\documentclass[12pt,a4paper,oneside]{book}

% Teoria de conjunts: una visió categòrica sense categories.
% Estructura comuna de la col·lecció aprendes / logcat (vegeu el volum de
% categories, versió 1.7): quatre parts (Teoria, Pràctica, Exercicis
% proposats, Tests), apèndixs, bibliografia, glossari i índex.

% Metadades pròpies del llibre: una única font de veritat.
% Dades pròpies del llibre. Es carrega abans del preàmbul compartit.
% Una sola font de veritat per a la portada i les metadades del PDF.
% Actualitzeu versió i data a cada versió publicada.
\newcommand{\llibretitol}{Teoria de conjunts}
\newcommand{\llibresubtitol}{Una visió categòrica sense categories}
\newcommand{\llibreautor}{Miquel Àngel Perelló}
\newcommand{\llibreversio}{1.7.1}
% Data editorial fixa (reproduïble), en lloc de \today.
\newcommand{\llibredata}{Octubre de 2026}
\newcommand{\llibreany}{2026}
% Els enllaços ja inclouen \url (no s'han d'embolcallar de nou).
\newcommand{\llibreweb}{\url{https://aprendes.com/logcat/}}
\newcommand{\llibrerepositori}{\url{https://github.com/maperello/aprendes}}
\newcommand{\llibreerrates}{\url{https://github.com/maperello/aprendes/issues}}

\title{\llibretitol\\[0.8em]\large\llibresubtitol}
\author{\llibreautor}
\date{\llibredata}

% S'ha de cridar després de carregar hyperref (és a preamble.tex).
\newcommand{\hypermetadata}{%
  \hypersetup{%
    pdftitle={\llibretitol: \llibresubtitol},
    pdfauthor={\llibreautor},
    pdfsubject={Teoria estructural de conjunts (ETCS) exposada en un llenguatge de primer ordre sobre aplicacions, sense pressuposar la teoria de categories: parts, construccions universals, exponencials, nombres naturals, cardinalitat, bons ordres, ordinals i cardinals. Versió \llibreversio\ (\llibredata). Llicència CC BY-SA 4.0},
    pdfkeywords={teoria de conjunts, ETCS, conjunts estructurals, Lawvere, topos ben apuntat, axioma d'elecció, ordinals, cardinals, hipòtesi del continu},
    bookmarksnumbered=true,
    bookmarksopen=true%
  }%
}


% Mode de sortida. Per defecte: pantalla/interactiu. El Makefile defineix
% \imprimible per generar una còpia pensada per imprimir: fons blanc,
% enllaços discrets i tests estàtics sense camps AcroForm.
\newif\ifimprimible
\ifdefined\imprimible\imprimibletrue\else\imprimiblefalse\fi
% Nom de la Part IV segons el mode: en paper no hi ha interactivitat.
\ifimprimible
  \newcommand{\NomTests}{Tests d'autoavaluació}
  \newcommand{\nomtests}{tests d'autoavaluació}
\else
  \newcommand{\NomTests}{Tests interactius}
  \newcommand{\nomtests}{tests interactius}
\fi

% Fitxers compartits de la col·lecció (idèntics als dels altres llibres).
\usepackage[sobri]{common/aprendes}
% ============================================================
%  Preàmbul compartit — Aprendes Universitat
%  Compartit pels llibres d'Aprendes Universitat.
%  Les dades pròpies de cada llibre (títol, autor i metadades PDF)
%  van al metadata.tex local de cada projecte, carregat ABANS
%  d'aquest fitxer. La macro \hypermetadata es crida DESPRÉS,
%  quan hyperref ja ha estat carregat.
% ============================================================

\usepackage{eurosym}
\usepackage{iftex}
\ifpdftex
  \usepackage[T1]{fontenc}
  \usepackage[utf8]{inputenc}
  \usepackage{lmodern}
\else
  % XeLaTeX / LuaLaTeX: Unicode natiu, res d'inputenc
  \usepackage{fontspec}
  % Carrega Latin Modern per nom de fitxer, que XeTeX/LuaTeX pot
  % resoldre via kpathsea fins i tot si fontconfig no l'ha registrat.
  \setmainfont{lmroman10-regular.otf}[
    BoldFont=lmroman10-bold.otf,
    ItalicFont=lmroman10-italic.otf,
    BoldItalicFont=lmroman10-bolditalic.otf,
    % Versaleta explícita (\textsc, \scshape): fontspec no la troba sola
    % quan la font es carrega per nom de fitxer.
    SmallCapsFont=lmromancaps10-regular.otf,
    SmallCapsFeatures={Letters=SmallCaps}
  ]
  \setsansfont{lmsans10-regular.otf}[
    BoldFont=lmsans10-bold.otf,
    ItalicFont=lmsans10-oblique.otf,
    BoldItalicFont=lmsans10-boldoblique.otf
  ]
  \setmonofont{lmmono10-regular.otf}[ItalicFont=lmmono10-italic.otf]
\fi
% microtype: protrusion i expansió completes només amb pdfLaTeX.
% Amb XeLaTeX/LuaLaTeX només hi ha protrusion i generaria avisos de
% «slot number»; per això només l'activem sota pdfTeX.
\ifpdftex
  \usepackage{microtype}
\fi
\usepackage[catalan, provide=*]{babel}
\usepackage[autostyle=true,babel=true]{csquotes}
% csquotes no porta estil català: el declarem perquè \enquote{...}
% doni guillemets «...» (exteriors) i cometes altes “...” (interiors),
% independentment de la versió de csquotes instal·lada.
\DeclareQuoteStyle{catalan}
  {\guillemotleft}{\guillemotright}
  {\quotedblbase}{\textquotedblleft}
\DeclareQuoteAlias{catalan}{catala}
\usepackage{amsmath,mathtools}% amssymb, amsthm: via aprendes.sty
\usepackage{geometry}
\usepackage{import}
\usepackage{graphicx}
\usepackage{tabularx}
\usepackage{xltabular}% tabularx que pagina (taula de notació)
\usepackage{tikz}
\usetikzlibrary{positioning,arrows.meta,matrix}
\usepackage[all]{xy}% diagrames de grafs (llaços i arestes múltiples) a l'apèndix de grafs
\usepackage{float}
\usepackage[section]{placeins}% FloatBarrier automàtic: cap figura no travessa una secció
\usepackage{pgfplots}
\pgfplotsset{compat=1.18}
\usepackage{imakeidx}
\usepackage{booktabs}

% La bibliografia es configura localment en cada projecte quan hi ha referències.

\setcounter{MaxMatrixCols}{10}

\graphicspath{{img/}}
\geometry{hmargin={3cm,3cm},vmargin={2.5cm,2.5cm},offset=0pt,footskip=1cm,headsep=1cm}


\setlist{itemsep=0.25em, topsep=0.25em}
\numberwithin{equation}{section}
\setcounter{secnumdepth}{3}
\setcounter{tocdepth}{2}

% Configuració dels paràgrafs
\usepackage{indentfirst}
\setlength{\parindent}{1.5em}
\setlength{\parskip}{0.3em}

% ============================================================
%  hyperref: ÚLTIM paquet, com recomana el seu manual.
%  Els colors dels enllaços es fixen aquí. Les metadades
%  pròpies del llibre ja s'han definit abans en metadata.tex;
%  la macro \hypermetadata s'executa després de carregar aquest fitxer.
% ============================================================
\usepackage[pdflang=ca-ES,pdfdisplaydoctitle=true]{hyperref}
% Color dels enllaços: el blau pissarra de la paleta d'aprendes,
% una mica enfosquit. Prou distint del vermell corporatiu dels
% títols, exemples i exercicis, i coherent amb la resta de colors.
\colorlet{EnllacBlau}{ResultatBlau!90!black}
\ifimprimible
  \hypersetup{
    colorlinks=true,
    linktoc=page,
    linkcolor=black,
    urlcolor=black,
    citecolor=black,
    filecolor=black,
    pdfborder={0 0 0},
    hypertexnames=false,
  }
\else
  \hypersetup{
    colorlinks=true,
    linktoc=page,
    linkcolor=EnllacBlau,
    urlcolor=EnllacBlau,
    citecolor=EnllacBlau,
    filecolor=EnllacBlau,
    hypertexnames=false,
  }
\fi
\usepackage{bookmark}% millora els marcadors del PDF; ha d'anar després d'hyperref


% ---- Paquets propis d'aquest llibre --------------------------------
\usepackage{stmaryrd}% \llbracket \rrbracket per a \seg; abans de common/macros.tex
\usepackage{longtable}% glossari de símbols
\usepackage[backend=biber,style=alphabetic,doi=true,url=true,isbn=false]{biblatex}
\addbibresource{bibliografia.bib}
\setlength{\emergencystretch}{3em}
\raggedbottom

% Macros matemàtiques i auxiliars compartides
\newcommand{\Qcb}[1]{#1}
\newcommand{\func}[1]{\operatorname{#1}}
\newcommand{\dbigcup}{\displaystyle\bigcup}
\newcommand{\dbigcap}{\displaystyle\bigcap}
\newcommand{\FRAME}[8]{%
\begin{center}
#5\par\smallskip
\includegraphics[width=#2,height=#3,keepaspectratio]{img/#7}
\end{center}}

\newcommand{\R}{\mathbb{R}}
\newcommand{\N}{\mathbb{N}}
\newcommand{\Z}{\mathbb{Z}}
\newcommand{\Q}{\mathbb{Q}}
\newcommand{\Natpos}{\mathbb{N}^{*}}
\newcommand{\Zstar}{\mathbb{Z}^{*}}
\newcommand{\mcd}{\operatorname{mcd}}
\newcommand{\mcm}{\operatorname{mcm}}
\newcommand{\abs}[1]{\left\lvert #1\right\rvert}
\newcommand{\classe}[1]{\left[#1\right]}
\newcommand{\divides}{\mid}
\newcommand{\notdivides}{\nmid}
\newcommand{\decimalperiodic}[2]{#1,\overline{#2}}
\DeclareMathOperator{\sgn}{sgn} 
% ------------------------------------------------------------
% Macros pròpies de teoria de categories
% ------------------------------------------------------------
\newcommand{\cat}[1]{\mathbf{#1}}          % noms de categories: \cat{Set}
\newcommand{\id}{\operatorname{id}}         % identitat
\providecommand{\Hom}{\operatorname{Hom}}   % conjunts de morfismes
\newcommand{\op}{^{\mathrm{op}}}            % categoria oposada
\newcommand{\Obj}{\operatorname{Obj}}       % col·lecció d'objectes
\newcommand{\comp}{\circ}                    % composició
\DeclareMathOperator{\dom}{dom}
\DeclareMathOperator{\cod}{cod}
\DeclareMathOperator{\Ob}{Ob}
\DeclareMathOperator{\Mor}{Mor}
\DeclareMathOperator{\Aut}{Aut}
\DeclareMathOperator{\Free}{Free}
\DeclareMathOperator{\Nat}{Nat}       % transformacions naturals com a homconjunt
\DeclareMathOperator{\Sub}{Sub}       % preordre de subobjectes
\DeclareMathOperator{\ev}{ev}         % morfisme d'avaluació (exponencials)
\DeclareMathOperator{\Imatge}{Im}     % imatge d'un morfisme (factorització imatge)
% Claudàtors semàntics [[-]] sense dependre de stmaryrd:
\providecommand{\llbracket}{\mathopen{[\![}}
\providecommand{\rrbracket}{\mathclose{]\!]}}
\newcommand{\yon}{\mathrm{y}}          % immersió de Yoneda
\DeclareMathOperator*{\colimop}{colim} % colímit (amb límits sota)

% ============================================================
%  Macros pròpies del volum de conjunts.
%  Es carreguen DESPRÉS de common/macros.tex, que ja defineix
%  \N, \Z, \Q, \R, \abs, \classe, \id, \dom, \cod, \Hom i \cat.
%  Requereix stmaryrd (per \llbracket), carregat a main.tex.
% ============================================================

% Vocabulari primitiu del llenguatge formal (capítol 1)
\DeclareMathOperator{\Dom}{Dom}
\DeclareMathOperator{\Cod}{Cod}
\DeclareMathOperator{\Com}{Com}
\providecommand{\Id}{\mathrm{Id}}

% Segment inicial de N (capítol 9)
\newcommand{\seg}[1]{\llbracket #1\rrbracket}

% Operadors de la part transfinita (capítol 15 endavant)
\DeclareMathOperator{\Hart}{Hart}
\DeclareMathOperator{\WO}{BonOrd}
\DeclareMathOperator{\camp}{camp}

% Tests interactius en PDF
\newcounter{testquestion}
\newcommand{\testprefix}{}
\newcommand{\testkeys}{}
\newcommand{\iniciTest}[2]{%
  \renewcommand{\testprefix}{#1}%
  \renewcommand{\testkeys}{#2}%
  \setcounter{testquestion}{0}%
  \gdef\testjustificacions{}%
  \label{test:#1}%
  \begin{center}
  \fbox{\begin{minipage}{0.9\textwidth}
  \ifimprimible
    \textbf{Instruccions.} Marca una sola resposta per pregunta. Comprova les respostes i les justificacions a l'apèndix de solucions.
  \else
    \textbf{Instruccions.} Marca una sola resposta per pregunta. En acabar, prem el bot\'{o} \textbf{Corregir}. El camp \textbf{Nota} mostrar\`{a} el resultat.\par\medskip
    \PushButton[name=#1corregir,width=3.2cm,height=0.8cm,bordercolor={0 0 1},onclick={var p='#1';var k='#2'.split(',');var opts=['a','b','c','d'];var s=0;for(var i=0;i<k.length;i++){var ans='';var n=0;for(var j=0;j<opts.length;j++){var q=this.getField(p+'q'+(i+1)+opts[j]);if(q && q.value!='Off'){ans=opts[j];n++;}}var f=this.getField(p+'f'+(i+1));if(n==1 && ans==k[i]){s++;if(f)f.value='Correcte';}else{if(f)f.value=(n==0?'Sense respondre. Correcta: '+k[i]:(n>1?'Marca una sola resposta. Correcta: '+k[i]:'Incorrecte. Correcta: '+k[i]));}}this.getField(p+'nota').value=s+' / '+k.length;app.alert('Nota: '+s+' / '+k.length);}]{Corregir}\quad
    \PushButton[name=#1netejar,width=3.2cm,height=0.8cm,bordercolor={1 0 0},onclick={var keys='#2'.split(',');var opts=['a','b','c','d'];this.getField('#1nota').value='';for(var i=1;i<=keys.length;i++){var f=this.getField('#1f'+i);if(f)f.value='';for(var j=0;j<opts.length;j++){var q=this.getField('#1q'+i+opts[j]);if(q)q.value='Off';}}}]{Netejar}\quad
    \textbf{Nota: }\TextField[name=#1nota,readonly=true,width=4cm,bordercolor={0 0 0}]{}
  \fi
  \end{minipage}}
  \end{center}
  \bigskip
}
\newcommand{\pregunta}{\stepcounter{testquestion}\item}
% Justificació breu de la resposta correcta d'una pregunta. No es mostra
% al costat de les opcions (per no donar pistes de longitud); s'acumula i
% \clauestatica la imprimeix després de la clau de respostes.
\newcommand{\testjustificacions}{}
\makeatletter
\newcommand{\justificacio}[1]{%
  \edef\@justnum{\thetestquestion}%
  \expandafter\g@addto@macro\expandafter\testjustificacions\expandafter{%
    \expandafter\justitem\expandafter{\@justnum}{#1}}%
}
\makeatother
\newcommand{\justitem}[2]{\item[\textbf{#1.}] #2}
\long\def\opcio#1#2{\item \ifimprimible$\square$\else\CheckBox[name=\testprefix{}q\thetestquestion#1,width=1.1em,height=1.1em,bordercolor={0 0 0},onclick={var opts=['a','b','c','d'];for(var j=0;j<opts.length;j++){if(opts[j]!='#1'){var q=this.getField('\testprefix{}q\thetestquestion'+opts[j]);if(q)q.value='Off';}}}]{}\fi\quad\textup{#1)}\enspace #2}
\newcommand{\feedback}{\ifimprimible\par\medskip\else\par\smallskip\textbf{Correcci\'{o}: }\TextField[name=\testprefix{}f\thetestquestion,readonly=true,width=9cm,bordercolor={0.6 0.6 0.6}]{}\par\medskip\fi}
% ---- Validació estructural de la clau (C7) ----
% \finalitzaTest compara el nombre de preguntes realment plantejades
% (\thetestquestion) amb el nombre de respostes de la clau, i emet un
% avís de compilació si no coincideixen. Cal cridar-la després de
% \end{Form}. Fem servir \@for, que recorre llistes separades per comes
% de manera segura, sense bucles \ifx artesanals.
\newcounter{testkeycount}
\makeatletter
\newcommand{\finalitzaTest}{%
  \setcounter{testkeycount}{0}%
  \expandafter\@for\expandafter\next\expandafter:\expandafter=\testkeys\do{\stepcounter{testkeycount}}%
  \ifnum\value{testquestion}=\value{testkeycount}\else
    \PackageWarning{interactive-tests}{El test '\testprefix' te \thetestquestion\space preguntes pero la clau en te \thetestkeycount\space respostes. Reviseu la crida a \string\iniciTest.}%
  \fi
}
\makeatother

% ============================================================
%  Claus estàtiques dels tests, reunides al final del llibre.
%
%  Els tests del llibre es corregeixen amb JavaScript d'AcroForm,
%  que no funciona en molts visors integrats, navegadors ni
%  aplicacions mòbils. Per a aquests casos, i per a l'ús en paper,
%  cada test té una clau de respostes i unes justificacions. Perquè
%  no apareguin just després de les preguntes, cosa que restaria
%  valor a l'autoavaluació, no s'imprimeixen al lloc del test:
%
%   - \clauestatica (cridada just abans de \finalitzaTest) desa la
%     clau i les justificacions del test en curs a \totesclaus i
%     imprimeix només una remissió a l'apèndix de solucions;
%   - \imprimeixsolucionstests, a caps/solucions-tests.tex, les
%     imprimeix totes al final del llibre, cadascuna amb un enllaç de
%     tornada al test.
% ============================================================

\makeatletter
\newcounter{clauestindex}
\gdef\totesclaus{}
\newcommand{\clauestatica}{%
  \begingroup
    \edef\@tmpclau{\noexpand\blocsoluciotest{\testprefix}{\thechapter}{\testkeys}{\unexpanded\expandafter{\testjustificacions}}}%
    \expandafter\g@addto@macro\expandafter\totesclaus\expandafter{\@tmpclau}%
  \endgroup
  \par\medskip
  \begin{center}
  \small\itshape
  \ifimprimible
  Les respostes i les justificacions d'aquest test són a l'apèndix \hyperref[cap:solucions-tests]{\textup{\enquote{Solucions dels tests}}}, pàgina~\pageref{sol:\testprefix}. Consulteu-les després d'haver respost tot el test.
  \else
  Si el vostre lector de PDF no admet el botó \textup{\textbf{Corregir}}, trobareu les respostes i les justificacions d'aquest test a l'apèndix \hyperref[cap:solucions-tests]{\textup{\enquote{Solucions dels tests}}}, pàgina~\pageref{sol:\testprefix}.
  \fi
  \end{center}
  \par\medskip
}
% #1 prefix del test, #2 número de capítol, #3 clau, #4 justificacions
\newcommand{\blocsoluciotest}[4]{%
  \section*{Test del capítol #2}\label{sol:#1}%
  \addcontentsline{toc}{section}{Test del capítol #2}%
  \noindent\textbf{Respostes.}\quad
  \setcounter{clauestindex}{0}%
  \@for\next:=#3\do{%
    \stepcounter{clauestindex}%
    \mbox{\textbf{\theclauestindex.}~\next}\hspace{1.2em plus 0.4em}%
  }%
  \par
  \def\@tmpjust{#4}%
  \ifx\@tmpjust\@empty\else
    \medskip
    \noindent\textbf{Justificacions.}
    \begin{list}{}{\setlength{\leftmargin}{2.2em}\setlength{\labelwidth}{1.8em}\setlength{\itemsep}{0.2em}\small}
    #4
    \end{list}
  \fi
  \par\smallskip
  \noindent{\small\hyperref[test:#1]{Tornar al test} (pàgina~\pageref{test:#1}).}
  \par\bigskip
}
\newcommand{\imprimeixsolucionstests}{\totesclaus}
\makeatother

% ============================================================
%  Entorns didàctics locals del volum de categories.
%  Definits aquí (no a l'estil compartit) per no acoblar-hi
%  el projecte. Reutilitzen tcolorbox, etoolbox i la paleta
%  ja carregats per aprendes.sty en mode «sobri».
%
%    \begin{objectius} ... \end{objectius}   -> bloc d'obertura
%    \begin{resumcap}   ... \end{resumcap}   -> bloc de tancament
%
%  Dins d'aquests blocs s'hi poden usar els encapçalaments
%  \apartat{...} per separar prerequisits/objectius/etc.
% ============================================================

% Color propi per als blocs didàctics: verd pissarra derivat de la
% paleta existent (cexer = verd), enfosquit per contrastar amb el fons.
\colorlet{DidacticVerd}{cexer!55!black}

% Estils de bloc: mateixa gramàtica visual que blocenunciat/blocobservacio.
\tcbset{
  blocdidactic/.style={
    enhanced, breakable, sharp corners,
    frame hidden, boxrule=0pt,
    colback=DidacticVerd!6,
    borderline west={2.0pt}{0pt}{DidacticVerd},
    left=8pt, right=4pt, top=5pt, bottom=5pt,
    before skip=1.0em, after skip=1.0em,
  },
}

% Encapçalament intern d'un apartat del bloc (Prerequisits, Objectius...).
\newcommand{\apartat}[1]{\par\smallskip{\color{DidacticVerd}\ifpdftex\scshape\else\sffamily\bfseries\fi #1}\par\nobreak\smallskip}

% Bloc d'obertura del capítol.
\newenvironment{objectius}{%
  \begin{tcolorbox}[blocdidactic]%
  \ignorespaces
}{\end{tcolorbox}}

% Bloc de tancament del capítol.
\newenvironment{resumcap}{%
  \begin{tcolorbox}[blocdidactic]%
  {\color{DidacticVerd}\ifpdftex\scshape\else\sffamily\bfseries\fi S\'{\i}ntesi del cap\'{\i}tol.}\hspace{0.6em}%
  \ignorespaces
}{\end{tcolorbox}}


\makeindex[title={Índex alfabètic},columns=2,intoc]

\begin{document}
\ifimprimible\pagecolor{white}\else\pagecolor{fons!20}\fi
\hypermetadata

\frontmatter
\maketitle
% Pàgina de crèdits (revers de la portada).
\thispagestyle{empty}
\vspace*{\fill}
\begin{flushleft}
\small
\textbf{\llibretitol}\\
\llibreautor\\[0.6em]
Versió \llibreversio\ · \llibredata\\
\llibreweb\\[1.2em]

\textbf{Llicència.} © \llibreany\ \llibreautor. Aquesta obra està subjecta a la llicència Creative Commons Reconeixement-CompartirIgual 4.0 Internacional (CC BY-SA 4.0): \url{https://creativecommons.org/licenses/by-sa/4.0/deed.ca}. Podeu copiar-la, redistribuir-la i adaptar-la, també amb finalitat comercial, sempre que en reconegueu l'autoria i distribuïu les obres derivades amb la mateixa llicència.\\[1em]

\textbf{Com citar-lo.} \llibreautor, \emph{\llibretitol}, versió \llibreversio, \llibreany. Projecte aprendes\,/\,logcat, \llibreweb.\\[1em]

\textbf{Errates i suggeriments.} Es poden comunicar obrint una incidència (\emph{issue}) al repositori del projecte, \llibreerrates, indicant la versió del llibre i la pàgina. Les correccions s'incorporen a les versions següents.
\end{flushleft}
\clearpage

\tableofcontents
\let\cleardoublepage\clearpage

\chapter*{Prefaci}
\addcontentsline{toc}{chapter}{Prefaci}

Aquest llibre forma part d'\textbf{aprendes}, una col·lecció de materials
d'estudi que comparteixen una mateixa manera de treballar: cada tema es
presenta en una part de \emph{teoria}, es reprèn en una part de
\emph{pràctica} amb problemes resolts, es consolida amb \emph{exercicis
proposats} i es comprova amb \emph{\nomtests}. Qui ja conegui
altres volums de la col·lecció hi reconeixerà l'estructura; qui s'hi
acosti per primera vegada la trobarà explicada a la guia de lectura.

El volum té un propòsit precís. Dins de la sèrie \textbf{logcat}, de
fonaments de la matemàtica i lògica categòrica, és el primer de quatre:
exposa la teoria de conjunts en la seva forma \emph{estructural}, el
segon, \emph{Una introducció a la teoria de categories}
\cite{perello-categories}, construeix el llenguatge general del qual
aquesta teoria és el cas fundacional, i els dos següents, dedicats a la
lògica categòrica, es basen en tots dos. Aquí els individus primitius
són les aplicacions, els conjunts són les identitats i cada construcció
---parts, sumes, productes, exponencials, nombres naturals--- es
defineix per la manera com es relaciona amb la resta mitjançant
aplicacions, i no per la naturalesa dels seus elements. És la teoria
elemental de la categoria dels conjunts de Lawvere, exposada en un
llenguatge de primer ordre que no pressuposa la teoria de categories:
per llegir aquest llibre no cal saber què és una categoria. Qui després
llegeixi el volum de categories hi reconeixerà cada construcció d'aquí
com el cas particular, en els conjunts, d'una noció general; l'apèndix
de correspondència entre els dos volums fa explícita aquesta lectura i
remet, per a cada construcció, al lloc del volum de categories on es
generalitza.

El recorregut va des dels axiomes sobre el domini, el codomini i la
composició fins a l'aritmètica dels ordinals i dels cardinals, la
cofinalitat i la hipòtesi del continu, passant per les construccions
universals, els valors de veritat, les relacions i l'elecció, les
imatges, els exponencials, el conjunt de parts i els nombres naturals.
A cada resultat s'hi indica de quins axiomes depèn, i en particular on
intervenen l'axioma d'elecció i, als darrers capítols, l'esquema de
substitució estructural que porta la teoria fins a la força de la teoria
clàssica. El capítol final fa balanç: situa els axiomes del llibre
respecte de la teoria de Zermelo--Fraenkel i de la formulació
categòrica de Lawvere.

Aquesta edició conté el volum complet: vint capítols de teoria, dinou
dels quals es desenvolupen també en les parts de pràctica, exercicis
proposats i tests, i els apèndixs, el glossari de símbols, la
bibliografia i l'índex.

Abans del primer capítol, la \emph{Introducció} explica què vol dir
mirar els conjunts des de les aplicacions i d'on ve aquesta idea.

\medskip
\noindent\textbf{Prerequisits.} El text pressuposa la maduresa matemàtica habitual d'un primer curs universitari de matemàtiques, informàtica o filosofia de la lògica: llegir fórmules de primer ordre amb quantificadors, seguir una demostració elemental i treballar informalment amb conjunts i funcions. No es pressuposa haver estudiat teoria axiomàtica de conjunts ni teoria de categories; la comparació amb la teoria clàssica es fa en observacions separades, que es poden saltar, i la terminologia categòrica apareix només com a nom de les construccions que ja s'han demostrat.

\medskip
\noindent\textbf{Errates i suggeriments.} Són benvinguts. Es poden comunicar obrint una incidència al repositori del projecte, \llibreerrates, indicant la versió (\llibreversio) i la pàgina.

% ------------------------------------------------------------
% Espai reservat per a l'autor (descomenteu i completeu):
%
% \medskip
% \noindent\textbf{Agraïments.} ...
% ------------------------------------------------------------

\phantomsection
\chapter*{Introducció}
\addcontentsline{toc}{chapter}{Introducció}
\label{cap:introduccio}

Aquest llibre desenvolupa una teoria estructural de conjunts en la qual les aplicacions són els únics individus primitius. Els conjunts, els elements i les construccions habituals es recuperen progressivament a partir dels dominis, els codominis i la composició.


\phantomsection
\section*{Una mirada diferent: el punt de vista estructural}
\addcontentsline{toc}{section}{Una mirada diferent: el punt de vista estructural}

En la teoria clàssica de conjunts ---la formulació tradicional de
Zermelo--Fraenkel--- els protagonistes indiscutibles són els
\textbf{elements}. Un conjunt es descriu visualment com un «sac» o una
«col·lecció» on guardem objectes dins, i la relació fonamental és la pertinença \( \in \), mitjançant la qual diem que un element \(x\) pertany al conjunt \(A\) escrivint \(x \in A\). Les aplicacions són actors secundaris, definides a posteriori com a conjunts de parells ordenats.

En aquest llibre canviarem la perspectiva: \textbf{els
protagonistes seran les aplicacions}.

Farem com en geometria: per entendre les figures sovint interessa més
saber \textbf{com es transformen o es relacionen entre elles} ---mitjançant
moviments, projeccions o simetries--- que no pas preguntar-se de què estan
fetes per dins. Adoptarem la mateixa actitud amb els conjunts: ens
interessarà com es relacionen mitjançant les aplicacions, més que no pas
la seva estructura interna de pertinença.

\begin{center}
\begin{tikzpicture}[scale=1,every node/.style={font=\small}]

%--------------------------------------------------
% Visió clàssica
%--------------------------------------------------

\node[font=\bfseries] at (1.5,3.2) {Visió clàssica (ZF)};

\draw[thick,rounded corners=14pt] (0,0) rectangle (3,2.5);

\node at (1.5,2.75) {$A$};

\fill (0.8,1.8) circle (2pt) node[right=2pt] {$x$};
\fill (2.2,1.0) circle (2pt) node[right=2pt] {$y$};
\fill (1.1,0.5) circle (2pt) node[right=2pt] {$z$};

\node at (1.5,-0.45)
{Contenidor d'elements: $x\in A$};

%--------------------------------------------------
% Fletxa central
%--------------------------------------------------

\draw[-{Stealth[scale=1.4]},line width=1.2pt]
(4.0,1.25)--(5.8,1.25);

%--------------------------------------------------
% Visió estructural
%--------------------------------------------------

\node[font=\bfseries] at (8.5,3.2)
{Visió estructural (aquest llibre)};

\node (A) at (7.2,1.8) {$A$};
\node (B) at (9.6,1.8) {$B$};
\node (C) at (8.4,0.4) {$C$};

% Aplicacions

\draw[-{Stealth[scale=1.2]},thick]
(A)--node[above]{$f$}(B);

\draw[-{Stealth[scale=1.2]},thick]
(A)--node[left]{$g$}(C);

\draw[-{Stealth[scale=1.2]},thick]
(C)--node[right]{$h$}(B);

%--------------------------------------------------
% Identitats
%--------------------------------------------------

% id_A
\draw[-{Stealth[scale=1]},thick]
(7.20,2.00)
.. controls (6.80,2.45) and (7.60,2.45) ..
(7.20,2.00);
\node[font=\scriptsize] at (6.55,2.48)
{$\id_A$};

% id_B
\draw[-{Stealth[scale=1]},thick]
(9.60,2.00)
.. controls (9.20,2.45) and (10.00,2.45) ..
(9.60,2.00);
\node[font=\scriptsize] at (10.25,2.48)
{$\id_B$};

% id_C
\draw[-{Stealth[scale=1]},thick]
(8.40,0.20)
.. controls (8.00,-0.25) and (8.80,-0.25) ..
(8.40,0.20);
\node[font=\scriptsize] at (9.05,-0.28)
{$\id_C$};

\node at (8.4,-0.85)
{Xarxa d'aplicacions};

\end{tikzpicture}
\end{center}

\subsection*{Una mica de context: d'on ve aquesta idea?}

Aquesta manera de reconsiderar els fonaments de la matemàtica \textbf{no és
nova}. La va proposar el matemàtic F.~William Lawvere durant la dècada de
1960~\cite{lawvere1964}, en adonar-se que la pràctica quotidiana ---en àlgebra, topologia o
anàlisi--- consisteix a estudiar com els objectes interactuen entre ells
mitjançant aplicacions, i no pas a registrar la naturalesa dels seus
elements interns.

Lawvere ho va formular en el llenguatge de la teoria de categories, en el
que es coneix com la \emph{teoria estructural de conjunts de Lawvere}
(ETCS)~\cite{lawvere1964,lawvere-etcs-reprint}, però \textbf{per seguir aquest llibre no cal conèixer la teoria de categories}. Una presentació moderna, propera a l'esperit d'aquest llibre, es troba a \cite{lawvere-rosebrugh}; per a la lectura des de la teoria de topos, vegeu \cite{maclane-moerdijk}.

El nostre objectiu és construir una teoria estructural de conjunts des de zero, amb un llenguatge formal de primer ordre elemental i amb els convenis habituals que permeten abreujar la notació.


\section*{Criteri d'escriptura formal}\label{not:criteri-formal}
Les definicions i els axiomes s'escriuran mitjançant fórmules del llenguatge formal o mitjançant abreviacions introduïdes explícitament. En les proposicions, els teoremes i les demostracions farem servir sistemàticament aquestes abreviacions, perquè la notació primitiva no amagui l'estructura matemàtica dels arguments.

Cada conveni introduït en un entorn de notació és eliminable: qualsevol expressió que el contingui es pot traduir a una fórmula construïda amb \(\Dom\), \(\Cod\), \(\Com\), la igualtat i els predicats definits anteriorment. Per tant, aquests convenis no amplien la teoria; només en simplifiquen l'escriptura.

La presentació alterna tres nivells. El nivell formal fixa les definicions i els axiomes; el nivell geomètric els representa mitjançant diagrames externs de fletxes; i les observacions de lectura clàssica expliquen com es recuperen els conjunts i les funcions habituals. Aquests tres nivells es mantindran diferenciats al llarg de l'obra.


\mainmatter

% ============================================================
% PART I — TEORIA
% Nivell I (capítols 1-18): ETCS, amb l'elecció a partir del capítol 4.
% Capítol 19: l'esquema SR com a principi adoptat. Capítol 20: balanç.
% ============================================================
\part{Teoria}\label{part:teoria}
\setcounter{chapter}{0}
\chapter{Aplicacions, dominis, codominis i composició}\label{cap:aplicacions}

\begin{objectius}
\apartat{Prerequisits} Llenguatge elemental de la lògica de primer ordre: variables, predicats, connectives, quantificadors i igualtat, i la distinció entre existència i unicitat. Familiaritat informal amb les funcions entre conjunts (domini, codomini, composició, identitat), que aquí es reconstrueix des de zero. No es pressuposa cap coneixement de teoria axiomàtica de conjunts ni de teoria de categories; les comparacions amb la teoria clàssica són només de lectura.
\apartat{Objectius} En acabar el capítol, el lector hauria de poder:
\begin{itemize}
  \item llegir i escriure fórmules amb els tres predicats primitius \(\Dom\), \(\Cod\) i \(\Com\), i traduir-hi la notació habitual \(f\colon A\to B\) i \(g\circ f\);
  \item enunciar els cinc axiomes estructurals i deduir-ne l'estabilitat del domini i del codomini i la unicitat de la composició;
  \item explicar per què un conjunt és, en aquesta teoria, una aplicació identitat, i fer servir el conveni de les majúscules per als conjunts;
  \item llegir un diagrama commutatiu com una igualtat de composicions;
  \item distingir elements generalitzats i elements globals, i aplicar el criteri d'igualtat d'aplicacions per elements generalitzats;
  \item definir isomorfisme, demostrar la unicitat de la inversa i les propietats de la isomorfia, i reconèixer una propietat estructural;
  \item definir conjunt terminal i conjunt inicial, demostrar-ne la unicitat llevat d'un únic isomorfisme i formular el principi de dualitat;
  \item explicar la funció de cadascun dels quatre axiomes d'existència i de separació: terminal, inicial, no-degeneració i separació global.
\end{itemize}
\end{objectius}

\section{El llenguatge formal i els axiomes inicials}

El nostre objectiu és construir axiomàticament una teoria estructural de conjunts, diferent de la teoria clàssica de conjunts. Per exposar-la usarem un llenguatge de primer ordre amb igualtat i, quan sigui convenient, la notació estàndard de la pràctica matemàtica com a abreviació del llenguatge formal.

Totes les variables del nostre llenguatge, majúscules i minúscules, recorren un mateix univers d'individus, que anomenarem \emph{aplicacions}. Fins que no hàgim definit els conjunts, no hi haurà cap distinció formal entre variables majúscules i minúscules. El vocabulari no lògic consta de dos predicats binaris, \(\Dom\) i \(\Cod\), i d'un predicat ternari, \(\Com\). La lectura prevista és la següent:
\begin{itemize}
  \item \(\Dom(f,a)\) es llegeix \enquote{\(a\) és el domini de \(f\)}.
  \item \(\Cod(f,b)\) es llegeix \enquote{\(b\) és el codomini de \(f\)}.
  \item \(\Com(f,g,h)\) es llegeix \enquote{\(h\) és el resultat de compondre primer \(f\) i després \(g\)}.
\end{itemize}

\section{Axiomes i propietats estructurals}\label{sec:axiomes-estructurals}

Primer hem de garantir que el domini, el codomini i la composició es comporten com esperem. Per això enunciem els cinc axiomes següents. 

\begin{axioma}[Existència i unicitat del domini i del codomini]\label{ax:dom-cod}
Tota aplicació té un domini i un codomini únics:
\[
\forall f\,\exists a\,\Bigl(\Dom(f,a)\land
  \forall a' \,\bigl(\Dom(f,a')\Rightarrow a=a'\bigr)\Bigr)
\]
i
\[
\forall f\,\exists b\,\Bigl(\Cod(f,b)\land
  \forall b'\,\bigl(\Cod(f,b')\Rightarrow b=b'\bigr)\Bigr).
\]
\end{axioma}

\begin{axioma}[Existència de la composició]\label{ax:composicio-existencia}
Dues aplicacions es poden compondre si i només si el codomini de la primera coincideix amb el domini de la segona:
\[
\forall f\,\forall g\,
\Bigl(
  \exists a\bigl(\Cod(f,a)\land\Dom(g,a)\bigr)
  \Leftrightarrow
  \exists h\,\Com(f,g,h)
\Bigr).
\]
\end{axioma}

Observeu que aquest axioma només afirma l'existència d'una aplicació composta. La seva unicitat es demostra a la \hyperref[prop:composicio-unica]{proposició~\ref*{prop:composicio-unica}} com una conseqüència dels axiomes.

\begin{axioma}[Domini i codomini de la composició]\label{ax:composicio-dom-cod}
L'aplicació composta conserva el domini de la primera aplicació i el codomini de la segona:
\[
\forall f\,\forall g\,\forall h\,\forall a\,\forall b\,
\Bigl(
  \Com(f,g,h)\land\Dom(f,a)\land\Cod(g,b)
  \Rightarrow
  \Dom(h,a)\land\Cod(h,b)
\Bigr).
\]
\end{axioma}

\begin{axioma}[Llei associativa]\label{ax:associativitat}
Si tres aplicacions es poden compondre consecutivament, el resultat no depèn de la manera d'agrupar-les:
\[
\begin{aligned}
\forall f\,\forall g\,\forall h\,\forall u\,\forall v\,\forall w\,\forall z\,
\Bigl(&\Com(f,g,u)\land\Com(g,h,v)\\
      &\land\Com(u,h,w)\land\Com(f,v,z)
      \Rightarrow w=z\Bigr).
\end{aligned}
\]
\end{axioma}

En aquesta fórmula, \(u\) és una composició de \(f\) seguida de \(g\), \(v\) és una composició de \(g\) seguida de \(h\), \(w\) és una composició d'\(u\) seguida d'\(h\), i \(z\) és una composició de \(f\) seguida de \(v\).

\begin{axioma}[Lleis de neutralitat]\label{ax:neutralitat}
El domini d'una aplicació actua com a neutre per la dreta en la composició:
\[
\forall f\,\forall a\,
\bigl(\Dom(f,a)\Rightarrow\Com(a,f,f)\bigr)
\]
i el codomini actua com a neutre per l'esquerra:
\[
\forall f\,\forall b\,
\bigl(\Cod(f,b)\Rightarrow\Com(f,b,f)\bigr).
\]
\end{axioma}

L'estabilitat del domini i del codomini, que inicialment es podria prendre com un axioma, és una conseqüència dels cinc axiomes anteriors.

\begin{proposicio}[Estabilitat del domini i del codomini]\label{prop:estabilitat}
Si una aplicació és domini o codomini d'una altra, aleshores és domini i codomini de si mateixa:
\[
\forall f \, \forall a \, \Bigl(\Dom(f,a)\Rightarrow\bigl(\Dom(a,a)\land\Cod(a,a)\bigr)\Bigr)
\]
i
\[
\forall f \, \forall b \, \Bigl(\Cod(f,b)\Rightarrow\bigl(\Dom(b,b)\land\Cod(b,b)\bigr)\Bigr).
\]
\end{proposicio}

\begin{prova}
Suposem primer que \(\Dom(f,a)\). Per la \hyperref[ax:neutralitat]{primera llei de neutralitat}, obtenim
\[
\Com(a,f,f).
\]
Com que aquesta composició existeix, l'\hyperref[ax:composicio-existencia]{axioma~\ref*{ax:composicio-existencia}} proporciona una aplicació \(x\) tal que
\[
\Cod(a,x)\land\Dom(f,x).
\]
La unicitat del domini de \(f\), juntament amb \(\Dom(f,a)\), implica \(x=a\). Per tant,
\begin{equation}\label{eq:codaa}
\Cod(a,a).
\end{equation}
Per l'\hyperref[ax:dom-cod]{axioma~\ref*{ax:dom-cod}}, existeixen aplicacions \(y\) i \(b\) tals que \(\Dom(a,y)\) i \(\Cod(f,b)\). L'\hyperref[ax:composicio-dom-cod]{axioma~\ref*{ax:composicio-dom-cod}}, aplicat a
\[
\Com(a,f,f),\qquad \Dom(a,y),\qquad \Cod(f,b),
\]
dona \(\Dom(f,y)\). Com que també \(\Dom(f,a)\), la unicitat del domini de \(f\) implica \(y=a\). Així,
\begin{equation}\label{eq:domaa}
\Dom(a,a).
\end{equation}
De \eqref{eq:codaa} i \eqref{eq:domaa} deduïm
\[
\Dom(a,a)\land\Cod(a,a).
\]
Considerem ara que \(\Cod(f,b)\). Per la \hyperref[ax:neutralitat]{segona llei de neutralitat},
\[
\Com(f,b,f).
\]
L'\hyperref[ax:composicio-existencia]{axioma~\ref*{ax:composicio-existencia}} implica que hi ha una aplicació \(x\) tal que
\[
\Cod(f,x)\land\Dom(b,x).
\]
Per la unicitat del codomini de \(f\), tenim \(x=b\), i per tant
\begin{equation}\label{eq:dombb}
\Dom(b,b).
\end{equation}
Finalment, per l'\hyperref[ax:dom-cod]{axioma~\ref*{ax:dom-cod}}, existeixen aplicacions \(y\) i \(c\) tals que \(\Cod(b,y)\) i \(\Dom(f,c)\). L'\hyperref[ax:composicio-dom-cod]{axioma~\ref*{ax:composicio-dom-cod}}, aplicat a
\[
\Com(f,b,f),\qquad \Dom(f,c),\qquad \Cod(b,y),
\]
dona \(\Cod(f,y)\). Com que també \(\Cod(f,b)\), la unicitat del codomini de \(f\) implica \(y=b\). Així,
\begin{equation}\label{eq:codbb}
\Cod(b,b).
\end{equation}
De \eqref{eq:dombb} i \eqref{eq:codbb} deduïm
\[
\Dom(b,b)\land\Cod(b,b).
\]
\end{prova}

També la unicitat de la composició es dedueix dels axiomes.

\begin{proposicio}[Unicitat de la composició]\label{prop:composicio-unica}
La composició de dues aplicacions que es poden compondre està unívocament determinada:
\[
\forall f \, \forall g \, \forall h \, \forall h' \, \bigl(\Com(f,g,h)\land\Com(f,g,h')\Rightarrow h=h'\bigr).
\]
\end{proposicio}

\begin{prova}
Suposem que
\[
\Com(f,g,h)
\qquad\text{i}\qquad
\Com(f,g,h').
\]
Per l'\hyperref[ax:dom-cod]{axioma~\ref*{ax:dom-cod}}, existeixen aplicacions \(a\) i \(b\) tals que \(\Dom(f,a)\) i \(\Cod(g,b)\). Per l'\hyperref[ax:composicio-dom-cod]{axioma~\ref*{ax:composicio-dom-cod}},
\[
\Dom(h,a)\land\Cod(h,b).
\]
Les \hyperref[ax:neutralitat]{lleis de neutralitat} donen
\[
\Com(g,b,g)
\qquad\text{i}\qquad
\Com(h,b,h).
\]
Apliquem ara l'\hyperref[ax:associativitat]{axioma~\ref*{ax:associativitat}} a les quatre relacions
\[
\Com(f,g,h),\qquad
\Com(g,b,g),\qquad
\Com(h,b,h),\qquad
\Com(f,g,h').
\]
En resulta \(h=h'\).
\end{prova}

\section{Identitats i conjunts}\label{sec:identitats-conjunts}

Definirem els conjunts com un tipus determinat d'aplicacions, però abans introduirem les aplicacions que actuaran com a identitats.

\begin{definicio}[Aplicació identitat]\index{Aplicació identitat}\label{def:identitat}
Una aplicació \(e\) és una \emph{identitat} si és el seu propi domini i el seu propi codomini:
\[
\operatorname{Id}(e)
\iff
\Dom(e,e)\land \Cod(e,e).
\]
\end{definicio}

L'\hyperref[ax:neutralitat]{axioma~\ref*{ax:neutralitat}} garanteix que qualsevol aplicació així definida actua com a neutre en totes les composicions compatibles. En efecte, si \(e\) és el domini d'una aplicació \(f\), és a dir, si \(\Dom(f,e)\), aleshores \(\Com(e,f,f)\); anàlogament, de \(\Cod(f,e)\) s'obté \(\Com(f,e,f)\). El que aporta la definició és el recíproc: les aplicacions neutres són exactament les que són el seu propi domini i codomini, com veurem tot seguit.

\begin{proposicio}\label{prop:dom-cod-identitats}
El domini i el codomini de qualsevol aplicació són identitats. A més, si una identitat \(e\) té \(a\) com a domini o com a codomini, aleshores \(e=a\).
\end{proposicio}

\begin{prova}
Si \(\Dom(f,a)\) o \(\Cod(f,a)\), la \hyperref[prop:estabilitat]{proposició~\ref*{prop:estabilitat}} dona
\[
\Dom(a,a)\land \Cod(a,a),
\]
i, per tant, \(a\) és una identitat.

Ara suposem que \(e\) és una identitat i que \(\Dom(e,a)\). Com que també \(\Dom(e,e)\), la unicitat del domini de \(e\) implica \(e=a\). El cas en què \(\Cod(e,a)\) és anàleg.
\end{prova}

\begin{definicio}[Conjunt]\index{Conjunt}\label{def:conjunt}
Un \emph{conjunt} és una aplicació que té per domini i codomini la mateixa aplicació:
\[
\operatorname{Conj}(a)
\iff
\Dom(a,a)\land \Cod(a,a).
\]
\end{definicio}

Si comparem la \hyperref[def:identitat]{definició~\ref*{def:identitat}} i la \hyperref[def:conjunt]{definició~\ref*{def:conjunt}} és obvi que tot conjunt és una identitat i a l'inrevés:
\[
\operatorname{Conj}(a)
\iff
\operatorname{Id}(a)
\iff
\Dom(a,a)\land \Cod(a,a).
\]

La \hyperref[prop:dom-cod-identitats]{proposició~\ref*{prop:dom-cod-identitats}} garanteix que el domini i el codomini de qualsevol aplicació són conjunts. També, si un conjunt \(a\) té \(b\) com a domini o com a codomini, aleshores \(a=b\).

\begin{notacio}[Conveni sobre les variables]\label{not:variables-conjunts-aplicacions}
A partir d'ara prendrem el conveni metalingüístic que permetrà reconèixer immediatament quines variables estan restringides als conjunts. Les lletres majúscules
\[
A,B,C,\ldots
\]
es reservaran per a variables restringides als conjunts. Així,
\[
\forall A\,\varphi(A)
\]
és una abreviació de
\[
\forall A\,\bigl(\operatorname{Conj}(A)\Rightarrow\varphi(A)\bigr),
\]
i
\[
\exists A\,\varphi(A)
\]
és una abreviació de
\[
\exists A\,\bigl(\operatorname{Conj}(A)\land\varphi(A)\bigr).
\]
La mateixa convenció s'aplica al quantificador d'existència única. Quan una majúscula aparegui lliure en una definició, un enunciat o una demostració, entendrem que designa un conjunt.

\bigskip

Les lletres minúscules
\[
f,g,h,x,y,z,\ldots
\]
denotaran aplicacions arbitràries. No se n'exclouen les identitats: una variable minúscula pot designar també un conjunt quan el context ho permeti.

Si \(A\) és un conjunt, escriurem \(\id_A\) quan vulguem destacar el seu paper d'identitat. Aquesta notació no designa una aplicació diferent de \(A\):
\[
\id_A=A.
\]
Totes aquestes convencions són eliminables i no modifiquen el llenguatge formal d'una sola mena. Per tant, després d'haver-les declarat, les podrem usar també en definicions i axiomes, sempre com a abreviacions del llenguatge primitiu.
\end{notacio}

\section{Notació i diagrames d'aplicacions}\label{sec:diagrames-aplicacions}

Un cop definits els conjunts i les identitats, podem introduir els convenis de notació que faran més llegible el desenvolupament posterior. Aquests convenis pertanyen al metallenguatge i són eliminables en el sentit establert pel \hyperref[not:criteri-formal]{criteri d'escriptura formal} de la introducció.

\begin{notacio}[Aplicacions, domini, codomini i composició]\label{not:aplicacions-tipades}
D'acord amb el conveni anterior, les expressions
\[
f:A\rightarrow B
\qquad\text{i}\qquad
A\xrightarrow{f}B
\]
són abreviacions de la fórmula
\[
\Dom(f,A)\land\Cod(f,B).
\]
No introduïm cap predicat primitiu nou per expressar aquest fet.
Per l'\hyperref[ax:dom-cod]{axioma~\ref*{ax:dom-cod}}, el domini i el codomini estan unívocament determinats. Per això escriurem
\[
A=\dom(f),
\qquad
B=\cod(f),
\qquad
f:\dom(f)\to\cod(f).
\]
Amb aquesta notació, l'estabilitat s'escriu
\[
\dom(\dom(f))=\dom(f),
\qquad
\cod(\dom(f))=\dom(f),
\]
\[
\dom(\cod(f))=\cod(f),
\qquad
\cod(\cod(f))=\cod(f).
\]
Si \(\cod(f)=\dom(g)\), l'existència i la unicitat de la composició determinen una única aplicació \(h\) tal que \(\Com(f,g,h)\). Escriurem
\[
h=g\circ f.
\]
Així, \(g\circ f\) és una abreviació metalingüística de l'única aplicació \(h\) que satisfà \(\Com(f,g,h)\); no és un terme primitiu nou del llenguatge.
\end{notacio}

\begin{notacio}[Diagrames d'aplicacions i commutativitat]
\label{not:diagrames-commutatius}

Un \emph{diagrama d'aplicacions} és una representació geomètrica
del metallenguatge en què els conjunts es representen mitjançant
vèrtexs i les aplicacions mitjançant fletxes. Un camí orientat
representa la composició successiva de les aplicacions que el
formen. Direm que un diagrama és \emph{commutatiu} quan els camins orientats que s'hi comparen,
amb el mateix origen i el mateix destí, determinen la mateixa aplicació composta.

Aquesta notació no introdueix nous individus, nous predicats ni
nous axiomes. Cada afirmació concreta de commutativitat és una
abreviació metalingüística d'una igualtat entre composicions i,
per tant, es pot traduir mitjançant el predicat \(\Com\) i la
igualtat. La noció general de diagrama finit, necessària per
estudiar els límits finits, s'introduirà quan correspongui.
\end{notacio}

Siguin
\[
f:A\to B,\qquad g:B\to C,\qquad h:A\to C.
\]
El triangle
\begin{center}
\begin{tikzpicture}[scale=1.1, every node/.style={font=\small}]
  \node (A) at (0,1.6) {\(A\)};
  \node (B) at (3.2,1.6) {\(B\)};
  \node (C) at (3.2,0) {\(C\)};
  \draw[-{Stealth[scale=1.15]}, thick] (A) -- node[above] {\(f\)} (B);
  \draw[-{Stealth[scale=1.15]}, thick] (B) -- node[right] {\(g\)} (C);
  \draw[-{Stealth[scale=1.15]}, thick] (A) -- node[below left] {\(h\)} (C);
\end{tikzpicture}
\end{center}
és \emph{commutatiu} si els dos camins d'\(A\) a \(C\) determinen la mateixa aplicació, és a dir, si
\[
h=g\circ f.
\]
En el llenguatge formal, aquesta igualtat s'expressa directament mitjançant
\[
\Com(f,g,h).
\]

De la mateixa manera, siguin
\[
f:A\to B,\qquad g:A\to C,\qquad h:B\to D,\qquad k:C\to D.
\]
El quadrat
\begin{center}
\begin{tikzpicture}[scale=1.1, every node/.style={font=\small}]
  \node (A) at (0,1.8) {\(A\)};
  \node (B) at (3.4,1.8) {\(B\)};
  \node (C) at (0,0) {\(C\)};
  \node (D) at (3.4,0) {\(D\)};
  \draw[-{Stealth[scale=1.15]}, thick] (A) -- node[above] {\(f\)} (B);
  \draw[-{Stealth[scale=1.15]}, thick] (A) -- node[left] {\(g\)} (C);
  \draw[-{Stealth[scale=1.15]}, thick] (B) -- node[right] {\(h\)} (D);
  \draw[-{Stealth[scale=1.15]}, thick] (C) -- node[below] {\(k\)} (D);
\end{tikzpicture}
\end{center}
és commutatiu quan els dos camins d'\(A\) a \(D\) determinen
la mateixa aplicació:
\[
h\circ f=k\circ g.
\]
Formalment, aquesta condició es pot escriure
\[
\exists u\,\bigl(\Com(f,h,u)\land\Com(g,k,u)\bigr).
\]
En efecte, aquesta fórmula afirma que els dos camins tenen una
mateixa aplicació composta.

En general, quan en un diagrama comparem diversos camins amb
el mateix origen i el mateix destí, la commutativitat significa
que totes les composicions indicades són iguals. Per exemple,
la llei associativa, l'\hyperref[ax:associativitat]{axioma~\ref*{ax:associativitat}}, afirma que, per a aplicacions
\[
A\xrightarrow{f}B\xrightarrow{g}C\xrightarrow{h}D,
\]
les dues maneres de reagrupar la composició determinen la mateixa
aplicació:
\[
h\circ(g\circ f)=(h\circ g)\circ f.
\]

\begin{figure}[H]
  \centering
  \begin{tikzpicture}[scale=1.2, every node/.style={font=\small}]
    \node (A) at (0,1.6) {\(A\)};
    \node (B) at (2.0,1.6) {\(B\)};
    \node (C) at (4.0,1.6) {\(C\)};
    \node (D) at (6.0,1.6) {\(D\)};
    \draw[-{Stealth[scale=1.2]}, thick] (A) -- node[above] {\(f\)} (B);
    \draw[-{Stealth[scale=1.2]}, thick] (B) -- node[above] {\(g\)} (C);
    \draw[-{Stealth[scale=1.2]}, thick] (C) -- node[above] {\(h\)} (D);
    \draw[-{Stealth[scale=1.2]}, thick] (A) to[bend right=20] node[below] {\(g \circ f\)} (C);
    \draw[-{Stealth[scale=1.2]}, thick] (B) to[bend left=20] node[above=6pt] {\(h \circ g\)} (D);
    \draw[-{Stealth[scale=1.2]}, thick] (A) to[bend right=38] node[below] {\(h\circ(g\circ f)=(h\circ g)\circ f\)} (D);
  \end{tikzpicture}
  \caption{La llei associativa: els dos camins agrupats de \(A\) a \(D\) són la mateixa aplicació.}
  \label{fig:associativa}
\end{figure}

Quan una aplicació aparegui com la fletxa que s'ha de trobar en una condició d'existència o d'unicitat, la podrem representar mitjançant una fletxa discontínua. Una vegada establertes aquestes condicions, la representarem mitjançant una fletxa contínua.

\begin{figure}[H]
  \centering
\begin{tikzpicture}[scale=1.1, every node/.style={font=\small}]
  \node (A) at (0,1.6) {\(A\)};
  \node (B) at (3.2,1.6) {\(B\)};
  \node (C) at (3.2,0) {\(C\)};
  \draw[-{Stealth[scale=1.15]}, thick] (A) -- node[above] {\(f\)} (B);
  \draw[-{Stealth[scale=1.15]}, thick] (B) -- node[right] {\(g\)} (C);
  \draw[-{Stealth[scale=1.15]}, thick, dashed] (A) -- node[below left] {\(g \circ f\)} (C);
\end{tikzpicture}
\caption{Diagrama de la composició: quan \(\cod(f)=\dom(g)\), la fletxa \(g \circ f\) existeix i és única.}
  \label{fig:composicio}
\end{figure}

Amb aquesta notació, l'\hyperref[ax:composicio-dom-cod]{axioma~\ref*{ax:composicio-dom-cod}} s'expressa així:
\[
\dom(g\circ f)=\dom(f)
\quad \text{i} \quad
\cod(g\circ f)=\cod(g).
\]

Finalment, l'\hyperref[ax:neutralitat]{axioma~\ref*{ax:neutralitat}} s'escriu així:

\[
f \circ \dom(f)=f,
\qquad
\cod(f) \circ f=f.
\]

\begin{notacio}[Identitat d'un conjunt]\label{not:identitat-conjunt}
Un conjunt \(A\) compleix \(\dom(A)=\cod(A)=A\), i per això podem escriure
\[
A:A\to A.
\]
Quan vulguem posar de relleu el seu paper neutral en la composició, escriurem
\[
\id_A.
\]
Formalment,
\[
\id_A=A;
\]
les expressions \(A\) i \(\id_A\) designen la mateixa aplicació, no dues aplicacions diferents.
\end{notacio}

Com a mostra d'aquesta notació, considerem una aplicació \(f:A\to B\). L'\hyperref[ax:neutralitat]{axioma~\ref*{ax:neutralitat}} s'expressa així:
\[
f\circ \id_A=f
\qquad\text{i}\qquad
\id_B\circ f=f.
\]
Els dos triangles commutatius següents posen de relleu el paper neutral dels conjunts:

\begin{figure}[H]
  \centering
  \begin{tikzpicture}[scale=1.05, every node/.style={font=\small}]
    \node (A1) at (0,1.6) {\(A\)};
    \node (A2) at (2.5,1.6) {\(A\)};
    \node (B1) at (2.5,0) {\(B\)};
    \draw[-{Stealth[scale=1.15]}, thick] (A1) -- node[above] {\(\id_A\)} (A2);
    \draw[-{Stealth[scale=1.15]}, thick] (A2) -- node[right] {\(f\)} (B1);
    \draw[-{Stealth[scale=1.15]}, thick] (A1) -- node[below left] {\(f\)} (B1);

    \node (A3) at (4.8,1.6) {\(A\)};
    \node (B2) at (7.3,1.6) {\(B\)};
    \node (B3) at (7.3,0) {\(B\)};
    \draw[-{Stealth[scale=1.15]}, thick] (A3) -- node[above] {\(f\)} (B2);
    \draw[-{Stealth[scale=1.15]}, thick] (B2) -- node[right] {\(\id_B\)} (B3);
    \draw[-{Stealth[scale=1.15]}, thick] (A3) -- node[below left] {\(f\)} (B3);
  \end{tikzpicture}
  \caption{Diagrames de les lleis de neutralitat: \(f\circ\id_A=f\) i \(\id_B\circ f=f\).}
  \label{fig:neutralitat}
\end{figure}

\section{Igualtat d'aplicacions}\label{sec:igualtat-aplicacions}

La pregunta que ens fem ara és: quan dues aplicacions són iguals? Els axiomes que hem introduït fins ara permeten el bon comportament de les aplicacions. De fet, en la \hyperref[sec:diagrames-aplicacions]{secció~\ref*{sec:diagrames-aplicacions}} hem vist com una aplicació \(f\) es podia escriure com 
\[
f:\dom(f) \rightarrow \cod(f),
\]
i, per tant, dues aplicacions només poden ser iguals si tenen el mateix domini i el mateix codomini; en aquest cas direm que són \emph{paral·leles}\index{Aplicacions paral·leles}. Ara bé, les aplicacions també es poden comparar per la manera com es componen amb altres aplicacions. Considerem, doncs, dues aplicacions \(f,g:A\to B\) i una aplicació \(x:X\to A\). L'\hyperref[ax:composicio-existencia]{axioma~\ref*{ax:composicio-existencia}} i l'\hyperref[ax:composicio-dom-cod]{axioma~\ref*{ax:composicio-dom-cod}} garanteixen que
\[
f\circ x:X\rightarrow B \quad \text{i} \quad g\circ x:X\rightarrow B.
\]
És natural esperar que si \(f\) i \(g\) són iguals, també ho siguin
\[
f\circ x=g\circ x.
\]

Aquesta propietat es formula en la proposició següent.

\begin{proposicio}[Aplicacions iguals]\label{prop:aplicacions-iguals}
Siguin \(f,g:A\to B\). Aleshores,
\[
f=g
\iff
\forall X \, \forall x \, \bigl( x:X \to A \Rightarrow f \circ x=g\circ x\bigr).
\]
\end{proposicio}

\begin{prova}
Si \(f=g\), la igualtat \(f\circ x=g\circ x\) se segueix per substitució per a qualsevol \(x:X\to A\) i qualsevol \(X\).

Recíprocament, suposem que \(f\circ x=g\circ x\) per a tot \(X\) i tot \(x:X\to A\). Prenem \(X=A\) i \(x=\id_A:A\to A\). Per l'\hyperref[ax:neutralitat]{axioma~\ref*{ax:neutralitat}}, obtenim
\[
f=f\circ\id_A=g\circ\id_A=g.
\]
\end{prova}

Observeu la importància de la frase \enquote{per a tot \(X\) i tot \(x:X\to A\)} en la prova anterior. Això és la raó per la qual introduïm la definició següent.

\begin{definicio}[Element generalitzat]\index{Element generalitzat}\label{def:element-generalitzat}
Un \emph{element generalitzat de \(A\) en l'estadi \(X\)}, o un \emph{element de \(A\) parametritzat per \(X\)}, és una aplicació amb domini \(X\) i codomini \(A\). Formalment,
\[
\operatorname{ElemGen}(x,X,A)
\iff
\Dom(x,X)\land\Cod(x,A).
\]
Les majúscules \(X\) i \(A\) estan restringides als conjunts per la convenció~\ref{not:variables-conjunts-aplicacions}.
\end{definicio}

\begin{notacio}[Elements generalitzats i avaluació]\label{not:elements-generalitzats}
Escriurem
\[
x\in_X A
\quad\Longleftrightarrow\quad
x:X\to A
\]
per indicar que \(x\) és un element generalitzat de \(A\) en l'estadi \(X\).

En el capítol següent escriurem també \(x\in_X m\) quan \(m\) sigui una part de \(A\); en aquest segon ús, la notació expressarà que \(x\) es factoritza a través de \(m\). El tipus de l'expressió situada a la dreta distingirà les dues lectures.

Si \(f:A\to B\) i \(x\in_X A\), anomenarem \emph{avaluació generalitzada de \(f\) en \(x\)} la composició \(f\circ x:X\to B\). Quan sigui útil, podrem escriure
\[
f_X(x)=f\circ x.
\]
\end{notacio}

Així, la \hyperref[def:element-generalitzat]{definició~\ref*{def:element-generalitzat}} i la \hyperref[not:elements-generalitzats]{notació~\ref*{not:elements-generalitzats}} permeten enunciar la \hyperref[prop:aplicacions-iguals]{proposició~\ref*{prop:aplicacions-iguals}} d'aquesta manera: dues aplicacions paral·leles són iguals si i només si les seves avaluacions coincideixen sobre tots els elements generalitzats del domini, sigui quin sigui l'estadi.

La \hyperref[prop:aplicacions-iguals]{proposició~\ref*{prop:aplicacions-iguals}} compara aplicacions paral·leles. En el cas de dos conjunts, considerats com les aplicacions identitat \(A:A\to A\) i \(B:B\to B\), la hipòtesi que siguin paral·leles ja exigeix \(A=B\). Per tant, aquesta proposició no proporciona una noció nova per comparar conjunts diferents. La qüestió adequada és quan dos conjunts tenen la mateixa estructura, i aquesta comparació es formularà mitjançant els isomorfismes en la \hyperref[sec:isomorfismes]{secció~\ref*{sec:isomorfismes}}.

\section{Isomorfismes}\label{sec:isomorfismes}

Introduïm ara el concepte que ens permetrà afirmar que dos conjunts són estructuralment iguals.

\begin{definicio}[Isomorfisme]\index{Isomorfisme}\label{def:isomorfisme}
Una aplicació \(f:A\to B\) és un \emph{isomorfisme} si existeix una aplicació \(g:B\to A\) tal que
\[
g\circ f=\id_A
\qquad\text{i}\qquad
f\circ g=\id_B.
\]
Formalment,
\[
\begin{aligned}
\operatorname{Iso}(f)
\iff {}&
\exists A\,\exists B\,\exists g\,
\Bigl(
    \Dom(f,A)\land \Cod(f,B)
\\
&\qquad\land \Dom(g,B)\land \Cod(g,A)
\\
&\qquad\land \Com(f,g,A)
\\
&\qquad\land \Com(g,f,B)
\Bigr).
\end{aligned}
\]
Aquí les majúscules quantificades són variables restringides als conjunts. Direm que dos conjunts \(A\) i \(B\) són \emph{isomorfs} si existeix un isomorfisme entre ells. Formalment,
\[
\operatorname{Isomorfs}(A,B)
\iff
\exists f\,\bigl(
\operatorname{Iso}(f)\land\Dom(f,A)\land\Cod(f,B)
\bigr).
\]
\end{definicio}

\begin{notacio}[Isomorfia de conjunts]\label{not:isomorfia}
Si \(f:A \to B\) és un isomorfisme escriurem abreujadament \(f:A \xrightarrow{\sim} B\). Si \(\operatorname{Isomorfs}(A,B)\), escriurem
\[
A\cong B.
\]
Si no són isomorfs, escriurem \(A\ncong B\).
\end{notacio}

\begin{figure}[H]
  \centering
  \begin{tikzpicture}[scale=1.2, every node/.style={font=\small}]
    \node (A) at (0,0) {\(A\)};
    \node (B) at (3,0) {\(B\)};
    \draw[-{Stealth[scale=1.2]}, thick] (A) edge[bend left=20] node[above] {\(f\)} (B);
    \draw[-{Stealth[scale=1.2]}, thick, dashed] (B) edge[bend left=20] node[below] {\(g\)} (A);
    \draw[-{Stealth[scale=1.2]}, thick] (A) to [out=135,in=225,looseness=6] node[left] {\(\id_A\)} (A);
    \draw[-{Stealth[scale=1.2]}, thick] (B) to [out=45,in=315,looseness=6] node[right] {\(\id_B\)} (B);
  \end{tikzpicture}
  \caption{Diagrama d'un isomorfisme: anar i tornar equival a no moure's, en els dos sentits.}
  \label{fig:isomorfisme}
\end{figure}

Quan \(f\) és un isomorfisme, diem que \(g\) és una \emph{inversa de} \(f\). A la \hyperref[prop:inversa-unica]{proposició~\ref*{prop:inversa-unica}} veurem que és única. 
\begin{proposicio}[Unicitat de la inversa]\label{prop:inversa-unica}
La inversa d'un isomorfisme, si existeix, és única.
\end{proposicio}

\begin{prova}
Suposem que \(g\) i \(g'\) són inverses de \(f\colon A \to B\). La llei associativa i les lleis de neutralitat donen
\[
g = g \circ \id_B = g \circ (f \circ g') = (g \circ f) \circ g' = \id_A \circ g' = g'.
\]
\end{prova}

\begin{notacio}[Inversa d'un isomorfisme]\label{not:inversa}
Per la unicitat demostrada, si \(f\) és un isomorfisme denotarem la seva única inversa per \(f^{-1}\).
\end{notacio}

Les propietats bàsiques dels isomorfismes es dedueixen ara mecànicament dels axiomes.

\begin{proposicio}\label{prop:isos-basics}
\begin{enumerate}
\item La identitat de tot conjunt \(A\) és un isomorfisme, amb \(A^{-1}=A\); en particular, \(A\cong A\).
\item Si \(f\) és un isomorfisme, \(f^{-1}\) també ho és, amb \({(f^{-1})}^{-1} = f\); en particular, \(A \cong B\) implica \(B \cong A\).
\item Si \(f\colon A \to B\) i \(g\colon B \to C\) són isomorfismes, \(g \circ f\) també ho és, amb \({(g\circ f)}^{-1}=f^{-1}\circ g^{-1}\); en particular, \(A \cong B\) i \(B \cong C\) impliquen \(A \cong C\).
\end{enumerate}
\end{proposicio}

\begin{prova}
(1) Un conjunt és una identitat, i \(A \circ A = A\) per l'\hyperref[ax:neutralitat]{axioma~\ref*{ax:neutralitat}}, de manera que \(\id_A\) és inversa d'ella mateixa. (2) Les dues equacions de la definició són simètriques en \(f\) i \(f^{-1}\). (3) Comprovem una de les dues equacions; l'altra és anàloga:
\[
(f^{-1} \circ g^{-1}) \circ (g \circ f)
= f^{-1} \circ (g^{-1} \circ g) \circ f
= f^{-1} \circ \id_B \circ f
= f^{-1} \circ f
= \id_A,
\]
on hem usat l'\hyperref[ax:associativitat]{axioma~\ref*{ax:associativitat}} per reagrupar sense ambigüitat.
\end{prova}

\begin{observacio}\label{obs:isomorfia-equivalencia}
La \hyperref[prop:isos-basics]{proposició~\ref*{prop:isos-basics}} diu que la isomorfia \(\cong\) es comporta com una relació d'equivalència entre els conjunts: reflexiva, simètrica i transitiva. Atenció a l'estatus d'aquesta afirmació: la relació \(\cong\) no forma part de la teoria (no hi ha, encara, cap noció de relació definida); és una relació del \emph{metallenguatge}, definida sobre l'univers per la fórmula d'existència d'un isomorfisme. De fet, el paper que té la isomorfia per als conjunts és el d'\emph{igualtat estructural}: per exemple, si un conjunt és terminal, aleshores tot conjunt isomorf també és terminal i, a l'inrevés. La propietat estructural de ser terminal la veurem a la \hyperref[sec:terminals]{secció~\ref*{sec:terminals}}.
\end{observacio}


\section{Conjunts terminals}\label{sec:terminals}

Les construccions que introduirem a partir d'ara es caracteritzen
mitjançant les aplicacions que hi arriben o que en surten i les
equacions que aquestes aplicacions satisfan. Abans de presentar-ne
la primera, precisem el significat de les expressions
\emph{propietat estructural} i \emph{propietat universal}. 

\begin{observacio}[Propietats estructurals i propietats universals]
\label{obs:propietats-estructurals-universals}

Una propietat \(\mathcal P\) dels conjunts és \emph{estructural}\index{Propietat estructural}
si només depèn de les aplicacions i de les seves composicions i es
conserva en substituir un conjunt per un conjunt isomorf. Això
significa que, si \(A\cong B\), aleshores
\[
\mathcal P(A)\Longleftrightarrow\mathcal P(B).
\]

Per tant, una propietat estructural no depèn de la presentació
particular d'un conjunt, sinó de la manera com aquest conjunt es
relaciona amb la resta de conjunts mitjançant les aplicacions.

Una \emph{propietat universal}\index{Propietat universal} és una propietat que caracteritza
un conjunt, sovint juntament amb unes aplicacions distingides,
mitjançant una condició d'existència i unicitat. Habitualment,
afirma que, per a qualsevol conjunt \(X\) i per a unes aplicacions
donades, existeix una única aplicació que fa commutar les
composicions corresponents, en el sentit de la
\hyperref[not:diagrames-commutatius]{notació~\ref*{not:diagrames-commutatius}}.

Les propietats universals que considerarem són estructurals.
A més, determinen el conjunt i les seves aplicacions distingides
llevat d'un únic isomorfisme compatible amb aquestes aplicacions.

La terminalitat és el primer exemple que donarem: un conjunt \(T\) és terminal
quan, per a cada conjunt \(X\), existeix una única aplicació
\(X\to T\), és a dir, és un conjunt al qual arriba una única aplicació des de qualsevol altre conjunt.
\end{observacio}

\begin{definicio}[Conjunt terminal]\index{Conjunt terminal}\label{def:terminal}
Un conjunt \(T\) és \emph{terminal} si, per a qualsevol conjunt \(X\), existeix una única aplicació \(X\to T\). Formalment,
\[
\operatorname{Terminal}(T)
\iff
\operatorname{Conj}(T)\land
\forall X\Bigl(
  \operatorname{Conj}(X)\Rightarrow
  \exists!f\bigl(\Dom(f,X)\land \Cod(f,T)\bigr)
\Bigr).
\]
\end{definicio}

\begin{figure}[H]
  \centering
  \begin{tikzpicture}[scale=1.2, every node/.style={font=\small}]
    \node (T) at (2.2,0) {\(T\)};
    \node (X) at (0,1.6) {\(X\)};
    \node (Y) at (2.2,1.9) {\(Y\)};
    \node (Z) at (4.4,1.6) {\(Z\)};
    \draw[-{Stealth[scale=1.2]}, thick, dashed] (X) -- (T);
    \draw[-{Stealth[scale=1.2]}, thick, dashed] (Y) -- (T);
    \draw[-{Stealth[scale=1.2]}, thick, dashed] (Z) -- (T);
  \end{tikzpicture}
  \caption{Diagrama de la terminalitat: des de cada conjunt arriba exactament una aplicació al conjunt terminal.}
  \label{fig:terminal}
\end{figure}

\begin{proposicio}[Invariància de la terminalitat per isomorfisme]\label{prop:terminal-invariant-iso}
Siguin \(T\) i \(T'\) conjunts isomorfs. Si \(T\) és terminal, aleshores \(T'\) també és terminal.
\end{proposicio}

\begin{prova}
Siguin \(u:T\xrightarrow{\sim}T'\) i \(X\) un conjunt qualsevol. Com que \(T\) és terminal, existeix una única aplicació \(t:X\to T\). Per tant,
\[
u\circ t:X\to T'
\]
proporciona una aplicació de \(X\) en \(T'\).

Si \(v:X\to T'\) és qualsevol altra aplicació, aleshores \(u^{-1}\circ v:X\to T\). Per la unicitat terminal,
\[
u^{-1}\circ v=t.
\]
Component per l'esquerra amb \(u\), obtenim \(v=u\circ t\). Així, existeix una única aplicació \(X\to T'\), i \(T'\) és terminal.
\end{prova}

\begin{proposicio}\label{prop:terminal-unic}
Si dos conjunts \(T_1\) i \(T_2\) són terminals, aleshores existeix un únic isomorfisme \(T_1\to T_2\). En particular, \(T_1\cong T_2\).
\end{proposicio}

\begin{prova}
Com que \(T_2\) és terminal, existeix una única aplicació
\[
f:T_1\to T_2.
\]
Com que \(T_1\) és terminal, existeix una única aplicació
\[
g:T_2\to T_1.
\]
La composició \(g\circ f:T_1\to T_1\) coincideix amb \(\id_{T_1}\), perquè la terminalitat de \(T_1\) garanteix que només hi ha una aplicació de \(T_1\) en \(T_1\). Anàlogament,
\[
f\circ g=\id_{T_2}.
\]
Per tant, \(f\) és un isomorfisme. La seva unicitat se segueix de la terminalitat de \(T_2\).
\end{prova}

\begin{axioma}[Existència d'un conjunt terminal]\label{ax:terminal-existencia}
Existeix almenys un conjunt terminal:
\[
\exists T\,\operatorname{Terminal}(T).
\]
\end{axioma}

\begin{notacio}[Conjunt terminal escollit]\label{not:terminal-escollit}
Per l'axioma i la proposició anteriors podem fixar un conjunt terminal, que denotarem per \(\mathbf1\). Per a cada conjunt \(X\), denotarem per
\[
{!}_X:X\to\mathbf1
\]
l'única aplicació determinada per la terminalitat de \(\mathbf1\). Qualsevol propietat estructural demostrada per a \(\mathbf1\)
la satisfarà qualsevol altre conjunt terminal, perquè tots
els conjunts terminals són isomorfs.

El símbol \(\mathbf1\) no és una constant del llenguatge formal, que només té variables: quan aparegui dins d'una fórmula, com passarà en dos dels axiomes d'aquest capítol, és una abreviació eliminable. Una fórmula \(\varphi(\mathbf1)\) abreuja
\[
\forall T\,\bigl(\operatorname{Terminal}(T)\Rightarrow\varphi(T)\bigr),
\]
i, com que els terminals existeixen i són tots isomorfs, per a les propietats estructurals \(\varphi\) aquesta fórmula equival a \(\exists T\,\bigl(\operatorname{Terminal}(T)\land\varphi(T)\bigr)\). El mateix conveni s'aplicarà a \(\mathbf0\) i, en general, a tots els conjunts que fixem per una propietat universal.
\end{notacio}

\section{Conjunts inicials i no-degeneració}\label{sec:inicials}

La noció dual de conjunt terminal és la de conjunt inicial.

\begin{observacio}[Nocions duals]\label{obs:dualitat}
Dues nocions són \emph{duals} quan la definició d'una s'obté de la definició de l'altra invertint sistemàticament el sentit de totes les aplicacions.

En el llenguatge formal, la transformació dual intercanvia domini i codomini i inverteix l'ordre de la composició:
\[
\Dom\longleftrightarrow\Cod,
\qquad
\Com(f,g,h)\longleftrightarrow\Com(g,f,h).
\]
Les connectives lògiques i els quantificadors no es modifiquen. Així, la fórmula abreujada \(f:A\to B\) es transforma en la fórmula que, en l'enunciat dual, s'escriu \(f:B\to A\). Això no significa que l'aplicació \(f\) de la teoria original hagi adquirit una inversa ni que n'hàgim capgirat realment el sentit: estem transformant una fórmula del metallenguatge en una altra fórmula.

La transformació és involutiva: el dual del dual és l'enunciat original. A més, els cinc axiomes estructurals de la secció~\ref{sec:axiomes-estructurals} es transformen, eventualment després de canviar el nom de variables quantificades, en els mateixos axiomes. Per inducció sobre la longitud d'una demostració, qualsevol resultat obtingut només a partir d'aquests axiomes dona lloc a un resultat dual. Aquest és un principi metateòric sobre el nostre sistema formal, no una operació nova sobre els individus de la teoria.

La noció dual de conjunt terminal és la de conjunt inicial. En la terminalitat, des de qualsevol conjunt surt una única aplicació que arriba al terminal:
\[
X\longrightarrow T.
\]
En dualitzar-la, obtenim que des de l'inicial surt una única aplicació cap a qualsevol conjunt:
\[
I\longrightarrow X.
\]
La proposició~\ref{prop:invariancia-inicialitat} és la dual de
la proposició~\ref{prop:terminal-invariant-iso}. Igualment, la
proposició~\ref{prop:inicial-unic} sobre la unicitat dels conjunts inicials llevat d'un
únic isomorfisme és la dual de la proposició~\ref{prop:terminal-unic} corresponent per als
conjunts terminals. Mantindrem, tanmateix, les demostracions
explícites perquè el text sigui autosuficient.

El principi de dualitat s'estén als axiomes d'existència del terminal i de l'inicial, que són duals l'un de l'altre, i a l'axioma de no-degeneració, que és autodual. No s'estén, en canvi, a l'axioma de separació global que introduirem a la secció~\ref{sec:lectura-classica}: el seu dual afirmaria que dues aplicacions paral·leles \(f,g\colon A\to B\) són iguals quan \(y\circ f=y\circ g\) per a tota \(y\colon B\to\mathbf0\), i com que cap conjunt no inicial no té aplicacions cap a \(\mathbf0\), això forçaria la igualtat de totes les aplicacions paral·leles. A partir d'aquell axioma, doncs, la teoria deixa de ser simètrica, i cada resultat que en depengui s'haurà de demostrar sense apel·lar a la dualitat.\index{Dualitat}
\end{observacio}

\begin{definicio}[Conjunt inicial]\index{Conjunt inicial}\label{def:inicial}
Un conjunt \(I\) és \emph{inicial} si, per a qualsevol conjunt \(Y\), existeix una única aplicació \(I\to Y\). Formalment,
\[
\operatorname{Inicial}(I)
\iff
\operatorname{Conj}(I)\land
\forall Y\Bigl(
  \operatorname{Conj}(Y)\Rightarrow
  \exists!f\bigl(\Dom(f,I)\land \Cod(f,Y)\bigr)
\Bigr).
\]
\end{definicio}

\begin{proposicio}[Invariància de la inicialitat per isomorfisme]
\label{prop:invariancia-inicialitat}
Siguin \(I\) i \(I'\) conjunts isomorfs. Si \(I\) és inicial,
aleshores \(I'\) també és inicial.
\end{proposicio}

\begin{prova}
Sigui
\[
u:I\xrightarrow{\sim}I'
\]
un isomorfisme i sigui \(X\) un conjunt. Com que \(I\) és
inicial, existeix una única aplicació
\[
f:I\to X.
\]
Per tant,
\[
f\circ u^{-1}:I'\to X
\]
proporciona una aplicació d'\(I'\) en \(X\).

Sigui ara \(v:I'\to X\) qualsevol altra aplicació. Aleshores
\[
v\circ u:I\to X.
\]
Per la unicitat determinada per la inicialitat d'\(I\),
\[
v\circ u=f.
\]
Component per la dreta amb \(u^{-1}\), obtenim
\[
v
=v\circ\id_{I'}
=v\circ(u\circ u^{-1})
=(v\circ u)\circ u^{-1}
=f\circ u^{-1}.
\]
Així, existeix una única aplicació \(I'\to X\), i \(I'\) és
inicial.
\end{prova}

\begin{proposicio}\label{prop:inicial-unic}
Si dos conjunts \(I_1\) i \(I_2\) són inicials, aleshores existeix un únic isomorfisme \(I_1\to I_2\). En particular, \(I_1\cong I_2\).
\end{proposicio}

\begin{prova}
Com que \(I_1\) és inicial, existeix una única aplicació
\[
f:I_1\to I_2.
\]
Com que \(I_2\) és inicial, existeix una única aplicació
\[
g:I_2\to I_1.
\]
La composició \(g\circ f:I_1\to I_1\) coincideix amb \(\id_{I_1}\), perquè la inicialitat de \(I_1\) garanteix que només hi ha una aplicació de \(I_1\) en \(I_1\). Anàlogament,
\[
f\circ g=\id_{I_2}.
\]
Per tant, \(f\) és un isomorfisme. La seva unicitat se segueix de la inicialitat de \(I_1\).
\end{prova}

\begin{axioma}[Existència d'un conjunt inicial]\label{ax:inicial-existencia}
Existeix almenys un conjunt inicial:
\[
\exists I\,\operatorname{Inicial}(I).
\]
\end{axioma}

\begin{notacio}[Conjunt inicial escollit]\label{not:inicial-escollit}
Per l'axioma i la proposició anteriors podem fixar un conjunt inicial, que denotarem per \(\mathbf0\). Per a cada conjunt \(X\), denotarem per
\[
\mathbf{0}_X:\mathbf0\to X
\]
l'única aplicació determinada per la inicialitat de \(\mathbf0\). Qualsevol propietat estructural demostrada per a \(\mathbf0\)
la satisfarà qualsevol altre conjunt inicial, perquè tots
els conjunts inicials són isomorfs.
\end{notacio}

\begin{axioma}[No-degeneració]\label{ax:no-degeneracio}
Cap conjunt inicial no és isomorf a cap conjunt terminal. Formalment,
\[
\forall I\,\forall T\,\Bigl(\operatorname{Inicial}(I)\land\operatorname{Terminal}(T)\Rightarrow\neg\operatorname{Isomorfs}(I,T)\Bigr).
\]
Amb els conjunts fixats \(\mathbf0\) i \(\mathbf1\) i la notació de la isomorfia, escriurem equivalentment
\[
\mathbf0\ncong\mathbf1.
\]
\end{axioma}

Aquest axioma exclou directament la identificació estructural
entre el conjunt inicial i el conjunt terminal. La seva formulació
només utilitza la noció d'isomorfia i no pressuposa la lectura
clàssica que introduirem a la secció següent.

\begin{proposicio}\label{prop:no-fletxa-u-zero}
No existeix cap aplicació
\[
x:\mathbf1\to\mathbf0.
\]
\end{proposicio}

\begin{prova}
Suposem que existeix \(x:\mathbf1\to\mathbf0\). Per la inicialitat de \(\mathbf0\), existeix una única aplicació
\[
y:\mathbf0\to\mathbf1.
\]
La composició \(x\circ y:\mathbf0\to\mathbf0\) coincideix amb \(\id_{\mathbf0}\), perquè de \(\mathbf0\) a \(\mathbf0\) només hi ha una aplicació. De la mateixa manera, la terminalitat de \(\mathbf1\) implica
\[
y\circ x=\id_{\mathbf1}.
\]
Per tant, \(x\) és un isomorfisme, en contradicció amb l'\hyperref[ax:no-degeneracio]{axioma de no-degeneració}.
\end{prova}

\begin{observacio}[Què passaria si \(\mathbf0\cong\mathbf1\)]\label{obs:zero-u-sense-separacio}
Suposem, només en aquesta observació, que existeixen isomorfismes inversos
\[
i:\mathbf1\to\mathbf0,
\qquad
j:\mathbf0\to\mathbf1.
\]
Per a qualsevol conjunt \(A\), la composició
\[
\mathbf1\xrightarrow{i}\mathbf0\xrightarrow{\mathbf{0}_A}A
\]
proporciona una aplicació \(\mathbf1\to A\). A més, si \(x,y:\mathbf1\to A\), aleshores
\[
x\circ j=y\circ j
\]
per la inicialitat de \(\mathbf0\), i component per la dreta amb \(i\) obtenim \(x=y\). Per tant, de \(\mathbf0\cong\mathbf1\) se segueix que existeix exactament una aplicació \(\mathbf1\to A\) per a cada conjunt \(A\).

També existeix almenys una aplicació entre dos conjunts qualssevol \(A\) i \(B\), obtinguda component
\[
A\xrightarrow{{!}_A}\mathbf1
\xrightarrow{i}\mathbf0
\xrightarrow{\mathbf{0}_B}B.
\]
Tanmateix, sense cap axioma addicional no se'n dedueix que totes les aplicacions paral·leles \(A\to B\) siguin iguals. Aquesta conclusió requerirà l'axioma de separació global que introduirem a la secció següent.
\end{observacio}

\section{Lectura clàssica mitjançant elements globals}\label{sec:lectura-classica}

Observeu que el desenvolupament fet fins ara és enterament estructural, en el
sentit de l'observació~\ref{obs:propietats-estructurals-universals}:
només hi intervenen aplicacions, composicions, identitats,
isomorfismes i propietats universals.

Ara introduirem una lectura addicional que permet recuperar part
del llenguatge habitual de la teoria clàssica de conjunts.
Aquesta lectura no substituirà els elements generalitzats en el
desenvolupament estructural.

\begin{definicio}[Element global o punt]\index{Element global o punt}\label{def:element-global}
Un \emph{element global}, o simplement un \emph{element} o un \emph{punt}, d'un conjunt \(A\) és una aplicació
\[
x:\mathbf1\to A.
\]
Formalment,
\[
\operatorname{Element}(x,A)
\iff
\Dom(x,\mathbf1)\land\Cod(x,A).
\]
\end{definicio}

\begin{notacio}[Pertinença global i valor d'una aplicació]\label{not:elements-globals}
La teoria no conté una relació primitiva de pertinença. En el metallenguatge escriurem
\[
x\in A
\quad\Longleftrightarrow\quad
x:\mathbf1\to A.
\]
Així, \(x\in A\) és el cas particular \(x\in_{\mathbf1}A\).

Si \(f:A\to B\) i \(x\in A\), escriurem
\[
f(x)=f\circ x
\]
per al valor de \(f\) en \(x\). Llavors \(f(x)\in B\). Si, a més, \(g:B\to C\), l'associativitat dona la regla habitual
\[
(g\circ f)(x)=g(f(x)).
\]
\end{notacio}

La primera igualtat es representa mitjançant el triangle següent:
\begin{figure}[H]
  \centering
  \begin{tikzpicture}[scale=1.2, every node/.style={font=\small}]
    \node (One) at (0,0) {\(\mathbf{1}\)};
    \node (A) at (2,1.5) {\(A\)};
    \node (B) at (4,0) {\(B\)};
    \draw[-{Stealth[scale=1.2]}, thick] (One) -- node[above left] {\(x\)} (A);
    \draw[-{Stealth[scale=1.2]}, thick] (A) -- node[above right] {\(f\)} (B);
    \draw[-{Stealth[scale=1.2]}, thick] (One) -- node[below] {\(f\circ x=f(x)\)} (B);
  \end{tikzpicture}
  \caption{El valor de \(f\) en \(x\) és una composició.}
  \label{fig:avaluacio_triangle}
\end{figure}

La compatibilitat amb la composició queda reflectida en el diagrama següent:
\begin{figure}[H]
  \centering
  \begin{tikzpicture}[scale=1.2, every node/.style={font=\small}]
    \node (One) at (0,0) {\(\mathbf{1}\)};
    \node (A) at (2,0) {\(A\)};
    \node (B) at (4,0) {\(B\)};
    \node (C) at (6,0) {\(C\)};
    \draw[-{Stealth[scale=1.2]}, thick] (One) -- node[above] {\(x\)} (A);
    \draw[-{Stealth[scale=1.2]}, thick] (A) -- node[above] {\(f\)} (B);
    \draw[-{Stealth[scale=1.2]}, thick] (B) -- node[above] {\(g\)} (C);
    \draw[-{Stealth[scale=1.2]}, thick] (One) to[bend left=35] node[above] {\((g\circ f)(x)\)} (C);
    \draw[-{Stealth[scale=1.2]}, thick] (One) to[bend right=35] node[below] {\(g(f(x))\)} (C);
  \end{tikzpicture}
  \caption{La regla \((g\circ f)(x)=g(f(x))\) és una conseqüència de l'associativitat.}
  \label{fig:associativitat_avaluacio_diagrama}
\end{figure}

\begin{definicio}[Aplicació constant]\index{Aplicació constant}\label{def:constant}
Una aplicació \(f:A\to B\) és \emph{constant} si existeix un element global \(b:\mathbf1\to B\) tal que
\[
f=b\circ {!}_A.
\]
Formalment,
\[
\begin{aligned}
\operatorname{Const}(f)
\iff {}&
\exists A\,\exists B\,\exists u\,\exists b\,
\Bigl(
    \Dom(f,A)\land \Cod(f,B)
\\
&\qquad\land \Dom(u,A)\land \Cod(u,\mathbf1)
\\
&\qquad\land \Dom(b,\mathbf1)\land \Cod(b,B)
\\
&\qquad\land \Com(u,b,f)
\Bigr).
\end{aligned}
\]
\end{definicio}

\begin{figure}[H]
  \centering
\begin{tikzpicture}[scale=1.1, every node/.style={font=\small}]
  \node (A) at (0,1.6) {\(A\)};
  \node (B) at (3.2,1.6) {\(B\)};
  \node (C) at (3.2,0) {\(C\)};
  \draw[-{Stealth[scale=1.15]}, thick] (A) -- node[above] {\(f\)} (B);
  \draw[-{Stealth[scale=1.15]}, thick] (C) -- node[right] {\(b\)} (B);
  \draw[-{Stealth[scale=1.15]}, thick] (A) -- node[below left] {\({!}_A\)} (C);
\end{tikzpicture}
\caption{La factorització per \(\mathbf1\) de l'aplicació constant: \(f=b\circ {!}_A\).}
  \label{fig:aplicacio-constant}
\end{figure}

\begin{observacio}[La convenció sobre les aplicacions constants]\label{obs:constants-domini-buit}
La definició anterior incorpora explícitament un valor \(b\in B\). Per tant, una aplicació amb domini \(\mathbf0\) és constant en aquest sentit només si el codomini té almenys un element global que permeti la factorització per \(\mathbf1\). En particular, \(\id_{\mathbf0}:\mathbf0\to\mathbf0\) no és constant segons aquesta convenció, perquè no existeix cap aplicació \(\mathbf1\to\mathbf0\).

Aquesta opció difereix de la convenció puntual de la teoria clàssica segons la qual tota funció amb domini buit és constant per vacuïtat. Adoptarem la definició estructural per factorització, i assenyalarem explícitament el cas buit quan sigui rellevant.
\end{observacio}

\begin{observacio}[Els punts de \(\mathbf1\) i de \(\mathbf0\)]\label{obs:punts-zero-u}
El conjunt terminal \(\mathbf1\) té exactament un element global: \(\id_{\mathbf1}\). En efecte, tota aplicació \(\mathbf1\to\mathbf1\) coincideix amb aquesta identitat per la terminalitat.

En canvi, la \hyperref[prop:no-fletxa-u-zero]{proposició~\ref*{prop:no-fletxa-u-zero}} afirma que \(\mathbf0\) no té cap element global. Tot i això, \(\mathbf0\) té un element generalitzat en l'estadi \(\mathbf0\), que és
\[
\id_{\mathbf0}:\mathbf0\to\mathbf0.
\]
Aquesta diferència mostra que els elements generalitzats contenen més informació estructural que els punts globals.
\end{observacio}

La \hyperref[prop:aplicacions-iguals]{proposició~\ref*{prop:aplicacions-iguals}} garanteix que la totalitat dels elements generalitzats separa sempre les aplicacions. No se segueix, en canvi, dels axiomes estructurals que els elements globals siguin suficients per distingir-les. Per recuperar el principi d'extensionalitat global d'aplicacions, habitual de la teoria clàssica de conjunts, introduïm l'axioma addicional següent.

\begin{axioma}[Separació d'aplicacions per elements globals]\label{ax:separacio-globals}
Dues aplicacions paral·leles que coincideixen sobre tots els elements globals són iguals. Formalment,
\[
\begin{aligned}
\forall f\,\forall g\,\forall A\,\forall B\,\Bigl(&
\Dom(f,A)\land\Cod(f,B)\land
\Dom(g,A)\land\Cod(g,B)
\\
&{}\land
\forall T\,\forall x\,\forall y\,\forall z\,
\Bigl(
\operatorname{Terminal}(T)\land\Dom(x,T)\land\Cod(x,A)
\\
&\qquad{}\land\Com(x,f,y)\land\Com(x,g,z)
\Rightarrow y=z
\Bigr)
\Rightarrow f=g
\Bigr).
\end{aligned}
\]
Amb els convenis de notació, aquesta fórmula s'escriu
\[
\left[
\forall x:\mathbf1\to A,\quad f\circ x=g\circ x
\right]
\Longrightarrow f=g.
\]
Equivalentment,
\[
f\ne g
\quad\Longrightarrow\quad
\exists x:\mathbf1\to A
\text{ tal que }
f\circ x\ne g\circ x.
\]
\end{axioma}

Aquest axioma és addicional: no és necessari per definir els conjunts, els isomorfismes, els conjunts inicials i terminals ni les construccions universals del capítol següent. La seva funció és assegurar que els elements globals determinen les aplicacions, tal com ocorre en la teoria clàssica de conjunts. En la terminologia de Lawvere, l'axioma diu que \(\mathbf1\) és un \emph{separador} i que la teoria és \emph{ben apuntada}\index{Ben apuntat}\index{Separador}: els punts globals separen les aplicacions. És un dels dos axiomes ---l'altre és l'elecció--- que distingeixen els conjunts constants de les altres estructures que satisfan la resta d'axiomes del llibre, com veurem al capítol~\ref{cap:models}.

\begin{proposicio}[Extensionalitat global de les aplicacions]\label{prop:extensionalitat-global}
Siguin \(f,g:A\to B\). Aleshores
\[
f=g
\quad\Longleftrightarrow\quad
\forall x:\mathbf1\to A,\quad f(x)=g(x).
\]
\end{proposicio}

\begin{prova}
Si \(f=g\), la igualtat dels valors se segueix per substitució. La implicació recíproca és exactament l'\hyperref[ax:separacio-globals]{axioma de separació d'aplicacions per elements globals}.
\end{prova}

\begin{observacio}[Separació global i no-degeneració]
\label{obs:separacio-no-degeneracio}

L'axioma de no-degeneració i l'axioma de separació global tenen
funcions diferents.

L'axioma de no-degeneració afirma directament que
\[
\mathbf0\ncong\mathbf1.
\]
Per tant, exclou que el conjunt inicial i el conjunt terminal
representin la mateixa estructura.

En canvi, l'axioma de separació global afirma que les aplicacions
amb domini \(\mathbf1\) són suficients per distingir dues
aplicacions paral·leles. Si \(f,g:A\to B\) i
\[
f\circ x=g\circ x
\]
per a tot \(x:\mathbf1\to A\), aleshores \(f=g\).

La separació global no substitueix la no-degeneració. En efecte,
suposem que
\[
\mathbf0\cong\mathbf1.
\]
L'\hyperref[obs:zero-u-sense-separacio]
{observació~\ref*{obs:zero-u-sense-separacio}} mostra que, en
aquest cas, per a cada conjunt \(A\) existiria exactament una
aplicació
\[
x_A:\mathbf1\to A.
\]
També existiria almenys una aplicació entre cada parell de
conjunts \(A\) i \(B\), obtinguda mitjançant la composició
\[
A\xrightarrow{{!}_A}\mathbf1
 \longrightarrow\mathbf0
 \xrightarrow{\mathbf{0}_B}B.
\]

Sense la separació global, d'això no se segueix que dues
aplicacions paral·leles \(f,g:A\to B\) hagin de ser iguals.

Si hi afegim la separació global, però, totes dues aplicacions
coincideixen sobre l'únic punt \(x_A\), perquè
\[
f\circ x_A
\quad\text{i}\quad
g\circ x_A
\]
són aplicacions \(\mathbf1\to B\), i només n'hi ha una. Per
separació global, \(f=g\).

Per tant, si s'acceptessin simultàniament
\(\mathbf0\cong\mathbf1\) i la separació global, entre cada parell
de conjunts existiria exactament una aplicació. Aquest és el
col·lapse estructural que l'axioma de no-degeneració exclou.
\end{observacio}

\begin{observacio}[El domini d'individus i la qüestió de mida]\label{obs:mida-univers}
Com en la teoria estructural de conjunts de Lawvere, la nostra teoria es defineix pels seus axiomes i no per referència a un domini de «tots els conjunts»: els conjunts són certes aplicacions, i la seva totalitat no és, ella mateixa, un objecte de la teoria. La qüestió de mida no es planteja, doncs, dins del sistema; apareix només quan volem \emph{interpretar-lo} en un model conjuntista, com farem tot seguit amb la interpretació clàssica.

En aquesta interpretació, els individus són les funcions entre conjunts, i aquí sí que cal una precaució: la col·lecció de totes les funcions entre tots els conjunts no és, en la teoria clàssica \textsf{ZFC}, un conjunt, sinó una totalitat massa gran per ser-ho. Per disposar d'un domini que sigui un conjunt honest ---i, per tant, d'una estructura de primer ordre legítima--- adoptem a la metateoria un únic axioma addicional, a la manera de Mac Lane~\cite{maclane-moerdijk,shulman-sets}: existeix un \emph{univers} \(\mathcal U\), un conjunt tancat per les operacions conjuntistes usuals i que conté un conjunt infinit. Els seus elements s'anomenen \emph{conjunts petits}, i la interpretació pren com a individus les funcions entre conjunts petits, que formen un subconjunt de \(\mathcal U\). L'existència d'un tal univers equival a la d'un cardinal inaccessible i no se segueix de \textsf{ZFC} sol; és un supòsit feble i estàndard, i n'hi ha prou amb un de sol per a tot el desenvolupament.

Podríem, alternativament, interpretar la teoria en una metateoria amb classes pròpies ---\textsf{NBG} o \textsf{ZFC} amb classes--- i prendre com a domini la classe de totes les funcions. Aquesta via és consistent, però importa la distinció conjunt/classe, basada en la pertinença, que el punt de vista estructural precisament evita: la noció de «classe pròpia» no té cap correlat intern en una teoria on els conjunts són identitats i no col·leccions. Preferim, doncs, la relativització a un univers, que expressa la mida dins de l'ontologia de conjunts, sense recórrer a cap noció primitiva de classe.

Convé subratllar que l'univers \(\mathcal U\) viu a la metateoria, no al llenguatge: dins la teoria segueix havent-hi només aplicacions i els predicats \(\Dom\), \(\Cod\) i \(\Com\), i no s'hi afegeix cap primitiva nova. L'univers és un artefacte del \emph{model}, no un objecte de la teoria. La contrapartida \emph{interna} de la mida ---la diferència entre les aplicacions entre dos conjunts fixats, que sí que formen un objecte, i la totalitat de les aplicacions, que no en forma cap--- apareixerà més endavant amb els \emph{exponencials}: l'objecte \(B^A\) de les aplicacions de \(A\) a \(B\) existeix per a cada parell de conjunts, mentre que no hi ha cap objecte de «totes les aplicacions». Aquesta serà la manifestació, dins la teoria, de la mateixa distinció que aquí expressem a la metateoria.
\end{observacio}

\begin{observacio}[Interpretació clàssica del conjunt terminal]\label{obs:model-classic-terminal}
Considerem la interpretació clàssica en què els individus són les funcions entre conjunts petits, en el sentit precisat a l'observació~\ref{obs:mida-univers}. Si \(f:S\to T\) és una funció clàssica, el seu domini i el seu codomini en la nostra teoria són les funcions identitat
\[
\id_S:S\to S,
\qquad
\id_T:T\to T,
\]
i \(\Com(f,g,h)\) expressa que \(h=g\circ f\) en el sentit habitual. Per tant, els conjunts de la nostra teoria es realitzen com les funcions identitat dels conjunts clàssics. Quan no hi hagi risc de confusió, identificarem metalingüísticament \(\id_S\) amb el conjunt clàssic subjacent \(S\).

Sota aquesta interpretació, els conjunts terminals són exactament les identitats dels conjunts unitaris. En efecte, si \(S=\{*\}\), per a qualsevol conjunt \(A\) existeix una única funció \(A\to S\). Recíprocament, si \(\id_T\) és terminal, hi ha una única funció \(\{*\}\to T\); com que aquestes funcions es corresponen amb els elements de \(T\), el conjunt \(T\) té exactament un element.

Fixat un conjunt unitari \(S=\{*\}\), per a qualsevol conjunt clàssic \(A\) hi ha una bijecció entre les funcions
\[
x:S\to A
\]
i els elements de \(A\): la funció \(x\) determina \(x(*)\), i cada \(a\in A\) determina l'única funció \(x_a:S\to A\) tal que \(x_a(*)=a\). Per això, una aplicació \(x:\mathbf1\to A\) representa, en la interpretació clàssica, un element del conjunt subjacent a \(A\).

Aquesta correspondència és externa. Els individus de la teoria continuen essent aplicacions, i \(\mathbf1\) és formalment una aplicació identitat. L'escriptura
\[
\mathbf1\in\mathbf1
\]
pertany a la notació metalingüística introduïda en aquest capítol: afirma que \(\id_{\mathbf1}:\mathbf1\to\mathbf1\) és l'únic element global de \(\mathbf1\); no introdueix una relació interna de pertinença.

Aquesta interpretació recupera la imatge ingènua d'un conjunt com una col·lecció, o un \enquote{sac}, de punts. Però és només una imatge externa. Dins la teoria, un conjunt es defineix com una aplicació identitat \(A:A\to A\), i els seus punts globals són aplicacions amb domini \(\mathbf1\) i codomini \(A\). El desenvolupament estructural treballarà també amb els elements generalitzats \(X\to A\), per a qualsevol estadi \(X\). L'axioma de separació global permet recuperar una altra característica de la teoria clàssica: una aplicació queda determinada pels valors que pren sobre tots els punts globals del seu domini.
\end{observacio}

\begin{resumcap}
\apartat{Idees centrals}
\begin{enumerate}
  \item El llenguatge té una sola mena d'individus, les aplicacions, i tres predicats primitius, \(\Dom\), \(\Cod\) i \(\Com\). Tota la notació habitual (\(f\colon A\to B\), \(g\circ f\), \(\id_A\), els diagrames) és una abreviació eliminable.
  \item Els cinc axiomes estructurals regulen el domini, el codomini i la composició; d'ells es dedueixen l'estabilitat del domini i del codomini i la unicitat de la composició.
  \item Un conjunt és una aplicació que és el seu propi domini i codomini, és a dir, una identitat. Els conjunts no es construeixen a partir d'elements: els elements, globals o generalitzats, són aplicacions cap al conjunt.
  \item Dues aplicacions paral·leles són iguals si i només si coincideixen sobre tots els elements generalitzats del domini; que n'hi hagi prou amb els elements globals és un axioma a part, la separació global, que fa la teoria ben apuntada.
  \item Un isomorfisme és una aplicació amb inversa, que és única. La isomorfia és la igualtat estructural: les propietats que estudiem són invariants per isomorfisme.
  \item Les propietats universals ---la terminalitat és la primera--- caracteritzen un conjunt per les aplicacions que hi arriben o que en surten, i el determinen llevat d'un únic isomorfisme.
  \item El principi de dualitat intercanvia domini i codomini i inverteix la composició; val per als axiomes estructurals i per als d'existència del terminal i de l'inicial, però no per a la separació global.
  \item Els axiomes de no-degeneració i de separació global tenen funcions diferents: el primer exclou \(\mathbf0\cong\mathbf1\); el segon diu que els punts globals determinen les aplicacions. Junts, exclouen el col·lapse en què entre cada parell de conjunts hi hauria una sola aplicació.
\end{enumerate}
\apartat{Errors típics} Pensar que \(\id_A\) és una aplicació diferent del conjunt \(A\) (són el mateix individu); confondre \enquote{\(A\) és un conjunt} amb \enquote{\(A\) té elements globals} (el conjunt \(\mathbf0\) no en té cap i és un conjunt); creure que la unicitat de la identitat vol dir que només hi ha una aplicació \(A\to A\); oblidar que \(g\circ f\) només està definida quan \(\cod(f)=\dom(g)\); llegir \(x\in A\) com una relació primitiva, quan és l'abreviació de \(x\colon\mathbf1\to A\); confondre el criteri d'igualtat per elements generalitzats, que és un teorema, amb la separació per elements globals, que és un axioma; pensar que una propietat universal determina un conjunt \emph{únic}, quan el determina llevat d'un únic isomorfisme; i aplicar el principi de dualitat a un resultat que usa la separació global.
\apartat{Lectura recomanada} \cite[cap.~1]{lawvere-rosebrugh} per a la presentació original dels conjunts abstractes, les aplicacions, la composició i els elements generalitzats, de la qual aquest capítol és una versió formalitzada; \cite[sessions~1--5]{lawvere-schanuel} per a una introducció molt elemental a la composició, les identitats, els isomorfismes i les seccions; \cite[pàg.~10--11]{bell2006} per a l'explicació filosòfica de per què els elements d'un conjunt abstracte han de ser aplicacions \(\mathbf1\to X\) i de \(\mathbf1\) com a separador; \cite{lawvere-etcs-reprint} per als axiomes originals d'ETCS, amb la introducció de McLarty; i, com a contrast, \cite[§§1--3]{halmos} per al punt de partida de la teoria clàssica, on els primers axiomes parlen d'elements i de pertinença.
\end{resumcap}


\chapter{Parts i construccions universals}\label{cap:parts}

\begin{objectius}
\apartat{Prerequisits} El capítol~\ref{cap:aplicacions} complet: el llenguatge formal i els convenis de notació, la composició i les lleis de neutralitat, els elements generalitzats i el criteri d'igualtat, els isomorfismes, el conjunt terminal i l'inicial, i l'axioma de separació global, que aquí només s'usa a les observacions de lectura clàssica. De la teoria clàssica de conjunts, n'hi ha prou amb la familiaritat informal amb subconjunts, funcions característiques, unions disjuntes i productes cartesians, que es reconstrueixen al capítol.
\apartat{Objectius} En acabar el capítol, el lector hauria de poder:
\begin{itemize}
  \item definir monomorfisme i epimorfisme per cancel·lació, i demostrar-ne les propietats bàsiques: identitats, composicions, isomorfismes, seccions i retraccions;
  \item explicar per què una part d'un conjunt és un monomorfisme, i no un conjunt, i treballar amb la inclusió i la pertinença generalitzada com a factoritzacions;
  \item reconèixer quan dues parts són equivalents i com recuperen, en el metallenguatge, la noció de subconjunt;
  \item enunciar la propietat universal d'una suma i d'un producte, escriure l'aplicació definida per casos i l'aparellament, i demostrar la unicitat llevat d'un únic isomorfisme;
  \item usar l'axioma de representació: associar a cada part la seva característica cap a \(\mathbf2\), i deduir-ne que els bimorfismes són isomorfismes;
  \item construir equalitzadors i pullbacks a partir dels productes i de la representació, sense cap axioma nou, i deduir-ne la imatge inversa i la intersecció de parts;
  \item situar la injectivitat, l'exhaustivitat i la bijectivitat clàssiques respecte dels monomorfismes, els epimorfismes i els isomorfismes, i saber quines equivalències depenen de la separació global.
\end{itemize}
\end{objectius}

Al \hyperref[cap:aplicacions]{capítol anterior} hem fixat el llenguatge de la teoria i hem introduït les nocions d'aplicació, conjunt, composició, isomorfisme, conjunt terminal i conjunt inicial. També hem establert els convenis de notació i la representació geomètrica del metallenguatge mitjançant diagrames d'aplicacions. En aquest capítol començarem a construir nous conjunts a partir d'altres conjunts.

El principi que guiarà totes les construccions serà el mateix. No descriurem un conjunt nou indicant de quins elements globals està format, sinó especificant les aplicacions que hi arriben o que en surten i les equacions que aquestes aplicacions han de satisfer. Les propietats que caracteritzen d'aquesta manera una construcció s'anomenen \emph{propietats universals}.

Abans d'introduir-les, necessitem dues classes d'aplicacions definides exclusivament mitjançant la composició: els monomorfismes i els epimorfismes. Els monomorfismes ens permetran definir les parts d'un conjunt sense recórrer als elements globals. Quan vulguem relacionar aquestes nocions estructurals amb la terminologia habitual de la teoria clàssica de conjunts, ho farem en observacions separades i usarem l'\hyperref[ax:separacio-globals]{axioma addicional de separació global}. Aquestes observacions no intervindran en el desenvolupament estructural del capítol.

\section{Monomorfismes i epimorfismes}\label{sec:monomorfismes-epimorfismes}

\begin{definicio}[Monomorfisme]\index{Monomorfisme}\label{def:monomorfisme}
Una aplicació \(m\) és un \emph{monomorfisme} si es pot cancel·lar per l'esquerra. Formalment,
\[
\begin{aligned}
\operatorname{Mono}(m)
\iff {}&
\exists A\,\exists B\,\Bigl(
\Dom(m,A)\land\Cod(m,B)
\\
&{}\land
\forall X\,\forall u\,\forall v\,
\Bigl(
\Dom(u,X)\land\Cod(u,A)
\\
&\qquad{}\land\Dom(v,X)\land\Cod(v,A)
\\
&\qquad{}\land\exists h\,
\bigl(\Com(u,m,h)\land\Com(v,m,h)\bigr)
\Rightarrow u=v
\Bigr)
\Bigr).
\end{aligned}
\]
\end{definicio}

\begin{notacio}[Fletxa monomòrfica]\label{not:mono}
Si \(m:A\to B\) i \(\operatorname{Mono}(m)\), escriurem
\[
m:A\rightarrowtail B.
\]
Amb aquesta notació, la condició de monomorfisme s'abreuja així: per a totes les aplicacions paral·leles \(u,v:X\to A\),
\[
m\circ u=m\circ v\Longrightarrow u=v.
\]
\end{notacio}

La situació descrita en la definició es representa mitjançant el diagrama extern

\begin{center}
	\begin{tikzpicture}[
		scale=1.1,
		every node/.style={font=\small},
		>=stealth
		]
		
		\node (X) at (0,0) {$X$};
		\node (A) at (3.4,0) {$A$};
		\node (B) at (6.8,0) {$B$};
		
		\draw[->,thick]
		([xshift=4pt,yshift=2pt]X.east) --
		node[above] {$u$}
		([xshift=-4pt,yshift=2pt]A.west);
		
		\draw[->,thick]
		([xshift=4pt,yshift=-2pt]X.east) --
		node[below] {$v$}
		([xshift=-4pt,yshift=-2pt]A.west);
		
		\draw[->,thick]
		([xshift=4pt]A.east) --
		node[above] {$m$}
		([xshift=-4pt]B.west);
		
	\end{tikzpicture}
\end{center}

La hipòtesi afirma que els dos camins de \(X\) a \(B\) coincideixen; la conclusió afirma que les dues aplicacions de \(X\) a \(A\) ja eren iguals. Equivalentment, pel contrarecíproc, tenim
\[
u\ne v\Longrightarrow m\circ u\ne m\circ v.
\]
Així, un monomorfisme no identifica dos elements generalitzats diferents de \(A\) que tinguin el mateix estadi \(X\).

\begin{definicio}[Epimorfisme]\index{Epimorfisme}\label{def:epimorfisme}
Una aplicació \(e\) és un \emph{epimorfisme} si es pot cancel·lar per la dreta. Formalment,
\[
\begin{aligned}
\operatorname{Epi}(e)
\iff {}&
\exists A\,\exists B\,\Bigl(
\Dom(e,A)\land\Cod(e,B)
\\
&{}\land
\forall Y\,\forall u\,\forall v\,
\Bigl(
\Dom(u,B)\land\Cod(u,Y)
\\
&\qquad{}\land\Dom(v,B)\land\Cod(v,Y)
\\
&\qquad{}\land\exists h\,
\bigl(\Com(e,u,h)\land\Com(e,v,h)\bigr)
\Rightarrow u=v
\Bigr)
\Bigr).
\end{aligned}
\]
\end{definicio}

\begin{notacio}[Fletxa epimòrfica]\label{not:epi}
Si \(e:A\to B\) i \(\operatorname{Epi}(e)\), escriurem
\[
e:A\twoheadrightarrow B.
\]
Amb aquesta notació, la condició d'epimorfisme s'abreuja així: per a totes les aplicacions paral·leles \(u,v:B\to Y\),
\[
u\circ e=v\circ e\Longrightarrow u=v.
\]
\end{notacio}

La situació dual de l'anterior és
\begin{center}
	\begin{tikzpicture}[
		scale=1.1,
		every node/.style={font=\small},
		>=stealth
		]
		
		\node (A) at (0,0) {$A$};
		\node (B) at (3.4,0) {$B$};
		\node (Y) at (6.8,0) {$Y$};
		
		\draw[->,thick]
		([xshift=4pt]A.east) --
		node[above] {$e$}
		([xshift=-4pt]B.west);
		
		\draw[->,thick]
		([xshift=4pt,yshift=2pt]B.east) --
		node[above] {$u$}
		([xshift=-4pt,yshift=2pt]Y.west);
		
		\draw[->,thick]
		([xshift=4pt,yshift=-2pt]B.east) --
		node[below] {$v$}
		([xshift=-4pt,yshift=-2pt]Y.west);
		
	\end{tikzpicture}
\end{center}

En aquest cas es comparen dues aplicacions que tenen el mateix domini \(B\) i, dualment, també el mateix codomini \(Y\). Si els dos camins d'\(A\) a \(Y\) coincideixen després de precompondre \(u\) i \(v\) amb \(e\), aleshores \(u=v\). Així, un epimorfisme determina les aplicacions que surten de \(B\) cap a un codomini fix.

\begin{definicio}[Bimorfisme]\index{Bimorfisme}\label{def:bimorfisme}
Una aplicació és un \emph{bimorfisme} si és alhora un monomorfisme i un epimorfisme. Formalment,
\[
\operatorname{Bimor}(f)
\iff
\operatorname{Mono}(f)\land\operatorname{Epi}(f).
\]
\end{definicio}

A la \hyperref[prop:bimorfisme-isomorfisme]{proposició~\ref*{prop:bimorfisme-isomorfisme}} demostrarem que, amb els axiomes d'aquest capítol, els bimorfismes són precisament els isomorfismes. Fins que no s'hagi demostrat aquest resultat, no identificarem les dues nocions.

\begin{proposicio}\label{prop:identitat-mono-epi}
Per a qualsevol conjunt \(A\), l'aplicació identitat
\[
\id_A:A\to A
\]
és un monomorfisme i un epimorfisme.
\end{proposicio}

\begin{prova}
Si \(\id_A\circ u=\id_A\circ v\), les lleis de neutralitat donen \(u=v\); per tant, \(\id_A\) és un monomorfisme. De la mateixa manera, si \(u\circ\id_A=v\circ\id_A\), obtenim \(u=v\), i \(\id_A\) és un epimorfisme.
\end{prova}

\begin{proposicio}\label{prop:composicio-mono-epi}
La composició de monomorfismes és un monomorfisme, i la composició d'epimorfismes és un epimorfisme.
\end{proposicio}

\begin{prova}
Siguin
\[
A\xrightarrow{m}B\xrightarrow{n}C
\]
dos monomorfismes. Si
\[
(n\circ m)\circ u=(n\circ m)\circ v,
\]
per l'associativitat tenim
\[
n\circ(m\circ u)=n\circ(m\circ v).
\]
Com que \(n\) és un monomorfisme, \(m\circ u=m\circ v\); i com que \(m\) també és un monomorfisme, \(u=v\). Per tant, \(n\circ m\) és un monomorfisme.

El cas dels epimorfismes és dual. Si \(e:A\to B\) i \(d:B\to C\) són epimorfismes i
\[
u\circ(d\circ e)=v\circ(d\circ e),
\]
l'associativitat dona
\[
(u\circ d)\circ e=(v\circ d)\circ e.
\]
Com que \(e\) és un epimorfisme, \(u\circ d=v\circ d\), i com que \(d\) és un epimorfisme, \(u=v\).
\end{prova}

\begin{proposicio}\label{prop:isomorfisme-mono-epi}
Tot isomorfisme és un monomorfisme i un epimorfisme.
\end{proposicio}

\begin{prova}
Sigui \(f:A\to B\) un isomorfisme i sigui \(f^{-1}:B\to A\) la seva inversa. Si \(f\circ u=f\circ v\), component per l'esquerra amb \(f^{-1}\) obtenim
\[
u=(f^{-1}\circ f)\circ u=(f^{-1}\circ f)\circ v=v.
\]
Per tant, \(f\) és un monomorfisme.

Si \(u\circ f=v\circ f\), component per la dreta amb \(f^{-1}\) obtenim
\[
u=u\circ(f\circ f^{-1})=v\circ(f\circ f^{-1})=v.
\]
Per tant, \(f\) és un epimorfisme.
\end{prova}

\begin{definicio}[Secció i retracció]\index{Secció i retracció}\label{def:seccio-retraccio}
Siguin \(s:A\to B\) i \(r:B\to A\). Direm que \(s\) és una \emph{secció} de \(r\), i que \(r\) és una \emph{retracció} de \(s\), si
\[
r\circ s=\id_A.
\]
Formalment, definim primer la relació entre les dues aplicacions:
\[
\begin{aligned}
\operatorname{SeccioDe}(s,r)
\iff
\exists A\,\exists B\,\Bigl(&
\Dom(s,A)\land\Cod(s,B)
\\
&{}\land\Dom(r,B)\land\Cod(r,A)
\land\Com(s,r,A)
\Bigr).
\end{aligned}
\]
A partir d'aquesta relació, les dues propietats unàries s'abreugen per
\[
\operatorname{Seccio}(s)
\iff
\exists r\,\operatorname{SeccioDe}(s,r),
\qquad
\operatorname{Retraccio}(r)
\iff
\exists s\,\operatorname{SeccioDe}(s,r).
\]
\end{definicio}

\begin{notacio}[Triangle de secció i retracció]\label{not:seccio-retraccio}
La condició \(r\circ s=\id_A\) es representa mitjançant el triangle commutatiu
\begin{center}
\begin{tikzpicture}[scale=1.1, every node/.style={font=\small}]
  \node (A1) at (0,1.7) {\(A\)};
  \node (B) at (3.3,1.7) {\(B\)};
  \node (A2) at (3.3,0) {\(A\)};
  \draw[-{Stealth[scale=1.15]}, thick] (A1) -- node[above] {\(s\)} (B);
  \draw[-{Stealth[scale=1.15]}, thick] (B) -- node[right] {\(r\)} (A2);
  \draw[-{Stealth[scale=1.15]}, thick] (A1) -- node[below left] {\(\id_A\)} (A2);
\end{tikzpicture}
\end{center}
\end{notacio}

\begin{proposicio}\label{prop:seccio-mono-retraccio-epi}
Tota secció és un monomorfisme i tota retracció és un epimorfisme.
\end{proposicio}

\begin{prova}
Suposem que \(r\circ s=\id_A\). Si \(s\circ u=s\circ v\), component per l'esquerra amb \(r\) obtenim
\[
u=(r\circ s)\circ u=(r\circ s)\circ v=v.
\]
Així, \(s\) és un monomorfisme.

Si \(u\circ r=v\circ r\), component per la dreta amb \(s\) obtenim
\[
u=u\circ(r\circ s)=v\circ(r\circ s)=v.
\]
Per tant, \(r\) és un epimorfisme.
\end{prova}

\begin{observacio}[Monomorfismes i epimorfismes escindits]\index{Monomorfisme escindit}\index{Epimorfisme escindit}\label{obs:escindits}
Siguin \(f:A\to B\) i \(g:B\to A\) tals que
\[
g\circ f=\id_A.
\]
Aleshores \(f\) és una secció de \(g\), i \(g\) és una retracció de \(f\). També direm que \(f\) és un \emph{monomorfisme escindit} i que \(g\) és un \emph{epimorfisme escindit}. Aquesta situació és intermèdia entre una aplicació qualsevol i un isomorfisme, perquè només s'exigeix una de les dues identitats d'inversió.

En la lectura clàssica, \(f\) insereix \(A\) en \(B\), i \(g\) permet recuperar cada element de \(A\) després d'aquesta inserció.

\begin{figure}[H]
  \centering
  \begin{tikzpicture}[scale=1.2, every node/.style={font=\small}]
    \node (A) at (0,0) {\(A\)};
    \node (B) at (3,0) {\(B\)};
    \draw[-{Stealth[scale=1.2]}, thick] (A) edge[bend left=20] node[above] {\(f\)} (B);
    \draw[-{Stealth[scale=1.2]}, thick] (B) edge[bend left=20] node[below] {\(g\)} (A);
    \draw[-{Stealth[scale=1.2]}, thick] (A) to [out=135,in=225,looseness=6] node[left] {\(\id_A\)} (A);
  \end{tikzpicture}
  \caption{Diagrama d'una secció i d'una retracció: la composició $g\circ f$ és la identitat d'$A$, però no s'exigeix que $f\circ g=\id_B$.}
  \label{fig:seccio-retraccio}
\end{figure}

\end{observacio}

Els epimorfismes escindits admeten una caracterització en termes d'elements generalitzats, que reforça la noció d'exhaustivitat i que serà el punt de partida de l'estudi de l'elecció.

\begin{definicio}[Aplicació exhaustiva sobre elements generalitzats]\index{Aplicació exhaustiva sobre elements generalitzats}\label{def:exhaustiva-generalitzada}
Una aplicació \(f\) és \emph{exhaustiva sobre elements generalitzats} si tot element generalitzat del seu codomini es factoritza a través de \(f\). Formalment,
\[
\begin{aligned}
\operatorname{ExhGen}(f)
\iff {}&
\exists A\,\exists B\,\Bigl(
\Dom(f,A)\land\Cod(f,B)
\\
&{}\land
\forall X\,\forall y\,\Bigl(
\Dom(y,X)\land\Cod(y,B)
\\
&\qquad{}\Rightarrow
\exists x\,\bigl(\Dom(x,X)\land\Cod(x,A)\land\Com(x,f,y)\bigr)
\Bigr)
\Bigr).
\end{aligned}
\]
\end{definicio}

\begin{notacio}[Exhaustivitat generalitzada]\label{not:exhaustiva-generalitzada}
En notació de fletxes, \(f:A\to B\) és exhaustiva sobre elements generalitzats si, per a tot estadi \(X\) i tot element generalitzat \(y:X\to B\), existeix un element generalitzat \(x:X\to A\) tal que
\[
f\circ x=y.
\]
És a dir, tot element generalitzat de \(B\) és una avaluació de \(f\): és de la forma \(f\circ x\) per a un element generalitzat \(x\) de \(A\) en el mateix estadi.
\end{notacio}

\begin{proposicio}[Exhaustivitat generalitzada i epimorfismes escindits]\label{prop:exhaustiva-generalitzada-escindit}
Sigui \(f:A\to B\). Són equivalents:
\begin{enumerate}
\item \(f\) és exhaustiva sobre elements generalitzats;
\item \(f\) té una secció, és a dir, existeix \(s:B\to A\) amb \(f\circ s=\id_B\);
\item \(f\) és un epimorfisme escindit.
\end{enumerate}
\end{proposicio}

\begin{prova}
L'equivalència entre (2) i (3) és la definició d'epimorfisme escindit.

\emph{(1) \(\Rightarrow\) (2).} Apliquem l'exhaustivitat generalitzada a l'estadi \(X=B\) i a l'element generalitzat \(\id_B:B\to B\). Obtenim una aplicació \(s:B\to A\) amb \(f\circ s=\id_B\), que és una secció de \(f\).

\emph{(2) \(\Rightarrow\) (1).} Sigui \(s\) una secció de \(f\), i sigui \(y:X\to B\) un element generalitzat qualsevol. Aleshores \(x:=s\circ y:X\to A\) compleix
\[
f\circ x=f\circ(s\circ y)=(f\circ s)\circ y=\id_B\circ y=y,
\]
de manera que \(y\) és l'avaluació \(f\circ x\). Per tant, \(f\) és exhaustiva sobre elements generalitzats.
\end{prova}

\begin{observacio}[Exhaustivitat generalitzada i exhaustivitat global]\label{obs:exhaustiva-generalitzada-global}
L'exhaustivitat sobre elements generalitzats és una condició més forta que l'exhaustivitat sobre els elements globals: la primera exigeix que tot element generalitzat \(y:X\to B\) sigui una avaluació de \(f\), per a qualsevol estadi \(X\), mentre que la segona només ho demana per als elements globals, és a dir, per a l'estadi terminal \(\mathbf1\). Per la proposició anterior, l'exhaustivitat generalitzada equival a l'existència d'una secció; és, doncs, una propietat almenys tan forta com ser epimorfisme, i possiblement més forta. La qüestió de si tot epimorfisme és exhaustiu sobre elements generalitzats ---és a dir, si tot epimorfisme és escindit--- no es decideix amb els axiomes actuals, i es reprendrà en estudiar l'axioma d'elecció.
\end{observacio}

\section{Parts d'un conjunt}\label{sec:parts}

Els monomorfismes permeten expressar la noció de part exclusivament mitjançant aplicacions.

\begin{definicio}[Part]\index{Part}\label{def:part}
Una \emph{part} d'un conjunt \(A\) és un monomorfisme amb codomini \(A\). Formalment,
\[
\operatorname{Part}(m,A)
\iff
\operatorname{Conj}(A)\land\operatorname{Mono}(m)\land\Cod(m,A).
\]
\end{definicio}

\begin{notacio}[Representació d'una part]\label{not:part}
Si \(\operatorname{Part}(m,A)\) i \(S=\dom(m)\), escriurem
\[
m:S\hookrightarrow A.
\]
Reservarem, per tant, la fletxa amb ganxo per indicar que un monomorfisme és considerat com una part del seu codomini. El conjunt \(S\) és el domini de la part i \(m\) especifica la manera com aquesta part està situada en \(A\).
\end{notacio}

Una part no és, per tant, només el seu domini. Dos monomorfismes amb el mateix domini i el mateix codomini poden situar el domini de maneres diferents dins del codomini.

\subsection{Inclusió de parts}\label{sec:inclusio-parts}

\begin{definicio}[Inclusió]\index{Inclusió}\label{def:inclusio}
Siguin \(m\) i \(n\) dues parts del mateix conjunt \(A\). Direm que \(m\) està \emph{inclosa} en \(n\) si es factoritza a través de \(n\). Formalment,
\[
\operatorname{Incl}(m,n,A)
\iff
\operatorname{Part}(m,A)\land\operatorname{Part}(n,A)
\land\exists k\,\Com(k,n,m).
\]
\end{definicio}

\begin{notacio}[Inclusió de parts]\label{not:inclusio}
Escriurem
\[
m\subseteq_A n
\]
com a abreviació de \(\operatorname{Incl}(m,n,A)\). Si \(m:S\hookrightarrow A\) i \(n:T\hookrightarrow A\), aquesta condició significa que existeix una aplicació \(k:S\to T\) tal que
\[
m=n\circ k.
\]
Geomètricament, la inclusió queda representada pel triangle commutatiu
\begin{center}
	\begin{tikzpicture}[scale=1.1, every node/.style={font=\small}, >=Stealth]
		\node (S) at (0,1.7) {$S$};
		\node (T) at (3.3,1.7) {$T$};
		\node (A) at (3.3,0) {$A$};
		
		\draw[->, thick, dashed]
		(S) -- node[above] {$k$} (T);
		
		\draw[{Hooks[right]}-{Stealth[scale=1.2]}, thick]
		(T) -- node[right] {$n$} (A);
		
		\draw[{Hooks[right]}-{Stealth[scale=1.2]}, thick]
		(S) -- node[below left] {$m$} (A);
	\end{tikzpicture}
\end{center}
La fletxa discontínua posa de relleu que \(k\) és determinada per la factorització.
\end{notacio}

\begin{proposicio}\label{prop:testimoni-inclusio}
Si \(m\subseteq_A n\), l'aplicació \(k\) que satisfà \(m=n\circ k\) és única. A més, \(k\) és un monomorfisme.
\end{proposicio}

\begin{prova}
Si \(m=n\circ k=n\circ k'\), com que \(n\) és un monomorfisme, tenim \(k=k'\).

Per demostrar que \(k\) és un monomorfisme, suposem que \(k\circ u=k\circ v\). Aleshores
\[
m\circ u=n\circ k\circ u=n\circ k\circ v=m\circ v.
\]
Com que \(m\) és un monomorfisme, \(u=v\).
\end{prova}

\section{Pertinença generalitzada}\label{sec:pertinenca-generalitzada}

La inclusió expressa la comparació entre dues parts. La pertinença generalitzada expressarà, també mitjançant una factorització, quan una aplicació amb codomini \(A\) passa a través d'una part de \(A\).

\begin{definicio}[Pertinença]\index{Pertinença}\label{def:pertinenca}
Un element generalitzat \(x\) pertany a una part \(m\) si es factoritza a través d'aquesta part. Formalment,
\[
\begin{aligned}
\operatorname{Pert}(x,X,m)
\iff {}&
\exists A\,\exists S\,\Bigl(
\operatorname{Part}(m,A)
\land\Dom(m,S)
\\
&{}\land\Dom(x,X)\land\Cod(x,A)
\\
&{}\land\exists \bar{x}\,\bigl(
\Dom(\bar{x},X)\land\Cod(\bar{x},S)
\land\Com(\bar{x},m,x)
\bigr)
\Bigr).
\end{aligned}
\]
\end{definicio}

\begin{notacio}[Pertinença a una part]\label{not:pertinenca-part}
Si \(m:S\hookrightarrow A\) i \(x:X\to A\), escriurem
\[
x\in_Xm
\]
com a abreviació de \(\operatorname{Pert}(x,X,m)\). Aquesta condició significa que existeix una aplicació \(\bar{x}:X\to S\) tal que
\[
x=m\circ\bar{x}.
\]
La factorització es representa mitjançant el triangle commutatiu
\begin{center}
\begin{tikzpicture}[scale=1.1, every node/.style={font=\small}, >=Stealth]
  \node (X) at (0,1.7) {$X$};
  \node (S) at (3.3,1.7) {$S$};
  \node (A) at (3.3,0) {$A$};

  \draw[->, thick, dashed]
    (X) -- node[above] {$\bar{x}$} (S);

  \draw[{Hooks[right]}-{Stealth[scale=1.2]}, thick]
    (S) -- node[right] {$m$} (A);

  \draw[->, thick]
    (X) -- node[below left] {$x$} (A);
\end{tikzpicture}
\end{center}
Quan l'estadi quedi determinat pel context, podrem escriure simplement \(x\in m\).
\end{notacio}

La pertinença no és una relació primitiva entre individus, sinó una propietat definida a partir de la composició d'aplicacions.

\begin{proposicio}\label{prop:testimoni-pertinenca}
Si \(x\in_Xm\), l'aplicació \(\bar{x}:X\to S\) que satisfà \(x=m\circ\bar{x}\) és única.
\end{proposicio}

\begin{prova}
Si
\[
x=m\circ\bar{x}=m\circ\bar{x}',
\]
com que \(m\) és un monomorfisme, \(\bar{x}=\bar{x}'\).
\end{prova}

\begin{proposicio}\label{prop:inclusio-pertinenca}
Siguin \(m:S\hookrightarrow A\) i \(n:T\hookrightarrow A\) dues parts de \(A\). Aleshores
\[
m\subseteq_A n
\]
si i només si, per a tot conjunt \(X\) i tota aplicació \(x:X\to A\),
\[
x\in_Xm\quad\Longrightarrow\quad x\in_Xn.
\]
\end{proposicio}

\begin{prova}
Suposem primer que \(m\subseteq_A n\). Existeix una aplicació \(k:S\to T\) tal que \(m=n\circ k\). Si \(x\in_Xm\), existeix una aplicació \(\bar{x}:X\to S\) tal que \(x=m\circ\bar{x}\). Les dues hipòtesis queden expressades pels triangles commutatius següents:

\begin{figure}[H]
	\centering
	\begin{tikzpicture}[scale=1.2, every node/.style={font=\small}]
		\node (U) at (0,1.5) {\(X\)};
		\node (S) at (2,1.5) {\(S\)};
		\node (T) at (4,1.5) {\(T\)};
		\node (A) at (2,0) {\(A\)};
		
		\draw[-{Stealth[scale=1.2]}, thick, dashed] (U) -- node[above] {\(\bar{x}\)} (S);
		\draw[-{Stealth[scale=1.2]}, thick, dashed] (S) -- node[above] {\(k\)} (T);
		\draw[-{Stealth[scale=1.2]}, thick] (U) -- node[below left] {\(x\)} (A);
		\draw[{Hooks[right]}-{Stealth[scale=1.2]}, thick] (S) -- node[left] {\(m\)} (A);
		\draw[{Hooks[right]}-{Stealth[scale=1.2]}, thick] (T) -- node[below right] {\(n\)} (A);
	\end{tikzpicture}
	\caption{Si \(x \in_X m\) i \(m \subseteq_A n\), el testimoni de \(x \in_X n\) és \(k \circ \bar{x}\).}
	\label{fig:fonamental}
\end{figure}

En compondre les dues fletxes discontínues obtenim
\[
x=m\circ\bar{x}
  =(n\circ k)\circ\bar{x}
  =n\circ(k\circ\bar{x}).
\]
Per tant, el triangle
\begin{center}
\begin{tikzpicture}[scale=1.1, every node/.style={font=\small}, >=Stealth]

    \node (A3) at (0,1.6) {$X$};
    \node (B3) at (3.2,1.6) {$T$};
    \node (C3) at (1.6,0) {$A$};

    \draw[->, thick, dashed]
        (A3) -- node[above] {$k\circ\bar{x}$} (B3);

    \draw[->, thick]
        (A3) -- node[left] {$x$} (C3);

    \draw[{Hooks[right]}-{Stealth[scale=1.2]}, thick]
        (B3) -- node[right] {$n$} (C3);

\end{tikzpicture}
\end{center}
commuta, i això significa que \(x\in_Xn\).

Recíprocament, suposem que, per a tot \(X\) i tot \(x:X\to A\), de \(x\in_Xm\) se segueix \(x\in_Xn\). Prenem
\[
X=S,
\qquad
x=m:S\hookrightarrow A.
\]
El triangle de l'esquerra mostra que \(m\in_Sm\), amb testimoni \(\id_S\). La hipòtesi proporciona llavors una aplicació \(k:S\to T\) que fa commutar el triangle de la dreta:

\begin{figure}[H]
	\centering
	\begin{tikzpicture}[scale=1.2, every node/.style={font=\small}]
		\node (U) at (0,1.5) {\(S\)};
		\node (S) at (2,1.5) {\(S\)};
		\node (T) at (4,1.5) {\(T\)};
		\node (A) at (2,0) {\(A\)};
		
		\draw[-{Stealth[scale=1.2]}, thick] (U) -- node[above] {\(\id_S\)} (S);
		\draw[-{Stealth[scale=1.2]}, thick, dashed] (S) -- node[above] {\(k\)} (T);
		\draw[{Hooks[right]}-{Stealth[scale=1.2]}, thick] (U) -- node[below left] {\(m\)} (A);
		\draw[{Hooks[right]}-{Stealth[scale=1.2]}, thick] (S) -- node[left] {\(m\)} (A);
		\draw[{Hooks[right]}-{Stealth[scale=1.2]}, thick] (T) -- node[below right] {\(n\)} (A);
	\end{tikzpicture}
\end{figure}

Aquest triangle expressa
\[
m=n\circ k,
\]
i, per definició, \(m\subseteq_A n\).
\end{prova}

La \hyperref[prop:inclusio-pertinenca]{proposició~\ref*{prop:inclusio-pertinenca}} és la forma estructural del criteri d'inclusió: una part està inclosa en una altra exactament quan tots els seus elements generalitzats pertanyen també a l'altra.

\section{Parts equivalents i subconjunts}\label{sec:parts-equivalents}

Abans de comparar dues parts, observem que el domini d'una part es pot substituir per un conjunt isomorf sense perdre la propietat de ser una part.

\begin{proposicio}[Reparametrització isomorfa d'una part]\label{prop:reparametritzacio-part}
Siguin
\[
m:S\hookrightarrow A
\qquad\text{i}\qquad
t:T\xrightarrow{\sim}S.
\]
Aleshores
\[
m\circ t:T\hookrightarrow A
\]
és una part de \(A\).
\end{proposicio}

\begin{prova}
Tot isomorfisme és un monomorfisme per la \hyperref[prop:isomorfisme-mono-epi]{proposició~\ref*{prop:isomorfisme-mono-epi}}, i la composició de monomorfismes és un monomorfisme per la \hyperref[prop:composicio-mono-epi]{proposició~\ref*{prop:composicio-mono-epi}}. Per tant, \(m\circ t\) és un monomorfisme amb codomini \(A\), és a dir, una part de \(A\).

Geomètricament, la nova part és el camí exterior del triangle commutatiu
\begin{center}
\begin{tikzpicture}[scale=1.1, every node/.style={font=\small}, >=stealth]
  \node (T) at (0,1.7) {\(T\)};
  \node (S) at (3.3,1.7) {\(S\)};
  \node (A) at (3.3,0) {\(A\)};
  \draw[->, thick] (T) -- node[above] {\(\sim\)} node[below] {\(t\)} (S);
  \draw[{Hooks[right]}-{Stealth[scale=1.2]}, thick] (S) -- node[right] {\(m\)} (A);
  \draw[{Hooks[right]}-{Stealth[scale=1.2]}, thick] (T) -- node[below left] {\(m\circ t\)} (A);
\end{tikzpicture}
\end{center}
\end{prova}

\begin{definicio}[Equivalència de parts]\index{Equivalència de parts}\label{def:parts-equivalents}
Dues parts \(m\) i \(n\) de \(A\) són \emph{equivalents} si cadascuna es factoritza a través de l'altra. Formalment,
\[
\operatorname{EqPart}(m,n,A)
\iff
\operatorname{Incl}(m,n,A)\land\operatorname{Incl}(n,m,A).
\]
\end{definicio}

\begin{notacio}[Equivalència de parts]\label{not:equivalencia-parts}
Escriurem
\[
m\equiv_A n
\]
com a abreviació de \(\operatorname{EqPart}(m,n,A)\).
\end{notacio}

Siguin \(m:S\hookrightarrow A\) i \(n:T\hookrightarrow A\). Si \(m\equiv_A n\), existeixen aplicacions \(k:S\to T\) i \(t:T\to S\) tals que
\[
m=n\circ k,
\qquad
n=m\circ t.
\]
Les dues factoritzacions es representen simultàniament mitjançant el diagrama
\begin{center}
\begin{tikzpicture}[scale=1.1, every node/.style={font=\small}, >=Stealth]

  \node (S) at (0,1.8) {$S$};
  \node (T) at (4.0,1.8) {$T$};
  \node (A) at (2.0,0) {$A$};

  \draw[->, thick, dashed]
    ([yshift=2pt]S.east) --
    node[above] {$k$}
    ([yshift=2pt]T.west);

  \draw[->, thick, dashed]
    ([yshift=-2pt]T.west) --
    node[below] {$t$}
    ([yshift=-2pt]S.east);

  \draw[{Hooks[right]}-{Stealth[scale=1.2]}, thick]
    (S) -- node[below left] {$m$} (A);

  \draw[{Hooks[right]}-{Stealth[scale=1.2]}, thick]
    (T) -- node[below right] {$n$} (A);

\end{tikzpicture}
\end{center}

\begin{proposicio}\label{prop:parts-equivalents-isomorfisme}
Si \(m\equiv_A n\), les aplicacions \(k:S\to T\) i \(t:T\to S\) que testimonien les dues factoritzacions són isomorfismes inversos l'una de l'altra. En particular, existeix un únic isomorfisme \(k:S\to T\) tal que
\[
m=n\circ k.
\]
Recíprocament, l'existència d'un isomorfisme \(k:S\to T\) que satisfà aquesta igualtat implica \(m\equiv_A n\).
\end{proposicio}

\begin{prova}
Suposem que
\[
m=n\circ k,
\qquad
n=m\circ t.
\]
Aleshores
\[
n\circ(k\circ t)
=(n\circ k)\circ t
=m\circ t
=n
=n\circ\id_T.
\]
Com que \(n\) és un monomorfisme,
\[
k\circ t=\id_T.
\]
Anàlogament,
\[
m\circ(t\circ k)
=(m\circ t)\circ k
=n\circ k
=m
=m\circ\id_S,
\]
i, com que \(m\) és un monomorfisme,
\[
t\circ k=\id_S.
\]
Per tant, \(k\) i \(t\) són isomorfismes inversos. La unicitat de \(k\) és una conseqüència de la \hyperref[prop:testimoni-inclusio]{proposició~\ref*{prop:testimoni-inclusio}}.

Recíprocament, si \(k:S\to T\) és un isomorfisme i \(m=n\circ k\), tenim \(m\subseteq_A n\). Com que
\[
n=m\circ k^{-1},
\]
també \(n\subseteq_A m\), i per tant \(m\equiv_A n\).
\end{prova}

\begin{proposicio}\label{prop:equivalencia-parts}
Per a parts \(m,n,p\) d'un mateix conjunt \(A\), es compleixen les propietats següents:
\[
m\equiv_A m;
\]
\[
m\equiv_A n\Longrightarrow n\equiv_A m;
\]
\[
m\equiv_A n\land n\equiv_A p
\Longrightarrow
m\equiv_A p.
\]
\end{proposicio}

\begin{prova}
La primera propietat es testimonia amb l'aplicació identitat del domini de \(m\). La segona se segueix directament de la simetria de les dues factoritzacions que formen la definició de \(m\equiv_A n\).

Per a la tercera, la \hyperref[prop:parts-equivalents-isomorfisme]{proposició~\ref*{prop:parts-equivalents-isomorfisme}} proporciona isomorfismes
\[
k:\dom(m)\xrightarrow{\sim}\dom(n),
\qquad
\ell:\dom(n)\xrightarrow{\sim}\dom(p)
\]
compatibles amb les tres parts. La composició \(\ell\circ k\) és un isomorfisme i proporciona la factorització de \(m\) a través de \(p\). La seva inversa proporciona la factorització de \(p\) a través de \(m\). Per tant, \(m\equiv_A p\).
\end{prova}

\begin{observacio}[Recuperació metalingüística dels subconjunts]\label{obs:subconjunt-metallenguatge}
Fixem un conjunt \(A\). Una mateixa part estructural pot ser representada per monomorfismes
\[
m:S\hookrightarrow A
\]
amb dominis diferents però isomorfs. La notació \(m\equiv_A n\) expressa que dues d'aquestes representacions determinen la mateixa manera de situar una part dins de \(A\), llevat d'un canvi isomorf del domini.

En el metallenguatge, podem recuperar la noció clàssica de \emph{subconjunt de \(A\)} considerant la classe de totes les parts de \(A\) equivalents a un representant donat. Aquesta classe no és, de moment, un conjunt intern de la teoria. En els raonaments treballarem amb una part concreta com a representant, i les proposicions anteriors garantiran que les propietats estructurals no depenen del representant escollit.
\end{observacio}

\section{Sumes de conjunts}\label{sec:suma}

\begin{definicio}[Suma]\index{Suma}\label{def:suma-conjunts}
Un conjunt \(S\), juntament amb aplicacions
\[
i_0:A_0\to S,
\qquad
i_1:A_1\to S,
\]
és una \emph{suma} dels conjunts \(A_0\) i \(A_1\) si compleix la propietat universal següent. Formalment,
\[
\begin{aligned}
&\operatorname{Suma}(S,i_0,i_1,A_0,A_1)
\\
\iff {}&
\operatorname{Conj}(S)
\land\Dom(i_0,A_0)\land\Cod(i_0,S)
\land\Dom(i_1,A_1)\land\Cod(i_1,S)
\\
&{}\land
\forall B\,\forall f_0\,\forall f_1\,
\Bigl(
\operatorname{Conj}(B)
\land\Dom(f_0,A_0)\land\Cod(f_0,B)
\\
&\qquad{}\land\Dom(f_1,A_1)\land\Cod(f_1,B)
\\
&\qquad{}\Rightarrow
\exists!h\,
\bigl(
\Dom(h,S)\land\Cod(h,B)
\land\Com(i_0,h,f_0)\land\Com(i_1,h,f_1)
\bigr)
\Bigr).
\end{aligned}
\]
\end{definicio}

\begin{notacio}[Suma i aplicació definida per casos]\label{not:suma}
Quan \(\operatorname{Suma}(S,i_0,i_1,A_0,A_1)\), representarem les dues \emph{injeccions} de la suma per dues fletxes que conflueixen en \(S\):
\[
A_0\xrightarrow{i_0}S\xleftarrow{i_1}A_1.
\]
Les aplicacions \(i_0\) i \(i_1\) s'anomenen \emph{injeccions} de la suma (o \emph{injeccions canòniques}). Reservem aquest terme per a les fletxes estructurals \(i_0,i_1\) d'una suma, i el terme \emph{inclusió} per als subobjectes \(S\hookrightarrow A\) i l'ordre entre parts (definició~\ref{def:inclusio}); són nocions distintes. En aquest punt les injeccions són simplement les fletxes de la propietat universal: encara \emph{no} sabem que siguin monomorfismes, cosa que dependrà de la disjunció de les sumes (definició~\ref{def:suma-disjunta} i axioma~\ref{ax:sumes-disjuntes}, capítol~\ref{cap:valors}). Per a \(f_0:A_0\to B\) i \(f_1:A_1\to B\), denotarem l'única aplicació proporcionada per la propietat universal per
\[
[f_0,f_1]:S\to B.
\]
Així, les fórmules \(\Com(i_0,[f_0,f_1],f_0)\) i \(\Com(i_1,[f_0,f_1],f_1)\) s'escriuen
\[
[f_0,f_1]\circ i_0=f_0,
\qquad
[f_0,f_1]\circ i_1=f_1.
\]
La propietat universal queda representada pel diagrama commutatiu
\begin{center}
\begin{tikzpicture}[scale=1.05, every node/.style={font=\small}, >=Stealth]
  \node (A0) at (0,2.1) {\(A_0\)};
  \node (S)  at (3.4,2.1) {\(S\)};
  \node (A1) at (6.8,2.1) {\(A_1\)};
  \node (B)  at (3.4,0) {\(B\)};

  \draw[->, thick] (A0) -- node[above] {\(i_0\)} (S);
  \draw[->, thick] (A1) -- node[above] {\(i_1\)} (S);
  \draw[->, thick] (A0) -- node[below left] {\(f_0\)} (B);
  \draw[->, thick] (A1) -- node[below right] {\(f_1\)} (B);
  \draw[->, thick, dashed] (S) -- node[right] {\([f_0,f_1]\)} (B);
\end{tikzpicture}
\end{center}
La fletxa discontínua posa de relleu que l'aplicació \([f_0,f_1]\) existeix i és única.
\end{notacio}

\begin{proposicio}[Extensionalitat de la definició per casos]\label{prop:extensionalitat-suma}
Si
\[
A_0\xrightarrow{i_0}S\xleftarrow{i_1}A_1
\]
és una suma i \(h:S\to B\), aleshores
\[
h=[h\circ i_0,h\circ i_1].
\]
En particular, per a dues aplicacions \(h,k:S\to B\),
\[
h=k
\quad\Longleftrightarrow\quad
h\circ i_0=k\circ i_0
\ \land\ 
h\circ i_1=k\circ i_1.
\]
\end{proposicio}

\begin{prova}
Les aplicacions \(h\) i \([h\circ i_0,h\circ i_1]\) tenen les mateixes composicions amb \(i_0\) i \(i_1\). La unicitat de la propietat universal implica
\[
h=[h\circ i_0,h\circ i_1].
\]
La segona afirmació se segueix aplicant aquesta identitat a \(h\) i a \(k\).
\end{prova}



\begin{axioma}[Existència de sumes]\label{ax:sumes}
Tot parell de conjunts té una suma. Formalment,
\[
\begin{aligned}
\forall A_0\,\forall A_1\,\Bigl(&
\operatorname{Conj}(A_0)\land\operatorname{Conj}(A_1)
\\
&{}\Rightarrow
\exists S\,\exists i_0\,\exists i_1\,
\operatorname{Suma}(S,i_0,i_1,A_0,A_1)
\Bigr).
\end{aligned}
\]
\end{axioma}

\begin{notacio}[Suma escollida]\label{not:suma-escollida}
Fixarem una suma de \(A_0\) i \(A_1\), que denotarem per \(A_0+A_1\), i escriurem
\[
i_0:A_0\to A_0+A_1,
\qquad
i_1:A_1\to A_0+A_1
\]
per a les injeccions escollides.
\end{notacio}

\begin{observacio}[Lectura clàssica de la suma i definició per casos]\label{ex:suma-classica}
En la teoria clàssica de conjunts, una suma es pot construir mitjançant dues còpies disjuntes dels conjunts:
\[
A_0+A_1
=
\bigl(\{0\}\times A_0\bigr)
\cup
\bigl(\{1\}\times A_1\bigr).
\]
Aquesta expressió pertany al metallenguatge i no és la definició formal de la suma dins de la teoria. Les etiquetes \(0\) i \(1\) només serveixen per distingir les dues còpies.

Donades aplicacions \(f_0:A_0\to B\) i \(f_1:A_1\to B\), l'aplicació
\[
[f_0,f_1]:A_0+A_1\to B
\]
es defineix per casos:
\[
[f_0,f_1](0,a)=f_0(a),
\qquad
[f_0,f_1](1,b)=f_1(b).
\]
La propietat universal expressa estructuralment que aquesta és l'única aplicació que actua com \(f_0\) sobre la primera còpia i com \(f_1\) sobre la segona.
\end{observacio}

\begin{proposicio}[Unicitat de la suma]\label{prop:suma-unica}
Dues sumes dels mateixos conjunts són isomorfes mitjançant un únic isomorfisme que conserva les injeccions.
\end{proposicio}

\begin{prova}
Siguin
\[
A_0\xrightarrow{i_0}S\xleftarrow{i_1}A_1
\]
i
\[
A_0\xrightarrow{j_0}T\xleftarrow{j_1}A_1
\]
dues sumes. Per la propietat universal de \(S\), existeix una única aplicació \(u:S\to T\) tal que
\[
u\circ i_0=j_0,
\qquad
u\circ i_1=j_1.
\]
Anàlogament, existeix una única aplicació \(v:T\to S\) tal que
\[
v\circ j_0=i_0,
\qquad
v\circ j_1=i_1.
\]
Les aplicacions \(v\circ u\) i \(\id_S\) tenen les mateixes composicions amb \(i_0\) i \(i_1\). Per la unicitat de la propietat universal,
\[
v\circ u=\id_S.
\]
De la mateixa manera,
\[
u\circ v=\id_T.
\]
Per tant, \(u\) és un isomorfisme. La seva unicitat forma part de la mateixa propietat universal.
\end{prova}

\begin{definicio}[Cograf]\index{Cograf}\label{def:cograf}
El \emph{cograf} d'una aplicació \(f:A\to B\) és l'aplicació de la suma \(A+B\) en \(B\) determinada per \(f\) i la identitat de \(B\). Formalment,
\[
\begin{aligned}
\operatorname{Cograf}(k,f)
\iff {}&
\exists A\,\exists B\,\exists S\,\exists i_0\,\exists i_1\,\Bigl(
\\
&\quad\Dom(f,A)\land\Cod(f,B)
\\
&\quad{}\land\operatorname{Suma}(S,i_0,i_1,A,B)
\\
&\quad{}\land\Dom(k,S)\land\Cod(k,B)
\\
&\quad{}\land\Com(i_0,k,f)\land\Com(i_1,k,B)
\Bigr).
\end{aligned}
\]
\end{definicio}

\begin{notacio}[Cograf]\label{not:cograf}
Denotarem el cograf de \(f\) per
\[
\kappa_f=[f,\id_B]:A+B\to B.
\]
La construcció queda representada pel diagrama commutatiu
\begin{center}
\begin{tikzpicture}[scale=1.05, every node/.style={font=\small}, >=Stealth]
  \node (A)  at (0,2.1) {\(A\)};
  \node (AB) at (3.4,2.1) {\(A+B\)};
  \node (Bt) at (6.8,2.1) {\(B\)};
  \node (Bb) at (3.4,0) {\(B\)};

  \draw[->, thick] (A) -- node[above] {\(i_0\)} (AB);
  \draw[->, thick] (Bt) -- node[above] {\(i_1\)} (AB);
  \draw[->, thick] (A) -- node[below left] {\(f\)} (Bb);
  \draw[->, thick] (Bt) -- node[below right] {\(\id_B\)} (Bb);
  \draw[->, thick, dashed] (AB) -- node[right] {\(\kappa_f\)} (Bb);
\end{tikzpicture}
\end{center}
Així,
\[
\kappa_f\circ i_0=f,
\qquad
\kappa_f\circ i_1=\id_B.
\]
\end{notacio}

\begin{observacio}[Lectura clàssica del cograf]\label{obs:cograf-classic}
En la representació clàssica de \(A+B\) com dues còpies disjuntes, el cograf envia la còpia d'un element \(a\in A\) al seu valor \(f(a)\in B\) i deixa fixa la còpia de cada element \(b\in B\). D'aquesta manera, els conjunts \(A\) i \(B\) apareixen situats un al costat de l'altre, i la branca corresponent a \(A\) recupera el diagrama intern de l'aplicació \(f\): de cada element del domini surt una fletxa cap al seu valor.
\end{observacio}

\begin{proposicio}\label{prop:cograf-retraccio}
El cograf \(\kappa_f\) és una retracció de la segona injecció \(i_1:B\to A+B\). En particular, \(\kappa_f\) és un epimorfisme i \(i_1\) és un monomorfisme. Simètricament, si existeix alguna aplicació \(g\colon B\to A\), el cograf \([\id_A,g]\colon A+B\to A\) és una retracció de \(i_0\), que és aleshores un monomorfisme.
\end{proposicio}

\begin{prova}
Per la definició de \([f,\id_B]\),
\[
\kappa_f\circ i_1=\id_B.
\]
La conclusió se segueix de la \hyperref[prop:seccio-mono-retraccio-epi]{proposició~\ref*{prop:seccio-mono-retraccio-epi}}; el cas de \(i_0\) és idèntic amb \([\id_A,g]\circ i_0=\id_A\). Observeu que l'argument necessita una aplicació entre els dos sumands: que les injeccions siguin monomorfismes \emph{en general} dependrà de la disjunció de les sumes (capítol~\ref{cap:valors}).
\end{prova}

\subsection{El conjunt \texorpdfstring{\(\mathbf2\)}{2}}\label{dos-oficial}

\begin{notacio}[El conjunt \(\mathbf2\) i els valors distingits]\label{not:dos}
Denotarem per \(\mathbf2\) la suma escollida \(\mathbf1+\mathbf1\):
\[
\mathbf2=\mathbf1+\mathbf1.
\]
Denotarem les dues injeccions per
\[
0:\mathbf1\to\mathbf2,
\qquad
1:\mathbf1\to\mathbf2.
\]
L'aplicació \(1\) serà la part distingida que intervindrà en l'axioma de representació.
\end{notacio}

\begin{corollari}\label{prop:zero-u-mono}
Les aplicacions \(0\) i \(1\) són monomorfismes.
\end{corollari}

\begin{prova}
La propietat universal de la suma determina una aplicació
\[
[\id_{\mathbf1},\id_{\mathbf1}]:\mathbf2\to\mathbf1.
\]
Aquesta aplicació és una retracció tant de \(0\) com de \(1\), perquè
\[
[\id_{\mathbf1},\id_{\mathbf1}]\circ0=\id_{\mathbf1}
\]
i
\[
[\id_{\mathbf1},\id_{\mathbf1}]\circ1=\id_{\mathbf1}.
\]
Per la \hyperref[prop:seccio-mono-retraccio-epi]{proposició~\ref*{prop:seccio-mono-retraccio-epi}}, \(0\) i \(1\) són monomorfismes.
\end{prova}

\begin{observacio}[Una aplicació de la suma cap a \(\mathbf2\)]\label{ex:suma-cap-dos}
Per a qualssevol conjunts \(A\) i \(B\), la propietat universal de la suma determina l'aplicació
\[
[0\circ {!}_A,1\circ {!}_B]:A+B\to\mathbf2.
\]
Aquesta aplicació és constant amb valor \(0\) sobre el sumand \(A\) i constant amb valor \(1\) sobre el sumand \(B\):
\[
[0\circ {!}_A,1\circ {!}_B]\circ i_0=0\circ {!}_A,
\]
\[
[0\circ {!}_A,1\circ {!}_B]\circ i_1=1\circ {!}_B.
\]
En la lectura clàssica, indica de quin dels dos sumands prové cada element de \(A+B\).
\end{observacio}

\begin{observacio}[Aplicacions de \(\mathbf2\) en un conjunt]\label{obs:dos-cap-a}
Donar una aplicació \(h:\mathbf2\to A\) equival a donar dues aplicacions
\[
a_0:\mathbf1\to A,
\qquad
a_1:\mathbf1\to A.
\]
En efecte, la propietat universal de \(\mathbf2=\mathbf1+\mathbf1\) determina l'aplicació
\[
[a_0,a_1]:\mathbf2\to A
\]
caracteritzada per
\[
[a_0,a_1]\circ0=a_0,
\qquad
[a_0,a_1]\circ1=a_1.
\]
Així, en la lectura clàssica, una aplicació de \(\mathbf2\) en \(A\) selecciona un parell ordenat d'elements de \(A\), que no han de ser necessàriament diferents.
\end{observacio}


\section{Representació de les parts}\label{representacio}

La pertinença a una part es pot convertir en una aplicació cap a \(\mathbf2\). Aquesta possibilitat no es dedueix dels axiomes anteriors i s'introdueix com un axioma nou.

\begin{definicio}[Aplicació característica]\index{Aplicació característica}\index{Classificador de parts}\label{def:caracteristica}
Una aplicació \(\chi\) és \emph{característica} d'una part \(m\) de \(A\) si classifica exactament els elements generalitzats que es factoritzen a través de \(m\). Formalment,
\[
\begin{aligned}
\operatorname{Car}(\chi,m,A)
\iff {}&
\operatorname{Part}(m,A)\land\Dom(\chi,A)\land\Cod(\chi,\mathbf2)
\\
&{}\land
\forall X\,\forall x\,
\Bigl(
\operatorname{Conj}(X)\land\Dom(x,X)\land\Cod(x,A)
\\
&\qquad{}\Rightarrow
\bigl(
\operatorname{Pert}(x,X,m)
\Longleftrightarrow
\chi\circ x=1\circ {!}_X
\bigr)
\Bigr).
\end{aligned}
\]
\end{definicio}

\begin{axioma}[Representació de les parts]\label{ax:representacio}
Tota part té una única aplicació característica, i tota aplicació cap a \(\mathbf2\) és la característica d'alguna part. Formalment,
\[
\begin{aligned}
\forall A\,\forall m\,\Bigl(&
\operatorname{Part}(m,A)
\Rightarrow
\exists!\chi\,\operatorname{Car}(\chi,m,A)
\Bigr),
\end{aligned}
\]
i
\[
\begin{aligned}
\forall A\,\forall\varphi\,\Bigl(&
\operatorname{Conj}(A)\land\Dom(\varphi,A)\land\Cod(\varphi,\mathbf2)
\\
&{}\Rightarrow
\exists m\,\operatorname{Car}(\varphi,m,A)
\Bigr).
\end{aligned}
\]
\end{axioma}

\begin{observacio}[Tres unicitats que no s'han de confondre]\label{obs:tres-unicitats}
L'axioma parla de dues unicitats diferents, i convé separar-les d'una tercera que no afirma. (i)~\emph{Unicitat de la característica:} per a cada monomorfisme \(m\), l'aplicació \(\chi\) amb \(\operatorname{Car}(\chi,m,A)\) és única com a fletxa, perquè la condició \(x\in_Xm\iff\chi\circ x=1\circ{!}_X\) la determina sobre tots els elements generalitzats d'\(A\). (ii)~\emph{Existència, però no unicitat, de la part:} per a cada \(\varphi\colon A\to\mathbf2\) existeix algun \(m\) amb \(\operatorname{Car}(\varphi,m,A)\), però no és únic com a fletxa: si \(m\colon S\hookrightarrow A\) ho compleix i \(u\colon S'\to S\) és un isomorfisme, \(m\circ u\) també ho compleix, perquè té els mateixos elements generalitzats. El que és únic és la \emph{classe d'equivalència} de la part: dues parts amb la mateixa característica tenen els mateixos elements generalitzats i són equivalents (proposició~\ref{prop:inclusio-pertinenca} en els dos sentits), i aleshores els seus dominis són isomorfs per un únic isomorfisme compatible amb les dues fletxes (proposició~\ref{prop:parts-equivalents-isomorfisme}). (iii)~\emph{Igualtat de fletxes:} que \(m\equiv_An\) no vol dir \(m=n\). En resum: les aplicacions \(A\to\mathbf2\) estan en correspondència bijectiva amb les \emph{classes} de parts equivalents d'\(A\), no amb les parts com a fletxes; és la mateixa situació que amb els subconjunts clàssics, que corresponen a les classes i no a les inclusions concretes (observació~\ref{obs:subconjunt-metallenguatge}).
\end{observacio}

\begin{notacio}[Aplicació característica]\label{not:caracteristica}
Si \(\operatorname{Car}(\chi,m,A)\), denotarem l'única aplicació característica per
\[
\chi_m:A\to\mathbf2.
\]
La seva propietat característica s'escriu
\[
x\in_Xm
\quad\Longleftrightarrow\quad
\chi_m\circ x=1\circ {!}_X.
\]
\end{notacio}

La factorització que expressa la pertinença \(x\in_Xm\) es representa mitjançant el triangle commutatiu
\begin{center}
\begin{tikzpicture}[scale=1.1, every node/.style={font=\small}, >=Stealth]
  \node (X) at (0,1.8) {\(X\)};
  \node (S) at (3.5,1.8) {\(S\)};
  \node (A) at (3.5,0) {\(A\)};

  \draw[->, thick, dashed] (X) -- node[above] {\(\bar{x}\)} (S);
  \draw[->, thick] (X) -- node[below left] {\(x\)} (A);
  \draw[{Hooks[right]}-{Stealth[scale=1.2]}, thick] (S) -- node[right] {\(m\)} (A);
\end{tikzpicture}
\end{center}
on \(x=m\circ\bar{x}\). Com que \(m\) és un monomorfisme, el testimoni \(\bar{x}\), quan existeix, és únic.

Prenent \(x=m:S\to A\) i \(\bar{x}=\id_S\), obtenim
\[
\chi_m\circ m=1\circ {!}_S.
\]
Aquesta igualtat queda representada pel quadrat commutatiu
\begin{center}
\begin{tikzpicture}[scale=1.1, every node/.style={font=\small}, >=Stealth]
  \node (S) at (0,1.9) {\(S\)};
  \node (One) at (3.8,1.9) {\(\mathbf1\)};
  \node (A) at (0,0) {\(A\)};
  \node (Two) at (3.8,0) {\(\mathbf2\)};

  \draw[->, thick] (S) -- node[above] {\({!}_S\)} (One);
  \draw[{Hooks[right]}-{Stealth[scale=1.2]}, thick] (S) -- node[left] {\(m\)} (A);
  \draw[{Hooks[right]}-{Stealth[scale=1.2]}, thick] (One) -- node[right] {\(1\)} (Two);
  \draw[->, thick] (A) -- node[below] {\(\chi_m\)} (Two);
\end{tikzpicture}
\end{center}
La propietat característica afirma que una aplicació \(x:X\to A\) es factoritza a través de \(m\) si i només si la composició \(\chi_m\circ x\) coincideix amb \(1\circ {!}_X\). Al \hyperref[teo:quadrat-caracteristic]{teorema~\ref*{teo:quadrat-caracteristic}} veurem que aquesta propietat diu exactament que el quadrat anterior és un pullback.

\begin{observacio}[Distingir elements mitjançant una aplicació característica]\label{obs:caracteristica-distingeix}
Siguin \(x,y:X\to A\) dos elements generalitzats. Si
\[
x\in_Xm
\qquad\text{i}\qquad
y\notin_Xm,
\]
aleshores
\[
\chi_m\circ x\neq\chi_m\circ y.
\]
En efecte, la primera condició implica \(\chi_m\circ x=1\circ {!}_X\), mentre que la segona implica \(\chi_m\circ y\neq1\circ {!}_X\).

En particular, quan \(X=\mathbf1\), l'aplicació característica recupera la noció clàssica de distingir dos elements d'un conjunt segons si pertanyen o no a una part. Els axiomes establerts fins aquí només afirmen que un element que no pertany a la part no és enviat al valor distingit \(1\); encara no afirmen que hagi de ser enviat necessàriament a \(0\).
\end{observacio}

L'axioma es formula per a tots els elements generalitzats. No afirma només quins punts globals pertanyen a una part, sinó quines aplicacions \(X\to A\) es factoritzen a través d'ella, per a qualsevol estadi \(X\).

\begin{proposicio}\label{prop:caracteristica-determina-part}
Dues parts \(m\) i \(n\) de \(A\) són equivalents si i només si tenen la mateixa aplicació característica:
\[
m\equiv_A n
\quad\Longleftrightarrow\quad
\chi_m=\chi_n.
\]
\end{proposicio}

\begin{prova}
Suposem primer que \(\chi_m=\chi_n\). Per a qualsevol \(x:X\to A\),
\[
\begin{aligned}
x\in_X m
&\Longleftrightarrow \chi_m\circ x=1\circ {!}_X\\
&\Longleftrightarrow \chi_n\circ x=1\circ {!}_X\\
&\Longleftrightarrow x\in_X n.
\end{aligned}
\]
Per la \hyperref[prop:inclusio-pertinenca]{proposició~\ref*{prop:inclusio-pertinenca}}, tenim \(m\subseteq_A n\) i \(n\subseteq_A m\). Per tant, \(m\equiv_A n\).

Recíprocament, suposem que \(m\equiv_A n\). La mateixa proposició implica que, per a qualsevol \(x:X\to A\),
\[
x\in_Xm\quad\Longleftrightarrow\quad x\in_Xn.
\]
Per tant, \(\chi_n\) satisfà la condició característica de \(m\). La unicitat establerta per l'axioma de representació implica
\[
\chi_m=\chi_n.
\]
\end{prova}

\begin{proposicio}\label{prop:caracteristica-total}
La característica de la part total \(\id_A:A\hookrightarrow A\) és l'aplicació constant
\[
1\circ {!}_A:A\to\mathbf2.
\]
\end{proposicio}

\begin{prova}
Per a tota aplicació \(x:X\to A\), tenim \(x\in_X\id_A\), amb testimoni \(x\). A més,
\[
(1\circ {!}_A)\circ x
=1\circ({!}_A\circ x)
=1\circ {!}_X.
\]
La unicitat de l'aplicació característica dona el resultat.
\end{prova}

\begin{proposicio}\label{prop:bimorfisme-isomorfisme}
Tot bimorfisme és un isomorfisme.
\end{proposicio}

\begin{prova}
Sigui \(f:A\to B\) un monomorfisme i un epimorfisme. Com que és un monomorfisme, podem considerar-lo una part de \(B\). Sigui
\[
\chi_f:B\to\mathbf2
\]
la seva aplicació característica.

La mateixa aplicació \(f:A\to B\) pertany a la part \(f\), amb testimoni \(\id_A\). Per tant,
\[
\chi_f\circ f=1\circ {!}_A.
\]
D'altra banda,
\[
(1\circ {!}_B)\circ f=1\circ {!}_A.
\]
Com que \(f\) és un epimorfisme,
\[
\chi_f=1\circ {!}_B.
\]
Així,
\[
\chi_f\circ\id_B=1\circ {!}_B,
\]
i l'axioma de representació implica
\[
\id_B\in_B f.
\]
Per tant, existeix una aplicació \(g:B\to A\) tal que
\[
f\circ g=\id_B.
\]
Finalment,
\[
f\circ(g\circ f)=(f\circ g)\circ f=f=f\circ\id_A.
\]
Com que \(f\) és un monomorfisme,
\[
g\circ f=\id_A.
\]
Per tant, \(g\) és la inversa de \(f\), i \(f\) és un isomorfisme.
\end{prova}

\begin{corollari}\label{cor:isomorfisme-bimorfisme}
Una aplicació és un isomorfisme si i només si és un monomorfisme i un epimorfisme.
\end{corollari}

\begin{prova}
Una implicació és la \hyperref[prop:isomorfisme-mono-epi]{proposició~\ref*{prop:isomorfisme-mono-epi}} i l'altra és la \hyperref[prop:bimorfisme-isomorfisme]{proposició~\ref*{prop:bimorfisme-isomorfisme}}.
\end{prova}

\begin{observacio}[Injectivitat clàssica i monomorfismes]\label{obs:injectiva-mono}
En la teoria clàssica de conjunts, una aplicació \(f:A\to B\) és \emph{injectiva} si, per a tots els elements globals \(x,y:\mathbf1\to A\),
\[
f\circ x=f\circ y
\quad\Longrightarrow\quad
x=y.
\]
En aquesta observació, i en les següents, ometrem l'expressió \enquote{respecte dels elements globals} perquè queda incorporada a la definició clàssica.

Tot monomorfisme és injectiu. En efecte, només cal aplicar la definició de monomorfisme a les aplicacions \(x,y:\mathbf1\to A\).

Recíprocament, suposem que \(f\) és injectiva i que
\[
f\circ u=f\circ v
\]
per a dues aplicacions \(u,v:X\to A\). Per a tot element global \(x:\mathbf1\to X\), l'associativitat dona
\[
f\circ(u\circ x)=f\circ(v\circ x).
\]
La injectivitat de \(f\) implica
\[
u\circ x=v\circ x.
\]
Per l'\hyperref[ax:separacio-globals]{axioma de separació global}, \(u=v\). Per tant,
\[
\boxed{
 f\text{ és injectiva}
 \quad\Longleftrightarrow\quad
 f\text{ és un monomorfisme}.
}
\]
Aquesta equivalència explica per què els monomorfismes són la formulació estructural de la injectivitat.
\end{observacio}

\begin{observacio}[Exhaustivitat clàssica i epimorfismes]\label{obs:exhaustiva-epi}
En la teoria clàssica de conjunts, una aplicació \(f:A\to B\) és \emph{exhaustiva} si, per a tot element global \(y:\mathbf1\to B\), existeix un element global \(x:\mathbf1\to A\) tal que
\[
f\circ x=y.
\]

Sota l'\hyperref[ax:separacio-globals]{axioma de separació global}, tota aplicació exhaustiva és un epimorfisme. En efecte, suposem que \(u,v:B\to C\) compleixen
\[
u\circ f=v\circ f.
\]
Per a qualsevol element global \(y:\mathbf1\to B\), l'exhaustivitat proporciona \(x:\mathbf1\to A\) amb \(f\circ x=y\). Llavors
\[
\begin{aligned}
u\circ y
&=u\circ(f\circ x)\\
&=(u\circ f)\circ x\\
&=(v\circ f)\circ x\\
&=v\circ(f\circ x)\\
&=v\circ y.
\end{aligned}
\]
La separació global implica \(u=v\), i per tant \(f\) és un epimorfisme.

El recíproc ---que tot epimorfisme sigui exhaustiu sobre els elements globals--- no es pot demostrar amb els axiomes d'aquest capítol, perquè demana saber que un conjunt sense punts globals és inicial i que els dos valors \(0,1\colon\mathbf1\to\mathbf2\) són diferents, fets que no establirem fins al capítol~\ref{cap:valors}. Allà el demostrarem, amb la separació global i sense cap axioma d'elecció. Fins aleshores, i en tot el desenvolupament estructural, mantindrem separades les nocions d'epimorfisme i d'exhaustivitat sobre punts; la noció més forta, l'exhaustivitat sobre elements \emph{generalitzats}, sí que és estrictament més forta que ser epimorfisme sense l'elecció, com s'ha dit a l'observació~\ref{obs:exhaustiva-generalitzada-global}.
\end{observacio}

\begin{observacio}[Bijectivitat clàssica i isomorfismes]\label{obs:bijectiva-iso}
Una aplicació és \emph{bijectiva} en el sentit clàssic si és injectiva i exhaustiva. Sota l'\hyperref[ax:separacio-globals]{axioma de separació global}, una aplicació bijectiva és un monomorfisme per l'\hyperref[obs:injectiva-mono]{observació~\ref*{obs:injectiva-mono}} i és un epimorfisme per l'\hyperref[obs:exhaustiva-epi]{observació~\ref*{obs:exhaustiva-epi}}. Per la \hyperref[prop:bimorfisme-isomorfisme]{proposició~\ref*{prop:bimorfisme-isomorfisme}}, és un isomorfisme.

Recíprocament, si \(f:A\to B\) és un isomorfisme amb inversa \(f^{-1}:B\to A\), aleshores és injectiva perquè
\[
f\circ x=f\circ y
\quad\Longrightarrow\quad
x=f^{-1}\circ f\circ x=f^{-1}\circ f\circ y=y,
\]
i és exhaustiva perquè, per a qualsevol \(y:\mathbf1\to B\), l'element global
\[
x=f^{-1}\circ y:\mathbf1\to A
\]
satisfà \(f\circ x=y\). Per tant,
\[
\boxed{
 f\text{ és bijectiva}
 \quad\Longleftrightarrow\quad
 f\text{ és un isomorfisme}.
}
\]
Aquesta equivalència recupera la caracterització clàssica de les bijeccions, però no intervé en les definicions ni en les demostracions estructurals del capítol.
\end{observacio}

\section{Productes de conjunts}\label{sec:producte}

\begin{definicio}[Producte]\index{Producte}\label{def:producte}
Un conjunt \(P\), juntament amb aplicacions
\[
p_A:P\to A,
\qquad
p_B:P\to B,
\]
és un \emph{producte} dels conjunts \(A\) i \(B\) si compleix la propietat universal següent. Formalment,
\[
\begin{aligned}
&\operatorname{Producte}(P,p_A,p_B,A,B)
\\
\iff {}&
\operatorname{Conj}(P)
\land\Dom(p_A,P)\land\Cod(p_A,A)
\land\Dom(p_B,P)\land\Cod(p_B,B)
\\
&{}\land
\forall X\,\forall f\,\forall g\,
\Bigl(
\operatorname{Conj}(X)
\land\Dom(f,X)\land\Cod(f,A)
\\
&\qquad{}\land\Dom(g,X)\land\Cod(g,B)
\\
&\qquad{}\Rightarrow
\exists!h\,
\bigl(
\Dom(h,X)\land\Cod(h,P)
\land\Com(h,p_A,f)\land\Com(h,p_B,g)
\bigr)
\Bigr).
\end{aligned}
\]
\end{definicio}

\begin{notacio}[Producte i aparellament]\label{not:producte}
Quan \(\operatorname{Producte}(P,p_A,p_B,A,B)\), representarem les dues projeccions canòniques per dues fletxes que divergeixen de \(P\):
\[
A\xleftarrow{p_A}P\xrightarrow{p_B}B.
\]
Per a \(f:X\to A\) i \(g:X\to B\), denotarem l'única aplicació proporcionada per la propietat universal per
\[
\langle f,g\rangle:X\to P.
\]
Així,
\[
p_A\circ\langle f,g\rangle=f,
\qquad
p_B\circ\langle f,g\rangle=g.
\]
La propietat universal queda representada pel diagrama commutatiu
\begin{center}
\begin{tikzpicture}[scale=1.05, every node/.style={font=\small}, >=Stealth]
  \node (A) at (0,2.1) {\(A\)};
  \node (P) at (3.4,2.1) {\(P\)};
  \node (B) at (6.8,2.1) {\(B\)};
  \node (X) at (3.4,0) {\(X\)};

  \draw[->, thick] (P) -- node[above] {\(p_A\)} (A);
  \draw[->, thick] (P) -- node[above] {\(p_B\)} (B);
  \draw[->, thick] (X) -- node[below left] {\(f\)} (A);
  \draw[->, thick] (X) -- node[below right] {\(g\)} (B);
  \draw[->, thick, dashed] (X) -- node[right] {\(\langle f,g\rangle\)} (P);
\end{tikzpicture}
\end{center}
La fletxa discontínua posa de relleu que l'aplicació \(\langle f,g\rangle\) existeix i és única.
\end{notacio}

\begin{axioma}[Existència de productes]\label{ax:productes}
Tot parell de conjunts té un producte. Formalment,
\[
\begin{aligned}
\forall A\,\forall B\,\Bigl(&
\operatorname{Conj}(A)\land\operatorname{Conj}(B)
\\
&{}\Rightarrow
\exists P\,\exists p_A\,\exists p_B\,
\operatorname{Producte}(P,p_A,p_B,A,B)
\Bigr).
\end{aligned}
\]
\end{axioma}

\begin{notacio}[Producte escollit]\label{not:producte-escollit}
Fixarem un producte de \(A\) i \(B\), que denotarem per \(A\times B\), i escriurem
\[
p_A:A\times B\to A,
\qquad
p_B:A\times B\to B
\]
per a les projeccions canòniques escollides.
\end{notacio}

\begin{observacio}[Recuperació de les parelles ordenades]\label{obs:parelles-ordenades}
En la teoria clàssica de conjunts, el producte cartesià es presenta com el conjunt
\[
A\times B=\{(a,b)\mid a\in A\text{ i }b\in B\}.
\]
Aquesta descripció pertany al metallenguatge. La propietat universal permet recuperar-la a partir de les aplicacions.

En efecte, un element global
\[
x:\mathbf1\to A\times B
\]
determina els elements globals
\[
a=p_A\circ x:\mathbf1\to A,
\qquad
b=p_B\circ x:\mathbf1\to B.
\]
Recíprocament, per a cada parell d'elements globals \(a:\mathbf1\to A\) i \(b:\mathbf1\to B\), existeix una única aplicació
\[
\langle a,b\rangle:\mathbf1\to A\times B
\]
amb aquestes dues components. En la interpretació clàssica, aquesta aplicació correspon a la parella ordenada \((a,b)\).

Més generalment, un element generalitzat \(x:X\to A\times B\) queda determinat pel parell d'aplicacions
\[
p_A\circ x:X\to A,
\qquad
p_B\circ x:X\to B,
\]
i qualsevol parell d'aplicacions amb domini comú \(X\) determina un únic element generalitzat del producte.
\end{observacio}

\begin{proposicio}[Unicitat del producte]\label{prop:producte-unic}
Dos productes dels mateixos conjunts són isomorfs mitjançant un únic isomorfisme que conserva les projeccions canòniques.
\end{proposicio}

\begin{prova}
La demostració és dual de la de la \hyperref[prop:suma-unica]{proposició~\ref*{prop:suma-unica}}. Si \(P\) i \(Q\) són dos productes, les seves propietats universals determinen aplicacions úniques \(u:P\to Q\) i \(v:Q\to P\) que conserven les projeccions. Les composicions \(v\circ u\) i \(\id_P\) tenen les mateixes components, i per la unicitat de la propietat universal són iguals. Anàlogament, \(u\circ v=\id_Q\).
\end{prova}

\begin{proposicio}\label{prop:projeccions-separen}
Siguin \(u,v:X\to A\times B\). Si
\[
p_A\circ u=p_A\circ v
\]
i
\[
p_B\circ u=p_B\circ v,
\]
aleshores \(u=v\).
\end{proposicio}

\begin{prova}
Les aplicacions \(u\) i \(v\) satisfan les mateixes equacions amb les projeccions. La unicitat de la propietat universal del producte implica \(u=v\).
\end{prova}

\begin{definicio}[Graf]\index{Graf}\label{def:graf}
El \emph{graf} d'una aplicació \(f:A\to B\) és l'única aplicació \(\gamma:A\to A\times B\) les components de la qual són \(\id_A\) i \(f\). Formalment,
\[
\begin{aligned}
\operatorname{Graf}(\gamma,f)
\iff {}&
\exists A\,\exists B\,\exists P\,\exists p_A\,\exists p_B\,\Bigl(
\\
&\quad\Dom(f,A)\land\Cod(f,B)
\\
&\quad{}\land\operatorname{Producte}(P,p_A,p_B,A,B)
\\
&\quad{}\land\Dom(\gamma,A)\land\Cod(\gamma,P)
\\
&\quad{}\land\Com(\gamma,p_A,A)\land\Com(\gamma,p_B,f)
\Bigr).
\end{aligned}
\]
\end{definicio}

\begin{notacio}[Graf]\label{not:graf}
Denotarem el graf de \(f\) per
\[
\gamma_f=\langle\id_A,f\rangle:A\to A\times B.
\]
\end{notacio}

\begin{observacio}[La definició clàssica d'una aplicació mitjançant el graf]\label{obs:graf-classic}
En la teoria clàssica de conjunts, el graf d'una aplicació \(f:A\to B\) és el subconjunt
\[
\Gamma_f=\{(a,b)\in A\times B\mid b=f(a)\}.
\]
Equivalentment,
\[
\Gamma_f=\{(a,f(a))\mid a\in A\}.
\]
La nostra aplicació
\[
\gamma_f:A\to A\times B
\]
correspon a la funció \(a\mapsto(a,f(a))\), que identifica \(A\) amb aquest subconjunt.

Recíprocament, en la teoria clàssica, un subconjunt \(G\subseteq A\times B\) és el graf d'una aplicació \(A\to B\) si, per a cada \(a\in A\), existeix un únic \(b\in B\) tal que \((a,b)\in G\). Estructuralment, aquesta condició expressa que la primera projecció restringida al graf és un isomorfisme. Aquesta observació recupera la presentació clàssica d'una aplicació com un subconjunt especial d'un producte cartesià, però dins la nostra teoria l'aplicació i el seu graf continuen essent aplicacions diferents.
\end{observacio}

\begin{proposicio}\label{prop:graf-part}
El graf \(\gamma_f\) és una part de \(A\times B\).
\end{proposicio}

\begin{prova}
Per la propietat universal del producte,
\[
p_A\circ\gamma_f=\id_A.
\]
Per tant, \(\gamma_f\) és una secció de \(p_A\). La \hyperref[prop:seccio-mono-retraccio-epi]{proposició~\ref*{prop:seccio-mono-retraccio-epi}} implica que \(\gamma_f\) és un monomorfisme.
\end{prova}

\section{Equalitzadors}\label{equalitzadors}

Siguin \(f,g:A\to B\) dues aplicacions paral·leles, és a dir, dues aplicacions amb el mateix domini \(A\) i el mateix codomini \(B\). Les representarem per
\begin{center}
\begin{tikzpicture}[scale=1.05, every node/.style={font=\small}, >=Stealth]
  \node (A) at (0,0) {\(A\)};
  \node (B) at (4.2,0) {\(B\)};

  \draw[->, thick]
    ([xshift=4pt,yshift=2pt]A.east) --
    node[above] {\(f\)}
    ([xshift=-4pt,yshift=2pt]B.west);
  \draw[->, thick]
    ([xshift=4pt,yshift=-2pt]A.east) --
    node[below] {\(g\)}
    ([xshift=-4pt,yshift=-2pt]B.west);
\end{tikzpicture}
\end{center}
Construirem una part de \(A\) que representi els elements generalitzats sobre els quals les dues aplicacions coincideixen.

Considerem la diagonal
\[
\Delta_B=\langle\id_B,\id_B\rangle:B\to B\times B.
\]

\begin{proposicio}\label{prop:diagonal-part}
La diagonal \(\Delta_B\) és una part de \(B\times B\).
\end{proposicio}

\begin{prova}
Tenim
\[
p_B^{(1)}\circ\Delta_B=\id_B,
\]
on \(p_B^{(1)}:B\times B\to B\) és la primera projecció. Per tant, \(\Delta_B\) és una secció i, per la \hyperref[prop:seccio-mono-retraccio-epi]{proposició~\ref*{prop:seccio-mono-retraccio-epi}}, és un monomorfisme.
\end{prova}

Sigui
\[
\operatorname{eq}_B=\chi_{\Delta_B}:B\times B\to\mathbf2
\]
l'aplicació característica de la diagonal.

\begin{proposicio}\label{prop:diagonal-igualtat}
Per a tot conjunt \(X\) i totes les aplicacions \(u,v:X\to B\),
\[
\langle u,v\rangle\in_X\Delta_B
\quad\Longleftrightarrow\quad
u=v.
\]
\end{proposicio}

\begin{prova}
Si \(\langle u,v\rangle\in_X\Delta_B\), existeix una aplicació \(w:X\to B\) tal que
\[
\langle u,v\rangle=\Delta_B\circ w.
\]
Component amb les dues projeccions obtenim
\[
u=w
\qquad\text{i}\qquad
v=w,
\]
i, per tant, \(u=v\).

Recíprocament, si \(u=v\), aleshores
\[
\langle u,v\rangle
=\langle u,u\rangle
=\Delta_B\circ u,
\]
i \(\langle u,v\rangle\in_X\Delta_B\).
\end{prova}

Per l'\hyperref[ax:representacio]{axioma~\ref*{ax:representacio}}, l'aplicació
\[
\operatorname{eq}_B\circ\langle f,g\rangle:A\to\mathbf2
\]
és la característica d'una part
\[
e:E\hookrightarrow A.
\]

\begin{definicio}[Equalitzador]\index{Equalitzador}\label{def:equalitzador}
Una aplicació \(e:E\to A\) és un \emph{equalitzador} de dues aplicacions paral·leles \(f,g:A\to B\) si iguala \(f\) i \(g\), i si tota aplicació que les iguala es factoritza d'una manera única a través de \(e\). Formalment,
\[
\begin{aligned}
\operatorname{Equalitzador}(e,f,g)
\iff
\exists A\,\exists B\,\exists E\,\Bigl(&
f:A\to B\land g:A\to B\land e:E\to A
\\
&{}\land f\circ e=g\circ e
\\
&{}\land\forall X\,\forall h\,
\Bigl(
 h:X\to A\land f\circ h=g\circ h
\\
&\qquad{}\Rightarrow
\exists!k\,\bigl(k:X\to E\land e\circ k=h\bigr)
\Bigr)
\Bigr).
\end{aligned}
\]
\end{definicio}

\begin{notacio}[Equalitzador escollit]\label{not:equalitzador}
Per a les aplicacions paral·leles \(f,g:A\to B\), denotarem per
\[
e:E\hookrightarrow A
\]
l'equalitzador construït mitjançant l'aplicació característica \(\operatorname{eq}_B\circ\langle f,g\rangle\).

La propietat universal afirma que les aplicacions \(h:X\to A\) que igualen \(f\) i \(g\), és a dir, les que satisfan
\[
f\circ h=g\circ h,
\]
són exactament les que passen per l'equalitzador, i hi passen d'una única manera:
\[
h=e\circ\bar h.
\]
Aquesta propietat queda representada pel diagrama commutatiu
\begin{center}
\begin{tikzpicture}[scale=1.0, every node/.style={font=\small}, >=Stealth]

  \node (E) at (0,0) {$E$};
  \node (A) at (4,0) {$A$};
  \node (X) at (2,3.46) {$X$};
  \node (B) at (8,0) {$B$};

  \draw[->, thick, dashed]
    (X) -- node[left] {$\bar h$} (E);

  \draw[->, thick]
    (X) -- node[right] {$h$} (A);

  \draw[{Hooks[right]}-{Stealth[scale=1.2]}, thick]
    (E) -- node[above] {$e$} (A);

  \draw[->, thick]
    ([xshift=4pt,yshift=2pt]A.east) --
    node[above] {$f$}
    ([xshift=-4pt,yshift=2pt]B.west);

  \draw[->, thick]
    ([xshift=4pt,yshift=-2pt]A.east) --
    node[below] {$g$}
    ([xshift=-4pt,yshift=-2pt]B.west);

\end{tikzpicture}
\end{center}
La fletxa discontínua expressa l'existència i la unicitat de \(\bar h\).
\end{notacio}

\begin{teorema}[Existència de l'equalitzador]\label{teo:equalitzador-universal}
La part \(e:E\hookrightarrow A\) classificada per \(\operatorname{eq}_B\circ\langle f,g\rangle\) és un equalitzador de \(f\) i \(g\). És a dir:
\begin{enumerate}
\item \(f\circ e=g\circ e\).
\item Si \(h:X\to A\) satisfà \(f\circ h=g\circ h\), existeix una única aplicació \(\bar h:X\to E\) tal que
\[
h=e\circ\bar h.
\]
\end{enumerate}
\end{teorema}

\begin{prova}
Com que \(e\in_E e\), l'axioma de representació dona
\[
\operatorname{eq}_B\circ\langle f,g\rangle\circ e
=1\circ {!}_E.
\]
Això significa que
\[
\langle f\circ e,g\circ e\rangle\in_E\Delta_B.
\]
Per la \hyperref[prop:diagonal-igualtat]{proposició~\ref*{prop:diagonal-igualtat}},
\[
f\circ e=g\circ e.
\]

Ara sigui \(h:X\to A\) tal que \(f\circ h=g\circ h\). La mateixa proposició implica
\[
\langle f\circ h,g\circ h\rangle\in_X\Delta_B.
\]
Per l'associativitat i la propietat universal del producte,
\[
\langle f,g\rangle\circ h
=\langle f\circ h,g\circ h\rangle.
\]
Per tant,
\[
\operatorname{eq}_B\circ\langle f,g\rangle\circ h
=1\circ {!}_X.
\]
Com que aquesta aplicació és la característica de \(e\), obtenim \(h\in_X e\). Per la definició de pertinença, existeix \(\bar h:X\to E\) tal que
\[
h=e\circ\bar h.
\]
La unicitat de \(\bar h\) se segueix del fet que \(e\) és un monomorfisme.
\end{prova}

\section{Pullbacks o productes fibrats}\label{pullbacks}

Siguin
\[
A\xrightarrow{f}C\xleftarrow{g}B
\]
dues aplicacions amb el mateix codomini. Un conjunt \(P\), juntament amb aplicacions
\[
q_A:P\to A,
\qquad
q_B:P\to B,
\]
determina el quadrat
\begin{center}
\begin{tikzpicture}[scale=1.05, every node/.style={font=\small}, >=Stealth]
  \node (P) at (0,2.2) {\(P\)};
  \node (B) at (4.0,2.2) {\(B\)};
  \node (A) at (0,0) {\(A\)};
  \node (C) at (4.0,0) {\(C\)};

  \draw[->, thick] (P) -- node[above] {\(q_B\)} (B);
  \draw[->, thick] (P) -- node[left] {\(q_A\)} (A);
  \draw[->, thick] (B) -- node[right] {\(g\)} (C);
  \draw[->, thick] (A) -- node[below] {\(f\)} (C);
\end{tikzpicture}
\end{center}
Direm que el quadrat és commutatiu si
\[
f\circ q_A=g\circ q_B.
\]

\begin{definicio}[Pullback o producte fibrat]\index{Pullback o producte fibrat}\label{def:pullback}
El quadrat commutatiu anterior és \emph{cartesià}, i el conjunt \(P\) és un \emph{pullback} o \emph{producte fibrat} de \(f\) i \(g\), si compleix la propietat universal següent: per a qualsevol parell d'aplicacions \(u:X\to A\) i \(v:X\to B\) que fan commutar el quadrat exterior, existeix una única aplicació \(w:X\to P\) que té \(u\) i \(v\) com a components. Formalment,
\[
\begin{aligned}
&\operatorname{Pullback}(P,q_A,q_B,f,g)
\\
\iff {}&
\exists A\,\exists B\,\exists C\,\Bigl(
 f:A\to C\land g:B\to C
\land q_A:P\to A\land q_B:P\to B
\\
&{}\land f\circ q_A=g\circ q_B
\\
&{}\land\forall X\,\forall u\,\forall v\,
\Bigl(
 u:X\to A\land v:X\to B\land f\circ u=g\circ v
\\
&\qquad{}\Rightarrow
\exists!w\,\bigl(
 w:X\to P\land q_A\circ w=u\land q_B\circ w=v
\bigr)
\Bigr)
\Bigr).
\end{aligned}
\]
\end{definicio}

Considerem ara el producte \(A\times B\) i les dues aplicacions
\[
f\circ p_A,
\qquad
g\circ p_B
\]
de \(A\times B\) en \(C\). Sigui
\[
e:P\hookrightarrow A\times B
\]
un equalitzador d'aquestes dues aplicacions. Definim
\[
q_A=p_A\circ e,
\qquad
q_B=p_B\circ e.
\]

\begin{notacio}[Pullback escollit]\label{not:pullback}
Denotarem un pullback escollit de \(f:A\to C\) i \(g:B\to C\) per
\[
A\times_C B,
\]
amb projeccions
\[
A\xleftarrow{q_A}A\times_C B\xrightarrow{q_B}B.
\]
La propietat universal queda representada pel diagrama
\begin{center}
\begin{tikzpicture}[
    >=Stealth, % Tipus de fletxa similar al de LaTeX/AMS
    node distance=2.5cm and 3.5cm,
    auto
]
    % Nodes principals (el quadrat de baix a la dreta)
    \node (AxCB) {$A \times_C B$};
    \node (B) [right=of AxCB] {$B$};
    \node (A) [below=of AxCB] {$A$};
    \node (C) [below=of B] {$C$};

    % Node extern (X) a dalt i a la meitat de la distància a l'esquerra
    \node (X) [above left=2.2cm and 2.5cm of AxCB] {$X$};

    % Fletxes del quadrat
    \draw[->] (AxCB) -- node {$q_B$} (B);
    \draw[->] (AxCB) -- node {$q_A$} (A);
    \draw[->] (B) -- node {$g$} (C);
    \draw[->] (A) -- node [below] {$f$} (C);

    % Fletxes des de X
    \draw[->, dashed] (X) -- node {$w$} (AxCB);
    \draw[->] (X) -- node {$v$} (B);
    \draw[->] (X) -- node [below left] {$u$} (A);

\end{tikzpicture}
\end{center}
La fletxa discontínua expressa que \(w\) existeix i és única sempre que \(f\circ u=g\circ v\).
\end{notacio}

\begin{teorema}[Existència del pullback]\label{teo:pullback-universal}
El conjunt \(P\) construït a partir de l'equalitzador de \(f\circ p_A\) i \(g\circ p_B\), juntament amb les aplicacions \(q_A\) i \(q_B\), és un pullback de \(f:A\to C\) i \(g:B\to C\).
\end{teorema}

\begin{prova}
Per la definició de l'equalitzador,
\[
f\circ q_A=g\circ q_B.
\]
Siguin \(u:X\to A\) i \(v:X\to B\) aplicacions que satisfan
\[
f\circ u=g\circ v.
\]
La propietat universal del producte determina una única aplicació
\[
\langle u,v\rangle:X\to A\times B
\]
tal que
\[
p_A\circ\langle u,v\rangle=u,
\qquad
p_B\circ\langle u,v\rangle=v.
\]
La hipòtesi implica
\[
(f\circ p_A)\circ\langle u,v\rangle
=(g\circ p_B)\circ\langle u,v\rangle.
\]
Per la propietat universal de l'equalitzador, existeix una única aplicació \(w:X\to P\) tal que
\[
\langle u,v\rangle=e\circ w.
\]
Aleshores
\[
q_A\circ w=p_A\circ e\circ w=p_A\circ\langle u,v\rangle=u
\]
i, anàlogament,
\[
q_B\circ w=v.
\]

Si \(w'\) és una altra aplicació amb les mateixes propietats, aleshores \(e\circ w\) i \(e\circ w'\) tenen les mateixes dues components. Per la \hyperref[prop:projeccions-separen]{proposició~\ref*{prop:projeccions-separen}},
\[
e\circ w=e\circ w'.
\]
Com que \(e\) és un monomorfisme, \(w=w'\).
\end{prova}

\begin{observacio}[Productes, fibres i canvis de base]\label{obs:pullback-exemples}
Els pullbacks recuperen diverses construccions habituals.
\begin{enumerate}
\item Si \(C=\mathbf1\), \(f={!}_A:A\to\mathbf1\) i \(g={!}_B:B\to\mathbf1\), la condició de commutativitat és automàtica. La propietat universal del pullback coincideix amb la del producte. Per tant,
\[
A\times_{\mathbf1}B\cong A\times B.
\]
\item Donades \(f:A\to C\) i \(g:B\to C\), la projecció
\[
q_B:A\times_C B\to B
\]
s'anomena el \emph{canvi de base}\index{Canvi de base} de \(f\) al llarg de \(g\), i la denotarem també \(g^{*}(f)\). El pullback substitueix l'aplicació \(f\) sobre \(C\) per una aplicació sobre \(B\). Dualment dins del mateix quadrat, \(q_A\) és el canvi de base de \(g\) al llarg de \(f\).
\item Si \(c:\mathbf1\to C\) és un element global i \(f:A\to C\), el producte fibrat
\[
\mathbf1\times_C A
\]
és la \emph{fibra}\index{Fibra} de \(f\) sobre \(c\). En la interpretació clàssica, els seus punts corresponen als elements \(a\in A\) tals que \(f(a)=c\).
\end{enumerate}
Aquestes lectures pertanyen al metallenguatge; la definició interna continua essent la propietat universal del quadrat cartesià.
\end{observacio}

\section{Imatge inversa d'una part}\label{imatge-inversa}

Sigui
\[
f:A\to B
\]
una aplicació i sigui
\[
m:S\hookrightarrow B
\]
una part de \(B\). Considerem un pullback de \(f\) i \(m\), amb projeccions
\[
q_A:P\to A,
\qquad
q_S:P\to S.
\]

\begin{proposicio}\label{prop:pullback-mono}
L'aplicació \(q_A:P\to A\) és un monomorfisme.
\end{proposicio}

\begin{prova}
Suposem que \(q_A\circ u=q_A\circ v\). Com que el quadrat commuta,
\[
m\circ q_S\circ u
=f\circ q_A\circ u
=f\circ q_A\circ v
=m\circ q_S\circ v.
\]
Com que \(m\) és un monomorfisme,
\[
q_S\circ u=q_S\circ v.
\]
Les aplicacions \(u\) i \(v\) tenen, doncs, les mateixes composicions amb \(q_A\) i \(q_S\). La unicitat de la propietat universal del pullback implica \(u=v\).
\end{prova}

Per tant, \(q_A\) és una part de \(A\), i el quadrat cartesià es representa per
\begin{center}
\begin{tikzpicture}[scale=1.05, every node/.style={font=\small}, >=Stealth]
  \node (P) at (0,2.2) {\(P\)};
  \node (S) at (4.0,2.2) {\(S\)};
  \node (A) at (0,0) {\(A\)};
  \node (B) at (4.0,0) {\(B\)};

  \draw[->, thick] (P) -- node[above] {\(q_S\)} (S);
  \draw[{Hooks[right]}-{Stealth[scale=1.2]}, thick] (P) -- node[left] {\(q_A\)} (A);
  \draw[{Hooks[right]}-{Stealth[scale=1.2]}, thick] (S) -- node[right] {\(m\)} (B);
  \draw[->, thick] (A) -- node[below] {\(f\)} (B);
\end{tikzpicture}
\end{center}

\begin{definicio}[Imatge inversa]\index{Imatge inversa}\label{def:imatge-inversa}
Sigui \(m:S\hookrightarrow B\) una part i sigui \(f:A\to B\). Una part \(q_A:P\hookrightarrow A\) és una \emph{imatge inversa} de \(m\) per \(f\) si existeix una aplicació \(q_S:P\to S\) tal que el quadrat anterior és cartesià. Formalment,
\[
\begin{aligned}
\operatorname{ImInv}(q_A,f,m)
\iff{}&
\exists A\,\exists B\,\exists P\,\exists q_S
\\[-0.2em]
&{}\Bigl(
\operatorname{Part}(m,B)
\land\operatorname{Part}(q_A,A)
\\
&\hspace{2.5em}{}\land f:A\to B
\land\operatorname{Pullback}(P,q_A,q_S,f,m)
\Bigr).
\end{aligned}
\]
\end{definicio}

\begin{notacio}[Imatge inversa]\label{not:imatge-inversa}
Denotarem la imatge inversa escollida per
\[
f^{-1}(m):P\hookrightarrow A.
\]
\end{notacio}

\begin{proposicio}\label{prop:pertinenca-imatge-inversa}
Per a tot conjunt \(X\) i tota aplicació \(x:X\to A\),
\[
x\in_X f^{-1}(m)
\quad\Longleftrightarrow\quad
f\circ x\in_X m.
\]
\end{proposicio}

\begin{prova}
Suposem que \(x\in_X f^{-1}(m)\). Existeix una aplicació \(w:X\to P\) tal que
\[
x=q_A\circ w.
\]
Com que el quadrat commuta,
\[
f\circ x=f\circ q_A\circ w=m\circ q_S\circ w.
\]
Per tant, \(f\circ x\in_X m\).

Recíprocament, suposem que \(f\circ x\in_X m\). Existeix una aplicació \(y:X\to S\) tal que
\[
f\circ x=m\circ y.
\]
Per la propietat universal del pullback, existeix una aplicació \(w:X\to P\) tal que
\[
q_A\circ w=x,
\qquad
q_S\circ w=y.
\]
Així, \(x\in_X f^{-1}(m)\).
\end{prova}

\begin{proposicio}\label{prop:caracteristica-imatge-inversa}
L'aplicació característica de la imatge inversa és
\[
\chi_{f^{-1}(m)}=\chi_m\circ f.
\]
\end{proposicio}

\begin{prova}
Per a tot \(x:X\to A\),
\[
\begin{aligned}
x\in_X f^{-1}(m)
&\Longleftrightarrow f\circ x\in_X m\\
&\Longleftrightarrow \chi_m\circ f\circ x=1\circ {!}_X.
\end{aligned}
\]
La unicitat de l'aplicació característica dona
\[
\chi_{f^{-1}(m)}=\chi_m\circ f.
\]
\end{prova}

\subsection{Intersecció de parts}\label{sec:interseccio}

Siguin
\[
m:S\hookrightarrow A,
\qquad
n:T\hookrightarrow A
\]
dues parts de \(A\). Considerem un pullback de \(m\) i \(n\). Per la \hyperref[prop:pullback-mono]{proposició~\ref*{prop:pullback-mono}}, aplicada en les dues orientacions del quadrat, les projeccions \(q_S:P\to S\) i \(q_T:P\to T\) són monomorfismes. Per tant, el quadrat cartesià es representa per
\begin{center}
\begin{tikzpicture}[scale=1.05, every node/.style={font=\small}, >=Stealth]
  \node (P) at (0,2.2) {\(P\)};
  \node (T) at (4.0,2.2) {\(T\)};
  \node (S) at (0,0) {\(S\)};
  \node (A) at (4.0,0) {\(A\)};

  \draw[{Hooks[right]}-{Stealth[scale=1.2]}, thick] (P) -- node[above] {\(q_T\)} (T);
  \draw[{Hooks[right]}-{Stealth[scale=1.2]}, thick] (P) -- node[left] {\(q_S\)} (S);
  \draw[{Hooks[right]}-{Stealth[scale=1.2]}, thick] (T) -- node[right] {\(n\)} (A);
  \draw[{Hooks[right]}-{Stealth[scale=1.2]}, thick] (S) -- node[below] {\(m\)} (A);
\end{tikzpicture}
\end{center}
La composició
\[
m\circ q_S=n\circ q_T:P\to A
\]
és un monomorfisme i, per tant, una part de \(A\).

\begin{definicio}[Intersecció]\index{Intersecció}\label{def:interseccio}
Siguin \(m:S\hookrightarrow A\) i \(n:T\hookrightarrow A\). Una part \(i:P\hookrightarrow A\) és una \emph{intersecció} de \(m\) i \(n\) si existeixen aplicacions \(q_S:P\to S\) i \(q_T:P\to T\) que formen un pullback de \(m\) i \(n\), i \(i=m\circ q_S=n\circ q_T\). Formalment,
\[
\begin{aligned}
\operatorname{Inter}(i,m,n,A)
\iff
\exists P\,\exists q_S\,\exists q_T\,\Bigl(&
\operatorname{Pullback}(P,q_S,q_T,m,n)
\\
&{}\land\operatorname{Part}(i,A)
\land i=m\circ q_S
\land i=n\circ q_T
\Bigr).
\end{aligned}
\]
\end{definicio}

\begin{notacio}[Intersecció]\label{not:interseccio}
Denotarem una intersecció escollida per
\[
m\cap n:=m\circ q_S=n\circ q_T:P\hookrightarrow A.
\]
\end{notacio}

\begin{proposicio}\label{prop:pertinenca-interseccio}
Per a tot conjunt \(X\) i tota aplicació \(x:X\to A\),
\[
x\in_X(m\cap n)
\quad\Longleftrightarrow\quad
x\in_X m\ \land\ x\in_X n.
\]
\end{proposicio}

\begin{prova}
Si \(x\in_X(m\cap n)\), existeix \(w:X\to P\) tal que
\[
x=m\circ q_S\circ w=n\circ q_T\circ w.
\]
Per tant, \(x\in_Xm\) i \(x\in_Xn\).

Recíprocament, si \(x\in_Xm\) i \(x\in_Xn\), existeixen aplicacions \(u:X\to S\) i \(v:X\to T\) tals que
\[
x=m\circ u=n\circ v.
\]
Per la propietat universal del pullback, existeix \(w:X\to P\) tal que
\[
q_S\circ w=u,
\qquad
q_T\circ w=v.
\]
Aleshores
\[
x=m\circ q_S\circ w,
\]
i, per tant, \(x\in_X(m\cap n)\).
\end{prova}

\begin{resumcap}
\apartat{Idees centrals}
\begin{enumerate}
  \item Els monomorfismes i els epimorfismes es defineixen per cancel·lació, només amb la composició; les seccions són monomorfismes i les retraccions, epimorfismes. Que mono i epi alhora sigui isomorfisme no és una obvietat: és un teorema que necessita l'axioma de representació.
  \item Una \emph{part} d'un conjunt és un monomorfisme cap a ell, no un conjunt: dues parts amb el mateix domini poden situar-lo de maneres diferents. La inclusió i la pertinença són factoritzacions, amb testimoni únic.
  \item Dues parts són equivalents quan cadascuna es factoritza per l'altra; aleshores els dominis són isomorfs per un únic isomorfisme compatible. El \emph{subconjunt} clàssic és, en el metallenguatge, una classe de parts equivalents.
  \item Les sumes i els productes es postulen per la seva propietat universal; l'aplicació definida per casos \([f_0,f_1]\) i l'aparellament \(\langle f,g\rangle\) en són les dues cares, i la unicitat llevat d'un únic isomorfisme és el patró de totes les construccions universals.
  \item El conjunt \(\mathbf2=\mathbf1+\mathbf1\) no es postula: és una suma. L'axioma de representació el converteix en el classificador de les parts: cada part té una única característica \(\chi_m\colon A\to\mathbf2\), i cada aplicació cap a \(\mathbf2\) classifica alguna part.
  \item Equalitzadors i pullbacks no demanen cap axioma: es construeixen a partir de la diagonal, els productes i la representació. D'ells surten la imatge inversa d'una part, amb \(\chi_{f^{-1}(m)}=\chi_m\circ f\), i la intersecció.
  \item La injectivitat clàssica coincideix amb ser monomorfisme gràcies a la separació global; l'exhaustivitat clàssica implica ser epimorfisme, i el recíproc es demostrarà al capítol~\ref{cap:valors}. L'exhaustivitat sobre elements generalitzats és més forta: equival a tenir secció.
\end{enumerate}
\apartat{Errors típics} Identificar una part amb el seu domini (\(S\) no diu res de com és dins \(A\): la part és \(m\)); pensar que \(m\subseteq_A n\) demana que els dominis estiguin inclosos l'un en l'altre com a conjunts, quan demana una factorització; creure que mono $+$ epi és isomorfisme \enquote{per definició}, quan cal l'axioma de representació; afirmar que les injeccions d'una suma són monomorfismes abans de tenir-ne una retracció o la disjunció; confondre l'aparellament \(\langle f,g\rangle\) (cap a un producte) amb la definició per casos \([f,g]\) (des d'una suma); llegir \(\chi_m\circ x\ne1\circ{!}_X\) com \(\chi_m\circ x=0\circ{!}_X\), que encara no s'ha demostrat; i tractar un epimorfisme com si fos exhaustiu sobre punts, o escindit, sense les hipòtesis corresponents.
\apartat{Lectura recomanada} \cite[cap.~2 i~3]{lawvere-rosebrugh} per a les sumes, els monomorfismes, les parts i els límits finits inversos (productes, equalitzadors, pullbacks), que és l'ordre que aquest capítol segueix; \cite[article~II, sessions~6--10, i article~IV, sessions~19--25]{lawvere-schanuel} per a seccions i retraccions, i per a productes i sumes amb molts exemples elementals; \cite[§2.1, §§2.2--2.3, §3.1--3.3 i §5.1]{awodey} per a monomorfismes i epimorfismes, objectes inicials i terminals, productes, coproductes, equalitzadors i pullbacks en una categoria general; \cite[pàg.~11]{bell2006} per a la lectura de les parts com a subobjectes i de \(\mathbf2\) com a classificador; i \cite[§§2--6]{halmos} per a les mateixes construccions fetes amb elements: especificació, parells, unions, complements i parells ordenats.
\end{resumcap}

\chapter[Valors de veritat i construccions finites]{Valors de veritat, parts\texorpdfstring{\\}{ }i construccions finites}\label{cap:valors}

\begin{objectius}
\apartat{Prerequisits} Del capítol~\ref{cap:parts}: parts, inclusió i pertinença generalitzada, sumes i productes amb les seves propietats universals, l'axioma de representació i les aplicacions característiques, i la construcció d'equalitzadors, pullbacks i imatges inverses. Del capítol~\ref{cap:aplicacions}: el conjunt inicial, la no-degeneració i la separació global. De la lògica elemental, les taules de veritat de la negació, la conjunció i la disjunció.
\apartat{Objectius} En acabar el capítol, el lector hauria de poder:
\begin{itemize}
  \item enunciar l'axioma de disjunció i estabilitat de les sumes i explicar què hi afegeix respecte de la propietat universal;
  \item demostrar que \(\mathbf0\) és estricte, que \(0\ne1\) a \(\mathbf2\) i que un conjunt sense punts és inicial, i deduir-ne que tot punt d'una suma prové d'un sol sumand i que tot epimorfisme és exhaustiu sobre punts;
  \item treballar amb la part buida, la part total i el complement d'una part, i amb el quadrat característic com a pullback;
  \item definir la negació, la conjunció i la disjunció com a aplicacions sobre \(\mathbf2\), calcular-ne les taules i traduir les operacions entre parts a operacions entre característiques;
  \item demostrar les lleis booleanes de les parts per extensionalitat de les taules de veritat;
  \item definir con, límit, cocon i colímit d'un diagrama finit, i construir tots els límits finits a partir de productes i equalitzadors, i tots els colímits finits a partir de sumes i coequalitzadors;
  \item distingir quines construccions del capítol demanen un axioma nou i quines no.
\end{itemize}
\end{objectius}

L'axioma de les sumes garanteix que tot parell de conjunts té una suma, i
l'axioma de representació converteix
\(\mathbf2=\mathbf1+\mathbf1\) en el conjunt que classifica les parts.
Aquests axiomes, però, encara no expressen que les dues aplicacions
canòniques d'una suma siguin disjuntes ni que aquesta descomposició es
conservi en fer un canvi de base. En particular, fins ara no hem pogut
demostrar que els dos elements globals canònics
\[
0,1:\mathbf1\to\mathbf2
\]
siguin diferents.

En les cinc primeres seccions reforçarem el comportament de les sumes amb
un axioma nou. A partir d'aquest axioma obtindrem la distinció entre
\(0\) i \(1\), identificarem les parts buida i total, construirem el
complement d'una part i recuperarem, en el metallenguatge, les operacions
i les lleis de l'àlgebra booleana de les parts.

\section{Sumes disjuntes i estables}\label{sec:sumes-disjuntes}

Considerem una suma
\[
A_0\xrightarrow{i_0}S\xleftarrow{i_1}A_1.
\]
La propietat universal de la suma determina les aplicacions que surten de
\(S\). Ara introduirem dues propietats que descriuen com les dues
injeccions se situen dins de \(S\) i com es comporten en fer
un canvi de base.

\begin{definicio}[Suma disjunta]\index{Suma disjunta}\label{def:suma-disjunta}
Una suma
\[
A_0\xrightarrow{i_0}S\xleftarrow{i_1}A_1
\]
és \emph{disjunta} si les dues injeccions són
monomorfismes i el seu pullback és inicial. Formalment,
\[
\begin{aligned}
&\operatorname{SumaDisj}(S,i_0,i_1,A_0,A_1)
\\
\iff{}&
\operatorname{Suma}(S,i_0,i_1,A_0,A_1)
\land\operatorname{Mono}(i_0)
\land\operatorname{Mono}(i_1)
\\
&{}\land
\exists P\,\exists q_0\,\exists q_1\,
\bigl(
\operatorname{Pullback}(P,q_0,q_1,i_0,i_1)
\land\operatorname{Inicial}(P)
\bigr).
\end{aligned}
\]
\end{definicio}

Geomètricament, la disjunció afirma que el quadrat cartesià
\begin{center}
\begin{tikzpicture}[scale=1.05, every node/.style={font=\small}, >=Stealth]
  \node (P) at (0,2.2) {\(P\cong\mathbf0\)};
  \node (A1) at (4.0,2.2) {\(A_1\)};
  \node (A0) at (0,0) {\(A_0\)};
  \node (S) at (4.0,0) {\(S\)};

  \draw[{Hooks[right]}-{Stealth[scale=1.2]}, thick] (P) -- node[above] {\(q_1\)} (A1);
  \draw[{Hooks[right]}-{Stealth[scale=1.2]}, thick] (P) -- node[left] {\(q_0\)} (A0);
  \draw[{Hooks[right]}-{Stealth[scale=1.2]}, thick] (A1) -- node[right] {\(i_1\)} (S);
  \draw[{Hooks[right]}-{Stealth[scale=1.2]}, thick] (A0) -- node[below] {\(i_0\)} (S);
\end{tikzpicture}
\end{center}
té un conjunt inicial com a vèrtex superior esquerre. Això expressa que
les dues parts de \(S\) representades per \(i_0\) i \(i_1\) no tenen
cap part comuna.

\begin{definicio}[Suma estable]\index{Suma estable}\label{def:suma-estable}
Una suma
\[
A_0\xrightarrow{i_0}S\xleftarrow{i_1}A_1
\]
és \emph{estable} si, per a tota aplicació \(f:B\to S\), les imatges
inverses de les dues injeccions formen una suma amb vèrtex
\(B\). Formalment,
\[
\begin{aligned}
&\operatorname{SumaEst}(S,i_0,i_1,A_0,A_1)
\\
\iff{}&
\operatorname{Suma}(S,i_0,i_1,A_0,A_1)
\\
&{}\land
\forall B\,\forall f\,
\Bigl(
\operatorname{Conj}(B)\land f:B\to S
\\
&\qquad{}\Rightarrow
\exists B_0\,\exists B_1\,
\exists j_0\,\exists j_1\,
\exists k_0\,\exists k_1\,
\Bigl(
\operatorname{Pullback}(B_0,j_0,k_0,f,i_0)
\\
&\qquad\qquad{}\land
\operatorname{Pullback}(B_1,j_1,k_1,f,i_1)
\land
\operatorname{Suma}(B,j_0,j_1,B_0,B_1)
\Bigr)
\Bigr).
\end{aligned}
\]
\end{definicio}

La definició queda representada pels dos quadrats cartesians
\begin{center}
\begin{tikzpicture}[scale=0.95, every node/.style={font=\small}, >=Stealth]
  \node (B0) at (0,2.2) {\(B_0\)};
  \node (A0) at (3.2,2.2) {\(A_0\)};
  \node (B) at (0,0) {\(B\)};
  \node (S0) at (3.2,0) {\(S\)};

  \node (B1) at (5.2,2.2) {\(B_1\)};
  \node (A1) at (8.4,2.2) {\(A_1\)};
  \node (Bbis) at (5.2,0) {\(B\)};
  \node (S1) at (8.4,0) {\(S\)};

  \draw[->, thick] (B0) -- node[above] {\(k_0\)} (A0);
  \draw[{Hooks[right]}-{Stealth[scale=1.2]}, thick] (B0) -- node[left] {\(j_0\)} (B);
  \draw[{Hooks[right]}-{Stealth[scale=1.2]}, thick] (A0) -- node[right] {\(i_0\)} (S0);
  \draw[->, thick] (B) -- node[below] {\(f\)} (S0);

  \draw[->, thick] (B1) -- node[above] {\(k_1\)} (A1);
  \draw[{Hooks[right]}-{Stealth[scale=1.2]}, thick] (B1) -- node[left] {\(j_1\)} (Bbis);
  \draw[{Hooks[right]}-{Stealth[scale=1.2]}, thick] (A1) -- node[right] {\(i_1\)} (S1);
  \draw[->, thick] (Bbis) -- node[below] {\(f\)} (S1);
\end{tikzpicture}
\end{center}
i per l'afirmació que les dues fletxes obtingudes tornen a confluir en
una suma:
\[
B_0\xrightarrow{j_0}B\xleftarrow{j_1}B_1.
\]
Així, una descomposició de \(S\) en dues parts es transporta enrere al
llarg de qualsevol aplicació \(f:B\to S\).

\begin{axioma}[Disjunció i estabilitat de les sumes]\label{ax:sumes-disjuntes}
Tota suma és disjunta i estable. Formalment,
\[
\begin{aligned}
\forall S\,\forall i_0\,\forall i_1\,
\forall A_0\,\forall A_1\,
\Bigl(&
\operatorname{Suma}(S,i_0,i_1,A_0,A_1)
\\
&{}\Rightarrow
\operatorname{SumaDisj}(S,i_0,i_1,A_0,A_1)
\land
\operatorname{SumaEst}(S,i_0,i_1,A_0,A_1)
\Bigr).
\end{aligned}
\]
\end{axioma}

\begin{observacio}[Informació aportada per l'axioma]\label{obs:necessitat-axioma}
La propietat universal de la suma determina com es defineix una
aplicació que surt de \(A_0+A_1\), però no proporciona per si sola una
construcció que demostri que les dues injeccions són
disjuntes ni que els seus canvis de base continuen formant una suma.
L'axioma~\ref{ax:sumes-disjuntes} incorpora aquestes dues propietats i
fa que les sumes es comportin estructuralment com les unions disjuntes
de la teoria clàssica de conjunts.
\end{observacio}

\section{La distinció entre \texorpdfstring{\(0\) i \(1\)}{0 i 1}}\label{sec:zero-u}

Abans de distingir els dos elements globals de \(\mathbf2\), obtindrem
una conseqüència estructural de l'estabilitat: si existeix una aplicació
d'un conjunt cap a \(\mathbf0\), aquest conjunt és inicial.

\begin{definicio}[Conjunt inicial estricte]\index{Conjunt inicial estricte}\label{def:inicial-estricte}
Un conjunt inicial \(I\) és \emph{estricte} si tota aplicació amb
codomini \(I\) és un isomorfisme. Formalment,
\[
\begin{aligned}
\operatorname{InicialEstricte}(I)
&\iff
\operatorname{Inicial}(I)
\\[-0.2em]
&\quad{}\land
\forall X\,\forall f\,
\Bigl(
\operatorname{Conj}(X)
\land\Dom(f,X)\land\Cod(f,I)
\Rightarrow\operatorname{Iso}(f)
\Bigr).
\end{aligned}
\]
\end{definicio}

\begin{proposicio}[Estrictesa del conjunt inicial]\label{prop:inicial-estricte}
El conjunt inicial escollit \(\mathbf0\) és estricte.
\end{proposicio}

\begin{prova}
Sigui \(f:X\to\mathbf0\). Les dues identitats
\[
\mathbf0\xrightarrow{\id_{\mathbf0}}\mathbf0
\xleftarrow{\id_{\mathbf0}}\mathbf0
\]
formen una suma de \(\mathbf0\) i \(\mathbf0\): donades dues
aplicacions amb domini \(\mathbf0\), totes dues són l'única aplicació
que surt de l'inicial.

Fem el canvi de base d'aquesta suma al llarg de \(f\). El pullback de
\(\id_{\mathbf0}\) al llarg de \(f\) és \(X\), amb projecció
\(\id_X:X\to X\). Per l'estabilitat, les dues identitats formen una
suma
\[
X\xrightarrow{\id_X}X\xleftarrow{\id_X}X.
\]
Aquesta suma és disjunta. El pullback de les dues identitats és el
mateix conjunt \(X\), i la disjunció implica que \(X\) és inicial.
Per la unicitat dels conjunts inicials llevat d'un únic isomorfisme, existeix un isomorfisme \(u:X\to\mathbf0\). Com que \(X\) és inicial, només hi ha una aplicació de \(X\) en \(\mathbf0\); per tant, \(f=u\) i \(f\) és un isomorfisme.
\end{prova}

\begin{teorema}[Distinció dels valors de veritat]\label{teo:zero-diferent-u}
Els dos elements globals canònics
\[
0,1:\mathbf1\to\mathbf2
\]
són diferents.
\end{teorema}

\begin{prova}
La suma
\[
\mathbf1\xrightarrow{0}\mathbf2\xleftarrow{1}\mathbf1
\]
és disjunta. Per tant, un pullback de \(0\) i \(1\) és inicial.
Suposem que \(0=1\). Aleshores el quadrat
\begin{center}
\begin{tikzpicture}[scale=1.05, every node/.style={font=\small}, >=Stealth]
  \node (OneA) at (0,2.0) {\(\mathbf1\)};
  \node (OneB) at (3.8,2.0) {\(\mathbf1\)};
  \node (OneC) at (0,0) {\(\mathbf1\)};
  \node (Two) at (3.8,0) {\(\mathbf2\)};

  \draw[->, thick] (OneA) -- node[above] {\(\id_{\mathbf1}\)} (OneB);
  \draw[->, thick] (OneA) -- node[left] {\(\id_{\mathbf1}\)} (OneC);
  \draw[{Hooks[right]}-{Stealth[scale=1.2]}, thick] (OneB) -- node[right] {\(1\)} (Two);
  \draw[{Hooks[right]}-{Stealth[scale=1.2]}, thick] (OneC) -- node[below] {\(0\)} (Two);
\end{tikzpicture}
\end{center}
seria cartesià: com que aquesta aplicació és un monomorfisme, el
pullback d'ella mateixa amb ella mateixa és el seu domini. Per la unicitat dels
pullbacks llevat d'isomorfisme, \(\mathbf1\) seria isomorf al pullback
inicial de \(0\) i \(1\). Així,
\[
\mathbf1\cong\mathbf0,
\]
en contradicció amb l'axioma de no-degeneració. Per tant,
\(0\neq1\).
\end{prova}

La conseqüència principal de l'estabilitat no és només una afirmació
sobre els punts globals d'una suma. És una descomposició de qualsevol
element generalitzat.

\begin{teorema}[Descomposició per canvi de base]\label{teo:descomposicio-canvi-base}
Sigui
\[
A_0\xrightarrow{i_0}A_0+A_1\xleftarrow{i_1}A_1
\]
una suma i sigui \(x:X\to A_0+A_1\). Considerem els pullbacks
\[
X_0:=X\times_{A_0+A_1}A_0,
\qquad
X_1:=X\times_{A_0+A_1}A_1,
\]
amb projeccions \(j_0:X_0\to X\), \(j_1:X_1\to X\) i aplicacions \(k_0:X_0\to A_0\), \(k_1:X_1\to A_1\). Aleshores
\[
X_0\xrightarrow{j_0}X\xleftarrow{j_1}X_1
\]
és una suma disjunta. En particular, l'aplicació induïda
\[
[j_0,j_1]:X_0+X_1\xrightarrow{\sim}X
\]
és un isomorfisme.
\end{teorema}

\begin{prova}
Per l'estabilitat de la suma, les dues projeccions
\(j_0\) i \(j_1\) formen una suma amb vèrtex \(X\). Per
l'axioma~\ref{ax:sumes-disjuntes}, aquesta suma és també disjunta.
Finalment, \(X_0+X_1\) i \(X\) són dos representants d'una mateixa
propietat universal; la unicitat de les sumes llevat d'un únic
isomorfisme dona el resultat.
\end{prova}

\begin{observacio}[Lectura dels elements generalitzats d'una suma]\label{obs:elements-generalitzats-suma}
Un element generalitzat
\[
x:X\to A_0+A_1
\]
no ha de provenir íntegrament d'un sol sumand. El
teorema~\ref{teo:descomposicio-canvi-base} afirma, en canvi, que el seu
estadi es descompon en una suma
\[
X\cong X_0+X_1,
\]
de manera que la restricció de \(x\) a \(X_0\) es factoritza a través
de \(i_0\), i la restricció a \(X_1\) es factoritza a través de
\(i_1\). Aquesta és la formulació estructural de l'exhaustivitat de les
sumes.
\end{observacio}

El teorema anterior té dues conseqüències sobre els elements globals. La primera usa, a més, l'axioma de separació global i caracteritza els conjunts sense elements globals.

\begin{lema}[Lema del buit]\label{lema:buit}
Tot conjunt sense elements globals és inicial.
\end{lema}

\begin{prova}
Sigui \(X\) un conjunt sense elements globals. Les aplicacions
\[
0\circ {!}_X,
\qquad
1\circ {!}_X
\]
de \(X\) a \(\mathbf2\) coincideixen sobre tots els elements globals de \(X\) ---no n'hi ha cap--- i, per l'\hyperref[ax:separacio-globals]{axioma de separació global}, són iguals. Aquesta igualtat diu que el parell \(({!}_X,{!}_X)\) satisfà l'equació de compatibilitat del pullback de \(0\) i \(1\). La propietat universal proporciona una aplicació \(X\to P\), on \(P\) és un pullback d'aquestes dues aplicacions; com que \(P\) és inicial, obtenim per composició una aplicació \(X\to\mathbf0\). Per l'estrictesa del conjunt inicial (proposició~\ref{prop:inicial-estricte}), \(X\cong\mathbf0\).
\end{prova}

La segona és l'exhaustivitat de les sumes sobre els elements globals, que ara es dedueix del teorema de descomposició per canvi de base.

\begin{teorema}[Exhaustivitat de les sumes]\label{teo:exhaustivitat-sumes}
Sigui \(A_0\xrightarrow{i_0}A_0+A_1\xleftarrow{i_1}A_1\) una suma. Tot element global \(x:\mathbf1\to A_0+A_1\) es factoritza a través d'exactament una de les dues injeccions.
\end{teorema}

\begin{prova}
Apliquem el teorema~\ref{teo:descomposicio-canvi-base} a \(x\), amb estadi \(X=\mathbf1\). Obtenim una suma
\[
X_0\longrightarrow\mathbf1\longleftarrow X_1.
\]
Els dos conjunts \(X_0,X_1\) no poden ser tots dos inicials, perquè la seva suma és \(\mathbf1\), que no és inicial. Per tant, almenys un dels dos no és inicial. Per la contraposició del lema~\ref{lema:buit}, aquest conjunt té un element global; suposem que és \(s:\mathbf1\to X_0\). Com que \(j_0\circ s=\id_{\mathbf1}\), tenim
\[
j_0\circ(s\circ j_0)=(j_0\circ s)\circ j_0=j_0=j_0\circ\id_{X_0}.
\]
Com que \(j_0\) és un monomorfisme, \(s\circ j_0=\id_{X_0}\); així, \(j_0\) és un isomorfisme. La commutativitat del quadrat cartesià dona
\[
x=x\circ j_0\circ s=i_0\circ k_0\circ s,
\]
de manera que \(x\) es factoritza a través d'\(i_0\). La factorització és única perquè \(i_0\) és un monomorfisme.

Els dos conjunts no poden tenir simultàniament un element global. En efecte, dos punts \(s_0:\mathbf1\to X_0\) i \(s_1:\mathbf1\to X_1\) satisfarien \(j_0\circ s_0=\id_{\mathbf1}=j_1\circ s_1\), i la propietat universal del pullback de \(j_0\) i \(j_1\) produiria una aplicació \(\mathbf1\to\mathbf0\), en contradicció amb la proposició~\ref{prop:no-fletxa-u-zero}. Per tant, \(x\) es factoritza a través d'una sola de les dues injeccions.
\end{prova}

\begin{corollari}[Els elements de \(\mathbf2\)]\label{cor:elements-de-2}
Tot element global de \(\mathbf2\) és \(0\) o \(1\). Com que el teorema~\ref{teo:zero-diferent-u} afirma que són diferents, \(\mathbf2\) té exactament dos elements globals.
\end{corollari}

\begin{prova}
És el cas \(A_0=A_1=\mathbf1\) del teorema~\ref{teo:exhaustivitat-sumes}: factoritzar-se a través de l'injecció \(0\) vol dir ser \(0\), i a través d'\(1\), ser \(1\).
\end{prova}

\begin{observacio}[Lectura clàssica dels punts d'una suma]\label{obs:punts-suma}
En el model clàssic dels conjunts, la suma \(A_0+A_1\) és la unió disjunta, i els seus punts són els punts d'\(A_0\) juntament amb els d'\(A_1\), sense solapament. El teorema~\ref{teo:exhaustivitat-sumes} i la disjunció de l'axioma~\ref{ax:sumes-disjuntes} recuperen aquesta imatge dins de la teoria: tot punt global prové d'un sumand o de l'altre, i mai de tots dos alhora.
\end{observacio}

\section{La part buida i la part total}\label{sec:buida-total}

El conjunt inicial estricte permet construir una part buida de qualsevol
conjunt. La part total ja havia aparegut al capítol anterior, representada
per la identitat.

\begin{definicio}[Part buida i part total]\index{Part buida i part total}\label{def:parts-extremals}
Una part \(m\) de \(A\) és \emph{buida} si el seu domini és inicial, i
és \emph{total} si \(m\) és un isomorfisme. Formalment,
\[
\begin{aligned}
\operatorname{Buida}(m,A)
&\iff
\operatorname{Part}(m,A)
\land
\exists I\,
\bigl(
\operatorname{Inicial}(I)\land\Dom(m,I)
\bigr),
\\[0.5ex]
\operatorname{Total}(m,A)
&\iff
\operatorname{Part}(m,A)
\land\operatorname{Iso}(m).
\end{aligned}
\]
\end{definicio}

\begin{proposicio}\label{prop:parts-extremals-existeixen}
Per a tot conjunt \(A\), l'única aplicació
\[
u_A:\mathbf0\to A
\]
és una part buida, i la identitat
\[
\id_A:A\to A
\]
és una part total.
\end{proposicio}

\begin{prova}
La identitat és un isomorfisme i, per tant, un monomorfisme.

Vegem que \(u_A\) és un monomorfisme. Siguin \(r,s:X\to\mathbf0\)
amb \(u_A\circ r=u_A\circ s\). L'existència de \(r\) implica, per
l'estrictesa de \(\mathbf0\), que \(X\) és inicial. Per tant, només hi
ha una aplicació de \(X\) en \(\mathbf0\), i \(r=s\). Així,
\(u_A\) és un monomorfisme amb codomini \(A\).
\end{prova}

\begin{notacio}[Parts extremals]\label{not:parts-extremals}
Fixarem com a representants de les parts buida i total de \(A\)
\[
\bot_A:=u_A:\mathbf0\hookrightarrow A,
\qquad
\top_A:=\id_A:A\hookrightarrow A.
\]
\end{notacio}

\begin{proposicio}[Unicitat de les parts extremals]\label{prop:parts-extremals-uniques}
Qualsevol part buida de \(A\) és equivalent a \(\bot_A\), i qualsevol
part total de \(A\) és equivalent a \(\top_A\).
\end{proposicio}

\begin{prova}
Sigui \(m:I\hookrightarrow A\) una part buida. Els conjunts \(I\) i
\(\mathbf0\) són inicials i, per tant, existeix un únic isomorfisme
\(t:I\to\mathbf0\). Com que només hi ha una aplicació d'un conjunt
inicial cap a \(A\),
\[
m=\bot_A\circ t.
\]
Si \(t^{-1}:\mathbf0\to I\) és la inversa, també
\[
\bot_A=m\circ t^{-1}.
\]
Així, \(m\equiv_A\bot_A\).

Sigui ara \(n:T\hookrightarrow A\) una part total. Com que \(n\) és un
isomorfisme,
\[
n=\top_A\circ n,
\qquad
\top_A=n\circ n^{-1}.
\]
Per tant, \(n\equiv_A\top_A\).
\end{prova}

A la \hyperref[prop:caracteristica-total]{proposició~\ref*{prop:caracteristica-total}} hem demostrat que la característica de la part total és
\[
\chi_{\top_A}=1\circ {!}_A.
\]
La característica de la part buida és la constant complementària.

\begin{proposicio}[Característica de la part buida]\label{prop:caracteristica-buida}
La característica de \(\bot_A\) és
\[
\chi_{\bot_A}=0\circ {!}_A.
\]
\end{proposicio}

\begin{prova}
Sigui \(x:X\to A\). Comprovarem que
\[
x\in_X\bot_A
\quad\Longleftrightarrow\quad
0\circ {!}_X=1\circ {!}_X.
\]

Suposem primer que \(x\in_X\bot_A\). Existeix una aplicació
\(w:X\to\mathbf0\). Per l'estrictesa de \(\mathbf0\), el conjunt
\(X\) és inicial; per tant, les dues aplicacions de \(X\) en
\(\mathbf2\) són iguals i
\[
0\circ {!}_X=1\circ {!}_X.
\]

Recíprocament, suposem aquesta igualtat. El parell
\(({!}_X,{!}_X)\) satisfà aleshores l'equació de compatibilitat per al pullback de
\(0,1:\mathbf1\to\mathbf2\). La seva propietat universal proporciona una aplicació \(X\to P\), on \(P\) és un pullback de \(0\) i \(1\). Com que \(P\) és inicial per la disjunció de la suma, obtenim per composició una aplicació \(X\to\mathbf0\). L'estrictesa implica que \(X\) és inicial. En
conseqüència, \(x:X\to A\) coincideix amb la composició de l'única
aplicació \(X\to\mathbf0\) amb \(\bot_A\), i
\(x\in_X\bot_A\).

Per la unicitat de l'aplicació característica,
\(\chi_{\bot_A}=0\circ {!}_A\).
\end{prova}

La propietat característica es pot expressar ara mitjançant un quadrat
cartesià.

\begin{teorema}[El quadrat característic]\label{teo:quadrat-caracteristic}
Sigui \(m:S\hookrightarrow A\) una part. El quadrat
\begin{center}
\begin{tikzpicture}[scale=1.05, every node/.style={font=\small}, >=Stealth]
  \node (S) at (0,2.1) {\(S\)};
  \node (One) at (4.0,2.1) {\(\mathbf1\)};
  \node (A) at (0,0) {\(A\)};
  \node (Two) at (4.0,0) {\(\mathbf2\)};

  \draw[->, thick] (S) -- node[above] {\({!}_S\)} (One);
  \draw[{Hooks[right]}-{Stealth[scale=1.2]}, thick] (S) -- node[left] {\(m\)} (A);
  \draw[{Hooks[right]}-{Stealth[scale=1.2]}, thick] (One) -- node[right] {\(1\)} (Two);
  \draw[->, thick] (A) -- node[below] {\(\chi_m\)} (Two);
\end{tikzpicture}
\end{center}
és cartesià. Per tant, \(m\) és una imatge inversa de la part
\(1:\mathbf1\hookrightarrow\mathbf2\) per \(\chi_m\).
\end{teorema}

\begin{prova}
El quadrat commuta perquè
\[
\chi_m\circ m=1\circ {!}_S.
\]
Siguin \(x:X\to A\) i \(t:X\to\mathbf1\) tals que
\[
\chi_m\circ x=1\circ t.
\]
Per la terminalitat de \(\mathbf1\), \(t={!}_X\). La definició
d'aplicació característica dona
\[
x\in_Xm.
\]
Per tant, existeix una aplicació \(\bar{x}:X\to S\) tal que
\[
m\circ\bar{x}=x.
\]
A més, \({!}_S\circ\bar{x}=t\) per la unicitat de les aplicacions cap
a \(\mathbf1\). L'aplicació \(\bar{x}\) és única perquè \(m\) és un
monomorfisme. Això és exactament la propietat universal del pullback.
\end{prova}

\begin{observacio}[Lectura clàssica de les parts extremals]\label{obs:parts-extremals-classiques}
En la teoria clàssica de conjunts, \(\bot_A\) representa el subconjunt
buit i \(\top_A\) representa tot el conjunt \(A\). Les seves
característiques són, respectivament, les aplicacions constants amb
valor \(0\) i amb valor \(1\). Dins la teoria, aquesta interpretació es
recupera a partir de les propietats universals del conjunt inicial, de
la identitat i del classificador de parts.
\end{observacio}

\section{La part complementària}\label{sec:complementaria}

La suma \(\mathbf2=\mathbf1+\mathbf1\) permet intercanviar els dos
valors de veritat.

\begin{definicio}[Negació]\index{Negació}\label{def:negacio}
La \emph{negació} és l'única aplicació
\[
\neg=[1,0]:\mathbf2\to\mathbf2
\]
que compleix
\[
\neg\circ0=1,
\qquad
\neg\circ1=0.
\]
Formalment,
\[
\begin{aligned}
\operatorname{Neg}(t)
\iff{}&
\Dom(t,\mathbf2)\land\Cod(t,\mathbf2)
\\
&{}\land\Com(0,t,1)
\land\Com(1,t,0).
\end{aligned}
\]
\end{definicio}

\begin{proposicio}\label{prop:negacio-involutiva}
La negació és una involució:
\[
\neg\circ\neg=\id_{\mathbf2}.
\]
En particular, és un isomorfisme invers de si mateix.
\end{proposicio}

\begin{prova}
Les aplicacions \(\neg\circ\neg\) i \(\id_{\mathbf2}\) tenen les
mateixes composicions amb les dues injeccions de la suma:
\[
(\neg\circ\neg)\circ0=\neg\circ1=0,
\qquad
(\neg\circ\neg)\circ1=\neg\circ0=1.
\]
L'extensionalitat de la definició per casos dona
\(\neg\circ\neg=\id_{\mathbf2}\).
\end{prova}

\begin{proposicio}[Característica del valor \(0\)]\label{prop:caracteristica-zero}
Considerada com una part
\[
0:\mathbf1\hookrightarrow\mathbf2,
\]
l'aplicació característica de \(0\) és la negació:
\[
\chi_0=\neg.
\]
\end{proposicio}

\begin{prova}
Sigui \(x:X\to\mathbf2\). Com que l'única aplicació
\(X\to\mathbf1\) és \({!}_X\), tenim
\[
x\in_X0
\quad\Longleftrightarrow\quad
x=0\circ {!}_X.
\]
Si aquesta igualtat es compleix, aleshores
\[
\neg\circ x
=\neg\circ0\circ {!}_X
=1\circ {!}_X.
\]
Recíprocament, si \(\neg\circ x=1\circ {!}_X\), componem amb
\(\neg\) i usem la involució:
\[
x
=\neg\circ\neg\circ x
=\neg\circ1\circ {!}_X
=0\circ {!}_X.
\]
Així, \(\neg\) compleix la propietat característica de la part
\(0\), i la unicitat dona el resultat.
\end{prova}

\begin{definicio}[Part complementària]\index{Part complementària}\label{def:complementaria}
Sigui \(m\) una part de \(A\). Una part \(n\) de \(A\) és
\emph{complementària} de \(m\) si és una imatge inversa de la part
\(0:\mathbf1\hookrightarrow\mathbf2\) per l'aplicació característica
de \(m\). Formalment,
\[
\begin{aligned}
\operatorname{Compl}(n,m,A)
\iff{}&
\operatorname{Part}(m,A)
\\
&{}\land
\exists\chi\,
\bigl(
\operatorname{Car}(\chi,m,A)
\land\operatorname{ImInv}(n,\chi,0)
\bigr).
\end{aligned}
\]
\end{definicio}

\begin{notacio}[Complement d'una part]\label{not:complementaria}
Fixarem una part complementària de \(m:S\hookrightarrow A\) i
l'escriurem
\[
m^c:S^c\hookrightarrow A.
\]
Geomètricament, \(m^c\) queda determinada pel quadrat cartesià
\begin{center}
\begin{tikzpicture}[scale=1.05, every node/.style={font=\small}, >=Stealth]
  \node (Sc) at (0,2.1) {\(S^c\)};
  \node (One) at (4.0,2.1) {\(\mathbf1\)};
  \node (A) at (0,0) {\(A\)};
  \node (Two) at (4.0,0) {\(\mathbf2\)};

  \draw[->, thick] (Sc) -- node[above] {\({!}_{S^c}\)} (One);
  \draw[{Hooks[right]}-{Stealth[scale=1.2]}, thick] (Sc) -- node[left] {\(m^c\)} (A);
  \draw[{Hooks[right]}-{Stealth[scale=1.2]}, thick] (One) -- node[right] {\(0\)} (Two);
  \draw[->, thick] (A) -- node[below] {\(\chi_m\)} (Two);
\end{tikzpicture}
\end{center}
\end{notacio}

\begin{proposicio}[Característica del complement]\label{prop:caracteristica-complement}
Per a tota part \(m\) de \(A\),
\[
\chi_{m^c}=\neg\circ\chi_m.
\]
\end{proposicio}

\begin{prova}
Per la proposició sobre la característica d'una imatge inversa,
\[
\chi_{m^c}=\chi_0\circ\chi_m.
\]
La proposició~\ref{prop:caracteristica-zero} identifica
\(\chi_0\) amb \(\neg\), i obtenim la fórmula anunciada.
\end{prova}

\begin{proposicio}[Involució del complement]\label{prop:complementari-involutiu}
Per a tota part \(m\) de \(A\),
\[
{(m^c)}^c\equiv_A m.
\]
\end{proposicio}

\begin{prova}
Per la proposició anterior i la involució de la negació,
\[
\chi_{{(m^c)}^c}
=\neg\circ\chi_{m^c}
=\neg\circ\neg\circ\chi_m
=\chi_m.
\]
La caracterització de les parts equivalents mitjançant les seves
aplicacions característiques dona
\({(m^c)}^c\equiv_A m\).
\end{prova}

\begin{teorema}[Descomposició en una part i el seu complement]\label{teo:part-complementari}
Per a tota part \(m:S\hookrightarrow A\), les dues parts
\[
m:S\hookrightarrow A,
\qquad
m^c:S^c\hookrightarrow A
\]
formen una suma disjunta. En particular,
\[
[m,m^c]:S+S^c\xrightarrow{\sim}A
\]
és un isomorfisme.
\end{teorema}

\begin{prova}
El quadrat característic del
teorema~\ref{teo:quadrat-caracteristic} identifica \(m\) amb el
pullback de \(1\) al llarg de \(\chi_m\). Per definició, \(m^c\) és el
pullback de \(0\) al llarg de la mateixa aplicació.

Les aplicacions
\[
\mathbf1\xrightarrow{1}\mathbf2\xleftarrow{0}\mathbf1
\]
formen també una suma, obtinguda de la suma canònica intercanviant els
dos sumands. Per l'estabilitat, els seus dos pullbacks al llarg de
\(\chi_m:A\to\mathbf2\) formen una suma
\[
S\xrightarrow{m}A\xleftarrow{m^c}S^c.
\]
Aquesta suma és disjunta per l'axioma~\ref{ax:sumes-disjuntes}. La
unicitat de les sumes llevat d'un únic isomorfisme dona
\([m,m^c]:S+S^c\xrightarrow{\sim}A\).
\end{prova}

\begin{observacio}[Lectura clàssica del complement]\label{obs:complement-classic}
Per a un element global \(x:\mathbf1\to A\), l'element
\(\chi_m\circ x\) és \(0\) o \(1\). La fórmula
\[
\chi_{m^c}=\neg\circ\chi_m
\]
implica, per tant,
\[
x\in m^c
\quad\Longleftrightarrow\quad
x\notin m.
\]
Aquesta equivalència és una lectura sobre punts globals. Per a un
element generalitzat \(x:X\to A\), l'estadi \(X\) pot descompondre's
en una part que es factoritza per \(m\) i una part que es factoritza
per \(m^c\); no cal que tota la família pertanyi íntegrament a una de
les dues parts.
\end{observacio}

Amb la part buida i la distinció entre \(0\) i \(1\) podem tancar la qüestió que el capítol anterior va deixar oberta (observació~\ref{obs:exhaustiva-epi}): els epimorfismes són exhaustius sobre els punts globals, i no cal cap axioma d'elecció per demostrar-ho.

\begin{proposicio}[Els epimorfismes són exhaustius sobre punts]\label{prop:epi-exhaustiu-punts}
Sigui \(f\colon A\to B\) un epimorfisme. Per a tot element global \(y\colon\mathbf1\to B\) existeix un element global \(x\colon\mathbf1\to A\) tal que \(f\circ x=y\). Juntament amb l'observació~\ref{obs:exhaustiva-epi}, una aplicació és un epimorfisme si i només si és exhaustiva sobre els elements globals.\par
\emph{Axiomes usats:} separació global, representació de les parts, disjunció i estabilitat de les sumes; \emph{cap elecció}.
\end{proposicio}

\begin{prova}
Tot element global \(y\colon\mathbf1\to B\) és un monomorfisme, perquè dues aplicacions paral·leles \(X\to\mathbf1\) sempre són iguals; és, doncs, una part de \(B\), amb característica \(\chi_y\). Sigui \(q\colon P\hookrightarrow A\) la imatge inversa \(f^{-1}(y)\), amb \(q'\colon P\to\mathbf1\) l'altra projecció del quadrat cartesià, de manera que \(f\circ q=y\circ q'\).

Si \(P\) té un element global \(s\colon\mathbf1\to P\), posem \(x:=q\circ s\); aleshores \(f\circ x=y\circ q'\circ s=y\), perquè \(q'\circ s=\id_{\mathbf1}\) per la terminalitat.

Si \(P\) no té cap element global, pel lema del buit (lema~\ref{lema:buit}) \(P\) és inicial, de manera que \(f^{-1}(y)\) és una part buida de \(A\) i, per la proposició~\ref{prop:caracteristica-buida} i la unicitat de les parts extremals, \(\chi_{f^{-1}(y)}=0\circ{!}_A\). D'altra banda, \(\chi_{f^{-1}(y)}=\chi_y\circ f\) (proposició~\ref{prop:caracteristica-imatge-inversa}). Així
\[
\chi_y\circ f=0\circ{!}_A=(0\circ{!}_B)\circ f,
\]
i com que \(f\) és un epimorfisme, \(\chi_y=0\circ{!}_B\). Però \(y\in y\) amb testimoni \(\id_{\mathbf1}\), de manera que \(\chi_y\circ y=1\); i també \(\chi_y\circ y=0\circ{!}_B\circ y=0\). Obtenim \(0=1\), contra el teorema~\ref{teo:zero-diferent-u}. Aquest cas és, doncs, impossible, i el primer dona la preimatge buscada.
\end{prova}

\begin{observacio}[Dues exhaustivitats, un sol epimorfisme]\label{obs:dues-exhaustivitats}
Tenim ara el quadre complet. Ser epimorfisme equival a ser exhaustiu sobre els \emph{punts globals}, i això es demostra amb els axiomes d'aquest capítol. Ser exhaustiu sobre els elements \emph{generalitzats} equival a tenir secció (proposició~\ref{prop:exhaustiva-generalitzada-escindit}) i és, en principi, més fort: que tot epimorfisme tingui secció és exactament l'axioma d'elecció, que s'introduirà al capítol~\ref{cap:relacions}. La diferència entre les dues exhaustivitats és la diferència entre \emph{saber que hi ha} una preimatge per a cada punt i \emph{tenir una regla} que en triï una per a cada element generalitzat alhora.
\end{observacio}

\section{Operacions booleanes amb parts}\label{sec:booleanes}

La intersecció de parts ja s'ha construït al capítol~2 mitjançant un
pullback. Ara introduirem la unió i la diferència, i mostrarem que les
parts d'un conjunt, considerades llevat d'equivalència, satisfan les
lleis booleanes.

\subsection{Producte d'aplicacions i distributivitat}

\begin{notacio}[Producte d'aplicacions]\label{not:producte-aplicacions}
Siguin
\[
f:A\to B,
\qquad
g:C\to D.
\]
Denotarem per
\[
f\times g:A\times C\to B\times D
\]
l'única aplicació determinada per
\[
p_B\circ(f\times g)=f\circ p_A,
\qquad
p_D\circ(f\times g)=g\circ p_C.
\]
Equivalentment,
\[
f\times g
=
\langle f\circ p_A,\,g\circ p_C\rangle.
\]
\end{notacio}

\begin{teorema}[Distributivitat del producte respecte de la suma]\label{teo:distributivitat}
Siguin
\[
B_0\xrightarrow{i_0}B_0+B_1\xleftarrow{i_1}B_1
\]
les injeccions d'una suma. Per a tot conjunt \(A\), les
dues aplicacions
\[
\id_A\times i_0:A\times B_0\to A\times(B_0+B_1),
\]
\[
\id_A\times i_1:A\times B_1\to A\times(B_0+B_1)
\]
formen una suma. En particular,
\[
A\times(B_0+B_1)
\cong
(A\times B_0)+(A\times B_1).
\]
\end{teorema}

\begin{prova}
Sigui
\[
p:A\times(B_0+B_1)\to B_0+B_1
\]
la segona projecció. El quadrat
\begin{center}
\begin{tikzpicture}[scale=1.0, every node/.style={font=\small}, >=Stealth]
  \node (AB0) at (0,2.1) {\(A\times B_0\)};
  \node (B0) at (4.2,2.1) {\(B_0\)};
  \node (ABS) at (0,0) {\(A\times(B_0+B_1)\)};
  \node (S) at (4.2,0) {\(B_0+B_1\)};

  \draw[->, thick] (AB0) -- node[above] {\(p_{B_0}\)} (B0);
  \draw[{Hooks[right]}-{Stealth[scale=1.2]}, thick] (AB0) -- node[left] {\(\id_A\times i_0\)} (ABS);
  \draw[{Hooks[right]}-{Stealth[scale=1.2]}, thick] (B0) -- node[right] {\(i_0\)} (S);
  \draw[->, thick] (ABS) -- node[below] {\(p\)} (S);
\end{tikzpicture}
\end{center}
és cartesià. En efecte, commuta per la definició del producte
d'aplicacions. Si \(u:X\to A\times(B_0+B_1)\) i \(v:X\to B_0\)
compleixen \(p\circ u=i_0\circ v\), l'aplicació
\[
w=\langle p_A\circ u,v\rangle:X\to A\times B_0
\]
compleix
\[
(\id_A\times i_0)\circ w=u,
\qquad
p_{B_0}\circ w=v.
\]
La unicitat se segueix de la propietat universal del producte. El cas
de \(i_1\) és anàleg.

Per tant, \(A\times B_0\) i \(A\times B_1\) són els dos canvis de base
de les injeccions al llarg de \(p\). L'estabilitat de la
suma implica que conflueixen en una suma amb vèrtex
\(A\times(B_0+B_1)\).
\end{prova}

\subsection{Les quatre injeccions de \texorpdfstring{\(\mathbf2\times\mathbf2\)}{2 per 2}}

Per a \(a,b\in\{0,1\}\), escriurem
\[
\delta_{ab}:=\langle a,b\rangle:
\mathbf1\to\mathbf2\times\mathbf2.
\]
Aquesta notació pertany al metallenguatge: els símbols \(a,b\) recorren
els dos elements globals canònics ja distingits.

\begin{proposicio}[Descomposició de \(\mathbf2\times\mathbf2\)]\label{prop:quatre-injeccions}
Les quatre aplicacions
\[
\delta_{00},\quad
\delta_{01},\quad
\delta_{10},\quad
\delta_{11}
\]
són les injeccions d'una suma iterada de quatre còpies de
\(\mathbf1\) amb vèrtex \(\mathbf2\times\mathbf2\).
\end{proposicio}

\begin{prova}
Apliquem el teorema~\ref{teo:distributivitat} a
\(\mathbf2=\mathbf1+\mathbf1\) en el segon factor. Obtenim una suma
\[
\mathbf2\times\mathbf1
\xrightarrow{\id_{\mathbf2}\times0}
\mathbf2\times\mathbf2
\xleftarrow{\id_{\mathbf2}\times1}
\mathbf2\times\mathbf1.
\]
La primera projecció
\[
p_{\mathbf2}:\mathbf2\times\mathbf1\to\mathbf2
\]
és un isomorfisme, amb inversa
\(\langle\id_{\mathbf2},{!}_{\mathbf2}\rangle\). Així, cadascun dels
dos sumands és una còpia de
\(\mathbf2=\mathbf1+\mathbf1\).

Seguim explícitament les quatre injeccions. Per a
\(a,b\in\{0,1\}\), la injecció corresponent és la composició
\[
\mathbf1
\xrightarrow{\langle a,\id_{\mathbf1}\rangle}
\mathbf2\times\mathbf1
\xrightarrow{\id_{\mathbf2}\times b}
\mathbf2\times\mathbf2.
\]
Les seves components són
\[
\begin{aligned}
p_1\circ
(\id_{\mathbf2}\times b)
\circ\langle a,\id_{\mathbf1}\rangle&=a,\\
p_2\circ
(\id_{\mathbf2}\times b)
\circ\langle a,\id_{\mathbf1}\rangle&=b.
\end{aligned}
\]
Per la unicitat de la propietat universal del producte, aquesta
composició és exactament \(\langle a,b\rangle=\delta_{ab}\). Això identifica les quatre còpies de \(\mathbf1\), en l'ordre produït
per aquesta descomposició, amb
\[
\delta_{00},\quad\delta_{10},\quad\delta_{01},\quad\delta_{11}.
\]
Per a les taules de veritat les enumerarem en l'ordre convencional
\(\delta_{00},\delta_{01},\delta_{10},\delta_{11}\); aquesta
reordenació no altera la propietat universal de la suma iterada.
\end{prova}

\begin{proposicio}[Disjunció de les quatre injeccions]\label{prop:delta-disjuntes}
Per a \(a,b,c,d\in\{0,1\}\), si \((a,b)\neq(c,d)\), el pullback de
\[
\delta_{ab},\delta_{cd}:\mathbf1\to\mathbf2\times\mathbf2
\]
és inicial.
\end{proposicio}

\begin{prova}
Sigui \(P\) un pullback de \(\delta_{ab}\) i \(\delta_{cd}\). Les dues aplicacions de \(P\) a \(\mathbf1\) són necessàriament \({!}_P\). En compondre la igualtat del pullback amb les dues projeccions de \(\mathbf2\times\mathbf2\), obtenim
\[
a\circ {!}_P=c\circ {!}_P,
\qquad
b\circ {!}_P=d\circ {!}_P.
\]
Com que els parells són diferents, almenys una d'aquestes igualtats és
\[
0\circ {!}_P=1\circ {!}_P.
\]
La propietat universal d'un pullback inicial \(R\) de \(0\) i \(1\) proporciona una aplicació \(P\to R\); per composició amb l'única aplicació \(R\to\mathbf0\), obtenim una aplicació \(P\to\mathbf0\). Per l'estrictesa de \(\mathbf0\), el conjunt \(P\) és inicial.
\end{prova}

\begin{proposicio}[Extensionalitat de les taules de veritat]\label{prop:extensionalitat-taules}
Siguin \(F,G:\mathbf2\times\mathbf2\to B\). Aleshores
\[
F=G
\quad\Longleftrightarrow\quad
F\circ\delta_{ab}=G\circ\delta_{ab}
\quad
\text{per a tots }a,b\in\{0,1\}.
\]

Anàlogament, dues aplicacions
\[
F,G:(\mathbf2\times\mathbf2)\times\mathbf2\to B
\]
són iguals si coincideixen sobre les vuit aplicacions
\[
\delta_{abc}:=\langle\delta_{ab},c\rangle,
\qquad a,b,c\in\{0,1\}.
\]
\end{proposicio}

\begin{prova}
La primera afirmació és l'extensionalitat de la definició per casos
aplicada a la suma de quatre còpies de \(\mathbf1\) de la proposició
anterior. Aplicant novament la distributivitat, ara a
\((\mathbf2\times\mathbf2)\times\mathbf2\), obtenim una suma de vuit
còpies de \(\mathbf1\); el mateix argument dona la segona afirmació.
\end{prova}

\subsection{Conjunció, disjunció i parts}

Les aplicacions \(\delta_{11}\) i \(\delta_{00}\) són
monomorfismes, perquè són injeccions d'una suma disjunta.
Per tant, representen parts de \(\mathbf2\times\mathbf2\).

\begin{definicio}[Conjunció i disjunció]\index{Conjunció i disjunció}\label{def:conjuncio-disjuncio}
La \emph{conjunció} és l'aplicació característica de la part
\(\delta_{11}\):
\[
\wedge:=\chi_{\delta_{11}}:
\mathbf2\times\mathbf2\to\mathbf2.
\]
La \emph{disjunció} és l'aplicació característica del complement de la
part \(\delta_{00}\):
\[
\vee:=\chi_{\delta_{00}^{c}}:
\mathbf2\times\mathbf2\to\mathbf2.
\]
\end{definicio}

\begin{proposicio}[Taules de veritat]\label{prop:taules-veritat}
Per a \(a,b\in\{0,1\}\),
\[
\wedge\circ\delta_{ab}
=
\begin{cases}
1,&(a,b)=(1,1),\\
0,&(a,b)\neq(1,1),
\end{cases}
\]
i
\[
\vee\circ\delta_{ab}
=
\begin{cases}
0,&(a,b)=(0,0),\\
1,&(a,b)\neq(0,0).
\end{cases}
\]
\end{proposicio}

\begin{prova}
Com que \(\wedge\) és la característica de \(\delta_{11}\), tenim
\[
\wedge\circ\delta_{11}=1.
\]
Si \((a,b)\neq(1,1)\), el pullback de \(\delta_{ab}\) i
\(\delta_{11}\) és inicial per la proposició~\ref{prop:delta-disjuntes}. Per tant, la imatge inversa de
\(\delta_{11}\) per \(\delta_{ab}\) és la part buida de
\(\mathbf1\). La fórmula de la característica d'una imatge inversa
dona
\[
\wedge\circ\delta_{ab}
=\chi_{\bot_{\mathbf1}}
=0.
\]
Això prova la taula de la conjunció.

El mateix argument, aplicat ara a la part \(\delta_{00}\), mostra
que la seva característica val \(1\) sobre \(\delta_{00}\) i
\(0\) sobre les altres tres injeccions. Com que
\(\delta_{00}^{c}\) és la part complementària,
\[
\vee=\chi_{\delta_{00}^{c}}
=\neg\circ\chi_{\delta_{00}}.
\]
Per tant, \(\vee\) val \(0\) sobre \(\delta_{00}\) i
\(1\) sobre les altres tres aplicacions. Aquesta és la taula de la
disjunció.
\end{prova}

\begin{proposicio}[Característica de la intersecció]\label{prop:caracteristica-interseccio}
Siguin \(m\) i \(n\) parts de \(A\). Aleshores
\[
\chi_{m\cap n}
=
\wedge\circ\langle\chi_m,\chi_n\rangle.
\]
\end{proposicio}

\begin{prova}
Posem
\[
c:=\langle\chi_m,\chi_n\rangle:
A\to\mathbf2\times\mathbf2.
\]
Com que \(\wedge\) és la característica de \(\delta_{11}\), la
composició \(\wedge\circ c\) és la característica de la imatge inversa
\(c^{-1}(\delta_{11})\).

Per a qualsevol \(x:X\to A\),
\[
\begin{aligned}
x\in_Xc^{-1}(\delta_{11})
&\Longleftrightarrow
c\circ x\text{ es factoritza a través de }\delta_{11}
\\
&\Longleftrightarrow
\chi_m\circ x=1\circ {!}_X
\ \land\
\chi_n\circ x=1\circ {!}_X
\\
&\Longleftrightarrow
x\in_Xm\ \land\ x\in_Xn
\\
&\Longleftrightarrow
x\in_X(m\cap n).
\end{aligned}
\]
La caracterització de la inclusió mitjançant els elements generalitzats
implica que \(c^{-1}(\delta_{11})\) i \(m\cap n\) són parts
equivalents. Per tant, tenen la mateixa aplicació característica, i
obtenim la fórmula anunciada.
\end{prova}

\begin{definicio}[Unió i diferència de parts]\index{Unió i diferència de parts}\label{def:operacions-parts}
Siguin \(m\) i \(n\) parts de \(A\). Una part \(u\) de \(A\) és una
\emph{unió} de \(m\) i \(n\) si la seva característica és
\[
\chi_u=\vee\circ\langle\chi_m,\chi_n\rangle.
\]
Formalment,
\[
\begin{aligned}
\operatorname{Unio}(u,m,n,A)
\iff{}&
\operatorname{Part}(u,A)
\land\operatorname{Part}(m,A)
\land\operatorname{Part}(n,A)
\\
&{}\land
\exists\alpha\,\exists\beta\,\exists\gamma\,
\Bigl(
\operatorname{Car}(\alpha,m,A)
\land\operatorname{Car}(\beta,n,A)
\\
&\qquad{}\land\operatorname{Car}(\gamma,u,A)
\land
\gamma=\vee\circ\langle\alpha,\beta\rangle
\Bigr).
\end{aligned}
\]

Una part \(d\) de \(A\) és una \emph{diferència} de \(m\) i \(n\) si
és una intersecció de \(m\) amb una part complementària de \(n\).
Formalment,
\[
\begin{aligned}
\operatorname{Dif}(d,m,n,A)
\iff
\exists c\,
\bigl(
\operatorname{Compl}(c,n,A)
\land\operatorname{Inter}(d,m,c,A)
\bigr).
\end{aligned}
\]
\end{definicio}

\begin{notacio}[Unió i diferència]\label{not:operacions-parts}
Fixarem representants i escriurem
\[
m\cup n,
\qquad
m\setminus n:=m\cap n^c.
\]
Aquestes parts són úniques llevat d'equivalència.
\end{notacio}

\begin{observacio}[Lectura clàssica de les operacions]\label{obs:lectura-classica-operacions}
Per a un element global \(x\in A\), les taules de veritat donen
\[
\begin{aligned}
x\in m\cap n
&\Longleftrightarrow x\in m\text{ i }x\in n,\\
x\in m\cup n
&\Longleftrightarrow x\in m\text{ o }x\in n,\\
x\in m\setminus n
&\Longleftrightarrow x\in m\text{ i }x\notin n.
\end{aligned}
\]
La primera equivalència ja és vàlida per a tots els elements
generalitzats, perquè la intersecció s'ha construït mitjançant un
pullback. Les altres dues són, en aquesta forma disjuntiva, lectures
sobre punts globals.
\end{observacio}

\subsection{Lleis booleanes}

\begin{teorema}[Lleis booleanes de les parts]\label{teo:lleis-booleanes}
Siguin \(m,n,p\) parts d'un mateix conjunt \(A\). Llevat
d'equivalència de parts, es compleixen les identitats següents.

\emph{Commutativitat:}
\[
m\cap n\equiv_A n\cap m,
\qquad
m\cup n\equiv_A n\cup m.
\]

\emph{Associativitat:}
\[
(m\cap n)\cap p\equiv_A m\cap(n\cap p),
\]
\[
(m\cup n)\cup p\equiv_A m\cup(n\cup p).
\]

\emph{Idempotència:}
\[
m\cap m\equiv_A m,
\qquad
m\cup m\equiv_A m.
\]

\emph{Elements neutres i absorbents:}
\[
m\cap\top_A\equiv_A m,
\qquad
m\cup\bot_A\equiv_A m,
\]
\[
m\cap\bot_A\equiv_A\bot_A,
\qquad
m\cup\top_A\equiv_A\top_A.
\]

\emph{Complementació:}
\[
m\cap m^c\equiv_A\bot_A,
\qquad
m\cup m^c\equiv_A\top_A,
\qquad
{(m^c)}^c\equiv_A m.
\]

\emph{Distributivitat:}
\[
m\cap(n\cup p)
\equiv_A
(m\cap n)\cup(m\cap p),
\]
\[
m\cup(n\cap p)
\equiv_A
(m\cup n)\cap(m\cup p).
\]

\emph{Absorció:}
\[
m\cap(m\cup n)\equiv_A m,
\qquad
m\cup(m\cap n)\equiv_A m.
\]

\emph{Lleis de De Morgan:}
\[
(m\cap n)^c\equiv_A m^c\cup n^c,
\]
\[
(m\cup n)^c\equiv_A m^c\cap n^c.
\]
\end{teorema}

\begin{prova}
Les fórmules de les característiques obtingudes fins ara són
\[
\chi_{m^c}=\neg\circ\chi_m,
\]
\[
\chi_{m\cap n}
=\wedge\circ\langle\chi_m,\chi_n\rangle,
\]
\[
\chi_{m\cup n}
=\vee\circ\langle\chi_m,\chi_n\rangle,
\]
i
\[
\chi_{\bot_A}=0\circ {!}_A,
\qquad
\chi_{\top_A}=1\circ {!}_A.
\]
Per tant, cada identitat de parts es redueix a una identitat entre
aplicacions construïdes amb
\(0,1,\neg,\wedge\) i \(\vee\).

Abreujarem \(\wedge\circ\langle a,b\rangle\) per \(a\wedge b\),
\(\vee\circ\langle a,b\rangle\) per \(a\vee b\), i les
composicions amb \(\neg\) de la manera habitual.
Per la proposició~\ref{prop:extensionalitat-taules}, les identitats de
dues variables es comproven sobre les quatre aplicacions
\(\delta_{ab}\), i les de tres variables sobre les vuit aplicacions
\(\delta_{abc}\). Les taules de veritat donen, per a
\(a,b,c\in\{0,1\}\),
\[
\begin{gathered}
a\wedge b=b\wedge a,
\qquad
a\vee b=b\vee a,\\
(a\wedge b)\wedge c=a\wedge(b\wedge c),
\qquad
(a\vee b)\vee c=a\vee(b\vee c),\\
a\wedge a=a,
\qquad
a\vee a=a,\\
a\wedge1=a,
\quad
a\vee0=a,
\quad
a\wedge0=0,
\quad
a\vee1=1,\\
a\wedge\neg a=0,
\qquad
a\vee\neg a=1,\\
a\wedge(b\vee c)
=(a\wedge b)\vee(a\wedge c),\\
a\vee(b\wedge c)
=(a\vee b)\wedge(a\vee c),\\
a\wedge(a\vee b)=a,
\qquad
a\vee(a\wedge b)=a,\\
\neg(a\wedge b)=\neg a\vee\neg b,
\qquad
\neg(a\vee b)=\neg a\wedge\neg b.
\end{gathered}
\]

Per exemple, la primera llei distributiva es verifica sense ometre cap
cas: si \(a=0\), tots dos membres valen \(0\); si \(a=1\), tots dos
membres valen \(b\vee c\). La primera llei de De Morgan també es
comprova directament: el membre esquerre val \(0\) únicament quan
\(a=b=1\), i el membre dret té exactament el mateix comportament. La
resta d'igualtats es llegeixen immediatament de les mateixes taules.
L'extensionalitat converteix aquestes verificacions finites en
igualtats d'aplicacions.

Finalment, dues parts són equivalents si i només si tenen la mateixa
aplicació característica. Substituint \(a,b,c\) per
\(\chi_m,\chi_n,\chi_p\) en les identitats anteriors obtenim totes les
equivalències enunciades. La involució del complement ja havia estat
demostrada a la proposició~\ref{prop:complementari-involutiu}.
\end{prova}

\begin{proposicio}[Inclusió i operacions booleanes]\label{prop:inclusio-booleana}
Per a dues parts \(m,n\) de \(A\),
\[
m\subseteq_A n
\quad\Longleftrightarrow\quad
m\cap n\equiv_A m
\quad\Longleftrightarrow\quad
m\cup n\equiv_A n.
\]
\end{proposicio}

\begin{prova}
Per la caracterització de la inclusió mitjançant els elements
generalitzats,
\(m\subseteq_A n\) és equivalent al fet que tota aplicació que es
factoritza per \(m\) també es factoritzi per \(n\). Juntament amb la
caracterització de la intersecció, això equival a
\[
x\in_X(m\cap n)
\quad\Longleftrightarrow\quad
x\in_Xm
\]
per a tot \(X\) i tot \(x:X\to A\), és a dir,
\(m\cap n\equiv_A m\).

Si \(m\cap n\equiv_A m\), la compatibilitat de la unió amb
l'equivalència i la llei d'absorció donen
\[
m\cup n
\equiv_A
(m\cap n)\cup n
\equiv_A n.
\]
Recíprocament, si \(m\cup n\equiv_A n\), aleshores
\[
m\cap n
\equiv_A
m\cap(m\cup n)
\equiv_A m
\]
per absorció. Això completa les equivalències.
\end{prova}

\begin{observacio}[Àlgebra booleana de les parts]\label{obs:algebra-booleana-parts}
Fixat un conjunt \(A\), les parts de \(A\) no formen encara un conjunt
intern de la teoria. En el metallenguatge, però, podem considerar les
seves classes d'equivalència. Les operacions
\[
\cap,\qquad\cup,\qquad{}^c,
\]
les classes de \(\bot_A\) i \(\top_A\), i l'ordre induït per
\(\subseteq_A\) converteixen aquestes classes en una àlgebra booleana.
Aquesta afirmació resumeix les lleis del
teorema~\ref{teo:lleis-booleanes}; no introdueix encara un conjunt
intern de totes les parts de \(A\).
\end{observacio}

\begin{observacio}[Quadre de dependències de les seccions~\ref{sec:sumes-disjuntes} a~\ref{sec:booleanes}]\label{obs:quadre-dependencies-cap3}
El quadre següent indica, per a cada resultat, quins dels axiomes següents intervenen a la seva demostració (directament o a través dels resultats que cita): \textbf{R}, representació de les parts (axioma~\ref{ax:representacio}, del qual depenen també els equalitzadors i els pullbacks); \textbf{D}, disjunció de les sumes, i \textbf{E}, estabilitat per canvi de base (les dues meitats de l'axioma~\ref{ax:sumes-disjuntes}); \textbf{N}, no-degeneració (axioma~\ref{ax:no-degeneracio}); \textbf{S}, separació global (axioma~\ref{ax:separacio-globals}). Els axiomes estructurals, el terminal, l'inicial, les sumes i els productes s'usen a tot arreu i no es consignen. Dues lectures del quadre: la columna \textbf{S} mostra que tota la lògica booleana de les parts es demostra \emph{sense} separació global, que només cal per al lema del buit i les seves conseqüències; i la columna \textbf{N} mostra que la no-degeneració intervé només per distingir \(0\) d'\(1\).
\begin{center}\small
\begin{tabular}{@{}>{\raggedright\arraybackslash}p{0.70\linewidth}ccccc@{}}
\toprule
Resultat & R & D & E & N & S\\
\midrule
Estrictesa de \(\mathbf0\) (prop.~\ref{prop:inicial-estricte}) & \(\bullet\) & \(\bullet\) & \(\bullet\) & & \\
\(0\ne1\) (teor.~\ref{teo:zero-diferent-u}) & \(\bullet\) & \(\bullet\) & & \(\bullet\) & \\
Descomposició per canvi de base (teor.~\ref{teo:descomposicio-canvi-base}) & \(\bullet\) & \(\bullet\) & \(\bullet\) & & \\
Lema del buit (lema~\ref{lema:buit}) & \(\bullet\) & \(\bullet\) & \(\bullet\) & & \(\bullet\) \\
Exhaustivitat de les sumes; elements de \(\mathbf2\) (teor.~\ref{teo:exhaustivitat-sumes}, cor.~\ref{cor:elements-de-2}) & \(\bullet\) & \(\bullet\) & \(\bullet\) & & \(\bullet\) \\
Unicitat de les parts extremals (prop.~\ref{prop:parts-extremals-uniques}) & \(\bullet\) & \(\bullet\) & \(\bullet\) & & \\
Característica de la part buida (prop.~\ref{prop:caracteristica-buida}) & \(\bullet\) & \(\bullet\) & \(\bullet\) & & \\
Quadrat característic (teor.~\ref{teo:quadrat-caracteristic}) & \(\bullet\) & & & & \\
Negació involutiva; \(\chi_0=\neg\); característica del complement (props.~\ref{prop:negacio-involutiva}, \ref{prop:caracteristica-zero}, \ref{prop:caracteristica-complement}) & \(\bullet\) & & & & \\
Descomposició \(A\cong S+S^c\) (teor.~\ref{teo:part-complementari}) & \(\bullet\) & & \(\bullet\) & & \\
Epimorfismes exhaustius sobre punts (prop.~\ref{prop:epi-exhaustiu-punts}) & \(\bullet\) & \(\bullet\) & \(\bullet\) & \(\bullet\) & \(\bullet\) \\
Distributivitat; quatre injeccions (teor.~\ref{teo:distributivitat}, prop.~\ref{prop:quatre-injeccions}) & \(\bullet\) & & \(\bullet\) & & \\
Disjunció de les \(\delta_{ab}\); taules de veritat (props.~\ref{prop:delta-disjuntes}, \ref{prop:taules-veritat}) & \(\bullet\) & \(\bullet\) & \(\bullet\) & & \\
Extensionalitat de les taules (prop.~\ref{prop:extensionalitat-taules}) & \(\bullet\) & & \(\bullet\) & & \\
Característica de la intersecció (prop.~\ref{prop:caracteristica-interseccio}) & \(\bullet\) & & & & \\
Lleis booleanes; inclusió i operacions (teor.~\ref{teo:lleis-booleanes}, prop.~\ref{prop:inclusio-booleana}) & \(\bullet\) & \(\bullet\) & \(\bullet\) & & \\
\bottomrule
\end{tabular}
\end{center}
Cap fila no cita un resultat posterior: l'ordre del quadre és l'ordre de les demostracions, de manera que no hi ha dependències circulars. En particular, \(0\ne1\) es demostra abans que qualsevol operació booleana, i cap llei booleana no l'usa.
\end{observacio}

\section{Diagrames finits i límits finits}\label{sec:limits-finits}

Els productes, els equalitzadors i els pullbacks responen a un mateix patró: un conjunt que resol universalment un sistema finit de compatibilitats entre aplicacions. Aquesta secció unifica aquest patró sota el nom de \emph{límit} i demostra que tots els límits finits existeixen sense afegir cap axioma.

\begin{observacio}[Els diagrames finits pertanyen al metallenguatge]\label{obs:diagrama-metateoric}
Un \emph{diagrama finit} consta d'una llista finita de conjunts
\[
A_1,\ldots,A_n
\]
i d'una llista finita d'aplicacions
\[
f_\alpha:A_{s(\alpha)}\to A_{t(\alpha)}.
\]
La finitud es compta en el metallenguatge. No hem codificat els diagrames com a conjunts interns de la teoria. Per tant, l'afirmació «tot diagrama finit té límit» és un metateorema: per a cada forma finita concreta proporciona una fórmula i una demostració dins de la teoria.
\end{observacio}

\begin{notacio}[Productes finits]\label{not:productes-finits}
Per a una llista finita de conjunts \(A_1,\ldots,A_n\), escriurem
\[
\prod_{i=1}^{n}A_i
\]
per un producte iterat escollit. El producte buit és \(\mathbf1\), el producte d'un sol conjunt és aquest mateix conjunt i els productes de longitud més gran es construeixen iterant els productes binaris.
\end{notacio}

\begin{proposicio}[Existència dels productes finits]\label{prop:productes-finits}
Tota llista finita de conjunts té un producte, únic llevat d'un únic isomorfisme compatible amb les projeccions.
\end{proposicio}

\begin{prova}
Per a la llista buida, la propietat universal és la del conjunt terminal \(\mathbf1\). Per a una llista d'un sol conjunt, la identitat proporciona el producte. Si ja existeix el producte de \(A_1,\ldots,A_n\), el producte binari
\[
\left(\prod_{i=1}^{n}A_i\right)\times A_{n+1}
\]
té la propietat universal del producte de la llista ampliada. La unicitat se segueix, com en el cas binari, aplicant les propietats universals en les dues direccions.
\end{prova}

Fixem ara un diagrama finit \(D\) amb conjunts \((A_i)_i\) i aplicacions \(f_\alpha:A_{s(\alpha)}\to A_{t(\alpha)}\).

\begin{definicio}[Con]\index{Con}\label{def:con}
Un \emph{con} sobre \(D\), amb vèrtex \(X\), és una família finita d'aplicacions
\[
c_i:X\to A_i
\]
tal que, per a cada aplicació \(f_\alpha\) del diagrama,
\[
f_\alpha\circ c_{s(\alpha)}=c_{t(\alpha)}.
\]

Per al diagrama fix \(D\), aquesta definició abrevia la fórmula finita
\[
\operatorname{Con}_D\bigl(X,(c_i)_i\bigr)
\iff
\operatorname{Conj}(X)
\land
\bigwedge_i\bigl(c_i:X\to A_i\bigr)
\land
\bigwedge_\alpha
\Com\bigl(c_{s(\alpha)},f_\alpha,c_{t(\alpha)}\bigr).
\]
Els símbols \(\bigwedge_i\) i \(\bigwedge_\alpha\) representen conjuncions finites del metallenguatge, que es poden desplegar en cada diagrama concret.
\end{definicio}

\begin{definicio}[Límit]\index{Límit}\label{def:limit}
Un con
\[
\bigl(L,(\ell_i)_i\bigr)
\]
és un \emph{límit} de \(D\) si, per a qualsevol altre con
\[
\bigl(X,(c_i)_i\bigr),
\]
existeix una única aplicació \(u:X\to L\) tal que
\[
\ell_i\circ u=c_i
\qquad\text{per a tot }i.
\]

Per a cada diagrama finit fix, aquesta condició és una abreviació d'una fórmula finita construïda amb els predicats ja definits.
\end{definicio}

\begin{notacio}[Con i límit]\label{not:limit}
La família \((\ell_i)_i\) s'anomena el \emph{con límit}. La seva propietat universal es representa afegint, des de qualsevol altre con, una única fletxa discontínua \(u:X\to L\) que fa commutar tots els triangles
\[
\ell_i\circ u=c_i.
\]
\end{notacio}

\begin{observacio}[Casos particulars]\label{obs:limits-casos}
El límit del diagrama buit és el conjunt terminal \(\mathbf1\). El límit de dos conjunts sense fletxes entre ells és el seu producte. El límit d'un parell d'aplicacions paral·leles és el seu equalitzador, i el límit d'un diagrama
\[
A\longrightarrow C\longleftarrow B
\]
és el pullback corresponent.
\end{observacio}

\begin{teorema}[Existència dels límits finits]\label{teo:limits-finits}
Tot diagrama finit té límit.
\end{teorema}

\begin{prova}
Sigui \(D\) un diagrama finit amb conjunts \((A_i)_i\) i aplicacions
\[
f_\alpha:A_{s(\alpha)}\to A_{t(\alpha)}.
\]
Formem els productes finits
\[
P=\prod_i A_i,
\qquad
Q=\prod_\alpha A_{t(\alpha)},
\]
amb projeccions \(p_i:P\to A_i\) i \(q_\alpha:Q\to A_{t(\alpha)}\). La propietat universal de \(Q\) determina dues aplicacions \(r,r':P\to Q\) mitjançant
\[
q_\alpha\circ r=p_{t(\alpha)},
\qquad
q_\alpha\circ r'=f_\alpha\circ p_{s(\alpha)}.
\]
Sigui
\[
\ell:L\hookrightarrow P
\]
l'equalitzador de \(r\) i \(r'\), i posem
\[
\ell_i:=p_i\circ\ell:L\to A_i.
\]

Com que \(r\circ\ell=r'\circ\ell\), en compondre amb cada projecció \(q_\alpha\) obtenim
\[
\ell_{t(\alpha)}=f_\alpha\circ\ell_{s(\alpha)}.
\]
Per tant, \((L,(\ell_i)_i)\) és un con.

Sigui ara \((X,(c_i)_i)\) un con qualsevol. La propietat universal de \(P\) proporciona una única aplicació \(c:X\to P\) tal que
\[
p_i\circ c=c_i
\qquad\text{per a tot }i.
\]
Per a cada \(\alpha\),
\[
\begin{aligned}
q_\alpha\circ r\circ c
&=p_{t(\alpha)}\circ c
=c_{t(\alpha)}\\
&=f_\alpha\circ c_{s(\alpha)}
=f_\alpha\circ p_{s(\alpha)}\circ c
=q_\alpha\circ r'\circ c.
\end{aligned}
\]
Com que les projeccions de \(Q\) separen aplicacions, \(r\circ c=r'\circ c\). La propietat universal de l'equalitzador determina una única \(u:X\to L\) amb
\[
\ell\circ u=c.
\]
Aleshores
\[
\ell_i\circ u=p_i\circ\ell\circ u=p_i\circ c=c_i.
\]
Si \(u'\) té les mateixes components, les projeccions de \(P\) donen \(\ell\circ u'=c\), i la unicitat de l'equalitzador implica \(u'=u\). Així, \((L,(\ell_i)_i)\) és un límit.
\end{prova}

\begin{observacio}[Balanç dels límits finits]\label{obs:tots-limits-finits}
La construcció només ha fet servir el conjunt terminal, els productes binaris i els equalitzadors. Per tant, els axiomes anteriors ja garanteixen tots els límits finits; no cal afegir cap axioma d'existència nou.
\end{observacio}

\section{Coequalitzadors}\label{sec:coequalitzadors}

Els equalitzadors resolen el problema d'igualar dues aplicacions paral·leles restringint-ne el domini. La construcció dual les iguala després de compondre-les amb una aplicació comuna.

\begin{definicio}[Coequalitzador]\index{Coequalitzador}\label{def:coequalitzador}
Una aplicació \(q\) és un \emph{coequalitzador} de dues aplicacions paral·leles \(f,g\) si
\[
q\circ f=q\circ g
\]
i si tota aplicació que també les coequalitza es factoritza de manera única a través de \(q\). Formalment,
\[
\begin{aligned}
\operatorname{Coeq}(q,f,g)
\iff
\exists A\,\exists B\,\exists Q\,\Bigl(&
 f:A\to B\land g:A\to B\land q:B\to Q
\\
&{}\land
\exists p\,
 \bigl(\Com(f,q,p)\land\Com(g,q,p)\bigr)
\\
&{}\land
\forall C\,\forall h\,\Bigl(
 h:B\to C
\\
&\qquad{}\land
 \exists p'\,
 \bigl(\Com(f,h,p')\land\Com(g,h,p')\bigr)
\\
&\qquad{}\Rightarrow
 \exists!\bar h\,
 \bigl(\bar h:Q\to C\land\Com(q,\bar h,h)\bigr)
\Bigr)
\Bigr).
\end{aligned}
\]
\end{definicio}

\begin{notacio}[Propietat universal del coequalitzador]\label{not:coequalitzador}
Per a \(f,g:A\to B\), escriurem un coequalitzador com
\[
A
\mathrel{\substack{\xrightarrow{\ f\ }\\[-0.6ex]\xrightarrow[\ g\ ]{}}}
B\xrightarrow{q}Q.
\]
La propietat universal afirma que, si \(h:B\to C\) satisfà
\[
h\circ f=h\circ g,
\]
existeix una única \(\bar h:Q\to C\) tal que
\[
\bar h\circ q=h.
\]
Es representa mitjançant el diagrama
\begin{center}
\begin{tikzpicture}[scale=1.1,every node/.style={font=\small},>=Stealth]
  \node (A) at (0,1.8) {\(A\)};
  \node (B) at (2.8,1.8) {\(B\)};
  \node (Q) at (5.6,1.8) {\(Q\)};
  \node (C) at (5.6,0) {\(C\)};
  \draw[->,thick] (A) to[bend left=22] node[above] {\(f\)} (B);
  \draw[->,thick] (A) to[bend right=22] node[below] {\(g\)} (B);
  \draw[->,thick] (B) -- node[above] {\(q\)} (Q);
  \draw[->,thick] (B) -- node[below left] {\(h\)} (C);
  \draw[->,thick,dashed] (Q) -- node[right] {\(\bar h\)} (C);
\end{tikzpicture}
\end{center}
La fletxa discontínua expressa simultàniament l'existència i la unicitat de la factorització.
\end{notacio}

\begin{proposicio}[Unicitat del coequalitzador]\label{prop:coequalitzador-unic}
Si \(q:B\to Q\) i \(q':B\to Q'\) són coequalitzadors de \(f,g:A\to B\), existeix un únic isomorfisme \(w:Q\to Q'\) tal que
\[
w\circ q=q'.
\]
\end{proposicio}

\begin{prova}
La propietat universal de \(q\), aplicada a \(q'\), proporciona una única \(w:Q\to Q'\) amb \(w\circ q=q'\). Simètricament, obtenim \(w':Q'\to Q\) amb \(w'\circ q'=q\). Aleshores
\[
(w'\circ w)\circ q=q=\id_Q\circ q.
\]
La unicitat de la factorització a través de \(q\) implica \(w'\circ w=\id_Q\). De manera anàloga, \(w\circ w'=\id_{Q'}\).
\end{prova}

\begin{proposicio}\label{prop:coequalitzador-epi}
Tot coequalitzador és un epimorfisme.
\end{proposicio}

\begin{prova}
Sigui \(q:B\to Q\) un coequalitzador de \(f,g\), i siguin \(y,y':Q\to C\) amb
\[
y\circ q=y'\circ q.
\]
L'aplicació comuna \(h:=y\circ q\) coequalitza \(f\) i \(g\). Tant \(y\) com \(y'\) són factoritzacions de \(h\) a través de \(q\); per la unicitat de la propietat universal, \(y=y'\).
\end{prova}

\begin{observacio}[Per què cal un axioma nou]\label{obs:necessitat-coequalitzadors}
Els equalitzadors s'han construït al capítol~\ref{cap:parts} a partir dels productes i de la representació de les parts. Les construccions disponibles fins ara no proporcionen una construcció dual dels coequalitzadors. Per garantir-ne l'existència, cal incorporar un axioma nou. Aquesta observació no afirma una independència lògica; només assenyala que els axiomes i les construccions establerts no n'han donat encara cap construcció.
\end{observacio}

\begin{axioma}[Existència de coequalitzadors]\label{ax:coequalitzadors}
Tot parell d'aplicacions paral·leles té un coequalitzador. Formalment,
\[
\forall A\,\forall B\,\forall f\,\forall g\,
\Bigl(
 f:A\to B\land g:A\to B
 \Rightarrow
 \exists q\,\operatorname{Coeq}(q,f,g)
\Bigr).
\]
\end{axioma}

\section{Pushouts}\label{sec:pushouts}

Els pushouts, o sumes amalgamades, són la construcció dual dels pullbacks. Es poden construir combinant una suma amb un coequalitzador.

\begin{definicio}[Pushout o suma amalgamada]\index{Pushout o suma amalgamada}\label{def:pushout}
Donades dues aplicacions amb domini comú
\[
f:C\to A,
\qquad
g:C\to B,
\]
un \emph{pushout} o una \emph{suma amalgamada} és un parell d'aplicacions
\[
j_A:A\to P,
\qquad
j_B:B\to P
\]
tal que
\[
j_A\circ f=j_B\circ g
\]
i que és universal amb aquesta propietat. Formalment,
\[
\begin{aligned}
&\operatorname{Pushout}(j_A,j_B,f,g)\\
&\iff
\exists A\,\exists B\,\exists C\,\exists P\,\Bigl(&
 f:C\to A\land g:C\to B
\\
&{}\land j_A:A\to P\land j_B:B\to P
\\
&{}\land
\exists d\,
 \bigl(\Com(f,j_A,d)\land\Com(g,j_B,d)\bigr)
\\
&{}\land
\forall X\,\forall a\,\forall b\,\Bigl(
 a:A\to X\land b:B\to X
\\
&\qquad{}\land
 \exists d'\,
 \bigl(\Com(f,a,d')\land\Com(g,b,d')\bigr)
\\
&\qquad{}\Rightarrow
 \exists!u\,\bigl(
 u:P\to X
\\
&\qquad\qquad{}\land
 \Com(j_A,u,a)\land\Com(j_B,u,b)
 \bigr)
\Bigr)
\Bigr).
\end{aligned}
\]
\end{definicio}

\begin{notacio}[Propietat universal del pushout]\label{not:pushout}
El quadrat
\begin{center}
\begin{tikzpicture}[scale=1.05,every node/.style={font=\small},>=Stealth]
  \node (C) at (0,2.2) {\(C\)};
  \node (B) at (4.0,2.2) {\(B\)};
  \node (A) at (0,0) {\(A\)};
  \node (P) at (4.0,0) {\(P\)};
  \node (X) at (6.8,1.1) {\(X\)};
  \draw[->,thick] (C) -- node[above] {\(g\)} (B);
  \draw[->,thick] (C) -- node[left] {\(f\)} (A);
  \draw[->,thick] (B) -- node[right] {\(j_B\)} (P);
  \draw[->,thick] (A) -- node[below] {\(j_A\)} (P);
  \draw[->,thick] (A) -- node[below] {\(a\)} (X);
  \draw[->,thick] (B) -- node[above] {\(b\)} (X);
  \draw[->,thick,dashed] (P) -- node[below right] {\(u\)} (X);
\end{tikzpicture}
\end{center}
és \emph{cocartesià} quan, per a totes les aplicacions \(a:A\to X\) i \(b:B\to X\) amb
\[
a\circ f=b\circ g,
\]
existeix una única \(u:P\to X\) tal que
\[
u\circ j_A=a,
\qquad
u\circ j_B=b.
\]
\end{notacio}

\begin{teorema}[Construcció del pushout]\label{teo:pushout-construccio}
Tot parell d'aplicacions \(f:C\to A\) i \(g:C\to B\) amb domini comú té un pushout. Es construeix com el coequalitzador de les aplicacions paral·leles
\[
i_A\circ f,
\qquad
i_B\circ g:C\to A+B,
\]
on \(i_A\) i \(i_B\) són les injeccions de la suma.
\end{teorema}

\begin{prova}
Sigui
\[
q:A+B\to P
\]
el coequalitzador de \(i_A\circ f\) i \(i_B\circ g\), i definim
\[
j_A:=q\circ i_A,
\qquad
j_B:=q\circ i_B.
\]
La igualtat del coequalitzador dona
\[
j_A\circ f=q\circ i_A\circ f
=q\circ i_B\circ g
=j_B\circ g.
\]

Siguin \(a:A\to X\) i \(b:B\to X\) amb \(a\circ f=b\circ g\). La propietat universal de la suma determina
\[
[a,b]:A+B\to X
\]
amb \([a,b]\circ i_A=a\) i \([a,b]\circ i_B=b\). A més,
\[
[a,b]\circ i_A\circ f
=a\circ f
=b\circ g
=[a,b]\circ i_B\circ g.
\]
Per la propietat universal del coequalitzador, existeix una única \(u:P\to X\) amb
\[
u\circ q=[a,b].
\]
Aleshores \(u\circ j_A=a\) i \(u\circ j_B=b\). Si \(u'\) compleix les mateixes igualtats, les aplicacions \(u'\circ q\) i \([a,b]\) coincideixen després de compondre-les amb \(i_A\) i \(i_B\); la unicitat en la propietat universal de la suma implica \(u'\circ q=[a,b]\), i la unicitat del coequalitzador dona \(u'=u\).
\end{prova}

\section{Colímits finits}\label{sec:colimits-finits}

La construcció dels colímits finits és dual de la dels límits finits. Ara sí que fa servir l'axioma d'existència de coequalitzadors.

\begin{notacio}[Sumes finites]\label{not:sumes-finits}
Per a una llista finita de conjunts \(A_1,\ldots,A_n\), escriurem
\[
\coprod_{i=1}^{n}A_i
\]
per una suma iterada escollida. La suma buida és \(\mathbf0\), la suma d'un sol conjunt és aquest mateix conjunt i les sumes de longitud més gran es construeixen iterant les sumes binàries.
\end{notacio}

\begin{proposicio}[Existència de les sumes finites]\label{prop:sumes-finits}
Tota llista finita de conjunts té una suma, única llevat d'un únic isomorfisme compatible amb les injeccions.
\end{proposicio}

\begin{prova}
La suma buida és el conjunt inicial \(\mathbf0\), i el cas d'un sol conjunt és immediat. El pas inductiu consisteix a sumar el conjunt següent a una suma iterada ja construïda. La unicitat se segueix de les propietats universals, exactament com en el cas binari.
\end{prova}

Fixem novament un diagrama finit \(D\) amb conjunts \((A_i)_i\) i aplicacions \(f_\alpha:A_{s(\alpha)}\to A_{t(\alpha)}\).

\begin{definicio}[Cocon]\index{Cocon}\label{def:cocon}
Un \emph{cocon} sobre \(D\), amb vèrtex \(X\), és una família finita d'aplicacions
\[
c_i:A_i\to X
\]
tal que, per a cada \(f_\alpha\),
\[
c_{t(\alpha)}\circ f_\alpha=c_{s(\alpha)}.
\]
Per al diagrama fix, aquesta condició és una abreviació d'una conjunció finita del llenguatge formal.
\end{definicio}

\begin{definicio}[Colímit]\index{Colímit}\label{def:colimit}
Un cocon
\[
\bigl(L,(\ell_i)_i\bigr)
\]
és un \emph{colímit} de \(D\) si, per a qualsevol altre cocon
\[
\bigl(X,(c_i)_i\bigr),
\]
existeix una única aplicació \(u:L\to X\) tal que
\[
u\circ\ell_i=c_i
\qquad\text{per a tot }i.
\]
Com en el cas dels límits, aquesta definició és un esquema finit del metallenguatge que es desplega en una fórmula per a cada diagrama concret.
\end{definicio}

\begin{teorema}[Existència dels colímits finits]\label{teo:colimits-finits}
Tot diagrama finit té colímit.
\end{teorema}

\begin{prova}
Formem les sumes finites
\[
S=\coprod_i A_i,
\qquad
R=\coprod_\alpha A_{s(\alpha)},
\]
amb injeccions \(u_i:A_i\to S\) i \(v_\alpha:A_{s(\alpha)}\to R\). La propietat universal de \(R\) determina dues aplicacions \(r,r':R\to S\) mitjançant
\[
r\circ v_\alpha=u_{s(\alpha)},
\qquad
r'\circ v_\alpha=u_{t(\alpha)}\circ f_\alpha.
\]
Sigui
\[
q:S\to L
\]
el coequalitzador de \(r\) i \(r'\), i posem
\[
\ell_i:=q\circ u_i:A_i\to L.
\]
Com que \(q\circ r=q\circ r'\), en compondre amb cada \(v_\alpha\) obtenim
\[
\ell_{s(\alpha)}=\ell_{t(\alpha)}\circ f_\alpha.
\]
Per tant, \((L,(\ell_i)_i)\) és un cocon.

Sigui \((X,(c_i)_i)\) un cocon qualsevol. La propietat universal de \(S\) proporciona una única aplicació \(c:S\to X\) amb
\[
c\circ u_i=c_i.
\]
Per a cada \(\alpha\),
\[
\begin{aligned}
c\circ r\circ v_\alpha
&=c_{s(\alpha)}
=c_{t(\alpha)}\circ f_\alpha\\
&=c\circ u_{t(\alpha)}\circ f_\alpha
=c\circ r'\circ v_\alpha.
\end{aligned}
\]
La unicitat en la propietat universal de la suma \(R\) implica \(c\circ r=c\circ r'\). Per la propietat universal del coequalitzador, existeix una única \(u:L\to X\) amb
\[
u\circ q=c.
\]
Aleshores
\[
u\circ\ell_i=u\circ q\circ u_i=c_i.
\]
Si \(u'\) té les mateixes components, les aplicacions \(u'\circ q\) i \(c\) coincideixen sobre totes les injeccions \(u_i\); per la unicitat en la propietat universal de la suma \(S\), són iguals, i la unicitat del coequalitzador dona \(u'=u\).
\end{prova}

\begin{resumcap}
\apartat{Idees centrals}
\begin{enumerate}
  \item L'axioma de disjunció i estabilitat diu que les sumes es comporten com unions disjuntes i que aquesta descomposició es conserva per canvi de base. És la primera propietat \emph{global} de les construccions: la propietat universal sola no la dona.
  \item Tres conseqüències en cadena: \(\mathbf0\) és estricte (tota aplicació cap a \(\mathbf0\) és isomorfisme), \(0\ne1\) a \(\mathbf2\), i un conjunt sense punts és inicial. Amb elles, tot punt d'una suma ve d'un sol sumand i tot epimorfisme és exhaustiu sobre punts, sense elecció.
  \item La part buida \(\bot_A\) i la part total \(\top_A\) tenen característiques \(0\circ{!}_A\) i \(1\circ{!}_A\); el quadrat característic és un pullback, de manera que tota part és la imatge inversa de \(1\colon\mathbf1\hookrightarrow\mathbf2\) per la seva característica. El complement \(m^c\) és la imatge inversa de \(0\), i \(m\) i \(m^c\) descomponen \(A\) en una suma.
  \item La negació, la conjunció i la disjunció són aplicacions \(\neg\colon\mathbf2\to\mathbf2\) i \(\wedge,\vee\colon\mathbf2\times\mathbf2\to\mathbf2\); les operacions entre parts es tradueixen en operacions entre característiques, i les lleis booleanes es comproven sobre els quatre (o vuit) punts de \(\mathbf2\times\mathbf2\) per extensionalitat de les taules de veritat.
  \item Un límit és un con universal sobre un diagrama finit. Tots els límits finits existeixen a partir de \(\mathbf1\), els productes binaris i els equalitzadors, sense cap axioma nou; els colímits finits existeixen a partir de \(\mathbf0\), les sumes i els coequalitzadors, que sí que es postulen.
  \item Els diagrames finits, els cons i els límits pertanyen al metallenguatge: per a cada forma finita concreta, l'enunciat és una fórmula i la demostració, una demostració de la teoria.
\end{enumerate}
\apartat{Errors típics} Creure que la propietat universal de la suma ja implica que les injeccions són disjuntes o monomorfismes; pensar que \(0\ne1\) és evident, quan és un teorema que depèn de la disjunció i de la no-degeneració; suposar que un conjunt \enquote{sense elements} és inicial per definició, quan és el lema del buit i usa la separació global; llegir \(x\in m^c\iff x\notin m\) per a elements generalitzats, quan només val per a punts globals (un element generalitzat es pot repartir entre \(m\) i \(m^c\)); confondre la unió, definida per la característica, amb una construcció universal com la suma; i aplicar la dualitat per obtenir els colímits dels límits, quan la construcció dels límits no usa cap axioma i la dels colímits usa el de coequalitzadors.
\apartat{Lectura recomanada} \cite[cap.~3 i~4]{lawvere-rosebrugh} per als límits finits, les sumes disjuntes i els colímits, i \cite[apèndix~A]{lawvere-rosebrugh} per a la lògica com a àlgebra de parts, que és el contingut de la secció~\ref{sec:booleanes}; \cite[article~VI, sessions~32--33]{lawvere-schanuel} per a una introducció molt elemental als subobjectes i a l'objecte de valors de veritat; \cite[§§3.2--3.4 i §5.4]{awodey} per a coequalitzadors, pushouts i colímits en una categoria general; \cite[pàg.~11--12]{bell2006} per al paper de \(\mathbf2\) com a objecte de valors de veritat; i \cite[§§4--5]{halmos} per a les operacions booleanes fetes amb elements.
\end{resumcap}

\chapter{Relacions, quocients i elecció}\label{cap:relacions}

\begin{objectius}
\apartat{Prerequisits} Del capítol~\ref{cap:parts}: parts, pertinença generalitzada, productes i aparellament, pullbacks i la caracterització dels monomorfismes; seccions, retraccions i exhaustivitat sobre elements generalitzats. Del capítol~\ref{cap:valors}: coequalitzadors i la seva unicitat, i la proposició que els epimorfismes són exhaustius sobre punts. De la teoria clàssica, la noció informal de relació d'equivalència, classe i partició.
\apartat{Objectius} En acabar el capítol, el lector hauria de poder:
\begin{itemize}
  \item definir una relació com a part d'un producte i llegir-la per elements generalitzats, amb el cas global com a cas particular;
  \item construir la relació inversa i la diagonal, i formular estructuralment la reflexivitat, la simetria i la transitivitat mitjançant les aplicacions \(e\), \(s\) i \(c\);
  \item demostrar que el parell nucli d'una aplicació qualsevol és una relació d'equivalència;
  \item definir el quocient com a coequalitzador de les projeccions i l'epimorfisme regular, i demostrar que tot quocient és un epimorfisme regular;
  \item enunciar l'axioma d'efectivitat, explicar per què cal, i demostrar la correspondència entre relacions d'equivalència i epimorfismes regulars;
  \item identificar relacions d'equivalència, particions i quocients com a tres descripcions del mateix objecte;
  \item enunciar l'axioma d'elecció en forma estructural, distingir-ne les tres lectures, i deduir-ne que tot epimorfisme és regular, escindit i exhaustiu sobre elements generalitzats;
  \item distingir el que depèn de l'elecció del que ja valia sense ella: en particular, que l'exhaustivitat sobre punts no en depèn.
\end{itemize}
\end{objectius}

Amb els productes disposem d'un lloc on viuen les parelles, i amb les parts, d'una manera d'assenyalar-ne subconjunts. La combinació de totes dues nocions és la de \emph{relació}. Aquest capítol la desenvolupa fins a les seves conseqüències més fortes: les relacions d'equivalència, formulades estructuralment; els quocients que en resulten, entesos com a coequalitzadors; la correspondència exacta entre relacions d'equivalència i epimorfismes regulars; i, finalment, l'axioma d'elecció i les seves conseqüències. Com als capítols anteriors, definirem les nocions primer sobre elements generalitzats i en recuperarem la lectura clàssica sobre punts globals com a cas particular.

\section{Relacions com a parts de productes}\label{sec:relacions}

\begin{definicio}[Relació]\index{Relació}\label{def:relacio}
Una \emph{relació} d'un conjunt \(A\) cap a un conjunt \(B\) és una part d'un producte de \(A\) i \(B\). Escrivint \(\operatorname{Part}(r,P)\) per a «\(r\) és una part de \(P\)», i recordant que el predicat \(\operatorname{Producte}\) (definició~\ref{def:producte}) porta com a arguments el vèrtex i les dues projeccions,
\[
\operatorname{Rel}(r,A,B)
\iff
\exists P\,\exists p_A\,\exists p_B\,
\bigl(\operatorname{Producte}(P,p_A,p_B,A,B)\land\operatorname{Part}(r,P)\bigr).
\]
\end{definicio}

Sigui \(r\colon R\hookrightarrow A\times B\) una relació, amb projeccions del producte \(p_A\colon A\times B\to A\) i \(p_B\colon A\times B\to B\). La composició amb \(r\) dona les dues \emph{projeccions de la relació}
\[
r_A:=p_A\circ r\colon R\to A,
\qquad
r_B:=p_B\circ r\colon R\to B.
\]

\begin{notacio}[Pertinença generalitzada a una relació]\label{not:relacio}
Siguin \(x\colon X\to A\) i \(y\colon X\to B\) dos elements generalitzats en un mateix estadi \(X\). Direm que \emph{estan relacionats en l'estadi \(X\)}, i ho escriurem \(x\mathrel{r_X}y\), quan la parella \(\langle x,y\rangle\colon X\to A\times B\) és un element generalitzat de la part \(r\):
\[
x\mathrel{r_X}y
\quad\overset{\text{def}}{\Longleftrightarrow}\quad
\langle x,y\rangle\in_X r.
\]
El cas global és el de l'estadi terminal \(X=\mathbf1\): per a punts \(a\colon\mathbf1\to A\) i \(b\colon\mathbf1\to B\), escriurem \(a\mathrel{r}b\) en lloc de \(a\mathrel{r_{\mathbf1}}b\).
\end{notacio}

\begin{observacio}[Lectura clàssica]\label{obs:relacio-classica}
En el model clàssic, una relació entre \(A\) i \(B\) és un conjunt de parelles ordenades, un subconjunt de \(A\times B\); la nostra definició hi coincideix, amb l'única diferència que el subconjunt es presenta com un monomorfisme. La pertinença generalitzada \(x\mathrel{r_X}y\) recupera, en l'estadi global \(\mathbf1\), la pertinença habitual \((a,b)\in r\).
\end{observacio}

\section{La relació inversa i la diagonal}\label{sec:inversa-diagonal}

Dues relacions es construeixen des de zero amb les eines del capítol~\ref{cap:parts}: la que inverteix l'ordre dels arguments i la que identifica cada element amb ell mateix.

\begin{definicio}[Relació inversa]\index{Relació inversa}\label{def:relacio-inversa}
Sigui \(r\colon R\hookrightarrow A\times B\) una relació. La seva \emph{inversa} és la part de \(B\times A\) obtinguda component \(r\) amb l'isomorfisme de commutació \(\sigma\colon A\times B\to B\times A\):
\[
\operatorname{Inv}(r)\;:=\;\sigma\circ r\colon R\hookrightarrow B\times A.
\]
\end{definicio}

\begin{proposicio}\label{prop:inversa-pertinenca}
Per a tot estadi \(X\) i tots els elements generalitzats \(x\colon X\to A\), \(y\colon X\to B\), es té \(y\mathrel{\operatorname{Inv}(r)_X}x\) si i només si \(x\mathrel{r_X}y\). En particular, \(\operatorname{Inv}(\operatorname{Inv}(r))\equiv_{A\times B}r\).
\end{proposicio}

\begin{prova}
L'isomorfisme de commutació compleix \(\sigma\circ\langle x,y\rangle=\langle y,x\rangle\). Per tant, \(\langle y,x\rangle\) es factoritza a través de \(\sigma\circ r\) si i només si \(\langle x,y\rangle\) es factoritza a través de \(r\), que és el que afirma l'equivalència. Com que \(\sigma\circ\sigma=\id\), aplicar la inversa dues vegades retorna la relació original.
\end{prova}

\begin{definicio}[Diagonal]\index{Diagonal}\label{def:diagonal}
La \emph{diagonal} d'un conjunt \(A\) és la part de \(A\times A\) donada per l'aplicació d'aparellament de la identitat amb ella mateixa:
\[
\Delta_A\;:=\;\langle\id_A,\id_A\rangle\colon A\hookrightarrow A\times A.
\]
\end{definicio}

\begin{proposicio}\label{prop:diagonal-mono}
La diagonal \(\Delta_A\) és un monomorfisme. Per a tot estadi \(X\) i tots \(x,x'\colon X\to A\), es té \(x\mathrel{(\Delta_A)_X}x'\) si i només si \(x=x'\). A més, \(\operatorname{Inv}(\Delta_A)\equiv_{A\times A}\Delta_A\).
\end{proposicio}

\begin{prova}
Component \(\Delta_A\) amb la primera projecció s'obté \(p_1\circ\langle\id_A,\id_A\rangle=\id_A\), que és un monomorfisme; per tant \(\Delta_A\) és un monomorfisme. Ara, \(\langle x,x'\rangle\) es factoritza a través de \(\Delta_A\) si i només si existeix \(w\colon X\to A\) amb \(\langle x,x'\rangle=\langle\id_A,\id_A\rangle\circ w=\langle w,w\rangle\), és a dir, si i només si \(x=w=x'\). Finalment, la simetria de la diagonal se segueix de \(\sigma\circ\langle\id_A,\id_A\rangle=\langle\id_A,\id_A\rangle\).
\end{prova}

\section{Relacions d'equivalència}\label{sec:equivalencia}

Les tres propietats clàssiques ---reflexiva, simètrica i transitiva--- s'han de formular estructuralment, és a dir, mitjançant aplicacions entre \(R\), \(A\) i els productes fibrats corresponents, i no només com a condicions sobre els punts globals. Una condició puntual no proporciona, en general, les aplicacions que realitzen internament les operacions sobre els parells relacionats; per això prenem com a definició la forma estructural, i en recuperem la lectura puntual en una observació.

Sigui \(r\colon R\hookrightarrow A\times A\) una relació sobre \(A\), amb projeccions \(r_1:=p_1\circ r\) i \(r_2:=p_2\circ r\) de \(R\) cap a \(A\).

\begin{definicio}[Relació d'equivalència]\index{Relació d'equivalència}\label{def:equivalencia}
Una relació \(r\colon R\hookrightarrow A\times A\) sobre \(A\) és una \emph{relació d'equivalència} si compleix les tres condicions estructurals següents.
\begin{enumerate}
\item \emph{Reflexivitat}: existeix una aplicació \(e\colon A\to R\) tal que
\[
r_1\circ e=\id_A,
\qquad
r_2\circ e=\id_A.
\]
\item \emph{Simetria}: existeix una aplicació \(s\colon R\to R\) tal que
\[
r_1\circ s=r_2,
\qquad
r_2\circ s=r_1.
\]
\item \emph{Transitivitat}: formant el producte fibrat \(R\times_A R\) de \(r_2\colon R\to A\) amb \(r_1\colon R\to A\), amb projeccions \(p,q\colon R\times_A R\to R\), existeix una aplicació \(c\colon R\times_A R\to R\) tal que
\[
r_1\circ c=r_1\circ p,
\qquad
r_2\circ c=r_2\circ q.
\]
\end{enumerate}
\end{definicio}

\begin{notacio}[El predicat d'equivalència]\label{not:equiv-predicat}
Escriurem \(\operatorname{Equiv}(r,A)\) per abreujar la conjunció de les tres condicions anteriors. En desplegar-la cal recordar que ni el producte \(A\times A\), ni el seu domini, ni les projeccions són termes primitius: s'han de quantificar amb el predicat de producte \(\operatorname{Producte}\) (definició~\ref{def:producte}) i el de pullback \(\operatorname{Pullback}\) (definició~\ref{def:pullback}). Abreujarem amb \(\operatorname{Prod}(P,\pi_1,\pi_2;A,A):=\operatorname{Producte}(P,\pi_1,\pi_2,A,A)\) el fet que \((P,\pi_1,\pi_2)\) és un producte de \(A\) amb ell mateix ---la mateixa abreviació \(\operatorname{Prod}(P,\pi_1,\pi_2;X,Y)\) que reprendrem al capítol~\ref{cap:exponencials}. A diferència de \(\operatorname{Rel}\), que amaga el seu propi producte sota un existencial intern, aquí necessitem que \(P\), les seves projeccions \(\pi_1,\pi_2\) i el domini \(R\) de la part \(r\) restin \emph{visibles} a tota la fórmula, perquè les tres condicions s'hi refereixen; per això els introduïm tots amb un únic bloc de quantificadors al capdavant i els mantenim dins d'un sol abast:
\[
\begin{aligned}
\operatorname{Equiv}(r,A)
\iff{}&
\exists P\,\exists\pi_1\,\exists\pi_2\,\exists R\,\exists r_1\,\exists r_2\,\Bigl(
\operatorname{Prod}(P,\pi_1,\pi_2;A,A)\land\operatorname{Part}(r,P)
\\
&\qquad{}\land\,\Dom(r,R)\land r_1=\pi_1\circ r\land r_2=\pi_2\circ r
\\
&\qquad{}\land\exists e\,\bigl(e\colon A\to R\land r_1\circ e=\id_A\land r_2\circ e=\id_A\bigr)
\\
&\qquad{}\land\exists s\,\bigl(s\colon R\to R\land r_1\circ s=r_2\land r_2\circ s=r_1\bigr)
\\
&\qquad{}\land\exists F\,\exists p\,\exists q\,\exists c\,\bigl(
\operatorname{Pullback}(F,p,q,r_2,r_1)
\\
&\qquad\qquad{}\land\,c\colon F\to R
\land r_1\circ c=r_1\circ p\land r_2\circ c=r_2\circ q
\bigr)
\Bigr),
\end{aligned}
\]
on ara \(P\) és el producte \(A\times A\) amb projeccions \(\pi_1,\pi_2\); \(r\colon R\hookrightarrow P\) és la part que representa la relació, amb domini \(R\); \(r_1,r_2\colon R\to A\) són les seves projeccions, \emph{definides dins la fórmula} per les equacions \(r_1=\pi_1\circ r\), \(r_2=\pi_2\circ r\) i no preses com a primitives; i \(F\) és el domini del producte fibrat \(R\times_A R\) de \(r_2\) amb \(r_1\). Tots ells queden lligats pels existencials del capdavant i determinats llevat d'isomorfisme pels predicats \(\operatorname{Prod}\) i \(\operatorname{Pullback}\). El parèntesi extern no es tanca fins després de l'última clàusula, de manera que cap condició no fa servir \(R\), \(r_1\), \(r_2\) ni \(F\) fora del seu abast. Aquí el desplegament arriba fins als predicats \(\operatorname{Prod}\), \(\operatorname{Part}\) i \(\operatorname{Pullback}\); les operacions que hi resten ---la composició \(\circ\), la igualtat d'aplicacions i les fletxes tipades--- són convenis eliminables segons el \hyperref[not:criteri-formal]{criteri d'escriptura formal} de la introducció. A partir d'ara la usarem abreujada amb \(A\times A\), \(R\times_A R\), \(r_1\), \(r_2\), \(\id_A\) i les fletxes tipades, entès que cada aparició s'elimina pels predicats corresponents ---i, en darrera instància, per \(\Dom\), \(\Cod\) i \(\Com\)--- quan els tipus són inequívocs.
\end{notacio}

\begin{notacio}[Lectura de les tres condicions]\label{not:equivalencia}
Les tres aplicacions es poden llegir en termes de la pertinença generalitzada. La reflexivitat equival a la inclusió de parts \(\Delta_A\subseteq_{A\times A}r\); la simetria, a \(\operatorname{Inv}(r)\equiv_{A\times A}r\). Per a la transitivitat, el producte fibrat \(R\times_A R\) recull els parells \((\rho,\rho')\) de testimonis amb \(r_2\circ\rho=r_1\circ\rho'\) ---és a dir, tals que el segon component de \(\rho\) coincideix amb el primer de \(\rho'\)---, i l'aplicació \(c\) en produeix un tercer testimoni que enllaça el primer component de \(\rho\) amb el segon de \(\rho'\).
\end{notacio}

\begin{observacio}[Lectura clàssica sobre punts globals]\label{obs:equivalencia-classica}
En l'estadi global, les tres condicions recuperen les propietats habituals. La reflexivitat dona \(a\mathrel{r}a\) per a tot punt \(a\colon\mathbf1\to A\), component \(e\) amb \(a\). La simetria dona \(a\mathrel{r}b\Rightarrow b\mathrel{r}a\), aplicant \(s\) al testimoni. I la transitivitat dona \(a\mathrel{r}b\land b\mathrel{r}c\Rightarrow a\mathrel{r}c\): dos testimonis amb el punt \(b\) comú determinen un element global de \(R\times_A R\), i \(c\) en produeix el testimoni de \(a\mathrel{r}c\). Convé remarcar que aquestes implicacions puntuals no bastarien, en canvi, per \emph{definir} l'equivalència: sense les aplicacions \(e\), \(s\) i \(c\) no disposaríem de les operacions internes que el desenvolupament posterior necessita.
\end{observacio}

\section{El parell nucli d'una aplicació}\label{sec:kernel-pair}

L'exemple fonamental de relació d'equivalència ---i el que farà de pont cap als quocients--- prové de qualsevol aplicació.

\begin{definicio}[Parell nucli]\index{Parell nucli}\label{def:kernel-pair}
Sigui \(f\colon A\to B\). El seu \emph{parell nucli} és la relació sobre \(A\) donada pel pullback de \(f\) amb ella mateixa:
\[
\operatorname{Nuc}(f)\;:=\;\langle k_1,k_2\rangle\colon A\times_B A\hookrightarrow A\times A,
\]
on \(k_1,k_2\colon A\times_B A\to A\) són les dues projeccions del pullback, que compleixen \(f\circ k_1=f\circ k_2\).
\end{definicio}

\begin{proposicio}[El parell nucli és un monomorfisme]\label{prop:kernel-mono}
L'aplicació \(\langle k_1,k_2\rangle\colon A\times_B A\to A\times A\) és un monomorfisme.
\end{proposicio}

\begin{prova}
Siguin \(u,v\colon X\to A\times_B A\) amb \(\langle k_1,k_2\rangle\circ u=\langle k_1,k_2\rangle\circ v\). Component aquesta igualtat amb les dues projeccions \(p_1,p_2\) del producte \(A\times A\), obtenim
\[
k_1\circ u=k_1\circ v,
\qquad
k_2\circ u=k_2\circ v.
\]
Ara bé, per la propietat universal del pullback, dues aplicacions amb codomini \(A\times_B A\) coincideixen tan bon punt coincideixen les seves composicions amb \(k_1\) i \(k_2\). Per tant, les igualtats anteriors impliquen \(u=v\), i \(\langle k_1,k_2\rangle\) és un monomorfisme.
\end{prova}

\begin{observacio}[La prova no usa que \(f\) sigui mono]\label{obs:kernel-no-mono}
Notem que la demostració anterior no suposa en cap moment que \(f\) sigui un monomorfisme ---no ho és, en general. El que fa mono el parell nucli és exclusivament la unicitat de la propietat universal del pullback, no cap propietat de \(f\).
\end{observacio}

\begin{proposicio}[El parell nucli és una relació d'equivalència]\label{prop:kernel-equivalencia}
Per a tota aplicació \(f\colon A\to B\), el parell nucli \(\operatorname{Nuc}(f)\) és una relació d'equivalència sobre \(A\).
\end{proposicio}

\begin{prova}
Escrivim \(K=A\times_B A\) amb projeccions \(k_1,k_2\), de manera que \(f\circ k_1=f\circ k_2\). Construïm les tres aplicacions estructurals.

\emph{Reflexivitat.} L'aplicació \(\id_A\) compleix \(f\circ\id_A=f\circ\id_A\), de manera que la propietat universal del pullback proporciona una única \(e\colon A\to K\) amb \(k_1\circ e=\id_A\) i \(k_2\circ e=\id_A\).

\emph{Simetria.} Les projeccions intercanviades compleixen \(f\circ k_2=f\circ k_1\), de manera que la propietat universal proporciona una única \(s\colon K\to K\) amb \(k_1\circ s=k_2\) i \(k_2\circ s=k_1\).

\emph{Transitivitat.} Formem \(K\times_A K\), el pullback de \(k_2\colon K\to A\) amb \(k_1\colon K\to A\), amb projeccions \(p,q\colon K\times_A K\to K\); per construcció \(k_2\circ p=k_1\circ q\). Aleshores
\[
f\circ(k_1\circ p)=f\circ(k_2\circ p)=f\circ(k_1\circ q)=f\circ(k_2\circ q),
\]
on la primera i la tercera igualtats són \(f\circ k_1=f\circ k_2\) i la central és \(k_2\circ p=k_1\circ q\). En particular \(f\circ(k_1\circ p)=f\circ(k_2\circ q)\), i la propietat universal del pullback \(K\) proporciona una única \(c\colon K\times_A K\to K\) amb \(k_1\circ c=k_1\circ p\) i \(k_2\circ c=k_2\circ q\). Aquestes són exactament les equacions de la transitivitat.
\end{prova}

\begin{observacio}[Lectura clàssica]\label{obs:kernel-classica}
En el model clàssic, el parell nucli de \(f\) és la relació «\(a\) i \(a'\) tenen la mateixa avaluació per \(f\)», és a dir, \(f\circ a=f\circ a'\) sobre punts globals. És l'exemple canònic de relació d'equivalència, i el teorema de correspondència d'aquest capítol dirà que, per als epimorfismes regulars, tota relació d'equivalència n'és un: la del seu propi quocient.
\end{observacio}

\section{Quocients i epimorfismes regulars}\label{sec:quocients}

Donada una relació d'equivalència, volem el conjunt que resulta d'identificar els elements relacionats. La seva propietat universal és la d'un coequalitzador de les dues projeccions.

\begin{definicio}[Quocient]\index{Quocient}\label{def:quocient}
Sigui \(r\colon R\hookrightarrow A\times A\) una relació sobre \(A\), amb projeccions \(r_1,r_2\colon R\to A\). Un \emph{quocient} de \(A\) per \(r\) és un coequalitzador de \(r_1\) i \(r_2\):
\[
\operatorname{Quoc}(q,r)
\iff
\operatorname{Coeq}(q,r_1,r_2).
\]
\end{definicio}

\begin{notacio}[Quocient]\label{not:quocient}
Denotarem el quocient per \(q\colon A\twoheadrightarrow A/r\). La seva propietat universal diu que \(q\) identifica els elements relacionats ---\(q\circ r_1=q\circ r_2\)--- i que és universal amb aquesta propietat: tota aplicació \(h\) amb \(h\circ r_1=h\circ r_2\) es factoritza de manera única a través de \(q\).
\end{notacio}

\begin{proposicio}[Existència i unicitat del quocient]\label{prop:quocient-existeix}
Tota relació sobre \(A\) té quocient, únic llevat d'un únic isomorfisme compatible, i l'aplicació quocient \(q\colon A\twoheadrightarrow A/r\) és un epimorfisme.
\end{proposicio}

\begin{prova}
L'existència és l'axioma de coequalitzadors (axioma~\ref{ax:coequalitzadors}) aplicat a \(r_1,r_2\); la unicitat llevat d'un únic isomorfisme és la unicitat dels coequalitzadors (proposició~\ref{prop:coequalitzador-unic}); i que \(q\) sigui epimorfisme, la proposició~\ref{prop:coequalitzador-epi}.
\end{prova}

Els quocients són coequalitzadors, i els coequalitzadors són un tipus especial d'epimorfisme, que aïllem amb un nom propi.

\begin{definicio}[Epimorfisme regular]\index{Epimorfisme regular}\label{def:epi-regular}
Una aplicació \(q\) és un \emph{epimorfisme regular} si és el coequalitzador d'algun parell d'aplicacions paral·leles:
\[
\operatorname{EpiReg}(q)
\iff
\exists f\,\exists g\,\operatorname{Coeq}(q,f,g).
\]
\end{definicio}

\begin{proposicio}[Tot quocient és un epimorfisme regular]\label{prop:quocient-regular}
Per a tota relació d'equivalència \(r\), l'aplicació quocient \(q\colon A\twoheadrightarrow A/r\) és un epimorfisme regular.
\end{proposicio}

\begin{prova}
Per definició, \(q\) és el coequalitzador de les projeccions \(r_1,r_2\), de manera que és, en particular, coequalitzador d'un parell d'aplicacions paral·leles.
\end{prova}

\begin{observacio}[Regular i escindit]\label{obs:regular-escindit}
Tot epimorfisme regular és un epimorfisme, per la proposició~\ref{prop:coequalitzador-epi}. El recíproc ---que tot epimorfisme sigui regular--- no el demostrarem amb els axiomes actuals; el retrobarem al final del capítol com a conseqüència de l'axioma d'elecció, en veure que tot epimorfisme escindit és regular, i el capítol~\ref{cap:imatges} en donarà una altra lectura a través de la factorització per la imatge.
\end{observacio}

\section{L'axioma d'efectivitat}\label{sec:exactitud}

Tenim dos camins entre aplicacions i relacions d'equivalència. D'una aplicació en surt una relació d'equivalència ---el seu parell nucli (proposició~\ref{prop:kernel-equivalencia}). I d'una relació d'equivalència en surt una aplicació ---el seu quocient. Manca saber que aquests dos camins es corresponen: que el parell nucli del quocient recupera la relació de partida. Això és el que la teoria, tal com la tenim, encara no garanteix.

\begin{definicio}[Parell nucli d'una aplicació, com a predicat]\index{Parell nucli d'una aplicació, com a predicat}\label{def:esparellnucli}
Direm que una relació \(n\) \emph{és un parell nucli de \(q\)}, i ho escriurem \(\operatorname{EsNuc}(n,q)\), quan \(n\) és la part de \(A\times A\) determinada per un pullback de \(q\) amb ella mateixa. Formalment,
\[
\begin{aligned}
\operatorname{EsNuc}(n,q)
\iff{}&
\exists A\,\exists B\,\exists K\,\exists k_1\,\exists k_2\,\Bigl(
q\colon A\to B
\\
&{}\land\operatorname{Pullback}(K,k_1,k_2,q,q)
\land n=\langle k_1,k_2\rangle
\land\operatorname{Part}(n,A\times A)
\Bigr).
\end{aligned}
\]
\end{definicio}

\begin{definicio}[Relació efectiva]\index{Relació efectiva}\label{def:efectiva}
Una relació d'equivalència \(r\) sobre \(A\) és \emph{efectiva} si existeix un quocient \(q\) de \(r\) i un parell nucli \(n\) de \(q\) tal que \(n\) i \(r\) són parts equivalents:
\[
\operatorname{Efectiva}(r,A)
\iff
\exists q\,\exists n\,\bigl(
\operatorname{Quoc}(q,r)\land\operatorname{EsNuc}(n,q)\land n\equiv_{A\times A}r
\bigr).
\]
\end{definicio}

\begin{axioma}[Efectivitat de les relacions d'equivalència]\label{ax:exactitud}
Tota relació d'equivalència és efectiva:
\[
\forall A\,\forall r\,\bigl(
\operatorname{Equiv}(r,A)
\;\Rightarrow\;
\operatorname{Efectiva}(r,A)
\bigr).
\]
Desplegant \(\operatorname{Efectiva}\), l'axioma afirma que per a tota relació d'equivalència \(r\) existeixen un quocient \(q\) i un parell nucli \(n\) de \(q\) amb \(n\equiv_{A\times A}r\). Tots els objectes ---quocient, pullback, producte--- estan quantificats existencialment; l'axioma no pressuposa cap tria canònica de cap d'ells.
\end{axioma}

\begin{observacio}[Sobre el nom i la formulació de l'axioma]\label{obs:exactitud-formulacio}
La formulació quantifica explícitament el quocient \(q\) en lloc d'usar un terme funcional \(q_r\). El motiu és que els quocients només estan determinats llevat d'un únic isomorfisme compatible: no disposem d'una elecció canònica de quocient per a cada relació, i per tant \(q_r\) no designaria un individu únic de la teoria. En quantificar \(q\), l'axioma afirma el que ha d'afirmar ---l'existència d'un quocient el parell nucli del qual recupera \(r\)--- sense pressuposar cap tria. La independència del representant se segueix de la unicitat del quocient llevat d'isomorfisme.

Quant al nom, hem preferit \emph{efectivitat de les relacions d'equivalència} a \emph{exactitud}. En la terminologia categòrica, una categoria regular té límits finits, admet factoritzacions (epimorfisme regular)--(monomorfisme) i els seus epimorfismes regulars són estables per canvi de base; una categoria exacta és una categoria regular en què, a més, tota relació d'equivalència és efectiva. Aquest axioma postula precisament aquesta condició addicional ---l'efectivitat---, damunt de la infraestructura regular que la teoria ja té; no forma part de la definició de regularitat que tot epimorfisme sigui regular, i de fet això no es demostra fins que s'afegeix l'axioma d'elecció (proposició~\ref{prop:epi-regular-eleccio}). El nom triat evita, doncs, atribuir a la teoria més estructura de la que aquest axioma li dona per si sol.

Com l'axioma de coequalitzadors (observació~\ref{obs:necessitat-coequalitzadors}), aquest és un axioma \emph{de treball}: en la formulació categòrica de la teoria, l'efectivitat de les relacions d'equivalència és un teorema que es demostra a partir del conjunt de parts, i per tant no figura entre els axiomes d'ETCS. El postulem aquí perquè el necessitem abans de disposar del conjunt de parts (capítol~\ref{cap:parts-interna}); el capítol~\ref{cap:models} en farà el balanç.
\end{observacio}

\begin{observacio}[Lectura clàssica]\label{obs:exactitud-classica}
En el model clàssic, l'efectivitat expressa que les classes d'equivalència parteixen el conjunt sense solapaments ni pèrdues: si dos elements tenen la mateixa classe en el quocient, ja estaven relacionats. Els coequalitzadors, garantits al capítol~\ref{cap:valors}, proporcionen el quocient; però res no assegura, sense aquest axioma, que en passar al quocient i tornar a prendre el parell nucli es recuperi la relació original en lloc d'una de més grossa.
\end{observacio}

\begin{lema}[Invariància del parell nucli sota isomorfisme del codomini]\label{lema:nucli-invariant}
Siguin \(q\colon A\to B\) i \(w\colon B\to B'\) un isomorfisme, i sigui \(q'=w\circ q\). Aleshores tot parell nucli de \(q\) és equivalent, com a part de \(A\times A\), a tot parell nucli de \(q'\).
\end{lema}

\begin{prova}
Siguin \(K\) un pullback de \(q\) amb ella mateixa, amb projeccions \(k_1,k_2\), i \(K'\) un pullback de \(q'\) amb ella mateixa, amb projeccions \(k_1',k_2'\). Com que \(w\) és un isomorfisme, per a qualsevol parell d'aplicacions \(u,v\) amb codomini \(A\) es té
\[
q\circ u=q\circ v
\quad\Longleftrightarrow\quad
q'\circ u=q'\circ v,
\]
ja que \(q'=w\circ q\) i \(w\) es pot cancel·lar. Per tant, un con sobre el diagrama de pullback de \(q\) ho és exactament sobre el de \(q'\), i les dues propietats universals determinen el mateix objecte llevat d'un únic isomorfisme: existeix un isomorfisme \(\theta\colon K\to K'\) amb \(k_1'\circ\theta=k_1\) i \(k_2'\circ\theta=k_2\). Component amb l'aparellament, \(\langle k_1',k_2'\rangle\circ\theta=\langle k_1,k_2\rangle\), de manera que les parts \(\langle k_1,k_2\rangle\) i \(\langle k_1',k_2'\rangle\) de \(A\times A\) es factoritzen l'una per l'altra a través d'un isomorfisme, i són equivalents.
\end{prova}

\begin{teorema}[Correspondència entre relacions d'equivalència i epimorfismes regulars]\label{teo:correspondencia}
Amb l'axioma d'efectivitat, els dos camins següents són inversos l'un de l'altre, llevat d'isomorfisme:
\begin{enumerate}
\item d'una relació d'equivalència \(r\) sobre \(A\), prendre un quocient \(q\colon A\twoheadrightarrow A/r\);
\item d'un epimorfisme regular \(q\colon A\twoheadrightarrow Q\), prendre el seu parell nucli \(\operatorname{Nuc}(q)\).
\end{enumerate}
És a dir: per a tota relació d'equivalència \(r\), \(\operatorname{Nuc}(q)\equiv_{A\times A} r\); i per a tot epimorfisme regular \(q\colon A\twoheadrightarrow Q\), el quocient de \(\operatorname{Nuc}(q)\) és isomorf a \(Q\) de manera compatible amb \(q\).
\end{teorema}

\begin{prova}
\emph{Camí (1) seguit de (2).} Sigui \(r\) una relació d'equivalència i \(q\) un quocient \emph{qualsevol} seu. L'axioma d'efectivitat (axioma~\ref{ax:exactitud}) proporciona \emph{algun} quocient \(q_0\) i un parell nucli seu equivalent a \(r\). Ara bé, dos quocients qualssevol de \(r\) són isomorfs de manera compatible, per la unicitat dels coequalitzadors (proposició~\ref{prop:coequalitzador-unic}): existeix un isomorfisme \(w\) amb \(q=w\circ q_0\). Pel lema~\ref{lema:nucli-invariant} anterior, el parell nucli és invariant sota postcomposició amb un isomorfisme, de manera que un parell nucli de \(q\) també és equivalent a \(r\).

\emph{Camí (2) seguit de (1).} Sigui \(q\colon A\twoheadrightarrow Q\) un epimorfisme regular, coequalitzador d'un parell \(f,g\colon C\to A\). El parell nucli \(\operatorname{Nuc}(q)\) té projeccions \(k_1,k_2\) amb \(q\circ k_1=q\circ k_2\), de manera que \(q\) coiguala \(k_1,k_2\) i, per la propietat universal del quocient \(A/\operatorname{Nuc}(q)\), es factoritza com \(q=\bar q\circ q'\), on \(q'\colon A\twoheadrightarrow A/\operatorname{Nuc}(q)\) és el quocient i \(\bar q\colon A/\operatorname{Nuc}(q)\to Q\). Vegem que \(\bar q\) és un isomorfisme.

\emph{\(\bar q\) és un epimorfisme}: com que \(q=\bar q\circ q'\) és un epimorfisme, \(\bar q\) també ho és.

\emph{\(\bar q\) és un isomorfisme}: com que \(q\) és coequalitzador de \(f,g\), tenim \(q\circ f=q\circ g\); per la propietat universal del pullback, el parell \((f,g)\) es factoritza a través de \(\operatorname{Nuc}(q)\), de manera que \(q'\circ f=q'\circ g\). Aleshores \(q'\) coiguala \(f,g\), i com que \(q\) és el coequalitzador de \(f,g\), existeix una única \(t\colon Q\to A/\operatorname{Nuc}(q)\) amb \(t\circ q=q'\). Ara \(\bar q\circ t\circ q=\bar q\circ q'=q=\id_Q\circ q\); com que \(q\) és epimorfisme, \(\bar q\circ t=\id_Q\). I \(t\circ\bar q\circ q'=t\circ q=q'\); com que \(q'\) és epimorfisme, \(t\circ\bar q=\id\). Per tant \(\bar q\) és un isomorfisme, invers de \(t\), compatible amb les aplicacions quocient.
\end{prova}

\section{Particions}\label{sec:particions}

La imatge geomètrica d'una relació d'equivalència és una \emph{partició}: una descomposició del conjunt en peces. En la teoria clàssica, una partició es descriu com una família de subconjunts no buits, disjunts i exhaustius; però encara no hem construït cap noció de família arbitrària de parts, i, a més, una relació d'equivalència pot tenir un nombre arbitrari de classes. Adoptem, doncs, la formulació estructural, que no depèn de cap noció de família ni es limita al cas finit.

\begin{definicio}[Partició]\index{Partició}\label{def:particio}
Una \emph{partició} d'un conjunt \(A\) és un epimorfisme regular amb domini \(A\),
\[
q\colon A\twoheadrightarrow Q,
\]
considerat llevat d'un isomorfisme compatible del codomini: dues particions \(q\colon A\twoheadrightarrow Q\) i \(q'\colon A\twoheadrightarrow Q'\) es consideren la mateixa si existeix un isomorfisme \(w\colon Q\to Q'\) amb \(w\circ q=q'\).
\end{definicio}

\begin{teorema}[Triangle partició--relació--quocient]\label{teo:triangle}
Amb l'axioma d'efectivitat, per a un conjunt \(A\) hi ha una correspondència biunívoca, llevat d'isomorfisme, entre:
\begin{enumerate}
\item les relacions d'equivalència sobre \(A\);
\item les particions de \(A\), és a dir, els epimorfismes regulars amb domini \(A\);
\item els quocients de \(A\), és a dir, les aplicacions \(q\colon A\twoheadrightarrow Q\) que són coequalitzador d'alguna relació d'equivalència sobre \(A\), considerades llevat d'isomorfisme compatible del codomini.
\end{enumerate}
\end{teorema}

\begin{prova}
La correspondència entre (1) i (2) és el teorema~\ref{teo:correspondencia}: a cada relació d'equivalència li correspon el seu quocient, que és un epimorfisme regular amb domini \(A\) (proposició~\ref{prop:quocient-regular}), i a cada epimorfisme regular el seu parell nucli, que és una relació d'equivalència (proposició~\ref{prop:kernel-equivalencia}); l'efectivitat fa que aquestes dues assignacions siguin inverses. Per a la identificació entre (2) i (3), tot quocient d'\(A\) per una relació d'equivalència és, per definició, un coequalitzador, i per tant un epimorfisme regular. Recíprocament, un epimorfisme regular \(q\) és el coequalitzador d'\emph{algun} parell d'aplicacions paral·leles, que no ha de ser cap relació d'equivalència; però el camí (2) seguit de (1) del teorema~\ref{teo:correspondencia} mostra que \(q=\bar q\circ q'\), amb \(q'\) el quocient d'\(A\) per \(\operatorname{Nuc}(q)\) i \(\bar q\) un isomorfisme. Per tant \(q\) és ell mateix un coequalitzador del parell nucli \(\operatorname{Nuc}(q)\), que és una relació d'equivalència (proposició~\ref{prop:kernel-equivalencia}), és a dir, un quocient d'\(A\) en el sentit de (3).
\end{prova}

\begin{observacio}[Partició estructural i partició clàssica]\label{obs:particio-classica}
Clàssicament, una partició d'\(A\) és una família de subconjunts no buits, dos a dos disjunts, la unió dels quals és \(A\). La definició~\ref{def:particio} no diu això literalment: diu que una partició és un epimorfisme regular \(q\colon A\twoheadrightarrow Q\). Les dues descripcions es connecten per les \emph{fibres} de \(q\), i la proposició següent demostra, dins la teoria, que les fibres tenen exactament les tres propietats clàssiques. El recíproc ---reconstruir \(q\) a partir d'una família de parts--- demana parlar de famílies indexades per un conjunt, cosa que farem al capítol~\ref{cap:families}; allà es veurà que una família \(E\to Q\) de parts d'\(A\) disjuntes i que cobreixen \(A\) és el mateix que una partició, i que el conjunt de les peces és el quocient \(Q\).
\end{observacio}

\begin{proposicio}[Les fibres d'una partició són les peces]\label{prop:fibres-particio}
Sigui \(q\colon A\twoheadrightarrow Q\) una partició. Per a cada punt \(y\colon\mathbf1\to Q\), la fibra \(q^{-1}(y)\hookrightarrow A\) és una part d'\(A\), i:
\begin{enumerate}
\item cap fibra no és buida: \(q^{-1}(y)\) té algun punt;
\item dues fibres sobre punts diferents són disjuntes: si \(y\ne y'\), aleshores \(q^{-1}(y)\cap q^{-1}(y')\) és la part buida;
\item les fibres cobreixen \(A\) sobre punts: tot punt \(a\colon\mathbf1\to A\) pertany a la fibra \(q^{-1}(q\circ a)\), i només a aquesta.
\end{enumerate}
La relació d'equivalència corresponent a \(q\) relaciona dos punts d'\(A\) si i només si són a la mateixa fibra.\par
\emph{Axiomes usats:} separació global, representació de les parts i disjunció i estabilitat de les sumes (a través del lema del buit i de l'exhaustivitat dels epimorfismes sobre punts); \emph{cap elecció}.
\end{proposicio}

\begin{prova}
Tot punt \(y\colon\mathbf1\to Q\) és un monomorfisme, de manera que \(q^{-1}(y)\) és la imatge inversa d'una part (definició~\ref{def:imatge-inversa}), i és una part d'\(A\).

\emph{(1)} Com que \(q\) és un epimorfisme, és exhaustiu sobre punts (proposició~\ref{prop:epi-exhaustiu-punts}): hi ha \(a\colon\mathbf1\to A\) amb \(q\circ a=y\), i per la proposició~\ref{prop:pertinenca-imatge-inversa}, \(a\in q^{-1}(y)\).

\emph{(2)} Un punt de la intersecció \(q^{-1}(y)\cap q^{-1}(y')\) és un punt \(a\) amb \(q\circ a=y\) i \(q\circ a=y'\) (proposició~\ref{prop:pertinenca-interseccio}), de manera que \(y=y'\). Si \(y\ne y'\), la intersecció no té cap punt i, pel lema del buit (lema~\ref{lema:buit}), és inicial: és la part buida.

\emph{(3)} Per a tot \(a\), \(q\circ a\) és un punt de \(Q\) i \(a\in q^{-1}(q\circ a)\) trivialment; i si \(a\in q^{-1}(y)\), aleshores \(y=q\circ a\), de manera que la fibra és única.

Finalment, dos punts \(a,a'\) estan relacionats pel parell nucli de \(q\) si i només si \(q\circ a=q\circ a'\), és a dir, si i només si són a la mateixa fibra; i el parell nucli de \(q\) és la relació d'equivalència que el teorema~\ref{teo:triangle} li assigna.
\end{prova}

\begin{observacio}[Per què només sobre punts]\label{obs:particio-punts}
Les tres propietats s'han enunciat sobre punts globals perquè és allà on tenen la forma clàssica. Sobre elements generalitzats la situació és la de sempre: un \(x\colon X\to A\) pot no caure dins de cap fibra sencera (si \(q\circ x\) no és constant), i aleshores no és a cap peça, sinó repartit entre diverses. És la mateixa diferència que separa el complement d'una part (capítol~\ref{cap:valors}) de la negació puntual. Les particions en un nombre finit de peces es podran definir quan s'hagin construït els nombres naturals i els conjunts finits.
\end{observacio}

\section{Epimorfismes escindits i quocients}\label{sec:escindits-quocients}

Les seccions i les retraccions (definició~\ref{def:seccio-retraccio}), els monomorfismes i els epimorfismes escindits (observació~\ref{obs:escindits}), i l'exhaustivitat sobre elements generalitzats (definició~\ref{def:exhaustiva-generalitzada}) ja s'han introduït al capítol~\ref{cap:parts}. No els redefinim; en recordem la terminologia i n'establim el resultat nou que fa de pont cap a l'elecció.

Recordem que una aplicació \(f\) és un \emph{epimorfisme escindit} si té una secció, és a dir, si existeix \(s\) amb \(f\circ s=\id\) (definició~\ref{def:seccio-retraccio}), i que tota secció és un monomorfisme i tota retracció un epimorfisme (proposició~\ref{prop:seccio-mono-retraccio-epi}).

\begin{proposicio}[Tot epimorfisme escindit és regular]\label{prop:escindit-regular}
Si \(q\colon A\to B\) és un epimorfisme escindit, aleshores \(q\) és el coequalitzador del seu parell nucli; en particular, és un epimorfisme regular.
\end{proposicio}

\begin{prova}
Sigui \(s\colon B\to A\) una secció de \(q\), amb \(q\circ s=\id_B\), i sigui \(\operatorname{Nuc}(q)\) el parell nucli, amb projeccions \(k_1,k_2\colon K\to A\) i \(q\circ k_1=q\circ k_2\). Comprovem que \(q\) té la propietat universal del coequalitzador de \(k_1,k_2\).

Ja tenim \(q\circ k_1=q\circ k_2\). Sigui ara \(h\colon A\to C\) amb \(h\circ k_1=h\circ k_2\); hem de trobar una única \(\bar h\colon B\to C\) amb \(\bar h\circ q=h\). Definim \(\bar h:=h\circ s\). Per veure que \(\bar h\circ q=h\), considerem les dues aplicacions \(\id_A\) i \(s\circ q\), totes dues de \(A\) a \(A\). Compleixen \(q\circ\id_A=q=q\circ(s\circ q)\), ja que \(q\circ s\circ q=\id_B\circ q=q\). Per la propietat universal del pullback, el parell \((\id_A,\,s\circ q)\) es factoritza a través de \(\operatorname{Nuc}(q)\): existeix \(d\colon A\to K\) amb \(k_1\circ d=\id_A\) i \(k_2\circ d=s\circ q\). Aleshores
\[
h=h\circ\id_A=h\circ k_1\circ d=h\circ k_2\circ d=h\circ s\circ q=\bar h\circ q,
\]
on hem usat \(h\circ k_1=h\circ k_2\) al pas central. La unicitat de \(\bar h\) se segueix del fet que \(q\) és un epimorfisme. Per tant, \(q\) és el coequalitzador de \(k_1,k_2\), i és un epimorfisme regular.
\end{prova}

\begin{observacio}[Una secció com a sistema de representants]\label{obs:seccio-representants}
En la lectura clàssica, una secció \(s\colon B\to A\) d'un quocient \(q\colon A\twoheadrightarrow B\) tria, per a cada peça, un element que la representa: com que \(q\circ s=\id_B\), el punt \(s(b)\) pertany a la fibra sobre \(b\), és a dir, a la classe corresponent. Una secció és, doncs, un \emph{sistema de representants} de les classes. Que tota partició admeti un sistema de representants ---que tot quocient tingui secció--- no se segueix, però, dels axiomes actuals: és el contingut de l'axioma d'elecció, que introduïm tot seguit.
\end{observacio}

\section{L'axioma d'elecció}\label{sec:eleccio}

\begin{axioma}[Elecció]\label{ax:eleccio}
Tot epimorfisme és escindit:
\[
\begin{aligned}
\forall A\,\forall B\,\forall q\,\Bigl(
&\operatorname{Epi}(q)\land\Dom(q,A)\land\Cod(q,B)
\\
&{}\Rightarrow
\exists s\,\bigl(\Dom(s,B)\land\Cod(s,A)\land\Com(s,q,\id_B)\bigr)
\Bigr).
\end{aligned}
\]
\end{axioma}

\begin{notacio}[Elecció]\label{not:eleccio}
En notació de fletxes: per a tot epimorfisme \(q\colon A\twoheadrightarrow B\) existeix una aplicació \(s\colon B\to A\) amb \(q\circ s=\id_B\).
\end{notacio}

\begin{observacio}[Les tres lectures de l'elecció]\label{obs:eleccio-lectures}
L'axioma admet tres formulacions que convé mantenir distingides. La \emph{formulació estructural} és la que hem donat: tot epimorfisme té secció. La \emph{lectura clàssica} diu que, donada una família de conjunts no buits ---les fibres d'un epimorfisme---, se'n pot triar simultàniament un element de cadascuna; és la forma més intuïtiva, però parla de fibres i eleccions, nocions que aquí són externes. I la \emph{formulació en termes de famílies} ---tota família de conjunts no buits té una funció d'elecció--- és, de moment, també externa, perquè no hem construït encara els objectes de famílies arbitràries. Adoptem la formulació estructural com a axioma, i reservem les altres dues per a la lectura.
\end{observacio}

\section{Conseqüències de l'elecció}\label{sec:consequencies-eleccio}

L'axioma d'elecció, combinat amb els resultats anteriors, tanca diverses qüestions que havien quedat obertes.

\begin{proposicio}[Tot epimorfisme és regular]\label{prop:epi-regular-eleccio}
Amb l'axioma d'elecció, tot epimorfisme és un epimorfisme regular.
\end{proposicio}

\begin{prova}
Per l'axioma d'elecció (axioma~\ref{ax:eleccio}), tot epimorfisme és escindit; i per la proposició~\ref{prop:escindit-regular}, tot epimorfisme escindit és regular. Avancem que al capítol~\ref{cap:imatges} (corol·lari~\ref{cor:epi-regular-sense-eleccio}) es veurà que la conclusió val sense l'elecció, per una altra via: la factorització imatge.
\end{prova}

\begin{corollari}[La correspondència s'estén a tots els epimorfismes]\label{cor:correspondencia-tots-epis}
Amb els axiomes d'efectivitat i d'elecció, la correspondència del teorema~\ref{teo:correspondencia} val entre les relacions d'equivalència sobre \(A\) i \emph{tots} els epimorfismes amb domini \(A\), no només els regulars.
\end{corollari}

\begin{prova}
Per la proposició~\ref{prop:epi-regular-eleccio}, tot epimorfisme és regular; per tant, la classe dels epimorfismes amb domini \(A\) coincideix amb la dels epimorfismes regulars amb domini \(A\), i el teorema~\ref{teo:correspondencia} s'hi aplica.
\end{prova}

\begin{proposicio}[Exhaustivitat generalitzada dels epimorfismes]\label{prop:epi-exhaustiu-generalitzat}
Amb l'axioma d'elecció, tot epimorfisme \(q\colon A\twoheadrightarrow B\) és exhaustiu sobre elements generalitzats: per a tot estadi \(X\) i tot \(y\colon X\to B\), existeix \(x\colon X\to A\) amb \(q\circ x=y\).
\end{proposicio}

\begin{prova}
Per l'elecció, \(q\) té una secció \(s\colon B\to A\). Per la proposició~\ref{prop:exhaustiva-generalitzada-escindit}, tenir secció equival a l'exhaustivitat sobre elements generalitzats; explícitament, donat \(y\colon X\to B\), l'element \(x:=s\circ y\) compleix \(q\circ x=q\circ s\circ y=y\).
\end{prova}

\begin{proposicio}[Epimorfisme i exhaustivitat global; sense elecció]\label{prop:epi-exhaustiu-global}
Amb l'axioma de separació global, i sense l'axioma d'elecció, una aplicació és un epimorfisme si i només si és exhaustiva sobre els elements globals.
\end{proposicio}

\begin{prova}
Si \(q\) és exhaustiva sobre els elements globals, és un epimorfisme per l'observació~\ref{obs:exhaustiva-epi} del capítol~\ref{cap:parts}, que usa la separació global. Recíprocament, si \(q\) és un epimorfisme, és exhaustiva sobre punts per la proposició~\ref{prop:epi-exhaustiu-punts} del capítol~\ref{cap:valors}, que usa la separació global i la disjunció de les sumes.
\end{prova}

\begin{observacio}[El que l'elecció no afegeix]\label{obs:eleccio-no-afegeix}
La proposició anterior també se segueix de la proposició~\ref{prop:epi-exhaustiu-generalitzat} prenent l'estadi \(X=\mathbf1\), però per aquesta via s'hi arriba amb l'elecció, cosa que és innecessària. Convé no atribuir a l'elecció més del que fa: el que aporta és el pas dels punts globals a \emph{tots} els estadis alhora, és a dir, la secció.
\end{observacio}

\begin{proposicio}[Representants de les classes]\label{prop:representants}
Amb l'axioma d'elecció, tot quocient \(q\colon A\twoheadrightarrow A/r\) admet una secció. En la interpretació clàssica, això significa que es pot triar un representant de cada classe d'equivalència.
\end{proposicio}

\begin{prova}
El quocient \(q\) és un epimorfisme (proposició~\ref{prop:quocient-existeix}); per l'elecció, té secció. La lectura com a sistema de representants és la de l'observació~\ref{obs:seccio-representants}.
\end{prova}

\begin{resumcap}
\apartat{Idees centrals}
\begin{enumerate}
  \item Una relació és una part d'un producte, i es llegeix per elements generalitzats: \(x\mathrel{r_X}y\) vol dir \(\langle x,y\rangle\in_X r\). Les relacions clàssiques són el cas de l'estadi \(\mathbf1\).
  \item Reflexivitat, simetria i transitivitat es formulen amb aplicacions \(e\colon A\to R\), \(s\colon R\to R\) i \(c\colon R\times_AR\to R\), no amb condicions sobre punts: les condicions puntuals en són conseqüència, però no bastarien per definir-les.
  \item Tota aplicació genera una relació d'equivalència, el seu parell nucli, i tota relació genera una aplicació, el seu quocient, que és un coequalitzador i, per tant, un epimorfisme regular.
  \item L'axioma d'efectivitat tanca el cercle: el parell nucli del quocient recupera la relació. Amb ell, relacions d'equivalència, epimorfismes regulars amb domini \(A\) i particions d'\(A\) són tres descripcions del mateix objecte.
  \item L'axioma d'elecció diu que tot epimorfisme té secció. En conseqüència, tot epimorfisme és exhaustiu sobre elements generalitzats i tot quocient admet un sistema de representants; també en resulta que tot epimorfisme és regular, encara que el capítol~\ref{cap:imatges} ho demostrarà sense elecció.
  \item L'exhaustivitat sobre \emph{punts} no és una conseqüència de l'elecció: ja valia sense ella. L'elecció és el pas de tenir una preimatge per a cada punt a tenir-les totes alhora.
\end{enumerate}
\apartat{Errors típics} Definir una relació d'equivalència per les tres condicions sobre punts globals, que no donen les aplicacions internes; pensar que el parell nucli és mono perquè \(f\) ho és (ho és sempre, per la unicitat del pullback); creure que el quocient només existeix amb l'elecció (existeix pels coequalitzadors); suposar que el parell nucli del quocient recupera la relació \enquote{per definició}, quan és l'axioma d'efectivitat; confondre epimorfisme regular (coequalitzador d'algun parell) amb epimorfisme escindit (amb secció), que és més fort sense l'elecció; i atribuir a l'elecció l'equivalència entre epimorfisme i exhaustivitat sobre punts.
\apartat{Lectura recomanada} \cite[cap.~4 i apèndix~B]{lawvere-rosebrugh} per als coequalitzadors, els epimorfismes i l'axioma d'elecció en la forma \enquote{tot epimorfisme té secció}, i per a la discussió de les formulacions equivalents de l'elecció; \cite[article~V, sessions~28--29]{lawvere-schanuel} per a una presentació elemental de les seccions i l'elecció; \cite[§§3.4, 5.5 i 9.8]{awodey} per a coequalitzadors, parells nuclis i categories regulars i exactes, que és el marc categòric d'aquest capítol; \cite[pàg.~12]{bell2006} per a la formulació de l'elecció en termes d'epis i seccions i per què és correcta per als conjunts abstractes; i \cite[§§7, 15--16]{halmos} per a les relacions d'equivalència, les particions i l'elecció en la teoria clàssica.
\end{resumcap}

\chapter{Imatges i factoritzacions}\label{cap:imatges}

\begin{objectius}
\apartat{Prerequisits} Del capítol~\ref{cap:relacions}: parell nucli, quocient, epimorfisme regular i l'axioma d'efectivitat. Del capítol~\ref{cap:parts}: bimorfismes, imatge inversa, el graf d'una aplicació i l'equivalència entre injectivitat i monomorfisme. Del capítol~\ref{cap:valors}: l'exhaustivitat dels epimorfismes sobre punts. L'axioma d'elecció no s'usa en aquest capítol.
\apartat{Objectius} En acabar el capítol, el lector hauria de poder:
\begin{itemize}
  \item definir la imatge d'una aplicació com a part mínima per la qual es factoritza, construir-la com a quocient del parell nucli i demostrar-ne la unicitat;
  \item demostrar, sense elecció, que el factor monomòrfic de la factorització és un monomorfisme, i que tot epimorfisme és regular;
  \item demostrar que els epimorfismes són estables per canvi de base i explicar en quin sentit la teoria és regular;
  \item caracteritzar monomorfismes, epimorfismes i isomorfismes per la degeneració d'un dels dos factors de la factorització imatge;
  \item definir la imatge directa d'una part i demostrar l'adjunció amb la imatge inversa;
  \item compondre relacions i llegir la composició per elements generalitzats, amb el paper dels recobriments epimòrfics;
  \item identificar les aplicacions amb les relacions funcionals i totals, i la composició de relacions amb la d'aplicacions.
\end{itemize}
\end{objectius}

El capítol precedent va lligar les relacions d'equivalència amb els
epimorfismes regulars, i va mostrar que tot epimorfisme regular és el
quocient del seu parell nucli. Aquestes eines tenen una conseqüència
immediata que encara no hem explotat: tota aplicació, sigui quina sigui,
es descompon de manera essencialment única en una part exhaustiva seguida
d'una part injectiva. La part injectiva és la \emph{imatge} de
l'aplicació.

Aquest és un capítol de consolidació. No hi afegim cap axioma: tot el que
segueix es dedueix dels productes, els pullbacks i els coequalitzadors ja
disponibles, i de l'axioma d'efectivitat del
capítol~\ref{cap:relacions}. Tampoc no hi fem servir l'axioma d'elecció:
en el punt on seria més temptador ---la prova que el factor monomòrfic de
la factorització imatge és un monomorfisme--- n'hi ha prou amb la
separació global i amb l'exhaustivitat dels epimorfismes sobre punts
(observació~\ref{obs:factoritzacio-eleccio}). Una conseqüència
d'aquesta via és que tot epimorfisme resulta regular sense elecció, cosa
que al capítol anterior només havíem obtingut amb ella. El resultat és
una manera uniforme de parlar
de la imatge d'una aplicació i de la imatge directa d'una part, que
completa la imatge inversa introduïda al capítol~\ref{cap:parts} i precisa
en quin sentit la teoria obtinguda és regular. Com sempre, definirem les
nocions sobre elements generalitzats i en recuperarem la lectura clàssica
com a cas particular.

\section{La imatge d'una aplicació}\label{sec:imatge}

Sigui \(f\colon A\to B\) una aplicació. El seu parell nucli
\(\operatorname{Nuc}(f)\) (definició~\ref{def:kernel-pair}) és una relació
d'equivalència sobre \(A\) (proposició~\ref{prop:kernel-equivalencia}):
identifica els elements generalitzats de \(A\) que \(f\) envia al mateix
lloc. És natural, doncs, prendre'n el quocient i comparar-lo amb \(B\).

\begin{definicio}[Imatge d'una aplicació]\index{Imatge d'una aplicació}\label{def:imatge}
Sigui \(f\colon A\to B\). Una \emph{imatge} de \(f\) és una part
\(m\colon I\hookrightarrow B\) per la qual \(f\) es factoritza com
\[
f=m\circ e,
\qquad
e\colon A\twoheadrightarrow I,
\]
amb \(e\) un epimorfisme, i que és \emph{mínima} amb aquesta propietat: si
\(m'\colon I'\hookrightarrow B\) és una part i \(f\) es factoritza a
través de \(m'\), aleshores \(m\) es factoritza a través de \(m'\).
\end{definicio}

\begin{notacio}[La imatge]\label{not:imatge}
Escriurem \(\operatorname{Im}(f)\) per a la part \(m\colon
I\hookrightarrow B\) de la definició anterior, i
\(\operatorname{Imatge}(m,e,f)\) per al predicat, eliminable, que afirma
que \(m\) és una imatge de \(f\) amb aplicació exhaustiva associada
\(e\). Desplegat, aquest predicat quantifica el domini \(I\), l'aplicació
\(e\), i les condicions que \(m\) és un monomorfisme, que \(e\) és un
epimorfisme, que \(f=m\circ e\) i la minimalitat. Com és habitual, un cop
introduïdes les abreviacions les farem servir en les proves; d'ara
endavant ometrem aquest desplegament quan els tipus siguin inequívocs.
\end{notacio}

La definició demana la minimalitat entre totes les parts a través de les
quals passa \(f\). El resultat principal del capítol és que aquesta part
mínima existeix sempre i que es pot construir explícitament: és el
quocient de \(A\) pel parell nucli de \(f\).

\begin{teorema}[Existència i factorització]\label{teo:factoritzacio}
Tota aplicació \(f\colon A\to B\) admet una imatge. A més, existeix una
factorització
\[
A\overset{e}{\twoheadrightarrow}\operatorname{Im}(f)\overset{m}{\hookrightarrow}B,
\]
on \(e\) és un epimorfisme regular i \(m\) un monomorfisme, i aquesta
factorització és única llevat d'un únic isomorfisme compatible.
\end{teorema}

\begin{prova}
\emph{Construcció.} Prenem el parell nucli \(n=\operatorname{Nuc}(f)\),
que és una relació d'equivalència sobre \(A\)
(proposició~\ref{prop:kernel-equivalencia}), i el seu quocient
\(e\colon A\twoheadrightarrow A/n\)
(proposició~\ref{prop:quocient-existeix}), un epimorfisme regular
(proposició~\ref{prop:quocient-regular}). Per definició de parell nucli,
\(f\circ k_1=f\circ k_2\), on \(k_1,k_2\) són les projeccions de \(n\); és
a dir, \(f\) coiguala el parell \((k_1,k_2)\). Com que \(e\) és el
coequalitzador d'aquest parell, la seva propietat universal
(definició~\ref{def:coequalitzador}) proporciona una única aplicació
\(m\colon A/n\to B\) amb
\[
f=m\circ e.
\]
Anomenem \(I:=A/n\).

\emph{\(m\) és un monomorfisme.} Per l'observació~\ref{obs:injectiva-mono}
del capítol~\ref{cap:parts}, que usa la separació global, n'hi ha prou amb
veure que \(m\) és injectiva sobre els punts globals. Siguin
\(y,y'\colon\mathbf1\to I\) amb \(m\circ y=m\circ y'\). Com que \(e\) és un
epimorfisme, és exhaustiu sobre punts
(proposició~\ref{prop:epi-exhaustiu-punts}): existeixen
\(a,a'\colon\mathbf1\to A\) amb \(e\circ a=y\) i \(e\circ a'=y'\).
Aleshores
\[
f\circ a=m\circ e\circ a=m\circ y=m\circ y'=m\circ e\circ a'=f\circ a',
\]
de manera que el parell \((a,a')\) es factoritza pel parell nucli \(n\)
de \(f\): hi ha \(d\colon\mathbf1\to A\times_BA\) amb \(k_1\circ d=a\) i
\(k_2\circ d=a'\). Com que \(e\) coiguala \(k_1\) i \(k_2\),
\[
y=e\circ a=e\circ k_1\circ d=e\circ k_2\circ d=e\circ a'=y'.
\]
Per tant \(m\) és injectiva sobre punts i, doncs, un monomorfisme.

\emph{Minimalitat.} Suposem que \(f\) es factoritza també a través d'una
part \(m'\colon I'\hookrightarrow B\), diguem \(f=m'\circ g\) amb
\(g\colon A\to I'\). Aleshores \(m'\circ g\) coiguala el parell nucli de
\(f\): en efecte, \(m'\circ g\circ k_1=f\circ k_1=f\circ k_2=m'\circ
g\circ k_2\), i com que \(m'\) és un monomorfisme, \(g\circ k_1=g\circ
k_2\). Per la propietat universal del coequalitzador \(e\), \(g\) es
factoritza de manera única com \(g=t\circ e\) per a alguna
\(t\colon I\to I'\). Aleshores \(m'\circ t\circ e=m'\circ g=f=m\circ e\);
com que \(e\) és un epimorfisme, \(m'\circ t=m\). Això mostra que \(m\) es
factoritza a través de \(m'\), que és la minimalitat.

\emph{Unicitat.} Siguin \(m_1\colon I_1\hookrightarrow B\) i
\(m_2\colon I_2\hookrightarrow B\) dues imatges de \(f\), amb aplicacions
exhaustives \(e_1,e_2\). Per la minimalitat de cadascuna respecte de
l'altra, \(m_1\) es factoritza a través de \(m_2\) i recíprocament: hi ha
\(u\colon I_1\to I_2\) i \(v\colon I_2\to I_1\) amb \(m_2\circ u=m_1\) i
\(m_1\circ v=m_2\). Aleshores \(m_1\circ v\circ u=m_2\circ u=m_1\); com
que \(m_1\) és un monomorfisme, \(v\circ u=\id_{I_1}\), i simètricament
\(u\circ v=\id_{I_2}\). Per tant \(u\) és un isomorfisme, i és compatible
amb les factoritzacions perquè \(u\circ e_1=e_2\) se'n dedueix component
\(m_2\circ u\circ e_1=m_1\circ e_1=f=m_2\circ e_2\) i cancel·lant el
monomorfisme \(m_2\). La unicitat d'\(u\) és clara pel mateix argument de
cancel·lació.

\emph{Tota factorització en epimorfisme regular i monomorfisme és una
imatge.} Falta veure que la unicitat s'aplica a \emph{qualsevol}
factorització \(f=m''\circ e''\) amb \(e''\colon A\twoheadrightarrow I''\)
epimorfisme regular i \(m''\colon I''\hookrightarrow B\) monomorfisme, i no
només a les que són mínimes per hipòtesi. Primer, \(e''\) és el
coequalitzador del seu propi parell nucli: si \(e''\) és el coequalitzador
d'un parell \(f',g'\colon C\to A\), aleshores \(e''\circ f'=e''\circ g'\) i
el parell \((f',g')\) es factoritza pel parell nucli de \(e''\) (propietat
universal del pullback); per tant tota aplicació que coiguala el parell
nucli coiguala \((f',g')\) i es factoritza de manera única per \(e''\), i
\(e''\) mateix coiguala el seu parell nucli. Segon, com que \(m''\) és un
monomorfisme, \(e''\circ u=e''\circ v\) si i només si \(f\circ u=f\circ v\),
de manera que el pullback de \(f\) amb ella mateixa ho és també de \(e''\)
amb ella mateixa: \(e''\) és un coequalitzador de les projeccions
\(k_1,k_2\) del parell nucli de \(f\). L'argument de la minimalitat,
repetit amb \(e''\) en lloc de \(e\), mostra aleshores que \(m''\) es
factoritza per tota part per la qual es factoritza \(f\): \(m''\) és una
imatge de \(f\), i el paràgraf anterior dona l'únic isomorfisme compatible
amb \((m,e)\).
\end{prova}

\begin{observacio}[Sense elecció]\label{obs:factoritzacio-eleccio}
La construcció de \(e\) i \(m\), la minimalitat i la unicitat no depenen
de cap axioma més enllà dels coequalitzadors i l'efectivitat. El pas que
podria semblar necessitar l'elecció ---que \(m\) sigui un monomorfisme---
s'ha fet amb la separació global, reduint-lo als punts globals, i amb
l'exhaustivitat dels epimorfismes sobre punts, que és un teorema del
capítol~\ref{cap:valors}. Amb l'elecció l'argument és encara més curt:
una secció \(s\) de \(e\) dona directament, per a \(x,x'\colon X\to I\)
amb \(m\circ x=m\circ x'\), que \(s\circ x\) i \(s\circ x'\) estan
relacionats pel parell nucli i, doncs, \(x=e\circ s\circ x=e\circ s\circ
x'=x'\); però no cal. En les teories sense separació global (com les
dels conjunts variables del capítol~\ref{cap:models}) cap de les dues
vies no és disponible, i la prova passa per elements generalitzats i
recobriments epimòrfics; és el preu de no tenir prou punts.
\end{observacio}

\begin{corollari}[Tot epimorfisme és regular; sense elecció]\label{cor:epi-regular-sense-eleccio}
Tota aplicació \(f\colon A\to B\) que és un epimorfisme és un epimorfisme
regular: és el coequalitzador del seu parell nucli.
\end{corollari}

\begin{prova}
Factoritzem \(f=m\circ e\) pel teorema~\ref{teo:factoritzacio}. Com que
\(f\) és un epimorfisme i \(f=m\circ e\), també \(m\) ho és; i \(m\) és un
monomorfisme. Per tant \(m\) és un bimorfisme i, per la
proposició~\ref{prop:bimorfisme-isomorfisme} del capítol~\ref{cap:parts},
un isomorfisme. Així \(f=m\circ e\) és el compost d'un coequalitzador amb
un isomorfisme, i és un coequalitzador del mateix parell: el parell
nucli de \(f\), que coincideix amb el de \(e\) pel
lema~\ref{lema:nucli-invariant}.
\end{prova}

\begin{observacio}[Què afegeix l'elecció, revisat]\label{obs:epi-regular-revisat}
Al capítol~\ref{cap:relacions} vam obtenir que tot epimorfisme és regular
com a conseqüència de l'elecció (proposició~\ref{prop:epi-regular-eleccio}).
El corol·lari anterior mostra que l'elecció hi era supèrflua: la
factorització imatge i el teorema del bimorfisme basten. El que l'elecció
continua afegint, i que no se segueix d'aquí, és que tot epimorfisme
sigui \emph{escindit}. Tenim, doncs, tres nocions en ordre de força:
epimorfisme, que amb aquest corol·lari coincideix amb epimorfisme regular;
i epimorfisme escindit, estrictament més fort sense l'elecció i
equivalent als anteriors amb ella.
\end{observacio}

\begin{observacio}[Lectura clàssica]\label{obs:imatge-classica}
En el model clàssic, la imatge de \(f\colon A\to B\) és el subconjunt
\(\{\,f(a):a\in A\,\}\subseteq B\) dels valors efectivament assolits. La
factorització \(A\twoheadrightarrow\operatorname{Im}(f)\hookrightarrow B\)
és la descomposició habitual d'una funció en la sobrejecció sobre la seva
imatge seguida de la inclusió de la imatge dins del codomini. El parell
nucli identifica els elements amb la mateixa imatge, i el quocient és el
conjunt de valors assolits, cadascun comptat una sola vegada.
\end{observacio}

\begin{lema}[Estabilitat dels epimorfismes per canvi de base]\label{lema:estabilitat-epi-regular}
Si \(e\colon A\twoheadrightarrow C\) és un epimorfisme i \(g\colon D\to C\)
és una aplicació qualsevol, aleshores la projecció
\(e'\colon D\times_C A\to D\) del pullback de \(e\) al llarg de \(g\) és
també un epimorfisme (i, pel corol·lari~\ref{cor:epi-regular-sense-eleccio},
un epimorfisme regular). No s'usa l'elecció.
\end{lema}

\begin{prova}
Considerem el pullback
\[
\begin{aligned}
D\times_C A &\xrightarrow{\;\pi_A\;} A\\
{\scriptstyle e'}\big\downarrow\ \ &\qquad\big\downarrow{\scriptstyle e}\\
D &\xrightarrow{\;\;g\;\;} C.
\end{aligned}
\]
Per la proposició~\ref{prop:epi-exhaustiu-global} del
capítol~\ref{cap:relacions}, n'hi ha prou amb veure que \(e'\) és
exhaustiva sobre punts. Sigui \(d\colon\mathbf1\to D\). El punt
\(g\circ d\colon\mathbf1\to C\) té una preimatge per \(e\), que és
exhaustiu sobre punts: \(a\colon\mathbf1\to A\) amb \(e\circ a=g\circ d\).
El parell \((d,a)\) és compatible, i la propietat universal del pullback
dona \(w\colon\mathbf1\to D\times_CA\) amb \(e'\circ w=d\). Per tant tot
punt de \(D\) té preimatge per \(e'\).

Amb l'elecció hi ha una prova que no passa pels punts: una secció \(s\)
de \(e\) dona la secció \(\langle\id_D,s\circ g\rangle\) de \(e'\), de
manera que \(e'\) és fins i tot escindit.
\end{prova}

\begin{observacio}[Una teoria regular]\label{obs:regularitat}
La factorització del teorema~\ref{teo:factoritzacio} és el que en la
terminologia categòrica es demana a una categoria \emph{regular}: tota
aplicació es descompon com un epimorfisme regular seguit d'un
monomorfisme, i, pel lema~\ref{lema:estabilitat-epi-regular}, els
epimorfismes regulars són estables per canvi de base. Amb l'axioma
d'efectivitat del capítol~\ref{cap:relacions} ja teníem les relacions
d'equivalència efectives; ara veiem que la factorització imatge n'és una
conseqüència i que, un cop demostrada l'estabilitat, la nostra teoria
disposa de la infraestructura d'una categoria regular~\cite{borceux-vol2},
i, com que tot epimorfisme és regular, d'una en què els epimorfismes i
els epimorfismes regulars coincideixen. Tot això s'ha obtingut sense
l'axioma d'elecció, però no sense la separació global, que és el que ha
permès reduir les proves als punts globals; la
secció~\ref{sec:imatge-directa} en complementa la lectura amb l'adjunció
entre imatge directa i imatge inversa.
\end{observacio}

\section{Caracteritzacions per la factorització}\label{sec:caracteritzacions}

La factorització imatge unifica en un sol enunciat les nocions de
monomorfisme, epimorfisme i isomorfisme: cadascuna correspon a una
degeneració d'un dels dos factors.

\begin{proposicio}[Monomorfismes i parell nucli trivial]\label{prop:mono-nucli-trivial}
Una aplicació \(f\colon A\to B\) és un monomorfisme si i només si el seu
parell nucli és la diagonal, és a dir,
\(\operatorname{Nuc}(f)\equiv_{A\times A}\Delta_A\). En aquest cas
l'aplicació exhaustiva de la factorització és un isomorfisme
\(A\xrightarrow{\sim}\operatorname{Im}(f)\).
\end{proposicio}

\begin{prova}
Si \(f\) és un monomorfisme, aleshores \(f\circ k_1=f\circ k_2\) implica
\(k_1=k_2\), de manera que el parell nucli és la diagonal
(proposició~\ref{prop:diagonal-mono}). Recíprocament, si
\(\operatorname{Nuc}(f)\equiv_{A\times A}\Delta_A\) i \(x,x'\colon X\to A\)
compleixen \(f\circ x=f\circ x'\), aleshores \(\langle x,x'\rangle\) es
factoritza a través del parell nucli, per tant a través de la diagonal, i
això dona \(x=x'\) (proposició~\ref{prop:diagonal-mono}); per tant \(f\) és
un monomorfisme. En aquest cas el quocient de \(A\) per la diagonal és el
mateix \(A\), de manera que \(e\colon A\to\operatorname{Im}(f)\) és un
isomorfisme.
\end{prova}

\begin{proposicio}[Epimorfismes i imatge total]\label{prop:epi-imatge-total}
Una aplicació \(f\colon A\to B\) és un epimorfisme (equivalentment, un
epimorfisme regular) si i només si la seva imatge és tot el codomini, és
a dir, si el monomorfisme \(m\colon\operatorname{Im}(f)\hookrightarrow B\)
és un isomorfisme.
\end{proposicio}

\begin{prova}
Si \(m\) és un isomorfisme, \(f=m\circ e\) és el compost d'un epimorfisme
regular amb un isomorfisme, i per tant un epimorfisme regular.
Recíprocament, si \(f\) és un epimorfisme, \(m\) és un epimorfisme (perquè
\(f=m\circ e\)) i un monomorfisme, i per tant un isomorfisme pel teorema
del bimorfisme (proposició~\ref{prop:bimorfisme-isomorfisme}); és
l'argument del corol·lari~\ref{cor:epi-regular-sense-eleccio}.
\end{prova}

\begin{corollari}[Bimorfisme i isomorfisme]\label{cor:bimorfisme-iso}
Una aplicació és un isomorfisme si i només si és alhora un monomorfisme i
un epimorfisme. Equivalentment, si i només si tots dos factors de
la seva factorització imatge són isomorfismes.
\end{corollari}

\begin{prova}
Si \(f\) és un isomorfisme, és tant mono com epi. Recíprocament,
si \(f\) és un monomorfisme, l'aplicació exhaustiva \(e\) és un
isomorfisme (proposició~\ref{prop:mono-nucli-trivial}); si a més és un
epimorfisme, el monomorfisme \(m\) també ho és
(proposició~\ref{prop:epi-imatge-total}). Aleshores \(f=m\circ e\) és
composició de dos isomorfismes, i és un isomorfisme. El resultat ja es va
demostrar al capítol~\ref{cap:parts} amb el classificador de parts; el que
oferim aquí és la seva lectura a través de la factorització imatge: cada
factor recull una de les dues hipòtesis.
\end{prova}

\section{Imatge directa d'una part}\label{sec:imatge-directa}

La factorització imatge s'aplica en particular a la composició d'una part
amb una aplicació, i produeix la noció d'\emph{imatge directa}, dual de la
imatge inversa del capítol~\ref{cap:parts}.

\begin{definicio}[Imatge directa]\index{Imatge directa}\label{def:imatge-directa}
Siguin \(f\colon A\to B\) una aplicació i \(s\colon S\hookrightarrow A\)
una part de \(A\). La \emph{imatge directa} de \(s\) per \(f\) és la
imatge de la composició \(f\circ s\colon S\to B\):
\[
f_{*}(s):=\operatorname{Im}(f\circ s)\;\colon\; f_{*}S\hookrightarrow B.
\]
\end{definicio}

\begin{notacio}[Imatge directa]\label{not:imatge-directa}
Escriurem \(f_{*}(s)\) o \(f_{*}S\) per a la part de \(B\) obtinguda, i
reservarem \(f^{*}\) per a la imatge inversa (definició~\ref{def:imatge-inversa}).
El predicat eliminable corresponent és
\(\operatorname{Imatge}(m,e,f\circ s)\) de la notació~\ref{not:imatge}
aplicat a \(f\circ s\).
\end{notacio}

\begin{proposicio}[Adjunció imatge directa\,--\,imatge inversa]\label{prop:adjuncio-imatges}
Siguin \(f\colon A\to B\), una part \(s\colon S\hookrightarrow A\) i una
part \(t\colon T\hookrightarrow B\). Aleshores
\[
f_{*}S\subseteq_B T
\qquad\text{si i només si}\qquad
S\subseteq_A f^{*}T.
\]
\end{proposicio}

\begin{prova}
Suposem \(f_{*}S\subseteq_B T\). Sobre elements generalitzats, un
\(x\colon X\to A\) es factoritza a través de \(s\) precisament quan
pertany a \(S\); en aquest cas \(f\circ x\) es factoritza a través de
\(f\circ s\) i, per la definició d'imatge, a través de \(f_{*}S\), i per
hipòtesi a través de \(t\). Per tant \(f\circ x\) pertany a \(T\), és a
dir, \(x\) pertany a la imatge inversa \(f^{*}T\)
(proposició~\ref{prop:pertinenca-imatge-inversa}); això és \(S\subseteq_A
f^{*}T\). Recíprocament, si \(S\subseteq_A f^{*}T\), tot element
generalitzat de \(S\) té imatge per \(f\) dins de \(T\); com que
\(f_{*}S\) és la part mínima de \(B\) a través de la qual passa \(f\circ
s\), i \(t\) és una part a través de la qual passa \(f\circ s\), la
minimalitat dona \(f_{*}S\subseteq_B T\).
\end{prova}

\begin{observacio}[Lectura clàssica]\label{obs:imatge-directa-classica}
En conjunts, la imatge directa és \(f_{*}(S)=\{\,f(a):a\in S\,\}\) i la
inversa és \(f^{*}(T)=\{\,a:f(a)\in T\,\}\). La
proposició~\ref{prop:adjuncio-imatges} és l'equivalència elemental
\(f(S)\subseteq T\iff S\subseteq f^{-1}(T)\), que expressa que prendre
imatge directa i prendre imatge inversa són operacions adjuntes entre les
parts de \(A\) i les de \(B\). Estructuralment, aquesta adjunció no
necessita cap axioma nou: només la factorització imatge i la propietat
universal de la imatge inversa.
\end{observacio}

\begin{corollari}[La imatge total és la imatge de \(f\)]\label{cor:estabilitat-exhaustiva}
La imatge directa de la part total de \(A\) és la imatge de \(f\):
\(f_{*}A=\operatorname{Im}(f)\). En conseqüència, \(f\) és un epimorfisme
si i només si \(f_{*}A\equiv_B\id_B\).
\end{corollari}

\begin{prova}
La part total de \(A\) és \(\id_A\), i \(f\circ\id_A=f\), de manera que
\(f_{*}A=\operatorname{Im}(f)\). La segona afirmació és la
proposició~\ref{prop:epi-imatge-total}.
\end{prova}

\section{Composició de relacions}\label{sec:composicio-relacions}

Les eines d'aquest capítol permeten compondre relacions, i no només
aplicacions. Una relació de \(A\) cap a \(B\) i una relació de \(B\) cap a
\(C\) es componen prenent la imatge de la relació producte a través de la
projecció adequada. Recuperem així, estructuralment, la composició
clàssica de relacions.

\begin{definicio}[Composició de relacions]\index{Composició de relacions}\label{def:composicio-relacions}
Siguin \(r\colon R\hookrightarrow A\times B\) una relació de \(A\) cap a
\(B\) i \(r'\colon R'\hookrightarrow B\times C\) una relació de \(B\) cap
a \(C\). El seu \emph{producte fibrat sobre \(B\)} és el pullback de la
segona projecció de \(r\) amb la primera de \(r'\):
\[
R\times_B R'\hookrightarrow R\times R',
\]
que recull els parells de testimonis \((\rho,\rho')\) amb el segon
component de \(\rho\) igual al primer de \(\rho'\); és a dir, el pullback
de \(r_2\colon R\to B\) amb \(r'_1\colon R'\to B\), amb projeccions
\(\pi\colon R\times_B R'\to R\) i \(\pi'\colon R\times_B R'\to R'\) que
compleixen \(r_2\circ\pi=r'_1\circ\pi'\). La \emph{composició}
\(r'\circ r\) és la imatge de l'aplicació
\[
g:=\langle r_1\circ\pi,\;r'_2\circ\pi'\rangle\colon
R\times_B R'\longrightarrow A\times C,
\]
que a cada testimoni li assigna el primer component de \(\rho\) i el segon
de \(\rho'\).
\end{definicio}

Abans de llegir la composició per elements generalitzats necessitem un fet
que s'usa sovint quan es treballa amb recobriments: la pertinença a una
part es pot comprovar després d'un epimorfisme.

\begin{lema}[Pertinença després d'un recobriment]\label{lema:pertinenca-recobriment}
Siguin \(m\colon S\hookrightarrow B\) una part, \(y\colon X\to B\) un
element generalitzat i \(p\colon X'\twoheadrightarrow X\) un epimorfisme.
Si \(y\circ p\in_{X'}m\), aleshores \(y\in_Xm\).
\end{lema}

\begin{prova}
Sigui \(m'\colon P\hookrightarrow X\) la imatge inversa \(y^{-1}(m)\), amb
l'altra projecció \(q\colon P\to S\), de manera que \(y\circ m'=m\circ q\).
Per la proposició~\ref{prop:pertinenca-imatge-inversa}, la hipòtesi
\(y\circ p\in_{X'}m\) equival a \(p\in_{X'}m'\): existeix
\(w\colon X'\to P\) amb \(m'\circ w=p\). Com que \(p\) és un epimorfisme i
es factoritza per \(m'\), també \(m'\) és un epimorfisme; i és un
monomorfisme. Pel teorema del bimorfisme
(proposició~\ref{prop:bimorfisme-isomorfisme}), \(m'\) és un isomorfisme,
de manera que \(\id_X\in_Xm'\) i, de nou per la
proposició~\ref{prop:pertinenca-imatge-inversa}, \(y\in_Xm\).
\end{prova}

\begin{proposicio}[La composició és una relació]\label{prop:composicio-es-relacio}
Amb les notacions anteriors, \(r'\circ r\) és una relació de \(A\) cap a
\(C\). Sobre elements generalitzats, \(x\mathrel{(r'\circ r)_X}z\) es
compleix si i només si existeix un recobriment epimòrfic
\(p\colon X'\twoheadrightarrow X\) i un element \(y\colon X'\to B\) amb
\[
(x\circ p)\mathrel{r_{X'}}y
\qquad\text{i}\qquad
y\mathrel{r'_{X'}}(z\circ p).
\]
Amb l'axioma d'elecció, el recobriment \(p\) és escindit, de manera que el
testimoni es pot recollir ja en el mateix estadi \(X\).
\end{proposicio}

\begin{prova}
Que \(r'\circ r\) és una part de \(A\times C\) és immediat de la
definició d'imatge, que produeix sempre un monomorfisme cap al codomini
\(A\times C\) (teorema~\ref{teo:factoritzacio}). Provem la caracterització
sobre elements generalitzats construint explícitament el testimoni.

Factoritzem \(g=m\circ e\) amb \(e\colon R\times_B R'\twoheadrightarrow
\operatorname{Im}(g)\) epimorfisme regular i
\(m\colon\operatorname{Im}(g)\hookrightarrow A\times C\) monomorfisme, de
manera que \(r'\circ r=m\).

\emph{(\(\Leftarrow\)) De les components al testimoni de la composició.}
Suposem que hi ha un recobriment \(p\colon X'\twoheadrightarrow X\) i un
\(y\colon X'\to B\) amb \((x\circ p)\mathrel{r_{X'}}y\) i
\(y\mathrel{r'_{X'}}(z\circ p)\). Els dos testimonis són fletxes
\(\rho\colon X'\to R\) i \(\rho'\colon X'\to R'\) amb
\(r_1\circ\rho=x\circ p\), \(r_2\circ\rho=y\), \(r'_1\circ\rho'=y\) i
\(r'_2\circ\rho'=z\circ p\). Com que \(r_2\circ\rho=y=r'_1\circ\rho'\), la
propietat universal del pullback \(R\times_B R'\) (definició~\ref{def:pullback})
proporciona una única \(w\colon X'\to R\times_B R'\) amb \(\pi\circ w=\rho\)
i \(\pi'\circ w=\rho'\). Aleshores
\[
g\circ w=\langle r_1\circ\pi\circ w,\;r'_2\circ\pi'\circ w\rangle
=\langle r_1\circ\rho,\;r'_2\circ\rho'\rangle
=\langle x\circ p,\;z\circ p\rangle=\langle x,z\rangle\circ p,
\]
de manera que \(\langle x,z\rangle\circ p\) es factoritza a través de
\(g\), i per tant a través de \(m=r'\circ r\). Com que \(p\) és un
epimorfisme, el lema~\ref{lema:pertinenca-recobriment} dona que
\(\langle x,z\rangle\) pertany a la part \(m\), és a dir,
\(x\mathrel{(r'\circ r)_X}z\).

\emph{(\(\Rightarrow\)) Del testimoni de la composició a les components.}
Suposem \(x\mathrel{(r'\circ r)_X}z\), és a dir, que
\(\langle x,z\rangle\colon X\to A\times C\) es factoritza a través de
\(m\). Sense elecció, \(e\) no té per què ser exhaustiu sobre elements
generalitzats (proposició~\ref{prop:exhaustiva-generalitzada-escindit});
el que sí que val és que l'element
generalitzat \(\bar x\colon X\to\operatorname{Im}(g)\) que testimonia
aquesta factorització ---l'únic amb \(m\circ\bar x=\langle x,z\rangle\),
únic perquè \(m\) és un monomorfisme--- es realitza, sobre algun
recobriment epimòrfic \(p\colon X'\twoheadrightarrow X\), com
la imatge per \(e\) d'un element generalitzat
\(w\colon X'\to R\times_B R'\): basta prendre el pullback de \(e\) al
llarg de \(\bar x\), amb projeccions \(p\) ---un epimorfisme pel
lema~\ref{lema:estabilitat-epi-regular}--- i \(w\). És a dir,
\(e\circ w=\bar x\circ p\), d'on
\(g\circ w=m\circ e\circ w=m\circ\bar x\circ p=\langle x,z\rangle\circ
p\). Posem \(\rho:=\pi\circ w\),
\(\rho':=\pi'\circ w\) i \(y:=r_2\circ\rho\). La commutativitat del
pullback dona \(r'_1\circ\rho'=r'_1\circ\pi'\circ w=r_2\circ\pi\circ
w=r_2\circ\rho=y\). D'altra banda, la definició de \(g\) dona
\(r_1\circ\rho=x\circ p\) i \(r'_2\circ\rho'=z\circ p\). Per tant \(\rho\)
testimonia \((x\circ p)\mathrel{r_{X'}}y\) i \(\rho'\) testimonia
\(y\mathrel{r'_{X'}}(z\circ p)\), com volíem.

En tots dos casos, amb l'axioma d'elecció (capítol~\ref{cap:relacions}) tot
epimorfisme és escindit; prenent una secció de \(p\) es transporta el
testimoni al mateix estadi \(X\) i el recobriment es pot ometre.
\end{prova}

\begin{observacio}[Lectura clàssica]\label{obs:composicio-relacions-classica}
En conjunts, la composició de relacions és la relació habitual
\[
r'\circ r=\{\,(a,c):\exists b\,\bigl((a,b)\in r\text{ i }(b,c)\in
r'\bigr)\,\}.
\]
El quantificador existencial sobre \(b\) és exactament el que la imatge
introdueix: la projecció del producte fibrat recull tots els encadenaments
possibles, i la imatge els comprimeix en el conjunt de parells
\((a,c)\) efectivament connectats. En el llenguatge global, l'existència
del testimoni \(b\) per a cada parell pot requerir l'elecció; sobre
elements generalitzats, en canvi, el testimoni existeix sempre després
d'un canvi d'estadi.
\end{observacio}

\begin{definicio}[Relació funcional, total, i funcional i total]\index{Relació funcional, total, i funcional i total}\label{def:relacio-funcional}
Sigui \(r\colon R\hookrightarrow A\times B\) una relació de \(A\) cap a
\(B\), amb projeccions \(r_1\colon R\to A\) i \(r_2\colon R\to B\). Diem
que \(r\) és:
\begin{enumerate}
\item \emph{funcional} si la primera projecció \(r_1\colon R\to A\) és un
monomorfisme (com a molt un testimoni per a cada element de \(A\));
\item \emph{total} si la primera projecció \(r_1\colon R\to A\) és un
epimorfisme regular (com a mínim un testimoni per a cada element de
\(A\));
\item \emph{funcional i total} si és alhora funcional i total.
\end{enumerate}
\end{definicio}

Les dues condicions són independents i, juntes, equivalen a demanar que
\(r_1\) sigui un isomorfisme: un monomorfisme que és també un epimorfisme
regular és un isomorfisme (corol·lari~\ref{cor:bimorfisme-iso}). Convé
mantenir-les separades, perquè les relacions funcionals que no són totals
---les aplicacions \emph{parcials}--- apareixen de manera natural, i la
correspondència amb les aplicacions (proposició~\ref{prop:aplicacio-funcional})
usa precisament la totalitat com a condició addicional sobre la
funcionalitat.

Cadascuna de les dues meitats té la seva lectura sobre elements
generalitzats. La funcionalitat és la unicitat del testimoni; la
totalitat, la seva existència. Ho recollim per separat i després en
combinem les conseqüències.

\begin{proposicio}[Lectura de la funcionalitat i la totalitat]\label{prop:funcional-lectura}
Sigui \(r\colon R\hookrightarrow A\times B\) una relació. Aleshores:
\begin{enumerate}
\item \(r\) és funcional si i només si per a tot element generalitzat
\(x\colon X\to A\) hi ha \emph{com a molt} un \(y\colon X\to B\) amb
\(x\mathrel{r_X}y\);
\item \(r\) és total si i només si per a tot element generalitzat
\(x\colon X\to A\) existeix \emph{almenys} un \(y\colon X\to B\) amb
\(x\mathrel{r_X}y\), possiblement després d'un canvi d'estadi;
\item \(r\) és funcional i total si i només si per a tot \(x\colon X\to A\)
existeix un \emph{únic} \(y\colon X\to B\) amb \(x\mathrel{r_X}y\), ja en
el mateix estadi \(X\).
\end{enumerate}
\end{proposicio}

\begin{prova}
Recordem que \(x\mathrel{r_X}y\) (definició~\ref{def:relacio}) significa que la parella
\(\langle x,y\rangle\colon X\to A\times B\) es factoritza a través de la
part \(r\), és a dir, que existeix \(\rho\colon X\to R\) amb
\(r\circ\rho=\langle x,y\rangle\); aquest \(\rho\), quan existeix, és el
\emph{testimoni} de la relació. Component amb les projeccions del
producte, això equival a \(r_1\circ\rho=x\) i \(r_2\circ\rho=y\).

\emph{Funcionalitat.} Suposem \(r_1\) mono. Si \(x\mathrel{r_X}y\) i
\(x\mathrel{r_X}y'\), tenim testimonis \(\rho,\rho'\colon X\to R\) amb
\(r_1\circ\rho=x=r_1\circ\rho'\); com que \(r_1\) és un monomorfisme,
\(\rho=\rho'\), i per tant \(y=r_2\circ\rho=r_2\circ\rho'=y'\). Hi ha,
doncs, com a molt un \(y\). Recíprocament, suposem la unicitat del \(y\) i
siguin \(\rho,\rho'\colon X\to R\) amb \(r_1\circ\rho=r_1\circ\rho'=:x\).
Posant \(y:=r_2\circ\rho\) i \(y':=r_2\circ\rho'\), tots dos són
testimonis de \(x\), d'on \(y=y'\); llavors \(\rho\) i \(\rho'\) tenen les
dues components iguals, i com que \(r=\langle r_1,r_2\rangle\) és un
monomorfisme, \(\rho=\rho'\). Això prova que \(r_1\) és un monomorfisme.

\emph{Totalitat.} Suposem \(r_1\) epimorfisme, i sigui \(x\colon X\to A\).
Prenem el pullback de \(r_1\) al llarg de \(x\), amb projeccions
\(u\colon X'\to X\) i \(\rho\colon X'\to R\), de manera que
\(r_1\circ\rho=x\circ u\); pel lema~\ref{lema:estabilitat-epi-regular},
\(u\) és un epimorfisme. Posant \(y:=r_2\circ\rho\colon X'\to B\), el
testimoni \(\rho\) dona \((x\circ u)\mathrel{r_{X'}}y\): aquesta és
l'existència \emph{a menys del canvi d'estadi} \(u\). Recíprocament, si
per a \(x=\id_A\) existeix un testimoni després d'un recobriment
\(u\colon X'\twoheadrightarrow A\), tenim \(\rho\colon X'\to R\) amb
\(r_1\circ\rho=u\); com que \(u\) és un epimorfisme que es factoritza per
\(r_1\), també \(r_1\) és un epimorfisme, i regular pel
corol·lari~\ref{cor:epi-regular-sense-eleccio}.

\emph{Funcional i total.} És la conjunció de les dues anteriors: la
totalitat dona l'existència d'un testimoni i la funcionalitat en dona la
unicitat, de manera que el \(y\) existeix i és únic. Quan totes dues es
compleixen, \(r_1\) és un isomorfisme
(corol·lari~\ref{cor:bimorfisme-iso}) i el testimoni és
\(\rho=r_1^{-1}\circ x\), definit ja en el mateix estadi \(X\) sense
passar a cap recobriment; llavors \(y=r_2\circ r_1^{-1}\circ x\).
\end{prova}

\begin{proposicio}[Aplicacions com a relacions funcionals]\label{prop:aplicacio-funcional}
Assignar a cada aplicació \(f\colon A\to B\) el seu graf
\(\langle\id_A,f\rangle\colon A\hookrightarrow A\times B\) estableix una
correspondència biunívoca, llevat d'equivalència de parts, entre les
aplicacions de \(A\) cap a \(B\) i les relacions funcionals de \(A\) cap a
\(B\) que a més són totals ---és a dir, funcionals i totals en el sentit
de la definició~\ref{def:relacio-funcional}, amb primera projecció alhora
monomorfisme i epimorfisme regular---.
Sota aquesta correspondència, la composició de relacions de la
definició~\ref{def:composicio-relacions} coincideix amb la composició
d'aplicacions.
\end{proposicio}

\begin{prova}
El graf d'una aplicació és un monomorfisme perquè la primera projecció
del producte el reparametritza com \(\id_A\)
(proposició~\ref{prop:graf-part}); de fet la seva primera projecció
\emph{és} \(\id_A\), un isomorfisme, de manera que el graf és funcional i
total en el sentit de la definició~\ref{def:relacio-funcional}, amb
\(y=f\circ x\) l'únic element relacionat amb cada \(x\).
Recíprocament, una relació funcional i total té primera projecció
\(r_1\colon R\to A\) alhora monomorfisme i epimorfisme, i per tant un
isomorfisme (corol·lari~\ref{cor:bimorfisme-iso}); doncs
\(f:=r_2\circ r_1^{-1}\colon A\to B\) és una aplicació el graf de la qual
és equivalent a \(r\). Aquestes dues assignacions són mútuament inverses
llevat d'equivalència de parts. Que la composició de relacions es
correspon amb la d'aplicacions es comprova directament. Siguin
\(f\colon A\to B\) i \(g\colon B\to C\), amb grafs
\(r=\langle\id_A,f\rangle\) i \(r'=\langle\id_B,g\rangle\), de manera que
\(r_1=\id_A\), \(r_2=f\), \(r'_1=\id_B\) i \(r'_2=g\). El pullback de
\(r_2=f\) amb \(r'_1=\id_B\) és \(A\) mateix, amb projeccions
\(\pi=\id_A\) i \(\pi'=f\): si \(u\colon X\to A\) i \(v\colon X\to B\)
compleixen \(f\circ u=v\), l'única \(w\colon X\to A\) amb \(\id_A\circ
w=u\) i \(f\circ w=v\) és \(w=u\). Aleshores l'aplicació de la
definició~\ref{def:composicio-relacions} és
\(\langle r_1\circ\pi,\,r'_2\circ\pi'\rangle=\langle\id_A,g\circ f\rangle\),
que ja és un monomorfisme (proposició~\ref{prop:graf-part}); la seva
imatge és ella mateixa (proposició~\ref{prop:mono-nucli-trivial}), és a
dir, \(r'\circ r\) és el graf de \(g\circ f\).
\end{prova}

\begin{resumcap}
\apartat{Idees centrals}
\begin{enumerate}
  \item La imatge d'una aplicació és la part mínima del codomini per la qual es factoritza, i existeix sempre: és el quocient del domini pel parell nucli. La factorització \(f=m\circ e\), epimorfisme seguit de monomorfisme, és única llevat d'un únic isomorfisme.
  \item Que \(m\) sigui un monomorfisme es demostra sobre punts, amb la separació global i l'exhaustivitat dels epimorfismes sobre punts: no cal l'elecció.
  \item Conseqüència: tot epimorfisme és regular (coequalitzador del seu parell nucli), sense elecció. L'elecció només afegeix que sigui escindit. Els epimorfismes són estables per canvi de base, i la teoria és regular.
  \item Mono, epi i iso es llegeixen a la factorització: \(f\) és mono si \(e\) és iso (parell nucli diagonal), epi si \(m\) és iso (imatge total), iso si ho són tots dos.
  \item La imatge directa \(f_*S=\operatorname{Im}(f\circ s)\) és adjunta a la imatge inversa: \(f_*S\subseteq T\iff S\subseteq f^*T\).
  \item Les relacions es componen prenent la imatge del producte fibrat; sobre elements generalitzats, el testimoni intermedi existeix després d'un recobriment epimòrfic. Les aplicacions són exactament les relacions funcionals i totals, i la composició de relacions estén la d'aplicacions.
\end{enumerate}
\apartat{Errors típics} Pensar que la imatge és \enquote{el conjunt dels valors} com si fos un conjunt a part (és una part del codomini, un monomorfisme); creure que cal l'elecció per a la factorització imatge o per a \enquote{tot epi és regular}, quan basta la separació global; oblidar que la minimalitat és part de la definició d'imatge, i no només la factorització; confondre \(f_*\) (imatge directa, que necessita la factorització) amb \(f^*\) (imatge inversa, un pullback); suposar que el testimoni \(y\) de la composició \(r'\circ r\) existeix en el mateix estadi \(X\) sense elecció; i identificar \enquote{relació funcional} amb \enquote{aplicació}, quan falta la totalitat.
\apartat{Lectura recomanada} \cite[§§4.3--4.5 i cap.~5]{lawvere-rosebrugh} per a la factorització imatge, les imatges directes i la lectura de les relacions com a parts de productes; \cite[article~IV, sessions~22--23]{lawvere-schanuel} per a la imatge d'una aplicació amb exemples elementals; \cite[§§5.3--5.4 i §§9.6--9.7]{awodey} per a les categories regulars, les factoritzacions i la seva relació amb els límits; \cite[cap.~2]{borceux-vol2} per a la teoria general de les categories regulars i exactes, que és el marc del capítol; i \cite[§§7--10]{halmos} per a les relacions, les funcions, les famílies i les inverses en la teoria clàssica.
\end{resumcap}

\chapter{Exponencials}\label{cap:exponencials}

\begin{objectius}
\apartat{Prerequisits} Del capítol~\ref{cap:parts}: productes, aparellament, producte d'aplicacions \(f\times g\) i la unicitat de les construccions universals; les aplicacions característiques i el classificador de parts. Del capítol~\ref{cap:valors}: la distributivitat del producte sobre la suma i l'estrictesa de \(\mathbf0\). Del capítol~\ref{cap:aplicacions}: elements generalitzats i separació global.
\apartat{Objectius} En acabar el capítol, el lector hauria de poder:
\begin{itemize}
  \item enunciar la propietat universal de l'exponencial amb l'avaluació i la transposició, i demostrar que currificar i descurrificar són inverses;
  \item demostrar la unicitat de l'exponencial llevat d'un únic isomorfisme compatible amb l'avaluació, pel patró habitual;
  \item explicar en quin sentit les aplicacions \(A\to B\) \enquote{són} els punts de \(B^A\), i per què un element generalitzat \(X\to B^A\) és una família d'aplicacions i no una sola;
  \item calcular \(A^{\mathbf1}\), \(A^{\mathbf0}\), \(\mathbf1^A\) i \(B^{A+C}\) identificant propietats universals, i definir \(P(A):=\mathbf2^A\);
  \item construir la composició i la identitat com a aplicacions internes, i saber quin axioma en garanteix la unicitat;
  \item llegir la naturalitat de la transposició com la coherència de les famílies d'aplicacions sota reindexació.
\end{itemize}
\end{objectius}

Fins ara les aplicacions han estat els individus primitius de la teoria:
tot el llenguatge parla d'aplicacions, i els conjunts no en són més que un
cas particular, les identitats. Aquest capítol fa el gir que dona nom al
llibre. Mostrarem que les aplicacions d'un conjunt \(A\) cap a un conjunt
\(B\) es poden reunir en un conjunt nou, l'\emph{exponencial} \(B^A\), els
punts del qual són precisament aquelles aplicacions. Les fletxes, que eren
la matèria primera, esdevenen així també punts d'un conjunt, i poden ser
manipulades com qualsevol altre element.

Aquesta reunió no se segueix dels axiomes anteriors: cal postular-la. El
capítol introdueix, doncs, un axioma nou ---l'existència d'exponencials---
i en treu les conseqüències immediates: la correspondència entre
aplicacions i punts globals de l'exponencial, el càlcul dels exponencials
dels conjunts \(\mathbf0\), \(\mathbf1\) i \(\mathbf2\), la representació
interna de la composició i la identitat, i la descripció de les famílies
d'aplicacions parametritzades per un estadi. Com sempre, formularem la
propietat universal sobre elements generalitzats i en recuperarem la
lectura clàssica com a cas particular.

\section{L'avaluació i la propietat universal}\label{sec:exponencial-universal}

La propietat que ha de caracteritzar l'exponencial és la mateixa que
regeix, en conjunts, el pas d'una funció de dues variables a una funció
que en retorna una altra. Si es té una regla que a cada parell
\((x,a)\) li assigna un element de \(B\), fixant la primera coordenada
s'obté, per a cada \(x\), una aplicació \(A\to B\); i recíprocament. El
conjunt \(B^A\) és el codomini del nom intern resultant \(x\mapsto f_x\),
és a dir, el conjunt de les aplicacions \(A\to B\), i l'avaluació és la
regla que en recupera el valor.

\begin{definicio}[Exponencial]\index{Exponencial}\label{def:exponencial}
Siguin \(A\) i \(B\) dos conjunts. Un \emph{exponencial} de \(B\) per
\(A\) és un conjunt \(E\), juntament amb una aplicació d'\emph{avaluació}
\[
\operatorname{ev}\colon E\times A\to B,
\]
que compleix la propietat universal següent: per a tot conjunt \(X\) i
tota aplicació \(h\colon X\times A\to B\), existeix una única aplicació
\[
\widehat h\colon X\to E
\qquad\text{tal que}\qquad
\operatorname{ev}\circ(\widehat h\times\id_A)=h.
\]
\end{definicio}

\begin{notacio}[L'exponencial i la transposició]\label{not:exponencial}
Escriurem \(B^A\) per a l'exponencial \(E\) de la definició anterior i
\(\operatorname{Exp}(E,\operatorname{ev},A,B)\) per al predicat,
eliminable, que l'afirma. En desplegar-lo cal recordar que ni el producte
ni les projeccions són termes primitius: s'han de quantificar amb el predicat
de producte \(\operatorname{Producte}\) (definició~\ref{def:producte}), que fixa
un objecte producte i les seves projeccions llevat d'isomorfisme. Reprenem
l'abreviació del capítol~\ref{cap:relacions}, que manté les projeccions
\emph{visibles} com a arguments del predicat:
\[
\operatorname{Prod}(P,\pi_1,\pi_2;X,Y):=\operatorname{Producte}(P,\pi_1,\pi_2,X,Y),
\]
que diu que \((P,\pi_1,\pi_2)\) és un producte de \(X\) i \(Y\). Cada producte que
hi intervé
---el de \(E\) amb \(A\) i, per a cada estadi de prova, el de \(X\) amb
\(A\)--- s'introdueix amb el seu propi bloc d'existencials, de manera que
totes les projeccions queden lligades dins l'abast on s'usen:
\[
\begin{aligned}
\operatorname{Exp}(E,\operatorname{ev},A,B)
\iff{}&
\operatorname{Conj}(E)
\\
&{}\land\exists P\,\exists\pi_1\,\exists\pi_2\,\bigl(
\operatorname{Prod}(P,\pi_1,\pi_2;E,A)\land\operatorname{ev}\colon P\to B\bigr)
\\
&{}\land\forall X\,\forall Q\,\forall q_1\,\forall q_2\,\forall h\,\Bigl(
\operatorname{Prod}(Q,q_1,q_2;X,A)\land h\colon Q\to B
\\
&\hphantom{{}\land\forall X\,\forall Q\,}
\Rightarrow
\exists!\,\widehat h\,\exists m\,\bigl(\widehat h\colon X\to E
\land m\colon Q\to P
\\
&\hphantom{{}\land\forall X\,\forall Q\,\Rightarrow\exists!\,\widehat h\,\exists m\,\bigl(}
{}\land\pi_1\circ m=\widehat h\circ q_1
\land\pi_2\circ m=q_2
\\
&\hphantom{{}\land\forall X\,\forall Q\,\Rightarrow\exists!\,\widehat h\,\exists m\,\bigl(}
{}\land\operatorname{ev}\circ m=h\bigr)\Bigr),
\end{aligned}
\]
on \(m\colon Q\to P\) és l'aplicació \(\widehat h\times\id_A\), definida
dins la fórmula per les dues equacions que la propietat universal de
\((P,\pi_1,\pi_2)\) associa a la parella \((\widehat h\circ q_1,q_2)\); com
que aquesta \(m\) és única un cop fixada \(\widehat h\), el quantificador
\(\exists m\) no afegeix cap indeterminació, i l'\(\exists!\) recau només
sobre \(\widehat h\). Ara les
projeccions \(q_1,q_2\) del producte variable \(Q\) queden lligades pels
seus propis quantificadors \(\forall q_1\,\forall q_2\), i no apareix cap
símbol lliure: el desplegament arriba fins al predicat de producte, i les
operacions que hi resten ---el producte \(\times\) d'aplicacions, la
composició \(\circ\), la parella \(\langle\cdot,\cdot\rangle\) i el
quantificador \(\exists!\)--- són convenis eliminables segons el
\hyperref[not:criteri-formal]{criteri d'escriptura formal} de la
introducció, que els tradueix en darrera instància a \(\Dom\), \(\Cod\) i
\(\Com\). A partir d'ara la usarem abreujada amb \(E\times A\), \(\id_A\) i
\(\times\), entès que cada aparició s'elimina pel predicat de producte i
pel criteri esmentat quan els tipus són inequívocs.
Anomenarem \(\widehat h\) la \emph{transposada} o \emph{currificació} de
\(h\), i escriurem \(\lambda(h):=\widehat h\). L'operació inversa, que a
una aplicació \(g\colon X\to B^A\) li associa
\(\operatorname{ev}\circ(g\times\id_A)\colon X\times A\to B\), l'anomenarem
\emph{descurrificació}.
\end{notacio}

L'existència d'exponencials no es dedueix dels axiomes anteriors. La
postulem.

\begin{axioma}[Existència d'exponencials]\label{ax:exponencials}
Tot parell de conjunts té un exponencial. Formalment,
\[
\forall A\,\forall B\,
\Bigl(
\operatorname{Conj}(A)\land\operatorname{Conj}(B)
\Rightarrow
\exists E\,\exists e\,
\operatorname{Exp}(E,e,A,B)
\Bigr),
\]
on \(e\) és una variable d'aplicació ordinària; un cop fixada, adoptem la
convenció \(\operatorname{ev}:=e\) per a l'avaluació de l'exponencial
triat.
\end{axioma}

\begin{proposicio}[Unicitat de l'exponencial]\label{prop:exponencial-unic}
L'exponencial d'un parell de conjunts és únic llevat d'un únic
isomorfisme compatible amb l'avaluació. És a dir, si \((E,\operatorname{ev})\)
i \((E',\operatorname{ev}')\) són tots dos exponencials de \(B\) per \(A\),
existeix un únic isomorfisme \(w\colon E\to E'\) amb
\(\operatorname{ev}'\circ(w\times\id_A)=\operatorname{ev}\).
\end{proposicio}

\begin{prova}
Apliquem la propietat universal de \((E',\operatorname{ev}')\) a
l'aplicació \(\operatorname{ev}\colon E\times A\to B\): hi ha una única
\(w\colon E\to E'\) amb \(\operatorname{ev}'\circ(w\times\id_A)=\operatorname{ev}\).
Simètricament, hi ha una única \(w'\colon E'\to E\) amb
\(\operatorname{ev}\circ(w'\times\id_A)=\operatorname{ev}'\). Component,
\(\operatorname{ev}\circ((w'\circ w)\times\id_A)=\operatorname{ev}\); però
també \(\operatorname{ev}\circ(\id_E\times\id_A)=\operatorname{ev}\), i la
unicitat de la transposada de \(\operatorname{ev}\) per la propietat
universal de \((E,\operatorname{ev})\) dona \(w'\circ w=\id_E\).
Simètricament \(w\circ w'=\id_{E'}\). Per tant \(w\) és un isomorfisme, i
és l'únic compatible amb l'avaluació per construcció.
\end{prova}

\begin{observacio}[Lectura clàssica]\label{obs:exponencial-classica}
En conjunts, \(B^A\) és el conjunt de totes les funcions de \(A\) en
\(B\), i l'avaluació és \(\operatorname{ev}(f,a)=f(a)\). Donada una funció
de dues variables \(h(x,a)\), la seva transposada \(\widehat h(x)\) és la
funció que a cada \(x\) li assigna la funció parcial
\(a\mapsto h(x,a)\); la condició
\(\operatorname{ev}\circ(\widehat h\times\id_A)=h\) diu simplement que
\(\widehat h(x)(a)=h(x,a)\). Aquest és el mecanisme habitual de fixar una
variable per obtenir una família de funcions de l'altra.
\end{observacio}

\section{Aplicacions com a punts}\label{sec:aplicacions-punts}

El cas \(X=\mathbf1\) de la propietat universal és el que justifica el
nom del capítol: identifica els punts globals de \(B^A\) amb les
aplicacions de \(A\) cap a \(B\).

\begin{teorema}[Punts globals de l'exponencial]\label{teo:punts-exponencial}
Amb l'axioma d'exponencials, per a cada parell de conjunts \(A\) i \(B\)
hi ha una correspondència biunívoca entre els punts globals
\(\mathbf1\to B^A\) i les aplicacions \(A\to B\). La correspondència
associa a un punt \(p\colon\mathbf1\to B^A\) l'aplicació
\(\operatorname{ev}\circ(p\times\id_A)\colon\mathbf1\times A\to B\), llegida
com una aplicació \(A\to B\) via l'isomorfisme
\(\mathbf1\times A\cong A\).
\end{teorema}

\begin{prova}
Fixem l'isomorfisme canònic \(\iota\colon A\xrightarrow{\sim}\mathbf1\times
A\) donat per la propietat universal del producte amb el terminal
(la primera projecció és l'única aplicació cap a \(\mathbf1\) i la segona
és \(\id_A\)). A un punt \(p\colon\mathbf1\to B^A\) li associem
\(\Phi(p):=\operatorname{ev}\circ(p\times\id_A)\circ\iota\colon A\to B\).
Recíprocament, a una aplicació \(f\colon A\to B\) li associem la
transposada de \(f\circ\iota^{-1}\colon\mathbf1\times A\to B\), és a dir,
\(\Psi(f):=\lambda(f\circ\iota^{-1})\colon\mathbf1\to B^A\). Que \(\Phi\) i
\(\Psi\) són mútuament inverses és exactament la unicitat de la transposada
de la definició~\ref{def:exponencial}: \(\Phi(\Psi(f))=f\) perquè
\(\operatorname{ev}\circ(\lambda(f\circ\iota^{-1})\times\id_A)=f\circ\iota^{-1}\),
i \(\Psi(\Phi(p))=p\) perquè tots dos tenen la mateixa descurrificació i la
transposada és única.
\end{prova}

\begin{observacio}[Lectura clàssica]\label{obs:punts-exponencial-classica}
El teorema diu que un element de \(B^A\) «és» una funció \(A\to B\), tal
com esperem del conjunt de funcions. La distinció que el llenguatge manté
---entre el \emph{punt} \(p\colon\mathbf1\to B^A\), que és una dada dins
d'un conjunt, i l'\emph{aplicació} \(f\colon A\to B\), que és una fletxa
del metallenguatge--- desapareix en la lectura clàssica, on totes dues coses
s'identifiquen. En el desenvolupament estructural convé, en canvi,
conservar-la: és el que permet tractar una aplicació com un objecte sense
confondre-la amb la fletxa que representa.
\end{observacio}

\begin{observacio}[Elements generalitzats de l'exponencial]\label{obs:generalitzats-exponencial}
Més enllà del cas global, la propietat universal identifica els elements
generalitzats \(X\to B^A\) amb les aplicacions \(X\times A\to B\). Un
element generalitzat de \(B^A\) a l'estadi \(X\) és, doncs, una família
d'aplicacions \(A\to B\) parametritzada per \(X\), i no una sola
aplicació. Aquesta lectura es desenvolupa a la
secció~\ref{sec:families}.
\end{observacio}

\section{\texorpdfstring{Els exponencials de \(\mathbf0\), \(\mathbf1\) i \(\mathbf2\)}{Els exponencials de 0, 1 i 2}}\label{sec:exponencials-basics}

Els conjunts distingits dels capítols anteriors donen exponencials que es
poden identificar explícitament, i que retornen nocions ja conegudes.

\begin{proposicio}[Exponencials amb base o exponent distingits]\label{prop:exponencials-basics}
Per a tot conjunt \(A\):
\begin{enumerate}
\item \(A^{\mathbf1}\cong A\): les aplicacions \(\mathbf1\to A\) són els
punts de \(A\).
\item \(A^{\mathbf0}\cong\mathbf1\): hi ha exactament una aplicació
\(\mathbf0\to A\).
\item \(\mathbf1^{A}\cong\mathbf1\): hi ha exactament una aplicació
\(A\to\mathbf1\).
\end{enumerate}
\end{proposicio}

\begin{prova}
\emph{(1)} Construïm les dues aplicacions inverses. D'una banda, sigui
\(\rho\colon A\times\mathbf1\xrightarrow{\sim}A\) l'isomorfisme
canònic de projecció. Definim \(\alpha\colon A\to A^{\mathbf1}\) com la
transposada de la projecció \(\rho\colon A\times\mathbf1\to A\); és a dir,
\(\alpha\) és l'única aplicació tal que
\(\operatorname{ev}\circ(\alpha\times\id_{\mathbf1})=\rho\). De l'altra, l'avaluació al punt únic de
\(\mathbf1\) dona \(\beta:=\operatorname{ev}\circ\langle\id,{!}\rangle
\colon A^{\mathbf1}\to A\), on \({!}\colon A^{\mathbf1}\to\mathbf1\) és l'única aplicació al
terminal. Comprovem les dues composicions. D'una banda,
\(\beta\circ\alpha=\operatorname{ev}\circ\langle\alpha,{!}_A\rangle
=\operatorname{ev}\circ(\alpha\times\id_{\mathbf1})\circ\langle\id_A,{!}_A\rangle
=\rho\circ\langle\id_A,{!}_A\rangle=\id_A\). De l'altra, \(\alpha\circ\beta\)
és una aplicació cap a l'exponencial, i basta comparar descurrificacions:
\(\operatorname{ev}\circ((\alpha\circ\beta)\times\id_{\mathbf1})
=\rho\circ(\beta\times\id_{\mathbf1})=\beta\circ p_1
=\operatorname{ev}\circ\langle p_1,{!}\rangle=\operatorname{ev}\),
perquè sobre \(A^{\mathbf1}\times\mathbf1\) la segona projecció és
l'única aplicació al terminal; per la unicitat de la transposada,
\(\alpha\circ\beta=\id_{A^{\mathbf1}}\). Per
tant \(A^{\mathbf1}\cong A\), i sota aquest isomorfisme els punts de
\(A^{\mathbf1}\) són els punts de \(A\).

\emph{(2)} Una aplicació \(X\times\mathbf0\to A\) surt del conjunt inicial:
\(X\times\mathbf0\cong\mathbf0\), perquè la projecció
\(X\times\mathbf0\to\mathbf0\) és un isomorfisme per l'estrictesa de
\(\mathbf0\) (proposició~\ref{prop:inicial-estricte}); de manera que n'hi ha exactament una; per la
propietat universal, hi ha exactament una aplicació \(X\to A^{\mathbf0}\)
per a cada \(X\), i això caracteritza \(A^{\mathbf0}\) com a terminal.

\emph{(3)} Una aplicació \(X\times A\to\mathbf1\) és única perquè
\(\mathbf1\) és terminal; per la propietat universal, hi ha exactament una
aplicació \(X\to\mathbf1^{A}\) per a cada \(X\), i \(\mathbf1^{A}\) és
terminal.
\end{prova}

Un darrer cas mereix un tractament a part. Pel
teorema~\ref{teo:punts-exponencial}, els punts de \(\mathbf2^{A}\) es
corresponen amb les aplicacions \(A\to\mathbf2\), que són les aplicacions
característiques del capítol~\ref{cap:parts}; pel classificador de parts,
aquestes es corresponen al seu torn amb les parts de \(A\) llevat
d'equivalència. L'exponencial \(\mathbf2^{A}\) és, doncs, l'objecte els
punts globals del qual classifiquen les parts de \(A\). Aquesta propietat
justifica definir-lo com el conjunt de parts.

\begin{definicio}[Conjunt de parts]\index{Conjunt de parts}\label{def:conjunt-parts}
Per a tot conjunt \(A\), definim el \emph{conjunt de parts} de \(A\) com
l'exponencial
\[
P(A):=\mathbf2^{A}.
\]
Els punts de \(P(A)\) es corresponen, pel classificador de parts del
capítol~\ref{cap:parts}, amb les parts de \(A\) llevat d'equivalència.
Amb aquesta definició, l'existència del conjunt de parts no requereix cap
axioma nou més enllà del d'exponencials: és el cas particular \(B=\mathbf2\).
\end{definicio}

\begin{observacio}[El que resta obert sobre \(P(A)\)]\label{obs:parts-obert}
Definir \(P(A)\) com a objecte no exhaureix el seu estudi. Queden per
desenvolupar la relació de \emph{pertinença interna}
\(\in_A\hookrightarrow A\times P(A)\), les operacions booleanes internes
sobre \(P(A)\) i l'ordre d'inclusió com a estructura sobre aquest conjunt.
Aquestes qüestions no s'aborden en aquest estat de la teoria.
\end{observacio}

\begin{observacio}[Lectura clàssica]\label{obs:exponencials-basics-classica}
Els quatre casos reprodueixen igualtats familiars de la combinatòria de
conjunts: \(A^1=A\), \(A^0=1\), \(1^A=1\) i \(2^A=\mathcal P(A)\). L'últim
és el que motiva la definició \(P(A):=\mathbf2^{A}\) d'aquest capítol, que
el següent desenvoluparà:
triar una part de \(A\) és el mateix que assignar a cada element un dels
dos valors de veritat, és a dir, donar una aplicació característica.
\end{observacio}

\section{L'exponencial d'una suma}\label{sec:exponencial-suma}

Un darrer isomorfisme, que farem servir en comptar aplicacions sobre
conjunts finits, tradueix una suma a l'exponent en un producte
d'exponencials. És la forma interna de la llei \(b^{a+c}=b^{a}\cdot
b^{c}\).

\begin{proposicio}[L'exponencial d'una suma]\label{prop:exponencial-suma}
Per a tots els conjunts \(A\), \(C\) i \(B\), hi ha un isomorfisme canònic
\[
B^{A+C}\cong B^{A}\times B^{C},
\]
on \(A+C\) és la suma de \(A\) i \(C\) (definició~\ref{def:suma-conjunts}).
\end{proposicio}

\begin{prova}
N'hi ha prou de comprovar que \(B^{A}\times B^{C}\) satisfà la propietat
universal de l'exponencial \(B^{A+C}\); la unicitat
(proposició~\ref{prop:exponencial-unic}) donarà l'isomorfisme canònic.
Per a tot conjunt \(X\), les aplicacions \(X\to B^{A+C}\) es corresponen,
per la propietat universal de l'exponencial, amb les aplicacions
\(X\times(A+C)\to B\). Ara bé, el producte distribueix sobre la suma
(teorema~\ref{teo:distributivitat}),
\[
X\times(A+C)\cong(X\times A)+(X\times C),
\]
de manera que, per la propietat universal de la suma
(definició~\ref{def:suma-conjunts}), una aplicació des d'aquest conjunt
cap a \(B\) és el mateix que una parella d'aplicacions \(X\times A\to B\) i
\(X\times C\to B\). Transposant de nou cada component, aquestes
corresponen a una aplicació \(X\to B^{A}\) i una aplicació \(X\to B^{C}\),
és a dir, a una única aplicació \(X\to B^{A}\times B^{C}\) per la propietat
universal del producte. Perquè l'argument sigui una comprovació de la
propietat universal i no només un recompte de correspondències, cal
exhibir l'avaluació: és l'aplicació
\[
\begin{aligned}
\operatorname{ev}'\colon(B^A\times B^C)\times(A+C)
&\;\cong\;(B^A\times B^C)\times A+(B^A\times B^C)\times C\\
&\xrightarrow{\ [\operatorname{ev}_A\circ(p_1\times\id_A),\ \operatorname{ev}_C\circ(p_2\times\id_C)]\ }B,
\end{aligned}
\]
i la cadena de correspondències anterior diu precisament que tota
\(h\colon X\times(A+C)\to B\) és \(\operatorname{ev}'\circ(\widehat h\times\id_{A+C})\)
per a una única \(\widehat h=\langle\lambda(h_A),\lambda(h_C)\rangle\), on
\(h_A,h_C\) són les dues components de \(h\) sota la distributivitat. Les
correspondències són naturals en \(X\) perquè cada pas ho és
(proposició~\ref{prop:families} i unicitat de les definicions per casos), i
\((B^A\times B^C,\operatorname{ev}')\) és, doncs, un exponencial de \(B\) per \(A+C\).
\end{prova}

\begin{observacio}[Lectura clàssica]\label{obs:exponencial-suma-classica}
Una aplicació definida sobre una unió disjunta \(A+C\) és exactament una
aplicació sobre \(A\) juntament amb una aplicació sobre \(C\), triades
independentment; el nombre de maneres de fer-ho és el producte dels
nombres respectius. En el cas particular \(A=C\) es recupera
\(B^{A+A}\cong B^{A}\times B^{A}\); i, iterant, l'exponencial d'un ordinal
finit es descompon en un producte de còpies, que és el que farà el
recompte d'aplicacions del capítol~\ref{cap:finits}.
\end{observacio}

\section{Composició i identitat internes}\label{sec:composicio-interna}

L'avaluació permet expressar dins de la teoria operacions que fins ara
eren metateòriques. La composició d'aplicacions, en particular, es
realitza com una aplicació entre exponencials.

\begin{proposicio}[Composició interna]\label{prop:composicio-interna}
Siguin \(A\), \(B\) i \(C\) conjunts. Existeix una única aplicació
\[
c\colon C^{B}\times B^{A}\to C^{A}
\]
tal que, sobre punts globals, \(c\) envia el parell \((g,f)\) ---amb
\(f\colon A\to B\) i \(g\colon B\to C\)--- al punt corresponent a
\(g\circ f\colon A\to C\). Anàlogament, existeix una única
\(\jmath\colon\mathbf1\to A^{A}\) que correspon a \(\id_A\).
\end{proposicio}

\begin{prova}
Per definir \(c\) n'hi ha prou de donar, per la propietat universal de
\(C^{A}\), una aplicació \((C^{B}\times B^{A})\times A\to C\). La prenem
com la composició
\[
(C^{B}\times B^{A})\times A
\longrightarrow
C^{B}\times(B^{A}\times A)
\xrightarrow{\id\times\operatorname{ev}}
C^{B}\times B
\xrightarrow{\operatorname{ev}}
C,
\]
on la primera fletxa és l'associativitat del producte i les altres dues
són avaluacions. La transposada d'aquesta aplicació és la \(c\) buscada.
Sobre punts globals, seguint les avaluacions, \((g,f)\) es transforma en
l'aplicació \(a\mapsto g(f(a))\), és a dir, \(g\circ f\): si \(\bar g\) i
\(\bar f\) són els punts que representen \(g\) i \(f\), aleshores
\(\operatorname{ev}\circ((c\circ\langle\bar g,\bar f\rangle)\times\id_A)\)
és \(\operatorname{ev}\circ(\bar g\times\id_B)\circ\iota_B\circ\operatorname{ev}\circ(\bar f\times\id_A)\),
on \(\iota_B\colon B\xrightarrow{\sim}\mathbf1\times B\) és l'isomorfisme
canònic; llevat de l'isomorfisme \(\mathbf1\times A\cong A\), això és
\(g\circ f\).
La \emph{unicitat} de \(c\) amb aquesta propietat no ve de la transposada
---la condició només en fixa els valors sobre els punts--- sinó de la
separació global: dues aplicacions \(C^B\times B^A\to C^A\) que
coincideixen sobre tots els punts globals són iguals, i tot punt global
del producte \(C^B\times B^A\) és un parell \(\langle\bar g,\bar f\rangle\)
de punts, per la propietat universal del producte. Per a \(\jmath\), la
transposada de la projecció \(\mathbf1\times A\to A\) dona el punt que
correspon a \(\id_A\), únic perquè els punts de \(A^A\) són les aplicacions
\(A\to A\).
\end{prova}

\begin{observacio}[Lectura clàssica]\label{obs:composicio-interna-classica}
Internament, el conjunt d'aplicacions \(A^{A}\) amb la composició
\(c\colon A^{A}\times A^{A}\to A^{A}\) i l'element distingit
\(\jmath\colon\mathbf1\to A^{A}\) forma el que en el metallenguatge és el
monoide de les aplicacions de \(A\) en si mateix. El capítol no
desenvolupa aquesta estructura algebraica; només fa notar que la
composició, que fins ara era una operació sobre fletxes, s'ha tornat una
aplicació entre conjunts, susceptible de ser tractada com qualsevol
altra.
\end{observacio}

\section{Famílies d'aplicacions}\label{sec:families}

L'observació~\ref{obs:generalitzats-exponencial} identificava un element
generalitzat \(X\to B^A\) amb una aplicació \(X\times A\to B\). Aquesta
lectura és la noció d'\emph{família d'aplicacions parametritzada per un
estadi}, i és la forma en què els exponencials intervenen en les
construccions posteriors.

\begin{proposicio}[Famílies parametritzades]\label{prop:families}
Per a conjunts \(X\), \(A\) i \(B\), la transposició estableix una
correspondència biunívoca, natural en \(X\), entre:
\begin{enumerate}
\item les aplicacions \(X\times A\to B\);
\item les aplicacions \(X\to B^A\).
\end{enumerate}
Una aplicació del primer tipus es llegeix com una família
\((f_x)_{x}\) d'aplicacions \(A\to B\) indexada pels elements generalitzats
\(x\) de \(X\); la del segon tipus n'és el nom intern.
\end{proposicio}

\begin{prova}
La correspondència és la transposició de la definició~\ref{def:exponencial},
que és biunívoca per la unicitat de \(\widehat h\). La naturalitat en
\(X\) vol dir que, per a tota aplicació \(t\colon X'\to X\),
transposar i després precompondre amb \(t\) dona el mateix que
precompondre amb \(t\times\id_A\) i després transposar; això se segueix de
la unicitat de la transposada aplicada a
\(h\circ(t\times\id_A)\), les dues transposicions de la qual coincideixen.
La lectura com a família s'obté fixant, per a cada element generalitzat
\(x\colon X'\to X\), l'aplicació \(f_x:=\operatorname{ev}\circ((\widehat
h\circ x)\times\id_A)\circ\iota\).
\end{prova}

\begin{observacio}[Lectura clàssica]\label{obs:families-classica}
En conjunts, una família d'aplicacions \((f_x)_{x\in X}\) de \(A\) en
\(B\) és el mateix que una funció de dues variables \(X\times A\to B\), o
que una funció \(X\to B^A\) que a cada índex li assigna un element del
conjunt de funcions. Les tres presentacions ---família indexada, funció de
dues variables i funció amb valors en un conjunt de funcions--- són la
mateixa dada vista de tres maneres. Els estadis del desenvolupament
estructural fan aquest paper d'índexs, ara sense sortir del llenguatge.
\end{observacio}

\begin{resumcap}
\apartat{Idees centrals}
\begin{enumerate}
  \item L'exponencial \(B^A\) es caracteritza per una propietat universal: tota \(h\colon X\times A\to B\) té una única transposada \(\widehat h\colon X\to B^A\) amb \(\operatorname{ev}\circ(\widehat h\times\id_A)=h\). Currificar i descurrificar són inverses, i l'exponencial és únic llevat d'un únic isomorfisme compatible amb l'avaluació.
  \item L'existència d'exponencials és un axioma nou; no se segueix dels anteriors.
  \item Els punts globals de \(B^A\) són les aplicacions \(A\to B\) (cas \(X=\mathbf1\)): les fletxes, individus primitius, es tornen punts d'un conjunt. Un element generalitzat \(X\to B^A\) és una família d'aplicacions parametritzada per \(X\), i la naturalitat de la transposició en garanteix la coherència sota reindexació.
  \item \(A^{\mathbf1}\cong A\), \(A^{\mathbf0}\cong\mathbf1\), \(\mathbf1^A\cong\mathbf1\), \(B^{A+C}\cong B^A\times B^C\): cada cas s'obté identificant una propietat universal. En particular \(P(A):=\mathbf2^A\) és el conjunt de parts, sense axioma nou.
  \item La composició i la identitat es realitzen com a aplicacions internes \(c\colon C^B\times B^A\to C^A\) i \(\jmath\colon\mathbf1\to A^A\); la seva unicitat és la separació global.
\end{enumerate}
\apartat{Errors típics} Confondre el punt \(p\colon\mathbf1\to B^A\) amb l'aplicació \(A\to B\) que representa (en la lectura clàssica s'identifiquen; en la teoria són individus diferents lligats per \(\operatorname{ev}\)); creure que un element generalitzat \(X\to B^A\) és una aplicació \(A\to B\), quan és una família; oblidar que \(\widehat h\) és única (la unicitat és el que fa funcionar totes les proves del capítol); pensar que \(B^A\) es construeix a partir dels axiomes anteriors; i deduir la unicitat de la composició interna de la transposada en lloc de la separació global.
\apartat{Lectura recomanada} \cite[cap.~5 i §§7.1--7.2]{lawvere-rosebrugh} per als conjunts d'aplicacions, l'avaluació, la transposició i les primeres conseqüències; \cite[article~V, sessions~30--31]{lawvere-schanuel} per a una introducció molt elemental als objectes d'aplicacions, amb els exemples de \(\mathbf2^A\) i de les lleis d'exponents; \cite[cap.~6]{awodey} per als exponencials i les categories cartesianament tancades en general; \cite[pàg.~11--12]{bell2006} per a la definició de l'exponencial per adjunció i per a \(2^X\) com a conjunt de parts; i \cite[§8]{halmos} per a les funcions i el conjunt de funcions en la teoria clàssica.
\end{resumcap}

\chapter{El conjunt de parts i la pertinença interna}\label{cap:parts-interna}

\begin{objectius}
\apartat{Prerequisits} Del capítol~\ref{cap:exponencials}: exponencials, avaluació, transposició i la definició \(P(A):=\mathbf2^A\), amb els punts de \(P(A)\) com a parts d'\(A\). Del capítol~\ref{cap:valors}: les operacions booleanes sobre \(\mathbf2\) i sobre les parts, i el complement. Del capítol~\ref{cap:imatges}: la imatge directa, la seva adjunció amb la imatge inversa, l'estabilitat dels epimorfismes i el lema de pertinença després d'un recobriment.
\apartat{Objectius} En acabar el capítol, el lector hauria de poder:
\begin{itemize}
  \item distingir la part \(m\), el seu nom \(\ulcorner m\urcorner\colon\mathbf1\to P(A)\) i el conjunt \(P(A)\), i construir la pertinença interna \(\in_A\) com la part classificada per l'avaluació;
  \item identificar les aplicacions \(X\to P(A)\) amb les parts de \(X\times A\), i llegir un element generalitzat de \(P(A)\) com una família de parts;
  \item definir les operacions booleanes internes per transposició, calcular-les sobre noms de parts i saber per a quines d'elles la lectura puntual val a tots els estadis;
  \item construir la imatge inversa interna \(P(f)\) i els quantificadors \(\exists_f\) i \(\forall_f\) com a aplicacions, i demostrar que \(\exists_f\dashv P(f)\dashv\forall_f\);
  \item demostrar que la imatge directa commuta amb el canvi de base, i usar-ho per llegir \(\exists_f\) sobre elements generalitzats;
  \item definir la inclusió interna sobre \(P(A)\) amb el quantificador universal i demostrar que és un ordre amb \(\cap\) i \(\cup\) com a ínfim i suprem.
\end{itemize}
\end{objectius}

El capítol anterior va definir el \emph{conjunt de parts} d'un conjunt
\(A\) com l'exponencial
\[
P(A):=\mathbf2^{A}
\]
(definició~\ref{def:conjunt-parts}), i en va identificar els punts globals
amb les aplicacions característiques \(A\to\mathbf2\), és a dir, amb les
parts de \(A\) llevat d'equivalència. Amb això, el conjunt de parts va
passar a existir com a \emph{objecte} de la teoria. Però definir un
objecte no és el mateix que estudiar-lo: quedaven obertes la pertinença
interna, les operacions booleanes sobre \(P(A)\) i l'ordre d'inclusió
(observació~\ref{obs:parts-obert}). Aquest capítol tanca aquestes
qüestions.

No hi introduirem cap axioma. Tot el que segueix és conseqüència de
l'existència d'exponencials del capítol~\ref{cap:exponencials}, del
classificador de parts del capítol~\ref{cap:parts} i de les operacions
booleanes sobre \(\mathbf2\) del capítol~\ref{cap:valors}. És, doncs, un
capítol de consolidació: fa interns, com a estructura sobre un únic
conjunt \(P(A)\), els fenòmens que fins ara vivien dispersos entre el
metallenguatge i el conjunt de valors de veritat.

Cal des del principi mantenir separades tres coses que la lectura clàssica
tendeix a confondre. Primer, una \emph{part} és una fletxa
\(m\colon S\hookrightarrow A\), un monomorfisme; és una dada del
metallenguatge. Segon, aquesta part té un únic \emph{punt classificador}
\[
\ulcorner m\urcorner\colon\mathbf1\to P(A),
\]
el nom intern de la seva aplicació característica \(\chi_m\colon A\to\mathbf2\);
és un element global d'un conjunt. Tercer, \(P(A)\) és el \emph{conjunt}
que reuneix tots aquests noms. La distinció entre la part, el seu nom i el
conjunt que els conté és el fil conductor del capítol.

\section{La relació de pertinença interna}\label{sec:pertinenca-interna}

L'avaluació de l'exponencial \(\mathbf2^{A}\) és una aplicació
\[
\operatorname{ev}\colon P(A)\times A\to\mathbf2,
\]
és a dir, una aplicació cap al conjunt de valors de veritat. Pel
classificador de parts (axioma~\ref{ax:representacio}), tota aplicació cap
a \(\mathbf2\) és la característica d'una part del seu domini. La part que
classifica l'avaluació és la relació de pertinença que buscàvem.

\begin{definicio}[Pertinença interna]\index{Pertinença interna}\label{def:pertinenca-interna}
Sigui \(A\) un conjunt. La \emph{relació de pertinença interna} de \(A\)
és la part de \(A\times P(A)\)
\[
{\in_A}\;\colon\;{\in_A}\hookrightarrow A\times P(A)
\]
classificada per l'avaluació, llegida sobre \(A\times P(A)\) via
l'isomorfisme de commutació \(\sigma\colon P(A)\times A\to A\times P(A)\).
És a dir, \(\in_A\) és la part de \(A\times P(A)\) la característica de la
qual és
\[
\chi_{\in_A}=\operatorname{ev}\circ\sigma^{-1}\colon A\times P(A)\to\mathbf2 .
\]
\end{definicio}

\begin{notacio}[Pertinença interna]\label{not:pertinenca-interna}
Escriurem \(\in_A\) tant per a la part com per al seu domini, seguint la
convenció de la relació inversa i la diagonal del
capítol~\ref{cap:relacions}. El predicat eliminable corresponent és
\(\operatorname{Car}(\operatorname{ev}\circ\sigma^{-1},{\in_A},A\times
P(A))\) de la definició~\ref{def:caracteristica}, amb \(\sigma\)
introduïda pel predicat de producte \(\operatorname{Prod}\). Per a un
parell d'elements generalitzats \(a\colon X\to A\) i \(\alpha\colon X\to
P(A)\), escriurem
\[
a\mathrel{\varepsilon}\alpha
\quad\text{en lloc de}\quad
\langle a,\alpha\rangle\in_X{\in_A},
\]
i llegirem «\(a\) pertany internament a \(\alpha\)».
\end{notacio}

La definició no fa res més que aplicar el classificador a una aplicació ja
existent. El seu contingut apareix quan la traduïm a la pertinença
generalitzada del capítol~\ref{cap:parts}.

\begin{proposicio}[Lectura de la pertinença interna]\label{prop:pertinenca-interna-lectura}
Siguin \(a\colon X\to A\) i \(\alpha\colon X\to P(A)\) dos elements
generalitzats a un mateix estadi \(X\). Aleshores
\[
a\mathrel{\varepsilon}\alpha
\quad\Longleftrightarrow\quad
\operatorname{ev}\circ\langle\alpha,a\rangle=1\circ{!}_X .
\]
En particular, si \(\alpha=\ulcorner m\urcorner\circ{!}_X\) és constant amb
valor el nom d'una part \(m\colon S\hookrightarrow A\), aleshores
\[
a\mathrel{\varepsilon}\alpha
\quad\Longleftrightarrow\quad
a\in_X m .
\]
\end{proposicio}

\begin{prova}
Per la propietat característica (notació~\ref{not:caracteristica}),
\(\langle a,\alpha\rangle\in_X{\in_A}\) equival a
\(\chi_{\in_A}\circ\langle a,\alpha\rangle=1\circ{!}_X\). Ara bé,
\(\chi_{\in_A}=\operatorname{ev}\circ\sigma^{-1}\) i
\(\sigma^{-1}\circ\langle a,\alpha\rangle=\langle\alpha,a\rangle\), d'on
la primera equivalència.

Per a la segona, sigui \(\alpha=\ulcorner m\urcorner\circ{!}_X\). El punt
classificador \(\ulcorner m\urcorner\) és, per definició, el nom intern de
\(\chi_m\), de manera que descurrificar-lo recupera \(\chi_m\): la
transposada de \(\operatorname{ev}\circ(\ulcorner m\urcorner\times\id_A)\)
és \(\ulcorner m\urcorner\), i per la unicitat de la transposada
(definició~\ref{def:exponencial}) es té
\(\operatorname{ev}\circ(\ulcorner m\urcorner\times\id_A)=\chi_m\circ\rho\),
on \(\rho\colon\mathbf1\times A\to A\) és la projecció. Component amb
\(\langle{!}_X,a\rangle\) i usant \(\alpha=\ulcorner m\urcorner\circ{!}_X\),
\[
\operatorname{ev}\circ\langle\alpha,a\rangle
=\chi_m\circ a .
\]
Així \(a\mathrel{\varepsilon}\alpha\) equival a
\(\chi_m\circ a=1\circ{!}_X\), que és exactament \(a\in_X m\).
\end{prova}

\begin{observacio}[La part, el seu nom i el conjunt]\label{obs:tres-nivells}
La proposició~\ref{prop:pertinenca-interna-lectura} és el lloc on es
paguen els dividends de mantenir separats els tres nivells. La
\emph{part} \(m\colon S\hookrightarrow A\) és una fletxa; el seu
\emph{nom} \(\ulcorner m\urcorner\colon\mathbf1\to P(A)\) és un punt del
conjunt \(P(A)\); i la pertinença \(a\in_X m\), que abans era una
factorització a través de \(m\), es converteix ara en la pertinença
interna \(a\mathrel{\varepsilon}\ulcorner m\urcorner\), una condició
expressada íntegrament amb aplicacions cap a \(P(A)\) i \(\mathbf2\). La
relació primitiva \(\in\) de la teoria clàssica, que allà és una noció no
definida entre individus, apareix aquí com una \emph{part construïda},
\(\in_A\hookrightarrow A\times P(A)\), classificada per l'avaluació.
\end{observacio}

\begin{observacio}[Lectura clàssica]\label{obs:pertinenca-interna-classica}
En conjunts, \(P(A)\) és el conjunt de subconjunts de \(A\), un punt
\(\alpha\in P(A)\) és un subconjunt concret, i
\(\operatorname{ev}(\alpha,a)\) val \(1\) exactament quan \(a\in\alpha\).
La part \(\in_A\subseteq A\times P(A)\) és, doncs, el graf de la relació
de pertinença: el conjunt de parells \((a,\alpha)\) amb \(a\in\alpha\). El
que en la presentació habitual és l'axioma que dona sentit a l'expressió
\(x\in y\), aquí és un teorema sobre una part determinada per l'avaluació.
\end{observacio}

\section{\texorpdfstring{Aplicacions cap a \(P(A)\) i parts de \(X\times A\)}{Aplicacions cap a P(A) i parts de X per A}}\label{sec:parts-parametritzades}

El teorema dels punts de l'exponencial (teorema~\ref{teo:punts-exponencial})
identificava els punts globals \(\mathbf1\to P(A)\) amb les parts de
\(A\). El cas general d'un estadi \(X\) qualsevol dona la versió
parametritzada: una aplicació \(X\to P(A)\) és una \emph{família de parts
de \(A\) indexada per \(X\)}, o equivalentment una part de \(X\times A\).

\begin{teorema}[Parts parametritzades]\label{teo:parts-parametritzades}
Per a conjunts \(X\) i \(A\), la transposició estableix una
correspondència biunívoca, natural en \(X\), entre:
\begin{enumerate}
\item les aplicacions \(\alpha\colon X\to P(A)\);
\item les aplicacions \(X\times A\to\mathbf2\);
\item les parts de \(X\times A\) llevat d'equivalència.
\end{enumerate}
La correspondència entre \((1)\) i \((2)\) és la descurrificació; la
correspondència entre \((2)\) i \((3)\) és el classificador de parts.
\end{teorema}

\begin{prova}
L'equivalència entre \((1)\) i \((2)\) és la propietat universal de
l'exponencial \(P(A)=\mathbf2^{A}\)
(definició~\ref{def:exponencial}) aplicada amb codomini \(\mathbf2\): a
cada \(\alpha\colon X\to P(A)\) li correspon
\(\operatorname{ev}\circ(\alpha\times\id_A)\colon X\times A\to\mathbf2\), i
la transposició és biunívoca per la unicitat de \(\widehat h\). La
naturalitat en \(X\) és la de la proposició~\ref{prop:families},
particularitzada a \(B=\mathbf2\).

L'equivalència entre \((2)\) i \((3)\) és l'axioma de representació
(axioma~\ref{ax:representacio}): tota aplicació \(X\times A\to\mathbf2\) és
la característica d'una única part de \(X\times A\) llevat d'equivalència,
i tota part de \(X\times A\) té una única característica. Componen les dues
correspondències.
\end{prova}

\begin{observacio}[Elements generalitzats de \(P(A)\)]\label{obs:generalitzats-parts}
El teorema diu que un element generalitzat de \(P(A)\) a l'estadi \(X\)
---una aplicació \(X\to P(A)\)--- no és una sola part de \(A\), sinó una
part de \(X\times A\), és a dir, una família de parts de \(A\)
parametritzada per \(X\). El cas global \(X=\mathbf1\) recupera una única
part de \(\mathbf1\times A\cong A\)
(teorema~\ref{teo:punts-exponencial}). Aquesta és la mateixa disciplina
que regeix tot el desenvolupament: les nocions es defineixen sobre
elements generalitzats, i el cas global n'és el particular \(X=\mathbf1\).
\end{observacio}

\begin{corollari}[La pertinença interna és universal]\label{cor:pertinenca-universal}
Sigui \(\alpha\colon X\to P(A)\) i sigui \(s_\alpha\hookrightarrow
X\times A\) la part que li correspon pel
teorema~\ref{teo:parts-parametritzades}. Aleshores, per a tot estadi
\(Y\) i tot parell \(a\colon Y\to A\), \(x\colon Y\to X\),
\[
a\mathrel{\varepsilon}(\alpha\circ x)
\quad\Longleftrightarrow\quad
\langle x,a\rangle\in_Y s_\alpha .
\]
En particular, la part \(\in_A\hookrightarrow A\times P(A)\) és la part
associada a la identitat \(\id_{P(A)}\colon P(A)\to P(A)\).
\end{corollari}

\begin{prova}
Per la proposició~\ref{prop:pertinenca-interna-lectura},
\(a\mathrel{\varepsilon}(\alpha\circ x)\) equival a
\(\operatorname{ev}\circ\langle\alpha\circ x,a\rangle=1\circ{!}_Y\). Ara
\(\langle\alpha\circ x,a\rangle=(\alpha\times\id_A)\circ\langle x,a\rangle\),
de manera que la condició és
\(\operatorname{ev}\circ(\alpha\times\id_A)\circ\langle x,a\rangle
=1\circ{!}_Y\); i \(\operatorname{ev}\circ(\alpha\times\id_A)\) és, per
construcció, la característica de \(s_\alpha\). Això és
\(\langle x,a\rangle\in_Y s_\alpha\). Prenent \(X=P(A)\) i
\(\alpha=\id_{P(A)}\), la part associada té característica
\(\operatorname{ev}\), que sobre \(A\times P(A)\) és la de \(\in_A\).
\end{prova}

\begin{observacio}[Lectura clàssica]\label{obs:parts-parametritzades-classica}
Les tres presentacions són les que un curs elemental usa
indistintament: una funció \(\alpha\colon X\to\mathcal P(A)\) que a cada
índex assigna un subconjunt; una relació \(R\subseteq X\times A\); i una
funció \((x,a)\mapsto\chi_R(x,a)\) amb valors de veritat. Que siguin la
mateixa dada és el que legitima escriure indistintament \(a\in\alpha_x\),
\((x,a)\in R\) o \(\chi_R(x,a)=1\). El teorema
dona a aquesta equivalència, habitualment implícita, un estatut
estructural.
\end{observacio}

\section{\texorpdfstring{Operacions booleanes internes sobre \(P(A)\)}{Operacions booleanes internes sobre P(A)}}\label{sec:booleanes-internes}

Al capítol~\ref{cap:valors} les operacions booleanes es van definir com a
aplicacions sobre el conjunt de valors de veritat: la conjunció i la
disjunció \(\wedge,\vee\colon\mathbf2\times\mathbf2\to\mathbf2\)
(definició~\ref{def:conjuncio-disjuncio}) i la negació
\(\neg\colon\mathbf2\to\mathbf2\) (definició~\ref{def:negacio}). Aquestes
operacions actuen sobre valors de veritat, un a un. Ara les elevarem a
operacions sobre \(P(A)\), on cada punt reuneix una part sencera. La clau
és que les operacions sobre \(\mathbf2\) actuen «puntualment» en cada
element de \(A\), i la transposició les converteix en operacions sobre el
conjunt de parts.

\begin{definicio}[Operacions booleanes internes]\index{Operacions booleanes internes}\label{def:booleanes-internes}
Sigui \(A\) un conjunt. Definim les aplicacions internes següents sobre
\(P(A)\), com a transposades de les operacions puntuals corresponents:
\begin{itemize}
\item la \emph{intersecció interna}
\(\cap_A\colon P(A)\times P(A)\to P(A)\), transposada de
\[
P(A)\times P(A)\times A
\xrightarrow{\;\langle\operatorname{ev}_1,\operatorname{ev}_2\rangle\;}
\mathbf2\times\mathbf2
\xrightarrow{\;\wedge\;}
\mathbf2 ;
\]
\item la \emph{unió interna}
\(\cup_A\colon P(A)\times P(A)\to P(A)\), transposada de la mateixa
aplicació amb \(\vee\) en lloc de \(\wedge\);
\item el \emph{complement intern}
\(\neg_A\colon P(A)\to P(A)\), transposada de
\(\neg\circ\operatorname{ev}\colon P(A)\times A\to\mathbf2\).
\end{itemize}
Aquí \(\operatorname{ev}_1\) i \(\operatorname{ev}_2\) denoten les
avaluacions de la primera i la segona còpia de \(P(A)\) sobre \(A\),
obtingudes component \(\operatorname{ev}\) amb les projeccions adequades.
Anàlogament, els \emph{punts distingits} \(\varnothing_A,\;A_A\colon
\mathbf1\to P(A)\) són els noms de les aplicacions constants
\(0\circ{!}_A\) i \(1\circ{!}_A\) (proposició~\ref{prop:caracteristica-total}),
és a dir, de la part buida i la part total del capítol~\ref{cap:valors}.
\end{definicio}

\begin{notacio}[Operacions internes]\label{not:booleanes-internes}
Per a elements generalitzats \(\alpha,\beta\colon X\to P(A)\) escriurem
\(\alpha\cap_A\beta\), \(\alpha\cup_A\beta\) i \(\neg_A\alpha\) per a les
composicions de \(\langle\alpha,\beta\rangle\) o \(\alpha\) amb les
aplicacions de la definició~\ref{def:booleanes-internes}. Quan el conjunt
\(A\) sigui inequívoc, ometrem el subíndex.
\end{notacio}

El sentit d'aquestes operacions és que calculen, punt a punt, la
pertinença interna al resultat. Aquesta és la propietat que les
caracteritza.

\begin{proposicio}[Càlcul de les operacions internes]\label{prop:booleanes-puntuals}
Siguin \(\alpha,\beta\colon X\to P(A)\) i \(a\colon X\to A\). Les
operacions internes es calculen component les avaluacions amb les
operacions sobre \(\mathbf2\):
\begin{align*}
\operatorname{ev}\circ\langle\alpha\cap\beta,a\rangle
&=\wedge\circ\langle\operatorname{ev}\circ\langle\alpha,a\rangle,
\operatorname{ev}\circ\langle\beta,a\rangle\rangle,\\
\operatorname{ev}\circ\langle\alpha\cup\beta,a\rangle
&=\vee\circ\langle\operatorname{ev}\circ\langle\alpha,a\rangle,
\operatorname{ev}\circ\langle\beta,a\rangle\rangle,\\
\operatorname{ev}\circ\langle\neg\alpha,a\rangle
&=\neg\circ\operatorname{ev}\circ\langle\alpha,a\rangle .
\end{align*}
En conseqüència, a tot estadi \(X\),
\[
a\mathrel{\varepsilon}(\alpha\cap\beta)
\quad\Longleftrightarrow\quad
a\mathrel{\varepsilon}\alpha\ \text{ i }\ a\mathrel{\varepsilon}\beta ;
\]
i sobre punts globals (\(X=\mathbf1\)), a més,
\(a\mathrel{\varepsilon}(\alpha\cup\beta)\) si i només si
\(a\mathrel{\varepsilon}\alpha\) o \(a\mathrel{\varepsilon}\beta\), i
\(a\mathrel{\varepsilon}(\neg\alpha)\) si i només si no
\(a\mathrel{\varepsilon}\alpha\). Finalment, sobre els noms de dues parts
\(m,n\) de \(A\), les operacions internes recuperen les operacions sobre
parts del capítol~\ref{cap:valors}:
\[
\ulcorner m\urcorner\cap\ulcorner n\urcorner=\ulcorner m\cap n\urcorner,
\qquad
\ulcorner m\urcorner\cup\ulcorner n\urcorner=\ulcorner m\cup n\urcorner,
\qquad
\neg\ulcorner m\urcorner=\ulcorner m^{c}\urcorner .
\]
\end{proposicio}

\begin{prova}
Les tres identitats són la descurrificació de la
definició~\ref{def:booleanes-internes}: per la propietat universal de
l'exponencial,
\(\operatorname{ev}\circ((\alpha\cap\beta)\times\id_A)
=\wedge\circ\langle\operatorname{ev}_1,\operatorname{ev}_2\rangle\circ(\langle\alpha,\beta\rangle\times\id_A)\),
i component amb \(\langle\id_X,a\rangle\) s'obté la primera; les altres
dues són idèntiques.

Per a la lectura de \(\cap\): \(\wedge\) és la característica de
\(\delta_{11}=\langle1,1\rangle\colon\mathbf1\hookrightarrow\mathbf2\times\mathbf2\)
(definició~\ref{def:conjuncio-disjuncio}), de manera que
\(\wedge\circ\langle u,v\rangle=1\circ{!}_X\) si i només si
\(\langle u,v\rangle\) es factoritza per \(\delta_{11}\), és a dir, si i
només si \(u=1\circ{!}_X\) i \(v=1\circ{!}_X\); amb
\(u=\operatorname{ev}\circ\langle\alpha,a\rangle\) i
\(v=\operatorname{ev}\circ\langle\beta,a\rangle\), la
proposició~\ref{prop:pertinenca-interna-lectura} dona l'equivalència, a
qualsevol estadi.

Per a \(\cup\) i \(\neg\) la lectura amb \enquote{o} i \enquote{no} només
val sobre punts globals: allà cada avaluació \(u\) i \(v\) és el punt
\(0\) o el punt \(1\) (corol·lari~\ref{cor:elements-de-2}), i les taules
de veritat (proposició~\ref{prop:taules-veritat}) decideixen. Sobre un estadi
general no val: per exemple, amb \(X=\mathbf2\), \(a\) tal que
\(\operatorname{ev}\circ\langle\alpha,a\rangle=\id_{\mathbf2}\) i
\(\operatorname{ev}\circ\langle\beta,a\rangle=\neg\), es té
\(a\mathrel{\varepsilon}(\alpha\cup\beta)\) sense que
\(a\mathrel{\varepsilon}\alpha\) ni \(a\mathrel{\varepsilon}\beta\).
És la mateixa frontera que separa el complement d'una part de la negació
puntual (observació~\ref{obs:complement-classic}).

Per a la darrera afirmació no cal cap lectura puntual: les parts
\(m\cap n\), \(m\cup n\) i \(m^c\) es van definir per les seves
característiques, \(\wedge\circ\langle\chi_m,\chi_n\rangle\),
\(\vee\circ\langle\chi_m,\chi_n\rangle\) i \(\neg\circ\chi_m\)
(proposició~\ref{prop:caracteristica-interseccio}, definició~\ref{def:operacions-parts}, proposició~\ref{prop:caracteristica-complement}), i les tres identitats
anteriors amb \(X=\mathbf1\), \(\alpha=\ulcorner m\urcorner\) i
\(\beta=\ulcorner n\urcorner\) diuen que la descurrificació de
\(\ulcorner m\urcorner\cap\ulcorner n\urcorner\) és exactament
\(\wedge\circ\langle\chi_m,\chi_n\rangle\), i anàlogament per a les
altres dues. Per la unicitat del nom (teorema~\ref{teo:punts-exponencial}),
els punts coincideixen.
\end{prova}

\begin{observacio}[Lectura clàssica]\label{obs:booleanes-internes-classica}
Les operacions internes són les operacions de conjunts habituals,
\(\alpha\cap\beta\), \(\alpha\cup\beta\) i \(A\setminus\alpha\), vistes ara
com a \emph{aplicacions} \(P(A)\times P(A)\to P(A)\) i \(P(A)\to P(A)\), i
no com a operacions definides element a element de manera externa. El
guany estructural és que l'àlgebra de subconjunts de \(A\), que en el
metallenguatge és una àlgebra booleana sobre les classes d'equivalència de
parts, es realitza ara com una estructura algebraica \emph{sobre un únic
conjunt} \(P(A)\), amb operacions que són aplicacions de la teoria.
\end{observacio}

\begin{observacio}[Les lleis booleanes es transporten]\label{obs:lleis-transportades}
Les lleis booleanes de les parts (teorema~\ref{teo:lleis-booleanes})
---commutativitat, associativitat, distributivitat, lleis de De~Morgan,
complementació--- valen també per a les operacions internes de
\(P(A)\). La raó és que cada llei és una igualtat entre aplicacions cap a
\(\mathbf2\) que ja es compleix puntualment sobre \(\mathbf2\)
(les taules de veritat), i la transposició, en ser biunívoca, transporta
la igualtat a la corresponent entre aplicacions cap a \(P(A)\). Aquest
capítol no reescriu les demostracions: n'hi ha prou de fer notar que la
unicitat de la transposada converteix cada identitat booleana puntual en
la identitat interna corresponent.
\end{observacio}

\section{Imatge inversa, imatge directa i quantificadors}\label{sec:quantificadors}

Una aplicació \(f\colon A\to B\) indueix operacions entre els conjunts de
parts \(P(A)\) i \(P(B)\). La imatge inversa del capítol~\ref{cap:parts} i
la imatge directa del capítol~\ref{cap:imatges}, que allà eren operacions
sobre parts individuals, es realitzen ara com a aplicacions internes; i
l'adjunció entre elles s'estén a una parella de quantificadors associats
a \(f\).

\begin{definicio}[Imatge inversa interna]\index{Imatge inversa interna}\label{def:imatge-inversa-interna}
Sigui \(f\colon A\to B\). La \emph{imatge inversa interna} és l'aplicació
\[
P(f)\colon P(B)\to P(A)
\]
transposada de
\(\operatorname{ev}\circ(\id_{P(B)}\times f)\colon P(B)\times A\to\mathbf2\).
Equivalentment, \(P(f)\) és l'única aplicació tal que, per a tot
\(\beta\colon X\to P(B)\) i tot \(a\colon X\to A\),
\[
a\mathrel{\varepsilon}\bigl(P(f)\circ\beta\bigr)
\quad\Longleftrightarrow\quad
(f\circ a)\mathrel{\varepsilon}\beta .
\]
\end{definicio}

\begin{proposicio}[La imatge inversa interna estén \(f^{-1}\)]\label{prop:imatge-inversa-interna}
Sobre els noms de parts, \(P(f)\) recupera la imatge inversa del
capítol~\ref{cap:parts}: per a tota part \(n\) de \(B\),
\[
P(f)\circ\ulcorner n\urcorner=\ulcorner f^{-1}(n)\urcorner .
\]
Més generalment, si \(\beta\colon X\to P(B)\) correspon a la part
\(s_\beta\) de \(X\times B\) (teorema~\ref{teo:parts-parametritzades}),
aleshores \(P(f)\circ\beta\) correspon a la part
\((\id_X\times f)^{-1}(s_\beta)\) de \(X\times A\).
\end{proposicio}

\begin{prova}
La descurrificació de \(P(f)\circ\beta\) és
\(\operatorname{ev}\circ(\id_{P(B)}\times f)\circ(\beta\times\id_A)
=\operatorname{ev}\circ(\beta\times\id_B)\circ(\id_X\times f)
=\chi_{s_\beta}\circ(\id_X\times f)\),
que és la característica de \((\id_X\times f)^{-1}(s_\beta)\)
(proposició~\ref{prop:caracteristica-imatge-inversa}). El cas dels noms és
\(X=\mathbf1\), amb \(\mathbf1\times A\cong A\): la descurrificació de
\(P(f)\circ\ulcorner n\urcorner\) és \(\chi_n\circ f=\chi_{f^{-1}(n)}\).
\end{prova}

\begin{proposicio}[La imatge inversa commuta amb el complement]\label{prop:imatge-inversa-complement}
Sigui \(f\colon A\to B\). Aleshores \(P(f)\circ\neg_B=\neg_A\circ P(f)\).
\end{proposicio}

\begin{prova}
Comparem descurrificacions, com a aplicacions \(P(B)\times A\to\mathbf2\).
D'una banda,
\(\operatorname{ev}\circ((P(f)\circ\neg_B)\times\id_A)
=\operatorname{ev}\circ(\id\times f)\circ(\neg_B\times\id_A)
=\operatorname{ev}\circ(\neg_B\times\id_B)\circ(\id\times f)
=\neg\circ\operatorname{ev}\circ(\id\times f)\).
De l'altra,
\(\operatorname{ev}\circ((\neg_A\circ P(f))\times\id_A)
=\neg\circ\operatorname{ev}\circ(P(f)\times\id_A)
=\neg\circ\operatorname{ev}\circ(\id\times f)\).
Per la unicitat de la transposada, les dues aplicacions coincideixen.
\end{prova}

Per construir el quantificador existencial com una \emph{aplicació}
\(P(A)\to P(B)\) necessitem la imatge directa del capítol~\ref{cap:imatges}
aplicada no a una part d'\(A\), sinó a la part universal \(\in_A\), i un
lema que diu que la imatge directa es comporta bé sota canvi de base. És
la propietat que fa que la lectura sobre elements generalitzats de
\(\exists_f\) sigui la que s'espera.

\begin{lema}[La imatge directa commuta amb el canvi de base]\label{lema:beck-chevalley}
Sigui
\[
\begin{aligned}
A' &\xrightarrow{\;h\;} A\\
{\scriptstyle f'}\big\downarrow\ \ &\qquad\big\downarrow{\scriptstyle f}\\
B' &\xrightarrow{\;h'\;} B
\end{aligned}
\]
un quadrat cartesià. Per a tota part \(s\colon S\hookrightarrow A\),
\[
(h')^{-1}\bigl(f_{*}s\bigr)\;\equiv_{B'}\;f'_{*}\bigl(h^{-1}(s)\bigr).
\]
\emph{Axiomes usats:} els del capítol~\ref{cap:imatges} (en particular
l'estabilitat dels epimorfismes, que usa la separació global); cap elecció.
\end{lema}

\begin{prova}
Factoritzem \(f\circ s=m\circ e\) amb \(e\colon S\twoheadrightarrow I\)
epimorfisme i \(m\colon I\hookrightarrow B\) la imatge \(f_*s\)
(teorema~\ref{teo:factoritzacio}). Fem el pullback de tota aquesta
factorització al llarg de \(h'\): obtenim \(I':=(h')^{-1}(m)\), una part
de \(B'\), i \(S':=B'\times_B S\) amb una aplicació \(e'\colon S'\to I'\)
que és el pullback de \(e\) al llarg de \(I'\to I\), i per tant un
epimorfisme (lema~\ref{lema:estabilitat-epi-regular}). D'altra banda,
com que el quadrat de l'enunciat és cartesià, \(S'=B'\times_BS\) coincideix
amb \(A'\times_AS\) (dos quadrats cartesians encaixats formen un quadrat
cartesià), que és el domini de \(h^{-1}(s)\); i la composició
\(f'\circ h^{-1}(s)\colon S'\to B'\) és \(I'\hookrightarrow B'\) precedida
de \(e'\). Tenim, doncs, una factorització de \(f'\circ h^{-1}(s)\) com a
epimorfisme seguit de monomorfisme, i per la unicitat de la factorització
imatge, \(f'_*(h^{-1}(s))\equiv I'=(h')^{-1}(f_*s)\).
\end{prova}

\begin{definicio}[Quantificadors associats a una aplicació]\index{Quantificadors associats a una aplicació}\label{def:quantificadors}
Sigui \(f\colon A\to B\). Sigui \(E_A\hookrightarrow P(A)\times A\) la
part classificada per l'avaluació (la pertinença interna amb els factors
intercanviats: \(E_A=\sigma^{-1}\circ{\in_A}\)). El \emph{quantificador
existencial} associat a \(f\) és l'aplicació
\[
\exists_f\colon P(A)\to P(B)
\]
transposada de la característica de la imatge directa
\((\id_{P(A)}\times f)_{*}E_A\hookrightarrow P(A)\times B\). El
\emph{quantificador universal} és
\[
\forall_f:=\neg_B\circ\exists_f\circ\neg_A\colon P(A)\to P(B).
\]
Escriurem \(\exists_A:=\exists_{{!}_A}\) i \(\forall_A:=\forall_{{!}_A}\)
per a l'única aplicació \({!}_A\colon A\to\mathbf1\); són aplicacions
\(P(A)\to P(\mathbf1)\cong\mathbf2\).
\end{definicio}

\begin{proposicio}[Lectura dels quantificadors]\label{prop:quantificadors-lectura}
Sigui \(f\colon A\to B\) i sigui \(\alpha\colon X\to P(A)\), amb part
associada \(s_\alpha\hookrightarrow X\times A\). Aleshores:
\begin{enumerate}
\item \(\exists_f\circ\alpha\) correspon a la part
\((\id_X\times f)_{*}(s_\alpha)\) de \(X\times B\), i
\(\forall_f\circ\alpha\) a la part
\(\bigl((\id_X\times f)_{*}(s_\alpha^{c})\bigr)^{c}\). En particular, sobre
noms, \(\exists_f\ulcorner m\urcorner=\ulcorner f_*m\urcorner\).
\item Per a \(b\colon X\to B\), \(b\mathrel{\varepsilon}(\exists_f\circ\alpha)\)
si i només si existeixen un recobriment epimòrfic
\(p\colon X'\twoheadrightarrow X\) i un \(a\colon X'\to A\) amb
\(f\circ a=b\circ p\) i \(a\mathrel{\varepsilon}(\alpha\circ p)\). Sobre
punts globals, \(b\mathrel{\varepsilon}(\exists_f\circ\alpha)\) si i només
si hi ha un punt \(a\) amb \(f\circ a=b\) i \(a\mathrel{\varepsilon}\alpha\).
\item Per a \(b\colon X\to B\), \(b\mathrel{\varepsilon}(\forall_f\circ\alpha)\)
si i només si, per a tot \(x\colon Y\to X\) i tot \(a\colon Y\to A\) amb
\(f\circ a=b\circ x\), es té \(a\mathrel{\varepsilon}(\alpha\circ x)\).
\end{enumerate}
Amb l'axioma d'elecció, el recobriment de~(2) és escindit i el testimoni
\(a\) es pot prendre ja sobre \(X\).
\end{proposicio}

\begin{prova}
\emph{(1)} La descurrificació de \(\exists_f\circ\alpha\) és
\(\chi_{(\id\times f)_*E_A}\circ(\alpha\times\id_B)\), la característica
de \((\alpha\times\id_B)^{-1}\bigl((\id_{P(A)}\times f)_*E_A\bigr)\). El
quadrat format per \(\alpha\times\id_A\colon X\times A\to P(A)\times A\),
\(\id_X\times f\), \(\id_{P(A)}\times f\) i \(\alpha\times\id_B\) és
cartesià (és el producte del quadrat cartesià trivial de vèrtex \(X\)
sobre \(P(A)\) amb el de \(f\) sobre ell mateix). Pel
lema~\ref{lema:beck-chevalley},
\((\alpha\times\id_B)^{-1}\bigl((\id\times f)_*E_A\bigr)
\equiv(\id_X\times f)_*\bigl((\alpha\times\id_A)^{-1}E_A\bigr)\), i
\((\alpha\times\id_A)^{-1}E_A\) té característica
\(\operatorname{ev}\circ(\alpha\times\id_A)\), que és la de \(s_\alpha\).
Per a \(\forall_f\), apliquem el que acabem de veure a \(\neg_A\circ\alpha\),
que correspon a \(s_\alpha^c\) (proposició~\ref{prop:booleanes-puntuals}),
i després \(\neg_B\). El cas dels noms és \(X=\mathbf1\).

\emph{(2)} Per~(1), \(b\mathrel{\varepsilon}(\exists_f\circ\alpha)\) vol dir
\(\langle\id_X,b\rangle\in_X(\id_X\times f)_*(s_\alpha)\). Si hi ha
\(p\) i \(a\) com a l'enunciat, aleshores
\(\langle p,a\rangle\in_{X'}s_\alpha\) (corol·lari~\ref{cor:pertinenca-universal})
i \((\id_X\times f)\circ\langle p,a\rangle=\langle p,f\circ a\rangle
=\langle\id_X,b\rangle\circ p\), de manera que
\(\langle\id_X,b\rangle\circ p\) pertany a la imatge directa, i pel
lema~\ref{lema:pertinenca-recobriment} també \(\langle\id_X,b\rangle\).
Recíprocament, si \(\langle\id_X,b\rangle\) pertany a
\((\id_X\times f)_*(s_\alpha)\), prenem el pullback de l'epimorfisme
\(s_\alpha\twoheadrightarrow\operatorname{Im}\) al llarg del testimoni;
dona un recobriment \(p\colon X'\twoheadrightarrow X\)
(lema~\ref{lema:estabilitat-epi-regular}) i un \(w\colon X'\to s_\alpha\),
i \(a:=\pi_A\circ s_\alpha\circ w\) compleix el que cal. Sobre punts
globals, un recobriment \(p\colon X'\twoheadrightarrow\mathbf1\) té un punt
(lema del buit i no-degeneració), i component-hi s'obté el testimoni
sobre \(\mathbf1\).

\emph{(3)} Usarem dos fets auxiliars. \emph{(a)} Per a una part
\(M\hookrightarrow Z\) i un element generalitzat \(z\colon Y\to Z\),
\(z\in_Y M^c\) si i només si la imatge inversa \(z^{-1}(M)\) és buida. En
efecte, \(\chi_{M^c}\circ z=\neg\circ\chi_M\circ z\)
(proposició~\ref{prop:caracteristica-complement}), que val \(1\circ{!}_Y\) si i
només si \(\chi_M\circ z=0\circ{!}_Y\) (componeu amb la involució \(\neg\)), és
a dir, si i només si \(\chi_{z^{-1}(M)}=\chi_{\bot_Y}\)
(proposicions~\ref{prop:caracteristica-imatge-inversa}
i~\ref{prop:caracteristica-buida}), que equival a \(z^{-1}(M)\equiv_Y\bot_Y\)
(proposició~\ref{prop:caracteristica-determina-part}). \emph{(b)} Si
\(p\colon T'\twoheadrightarrow T\) és un epimorfisme i \(T'\) és inicial,
\(T\) és inicial: \(0\circ{!}_T\circ p=1\circ{!}_T\circ p\) perquè \(T'\) és
inicial, d'on \(0\circ{!}_T=1\circ{!}_T\), i l'argument de la
proposició~\ref{prop:caracteristica-buida} dona una aplicació \(T\to\mathbf0\).

Escrivim \(K:=s_\alpha\hookrightarrow X\times A\) i
\(N:=(\id_X\times f)_*(K^c)\hookrightarrow X\times B\). Per~(1),
\(b\mathrel{\varepsilon}(\forall_f\circ\alpha)\) vol dir
\(\langle\id_X,b\rangle\in_X N^c\); i, pel
corol·lari~\ref{cor:pertinenca-universal},
\(a\mathrel{\varepsilon}(\alpha\circ x)\) vol dir \(\langle x,a\rangle\in_Y K\).

\(\Rightarrow\)) Suposem \(\langle\id_X,b\rangle\in_X N^c\), i siguin
\(x\colon Y\to X\), \(a\colon Y\to A\) amb \(f\circ a=b\circ x\). Sigui
\(t\colon T\hookrightarrow Y\) la imatge inversa \(\langle x,a\rangle^{-1}(K^c)\).
Aleshores \(\langle x,a\rangle\circ t\in K^c\) i, per tant,
\((\id_X\times f)\circ\langle x,a\rangle\circ t=\langle\id_X,b\rangle\circ x\circ t\)
pertany a la imatge directa \(N\). Però també pertany a \(N^c\), perquè
\(\langle\id_X,b\rangle\) hi pertany. Com que \(N\cap N^c\equiv\bot\), \(T\)
admet una aplicació cap a \(\mathbf0\) i és inicial (estrictesa). Per~(a),
\(\langle x,a\rangle\in_Y(K^c)^c\equiv K\), és a dir,
\(a\mathrel{\varepsilon}(\alpha\circ x)\).

\(\Leftarrow\)) Suposem la condició de l'enunciat, i sigui
\(t\colon T\hookrightarrow X\) la imatge inversa \(\langle\id_X,b\rangle^{-1}(N)\),
de manera que \(\langle t,b\circ t\rangle\in_T N\). Com a la prova de~(2),
el pullback de l'epimorfisme \(K^c\twoheadrightarrow N\) al llarg del
testimoni d'aquesta pertinença dona un epimorfisme
\(p\colon T'\twoheadrightarrow T\) i un element \(\langle x',a'\rangle\colon
T'\to X\times A\) de \(K^c\) amb \((\id_X\times f)\circ\langle x',a'\rangle
=\langle t,b\circ t\rangle\circ p\); és a dir, \(x'=t\circ p\) i
\(f\circ a'=b\circ x'\). La hipòtesi dona \(\langle x',a'\rangle\in_{T'}K\), i
com que també és a \(K^c\), \(T'\) és inicial. Per~(b), \(T\) és inicial, i
per~(a), \(\langle\id_X,b\rangle\in_X N^c\), és a dir,
\(b\mathrel{\varepsilon}(\forall_f\circ\alpha)\).
\end{prova}

\begin{proposicio}[Els quantificadors com a adjunts]\label{prop:quantificadors-adjuncio}
Sigui \(f\colon A\to B\). Per a tot estadi \(X\), tota
\(\alpha\colon X\to P(A)\) i tota \(\beta\colon X\to P(B)\), amb parts
associades \(s_\alpha\hookrightarrow X\times A\) i
\(s_\beta\hookrightarrow X\times B\), i escrivint \(\subseteq\) per a la
inclusió de parts:
\[
s_{\exists_f\alpha}\subseteq s_\beta
\ \Longleftrightarrow\
s_\alpha\subseteq s_{P(f)\beta},
\qquad
s_{P(f)\beta}\subseteq s_\alpha
\ \Longleftrightarrow\
s_\beta\subseteq s_{\forall_f\alpha} .
\]
És a dir, \(\exists_f\) és adjunt per l'esquerra i \(\forall_f\) adjunt
per la dreta de la imatge inversa interna \(P(f)\), a cada estadi.
\end{proposicio}

\begin{prova}
Per les proposicions~\ref{prop:imatge-inversa-interna}
i~\ref{prop:quantificadors-lectura}, \(s_{\exists_f\alpha}=(\id_X\times f)_*s_\alpha\)
i \(s_{P(f)\beta}=(\id_X\times f)^{-1}s_\beta\); la primera equivalència és,
doncs, l'adjunció imatge directa\,--\,imatge inversa del
capítol~\ref{cap:imatges} (proposició~\ref{prop:adjuncio-imatges}) per a
l'aplicació \(\id_X\times f\). La segona se'n dedueix per dualitat
booleana:
\(s_\beta\subseteq s_{\forall_f\alpha}=\bigl((\id\times f)_*(s_\alpha^c)\bigr)^c\)
equival, per l'antimonotonia del complement, a
\((\id\times f)_*(s_\alpha^c)\subseteq s_\beta^c\), i per la primera
adjunció a \(s_\alpha^c\subseteq(\id\times f)^{-1}(s_\beta^c)
=\bigl((\id\times f)^{-1}s_\beta\bigr)^c\) (la imatge inversa commuta amb
el complement, proposició~\ref{prop:imatge-inversa-complement}), que per
l'antimonotonia és \((\id\times f)^{-1}s_\beta\subseteq s_\alpha\).
\end{prova}

\begin{observacio}[Lectura clàssica]\label{obs:quantificadors-classica}
Per a \(f\colon A\to B\), la imatge inversa és
\(P(f)(\beta)=f^{-1}(\beta)=\{\,a:f(a)\in\beta\,\}\); el quantificador
existencial és la imatge directa
\(\exists_f(\alpha)=\{\,b:\exists a\,(f(a)=b\ \text{i}\ a\in\alpha)\,\}\);
i l'universal és
\(\forall_f(\alpha)=\{\,b:\forall a\,(f(a)=b\Rightarrow a\in\alpha)\,\}\).
Els noms «existencial» i «universal» no són metafòrics: quan
\(f={!}_A\colon A\to\mathbf1\), \(\exists_A\) diu si hi ha algun element de
\(A\) que pertany a la part, i \(\forall_A\) si tots hi pertanyen. La
lògica de predicats sobre \(A\), habitualment presa com a llenguatge
extern, apareix aquí com l'àlgebra de \(P(A)\) amb aquests adjunts. El
preu de llegir-ho sobre elements generalitzats és el recobriment de la
proposició~\ref{prop:quantificadors-lectura}: la mateixa frontera que a
la composició de relacions.
\end{observacio}

\section{\texorpdfstring{La inclusió com a ordre sobre \(P(A)\)}{La inclusió com a ordre sobre P(A)}}\label{sec:inclusio-ordre}

La inclusió de parts \(\subseteq_A\) del capítol~\ref{cap:parts} era una
relació del metallenguatge entre monomorfismes. Ara la realitzarem com una
\emph{part} de \(P(A)\times P(A)\), és a dir, com una relació interna
sobre el conjunt de parts, i mostrarem que és un ordre. L'eina és el
quantificador universal al llarg d'una projecció.

\begin{definicio}[Inclusió interna]\index{Inclusió interna}\label{def:inclusio-interna}
Sigui \(\Rightarrow:=\vee\circ(\neg\times\id_{\mathbf2})\colon
\mathbf2\times\mathbf2\to\mathbf2\) la implicació booleana, i sigui
\(I_A\hookrightarrow P(A)\times P(A)\times A\) la part classificada per
\(\Rightarrow\circ\langle\operatorname{ev}_1,\operatorname{ev}_2\rangle\).
La \emph{relació d'inclusió interna} sobre \(P(A)\) és la part
\[
{\subseteq_A}\;:=\;\forall_\pi(I_A)\;\hookrightarrow\;P(A)\times P(A),
\]
on \(\pi\colon P(A)\times P(A)\times A\to P(A)\times P(A)\) és la
projecció i \(\forall_\pi\) s'aplica al nom de la part \(I_A\) (és a dir,
\(\subseteq_A\) és la part de \(P(A)\times P(A)\) el nom de la qual és
\(\forall_\pi\ulcorner I_A\urcorner\)).
\end{definicio}

\begin{proposicio}[Lectura de la inclusió interna]\label{prop:inclusio-interna-lectura}
Siguin \(\alpha,\beta\colon X\to P(A)\), amb parts associades
\(s_\alpha,s_\beta\hookrightarrow X\times A\). Aleshores
\[
\begin{aligned}
\langle\alpha,\beta\rangle\in_X{\subseteq_A}
\ &\Longleftrightarrow\
s_\alpha\subseteq s_\beta\\
&\Longleftrightarrow\
\text{per a tot } x\colon Y\to X \text{ i tot } a\colon Y\to A,\\
&\qquad\qquad
a\mathrel{\varepsilon}(\alpha\circ x)\Rightarrow a\mathrel{\varepsilon}(\beta\circ x).
\end{aligned}
\]
En particular, sobre els noms de dues parts \(m,n\) de \(A\),
\(\langle\ulcorner m\urcorner,\ulcorner n\urcorner\rangle\in_{\mathbf1}{\subseteq_A}\)
si i només si \(m\subseteq_A n\).
\end{proposicio}

\begin{prova}
Per la lectura del quantificador universal
(proposició~\ref{prop:quantificadors-lectura}(3)) amb \(f=\pi\),
\(\langle\alpha,\beta\rangle\in_X\forall_\pi(I_A)\) vol dir que, per a tot
\(x\colon Y\to X\) i tot \(c\colon Y\to P(A)\times P(A)\times A\) amb
\(\pi\circ c=\langle\alpha,\beta\rangle\circ x\) ---és a dir, per a tot
\(c=\langle\alpha\circ x,\beta\circ x,a\rangle\) amb \(a\colon Y\to A\)---,
\(c\in_YI_A\), és a dir,
\(\Rightarrow\circ\langle u,v\rangle=1\circ{!}_Y\) amb
\(u=\operatorname{ev}\circ\langle\alpha\circ x,a\rangle\) i
\(v=\operatorname{ev}\circ\langle\beta\circ x,a\rangle\). Ara bé,
\(\Rightarrow\circ\langle u,v\rangle=\neg u\vee v\) és la característica
de la part \(T^c\cup U\) de \(Y\), on \(T,U\) són les parts classificades
per \(u,v\); i \(T^c\cup U\) és total si i només si \(T\subseteq U\)
(lleis booleanes, teorema~\ref{teo:lleis-booleanes}). Per la
proposició~\ref{prop:pertinenca-interna-lectura}, \(T\) i \(U\) són les
imatges inverses \(\langle x,a\rangle^{-1}s_\alpha\) i
\(\langle x,a\rangle^{-1}s_\beta\). Que la condició valgui per a tots els
\(x,a\) és, prenent \(Y=s_\alpha\) i \(\langle x,a\rangle\) la seva
inclusió, exactament \(s_\alpha\subseteq s_\beta\); i recíprocament,
\(s_\alpha\subseteq s_\beta\) la implica per monotonia de la imatge
inversa. Això dona les dues equivalències (la tercera és la segona,
llegida amb el corol·lari~\ref{cor:pertinenca-universal}). Per als noms,
\(X=\mathbf1\), i \(s_{\ulcorner m\urcorner}\) és \(m\) vista dins
\(\mathbf1\times A\cong A\).
\end{prova}

\begin{teorema}[La inclusió interna és un ordre parcial]\label{teo:inclusio-ordre}
La relació \(\subseteq_A\) sobre \(P(A)\) és reflexiva, antisimètrica i
transitiva, en el sentit intern següent, per a tots
\(\alpha,\beta,\gamma\colon X\to P(A)\):
\begin{enumerate}
\item \(\langle\alpha,\alpha\rangle\in_X{\subseteq_A}\);
\item si \(\langle\alpha,\beta\rangle\in_X{\subseteq_A}\) i
\(\langle\beta,\alpha\rangle\in_X{\subseteq_A}\), aleshores
\(\alpha=\beta\);
\item si \(\langle\alpha,\beta\rangle\in_X{\subseteq_A}\) i
\(\langle\beta,\gamma\rangle\in_X{\subseteq_A}\), aleshores
\(\langle\alpha,\gamma\rangle\in_X{\subseteq_A}\).
\end{enumerate}
A més, la intersecció i la unió internes són l'ínfim i el suprem respecte
d'aquest ordre: \(\langle\gamma,\alpha\cap\beta\rangle\in_X{\subseteq_A}\)
si i només si \(\langle\gamma,\alpha\rangle\in_X{\subseteq_A}\) i
\(\langle\gamma,\beta\rangle\in_X{\subseteq_A}\), i dualment per a
\(\cup\).
\end{teorema}

\begin{prova}
Per la proposició~\ref{prop:inclusio-interna-lectura}, tot es redueix a la
inclusió de parts de \(X\times A\): \(s_\alpha\subseteq s_\alpha\); si
\(s_\alpha\subseteq s_\beta\) i \(s_\beta\subseteq s_\alpha\), les parts
són equivalents i tenen la mateixa característica, de manera que
\(\alpha=\beta\) per la unicitat de la transposada; i la transitivitat de
la inclusió de parts és immediata de la
proposició~\ref{prop:inclusio-pertinenca}. Per a
l'ínfim, \(s_{\alpha\cap\beta}=s_\alpha\cap s_\beta\) per la
proposició~\ref{prop:booleanes-puntuals} (les característiques
coincideixen), i una part és inclosa en una intersecció si i només si ho
és en cadascun dels factors (proposició~\ref{prop:pertinenca-interseccio});
dualment, \(s_{\alpha\cup\beta}=s_\alpha\cup s_\beta\) i la unió és el
suprem per les lleis booleanes (proposició~\ref{prop:inclusio-booleana}).
\end{prova}

\begin{observacio}[Lectura clàssica]\label{obs:inclusio-ordre-classica}
El teorema diu que \((P(A),\subseteq)\) és un reticle booleà: un conjunt
parcialment ordenat per la inclusió, amb la intersecció com a ínfim, la
unió com a suprem, la part buida i la total com a mínim i màxim, i el
complement com a operació de complementació. És l'estructura de reticle de
potència de la teoria clàssica, ara construïda com una relació interna
sobre l'objecte \(P(A)\) en lloc de descrita externament sobre la
col·lecció de subconjunts.
\end{observacio}

\begin{resumcap}
\apartat{Idees centrals}
\begin{enumerate}
  \item Tres nivells: la part \(m\colon S\hookrightarrow A\) (una fletxa), el seu nom \(\ulcorner m\urcorner\colon\mathbf1\to P(A)\) (un punt) i el conjunt \(P(A)\). La pertinença interna \(\in_A\hookrightarrow A\times P(A)\) és la part classificada per l'avaluació; sobre noms recupera la pertinença generalitzada.
  \item Una aplicació \(X\to P(A)\) és el mateix que una part de \(X\times A\): una família de parts d'\(A\) indexada per \(X\). Tot el capítol llegeix les construccions sobre \(P(A)\) a través d'aquesta correspondència.
  \item Les operacions booleanes internes es defineixen per transposició i es calculen component avaluacions; la lectura amb \enquote{i} val a tots els estadis, les lectures amb \enquote{o} i \enquote{no} només sobre punts. Sobre noms recuperen les operacions de parts, i les lleis booleanes es transporten.
  \item \(P(f)\) és la imatge inversa interna; \(\exists_f\) és la transposada de la imatge directa de la pertinença al llarg de \(\id\times f\), i \(\forall_f:=\neg\exists_f\neg\). La imatge directa commuta amb el canvi de base (Beck--Chevalley), i això dona la lectura de \(\exists_f\) sobre elements generalitzats (amb recobriment) i l'adjunció \(\exists_f\dashv P(f)\dashv\forall_f\) a cada estadi.
  \item La inclusió interna és \(\forall_\pi\) de la implicació puntual; es llegeix com la inclusió de les parts associades, i és un ordre amb \(\cap\) i \(\cup\) com a ínfim i suprem.
  \item Cap axioma nou: tot surt dels exponencials, el classificador i la factorització imatge.
\end{enumerate}
\apartat{Errors típics} Identificar la part \(m\) amb el seu nom \(\ulcorner m\urcorner\) o amb el conjunt \(S\); llegir \(a\mathrel{\varepsilon}(\alpha\cup\beta)\) com \enquote{\(a\mathrel{\varepsilon}\alpha\) o \(a\mathrel{\varepsilon}\beta\)} per a un element generalitzat; creure que \(\exists_f\) \enquote{existeix} per la seva lectura puntual, quan és una aplicació que cal construir amb la imatge directa; esperar el testimoni de \(\exists_f\) en el mateix estadi sense elecció; definir la inclusió interna amb \(\forall_A\) (que quantifica cap a \(\mathbf2\)) en lloc de \(\forall_\pi\); i oblidar que \(\forall_f=\neg\exists_f\neg\) depèn del caràcter booleà de la teoria.
\apartat{Lectura recomanada} \cite[cap.~7--8]{lawvere-rosebrugh} per al conjunt de parts, la pertinença interna, els quantificadors com a adjunts i el reticle de parts; \cite[article~VI, sessions~32--33]{lawvere-schanuel} per a una presentació elemental de l'objecte de parts; \cite[cap.~9]{awodey} per a les adjuncions en general, amb els quantificadors com a exemple; \cite[cap.~IV]{maclane-moerdijk} per als objectes de parts en un topos i la condició de Beck--Chevalley; i \cite[§§5--9]{halmos} per al conjunt de parts i les operacions amb elements.
\end{resumcap}

\chapter{L'objecte dels nombres naturals}\label{cap:naturals}

\begin{objectius}
\apartat{Prerequisits} Del capítol~\ref{cap:exponencials}: exponencials, transposició, noms d'aplicacions i composició interna. Del capítol~\ref{cap:parts}: parts, equalitzadors, seccions i retraccions, i el teorema del bimorfisme. Del capítol~\ref{cap:valors}: \(0\ne1\), el lema del buit i la lectura de les sumes i les unions sobre punts. Del capítol~\ref{cap:parts-interna}: quantificadors i operacions booleanes internes, per a la part auxiliar del teorema del bon ordre.
\apartat{Objectius} En acabar el capítol, el lector hauria de poder:
\begin{itemize}
  \item enunciar l'axioma de l'objecte dels nombres naturals com una propietat universal de recursió i demostrar-ne la unicitat llevat d'un únic isomorfisme;
  \item construir la iterada d'una autoaplicació i la recursió amb paràmetres a partir de l'axioma, pujant als exponencials;
  \item demostrar el principi d'inducció i la unicitat per inducció, amb i sense paràmetres, i explicar per què són la cara d'unicitat de la recursió;
  \item deduir les equacions de Peano, en particular la injectivitat del successor mitjançant el predecessor construït sobre \(\N\times\N\), i que tot punt és zero o successor;
  \item definir suma, producte i exponenciació per recursió i demostrar-ne les lleis per inducció amb paràmetres;
  \item construir l'ordre dels naturals com una part de \(\N\times\N\), demostrar que és total sobre punts i que \(\N\) és ben ordenat, i saber per què aquestes propietats es comproven sobre punts globals.
\end{itemize}
\end{objectius}

Els capítols precedents han completat el que es podria anomenar la part
\emph{finita} de la teoria: aplicacions, parts, construccions booleanes,
relacions, quocients, imatges, exponencials i el conjunt de parts. Cap
d'aquestes construccions produeix, per si sola, un conjunt infinit; els
axiomes acumulats permeten formar sumes, productes i exponencials de
conjunts ja donats, però mai no forcen l'existència d'un conjunt amb
infinits elements distingibles. Aquest capítol introdueix l'últim axioma
estructural essencial del llibre, que trenca aquesta limitació:
l'existència d'un \emph{objecte dels nombres naturals}.

La forma d'aquest axioma segueix la disciplina de tot el
desenvolupament. No postularem una llista de propietats de \(\N\) ---les
equacions de Peano--- sinó una única \emph{propietat universal}: la
capacitat de definir aplicacions per recursió. D'aquesta propietat en
derivarem, com a teoremes, el principi d'inducció, les equacions de Peano
i tota l'aritmètica elemental. Aquesta és la lliçó del capítol: el que en
la presentació habitual són cinc axiomes independents sobre el zero i el
successor apareix aquí com a conseqüència d'una sola condició de
definibilitat, coherent amb la manera com el llibre ha tractat el
producte, l'exponencial i el conjunt de parts.

\section{La propietat universal de recursió}\label{sec:nno-universal}

En conjunts, definir una successió \(u_0,u_1,u_2,\dots\) per recursió vol
dir donar un valor inicial \(u_0\) i una regla que, a partir de \(u_n\),
en produeix el següent \(u_{n+1}\). Perquè aquesta descripció determini
una única successió cal que el conjunt d'índexs estigui generat pel zero i
el pas al successor sense cap altra relació. La propietat universal que
segueix diu exactament això: que a partir d'un punt inicial i una
autoaplicació de qualsevol conjunt hi ha una i una sola manera de
respectar el zero i el successor.

\begin{axioma}[Existència d'un objecte dels nombres naturals]\label{ax:naturals}
Existeixen un conjunt \(\N\), un element global \(0\colon\mathbf1\to\N\) i
una aplicació \(s\colon\N\to\N\), el \emph{successor}, amb la propietat
universal següent: per a tot conjunt \(A\), tot element global
\(a\colon\mathbf1\to A\) i tota aplicació \(t\colon A\to A\), existeix una
\emph{única} aplicació
\[
u\colon\N\to A
\qquad\text{tal que}\qquad
u\circ 0=a
\quad\text{i}\quad
u\circ s=t\circ u .
\]
\end{axioma}

\begin{notacio}[L'objecte dels nombres naturals]\label{not:nno}
Anomenarem \(\N\) l'\emph{objecte dels nombres naturals}, i
\((\N,0,s)\) la terna que el presenta. L'aplicació \(u\) de
l'axioma~\ref{ax:naturals} és la \emph{recursora} de les dades
\((a,t)\), i escriurem \(u=\operatorname{rec}(a,t)\) quan calgui
distingir-la. Les dues condicions
\[
u\circ 0=a,
\qquad
u\circ s=t\circ u
\]
diuen que \(u\) fa commutar el diagrama
\[
\mathbf1\xrightarrow{\;0\;}\N\xrightarrow{\;s\;}\N,
\qquad
\mathbf1\xrightarrow{\;a\;}A\xrightarrow{\;t\;}A,
\]
amb \(u\) vertical; una terna \((A,a,t)\) formada per un conjunt, un punt
i una autoaplicació l'anomenarem un \emph{sistema recursiu}. El predicat
eliminable corresponent, \(\operatorname{Nno}(\N,0,s)\), s'escriu com la
propietat universal, quantificant els punts globals amb el predicat de
conjunt terminal del capítol~\ref{cap:aplicacions} i les composicions amb
el predicat \(\operatorname{Com}\); no el desplegarem, perquè segueix el
mateix patró que \(\operatorname{Exp}\)
(notació~\ref{not:exponencial}).
\end{notacio}

Com en els altres axiomes d'existència, la terna \((\N,0,s)\) és única
llevat d'un isomorfisme compatible. Aquesta unicitat és el primer teorema,
i la seva demostració és el patró de tots els arguments per recursió que
seguiran.

\begin{teorema}[Unicitat de l'objecte dels nombres naturals]\label{teo:nno-unicitat}
Siguin \((\N,0,s)\) i \((\N',0',s')\) dues ternes que compleixen la
propietat universal de l'axioma~\ref{ax:naturals}. Aleshores existeix un
únic isomorfisme \(\varphi\colon\N\to\N'\) amb
\[
\varphi\circ 0=0'
\qquad\text{i}\qquad
\varphi\circ s=s'\circ\varphi .
\]
\end{teorema}

\begin{prova}
Apliquem la propietat universal de \((\N,0,s)\) al sistema recursiu
\((\N',0',s')\): hi ha una única \(\varphi\colon\N\to\N'\) amb
\(\varphi\circ 0=0'\) i \(\varphi\circ s=s'\circ\varphi\). Simètricament,
la propietat de \((\N',0',s')\) aplicada a \((\N,0,s)\) dona una única
\(\psi\colon\N'\to\N\) amb \(\psi\circ 0'=0\) i
\(\psi\circ s'=s\circ\psi\). La composició \(\psi\circ\varphi\colon\N\to\N\)
compleix \((\psi\circ\varphi)\circ 0=0\) i
\((\psi\circ\varphi)\circ s=s\circ(\psi\circ\varphi)\); però la identitat
\(\id_\N\) compleix les mateixes dues condicions, i la propietat universal
de \((\N,0,s)\) aplicada al sistema \((\N,0,s)\) mateix diu que hi ha una
sola aplicació així. Per tant \(\psi\circ\varphi=\id_\N\, \); igualment
\(\varphi\circ\psi=\id_{\N'}\). Així \(\varphi\) és un isomorfisme, i és
l'únic amb les propietats requerides per la unicitat de la recursora.
\end{prova}

\begin{observacio}[El paper de la unicitat en la recursió]\label{obs:unicitat-recursio}
La demostració anterior conté, en miniatura, l'estratègia que farem servir
constantment: per veure que dues aplicacions \(\N\to A\) coincideixen,
n'hi ha prou de comprovar que totes dues resolen un mateix problema de
recursió, perquè la solució és única. Aquest «principi d'unicitat de la
recursora» és, com veurem a la secció~\ref{sec:induccio}, una altra cara
del principi d'inducció.
\end{observacio}

\begin{observacio}[Lectura clàssica]\label{obs:nno-classica}
En conjunts, \(\N=\{0,1,2,\dots\}\), \(0\) és el zero i \(s(n)=n+1\). La
propietat universal és el \emph{teorema de recursió} de Dedekind: donats
un conjunt \(A\), un element \(a\in A\) i una funció \(t\colon A\to A\),
hi ha una única successió \(u\colon\N\to A\) amb \(u(0)=a\) i
\(u(n+1)=t(u(n))\). El que habitualment es demostra a partir dels axiomes
de Peano ---que la recursió està ben definida--- és aquí l'axioma mateix,
i els axiomes de Peano en seran conseqüència
(secció~\ref{sec:peano}).
\end{observacio}

\section{Iteració i el teorema de recursió}\label{sec:recursio}

El cas més simple de la propietat universal és aquell en què el valor
inicial és un punt genèric i la regla no depèn de l'índex: la
\emph{iteració} d'una aplicació. En separem l'enunciat perquè és la forma
en què la recursió apareix més sovint i perquè fixa la notació per a les
potències d'una autoaplicació.

\begin{definicio}[Iterada d'una autoaplicació]\index{Iterada d'una autoaplicació}\label{def:iterada}
Sigui \(t\colon A\to A\) una autoaplicació d'un conjunt \(A\). La
\emph{iterada} de \(t\) és l'única aplicació
\[
\operatorname{it}(t)\colon\N\times A\to A
\]
tal que, escrivint \(t^{n}(x):=\operatorname{it}(t)\circ\langle n,x\rangle\)
per a elements generalitzats \(n\colon X\to\N\) i \(x\colon X\to A\),
\[
t^{0}(x)=x
\qquad\text{i}\qquad
t^{\,s\circ n}(x)=t\bigl(t^{n}(x)\bigr).
\]
\end{definicio}

\begin{proposicio}[Existència i unicitat de la iterada]\label{prop:iterada}
Per a tota autoaplicació \(t\colon A\to A\), l'aplicació
\(\operatorname{it}(t)\) de la definició~\ref{def:iterada} existeix i és
única.
\end{proposicio}

\begin{prova}
Considerem l'exponencial \(A^{A}\) del capítol~\ref{cap:exponencials} i
el sistema recursiu \((A^{A},\ulcorner\id_A\urcorner,C_t)\), on
\(\ulcorner\id_A\urcorner\colon\mathbf1\to A^{A}\) és el nom de la
identitat i \(C_t\colon A^{A}\to A^{A}\) és la composició interna per
\(t\) per l'esquerra: la transposada de
\(t\circ\operatorname{ev}\colon A^A\times A\to A\), que sobre punts envia
\(g\) a \(t\circ g\) (és la composició interna del
capítol~\ref{cap:exponencials} amb el primer argument fixat).
Per l'axioma~\ref{ax:naturals} hi ha una única
\(\pi\colon\N\to A^{A}\) amb \(\pi\circ 0=\ulcorner\id_A\urcorner\) i
\(\pi\circ s=C_t\circ\pi\); intuïtivament \(\pi(n)=\ulcorner
t^{n}\urcorner\). Descurrificant \(\pi\) s'obté
\(\operatorname{it}(t):=\operatorname{ev}\circ(\pi\times\id_A)
\colon\N\times A\to A\), i les dues equacions de \(\pi\), llegides a
través de l'avaluació, són exactament \(t^{0}(x)=x\) i
\(t^{\,s\circ n}(x)=t(t^{n}(x))\). La unicitat de \(\operatorname{it}(t)\)
es redueix, per la unicitat de la transposada
(definició~\ref{def:exponencial}), a la de \(\pi\), que és la de la
recursora.
\end{prova}

\begin{observacio}[La recursió general via iteració]\label{obs:recursio-iteracio}
La iteració és, en realitat, tan general com la recursió: donada una regla
\(t\colon A\to A\) i un valor inicial \(a\colon\mathbf1\to A\), la
recursora \(u=\operatorname{rec}(a,t)\) de l'axioma~\ref{ax:naturals}
s'obté iterant \(t\) a partir de \(a\), és a dir,
\(u\circ n=t^{n}(a)\) per a tot \(n\colon X\to\N\). Aquesta és la lectura
que justifica escriure la recursió com «aplicar \(n\) vegades la regla a
la llavor». La proposició~\ref{prop:iterada} mostra que aquest «aplicar
\(n\) vegades» és una construcció legítima de la teoria, no una
abreviatura metalingüística: viu com l'aplicació \(\operatorname{it}(t)\).
\end{observacio}

\section{El principi d'inducció}\label{sec:induccio}

La propietat universal és una afirmació sobre \emph{definició} d'aplicacions.
El principi d'inducció n'és la contrapartida sobre \emph{demostració}:
diu que \(\N\) no té cap part pròpia que contingui el zero i sigui tancada
pel successor. Els dos enunciats estan estretament lligats: el principi
d'inducció equival a la part d'\emph{unicitat} de la propietat universal
---la que garanteix que dues aplicacions definides per la mateixa recursió
coincideixen---, mentre que la part d'\emph{existència}, que produeix
efectivament les aplicacions recursives, no se'n dedueix. Aquesta
correspondència, que precisarem a
l'observació~\ref{obs:recursio-i-induccio}, és el cor del capítol.

\begin{teorema}[Principi d'inducció]\label{teo:induccio}
Sigui \(m\colon S\hookrightarrow\N\) una part de \(\N\) tal que
\begin{enumerate}
\item \(0\) es factoritza per \(m\) (és a dir, \(0\in_{\mathbf1}m\)); i
\item \(s\) preserva \(m\): existeix una aplicació \(s_S\colon S\to S\)
tal que \(m\circ s_S=s\circ m\); és a dir, \(s\circ m\) es factoritza per
\(m\).
\end{enumerate}
Aleshores \(m\) és un isomorfisme; equivalentment, \(S=\N\) com a parts de
\(\N\).
\end{teorema}

\begin{prova}
Les hipòtesis diuen que \((S,0_S,s_S)\) és un sistema recursiu, on
\(0_S\colon\mathbf1\to S\) és la factorització de \(0\) per \(m\)
(condició~1) i \(s_S\colon S\to S\) és l'aplicació donada per la
condició~2. Aquestes compleixen
\[
m\circ 0_S=0,
\qquad
m\circ s_S=s\circ m .
\]
Així \(m\colon S\to\N\) és un morfisme de sistemes recursius. Per la
propietat universal de \((\N,0,s)\) aplicada al sistema \((S,0_S,s_S)\),
hi ha una única \(u\colon\N\to S\) amb \(u\circ 0=0_S\) i
\(u\circ s=s_S\circ u\). Llavors \(m\circ u\colon\N\to\N\) compleix
\(m\circ u\circ 0=m\circ 0_S=0\) i
\(m\circ u\circ s=m\circ s_S\circ u=s\circ m\circ u\); com que \(\id_\N\)
compleix les mateixes condicions, la unicitat de la recursora dona
\(m\circ u=\id_\N\). Aleshores \(m\) és una retracció i, per tant, un
epimorfisme (proposició~\ref{prop:seccio-mono-retraccio-epi}); però \(m\)
és també un monomorfisme (és una part). Una aplicació que és alhora
monomorfisme i epimorfisme és un isomorfisme
(corol·lari~\ref{cor:isomorfisme-bimorfisme}). Per tant \(m\) és un
isomorfisme.
\end{prova}

\begin{observacio}[Inducció sobre elements generalitzats]\label{obs:induccio-generalitzada}
La versió del teorema~\ref{teo:induccio} en termes de parts és la més
directa d'usar, però la seva força pràctica apareix en la formulació
predicativa. Sigui \(\varphi\) una propietat d'elements generalitzats de
\(\N\), és a dir, una part \(m_\varphi\hookrightarrow\N\). Aleshores el
teorema diu:
\[
\bigl(0\in m_\varphi\ \text{ i }\ (n\in m_\varphi\Rightarrow s\circ
n\in m_\varphi)\bigr)
\quad\Longrightarrow\quad
n\in m_\varphi\ \text{ per a tot } n\colon X\to\N .
\]
En aquesta forma, l'inducció és la regla de demostració habitual: cas
base al zero, pas inductiu del successor, conclusió per a tot natural.
Notem que la propietat \(\varphi\) ha de ser una \emph{part} de \(\N\),
és a dir, ha de correspondre a una aplicació \(\N\to\mathbf2\) via el
classificador (axioma~\ref{ax:representacio}); no tota afirmació del
metallenguatge defineix una part, i aquesta restricció és el que
distingeix la inducció interna de l'esquema d'inducció clàssic
(observació~\ref{obs:induccio-classica}).
\end{observacio}

\begin{proposicio}[Unicitat per inducció]\label{prop:unicitat-induccio}
Siguin \(f,g\colon\N\to A\) dues aplicacions. Si \(f\circ 0=g\circ 0\) i
\(f\circ s=g\circ s\) sempre que \(f\) i \(g\) coincideixin ---més
precisament, si la part on \(f\) i \(g\) coincideixen conté el zero i és
tancada pel successor--- aleshores \(f=g\).
\end{proposicio}

\begin{prova}
Sigui \(e\colon E\hookrightarrow\N\) l'equalitzador de \(f\) i \(g\)
(capítol~\ref{cap:parts}), la part de \(\N\) on totes dues coincideixen.
La hipòtesi diu que \(E\) conté el zero i és tancada pel successor, de
manera que pel teorema~\ref{teo:induccio} \(e\) és un isomorfisme; per
tant \(f\) i \(g\) coincideixen a tot \(\N\), és a dir, \(f=g\).
\end{prova}

Les demostracions per inducció de la resta del capítol tenen, gairebé
sempre, \emph{paràmetres}: es vol veure que dues aplicacions
\(X\times\N\to A\) coincideixen, fent inducció sobre la variable de \(\N\)
amb la de \(X\) fixada. La part de \(X\times\N\) on coincideixen no és
una part de \(\N\), i el teorema~\ref{teo:induccio} no s'hi aplica
directament; els exponencials resolen la dificultat.

\begin{corollari}[Unicitat per inducció amb paràmetres]\label{cor:induccio-parametres}
Siguin \(f,g\colon X\times\N\to A\). Si
\(f\circ(\id_X\times0)=g\circ(\id_X\times0)\) (com a aplicacions
\(X\times\mathbf1\to A\)) i, per a tot estadi \(Y\) i tots
\(x\colon Y\to X\), \(n\colon Y\to\N\),
\[
f\circ\langle x,n\rangle=g\circ\langle x,n\rangle
\quad\Longrightarrow\quad
f\circ\langle x,s\circ n\rangle=g\circ\langle x,s\circ n\rangle,
\]
aleshores \(f=g\).
\end{corollari}

\begin{prova}
Transposem respecte de \(X\): \(f\) i \(g\) corresponen a
\(\widehat f,\widehat g\colon\N\to A^{X}\), amb
\(f=\operatorname{ev}\circ(\widehat f\times\id_X)\circ\sigma\), on
\(\sigma\colon X\times\N\to\N\times X\) és la commutació. Sigui
\(e\colon E\hookrightarrow\N\) l'equalitzador de \(\widehat f\) i
\(\widehat g\). La primera hipòtesi diu que
\(\widehat f\circ0\) i \(\widehat g\circ0\) tenen la mateixa
descurrificació, de manera que \(0\in E\). Si \(n\colon Y\to\N\) pertany a
\(E\), aleshores \(\widehat f\circ n=\widehat g\circ n\) i, descurrificant,
\(f\circ\langle x,n\rangle=g\circ\langle x,n\rangle\) per a tot
\(x\colon Y\to X\); per la segona hipòtesi,
\(f\circ\langle x,s\circ n\rangle=g\circ\langle x,s\circ n\rangle\) per a
tot \(x\), i com que els \(\langle x,s\circ n\rangle\) amb \(x\) arbitrari
determinen la restricció a \(X\times Y\) (prenent \(x\) la projecció
\(X\times Y\to X\) i \(n\) composta amb l'altra projecció), les
transposades coincideixen sobre \(s\circ n\): \(s\circ n\in E\). Per
inducció (teorema~\ref{teo:induccio}), \(E=\N\), d'on
\(\widehat f=\widehat g\) i \(f=g\).
\end{prova}

En la pràctica escriurem aquestes induccions amb elements generalitzats
\(m,n\colon X\to\N\), \enquote{amb \(m\) fixat i inducció sobre \(n\)}: el
corol·lari és el que legitima aquesta manera de parlar.

\begin{observacio}[Recursió i inducció]\label{obs:recursio-i-induccio}
La proposició~\ref{prop:unicitat-induccio} explicita el que
l'observació~\ref{obs:unicitat-recursio} anunciava: la part
d'\emph{unicitat} de la propietat universal és exactament el principi
d'inducció. La part d'\emph{existència} de la recursora, en canvi, és el
que garanteix que les definicions per recursió produeixen aplicacions
efectives. Existència i unicitat de la recursora són, doncs, la definició
per recursió i la demostració per inducció, respectivament.
\end{observacio}

\section{Les equacions de Peano}\label{sec:peano}

Amb la inducció disponible podem derivar, com a teoremes, les propietats
que en la presentació axiomàtica de Peano es prenen com a punt de partida:
que el zero no és cap successor i que el successor és injectiu. La
coherència amb el llibre demana obtenir-les de la propietat universal, i
no a l'inrevés.

\begin{teorema}[Equacions de Peano]\label{teo:peano}
En tota terna \((\N,0,s)\) que compleixi l'axioma~\ref{ax:naturals}:
\begin{enumerate}
\item el successor \(s\colon\N\to\N\) és un monomorfisme;
\item el zero no és cap successor: no hi ha cap \(n\colon X\to\N\) amb
\(X\) no inicial tal que \(s\circ n=0\circ{!}_X\). En termes clàssics,
això diu que la imatge del successor i el zero són disjunts.
\end{enumerate}
\end{teorema}

\begin{prova}
\emph{(2) El zero no és cap successor.} Definim per recursió una aplicació
\(d\colon\N\to\mathbf2\) amb \(d\circ 0=1\) i
\(d\circ s=0\circ{!}_\N\), és a dir, la recursora del sistema
\((\mathbf2,1,0\circ{!}_{\mathbf2})\) ---la regla constant \(0\)---. Si fos
\(s\circ n=0\circ{!}_X\) per a algun \(n\colon X\to\N\) amb \(X\) no
inicial, component amb \(d\) obtindríem
\(d\circ s\circ n=d\circ 0\circ{!}_X\), és a dir,
\(0\circ{!}_X=1\circ{!}_X\). Com que \(X\) no és inicial, té un punt
\(x\) (lema del buit, lema~\ref{lema:buit}), i component-hi obtenim
\(0=1\), contra el teorema~\ref{teo:zero-diferent-u}. Per tant cap
successor no és el zero.

\emph{(1) El successor és injectiu.} Construïm un \emph{predecessor}
\(p\colon\N\to\N\) amb \(p\circ s=\id_\N\). No es pot fer amb una recursió
directa sobre \(\N\): la condició \(p\circ s=\id_\N\) no té la forma
\(u\circ s=t\circ u\) de l'axioma (una recursora amb \(t=\id\) compleix
\(u\circ s=u\), i és constant). El recurs habitual és recórrer parells
consecutius. Sigui \(q\colon\N\to\N\times\N\) la recursora del sistema
\((\N\times\N,\ \langle0,0\rangle,\ \langle p_2,\,s\circ p_2\rangle)\):
\[
q\circ0=\langle0,0\rangle,
\qquad
q\circ s=\langle p_2,s\circ p_2\rangle\circ q .
\]
Intuïtivament \(q(n)=(n-1,n)\) per a \(n\ge1\) i \(q(0)=(0,0)\). Vegem
primer que \(p_2\circ q=\id_\N\): es compleix
\(p_2\circ q\circ0=0\) i
\(p_2\circ q\circ s=p_2\circ\langle p_2,s\circ p_2\rangle\circ q
=s\circ(p_2\circ q)\), i \(\id_\N\) compleix \(\id\circ0=0\) i
\(\id\circ s=s\circ\id\); per la unicitat de la recursora del sistema
\((\N,0,s)\), \(p_2\circ q=\id_\N\). Definim \(p:=p_1\circ q\). Aleshores
\[
p\circ s=p_1\circ q\circ s=p_1\circ\langle p_2,s\circ p_2\rangle\circ q
=p_2\circ q=\id_\N .
\]
Així \(s\) és una secció de \(p\) i, per tant, un monomorfisme
(proposició~\ref{prop:seccio-mono-retraccio-epi}): si
\(s\circ n=s\circ n'\), component amb \(p\) resulta \(n=n'\). Conservarem
l'aplicació \(p\), el \emph{predecessor}, que compleix \(p\circ0=0\) i
\(p\circ s=\id_\N\).
\end{prova}

\begin{observacio}[Per què cal el parell]\label{obs:predecessor-parell}
L'argument anterior és el primer exemple d'una tècnica general: una
funció que \enquote{mira enrere} (el predecessor) no és una recursió
simple, però esdevé recursió simple si se'n recorda prou informació (el
parell \((n-1,n)\)). És la mateixa idea que, en programació, converteix
una recursió sobre el valor anterior en una iteració amb estat. Amb la
recursió amb paràmetres del final del capítol
(teorema~\ref{teo:recursio-parametres}) es pot fer el mateix per a
regles que depenen de l'índex.
\end{observacio}

\begin{lema}[Tot punt és zero o successor]\label{lema:zero-o-successor}
Sigui \(U\hookrightarrow\N\) la unió de la part \(0\colon\mathbf1\hookrightarrow\N\)
i de la imatge \(s_*\N\hookrightarrow\N\) del successor. Aleshores \(U\) és
total; en conseqüència, tot punt \(l\colon\mathbf1\to\N\) és \(0\) o és de
la forma \(s\circ l'\) per a un únic punt \(l'\), a saber \(l'=p\circ l\).
\end{lema}

\begin{prova}
\(0\in U\) perquè \(0\) pertany al primer sumand de la unió. Si \(n\in U\),
aleshores \(s\circ n\in s_*\N\subseteq U\) (tot \(s\circ n\) es factoritza
per la imatge de \(s\)). Per inducció (teorema~\ref{teo:induccio}), \(U\)
és total. Un punt \(l\) de \(\N\) pertany, doncs, a la unió
\(0\cup s_*\N\), i un punt d'una unió pertany a un dels dos membres
(la característica de la unió és \(\vee\) de les característiques, i
sobre un punt els valors són \(0\) o \(1\), corol·lari~\ref{cor:elements-de-2}).
Si \(l\in s_*\N\), hi ha \(l'\) amb \(s\circ l'=l\) (exhaustivitat de
l'epimorfisme \(\N\twoheadrightarrow s_*\N\) sobre punts,
proposició~\ref{prop:epi-exhaustiu-punts}), únic per la injectivitat de
\(s\), i \(p\circ l=p\circ s\circ l'=l'\).
\end{prova}

\begin{observacio}[Peano complet]\label{obs:peano-complet}
Amb el teorema~\ref{teo:peano} i el principi d'inducció
(teorema~\ref{teo:induccio}) es tenen les cinc afirmacions de Peano:
\(0\) és un natural; el successor és una aplicació \(\N\to\N\); el zero no
és successor de cap natural; el successor és injectiu; i val la inducció.
La diferència amb la presentació habitual és l'ordre lògic: aquí no són
axiomes sinó conseqüències d'una única propietat universal, que a més
garanteix ---cosa que els axiomes de Peano de primer ordre no fan--- que
les definicions per recursió estan sempre ben definides.
\end{observacio}

\begin{observacio}[Lectura clàssica]\label{obs:induccio-classica}
En la teoria clàssica de primer ordre, la inducció és un \emph{esquema}:
una còpia de l'axioma per a cada fórmula del llenguatge. Aquí la inducció
és un únic teorema, però quantifica sobre \emph{parts} de \(\N\), és a
dir, sobre aplicacions \(\N\to\mathbf2\). Com que, per l'axioma de
representació (axioma~\ref{ax:representacio}), les parts de \(\N\) es
corresponen amb els punts de \(P(\N)\), l'inducció d'aquest capítol és una
inducció de «segon ordre» respecte del llenguatge intern: abraça totes
les propietats que defineixen una part, no només les expressables per una
fórmula prefixada. Aquesta és la font de la categoricitat de \(\N\)
(teorema~\ref{teo:nno-unicitat}), que l'aritmètica de primer ordre no té.
\end{observacio}

\section{Suma, producte i exponenciació}\label{sec:aritmetica}

Tota l'aritmètica elemental es defineix ara per recursió. Definirem cada
operació fixant un argument i iterant sobre l'altre, i en derivarem les
propietats per inducció. Presentem la suma amb detall i indiquem el
producte i l'exponenciació, que segueixen el mateix patró.

\begin{definicio}[Suma]\index{Suma}\label{def:suma-naturals}
La \emph{suma} és l'aplicació \(+\colon\N\times\N\to\N\) definida per
recursió sobre el segon argument: és la iterada del successor,
\[
m+n:=s^{n}(m)=\operatorname{it}(s)\circ\langle n,m\rangle,
\]
de manera que, per a elements generalitzats \(m,n\colon X\to\N\),
\[
m+0=m
\qquad\text{i}\qquad
m+(s\circ n)=s\circ(m+n).
\]
\end{definicio}

\begin{proposicio}[Propietats de la suma]\label{prop:suma}
La suma és associativa, commutativa i té el zero com a element neutre:
per a tots \(k,m,n\colon X\to\N\),
\[
(k+m)+n=k+(m+n),
\qquad
m+n=n+m,
\qquad
0+n=n=n+0 .
\]
A més, és \emph{cancel·lativa}: si \(m+k=m+l\), aleshores \(k=l\); i una
suma és nul·la només trivialment: si \(k+l=0\), aleshores \(k=l=0\).
\end{proposicio}

\begin{prova}
Totes les igualtats són igualtats entre aplicacions \(\N\times\N\to\N\)
(o \(\N^3\to\N\)), i es demostren per inducció sobre un dels arguments amb
els altres com a paràmetres, és a dir, pel
corol·lari~\ref{cor:induccio-parametres}; escrivim els arguments com a
elements generalitzats \(k,m,n\colon X\to\N\). La definició dona
\(n+0=n\) directament; provem primer les dues igualtats auxiliars
\[
0+n=n,
\qquad
(s\circ m)+n=s\circ(m+n),
\]
i després en deduïm la neutralitat, l'associativitat i la commutativitat.

\emph{Neutralitat per l'esquerra \(0+n=n\).} La part de \(\N\) on
\(0+n=n\) conté el zero, perquè \(0+0=0\); i és tancada pel successor,
perquè si \(0+n=n\) aleshores \(0+(s\circ n)=s\circ(0+n)=s\circ n\). Per
inducció, val a tot \(\N\).

\emph{Segona igualtat auxiliar \((s\circ m)+n=s\circ(m+n)\).} Per inducció
sobre \(n\). Per a \(n=0\): \((s\circ m)+0=s\circ m=s\circ(m+0)\). Pas de
\(n\) a \(s\circ n\): \((s\circ m)+(s\circ n)=s\circ((s\circ m)+n)
=s\circ s\circ(m+n)=s\circ(m+(s\circ n))\), on la primera i la darrera
igualtats són la definició de la suma i la central és la hipòtesi
d'inducció.

\emph{Associativitat.} Per inducció sobre \(n\). Per a \(n=0\):
\((k+m)+0=k+m=k+(m+0)\). Pas de \(n\) a \(s\circ n\):
\((k+m)+(s\circ n)=s\circ((k+m)+n)=s\circ(k+(m+n))=k+(m+(s\circ n))\),
usant la definició de la suma i la hipòtesi d'inducció.

\emph{Commutativitat.} Per inducció sobre \(n\), amb les dues igualtats
auxiliars. Per a \(n=0\): \(m+0=m=0+m\). Pas de \(n\) a \(s\circ n\):
\(m+(s\circ n)=s\circ(m+n)=s\circ(n+m)=(s\circ n)+m\), on la primera
igualtat és la definició, la segona la hipòtesi d'inducció i la tercera la
segona igualtat auxiliar.

\emph{Cancel·lació.} Aquí la propietat que es propaga no és una igualtat
entre dues aplicacions sinó una implicació quantificada,
\(\forall k\,\forall l\,(m+k=m+l\Rightarrow k=l)\); la part de \(\N\) que
defineix existeix pels quantificadors interns i la implicació booleana
(capítol~\ref{cap:parts-interna}), i la seva pertinença es llegeix, per la
proposició~\ref{prop:quantificadors-lectura}(3), així: \(m\colon X\to\N\) hi
pertany si per a tot \(y\colon Y\to X\) i tots \(k,l\colon Y\to\N\), de
\(m\circ y+k=m\circ y+l\) se segueix \(k=l\). Els dos passos següents
comproven aquesta condició a qualsevol estadi, i el
teorema~\ref{teo:induccio} s'hi aplica. Per inducció sobre \(m\), doncs. Per a \(m=0\): de
\(0+k=0+l\) i la neutralitat per l'esquerra surt \(k=l\). Pas de \(m\) a
\(s\circ m\): si \((s\circ m)+k=(s\circ m)+l\), aleshores per la segona
igualtat auxiliar \(s\circ(m+k)=s\circ(m+l)\); la injectivitat del
successor (teorema~\ref{teo:peano}(1)) dona
\(m+k=m+l\), i la hipòtesi d'inducció \(k=l\).

\emph{Suma nul·la.} Aquesta afirmació la provem per a punts globals
\(k,l\colon\mathbf1\to\N\), que és com es farà servir. Pel
lema~\ref{lema:zero-o-successor}, \(l=0\) o \(l=s\circ l'\) per a algun
punt \(l'\). En el segon cas,
\(k+l=k+(s\circ l')=s\circ(k+l')\ne0\), perquè el zero no és cap successor
(teorema~\ref{teo:peano}(2)); per tant \(k+l=0\) força \(l=0\), i llavors
\(k=k+0=0\).
\end{prova}

\begin{definicio}[Producte i exponenciació]\index{Producte i exponenciació}\label{def:producte-exp}
El \emph{producte} \(\cdot\colon\N\times\N\to\N\) es defineix per recursió
sobre el segon argument iterant la suma per \(m\):
\[
m\cdot 0=0,
\qquad
m\cdot(s\circ n)=(m\cdot n)+m .
\]
L'\emph{exponenciació} \(\N\times\N\to\N\), \((m,n)\mapsto m^{n}\), es
defineix iterant el producte per \(m\):
\[
m^{0}=s\circ 0\;(=1),
\qquad
m^{\,s\circ n}=(m^{n})\cdot m .
\]
\end{definicio}

\begin{proposicio}[Les lleis de l'aritmètica]\label{prop:lleis-aritmetica}
Per a tots \(k,m,n\colon X\to\N\), el producte és associatiu, commutatiu,
distributiu sobre la suma i té \(1:=s\circ 0\) com a element neutre, i
l'exponenciació satisfà les lleis dels exponents:
\[
m\cdot(n+k)=m\cdot n+m\cdot k,
\qquad
m^{\,n+k}=m^{n}\cdot m^{k},
\qquad
m^{\,n\cdot k}=(m^{n})^{k}.
\]
\end{proposicio}

\begin{prova}
Demostrem la distributivitat per l'esquerra
\(m\cdot(n+k)=m\cdot n+m\cdot k\) com a model; és la llei de la qual
depenen l'associativitat i la commutativitat del producte, i la primera
llei dels exponents. Fem inducció sobre \(k\)
(teorema~\ref{teo:induccio}), amb \(m,n\) fixats, usant la unicitat per
inducció (proposició~\ref{prop:unicitat-induccio}) per garantir que la
part on val la igualtat és tot \(\N\).

\emph{Cas base \(k=0\).} Per la definició de suma i de producte
(definicions~\ref{def:suma-naturals} i~\ref{def:producte-exp}),
\(m\cdot(n+0)=m\cdot n=m\cdot n+0=m\cdot n+m\cdot 0\).

\emph{Pas de \(k\) a \(s\circ k\).}
\[
m\cdot\bigl(n+(s\circ k)\bigr)
=m\cdot\bigl(s\circ(n+k)\bigr)
=m\cdot(n+k)+m
=(m\cdot n+m\cdot k)+m,
\]
on la primera igualtat és la definició de suma, la segona la definició de
producte, i la tercera la hipòtesi d'inducció; i, reordenant amb
l'associativitat de la suma (proposició~\ref{prop:suma}),
\[
(m\cdot n+m\cdot k)+m=m\cdot n+(m\cdot k+m)=m\cdot n+m\cdot(s\circ k),
\]
usant altre cop la definició de producte. Per inducció, la igualtat val a
tot \(\N\).

Les altres lleis es demostren totes per inducció seguint aquest mateix
esquema ---inducció sobre un argument, amb la distributivitat i la
proposició~\ref{prop:suma} com a lemes---, exactament com la prova
anterior i com la de la proposició~\ref{prop:suma}. Cap no requereix cap
axioma nou: totes són conseqüència de l'axioma~\ref{ax:naturals}, de la
definició recursiva de les operacions (definicions~\ref{def:suma-naturals}
i~\ref{def:producte-exp}) i de les propietats de la suma ja demostrades
(proposició~\ref{prop:suma}).
\end{prova}

\begin{observacio}[Lectura clàssica]\label{obs:lleis-classica}
L'estructura \((\N,0,s,+,\cdot)\) és el semianell dels nombres naturals:
les lleis de la proposició~\ref{prop:lleis-aritmetica} són precisament els
axiomes de semianell commutatiu, ara \emph{teoremes} d'inducció en lloc de
postulats. Les lleis dels exponents \(m^{p+q}=m^{p}\cdot m^{q}\) i
\(m^{p\cdot q}=(m^{p})^{q}\) són les que, al capítol següent, donaran els
recomptes \(|B^{A}|=m^{n}\) i la seva compatibilitat amb la suma i el
producte de cardinals.
\end{observacio}

\section{L'ordre dels naturals}\label{sec:ordre-naturals}

Fins ara hem tractat \(\N\) com un objecte algebraic ---amb el zero, el
successor i les operacions---, però encara no hem construït la relació
d'ordre que en fa una successió \(0<1<2<\dots\). La definim aquí, a partir
de la suma, i en derivem la propietat que el capítol següent farà servir
constantment: \(\N\) és \emph{ben ordenat}, és a dir, tota part no buida en
té un mínim. Com totes les relacions del llibre, l'ordre no és una fórmula
externa sinó una \emph{part} d'un producte
(capítol~\ref{cap:relacions}).

\begin{observacio}[Com es demostren les propietats de l'ordre]\label{obs:ordre-sobre-punts}
Les propietats de l'ordre que segueixen ---totalitat, bon ordre--- són
afirmacions sobre \emph{parts} de \(\N\) o de \(\N\times\N\): per exemple,
la totalitat diu que la unió \(\operatorname{Leq}\cup\operatorname{Inv}(\operatorname{Leq})\)
és la part total de \(\N\times\N\). Per la separació global i el lema del
buit, dues parts d'un mateix conjunt amb els mateixos punts globals són
equivalents (capítol~\ref{cap:relacions}, pràctica), de manera que n'hi ha
prou amb comprovar aquestes afirmacions \emph{sobre punts globals}; i
sobre punts, la pertinença a una unió és \enquote{o} i la pertinença a
una imatge és \enquote{existeix} (capítol~\ref{cap:valors},
corol·lari~\ref{cor:elements-de-2}; capítol~\ref{cap:imatges}). Per això,
en aquesta secció, \(m,n,k\) denotaran punts globals de \(\N\) llevat que
es digui el contrari. Les igualtats aritmètiques de la secció anterior,
en canvi, són igualtats d'aplicacions i s'han demostrat sobre elements
generalitzats.
\end{observacio}

\begin{definicio}[Ordre dels naturals]\index{Ordre dels naturals}\label{def:ordre-naturals}
La relació \emph{menor o igual} sobre \(\N\) és la part
\(\operatorname{Leq}\hookrightarrow\N\times\N\) donada per la imatge
(teorema~\ref{teo:factoritzacio}) de l'aplicació
\[
\langle\pi_1,+\rangle\colon\N\times\N\longrightarrow\N\times\N,
\qquad
(m,k)\longmapsto(m,\,m+k).
\]
Escriurem \(m\le n\) per «\(\langle m,n\rangle\) es factoritza per
\(\operatorname{Leq}\)», és a dir, per la pertinença generalitzada
\(\langle m,n\rangle\in\operatorname{Leq}\). Sobre punts globals, per
l'exhaustivitat dels epimorfismes sobre punts,
\[
m\le n
\quad\Longleftrightarrow\quad
\exists k\colon\mathbf1\to\N\ \text{ amb }\ m+k=n
\]
(sobre elements generalitzats, el \(k\) existeix després d'un
recobriment, com sempre amb les imatges).
L'ordre \emph{estricte} és la part
\(\operatorname{Lt}\hookrightarrow\N\times\N\) amb
\(m<n\Leftrightarrow(m\le n\wedge m\ne n)\), equivalentment
\(\exists k\colon m+(s\circ k)=n\).
\end{definicio}

\begin{observacio}[L'ordre és una part, no un esquema]\label{obs:ordre-part}
La definició no introdueix cap primitiva nova: \(\operatorname{Leq}\) és
una part concreta de \(\N\times\N\), obtinguda per factorització imatge
d'una aplicació ja disponible, i \(m\le n\) n'és la pertinença
generalitzada en el sentit del capítol~\ref{cap:relacions}. El
quantificador \(\exists k\) de la lectura equival, internament, a
travessar aquesta imatge; no és una quantificació metateòrica. Aquesta és
la diferència amb els «comentaris externs» sobre finitud que fèiem abans
de tenir \(\N\): l'ordre, ara, viu dins la teoria.
\end{observacio}

\begin{proposicio}[L'ordre és un ordre total]\label{prop:ordre-total}
La relació \(\le\) sobre \(\N\) és reflexiva, antisimètrica, transitiva i
total; és a dir, per a tots els punts \(k,m,n\colon\mathbf1\to\N\):
\begin{enumerate}
\item \(n\le n\) per a tot \(n\);
\item si \(m\le n\) i \(n\le m\), aleshores \(m=n\);
\item si \(k\le m\) i \(m\le n\), aleshores \(k\le n\);
\item per a tots \(m,n\), \(m\le n\) o bé \(n\le m\).
\end{enumerate}
A més, \(0\le n\) sempre, i \(m\le n\) si i només si \(s\circ m\le s\circ
n\).
\end{proposicio}

\begin{prova}
\emph{Reflexivitat.} \(n+0=n\) (definició~\ref{def:suma-naturals}), d'on
\(n\le n\) prenent \(k=0\).

\emph{Transitivitat.} Si \(k+a=m\) i \(m+b=n\), aleshores
\(k+(a+b)=(k+a)+b=m+b=n\) per l'associativitat de la suma
(proposició~\ref{prop:suma}), d'on \(k\le n\) prenent l'increment
\(a+b\).

\emph{\(0\le n\) i el pas del successor.} Com que \(0+n=n\)
(proposició~\ref{prop:suma}), \(0\le n\) per a tot \(n\). I
\((s\circ m)+k=s\circ(m+k)\) per la igualtat auxiliar de la
proposició~\ref{prop:suma}, i \(s\circ(m+k)=s\circ n\) equival a \(m+k=n\)
per la injectivitat del successor (teorema~\ref{teo:peano}(1)); així
\(m+k=n\Leftrightarrow(s\circ m)+k=s\circ n\), i per tant \(m\le n\Leftrightarrow s\circ m\le
s\circ n\).

\emph{Antisimetria.} Suposem \(m+k=n\) i \(n+l=m\). Substituint,
\(m+(k+l)=(m+k)+l=n+l=m=m+0\), i per la cancel·lació de la suma
(proposició~\ref{prop:suma}) \(k+l=0\); com que una suma nul·la té tots dos
sumands nuls (proposició~\ref{prop:suma}), \(k=l=0\), d'on \(n=m+k=m\).

\emph{Totalitat.} Fixat el punt \(m\), sigui \(T_m\hookrightarrow\N\) la
part dels \(n\) amb \(m\le n\) o \(n\le m\): és la imatge inversa, per
\(\langle m\circ{!},\id\rangle\colon\N\to\N\times\N\), de la unió
\(\operatorname{Leq}\cup\operatorname{Inv}(\operatorname{Leq})\). Vegem
per inducció (teorema~\ref{teo:induccio}) que \(T_m\) és total, i
aleshores tot punt \(n\) hi pertany. Per a \(n=0\): \(0\le m\) sempre, així que
\(m\le0\vee0\le m\). Pas de \(n\) a \(s\circ n\): sigui \(n\) un punt de \(T_m\), de manera que
\(m\le n\) o \(n\le m\) (pertinença d'un punt a una unió). Si \(m\le n\), aleshores \(m\le n\le
s\circ n\) per transitivitat (car \(n\le s\circ n\), ja que
\(n+s\circ0=s\circ n\)). Si \(n\le m\), distingim: o bé \(n=m\), i
aleshores \(m=n\le s\circ n\); o bé \(n<m\), i aleshores \(s\circ n\le m\)
perquè de \(n+(s\circ k)=m\) surt \(s\circ n+k=m\) amb el pas del
successor. En tots els casos \(m\le s\circ n\) o \(s\circ n\le m\), és a dir,
\(s\circ n\in T_m\). Hem comprovat el cas base i el pas sobre punts;
com que \(T_m\) és una part de \(\N\) i les parts es determinen pels
punts (observació~\ref{obs:ordre-sobre-punts}), la part dels \(n\) amb
\(n\in T_m\Rightarrow s\circ n\in T_m\) és total, i el
teorema~\ref{teo:induccio} s'aplica.
\end{prova}

\begin{teorema}[\(\N\) és ben ordenat]\label{teo:bon-ordre-N}
Tota part no inicial \(P\hookrightarrow\N\) té un \emph{element mínim}: un
punt \(m_0\colon\mathbf1\to\N\) que es factoritza per \(P\) i compleix
\(m_0\le k\) per a tot element generalitzat \(k\) de \(P\). El mínim, quan
existeix, és únic.
\end{teorema}

\begin{prova}
La unicitat és immediata per l'antisimetria
(proposició~\ref{prop:ordre-total}): dos mínims \(m_0,m_0'\) compleixen
\(m_0\le m_0'\) i \(m_0'\le m_0\), d'on \(m_0=m_0'\).

Per a l'existència provem el contrarecíproc: si \(P\) no té element mínim,
aleshores \(P\) és inicial. Tot el raonament és sobre punts globals, en el
sentit de l'observació~\ref{obs:ordre-sobre-punts}: les parts de \(\N\)
queden determinades pels seus punts (separació global i lema del buit),
de manera que una part conté el zero, o és tancada pel successor, o és
total, si i només si ho és sobre punts; i sobre punts, \enquote{o},
\enquote{existeix} i la pertinença a una imatge tenen el sentit clàssic.
Introduïm la part auxiliar
\(R\hookrightarrow\N\) classificada per la fórmula interna
\[
R(n)\ :\Longleftrightarrow\ \forall k\ (k\le n\Rightarrow \neg(k\mathrel{\varepsilon}P)),
\]
és a dir, «cap natural \(\le n\) no pertany a \(P\)». Aquesta part existeix
pels quantificadors interns i les operacions booleanes sobre \(P(\N)\)
(capítol~\ref{cap:parts-interna}), i el quantificador acotat \(\forall
k(k\le n\Rightarrow\cdots)\) és intern perquè \(\le\) és una part
(definició~\ref{def:ordre-naturals}).

Vegem que \(R=\N\) per inducció (teorema~\ref{teo:induccio}), sota la
hipòtesi que \(P\) no té mínim.

\emph{Cas base.} \(R(0)\) diu que cap \(k\le0\) no és a \(P\); com que
\(k\le0\) força \(k=0\) (de \(k+a=0\) surt \(k=0\) perquè una suma nul·la té els
dos sumands nuls, proposició~\ref{prop:suma}), \(R(0)\) equival a \(0\notin P\). Ara
bé, si \(0\in P\), aleshores \(0\) seria el mínim de \(P\), perquè \(0\le
k\) per a tot \(k\) (proposició~\ref{prop:ordre-total}); com que \(P\) no
té mínim, \(0\notin P\), i doncs \(R(0)\).

\emph{Pas de \(n\) a \(s\circ n\).} Suposem \(R(n)\): cap \(k\le n\) no és
a \(P\). Els naturals \(\le s\circ n\) són els \(\le n\) i el mateix
\(s\circ n\) ---perquè \(k\le s\circ n\) i \(k\ne s\circ n\) impliquen
\(k\le n\) pel pas del successor (proposició~\ref{prop:ordre-total}). Per
la hipòtesi \(R(n)\), cap dels primers no és a \(P\); resta veure que
\(s\circ n\notin P\). Però si \(s\circ n\in P\), aleshores \(s\circ n\)
seria el \emph{mínim} de \(P\): per a tot \(j\in P\), com que cap \(k\le
n\) no és a \(P\) (per \(R(n)\)), ha de ser \(n<j\), és a dir \(s\circ
n\le j\) (pas del successor). Això contradiu que \(P\) no tingui mínim.
Per tant \(s\circ n\notin P\), i \(R(s\circ n)\).

Hem comprovat sobre punts que \(R\) conté el zero i que tot punt de \(R\)
té el successor a \(R\); això vol dir que la part \(R\cap s^{-1}(R^c)\)
no té punts i és buida, és a dir, \(s\circ r\) es factoritza per \(R\),
que és la hipòtesi del teorema~\ref{teo:induccio}. Per inducció,
\(R=\N\). Però \(R(n)\) per a tot punt \(n\) diu, en particular prenent
\(k=n\), que cap punt de \(\N\) no pertany a \(P\); per tant \(P\) no té
punts i, pel lema del buit (lema~\ref{lema:buit}), és inicial. Això prova
el contrarecíproc, i per tant tota part no inicial de \(\N\) té un mínim.
Finalment, que \(m_0\) sigui mínim per a tot element generalitzat \(k\) de
\(P\), i no només per als punts, és la inclusió de parts
\(P\subseteq\{k:m_0\le k\}\), que de nou es comprova sobre punts.
\end{prova}

\begin{observacio}[Lectura clàssica i ús posterior]\label{obs:bon-ordre-classica}
El teorema~\ref{teo:bon-ordre-N} és el \emph{principi del bon ordre} dels
naturals, equivalent en la matemàtica clàssica al principi d'inducció: tota
col·lecció no buida de naturals té un mínim. El farem servir de manera
essencial al capítol següent ---per exemple, per triar canònicament un
element de cada part no buida de \(\N\) sense recórrer a cap axioma
d'elecció (el mínim és únic, i per tant una tria \emph{definida})--- i és
la base de tota la comparació de segments inicials que sosté la noció de
finitud. La restricció a \(\N\) és deliberada: la noció general de conjunt
\emph{ben ordenat} i el seu paper en la teoria de l'ordre es
desenvoluparan més endavant, quan calguin per als ordinals.
\end{observacio}

\begin{definicio}[Successió]\index{Successió}\label{def:successio}
Una \emph{successió} d'elements d'un conjunt \(A\) és una aplicació
\(a\colon\N\to A\). El conjunt de totes les successions d'\(A\) és
l'exponencial \(A^{\N}\). Un element global de \(A^{\N}\) és una successió
en el sentit clàssic \((a_0,a_1,a_2,\dots)\), amb \(a_n:=a\circ n\); un
element generalitzat \(X\to A^{\N}\) és una família de successions
parametritzada per \(X\), o equivalentment una aplicació
\(X\times\N\to A\) (proposició~\ref{prop:families}).
\end{definicio}

La recursió de l'axioma~\ref{ax:naturals} defineix aplicacions
\(\N\to A\) a partir de dades \emph{globals} \((a,t)\). Sovint, però, cal
definir una successió que depèn d'un paràmetre \(x\) que recorre un estadi
\(X\), amb un valor inicial \(a(x)\) i una regla \(t(x,-)\) que també en
depèn. Els exponencials permeten reduir aquest cas al de
l'axioma~\ref{ax:naturals}.

\begin{teorema}[Recursió amb paràmetres]\label{teo:recursio-parametres}
Siguin \(A\) i \(X\) conjunts, \(a\colon X\to A\) una aplicació i
\(t\colon X\times A\to A\) una aplicació. Aleshores existeix una única
aplicació
\[
u\colon X\times\N\to A
\]
tal que, per a tot \(x\colon Y\to X\) i tot \(n\colon Y\to\N\),
\[
u\circ\langle x,0\circ{!}_Y\rangle=a\circ x
\qquad\text{i}\qquad
u\circ\langle x,s\circ n\rangle=t\circ\langle x,\,u\circ\langle
x,n\rangle\rangle .
\]
\end{teorema}

\begin{prova}
Transposem respecte de \(X\). Pel capítol~\ref{cap:exponencials}, una
aplicació \(u\colon X\times\N\to A\) equival a la seva transposada
\(\widehat u\colon\N\to A^{X}\), i les dues condicions sobre \(u\) es
tradueixen, per la unicitat de la transposada, en condicions sobre
\(\widehat u\). El valor inicial \(a\colon X\to A\) té nom
\(\ulcorner a\urcorner\colon\mathbf1\to A^{X}\), i la regla
\(t\colon X\times A\to A\) determina una autoaplicació
\(T\colon A^{X}\to A^{X}\): la transposada de
\(t\circ\langle\pi_X,\operatorname{ev}\rangle\colon A^X\times X\to A\),
que sobre punts envia \(g\) a \(x\mapsto t(x,g(x))\). Aplicant
l'axioma~\ref{ax:naturals} al sistema recursiu
\((A^{X},\ulcorner a\urcorner,T)\) s'obté una única
\(\widehat u\colon\N\to A^{X}\) amb
\(\widehat u\circ 0=\ulcorner a\urcorner\) i
\(\widehat u\circ s=T\circ\widehat u\). La seva descurrificada \(u\) és
l'aplicació buscada, i és única perquè ho és \(\widehat u\).
\end{prova}

\begin{observacio}[El paper dels exponencials]\label{obs:exponencials-recursio}
El teorema~\ref{teo:recursio-parametres} és la raó tècnica per la qual
convé haver introduït els exponencials abans que \(\N\). L'axioma de
recursió, tal com s'enuncia, només parla de dades globals sobre un
conjunt \(A\); la recursió amb paràmetres, que és la que fa falta a la
pràctica, s'hi redueix pujant a l'exponencial \(A^{X}\) i tornant a baixar
per descurrificació. Sense exponencials caldria un axioma de recursió més
fort, amb paràmetres incorporats; amb ells, la forma simple de
l'axioma~\ref{ax:naturals} és suficient. Aquesta economia ---un axioma
mínim que la resta de l'estructura amplifica--- és la mateixa que hem vist
amb el producte i el conjunt de parts.
\end{observacio}

\begin{observacio}[Lectura clàssica]\label{obs:recursio-parametres-classica}
La recursió amb paràmetres és la forma en què es defineixen en la
pràctica la major part de les successions: el terme \(u(x,n+1)\) es calcula
a partir de \(u(x,n)\) i del paràmetre \(x\), com en la definició de la
successió de Fibonacci, dels factorials, o de qualsevol iteració amb dades
d'entrada. El teorema garanteix que totes aquestes definicions produeixen
funcions ben definides, i el fa reduint-les al teorema de recursió simple
a través del conjunt de funcions \(A^{X}\).
\end{observacio}

\begin{resumcap}
\apartat{Idees centrals}
\begin{enumerate}
  \item L'objecte dels nombres naturals es postula per una sola propietat universal, la recursió: per a tot sistema recursiu \((A,a,t)\) hi ha una única \(u\colon\N\to A\) amb \(u\circ0=a\) i \(u\circ s=t\circ u\). És únic llevat d'un únic isomorfisme.
  \item La unicitat de la recursora és el principi d'inducció; l'existència és la definició per recursió. Els exponencials estenen totes dues coses als paràmetres: recursió amb paràmetres i inducció amb paràmetres, per transposició.
  \item Les equacions de Peano són teoremes. El successor és injectiu perquè té una retracció, el predecessor, que no és una recursió simple sinó la primera projecció d'una recursió sobre \(\N\times\N\). Tot punt és zero o successor.
  \item Suma, producte i exponenciació es defineixen iterant, i les seves lleis es demostren per inducció amb paràmetres, com a igualtats d'aplicacions \(\N^k\to\N\).
  \item L'ordre \(\le\) és una part de \(\N\times\N\), la imatge de \((m,k)\mapsto(m,m+k)\). Les seves propietats (total, bon ordre) són afirmacions sobre parts, i per la separació global es comproven sobre punts globals, on \enquote{o} i \enquote{existeix} tenen el sentit clàssic.
  \item \(\N\) és ben ordenat: tota part no buida té un mínim, únic i, per tant, definible sense elecció.
\end{enumerate}
\apartat{Errors típics} Definir el predecessor com la recursora de \((\N,0,\id)\), que és constant; fer inducció sobre \(n\) amb un paràmetre \(m\) generalitzat sense passar per l'exponencial (la part on dues aplicacions \(X\times\N\to A\) coincideixen no és una part de \(\N\)); llegir \enquote{\(m\le n\) o \(n\le m\)} sobre elements generalitzats, quan només val sobre punts; prendre les cinc propietats de Peano com a axiomes, quan aquí són conseqüències; i oblidar que la inducció del llibre quantifica sobre parts de \(\N\), no sobre fórmules, cosa que la fa de \enquote{segon ordre} i dona la unicitat de \(\N\).
\apartat{Lectura recomanada} \cite[§§7.3--7.5]{lawvere-rosebrugh} per a l'objecte dels nombres naturals, la recursió i la inducció en conjunts abstractes; \cite[sessió~26 i article~V]{lawvere-schanuel} per a una presentació elemental dels sistemes dinàmics discrets i de la recursió; \cite[§§11--13]{halmos} per als axiomes de Peano, el teorema de recursió i l'aritmètica en la teoria clàssica, que és el que aquí s'ha derivat; \cite[§4.7]{maclane-moerdijk} per a l'objecte de nombres naturals en un topos i la recursió amb paràmetres; i \cite[pàg.~12]{bell2006} per a la definició de conjunt infinit de Dedekind i l'axioma d'infinitud.
\end{resumcap}

\chapter{Conjunts finits i cardinalitat finita}\label{cap:finits}

\begin{objectius}
\apartat{Prerequisits} Del capítol~\ref{cap:naturals}: l'objecte \(\N\), la inducció (amb i sense paràmetres), les equacions de Peano, l'aritmètica i l'ordre, amb la convenció que les propietats de l'ordre es comproven sobre punts. Del capítol~\ref{cap:valors}: sumes disjuntes, complements i distributivitat. Dels capítols~\ref{cap:imatges} i~\ref{cap:exponencials}: factorització imatge, exponencials i composició interna. L'axioma d'elecció només apareix en un punt, assenyalat.
\apartat{Objectius} En acabar el capítol, el lector hauria de poder:
\begin{itemize}
  \item distingir els ordinals finits \([k]\), definits al metanivell per a numerals estàndard, dels segments inicials \(\seg{n}\), que formen una família interna, i explicar per què la finitud es defineix amb els segons;
  \item construir el conjunt de cardinals \(\operatorname{Card}(A)\hookrightarrow\N\) codificant les bijeccions \(\seg{n}\to A\) dins un exponencial fix, i definir finitud i cardinal a partir d'ell;
  \item demostrar, per inducció interna sobre un predicat codificat, que un monomorfisme \(\seg{m}\hookrightarrow\seg{n}\) força \(m\le n\), i deduir-ne la unicitat del cardinal;
  \item demostrar el principi dels calaixos i l'equivalència entre mono, epi i iso en cardinal igual, i saber on intervé l'elecció;
  \item demostrar els recomptes \(|A+B|=|A|+|B|\), \(|A\times B|=|A|\cdot|B|\), \(|P(A)|=2^{|A|}\) i \(|B^A|=|B|^{|A|}\) per inducció.
\end{itemize}
\end{objectius}

Al llarg del llibre, la paraula \emph{finit} ha aparegut només com a
comentari extern informal: dèiem que una construcció produïa «un nombre
finit» de casos o que un diagrama tenia «finites» fletxes, sense que
aquesta finitud fos res precís. Ara que disposem de l'objecte dels nombres
naturals (capítol~\ref{cap:naturals}), podem convertir-la en una noció
\emph{precisa} i \emph{interna}: un conjunt \(A\) serà finit quan existeixi
un punt \(n\colon\mathbf1\to\N\) amb \(\seg{n}\approx A\), és a dir, quan
\(A\) sigui equipotent a algun dels segments inicials
(definició~\ref{def:segment-inicial}) que el successor genera. Aquesta és
una propietat interna, expressada dins el llenguatge mitjançant la família
uniforme dels segments \(n\mapsto\seg{n}\), i no requereix cap axioma nou.
Els ordinals finits \([k]\) que introduïm tot seguit
(definició~\ref{def:ordinal-finit}) no en són la definició, sinó només la
\emph{lectura externa} per als numerals estàndard
(observació~\ref{obs:ordinals-metanivell}): una notació de càlcul còmoda,
lligada als segments per la coincidència de
l'observació~\ref{obs:segment-numeral}. El capítol és, més aviat, una
primera aplicació extensa de tot el que s'ha acumulat ---sumes,
exponencials, conjunt de parts, factorització imatge i inducció--- i el
pont cap als exercicis computacionals de la part pràctica.

\section{Els ordinals finits}\label{sec:ordinals-finits}

La successió \(\mathbf1,\;\mathbf1+\mathbf1,\;\mathbf1+\mathbf1+\mathbf1,
\dots\) no és, de moment, res més que una família de conjunts descrita amb
punts suspensius. Comencem donant-li un nom precís \emph{al metanivell},
una recursió externa sobre els numerals estàndard; tot seguit
(secció~\ref{sec:segments-interns}) en construirem la versió \emph{interna}
i uniforme, els segments inicials, que és la que sostindrà tota la teoria.

\begin{definicio}[Ordinal finit]\index{Ordinal finit}\label{def:ordinal-finit}
Per a cada natural estàndard \(k\) del metallenguatge definim
\emph{externament} un conjunt \([k]\), l'\emph{ordinal finit}
associat a \(k\), per la recursió metateòrica
\[
[0]:=\mathbf0,
\qquad
[k+1]:=[k]+\mathbf1,
\]
on \(\mathbf0\) és el conjunt inicial (capítol~\ref{cap:aplicacions}),
\(+\) la suma de conjunts (definició~\ref{def:suma-conjunts}) i
\(\mathbf1\) el conjunt terminal. Aquesta notació assigna un conjunt
\([k]\), concret llevat d'isomorfisme, a cada numeral estàndard, un per a
cada \(k\) del metallenguatge; \emph{no} s'estén a un punt intern arbitrari
\(n\colon\mathbf1\to\N\). Per al pas de \([k]\) a \([k+1]\) escriurem
\(\iota_k\colon[k]\to[k+1]\) per a la injecció canònica \(i_0\) de la suma
\([k]+\mathbf1\), i \(\ast_k\colon\mathbf1\to[k+1]\) per a la injecció
\(i_1\) del sumand nou.
\end{definicio}

\begin{observacio}[Sobre l'estatut de la recursió d'ordinals]\label{obs:ordinals-metanivell}
La regla anterior és una recursió al metanivell, no una aplicació interna
cap a un «conjunt de tots els conjunts», que no existeix. A cada numeral
estàndard \(\bar k\colon\mathbf1\to\N\) ---el zero, el successor del zero,
i així successivament--- li correspon el conjunt \([k]\), concret llevat
d'isomorfisme, un per a cada \(k\) del metallenguatge. No pretenem que tot
punt global d'un \(\N\) intern sigui d'aquesta forma: en una teoria de
primer ordre no es pot identificar sense més una variable interna
\(n\colon\mathbf1\to\N\) amb un numeral construït externament, i per tant la
notació \([\,\cdot\,]\) no s'estén a un \(n\) intern arbitrari.

Per aquesta raó, la finitud i la cardinalitat internes \emph{no} es
basaran en la família metanivell \(k\mapsto[k]\), sinó en la família
\emph{interna} dels segments inicials \(n\mapsto\seg{n}\)
(definició~\ref{def:segment-inicial}), que prové d'un sol objecte, la
relació \(\operatorname{Lt}\hookrightarrow\N\times\N\). Els ordinals
\([k]\) es conservaran només com a \emph{notació de càlcul} per als
numerals estàndard, lligada als segments per la coincidència externa de
l'observació~\ref{obs:segment-numeral}.
\end{observacio}

Convé fixar de seguida les mides petites, que reapareixeran a cada
inducció.

\begin{proposicio}[Els primers ordinals]\label{prop:ordinals-petits}
Es tenen isomorfismes canònics \([1]\cong\mathbf1\) i
\([2]\cong\mathbf1+\mathbf1\cong\mathbf2\). A més, per a cada numeral
estàndard \(k\) la injecció \(\iota_k\colon[k]\hookrightarrow[k+1]\) és un
monomorfisme, i \(\ast_k\) és el seu complement dins \([k+1]\).
\end{proposicio}

\begin{prova}
El cas \(k=0\) de la definició dona \([1]=\mathbf0+\mathbf1\cong\mathbf1\)
per la neutralitat del conjunt inicial respecte de la suma
(capítol~\ref{cap:valors}); aplicant-ho un cop més,
\([2]=[1]+\mathbf1\cong\mathbf1+\mathbf1\), que és el conjunt de valors de
veritat \(\mathbf2\). Que \(\iota_k\) és un monomorfisme amb complement
\(\ast_k\) és la disjunció de la suma (definició~\ref{def:suma-disjunta}):
en una suma disjunta les dues injeccions són monomorfismes que
representen parts complementàries. Com que aquestes afirmacions són un
esquema sobre numerals estàndard, cada instància es comprova externament.
\end{prova}

\section{Els segments inicials com a família interna}\label{sec:segments-interns}

Els ordinals \([k]\) de la definició~\ref{def:ordinal-finit} pateixen la
limitació que l'observació~\ref{obs:ordinals-metanivell} ja assenyalava:
la notació \(k\mapsto[k]\) només està definida per als numerals estàndard
del metallenguatge i no és, per si sola, una família interna. Ara que
disposem de l'ordre de
\(\N\) com a \emph{part} \(\operatorname{Lt}\hookrightarrow\N\times\N\)
(definició~\ref{def:ordre-naturals}), podem donar una segona presentació
dels ordinals finits que \emph{sí} és interna i uniforme, i comprovar que
coincideix amb la primera. Aquesta presentació és la que farà de la
finitud una propietat interna.

\begin{definicio}[Segment inicial]\index{Segment inicial}\label{def:segment-inicial}
Per a cada punt \(n\colon\mathbf1\to\N\), el \emph{segment inicial}
\(\seg{n}\hookrightarrow\N\) és la fibra de la part
\(\operatorname{Lt}\hookrightarrow\N\times\N\) sobre \(n\): la part de
\(\N\) formada pels \(k\) amb \(k<n\), és a dir la imatge inversa
\[
\seg{n}:=\{\,k\colon\mathbf1\to\N \mid k<n\,\}
=\operatorname{Lt}\bigl|_{\N\times\{n\}} ,
\]
obtinguda com el producte fibrat de \(\operatorname{Lt}\hookrightarrow
\N\times\N\) al llarg de \(\langle\id_\N,n\circ{!}\rangle\colon\N\to
\N\times\N\) (definició~\ref{def:pullback}).
\end{definicio}

\begin{observacio}[La família de segments és interna i uniforme]\label{obs:segments-interns}
A diferència de \(k\mapsto[k]\), la regla \(n\mapsto\seg{n}\) \emph{sí} que és una única part interna: prové d'un sol
objecte, la relació \(\operatorname{Lt}\hookrightarrow\N\times\N\), de la
qual \(\seg{n}\) és la fibra sobre \(n\). Dit d'una altra
manera, \(\operatorname{Lt}\), vista com una aplicació
\(\operatorname{Lt}\to\N\) per la segona projecció, és la
\emph{família uniforme} \((\seg{n})_{n\in\N}\) de tots els
segments alhora, indexada internament per \(\N\)
(capítol~\ref{cap:relacions}). És aquesta uniformitat la que permetrà
quantificar sobre segments dins el llenguatge, cosa que amb \([k]\) sol no
era possible.
\end{observacio}

\begin{observacio}[Coincidència externa amb els numerals estàndard]\label{obs:segment-numeral}
Per a cada numeral estàndard \(k\) del metallenguatge, hi ha un isomorfisme
\(\seg{\bar k}\cong[k]\) entre el segment del numeral \(\bar
k\colon\mathbf1\to\N\) i l'ordinal finit \([k]\). Aquesta és una afirmació
\emph{externa}, un esquema amb una instància per a cada \(k\), i es
comprova per inducció metateòrica sobre \(k\):
\[
\seg{\bar 0}\cong\mathbf0=[0],
\qquad
\seg{\overline{k+1}}\cong\seg{\bar k}+\mathbf1\cong[k]+\mathbf1=[k+1],
\]
on la primera línia val perquè \(\{j:j<0\}\) és inicial, i el pas usa que
\(j<\overline{k+1}\) equival a \(j\le\bar k\), és a dir \(j<\bar k\) o
\(j=\bar k\) (proposició~\ref{prop:ordre-total}), cosa que descompon
\(\seg{\overline{k+1}}\) en la suma disjunta \(\seg{\bar k}+\mathbf1\).

No enunciem cap identificació \(\seg{n}\cong[n]\) per a un \(n\colon
\mathbf1\to\N\) \emph{intern} arbitrari: això requeriria una inducció
interna sobre un objecte, \([n]\), que no és intern. Els \([k]\) queden,
doncs, com a abreviatura còmoda per als numerals concrets, i tota
construcció interna del capítol es fa amb els segments \(\seg{n}\).
\end{observacio}

El que \emph{sí} és intern, i és tot el que la teoria de la finitud
necessita, és la manera com un segment creix en passar al successor. Ho
enunciem sobre una variable interna \(n\).

\begin{proposicio}[Descomposició successora del segment]\label{prop:segment-successor}
Per a tota variable interna \(n\colon\mathbf1\to\N\) hi ha un isomorfisme
\[
\seg{s\circ n}\ \cong\ \seg{n}+\mathbf1,
\]
en què el sumand \(\seg{n}\) és la part de \(\seg{s\circ n}\) dels
\(k<n\) i el sumand \(\mathbf1\) és el punt \(n\) mateix. A més
\(\seg{0}\cong\mathbf0\).
\end{proposicio}

\begin{prova}
\(\seg{0}=\{k:k<0\}\) és inicial, perquè \(k<0\) no es compleix per a cap
punt \(k\) (de \(k+s\circ j=0\) sortiria \(s\circ(k+j)=0\), contra el
fet que el zero no és cap successor, teorema~\ref{teo:peano}(2)).

Per al successor, un punt \(k\in\seg{s\circ n}\) compleix \(k<s\circ n\),
que equival a \(k\le n\), és a dir \(k<n\) o bé \(k=n\), i no totes dues
alhora perquè \(n\not<n\) (pas del successor de l'ordre,
proposició~\ref{prop:ordre-total}); com a tota la secció de l'ordre, el
raonament és sobre punts, i la separació global el converteix en una
igualtat de parts (observació~\ref{obs:ordre-sobre-punts}). Aquesta disjunció exclusiva
descompon la part \(\seg{s\circ n}\hookrightarrow\N\) en dues parts
complementàries: la dels \(k<n\), que és \(\seg{n}\), i la del punt únic
\(k=n\), que és una còpia de \(\mathbf1\). Una part i el seu complement
formen una suma disjunta (teorema~\ref{teo:part-complementari}), de manera
que \(\seg{s\circ n}\cong\seg{n}+\mathbf1\).
\end{prova}

\section{Equipotència i el conjunt de cardinals}\label{sec:equipotencia}

Per definir la finitud com una propietat \emph{interna} ---i, alhora, el
cardinal com una funció ben definida--- ens cal, per a un conjunt \(A\)
fixat, poder reunir en una part de \(\N\) tots els \(n\) tals que \(\seg{n}\)
és isomorf a \(A\). Aquesta part es construeix sense cap circularitat, ni
amb cap axioma nou, a partir de dues eines ja disponibles: el subconjunt
dels isomorfismes dins un exponencial, i la factorització imatge. La
construcció és el cas finit d'una noció més general ---l'\emph{equipotència}
de dos conjunts--- que reprendrem quan comparem cardinals infinits.

\begin{definicio}[El subconjunt dels isomorfismes]\index{El subconjunt dels isomorfismes}\label{def:iso-exponencial}
Siguin \(A\) i \(B\) dos conjunts. El \emph{subconjunt dels isomorfismes}
\(\operatorname{Isos}(A,B)\hookrightarrow B^{A}\) és la part de
l'exponencial \(B^{A}\) formada pels punts \(f\) que admeten una inversa
bilateral. Es construeix internament així. Per la composició interna
(proposició~\ref{prop:composicio-interna}) hi ha aplicacions
\[
c_{A}\colon A^{B}\times B^{A}\to A^{A},
\qquad
c_{B}\colon B^{A}\times A^{B}\to B^{B},
\]
que sobre punts envien \((g,f)\mapsto g\circ f\) i \((f,g)\mapsto f\circ
g\) respectivament, i identitats internes \(\jmath_A\colon\mathbf1\to
A^{A}\), \(\jmath_B\colon\mathbf1\to B^{B}\). Considerem l'aplicació
\[
\Phi\colon B^{A}\times A^{B}\longrightarrow A^{A}\times B^{B},
\qquad
(f,g)\longmapsto\bigl(c_A(g,f),\,c_B(f,g)\bigr)=(g\circ f,\ f\circ g),
\]
i l'equalitzador (definició~\ref{def:equalitzador}) de \(\Phi\) amb
l'aplicació constant de valor \((\jmath_A,\jmath_B)\): la part
\(P\hookrightarrow B^{A}\times A^{B}\) dels parells \((f,g)\) amb \(g\circ
f=\id_A\) i \(f\circ g=\id_B\). Aleshores \(\operatorname{Isos}(A,B)\) és la
imatge (teorema~\ref{teo:factoritzacio}) de \(P\) per la projecció a
\(B^{A}\): els \(f\) per als quals existeix un tal \(g\).
\end{definicio}

\begin{observacio}[Els punts de \(\operatorname{Isos}(A,B)\)]\label{obs:iso-punts}
Pels punts globals de l'exponencial (teorema~\ref{teo:punts-exponencial}),
un punt \(\mathbf1\to\operatorname{Isos}(A,B)\) és exactament una aplicació
\(f\colon A\to B\) que és un isomorfisme (definició~\ref{def:isomorfisme}).
Així, \(\operatorname{Isos}(A,B)\) és no inicial si i només si existeix un
isomorfisme \(A\cong B\). La part es defineix per a tot parell de conjunts
amb un únic bloc d'operacions internes, sense triar cap isomorfisme.
\end{observacio}

\begin{definicio}[Equipotència]\index{Equipotència}\label{def:equipotencia}
Dos conjunts \(A\) i \(B\) són \emph{equipotents}, i escrivim \(A\approx
B\), si \(\operatorname{Isos}(A,B)\) és no inicial ---equivalentment, si
existeix un isomorfisme \(A\cong B\). La relació \(\approx\) és reflexiva
(la identitat), simètrica (la inversa d'un isomorfisme) i transitiva (la
composició), és a dir, una relació d'equivalència. Cal entendre aquesta
afirmació \emph{metateòricament}: com que no hi ha cap conjunt de tots els
conjunts, \(\approx\) no és una part interna d'un producte, sinó una
relació entre conjunts formulada al metanivell; reflexivitat, simetria i
transitivitat s'hi comproven exhibint, en cada cas, l'isomorfisme
corresponent.
\end{definicio}

Ara fixem \(A\) i fem variar el segment. La família uniforme
\((\seg{n})_{n\in\N}\) (observació~\ref{obs:segments-interns}) permet
reunir en un sol objecte tots els isomorfismes \(\seg{n}\to A\), i
projectar-los sobre \(\N\) dona la part buscada.

\begin{definicio}[El conjunt de cardinals d'un conjunt]\index{El conjunt de cardinals d'un conjunt}\label{def:card-A}
Sigui \(A\) un conjunt. Cada segment \(\seg{n}\) és una part de \(\N\)
(definició~\ref{def:segment-inicial}), de manera que una aplicació
\(\varphi\colon\seg{n}\to A\) s'estén de manera única a una aplicació
\(\bar\varphi\colon\N\to A+\mathbf1\) que val \(\varphi\) sobre \(\seg{n}\)
i el punt afegit \(\ast\colon\mathbf1\to A+\mathbf1\) fora de \(\seg{n}\)
(definició per casos amb el complement, capítol~\ref{cap:valors}). Així
totes les aplicacions des de tots els segments viuen dins un \emph{únic}
exponencial fix, \((A+\mathbf1)^{\N}\), i podem reunir-les en un sol
objecte.

Definim una part d'un objecte \emph{fix}, sense exponencials de domini
variable:
\[
\mathcal E_A\ \hookrightarrow\ \N\times(A+\mathbf1)^{\N}\times\N^{A},
\]
els punts de la qual són els triples \((n,g,h)\) ---amb
\(g\colon\N\to A+\mathbf1\) i \(h\colon A\to\N\)--- que satisfan
internament les tres condicions següents, on \(i_A\colon A\to A+\mathbf1\)
i \(i_1\colon\mathbf1\to A+\mathbf1\) són les injeccions de la suma:
\begin{enumerate}
\item[(i)] per a tot \(k<n\), existeix (un únic) \(a\) amb \(g(k)=i_A(a)\)
i \(h(a)=k\);
\item[(ii)] per a tot \(a\in A\), \(h(a)<n\) i \(g(h(a))=i_A(a)\);
\item[(iii)] per a tot \(k\) amb \(k\not<n\), \(g(k)=i_1(\ast)\).
\end{enumerate}
Aquestes condicions només fan servir avaluacions dels exponencials
\emph{fixos} \((A+\mathbf1)^{\N}\) i \(\N^{A}\), l'ordre \(<\)
(definició~\ref{def:ordre-naturals}), igualtats i quantificadors acotats,
tots ja construïts; \(\mathcal E_A\) n'és la part corresponent, tallada per
equalitzadors i imatges inverses. Les condicions (i)--(ii) diuen que
\(g\) i \(h\) exhibeixen una bijecció entre \(\seg{n}\) i \(A\) (cadascuna
inversa de l'altra sobre el seu domini), i (iii) fixa \(g\) fora de
\(\seg{n}\), de manera que un tal triple existeix si i només si
\(\seg{n}\approx A\).

El \emph{conjunt de cardinals} de \(A\) és la imatge
(teorema~\ref{teo:factoritzacio}) de la projecció
\(\pi\colon\mathcal E_A\to\N\):
\[
\operatorname{Card}(A)\ :=\ \operatorname{Im}(\pi)\ =\
\{\,n\colon\mathbf1\to\N \mid \seg{n}\approx A\,\}
\ \hookrightarrow\ \N .
\]
\end{definicio}

\begin{proposicio}[\(\operatorname{Card}(A)\) és una part interna]\label{prop:card-interna}
Per a tot conjunt \(A\), \(\operatorname{Card}(A)\) és una part ben
definida de \(\N\), i un punt \(n\colon\mathbf1\to\N\) hi pertany si i
només si \(\seg{n}\approx A\). La construcció no fa servir cap axioma
d'elecció ni cap noció prèvia de finitud o cardinal.
\end{proposicio}

\begin{prova}
La part \(\mathcal E_A\hookrightarrow\N\times(A+\mathbf1)^{\N}\times\N^{A}\)
és interna per construcció (definició~\ref{def:card-A}): les condicions
(i)--(iii) són conjuncions d'equacions entre avaluacions d'exponencials
fixos i de condicions acotades per l'ordre \(<\), totes tallables amb
equalitzadors, imatges inverses i quantificadors acotats ja disponibles
(capítols~\ref{cap:exponencials} i~\ref{cap:parts-interna}). La imatge
d'una aplicació existeix sempre (teorema~\ref{teo:factoritzacio}), de
manera que \(\operatorname{Card}(A)=\operatorname{Im}(\pi)\) és una part de
\(\N\).

Un punt \(n\) es factoritza per aquesta imatge si i només si la fibra
\((\mathcal E_A)_n\) és no inicial (exhaustivitat de la coprojecció sobre
la imatge, capítol~\ref{cap:imatges}). Ara bé, un punt de \((\mathcal
E_A)_n\) és un parell \((g,h)\) que satisfà (i)--(iii); per (i) i (ii),
\(g\) restringida a \(\seg{n}\) i \(h\) són mútuament inverses, i
defineixen un isomorfisme \(\seg{n}\to A\), és a dir \(\seg{n}\approx A\).
Recíprocament, un isomorfisme \(\varphi\colon\seg{n}\to A\) dona el triple
\((n,\bar\varphi,\varphi^{-1})\), amb \(\bar\varphi\) l'extensió per
\(\ast\), que satisfà (i)--(iii). Així \(n\in\operatorname{Card}(A)\) si i
només si \(\seg{n}\approx A\). Cap pas no ha usat la finitud ni el
cardinal, sinó només exponencials fixos, l'ordre, equalitzadors i la
imatge.
\end{prova}

\begin{definicio}[Conjunt finit i cardinal]\index{Conjunt finit i cardinal}\label{def:conjunt-finit}
Un conjunt \(A\) és \emph{finit} si \(\operatorname{Card}(A)\) és no
inicial (definició~\ref{def:card-A}); és a dir, si existeix un punt
\(n\colon\mathbf1\to\N\) amb \(\seg{n}\approx A\). Quan \(A\) és finit,
provarem tot seguit que \(\operatorname{Card}(A)\) té un \emph{únic} punt
(teorema~\ref{teo:unicitat-cardinal}); aquest punt és el \emph{cardinal}
de \(A\), que denotarem \(|A|\colon\mathbf1\to\N\).
\end{definicio}

\begin{observacio}[La finitud és una propietat interna]\label{obs:finitud-propietat}
Aquesta definició \emph{sí} és interna, a diferència del que passaria amb
\([n]\). La clau és que \(\operatorname{Card}(A)\) s'ha construït
(definició~\ref{def:card-A}, proposició~\ref{prop:card-interna}) com una
part genuïna de \(\N\) ---la imatge d'una projecció---, sense recórrer a
cap noció prèvia de finitud ni de cardinal, sinó només a la família interna
de segments (observació~\ref{obs:segments-interns}), al subconjunt dels
isomorfismes d'un exponencial (definició~\ref{def:iso-exponencial}) i a la
factorització imatge (teorema~\ref{teo:factoritzacio}). Per tant «\(A\) és
finit» ---que \(\operatorname{Card}(A)\) sigui no inicial--- és una
propietat interna de \(A\), i no un esquema indexat per numerals com ho era
la caracterització amb \([n]\).

No hi ha circularitat entre finitud i cardinal: primer es construeix la
part \(\operatorname{Card}(A)\); després, «finit» n'afirma la no-inicialitat
i el «cardinal» n'és l'únic punt. Cap de les dues nocions no s'ha usat en
la construcció de \(\operatorname{Card}(A)\). Que aquest punt és únic ---i
per tant que \(|A|\) és una funció ben definida sobre els conjunts finits,
invariant per isomorfisme--- és el contingut del
teorema~\ref{teo:unicitat-cardinal}, que provem tot seguit.
\end{observacio}

\begin{observacio}[Relació amb la finitud estàndard]\label{obs:finitud-equivalencia}
La definició interna es pot comparar, \emph{externament}, amb la noció
ingènua de finitud «isomorf a algun \([k]\)». Per a cada numeral estàndard
\(k\), la coincidència \(\seg{\bar k}\cong[k]\)
(observació~\ref{obs:segment-numeral}) dona: si \(A\cong[k]\) per a algun
numeral \(k\), aleshores \(\seg{\bar k}\approx A\), de manera que
\(\bar k\in\operatorname{Card}(A)\) i \(A\) és finit en el sentit intern
(definició~\ref{def:conjunt-finit}). Recíprocament, si es disposa
\emph{externament} d'un punt global \(n\in\operatorname{Card}(A)\) que
sigui un numeral estàndard \(\bar k\), aleshores \(A\cong\seg{\bar
k}\cong[k]\).

Convé distingir tres nivells, que no s'han de confondre:
\begin{enumerate}
\item \(\operatorname{Card}(A)\hookrightarrow\N\) és una \emph{part
interna}, sempre ben definida (proposició~\ref{prop:card-interna});
\item «\(A\) és finit» és l'afirmació, \emph{externa}, que aquesta part és
no inicial ---és a dir, que té algun punt global;
\item un tal punt global és un \(n\colon\mathbf1\to\N\) que no té per què
ser un numeral estàndard; la identificació amb un \([k]\) concret només és
possible quan, externament, se sap que \(n\) és un numeral.
\end{enumerate}
Els resultats interns del capítol (comparació, unicitat, recomptes) no
requereixen cap identificació amb numerals: es demostren sobre els segments
\(\seg{n}\) amb \(n\) intern.
\end{observacio}

\begin{observacio}[Treballem sobre els segments interns]\label{obs:seg-vs-corxet}
En tot el que segueix, els resultats interns ---comparació de mides,
unicitat del cardinal, recomptes--- es demostren \emph{sobre els segments
\(\seg{n}\)} amb \(n\colon\mathbf1\to\N\) intern, i les inducions es fan
amb el principi d'inducció (teorema~\ref{teo:induccio}) usant la
descomposició successora \(\seg{s\circ n}\cong\seg{n}+\mathbf1\)
(proposició~\ref{prop:segment-successor}) allà on abans s'hauria escrit
\([s\circ n]=[n]+\mathbf1\). Així cap prova no depèn de la família
metanivell \([n]\): la finitud i el cardinal són interns
(definició~\ref{def:conjunt-finit}), i els \([k]\) queden reservats a la
lectura externa amb numerals concrets (observació~\ref{obs:segment-numeral}).
\end{observacio}

El nucli de tota la teoria és el fet que dos ordinals isomorfs són iguals.
El demostrarem per inducció (lema~\ref{lema:comparacio-ordinals}), i en derivarem alhora la unicitat del
cardinal i el principi dels calaixos.

Abans, però, cal assegurar que la inducció d'aquest lema \emph{sí} que és
interna. La dificultat és la mateixa que resolguérem per a
\(\operatorname{Card}(A)\): tant el domini \(\seg{m}\) com el codomini
\(\seg{n}\) varien amb els punts \(m,n\), de manera que «existeix un
monomorfisme \(\seg{m}\hookrightarrow\seg{n}\)» no és, escrit així, una
fórmula amb \(m\) lliure sobre la qual aplicar el principi d'inducció
(teorema~\ref{teo:induccio}). Ho resolem codificant tots els monomorfismes
entre segments dins un \emph{exponencial fix}, exactament com fèrem amb
\(\mathcal E_A\).

\begin{definicio}[Codificació uniforme dels monomorfismes entre segments]\index{Codificació uniforme dels monomorfismes entre segments}\label{def:mono-segments}
De la mateixa manera que a la definició~\ref{def:card-A}, una aplicació
\(\seg{m}\to\seg{n}\subseteq\N\) s'estén de manera única a una aplicació
\(\N\to\N+\mathbf1\) que hi val el punt afegit fora de \(\seg{m}\). Definim
una part d'un objecte \emph{fix},
\[
\mathcal M\ \hookrightarrow\ \N\times\N\times(\N+\mathbf1)^{\N},
\]
els punts \((m,n,g)\) de la qual ---amb \(g\colon\N\to\N+\mathbf1\), i
\(i_\N\colon\N\to\N+\mathbf1\), \(i_1\colon\mathbf1\to\N+\mathbf1\) les
injeccions de la suma--- satisfan internament:
\begin{enumerate}
\item[(m1)] per a tot \(k<m\), existeix un únic \(\ell<n\) amb
\(g(k)=i_\N(\ell)\);
\item[(m2)] per a tot \(k\) amb \(k\not<m\), \(g(k)=i_1(\ast)\);
\item[(m3)] per a tot \(k,k'<m\), si \(g(k)=g(k')\) aleshores \(k=k'\).
\end{enumerate}
Aquestes condicions només fan servir avaluacions de l'exponencial \emph{fix}
\((\N+\mathbf1)^{\N}\), l'ordre \(<\) (definició~\ref{def:ordre-naturals}),
igualtats i quantificadors acotats; \(\mathcal M\) n'és la part
corresponent, tallada per equalitzadors i imatges inverses. Les condicions
(m1)--(m2) diuen que \(g\) codifica una aplicació \(\seg{m}\to\seg{n}\)
(estesa per \(\ast\)), i (m3) que és injectiva. Així la fibra
\(\mathcal M_{m,n}\) sobre \((m,n)\) és no inicial si i només si existeix un
monomorfisme \(\seg{m}\hookrightarrow\seg{n}\).
\end{definicio}

Amb aquesta codificació, «existeix un monomorfisme
\(\seg{m}\hookrightarrow\seg{n}\)» esdevé la condició interna
«\(\mathcal M_{m,n}\) és no inicial», i podem formar el predicat intern
\[
P(m)\ :\Longleftrightarrow\ \forall n\colon\mathbf1\to\N\ \bigl(\mathcal
M_{m,n}\text{ és no inicial}\ \Rightarrow\ m\le n\bigr),
\]
que sí que és una fórmula amb \(m\) com a única variable lliure, apta per al
principi d'inducció (teorema~\ref{teo:induccio}). La prova del lema aplica
la inducció interna precisament a \(P\).

\begin{lema}[Comparació de segments]\label{lema:comparacio-ordinals}
Siguin \(m,n\colon\mathbf1\to\N\).
\begin{enumerate}
\item Si existeix un monomorfisme \(\seg{m}\hookrightarrow\seg{n}\),
aleshores \(m\le n\).
\item Si existeix un isomorfisme \(\seg{m}\cong\seg{n}\), aleshores \(m=n\).
\end{enumerate}
\end{lema}

\begin{prova}
La segona afirmació se segueix de la primera: un isomorfisme dona
monomorfismes \(\seg{m}\hookrightarrow\seg{n}\) i
\(\seg{n}\hookrightarrow\seg{m}\), d'on \(m\le n\) i \(n\le m\), i per
l'antisimetria (proposició~\ref{prop:ordre-total}) \(m=n\).

Provem la primera aplicant el principi d'inducció
(teorema~\ref{teo:induccio}) al predicat intern \(P(m)\) definit més amunt
(definició~\ref{def:mono-segments}); com que \(P\) quantifica internament
sobre \(n\), la inducció és uniforme en \(n\) sense sortir del llenguatge.
Per a \(P(0)\): si \(\mathcal M_{0,n}\) és no inicial, aleshores \(0\le n\)
sempre (proposició~\ref{prop:ordre-total}). Per al pas, suposem \(P(m)\) i
provem \(P(s\circ m)\): sigui \(n\) tal que \(\mathcal M_{s\circ m,\,n}\)
és no inicial, és a dir hi ha un monomorfisme
\(f\colon\seg{s\circ m}\hookrightarrow\seg{n}\). Com que \(\seg{s\circ m}\cong\seg{m}+\mathbf1\)
(proposició~\ref{prop:segment-successor}) és no inicial, també ho és
\(\seg{n}\); per tant \(n\ne0\) ---si \(n=0\), \(\seg{0}\cong\mathbf0\) no rebria cap
aplicació des d'un conjunt no inicial (estrictesa de \(\mathbf0\),
capítol~\ref{cap:valors})---, és a dir \(n=s\circ n'\) per a un (únic)
punt \(n'\) (lema~\ref{lema:zero-o-successor}), i
\(\seg{n}\cong\seg{n'}+\mathbf1\)
(proposició~\ref{prop:segment-successor}).

Escrivim \(\seg{s\circ m}\cong\seg{m}+\mathbf1\) amb les injeccions
\(\iota\colon\seg{m}\hookrightarrow\seg{s\circ m}\) i
\(\ast\colon\mathbf1\to\seg{s\circ m}\), i
\(\seg{n}\cong\seg{n'}+\mathbf1\) amb
\(\jmath\colon\seg{n'}\hookrightarrow\seg{n}\) i
\(\ast'\colon\mathbf1\to\seg{n}\). Volem un monomorfisme
\(g\colon\seg{m}\hookrightarrow\seg{n'}\). El pas essencial és decidir on
va a parar el punt nou \(f\circ\ast\colon\mathbf1\to\seg{n}\). Com que
\(\seg{n}\cong\seg{n'}+\mathbf1\) és una suma disjunta, aquest punt global
es factoritza per exactament una de les dues injeccions
(teorema~\ref{teo:exhaustivitat-sumes}: tot punt global d'una suma prové
d'un únic sumand); distingim els dos casos.

\emph{Cas 1: \(f\circ\ast=\ast'\).} Aleshores la imatge inversa
\((f\circ\iota)^{-1}(\ast')\hookrightarrow\seg{m}\) és buida: un punt
\(x\) seu compliria \(f\circ\iota\circ x=\ast'=f\circ\ast\), i la
injectivitat de \(f\) donaria \(\iota\circ x=\ast\), impossible perquè
\(\iota\) i \(\ast\) són parts complementàries de \(\seg{s\circ m}\) (la
seva intersecció és inicial i no té punts); sense punts, la imatge
inversa és inicial pel lema del buit (lema~\ref{lema:buit}). Per tant
\(f\circ\iota\) es factoritza pel complement de \(\ast'\) (descomposició
\(\seg{n}\cong\ast'+(\ast')^c\) i estabilitat), que és
\(\jmath\colon\seg{n'}\hookrightarrow\seg{n}\):
existeix un únic \(g\colon\seg{m}\to\seg{n'}\) amb \(\jmath\circ
g=f\circ\iota\). Aquest \(g\) és un monomorfisme perquè \(\jmath\circ g\)
ho és i \(\jmath\) és un monomorfisme.

\emph{Cas 2: \(f\circ\ast=\jmath\circ q\)} per a un únic punt
\(q\colon\mathbf1\to\seg{n'}\). Reduïm al Cas~1 permutant a \(\seg{n}\) el
punt \(\jmath\circ q\) amb el punt nou \(\ast'\). Definim
\(\tau\colon\seg{n}\to\seg{n}\) \emph{per casos}, usant la descomposició de
\(\seg{n'}\) en el singleton determinat per \(q\) i el seu complement
\(C\hookrightarrow\seg{n'}\) (tota part té complement,
capítol~\ref{cap:valors}), de manera que
\(\seg{n}\cong C+\mathbf1_q+\mathbf1_{\ast'}\), on \(\mathbf1_q\) és la
fibra de \(q\) i \(\mathbf1_{\ast'}\) el sumand nou. Sobre aquesta
descomposició, \(\tau\) actua com:
\[
\tau|_C=\id_C,
\qquad
\tau|_{\mathbf1_q}\colon\mathbf1_q\xrightarrow{\ \sim\ }\mathbf1_{\ast'},
\qquad
\tau|_{\mathbf1_{\ast'}}\colon\mathbf1_{\ast'}\xrightarrow{\ \sim\ }\mathbf1_q,
\]
és a dir, la identitat sobre \(C\) i l'intercanvi dels dos singletons.
Aquesta definició per casos (propietat universal de la suma,
notació~\ref{not:suma}) fa immediat que \(\tau\circ\tau=\id_{\seg{n}}\),
ja que intercanviar dues vegades és la identitat sobre cada sumand; per
tant \(\tau\) és una involució, i en particular un isomorfisme. Aleshores
\(\tau\circ f\colon\seg{s\circ m}\hookrightarrow\seg{n}\) és un
monomorfisme amb \(\tau\circ f\circ\ast=\tau\circ\jmath\circ
q=\ast'\), de manera que recau en el Cas~1; aplicant-hi el Cas~1 obtenim
el monomorfisme \(g\colon\seg{m}\hookrightarrow\seg{n'}\).

En tots dos casos hi ha un monomorfisme
\(g\colon\seg{m}\hookrightarrow\seg{n'}\), és a dir \(\mathcal M_{m,n'}\) és
no inicial. Per la hipòtesi d'inducció \(P(m)\), aplicada a \(n'\), obtenim
\(m\le n'\), i doncs \(s\circ m\le s\circ n'=n\) (pas del successor,
proposició~\ref{prop:ordre-total}). Com que \(n\) era arbitrari, això prova
\(P(s\circ m)\) i tanca la inducció.
\end{prova}

\begin{teorema}[Unicitat del cardinal finit]\label{teo:unicitat-cardinal}
Si \(A\) és finit, la part \(\operatorname{Card}(A)\hookrightarrow\N\) té un
\emph{únic} punt: el cardinal \(n\) amb \(\seg{n}\approx A\) és únic. En
conseqüència, la notació \(|A|\) està ben definida, i dos conjunts finits
són isomorfs si i només si tenen el mateix cardinal.
\end{teorema}

\begin{prova}
Treballem sobre els segments interns (observació~\ref{obs:seg-vs-corxet}).
Si \(\seg{m}\approx A\) i \(\seg{n}\approx A\), aleshores
\(\seg{m}\approx\seg{n}\) per transitivitat de l'equipotència
(definició~\ref{def:equipotencia}), és a dir hi ha un isomorfisme
\(\seg{m}\cong\seg{n}\); pel lema~\ref{lema:comparacio-ordinals}, \(m=n\).
Això prova que \(\operatorname{Card}(A)\) té com a molt un punt, i dona la
bona definició de \(|A|\). Per a la darrera afirmació, si \(A\cong B\) amb
\(\seg{|A|}\approx A\) i \(\seg{|B|}\approx B\), aleshores
\(\seg{|A|}\approx\seg{|B|}\) i \(|A|=|B|\); recíprocament, \(|A|=|B|\)
dona \(A\approx\seg{|A|}=\seg{|B|}\approx B\), és a dir \(A\cong B\).
\end{prova}

\section{El principi dels calaixos}\label{sec:calaixos}

La comparació d'ordinals conté ja, essencialment, el principi dels
calaixos. L'enunciem en la forma habitual i en derivem l'equivalència, per
a conjunts finits del mateix cardinal, entre injectivitat, exhaustivitat i
bijectivitat.

\begin{teorema}[Principi dels calaixos]\label{teo:calaixos}
Siguin \(A\) i \(B\) conjunts finits amb \(|A|>|B|\). Aleshores no hi ha
cap monomorfisme \(A\hookrightarrow B\). Dit altrament: tota aplicació
d'un conjunt finit cap a un de cardinal estrictament menor identifica
almenys dos elements.
\end{teorema}

\begin{prova}
Un monomorfisme \(A\hookrightarrow B\) donaria, component amb els
isomorfismes \(\seg{|A|}\approx A\) i \(B\approx\seg{|B|}\)
(definició~\ref{def:conjunt-finit}), un monomorfisme
\(\seg{|A|}\hookrightarrow\seg{|B|}\); pel
lema~\ref{lema:comparacio-ordinals} seria \(|A|\le|B|\), contra la
hipòtesi. La reformulació és la negació de la definició de monomorfisme
(capítol~\ref{cap:aplicacions}): si l'aplicació no és un monomorfisme,
dos elements generalitzats distints tenen la mateixa imatge.
\end{prova}

\begin{observacio}[Les induccions dels recomptes són internes]\label{obs:recomptes-interns}
Les proves que segueixen fan inducció sobre un punt \(n\) en afirmacions
com \enquote{\(A+\seg{n}\) és finit amb cardinal \(|A|+n\)}, és a dir,
\enquote{existeix un isomorfisme \(\seg{|A|+n}\to A+\seg{n}\)}, en què
tant el domini com el codomini varien amb \(n\). Perquè el
teorema~\ref{teo:induccio} s'hi pugui aplicar cal que la propietat sigui
una part de \(\N\); i ho és, per la mateixa tècnica de la
definició~\ref{def:card-A} i de la definició~\ref{def:mono-segments}:
totes les aplicacions \(\seg{j}\to A+\seg{n}\) s'estenen a aplicacions
\(\N\to(A+\N)+\mathbf1\), de manera que \enquote{existeix un isomorfisme}
és la no-inicialitat d'una fibra d'una part d'un exponencial fix, i és
una fórmula interna en \(n\). No repetirem la codificació a cada
recompte; el lector la pot reconstruir seguint el model de
\(\mathcal E_A\).
\end{observacio}

\begin{lema}[Recompte de la suma finita]\label{lema:recompte-suma}
Siguin \(A\) i \(B\) conjunts finits. Aleshores la suma \(A+B\) és finita i
\[
|A+B|=|A|+|B|,
\]
on el segon membre és la suma de naturals (definició~\ref{def:suma-naturals}).
En particular, si \(B\) no és inicial ---és a dir, \(|B|\ge 1\)--- aleshores
\(|A+B|>|A|\).
\end{lema}

\begin{prova}
Com que \(B\approx\seg{|B|}\), n'hi ha prou de provar, per inducció sobre
el punt \(n\) (teorema~\ref{teo:induccio}) i amb \(A\) fixat, que
\(A+\seg{n}\) és finit amb cardinal \(|A|+n\); la propietat és interna en
\(n\) per l'observació~\ref{obs:recomptes-interns}. Escrivim les dues
etapes en termes de \(B\). Per a \(|B|=0\), \(B\approx\seg{0}\cong\mathbf0\)
i la neutralitat del conjunt inicial respecte de la suma
(capítol~\ref{cap:valors}) dona \(A+B\cong A\), d'on
\(|A+B|=|A|=|A|+0\).

Pas de \(|B|=n\) a \(|B|=s\circ n\). Sigui \(B\approx\seg{s\circ
n}\cong\seg{n}+\mathbf1\) (proposició~\ref{prop:segment-successor}). Per
l'associativitat de la suma de conjunts (capítol~\ref{cap:valors}),
\[
A+B\cong A+(\seg{n}+\mathbf1)\cong (A+\seg{n})+\mathbf1 .
\]
Per la hipòtesi d'inducció, \(A+\seg{n}\) és finit amb
\(|A+\seg{n}|=|A|+n\); i afegir un sumand terminal incrementa el cardinal
en un, perquè \((A+\seg{n})+\mathbf1\approx\seg{|A|+n}+\mathbf1\cong
\seg{s\circ(|A|+n)}\) per la descomposició successora
(proposició~\ref{prop:segment-successor}). Així \(A+B\) és finit i
\[
|A+B|=s\circ(|A|+n)=|A|+(s\circ n)=|A|+|B|,
\]
on la penúltima igualtat és la definició recursiva de la suma
(definició~\ref{def:suma-naturals}). Finalment, si \(|B|\ge1\) aleshores
\(|B|=s\circ k\) per a algun \(k\) (lema~\ref{lema:zero-o-successor}) i
\(|A|+|B|=|A|+(s\circ k)\), que és \(>|A|\) per la definició de l'ordre
estricte (definició~\ref{def:ordre-naturals}).
\end{prova}

\begin{lema}[Tota part d'un conjunt finit és finita]\label{lema:part-finita}
Sigui \(S\hookrightarrow A\) una part d'un conjunt finit \(A\). Aleshores
\(S\) i el seu complement són finits, i \(|S|\le|A|\).
\end{lema}

\begin{prova}
Com que \(A\) és finit, \(A\approx\seg{n}\) amb \(n=|A|\); una part de
\(A\) es correspon, per aquest isomorfisme, amb una part de \(\seg{n}\),
així que n'hi ha prou de provar-ho per a \(A=\seg{n}\), per inducció sobre
\(n\) (teorema~\ref{teo:induccio}).

Per a \(n=0\), \(\seg{0}\cong\mathbf0\) i l'única part és la inicial, que
és finita amb cardinal \(0\).

Pas de \(n\) a \(s\circ n\). Per la descomposició successora,
\(\seg{s\circ n}\cong\seg{n}+\mathbf1\)
(proposició~\ref{prop:segment-successor}). Sigui \(S\hookrightarrow
\seg{s\circ n}\) una part. La imatge inversa de \(S\) per les dues
injeccions la descompon en dues parts complementàries: \(S_0\hookrightarrow
\seg{n}\), la part de \(S\) situada sobre \(\seg{n}\), i \(S_1\hookrightarrow
\mathbf1\), la part situada sobre el punt nou. Ara \(S\cong S_0+S_1\)
(les dues parts, complementàries dins \(S\), formen una suma,
capítol~\ref{cap:valors}). Per la hipòtesi d'inducció, \(S_0\) és finit amb
\(|S_0|\le n\). I \(S_1\hookrightarrow\mathbf1\) és inicial o terminal
---les úniques parts de \(\mathbf1\) (capítol~\ref{cap:valors})---, en tot
cas finit amb \(|S_1|\le1\). Pel recompte de la suma
(lema~\ref{lema:recompte-suma}), \(S\) és finit i
\[
|S|=|S_0|+|S_1|\le n+1=s\circ n=|\seg{s\circ n}| .
\]
El complement de \(S\) dins \(\seg{s\circ n}\) es tracta igual, canviant
\(S_0,S_1\) pels seus complements dins \(\seg{n}\) i \(\mathbf1\)
respectivament.
\end{prova}

\begin{teorema}[Injectiu, exhaustiu i bijectiu en el cas finit]\label{teo:endo-finit}
Sigui \(f\colon A\to B\) una aplicació entre conjunts finits amb
\(|A|=|B|\). Són equivalents:
\begin{enumerate}
\item \(f\) és un monomorfisme;
\item \(f\) és un epimorfisme;
\item \(f\) és un isomorfisme.
\end{enumerate}
\end{teorema}

\begin{prova}
Que (3) implica (1) i (2) és immediat. Vegem que (1) implica (3) i,
després, que (2) implica (3).

\emph{(1) \(\Rightarrow\) (3).} Suposem \(f\) un monomorfisme. Per la
factorització imatge (teorema~\ref{teo:factoritzacio}), \(f=m\circ e\) amb
\(e\colon A\twoheadrightarrow\operatorname{Im}(f)\) epimorfisme regular i
\(m\colon\operatorname{Im}(f)\hookrightarrow B\) monomorfisme; com que
\(f\) és un monomorfisme, \(e\) és un isomorfisme
(proposició~\ref{prop:mono-nucli-trivial}), i per tant
\(\operatorname{Im}(f)\cong A\) és finit amb
\(|\operatorname{Im}(f)|=|A|=|B|\). Ara \(m\) exhibeix
\(\operatorname{Im}(f)\) com una part de \(B\). Si \(m\) no fos exhaustiu,
el seu complement \(C\hookrightarrow B\) seria una part no inicial, finita
pel lema~\ref{lema:part-finita} (és una part del finit \(B\)), i \(B\) es
descompondria en la suma disjunta \(B\cong\operatorname{Im}(f)+C\)
(el complement d'una part i la seva part formen una suma,
capítol~\ref{cap:valors}). Com que tots dos sumands són finits, s'hi aplica
el recompte de la suma (lema~\ref{lema:recompte-suma}); i com que \(C\) no
és inicial, \(|B|=|\operatorname{Im}(f)|+|C|>|\operatorname{Im}(f)|=|B|\),
una contradicció. Per tant \(m\) és exhaustiu, i sent un monomorfisme
exhaustiu és un isomorfisme (proposició~\ref{prop:epi-imatge-total} i
corol·lari~\ref{cor:bimorfisme-iso}). Aleshores \(f=m\circ e\) és
composició de dos isomorfismes, i és un isomorfisme.

\emph{(2) \(\Rightarrow\) (3).} Suposem ara \(f\) un epimorfisme. Per
l'axioma d'elecció (axioma~\ref{ax:eleccio}), \(f\) és escindit: existeix
una secció
\(s\colon B\to A\) amb \(f\circ s=\id_B\). Tota secció és un monomorfisme
(proposició~\ref{prop:seccio-mono-retraccio-epi}), de manera que
\(s\colon B\to A\) és un monomorfisme entre conjunts finits amb
\(|B|=|A|\). Per la implicació (1)\(\Rightarrow\)(3) ja provada, \(s\) és un
isomorfisme. Aleshores, de \(f\circ s=\id_B\) se segueix \(f=s^{-1}\)
---component per la dreta amb \(s^{-1}\)---, i \(f\) és un isomorfisme.
Aquesta prova evita tota descomposició en fibres i no és circular: només
usa l'elecció, ja disponible, i la direcció (1)\(\Rightarrow\)(3). És
l'únic lloc del capítol on intervé l'elecció, i s'hi pot evitar: com que
\(A\approx\seg{n}\) és ben ordenat, una secció de \(f\) es pot
\emph{definir} enviant cada punt \(b\) al mínim de la fibra \(f^{-1}(b)\),
que és no buida perquè \(f\) és exhaustiu sobre punts; la construcció
interna d'aquesta secció és un exercici.
\end{prova}

\begin{observacio}[Lectura clàssica]\label{obs:calaixos-classica}
En conjunts, el teorema~\ref{teo:calaixos} és l'enunciat familiar segons
el qual, si es reparteixen \(|A|\) objectes en \(|B|<|A|\) calaixos, algun
calaix en rep més d'un. El teorema~\ref{teo:endo-finit} és el fet, també
familiar i propi \emph{només} del cas finit, que una aplicació d'un
conjunt finit en si mateix és injectiva exactament quan és exhaustiva.
Sobre conjunts infinits ---\(\N\) mateix, amb el successor--- això falla:
el successor és injectiu però no exhaustiu.
\end{observacio}

\section{Parts d'un conjunt finit}\label{sec:parts-finites}

El pas de \(A\) a \(A+\mathbf1\) duplica el nombre de parts: cada part de
\(A+\mathbf1\) o bé conté el punt nou o bé no. Aquesta observació,
convertida en inducció, dona el recompte de parts. Abans, però, ens cal el
recompte del producte, que la prova utilitza sota la forma de la
duplicació \(|A\times\mathbf2|=2\cdot|A|\).

\begin{lema}[Recompte del producte finit]\label{lema:recompte-producte}
Siguin \(A\) i \(B\) conjunts finits. Aleshores el producte \(A\times B\)
és finit i
\[
|A\times B|=|A|\cdot|B|,
\]
on el segon membre és el producte de naturals
(definició~\ref{def:producte-exp}). En particular
\(|A\times\mathbf2|=2\cdot|A|\).
\end{lema}

\begin{prova}
Fem inducció sobre \(|B|\), amb \(A\) fixat (teorema~\ref{teo:induccio}).
Per a \(|B|=0\), \(B\approx\seg{0}\cong\mathbf0\) i
\(A\times\mathbf0\cong\mathbf0\) ---el producte per l'inicial és inicial,
per la distributivitat del producte sobre la suma i l'estrictesa de
\(\mathbf0\) (capítol~\ref{cap:valors})---, d'on \(|A\times B|=0=|A|\cdot0\).

Pas de \(|B|=n\) a \(|B|=s\circ n\). Sigui \(B\approx\seg{s\circ
n}\cong\seg{n}+\mathbf1\) (proposició~\ref{prop:segment-successor}). Per la
distributivitat del producte sobre la suma
(teorema~\ref{teo:distributivitat}),
\[
A\times B\cong A\times(\seg{n}+\mathbf1)\cong (A\times\seg{n})+(A\times\mathbf1)
\cong (A\times\seg{n})+A,
\]
usant \(A\times\mathbf1\cong A\). Per la hipòtesi d'inducció
\(A\times\seg{n}\) és finit amb \(|A\times\seg{n}|=|A|\cdot n\); pel
recompte de la suma (lema~\ref{lema:recompte-suma}),
\[
|A\times B|=|A\times\seg{n}|+|A|=|A|\cdot n+|A|=|A|\cdot(s\circ n)=|A|\cdot|B|,
\]
on la penúltima igualtat és la definició recursiva del producte
(definició~\ref{def:producte-exp}). El cas \(B=\mathbf2\) (amb
\(|\mathbf2|=2\)) dona \(|A\times\mathbf2|=|A|\cdot2=2\cdot|A|\) per la
commutativitat del producte (proposició~\ref{prop:lleis-aritmetica}).
\end{prova}

\begin{teorema}[Nombre de parts d'un conjunt finit]\label{teo:nombre-parts}
Si \(A\) és finit amb \(|A|=n\), aleshores el conjunt de parts
\(\mathcal P(A)\) (definició~\ref{def:conjunt-parts}) és finit i
\[
|\mathcal P(A)|=2^{n},
\]
on \(2^{n}\) és l'exponenciació de naturals
(definició~\ref{def:producte-exp}).
\end{teorema}

\begin{prova}
Com que \(\mathcal P(A)\cong\mathbf2^{A}\)
(definició~\ref{def:conjunt-parts}) i \(A\approx\seg{n}\) amb \(n=|A|\),
n'hi ha prou de provar \(|\mathbf2^{\seg{n}}|=2^{n}\) per inducció sobre
\(n\). Per a \(n=0\), \(\mathbf2^{\seg{0}}\cong\mathbf2^{\mathbf0}\cong
\mathbf1\), i \(2^{0}=1\) (definició~\ref{def:producte-exp}). Pas de \(n\)
a \(s\circ n\): de \(\seg{s\circ n}\cong\seg{n}+\mathbf1\)
(proposició~\ref{prop:segment-successor}) i de la conversió d'un
exponencial d'una suma en un producte d'exponencials
(proposició~\ref{prop:exponencial-suma}),
\[
\mathbf2^{\seg{s\circ n}}\cong\mathbf2^{\seg{n}+\mathbf1}\cong
\mathbf2^{\seg{n}}\times\mathbf2^{\mathbf1}\cong\mathbf2^{\seg{n}}\times\mathbf2 .
\]
Per la hipòtesi d'inducció \(\mathbf2^{\seg{n}}\) és finit de cardinal
\(2^{n}\), i multiplicar per \(\mathbf2\) duplica el cardinal
(lema~\ref{lema:recompte-producte}), de manera que
\(\mathbf2^{\seg{s\circ n}}\) té cardinal \(2^{n}+2^{n}=2\cdot2^{n}=
2^{s\circ n}\), on la darrera igualtat és la definició recursiva de
l'exponenciació.
\end{prova}

\begin{observacio}[Els recomptes de la suma i del producte]\label{obs:recompte-producte-suma}
Els dos recomptes que aquestes proves fan servir ---\(|A+B|=|A|+|B|\) i
\(|A\times B|=|A|\cdot|B|\) per a conjunts finits--- s'han establert com a
lemes previs (lemes~\ref{lema:recompte-suma} i~\ref{lema:recompte-producte}),
tots dos per inducció sobre el cardinal del segon factor. Són la traducció,
al recompte de cardinals, de la suma disjunta i de l'enumeració fila per
fila del producte; la duplicació \(|A\times\mathbf2|=2\cdot|A|\) que hem
usat n'és el cas \(B=\mathbf2\).
\end{observacio}

\section{Recompte d'aplicacions}\label{sec:recompte-aplicacions}

L'exponencial \(B^{A}\) internalitza el conjunt de les aplicacions de
\(A\) cap a \(B\) (capítol~\ref{cap:exponencials}). Quan tots dos conjunts
són finits, en compta exactament les aplicacions.

\begin{teorema}[Nombre d'aplicacions entre conjunts finits]\label{teo:nombre-aplicacions}
Si \(A\) i \(B\) són finits amb \(|A|=n\) i \(|B|=m\), aleshores
l'exponencial \(B^{A}\) és finit i
\[
|B^{A}|=m^{n} .
\]
\end{teorema}

\begin{prova}
Fixats isomorfismes \(A\approx\seg{n}\) i \(B\approx\seg{m}\) (amb
\(n=|A|\), \(m=|B|\)), l'exponencial només depèn dels dos conjunts llevat
d'isomorfisme, de manera que n'hi ha prou de provar
\(|\seg{m}^{\seg{n}}|=m^{n}\) per inducció sobre \(n\), uniformement en
\(m\). Per a \(n=0\), \(\seg{m}^{\seg{0}}\cong\seg{m}^{\mathbf0}\cong
\mathbf1\), i \(m^{0}=1\). Pas de \(n\) a \(s\circ n\): de
\(\seg{s\circ n}\cong\seg{n}+\mathbf1\)
(proposició~\ref{prop:segment-successor}) i de la conversió de
l'exponencial d'una suma (proposició~\ref{prop:exponencial-suma}),
\[
\seg{m}^{\seg{s\circ n}}\cong\seg{m}^{\seg{n}+\mathbf1}\cong
\seg{m}^{\seg{n}}\times\seg{m}^{\mathbf1}\cong\seg{m}^{\seg{n}}\times\seg{m} .
\]
Per la hipòtesi d'inducció \(\seg{m}^{\seg{n}}\) té cardinal \(m^{n}\), i
pel recompte del producte (lema~\ref{lema:recompte-producte}) multiplicar
per \(\seg{m}\) multiplica el cardinal per \(m\); així
\(|\seg{m}^{\seg{s\circ n}}|=m^{n}\cdot m=m^{s\circ n}\), la definició
recursiva de l'exponenciació.
\end{prova}

\begin{observacio}[Lectura clàssica i enllaç amb la pràctica]\label{obs:recompte-classica}
Els quatre recomptes ---\(2^{n}\) parts, \(m^{n}\) aplicacions,
\(|A+B|=|A|+|B|\) i \(|A\times B|=|A|\cdot|B|\)--- són l'aritmètica
combinatòria elemental, ara obtinguda sense sortir de la teoria: cada
igualtat és un isomorfisme construït per inducció, no un recompte de
punts. Aquests enunciats són el punt de partida dels exercicis
computacionals de la part pràctica, on es calculen cardinals concrets i
es refà, en casos petits, la inducció general. En particular, el nombre
d'aplicacions injectives i el de bijeccions de \([k]\) en si mateix
---les permutacions--- es tracten allà com a aplicació directa del
principi dels calaixos i del teorema~\ref{teo:endo-finit}.
\end{observacio}

\begin{resumcap}
\apartat{Idees centrals}
\begin{enumerate}
  \item Dues famílies de \enquote{nombres finits}: els ordinals \([k]=\mathbf1+\dots+\mathbf1\), definits al metanivell per a numerals estàndard, i els segments \(\seg{n}=\{k<n\}\), fibres de \(\operatorname{Lt}\hookrightarrow\N\times\N\), que formen una família interna. Coincideixen externament sobre numerals; la finitud es defineix amb els segments.
  \item \(\operatorname{Card}(A)\hookrightarrow\N\) és la imatge de la projecció d'una part \(\mathcal E_A\) d'un exponencial fix que codifica les bijeccions \(\seg{n}\to A\); \(A\) és finit si \(\operatorname{Card}(A)\) té algun punt, i el cardinal és aquest punt, únic.
  \item La unicitat ve de la comparació de segments: un monomorfisme \(\seg{m}\hookrightarrow\seg{n}\) força \(m\le n\). La prova és una inducció interna sobre un predicat codificat, amb el truc de permutar dos punts per reduir al cas en què el punt nou va al punt nou.
  \item Principi dels calaixos; en cardinal igual, mono, epi i iso coincideixen. L'elecció només s'usa per escindir l'epimorfisme, i es pot evitar amb el bon ordre.
  \item Recomptes per inducció sobre el cardinal: \(|A+B|=|A|+|B|\), \(|A\times B|=|A|\cdot|B|\), \(|P(A)|=2^{|A|}\), \(|B^A|=|B|^{|A|}\); cadascun és un isomorfisme construït amb la descomposició successora, la distributivitat i l'exponencial d'una suma.
  \item Les induccions sobre afirmacions d'existència d'isomorfismes són internes perquè les aplicacions entre objectes variables es codifiquen dins exponencials fixos.
\end{enumerate}
\apartat{Errors típics} Definir la finitud com \enquote{isomorf a algun \([k]\)}, que és un esquema extern i no una propietat interna; fer inducció sobre \enquote{existeix un isomorfisme \(\seg{n}\cong\cdots\)} sense comprovar que és una part de \(\N\); creure que \(\operatorname{Card}(A)\) necessita triar un isomorfisme (es construeix com una imatge, sense cap tria); oblidar que \enquote{un punt d'una suma ve d'un sumand} val sobre punts globals i no sobre elements generalitzats, cosa essencial al lema de comparació; i pensar que mono \(\Rightarrow\) iso en cardinal igual val per a conjunts infinits (el successor de \(\N\) ho desmenteix).
\apartat{Lectura recomanada} \cite[§§7.6--7.7 i cap.~8]{lawvere-rosebrugh} per als conjunts finits, el principi dels calaixos i els recomptes en conjunts abstractes; \cite[article~V, sessió~30]{lawvere-schanuel} per al recompte d'aplicacions \(|B^A|=|B|^{|A|}\) amb exemples; \cite[§§12--13]{halmos} per als nombres naturals finits i l'aritmètica de cardinals finits en la teoria clàssica; i \cite[§§7.1--7.2]{awodey} per a l'objecte de nombres naturals i la seva propietat universal en una categoria, que és el que fa possible la codificació interna d'aquest capítol.
\end{resumcap}

\chapter{Infinitud, enumerabilitat i seqüències}\label{cap:infinitud}

\begin{objectius}
\apartat{Prerequisits} Del capítol~\ref{cap:finits}: finitud interna, cardinal, principi dels calaixos, l'equivalència mono/epi/iso en cardinal igual i la tècnica de codificar aplicacions entre objectes variables dins un exponencial fix. Del capítol~\ref{cap:naturals}: recursió (simple i amb paràmetres), inducció, aritmètica i bon ordre de \(\N\). Del capítol~\ref{cap:parts-interna}: quantificadors interns, per tallar parts definides per condicions.
\apartat{Objectius} En acabar el capítol, el lector hauria de poder:
\begin{itemize}
  \item distingir \enquote{infinit} (no finit) d'\enquote{infinit de Dedekind} (endomonomorfisme no exhaustiu), demostrar que \(\N\) és infinit de Dedekind i que Dedekind implica infinit sense cap axioma;
  \item demostrar que infinit implica Dedekind amb l'elecció numerable, construint primer una família de llistes injectives i després, sense més elecció, una successió injectiva;
  \item definir enumerable i numerable, i demostrar l'estabilitat per imatge, part, suma i producte sense elecció, amb un aparellament \(\N\cong\N\times\N\) construït per recursió;
  \item demostrar que tota part de \(\N\) és enumerable (el mínim substitueix l'elecció) i que tot enumerable infinit és isomorf a \(\N\);
  \item construir \(\operatorname{List}(A)\) com a part d'un exponencial fix i demostrar-ne el principi de recursió amb el truc de la diagonal sobre la llargada;
  \item situar l'elecció numerable: on cal (unió numerable d'enumerables, dues infinituds) i on no (parts de \(\N\), llistes).
\end{itemize}
\end{objectius}

El capítol anterior va internalitzar la finitud: un
conjunt és finit quan és isomorf a algun segment inicial \(\seg{n}\)
(definició~\ref{def:conjunt-finit}), i el seu tret distintiu és
l'equivalència, en cardinal igual, entre injectivitat, exhaustivitat i
bijectivitat (teorema~\ref{teo:endo-finit}). El present capítol pren
aquesta darrera equivalència com a frontissa: un conjunt és \emph{infinit}
quan simplement no és finit (definició~\ref{def:conjunt-infinit}), i el que
la fa manejable és que en un conjunt infinit aquella equivalència falla, de
manera que admet un monomorfisme estricte en si mateix. Aquesta darrera
propietat ---la infinitud de \emph{Dedekind}--- no és la definició
d'infinit, sinó una noció a priori distinta; la seva equivalència amb «no
finit» depèn de l'elecció, que explicitarem més endavant
(teorema~\ref{teo:infinit-implica-dedekind}). A partir d'aquí desenvolupem la
comparació de cardinals \emph{sense} la restricció de finitud ---els
conjunts enumerables, els subobjectes de \(\N\), les llistes finites i
les unions indexades numerablement--- amb un fil conductor constant:
distingir en cada pas què es demostra amb els axiomes ja assumits i què en
requereix un de nou. L'axioma nou que reapareixerà és l'\emph{elecció
numerable}; el senyalarem cada cop que intervingui.

\section{Conjunts infinits i infinitud de Dedekind}\label{sec:infinits}

La manera més directa de dir que un conjunt és infinit és negar que sigui
finit. Aquesta definició, però, és purament externa a l'estructura del
conjunt: no exhibeix cap testimoni de la infinitud. Dedekind en va donar
una de positiva ---un conjunt és infinit si es pot posar en bijecció amb
una part pròpia--- que sí que n'aporta un, i que resulta molt més
manejable. Introduïm totes dues i n'estudiem la relació.

\begin{definicio}[Conjunt infinit]\index{Conjunt infinit}\label{def:conjunt-infinit}
Un conjunt \(A\) és \emph{infinit} si no és finit; és a dir, si
\(\operatorname{Card}(A)\) és inicial (definició~\ref{def:conjunt-finit}):
no hi ha cap punt \emph{intern} \(n\colon\mathbf1\to\N\) amb
\(\seg{n}\approx A\).
\end{definicio}

\begin{observacio}[Lectura externa amb numerals estàndard]\label{obs:infinit-numeral}
En particular, si \(A\) és infinit, aleshores no és isomorf a cap ordinal
\([k]\) associat a un numeral estàndard \(k\): si ho fos, la coincidència
\(\seg{\bar k}\cong[k]\) (observació~\ref{obs:segment-numeral}) donaria
\(\bar k\in\operatorname{Card}(A)\), contra la definició. El recíproc, en
canvi, no val en general: un punt global intern \(n\colon\mathbf1\to\N\)
no ha de ser necessàriament un numeral estàndard, de manera que
\(\operatorname{Card}(A)\) pot ser no inicial ---i \(A\) finit en el sentit
intern--- sense que \(A\) sigui isomorf a cap \([k]\) estàndard. Només es
pot afirmar el recíproc sota una hipòtesi externa que garanteixi que el
testimoni intern de finitud és estàndard.
\end{observacio}

\begin{definicio}[Infinit de Dedekind]\index{Infinit de Dedekind}\label{def:dedekind-infinit}
Un conjunt \(A\) és \emph{infinit de Dedekind} si existeix un monomorfisme
\(f\colon A\hookrightarrow A\) que no és exhaustiu; equivalentment, si
existeix un monomorfisme \(A\hookrightarrow A\) la imatge del qual és una
part pròpia de \(A\).
\end{definicio}

\begin{observacio}[Les dues nocions i el cas finit]\label{obs:dues-infinituds}
Sobre un conjunt finit les dues definicions coincideixen amb la negació
l'una de l'altra: pel teorema~\ref{teo:endo-finit}, un monomorfisme d'un
conjunt finit en si mateix és sempre exhaustiu, de manera que cap conjunt
finit és infinit de Dedekind. La qüestió delicada és la recíproca ---si
tot conjunt infinit és infinit de Dedekind---, i és precisament on
apareixerà un axioma addicional. Convé, doncs, mantenir les dues nocions
separades fins que en tinguem la relació exacta.
\end{observacio}

El primer fet, i el que justifica tota la teoria, és que \(\N\) mateix és
infinit de Dedekind, i el testimoni és el successor.

\begin{proposicio}[\(\N\) és infinit de Dedekind]\label{prop:N-dedekind}
El successor \(s\colon\N\to\N\) és un monomorfisme no exhaustiu. Per tant
\(\N\) és infinit de Dedekind, i en conseqüència infinit.
\end{proposicio}

\begin{prova}
Que \(s\) és un monomorfisme és el teorema~\ref{teo:peano}(1): de
\(s\circ m=s\circ n\) se segueix \(m=n\). Que no és exhaustiu és el
teorema~\ref{teo:peano}(2): \(0\) no és de la forma \(s\circ n\)
per a cap \(n\), de manera que el punt \(0\colon\mathbf1\to\N\) no factoritza
per \(s\), i la imatge de \(s\) és la part pròpia complementària del zero.
Així \(s\) exhibeix \(\N\) com a infinit de Dedekind. Que \(\N\) és
infinit se segueix de l'observació~\ref{obs:dues-infinituds}: si fos
finit, \(s\) hauria de ser exhaustiu pel teorema~\ref{teo:endo-finit},
contra el que acabem de veure.
\end{prova}

Una direcció de l'equivalència entre les dues infinituds no costa res i
no fa servir cap axioma nou: qui és infinit de Dedekind és infinit.

\begin{proposicio}[Dedekind implica infinit]\label{prop:dedekind-implica-infinit}
Tot conjunt infinit de Dedekind és infinit. La implicació no usa cap
axioma més enllà dels ja assumits.
\end{proposicio}

\begin{prova}
És el contrarecíproc de la primera part de
l'observació~\ref{obs:dues-infinituds}. Si \(A\) fos finit i infinit de
Dedekind alhora, tindríem un monomorfisme \(A\hookrightarrow A\) no
exhaustiu; però pel teorema~\ref{teo:endo-finit} tot monomorfisme d'un
conjunt finit en si mateix és exhaustiu, una contradicció. Per tant un
conjunt infinit de Dedekind no pot ser finit.
\end{prova}

La direcció oposada ---tot conjunt infinit és infinit de Dedekind--- és
la que requereix construir, dins d'un conjunt del qual només sabem que no
és finit, una còpia de \(\N\); i triar-la exigeix una successió d'eleccions.

\begin{teorema}[Infinit implica Dedekind, amb elecció numerable]\label{teo:infinit-implica-dedekind}
Suposem l'\emph{axioma d'elecció numerable}
(axioma~\ref{ax:eleccio-numerable}). Aleshores tot conjunt infinit és
infinit de Dedekind. Equivalentment, sota aquest axioma un conjunt és
infinit si i només si admet un monomorfisme \(\N\hookrightarrow A\).
\end{teorema}

\begin{prova}
Sigui \(A\) infinit. La dificultat és que no podem anar \emph{triant}
\(a_{n+1}\) dins el complement \(A\setminus\{a_0,\dots,a_n\}\): aquest
complement depèn de les tries anteriors, de manera que la família de
conjunts no buits sobre la qual triaríem no està donada per endavant i
l'axioma~\ref{ax:eleccio-numerable} no s'hi aplica directament. Ho evitem
fixant \emph{abans} de triar una única família numerable, independent de
tota elecció.

\emph{La família.} Per a cada \(n\colon\mathbf1\to\N\), sigui
\[
I_n\ :=\ \operatorname{Inj}\bigl(\seg{s\circ n},A\bigr)\ \hookrightarrow\
A^{\seg{s\circ n}}
\]
la part de l'exponencial \emph{fix} \(A^{\seg{s\circ n}}\) formada pels
monomorfismes \(\seg{s\circ n}\hookrightarrow A\) ---les llistes
injectives de llargada \(n+1\)---, tallada per la condició interna
d'injectivitat (quantificadors acotats del capítol~\ref{cap:parts-interna}).
Cada \(I_n\) és \emph{no inicial}, és a dir, per a cada \(n\) existeix un
monomorfisme \(\seg{s\circ n}\hookrightarrow A\). Ho veiem per inducció
sobre \(n\), sense elecció. Per a \(n=0\), \(A\) és no inicial ---no és
finit, en particular no és buit---, de manera que hi ha un punt
\(\mathbf1\cong\seg{1}\hookrightarrow A\). Suposem un monomorfisme
\(g\colon\seg{s\circ n}\hookrightarrow A\); la seva imatge és una part
finita de \(A\) (cardinal \(n+1\)), i com que \(A\) no és finit, aquesta
imatge no és tot \(A\) (altrament \(\seg{s\circ n}\approx A\) i \(A\) seria
finit, definició~\ref{def:conjunt-finit}). El seu complement és, doncs, no
inicial i en conté algun punt \(x\colon\mathbf1\to A\); estenent \(g\) amb
\(x\) sobre el sumand nou de \(\seg{s\circ(s\circ n)}\cong\seg{s\circ
n}+\mathbf1\) (proposició~\ref{prop:segment-successor}) s'obté un
monomorfisme \(\seg{s\circ(s\circ n)}\hookrightarrow A\), perquè \(x\) no
és a la imatge de \(g\). Això tanca la inducció i prova que
\(I_n\) és no inicial per a tot \(n\). Noteu que aquesta inducció només
afirma \emph{existència} fibra a fibra: no proporciona una tria uniforme de
les \(g\), i és per obtenir-ne una que caldrà l'elecció numerable.

Aquestes parts es reuneixen en una \emph{sola} família interna. Codifiquem
cada \(f\in I_n\), com a la definició~\ref{def:card-A}, per la seva
extensió \(\bar f\colon\N\to A+\mathbf1\) que val \(f\) sobre
\(\seg{s\circ n}\) i el punt afegit fora; així totes les llistes
injectives, de qualsevol llargada, viuen dins l'exponencial \emph{fix}
\((A+\mathbf1)^{\N}\). Sigui
\[
E\ :=\ \bigl\{\,(n,\bar f)\ \mid\ n\colon\mathbf1\to\N,\ f\in I_n\,\bigr\}\
\hookrightarrow\ \N\times(A+\mathbf1)^{\N},
\]
la part corresponent (tallada per equalitzadors i condicions acotades), i
\(p\colon E\to\N\) la projecció, la fibra de la qual sobre \(n\) és
exactament \(I_n\). Com que cada \(I_n\) és no inicial, tot punt \(n\) té preimatge per \(p\),
i \(p\) és un epimorfisme (exhaustivitat sobre punts,
proposició~\ref{prop:epi-exhaustiu-global}), regular pel
corol·lari~\ref{cor:epi-regular-sense-eleccio}.
Aquesta família \emph{ja existeix} abans de qualsevol tria.

\emph{La tria.} L'axioma d'elecció numerable
(axioma~\ref{ax:eleccio-numerable}) dona una secció
\(\sigma\colon\N\to E\) de \(p\); és a dir, per a cada \(n\) un monomorfisme
\(f_n\colon\seg{s\circ n}\hookrightarrow A\) escollit \emph{d'un sol cop}
per a tota la família alhora, i no pas un a un. Escrivim
\(f_n(k)\in A\) per al valor de la \(k\)-èsima component,
\(k<n+1\).

\emph{El monomorfisme \(\N\hookrightarrow A\).} A partir de la successió
\((f_n)_n\) construïm \(a\colon\N\to A\) \emph{sense més elecció}. La
temptació és definir \(a_n\) com el primer valor de \(f_n\) que no és cap
dels \(a_0,\dots,a_{n-1}\); però això és una recursió que depèn de tota la
història anterior, i no té la forma de l'axioma~\ref{ax:naturals}. Ho
evitem fent servir llistes prou llargues perquè els valors ja usats no
puguin esgotar-les. Per a cada punt \(n\), sigui
\[
V_n\ \hookrightarrow\ A
\]
la part dels \(x\) tals que \(x=f_{2^m-1}(i)\) per a algun \(m<n\) i algun
\(i<2^m\): és la imatge de la part de \(\N\times\N\) dels parells
\((m,i)\) amb \(m<n\), \(i<2^m\) per l'aplicació
\((m,i)\mapsto f_{2^m-1}(i)\), i la regla \(n\mapsto V_n\) és una
aplicació \(\N\to P(A)\) (quantificadors interns,
capítol~\ref{cap:parts-interna}); aquí \(2^m\) és l'exponenciació de
naturals (definició~\ref{def:producte-exp}) i \(2^m-1\) el seu
predecessor. La part \(V_n\) és finita i \(|V_n|<2^n\). En efecte, el conjunt
d'índexs \(J_n:=\{(m,i):m<n,\ i<2^m\}\) és finit amb \(|J_n|+1=2^n\): per
inducció sobre \(n\), \(|J_0|+1=0+1=2^0\), i
\(J_{n+1}\cong J_n+\seg{2^n}\) (s'hi afegeix la fila \(m=n\)), de manera que,
pel recompte de la suma (lema~\ref{lema:recompte-suma}),
\(|J_{n+1}|+1=(|J_n|+1)+2^n=2^n+2^n=2^{n+1}\). A més, \(V_n\) és imatge de
\(J_n\): un epimorfisme \(J_n\twoheadrightarrow V_n\) té una secció ---enviant
cada punt al mínim índex que hi va a parar, ja que \(J_n\) és ben ordenat
com a part finita de \(\N\times\N\) ordenada lexicogràficament, o bé per
l'elecció---, que és un monomorfisme \(V_n\hookrightarrow J_n\); pel
lema~\ref{lema:part-finita}, \(V_n\) és finita amb
\(|V_n|\le|J_n|<2^n\). En canvi, la llista
\(f_{2^n-1}\) té \(2^n\) valors distints. Pel principi dels calaixos
(teorema~\ref{teo:calaixos}) algun valor de \(f_{2^n-1}\) no és a \(V_n\).
Definim
\[
a_n:=f_{2^n-1}(k_n),
\qquad
k_n:=\min\{\,k<2^n\ :\ f_{2^n-1}(k)\notin V_n\,\},
\]
on el mínim existeix pel bon ordre de \(\N\)
(teorema~\ref{teo:bon-ordre-N}) i és únic: \(a\) és una aplicació
\(\N\to A\) definida sense cap tria (internament, el seu graf és la part
de \(\N\times A\) tallada per la condició, funcional i total). Per
construcció \(a_n\notin V_n\), mentre que \(a_m\in V_n\) per a tot
\(m<n\), perquè \(a_m\) és un valor de \(f_{2^m-1}\); per tant
\(a_n\ne a_m\) per a \(m<n\), i \(a\) és injectiva sobre punts, és a dir,
un monomorfisme (separació global).

Finalment, sigui \(C\) el complement de la imatge de \(a\). El primer
isomorfisme de
\[
A\;\cong\;\N + C\;\xrightarrow{\;s+\id_C\;}\;\N+C\;\cong\;A
\]
descompon \(A\) en aquesta imatge (isomorfa a \(\N\) via \(a\)) i \(C\); la
composició reprodueix el desplaçament del successor en el sumand \(\N\) i
omet el punt corresponent a \(0\), de manera que és un monomorfisme
\(A\hookrightarrow A\) no exhaustiu. Així \(A\) és infinit de Dedekind. La
reformulació és immediata: \(a\colon\N\hookrightarrow A\) és el testimoni, i
recíprocament un tal monomorfisme dona, transportant el successor, un
endomorfisme no exhaustiu de \(A\).
\end{prova}

\begin{observacio}[Per què cal l'elecció aquí]\label{obs:necessitat-eleccio-dedekind}
La prova anterior no pot evitar l'elecció numerable amb els axiomes que
tenim. Fixeu-vos on entra i on no. La família \((I_n)_n\) dels monomorfismes
\(\seg{s\circ n}\hookrightarrow A\) està definida \emph{abans} de tota
elecció i cada \(I_n\) és no inicial; això és una afirmació d'existència
fibra a fibra, provada per inducció sense triar res. El que \emph{no}
tenim, sense elecció, és una \emph{secció} uniforme d'aquesta família ---una
tria simultània d'un \(f_n\) a cada fibra---: en cada \(I_n\) hi ha molts
monomorfismes i res no en distingeix cap de canònic. És exactament aquest
pas, i només aquest, el que consumeix l'axioma~\ref{ax:eleccio-numerable};
un cop tenim la successió \((f_n)_n\), la resta de la construcció (el pas al
mínim índex nou dins una llista prou llarga, el desplaçament del
successor) torna a ser canònica. Es
tracta d'un ús genuí i, en el marc habitual, inevitable: sense cap forma
d'elecció és consistent que existeixin conjunts infinits ---no isomorfs a
cap \([k]\) estàndard--- que \emph{no} siguin infinits de Dedekind, els
anomenats conjunts infinits \emph{amb finitud de Dedekind}. La direcció
fàcil (proposició~\ref{prop:dedekind-implica-infinit}), en canvi, no
necessita res. Aquesta asimetria és el motiu pel qual mantenim les dues
nocions d'infinitud amb noms diferents.
\end{observacio}

Enunciem l'axioma que hem fet servir, en la forma que ens caldrà després.

\begin{axioma}[Elecció numerable]\label{ax:eleccio-numerable}
Sigui \(p\colon E\twoheadrightarrow\N\) un epimorfisme regular ---una
família de conjunts no buits indexada per \(\N\), on la fibra sobre cada
\(n\) és la part no buida corresponent. Aleshores \(p\) admet una secció
\(\sigma\colon\N\to E\), és a dir, una aplicació amb \(p\circ\sigma=\id_\N\)
que tria un element de cada fibra.
\end{axioma}

\begin{observacio}[Elecció, elecció numerable i una nota terminològica]\label{obs:eleccio-global-vs-numerable}
L'\emph{axioma d'elecció} del llibre (capítol~\ref{cap:relacions}) diu que
\emph{tot} epimorfisme regular s'escindeix; l'elecció numerable
(axioma~\ref{ax:eleccio-numerable}) n'és el cas particular en què l'índex
és \(\N\), i per tant se'n dedueix. Si el llibre assumeix l'elecció plena,
l'axioma~\ref{ax:eleccio-numerable} no és una hipòtesi nova sinó una
instància de la que ja hi ha, i el
teorema~\ref{teo:infinit-implica-dedekind} és un teorema sense reserves.
L'enunciem a part perquè la major part de la teoria d'aquest capítol només
en necessita la versió feble, i perquè convé veure exactament on entra
alguna forma d'elecció: en tot el que segueix, els resultats que no
l'esmenten no en depenen.

Una advertència de vocabulari, que serà important en els capítols
transfinits. Anomenem aquest principi simplement \emph{axioma d'elecció};
evitem el terme «elecció global», que en teoria de conjunts i de classes
designa un principi més fort ---una funció d'elecció única sobre tota la
classe dels conjunts--- que aquí ni enunciem ni necessitem. Igualment, cal
no confondre l'elecció amb cap principi de \emph{substitució} o de reunió de
famílies: l'elecció tria elements dins les fibres d'una família \emph{que ja
existeix} (com \(p\colon E\to\N\) a la prova anterior), mentre que la
substitució serviria per garantir l'existència d'una família d'objectes
definida uniformement. Són funcions lògiques distintes, i en aquest capítol
només usem la primera.
\end{observacio}

\section{Conjunts enumerables i numerables}\label{sec:enumerables}

Un conjunt és enumerable quan els seus elements es poden \emph{llistar}
amb els naturals, potser amb repeticions. La manera estructural de dir-ho
és que hi ha una aplicació exhaustiva des de \(\N\).

\begin{definicio}[Conjunt enumerable]\index{Conjunt enumerable}\label{def:enumerable}
Un conjunt \(A\) és \emph{enumerable} si és inicial o bé existeix un
epimorfisme regular \(e\colon\N\twoheadrightarrow A\). Una tal \(e\) és una
\emph{enumeració} de \(A\): l'element \(a_n:=e\circ n\) és el terme
\(n\)-èsim, i l'exhaustivitat diu que tot element generalitzat de \(A\)
és, localment, algun \(a_n\).
\end{definicio}

\begin{definicio}[Conjunt numerable]\index{Conjunt numerable}\label{def:numerable}
Un conjunt és \emph{numerable} si és finit
(definició~\ref{def:conjunt-finit}) o enumerable infinit. Un conjunt
enumerable és, doncs, o bé finit o bé numerable infinit.
\end{definicio}

\begin{observacio}[Enumerable, sense elecció]\label{obs:enumerable-sense-eleccio}
La definició~\ref{def:enumerable} \emph{no} amaga cap elecció: demana una
única aplicació \(e\), no una successió de tries. Per això tota la
combinatòria d'enumerabilitat que en depèn ---imatges d'enumerables,
unions de dues enumeracions, producte de dues--- es fa sense cap axioma
nou. L'elecció numerable només reapareixerà quan calgui triar \emph{una
enumeració per a cadascun} d'una família infinita de conjunts, no quan en
tinguem una de sola.
\end{observacio}

Les operacions bàsiques preserven l'enumerabilitat. Reunim els fets
elementals, tots demostrables sense elecció.

\begin{proposicio}[Estabilitat de l'enumerabilitat]\label{prop:estabilitat-enumerable}
Siguin \(A\) i \(B\) enumerables.
\begin{enumerate}
\item Tota imatge d'\(A\) és enumerable: si \(g\colon A\to C\) és
exhaustiva, \(C\) és enumerable.
\item Tota part d'\(A\) és enumerable.
\item La suma \(A+B\) és enumerable.
\item El producte \(A\times B\) és enumerable.
\end{enumerate}
Cap dels quatre enunciats no requereix elecció.
\end{proposicio}

\begin{prova}
(1) Si \(e\colon\N\twoheadrightarrow A\) és una enumeració i \(g\) és
exhaustiva, la composició \(g\circ e\colon\N\twoheadrightarrow C\) ho és
(l'exhaustivitat es preserva per composició, capítol~\ref{cap:imatges}),
i és una enumeració de \(C\); el cas inicial és trivial.

(2) Sigui \(m\colon P\hookrightarrow A\) una part no inicial. Fixat un punt
\(p_0\colon\mathbf1\to P\) ---que existeix perquè \(P\) no és inicial---,
definim \(r\colon A\to P\) enviant els elements de \(P\) en si mateixos i
la resta a \(p_0\); això és una retracció de \(m\), construïda amb el cas
finit de decisió del complement (capítol~\ref{cap:valors}) i sense cap
tria infinita. Aleshores \(r\circ e\colon\N\twoheadrightarrow P\) és una
enumeració. El cas \(P\) inicial és immediat.

(3) De dues enumeracions \(e\colon\N\twoheadrightarrow A\) i
\(f\colon\N\twoheadrightarrow B\) i d'un isomorfisme
\(\N\cong\N+\N\) ---per exemple la partició en parells i senars, obtinguda
de l'aritmètica del capítol~\ref{cap:naturals}--- s'obté
\(\N\cong\N+\N\xrightarrow{e+f}A+B\), exhaustiva.

(4) N'hi ha prou amb una bijecció \(\N\cong\N\times\N\) ---un
\emph{aparellament}--- perquè aleshores
\(\N\cong\N\times\N\xrightarrow{e\times f}A\times B\) és exhaustiva. Un tal
aparellament es construeix explícitament dins la teoria \emph{sense divisió},
enumerant \(\N\times\N\) per diagonals: es recorre primer la diagonal
\(i+j=0\), després \(i+j=1\), i així successivament, i dins de cada diagonal
s'ordena per \(j\) creixent. Formalment, sigui \(\nu\colon\N\times\N\to\N\times\N\) l'aplicació
\enquote{següent parell en l'ordre per diagonals}, definida per casos
sobre el punt \(i\) (zero o successor, lema~\ref{lema:zero-o-successor}):
\(\nu(0,j)=(s\circ j,0)\) i \(\nu(s\circ i',j)=(i',s\circ j)\). La recursora
\(u\colon\N\to\N\times\N\) del sistema \((\N\times\N,\langle0,0\rangle,\nu)\)
recorre les diagonals, i és bijectiva: injectiva perquè \(\nu\) ho és i
\((0,0)\) no és a la seva imatge (el mateix argument que a les equacions
de Peano), i exhaustiva perquè la part dels \((i,j)\) assolits conté
\((0,0)\) i és tancada per \(\nu\), i tot parell s'obté de \((0,0)\)
iterant \(\nu\) (inducció sobre \(i+j\) i després sobre \(j\)). Els detalls
són un exercici de la part pràctica. No cal cap divisió ni la fórmula
tancada \(\tfrac12(i+j)(i+j+1)+j\), que és correcta però pressuposa
aritmètica que el llibre no ha desenvolupat. És, en tot cas, una
aplicació concreta, no una tria.
\end{prova}

\begin{observacio}[Lectura clàssica]\label{obs:enumerable-classica}
En termes clàssics, un conjunt no buit és enumerable exactament quan els
seus elements es poden disposar en una successió
\(a_0,a_1,a_2,\dots\) que els reculli tots ---amb repeticions permeses.
La proposició~\ref{prop:estabilitat-enumerable} recull els fets habituals:
la imatge d'un conjunt enumerable, qualsevol dels seus subconjunts, i la
unió i el producte de dos enumerables són enumerables; en particular
\(\Z\) i \(\Q\) ho són, com es comprova aparellant-los amb \(\N\). El punt
que la presentació estructural fa visible, i que sovint queda amagat, és
que aquests resultats no fan servir cap elecció: totes les enumeracions
involucrades es construeixen a partir d'una o dues enumeracions donades,
mitjançant aparellaments explícits.
\end{observacio}

\section{\texorpdfstring{Subobjectes de \(\N\)}{Subobjectes de N}}\label{sec:subobjectes-N}

Els subobjectes de \(\N\) ---les seves parts--- són l'exemple bàsic i
mereixen un tractament propi, perquè hi funciona una construcció que en
general falla: es pot triar sempre el \emph{mínim}. Això fa que tot
subobjecte de \(\N\) sigui enumerable sense recórrer a cap elecció.

\begin{teorema}[Els subobjectes de \(\N\) són enumerables]\label{teo:subobjectes-N}
Tota part \(P\hookrightarrow\N\) és enumerable, i la seva enumeració es
construeix sense elecció. A més, \(P\) és finit o bé isomorf a \(\N\); en
el segon cas, l'enumeració es pot prendre estrictament creixent i
bijectiva.
\end{teorema}

\begin{prova}
Si \(P\) és inicial, és enumerable per definició. Suposem \(P\) no
inicial. La clau és el \emph{principi del mínim}: tota part no inicial de
\(\N\) té un element mínim, la qual cosa és una reformulació de la
inducció (teorema~\ref{teo:induccio}) i no usa elecció, perquè el mínim,
quan existeix, és \emph{únic} i per tant canònic. Definim una autoaplicació \(t\colon P\to P\): \(t(x)\) és el mínim dels
elements de \(P\) estrictament més grans que \(x\) si n'hi ha, i \(x\)
altrament. És una aplicació genuïna, perquè el seu graf ---la part de
\(P\times P\) dels \((x,y)\) amb \(y\in P\), \(x<y\) i
\(\forall z\in P\,(x<z\Rightarrow y\le z)\), unida a la part dels \((x,x)\)
amb \(\neg\exists z\in P\,(x<z)\)--- és funcional (el mínim és únic) i
total (o hi ha elements més grans o no), i es talla amb els
quantificadors interns. La successió \(a\colon\N\to P\) és la recursora
del sistema \((P,\min P,t)\): \(a_0\) és el mínim de \(P\) i
\(a_{n+1}=t(a_n)\). És una recursió simple de l'axioma~\ref{ax:naturals},
sense cap tria.

Dos casos. Si per a algun \(n\) no queda cap element de \(P\) més gran que
\(a_n\), la successió s'estabilitza i la seva imatge
\(\{a_0<a_1<\dots<a_n\}\) és exactament \(P\), que és finit. Si en cada pas
sempre en queda un de més gran, \(a\) és estrictament creixent, per tant
un monomorfisme, i la seva imatge és tot \(P\) perquè cap element de \(P\)
pot ser omès. En efecte, \(a_i\ge i\) per a tot \(i\) (inducció: \(a_{i+1}>a_i\)),
de manera que, donat \(k\in P\), la part dels \(i\) amb \(k<a_i\) és no
inicial (conté \(s\circ k\)) i té un mínim \(j\); com que \(a_0=\min P\le k\),
\(j=i+1\) per a algun \(i\), amb \(a_i\le k<a_{i+1}\). Si fos \(a_i<k\), aleshores
\(k\) seria un element de \(P\) estrictament més gran que \(a_i\) i menor que
\(a_{i+1}=t(a_i)\), contra la definició de \(t\) com a mínim d'aquests
elements; per tant \(k=a_i\). Aleshores
\(a\colon\N\to P\) és una bijecció creixent, i \(P\cong\N\).
\end{prova}

\begin{corollari}[Caracterització dels enumerables infinits]\label{cor:enumerable-infinit-N}
Un conjunt és enumerable infinit si i només si és isomorf a \(\N\). En
conseqüència, tots els conjunts numerables infinits són isomorfs entre si,
i \(\N\) n'és el representant.
\end{corollari}

\begin{prova}
Si \(A\cong\N\), és clarament enumerable (per la identitat) i infinit
(proposició~\ref{prop:N-dedekind}). Recíprocament, sigui \(A\) enumerable
infinit i \(e\colon\N\twoheadrightarrow A\) una enumeració. Volem un
monomorfisme \(A\hookrightarrow\N\); un cop el tinguem, la seva imatge és
una part de \(\N\) isomorfa a \(A\), que pel teorema~\ref{teo:subobjectes-N}
és finita o isomorfa a \(\N\), i com que \(A\) és infinit ha de ser el
segon cas. Per construir el monomorfisme sense elecció, definim
\(r\colon A\to\N\) enviant cada element \(a\) al \emph{mínim} índex \(n\)
amb \(e_n=a\); aquest mínim existeix i és únic pel principi del mínim, de
manera que \(r\) està ben definida i és una secció de \(e\), en particular
un monomorfisme. Per tant \(A\cong\N\).
\end{prova}

\begin{observacio}[El mínim substitueix l'elecció]\label{obs:minim-vs-eleccio}
Els dos resultats anteriors s'obtenen sense elecció \emph{precisament}
perquè a \(\N\) hi ha una tria canònica ---el mínim--- disponible a cada
pas. És el mateix fenomen que fa que l'elecció sigui innecessària sobre
conjunts ben ordenats: quan cada part no buida té un element distingit, no
cal cap axioma per triar-lo. En un conjunt infinit qualsevol, sense un tal
ordre, aquesta comoditat desapareix, i és per això que el
teorema~\ref{teo:infinit-implica-dedekind} sí que necessita l'elecció
numerable mentre que el teorema~\ref{teo:subobjectes-N} no.
\end{observacio}

\section{Llistes finites i iteració}\label{sec:llistes}

Una \emph{llista finita} d'elements d'un conjunt \(A\) és una successió
finita \((a_0,\dots,a_{n-1})\) de llargada arbitrària però determinada.
Reunir totes les llistes de totes les llargades en un sol conjunt és la
construcció que fa d'\(A\) el generador d'un monoide lliure, i és el
prototip d'una unió indexada per \(\N\) que, en aquest cas, es fa sense
elecció.

\begin{definicio}[Conjunt de llistes finites]\index{Conjunt de llistes finites}\label{def:llistes}
El \emph{conjunt de les llistes finites} d'\(A\), denotat
\(\operatorname{List}(A)\), es defineix com una part de l'exponencial
\emph{fix} \((A+\mathbf1)^{\N}\) de les successions amb un valor de
terminació. Escrivim \(i_A\colon A\to A+\mathbf1\) i
\(\ast\colon\mathbf1\to A+\mathbf1\) per a les injeccions. Una successió
\(u\colon\N\to A+\mathbf1\) és una \emph{llista} quan el conjunt de
posicions \emph{ocupades} ---les \(k\) amb \(u(k)=i_A(a)\) per a algun
\(a\)--- és un segment inicial finit; és a dir, quan existeix
\(n\colon\mathbf1\to\N\) tal que
\begin{enumerate}
\item[(l1)] per a tot \(k<n\), \(u(k)=i_A(a)\) per a un únic \(a\in A\);
\item[(l2)] per a tot \(k\) amb \(k\not<n\), \(u(k)=\ast\).
\end{enumerate}
Aquesta \(n\) és única quan existeix (és la \emph{llargada}), i les dues
condicions es tallen amb quantificadors acotats i igualtats sobre
l'exponencial fix (capítol~\ref{cap:parts-interna}), de manera que
\[
\operatorname{List}(A)\ :=\ \bigl\{\,u\in(A+\mathbf1)^{\N}\ \mid\ u\text{
satisfà (l1)--(l2) per a algun }n\,\bigr\}\ \hookrightarrow\
(A+\mathbf1)^{\N}
\]
és una part interna ben definida. La successió constant \(\ast\)
(cas \(n=0\)) és la llista buida
\(\varepsilon\colon\mathbf1\to\operatorname{List}(A)\).
\end{definicio}

\begin{observacio}[La condició de segment inicial és essencial]\label{obs:llistes-segment}
No n'hi ha prou d'exigir que \(u\) sigui \emph{eventualment constant} igual
a \(\ast\): aquesta condició, més feble, admet successions com
\(a,\ast,b,\ast,\ast,\dots\), amb una posició de terminació seguida de nous
elements d'\(A\), que no representen cap llista. La condició (l1)--(l2)
demana que \emph{totes} les posicions ocupades formin un segment inicial
\(\seg{n}\): un cop apareix el primer \(\ast\), tots els termes següents
són \(\ast\). Amb aquesta condició, i només amb ella, la llargada està ben
definida i la lectura
\[
\operatorname{List}(A)\ \cong\ \sum_{n\colon\mathbf1\to\N}A^{\seg{n}}
\;=\;\coprod_{n\in\N}A^{\seg{n}}
\]
és correcta: cada llista de llargada \(n\) es correspon amb la seva
restricció \(\seg{n}\to A\), i \(A^{\seg{0}}\cong\mathbf1\) és la llista
buida. Aquesta igualtat és, doncs, un \emph{teorema} sobre la part just
construïda, no la definició: la suma indexada per \(\N\) del membre dret no
es pot formar directament, perquè els axiomes només donen sumes finites
(capítol~\ref{cap:valors}). Vegeu l'observació~\ref{obs:existencia-llistes}.
\end{observacio}

\begin{observacio}[Sobre l'existència i les dues lectures]\label{obs:existencia-llistes}
La definició~\ref{def:llistes} construeix \(\operatorname{List}(A)\)
\emph{directament}, com una part d'un exponencial fix, i per tant la seva
existència no demana cap suma numerable ni cap principi nou: només
l'exponencial \((A+\mathbf1)^{\N}\) i els quantificadors interns
(capítols~\ref{cap:exponencials} i~\ref{cap:parts-interna}). Aquesta és la
via segura. Les altres dues descripcions habituals ---la suma
\(\coprod_{n}A^{\seg{n}}\) i l'objecte inicial de les
\((\mathbf1+A\times{-})\)-àlgebres--- \emph{no} es prenen aquí com a
definicions: la primera perquè la suma numerable no és una construcció
finita (observació~\ref{obs:llistes-segment}), i la segona perquè afirmar
que aquesta àlgebra inicial \emph{existeix} és precisament el que s'ha de
demostrar. Totes dues es recuperen com a \emph{propietats} de la part ja
construïda: la lectura com a suma és l'observació~\ref{obs:llistes-segment},
i la inicialitat és el contingut de la
proposició~\ref{prop:recursio-llistes}, que provem tot seguit sobre
\(\operatorname{List}(A)\) equipat amb \(\varepsilon\) i \(\operatorname{cons}\).
\end{observacio}

\begin{proposicio}[Recursió i iteració sobre llistes]\label{prop:recursio-llistes}
El conjunt \(\operatorname{List}(A)\) ve equipat amb el punt
\(\varepsilon\colon\mathbf1\to\operatorname{List}(A)\) de la llista buida i
amb l'operació de prefixar un element,
\(\operatorname{cons}\colon A\times\operatorname{List}(A)\to
\operatorname{List}(A)\), i satisfà el principi de recursió corresponent:
per a tot conjunt \(B\) amb un punt \(b\colon\mathbf1\to B\) i una aplicació
\(t\colon A\times B\to B\), existeix una única
\(h\colon\operatorname{List}(A)\to B\) amb
\[
h\circ\varepsilon=b
\qquad\text{i}\qquad
h\circ\operatorname{cons}=t\circ(\id_A\times h).
\]
En particular, la llargada
\(\ell\colon\operatorname{List}(A)\to\N\), la concatenació i la iteració
d'una operació al llarg d'una llista es defineixen totes per aquesta
recursió, sense cap axioma nou.
\end{proposicio}

\begin{prova}
\emph{L'operació \(\operatorname{cons}\).} Sobre la representació de la
definició~\ref{def:llistes}, prefixar \(a\in A\) a una llista \(u\) de
llargada \(n\) és la successió \(\operatorname{cons}(a,u)\) de llargada
\(n+1\) que val \(i_A(a)\) a la posició \(0\) i \(u(k-1)\) ---és a dir, el
desplaçament de \(u\)--- a la posició \(k\) per a \(1\le k<n+1\), i \(\ast\)
a la resta. El desplaçament que reindexa el terme de la posició \(k\) cap a
la posició \(k+1\), i la reassignació de la posició \(0\), són aplicacions
internes sobre
\((A+\mathbf1)^{\N}\) (avaluació, predecessor acotat i definició per casos,
capítols~\ref{cap:exponencials} i~\ref{cap:parts-interna}), i preserven la
condició (l1)--(l2); per tant
\(\operatorname{cons}\colon A\times\operatorname{List}(A)\to
\operatorname{List}(A)\) està ben definida. La llargada compleix
\(\ell(\varepsilon)=0\) i \(\ell(\operatorname{cons}(a,u))=s(\ell(u))\).

\emph{Existència de \(h\).} Escrivim \(L:=\operatorname{List}(A)\). Les
aplicacions \(h_m\colon L\to B\) que \enquote{encerten} sobre les llistes
de llargada \(\le m\) no les construïm amb dominis variables, sinó com a
punts de l'exponencial fix \(B^{L}\), per recursió simple. Sigui
\(\Phi\colon B^{L}\to B^{L}\) la transposada de l'aplicació
\(B^{L}\times L\to B\) definida per casos sobre la descomposició
\(L\cong\mathbf1+A\times L\) (llista buida, o primer element i cua,
recuperats internament, amb \(\operatorname{cons}\) com a isomorfisme):
\[
(h,\varepsilon)\mapsto b,
\qquad
(h,\operatorname{cons}(a,u))\mapsto t\bigl(a,\operatorname{ev}(h,u)\bigr).
\]
Sigui \(H\colon\N\to B^{L}\) la recursora del sistema
\((B^{L},\ulcorner b\circ{!}_L\urcorner,\Phi)\), de manera que
\(H_{0}\) és la constant \(b\) i \(H_{m+1}=\Phi(H_m)\). Per inducció
sobre \(m\) (amb la llargada com a paràmetre), \(H_m\) coincideix amb
\(H_{m'}\) sobre les llistes de llargada \(\le m\) per a tot \(m'\ge m\):
per a \(m=0\), totes valen \(b\) sobre \(\varepsilon\); i si val per a
\(m\), aleshores sobre \(\operatorname{cons}(a,u)\) amb \(\ell(u)\le m\)
es té \(H_{m'+1}(\operatorname{cons}(a,u))=t(a,H_{m'}(u))=t(a,H_m(u))
=H_{m+1}(\operatorname{cons}(a,u))\). Definim \(h\) per la
\emph{diagonal}: \(h:=\operatorname{ev}\circ\langle H\circ\ell,\id_L\rangle\),
és a dir, \(h(u):=H_{\ell(u)}(u)\). Aleshores \(h(\varepsilon)=H_0(\varepsilon)=b\),
i \(h(\operatorname{cons}(a,u))=H_{\ell(u)+1}(\operatorname{cons}(a,u))
=t(a,H_{\ell(u)}(u))=t(a,h(u))\), que són les dues equacions.

\emph{Unicitat.} Si \(h'\) compleix les mateixes dues equacions, aleshores
\(h=h'\) sobre les llistes de llargada \(0\) (totes dues valen \(b\)), i si
coincideixen sobre les de llargada \(\le m\), l'equació del
\(\operatorname{cons}\) força la coincidència sobre les de llargada \(m+1\);
per inducció sobre \(m\) (teorema~\ref{teo:induccio}; la part dels \(m\)
sobre els quals coincideixen és interna, amb un quantificador acotat per
\(\ell\)), \(h=h'\) a tot arreu.
Això estableix la propietat universal ---la inicialitat de
\((\operatorname{List}(A),\varepsilon,\operatorname{cons})\) entre les
\((\mathbf1+A\times{-})\)-àlgebres--- \emph{a partir} de la construcció, i
no a l'inrevés. La llargada s'obté amb \(B=\N\), \(b=0\), \(t=s\circ\pi_2\);
la concatenació amb una llista fixada iterant \(\operatorname{cons}\); i
qualsevol plegat (\emph{fold}) d'una operació binària, amb el \(t\)
corresponent.
\end{prova}

\begin{observacio}[Enumerabilitat de \(\operatorname{List}(A)\)]\label{obs:llistes-enumerables}
Si \(A\) és enumerable, també ho és \(\operatorname{List}(A)\), i sense
elecció. L'argument és en dos passos. Primer, \(\operatorname{List}(\N)\)
és enumerable: la codificació
\(c\colon\operatorname{List}(\N)\to\N\) definida per recursió sobre llistes
(proposició~\ref{prop:recursio-llistes}) amb \(c(\varepsilon)=0\) i
\(c(\operatorname{cons}(k,u))=s\bigl(\pi(k,c(u))\bigr)\), on
\(\pi\colon\N\times\N\to\N\) és un aparellament bijectiu
(proposició~\ref{prop:estabilitat-enumerable}), és bijectiva: injectiva
perquè \(\pi\) i \(s\) ho són i \(0\) no és successor, i exhaustiva per
inducció sobre el codi. Segon, una enumeració \(e\colon\N\twoheadrightarrow A\)
indueix \(\operatorname{List}(e)\colon\operatorname{List}(\N)\to\operatorname{List}(A)\)
(aplicar \(e\) terme a terme, també per recursió sobre llistes), que és
exhaustiva sobre punts: tota llista \((a_0,\dots,a_{n-1})\) és la imatge
de la llista dels mínims índexs \(k_i\) amb \(e(k_i)=a_i\), que es
defineixen sense cap tria. Així
\(\operatorname{List}(e)\circ c^{-1}\colon\N\twoheadrightarrow\operatorname{List}(A)\)
és una enumeració. El que fa possible evitar l'elecció és que les
enumeracions de totes les llistes deriven \emph{uniformement} d'una sola
enumeració d'\(A\) i d'un sol aparellament; és l'observació clau que
separa aquest cas del de la secció següent.
\end{observacio}

\section{Unions indexades numerablement}\label{sec:unions-numerables}

L'últim tema és el que fa explícit el paper de l'elecció numerable en la
seva forma més habitual. Es tracta de reunir una família de conjunts
\((A_n)_{n\in\N}\) en una sola unió i de decidir quan aquesta unió és
enumerable.

\begin{definicio}[Família i unió numerables]\index{Família i unió numerables}\label{def:familia-numerable}
Una \emph{família numerable} de conjunts és una aplicació
\(p\colon E\to\N\); el conjunt \(A_n\) és la fibra
\(p^{-1}(n)\hookrightarrow E\) sobre el punt \(n\)
(capítol~\ref{cap:parts}). La \emph{unió} de la família és la imatge de
\(E\) per una aplicació cap a un conjunt ambient, o simplement \(E\) mateix
quan es pensa com la reunió disjunta \(\coprod_{n\in\N}A_n\); la
\emph{unió indexada} \(\bigcup_{n\in\N}A_n\) dins un conjunt \(U\) que
conté totes les fibres és la imatge (capítol~\ref{cap:imatges}) de \(E\to
U\).
\end{definicio}

\begin{teorema}[Unió numerable d'enumerables, amb elecció numerable]\label{teo:unio-numerable}
Suposem l'axioma d'elecció numerable
(axioma~\ref{ax:eleccio-numerable}). Sigui \((A_n)_{n\in\N}\) una família
numerable en què cada \(A_n\) és enumerable. Aleshores la unió disjunta
\(\coprod_{n\in\N}A_n\) és enumerable, i per tant també ho és qualsevol
imatge seva ---en particular la unió indexada dins un conjunt ambient.
\end{teorema}

\begin{prova}
Per hipòtesi, per a cada \(n\) el conjunt \(A_n=p^{-1}(n)\) és enumerable:
o bé és inicial, o bé \emph{existeix} un epimorfisme
\(\N\twoheadrightarrow A_n\). Reunim totes les enumeracions de totes les
fibres en una sola part d'un objecte fix:
\[
E':=\bigl\{\,(n,e)\in\N\times E^{\N}\ :\ p\circ e=n\circ{!}_\N
\ \text{ i }\ \forall x\in E\,\bigl(p(x)=n\Rightarrow\exists k\,(e(k)=x)\bigr)\bigr\},
\]
tallada amb els quantificadors interns (capítol~\ref{cap:parts-interna});
la fibra de \(E'\to\N\) sobre \(n\) és el conjunt de les enumeracions de
\(A_n\), no buida exactament quan \(A_n\) no és inicial. Sigui
\(N'\hookrightarrow\N\) la part dels \(n\) amb \(A_n\) no inicial (la imatge
de \(p\)). Si \(N'\) és inicial, la unió és inicial i no hi ha res a
provar. Altrament, \(N'\) és enumerable sense elecció
(teorema~\ref{teo:subobjectes-N}): sigui \(r\colon\N\twoheadrightarrow N'\)
una enumeració. La família \(E'\) reindexada per \(r\) és una família
\(E''\twoheadrightarrow\N\) amb totes les fibres no buides, i l'axioma
d'elecció numerable en dona una secció: una \emph{tria simultània} d'una
enumeració \(e_{r(j)}\colon\N\twoheadrightarrow A_{r(j)}\) per a cada
\(j\). Per a cada \(n\in N'\) prenem \(e_n:=e_{r(j_n)}\), on \(j_n\) és
el \emph{mínim} índex \(j\) amb \(r(j)=n\), que existeix perquè \(r\) és
exhaustiva i és únic pel bon ordre de \(\N\): així l'assignació
\(n\mapsto e_n\) és explícita i no amaga cap elecció addicional. Amb aquesta tria, l'aplicació
\[
\N\times\N\longrightarrow\coprod_{n\in\N}A_n,
\qquad
(j,k)\longmapsto e_{r(j)}(k),
\]
és exhaustiva sobre punts (tot punt de la unió és a alguna fibra
\(A_{r(j)}\), i allà és algun \(e_{r(j)}(k)\)), per tant epimorfisme, i component-la amb un aparellament
\(\N\cong\N\times\N\) s'obté una enumeració de la unió disjunta. Tota
imatge d'aquesta unió és enumerable per la
proposició~\ref{prop:estabilitat-enumerable}(1).
\end{prova}

\begin{observacio}[On és l'elecció, i on no]\label{obs:eleccio-unio}
El pas irreductible és la tria simultània de les \(e_n\). Sabem que
cada \(A_n\) enumerable \emph{té} alguna enumeració, però la definició
d'enumerabilitat n'afirma només l'existència; res no en selecciona cap de
canònica quan els \(A_n\) són conjunts arbitraris. Encadenar aquestes
infinites eleccions en una sola aplicació \(n\mapsto e_n\) és precisament
el contingut de l'axioma~\ref{ax:eleccio-numerable}. Contrasta amb dos
casos ja vistos on l'elecció \emph{no} feia falta: la unió de \emph{dos}
enumerables (proposició~\ref{prop:estabilitat-enumerable}), on hi ha només
dues enumeracions donades; i el conjunt de llistes
(observació~\ref{obs:llistes-enumerables}), on les enumeracions dels
sumands eren totes derivades \emph{uniformement} d'una de sola. La
diferència no és la mida de la família ---en tots dos casos infinita en el
segon--- sinó si les enumeracions vénen donades o s'han de triar. De fet,
sense cap elecció és consistent que una unió numerable de conjunts de dos
elements no sigui enumerable: la conclusió del
teorema~\ref{teo:unio-numerable} pot fallar, i per tant l'axioma hi és
essencial.
\end{observacio}

\begin{observacio}[El producte numerable]\label{obs:producte-numerable}
El producte d'una família numerable, \(\prod_{n\in\N}A_n\), planteja una
qüestió d'elecció encara més aguda. El producte existeix com a objecte
---és un límit, i els límits numerables es poden construir a partir dels
exponencials i de \(\N\)---, però la seva \emph{no-vacuïtat} equival, per a
famílies de conjunts no buits, a una forma d'elecció: dir que
\(\prod_{n\in\N}A_n\) té un punt global és exactament triar un element de
cada \(A_n\). Per aquest producte l'elecció numerable és, doncs, no una
comoditat de la prova sinó la substància de l'enunciat. És el motiu pel
qual, en teories sense elecció, s'ha de distingir amb cura entre
l'existència d'un producte com a objecte i l'existència d'elements seus.
\end{observacio}

\begin{resumcap}
\apartat{Idees centrals}
\begin{enumerate}
  \item \emph{Infinit} és la negació d'\emph{finit}; \emph{infinit de Dedekind} és admetre un endomonomorfisme no exhaustiu. \(\N\) ho és pel successor; Dedekind implica infinit pel cas finit del capítol anterior, sense cap axioma.
  \item Infinit implica Dedekind amb l'elecció numerable, i només per un pas: triar alhora una llista injectiva de cada llargada. La successió injectiva es construeix després sense tria, prenent en llistes de llargada \(2^n\) el primer valor no usat.
  \item Enumerable és ser imatge de \(\N\); l'estabilitat per imatge, part, suma i producte és sense elecció, amb un aparellament \(\N\cong\N\times\N\) definit per recursió sobre les diagonals.
  \item Tota part de \(\N\) és enumerable, finita o isomorfa a \(\N\), perquè el mínim és una tria canònica; d'aquí, tot enumerable infinit és isomorf a \(\N\).
  \item \(\operatorname{List}(A)\) és una part d'\((A+\mathbf1)^{\N}\); el seu principi de recursió es demostra pujant a \(B^{\operatorname{List}(A)}\) i prenent la diagonal sobre la llargada. És enumerable, uniformement, si \(A\) ho és.
  \item La unió numerable d'enumerables necessita l'elecció numerable (i és consistent que falli sense ella); el producte numerable la incorpora al seu enunciat.
\end{enumerate}
\apartat{Errors típics} Confondre els dos usos de la paraula: \emph{enumerable} i \emph{numerable} qualifiquen conjunts (imatge de \(\N\); finit o enumerable infinit), mentre que \enquote{família, unió o producte numerable} i \enquote{elecció numerable} volen dir \emph{indexat per \(\N\)}; identificar infinit i infinit de Dedekind sense elecció; definir una successió \enquote{pel primer element no usat} com si fos una recursió de l'axioma (depèn de tota la història); creure que la unió de dos enumerables i la de numerablement molts estan al mateix nivell axiomàtic; pensar que \(\operatorname{List}(A)\) és \enquote{la suma \(\coprod_nA^{\seg n}\)}, que els axiomes no formen directament; i confondre l'existència del producte numerable com a objecte amb l'existència dels seus punts.
\apartat{Lectura recomanada} \cite[§§7.6--7.8]{lawvere-rosebrugh} per a Dedekind-infinitud, enumerabilitat i el paper de l'elecció; \cite[§§13, 15--16 i 23]{halmos} per als conjunts infinits, l'elecció i els conjunts numerables en la teoria clàssica; \cite[apèndix~B]{lawvere-rosebrugh} per a les formes de l'elecció; i \cite[pàg.~12]{bell2006} per a la definició de Dedekind i l'axioma d'infinitud com a existència d'un conjunt \(X\cong X+1\).
\end{resumcap}

\chapter[Schröder--Bernstein i Cantor]{Comparació de conjunts,\texorpdfstring{\\}{ }Schröder--Bernstein i Cantor}\label{cap:cantor}

\begin{objectius}
\apartat{Prerequisits} Del capítol~\ref{cap:parts-interna}: conjunt de parts, pertinença i inclusió internes, el quantificador universal al llarg d'una projecció, i l'ordre \((P(A),\subseteq)\). Del capítol~\ref{cap:imatges}: imatge directa i la seva adjunció. Del capítol~\ref{cap:valors}: complements i definició per casos. Dels capítols~\ref{cap:parts} i~\ref{cap:valors}: que mono i epi es comproven sobre punts globals.
\apartat{Objectius} En acabar el capítol, el lector hauria de poder:
\begin{itemize}
  \item definir la comparació \(A\preceq B\) i l'equipotència com a relacions del metallenguatge, i demostrar que \(\preceq\) és un preordre invariant;
  \item demostrar Schröder--Bernstein construint el menor punt fix d'un operador creixent sobre \(P(A)\) amb el quantificador universal, i explicar per què no hi intervé cap elecció;
  \item demostrar que cap aplicació \(A\to P(A)\) és exhaustiva (argument diagonal) i deduir-ne el teorema de Cantor i la inexistència d'un conjunt universal;
  \item distingir \enquote{no hi ha exhaustiva} de \enquote{no hi ha isomorfisme}, i situar la paradoxa de Russell com el que el llibre evita en no tenir pertinença global.
\end{itemize}
\end{objectius}

Els deu capítols anteriors han construït tota la infraestructura dels
conjunts ---parts, productes i sumes, quocients, imatges, exponencials i
el conjunt de parts, els nombres naturals i la cardinalitat fins al
numerable--- i han comparat conjunts sobretot en el cas finit i en el
numerable. Ara comença la comparació \emph{general} de mides, i el primer
canvi és de punt de vista. Un cardinal no serà, d'entrada, un conjunt
privilegiat que representa una mida, sinó una manera de \emph{comparar}
conjunts: un conjunt és més petit o igual que un altre quan s'hi injecta, i
tots dos tenen la mateixa mida quan són isomorfs. Amb només aquest
vocabulari ---monomorfismes i isomorfismes, res més que el que ja tenim---
s'obtenen els dos teoremes bàsics de la teoria de Cantor. El primer,
\emph{Schröder--Bernstein}, diu que dues injeccions en sentits oposats
basten per construir un isomorfisme: la comparació de mides és
antisimètrica. El segon, el \emph{teorema de Cantor}, diu que el pas al
conjunt de parts sempre produeix una mida estrictament més gran, de manera
que cap univers de conjunts no queda esgotat per una sola mida. Cap dels
dos no necessita cap axioma nou; en particular, cap dels dos no necessita
una successió prèvia d'alephs ni l'elecció. Aquest és el capítol més curt de
la part, i el que fixa el llenguatge de tots els que segueixen.

\section{La comparació de conjunts}\label{sec:comparacio-conjunts}

Recordem que dos conjunts són \emph{equipotents}, \(A\approx B\), quan hi ha
un isomorfisme \(A\cong B\) (definició~\ref{def:equipotencia}). L'equipotència
és la igualtat de mida. La comparació que falta ---quan una mida és
\emph{menor o igual} que una altra--- s'expressa amb un monomorfisme.

\begin{definicio}[Comparació de conjunts]\index{Comparació de conjunts}\label{def:comparacio}
Siguin \(A\) i \(B\) conjunts.
\begin{itemize}
\item \(A\preceq B\) (\emph{\(A\) es domina per \(B\)}) si existeix un
monomorfisme \(A\hookrightarrow B\).
\item \(A\approx B\) (\emph{\(A\) és equipotent a \(B\)}) si existeix un
isomorfisme \(A\cong B\) (definició~\ref{def:equipotencia}).
\item \(A\prec B\) (\emph{\(A\) és estrictament menor que \(B\)}) si
\(A\preceq B\) però \(A\not\approx B\).
\end{itemize}
\end{definicio}

\begin{observacio}[Una comparació entre conjunts, no una relació interna]\label{obs:comparacio-metateoria}
Com l'equipotència (definició~\ref{def:equipotencia}), les relacions
\(\preceq\), \(\approx\) i \(\prec\) són afirmacions \emph{sobre} conjunts,
formulades al metallenguatge: diuen que existeix una certa aplicació entre
dos conjunts donats. No són parts internes d'un producte ni relacions sobre
un suposat «conjunt de tots els conjunts», que no existeix
(secció~\ref{sec:cantor}). Aquesta distinció, ja feta per a \(\approx\), és
la que permet parlar de mides sense convertir la pertinença global en una
noció primitiva; la mantindrem al llarg de tota la part transfinita.
\end{observacio}

La comparació és un preordre ---reflexiva i transitiva--- i respecta
l'equipotència. Són fets immediats, però convé enunciar-los perquè fixen que
\(\preceq\) mesura una mida i no una còpia concreta.

\begin{proposicio}[La comparació és un preordre invariant]\label{prop:preceq-preordre}
Per a conjunts qualssevol:
\begin{enumerate}
\item \(A\preceq A\) (reflexivitat).
\item Si \(A\preceq B\) i \(B\preceq C\), aleshores \(A\preceq C\)
(transitivitat).
\item Si \(A\approx A'\), \(B\approx B'\) i \(A\preceq B\), aleshores
\(A'\preceq B'\) (invariància per equipotència).
\item Si \(A\approx B\), aleshores \(A\preceq B\) i \(B\preceq A\).
\end{enumerate}
Cap enunciat no usa cap axioma més enllà dels de composició.
\end{proposicio}

\begin{prova}
(1) La identitat \(\id_A\colon A\to A\) és un monomorfisme
(definició~\ref{def:monomorfisme}). (2) La composició de dos monomorfismes és
un monomorfisme: si \(m\colon A\hookrightarrow B\) i \(n\colon
B\hookrightarrow C\) són monomorfismes i \((n\circ m)\circ f=(n\circ m)\circ
g\), aleshores \(n\circ(m\circ f)=n\circ(m\circ g)\) dona \(m\circ f=m\circ
g\) i, altre cop, \(f=g\). (3) Siguin \(u\colon A'\cong A\), \(v\colon B\cong
B'\) isomorfismes i \(m\colon A\hookrightarrow B\) un monomorfisme; aleshores
\(v\circ m\circ u\colon A'\to B'\) és un monomorfisme, perquè els
isomorfismes són monomorfismes (composició amb la inversa) i la classe dels
monomorfismes és tancada per composició. (4) Un isomorfisme és alhora un
monomorfisme i té una inversa, que també ho és; d'aquí \(A\preceq B\) i
\(B\preceq A\).
\end{prova}

L'enunciat que \emph{falta} a la proposició anterior és el recíproc del
punt~(4): si \(A\preceq B\) i \(B\preceq A\), és \(A\approx B\)? És a dir, la
comparació és \emph{antisimètrica} llevat d'equipotència? Que la resposta
sigui afirmativa ---i sense cap axioma d'elecció--- és el teorema de
Schröder--Bernstein, l'objecte de la secció següent.

\section{El teorema de Schröder--Bernstein}\label{sec:schroder-bernstein}

Dues injeccions en sentits oposats, \(f\colon A\hookrightarrow B\) i
\(g\colon B\hookrightarrow A\), no donen directament cap isomorfisme: no hi
ha cap raó perquè \(f\) sigui exhaustiva, ni \(g\). El que sí que es pot fer
és \emph{combinar-les} en una sola bijecció, decidint per a cada element
d'\(A\) si l'enviem per \(f\) o pel tros invertible de \(g\). La decisió es
pren identificant una part distingida d'\(A\) ---els «punts bons»--- sobre la
qual cal usar \(f\), i el seu complement, sobre el qual cal usar
\(g^{-1}\). Tota la construcció viu dins \(P(A)\) i no fa servir cap tria.

\begin{teorema}[Schröder--Bernstein]\index{Schröder--Bernstein, teorema de}\label{teo:schroder-bernstein}
Siguin \(A\) i \(B\) conjunts. Si \(A\preceq B\) i \(B\preceq A\), aleshores
\(A\approx B\). És a dir, de dos monomorfismes \(f\colon A\hookrightarrow B\)
i \(g\colon B\hookrightarrow A\) se'n construeix un isomorfisme
\(A\cong B\).\par
\emph{Axiomes usats:} conjunt de parts, inclusió i pertinença internes i el
quantificador universal \(\forall\) (capítols~\ref{cap:exponencials}
i~\ref{cap:parts-interna}), complement, implicació booleana i definició per
casos (capítol~\ref{cap:valors}), imatge directa i la seva adjunció
(capítol~\ref{cap:imatges}). Cap axioma nou; en particular, \emph{cap
elecció}.
\end{teorema}

\begin{prova}
Escrivim \(g(Y)\hookrightarrow A\) per a la imatge directa
(definició~\ref{def:imatge-directa}) d'una part \(Y\hookrightarrow B\), i
igualment \(f(X)\hookrightarrow B\) per a \(X\hookrightarrow A\). Considerem
l'operació sobre parts d'\(A\)
\[
\Phi\colon P(A)\longrightarrow P(A),
\qquad
\Phi(X)\ :=\ A\setminus g\bigl(B\setminus f(X)\bigr),
\]
composició d'imatges directes i complements
(capítols~\ref{cap:imatges} i~\ref{cap:valors}); és una aplicació interna
\(P(A)\to P(A)\) ben definida. Observem que \(\Phi\) és \emph{creixent} per
la inclusió: si \(X\subseteq X'\), aleshores \(f(X)\subseteq f(X')\), perquè
la imatge directa és monòtona ---és adjunta per l'esquerra de la imatge
inversa (proposició~\ref{prop:adjuncio-imatges})---; d'on \(B\setminus
f(X')\subseteq B\setminus f(X)\), perquè el complement inverteix la inclusió
(capítol~\ref{cap:valors}); d'on \(g(B\setminus f(X'))\subseteq g(B\setminus
f(X))\) altre cop per monotonia de \(g_{*}\); i finalment \(\Phi(X)\subseteq
\Phi(X')\) prenent complements. El nombre parell d'inversions del complement
retorna el sentit de la inclusió.

\emph{Un punt fix de \(\Phi\).} Diem que una part \(X\hookrightarrow A\) és
\emph{\(\Phi\)-tancada} si \(\Phi(X)\subseteq X\). Volem la intersecció de
totes les parts \(\Phi\)-tancades. Convé construir-la amb cura, perquè és un
ínfim \emph{arbitrari} ---indexat per una part de \(P(A)\)--- i el llibre
només té ínfims binaris directament; l'obtenim, però, dins conjunts ja
construïts, mitjançant el quantificador universal. Considerem sobre el
producte \(A\times P(A)\) el predicat intern
\[
\varphi(a,X)\ :\equiv\ \bigl(\Phi(X)\subseteq X\bigr)\Rightarrow
\bigl(a\mathrel{\varepsilon} X\bigr),
\]
on l'antecedent \(\Phi(X)\subseteq X\) és la inclusió interna
(definició~\ref{def:inclusio-interna}), avaluada component \(\Phi\) amb la
projecció, el consequent \(a\mathrel{\varepsilon} X\) és la pertinença
interna (definició~\ref{def:pertinenca-interna}), i la implicació és la
booleana interna del capítol~\ref{cap:valors}. Sigui
\(R_\varphi\hookrightarrow A\times P(A)\) la part que classifica \(\varphi\).
Aplicant a \(R_\varphi\) el quantificador universal \(\forall_{\pi}\) associat
a la projecció \(\pi\colon A\times P(A)\to A\)
(definició~\ref{def:quantificadors}), s'obté una part
\[
D\ :=\ \forall_{\pi}(R_\varphi)\ \hookrightarrow\ A,
\]
els elements \(a\) de la qual són, per la lectura de \(\forall_{\pi}\), els
que compleixen \(\varphi(a,X)\) \emph{per a tota} \(X\); és a dir, els \(a\)
que pertanyen a tota part \(\Phi\)-tancada. Escriurem informalment
\(D=\bigcap\{X\in P(A):\Phi(X)\subseteq X\}\), amb el benentès que aquesta
intersecció arbitrària és, per definició, la part \(\forall_{\pi}(R_\varphi)\)
tot just construïda dins \(A\). La col·lecció de parts \(\Phi\)-tancades és
no buida: la part total \(A\) ho és, perquè \(\Phi(A)\subseteq A\) sempre.
Vegem que \(\Phi(D)=D\).

Una precisió abans de seguir. Per la lectura de \(\forall_\pi\)
(proposició~\ref{prop:quantificadors-lectura}(3)), les \(X\) que intervenen a
la definició de \(D\) no són només parts globals d'\(A\), sinó elements
generalitzats \(X\colon Y\to P(A)\), és a dir, parts \(s_X\hookrightarrow
Y\times A\) a un estadi \(Y\). Els arguments següents valen a qualsevol estadi:
\(\Phi\) és una aplicació interna, composició de \(\exists_f\), \(\neg_B\),
\(\exists_g\) i \(\neg_A\), i sobre les parts associades actua per imatges
directes al llarg de \(\id_Y\times f\) i \(\id_Y\times g\) i per complements
(proposicions~\ref{prop:quantificadors-lectura}(1)
i~\ref{prop:booleanes-puntuals}). Per tant la monotonia que acabem de
provar val igualment a cada estadi, i podem raonar amb \(X\) arbitràries.
Primer, \(\Phi(D)\subseteq D\). Sigui \(X\) qualsevol part
\(\Phi\)-tancada. Com que \(D\subseteq X\) i \(\Phi\) és creixent,
\(\Phi(D)\subseteq\Phi(X)\subseteq X\). Així \(\Phi(D)\subseteq X\) per a
tota \(X\) \(\Phi\)-tancada, d'on \(\Phi(D)\subseteq D\), que és la
intersecció de totes. Per tant \(D\) és ella mateixa \(\Phi\)-tancada.
Segon, \(D\subseteq\Phi(D)\). Com que \(\Phi(D)\subseteq D\) i \(\Phi\) és
creixent, \(\Phi(\Phi(D))\subseteq\Phi(D)\); és a dir, \(\Phi(D)\) és també
\(\Phi\)-tancada. Per la definició de \(D\) com a ínfim, \(D\subseteq
\Phi(D)\). Les dues inclusions donen \(\Phi(D)=D\).

\emph{L'isomorfisme.} Despleguem la igualtat \(D=\Phi(D)=A\setminus
g(B\setminus f(D))\). Prenent complements a \(A\),
\[
A\setminus D\ =\ g\bigl(B\setminus f(D)\bigr).
\]
Aquesta és la clau: sobre \(D\) usarem \(f\), i sobre el complement
\(A\setminus D\) usarem la inversa de \(g\). Definim
\(h\colon A\to B\) per casos (capítol~\ref{cap:valors}), segons la
part \(D\) i el seu complement (teorema~\ref{teo:part-complementari}):
\[
h(a)\ :=\
\begin{cases}
f(a) & \text{si } a\in D,\\[1mm]
g^{-1}(a) & \text{si } a\in A\setminus D.
\end{cases}
\]
La segona clàusula té sentit, i és una aplicació genuïna: la
restricció de \(g\) a \(B\setminus f(D)\) és un monomorfisme, i tot
monomorfisme és un isomorfisme sobre la seva imatge (la seva factorització
imatge té el factor \(e\) isomorfisme, proposició~\ref{prop:mono-nucli-trivial});
la imatge és \(g(B\setminus f(D))=A\setminus D\), de manera que hi ha un
isomorfisme \(A\setminus D\cong B\setminus f(D)\), que és el que escrivim
\(g^{-1}\). Així \(h=[f|_D,\,g^{-1}]\colon D+(A\setminus D)\cong A\to B\) és la
definició per casos sobre la descomposició de \(A\)
(teorema~\ref{teo:part-complementari}).

Comprovem que \(h\) és un isomorfisme veient que és alhora un monomorfisme
i un epimorfisme. Per la separació global, n'hi ha prou amb comprovar-ho
sobre \emph{punts globals}: una aplicació injectiva sobre punts és mono
(observació~\ref{obs:injectiva-mono}) i una d'exhaustiva sobre punts és
epi (observació~\ref{obs:exhaustiva-epi}); i sobre punts, pertànyer a
\(D\) o al seu complement, i a \(f(D)\) o al seu complement, són
alternatives excloents (corol·lari~\ref{cor:elements-de-2}). En el que
segueix, doncs, \(a,a',b\) són punts.

\emph{Injectivitat.} Suposem \(h(a)=h(a')\). Si tots dos són a \(D\),
\(f(a)=f(a')\) dona \(a=a'\) perquè \(f\) és un monomorfisme. Si tots dos són
a \(A\setminus D\), \(g^{-1}(a)=g^{-1}(a')\) dona \(a=g(g^{-1}(a))=
g(g^{-1}(a'))=a'\). Queda descartar el cas mixt: \(a\in D\) i \(a'\in
A\setminus D\). Aleshores \(h(a)=f(a)\in f(D)\), mentre que
\(h(a')=g^{-1}(a')\) amb \(a'=g(b')\) i \(b'\in B\setminus f(D)\); és a dir
\(h(a')=b'\notin f(D)\). Per tant \(h(a)\in f(D)\) i \(h(a')\notin f(D)\), de
manera que \(h(a)\neq h(a')\), contra la hipòtesi. El cas mixt és, doncs,
impossible, i \(h\) és un monomorfisme.

\emph{Exhaustivitat.} Sigui \(b\) un punt de \(B\). Distingim si
\(b\in f(D)\) o no. Si \(b\in f(D)\), aleshores \(b=f(a)\) amb
\(a\in D\), i \(h(a)=f(a)=b\). Si \(b\in B\setminus f(D)\), aleshores
\(g(b)\in g(B\setminus f(D))=A\setminus D\), de manera que
\(h(g(b))=g^{-1}(g(b))=b\). En tots dos casos \(b\) té preimatge per \(h\);
\(h\) és exhaustiva sobre punts i, doncs, un epimorfisme.

Un bimorfisme és un isomorfisme
(proposició~\ref{prop:bimorfisme-isomorfisme}); per tant \(h\colon A\cong
B\) i \(A\approx B\).
\end{prova}

\begin{observacio}[Per què no cal elecció]\label{obs:sb-sense-eleccio}
La construcció de \(D\) no tria res: \(D\) és l'\emph{ínfim} d'una família de
parts ---les \(\Phi\)-tancades--- completament determinada per les dades
\(f\) i \(g\), i l'ínfim, quan existeix, és únic i canònic, igual que ho era
el mínim d'una part de \(\N\) (observació~\ref{obs:minim-vs-eleccio}). Tot
viu dins un conjunt de parts \(P(A)\) ja construït, i la partició d'\(A\) en
\(D\) i el seu complement és una decisió \emph{booleana} interna, no una
successió de tries. És el mateix fenomen del capítol anterior: quan hi ha un
element distingit canònic ---aquí, l'ínfim--- l'elecció és innecessària. La
prova alternativa, que itera \(A_0=A\setminus g(B)\), \(A_{n+1}=g(f(A_n))\) i
en pren la unió, també és correcta, però demana reunir una família indexada
per \(\N\) i comprovar que està uniformement definida; la via de l'ínfim
evita aquesta reunió i és, en aquest llibre, la més neta.
\end{observacio}

\begin{corollari}[Antisimetria de la comparació]\label{cor:preceq-antisimetric}
Per a conjunts qualssevol, \(A\preceq B\) i \(B\preceq A\) impliquen
\(A\approx B\). Per tant la comparació \(\preceq\) és un ordre parcial
\emph{llevat d'equipotència}: reflexiva, transitiva
(proposició~\ref{prop:preceq-preordre}) i antisimètrica mòdul \(\approx\).
\end{corollari}

\begin{prova}
És exactament el teorema~\ref{teo:schroder-bernstein}: de les dues
dominacions se'n treu un isomorfisme. La reflexivitat i la transitivitat són
la proposició~\ref{prop:preceq-preordre}.
\end{prova}

\begin{observacio}[Lectura clàssica]\label{obs:sb-classica}
En termes clàssics, Schröder--Bernstein diu que si hi ha una injecció de
\(A\) en \(B\) i una de \(B\) en \(A\), aleshores hi ha una bijecció entre
\(A\) i \(B\): \(\abs{A}\le\abs{B}\) i \(\abs{B}\le\abs{A}\) donen
\(\abs{A}=\abs{B}\). És el que fa que els cardinals estiguin \emph{totalment
ordenats} un cop se sap, a més, que dos conjunts qualssevol són comparables
---cosa que, aquella sí, requereix l'elecció i que veurem més endavant. Aquí
només hem establert l'antisimetria, que no en necessita cap. La demostració
per punts bons que hem donat és, essencialment, la de Banach; la formulació
amb l'operador creixent \(\Phi\) i el seu menor punt fix és el nucli del
teorema de Knaster--Tarski sobre reticles complets, del qual \(P(A)\) n'és un
exemple (teorema~\ref{teo:inclusio-ordre}).
\end{observacio}

\section{L'argument diagonal i el teorema de Cantor}\label{sec:cantor}

El segon teorema compara un conjunt amb el seu conjunt de parts. La
desigualtat fàcil ---que \(A\) s'injecta en \(P(A)\)--- ja fixa que
\(A\preceq P(A)\). El contingut és que la desigualtat és \emph{estricta}: no
hi pot haver cap isomorfisme \(A\cong P(A)\), ni tan sols cap aplicació
exhaustiva \(A\to P(A)\). La raó és un únic argument, la
\emph{diagonal} de Cantor, que treballa amb la pertinença interna del
capítol~\ref{cap:parts-interna}.

\begin{teorema}[Cap aplicació \(A\to P(A)\) no és exhaustiva]\label{teo:diagonal}
Sigui \(A\) un conjunt i \(f\colon A\to P(A)\) una aplicació qualsevol.
Aleshores \(f\) no és exhaustiva. Concretament, la part
\[
D\ :=\ \{\,a\in A\ \mid\ a\not\mathrel{\varepsilon} f(a)\,\}\ \hookrightarrow\ A
\]
---els elements que \emph{no} pertanyen internament a la part que se'ls
assigna--- no és el valor \(f(a)\) de cap \(a\in A\).\par
\emph{Axiomes usats:} conjunt de parts i pertinença interna
(capítol~\ref{cap:parts-interna}), complement (capítol~\ref{cap:valors}).
Cap axioma nou.
\end{teorema}

\begin{prova}
La part \(D\) està ben definida: la pertinença interna \({\in_A}\) és una
part de \(A\times P(A)\) (definició~\ref{def:pertinenca-interna}), i \(D\) és
la imatge inversa per \(\langle\id_A,f\rangle\colon A\to A\times P(A)\) del
seu complement; és a dir, \(D=\{a : \langle a,f(a)\rangle\notin{\in_A}\}\),
tallada amb el complement (capítol~\ref{cap:valors}) i una imatge inversa
(capítol~\ref{cap:parts}). El seu nom
\(\ulcorner D\urcorner\colon\mathbf1\to P(A)\) és un element de \(P(A)\).

Suposem, cap a una contradicció, que \(D\) fos un valor de \(f\): que
existís un element \(d\colon\mathbf1\to A\) amb \(f(d)=\ulcorner
D\urcorner\), és a dir amb \(f(d)=D\) com a parts d'\(A\). Preguntem-nos si
\(d\in D\). Per la definició de \(D\),
\[
d\in D
\quad\Longleftrightarrow\quad
d\not\mathrel{\varepsilon} f(d).
\]
Però \(f(d)=D\), de manera que \(d\mathrel{\varepsilon} f(d)\) equival a
\(d\in D\) (la pertinença interna a una part constant és la pertinença
generalitzada, proposició~\ref{prop:pertinenca-interna-lectura}). Substituint,
\[
d\in D
\quad\Longleftrightarrow\quad
d\notin D,
\]
que és una contradicció: la distinció dels dos valors de veritat
(teorema~\ref{teo:zero-diferent-u}) prohibeix que un element compleixi una
condició si i només si no la compleix. Per tant no existeix cap \(d\) amb
\(f(d)=D\), i \(D\) no és a la imatge de \(f\). Com que \(D\) és un element
de \(P(A)\) que \(f\) no assoleix, \(f\) no és exhaustiva.
\end{prova}

La desigualtat en l'altre sentit és immediata: cada element determina el seu
singletó, i això injecta \(A\) dins \(P(A)\).

\begin{proposicio}[El monomorfisme dels singletons]\label{prop:singletons}
L'aplicació \(\{-\}\colon A\to P(A)\) que envia cada \(a\in A\) al singletó
\(\{a\}\hookrightarrow A\) és un monomorfisme. Per tant \(A\preceq P(A)\).
\end{proposicio}

\begin{prova}
El singletó \(\{a\}\) és la imatge del punt \(a\colon\mathbf1\to A\); en
termes interns, \(\{-\}\colon A\to P(A)\) és la transposada de la
característica de la diagonal \(\Delta\hookrightarrow A\times A\)
(capítols~\ref{cap:relacions} i~\ref{cap:parts-interna}), i és per tant una
aplicació interna ben definida. És un monomorfisme perquè és injectiu
sobre punts (observació~\ref{obs:injectiva-mono}): si \(\{a\}=\{a'\}\) com
a parts d'\(A\) per a dos punts \(a,a'\), com que sempre
\(a\mathrel{\varepsilon}\{a\}\), resulta \(a\mathrel{\varepsilon}\{a'\}\), és a
dir \(a=a'\). Així \(A\preceq P(A)\).
\end{prova}

\begin{teorema}[Cantor]\index{Cantor, teorema de}\label{teo:cantor}
Per a tot conjunt \(A\),
\[
A\ \prec\ P(A).
\]
És a dir, \(A\preceq P(A)\) però \(A\not\approx P(A)\): el conjunt de parts
és estrictament més gran que el conjunt.\par
\emph{Axiomes usats:} els del teorema~\ref{teo:diagonal} i de la
proposició~\ref{prop:singletons}; conjunt de parts, pertinença interna i
complement. Cap axioma nou; cap elecció.
\end{teorema}

\begin{prova}
La dominació \(A\preceq P(A)\) és la proposició~\ref{prop:singletons}. Falta
veure \(A\not\approx P(A)\). Si hi hagués un isomorfisme \(A\cong P(A)\),
seria en particular una aplicació exhaustiva \(A\to P(A)\), cosa impossible
pel teorema~\ref{teo:diagonal}. Per tant no hi ha cap isomorfisme, \(A\not
\approx P(A)\), i combinant amb la dominació, \(A\prec P(A)\).
\end{prova}

\begin{observacio}[No exhaustiva és més fort que no isomorf]\label{obs:cantor-fort}
El teorema~\ref{teo:diagonal} diu una cosa més forta que el teorema de
Cantor: no és només que no hi hagi cap \emph{isomorfisme} \(A\cong P(A)\),
sinó que no hi ha cap aplicació \emph{exhaustiva} \(A\to P(A)\), sigui o no
injectiva. Aquesta és la formulació que convé recordar, perquè és la que es
transporta a altres contextos ---per exemple, per veure que \(\mathbf2^{\N}\)
no és enumerable: una enumeració seria una exhaustiva
\(\N\twoheadrightarrow\mathbf2^{\N}\approx P(\N)\), que el
teorema~\ref{teo:diagonal} prohibeix. Convé mantenir separades, doncs, les
dues negacions: «no hi ha exhaustiva» i «no hi ha monomorfisme» són
afirmacions distintes, i la diagonal dona la primera.
\end{observacio}

\begin{corollari}[No hi ha conjunt universal ni mida màxima]\label{cor:no-universal}
No existeix cap conjunt \(U\) tal que tot conjunt s'injecti en \(U\). En
particular, no hi ha cap conjunt de mida màxima: per a tot \(A\) existeix un
conjunt estrictament més gran, a saber \(P(A)\).
\end{corollari}

\begin{prova}
Si un tal \(U\) existís, aleshores \(P(U)\preceq U\), perquè \(P(U)\) també
seria un conjunt i s'hauria d'injectar en \(U\). Però la
proposició~\ref{prop:singletons} dona \(U\preceq P(U)\), de manera que
Schröder--Bernstein (teorema~\ref{teo:schroder-bernstein}) forçaria
\(U\approx P(U)\), contra el teorema de Cantor
(teorema~\ref{teo:cantor}). Per tant cap conjunt no domina tots els
conjunts. La segona afirmació és el teorema de Cantor mateix:
\(A\prec P(A)\) exhibeix, per a cada \(A\), un conjunt més gran.
\end{prova}

\begin{observacio}[Lectura clàssica i el lloc de la paradoxa]\label{obs:cantor-classica}
Clàssicament, el teorema de Cantor és \(\abs{A}<2^{\abs{A}}\): el cardinal
del conjunt de parts és estrictament més gran que el del conjunt, perquè
\(P(A)\approx\mathbf2^{A}\) (el conjunt de parts és l'exponencial de base
\(\mathbf2\), capítol~\ref{cap:exponencials}). D'aquí surt la torre infinita
de cardinals \(\abs{A}<2^{\abs{A}}<2^{2^{\abs{A}}}<\cdots\), que no s'atura
mai. El corol·lari~\ref{cor:no-universal} és la forma estructural del fet que
«el conjunt de tots els conjunts» no existeix: si existís, contindria el seu
propi conjunt de parts i violaria Cantor. Aquesta és, de fet, la manera
correcta de dir en aquest llibre el que en la teoria clàssica es formula amb
la paradoxa de Russell; aquí no hi ha cap paradoxa perquè mai no hem tractat
la pertinença global entre conjunts com una relació, sinó sempre la
pertinença \emph{a una part d'un conjunt fix}. La diagonal de Cantor
reapareixerà, en versió ordinal, com el teorema de Burali--Forti, i en versió
de famílies, com el teorema de König.
\end{observacio}

\begin{resumcap}
\apartat{Idees centrals}
\begin{enumerate}
  \item Comparar mides és exhibir aplicacions: \(A\preceq B\) és un monomorfisme, \(A\approx B\) un isomorfisme. Són relacions del metallenguatge entre conjunts, no parts d'un producte.
  \item Schröder--Bernstein: dues dominacions oposades donen un isomorfisme. La prova és un menor punt fix d'un operador creixent \(\Phi\) sobre \(P(A)\), obtingut com a ínfim amb el quantificador universal; cap elecció.
  \item L'isomorfisme \(h=[f|_D,g^{-1}]\) es defineix per casos sobre \(D\) i el seu complement, i es comprova que és bimorfisme sobre punts globals, on les alternatives són excloents.
  \item Argument diagonal: cap \(A\to P(A)\) és exhaustiva, perquè la part \(\{a:a\not\mathrel{\varepsilon}f(a)\}\) no és cap valor. Amb els singletons, \(A\prec P(A)\): Cantor.
  \item No hi ha conjunt universal ni mida màxima; la paradoxa de Russell no apareix perquè el llibre no té pertinença global, només pertinença a parts d'un conjunt fix.
\end{enumerate}
\apartat{Errors típics} Creure que dues injeccions oposades donen directament una bijecció (cal construir-la); intentar l'ínfim de les parts \(\Phi\)-tancades sense el quantificador universal (els axiomes només donen ínfims binaris); fer la distinció de casos de la prova de Schröder--Bernstein sobre elements generalitzats, on \enquote{o} no exclou; confondre \enquote{no isomorf} amb \enquote{no exhaustiu} (la diagonal dona la segona, més forta); i llegir el teorema de Cantor com una afirmació sobre un conjunt de tots els cardinals.
\apartat{Lectura recomanada} \cite[§§8.1--8.3]{lawvere-rosebrugh} per a l'argument diagonal en conjunts abstractes i el teorema del punt fix de Lawvere, del qual Cantor és un cas; \cite[sessió~29]{lawvere-schanuel} per a una presentació elemental de la diagonal; \cite[§§22 i 25]{halmos} per a Schröder--Bernstein i el teorema de Cantor en la teoria clàssica; i \cite[cap.~3]{jech2003} per a la demostració clàssica de Schröder--Bernstein i l'aritmètica cardinal que en resulta.
\end{resumcap}

\chapter{Famílies, fibres i operacions indexades}\label{cap:families}

\begin{objectius}
\apartat{Prerequisits} Del capítol~\ref{cap:parts}: pullbacks, imatge inversa i fibres. Del capítol~\ref{cap:relacions}: seccions, l'axioma d'elecció i les particions. Del capítol~\ref{cap:parts-interna}: conjunt de parts, pertinença interna i els quantificadors \(\exists_f\), \(\forall_f\) amb les seves lectures i adjuncions. Dels capítols~\ref{cap:infinitud} i~\ref{cap:cantor}: les famílies numerables i la construcció d'ínfims arbitraris amb \(\forall\).
\apartat{Objectius} En acabar el capítol, el lector hauria de poder:
\begin{itemize}
  \item representar una família de conjunts indexada per \(I\) com una aplicació \(p\colon E\to I\), i llegir-ne les fibres com a parts del total;
  \item identificar seccions amb funcions d'elecció, i situar l'elecció com l'existència de seccions per a famílies de fibres no buides;
  \item construir la suma (el total) i el producte (les seccions, dins \(E^I\)) d'una família, i distingir l'existència del producte com a objecte de l'existència dels seus punts;
  \item codificar una família de parts d'un conjunt fix com \(\alpha\colon I\to P(A)\), construir-ne la unió i la intersecció indexades amb els quantificadors, i demostrar-ne les propietats de suprem, ínfim i De Morgan;
  \item demostrar que una família de parts disjuntes que cobreixen \(A\) és una partició, tancant el pont del capítol~\ref{cap:relacions};
  \item distingir família interna de regla externa, aplicar la regla de treball, i saber què promet l'esquema SR i per què no és l'elecció.
\end{itemize}
\end{objectius}

Fins ara hem indexat col·leccions de conjunts d'una manera implícita: una
família \((A_n)_{n\in\N}\) apareixia al capítol~\ref{cap:infinitud} com una
aplicació \(p\colon E\to\N\) amb les fibres per components
(definició~\ref{def:familia-numerable}), i els arguments diagonals del
capítol~\ref{cap:cantor} només comparaven conjunts d'un en un. Els capítols
transfinits que vindran ---recursió sobre un bon ordre, suprems de tipus
ordinals, successions d'alephs--- necessiten manejar \emph{famílies} de
conjunts amb tota generalitat, i sobretot necessiten un criteri net per
distingir dues coses que en la pràctica clàssica es confonen: una família
que \emph{ja existeix} com a objecte, i una \emph{regla} que a cada índex
assigna un conjunt però que encara no s'ha reunit en cap objecte. Aquesta
distinció és el cor del capítol i la clau de tot el que segueix; ja es va
insinuar amb l'assignació \(n\mapsto\seg{n}\) del capítol~\ref{cap:naturals},
i reapareixerà a cada recursió transfinita.

La idea central és simple i no demana cap vocabulari nou. Una família de
conjunts indexada per un conjunt \(I\) es representa amb una aplicació
\(p\colon E\to I\); el conjunt \(E\) és el \emph{total}, i el conjunt indexat
per \(i\) és la \emph{fibra} de \(p\) sobre \(i\). Tot el capítol desplega
aquesta idea ---suma, producte, unions i interseccions indexades--- amb les
eines que ja tenim: pullbacks, exponencials, el conjunt de parts i els
quantificadors interns. No parlarem de functors ni de fibracions; no en
calen.

\section{Famílies com a aplicacions, i les seves fibres}\label{sec:families-aplicacions}

\begin{definicio}[Família de conjunts]\index{Família de conjunts}\label{def:familia}
Una \emph{família de conjunts indexada per \(I\)} és una aplicació
\(p\colon E\to I\). El conjunt \(I\) és l'\emph{índex}, \(E\) és el conjunt
\emph{total}, i per a cada element \(i\colon\mathbf1\to I\) la \emph{fibra}
de la família sobre \(i\) és el producte fibrat
\[
E_i\ :=\ \mathbf1\times_I E
\]
de \(i\colon\mathbf1\to I\) i \(p\colon E\to I\)
(definició~\ref{def:pullback}), amb la projecció
\(E_i\hookrightarrow E\) que n'és un monomorfisme
(proposició~\ref{prop:pullback-mono}). Escrivim també \(A_i:=E_i\) quan
volem pensar la família com \((A_i)_{i\in I}\).
\end{definicio}

\begin{observacio}[La fibra és una part del total]\label{obs:fibra-part}
La fibra \(E_i=\mathbf1\times_I E\) és exactament la imatge inversa
\(p^{-1}(i)\hookrightarrow E\) del punt \(i\) per \(p\)
(definició~\ref{def:imatge-inversa}, i
observació~\ref{obs:pullback-exemples}): en la lectura clàssica, els seus
punts són els \(e\in E\) amb \(p(e)=i\). Per tant una família \(p\colon E\to
I\) no és res més que el seu total \(E\) \emph{repartit} en parts disjuntes,
una per cada índex, i cada part etiquetada pel seu índex a través de \(p\).
Aquesta és la manera estructural de dir \enquote{una família de conjunts}
sense sortir dels conjunts ja construïts: tot viu dins \(E\), i les fibres
són parts seves.
\end{observacio}

L'exemple més senzill és el d'una família que no varia amb l'índex.

\begin{definicio}[Família constant]\index{Família constant}\label{def:familia-constant}
Donat un conjunt \(A\), la \emph{família constant} de valor \(A\) indexada
per \(I\) és la projecció
\[
\pi_I\colon I\times A\longrightarrow I .
\]
La seva fibra sobre cada \(i\) és \(\{i\}\times A\cong A\)
(observació~\ref{obs:pullback-exemples}), de manera que representa la família
\((A)_{i\in I}\) amb totes les fibres iguals a \(A\).
\end{definicio}

\begin{observacio}[Tota aplicació és una família, i tota part una família de valors booleans]\label{obs:familia-exemples}
La noció és prou àmplia per recobrir casos ja coneguts. Una aplicació
qualsevol \(p\colon E\to I\) es \emph{pot} llegir com a família, i triar com
llegir-la depèn de què vulguem: la mateixa \(p\colon E\to\N\) del
capítol~\ref{cap:infinitud} era la família numerable
(definició~\ref{def:familia-numerable}), cas particular amb \(I=\N\). Una
part \(m\colon S\hookrightarrow I\), vista com la seva característica
\(\chi_S\colon I\to\mathbf2\), és la família indexada per \(I\) amb fibra
\(\mathbf1\) sobre els \(i\in S\) i \(\mathbf0\) fora: una família de conjunts
\enquote{de com a molt un element}. Aquests exemples fixen la intuïció; el
contingut del capítol és el que se'n construeix.
\end{observacio}

\section{Seccions i funcions d'elecció}\label{sec:seccions}

Una família \(p\colon E\to I\) reparteix \(E\) en fibres; \emph{triar} un
element a cada fibra és exactament donar una secció de \(p\).

\begin{definicio}[Secció d'una família]\index{Secció d'una família}\label{def:seccio-familia}
Una \emph{secció} de la família \(p\colon E\to I\) és una aplicació
\(s\colon I\to E\) amb \(p\circ s=\id_I\). Per a cada índex \(i\), el valor
\(s(i)\) és un element de la fibra \(E_i\), perquè \(p(s(i))=i\); una secció
és, doncs, una \emph{tria simultània} d'un element a cada fibra ---una
\emph{funció d'elecció} per a la família.
\end{definicio}

\begin{observacio}[Quan existeix una secció: el lloc exacte de l'elecció]\label{obs:seccio-eleccio}
L'existència d'una secció \emph{no} és automàtica, i és precisament aquí on
apareix ---o no--- l'axioma d'elecció. (Diem indistintament que una fibra és
\emph{no inicial} o \emph{no buida}: per la separació global i el lema del
buit, un conjunt és inicial si i només si no té cap punt global, de manera
que les dues expressions coincideixen en aquest llibre.) Si totes les fibres són no inicials,
és a dir si \(p\) és un epimorfisme (exhaustivitat sobre punts,
capítol~\ref{cap:valors}; i tot epimorfisme és regular,
capítol~\ref{cap:imatges}), l'axioma d'elecció
(axioma~\ref{ax:eleccio}) diu exactament que \(p\) admet una secció
(proposició~\ref{prop:epi-regular-eleccio}). Sense elecció, en canvi, una
família de fibres totes no buides pot no tenir cap secció. Convé retenir la
distinció, que serà constant en els capítols transfinits: \emph{tenir
fibres no buides} és una propietat de la família ja donada; \emph{tenir una
secció} és una tria simultània que, en general, costa l'elecció. Quan les
fibres porten un element distingit canònic ---un mínim, un ínfim--- la secció
es construeix sense elecció, com vam veure amb els subobjectes de \(\N\)
(observació~\ref{obs:minim-vs-eleccio}) i amb la part \(D\) de
Schröder--Bernstein (observació~\ref{obs:sb-sense-eleccio}).
\end{observacio}

\section{Suma i producte d'una família}\label{sec:suma-producte}

Les dues operacions bàsiques sobre una família són reunir totes les fibres
en un sol conjunt ---la suma--- i formar el conjunt de totes les seccions
---el producte. Totes dues es construeixen amb el que ja tenim.

\begin{definicio}[Suma d'una família]\index{Suma d'una família}\label{def:suma-familia}
La \emph{suma} de la família \(p\colon E\to I\) és el seu propi total \(E\),
pensat com la reunió disjunta de les fibres,
\[
\sum_{i\in I}E_i\ :=\ E,
\]
en què cada fibra \(E_i\hookrightarrow E\) apareix marcada pel seu índex a
través de \(p\). L'aplicació \(p\) mateixa recorda, per a cada element del
total, a quina fibra pertany.
\end{definicio}

\begin{observacio}[La suma no demana res de nou]\label{obs:suma-existeix}
La suma d'una família \emph{existeix sempre}, sense cap axioma ni cap tria:
és el conjunt \(E\), que és una de les dades. Aquesta trivialitat aparent és
important. Recollir \enquote{totes les fibres alhora} és exactament tenir el
total \(E\); el que pot mancar, com veurem, no és la reunió sinó la
\emph{regla} que produeix \(E\) a partir d'una assignació índex-a-conjunt
quan aquesta assignació no ve ja d'una aplicació. Per a la família constant
\(I\times A\to I\), la suma és \(I\times A\), i recupera el producte
cartesià habitual llegit com \enquote{\(I\) còpies de \(A\)}.
\end{observacio}

\begin{definicio}[Producte d'una família]\index{Producte d'una família}\label{def:producte-familia}
El \emph{producte} de la família \(p\colon E\to I\) és el conjunt de les
seves seccions,
\[
\prod_{i\in I}E_i\ :=\ \{\,s\in E^{I}\ \mid\ p\circ s=\id_I\,\}
\ \hookrightarrow\ E^{I},
\]
la part de l'exponencial \(E^{I}\) (axioma~\ref{ax:exponencials}) tallada per
l'equació \(p\circ s=\id_I\) ---l'equalitzador de les dues aplicacions
paral·leles \(E^{I}\to I^{I}\) donades per \(s\mapsto p\circ s\)
(postcomposició interna amb \(p\)) i per la constant \(s\mapsto\id_I\). Els seus elements
globals són exactament les seccions de \(p\) (definició~\ref{def:seccio-familia}).
\end{definicio}

\begin{proposicio}[Existència del producte i el lloc de l'elecció]\label{prop:producte-existeix}
Per a tota família \(p\colon E\to I\), el producte \(\prod_{i\in I}E_i\)
\emph{existeix com a conjunt}: és una part interna ben definida de
\(E^{I}\). En canvi, que sigui \emph{no inicial} ---que tingui algun element
global--- equival a l'existència d'una secció de \(p\), i per tant, quan les
fibres són no buides, a una forma d'elecció.\par
\emph{Axiomes usats:} exponencials (axioma~\ref{ax:exponencials}),
equalitzadors i parts (capítols~\ref{cap:parts} i~\ref{cap:valors}); cap
axioma nou per a l'existència \emph{com a objecte}. La no-vacuïtat, per a
fibres no buides, usa l'elecció (axioma~\ref{ax:eleccio}).
\end{proposicio}

\begin{prova}
L'exponencial \(E^{I}\) existeix per l'axioma~\ref{ax:exponencials}, i les
dues aplicacions paral·leles \(E^{I}\to I^{I}\) donades per
\(s\mapsto p\circ s\) (postcomposició interna amb \(p\): la composició
interna del capítol~\ref{cap:exponencials} amb el primer argument fixat al
nom de \(p\)) i per \(s\mapsto\ulcorner\id_I\urcorner\) (constant) tenen un
equalitzador, que és una part de \(E^{I}\) (capítol~\ref{cap:parts}); aquesta part és, per construcció,
\(\prod_{i\in I}E_i\). Fins aquí, cap tria. Un element global del producte és
una \(s\colon\mathbf1\to E^{I}\) amb \(p\circ s=\id_I\), és a dir ---per
transposició (definició~\ref{def:exponencial})--- una aplicació
\(I\to E\) amb \(p\circ s=\id_I\): una secció. Per tant el producte és no
inicial si i només si \(p\) té una secció. Si les fibres són totes no buides
(\(p\) epimorfisme), l'existència d'una tal secció és exactament el
que afirma l'axioma d'elecció (axioma~\ref{ax:eleccio});
sense elecció pot fallar.
\end{prova}

\begin{observacio}[Objecte contra elements: una distinció que cal mantenir]\label{obs:objecte-vs-elements}
La proposició anterior separa dues qüestions que la notació
\(\prod_{i\in I}E_i\) tendeix a confondre: \emph{existeix el producte?} i
\emph{té elements?}. La primera es respon sempre que sí, amb exponencials i
equalitzadors. La segona, per a fibres no buides, és el contingut mateix de
l'elecció. Aquesta és la versió general del que vam veure amb el producte
numerable (observació~\ref{obs:producte-numerable}): l'objecte hi és sempre;
els seus punts, no necessàriament. La mateixa distinció ---objecte contra
punts globals--- reapareixerà amb les famílies definides per regles, a la
secció~\ref{sec:regla-vs-familia}.
\end{observacio}

\section{Famílies de parts d'un conjunt fix}\label{sec:families-parts}

Un cas particularment manejable, i el que més usarem, és el de les famílies
totes les fibres de les quals són parts d'un mateix conjunt fix \(A\).
Aleshores la família es codifica com una aplicació cap al conjunt de parts, i
totes les operacions es fan \emph{acotadament} dins \(A\) i \(P(A)\).

\begin{definicio}[Família de parts d'un conjunt fix]\index{Família de parts d'un conjunt fix}\label{def:familia-parts}
Una \emph{família de parts} de \(A\) indexada per \(I\) és una aplicació
\[
\alpha\colon I\longrightarrow P(A),
\]
que a cada índex \(i\) assigna una part \(\alpha(i)\hookrightarrow A\)
(el seu nom, definició~\ref{def:conjunt-parts}). La \emph{relació de
pertinença total} de la família és la part
\[
R\ :=\ \{\,(i,a)\in I\times A\ \mid\ a\mathrel{\varepsilon}\alpha(i)\,\}
\ \hookrightarrow\ I\times A,
\]
la imatge inversa de la pertinença interna \({\in_A}\)
(definició~\ref{def:pertinenca-interna}) per \(\id_I\times\alpha\), seguida
de la commutació. Escrivim \(R_i\hookrightarrow A\) per a la fibra de la
projecció \(R\to I\) sobre \(i\), que és exactament \(\alpha(i)\).
\end{definicio}

\begin{observacio}[Tot és acotat]\label{obs:familia-parts-acotada}
Aquesta és la situació ideal per als capítols transfinits, i la raó per la
qual la guia insisteix a \emph{acotar} les col·leccions. Una família de parts
de \(A\) és un únic objecte, \(\alpha\colon I\to P(A)\), que viu dins un
exponencial fix \(P(A)^{I}\); la seva relació de pertinença total \(R\) viu
dins \(I\times A\); i cada fibra és una part de \(A\). Res no s'escapa a un
ambient prèviament construït. Sempre que puguem expressar una col·lecció de
conjunts com a família de parts d'un fix, les seves unions, interseccions i
suprems es formaran sense cap principi nou, com veurem tot seguit. La
dificultat transfinita començarà exactament quan una col·lecció \emph{no} es
pugui allotjar així.
\end{observacio}

\begin{definicio}[Unió i intersecció indexades]\index{Unió i intersecció indexades}\label{def:unio-interseccio-indexada}
Sigui \(\alpha\colon I\to P(A)\) una família de parts de \(A\), amb relació
de pertinença total \(R\hookrightarrow I\times A\) i projecció
\(\pi_A\colon I\times A\to A\). La \emph{unió indexada} i la
\emph{intersecció indexada} de la família són les parts de \(A\)
\[
\bigcup_{i\in I}\alpha(i)\ :=\ \exists_{\pi_A}(R),
\qquad
\bigcap_{i\in I}\alpha(i)\ :=\ \forall_{\pi_A}(R),
\]
on \(\exists_{\pi_A}\) i \(\forall_{\pi_A}\) són els quantificadors interns
associats a la projecció (definició~\ref{def:quantificadors}), aplicats al
nom de la part \(R\). Sobre punts globals (proposició~\ref{prop:quantificadors-lectura}):
un punt \(a\) pertany a la unió quan existeix un punt \(i\) amb
\(a\mathrel{\varepsilon}\alpha(i)\), i a la intersecció quan
\(a\mathrel{\varepsilon}\alpha(i)\) per a tot \(i\); sobre elements
generalitzats, la lectura de la unió és amb recobriment, com sempre amb
\(\exists\).
\end{definicio}

\begin{observacio}[La lectura amb quantificadors]\label{obs:unio-interseccio-lectura}
La formulació anterior és la traducció directa de \enquote{existeix un
índex\ldots} i \enquote{per a tot índex\ldots} al llenguatge intern del
capítol~\ref{cap:parts-interna}. Per a la unió: \(a\in\bigcup_i\alpha(i)\)
diu que la fibra de \(R\) sobre \(a\) ---els índexs \(i\) amb
\(a\mathrel{\varepsilon}\alpha(i)\)--- és no buida, i això és exactament
\(\exists_{\pi_A}(R)\). Per a la intersecció, \(a\in\bigcap_i\alpha(i)\) diu
que \emph{tot} índex compleix \(a\mathrel{\varepsilon}\alpha(i)\), és a dir
\(\forall_{\pi_A}\) del predicat \enquote{\((i,a)\in R\)}: tots els punts de
la fibra de \(\pi_A\) sobre \(a\), que són els \((i,a)\), han de ser a \(R\).
La construcció usa \emph{només} els
quantificadors interns, ja disponibles; cap suprem infinit adjuntat, cap
axioma nou. És la mateixa tècnica amb què vam formar la part \(D\) de
Schröder--Bernstein com un ínfim arbitrari
(teorema~\ref{teo:schroder-bernstein}).
\end{observacio}

\begin{proposicio}[Propietats de la unió i la intersecció indexades]\label{prop:unio-interseccio-propietats}
Sigui \(\alpha\colon I\to P(A)\) una família de parts de \(A\). Aleshores,
dins \(P(A)\):
\begin{enumerate}
\item \(\bigcup_{i\in I}\alpha(i)\) és la menor part de \(A\) que conté totes
les \(\alpha(i)\); és a dir, el suprem de la família a \((P(A),\subseteq)\).
\item \(\bigcap_{i\in I}\alpha(i)\) és la major part de \(A\) continguda en
totes les \(\alpha(i)\); és a dir, l'ínfim de la família.
\item Es compleixen les lleis de De Morgan indexades:
\[
A\setminus\bigcup_{i\in I}\alpha(i)=\bigcap_{i\in I}\bigl(A\setminus\alpha(i)\bigr),
\qquad
A\setminus\bigcap_{i\in I}\alpha(i)=\bigcup_{i\in I}\bigl(A\setminus\alpha(i)\bigr).
\]
\end{enumerate}
\emph{Axiomes usats:} conjunt de parts, pertinença interna i quantificadors
(capítol~\ref{cap:parts-interna}), complement (capítol~\ref{cap:valors}); cap
axioma nou.
\end{proposicio}

\begin{prova}
(1) i (2) són les propietats universals de \(\exists_{\pi_A}\) i
\(\forall_{\pi_A}\) com a adjunts de la imatge inversa
(proposició~\ref{prop:quantificadors-adjuncio}, amb l'estadi \(X=\mathbf1\)):
\(\exists_{\pi_A}(R)\subseteq S\) equival a \(R\subseteq\pi_A^{-1}(S)=I\times S\),
és a dir, sobre punts, a \(\alpha(i)\subseteq S\) per a tot \(i\) ---la
propietat del suprem---; i dualment, \(S\subseteq\forall_{\pi_A}(R)\) equival
a \(I\times S\subseteq R\), és a dir, a \(S\subseteq\alpha(i)\) per a tot
\(i\) ---la de l'ínfim. (3) surt de les mateixes adjuncions per
complementació, com a la construcció de \(\forall_f=\neg\exists_f\neg\) de la
proposició~\ref{prop:quantificadors-adjuncio}, transportant les lleis de De
Morgan booleanes (observació~\ref{obs:lleis-transportades}) a la família.
L'argument no fa servir cap hipòtesi sobre \(I\): si \(I\) és inicial, la
unió és la part buida i la intersecció és tot \(A\), que són, en efecte, el
suprem i l'ínfim de la família buida.
\end{prova}

Amb aquest llenguatge podem tancar el pont que el capítol~\ref{cap:relacions}
va deixar obert: una partició estructural (un epimorfisme) té fibres
disjuntes que cobreixen el conjunt (proposició~\ref{prop:fibres-particio});
recíprocament, una família de parts amb aquestes propietats és una
partició.

\begin{proposicio}[Les famílies de parts disjuntes que cobreixen són particions]\label{prop:familia-particio}
Sigui \(\alpha\colon I\to P(A)\) una família de parts de \(A\) tal que,
sobre punts: cada \(\alpha(i)\) és no buida; dos índexs diferents donen
parts disjuntes, \(i\ne j\Rightarrow\alpha(i)\cap\alpha(j)=\varnothing_A\);
i la unió indexada és tot \(A\). Aleshores existeix un epimorfisme
\(q\colon A\twoheadrightarrow I\) ---una partició--- les fibres del qual són
exactament les \(\alpha(i)\), i és únic amb aquesta propietat.\par
\emph{Axiomes usats:} els del capítol~\ref{cap:parts-interna} i la separació
global; cap elecció.
\end{proposicio}

\begin{prova}
Sigui \(R\hookrightarrow I\times A\) la relació de pertinença total de la
família, amb projeccions \(r_I\colon R\to I\) i \(r_A\colon R\to A\). La
hipòtesi que les parts cobreixen \(A\) diu que \(\exists_{\pi_A}(R)\) és
total, és a dir, que \(r_A\) és exhaustiva sobre punts: tot punt \(a\)
pertany a alguna \(\alpha(i)\). La disjunció diu que \(r_A\) és injectiva
sobre punts: si \((i,a)\) i \((j,a)\) són punts de \(R\), aleshores
\(a\in\alpha(i)\cap\alpha(j)\), que és buida si \(i\ne j\); per tant \(i=j\).
Per la separació global, \(r_A\) és un bimorfisme i, doncs, un isomorfisme
(proposició~\ref{prop:bimorfisme-isomorfisme}). Posem
\(q:=r_I\circ r_A^{-1}\colon A\to I\). És exhaustiva sobre punts, perquè
cada \(\alpha(i)\) té un punt \(a\), i aleshores \((i,a)\in R\) i
\(q(a)=i\); per tant és un epimorfisme, i una partició. La fibra
\(q^{-1}(i)\) és, per construcció, la part dels \(a\) amb \((i,a)\in R\), és
a dir, \(\alpha(i)\). La unicitat: una altra \(q'\) amb les mateixes fibres
coincideix amb \(q\) sobre tots els punts, i per la separació global
\(q'=q\).
\end{prova}

\section{Regla externa contra família interna}\label{sec:regla-vs-familia}

Arribem al punt central del capítol, la distinció que la guia demana fer molt
explícita perquè governa tota la teoria transfinita. Cal separar dos usos de
la notació \enquote{\(i\mapsto X_i\)}.

\begin{definicio}[Família interna i regla externa]\index{Família interna i regla externa}\label{def:familia-vs-regla}
Una \emph{família interna} indexada per \(I\) és una aplicació \emph{ja
donada} \(p\colon E\to I\) (definició~\ref{def:familia}); les seves fibres
\(E_i\) són parts d'un total \(E\) que existeix com a objecte. En canvi, una
\emph{regla externa} és una prescripció que, al metallenguatge, associa a
cada element \(i\) d'\(I\) un conjunt \(X_i\), \emph{sense} que se n'hagi
exhibit cap aplicació \(E\to I\) les fibres de la qual siguin els \(X_i\).
Una regla externa \emph{es representa} per una família interna quan existeix
una tal \(p\colon E\to I\) amb \(E_i\cong X_i\) per a tot \(i\), de manera
compatible.
\end{definicio}

\begin{observacio}[El problema, i on ja el vam trobar]\label{obs:regla-externa-problema}
La diferència no és pedant: hi ha regles externes perfectament definides que
\emph{no} es poden representar per cap família interna amb els axiomes que
tenim. L'exemple prototípic ja va aparèixer al capítol~\ref{cap:naturals}: la
regla \(n\mapsto\seg{n}\), que a cada natural assigna el seu segment inicial,
és una prescripció clara, però reunir tots els \(\seg{n}\) en un sol total
\(E\to\N\) demana una construcció que no és cap dels axiomes bàsics ---de fet
es fa allotjant tots els \(\seg{n}\) dins un exponencial fix, com al
conjunt de llistes (definició~\ref{def:llistes}). Sempre que puguem allotjar
els \(X_i\) dins un conjunt ambient comú ---per exemple si tots són parts
d'un fix \(A\), i aleshores la regla és una aplicació \(I\to P(A)\)
(definició~\ref{def:familia-parts})--- la representació existeix i no cal res
nou. El problema genuí apareix quan els \(X_i\) \emph{creixen sense cota
comuna coneguda}: aleshores no hi ha cap ambient dins el qual tallar-los per
separació, i la reunió pot escapar-se dels axiomes.
\end{observacio}

\begin{observacio}[La regla de treball]\label{obs:regla-treball}
D'aquí surt una disciplina que mantindrem a tot el desenvolupament
transfinit: \textbf{cada vegada que aparegui una notació \(i\mapsto X_i\),
cal indicar quina aplicació \(E\to I\) representa la família; i si no
existeix encara, cal declarar-ho.} Escriure \(\bigcup_{i\in I}X_i\) o
\(\sup_i X_i\) només té sentit un cop s'ha exhibit el total \(E\), o un
ambient \(A\) que contingui tots els \(X_i\). Aquesta és la lliçó que fa
manejable la recursió transfinita: mentre els valors visquin dins un
codomini fix, la família és interna i tot és de Nivell~I; quan els valors
siguin conjunts nous sense codomini comú, cal un principi addicional, que
introduïm a la secció següent i no abans.
\end{observacio}

\section{Nota sobre la reunió de famílies no acotades (SR)}\label{sec:sr}

\begin{nota}
Aquesta secció és de \emph{Nivell~II}: enuncia, però \emph{no adopta},
un principi més fort que els axiomes del llibre. Els capítols de Nivell~I no
en depenen, i el text ho assenyalarà cada vegada que calgui.
\end{nota}

Fins ara, tota família que hem format era interna perquè els seus valors
vivien dins un conjunt ja construït: parts d'un fix \(A\), elements d'un
exponencial, fibres d'un total donat. Els capítols següents arribaran, però,
a construccions on aquesta comoditat desapareix: una recursió transfinita els
valors de la qual són conjunts \emph{nous} sense codomini comú, el pas d'un
bon ordre al seu representant ordinal, una successió \(\alpha\mapsto\aleph_\alpha\).
En tots aquests casos tindrem una \emph{regla externa} ben definida
\(i\mapsto X_i\) i voldrem reunir-la en una \emph{família interna} \(E\to I\),
però cap axioma de Nivell~I no ens ho garantirà. El principi que ho permet,
que designarem provisionalment per \textbf{SR} (\enquote{substitució
estructural}), es pot enunciar orientativament així.

\begin{observacio}[Esquema de reunió de famílies (SR), provisional]\label{obs:sr}
Suposem que una fórmula, per a cada element \(i\) d'un conjunt \(I\),
determina un conjunt \(X_i\) \emph{llevat d'un únic isomorfisme compatible
amb les dades}. Aleshores ---diu SR--- existeix una aplicació \(p\colon E\to
I\) la fibra de la qual sobre cada \(i\) representa \(X_i\). És a dir: tota
regla externa uniformement i invariantment definida es pot representar per
una família interna.

No prenem això com a axioma en aquesta forma. L'enunciat definitiu, com un
\emph{esquema} de fórmules amb la condició de funcionalitat a isomorfisme
únic i l'estructura que la garanteix explicitades, és el
principi~\ref{sr:esquema} del capítol~\ref{cap:alephs}, que també n'identifica
la formulació adoptada a la literatura; aquí només el deixem anunciat
perquè el lector sàpiga on i per què farà falta.
\end{observacio}

\begin{observacio}[SR no és l'elecció]\label{obs:sr-no-eleccio}
Convé, ja d'entrada, no confondre SR amb l'axioma d'elecció, perquè fan
feines lògiques diferents ---la mateixa distinció que ja vam marcar al
capítol~\ref{cap:infinitud}
(observació~\ref{obs:eleccio-global-vs-numerable}). L'elecció tria un element
dins cada fibra d'una família que \emph{ja existeix} \(p\colon E\to I\)
(una secció, secció~\ref{sec:seccions}); SR, en canvi, garanteix
l'\emph{existència} de la família \(E\to I\) a partir d'una regla que defineix
els seus valors. La primera opera sobre un total donat; la segona el
construeix. En els capítols que vénen, mantindrem separats els dos usos i,
seguint la disciplina de la caixa d'axiomes, marcarem \enquote{elecció} i
\enquote{SR} allà on cadascun intervingui, mai indistintament.
\end{observacio}

\begin{resumcap}
\apartat{Idees centrals}
\begin{enumerate}
  \item Una família de conjunts indexada per \(I\) és una aplicació \(p\colon E\to I\); les fibres \(E_i=p^{-1}(i)\) són parts del total. La família constant és una projecció.
  \item Una secció de \(p\) és una funció d'elecció; per a fibres no buides, l'existència de seccions és l'axioma d'elecció. Quan les fibres tenen un element canònic (mínim, ínfim), la secció es construeix sense elecció.
  \item La suma d'una família és el total \(E\); existeix sempre. El producte és la part de \(E^I\) de les seccions: existeix sempre com a objecte, i té punts exactament quan \(p\) té secció.
  \item Una família de parts d'un conjunt fix és \(\alpha\colon I\to P(A)\); la unió i la intersecció indexades són \(\exists_{\pi_A}(R)\) i \(\forall_{\pi_A}(R)\), suprem i ínfim a \((P(A),\subseteq)\), amb les lleis de De Morgan; tot acotat dins \(A\).
  \item Una família de parts no buides, disjuntes i que cobreixen \(A\) és una partició (un epimorfisme \(A\twoheadrightarrow I\)), i recíprocament.
  \item Família interna (una \(p\colon E\to I\) donada) contra regla externa (una prescripció \(i\mapsto X_i\)): la regla de treball exigeix dir sempre quina aplicació representa la família. Quan els valors no tenen ambient comú, cal SR, que no és l'elecció: l'elecció tria dins fibres d'una família que existeix, SR fa existir la família.
\end{enumerate}
\apartat{Errors típics} Pensar que \enquote{una família de conjunts} és una regla i no una aplicació; confondre l'existència del producte amb la seva no-vacuïtat; llegir la unió indexada sobre elements generalitzats amb un \enquote{existeix} en el mateix estadi; escriure \(\bigcup_iX_i\) sense haver exhibit un total o un ambient; i confondre SR (existència de la família) amb l'elecció (secció d'una família existent).
\apartat{Lectura recomanada} \cite[cap.~8 i §§9.1--9.2]{lawvere-rosebrugh} per a les famílies com a aplicacions, les fibres i el pas als conjunts variables; \cite[§§9--10]{halmos} per a les famílies, les unions i els productes indexats en la teoria clàssica, on la distinció entre regla i família queda amagada per l'axioma de reemplaçament; \cite[§§2--3]{shulman-comparing} per a la formulació estructural de les famílies i del reemplaçament; i \cite[cap.~I]{maclane-moerdijk} per a les categories de fletxes \(\mathbf{Set}/I\) com a marc general de les famílies indexades.
\end{resumcap}

\chapter{Ordres, bons ordres i segments}\label{cap:bons-ordres}

\begin{objectius}
\apartat{Prerequisits} Del capítol~\ref{cap:relacions}: relacions com a parts d'un producte, diagonal, relació inversa. Del capítol~\ref{cap:naturals}: l'ordre de \(\N\), el seu bon ordre, i la convenció que les propietats d'un ordre són afirmacions sobre parts que es comproven sobre punts globals (observació~\ref{obs:ordre-sobre-punts}). Del capítol~\ref{cap:valors}: complements i el lema del buit. Del capítol~\ref{cap:cantor}: la comparació de conjunts, per contrast amb la comparació d'ordres.
\apartat{Objectius} En acabar el capítol, el lector hauria de poder:
\begin{itemize}
  \item definir ordre estricte, ordre total i bon ordre com a parts de \(W\times W\) amb condicions, i saber quines es llegeixen sobre punts;
  \item demostrar que tot bon ordre és total i que el mínim és únic;
  \item definir els segments inicials \(W_{<w}\) i les parts inicials, i demostrar que en un bon ordre les parts inicials pròpies són exactament els segments, que són ben ordenats i determinen el seu punt;
  \item demostrar la rigidesa (cap punt baixa sota una creixent injectiva), la unicitat de l'isomorfisme d'ordre i que cap bon ordre és isomorf a un segment propi;
  \item definir la comparació \(\preceq_{\mathrm o}\) per segments, demostrar-ne l'antisimetria i la transitivitat, i la part d'exclusivitat de la tricotomia, sabent que la comparabilitat total demana recursió transfinita i no elecció.
\end{itemize}
\end{objectius}

Amb el llenguatge de les famílies a punt (capítol~\ref{cap:families}), la
part transfinita necessita ara la noció que en serà el fil conductor: la
d'\emph{ordre}, i sobretot la de \emph{bon ordre}. Un bon ordre és la forma
abstracta d'un procés que es recorre \enquote{de baix cap amunt} sense
saltar-se cap pas ni quedar-se encallat: tota subcol·lecció no buida té un
primer element. Els naturals ja en van ser el primer exemple
(teorema~\ref{teo:bon-ordre-N}); el que fem aquí és desprendre la propietat
de \(\N\) i estudiar-la sobre un conjunt qualsevol. La recompensa és doble.
D'una banda, els bons ordres tenen una \emph{rigidesa} notable ---entre dos
n'hi ha, com a molt, un isomorfisme d'ordre--- que els fa servir de
\emph{mesura de forma}: dir que un és \enquote{més curt} que un altre serà,
literalment, que és isomorf a un segment inicial de l'altre. De l'altra,
aquesta rigidesa és exactament el que permetrà, al capítol següent, definir
funcions per recursió al llarg d'un bon ordre. Aquest capítol estableix les
definicions ---sempre primer sobre parts, després sobre punts--- i els cinc
resultats bàsics sobre segments i comparació, sense cap axioma nou.

\section{Ordres estrictes, totals i bons ordres}\label{sec:ordres}

Com totes les relacions del llibre (capítol~\ref{cap:relacions}), un ordre
no serà una fórmula sinó una \emph{part} d'un producte; les lectures sobre
punts en són la pertinença generalitzada. Mantenim la convenció del
capítol~\ref{cap:naturals} (observació~\ref{obs:ordre-sobre-punts}): les
propietats d'un ordre són afirmacions sobre parts de \(W\times W\) o de
\(W\), i per la separació global i el lema del buit n'hi ha prou amb
comprovar-les sobre punts globals, on \enquote{o} i \enquote{existeix}
tenen el sentit clàssic. Les condicions que es formulen amb una implicació
(com la transitivitat) són inclusions de parts i valen a tots els estadis;
les que es formulen amb una disjunció (com la totalitat) són igualtats de
parts amb una unió, i sobre un estadi general un element pot quedar
repartit entre els membres.

\begin{definicio}[Ordre estricte]\index{Ordre estricte}\label{def:ordre-estricte}
Un \emph{ordre estricte} sobre un conjunt \(W\) és una part
\[
{<}\ \hookrightarrow\ W\times W
\]
que és \emph{irreflexiva} i \emph{transitiva}. Escrivim \(u<w\) per a la
pertinença generalitzada \((u,w)\mathrel{\varepsilon}{<}\)
(capítol~\ref{cap:relacions}). Les dues condicions, sobre elements
generalitzats \(u,v,w\colon X\to W\) a un estadi qualsevol, són:
\begin{itemize}
\item[(o1)] \emph{irreflexivitat}: mai \(w<w\); és a dir, la part \({<}\) és
disjunta de la diagonal \(\Delta\hookrightarrow W\times W\)
(capítol~\ref{cap:relacions});
\item[(o2)] \emph{transitivitat}: si \(u<v\) i \(v<w\), aleshores \(u<w\).
\end{itemize}
El parell \((W,<)\) és un \emph{conjunt ordenat}.
\end{definicio}

\begin{observacio}[L'ordre és una part, no un esquema]\label{obs:ordre-part-general}
Com per a \(\N\) (observació~\ref{obs:ordre-part}), \({<}\) és una part
concreta de \(W\times W\); \enquote{\(u<w\)} n'és la pertinença
generalitzada, i les condicions (o1)--(o2) es tallen amb la diagonal, la
composició de relacions i els quantificadors interns
(capítols~\ref{cap:relacions} i~\ref{cap:parts-interna}). De \({<}\) se'n
deriva la relació \emph{no estricta} \(\le\), la part
\({\le}:={<}\cup\Delta\hookrightarrow W\times W\), amb \(u\le w\Leftrightarrow
(u<w\ \vee\ u=w)\); és reflexiva, i és transitiva i antisimètrica quan
\({<}\) és un ordre estricte (la irreflexivitat de \({<}\) dona
l'antisimetria de \(\le\)). Treballarem indistintament amb \({<}\) i amb
\(\le\), segons convingui.
\end{observacio}

\begin{definicio}[Ordre total]\index{Ordre total}\label{def:ordre-total-gen}
Un ordre estricte \({<}\) sobre \(W\) és \emph{total} (o \emph{lineal}) si
la unió \(\Delta\cup{<}\cup\operatorname{Inv}(<)\) és la part total de
\(W\times W\); sobre punts globals, si dos punts qualssevol són
comparables:
\[
\text{per a tots } u,w\colon\mathbf1\to W,\qquad u=w\ \ \text{o}\ \ u<w\ \ \text{o}\ \ w<u,
\]
i aquestes tres possibilitats són mútuament excloents (\emph{tricotomia}
puntual), la darrera exclusió per la irreflexivitat i la transitivitat.
\end{definicio}

\begin{definicio}[Bon ordre]\index{Bon ordre}\label{def:bon-ordre}
Un ordre estricte \({<}\) sobre \(W\) és un \emph{bon ordre}, i \((W,<)\) un
conjunt \emph{ben ordenat}, si tota part no inicial de \(W\) té un element
mínim: per a tota part \(P\hookrightarrow W\) no inicial existeix un punt
\(m\colon\mathbf1\to W\) que es factoritza per \(P\) i compleix \(m\le k\)
per a tot element generalitzat \(k\) de \(P\). El mínim, quan existeix, és
únic.
\end{definicio}

\begin{lema}[Les parts es determinen pels punts]\label{lema:transferencia-punts}
Siguin \(R,S\hookrightarrow W\times W\) dues parts (per exemple, dues relacions
sobre \(W\)). Sota la separació global (teoria ben apuntada,
axioma~\ref{ax:separacio-globals}), si \(R\) i \(S\) tenen exactament els
mateixos \emph{punts globals} ---per a tot \(p\colon\mathbf1\to W\times W\),
\(p\) es factoritza per \(R\) si i només si es factoritza per \(S\)---,
aleshores \(R\equiv S\) com a parts, i per tant coincideixen sobre \emph{tot}
element generalitzat \(X\to W\times W\), sigui quin sigui l'estadi \(X\). En
particular, una relació que conté tots els punts globals de \(W\times W\)
és la part total.
\end{lema}

\begin{prova}
És la reducció als punts que ja vam fer al capítol~\ref{cap:relacions}
(pràctica) i al capítol~\ref{cap:naturals}: la part \(R\cap S^c\) no té cap
punt i, pel lema del buit (lema~\ref{lema:buit}), és buida, de manera que
\(R\subseteq S\) per la descomposició \(R\cong(R\cap S)+(R\cap S^c)\); i
simètricament. Dues parts incloses l'una en l'altra són equivalents, i la
pertinença d'un element generalitzat es decideix per factorització, igual
en totes dues. Observeu el que el lema \emph{no} diu: que una disjunció
verificada sobre punts (\enquote{\(u<w\) o \(w<u\)}) valgui, disjunt per
disjunt, sobre elements generalitzats; el que val és que la \emph{unió} de
les parts corresponents és total.
\end{prova}

\begin{proposicio}[Tot bon ordre és total]\label{prop:bon-ordre-total}
Si \({<}\) és un bon ordre sobre \(W\), aleshores és total. El mínim d'una
part no inicial és, a més, únic.
\end{proposicio}

\begin{prova}
La unicitat del mínim surt de l'antisimetria de \(\le\)
(observació~\ref{obs:ordre-part-general}): dos mínims \(m,m'\) compleixen
\(m\le m'\) i \(m'\le m\), d'on \(m=m'\). Per a la totalitat, provem primer la
comparabilitat sobre \emph{punts globals} i després la transferim a tots els
estadis. Siguin \(u,w\colon\mathbf1\to W\) dos punts distints i considerem la
part \(\{u,w\}\hookrightarrow W\) ---la imatge del parell
\(\langle u,w\rangle\colon\mathbf1+\mathbf1\to W\)
(teorema~\ref{teo:factoritzacio})---, que és no inicial. Pel bon ordre, té un
mínim, que és \(u\) o bé \(w\). Si el mínim és \(u\), aleshores \(u\le w\), i
com que \(u\neq w\), \(u<w\); si és \(w\), simètricament \(w<u\). Així la part
\(\Delta\cup{<}\cup\operatorname{Inv}(<)\hookrightarrow W\times W\) conté
tots els punts globals de \(W\times W\). Pel
lema~\ref{lema:transferencia-punts} (separació global), aquesta part és tot
\(W\times W\), que és la totalitat de la definició~\ref{def:ordre-total-gen}.
\end{prova}

\begin{observacio}[Lectura clàssica]\label{obs:bon-ordre-classic}
Un bon ordre és un ordre total en què tot subconjunt no buit té un primer
element. La condició, aparentment modesta, és molt forta: prohibeix les
cadenes descendents infinites \(w_0>w_1>w_2>\cdots\) ---una tal cadena tindria
per imatge una part sense mínim--- i és el que fa possible tant la inducció
com la recursió al llarg de \(W\). Els naturals amb el seu ordre habitual en
són l'exemple canònic (teorema~\ref{teo:bon-ordre-N}); en canvi \(\Z\), \(\Q\)
i \(\R\) amb l'ordre usual \emph{no} són ben ordenats, perquè tenen parts no
buides sense mínim (el mateix \(\Z\), o l'interval obert). Que \emph{tot}
conjunt admeti \emph{algun} bon ordre és una afirmació molt més forta ---el
teorema del bon ordre--- que requereix l'elecció i que veurem al
capítol~\ref{cap:zorn}.
\end{observacio}

\section{Segments inicials}\label{sec:segments}

La noció que organitza tota la comparació de bons ordres és la de segment
inicial: el tros de \(W\) que queda \enquote{estrictament per sota} d'un
punt.

\begin{definicio}[Segment inicial]\index{Segment inicial}\label{def:segment}
Sigui \((W,<)\) un conjunt ordenat i \(w\colon\mathbf1\to W\) un punt. El
\emph{segment inicial} determinat per \(w\) és la part
\[
W_{<w}\ :=\ \{\,u\in W\ \mid\ u<w\,\}\ \hookrightarrow\ W,
\]
la imatge inversa de la relació \({<}\) al llarg de
\(\langle\id_W,w\circ{!}_W\rangle\); en la lectura sobre punts, els elements
de \(W_{<w}\) són els \(u\) amb \(u<w\). Més generalment, una part
\(S\hookrightarrow W\) és un \emph{segment inicial} (o \emph{part inicial})
de \(W\) si és tancada cap avall: si \(u<w\) i \(w\in S\), aleshores \(u\in
S\).
\end{definicio}

\begin{observacio}[Els \(W_{<w}\) i les parts inicials]\label{obs:segments-inicials}
Tot segment \(W_{<w}\) és una part inicial: si \(u<w'\) i \(w'<w\), la
transitivitat dona \(u<w\). El recíproc, per a bons ordres, és gairebé cert i
molt útil: veurem (proposició~\ref{prop:parts-inicials}) que en un bon ordre
les úniques parts inicials pròpies són exactament els segments \(W_{<w}\).
Aquesta identificació ---parts inicials \(=\) segments--- és la que farà que
\enquote{ser un segment de} funcioni com una relació d'ordre entre bons
ordres.
\end{observacio}

\begin{proposicio}[Els segments d'un bon ordre són ben ordenats]\label{prop:segment-ben-ordenat}
Sigui \((W,<)\) un bon ordre i \(S\hookrightarrow W\) una part qualsevol.
Aleshores la restricció de \({<}\) a \(S\) ---la part
\({<}\cap(S\times S)\hookrightarrow S\times S\)--- és un bon ordre sobre
\(S\). En particular, tot segment inicial \(W_{<w}\) és ben ordenat.\par
\emph{Axiomes usats:} parts i interseccions (capítols~\ref{cap:parts}
i~\ref{cap:valors}); cap axioma nou.
\end{proposicio}

\begin{prova}
La restricció \({<}\!\restriction_S\,={<}\cap(S\times S)\) és una part de
\(S\times S\) (intersecció de parts, capítol~\ref{cap:valors}), i és
irreflexiva i transitiva perquè ho és \({<}\): les condicions (o1)--(o2) es
conserven en restringir a una part. Per al bon ordre, sigui
\(P\hookrightarrow S\) una part no inicial; via la inclusió
\(S\hookrightarrow W\), \(P\) és també una part no inicial de \(W\)
(la composició de monomorfismes és un monomorfisme,
proposició~\ref{prop:preceq-preordre}), que pel bon ordre de \(W\) té un
mínim \(m\). Com que \(m\in P\subseteq S\) i \(m\le k\) per a tot \(k\in
P\), aquest \(m\) és també el mínim de \(P\) dins \(S\). Per tant
\({<}\!\restriction_S\) és un bon ordre. El cas \(S=W_{<w}\) n'és una
instància.
\end{prova}

\begin{proposicio}[Les parts inicials pròpies són segments]\label{prop:parts-inicials}
Sigui \((W,<)\) un bon ordre i \(S\subsetneq W\) una part inicial pròpia.
Aleshores existeix un únic punt \(w\colon\mathbf1\to W\) tal que
\(S=W_{<w}\); concretament, \(w\) és el mínim del complement
\(W\setminus S\).
\end{proposicio}

\begin{prova}
Com que \(S\) és pròpia, el complement \(W\setminus S\)
(teorema~\ref{teo:part-complementari}) és no inicial, i pel bon ordre té un
mínim \(w\). Vegem \(S=W_{<w}\). Si \(u\in S\), aleshores \(u<w\): altrament
\(w\le u\), i com que \(S\) és inicial i \(u\in S\), de \(w\le u\) en sortiria
\(w\in S\) ---si \(w<u\) per inicialitat, i si \(w=u\) trivialment---, contra
\(w\in W\setminus S\); així \(u<w\), és a dir \(u\in W_{<w}\).
Recíprocament, si \(u<w\), aleshores \(u\in S\): com que \(w\) és el
\emph{mínim} de \(W\setminus S\), cap element estrictament menor que \(w\)
és al complement, de manera que \(u\in S\). Per tant \(S=W_{<w}\). La
unicitat de \(w\) és la del punt que determina un segment donat, que provem
tot seguit (proposició~\ref{prop:segment-determina-punt}).
\end{prova}

\begin{proposicio}[Un segment determina el seu punt]\label{prop:segment-determina-punt}
Sigui \((W,<)\) un bon ordre. Si \(W_{<v}=W_{<w}\) com a parts de \(W\),
aleshores \(v=w\). És a dir, dos segments inicials iguals ---i, més avall,
isomorfs--- corresponen al mateix punt.
\end{proposicio}

\begin{prova}
Suposem \(W_{<v}=W_{<w}\) i, cap a una contradicció, \(v\neq w\). Per la
totalitat (proposició~\ref{prop:bon-ordre-total}), \(v<w\) o \(w<v\); diguem
\(v<w\). Aleshores \(v\in W_{<w}=W_{<v}\), és a dir \(v<v\), contra la
irreflexivitat. Simètricament si \(w<v\). Per tant \(v=w\).
\end{prova}

\section{Aplicacions creixents i isomorfismes d'ordre}\label{sec:creixents}

\begin{definicio}[Aplicació creixent i isomorfisme d'ordre]\index{Aplicació creixent i isomorfisme d'ordre}\label{def:creixent}
Siguin \((W,<_W)\) i \((V,<_V)\) conjunts ordenats. Una aplicació
\(f\colon W\to V\) és \emph{creixent} (o \emph{monòtona}) si conserva
l'ordre estricte: \(u<_W w\) implica \(f(u)<_V f(w)\), per a tots els
elements generalitzats \(u,w\colon X\to W\); com a parts,
\({<_W}\subseteq(f\times f)^{-1}({<_V})\), una inclusió, que val a tots
els estadis si val sobre punts. És un \emph{isomorfisme
d'ordre} si és un isomorfisme de conjunts (definició~\ref{def:isomorfisme})
i la seva inversa també és creixent; equivalentment, si és una bijecció amb
\(u<_W w\Leftrightarrow f(u)<_V f(w)\). Escrivim \((W,<_W)\cong(V,<_V)\) quan
n'existeix un.
\end{definicio}

\begin{observacio}[Creixent i injectiva]\label{obs:creixent-injectiva}
Sobre un ordre total ---en particular un bon ordre--- tota aplicació
creixent és automàticament un monomorfisme: sobre punts, si \(f(u)=f(w)\)
amb \(u\neq w\), la totalitat dona \(u<w\) o \(w<u\), i la monotonia
\(f(u)<f(w)\) o \(f(w)<f(u)\), contra \(f(u)=f(w)\) i la irreflexivitat;
injectiva sobre punts és mono (observació~\ref{obs:injectiva-mono}). Per tant, entre
conjunts totalment ordenats, un isomorfisme d'ordre no és més que una
aplicació creixent, exhaustiva i amb inversa creixent; i, com veurem, entre
bons ordres n'hi ha com a molt un.
\end{observacio}

Els isomorfismes d'ordre entre bons ordres són \emph{rígids}: no en poden
existir dos diferents. Aquest és el resultat que fa dels bons ordres una
bona mesura de forma, i el que garantirà, al capítol següent, la unicitat de
les definicions per recursió transfinita.

\begin{lema}[Rigidesa: cap punt es mou cap avall]\label{lema:rigidesa}
Sigui \((W,<)\) un bon ordre i \(f\colon W\to W\) una aplicació creixent i
injectiva. Aleshores \(w\le f(w)\) per a tot \(w\colon\mathbf1\to W\); cap
element no baixa.
\end{lema}

\begin{prova}
Sigui \(P:=\{\,w\in W\mid f(w)<w\,\}\hookrightarrow W\) la part dels punts
que \emph{baixen} ---la imatge inversa per \(\langle f,\id_W\rangle\) de la
relació \({<}\) (capítols~\ref{cap:parts} i~\ref{cap:relacions}). Volem
\(P\) inicial. Si no ho fos, pel bon ordre tindria un mínim \(w_0\), amb
\(f(w_0)<w_0\). Aplicant \(f\), que és creixent, a la desigualtat
\(f(w_0)<w_0\) obtenim \(f(f(w_0))<f(w_0)\); és a dir, el punt
\(w_1:=f(w_0)\) també baixa, \(w_1\in P\), i \(w_1=f(w_0)<w_0\) és
estrictament menor que el mínim \(w_0\) de \(P\), una contradicció. Per tant
\(P\) és inicial: cap punt baixa, i \(w\le f(w)\) per a tot \(w\).
\end{prova}

\begin{teorema}[Unicitat de l'isomorfisme d'ordre]\label{teo:unicitat-iso}
Entre dos bons ordres \((W,<_W)\) i \((V,<_V)\) hi ha, com a molt, un
isomorfisme d'ordre. En particular, l'únic automorfisme d'ordre d'un bon
ordre és la identitat.\par
\emph{Axiomes usats:} bon ordre, parts, imatges inverses
(capítols~\ref{cap:parts} i~\ref{cap:imatges}); cap axioma nou; cap elecció.
\end{teorema}

\begin{prova}
Siguin \(f,g\colon W\to V\) dos isomorfismes d'ordre. Aleshores
\(g^{-1}\circ f\colon W\to W\) és un automorfisme d'ordre (composició
d'isomorfismes d'ordre), i també ho és \(f^{-1}\circ g\), la seva inversa.
Pel lema~\ref{lema:rigidesa} aplicat a \(g^{-1}\circ f\), \(w\le
(g^{-1}\circ f)(w)\) per a tot \(w\); aplicat a \(f^{-1}\circ g\),
\(w\le(f^{-1}\circ g)(w)\), i component amb \(g^{-1}\circ f\)
(creixent) s'obté \((g^{-1}\circ f)(w)\le w\). Les dues desigualtats donen
\((g^{-1}\circ f)(w)=w\) per a tot \(w\), és a dir \(g^{-1}\circ f=\id_W\),
d'on \(f=g\). El cas \(W=V\) dona que l'únic automorfisme és la identitat.
\end{prova}

\begin{corollari}[Un bon ordre no és isomorf a cap segment propi]\label{cor:no-segment-propi}
Cap bon ordre \((W,<)\) no és isomorf a cap dels seus segments inicials
propis \(W_{<w}\). Més en general, cap isomorfisme d'ordre no aplica \(W\)
sobre una part inicial pròpia.
\end{corollari}

\begin{prova}
Suposem un isomorfisme d'ordre \(f\colon W\cong W_{<w}\) per a algun \(w\).
Compost amb la inclusió \(W_{<w}\hookrightarrow W\), dona una aplicació
creixent i injectiva \(f\colon W\to W\) amb \(f(w)\in W_{<w}\), és a dir
\(f(w)<w\). Però això contradiu el lema~\ref{lema:rigidesa}, segons el qual
\(w\le f(w)\). Per tant no existeix cap tal isomorfisme. El cas d'una part
inicial pròpia qualsevol es redueix a un segment
(proposició~\ref{prop:parts-inicials}).
\end{prova}

\section{Comparació de bons ordres}\label{sec:comparacio-bons-ordres}

Els resultats anteriors permeten \emph{comparar} bons ordres per la relació
\enquote{ser isomorf a un segment inicial de}. Aquesta relació es comporta
com un ordre, i és la que, al capítol dels tipus ordinals
(capítol~\ref{cap:ordinals}), es convertirà en la relació \(<\) entre ordinals.

\begin{definicio}[Un bon ordre com a segment d'un altre]\index{Un bon ordre com a segment d'un altre}\label{def:mes-curt}
Direm que \((W,<_W)\) és \emph{(isomorf a) un segment inicial} de
\((V,<_V)\), i escriurem \(W\preceq_{\mathrm o} V\), si existeix un
isomorfisme d'ordre de \(W\) sobre un segment inicial de \(V\) ---sobre tot
\(V\), o sobre algun \(V_{<v}\). Escrivim \(W\prec_{\mathrm o}V\) quan és
sobre un segment \emph{propi} \(V_{<v}\), i \(W\cong V\) quan és sobre tot
\(V\).
\end{definicio}

\begin{proposicio}[La comparació per segments és antisimètrica i transitiva]\label{prop:comparacio-ordres}
Sobre bons ordres:
\begin{enumerate}
\item si \(W\preceq_{\mathrm o}V\) i \(V\preceq_{\mathrm o}W\), aleshores
\(W\cong V\);
\item si \(W\preceq_{\mathrm o}V\) i \(V\preceq_{\mathrm o}U\), aleshores
\(W\preceq_{\mathrm o}U\);
\item l'isomorfisme d'ordre de qualsevol de les situacions és \emph{únic}
(teorema~\ref{teo:unicitat-iso}).
\end{enumerate}
\emph{Axiomes usats:} els dels resultats anteriors; cap axioma nou; cap
elecció.
\end{proposicio}

\begin{prova}
(2) Si \(f\colon W\cong V_{\le}\) sobre un segment de \(V\) i \(g\colon
V\cong U_{\le}\) sobre un segment de \(U\), aleshores \(g\) aplica el segment
\(V_{\le}\) sobre un segment de \(U\). En efecte, sigui \(T\) la part
inicial de \(U\) sobre la qual \(g\) és un isomorfisme d'ordre, i \(S\) una
part inicial de \(V\). Si \(u<_U g(s)\) amb \(s\in S\), aleshores \(u\in T\)
perquè \(T\) és inicial, de manera que \(u=g(v)\) per a un \(v\in V\); com que
\(g\) també \emph{reflecteix} l'ordre (la seva inversa és creixent),
\(v<_V s\), i \(v\in S\) perquè \(S\) és inicial. Així \(g(S)\) és una part
inicial de \(U\) (proposició~\ref{prop:parts-inicials}). Per tant \(g\circ f\) aplica
\(W\) isomòrficament sobre un segment de \(U\); així \(W\preceq_{\mathrm
o}U\). (1) Si \(W\preceq_{\mathrm o}V\preceq_{\mathrm o}W\), per (2)
\(W\preceq_{\mathrm o}W\) via un isomorfisme d'ordre de \(W\) sobre un
segment inicial de \(W\); pel corol·lari~\ref{cor:no-segment-propi} aquest
segment no pot ser propi, de manera que és tot \(W\) i la composició és un
automorfisme ---la identitat (teorema~\ref{teo:unicitat-iso}). Desfent-la,
cadascun dels dos isomorfismes és sobre tot l'altre, d'on \(W\cong V\). (3)
és el teorema~\ref{teo:unicitat-iso}.
\end{prova}

Falta el resultat central de comparació: que dos bons ordres qualssevol
\emph{sempre} són comparables ---un dels dos és isomorf a un segment de
l'altre. És la \emph{tricotomia}, i és la peça que, en el llenguatge
estructural, substitueix moltes manipulacions clàssiques amb ordinals de von
Neumann. La seva prova, però, construeix l'isomorfisme de comparació
\emph{punt per punt al llarg del bon ordre}, i això és una recursió
transfinita, que encara no tenim. Enunciem la tricotomia aquí, amb la seva
dependència declarada, i en donem la part que \emph{no} depèn de la recursió;
la construcció de l'isomorfisme es completa al capítol següent.

\begin{teorema}[Tricotomia de la comparació; part sense recursió]\label{teo:tricotomia}
Siguin \((W,<_W)\) i \((V,<_V)\) dos bons ordres. Aleshores \emph{com a molt
una} de les tres situacions es dona:
\[
W\prec_{\mathrm o}V,
\qquad
W\cong V,
\qquad
V\prec_{\mathrm o}W,
\]
i quan es dona, l'isomorfisme d'ordre corresponent és únic. Que \emph{almenys
una} es dona sempre ---la comparabilitat total--- es demostra construint un
isomorfisme d'ordre entre \(W\) i un segment de \(V\), o a l'inrevés, per
recursió transfinita; ho establirem al teorema~\ref{teo:comparabilitat} del
capítol~\ref{cap:recursio}.\par
\emph{Axiomes usats (per a l'exclusivitat):} unicitat de l'isomorfisme
d'ordre i no-isomorfisme amb segments propis
(teorema~\ref{teo:unicitat-iso}, corol·lari~\ref{cor:no-segment-propi}); cap
elecció. La comparabilitat total usa la recursió
transfinita del capítol~\ref{cap:recursio}, també sense elecció.
\end{teorema}

\begin{prova}[Exclusivitat mútua]
Provem que les tres situacions són incompatibles dues a dues; l'existència
d'almenys una es remet al capítol~\ref{cap:recursio}.
Suposem \(W\cong V\) i
alhora \(W\prec_{\mathrm o}V\); aleshores, component, \(V\cong W\prec_{\mathrm
o}V\) dona \(V\) isomorf a un segment propi de si mateix, contra el
corol·lari~\ref{cor:no-segment-propi}. Igual amb \(W\cong V\) i
\(V\prec_{\mathrm o}W\). Finalment, si \(W\prec_{\mathrm o}V\) i
\(V\prec_{\mathrm o}W\) alhora, la transitivitat
(proposició~\ref{prop:comparacio-ordres}) dona \(W\prec_{\mathrm o}W\), un
isomorfisme de \(W\) sobre un segment propi, altre cop contra el
corol·lari~\ref{cor:no-segment-propi}. Per tant les tres són mútuament
excloents, i la unicitat de l'isomorfisme és el
teorema~\ref{teo:unicitat-iso}.
\end{prova}

\begin{observacio}[Per què la comparabilitat necessita recursió, i per què no elecció]\label{obs:tricotomia-recursio}
La direcció que aquí deixem pendent ---que almenys una de les tres situacions
es dona--- no és una mancança de la nostra prova, sinó un fet estructural:
per decidir quin dels dos bons ordres és \enquote{més curt} cal
\emph{aparellar-los} element a element començant pels mínims respectius, i
aquest aparellament es defineix per recursió al llarg d'un dels dos ordres.
No hi intervé cap \emph{tria}: a cada pas, l'element que toca és el
\emph{mínim} del que queda, que és únic i canònic (com als subobjectes de
\(\N\), observació~\ref{obs:minim-vs-eleccio}); per això la comparabilitat de
bons ordres, a diferència de la comparabilitat de conjunts \emph{qualssevol}
(que veurem amb el teorema del bon ordre), \emph{no} usa l'elecció. El que necessita és
únicament el principi de recursió transfinita amb codomini fix, que
construïm al capítol~\ref{cap:recursio}; per això n'hem separat l'enunciat.
Aquesta és, també, la disciplina que la part transfinita mantindrà: enunciar
un resultat només quan les seves eines ja existeixen, o declarar
explícitament de quina depèn.
\end{observacio}

\begin{resumcap}
\apartat{Idees centrals}
\begin{enumerate}
  \item Un ordre estricte és una part \({<}\hookrightarrow W\times W\) irreflexiva i transitiva; total quan \(\Delta\cup{<}\cup\operatorname{Inv}(<)\) és tot \(W\times W\); bon ordre quan tota part no inicial té mínim. Les propietats es comproven sobre punts, i les disjunctives s'entenen com a unions.
  \item Tot bon ordre és total i el mínim és únic; els segments \(W_{<w}\) són les parts inicials pròpies, són ben ordenats i determinen el seu punt.
  \item Rigidesa: sota una creixent injectiva d'un bon ordre en si mateix cap punt baixa. D'aquí, entre dos bons ordres hi ha com a molt un isomorfisme d'ordre, l'únic automorfisme és la identitat, i cap bon ordre és isomorf a un segment propi.
  \item \(W\preceq_{\mathrm o}V\) (\(W\) isomorf a un segment inicial de \(V\)) és antisimètrica llevat d'isomorfisme i transitiva, i les tres situacions \(W\prec_{\mathrm o}V\), \(W\cong V\), \(V\prec_{\mathrm o}W\) s'exclouen. Que almenys una es doni demana recursió transfinita, no elecció.
\end{enumerate}
\apartat{Errors típics} Llegir la totalitat com \enquote{\(u<w\) o \(w<u\)} per a elements generalitzats (val sobre punts; com a parts és una unió total); creure que un bon ordre pot tenir automorfismes no trivials (com els ordres totals densos); confondre \(\preceq_{\mathrm o}\) (segments, forma) amb \(\preceq\) (monomorfismes, mida) del capítol~\ref{cap:cantor}; i pensar que la comparabilitat de bons ordres usa l'elecció, quan el que usa és el mínim.
\apartat{Lectura recomanada} \cite[§§14 i 17]{halmos} per als conjunts ordenats, els bons ordres, els segments i la comparabilitat en la teoria clàssica, que és el que aquí es reconstrueix sobre parts; \cite[cap.~2]{jech2003} per a la teoria clàssica dels bons ordres i dels ordinals; \cite[cap.~I.6--I.7]{kunen2011} per a la presentació axiomàtica estàndard dels bons ordres i la recursió; i \cite[§§8.4--8.5]{lawvere-rosebrugh} per a l'ordre i els bons ordres en conjunts abstractes.
\end{resumcap}

\chapter{Inducció i recursió transfinites}\label{cap:recursio}

\begin{objectius}
\apartat{Prerequisits} Del capítol~\ref{cap:bons-ordres}: bons ordres, segments inicials, parts inicials, rigidesa i la comparació \(\preceq_{\mathrm o}\). Del capítol~\ref{cap:families}: famílies de parts d'un conjunt fix, unions indexades, i la distinció família interna\,--\,regla externa. Dels capítols~\ref{cap:exponencials} i~\ref{cap:parts-interna}: exponencials i parts tallades per condicions internes. Del capítol~\ref{cap:imatges}: el graf d'una relació funcional i total.
\apartat{Objectius} En acabar el capítol, el lector hauria de poder:
\begin{itemize}
  \item enunciar i demostrar la inducció transfinita pel mínim del complement, i explicar per què no cal cas base;
  \item formalitzar aplicacions parcials, històries i regles de recursió dins l'exponencial fix \((X+\mathbf1)^W\);
  \item demostrar la recursió transfinita amb codomini fix: compatibilitat de solucions parcials, existència fins a cada punt, reunió global i unicitat, sense elecció;
  \item demostrar la comparabilitat total dels bons ordres aparellant pel mínim lliure, i situar-la respecte de la comparabilitat de conjunts (que sí que demana elecció);
  \item reconèixer definicions per recursió amb codomini fix (sumes iterades, successions de parts, tancaments) i saber quan la separació en casos successor/límit és legítima;
  \item identificar la frontera: quan els valors són conjunts nous sense codomini comú, la recursió d'aquest capítol no basta i cal SR.
\end{itemize}
\end{objectius}

El capítol anterior va establir la rigidesa dels bons ordres i en va deixar
una peça pendent: que dos bons ordres qualssevol són comparables
(teorema~\ref{teo:tricotomia}). La comparació es construeix aparellant els
dos ordres començant pels mínims, i aquest aparellament és una \emph{recursió}
al llarg d'un bon ordre. Aquest capítol construeix l'eina que faltava. Primer
la \emph{inducció transfinita}, que generalitza a un bon ordre qualsevol el
principi d'inducció de \(\N\) (teorema~\ref{teo:induccio}): per provar que una
propietat val a tot \(W\), n'hi ha prou de provar-la a cada \(w\) suposant-la
a tot el segment inicial \(W_{<w}\). Després la \emph{recursió transfinita amb
codomini fix}: quan els valors d'una definició viuen tots dins un conjunt
\(X\) prefixat, una regla que digui com determinar \(F(w)\) a partir de la
\enquote{història} \(F|_{W_{<w}}\) defineix una única aplicació \(F\colon W\to
X\). Tot això és de Nivell~I i no costa cap axioma nou: viu dins exponencials
i conjunts de parts ja construïts. Amb la recursió a la mà, tanquem la
comparabilitat total dels bons ordres. Finalment marquem \emph{la frontera}:
quan els valors \(F(w)\) han de ser conjunts \emph{nous} sense codomini comú,
la recursió d'aquest capítol no basta i cal la substitució estructural SR
---el pas a Nivell~II---, que il·lustrem amb un exemple.

\section{Inducció transfinita}\label{sec:induccio-transfinita}

\begin{teorema}[Inducció transfinita]\label{teo:induccio-transfinita}
Sigui \((W,<)\) un bon ordre i \(S\hookrightarrow W\) una part que compleix
la condició de \emph{progressivitat}
\[
\bigl(W_{<w}\subseteq S\bigr)\ \Longrightarrow\ w\in S
\qquad\text{per a tot } w\colon\mathbf1\to W .
\]
Aleshores \(S=W\).\par
\emph{Axiomes usats:} bon ordre (definició~\ref{def:bon-ordre}), complement
(capítol~\ref{cap:valors}); cap axioma nou; cap elecció.
\end{teorema}

\begin{prova}
Suposem, cap a una contradicció, que \(S\neq W\); aleshores el complement
\(W\setminus S\) (teorema~\ref{teo:part-complementari}) és no inicial. Pel bon
ordre (definició~\ref{def:bon-ordre}) té un element mínim \(w_0\). Com que
\(w_0\) és el mínim del complement, tot element estrictament menor que \(w_0\)
és fora del complement, és a dir dins \(S\): per tant \(W_{<w_0}\subseteq S\).
Però aleshores la progressivitat, aplicada a \(w_0\), dona \(w_0\in S\),
contra \(w_0\in W\setminus S\). La contradicció mostra que \(W\setminus S\) és
inicial, és a dir \(S=W\).
\end{prova}

\begin{observacio}[La forma habitual, sobre propietats]\label{obs:induccio-propietats}
En la pràctica \(S\) és la part on val una propietat interna
\(\varphi\) ---la part classificada per \(\varphi\)
(capítol~\ref{cap:parts-interna})---, i el teorema es llegeix així: si de
\enquote{\(\varphi\) val a tot \(u<w\)} se'n dedueix \enquote{\(\varphi\) val
a \(w\)}, per a tot \(w\), aleshores \(\varphi\) val a tot \(W\). No cal cap
\emph{cas base} separat: el cas del mínim global \(\bot\) de \(W\) hi és
inclòs, perquè \(W_{<\bot}\) és inicial i la hipòtesi
\enquote{\(W_{<\bot}\subseteq S\)} es compleix buidament, de manera que la
progressivitat força \(\bot\in S\). Aquesta és la diferència amb la inducció
de \(\N\) (teorema~\ref{teo:induccio}), on el zero i el successor se separaven:
en un bon ordre general no tot element és un successor, i la formulació
\enquote{tot el segment inicial} unifica els dos casos ---i els límits--- en
un de sol.
\end{observacio}

\begin{observacio}[Lectura clàssica]\label{obs:induccio-transf-classica}
És la \emph{inducció ben fundada} o \emph{inducció transfinita} habitual: per
demostrar que una propietat val per a tot element d'un conjunt ben ordenat,
es demostra per a cada element suposant-la per a tots els anteriors. Sobre
\(\N\) recupera la inducció forta (o \emph{completa}). La bona fundació ---la
inexistència de cadenes descendents infinites, garantida pel bon ordre
(observació~\ref{obs:bon-ordre-classic})--- és exactament el que fa vàlid el
principi: sense ella, el \enquote{mínim del complement} podria no existir.
Convé mantenir aquesta implicació en la direcció \emph{segura}: tot bon ordre
exclou tota cadena descendent (fins i tot numerable). El \emph{recíproc}
---que un ordre sense cadenes descendents numerables sigui ben fundat, o ben
ordenat--- \emph{no} és gratuït: la seva demostració habitual construeix una
cadena descendent triant successivament testimonis, i per tant pot requerir
alguna forma d'elecció (típicament l'elecció dependent). No l'usem en aquesta
direcció sense declarar-ho.
\end{observacio}

\section{Recursió transfinita amb codomini fix}\label{sec:recursio-fix}

Passem de \emph{demostrar} per inducció a \emph{definir} per recursió. La
situació és: un conjunt \emph{fix} \(X\), i una regla que, coneguts els valors
de \(F\) sobre tot el segment \(W_{<w}\), determina el valor \(F(w)\). Volem
que aquesta regla defineixi una única \(F\colon W\to X\). La formalització de
\enquote{els valors sobre \(W_{<w}\)} és una \emph{història parcial}.

\begin{definicio}[Aplicació parcial i història]\index{Aplicació parcial i història}\label{def:historia}
Una \emph{aplicació parcial} de \(W\) cap a \(X\) és una aplicació
\(\bar h\colon W\to X+\mathbf1\) cap a \(X\) ampliat amb un punt de
\enquote{no definit} \(\ast\colon\mathbf1\to X+\mathbf1\); el seu
\emph{domini de definició} és la part \(\{w\in W: \bar h(w)\neq\ast\}\). Tota
aplicació parcial viu dins l'exponencial \emph{fix} \((X+\mathbf1)^{W}\)
(axioma~\ref{ax:exponencials}), i la seva restricció a una part és de nou una
aplicació parcial. Una \emph{història en \(w\)} és una aplicació
\(h\colon W_{<w}\to X\) ---una aplicació parcial amb domini de definició
exactament el segment \(W_{<w}\).
\end{definicio}

\begin{definicio}[Regla de recursió]\index{Regla de recursió}\label{def:regla-recursio}
Una \emph{regla de recursió} amb valors a \(X\) és una manera interna
d'assignar, a cada punt \(w\colon\mathbf1\to W\) i cada història
\(h\colon W_{<w}\to X\), un element \(\Phi(w,h)\in X\). Formalment, \(\Phi\)
és una aplicació interna del conjunt de parells \((w,h)\) ---una part de
\(W\times(X+\mathbf1)^{W}\), tallada per la condició \enquote{\(h\) és una
història en \(w\)}, que és interna (el domini de definició de \(h\) és
\(W_{<w}\): igualtats i quantificadors acotats sobre l'exponencial fix,
capítols~\ref{cap:exponencials} i~\ref{cap:parts-interna})--- cap a \(X\).
\end{definicio}

\begin{teorema}[Recursió transfinita amb codomini fix]\label{teo:recursio-transfinita}
Siguin \((W,<)\) un bon ordre, \(X\) un conjunt i \(\Phi\) una regla de
recursió amb valors a \(X\) (definició~\ref{def:regla-recursio}). Aleshores
existeix una \emph{única} aplicació \(F\colon W\to X\) tal que
\[
F(w)=\Phi\bigl(w,\ F|_{W_{<w}}\bigr)
\qquad\text{per a tot } w\colon\mathbf1\to W .
\]
\emph{Axiomes usats:} exponencials i conjunt de parts
(capítols~\ref{cap:exponencials} i~\ref{cap:parts-interna}), inducció
transfinita (teorema~\ref{teo:induccio-transfinita}), imatges i unions de
parts (capítols~\ref{cap:imatges} i~\ref{cap:valors}); cap axioma nou; cap
elecció.
\end{teorema}

\begin{prova}
Diem que una aplicació parcial \(a\colon D\to X\), amb domini de definició una
part inicial \(D\hookrightarrow W\), és una \emph{solució parcial} (o
\emph{aproximació}) si compleix l'equació de recursió allà on és definida:
\[
a(w)=\Phi\bigl(w,\ a|_{W_{<w}}\bigr)
\qquad\text{per a tot } w\in D .
\]
Aquesta condició té sentit perquè, si \(D\) és inicial i \(w\in D\), aleshores
\(W_{<w}\subseteq D\) i la restricció \(a|_{W_{<w}}\) és una història en \(w\).
El conjunt de les solucions parcials és una part interna \(\mathcal
A\hookrightarrow(X+\mathbf1)^{W}\), tallada per la condició
\enquote{el domini és inicial} i per l'equació anterior (quantificadors
acotats i igualtats sobre l'exponencial fix,
capítols~\ref{cap:exponencials} i~\ref{cap:parts-interna}). Tot el que segueix
viu dins aquest exponencial fix i dins \(P(W)\).

\emph{Compatibilitat.} Dues solucions parcials coincideixen allà on totes
dues són definides: si \(a\colon D\to X\) i \(b\colon E\to X\) són solucions
parcials, aleshores \(a=b\) sobre \(D\cap E\). En efecte, \(D\cap E\) és una
part inicial (intersecció de parts inicials), i apliquem la inducció
transfinita (teorema~\ref{teo:induccio-transfinita}) sobre \(W\) a la part
\[
S=\{\,w\in W\ \mid\ w\notin D\cap E\ \text{ o bé }\ a(w)=b(w)\,\}.
\]
Sigui \(w\) amb \(W_{<w}\subseteq S\); volem \(w\in S\). Si \(w\notin D\cap
E\), ja hi és. Si \(w\in D\cap E\), aleshores \(W_{<w}\subseteq D\cap E\)
(inicialitat), i per la hipòtesi \(W_{<w}\subseteq S\) tenim \(a=b\) sobre
\(W_{<w}\); per tant \(a|_{W_{<w}}=b|_{W_{<w}}\), i l'equació de recursió dona
\[
a(w)=\Phi(w,a|_{W_{<w}})=\Phi(w,b|_{W_{<w}})=b(w),
\]
és a dir \(w\in S\). Per inducció \(S=W\), és a dir \(a=b\) sobre \(D\cap E\).

\emph{Existència de solucions prou llargues.} Vegem, també per inducció, que
per a cada \(w\colon\mathbf1\to W\) existeix una solució parcial el domini de
la qual conté \(w\) ---de fet conté el segment \(W_{\le w}=W_{<w}\cup\{w\}\).
Sigui
\[
T=\{\,w\in W\ \mid\ \text{existeix una solució parcial definida a } W_{\le w}\,\}.
\]
Prenem \(w\) amb \(W_{<w}\subseteq T\). Per a cada \(u<w\) hi ha, doncs, una
solució parcial definida fins a \(u\); per la compatibilitat, totes aquestes
solucions coincideixen a les seves interseccions, de manera que la seva
\emph{reunió} és una única aplicació ben definida \(g\colon W_{<w}\to X\)
---el graf és la unió indexada dels grafs, funcional per la compatibilitat
(capítol~\ref{cap:families}; vegeu la reunió global, més avall)---, i \(g\) és una solució parcial amb domini el
segment inicial \(W_{<w}\) (satisfà l'equació a cada \(u<w\) perquè hi
coincideix amb alguna solució parcial que sí que la satisfà). Ara estenem
\(g\) a \(W_{\le w}\) posant
\[
g^{+}(u)=g(u)\ \text{ per a } u<w,
\qquad
g^{+}(w)=\Phi\bigl(w,\ g\bigr),
\]
que és una aplicació parcial amb domini inicial \(W_{\le w}\); satisfà
l'equació a tot \(u<w\) (com \(g\)) i a \(w\) per construcció, ja que
\(g^{+}|_{W_{<w}}=g\). Per tant \(g^{+}\) és una solució parcial definida a
\(W_{\le w}\), i \(w\in T\). Per inducció transfinita \(T=W\).

\emph{La reunió global.} Per la compatibilitat, totes les solucions parcials
formen una família compatible, i es poden reunir en una sola aplicació.
Concretament, sigui \(G\hookrightarrow W\times X\) la relació
\(\{(w,x):\exists a\in\mathcal A\ (a(w)=x)\}\), una unió indexada per
\(\mathcal A\) (capítol~\ref{cap:families}); és funcional per la
compatibilitat, i el seu domini \(W'=\bigcup\{\,D\mid a\colon D\to X\
\text{solució parcial}\,\}\) és una part inicial de \(W\). Per la
proposició~\ref{prop:aplicacio-funcional}, \(G\) és el graf d'una única
aplicació \(F\colon W'\to X\), que coincideix amb cada solució parcial sobre
el seu domini; i \(F\) mateixa és una solució parcial. Com que \(T=W\), per a cada \(w\) hi ha una solució parcial
definida a \(W_{\le w}\), de manera que \(w\in W'\); per tant \(W'=W\) i
\(F\colon W\to X\) és total. Satisfà l'equació de recursió a tot \(w\) perquè
és una solució parcial amb domini \(W\).

\emph{Unicitat.} Si \(F,F'\colon W\to X\) satisfan totes dues l'equació, són
solucions parcials amb domini \(W\); per la compatibilitat coincideixen sobre
\(W\cap W=W\), és a dir \(F=F'\).
\end{prova}

\begin{observacio}[On és cada axioma, i on no hi ha elecció]\label{obs:recursio-sense-eleccio}
La construcció no fa cap tria. Cada solució parcial que apareix està
\emph{determinada} per l'equació de recursió ---no se'n tria cap d'un conjunt
d'alternatives---, i la reunió de solucions compatibles és una operació
canònica sobre parts (capítol~\ref{cap:families}), no una selecció. Tot viu
dins l'exponencial fix \((X+\mathbf1)^{W}\) i dins \(P(W)\): és la
manifestació, aquí, del principi de la guia segons el qual mentre els valors
visquin en un codomini fix la família és interna i tot és de Nivell~I. L'única
cosa que fa el bon ordre és garantir, via la inducció transfinita, la
compatibilitat i l'existència fins a cada punt; en un ordre no ben fundat la
prova s'ensorraria, perquè el \enquote{mínim del complement} podria no
existir. Compareu amb la comparabilitat de conjunts \emph{qualssevol}, que sí
que necessita elecció (capítol~\ref{cap:zorn}): la diferència és, altre cop,
que aquí cada pas té una resposta canònica.
\end{observacio}

\begin{observacio}[La recursió de \(\N\) com a cas particular]\label{obs:recursio-N-cas}
Sobre \(W=\N\) amb el seu bon ordre (teorema~\ref{teo:bon-ordre-N}), una regla
que només depèn del valor anterior ---\(\Phi(w,h)\) donat per \(h(\text{el
predecessor de }w)\), i un valor inicial a \(0\)--- és exactament la recursió
primitiva amb paràmetres (teorema~\ref{teo:recursio-parametres}). La recursió
transfinita la generalitza en dues direccions: el pas pot mirar \emph{tota} la
història, no només el terme anterior, i l'índex pot ser un bon ordre qualsevol,
amb punts límit que no tenen predecessor immediat. És per als límits que la
formulació \enquote{història sobre tot \(W_{<w}\)} és imprescindible: allà no
hi ha cap \enquote{terme anterior} d'on partir.
\end{observacio}

\section{La comparabilitat total dels bons ordres}\label{sec:comparabilitat}

Amb la recursió podem tancar la comparació de bons ordres que el capítol
anterior va deixar oberta: dos bons ordres qualssevol són comparables. La
idea és definir, per recursió, l'única aplicació que aparella \(W\) amb un
segment inicial de \(V\) enviant cada element al \enquote{primer lloc lliure}.

\begin{teorema}[Comparabilitat total dels bons ordres]\label{teo:comparabilitat}
Siguin \((W,<_W)\) i \((V,<_V)\) dos bons ordres. Aleshores es dona
\emph{almenys una} ---i per tant, amb el teorema~\ref{teo:tricotomia},
\emph{exactament una}--- de les tres situacions:
\[
W\prec_{\mathrm o}V,
\qquad
W\cong V,
\qquad
V\prec_{\mathrm o}W .
\]
En conseqüència, la comparació per segments és un ordre total (llevat
d'isomorfisme d'ordre) sobre els bons ordres.\par
\emph{Axiomes usats:} recursió transfinita
(teorema~\ref{teo:recursio-transfinita}), bon ordre i segments
(capítol~\ref{cap:bons-ordres}); cap axioma nou; \emph{cap elecció}.
\end{teorema}

\begin{prova}
Adjuntem a \(V\) un punt exterior \(\top\), formant \(V^{+}=V+\mathbf1\), i
definim per recursió transfinita (teorema~\ref{teo:recursio-transfinita})
sobre \(W\) una aplicació \(F\colon W\to V^{+}\) amb la regla següent: donada
la història \(F|_{W_{<w}}\), sigui
\[
F(w)\ :=\
\begin{cases}
\min\bigl(V\setminus F(W_{<w})\bigr) & \text{si } V\setminus F(W_{<w})\neq\varnothing,\\[1mm]
\top & \text{altrament,}
\end{cases}
\]
on \(F(W_{<w})\hookrightarrow V\) és la imatge de la història
(capítol~\ref{cap:imatges}), \(V\setminus F(W_{<w})\) el seu complement a
\(V\), i el mínim és el del bon ordre de \(V\) ---únic i canònic, de manera
que la regla és interna i \emph{no} usa elecció. Pel
teorema~\ref{teo:recursio-transfinita}, \(F\) existeix i és única.

\emph{\(F\) és creixent mentre no val \(\top\).} Si \(u<_W w\) i \(F(w)\neq
\top\), aleshores \(F(u)\neq\top\) i \(F(u)<_V F(w)\). Llegim \(F(W_{<w})\)
com la part de \(V\) formada pels valors de la història diferents de
\(\top\); com que \(W_{<u}\subseteq W_{<w}\), es té
\(F(W_{<u})\subseteq F(W_{<w})\). Primer, \(F(u)\neq\top\): si fos
\(F(u)=\top\), seria \(F(W_{<u})=V\), d'on \(F(W_{<w})=V\) i \(F(w)=\top\),
contra la hipòtesi. Així \(F(u)=\min(V\setminus F(W_{<u}))\), i \(F(u)\in
F(W_{<w})\) perquè \(u<_W w\). En canvi, \(F(w)=\min(V\setminus F(W_{<w}))\)
no és a \(F(W_{<w})\); en particular \(F(w)\neq F(u)\), i \(F(w)\) tampoc no
és a \(F(W_{<u})\), és a dir, \(F(w)\in V\setminus F(W_{<u})\). Per la
\emph{minimalitat de \(F(u)\)} en aquest complement, \(F(u)\le_V F(w)\), i com
que són diferents, \(F(u)<_V F(w)\).
En particular, sobre la part \(W_0:=\{w: F(w)\neq\top\}\), que és inicial
(si \(F(w)\neq\top\) i \(u<_W w\), aleshores \(F(u)\neq\top\)), \(F\) és
creixent i injectiva, i la seva imatge és una part inicial de \(V\): si \(v<_VF(w)\) amb
\(w\in W_0\), aleshores \(v\in F(W_{<w})\), perquè \(F(w)\) és el mínim del
complement de \(F(W_{<w})\) i \(v\) no pot ser-hi; així \(v\) és a la
imatge.

\emph{Els tres casos.} Distingim segons on s'atura \(F\).
\begin{itemize}
\item Si \(F\) no pren mai el valor \(\top\) (és a dir \(W_0=W\)), aleshores
\(F\colon W\to V\) és un isomorfisme d'ordre sobre un segment inicial de
\(V\): creixent, injectiva, amb imatge inicial. Si la imatge és tot \(V\),
\(W\cong V\); si és un segment propi, \(W\prec_{\mathrm o}V\).
\item Si \(F\) pren el valor \(\top\), sigui \(w_0\) el mínim amb
\(F(w_0)=\top\). Aleshores \(V\setminus F(W_{<w_0})=\varnothing\), és a dir
\(F(W_{<w_0})=V\): la restricció \(F|_{W_{<w_0}}\colon W_{<w_0}\to V\) és
creixent, injectiva i exhaustiva, per tant un isomorfisme d'ordre
\(W_{<w_0}\cong V\). Així \(V\) és isomorf al segment \(W_{<w_0}\) de \(W\), que és propi
(\(w_0\notin W_{<w_0}\)), d'on \(V\prec_{\mathrm o}W\).
\end{itemize}
En tots els casos es dona una de les tres situacions. L'exclusivitat i la
unicitat de l'isomorfisme són el teorema~\ref{teo:tricotomia}. Cap pas no ha
usat elecció: l'única selecció a cada etapa és el mínim, canònic.
\end{prova}

\begin{observacio}[Comparabilitat de bons ordres vs. d'ordres qualssevol]\label{obs:comparabilitat-vs-eleccio}
Aquest teorema completa el que l'observació~\ref{obs:tricotomia-recursio}
anunciava: la comparabilitat \emph{total} dels bons ordres és un teorema
\emph{sense elecció}, perquè l'aparellament es fa prenent a cada pas el mínim
disponible. És un contrast important amb la comparabilitat de conjunts
\emph{qualssevol} ---donats \(A\) i \(B\), \(A\preceq B\) o \(B\preceq A\)---,
que sí que requereix l'elecció i que provarem via el teorema del bon ordre (capítol~\ref{cap:zorn}). La moralitat: un bon ordre porta incorporada, a
cada pas, la tria canònica que en un conjunt nu s'ha de postular.
\end{observacio}

\section{Exemples de definició per recursió}\label{sec:exemples-recursio}

La recursió transfinita amb codomini fix cobreix una gran varietat de
definicions. En donem tres, totes de Nivell~I ---el codomini és fix en cada
cas--- i una nota sobre la separació en casos successor i límit.

\begin{observacio}[Sumes iterades i tancaments]\label{obs:exemples-recursio}
\emph{Sumes iterades en un monoide.} Si \((M,\cdot,e)\) és un monoide intern i
\((a_v)_{v\in W}\) una família d'elements de \(M\) indexada per un bon ordre
\(W\) ---una aplicació \(W\to M\)---, la recursió amb \(X=M\) i
\(\Phi(w,h)=\bigl(\prod_{u<w}a_u\bigr)\cdot a_w\), és a dir \enquote{el
producte de la història, multiplicat pel terme actual}, defineix les sumes
parcials. El codomini \(M\) és fix.

\emph{Successió transfinita de parts d'un fix.} Amb \(X=P(A)\) per a un
conjunt fix \(A\), qualsevol regla que construeixi \(S_w\hookrightarrow A\) a
partir de les parts anteriors \((S_u)_{u<w}\) ---per exemple
\(S_w=\bigcup_{u<w}g(S_u)\) per a una operació \(g\colon P(A)\to P(A)\)---
defineix una successió transfinita \((S_w)_{w\in W}\) de parts d'\(A\). Com que
totes viuen dins el fix \(P(A)\), el codomini és fix i n'hi ha prou amb el
teorema~\ref{teo:recursio-transfinita}.

\emph{Tancament d'una part sota una operació.} Donada \(g\colon P(A)\to P(A)\)
monòtona, el \emph{tancament} d'una part \(S_0\) sota \(g\) s'obté iterant
\(g\) transfinitament fins a l'estabilitat; cada iterada és una part del fix
\(A\), de manera que tota la construcció és de Nivell~I. És la mateixa tècnica
del menor punt fix que ja vam usar a Schröder--Bernstein
(teorema~\ref{teo:schroder-bernstein}), ara desplegada com a recursió.
\end{observacio}

\begin{observacio}[Successor i límit: separar només quan és exhaustiu]\label{obs:successor-limit}
És còmode definir una recursió separant tres casos ---el mínim \(\bot\), els
punts \emph{successors} (els que tenen un predecessor immediat) i els punts
\emph{límit} (els que no en tenen)--- i donant una clàusula per a cadascun.
Aquesta separació és \emph{legítima} només si s'ha comprovat que és
\emph{exhaustiva} per al bon ordre considerat: que tot element és \(\bot\), o
successor, o límit. Convé demostrar-ho abans d'usar-la, perquè la nostra
formulació general (teorema~\ref{teo:recursio-transfinita}) \emph{no} la
requereix ---la regla \(\Phi(w,h)\) mira tota la història i cobreix els tres
casos uniformement. La separació és una comoditat de presentació, no una
necessitat; i escriure-la sense provar-ne l'exhaustivitat és una de les fonts
d'error més habituals en els arguments transfinits.
\end{observacio}

\section{La frontera: quan el codomini no és fix}\label{sec:frontera-sr}

\begin{nota}
Aquesta secció torna a tocar el \emph{Nivell~II}: descriu una situació
que la recursió d'aquest capítol \emph{no} cobreix, i per a la qual cal la
substitució estructural SR, ja anunciada al capítol~\ref{cap:families}.
\end{nota}

Tot el que hem fet suposa un codomini \emph{fix} \(X\): els valors \(F(w)\)
són elements d'un conjunt donat per endavant. Hi ha, però, recursions on el
valor \(F(w)\) ha de ser un \emph{conjunt nou}, que no viu dins cap codomini
prefixat.

\begin{observacio}[Un exemple on el codomini no és fix]\label{obs:frontera-exemple}
Considerem la regla que a cada natural \(n\) assigna la iterada del conjunt de
parts,
\[
V_0=A,
\qquad
V_{n+1}=P(V_n),
\]
i que en un bon ordre amb límits assignaria als límits la reunió dels
anteriors. Cada \(V_{n}\) és estrictament més gran que l'anterior pel teorema
de Cantor (teorema~\ref{teo:cantor}), i cap dels axiomes de Nivell~I no
proporciona un conjunt \(X\) dins del qual tots els \(V_n\) s'allotgin com
a parts d'una manera uniforme: les construccions finites a partir d'\(A\)
donen conjunts de mida \(P^k(A)\) per a un \(k\) concret, i no hi ha cap
construcció que els recorri tots. La successió \((V_n)_{n\in\N}\) és, doncs,
una \emph{regla externa} (definició~\ref{def:familia-vs-regla}) sense
codomini fix conegut, i el teorema~\ref{teo:recursio-transfinita} no s'hi
aplica. (No és que un tal \(X\) sigui \emph{impossible}: en models més
grans n'hi ha; és que els axiomes de Nivell~I no en garanteixen cap, i de
fet hi ha models de tots ells on la reunió no existeix.) Reunir aquesta successió en una família interna
\(n\mapsto V_n\) ---i, més encara, prolongar-la transfinitament--- és
exactament el que fa la substitució estructural SR
(observació~\ref{obs:sr}): una recursió els valors de la qual són conjunts
variables. Sense SR, la successió dels \(V_n\) queda com a regla externa,
definible terme a terme però no reunible en un sol objecte.
\end{observacio}

\begin{observacio}[La línia divisòria, i el que en depèn]\label{obs:frontera-balanc}
La distinció és neta i val la pena tenir-la ben present. Si els valors d'una
recursió \emph{viuen tots dins un conjunt fix} ---elements d'un \(X\), parts
d'un \(A\), aplicacions d'un exponencial---, la recursió és de Nivell~I i la
dona el teorema~\ref{teo:recursio-transfinita} sense cap axioma nou. Si els
valors són \emph{conjunts nous sense codomini comú}, cal SR (Nivell~II). En
els capítols que vénen, el pas d'un bon ordre al seu representant ordinal
escollit, la successió \(\alpha\mapsto\aleph_\alpha\) i les jerarquies
iterades com la dels \(V_n\) cauran del costat de Nivell~II; en canvi, tota la
teoria dels tipus ordinals \emph{acotats} ---els que viuen sobre parts d'un
conjunt fix--- es podrà fer amb la recursió d'aquest capítol. Mantindrem la
disciplina de la caixa d'axiomes: marcarem \enquote{SR} exactament allà on el
codomini deixi de ser fix, i enlloc més.
\end{observacio}

\begin{resumcap}
\apartat{Idees centrals}
\begin{enumerate}
  \item Inducció transfinita: en un bon ordre, una part progressiva (\(W_{<w}\subseteq S\Rightarrow w\in S\)) és tot \(W\). La prova és el mínim del complement; el cas del mínim de \(W\) queda inclòs.
  \item Recursió amb codomini fix: una regla \(\Phi(w,h)\) que llegeix la història determina una única \(F\colon W\to X\). Es construeix dins \((X+\mathbf1)^W\): les solucions parcials són compatibles (inducció), n'hi ha fins a cada punt (inducció) i es reuneixen en una (graf funcional i total). Cap elecció.
  \item Comparabilitat total dels bons ordres: l'aparellament \(F(w)=\min(V\setminus F(W_{<w}))\), amb un punt \(\top\) quan \(V\) s'esgota, és creixent amb imatge inicial, i els tres casos surten d'on s'atura. Sense elecció: el mínim és la tria.
  \item Sumes iterades, successions de parts d'un fix i tancaments són recursions amb codomini fix. La separació successor/límit és una comoditat, legítima només si és exhaustiva.
  \item Frontera: quan els valors són conjunts nous sense codomini comú (la torre \(P^n(A)\)), els axiomes de Nivell~I no garanteixen la família, i cal SR.
\end{enumerate}
\apartat{Errors típics} Afegir un cas base a la inducció transfinita (ja hi és); creure que la recursió transfinita demana elecció (cada pas està determinat); confondre la comparabilitat de bons ordres (sense elecció) amb la de conjunts (amb elecció); usar la separació successor/límit sense haver provat que és exhaustiva; i afirmar que la torre \(P^n(A)\) no cap en cap conjunt, quan el que val és que cap axioma de Nivell~I no en dona un.
\apartat{Lectura recomanada} \cite[§§17--18]{halmos} per a la inducció i la recursió transfinites i la comparabilitat dels bons ordres en la teoria clàssica; \cite[cap.~2]{jech2003} i \cite[cap.~I.9]{kunen2011} per al teorema de recursió transfinita i el seu ús amb el reemplaçament, que és el que aquí es marca com a frontera; i \cite[§§3--4]{shulman-comparing} per a la recursió transfinita en teoria estructural de conjunts i el paper del reemplaçament de contextos.
\end{resumcap}

\chapter{Zorn, Hartogs i el teorema del bon ordre}\label{cap:zorn}

\begin{objectius}
\apartat{Prerequisits} Del capítol~\ref{cap:bons-ordres}: bons ordres, segments, rigidesa i comparació. Del capítol~\ref{cap:recursio}: inducció i recursió transfinites amb codomini fix, i la comparabilitat total. Del capítol~\ref{cap:relacions}: l'axioma d'elecció, els quocients i l'efectivitat. Del capítol~\ref{cap:families}: seccions, funcions d'elecció i unions indexades.
\apartat{Objectius} En acabar el capítol, el lector hauria de poder:
\begin{itemize}
  \item codificar un bon ordre sobre una part d'\(A\) com un parell \((S,R)\) dins \(P(A)\times P(A\times A)\), i explicar per què el suport ha de ser una dada;
  \item construir \(\Hart(A)\) com a quocient de \(\WO(A)\) per isomorfisme, demostrar que és un bon ordre i que el segment anterior a un tipus té aquell tipus;
  \item demostrar el teorema de Hartogs, sense elecció, i veure'l com l'anàleg ordinal del teorema de Cantor;
  \item demostrar el teorema del bon ordre amb una recursió \emph{acotada per \(\Hart(A)\)}, i explicar per què la cota evita la recursió sobre tots els ordinals;
  \item distingir maximal de màxim, demostrar el lema de Zorn amb la mateixa cota i aplicar-lo (existència de bases);
  \item demostrar el cercle d'equivalències entre elecció, bon ordre i Zorn, marcant quina direcció introdueix l'elecció.
\end{itemize}
\end{objectius}

Els dos capítols anteriors han donat la teoria dels bons ordres i la
recursió transfinita \emph{sense} usar l'elecció. Aquest capítol hi torna a
introduir l'elecció (axioma~\ref{ax:eleccio}) i en treu les tres grans
conseqüències que, en la teoria clàssica, apareixen encadenades: el
\emph{lema de Zorn}, el \emph{teorema del bon ordre} i, com a frontissa
tècnica entre tots dos, el \emph{teorema de Hartogs}. La novetat respecte del
tractament habitual és doble. Primer, el punt de vista estructural manté totes
les construccions \emph{acotades}: en lloc de parlar de \enquote{tots els
ordinals}, treballarem amb parells \((S,R)\), on \(S\) és una part d'un
conjunt fix \(A\) i \(R\) un bon ordre sobre \(S\); aquests parells formen una
part de \(P(A)\times P(A\times A)\), i els seus tipus d'ordre s'obtenen
per un quocient efectiu. Segon, la construcció de Hartogs ---un bon ordre
\(\Hart(A)\) que no s'injecta en \(A\)--- proporciona la \emph{cota} que atura
les recursions: gràcies a ella, el teorema del bon ordre no necessita cap
recursió \enquote{sobre tots els ordinals}, sinó una de sola, acotada per
\(\Hart(A)\). Al final del capítol tancarem el cercle clàssic mostrant que
elecció, teorema del bon ordre i lema de Zorn són equivalents.


\section{El teorema de Hartogs}\label{sec:hartogs}

Comencem per Hartogs, perquè no usa l'elecció i perquè és la peça que fa
acotades les construccions posteriors. La idea és reunir tots els
\emph{tipus} de bons ordres que es poden posar sobre parts d'\(A\), i veure
que aquesta col·lecció, ordenada, és un bon ordre massa gran per cabre dins
\(A\).

\begin{definicio}[Bons ordres sobre parts d'un conjunt fix]\index{Bons ordres sobre parts d'un conjunt fix}\label{def:bons-ordres-parts}
Sigui \(A\) un conjunt. Un \emph{bon ordre sobre una part d'\(A\)} és un
\emph{parell} \((S,R)\) format per un \emph{suport} \(S\hookrightarrow A\)
---un element de \(P(A)\)--- i una relació \(R\hookrightarrow S\times S\)
---un element de \(P(S\times S)\hookrightarrow P(A\times A)\)
(definició~\ref{def:relacio})--- tal que \(R\) és un bon ordre estricte
(definició~\ref{def:bon-ordre}) sobre \emph{tot} \(S\). Codifiquem el parell
com un element del conjunt fix
\[
(S,R)\ \in\ P(A)\times P(A\times A).
\]
La col·lecció d'aquests bons ordres és una part interna
\[
\WO(A)\ \hookrightarrow\ P(A)\times P(A\times A),
\]
tallada per les condicions internes \enquote{\(R\hookrightarrow S\times S\), i
\(R\) és irreflexiva, transitiva, total i ben fundada sobre \(S\)} ---totes
expressables amb igualtats i quantificadors acotats sobre
\(P(A)\times P(A\times A)\) (capítols~\ref{cap:relacions}
i~\ref{cap:parts-interna}). Escrivim \(\camp(S,R):=S\) per al suport, que ara
és una dada del parell i no una quantitat derivada.
\end{definicio}

\begin{observacio}[Per què el parell, i no la sola relació]\label{obs:parell-suport}
Codificar el bon ordre només per la relació \(R\) i recuperar-ne el suport
com \(\camp(R)=\{a: \exists b\,(a\mathrel R b\vee b\mathrel R a)\}\) falla en
els casos degenerats. Sobre un singletó \(\{a\}\), un bon ordre estricte és la
relació \emph{buida}, i aleshores \(\camp(R)=\varnothing\), de manera que la
relació buida sola no distingeix l'ordre buit sobre \(\varnothing\) dels
ordres sobre subconjunts d'un element. El parell \((S,R)\) fixa el suport
explícitament i elimina l'ambigüitat: el bon ordre buit sobre \(\{a\}\) queda
representat per \((\{a\},\varnothing)\), diferent de
\((\varnothing,\varnothing)\). Amb la codificació antiga, en canvi,
\(\Hart(A)\) d'un singletó \(A\) podria contenir un únic tipus que
\emph{sí} s'injecta en \(A\), contradient l'enunciat de Hartogs
(teorema~\ref{teo:hartogs}); el parell ho evita.
\end{observacio}

\begin{observacio}[Tot és acotat]\label{obs:hartogs-acotat}
Aquesta és la clau del tractament estructural de Hartogs. Un bon ordre sobre
una part d'\(A\) \emph{no} és un objecte extern: és un element del conjunt fix
\(P(A)\times P(A\times A)\), i \(\WO(A)\) n'és una part. No parlem de
\enquote{tots els bons ordres} en abstracte, sinó dels que caben sobre parts
d'un conjunt donat, que formen un conjunt ja construït. Tota la resta de la
secció ---la identificació per isomorfisme, el quocient, l'ordre--- passa dins
aquest ambient, sense cap classe global.
\end{observacio}

\begin{definicio}[Tipus d'ordre; el conjunt de Hartogs]\index{Tipus d'ordre; el conjunt de Hartogs}\label{def:hartogs}
Sobre \(\WO(A)\), la relació \enquote{\((S,R)\) i \((S',R')\) són isomorfs com
a bons ordres} ---és a dir, hi ha una bijecció \(S\to S'\) que transporta
\(R\) en \(R'\)--- és una relació d'equivalència interna \(\cong\) (reflexiva,
simètrica i transitiva, per composició d'isomorfismes d'ordre). Per
l'efectivitat de les relacions d'equivalència (axioma~\ref{ax:exactitud})
existeix el quocient
\[
\Hart(A)\ :=\ \WO(A)/\cong,
\]
el \emph{conjunt dels tipus d'ordre} dels bons ordres sobre parts d'\(A\). La
classe \([(S,R)]\in\Hart(A)\) és el \emph{tipus d'ordre} del bon ordre
\((S,R)\); quan el suport se sobreentén escrivim simplement \([R]\).
Ordenem \(\Hart(A)\) per la relació de segment inicial: per a dos tipus
\(\alpha=[(S,R)]\) i \(\beta=[(S',R')]\),
\[
\alpha<\beta
\quad:\Longleftrightarrow\quad
(S,R)\ \text{és isomorf a un segment inicial propi de}\ (S',R')
\]
(definició~\ref{def:mes-curt}), que no depèn dels representants:
component isomorfismes d'ordre, si \((S,R)\cong(S_1,R_1)\) i
\((S',R')\cong(S'_1,R'_1)\), un isomorfisme de \((S,R)\) sobre un segment
propi de \((S',R')\) dona un isomorfisme de \((S_1,R_1)\) sobre un segment
propi de \((S'_1,R'_1)\).
\end{definicio}

\begin{proposicio}[\(\Hart(A)\) és un bon ordre]\label{prop:hartogs-bon-ordre}
La relació \(<\) sobre \(\Hart(A)\) de la definició~\ref{def:hartogs} és un
bon ordre.\par
\emph{Axiomes usats:} conjunt de parts, relacions i quocients efectius
(capítols~\ref{cap:relacions} i~\ref{cap:parts-interna},
axioma~\ref{ax:exactitud}), comparabilitat de bons ordres
(teorema~\ref{teo:comparabilitat}); \emph{cap elecció}.
\end{proposicio}

\begin{prova}
Que \(<\) és un ordre estricte total és la comparació de bons ordres del
capítol~\ref{cap:recursio}: dos tipus qualssevol són comparables
(teorema~\ref{teo:comparabilitat}), i l'exclusivitat i la transitivitat són
la tricotomia i la proposició~\ref{prop:comparacio-ordres}; la
irreflexivitat és que cap bon ordre no és isomorf a un segment propi seu
(corol·lari~\ref{cor:no-segment-propi}). Falta el bon ordre: tota part no
inicial \(T\hookrightarrow\Hart(A)\) té mínim. Sigui
\(\alpha_0=[(S_0,R_0)]\in T\) qualsevol. Si \(\alpha_0\) ja és el mínim, fet.
Si no, els tipus de \(T\) estrictament menors que \(\alpha_0\) corresponen,
per definició de \(<\), a segments inicials propis de \((S_0,R_0)\); i els
segments inicials de \((S_0,R_0)\) estan en bijecció creixent amb els punts
del suport \(S_0\) (proposició~\ref{prop:parts-inicials}), que és un bon
ordre. Per tant \(\{\beta\in T:\beta<\alpha_0\}\) es correspon amb una part no
inicial de \(S_0\), que té mínim pel bon ordre \(R_0\); aquest mínim dona el
tipus mínim de \(T\). Així \(\Hart(A)\) és ben
ordenat.
\end{prova}

\begin{lema}[El segment anterior a un tipus té aquell tipus]\label{lema:segment-tipus}
Sigui \(\alpha=[(S,R)]\in\Hart(A)\). El segment inicial propi
\[
\Hart(A)_{<\alpha}\ :=\ \{\,\beta\in\Hart(A):\beta<\alpha\,\}
\]
amb l'ordre \(<\) heretat és un bon ordre de tipus exactament \(\alpha\); és a
dir, \((\Hart(A)_{<\alpha},<)\) és isomorf a \((S,R)\).\par
\emph{Axiomes usats:} els de la proposició~\ref{prop:hartogs-bon-ordre};
\emph{cap elecció}.
\end{lema}

\begin{prova}
Per definició de \(<\), un tipus \(\beta=[(S',R')]\) compleix \(\beta<\alpha\)
si i només si \((S',R')\) és isomorf a un segment inicial propi de \((S,R)\).
Ara bé, els segments inicials propis de \((S,R)\) estan en bijecció creixent
amb els punts de \(S\): a cada \(s\in S\) li correspon el segment
\(S_{<s}=\{t\in S: t\mathrel R s\}\), i tot segment inicial propi és d'aquesta
forma (proposició~\ref{prop:parts-inicials}). Dos punts diferents donen
segments no isomorfs, perquè cap bon ordre no és isomorf a un segment propi
seu (corol·lari~\ref{cor:no-segment-propi}). Així l'aplicació
\[
s\ \longmapsto\ [\,(S_{<s},\,R\!\restriction\!S_{<s})\,]
\]
és una bijecció de \(S\) sobre \(\Hart(A)_{<\alpha}\), i és creixent perquè
\(s\mathrel R s'\) equival a \(S_{<s}\) ser segment inicial propi de
\(S_{<s'}\). Per tant és un isomorfisme d'ordre \((S,R)\cong
\Hart(A)_{<\alpha}\), i \(\Hart(A)_{<\alpha}\) té tipus \(\alpha\).
\end{prova}

\begin{teorema}[Hartogs]\index{Hartogs, teorema de}\label{teo:hartogs}
Per a tot conjunt \(A\), el bon ordre \(\Hart(A)\) \emph{no} admet cap
monomorfisme \(\Hart(A)\hookrightarrow A\). És a dir, \(\Hart(A)\not\preceq A\):
hi ha un bon ordre estrictament \enquote{massa llarg} per a \(A\).\par
\emph{Axiomes usats:} els de la proposició~\ref{prop:hartogs-bon-ordre}; cap
axioma nou; \emph{cap elecció}.
\end{teorema}

\begin{prova}
Suposem, cap a una contradicció, un monomorfisme \(m\colon\Hart(A)\hookrightarrow
A\). Transportant l'ordre de \(\Hart(A)\) per \(m\), la imatge
\(S_{*}:=m(\Hart(A))\hookrightarrow A\) és una part d'\(A\) dotada d'un bon
ordre \(R_{*}\hookrightarrow S_{*}\times S_{*}\) isomorf a \(\Hart(A)\) ---és a
dir, \((S_{*},R_{*})\in\WO(A)\), i el seu tipus
\([(S_{*},R_{*})]\in\Hart(A)\) és el tipus del bon ordre \(\Hart(A)\) mateix.
Ara bé, per construcció, \(\Hart(A)\) ordena \emph{tots} els tipus de bons
ordres sobre parts d'\(A\), inclòs \([(S_{*},R_{*})]\); i pel
lema~\ref{lema:segment-tipus} el segment inicial de \(\Hart(A)\) determinat per
un tipus \(\alpha\) és isomorf a un bon ordre de tipus \(\alpha\). En
particular, \(\Hart(A)\) seria isomorf al segment inicial de \(\Hart(A)\)
determinat per \([(S_{*},R_{*})]\), un segment inicial \emph{propi} de si
mateix ---perquè \([(S_{*},R_{*})]\) és un element de \(\Hart(A)\). Això
contradiu que cap bon ordre no és isomorf a cap segment propi seu
(corol·lari~\ref{cor:no-segment-propi}). Per tant no existeix cap tal \(m\), i
\(\Hart(A)\not\preceq A\).
\end{prova}

\begin{observacio}[Hartogs com a paral·lel de Cantor]\label{obs:hartogs-cantor}
El teorema de Hartogs és, per als \emph{ordres}, l'anàleg del que el teorema
de Cantor (teorema~\ref{teo:cantor}) és per a les \emph{mides}: així com
\(P(A)\) produeix una mida estrictament més gran que \(A\), \(\Hart(A)\)
produeix un bon ordre estrictament més llarg que qualsevol que càpiga dins
\(A\). Cap dels dos no usa elecció, i tots dos són \enquote{diagonals} en un
sentit ampli: la contradicció surt de fer actuar la construcció sobre si
mateixa. La diferència de funció és que \(P\) mesura cardinalitat i
\(\Hart\) mesura longitud d'ordre; per això \(\Hart(A)\) serà la cota natural
de les recursions que segueixen, on el que importa és la longitud.
\end{observacio}

\section{El teorema del bon ordre}\label{sec:teorema-bon-ordre}

Amb l'elecció i la cota de Hartogs podem provar que tot conjunt admet un bon
ordre. La idea és anar exhaurint \(A\) triant successivament un element del
que queda; el que garanteix que el procés \emph{acaba} ---i per tant que és
una recursió acotada, no una \enquote{sobre tots els ordinals}--- és
\(\Hart(A)\).

\begin{teorema}[Teorema del bon ordre (Zermelo)]\label{teo:bon-ordre}
Suposem l'axioma d'elecció (axioma~\ref{ax:eleccio}). Aleshores tot conjunt
\(A\) admet un bon ordre: existeix una relació \(R\hookrightarrow A\times A\)
que és un bon ordre sobre tot \(A\).\par
\emph{Axiomes usats:} elecció (axioma~\ref{ax:eleccio}), Hartogs
(teorema~\ref{teo:hartogs}), recursió transfinita amb codomini fix
(teorema~\ref{teo:recursio-transfinita}); cap recursió \enquote{sobre tots els
ordinals} ---la construcció és acotada per \(\Hart(A)\).
\end{teorema}

\begin{prova}
Per l'elecció (axioma~\ref{ax:eleccio}), la família de les parts no inicials
d'\(A\) admet una \emph{funció d'elecció}: una aplicació
\(c\colon P(A)\setminus\{\varnothing\}\to A\) amb \(c(S)\in S\) per a tota
part no inicial \(S\) ---una secció de la família corresponent
(definició~\ref{def:seccio-familia}, proposició~\ref{prop:epi-regular-eleccio}).
Afegim un punt exterior \(\top\), formant \(A^{+}=A+\mathbf1\), i definim per
recursió transfinita (teorema~\ref{teo:recursio-transfinita}) sobre el bon
ordre \(\Hart(A)\) una aplicació \(F\colon\Hart(A)\to A^{+}\) amb la regla:
donada la història \(F|_{\Hart(A)_{<\alpha}}\),
\[
F(\alpha)\ :=\
\begin{cases}
c\bigl(A\setminus F(\Hart(A)_{<\alpha})\bigr) & \text{si } A\setminus F(\Hart(A)_{<\alpha})\neq\varnothing,\\[1mm]
\top & \text{altrament,}
\end{cases}
\]
és a dir, a cada etapa es tria ---amb \(c\)--- un element d'\(A\) encara no
fet servir, i \(\top\) quan ja no en queda cap. El codomini \(A^{+}\) és
\emph{fix}, de manera que el teorema~\ref{teo:recursio-transfinita} s'aplica i
\(F\) existeix.

\emph{\(F\) s'atura.} Sobre la part \(H_0=\{\alpha:F(\alpha)\neq\top\}\), la
regla fa \(F\) \emph{injectiva} cap a \(A\) (encara no hi ha cap ordre fixat
sobre \(A\), de manera que no té sentit dir-la \enquote{creixent}): si
\(\beta<\alpha\) i \(F(\alpha)\neq\top\), aleshores
\(F(\alpha)\notin F(\Hart(A)_{<\alpha})\), en particular \(F(\alpha)\neq
F(\beta)\). Si \(F\) no prengués mai \(\top\),
\(F\colon\Hart(A)\hookrightarrow A\) seria un monomorfisme, contra el teorema
de Hartogs (teorema~\ref{teo:hartogs}). Per tant existeix algun \(\alpha\) amb
\(F(\alpha)=\top\); sigui \(\alpha_0\) el mínim
(bon ordre de \(\Hart(A)\), proposició~\ref{prop:hartogs-bon-ordre}).

\emph{El bon ordre sobre \(A\).} Per la minimalitat de \(\alpha_0\),
\(F(\alpha)\neq\top\) per a tot \(\alpha<\alpha_0\); i per a
\(\alpha>\alpha_0\) també \(F(\alpha)=\top\), perquè
\(F(\Hart(A)_{<\alpha})\supseteq F(\Hart(A)_{<\alpha_0})=A\) i el complement
és buit. Així \(F(\alpha)=\top\) exactament per a \(\alpha\ge\alpha_0\); i
\(F(\alpha_0)=\top\) vol dir \(A\setminus F(\Hart(A)_{<\alpha_0})=\varnothing\),
és a dir \(F(\Hart(A)_{<\alpha_0})=A\). Així la restricció
\[
F|_{\Hart(A)_{<\alpha_0}}\colon\ \Hart(A)_{<\alpha_0}\ \longrightarrow\ A
\]
és injectiva i exhaustiva: una bijecció entre el segment inicial
\(\Hart(A)_{<\alpha_0}\) ---un bon ordre (proposició~\ref{prop:segment-ben-ordenat})---
i \(A\) (encara sense ordre). Transportant l'ordre del segment per aquesta
bijecció s'obté un bon ordre \(R\hookrightarrow A\times A\) sobre tot \(A\), i
respecte d'aquest ordre transportat la bijecció esdevé, ja sí, un isomorfisme
d'ordre creixent. El procés ha estat
\emph{acotat} per \(\Hart(A)\): no hem recorregut \enquote{tots els
ordinals}, sinó el bon ordre concret \(\Hart(A)\), i ens hem aturat en
\(\alpha_0<\) la seva longitud.
\end{prova}

\begin{observacio}[Per què cal la cota de Hartogs]\label{obs:cota-hartogs}
La construcció \enquote{anar triant un element del que queda} és la prova
clàssica de Zermelo, però tal com se sol enunciar amaga una recursió sobre
una col·lecció d'ordinals que no és, a priori, un conjunt. El tractament
estructural ho evita fent la recursió sobre el bon ordre \emph{concret}
\(\Hart(A)\), que sí que és un conjunt (definició~\ref{def:hartogs}): la
recursió transfinita amb codomini fix (teorema~\ref{teo:recursio-transfinita})
s'hi aplica directament, i el teorema de Hartogs garanteix que abans
d'esgotar \(\Hart(A)\) ja s'han esgotat els elements d'\(A\). Aquesta és la
raó tècnica per la qual Hartogs precedeix el teorema del bon ordre: és la cota
que converteix una recursió \enquote{potencialment sense fi} en una recursió
acotada de Nivell~I més elecció.
\end{observacio}

\section{El lema de Zorn}\label{sec:zorn}

La tercera forma de maximalitat és el lema de Zorn, que en la pràctica
matemàtica és la més usada. El demostrarem a partir del teorema del bon
ordre; primer, el vocabulari.

\begin{definicio}[Cadena, cota superior, maximal]\index{Cadena, cota superior, maximal}\label{def:zorn-nocions}
Sigui \((P,\le)\) un conjunt parcialment ordenat (un ordre reflexiu,
antisimètric i transitiu, capítol~\ref{cap:bons-ordres}).
\begin{itemize}
\item Una \emph{cadena} de \(P\) és una part \(C\hookrightarrow P\) totalment
ordenada per \(\le\): dos elements qualssevol de \(C\) són comparables.
\item Una \emph{cota superior} d'una part \(S\hookrightarrow P\) és un
element \(b\in P\) amb \(s\le b\) per a tot \(s\in S\).
\item Un element \(m\in P\) és \emph{maximal} si cap element no el supera
estrictament: \(m\le x\) implica \(m=x\).
\end{itemize}
\end{definicio}

\begin{observacio}[Maximal no és màxim]\label{obs:maximal-maxim}
Convé no confondre \emph{maximal} amb \emph{màxim}. Un \emph{màxim} és un
element \(\ge\) que \emph{tots} els altres; un \emph{maximal} és només un
element que \emph{cap} altre no supera. En un ordre total coincideixen, però
en un ordre parcial pot haver-hi diversos maximals i cap màxim: per exemple,
en les parts pròpies d'un conjunt de dos elements ordenades per inclusió, els
dos singletons són tots dos maximals, i no hi ha màxim. El lema de Zorn dona
un \emph{maximal}, no un màxim; afirmar que dona un màxim seria fals.
\end{observacio}

\begin{teorema}[Lema de Zorn]\index{Zorn, lema de}\label{teo:zorn}
Suposem l'axioma d'elecció (axioma~\ref{ax:eleccio}). Sigui \((P,\le)\) un
conjunt parcialment ordenat no inicial en què \emph{tota cadena té una cota
superior}. Aleshores \(P\) té almenys un element maximal.\par
\emph{Axiomes usats:} elecció (axioma~\ref{ax:eleccio}), teorema del bon
ordre (teorema~\ref{teo:bon-ordre}) i recursió transfinita
(teorema~\ref{teo:recursio-transfinita}); tot acotat, cap recursió sobre una
classe.
\end{teorema}

\begin{prova}
Suposem, cap a una contradicció, que \(P\) no té cap maximal: per a tot
\(x\in P\) existeix \(x'\) amb \(x<x'\). Aleshores tota cadena \(C\) té, a
més d'una cota superior, una cota superior \emph{estricta} ---un element
estrictament més gran que tots els de \(C\)---: prenent una cota superior
\(b\) de \(C\) i un \(b'>b\) (que existeix perquè \(b\) no és maximal),
\(b'\) supera estrictament tot \(C\).

Considerem la família indexada per les parts de \(P\) que té, sobre cada
part \(C\), la part de \(P\) de les seves cotes superiors estrictes
(una part de \(P(P)\times P\), tallada amb quantificadors interns). Les
fibres sobre les cadenes són no buides pel paràgraf anterior; restringint
la família a la part de les cadenes i aplicant l'elecció
(axioma~\ref{ax:eleccio}) obtenim una aplicació \(c\) que a cada cadena
\(C\) assigna una cota superior estricta \(c(C)\notin C\). Fixem també un
punt qualsevol \(p_0\in P\) (\(P\) no és inicial). Definim per recursió
transfinita (teorema~\ref{teo:recursio-transfinita}) sobre \(\Hart(P)\) una
aplicació \(F\colon\Hart(P)\to P\) amb la regla \emph{total}
\[
F(\alpha):=
\begin{cases}
c\bigl(F(\Hart(P)_{<\alpha})\bigr) & \text{si } F(\Hart(P)_{<\alpha})\text{ és una cadena},\\[1mm]
p_0 & \text{altrament},
\end{cases}
\]
on la condició \enquote{ser una cadena} és interna (quantificadors acotats
sobre \(P\times P\)); la segona clàusula no s'arribarà a usar mai, però fa
que la regla estigui definida per a tota història, com el
teorema~\ref{teo:recursio-transfinita} exigeix. Per inducció transfinita
(teorema~\ref{teo:induccio-transfinita}) es comprova alhora que la imatge
\(F(\Hart(P)_{<\alpha})\) és sempre una cadena i que \(F\) és estrictament
creixent: si ho és per a tot \(\beta<\alpha\), els valors \(F(\beta)\)
estan totalment ordenats, de manera que la imatge és una cadena, el primer
cas de la regla s'aplica, i \(F(\alpha)=c(\cdot)\) supera estrictament tots
els \(F(\beta)\) amb \(\beta<\alpha\). Per tant \(F\colon\Hart(P)\to P\) és un
monomorfisme (una aplicació estrictament creixent des d'un ordre total és
injectiva), és a dir \(\Hart(P)\preceq P\). Però això contradiu el teorema de
Hartogs (teorema~\ref{teo:hartogs}), \(\Hart(P)\not\preceq P\). La
contradicció ve de suposar que \(P\) no té maximal; per tant \(P\) en té un.
\end{prova}

\begin{observacio}[La cota de Hartogs, altre cop]\label{obs:zorn-hartogs}
La prova de Zorn és, estructuralment, la mateixa que la del teorema del bon
ordre: una recursió que tria a cada pas un element \enquote{nou} (aquí, una
cota estricta) i que \emph{no pot continuar indefinidament} perquè
\(\Hart(P)\) no s'injecta en \(P\). És el teorema de Hartogs, i no cap
recursió sobre una classe d'ordinals, el que atura el procés i dona la
contradicció. Aquesta és la manera estructural, i acotada, de fer el que en
la prova clàssica es fa \enquote{iterant fins que ja no es pot}.
\end{observacio}

\subsection{Una aplicació: tot espai vectorial té una base}\label{sub:base}

Perquè Zorn no quedi com un enunciat purament fundacional, en donem una
aplicació matemàtica genuïna: l'existència de bases. Es formula amb el
vocabulari d'aquest llibre prenent un espai vectorial com un conjunt amb les
operacions internes corresponents; el que importa aquí és l'estructura de
l'argument de Zorn.

\begin{proposicio}[Existència de base]\label{prop:base}
Suposem l'elecció. Tot espai vectorial \(V\) sobre un cos \(K\) té una base:
una part \(B\hookrightarrow V\) linealment independent i generadora.\par
\emph{Axiomes usats:} lema de Zorn (teorema~\ref{teo:zorn}); per tant elecció.
\end{proposicio}

\begin{prova}
Sigui \(P\) el conjunt de les parts linealment independents de \(V\),
ordenat per inclusió ---una part de \(P(V)\) tallada per la condició interna
d'independència lineal (cap combinació lineal finita no trivial no s'anul·la),
que és un conjunt parcialment ordenat no inicial (la part buida hi és, o
qualsevol singletó no nul). Vegem que tota cadena \(C\hookrightarrow P\) té
cota superior: la unió \(U=\bigcup C\) (unió indexada de parts,
capítol~\ref{cap:families}) és linealment independent, perquè tota
combinació lineal involucra només un nombre \emph{finit} de vectors, que per
ser \(C\) una cadena viuen tots en un mateix membre de la cadena, on són
independents; i \(U\) conté tots els membres de \(C\), és a dir n'és cota
superior. Pel lema de Zorn (teorema~\ref{teo:zorn}), \(P\) té un element
maximal \(B\). Aquesta \(B\) és una base: és independent per ser a \(P\), i
és generadora perquè si algun \(v\in V\) no fos combinació lineal de \(B\),
aleshores \(B\cup\{v\}\) seria independent i estrictament més gran, contra la
maximalitat de \(B\).
\end{prova}

\begin{observacio}[Altres aplicacions]\label{obs:zorn-aplicacions}
El mateix esquema ---ordenar per inclusió, veure que les unions de cadenes
resten dins la classe, i prendre un maximal--- dona molts resultats
d'existència no constructius: l'extensió d'un filtre a un ultrafiltre, la
clausura algebraica d'un cos, l'existència d'ideals maximals en un anell amb
unitat, o el teorema de Hahn--Banach en la seva forma algebraica. En tots,
Zorn substitueix una tria \enquote{fins a l'infinit} per un maximal; i en
tots, la propietat de tancament per unions de cadenes és el que cal
comprovar.
\end{observacio}

\section{Les tres equivalències}\label{sec:equivalencies}

Tanquem el capítol amb el cercle clàssic: sobre la base dels altres axiomes
del llibre, les tres afirmacions són equivalents.

\begin{teorema}[Elecció, bon ordre i Zorn són equivalents]\label{teo:equivalencies}
Sobre els axiomes del llibre \emph{sense} l'elecció, les tres afirmacions
següents són equivalents:
\[
\begin{aligned}
\text{(E) elecció}
\quad&\Longleftrightarrow\quad
\text{(BO) tot conjunt admet un bon ordre}\\
&\Longleftrightarrow\quad
\text{(Z) lema de Zorn.}
\end{aligned}
\]
\emph{Axiomes usats:} recursió transfinita i Hartogs
(capítols~\ref{cap:recursio} i aquest), que no depenen de l'elecció; les
implicacions que \emph{introdueixen} l'elecció es marquen a la prova.
\end{teorema}

\begin{prova}
\emph{(E) \(\Rightarrow\) (BO)} és el teorema del bon ordre
(teorema~\ref{teo:bon-ordre}); \emph{(E) \(\Rightarrow\) (Z)} és el lema de
Zorn (teorema~\ref{teo:zorn}). Falten dues direccions que \emph{retornen}
l'elecció.

\emph{(BO) \(\Rightarrow\) (E).} Sigui \(p\colon E\twoheadrightarrow I\) una
família de parts no inicials (un epimorfisme regular). Per (BO), fixem un bon
ordre sobre \(E\). Aleshores cada fibra \(E_i=p^{-1}(i)\) és una part no
inicial d'\(E\), i té un \emph{mínim} per al bon ordre
(definició~\ref{def:bon-ordre}); l'aplicació \(s\colon I\to E\) que a cada
\(i\) assigna \(\min(E_i)\) és una secció de \(p\) ---està ben definida
perquè el mínim és únic i canònic, sense cap tria addicional. Per tant tota
família de parts no buides té secció, que és l'elecció (axioma~\ref{ax:eleccio}).

\emph{(Z) \(\Rightarrow\) (BO)}, i per tant \((Z)\Rightarrow(E)\). Sigui \(A\)
un conjunt; volem ben ordenar-lo. Considerem el conjunt \(P\) dels
\emph{bons ordres parcials} d'\(A\): els parells \((S,R)\) amb
\(S\hookrightarrow A\) i \(R\) un bon ordre sobre \(S\) ---una part de
\(P(A)\times P(A\times A)\). Ordenem \(P\) per \emph{extensió inicial}:
\((S,R)\le(S',R')\) si \(S\hookrightarrow S'\), \(R\) és la restricció de
\(R'\) a \(S\), i \(S\) és un \emph{segment inicial} de \((S',R')\). Tota cadena \(C\) té cota superior: siguin
\(S_{C}=\bigcup\{S:(S,R)\in C\}\) i \(R_{C}=\bigcup\{R:(S,R)\in C\}\) les
unions indexades (capítol~\ref{cap:families}). \(R_{C}\) és un ordre
estricte total sobre \(S_{C}\), perquè dos punts qualssevol són en un mateix
membre de la cadena (que està totalment ordenada per extensió inicial) i
allà són comparables; i és un bon ordre: si \(T\hookrightarrow S_{C}\) és
una part no buida, prenem un punt \(t\in T\) i un membre \((S,R)\in C\) amb
\(t\in S\); el mínim de \(T\cap S\) per \(R\) és el mínim de tot \(T\),
perquè qualsevol altre punt de \(T\) o bé és a \(S\), o bé és en un membre
que estén \((S,R)\) \emph{inicialment}, i aleshores és \(\ge\) tots els
punts de \(S\). Finalment \((S_{C},R_{C})\) estén inicialment cada membre
de \(C\), de manera que n'és cota superior. (La cadena buida té per cota
\((\varnothing,\varnothing)\in P\).) Pel lema de Zorn
(teorema~\ref{teo:zorn}), \(P\) té un maximal \((S_{*},R_{*})\). Ha de ser
\(S_{*}=A\): si quedés un element \(a\in A\setminus S_{*}\), afegint-lo al
final de \(R_{*}\) obtindríem un bon ordre parcial estrictament més gran,
contra la maximalitat. Per tant \(R_{*}\) és un bon ordre sobre tot \(A\), i
val (BO).

Les tres afirmacions són, doncs, equivalents. Cap direcció no ha exigit una
recursió més forta que la del capítol~\ref{cap:recursio}: la construcció de
\(\Hart\) i les recursions són totes acotades i de codomini fix.
\end{prova}

\begin{resumcap}
\apartat{Idees centrals}
\begin{enumerate}
  \item Un bon ordre sobre una part d'\(A\) és un parell \((S,R)\) dins \(P(A)\times P(A\times A)\); el suport ha de ser una dada perquè la relació buida no el determina. \(\WO(A)\) és una part interna i \(\Hart(A)=\WO(A)/\!\cong\) existeix per l'efectivitat.
  \item \(\Hart(A)\), ordenat per segment inicial, és un bon ordre, i el segment anterior a un tipus \(\alpha\) té tipus \(\alpha\). Teorema de Hartogs: \(\Hart(A)\not\preceq A\), sense elecció; és l'anàleg ordinal de Cantor.
  \item Teorema del bon ordre: amb una funció d'elecció sobre les parts no buides d'\(A\) i una recursió \emph{acotada per \(\Hart(A)\)}, s'exhaureix \(A\) abans d'esgotar \(\Hart(A)\), i la bijecció resultant transporta el bon ordre del segment.
  \item Lema de Zorn amb la mateixa cota: si no hi hagués maximal, una recursió que tria cotes estrictes donaria \(\Hart(P)\preceq P\), contra Hartogs. Zorn dona un maximal, no un màxim.
  \item Elecció, bon ordre i Zorn són equivalents; (BO)\(\Rightarrow\)(E) pel mínim de cada fibra, (Z)\(\Rightarrow\)(BO) per maximalitat en els bons ordres parcials.
\end{enumerate}
\apartat{Errors típics} Codificar un bon ordre només per la relació, perdent el suport; parlar de \enquote{tots els ordinals} en lloc de \(\Hart(A)\); donar una regla de recursió definida només en el cas bo (ha de ser total); confondre maximal amb màxim, o cota superior amb cota estricta; i oblidar que la unió d'una cadena de bons ordres parcials és ben ordenada \emph{perquè} les extensions són inicials.
\apartat{Lectura recomanada} \cite[§§15--17]{halmos} per a l'elecció, Zorn i el teorema del bon ordre en la teoria clàssica; \cite[cap.~5]{jech2003} per al teorema de Hartogs, l'aleph de Hartogs i les equivalències de l'elecció; \cite[cap.~I.12]{kunen2011} per a la presentació axiomàtica d'aquests resultats; i \cite[apèndix~B]{lawvere-rosebrugh} per a la formulació estructural de l'elecció i els seus equivalents.
\end{resumcap}

\chapter{Tipus ordinals}\label{cap:ordinals}

\begin{objectius}
\apartat{Prerequisits} Del capítol~\ref{cap:bons-ordres}: bons ordres, segments, rigidesa, comparació per segments i tricotomia. Del capítol~\ref{cap:recursio}: la comparabilitat total. Del capítol~\ref{cap:zorn}: el conjunt de Hartogs \(\Hart(A)\) i el fet que és un bon ordre. Del capítol~\ref{cap:families}: unions indexades de parts. Del capítol~\ref{cap:cantor}: el teorema de Cantor, com a model de Burali--Forti.
\apartat{Objectius} En acabar el capítol, el lector hauria de poder:
\begin{itemize}
  \item definir un tipus ordinal com un bon ordre llevat d'isomorfisme, sense representant canònic, i comprovar la independència del representant de cada definició;
  \item demostrar que els tipus estan totalment ordenats per segments, i ben ordenats dins \(\Hart(A)\);
  \item explicar per què \(\alpha\in\beta\) no és una notació adequada en aquest llibre;
  \item definir zero, successor i límit, i demostrar que són exhaustius, excloents i invariants;
  \item construir el suprem d'una família representada dins un bon ordre com el tipus de la reunió dels segments oberts, i saber per què cal l'acotació;
  \item demostrar el teorema de Burali--Forti i llegir-lo com el paral·lel ordinal del teorema de Cantor.
\end{itemize}
\end{objectius}

Els bons ordres mesuren la \emph{forma} d'un procés que es recorre sense
saltar-se cap pas; però dos bons ordres de la mateixa forma ---isomorfs com a
ordres--- no s'han de distingir. El \emph{tipus ordinal} és, precisament, un
bon ordre considerat llevat d'isomorfisme d'ordre: el que queda quan
s'oblida la naturalesa concreta dels elements i es reté només la manera com
estan ordenats. Aquest capítol desenvolupa els tipus ordinals amb la
disciplina estructural del llibre: es defineixen les nocions
\emph{sobre representants} ---bons ordres concrets--- i després es demostra que
no depenen del representant, gràcies a la unicitat de l'isomorfisme d'ordre
(teorema~\ref{teo:unicitat-iso}). L'ordre entre tipus és la comparació per
segments inicials que ja coneixem; el zero, el successor i els punts límit
s'obtenen de construccions elementals sobre bons ordres; i els suprems es
formen \emph{acotadament}, dins un bon ordre donat, sense parlar mai del
\enquote{conjunt de tots els ordinals} ---que, com veurem amb el teorema de
Burali--Forti, no existeix. El paral·lelisme amb el capítol de Cantor és
deliberat: així com no hi ha una mida màxima, tampoc no hi ha una longitud
d'ordre màxima.

\section{Tipus ordinals i el seu ordre}\label{sec:tipus-ordinals}

\begin{definicio}[Tipus ordinal]\index{Tipus ordinal}\label{def:tipus-ordinal}
Un \emph{tipus ordinal} és un conjunt ben ordenat considerat llevat
d'isomorfisme d'ordre (definició~\ref{def:creixent}). Diem que un bon ordre
\((W,<)\) \emph{representa} el tipus \(\alpha\), i escrivim
\(\alpha=\operatorname{tipus}(W)\), quan \(\alpha\) és la seva classe
d'isomorfisme d'ordre. Dos bons ordres representen el mateix tipus si i només
si són isomorfs com a ordres. \emph{No} escollim cap representant canònic:
treballem amb bons ordres qualssevol i comprovem, cada vegada, que el que
definim no depèn del representant triat.
\end{definicio}

\begin{observacio}[Tipus acotats i teoria global]\label{obs:tipus-acotats}
Com a la comparació de conjunts (observació~\ref{obs:comparacio-metateoria}),
\enquote{tipus ordinal} és, en general, una noció metateòrica: parla de bons
ordres llevat d'isomorfisme, no d'elements d'un conjunt fix. Ara bé, per a un
conjunt \(A\) donat, els tipus dels bons ordres sobre parts d'\(A\) sí que
formen un conjunt ---és el conjunt de Hartogs \(\Hart(A)\)
(definició~\ref{def:hartogs})---, i tota la teoria \emph{acotada} d'aquest
capítol es pot llegir dins un \(\Hart(A)\) prou gran. La distinció importa: els
enunciats sobre \enquote{un tipus ordinal} són de Nivell~I quan es poden
situar dins un \(\Hart(A)\); una teoria \emph{global} de tots els tipus
ordinals alhora seria de Nivell~II i requeriria SR
(observació~\ref{obs:sr}). En aquest capítol ens mantenim, tret d'on
s'indiqui, en el terreny acotat.
\end{observacio}

L'ordre entre tipus és la comparació per segments inicials del
capítol~\ref{cap:bons-ordres}, ara llegida sobre tipus.

\begin{definicio}[Ordre dels tipus ordinals]\index{Ordre dels tipus ordinals}\label{def:ordre-ordinals}
Siguin \(\alpha=\operatorname{tipus}(W)\) i \(\beta=\operatorname{tipus}(V)\).
Definim
\[
\alpha<\beta
\quad\Longleftrightarrow\quad
W\ \text{és isomorf a un segment inicial propi de}\ V
\]
(definició~\ref{def:mes-curt}), i \(\alpha\le\beta\) si \(\alpha<\beta\) o
\(\alpha=\beta\). La unicitat de l'isomorfisme d'ordre
(teorema~\ref{teo:unicitat-iso}) garanteix que la definició \emph{no depèn dels
representants}: si \(W\cong W'\) i \(V\cong V'\), aleshores \(W\) és isomorf a
un segment propi de \(V\) si i només si \(W'\) ho és de \(V'\).
\end{definicio}

\begin{proposicio}[Els tipus ordinals estan totalment ordenats]\label{prop:ordinals-total}
La relació \(<\) sobre tipus ordinals és un ordre estricte total: per a
tipus \(\alpha,\beta\) qualssevol es dona exactament una de \(\alpha<\beta\),
\(\alpha=\beta\), \(\beta<\alpha\). A més, restringida als tipus dels bons
ordres sobre parts d'un conjunt fix \(A\) ---és a dir, sobre \(\Hart(A)\)---,
és un bon ordre.\par
\emph{Axiomes usats:} comparabilitat i tricotomia de bons ordres
(teoremes~\ref{teo:comparabilitat} i~\ref{teo:tricotomia}), Hartogs
(proposició~\ref{prop:hartogs-bon-ordre}); cap axioma nou; \emph{cap elecció}.
\end{proposicio}

\begin{prova}
La tricotomia i la comparabilitat dels bons ordres
(teoremes~\ref{teo:tricotomia} i~\ref{teo:comparabilitat}) diuen exactament
que, per a dos representants \(W\) i \(V\), es dona una i una sola de: \(W\)
isomorf a un segment propi de \(V\) (és a dir \(\alpha<\beta\)), \(W\cong V\)
(és a dir \(\alpha=\beta\)), o \(V\) isomorf a un segment propi de \(W\) (és a
dir \(\beta<\alpha\)). La transitivitat és la de la comparació per segments
(proposició~\ref{prop:comparacio-ordres}), i la irreflexivitat és que cap bon
ordre no és isomorf a un segment propi seu
(corol·lari~\ref{cor:no-segment-propi}). Que sobre \(\Hart(A)\) és un bon
ordre és la proposició~\ref{prop:hartogs-bon-ordre}.
\end{prova}

\begin{observacio}[No s'ha d'escriure \(\alpha\in\beta\)]\label{obs:no-pertinença-ordinal}
Hi ha una diferència important amb la teoria clàssica de von Neumann, on un
ordinal \emph{és} el conjunt dels ordinals menors i \enquote{\(\alpha<\beta\)}
es defineix com \enquote{\(\alpha\in\beta\)}. En aquest llibre la pertinença
és sempre \emph{relativa a una part d'un conjunt fix}
(capítol~\ref{cap:parts-interna}), i no hi ha cap relació de pertinença global
entre conjunts; per tant \(\alpha\in\beta\) \emph{no} és una notació adequada
per a l'ordre ordinal. La relació correcta és \(\alpha<\beta\), definida
mitjançant segments inicials (definició~\ref{def:ordre-ordinals}). Aquesta no
és una qüestió de gust: escriure \(\alpha\in\beta\) suggeriria una jerarquia
material de pertinença que el llibre no assumeix, i que precisament el punt de
vista estructural vol evitar.
\end{observacio}

\section{Zero, successor i límit}\label{sec:zero-successor-limit}

Les tres classes bàsiques de tipus ordinals es defineixen sobre
representants, amb construccions elementals de bons ordres, i després es
comprova la invariància.

\begin{definicio}[Zero i successor]\index{Zero i successor}\label{def:zero-successor}
El \emph{tipus zero}, \(0\), és el tipus del bon ordre \emph{inicial}
\(\mathbf0\) (l'ordre buit): el menor de tots els tipus, perquè \(\mathbf0\) és
isomorf a un segment inicial de qualsevol bon ordre. Donat un bon ordre
\((W,<)\), el seu \emph{successor} és el tipus del bon ordre \(W+\mathbf1\)
---la suma disjunta (definició~\ref{def:suma-disjunta}) amb un punt nou
\(\ast\) col·locat \emph{després} de tots els elements de \(W\): l'ordre en
què \(u<\ast\) per a tot \(u\in W\) i la restricció a \(W\) és l'original. Si
\(\alpha=\operatorname{tipus}(W)\), el successor es denota \(\alpha+1\).
\end{definicio}

\begin{proposicio}[Invariància i caracterització del successor]\label{prop:successor-invariant}
El successor no depèn del representant: si \(W\cong V\), aleshores
\(W+\mathbf1\cong V+\mathbf1\). A més, per a tot bon ordre \(W\):
\begin{enumerate}
\item \(\alpha<\alpha+1\), i \(\alpha+1\) és el \emph{menor} tipus
estrictament més gran que \(\alpha\): no hi ha cap tipus \(\gamma\) amb
\(\alpha<\gamma<\alpha+1\);
\item \(\alpha+1\) té un element \emph{màxim} (el punt nou \(\ast\)), i el
segment inicial que aquest determina és, precisament, \(\alpha\).
\end{enumerate}
\emph{Axiomes usats:} suma disjunta (capítol~\ref{cap:valors}), unicitat de
l'isomorfisme d'ordre (teorema~\ref{teo:unicitat-iso}); cap elecció.
\end{proposicio}

\begin{prova}
\emph{Invariància.} Un isomorfisme d'ordre \(f\colon W\cong V\) s'estén a
\(f+\id\colon W+\mathbf1\cong V+\mathbf1\) enviant \(\ast\) a \(\ast\); conserva
l'ordre perquè els punts nous ocupen la mateixa posició final. \emph{(1)} El
segment inicial de \(W+\mathbf1\) determinat pel punt \(\ast\) és exactament
\(W\), de manera que \(W\) és isomorf a un segment propi de \(W+\mathbf1\), és a
dir \(\alpha<\alpha+1\). Si hi hagués \(\gamma\) amb \(\alpha<\gamma<\alpha+1\),
un representant \(G\) de \(\gamma\) seria isomorf a un segment propi de
\(W+\mathbf1\); els segments propis de \(W+\mathbf1\) són \(W\) mateix i els
segments de \(W\) (proposició~\ref{prop:parts-inicials}), de manera que
\(\gamma\le\alpha\), contra \(\alpha<\gamma\). \emph{(2)} El punt \(\ast\) és
màxim per construcció, i el seu segment inicial és \(W\), que representa
\(\alpha\) (proposició~\ref{prop:segment-determina-punt}).
\end{prova}

\begin{definicio}[Tipus límit]\index{Tipus límit}\label{def:tipus-limit}
Un tipus ordinal \(\lambda\) és \emph{successor} si \(\lambda=\alpha+1\) per a
algun \(\alpha\); equivalentment, si un (tot) representant seu té un element
màxim. És \emph{límit} si no és zero ni successor: és a dir, si és no inicial i
cap representant seu no té element màxim ---per a tot element hi ha un de
posterior.
\end{definicio}

\begin{proposicio}[Zero, successor i límit són exhaustius i invariants]\label{prop:tricotomia-successor}
Tot tipus ordinal és exactament un de: zero, successor o límit; i la classe a
què pertany no depèn del representant.\par
\emph{Axiomes usats:} bon ordre (definició~\ref{def:bon-ordre}), unicitat de
l'isomorfisme d'ordre (teorema~\ref{teo:unicitat-iso}); cap elecció.
\end{proposicio}

\begin{prova}
Sigui \(W\) un representant de \(\lambda\). Si \(W\) és inicial, \(\lambda=0\).
Si no, \(W\) té un element mínim però pot tenir o no un element màxim. Tenir
element màxim és una propietat invariant per isomorfisme d'ordre (un
isomorfisme envia màxim a màxim), de manera que no depèn del representant. Si
\(W\) té màxim \(m\), aleshores \(W\cong W_{<m}+\mathbf1\) ---el segment
\(W_{<m}\) seguit de \(m\)---, de manera que \(\lambda=\operatorname{tipus}(W_{<m})+1\)
és successor. Si \(W\) no té màxim i no és inicial, \(\lambda\) no és zero (no
inicial) ni successor (un successor té màxim, proposició~\ref{prop:successor-invariant}),
i per tant és límit. Les tres classes són mútuament excloents ---zero és
inicial, successor té màxim, límit no--- i exhaustives.
\end{prova}

\begin{observacio}[Lectura clàssica; el primer límit]\label{obs:limit-classic}
En la notació habitual, \(0,1,2,\dots\) són els tipus dels bons ordres finits
(els \(\seg{n}\), capítol~\ref{cap:finits}), tots zero o successors; el primer
tipus \emph{límit} és \(\omega=\operatorname{tipus}(\N)\), el tipus del bon
ordre dels naturals (teorema~\ref{teo:bon-ordre-N}), que no té element màxim.
Després venen \(\omega+1,\omega+2,\dots\), successors, i el següent límit
\(\omega+\omega=\omega\cdot2\), i així successivament. La distinció
successor/límit és la que governa les definicions per recursió transfinita:
als successors es fa un pas, als límits es pren un suprem
(capítol~\ref{cap:recursio}). És per als límits, on no hi ha \enquote{element
anterior}, que la recursió ha de mirar tota la història i no només el terme
previ.
\end{observacio}

\section{Suprems acotats}\label{sec:suprems}

Una família de tipus ordinals, si està \emph{representada dins un bon ordre}
comú, té suprem; i el suprem es forma acotadament, sense sortir d'aquell bon
ordre. Aquesta és la manera estructural de reunir tipus sense invocar el
\enquote{conjunt de tots els ordinals}.

\begin{definicio}[Família de tipus representada dins un bon ordre]\index{Família de tipus representada dins un bon ordre}\label{def:familia-tipus}
Una família de tipus ordinals \((\alpha_i)_{i\in I}\) està \emph{representada
dins} un bon ordre \(W\) si existeix una aplicació \(r\colon I\to W\) tal que,
per a cada \(i\), el tipus \(\alpha_i\) és el del segment inicial
\(W_{<r(i)}\). És a dir, cada \(\alpha_i\) apareix com el tipus d'un segment de
\(W\), i la família és la imatge d'una aplicació al conjunt d'índexs
(capítol~\ref{cap:families}).
\end{definicio}

\begin{proposicio}[Suprem acotat]\label{prop:suprem}
Sigui \((\alpha_i)_{i\in I}\) una família de tipus representada dins un bon
ordre \(W\) per \(r\colon I\to W\), és a dir, amb
\(\alpha_i=\operatorname{tipus}(W_{<r(i)})\). Aleshores existeix el
\emph{suprem} \(\sup_i\alpha_i\): el menor tipus ordinal \(\ge\) que tots els
\(\alpha_i\). Concretament, és el tipus de la reunió
\[
S\;=\;\bigcup_{i\in I}W_{<r(i)}
\]
dels segments inicials \emph{oberts}, que és una part inicial de \(W\). Si
\(I\) és inicial, \(S\) és buida i el suprem és \(0\).\par
\emph{Axiomes usats:} unions indexades de parts (capítol~\ref{cap:families}),
bon ordre i segments (capítol~\ref{cap:bons-ordres}); cap axioma nou; \emph{cap
elecció}.
\end{proposicio}

\begin{prova}
Sigui \(S=\bigcup_{i\in I}W_{<r(i)}\hookrightarrow W\) la unió indexada
(capítol~\ref{cap:families}). Com a unió de parts inicials, \(S\) és una part
inicial de \(W\), i per tant un segment \(W_{<w_0}\) o bé tot \(W\)
(proposició~\ref{prop:parts-inicials}); en qualsevol cas és ben ordenat
(proposició~\ref{prop:segment-ben-ordenat}). Sigui
\(\sigma=\operatorname{tipus}(S)\).

\emph{\(\sigma\) és cota superior.} Cada \(W_{<r(i)}\) és una part inicial de
\(S\) (és inicial a \(W\) i està continguda a \(S\)), de manera que és tot
\(S\) o un segment propi de \(S\); en tots dos casos \(\alpha_i\le\sigma\).

\emph{\(\sigma\) és la menor.} Sigui \(\beta\) amb \(\beta\ge\alpha_i\) per a
tot \(i\), amb representant \(V\). Per a cada \(i\) hi ha un únic
isomorfisme d'ordre \(f_i\) de \(W_{<r(i)}\) sobre una part inicial de \(V\)
(teorema~\ref{teo:unicitat-iso}). Aquests isomorfismes són compatibles: si
\(r(i)\le r(j)\), aleshores \(W_{<r(i)}\) és una part inicial de
\(W_{<r(j)}\), la restricció de \(f_j\) és un isomorfisme d'ordre de
\(W_{<r(i)}\) sobre una part inicial de \(V\) (la imatge d'una part inicial
per un isomorfisme sobre una part inicial és inicial), i per la unicitat
coincideix amb \(f_i\); i per la totalitat de \(W\), dos índexs sempre són
en un dels dos casos. La reunió dels grafs de les \(f_i\) és, doncs, una
relació funcional i total sobre \(S\) (proposició~\ref{prop:aplicacio-funcional}),
és a dir, una aplicació \(f\colon S\to V\) que és creixent, injectiva i té
imatge inicial, perquè cada punt de \(S\) i cada desigualtat entre dos punts
de \(S\) viuen dins un mateix \(W_{<r(j)}\), on \(f\) és \(f_j\). Així \(S\)
és isomorf a una part inicial de \(V\), i \(\sigma\le\beta\). Per tant
\(\sigma=\sup_i\alpha_i\).
\end{prova}

\begin{observacio}[Per què els segments oberts]\label{obs:suprem-oberts}
La reunió dels segments \emph{tancats} \(W_{\le r(i)}\) dona en general un
tipus massa gran: si \(I=\mathbf1\) i \(r\) és el mínim de \(W\), aleshores
\(\alpha=\operatorname{tipus}(W_{<r})=0\), mentre que \(W_{\le r}\) té un
element i tipus \(1\). El suprem d'una família amb màxim és aquest màxim, i
la fórmula amb segments oberts ho respecta: si \(r(i_0)\) és el màxim dels
\(r(i)\), aleshores \(S=W_{<r(i_0)}\) i \(\sigma=\alpha_{i_0}\).
\end{observacio}

\begin{observacio}[Per què el suprem ha de ser acotat]\label{obs:suprem-acotat}
La hipòtesi \enquote{representada dins un bon ordre \(W\)} és essencial i és el
que manté la construcció de Nivell~I. Reunir una família \emph{arbitrària} de
tipus ordinals ---sense un \(W\) comú que els contingui tots com a segments---
no es pot fer amb els axiomes actuals: caldria formar primer un bon ordre prou
llarg, i si els tipus no estan acotats cap axioma de Nivell~I no en garanteix
cap; i si la família és la de \emph{tots} els tipus, un tal bon ordre és
impossible, que és el contingut del teorema de Burali--Forti que ve tot
seguit. Per això, cada vegada que prenem un suprem, hem d'exhibir abans el bon
ordre \(W\) dins el qual la família viu ---sovint un \(\Hart(A)\) prou gran. És
la mateixa disciplina de la guia: no escriure \(\sup_i\alpha_i\) sense haver
construït primer la família acotada que es reuneix.
\end{observacio}

\section{El teorema de Burali--Forti}\label{sec:burali-forti}

Així com el teorema de Cantor prohibeix una mida màxima, el teorema de
Burali--Forti prohibeix un \enquote{conjunt de tots els tipus ordinals}: la
col·lecció de tots ells, si fos un conjunt ben ordenat, tindria un tipus
estrictament menor que ell mateix.

\begin{teorema}[Burali--Forti]\label{teo:burali-forti}
No existeix cap conjunt \(\Omega\) que contingui un representant de
\emph{cada} tipus ordinal i que estigui ben ordenat per la relació de segment
inicial de manera que el tipus del seu segment \(\Omega_{<\gamma}\) sigui,
per a cada \(\gamma\in\Omega\), el tipus representat per \(\gamma\).\par
\emph{Axiomes usats:} bon ordre i segments, no-isomorfisme amb segments propis
(corol·lari~\ref{cor:no-segment-propi}); cap axioma nou; \emph{cap elecció}.
\end{teorema}

\begin{prova}
Suposem que un tal \(\Omega\) existís. Aleshores \((\Omega,<)\) seria un bon
ordre, i tindria ell mateix un tipus \(\omega_\Omega=\operatorname{tipus}(\Omega)\).
Com que \(\Omega\) conté un representant de \emph{cada} tipus ordinal, conté un
element \(\gamma\) que representa \(\omega_\Omega\): és a dir, per la hipòtesi
sobre \(\Omega\), \(\operatorname{tipus}(\Omega_{<\gamma})=\omega_\Omega=\operatorname{tipus}(\Omega)\).
Però \(\Omega_{<\gamma}\) és un segment inicial \emph{propi} de \(\Omega\)
(propi perquè \(\gamma\in\Omega\)), i acabem de dir que és isomorf a \(\Omega\)
sencer. Això contradiu que cap bon ordre no és isomorf a cap segment inicial
propi seu (corol·lari~\ref{cor:no-segment-propi}). Per tant no existeix cap tal
\(\Omega\).
\end{prova}

\begin{observacio}[Què diu exactament el teorema]\label{obs:burali-hipotesi}
L'enunciat és precís i convé no ampliar-lo en llegir-lo informalment. Afirma
que no hi ha cap conjunt \(\Omega\) que \emph{alhora} (i)~contingui un
representant de cada tipus, (ii)~estigui ben ordenat per segments, i
(iii)~tingui la propietat que el segment \(\Omega_{<\gamma}\) sigui del tipus
que \(\gamma\) representa. La frase \enquote{no hi ha conjunt de tots els
ordinals} n'és la lectura habitual, i és correcta en aquest sentit: no hi ha
cap conjunt ben ordenat que \emph{indexi coherentment} tots els tipus. El que
el teorema no diu és res sobre col·leccions de representants sense aquesta
estructura d'ordre, que en aquest llibre ni tan sols es poden formular sense
un conjunt ambient. Com que els tipus no són elements de cap conjunt fix
(observació~\ref{obs:tipus-acotats}), la condició (iii) és la manera de dir
\enquote{\(\Omega\) és el conjunt dels tipus} sense pertinença global.
\end{observacio}

\begin{observacio}[Burali--Forti com a paral·lel de Cantor]\label{obs:burali-cantor}
El teorema és l'anàleg ordinal exacte del teorema de Cantor
(teorema~\ref{teo:cantor}), i comparteix amb ell l'estructura diagonal: la
contradicció surt de fer que la construcció s'apliqui a si mateixa. Cantor:
\(A\prec P(A)\), de manera que no hi ha conjunt universal ni mida màxima
(corol·lari~\ref{cor:no-universal}). Burali--Forti: no hi ha un bon ordre de
tots els tipus ordinals, de manera que no hi ha longitud d'ordre màxima. En
tots dos casos, la moralitat és que la col·lecció de \enquote{tots els}
objectes d'una mena ---totes les parts, tots els tipus--- és massa gran per
ser un conjunt del mateix tipus. És per això que la teoria \emph{global} dels
ordinals és de Nivell~II (observació~\ref{obs:tipus-acotats}): reunir tots els
tipus demanaria un principi ---SR--- que els axiomes de Nivell~I no donen, i
que ni tan sols amb SR produiria un \emph{conjunt} ben ordenat de tots ells,
sinó una família estratificada.
\end{observacio}

\begin{resumcap}
\apartat{Idees centrals}
\begin{enumerate}
  \item Un tipus ordinal és la classe d'un bon ordre llevat d'isomorfisme d'ordre; un bon ordre concret de la classe n'és un \emph{representant}, i el seu conjunt subjacent, el \emph{suport}. No hi ha representant canònic, i cada definició es comprova independent del representant gràcies a la unicitat de l'isomorfisme d'ordre.
  \item \(\alpha<\beta\) vol dir que un representant de \(\alpha\) és isomorf a un segment propi d'un de \(\beta\); l'ordre és total, i dins \(\Hart(A)\) és un bon ordre. No s'escriu \(\alpha\in\beta\): no hi ha pertinença global.
  \item El successor \(\alpha+1\) és el tipus de \(W+\mathbf1\), el menor tipus més gran que \(\alpha\). Tot tipus és exactament zero (representant buit), successor (representant amb màxim) o límit (representant no buit i sense màxim); \(\omega=\operatorname{tipus}(\N)\) és el primer límit.
  \item El suprem d'una família representada dins \(W\) per \(r\colon I\to W\) és el tipus de \(\bigcup_iW_{<r(i)}\), la reunió dels segments \emph{oberts}; l'acotació és imprescindible.
  \item Burali--Forti: no hi ha cap bon ordre \(\Omega\) que indexi coherentment tots els tipus (cada \(\gamma\in\Omega\) representa el tipus del segment \(\Omega_{<\gamma}\)). És el paral·lel ordinal de Cantor, i és per això que la teoria global dels ordinals queda fora del Nivell~I.
\end{enumerate}
\apartat{Errors típics} Escollir un representant canònic i oblidar comprovar la independència; escriure \(\alpha\in\beta\); confondre \(\mathbf1+W\) (punt nou al principi, mateix tipus que \(W\) si \(W\) és infinit sense màxim, com \(\mathbf1+\N\cong\N\)) amb \(W+\mathbf1\); prendre el suprem amb segments tancats \(W_{\le r(i)}\), que dona un tipus massa gran; i parlar del \enquote{conjunt de tots els ordinals}.
\apartat{Lectura recomanada} \cite[§§19--21]{halmos} per als ordinals, la seva aritmètica inicial i el teorema de Burali--Forti en la presentació de von Neumann; \cite[cap.~2]{jech2003} i \cite[cap.~I.7--I.8]{kunen2011} per a la teoria clàssica dels ordinals; i \cite[§§8.4--8.5]{lawvere-rosebrugh} per a la lectura estructural dels bons ordres com a tipus.
\end{resumcap}

\chapter{Aritmètica ordinal}\label{cap:aritmetica-ordinal}

\begin{objectius}
\apartat{Prerequisits} Del capítol~\ref{cap:ordinals}: tipus ordinals, el seu ordre, successor, límit i suprem acotat. Del capítol~\ref{cap:bons-ordres}: rigidesa, segments i comparació. Del capítol~\ref{cap:recursio}: recursió transfinita amb codomini fix. Del capítol~\ref{cap:zorn}: el conjunt de Hartogs com a cota. Dels capítols~\ref{cap:valors} i~\ref{cap:parts}: sumes i productes de conjunts.
\apartat{Objectius} En acabar el capítol, el lector hauria de poder:
\begin{itemize}
  \item definir la suma i el producte ordinals sobre representants (concatenació; \(\beta\) còpies d'\(\alpha\) amb l'ordre antilexicogràfic), demostrar que són bons ordres ben definits i associatius;
  \item calcular \(1+\omega\), \(\omega+1\), \(2\cdot\omega\), \(\omega\cdot2\), i demostrar la no-commutativitat i la distributivitat només per l'esquerra;
  \item definir l'exponenciació per funcions de suport finit, deduir-ne les equacions recursives per a base no nul·la, i explicar per què la base \(0\) es tracta a part;
  \item demostrar que una immersió d'ordre entre bons ordres compara els tipus, i deduir-ne la monotonia de les tres operacions;
  \item distingir l'aritmètica ordinal (recorda l'ordre) de la cardinal (l'oblida), amb \(2^{\omega}=\omega\) com a exemple clau.
\end{itemize}
\end{objectius}

Els tipus ordinals mesuren la forma d'un bon ordre; ara els operem. La suma,
el producte i l'exponenciació ordinals es defineixen ---seguint la disciplina
del capítol anterior--- \emph{sobre representants}, amb construccions
concretes de bons ordres, i després es comprova que el resultat no depèn del
representant. Dues coses fan aquesta aritmètica diferent de la dels naturals,
i val la pena anunciar-les d'entrada. Primer, \emph{no és commutativa}: l'ordre
en què es concatenen o s'iteren els bons ordres importa, i \(1+\omega\) no és
\(\omega+1\). Segon, recorda \emph{l'ordre} i no la \emph{mida}: dos bons
ordres poden ser equipotents ---tenir la mateixa cardinalitat--- i tenir tipus
ordinals ben diferents; l'aritmètica ordinal distingeix el que la cardinal
(capítol següent) oblidarà. La suma i el producte són operacions binàries de
Nivell~I, construïdes directament sobre representants. L'exponenciació es
defineix per recursió transfinita (capítol~\ref{cap:recursio}), i la seva
clàusula límit ---un suprem--- és el primer lloc on cal vigilar l'acotació:
mentre la base i l'exponent visquin dins un \(\Hart(A)\) prou gran, tot és de
Nivell~I; una teoria completament general demanaria SR.

\section{Suma ordinal}\label{sec:suma-ordinal}

\begin{definicio}[Suma ordinal]\index{Suma ordinal}\label{def:suma-ordinal}
Siguin \(\alpha=\operatorname{tipus}(A)\) i \(\beta=\operatorname{tipus}(B)\)
tipus ordinals amb representants \((A,<_A)\) i \((B,<_B)\). La \emph{suma
ordinal} \(\alpha+\beta\) és el tipus del bon ordre sobre la suma disjunta
\(A+B\) (definició~\ref{def:suma-disjunta}) en què \emph{tots els elements
d'\(A\) precedeixen tots els de \(B\)}, i dins de cada sumand es manté l'ordre
original:
\[
x<y
\ \Longleftrightarrow\
\begin{cases}
x<_A y & \text{si } x,y\in A,\\
x<_B y & \text{si } x,y\in B,\\
\text{cert} & \text{si } x\in A,\ y\in B,\\
\text{fals} & \text{si } x\in B,\ y\in A.
\end{cases}
\]
\end{definicio}

\begin{proposicio}[La suma és un bon ordre, invariant]\label{prop:suma-bon-ordre}
L'ordre de la definició~\ref{def:suma-ordinal} és un bon ordre sobre \(A+B\),
i el seu tipus no depèn dels representants: si \(A\cong A'\) i \(B\cong B'\)
com a ordres, aleshores \(A+B\cong A'+B'\). Per tant \(\alpha+\beta\) està ben
definida.\par
\emph{Axiomes usats:} suma disjunta (capítol~\ref{cap:valors}), bon ordre
(capítol~\ref{cap:bons-ordres}), unicitat de l'isomorfisme d'ordre
(teorema~\ref{teo:unicitat-iso}); cap elecció.
\end{proposicio}

\begin{prova}
\emph{És un ordre estricte total}: la irreflexivitat i la transitivitat es
comproven per casos segons els sumands (els casos mixtos són consistents
perquè \enquote{\(A\) abans de \(B\)} és una decisió global), i la totalitat
val perquè dos elements qualssevol o són al mateix sumand ---on \(<_A\) o
\(<_B\) els compara--- o a sumands diferents ---on la clàusula mixta els
compara. \emph{És ben ordenat}: sigui \(P\hookrightarrow A+B\) una part no
inicial. Si \(P\cap A\) és no inicial, el seu mínim (pel bon ordre d'\(A\)) és
el mínim de \(P\), perquè tot element d'\(A\) precedeix tot element de \(B\).
Si \(P\cap A\) és inicial, aleshores \(P\subseteq B\) i el mínim de \(P\) és el
seu mínim pel bon ordre de \(B\). \emph{Invariància}: dos isomorfismes d'ordre
\(f\colon A\cong A'\) i \(g\colon B\cong B'\) donen \(f+g\colon A+B\cong A'+B'\),
que conserva l'ordre perquè respecta cada sumand i la relació
\enquote{\(A\) abans de \(B\)}.
\end{prova}

\begin{proposicio}[Associativitat de la suma]\label{prop:suma-associativa}
Per a tipus ordinals qualssevol, \((\alpha+\beta)+\gamma=\alpha+(\beta+\gamma)\).
\end{proposicio}

\begin{prova}
Els dos costats són el tipus del bon ordre sobre \(A+B+C\) en què tot \(A\)
precedeix tot \(B\), que precedeix tot \(C\), amb els ordres originals dins de
cada sumand. L'associativitat de la suma disjunta de conjunts
(capítol~\ref{cap:valors}) dona un isomorfisme \((A+B)+C\cong A+(B+C)\) que
conserva aquest ordre; per la unicitat de l'isomorfisme d'ordre
(teorema~\ref{teo:unicitat-iso}), els tipus coincideixen.
\end{prova}

\begin{proposicio}[La suma no és commutativa]\label{prop:suma-no-commutativa}
En general \(\alpha+\beta\neq\beta+\alpha\). Concretament,
\[
1+\omega=\omega,
\qquad\text{però}\qquad
\omega+1\neq\omega .
\]
\end{proposicio}

\begin{prova}
\emph{\(1+\omega=\omega\)}: \(1+\omega\) és el tipus d'un punt seguit d'una
còpia de \(\N\); l'aplicació que envia el punt afegit a \(0\) i el natural
\(n\) a \(n+1\) és un isomorfisme d'ordre amb \(\N\) ---simplement desplaça la
numeració---, de manera que \(1+\omega=\operatorname{tipus}(\N)=\omega\).
\emph{\(\omega+1\neq\omega\)}: \(\omega+1\) és el tipus d'una còpia de \(\N\)
seguida d'un punt \(\ast\) col·locat \emph{després} de tots els naturals.
Aquest bon ordre té un element \emph{màxim} (el punt \(\ast\)), mentre que
\(\N\) no en té cap; com que tenir màxim és invariant per isomorfisme d'ordre,
\(\omega+1\not\cong\omega\), és a dir \(\omega+1\neq\omega\). De fet
\(\omega+1\) és el successor de \(\omega\) (definició~\ref{def:zero-successor}),
estrictament més gran que \(\omega\), mentre que \(1+\omega=\omega\): la suma
per l'esquerra d'un punt \enquote{s'absorbeix} i per la dreta no.
\end{prova}

\begin{observacio}[Lectura clàssica: concatenar ordres]\label{obs:suma-classica}
La suma ordinal és la \emph{concatenació} de bons ordres: \(\alpha+\beta\) és
\enquote{primer tot \(\alpha\), després tot \(\beta\)}. La no-commutativitat
és intuïtiva des d'aquest punt de vista: afegir un element \emph{abans} d'una
successió infinita ascendent no canvia la forma (es pot reabsorbir desplaçant
els índexs), però afegir-lo \emph{després} crea un element màxim que abans no
hi era, i canvia la forma. És el primer exemple, i el més net, de la
diferència entre ordre i mida: \(1+\omega\) i \(\omega+1\) són equipotents
---tots dos numerables--- però tenen tipus ordinals diferents.
\end{observacio}

\section{Producte ordinal}\label{sec:producte-ordinal}

\begin{definicio}[Producte ordinal]\index{Producte ordinal}\label{def:producte-ordinal}
Siguin \(\alpha=\operatorname{tipus}(A)\) i \(\beta=\operatorname{tipus}(B)\).
El \emph{producte ordinal} \(\alpha\cdot\beta\) és el tipus del bon ordre
sobre el producte \(A\times B\) ordenat \emph{primer per la coordenada de
\(B\) i després per la d'\(A\)} ---l'ordre \emph{antilexicogràfic}:
\[
(a,b)<(a',b')
\quad\Longleftrightarrow\quad
b<_B b'\ \ \text{o bé}\ \ (b=b'\ \text{i}\ a<_A a').
\]
Intuïtivament, \(\alpha\cdot\beta\) són \(\beta\) còpies d'\(\alpha\)
col·locades una rere l'altra: per a cada \(b\in B\) hi ha una còpia sencera
d'\(A\), i les còpies se succeeixen segons l'ordre de \(B\). Fixem aquesta
convenció ---la coordenada de \(B\) mana--- abans de qualsevol càlcul.
\end{definicio}

\begin{proposicio}[El producte és un bon ordre, invariant i associatiu]\label{prop:producte-bon-ordre}
L'ordre antilexicogràfic de la definició~\ref{def:producte-ordinal} és un bon
ordre sobre \(A\times B\); el seu tipus no depèn dels representants; i el
producte és associatiu, \((\alpha\cdot\beta)\cdot\gamma=\alpha\cdot(\beta\cdot\gamma)\).\par
\emph{Axiomes usats:} producte (capítol~\ref{cap:parts}), bon ordre, unicitat
de l'isomorfisme d'ordre (teorema~\ref{teo:unicitat-iso}); cap elecció.
\end{proposicio}

\begin{prova}
\emph{Bon ordre}: és un ordre estricte total (l'antilexicogràfic d'ordres
totals ho és). Per al bon ordre, sigui \(P\hookrightarrow A\times B\) no
inicial; la projecció \(\pi_B(P)\hookrightarrow B\) és no inicial i té un mínim
\(b_0\); dins la fibra \(\{(a,b_0)\in P\}\), la projecció a \(A\) és no
inicial i té un mínim \(a_0\); aleshores \((a_0,b_0)\) és el mínim de \(P\),
perquè cap element amb segona coordenada \(>b_0\) no el precedeix, i dins
\(b_0\) mana \(a_0\). \emph{Invariància}: \(f\times g\) transporta l'ordre.
\emph{Associativitat}: tots dos costats són el tipus de \(A\times B\times C\)
amb l'ordre \enquote{la coordenada de \(C\) primer, després la de \(B\),
després la d'\(A\)}; l'isomorfisme canònic \((A\times B)\times C\cong A\times(B\times C)\)
el conserva, i la unicitat (teorema~\ref{teo:unicitat-iso}) tanca.
\end{prova}

\begin{proposicio}[El producte no és commutatiu; distributivitat per l'esquerra]\label{prop:producte-propietats}
En general \(\alpha\cdot\beta\neq\beta\cdot\alpha\); concretament
\[
2\cdot\omega=\omega,
\qquad\text{però}\qquad
\omega\cdot2=\omega+\omega\neq\omega .
\]
A més val la distributivitat \emph{per l'esquerra}:
\(\alpha\cdot(\beta+\gamma)=\alpha\cdot\beta+\alpha\cdot\gamma\); en canvi la
distributivitat per la dreta \emph{falla} en general.
\end{proposicio}

\begin{prova}
\emph{\(2\cdot\omega=\omega\)}: són \(\omega\) còpies del bon ordre de dos
elements, és a dir la successió \(0<1<2<3<\cdots\) agrupada de dos en dos
\((0,1),(2,3),\dots\); reindexant, és isomorfa a \(\N\), d'on
\(2\cdot\omega=\omega\). \emph{\(\omega\cdot2=\omega+\omega\)}: són \(2\)
còpies de \(\N\), l'una després de l'altra ---exactament la concatenació
\(\omega+\omega\), que té un element que no és el mínim i no té predecessor
immediat (el primer de la segona còpia); \(\N\) no en té cap de tal, perquè
tot punt no nul és un successor (lema~\ref{lema:zero-o-successor}), i la
propietat és invariant per isomorfisme d'ordre. Per tant
\(\omega+\omega\not\cong\omega\). Per
tant \(\omega\cdot2\neq2\cdot\omega\). \emph{Distributivitat esquerra}:
\(\alpha\cdot(\beta+\gamma)\) és \((\beta+\gamma)\) còpies d'\(\alpha\), és a
dir \(\beta\) còpies seguides de \(\gamma\) còpies, que és
\(\alpha\cdot\beta+\alpha\cdot\gamma\); l'isomorfisme d'ordre és el canònic
\(A\times(B+C)\cong(A\times B)+(A\times C)\), que conserva l'antilexicogràfic.
\emph{La dreta falla}: \((\ 1+1)\cdot\omega=2\cdot\omega=\omega\), però
\(1\cdot\omega+1\cdot\omega=\omega+\omega\neq\omega\).
\end{prova}

\begin{observacio}[Lectura clàssica: substituir cada punt per una còpia]\label{obs:producte-classic}
El producte \(\alpha\cdot\beta\) s'obté prenent el bon ordre \(\beta\) i
substituint \emph{cada} un dels seus punts per una còpia sencera d'\(\alpha\).
La convenció \enquote{la coordenada de \(\beta\) mana} és la que fa que això
sigui \(\beta\) còpies d'\(\alpha\) i no a l'inrevés; amb la convenció
oposada s'obtindria \(\beta\cdot\alpha\), i com que el producte no és
commutatiu, cal fixar-la. La distributivitat només val per l'esquerra
perquè substituir cada punt d'una concatenació \(\beta+\gamma\) per una còpia
d'\(\alpha\) dona la concatenació de les substitucions; en canvi
\enquote{repartir} una còpia d'\(\alpha\) sobre una suma a la dreta trenca
l'estructura.
\end{observacio}

\section{Exponenciació ordinal}\label{sec:exponenciacio-ordinal}

L'exponenciació no es defineix per una construcció directa sinó \emph{per
recursió} sobre l'exponent, amb un suprem als límits. És on la disciplina de
l'acotació esdevé imprescindible.

Per evitar una recursió sobre \enquote{tots els tipus}, que demanaria un
codomini fix i una cota per a cada suprem, donem una construcció
\emph{directa} sobre representants, com hem fet amb la suma i el producte,
i en deduïm les equacions recursives com a teorema.

\begin{definicio}[Funcions de suport finit]\index{Funcions de suport finit}\label{def:suport-finit}
Siguin \((A,<_A)\) i \((B,<_B)\) bons ordres, amb \(A\) no buit i mínim
\(0_A\). Una aplicació \(f\colon B\to A\) té \emph{suport finit} si la part
\(\operatorname{supp}(f):=\{b:f(b)\neq0_A\}\hookrightarrow B\) és finita
(capítol~\ref{cap:finits}). Les aplicacions de suport finit formen una part
\[
E(B,A)\ \hookrightarrow\ A^{B},
\]
tallada per la condició interna \enquote{la part classificada per
\(\neg\circ\chi_{0_A}\circ\operatorname{ev}\) és finita}. L'ordenem
\emph{antilexicogràficament}: \(f<g\) si \(f\neq g\) i, al \emph{màxim}
punt \(b\) on difereixen, \(f(b)<_Ag(b)\). El màxim existeix perquè els
punts on difereixen formen una part finita no buida
(\(\subseteq\operatorname{supp}f\cup\operatorname{supp}g\)) d'un ordre
total.
\end{definicio}

\begin{definicio}[Exponenciació ordinal]\index{Exponenciació ordinal}\label{def:exponenciacio-ordinal}
Siguin \(\alpha=\operatorname{tipus}(A)\) i \(\beta=\operatorname{tipus}(B)\).
Si \(\alpha>0\), definim
\[
\alpha^{\beta}:=\operatorname{tipus}\bigl(E(B,A),<\bigr).
\]
Si \(\alpha=0\), és a dir, si \(A\) és inicial, no hi ha mínim \(0_A\) i
\(E(B,A)\) no es pot formar; definim aleshores
\[
0^{\beta}:=\operatorname{tipus}\bigl(A^{B}\bigr)
=\begin{cases}1&\text{si }\beta=0,\\0&\text{si }\beta>0,\end{cases}
\]
on \(A^{B}\), que té com a molt un punt, porta l'únic ordre possible. Els
dos valors es justifiquen dins la teoria: si \(B\) és inicial,
\(A^{B}\cong\mathbf1\) (capítol~\ref{cap:exponencials}); si \(B\) no és
inicial, un punt de \(A^{B}\) seria una aplicació \(B\to\mathbf0\), que per
l'estrictesa de \(\mathbf0\) (proposició~\ref{prop:inicial-estricte}) faria
\(B\) inicial, de manera que \(A^{B}\) no té cap punt i és inicial (lema
del buit).
\end{definicio}

\begin{teorema}[L'exponenciació està ben definida i satisfà la recursió]\label{teo:exponenciacio-recursio}
Sigui \(\alpha>0\). Per a tot bon ordre \(B\):
\begin{enumerate}
\item \(E(B,A)\) és un bon ordre, i el seu tipus no depèn dels representants
d'\(\alpha\) i \(\beta\);
\item si \(B'\hookrightarrow B\) és una part inicial, l'extensió per \(0_A\)
fora de \(B'\) identifica \(E(B',A)\) amb una part inicial de \(E(B,A)\);
\item \(\alpha^{0}=1\), \(\alpha^{\beta+1}=\alpha^{\beta}\cdot\alpha\), i, si
\(\lambda\) és límit, \(\alpha^{\lambda}=\sup_{\beta<\lambda}\alpha^{\beta}\),
on el suprem és el de la proposició~\ref{prop:suprem}.
\end{enumerate}
\emph{Axiomes usats:} exponencials, finitud (capítol~\ref{cap:finits}),
bons ordres, inducció transfinita; cap elecció; Nivell~I.
\end{teorema}

\begin{prova}
\emph{(2)} Sigui \(\iota\colon E(B',A)\to E(B,A)\) l'extensió per \(0_A\).
És estrictament creixent (els punts on difereixen dues extensions són els
mateixos). La imatge és inicial: si \(g<\iota(f)\) i \(b\) és el màxim punt
on difereixen, aleshores \(b\in B'\) ---si no, \(\iota(f)(b)=0_A\) i
\(g(b)<0_A\) seria impossible---, i per sobre de \(b\) les dues funcions
coincideixen, de manera que \(g\) també s'anul·la fora de \(B'\) (que és
inicial i conté \(b\), de manera que tot punt fora de \(B'\) és \(>b\)).

\emph{(3), cas zero i successor.} Si \(B\) és buit, \(A^{B}\) té un sol
punt, de suport buit: \(\alpha^0=1\). Si \(B=C+\mathbf1\) amb punt final
\(\ast\), l'aplicació \(f\mapsto(f|_C,f(\ast))\) és una bijecció
\(E(B,A)\cong E(C,A)\times A\) (una funció de suport finit sobre \(B\) és
una sobre \(C\) més un valor a \(\ast\)), i l'ordre antilexicogràfic, que
mira primer el punt \(\ast\), es correspon exactament amb l'ordre del
producte ordinal, on mana la coordenada d'\(A\)
(definició~\ref{def:producte-ordinal}): \(\alpha^{\beta+1}=\alpha^{\beta}\cdot\alpha\).

\emph{(3), cas límit.} Si \(B\) no té màxim, tota \(f\in E(B,A)\) té suport
finit, acotat per algun punt \(b\) de \(B\) (si el suport és no buit, el seu
màxim té un posterior, perquè \(B\) no té màxim); per tant
\(E(B,A)=\bigcup_{b\in B}E(B_{<b},A)\), una unió de parts inicials
encadenades, per~(2). Pel mateix argument de la proposició~\ref{prop:suprem}
(isomorfismes únics compatibles sobre una reunió de parts inicials), el tipus
de la reunió és el suprem dels tipus \(\alpha^{\operatorname{tipus}(B_{<b})}\),
és a dir, \(\sup_{\beta<\lambda}\alpha^{\beta}\).

\emph{(1)} Per inducció transfinita sobre els segments de \(B\) i \(B\)
mateix (teorema~\ref{teo:induccio-transfinita}): el cas buit és trivial; el
successor és un producte ordinal de bons ordres, que és un bon ordre
(proposició~\ref{prop:producte-bon-ordre}); i al límit, una reunió
encadenada de parts inicials ben ordenades és ben ordenada (el mínim d'una
part no buida es troba ja dins de qualsevol membre que la talli, com a la
prova del lema de Zorn). La independència dels representants: isomorfismes
d'ordre \(A\cong A'\) i \(B\cong B'\) indueixen, per composició, un
isomorfisme d'ordre \(E(B,A)\cong E(B',A')\) que respecta el mínim i el
suport.
\end{prova}

\begin{observacio}[Per què no definir-la directament per recursió]\label{obs:exponenciacio-acotacio}
La presentació habitual defineix \(\alpha^{\beta}\) per recursió transfinita
sobre l'exponent, amb un suprem als límits. Aquí això tindria dos
inconvenients. Primer, una recursió necessita un codomini fix i cada suprem
una família representada dins un bon ordre (observació~\ref{obs:suprem-acotat});
caldria exhibir-los. La construcció per funcions de suport finit els
proporciona d'entrada: totes les potències \(\alpha^{\beta}\) amb
\(\beta\le\lambda\) viuen, per~(2), com a parts inicials d'un sol bon ordre
\(E(L,A)\), sense cap cota externa ni cap apel·lació a SR. Segon, per a la base \(0\) les equacions del zero i del successor
continuen valent (\(0^{0}=1\) i \(0^{\beta+1}=0^{\beta}\cdot0=0\)), però la
del límit \emph{no}: amb la clàusula límit literal,
\(0^{\omega}=\sup\{0^0,0^1,0^2,\dots\}=\sup\{1,0,0,\dots\}=1\), mentre que
la definició dona \(0^{\omega}=0\), d'acord amb el fet que no hi ha cap
aplicació \(\N\to\mathbf0\). La raó és que la família \(1,0,0,\dots\) no és
creixent, i un suprem només recull bé una família creixent. Per això la
base \(0\) es defineix a part, i la clàusula límit del
teorema~\ref{teo:exponenciacio-recursio} només s'afirma per a
\(\alpha>0\), on la família \((\alpha^\beta)_\beta\) és creixent
(proposició~\ref{prop:monotonia-ordinal}). Una exponenciació amb bases i exponents variant en famílies
sense un conjunt ambient comú, en canvi, sí que demanaria SR
(observació~\ref{obs:sr}).
\end{observacio}

\begin{proposicio}[Càlculs bàsics amb l'exponenciació]\label{prop:exponenciacio-calculs}
Amb la definició~\ref{def:exponenciacio-ordinal} i el teorema~\ref{teo:exponenciacio-recursio}, per a \(\alpha>0\):
\[
\alpha^{1}=\alpha,
\qquad
\alpha^{2}=\alpha\cdot\alpha,
\qquad
2^{\omega}=\omega,
\qquad
\omega^{2}=\omega\cdot\omega .
\]
En particular \(2^{\omega}=\omega\neq\mathfrak c\): l'exponenciació
\emph{ordinal} \(2^{\omega}\) és numerable, ben diferent de l'exponenciació
\emph{cardinal} \(2^{\aleph_0}\), que és el continu.
\end{proposicio}

\begin{prova}
\(\alpha^{1}=\alpha^{0+1}=\alpha^{0}\cdot\alpha=1\cdot\alpha=\alpha\), i
\(\alpha^{2}=\alpha^{1}\cdot\alpha=\alpha\cdot\alpha\). Per a \(2^{\omega}\):
és el suprem dels \(2^{n}\) amb \(n<\omega\), tots finits, i el suprem d'una
successió creixent de tipus finits no acotada és \(\omega\); per tant
\(2^{\omega}=\omega\). I \(\omega^{2}=\omega^{1}\cdot\omega=\omega\cdot\omega\).
La comparació amb la cardinal se subratlla perquè és el nucli de la distinció
del capítol següent: \(2^{\omega}=\omega\) com a \emph{tipus ordinal} (una
longitud numerable), mentre que com a \emph{cardinal} \(2^{\aleph_0}\) és no
numerable (teorema~\ref{teo:cantor}); les dues operacions, escrites igual,
signifiquen coses distintes.
\end{prova}

\section{Monotonia i ordre contra mida}\label{sec:monotonia-ordre-mida}

\begin{lema}[Les immersions d'ordre comparen tipus]\label{lema:immersio-ordre}
Siguin \(W\) i \(V\) bons ordres i \(f\colon W\to V\) una aplicació
estrictament creixent. Aleshores \(\operatorname{tipus}(W)\le\operatorname{tipus}(V)\).
\end{lema}

\begin{prova}
Per la comparabilitat (teorema~\ref{teo:comparabilitat}), si no fos
\(\operatorname{tipus}(W)\le\operatorname{tipus}(V)\) tindríem
\(V\cong W_{<w}\) per a algun \(w\); compost amb \(f\), donaria una
aplicació estrictament creixent \(W\to W_{<w}\hookrightarrow W\) amb valor a
\(w\) estrictament menor que \(w\), contra la rigidesa
(lema~\ref{lema:rigidesa}).
\end{prova}

\begin{proposicio}[Monotonia]\label{prop:monotonia-ordinal}
Les tres operacions són \emph{estrictament creixents per la dreta} i
\emph{no decreixents per l'esquerra}. En particular, si \(\beta<\gamma\):
\[
\alpha+\beta<\alpha+\gamma,
\qquad
\alpha\cdot\beta<\alpha\cdot\gamma\ (\alpha>0),
\qquad
\alpha^{\beta}<\alpha^{\gamma}\ (\alpha>1),
\]
mentre que per l'esquerra només val \(\le\): de \(\beta<\gamma\) se'n
segueix \(\beta+\alpha\le\gamma+\alpha\), i pot fallar l'estrictesa
(\(1<2\) però \(1+\omega=2+\omega=\omega\)).
\end{proposicio}

\begin{prova}
Siguin \(A,B,C\) representants d'\(\alpha,\beta,\gamma\), amb
\(B\cong C_{<c}\) per a un punt \(c\) de \(C\) (perquè \(\beta<\gamma\)).

\emph{Suma per la dreta.} \(A+C_{<c}\) és el segment de \(A+C\) determinat
pel punt \(c\) del segon sumand, que és propi; d'on
\(\alpha+\beta<\alpha+\gamma\).

\emph{Producte per la dreta.} Si \(\alpha>0\), sigui \(a_0\) el mínim d'\(A\).
En l'ordre antilexicogràfic, \(A\times C_{<c}\) és el segment de
\(A\times C\) determinat per \((a_0,c)\), que és propi; d'on
\(\alpha\cdot\beta<\alpha\cdot\gamma\).

\emph{Exponenciació per la dreta.} Si \(\alpha>1\), el representant \(A\)
té, a més del mínim \(0_A\), algun punt \(a_1\neq0_A\). Per la invariància
(teorema~\ref{teo:exponenciacio-recursio}(1)),
\(\alpha^{\beta}=\operatorname{tipus}(E(C_{<c},A))\), i per
l'apartat~(2) d'aquell teorema l'extensió per \(0_A\) identifica
\(E(C_{<c},A)\) amb una part inicial de \(E(C,A)\). Aquesta part és
\emph{pròpia}: l'aplicació \(g\colon C\to A\) que val \(a_1\) al punt \(c\) i
\(0_A\) a la resta (definició per casos sobre el singletó \(\{c\}\) i el seu
complement) té suport \(\{c\}\), finit, i no és l'extensió de cap funció
sobre \(C_{<c}\), perquè \(g(c)\neq0_A\). Per la
proposició~\ref{prop:parts-inicials}, \(E(C_{<c},A)\) és isomorf a un segment
propi de \(E(C,A)\), és a dir, \(\alpha^{\beta}<\alpha^{\gamma}\). No cal cap
inducció sobre els tipus, que no formen cap conjunt: l'argument passa
íntegrament dins el bon ordre fix \(C\).

\emph{Per l'esquerra.} Si \(\beta\le\gamma\), hi ha una immersió estrictament
creixent \(B\to C\) (sobre un segment); en resulten immersions estrictament
creixents \(B+A\to C+A\) (per a \(\beta+\alpha\le\gamma+\alpha\)) i
\(B\times A\to C\times A\) en l'ordre antilexicogràfic (per a
\(\beta\cdot\alpha\le\gamma\cdot\alpha\)),
i pel lema~\ref{lema:immersio-ordre} els tipus es comparen. Per a
l'exponenciació, sigui \(D\) un representant de \(\delta\); volem
\(\beta^{\delta}\le\gamma^{\delta}\). Si \(\beta=0\), \(\beta^{\delta}\) val
\(1\) o \(0\) segons si \(\delta=0\) o no
(definició~\ref{def:exponenciacio-ordinal}), i en el primer cas també
\(\gamma^{0}=1\); la desigualtat és immediata. Si \(\beta>0\), la immersió
\(j\colon B\to C\) és un isomorfisme d'ordre sobre una part inicial no
inicial de \(C\), que conté el mínim de \(C\); per tant \(j(0_B)=0_C\). La
postcomposició \(f\mapsto j\circ f\) envia \(E(D,B)\) dins \(E(D,C)\): conserva
el suport, perquè \(j\) és injectiva i \(j(0_B)=0_C\), i és estrictament
creixent, perquè \(f\) i \(f'\) difereixen exactament als mateixos punts que
\(j\circ f\) i \(j\circ f'\), i al màxim d'aquests punts \(j\) conserva la
desigualtat. Pel lema~\ref{lema:immersio-ordre},
\(\beta^{\delta}\le\gamma^{\delta}\). L'exemple \(1+\omega=2+\omega=\omega\)
(proposició~\ref{prop:suma-no-commutativa}) mostra que per l'esquerra
l'estrictesa pot fallar.
\end{prova}

\begin{observacio}[Ordre contra mida: la lliçó de l'aritmètica ordinal]\label{obs:ordre-vs-mida}
L'aritmètica ordinal \emph{recorda l'ordre}; la cardinal, que veurem tot
seguit, l'\emph{oblidarà}. Els càlculs d'aquest capítol ho fan tangible.
Ordinalment, \(1+\omega=\omega\) però \(\omega+1>\omega\);
\(2\cdot\omega=\omega\) però \(\omega\cdot2>\omega\); \(2^{\omega}=\omega\).
Els tipus \(\omega\), \(\omega+1\), \(\omega\cdot2\), \(\omega^{2}\) són
\emph{diferents} dos a dos com a formes d'ordre; en canvi \(2^{\omega}=\omega\)
(exponenciació \emph{ordinal}!), de manera que \(2^{\omega}\) no és un tipus
nou sinó el mateix \(\omega\). Tots ells són a més \emph{equipotents}: tots
numerables, tots de cardinal \(\aleph_0\). Com a \emph{mides} són el mateix. Aquesta és la raó de mantenir separades les dues
aritmètiques: la ordinal distingeix \(\omega+1\) de \(\omega\), cosa que la
cardinal no pot fer, perquè oblidar l'ordre és precisament identificar tots
els bons ordres numerables. El capítol següent construeix l'aritmètica
cardinal, on \enquote{\(+\)}, \enquote{\(\cdot\)} i \enquote{\(\,\hat{}\,\)}
signifiquen suma disjunta, producte i exponencial de conjunts \emph{sense}
ordre, i on \(2^{\aleph_0}\neq\aleph_0\) per Cantor.
\end{observacio}

\begin{resumcap}
\apartat{Idees centrals}
\begin{enumerate}
  \item Siguin \(A\) i \(B\) representants dels tipus \(\alpha\) i \(\beta\). La suma \(\alpha+\beta\) és el tipus de la concatenació \(A+B\) (tot \(A\) abans de tot \(B\)). El producte \(\alpha\cdot\beta\) és el tipus de \(A\times B\) amb l'ordre antilexicogràfic: una còpia d'\(A\) per cada punt de \(B\). Totes dues operacions no depenen dels representants i són associatives.
  \item No són commutatives: \(1+\omega=\omega<\omega+1\), \(2\cdot\omega=\omega<\omega\cdot2\). El producte distribueix només per l'esquerra.
  \item Per a \(\alpha>0\), l'exponenciació \(\alpha^\beta\) és el tipus del conjunt \(E(B,A)\) de les aplicacions de suport finit \(B\to A\), amb l'ordre antilexicogràfic; satisfà \(\alpha^0=1\), \(\alpha^{\beta+1}=\alpha^\beta\cdot\alpha\) i el suprem als límits, sense cap cota externa. Per a \(\alpha=0\), \(0^\beta\) és el tipus de \(A^B\): \(1\) si \(\beta=0\) i \(0\) altrament; la clàusula límit no val per a aquesta base. Així \(2^\omega=\omega\) ordinalment.
  \item Una immersió estrictament creixent compara tipus. Les tres operacions són estrictament creixents per la dreta i no decreixents per l'esquerra.
  \item L'aritmètica ordinal recorda l'ordre; la cardinal, que ve tot seguit, l'oblida: \(\omega\), \(\omega+1\), \(\omega\cdot2\), \(\omega^2\) són tipus diferents i la mateixa mida.
\end{enumerate}
\apartat{Errors típics} Invertir la convenció del producte (\(\omega\cdot2\) són dues còpies de \(\omega\), no \(\omega\) còpies de \(2\)); pensar que \(\omega+1=1+\omega\); confondre l'exponenciació ordinal \(2^\omega=\omega\) amb la cardinal \(2^{\aleph_0}\); aplicar la clàusula límit a la base \(0\); i creure que la monotonia per l'esquerra és estricta.
\apartat{Lectura recomanada} \cite[§§19--21]{halmos} per a l'aritmètica ordinal en la teoria clàssica; \cite[cap.~2]{jech2003} per a la suma, el producte i l'exponenciació ordinals, amb la forma normal de Cantor; i \cite[cap.~I.9]{kunen2011} per a la definició recursiva de les operacions i les seves propietats.
\end{resumcap}

\chapter{Cardinals i aritmètica cardinal}\label{cap:cardinals}

\begin{objectius}
\apartat{Prerequisits} Del capítol~\ref{cap:cantor}: comparació i equipotència, Schröder--Bernstein, Cantor i la diagonal. Del capítol~\ref{cap:zorn}: Hartogs, el teorema del bon ordre i les equivalències de l'elecció. Del capítol~\ref{cap:families}: sumes i productes de famílies. Dels capítols~\ref{cap:valors}, \ref{cap:parts} i~\ref{cap:exponencials}: les construccions i els isomorfismes canònics que transporten les lleis.
\apartat{Objectius} En acabar el capítol, el lector hauria de poder:
\begin{itemize}
  \item tractar el cardinal \(\abs{A}\) com a classe d'equipotència, sense representant, i definir-ne la suma, el producte i l'exponenciació comprovant-ne la invariància;
  \item demostrar les lleis de l'aritmètica cardinal com a isomorfismes canònics, i la monotonia en l'exponent només per a base no nul·la;
  \item demostrar Cantor en forma cardinal i distingir \(2^{\aleph_0}\) de \(2^{\omega}\);
  \item demostrar que la comparabilitat total dels cardinals equival a l'elecció;
  \item construir el cardinal inicial d'un conjunt dins \(\Hart(A)\) i explicar la prudència amb \enquote{el cardinal és un conjunt};
  \item demostrar el teorema de König, identificant cada ús de l'elecció, i deduir-ne Cantor.
\end{itemize}
\end{objectius}

L'aritmètica ordinal recordava l'ordre; la cardinal l'oblida. Un
\emph{cardinal} és el que queda d'un conjunt quan se n'oblida qualsevol
ordre i es reté només la \emph{mida} ---la classe d'equipotència. Aquest
capítol desenvolupa l'aritmètica de les mides amb la mateixa disciplina de
sempre: les operacions \(+\), \(\cdot\) i l'exponenciació es defineixen
\emph{sobre representants} ---suma disjunta, producte i exponencial de
conjunts--- i es comprova que no depenen del representant, és a dir que estan
ben definides sobre classes d'equipotència. En aquest primer nivell no cal
cap \enquote{nombre cardinal} com a objecte: n'hi ha prou amb la relació
d'equipotència i la comparació per monomorfismes del capítol~\ref{cap:cantor}.
El teorema de Cantor hi reapareix com la desigualtat fonamental
\(\abs{A}<2^{\abs{A}}\). Després, amb l'elecció, veurem que dos conjunts
qualssevol són comparables, de manera que els cardinals queden totalment
ordenats; i que cada conjunt té un \emph{cardinal inicial} ---el menor tipus
d'ordre de la seva mida---, amb la prudència que aquest és un tipus, no un
conjunt escollit. Tanquem amb el teorema de König, un refinament diagonal de
Cantor per a famílies.

\section{Cardinalitat sense nombres cardinals}\label{sec:cardinalitat-relacional}

Recordem del capítol~\ref{cap:cantor} la comparació de conjunts: \(A\preceq B\)
si hi ha un monomorfisme \(A\hookrightarrow B\), \(A\approx B\) si hi ha un
isomorfisme (equipotència, definició~\ref{def:equipotencia}), i \(A\prec B\)
si \(A\preceq B\) però \(A\not\approx B\) (definició~\ref{def:comparacio}). El
\emph{cardinal} d'\(A\), que escrivim \(\abs{A}\), és \(A\) considerat llevat
d'equipotència; \(\abs{A}=\abs{B}\) vol dir \(A\approx B\). Com amb els tipus
ordinals, no cal escollir cap representant: treballem amb conjunts qualssevol
i comprovem la invariància.

\begin{definicio}[Operacions cardinals]\index{Operacions cardinals}\label{def:operacions-cardinals}
Per a conjunts \(A\) i \(B\), es defineixen la \emph{suma}, el \emph{producte}
i l'\emph{exponenciació} cardinals per
\[
\abs{A}+\abs{B}:=\abs{A+B},
\qquad
\abs{A}\cdot\abs{B}:=\abs{A\times B},
\qquad
\abs{B}^{\abs{A}}:=\abs{B^{A}},
\]
és a dir, com els cardinals de la suma disjunta
(definició~\ref{def:suma-disjunta}), del producte
(definició~\ref{def:producte}) i de l'exponencial
(definició~\ref{def:exponencial}), respectivament. Aquestes són operacions
\emph{sobre classes d'equipotència}, no sobre cap conjunt global de cardinals.
\end{definicio}

\begin{proposicio}[Les operacions estan ben definides]\label{prop:operacions-invariants}
La suma, el producte i l'exponenciació cardinals no depenen dels
representants: si \(A\approx A'\) i \(B\approx B'\), aleshores
\[
A+B\approx A'+B',
\qquad
A\times B\approx A'\times B',
\qquad
B^{A}\approx B'^{A'} .
\]
\emph{Axiomes usats:} suma, producte, exponencials
(capítols~\ref{cap:valors},~\ref{cap:parts} i~\ref{cap:exponencials}); cap
elecció.
\end{proposicio}

\begin{prova}
Siguin \(f\colon A\cong A'\) i \(g\colon B\cong B'\) isomorfismes. Aleshores
\(f+g\colon A+B\cong A'+B'\) (functorialitat de la suma disjunta,
capítol~\ref{cap:valors}) i \(f\times g\colon A\times B\cong A'\times B'\)
(functorialitat del producte, capítol~\ref{cap:parts}) són isomorfismes. Per a
l'exponencial, \(h\mapsto g\circ h\circ f^{-1}\) és una bijecció
\(B^{A}\cong B'^{A'}\), amb inversa \(k\mapsto g^{-1}\circ k\circ f\);
internament, és la transposada de
\(g\circ\operatorname{ev}\circ(\id_{B^A}\times f^{-1})\colon B^A\times A'\to B'\)
(propietat universal de l'exponencial, definició~\ref{def:exponencial}). Les
tres construccions són isomorfismes de conjunts, de manera que les mides
coincideixen.
\end{prova}

\begin{proposicio}[Lleis de l'aritmètica cardinal]\label{prop:lleis-cardinals}
Per a cardinals qualssevol \(\kappa=\abs{A}\), \(\mu=\abs{B}\),
\(\nu=\abs{C}\):
\begin{enumerate}
\item suma i producte són \emph{commutatius} i \emph{associatius}, amb neutres
\(\abs{\mathbf0}=0\) per a la suma i \(\abs{\mathbf1}=1\) per al producte;
\item distributivitat: \(\kappa\cdot(\mu+\nu)=\kappa\cdot\mu+\kappa\cdot\nu\);
\item lleis d'exponents: \(\kappa^{\mu+\nu}=\kappa^{\mu}\cdot\kappa^{\nu}\),
\((\kappa\cdot\mu)^{\nu}=\kappa^{\nu}\cdot\mu^{\nu}\),
\((\kappa^{\mu})^{\nu}=\kappa^{\mu\cdot\nu}\);
\item monotonia: si \(\kappa\le\mu\) (és a dir \(A\preceq B\)), aleshores
\(\kappa+\nu\le\mu+\nu\), \(\kappa\cdot\nu\le\mu\cdot\nu\) i
\(\kappa^{\nu}\le\mu^{\nu}\) (monotonia en la \emph{base}); a l'\emph{exponent},
si a més \(\nu\neq 0\), aleshores \(\nu^{\kappa}\le\nu^{\mu}\);
\item compatibilitat amb els recomptes finits: sobre conjunts finits, aquestes
operacions coincideixen amb la suma, el producte i la potència de naturals
(capítol~\ref{cap:finits}).
\end{enumerate}
\emph{Axiomes usats:} propietats universals de suma, producte i exponencial;
cap elecció.
\end{proposicio}

\begin{prova}
Cada llei és un isomorfisme canònic de conjunts, transportat a mides per la
proposició~\ref{prop:operacions-invariants}. \emph{(1)} La commutativitat de
la suma és \(A+B\cong B+A\) i la del producte \(A\times B\cong B\times A\); les
associativitats són els isomorfismes canònics corresponents; els neutres són
\(A+\mathbf0\cong A\) i \(A\times\mathbf1\cong A\). \emph{(2)} És la
distributivitat \(A\times(B+C)\cong(A\times B)+(A\times C)\)
(capítol~\ref{cap:valors}). \emph{(3)} Són els isomorfismes exponencials, que alineem lletra a lletra amb
l'enunciat. Per a \(\kappa^{\mu+\nu}=\kappa^{\mu}\cdot\kappa^{\nu}\), amb
\(\kappa=\abs C\), \(\mu=\abs B\), \(\nu=\abs A\): \(C^{B+A}\cong C^{B}\times
C^{A}\) (una aplicació d'una suma és un parell d'aplicacions). Per a
\((\kappa\cdot\mu)^{\nu}=\kappa^{\nu}\cdot\mu^{\nu}\), amb \(\kappa=\abs C\),
\(\mu=\abs B\), \(\nu=\abs A\): \((C\times B)^{A}\cong C^{A}\times B^{A}\). Per
a \((\kappa^{\mu})^{\nu}=\kappa^{\mu\cdot\nu}\), amb \(\kappa=\abs C\),
\(\mu=\abs B\), \(\nu=\abs A\): \((C^{B})^{A}\cong C^{B\times A}\)
(currificació, definició~\ref{def:exponencial}). En els tres casos les lletres
dels representants s'assignen igual que a l'enunciat, i l'ordre del producte
\(B\times A\) és el que correspon a \(\mu\cdot\nu\) (indiferent per
commutativitat).
\emph{(4)} La monotonia \emph{en la base} surt de compondre un monomorfisme
\(A\hookrightarrow B\) amb les operacions, que preserven monomorfismes: de
\(A\hookrightarrow B\) s'obté \(A+C\hookrightarrow B+C\),
\(A\times C\hookrightarrow B\times C\) i \(A^{C}\hookrightarrow B^{C}\) per
postcomposició. Per a l'\emph{exponent}, en canvi, un monomorfisme
\(A\hookrightarrow B\) indueix per \emph{precomposició} una aplicació
\(C^{B}\to C^{A}\) ---en la direcció contrària---, i per anar de \(C^{A}\) cap
a \(C^{B}\) cal estendre cada \(f\colon A\to C\) a \(B\to C\). Amb
\(\nu=\abs C\neq 0\) fixem un punt \(c_0\in C\) i estenem per casos: sobre la
imatge de \(A\) copiem \(f\), i sobre el complement de \(A\) dins \(B\)
posem el valor constant \(c_0\); això dona un monomorfisme \(C^{A}\hookrightarrow
C^{B}\), és a dir \(\nu^{\kappa}\le\nu^{\mu}\). Quan \(\nu=0\) la construcció
falla i l'enunciat és fals: \(0^{0}=1\not\le 0=0^{1}\), tot i que \(0\le 1\).
Per això la monotonia en l'exponent només s'afirma per a \(\nu\neq 0\).
\emph{(5)} Sobre segments
\(\seg{m}\) i \(\seg{n}\), \(\seg{m}+\seg{n}\approx\seg{m+n}\),
\(\seg{m}\times\seg{n}\approx\seg{m\cdot n}\) i
\(\seg{n}^{\seg{m}}\approx\seg{n^{m}}\) són els recomptes finits del
capítol~\ref{cap:finits}.
\end{prova}

\begin{teorema}[Cantor, en forma cardinal]\label{teo:cantor-cardinal}
Per a tot conjunt \(A\),
\[
\abs{A}\ <\ 2^{\abs{A}} .
\]
És a dir, \(\abs{A}\le 2^{\abs{A}}\) i \(\abs{A}\neq 2^{\abs{A}}\), on
\(2^{\abs{A}}=\abs{\mathbf2^{A}}=\abs{P(A)}\).\par
\emph{Axiomes usats:} teorema de Cantor (teorema~\ref{teo:cantor}), exponencial
\(P(A)\cong\mathbf2^{A}\) (capítol~\ref{cap:exponencials}); cap elecció.
\end{teorema}

\begin{prova}
Com que \(P(A)\cong\mathbf2^{A}\) (el conjunt de parts és l'exponencial de
base \(\mathbf2\), capítol~\ref{cap:exponencials}), es té
\(2^{\abs{A}}=\abs{\mathbf2^{A}}=\abs{P(A)}\). El teorema de Cantor
(teorema~\ref{teo:cantor}) diu \(A\prec P(A)\), és a dir
\(\abs{A}\le\abs{P(A)}\) i \(\abs{A}\neq\abs{P(A)}\). Per tant
\(\abs{A}<2^{\abs{A}}\).
\end{prova}

\begin{observacio}[Ordinal contra cardinal, un cop més]\label{obs:cardinal-vs-ordinal}
L'exponenciació cardinal \(2^{\abs{A}}\) és radicalment diferent de l'ordinal
del capítol anterior. Ordinalment, \(2^{\omega}=\omega\)
(proposició~\ref{prop:exponenciacio-calculs}): una longitud numerable.
Cardinalment, \(2^{\aleph_0}=\abs{P(\N)}>\aleph_0\) pel teorema de Cantor: una
mida no numerable, el continu. La mateixa notació \enquote{\(2\) elevat a}
significa coses oposades segons si recorda l'ordre (ordinal) o l'oblida
(cardinal). És per evitar aquesta confusió que hem construït les dues
aritmètiques per separat, i que escriurem \(\aleph_0\) per al cardinal de
\(\N\) i \(\omega\) per al seu tipus d'ordre, tot i que \enquote{són} el
mateix conjunt.
\end{observacio}

\section{Comparabilitat i ordre total dels cardinals}\label{sec:comparabilitat-cardinals}

En el primer nivell, la comparació \(\preceq\) és un preordre invariant
(proposició~\ref{prop:preceq-preordre}) i, per Schröder--Bernstein, és
antisimètrica mòdul equipotència (corol·lari~\ref{cor:preceq-antisimetric}).
El que \emph{no} tenim sense elecció és que sigui \emph{total}: que dos
conjunts qualssevol siguin comparables. Això ho dona el teorema del bon ordre.

\begin{teorema}[Comparabilitat total dels cardinals]\label{teo:comparabilitat-cardinals}
Suposem l'axioma d'elecció (axioma~\ref{ax:eleccio}). Aleshores dos conjunts
qualssevol són comparables:
\[
A\preceq B
\qquad\text{o bé}\qquad
B\preceq A .
\]
Combinat amb Schröder--Bernstein (teorema~\ref{teo:schroder-bernstein}), la
comparació de cardinals és un \emph{ordre total llevat d'equipotència}:
reflexiu, transitiu, antisimètric mòdul \(\approx\) i total.\par
\emph{Axiomes usats:} teorema del bon ordre (teorema~\ref{teo:bon-ordre}) ---i
per tant elecció---, comparabilitat de bons ordres
(teorema~\ref{teo:comparabilitat}), Schröder--Bernstein
(teorema~\ref{teo:schroder-bernstein}).
\end{teorema}

\begin{prova}
Pel teorema del bon ordre (teorema~\ref{teo:bon-ordre}), \(A\) i \(B\) admeten
bons ordres; siguin \(W_A\) i \(W_B\) representants. Per la comparabilitat de
bons ordres (teorema~\ref{teo:comparabilitat}), un dels dos és isomorf a un
segment inicial de l'altre ---diguem \(W_A\) isomorf a un segment de \(W_B\).
Aquest isomorfisme d'ordre és en particular un monomorfisme de conjunts
\(A\hookrightarrow B\), d'on \(A\preceq B\). Simètricament en l'altre cas. Per
tant \(A\preceq B\) o \(B\preceq A\). La resta ---reflexivitat, transitivitat,
antisimetria mòdul \(\approx\)--- és la proposició~\ref{prop:preceq-preordre} i
el corol·lari~\ref{cor:preceq-antisimetric}.
\end{prova}

\begin{observacio}[La comparabilitat és equivalent a l'elecció]\label{obs:comparabilitat-eleccio}
La comparabilitat total dels cardinals no és només \emph{conseqüència} de
l'elecció: n'és \emph{equivalent}. Si dos conjunts qualssevol són comparables,
aleshores en particular tot conjunt \(A\) és comparable amb \(\Hart(A)\); com
que \(\Hart(A)\not\preceq A\) (teorema~\ref{teo:hartogs}), ha de ser
\(A\preceq\Hart(A)\), és a dir \(A\) s'injecta en un bon ordre, i per tant
hereta un bon ordre ---d'on el teorema del bon ordre, equivalent a l'elecció
(teorema~\ref{teo:equivalencies}). Aquesta és una altra manifestació del
paper de Hartogs: la seva cota converteix la comparabilitat en bon ordre.
\end{observacio}

\section{Cardinals inicials}\label{sec:cardinals-inicials}

Fins ara \enquote{\(\abs{A}\)} ha estat una classe d'equipotència, no un
objecte. Amb l'elecció podem, si volem, triar per a cada mida un
\emph{representant} distingit: el menor tipus d'ordre d'aquella mida. Cal
fer-ho amb la prudència que la guia recomana.

\begin{definicio}[Bon ordre inicial i cardinal inicial]\index{Bon ordre inicial i cardinal inicial}\label{def:cardinal-inicial}
Un bon ordre \(W\) és \emph{inicial} si cap bon ordre de tipus estrictament
menor no és equipotent a \(W\): per a tot bon ordre \(V\) amb
\(\operatorname{tipus}(V)<\operatorname{tipus}(W)\), es té \(V\not\approx W\).
Amb l'elecció, cada conjunt \(A\) és equipotent a algun bon ordre
(teorema~\ref{teo:bon-ordre}); el \emph{cardinal inicial} de \(A\) és el
\emph{menor} tipus d'ordre \(\alpha\) tal que un (tot) representant seu és
equipotent a \(A\). Aquest tipus és inicial i únic.
\end{definicio}

\begin{proposicio}[Existència i unicitat del cardinal inicial]\label{prop:cardinal-inicial}
Suposem l'elecció. Per a tot conjunt \(A\) existeix un únic cardinal inicial
\(\alpha\), i és un tipus d'ordre inicial. La cerca està \emph{acotada}: el
menor tipus equipotent a \(A\) és un element de \(\Hart(A)\), sense recórrer
cap classe global.\par
\emph{Axiomes usats:} teorema del bon ordre (teorema~\ref{teo:bon-ordre}),
Hartogs i el bon ordre de \(\Hart\) (proposició~\ref{prop:hartogs-bon-ordre});
elecció.
\end{proposicio}

\begin{prova}
Sigui \(E\hookrightarrow\Hart(A)\) la part dels tipus \([(S,R)]\) amb
\(S\approx A\) (una condició interna sobre \(\WO(A)\), invariant per
isomorfisme, que baixa al quocient). És no inicial: pel teorema del bon ordre
\(A\) admet un bon ordre \(R\), i \([(A,R)]\in E\). Per ser \(\Hart(A)\) un bon
ordre (proposició~\ref{prop:hartogs-bon-ordre}), \(E\) té un mínim \(\alpha\),
únic. És inicial: si un bon ordre \(V\) amb \(\operatorname{tipus}(V)<\alpha\)
fos equipotent a \(A\), aleshores \(V\) seria isomorf a un segment propi d'un
representant \((S,R)\) de \(\alpha\), és a dir, a un bon ordre sobre una part
de \(S\subseteq A\); el seu tipus seria un element de \(E\) menor que
\(\alpha\), contra la minimalitat. Tota la cerca ha passat dins \(\Hart(A)\),
acotadament.
\end{prova}

\begin{observacio}[\enquote{El cardinal és un conjunt}: una prudència]\label{obs:cardinal-prudencia}
Convé ser prudent amb la frase \enquote{el cardinal de \(A\) és un conjunt}.
En el tractament estructural, el cardinal inicial de \(A\) és un \emph{tipus
d'ordre} inicial ---únic llevat d'un únic isomorfisme d'ordre
(teorema~\ref{teo:unicitat-iso})---, no un conjunt privilegiat escollit
d'entrada. Triar-ne un representant concret (un bon ordre determinat d'aquell
tipus) és una \emph{tria addicional}, legítima però no canònica. Aquesta és la
raó per la qual, en el primer nivell, hem treballat amb \(\abs{A}\) com a
classe d'equipotència i no com a objecte: la teoria relacional no necessita
representants, i només els introdueix ---amb l'elecció--- qui vulgui una
successió distingida de cardinals. Una \emph{successió global} de tots els
cardinals inicials, d'altra banda, seria de Nivell~II: reunir-los tots
demanaria SR (observació~\ref{obs:sr}), igual que la successió dels
\(\aleph_\alpha\) que veurem al capítol opcional.
\end{observacio}

\section{El teorema de König}\label{sec:konig}

Tanquem amb un refinament de Cantor per a famílies: si cada terme d'una
família és estrictament més petit que el corresponent d'una altra, la suma
dels primers és estrictament menor que el producte dels segons. La prova és,
en el fons, la diagonal de Cantor.

\begin{teorema}[König]\index{König, teorema de}\label{teo:konig}
Suposem l'elecció. Siguin \((A_i)_{i\in I}\) i \((B_i)_{i\in I}\) dues famílies
de conjunts, indexades pel mateix \(I\), amb \(A_i\prec B_i\) per a cada \(i\).
Aleshores
\[
\sum_{i\in I}\abs{A_i}\ <\ \prod_{i\in I}\abs{B_i},
\]
on la suma és el cardinal de la suma disjunta \(\coprod_i A_i\) i el producte
el del producte \(\prod_i B_i\) (capítol~\ref{cap:families}).\par
\emph{Axiomes usats:} elecció (axioma~\ref{ax:eleccio}), famílies, suma i
producte (capítol~\ref{cap:families}), la diagonal de Cantor
(teorema~\ref{teo:diagonal}).
\end{teorema}

\begin{prova}
L'elecció s'usa en tres llocs, que assenyalem: per triar alhora un
monomorfisme \(m_i\) i un punt fora de la seva imatge per a cada \(i\); per
deduir de \(F_i\) exhaustiva que \(B_i\preceq A_i\) (una secció); i per triar
alhora els punts \(d_i\) de la diagonal. Les tres tries són seccions de
famílies internes indexades per \(I\), construïdes amb exponencials
(capítol~\ref{cap:families}).

\emph{\(\le\).} Hi ha un monomorfisme \(\coprod_i A_i\hookrightarrow\prod_i B_i\):
com que \(A_i\prec B_i\), en particular \(A_i\preceq B_i\), i triant (per
l'elecció) un monomorfisme \(m_i\colon A_i\hookrightarrow B_i\) i un punt
\(b_i\in B_i\setminus\operatorname{Im}(m_i)\) ---no buit perquè \(m_i\) no és
exhaustiu, ja que \(A_i\not\approx B_i\)---, enviem un element \(a\) de la
fibra \(A_i\) a la família \((x_j)_{j\in I}\) amb \(x_i=m_i(a)\) i \(x_j=b_j\)
per a \(j\neq i\). És injectiva perquè la coordenada \(i\) i la posició del
\enquote{no-\(b\)} determinen \(a\) i \(i\).

\emph{\(\neq\).} Vegem que cap aplicació \(F\colon\coprod_i A_i\to\prod_i B_i\)
no és exhaustiva ---la diagonal. Per a cada \(i\), la restricció de \(F\) a la
fibra \(A_i\), seguida de la projecció \(\pi_i\colon\prod_j B_j\to B_i\), dona
una aplicació \(F_i\colon A_i\to B_i\); com que \(A_i\prec B_i\), \(F_i\) no és
exhaustiva (si ho fos, una secció seva, que existeix per l'elecció, donaria
\(B_i\preceq A_i\), i per Schröder--Bernstein \(A_i\approx B_i\)), de manera que existeix
\(d_i\in B_i\setminus\operatorname{Im}(F_i)\) ---triem-ne un per l'elecció. La
família \((d_i)_{i\in I}\in\prod_i B_i\) no és cap valor de \(F\): si fos
\(F(a)=(d_i)_i\) amb \(a\) a la fibra \(A_{i_0}\), aleshores la coordenada
\(i_0\) donaria \(F_{i_0}(a)=d_{i_0}\), contra \(d_{i_0}\notin\operatorname{Im}(F_{i_0})\).
Per tant \(F\) no és exhaustiva, cap \(\coprod_i A_i\to\prod_i B_i\) no ho és,
i \(\sum_i\abs{A_i}\neq\prod_i\abs{B_i}\). Amb la desigualtat \(\le\), s'obté
l'estricta.
\end{prova}

\begin{observacio}[König conté Cantor]\label{obs:konig-cantor}
El teorema de Cantor és el cas particular de König amb \(I=A\), \(A_i=\mathbf1\)
i \(B_i=\mathbf2\) per a tot \(i\): aleshores \(A_i=\mathbf1\prec\mathbf2=B_i\),
i König dona
\[
\abs{A}=\sum_{i\in A}1<\prod_{i\in A}2=2^{\abs{A}},
\]
que és exactament \(\abs{A}<2^{\abs{A}}\) (teorema~\ref{teo:cantor-cardinal}).
La \enquote{idea diagonal} que apareix a la prova de König ---construir una
família \((d_i)\) que difereix de cada valor en alguna coordenada--- és
literalment la diagonal de Cantor, ara repartida coordenada a coordenada. Per
això König es pot veure com la forma \emph{indexada} del teorema de Cantor.
\end{observacio}

\begin{resumcap}
\apartat{Idees centrals}
\begin{enumerate}
  \item El cardinal \(\abs{A}\) és la classe del conjunt \(A\) llevat d'equipotència (no un tipus d'ordre: no hi ha ordre); suma, producte i exponenciació són els cardinals de \(A+B\), \(A\times B\) i \(B^A\), i les lleis són isomorfismes canònics. La monotonia en l'exponent demana base no nul·la: \(0^0=1\not\le0=0^1\).
  \item Cantor: \(\abs{A}<2^{\abs{A}}\). L'exponenciació cardinal \(2^{\aleph_0}\) no té res a veure amb l'ordinal \(2^\omega=\omega\).
  \item Sense elecció, \(\preceq\) és un preordre antisimètric mòdul \(\approx\); la totalitat és equivalent a l'elecció (Hartogs fa el pas de comparabilitat a bon ordre).
  \item Amb l'elecció, el cardinal inicial d'\(A\) és el menor tipus d'ordre de \(\Hart(A)\) el suport d'un representant del qual és equipotent a \(A\); és un tipus, no un conjunt, i triar-ne un representant és una tria més.
  \item König: si \(A_i\prec B_i\) per a tot \(i\), \(\sum\abs{A_i}<\prod\abs{B_i}\). La prova és una diagonal coordenada a coordenada, i usa l'elecció tres vegades. Conté Cantor com a cas particular.
\end{enumerate}
\apartat{Errors típics} Escriure \(\abs{A}\) com si fos un conjunt fix; aplicar la monotonia en l'exponent amb base \(0\); confondre \(2^\omega\) ordinal amb \(2^{\aleph_0}\) cardinal; creure que la comparabilitat de cardinals és gratuïta; i passar de \enquote{exhaustiva} a \enquote{domina} sense elecció.
\apartat{Lectura recomanada} \cite[§§22--25]{halmos} per a l'aritmètica cardinal clàssica; \cite[cap.~3 i~5]{jech2003} per als cardinals, König i el paper de l'elecció; \cite[cap.~I.10]{kunen2011} per als cardinals inicials; i \cite[§§8.5--8.6]{lawvere-rosebrugh} per a la mida dels conjunts abstractes.
\end{resumcap}

\chapter{Alephs, cofinalitat i continu}\label{cap:alephs}

\begin{objectius}
\apartat{Prerequisits} Del capítol~\ref{cap:cardinals}: aritmètica cardinal, comparabilitat amb l'elecció, cardinals inicials i König. Del capítol~\ref{cap:zorn}: Hartogs. Dels capítols~\ref{cap:ordinals} i~\ref{cap:recursio}: tipus ordinals, suprems acotats i recursió transfinita amb codomini fix. Del capítol~\ref{cap:families}: famílies internes i regles externes. Del capítol~\ref{cap:infinitud}: \(\N\cong\N\times\N\).
\apartat{Objectius} En acabar el capítol, el lector hauria de poder:
\begin{itemize}
  \item enunciar l'esquema SR, explicar la condició d'isomorfisme únic i distingir SR de l'elecció;
  \item demostrar, al Nivell~I, que \(\kappa^{+}=\abs{\Hart(K)}\), i situar el pas de cada \(\kappa\) al seu successor fora de SR;
  \item construir, sota SR, la família \((\aleph_\beta)_{\beta<\lambda}\) per a un \(\lambda\) fix, identificant la instància de l'esquema, i demostrar que tot cardinal infinit és un aleph;
  \item definir cofinalitat, regularitat i singularitat, i justificar els exemples \(\aleph_0\), \(\aleph_{\alpha+1}\) i \(\aleph_\omega\), sabent quin resultat clàssic se cita;
  \item definir els beths i el continu, i formular HC i HCG, separant el que és de Nivell~I del que demana SR;
  \item explicar en quin sentit HC és independent, i sota quina hipòtesi de consistència.
\end{itemize}
\end{objectius}

Aquest capítol és opcional i de \emph{Nivell~II}: a diferència de tots els
anteriors, fa servir un principi més fort que els axiomes bàsics del llibre
---la substitució estructural SR, fins ara només anunciada
(observació~\ref{obs:sr})---, i el declara explícitament a cada construcció que
el necessita. La raó és estructural i ja la coneixem: reunir en un sol objecte
una successió transfinita de cardinals que \emph{creixen sense cota comuna}
---els alephs \(\aleph_\alpha\), els beths \(\beth_\alpha\)--- és exactament el
tipus de regla externa que els axiomes de Nivell~I no poden representar
(observació~\ref{obs:frontera-exemple}). Amb SR, sí. El capítol recorre la
jerarquia dels cardinals infinits: \(\aleph_0\) i el cardinal successor
\(\kappa^{+}\) (que encara és de Nivell~I, gràcies a Hartogs), la successió
\(\alpha\mapsto\aleph_\alpha\) (ja de Nivell~II), la cofinalitat i la distinció
regular\slash singular, els nombres \(\beth\), i el continu
\(\mathfrak c=2^{\aleph_0}\). Es tanca amb la hipòtesi del continu, presentada
com el que és ---una qüestió \emph{independent} dels axiomes--- i no com un
problema pendent de solució elemental. No entrem en forcing ni en grans
cardinals; l'objectiu és formular amb precisió què pregunta la hipòtesi del
continu i per què la teoria desenvolupada permet plantejar-la.

\section{La substitució estructural com a esquema adoptat}\label{sec:sr-adoptat}

Fins ara SR ha estat un \emph{esquema provisional}
(observació~\ref{obs:sr}): l'enunciàvem per assenyalar on faria falta, però no
l'assumíem. En aquest capítol el prenem com a principi de treball, en la forma
d'esquema, i el marcarem a la caixa d'axiomes de cada resultat que en depengui.

\begin{principi}[Esquema de reunió de famílies (substitució estructural, SR)]\label{sr:esquema}
Sigui \(\varphi(i,X,R,\vec p)\) una fórmula del llenguatge del llibre en què
\(i\) és una variable d'element, \(X\) una variable de conjunt, \(R\) una
variable d'\emph{estructura sobre \(X\)} (una part de \(X\), de \(X\times X\),
una aplicació amb \(X\) com a domini o codomini, o una llista finita
d'estructures d'aquests tipus) i \(\vec p\) una llista de
paràmetres fixats. Suposem que, per a un conjunt \(I\) donat, \(\varphi\) és
\emph{funcional a isomorfisme únic}: per a cada element \(i\colon\mathbf1\to I\),
\begin{enumerate}
  \item existeix algun parell \((X,R)\) tal que \(\varphi(i,X,R,\vec p)\), i
  \item si \(\varphi(i,X,R,\vec p)\) i \(\varphi(i,X',R',\vec p)\), hi ha un
    \emph{únic} isomorfisme \(u\colon X\to X'\) que transporta \(R\) en \(R'\).
\end{enumerate}
Aleshores existeixen una aplicació \(p\colon E\to I\) i una estructura
\(S\) sobre \(E\) que es restringeix a cada fibra, tals que, per a tot
\(i\colon\mathbf1\to I\), la fibra \(E_i=p^{-1}(i)\) amb l'estructura
restringida \(S_i\) satisfà \(\varphi(i,E_i,S_i,\vec p)\).
\end{principi}

En paraules: una regla externa \(i\mapsto(X_i,R_i)\) que determina cada
\((X_i,R_i)\) llevat d'un \emph{únic} isomorfisme es representa per una família
interna \(E\to I\). La unicitat de l'isomorfisme a la clàusula~2 és essencial, i
és el que diferencia \enquote{determinat llevat d'isomorfisme únic} de
\enquote{determinat llevat d'isomorfisme}: un conjunt infinit determinat només
llevat d'isomorfisme té molts automorfismes, i no hi ha cap manera canònica
d'enganxar les fibres; un bon ordre, o un cardinal inicial amb el seu ordre, en
té exactament un, la identitat, i les fibres s'enganxen sense ambigüitat.

\begin{observacio}[SR és un esquema, no un axioma més]\label{obs:sr-esquema-opcional}
Convé ser precisos sobre l'estatus de SR, perquè no és del mateix tipus lògic
que els axiomes que hem fet servir fins ara. L'\emph{elecció}
(axioma~\ref{ax:eleccio}) o l'\emph{efectivitat} (axioma~\ref{ax:exactitud})
són \emph{sentències}: cadascuna, un únic enunciat afegit a la teoria. SR, en
canvi, és un \emph{esquema} ---una família d'enunciats, una instància per a
cada fórmula funcional i invariant per isomorfisme---, de la mateixa manera que
el reemplaçament de la teoria clàssica és un esquema i no una sentència. Per
això no el presentem com un axioma més a l'estil de l'elecció, sinó com un
\emph{principi} a part, i per això, a la caixa de dependències, \enquote{SR} no
figura com una hipòtesi al mateix nivell que \enquote{elecció}: assenyala que
s'ha usat una instància de l'esquema de reunió, i indicarem \emph{quina} cada
vegada. SR és, a més, \emph{opcional}: no forma part dels axiomes de Nivell~I
del llibre, i tots els capítols anteriors en són independents.
\end{observacio}

\begin{observacio}[Què afegeix SR, i què no]\label{obs:sr-abast}
SR converteix una \emph{regla externa} \(i\mapsto X_i\) en una \emph{família
interna} \(E\to I\) (definició~\ref{def:familia-vs-regla}), sempre que la regla
sigui uniforme i invariant. És exactament el que mancava per reunir famílies no
acotades: amb SR, l'objecte total \(E\) existeix, i amb ell la suma
\(\coprod_i X_i\), el suprem d'una família de tipus no prèviament acotada, i les
successions transfinites de conjunts variables. El que SR \emph{no} és ---cal
repetir-ho--- és una forma d'elecció (observació~\ref{obs:sr-no-eleccio}):
l'elecció tria dins fibres d'una família \emph{ja existent}; SR fa \emph{existir}
la família. En aquest capítol usarem tots dos, i els marcarem per separat.
\end{observacio}

\begin{observacio}[Quina formulació adoptem, i què en citem]\label{obs:sr-rigor}
El principi~\ref{sr:esquema} és la versió, en el llenguatge d'aquest llibre, del
\emph{reemplaçament estructural} que Shulman formula i estudia a
\cite{shulman-comparing}: una instància de l'esquema per a cada fórmula
\(\varphi\) del llenguatge, amb la condició de funcionalitat a isomorfisme únic
i amb els paràmetres \(\vec p\) i el conjunt d'índexs \(I\) com a dades fixades.
Tres precisions. Primera, l'\emph{estructura} \(R\) és la que fa que l'isomorfisme
de la clàusula~2 sigui únic: en les instàncies que usarem, \(R\) serà sempre un bon ordre sobre \(X\), o una aplicació
\(X\to I'\) amb un bon ordre sobre cada fibra, de manera que la unicitat és
la rigidesa dels bons ordres
(teorema~\ref{teo:unicitat-iso}). Segona, la formulació de McLarty a
\cite{mclarty-catset} demana només \enquote{llevat d'isomorfisme}, sense
unicitat; en presència de l'elecció (axioma~\ref{ax:eleccio}) les dues versions
són equivalents, perquè l'elecció permet fixar representants i, un cop fixats,
l'esquema amb unicitat reconstrueix la família; aquí preferim la versió amb
unicitat perquè no barreja SR amb l'elecció. Tercera, SR és un esquema sobre
fórmules del llenguatge, no sobre \enquote{regles} arbitràries del
metallenguatge: una regla \(i\mapsto X_i\) que no es pugui escriure com una
fórmula \(\varphi\) no en dona cap instància.

Amb aquesta formulació, l'afirmació que ETCS${}+{}$SR \enquote{té la força de
ZFC} és un \emph{teorema citat}, no un resultat del llibre: Osius
\cite{osius1974} va demostrar que ETCS més reemplaçament és equivalent a ZFC,
i Shulman \cite{shulman-comparing} en dona la forma precisa com a
bi-interpretabilitat entre ETCS${}+{}$reemplaçament estructural i ZFC. El
que sí que fem aquí és marcar, a cada resultat que en depèn, \emph{quina}
instància de l'esquema s'usa (quina \(\varphi\), quin \(I\), quina estructura
\(R\)), de manera que el lector pugui comprovar que compleix les dues clàusules.
\end{observacio}

\begin{lema}[Separació completa sota SR]\label{lema:separacio-sr}
Suposem SR. Sigui \(\psi(i,\vec p)\) una fórmula qualsevol del llenguatge del
llibre ---amb quantificadors no acotats, si cal--- i \(I\) un conjunt. Aleshores
existeix una part \(m\colon S\hookrightarrow I\) tal que, per a tot punt
\(i\colon\mathbf1\to I\), \(i\) es factoritza per \(m\) si i només si
\(\psi(i,\vec p)\). En particular, la inducció transfinita
(teorema~\ref{teo:induccio-transfinita}) s'aplica a propietats
\(\psi\) dels punts d'un bon ordre definides per fórmules arbitràries.\par
\emph{Axiomes usats:} SR, amb la instància indicada a la prova.
\end{lema}

\begin{prova}
Considerem la fórmula
\[
\varphi(i,X,R,\vec p)\ :\equiv\
\bigl(\psi(i,\vec p)\land\operatorname{Terminal}(X)\bigr)\lor
\bigl(\neg\psi(i,\vec p)\land\operatorname{Inicial}(X)\bigr),
\]
amb \(R:=\id_X\) com a estructura trivial. És funcional a isomorfisme únic:
per a cada \(i\) es dona \(\psi(i,\vec p)\) o la seva negació, i hi ha un \(X\)
adequat (\(\mathbf1\) o \(\mathbf0\)); i dos \(X,X'\) que compleixen
\(\varphi\) per al mateix \(i\) són tots dos terminals o tots dos inicials, de
manera que hi ha un únic isomorfisme entre ells
(proposicions~\ref{prop:terminal-unic} i~\ref{prop:inicial-unic}). SR dona una
família \(p\colon E\to I\) amb fibra terminal sobre els \(i\) que compleixen
\(\psi\) i inicial sobre els altres. L'aplicació \(p\) és un monomorfisme: és
injectiva sobre punts, perquè dos punts de \(E\) amb la mateixa imatge \(i\)
són a la fibra \(E_i\), que és terminal o inicial, i per tant coincideixen
(observació~\ref{obs:injectiva-mono}). Un punt \(i\) es factoritza per \(p\)
si i només si la fibra \(E_i\) té algun punt, és a dir, si i només si és
terminal, és a dir, si i només si \(\psi(i,\vec p)\). Prenem \(m:=p\). La
darrera afirmació és immediata, perquè el teorema~\ref{teo:induccio-transfinita}
només demana que la propietat defineixi una part i la hipòtesi de
progressivitat sobre punts.
\end{prova}

\section{\texorpdfstring{\(\aleph_0\)}{aleph 0} i el cardinal successor}\label{sec:aleph0-successor}

Comencem pel més petit dels cardinals infinits i per l'operació que fa pujar un
graó. Aquesta secció és, de fet, de \emph{Nivell~I}: no necessita SR, perquè el
cardinal successor està acotat per Hartogs.

\begin{definicio}[\(\aleph_0\)]\index{aleph zero@\(\aleph_0\)}\label{def:aleph0}
\(\aleph_0:=\abs{\N}\), el cardinal del conjunt dels naturals ---el menor
cardinal infinit, perquè tot conjunt infinit conté una còpia de \(\N\)
(capítol~\ref{cap:infinitud}) i tot enumerable infinit és equipotent a \(\N\)
(corol·lari~\ref{cor:enumerable-infinit-N}). Com a \emph{tipus d'ordre}, el
mateix conjunt és \(\omega\); com a \emph{cardinal}, és \(\aleph_0\). Escrivim
\(\aleph_0\) quan pensem la mida i \(\omega\) quan pensem la forma d'ordre.
\end{definicio}

\begin{definicio}[Cardinal successor]\index{Cardinal successor}\label{def:cardinal-successor}
Sigui \(\kappa\) un cardinal, representat per un conjunt \(K\). El \emph{cardinal
successor} \(\kappa^{+}\) és el menor cardinal estrictament més gran que
\(\kappa\): el menor cardinal que \emph{no} s'injecta en \(K\).
\end{definicio}

\begin{proposicio}[Existència del cardinal successor; Nivell I]\label{prop:cardinal-successor}
Per a tot cardinal \(\kappa=\abs{K}\) existeix \(\kappa^{+}\), i és
exactament el cardinal del bon ordre de Hartogs:
\[
\kappa^{+}=\abs{\Hart(K)}.
\]
A més, \(\Hart(K)\) és un bon ordre inicial, de manera que \(\kappa^{+}\) és un
cardinal inicial. L'existència \emph{no} usa SR.\par
\emph{Axiomes usats:} Hartogs (teorema~\ref{teo:hartogs},
proposició~\ref{prop:hartogs-bon-ordre}), teorema del bon ordre
(teorema~\ref{teo:bon-ordre}), comparabilitat de bons ordres
(teorema~\ref{teo:comparabilitat}); elecció (només a través del teorema del
bon ordre); \emph{cap SR}.
\end{proposicio}

\begin{prova}
Escrivim \(H=\Hart(K)\). Hem de veure dues coses: que \(\abs{H}>\kappa\), i que
tot conjunt \(X\) amb \(\abs{X}>\kappa\) satisfà \(\abs{H}\le\abs{X}\).

La primera és el teorema de Hartogs: \(H\not\preceq K\), és a dir,
\(\abs{H}\not\le\kappa\), i per la comparabilitat dels cardinals
(teorema~\ref{teo:comparabilitat-cardinals}) \(\abs{H}>\kappa\).

Per a la segona, sigui \(X\) un conjunt amb \(\abs{X}>\kappa\), és a dir,
\(X\not\preceq K\). Pel teorema del bon ordre, \(X\) admet un bon ordre; sigui
\(W\) aquest bon ordre sobre \(X\). Per la comparabilitat dels bons ordres,
o bé \(W\) és isomorf a un segment inicial propi de \(H\), o bé \(H\) és
isomorf a un segment inicial (propi o no) de \(W\). El primer cas és impossible:
el segment inicial propi de \(H\) determinat per un tipus \(\alpha=[(S,R)]\) té
precisament el tipus \(\alpha\) (lema~\ref{lema:segment-tipus}), és a dir, és
isomorf a un bon ordre \((S,R)\) sobre una part \(S\) de \(K\); així \(W\)
s'injectaria en \(K\) i tindríem \(X\preceq K\), contra la hipòtesi. Queda el segon cas, que dona
precisament un monomorfisme \(H\hookrightarrow X\), és a dir, \(\abs{H}\le\abs{X}\).

Finalment, \(H\) és inicial: un bon ordre \(V\) de tipus menor que el de
\(H\) és isomorf a un segment inicial propi de \(H\), que pel mateix
lema~\ref{lema:segment-tipus} s'injecta en \(K\); si \(V\) fos equipotent a
\(H\) tindríem \(H\approx V\preceq K\), contra Hartogs. Tota la cerca ha quedat dins de \(H\), un conjunt ja construït;
per això no cal SR.
\end{prova}

\begin{observacio}[Per què el successor és Hartogs]\label{obs:successor-hartogs}
La igualtat \(\kappa^{+}=\abs{\Hart(K)}\) és la lectura cardinal del teorema de
Hartogs: \(\Hart(K)\) és el menor bon ordre que no cap dins de \(K\), i amb
l'elecció, que fa que tot conjunt sigui ben ordenable, \enquote{el menor bon
ordre que no cap} i \enquote{el menor conjunt que no cap} coincideixen. Sense
elecció, \(\abs{\Hart(K)}\) continua essent el menor cardinal \emph{ben
ordenable} que no és \(\le\kappa\), però pot haver-hi conjunts \(X\) amb
\(\abs{X}>\kappa\) en què \(\Hart(K)\) no s'injecta, i aleshores
\enquote{el menor cardinal \(>\kappa\)} deixa de tenir sentit.
\end{observacio}

\begin{observacio}[On encara no cal SR]\label{obs:successor-sense-sr}
Observem que \(\aleph_0\) i el pas \(\kappa\mapsto\kappa^{+}\) són de
Nivell~I: cada \(\kappa^{+}\) individual s'obté acotadament per Hartogs. SR
només farà falta quan vulguem reunir un \emph{tram} de la successió
\(\aleph_0,\aleph_1,\aleph_2,\dots\) ---el que va fins a un ordinal fix
\(\lambda\)--- en una família interna i prolongar-lo als límits; és a dir,
quan passem d'un pas a una família. És la mateixa frontera de sempre: una
operació binària o un successor són acotats; una successió transfinita no
acotada demana SR.
\end{observacio}

\section{La successió dels alephs}\label{sec:alephs}

Ara sí, reunim els cardinals infinits en successions. Un punt de rigor
important: la col·lecció de \emph{tots} els tipus ordinals és pròpia, no un
conjunt, i per tant no pot ser el conjunt d'índexs d'una família interna
\(E\to I\). Per això no construïm cap objecte intern
\enquote{\(\alpha\mapsto\aleph_\alpha\) per a tots els \(\alpha\)}; construïm,
per a cada bon ordre \emph{fix} \(\lambda\), el tram
\((\aleph_\beta)_{\beta<\lambda}\) indexada pel \emph{conjunt} \(\lambda\).
La notació global \(\alpha\mapsto\aleph_\alpha\) és aleshores una regla del
metallenguatge (una funció-classe), no un objecte del sistema. Aquest és el
primer ús essencial de SR.

\begin{definicio}[Els alephs fins a un ordinal fix]\index{Els alephs fins a un ordinal fix}\label{def:alephs}
Sigui \(\lambda\) un bon ordre (un tipus ordinal fix). La família
\((\aleph_\beta)_{\beta<\lambda}\), indexada pel conjunt \(\lambda\), es
defineix per recursió transfinita sobre \(\lambda\)
(teorema~\ref{teo:recursio-transfinita}):
\[
\aleph_0=\abs{\N},
\qquad
\aleph_{\beta+1}=(\aleph_\beta)^{+},
\qquad
\aleph_\mu=\sup_{\beta<\mu}\aleph_\beta\quad(\mu<\lambda\ \text{límit}).
\]
Els casos zero i successor són de Nivell~I (definició~\ref{def:cardinal-successor});
la clàusula límit \emph{i} la reunió del tram en una família indexada pel
conjunt \(\lambda\) requereixen SR, amb la instància que es descriu a la
prova de la proposició~\ref{prop:alephs}. La notació
\(\alpha\mapsto\aleph_\alpha\) sense cota superior és només la regla de
classe subjacent, no un conjunt.
\end{definicio}

\begin{proposicio}[Els alephs, sota SR]\label{prop:alephs}
Suposem SR. Aleshores:
\begin{enumerate}
\item per a cada bon ordre fix \(\lambda\), la família
\((\aleph_\beta)_{\beta<\lambda}\) existeix com a família interna indexada pel
\emph{conjunt} \(\lambda\);
\item és estrictament creixent: \(\beta<\gamma<\lambda\) implica
\(\aleph_\beta<\aleph_\gamma\);
\item \emph{tot} cardinal infinit és un aleph, en sentit metateòric: per a
tot conjunt infinit \(A\) existeix algun tipus ordinal \(\alpha\) amb
\(\abs{A}=\aleph_\alpha\), i aquest \(\alpha\) és únic. (No s'afirma cap
objecte intern que contingui tots els \(\aleph_\alpha\) alhora; l'enunciat és
un esquema sobre \(A\), no una igualtat entre famílies.)
\end{enumerate}
\emph{Axiomes usats:} SR (principi~\ref{sr:esquema}) per a la família sobre
\(\lambda\), amb la instància \enquote{\((X,R)\) és una solució fins a
\(\gamma\)} als límits; elecció (proposició~\ref{prop:cardinal-inicial}) per
a~(3); Hartogs per als successors.
\end{proposicio}

\begin{prova}
\emph{(1)} Fixem un representant \(L\) de \(\lambda\). Anomenem \emph{solució
fins a \(\beta\)} (per a \(\beta\in L\)) una família \(q\colon X\to L_{\le\beta}\)
amb un bon ordre sobre cada fibra (codificat com una part de \(X\times X\)),
tal que la fibra sobre \(0\) té tipus \(\aleph_0\), la fibra sobre
\(\gamma+1\) té per tipus el successor cardinal del de la fibra sobre
\(\gamma\), i la fibra sobre un límit \(\mu\) té per tipus el suprem dels de
les fibres sobre \(\gamma<\mu\). Les fibres són bons ordres inicials, de
manera que dues solucions fins a \(\beta\) estan lligades per un \emph{únic}
isomorfisme compatible (rigidesa, teorema~\ref{teo:unicitat-iso}, fibra a
fibra). Per inducció transfinita sobre \(\beta\) es veu que n'hi ha una per
a cada \(\beta\in L\); la propietat \enquote{existeix una solució fins a
\(\beta\)} quantifica sobre conjunts arbitraris i no és, en principi, una part
de \(L\), però ho és pel lema~\ref{lema:separacio-sr}, de manera que la
inducció és legítima. Pas a pas: al successor s'afegeix una fibra
(proposició~\ref{prop:cardinal-successor}, Nivell~I); al límit \(\mu\),
l'esquema SR (principi~\ref{sr:esquema}) aplicat al conjunt d'índexs
\(L_{<\mu}\) i a la fórmula \(\varphi(\gamma,X,R)\): \enquote{\((X,R)\) és
una solució fins a \(\gamma\)}, amb \(R\) la llista formada per la
projecció i l'ordre de les fibres, dona una família \(F\to L_{<\mu}\) la fibra \(F_\gamma\) de la qual és una
solució fins a \(\gamma\). Aquestes solucions són compatibles en el sentit
següent. Si \(\gamma\le\delta<\mu\), la \emph{restricció} de la solució
\(F_\delta\) a \(L_{\le\gamma}\) ---la imatge inversa de \(L_{\le\gamma}\) per
la seva projecció, amb els ordres de les fibres restringits--- és de nou una
solució fins a \(\gamma\), perquè les condicions de la recursió sobre
\(\gamma'\le\gamma\) només miren fibres sobre \(L_{\le\gamma}\). Per la
clàusula d'unicitat de SR hi ha, doncs, un \emph{únic} isomorfisme
\(u_{\gamma\delta}\) de \(F_\gamma\) sobre aquesta restricció que commuta amb
les projeccions cap a \(L_{\le\gamma}\) i conserva l'ordre de cada fibra; i
per la unicitat, aquests isomorfismes són coherents
(\(u_{\gamma\varepsilon}=u_{\delta\varepsilon}|\circ u_{\gamma\delta}\) per
a \(\gamma\le\delta\le\varepsilon\)), perquè els dos costats són isomorfismes
del mateix tipus entre les mateixes solucions. Prenem ara, per a cada
\(\gamma<\mu\), la fibra de \(F_\gamma\) sobre el punt final \(\gamma\)
(la \enquote{diagonal}): és la part de \(F\) dels elements que es projecten
a \((\gamma,\gamma)\), i com a \(\gamma\) varia forma una família interna
\(D\to L_{<\mu}\), tallada dins \(F\) per una condició sobre les dues
projeccions. Per la coherència, la restricció de \(D\) a \(L_{\le\gamma}\) és
isomorfa a \(F_\gamma\) mateixa, de manera que \(D\) satisfà la recursió a
tots els \(\gamma<\mu\). El suprem d'aquests tipus és el
tipus del límit directe dels representants (cada fibra s'immergeix com a
part inicial de les posteriors, de manera única), per l'argument de la
proposició~\ref{prop:suprem}; i és un cardinal inicial, perquè un tipus
menor és menor que algun \(\aleph_\gamma\) i, per tant, de mida
\(\le\aleph_\gamma<\aleph_{\gamma+1}\le\) la del suprem. Afegint aquesta
fibra s'obté la solució fins a \(\mu\). Una darrera aplicació del mateix
argument sobre tot \(L\) dona la família \(E\to\lambda\).

\emph{(2)} Per inducció transfinita (teorema~\ref{teo:induccio-transfinita}):
al successor, \(\aleph_{\gamma+1}=(\aleph_\gamma)^{+}>\aleph_\gamma\); al
límit \(\mu\), \(\aleph_\mu\) és cota superior de tots els \(\aleph_\gamma\),
\(\gamma<\mu\), i per a cada \(\gamma<\mu\) també \(\gamma+1<\mu\), de manera
que \(\aleph_\mu\ge\aleph_{\gamma+1}>\aleph_\gamma\).

\emph{(3)} Per inducció, \(\beta\le\aleph_\beta\) com a tipus (al successor,
\(\aleph_{\beta+1}>\aleph_\beta\ge\beta\), i és un tipus estrictament més gran;
al límit, el suprem domina cada terme). Sigui \(A\) infinit i \(\kappa\) el
seu cardinal inicial (proposició~\ref{prop:cardinal-inicial}, amb elecció).
Prenem \(\lambda=\kappa+1\) i la família de (1). Com que
\(\aleph_\kappa\ge\kappa\), la part dels \(\beta\le\kappa\) amb
\(\aleph_\beta\ge\kappa\) és no buida; sigui \(\alpha\) el seu mínim. Si
\(\alpha=0\), \(\kappa\le\aleph_0\) i, com que \(A\) és infinit,
\(\kappa=\aleph_0\). Si \(\alpha=\gamma+1\), \(\aleph_\gamma<\kappa\), d'on
\(\aleph_\alpha=(\aleph_\gamma)^{+}\le\kappa\le\aleph_\alpha\). Si \(\alpha\)
és límit, \(\aleph_\gamma<\kappa\) per a tot \(\gamma<\alpha\), d'on
\(\aleph_\alpha=\sup_{\gamma<\alpha}\aleph_\gamma\le\kappa\le\aleph_\alpha\). En
tots els casos \(\aleph_\alpha=\kappa=\abs{A}\). La unicitat és la
monotonia estricta de (2).
\end{prova}

\begin{observacio}[Lectura clàssica]\label{obs:alephs-classica}
Els alephs enumeren \emph{tots} els cardinals infinits en ordre creixent:
\(\aleph_0<\aleph_1<\aleph_2<\cdots<\aleph_\omega<\aleph_{\omega+1}<\cdots\),
un per a cada tipus ordinal. La part (3) ---que no hi ha cap cardinal infinit
\enquote{entre els alephs}--- és el que fa que la jerarquia sigui exhaustiva, i
depèn de l'elecció (via el cardinal inicial). Sense elecció, hi pot haver
conjunts infinits que no siguin equipotents a cap aleph. Els trams de la
successió, com a objectes, són el primer lloc del llibre on SR és
inevitable: cada \(\aleph_\alpha\) individual existeix sense SR, però un tram
\((\aleph_\beta)_{\beta<\lambda}\) com a família interna, no. La família de
\emph{tots} els alephs no existeix com a objecte ni tan sols amb SR: els
tipus ordinals no formen cap conjunt d'índexs (Burali--Forti).
\end{observacio}

\section{Cofinalitat; cardinals regulars i singulars}\label{sec:cofinalitat}

La cofinalitat mesura \enquote{amb quina rapidesa} es pot arribar al final d'un
bon ordre, i distingeix dues menes de cardinals límit.

\begin{definicio}[Cofinalitat]\index{Cofinalitat}\label{def:cofinalitat}
Sigui \((W,<)\) un bon ordre sense element màxim. Una part \(C\hookrightarrow W\)
és \emph{cofinal} si per a tot \(w\in W\) existeix \(c\in C\) amb \(w\le c\)
---\(C\) \enquote{arriba fins al final} de \(W\). La \emph{cofinalitat}
\(\operatorname{cf}(W)\) és el menor tipus ordinal d'una part cofinal de \(W\),
ben ordenada per l'ordre restringit. Per a un cardinal límit \(\kappa\),
\(\operatorname{cf}(\kappa)\) és la cofinalitat d'un bon ordre inicial que el
representa.
\end{definicio}

\begin{definicio}[Regular i singular]\index{Regular i singular}\label{def:regular-singular}
Un cardinal infinit \(\kappa\) és \emph{regular} si \(\operatorname{cf}(\kappa)=\kappa\)
---no es pot arribar al seu final amb menys de \(\kappa\) passos---, i
\emph{singular} si \(\operatorname{cf}(\kappa)<\kappa\).
\end{definicio}

\begin{teorema}[Hessenberg; citat]\label{teo:hessenberg}
Per a tot cardinal infinit \(\kappa\), \(\kappa\cdot\kappa=\kappa\). En
conseqüència, amb l'elecció, la reunió de \(\le\kappa\) conjunts de cardinal
\(\le\kappa\) té cardinal \(\le\kappa\).\par
\emph{No el demostrem}: la prova clàssica bon-ordena \(\kappa\times\kappa\)
per l'ordre de Gödel (primer pel màxim de les dues coordenades) i és de
Nivell~I amb l'elecció, però llarga; vegeu~\cite[teorema~3.5]{jech2003}. El
cas \(\kappa=\aleph_0\) és \(\N\cong\N\times\N\), que sí que hem demostrat
(proposició~\ref{prop:estabilitat-enumerable}).
\end{teorema}

\begin{proposicio}[Exemples de regularitat, sota elecció i SR]\label{prop:regular-exemples}
Suposem elecció i SR.
\begin{enumerate}
\item \(\aleph_0\) és regular: cap part finita de \(\N\) no és cofinal.
\item Tot cardinal successor \(\aleph_{\alpha+1}\) és regular.
\item \(\aleph_\omega\) és singular: \(\operatorname{cf}(\aleph_\omega)=\aleph_0<\aleph_\omega\),
perquè la successió \((\aleph_n)_{n<\omega}\) és cofinal en \(\aleph_\omega\) i
té longitud \(\omega\).
\end{enumerate}
\emph{Axiomes usats:} elecció, SR (per a \(\aleph_\omega\) i la successió dels
alephs); cofinalitat (definició~\ref{def:cofinalitat}).
\end{proposicio}

\begin{prova}
\emph{(1)} Una part finita de \(\N\) té un màxim, de manera que no arriba fins al
final; una part cofinal és, doncs, infinita, i tota part infinita de \(\N\)
és isomorfa a \(\N\) com a ordre (teorema~\ref{teo:subobjectes-N}), de tipus
\(\omega\).
\emph{(2)} Si \(C\) fos una part cofinal de \(\aleph_{\alpha+1}\) de cardinal
\(\le\aleph_\alpha\), aleshores \(\aleph_{\alpha+1}\) seria la unió de \(\le\aleph_\alpha\)
conjunts cadascun de cardinal \(\le\aleph_\alpha\) (els segments determinats pels
elements de \(C\)); però una tal unió té cardinal \(\le\aleph_\alpha\cdot\aleph_\alpha=\aleph_\alpha<\aleph_{\alpha+1}\)
(teorema~\ref{teo:hessenberg}), contradicció.
Per tant tota part cofinal té cardinal \(\aleph_{\alpha+1}\), i el successor és
regular. \emph{(3)} La successió \((\aleph_n)_{n<\omega}\) és estrictament
creixent (proposició~\ref{prop:alephs}) i el seu suprem és \(\aleph_\omega\)
(definició~\ref{def:alephs}); per tant \(\{\aleph_n:n<\omega\}\) és cofinal en
\(\aleph_\omega\) i té tipus \(\omega\), d'on \(\operatorname{cf}(\aleph_\omega)=\aleph_0<\aleph_\omega\).
\end{prova}

\begin{observacio}[Per què importa la cofinalitat]\label{obs:cofinalitat-importa}
La distinció regular\slash singular és la que governa quins cardinals es poden
\enquote{assolir des de baix}. Un cardinal regular no és el suprem de menys
cardinals més petits; un de singular, sí. Aquesta és, també, la raó per la qual
el teorema de König (teorema~\ref{teo:konig}) té conseqüències sobre la
cofinalitat: implica \(\operatorname{cf}(2^{\aleph_0})>\aleph_0\): si \(2^{\aleph_0}\) fos la
reunió d'una successió \((A_n)_{n\in\N}\) de parts de cardinal
\(<2^{\aleph_0}\), König donaria
\(2^{\aleph_0}\le\sum_n\abs{A_n}<\prod_n2^{\aleph_0}=(2^{\aleph_0})^{\aleph_0}=2^{\aleph_0\cdot\aleph_0}=2^{\aleph_0}\),
absurd. És a dir, el continu no pot ser un cardinal de cofinalitat numerable ---una de les
poques coses que la teoria \emph{sí} diu sobre \(2^{\aleph_0}\), a diferència del
seu valor exacte.
\end{observacio}

\section{Els nombres beth i el continu}\label{sec:beths-continu}

Els alephs pugen \enquote{d'un en un} pel cardinal successor; els beths pugen
\enquote{de cop} pel conjunt de parts. La diferència entre les dues jerarquies
és, precisament, el contingut de la hipòtesi del continu.

\begin{definicio}[Els beths fins a un ordinal fix]\index{Els beths fins a un ordinal fix}\label{def:beths}
Com amb els alephs, per a cada bon ordre fix \(\lambda\) la família
\((\beth_\beta)_{\beta<\lambda}\), indexada pel conjunt \(\lambda\), es
defineix per recursió transfinita sobre \(\lambda\), sota SR:
\[
\beth_0=\aleph_0,
\qquad
\beth_{\beta+1}=2^{\beth_\beta},
\qquad
\beth_\mu=\sup_{\beta<\mu}\beth_\beta\quad(\mu<\lambda\ \text{límit}).
\]
La notació global \(\alpha\mapsto\beth_\alpha\) és, de nou, una regla de
classe del metallenguatge, no un objecte intern indexat per tots els
ordinals. El \emph{continu} és
\[
\mathfrak c\ :=\ \abs{P(\N)}\ =\ 2^{\aleph_0}\ =\ \beth_1 ,
\]
i, notablement, ja és de Nivell~I: \(P(\N)\) està construït i \(\beth_1\) és un
sol pas de conjunt de parts, sense reunir cap família no acotada.
\emph{Axiomes usats:} SR (per a la família sobre \(\lambda\) i els suprems
límit), Cantor cardinal (teorema~\ref{teo:cantor-cardinal}) per a l'estricta
pujada.
\end{definicio}

\begin{proposicio}[Els beths pugen estrictament; \(\aleph_\alpha\le\beth_\alpha\)]\label{prop:beths}
Suposem SR. La successió dels beths és estrictament creixent
---\(\beth_\alpha<\beth_{\alpha+1}=2^{\beth_\alpha}\) pel teorema de Cantor
(teorema~\ref{teo:cantor-cardinal})--- i domina els alephs terme a terme:
\(\aleph_\alpha\le\beth_\alpha\) per a tot \(\alpha\), amb igualtat a
\(\alpha=0\).\par
\emph{Axiomes usats:} SR, Cantor cardinal, monotonia de l'exponenciació
cardinal (proposició~\ref{prop:lleis-cardinals}); elecció, a través del
cardinal successor i de la comparabilitat dels cardinals.
\end{proposicio}

\begin{prova}
L'estricta pujada al successor és \(\beth_\alpha<2^{\beth_\alpha}\)
(teorema~\ref{teo:cantor-cardinal}); als límits, el suprem d'una successió
estrictament creixent és més gran que cada terme. Per a
\(\aleph_\alpha\le\beth_\alpha\), inducció transfinita
(teorema~\ref{teo:induccio-transfinita}): a \(\alpha=0\), \(\aleph_0=\beth_0\); al
successor, \(\aleph_{\alpha+1}=(\aleph_\alpha)^{+}\le 2^{\aleph_\alpha}\le
2^{\beth_\alpha}=\beth_{\alpha+1}\) ---la primera desigualtat perquè
\((\aleph_\alpha)^{+}\) és el menor cardinal \(>\aleph_\alpha\) i
\(2^{\aleph_\alpha}>\aleph_\alpha\) (Cantor), la segona per monotonia i la
hipòtesi d'inducció; a un límit \(\mu\),
\(\aleph_\mu=\sup_{\beta<\mu}\aleph_\beta\le\sup_{\beta<\mu}\beth_\beta=\beth_\mu\).
Totes dues induccions són sobre el bon ordre fix que indexa les famílies, amb
propietats que són parts pel lema~\ref{lema:separacio-sr}.
\end{prova}

\section{La hipòtesi del continu}\label{sec:hipotesi-continu}

Amb els alephs i els beths a la mà podem formular amb precisió la pregunta que
Cantor va deixar oberta.

\begin{observacio}[La hipòtesi del continu]\label{obs:hipotesi-continu}
El continu és \(\mathfrak c=2^{\aleph_0}=\beth_1\). Sabem, pel teorema de Cantor,
que \(\mathfrak c>\aleph_0\) (teorema~\ref{teo:cantor-cardinal}); per tant, sota
l'elecció, \(\mathfrak c\ge\aleph_1\), el menor cardinal no numerable. La
\emph{hipòtesi del continu} (HC) és l'afirmació que no hi ha cap cardinal entre
\(\aleph_0\) i \(\mathfrak c\):
\[
2^{\aleph_0}=\aleph_1,
\qquad\text{equivalentment}\qquad
\beth_1=\aleph_1 .
\]
La \emph{hipòtesi del continu generalitzada} (HCG) afirma \(2^{\aleph_\alpha}=\aleph_{\alpha+1}\)
per a tot \(\alpha\), és a dir \(\beth_\alpha=\aleph_\alpha\) per a tot \(\alpha\):
que les dues jerarquies, la dels successors i la dels conjunts de parts,
coincideixen.
\end{observacio}

\begin{observacio}[HC no necessita SR]\label{obs:hc-sense-sr}
Remarquem que la sola \emph{formulació} d'HC ---i la seva veritat o
falsedat--- no depèn de l'esquema SR, tot i que aquest capítol és, en conjunt,
de Nivell~II. En efecte, els dos membres de \(2^{\aleph_0}=\aleph_1\) existeixen
ja al Nivell~I: el continu és \(\mathfrak c=\abs{P(\N)}\), i \(P(\N)\) està
construït pel conjunt de parts; i \(\aleph_1=(\aleph_0)^{+}\) existeix
individualment per Hartogs (definició~\ref{def:cardinal-successor}), sense
reunir cap família no acotada de cardinals. El que \emph{sí} requereix SR és
disposar d'un tram \((\aleph_\beta)_{\beta<\lambda}\) o
\((\beth_\beta)_{\beta<\lambda}\) com a família interna ---per exemple per
enunciar HCG uniformement sobre tots els \(\beta<\lambda\); HCG per a
\emph{tots} els ordinals és un esquema d'enunciats, un per a cada
\(\lambda\)---, no pas l'equació d'HC en si. Per
tant, associar \enquote{el continu} al Nivell~II és imprecís: el continu i HC
són de Nivell~I; només la jerarquia global dels alephs és de Nivell~II.
\end{observacio}

\begin{observacio}[Què pregunta HC, i per què Cantor no la resol]\label{obs:hc-independencia}
El punt essencial ---i el tancament del capítol--- és que HC \emph{no és un
teorema pendent de demostració elemental}, sinó una qüestió \emph{independent}
dels axiomes habituals. La teoria desenvolupada en aquest llibre permet
\emph{formular} HC amb tota precisió: distingeix la pujada per successor
(\(\aleph\)) de la pujada per conjunt de parts (\(\beth\)), i HC pregunta si
coincideixen al primer graó. Però la teoria \emph{no} determina la resposta.
El que Cantor va establir ---i és tot el que un argument diagonal pot
establir--- és la \emph{desigualtat} \(\aleph_0<2^{\aleph_0}\)
(teorema~\ref{teo:cantor}); saber \emph{on cau} exactament \(2^{\aleph_0}\) dins
la successió dels alephs és una altra qüestió. Els resultats de Gödel (1938)~\cite{godel1940} i
Cohen (1963)~\cite{cohen1966} mostren que, \emph{si ZFC és consistent}, ni HC ni la seva negació no es
poden demostrar a ZFC: HC és \emph{independent de ZFC}, relativament a la
seva consistència. Com que ETCS${}+{}$SR i ZFC són bi-interpretables
(observació~\ref{obs:sr-rigor}), la independència es transporta a la teoria
d'aquest llibre amb SR; i a fortiori a la teoria de Nivell~I, que és més
feble. Cap d'aquests resultats no es demostra aquí: són metateoremes citats. La teoria fixa el que
és forçós ---\(\mathfrak c>\aleph_0\), i \(\operatorname{cf}(\mathfrak c)>\aleph_0\)
per König (observació~\ref{obs:cofinalitat-importa})--- i deixa lliure la resta.
No entrem aquí en els mètodes ---forcing, models interns--- que estableixen
aquesta independència; només subratllem que la pregunta està ben plantejada i
que la seva resposta no és accessible pels mitjans d'aquest volum, ni per cap
refinament elemental d'ells. Per a un tractament complet d'aquests mètodes i de
l'aritmètica cardinal clàssica, vegeu els manuals estàndard~\cite{jech2003,kunen2011}.
\end{observacio}

\begin{resumcap}
\apartat{Idees centrals}
\begin{enumerate}
  \item SR és un esquema: una regla externa que determina cada \((X_i,R_i)\) llevat d'un \emph{únic} isomorfisme es representa per una família interna \(E\to I\). Fa existir famílies; l'elecció tria dins famílies que ja existeixen. És opcional, i porta la teoria a la força de ZFC.
  \item \(\aleph_0=\abs{\N}\) i \(\kappa^{+}=\abs{\Hart(K)}\) són de Nivell~I. La família \((\aleph_\beta)_{\beta<\lambda}\), per a un \(\lambda\) fix, demana SR als límits; la instància és \enquote{ser una solució de la recursió fins a \(\gamma\)}, rígida perquè les fibres són bons ordres.
  \item Amb l'elecció, tot cardinal infinit és un aleph. Cofinalitat: \(\aleph_0\) i els successors són regulars (els successors, amb el teorema de Hessenberg, citat); \(\aleph_\omega\) és singular.
  \item Els beths pugen pel conjunt de parts i dominen els alephs. El continu \(\mathfrak c=\beth_1=2^{\aleph_0}\) i la hipòtesi del continu són de Nivell~I; els trams de la successió com a famílies, de Nivell~II.
  \item HC és independent de ZFC si ZFC és consistent (Gödel, Cohen); per la bi-interpretabilitat, també d'aquesta teoria amb SR. König en diu una cosa forçosa: \(\operatorname{cf}(2^{\aleph_0})>\aleph_0\).
\end{enumerate}
\apartat{Errors típics} Parlar de la família de \emph{tots} els alephs com d'un objecte (ni amb SR no existeix: només els trams fins a un ordinal fix); creure que SR és una forma d'elecció; aplicar SR a regles determinades només llevat d'isomorfisme (sense unicitat); atribuir HC al Nivell~II; afirmar que HC és \enquote{independent} sense la hipòtesi de consistència; i usar \(\kappa\cdot\kappa=\kappa\) com si s'hagués demostrat al llibre.
\apartat{Lectura recomanada} \cite[§§2--4 i~7]{shulman-comparing} per al reemplaçament estructural i la comparació amb ZFC; \cite{osius1974} per al resultat original; \cite[cap.~3, 5 i~12--15]{jech2003} per als alephs, la cofinalitat, Hessenberg, König i la independència d'HC; \cite[cap.~I.10 i~VII]{kunen2011} per al forcing; \cite{godel1940,cohen1966} per als resultats originals d'independència; i \cite[§25]{halmos} per a una primera lectura dels alephs i del continu.
\end{resumcap}

\chapter{Models, mida i comparació amb \texorpdfstring{ZF\slash ETCS}{ZF/ETCS}}\label{cap:models}

\begin{objectius}
\apartat{Prerequisits} Tot el llibre, i en particular els axiomes tal com s'han anat introduint, l'esquema SR del capítol~\ref{cap:alephs} i les observacions de lectura clàssica. Familiaritat informal amb els axiomes de ZFC, que es recorden a la taula de la secció~\ref{sec:taula-comparativa}.
\apartat{Objectius} En acabar el capítol, el lector hauria de poder:
\begin{itemize}
  \item recórrer els axiomes del llibre per grups i explicar-ne les dependències;
  \item identificar l'estructura obtinguda com un topos ben apuntat amb objecte dels nombres naturals i elecció, és a dir, un model d'ETCS;
  \item comparar, axioma per axioma, aquesta presentació amb ZFC i amb l'ETCS de Lawvere, distingint els axiomes equivalents, els derivats i els addicionals;
  \item situar la força de la teoria (Nivell~I, amb i sense SR) respecte de la teoria de Zermelo acotada i de ZFC;
  \item explicar com es recupera la pertinença, què se'n perd i què se'n guanya.
\end{itemize}
\end{objectius}

Els capítols anteriors han bastit, a partir d'un llenguatge en què els
únics individus són les aplicacions, tota la infraestructura que
habitualment associem als conjunts: parts i pertinença, productes i sumes,
quocients, imatges, exponencials, el conjunt de parts, els nombres
naturals i la cardinalitat fins al numerable. I, en una segona part, han
portat aquesta construcció fins a la teoria clàssica de l'infinit: la
comparació general de mides amb Schröder--Bernstein i el teorema de Cantor
(capítol~\ref{cap:cantor}), els bons ordres i la recursió transfinita
(capítols~\ref{cap:bons-ordres} i~\ref{cap:recursio}), Zorn, Hartogs i el
teorema del bon ordre (capítol~\ref{cap:zorn}), els tipus ordinals i la seva
aritmètica (capítols~\ref{cap:ordinals} i~\ref{cap:aritmetica-ordinal}), els
cardinals i la seva aritmètica (capítol~\ref{cap:cardinals}) i ---amb un principi
addicional--- la jerarquia sencera dels alephs i dels beths
(capítol~\ref{cap:alephs}). En cap
moment hem necessitat la relació de pertinença \(\in\) com a noció primitiva.
És natural que, un cop acabada la construcció, el lector es faci tres
preguntes: si això que hem obtingut és \emph{realment} una teoria de
conjunts; què és capaç d'expressar i què no; i quina relació guarda amb la
teoria clàssica de Zermelo--Fraenkel que va aprendre primer. Aquest capítol
final hi respon. No conté cap axioma nou ni cap construcció nova: és un
capítol metateòric, escrit per situar l'edifici sencer ---inclosa la part
transfinita--- i per assenyalar-ne els límits amb honestedat.

\section{Recapitulació dels axiomes}\label{sec:recapitulacio-axiomes}

Convé, abans de res, tenir a la vista tots els axiomes que hem assumit i
l'ordre en què depenen els uns dels altres. Els agrupem pel paper que fan.

El \textbf{grup del llenguatge de composició} fixa què és una aplicació i
com es componen: l'existència i unicitat de domini i codomini
(axioma~\ref{ax:dom-cod}), l'existència de la composició quan els tipus
encaixen (axioma~\ref{ax:composicio-existencia}), el seu domini i codomini
(axioma~\ref{ax:composicio-dom-cod}), l'associativitat
(axioma~\ref{ax:associativitat}) i les lleis de neutralitat
(axioma~\ref{ax:neutralitat}). Aquests cinc axiomes són la base sobre la
qual descansa tota la resta; en el vocabulari que hem evitat, diuen que les
aplicacions i els conjunts formen una categoria.

El \textbf{grup dels objectes distingits i la separació} dona els punts de
partida per construir: l'existència d'un conjunt terminal
(axioma~\ref{ax:terminal-existencia}) i d'un d'inicial
(axioma~\ref{ax:inicial-existencia}), la no-degeneració que els separa,
\(\mathbf0\ncong\mathbf1\) (axioma~\ref{ax:no-degeneracio}), i la separació
de les aplicacions pels elements globals
(axioma~\ref{ax:separacio-globals}), que és la que fa que els punts
\(\mathbf1\to A\) siguin prou nombrosos per distingir aplicacions ---la
propietat que més endavant expressarem dient que la teoria és \emph{ben
apuntada}.

El \textbf{grup de les construccions finites} postula les operacions
universals bàsiques: l'existència de sumes (axioma~\ref{ax:sumes}), la seva
disjunció i estabilitat (axioma~\ref{ax:sumes-disjuntes}), l'existència de
productes (axioma~\ref{ax:productes}), la representació de les parts per
aplicacions cap a \(\mathbf2\) (axioma~\ref{ax:representacio}) ---el nucli
lògic de tot el llibre--- i l'existència de coequalitzadors
(axioma~\ref{ax:coequalitzadors}). Amb ells s'obtenen tots els límits i
colímits finits i tota l'àlgebra booleana de parts.

El \textbf{grup de la regularitat i l'exactitud} controla les imatges i
els quocients: l'efectivitat de les relacions d'equivalència
(axioma~\ref{ax:exactitud}), que casa quocients i parells nucli. Combinat
amb els coequalitzadors i la separació global, dona la factorització
imatge, i amb ella que tot epimorfisme és regular i estable per canvi de
base, sense elecció (capítol~\ref{cap:imatges}).

El \textbf{grup de l'ordre superior} fa la teoria capaç de parlar de si
mateixa: l'existència d'exponencials (axioma~\ref{ax:exponencials}), de la
qual surt, prenent \(\mathbf2^{A}\), el conjunt de parts com a objecte, la
pertinença interna i els quantificadors.

El \textbf{grup de l'infinit} afegeix l'única font d'infinitud: l'objecte
dels nombres naturals amb la seva propietat universal de recursió
(axioma~\ref{ax:naturals}), del qual pengen la inducció, l'aritmètica, la
finitud interna i tota la cardinalitat numerable.

Finalment, els \textbf{axiomes d'elecció}, que hem mantingut separats
perquè molts resultats no en depenen: l'elecció global ---tot epimorfisme
s'escindeix--- (axioma~\ref{ax:eleccio}) i, com a cas particular seu,
l'elecció numerable (axioma~\ref{ax:eleccio-numerable}), que va ser
necessària per identificar les dues nocions d'infinitud
(teorema~\ref{teo:infinit-implica-dedekind}) i per la unió numerable
d'enumerables (teorema~\ref{teo:unio-numerable}).

Fora d'aquesta llista, i amb un estatus deliberadament diferent, hi ha
l'\textbf{esquema de reunió de famílies} o \emph{substitució estructural}
(principi~\ref{sr:esquema}), que la part transfinita ha necessitat només al
capítol opcional dels alephs (capítol~\ref{cap:alephs}). Insistim que SR \emph{no} és un axioma més al costat dels anteriors: és un
\emph{esquema} ---una instància per a cada fórmula funcional i invariant per
isomorfisme---, no una sentència, exactament com el reemplaçament de la
teoria clàssica és un esquema i no un axioma únic. Per això l'hem presentat
com un \emph{principi} a part i no amb l'entorn dels axiomes, i per això tots
els capítols fins al~18 en són independents: la teoria de Nivell~I ---tota
la que va del llenguatge d'aplicacions fins als cardinals--- no el fa servir.
Només la reunió de famílies transfinites no acotades ---la successió dels
alephs, els beths--- el requereix. Aquesta és la línia divisòria que
recorrerem amb detall a la secció~\ref{sec:etcs-sr}.

\begin{observacio}[El mapa de dependències]\label{obs:mapa-dependencies}
La lectura del llibre segueix, de fet, aquest ordre de dependència. El
llenguatge de composició no pressuposa res; els objectes distingits el
necessiten; les construccions finites necessiten els objectes distingits;
la regularitat necessita sumes i coequalitzadors; els exponencials
necessiten els productes; l'objecte dels nombres naturals necessita els
exponencials per a la recursió amb paràmetres
(observació~\ref{obs:exponencials-recursio}); i l'elecció és transversal,
invocada només on s'assenyala. Cap axioma no es fa servir abans
d'enunciar-se, i cap grup no depèn d'un de posterior. Aquesta és la raó
per la qual el llibre es pot llegir linealment: l'estructura lògica i
l'estructura expositiva coincideixen.
\end{observacio}

\section{Què és, categòricament, l'estructura obtinguda}\label{sec:caracteritzacio}

Reunits tots els axiomes, es pot dir en una sola frase quina mena
d'estructura formen les aplicacions i els conjunts d'aquest llibre. En el
vocabulari de la teoria de categories ---que no hem fet servir en cap
demostració, però que aquí, en un capítol de balanç, resulta el més
econòmic per anomenar el conjunt---, l'estructura és un \emph{topos ben
apuntat, amb objecte dels nombres naturals i elecció}. Desglossem-ho amb
paraules planes, perquè cada peça correspon exactament a un grup
d'axiomes.

Ser un \emph{topos elemental} vol dir tres coses, que són les que hem
postulat: que hi ha tots els límits finits ---terminal, productes,
equalitzadors i les seves combinacions, que vam obtenir al
capítol~\ref{cap:valors}---; que hi ha exponencials
(axioma~\ref{ax:exponencials}); i que les parts d'un conjunt qualsevol
estan representades per les aplicacions cap a un objecte fix \(\mathbf2\),
el \emph{classificador de subobjectes} (axioma~\ref{ax:representacio}).
Aquesta darrera és la propietat decisiva: fa que la lògica dels
subconjunts sigui interna a la teoria, amb el conjunt de parts, la
pertinença i els quantificadors com a objectes i aplicacions de ple dret,
tal com vam veure al capítol~\ref{cap:parts-interna}.

\emph{Ben apuntat} (\emph{well-pointed}) és el que afegeix l'axioma de separació pels elements
globals (axioma~\ref{ax:separacio-globals}) juntament amb la
no-degeneració (axioma~\ref{ax:no-degeneracio}): els punts
\(\mathbf1\to A\) basten per distingir aplicacions, i \(\mathbf1\) no
coincideix amb \(\mathbf0\). És el que permet raonar «element a element»,
com fem contínuament amb els elements globals i generalitzats, sense
trair el punt de vista estructural. Un topos qualsevol no té per què ser
ben apuntat; el nostre sí, i això l'acosta molt a la categoria dels
conjunts ordinaris.

L'\emph{objecte dels nombres naturals} (axioma~\ref{ax:naturals}) i
l'\emph{elecció} (axioma~\ref{ax:eleccio}) són els dos ingredients que
completen el quadre. Un topos ben apuntat amb aquests dos afegits és,
precisament, un model de la teoria estructural de conjunts de Lawvere,
l'ETCS~\cite{lawvere1964,lawvere-etcs-reprint}. Dit d'una altra manera:

\begin{observacio}[La tesi del llibre, en una frase]\label{obs:tesi-etcs}
Els axiomes d'aquest llibre són una axiomatització d'ETCS presentada en un
llenguatge de primer ordre d'aplicacions, sense el vocabulari de la teoria
de categories. Tot el que hem demostrat és, doncs, un teorema d'ETCS; i
qualsevol model d'ETCS ---en particular la categoria dels conjunts
clàssics--- satisfà els nostres axiomes. Aquesta és la resposta curta a la
pregunta «és realment una teoria de conjunts?»: ho és tant com ETCS, que
és una axiomatització \emph{autònoma} de la categoria dels conjunts,
proposada com a fonament alternatiu de la matemàtica~\cite{lawvere-rosebrugh,mclarty-catset}.
Evitem l'expressió \enquote{axiomatització completa}, que es podria confondre
amb la \emph{completesa lògica} (que tot enunciat o la seva negació és
demostrable) o amb la \emph{categoricitat} (que hi ha un únic model llevat
d'isomorfisme): cap de les dues no val aquí ---ETCS és una teoria de primer
ordre, subjecta a la incompletesa i amb molts models no isomorfs---, ni cap de
les dues no es demostra en aquest llibre. \enquote{Autònoma} vol dir només que
els axiomes es basten a si mateixos com a presentació dels conjunts, sense dependre d'una teoria material prèvia.
\end{observacio}

\section{La part transfinita dins el balanç}\label{sec:transfinit-balanc}

El capítol que en la primera versió d'aquest llibre tancava el volum
s'aturava a la cardinalitat numerable. La part transfinita
(capítols~\ref{cap:cantor}--\ref{cap:alephs}) ha estès el desenvolupament
fins a la teoria clàssica de l'infinit, i convé situar-la dins aquest
balanç, perquè il·lumina exactament què es pot fer amb els axiomes de
Nivell~I i què demana un principi addicional.

El primer bloc és \emph{purament diagonal} i no costa cap axioma nou. El
teorema de \textbf{Schröder--Bernstein} (teorema~\ref{teo:schroder-bernstein})
---dues injeccions oposades donen un isomorfisme--- es prova pel menor punt
fix d'un operador creixent sobre \(P(A)\), sense elecció. El teorema de
\textbf{Cantor} (teorema~\ref{teo:cantor}) ---\(A\prec P(A)\)--- és
l'argument diagonal, i tampoc no necessita res de nou; d'ell surt que no hi
ha conjunt universal ni mida màxima. Aquests dos resultats són, en el
llenguatge estructural, teoremes elementals sobre parts i pertinença interna.

El segon bloc és la \emph{teoria de l'ordre}, també de Nivell~I. Els
\textbf{bons ordres} (capítol~\ref{cap:bons-ordres}) tenen la rigidesa que
els fa mesura de forma ---entre dos n'hi ha com a molt un isomorfisme
d'ordre---, i sobre ells es construeixen la \textbf{inducció} i la
\textbf{recursió transfinites} amb codomini fix (capítol~\ref{cap:recursio}),
la reunió de solucions parcials dins un exponencial fix, \emph{sense}
elecció. Amb la recursió es tanca la comparabilitat de bons ordres
(teorema~\ref{teo:comparabilitat}). El teorema de \textbf{Hartogs}
(teorema~\ref{teo:hartogs}) ---un bon ordre que no s'injecta en \(A\)--- és
el paral·lel ordinal de Cantor, i és igualment de Nivell~I: es construeix
tallant els bons ordres sobre parts d'\(A\) ---parells \((S,R)\) dins
\(P(A)\times P(A\times A)\)--- i quocientant-los per isomorfisme.

El tercer bloc introdueix l'\emph{elecció}. El teorema del \textbf{bon
ordre} (teorema~\ref{teo:bon-ordre}) i el \textbf{lema de Zorn}
(capítol~\ref{cap:zorn}) es demostren amb l'elecció i la cota de Hartogs
---que evita, precisament, la recursió \enquote{sobre tots els ordinals}---,
i resulten equivalents a l'elecció. Sobre aquesta base, els \textbf{tipus
ordinals} i la seva aritmètica (capítols~\ref{cap:ordinals}
i~\ref{cap:aritmetica-ordinal}) i els \textbf{cardinals} i la seva
(capítol~\ref{cap:cardinals}) es desenvolupen amb el teorema de Cantor
reaparegut com \(\abs{A}<2^{\abs{A}}\), la comparabilitat total dels
cardinals i el teorema de Burali--Forti (teorema~\ref{teo:burali-forti}) com
a segon gran resultat diagonal ---no hi ha longitud d'ordre màxima.

Tot això és de Nivell~I o de Nivell~I més elecció: viu dins conjunts ja
construïts, acotat per un \(\Hart\) sempre que cal reunir tipus. \emph{Només}
el capítol final dels alephs (capítol~\ref{cap:alephs}) surt d'aquest terreny,
i ho fa exactament en un punt: reunir la successió transfinita no acotada
\(\alpha\mapsto\aleph_\alpha\) ---i la dels beths--- en una família. Aquest
és l'únic lloc de tot el llibre on cal SR, i és el tema de la
secció~\ref{sec:etcs-sr}.

\section{Ordinal material i tipus estructural}\label{sec:vonneumann}

La part transfinita permet, a més, precisar una comparació que la
introducció només podia insinuar: la que hi ha entre l'\emph{ordinal
material} de la teoria clàssica i el \emph{tipus d'ordre estructural}
d'aquest llibre.

En la teoria de von Neumann, un ordinal \emph{és} un conjunt: un conjunt
transitiu ben ordenat per la pertinença, de manera que cada ordinal és
literalment el conjunt dels ordinals menors ---\(0=\varnothing\),
\(1=\{0\}\), \(2=\{0,1\}\), \(\omega=\{0,1,2,\dots\}\)--- i
\enquote{\(\alpha<\beta\)} es defineix com \enquote{\(\alpha\in\beta\)}. És
una construcció elegant i canònica: cada tipus de bon ordre té \emph{un}
representant distingit, l'ordinal de von Neumann, i la relació d'ordre és la
pertinença mateixa. El preu és que descansa sobre la pertinença global i la
jerarquia acumulativa: un ordinal arrossega tota la seva història de
membres, i la seva construcció ---per recursió al llarg de tots els
ordinals--- fa servir el reemplaçament.

En aquest llibre un \emph{tipus ordinal} no és cap conjunt privilegiat, sinó
un bon ordre \emph{considerat llevat d'isomorfisme d'ordre}
(definició~\ref{def:tipus-ordinal}), i \enquote{\(\alpha<\beta\)} es defineix
per segments inicials, no per pertinença ---de fet vam advertir explícitament
que \(\alpha\in\beta\) no és una notació adequada. Els dos enfocaments
descriuen la mateixa matemàtica, i la traducció és neta: el tipus estructural
d'un bon ordre \(W\) correspon a la classe d'isomorfisme d'ordre de \(W\),
i l'ordinal material de von Neumann n'és el representant canònic que la teoria
clàssica \emph{tria} dins aquesta classe. On l'estructural diu \enquote{el
tipus de \(W\), únic llevat d'un únic isomorfisme}, el material diu
\enquote{l'ordinal \(\in\)-transitiu isomorf a \(W\)}. La diferència no és de
contingut sinó de si es fa o no una tria canònica de representant: el
tractament estructural la evita perquè no la necessita, i perquè triar-la
demanaria la mateixa pertinença global que hem volgut no fer primitiva.

\begin{observacio}[Per què l'estructural no tria representant]\label{obs:vonneumann-representant}
Es podria pensar que l'ordinal de von Neumann és \enquote{millor} perquè és
canònic. Però la seva canonicitat és exactament el que costa la pertinença
global: per triar, dins cada classe d'isomorfisme de bons ordres, l'ordinal
\(\in\)-transitiu, cal disposar de la relació \(\in\) i de la recursió que el
construeix ---i, per als ordinals grans, del reemplaçament. El tractament
estructural prefereix treballar amb el tipus \emph{sense} representant, com fa
la pràctica matemàtica quan diu \enquote{el tipus d'ordre de \(\Q\)} sense
identificar-lo amb cap conjunt concret. Quan es vol un representant ---per
exemple per fer aritmètica còmoda--- se'n pot triar un amb l'elecció (el
cardinal inicial n'és un exemple, capítol~\ref{cap:cardinals}), però és una
tria \emph{addicional}, no un tret primitiu. Aquesta és la mateixa filosofia
que separa \(\abs{A}\) com a classe d'equipotència del cardinal inicial
escollit: l'estructura reté el que és invariant, i deixa la tria de
representant per a qui la vulgui.
\end{observacio}

\section{Interpretació en conjunts petits}\label{sec:conjunts-petits}

Una manera de convèncer-se que els axiomes descriuen conjunts és
interpretar-los en un univers ja conegut i comprovar que s'hi compleixen.
Fixem, doncs, un univers de conjunts clàssics ---els conjunts d'una teoria
de Zermelo--Fraenkel, posem--- i llegim-hi cada noció del llibre.

Un \emph{conjunt} del llibre és, en aquesta lectura, un conjunt ordinari.
Una \emph{aplicació} és una funció en el sentit habitual, amb el seu
domini i el seu codomini fixats ---no un simple graf, sinó una funció amb
els seus dos extrems marcats, que és el que demana l'axioma~\ref{ax:dom-cod}.
La \emph{composició} és la composició de funcions; les \emph{identitats}
són les funcions identitat. El conjunt \emph{terminal} \(\mathbf1\) és
qualsevol singletó \(\{\ast\}\); l'\emph{inicial} \(\mathbf0\) és el
conjunt buit; la \emph{no-degeneració} és el fet que \(\varnothing\) i
\(\{\ast\}\) no són equipotents. Un \emph{element global} \(\mathbf1\to A\)
és, exactament, un element de \(A\) en el sentit ordinari, perquè triar
una funció des d'un singletó és triar un valor. El conjunt \(\mathbf2\) és
qualsevol parell \(\{0,1\}\), i una \emph{part} de \(A\) es correspon amb
un subconjunt ordinari via la seva funció característica
(axioma~\ref{ax:representacio}). Les \emph{sumes} són unions disjuntes; els
\emph{productes}, productes cartesians; els \emph{coequalitzadors},
quocients per relacions d'equivalència; els \emph{exponencials}
\(B^{A}\), el conjunt de totes les funcions de \(A\) a \(B\); l'\emph{objecte
dels nombres naturals}, el conjunt \(\{0,1,2,\dots\}\) amb el zero i el
successor.

\begin{observacio}[La lectura és fidel, però no és l'única]\label{obs:interpretacio-fidel}
Sota aquesta interpretació, cadascun dels divuit axiomes ---els disset
estructurals més l'elecció global, sense comptar-hi l'elecció numerable,
que n'és un cas particular--- és un teorema
elemental de la teoria clàssica de conjunts: que tota funció té domini i
codomini, que el producte cartesià existeix, que el conjunt de funcions
existeix, i així successivament. Per això la categoria dels conjunts
clàssics és un model dels nostres axiomes, i qualsevol cosa que hàgim
demostrat hi val. Ara bé ---i això és essencial---, \emph{no és l'únic}
model. Els mateixos axiomes se satisfan en universos molt més petits o de
mides diferents: per exemple, la col·lecció dels conjunts \emph{finits}, amb
les funcions entre ells, compleix tots els axiomes \emph{llevat} del de
l'infinit (axioma~\ref{ax:naturals}), i restringint-se a un univers de
Grothendieck (un univers de conjunts de rang menor que un cardinal
inaccessible) s'obtenen models complets tancats sota totes les operacions,
inclòs el conjunt de parts. La teoria no fixa un univers; en descriu una
\emph{forma}, i molts universos de mides diferents la tenen.
\end{observacio}

\begin{observacio}[Per què no serveixen els hereditàriament numerables]\label{obs:no-hcountable}
Podria temptar prendre com a model \enquote{petit} la col·lecció dels conjunts
\emph{hereditàriament numerables}. No funciona: aquesta col·lecció \emph{no}
és tancada sota el conjunt de parts. En efecte, \(\N\) és hereditàriament
numerable, però \(P(\N)\) no ho és ---és no numerable---, de manera que
l'axioma del conjunt de parts falla dins d'aquell univers. Un model complet
dels axiomes ha de ser tancat sota totes les operacions que aquests postulen;
per això cal, o bé un univers de Grothendieck (tancat sota parts amb el marge
de mida corresponent), o bé un univers amb hipòtesis de mida explícites que
garanteixin la clausura. La numerabilitat hereditària, tota sola, no la dona.
\end{observacio}

Aquesta pluralitat de models no és un defecte sinó el senyal distintiu
d'una axiomatització estructural: descriu els conjunts \emph{llevat
d'isomorfisme i llevat d'equivalència de la categoria sencera}, sense
comprometre's amb cap realització particular. És el mateix que passa amb
els axiomes de grup, que no descriuen \emph{un} grup sinó la forma comuna
a tots.

\section{Pertinença primitiva i pertinença recuperada}\label{sec:pertinenca}

La diferència més visible amb la teoria clàssica és el paper de la
pertinença. Val la pena precisar-la, perquè és on es concentra bona part
del canvi de perspectiva anunciat a la introducció.

En la teoria de Zermelo--Fraenkel, \(\in\) és una relació \emph{primitiva}
i \emph{global}: té sentit preguntar-se, per a dos conjunts qualssevol
\(x\) i \(A\), si \(x\in A\), i la resposta forma part de les dades del
món. Que \(3\in 5\) (en la codificació de von Neumann els naturals són
conjunts, i \(5=\{0,1,2,3,4\}\)) és un fet; que \(3\in\pi\), amb una certa
codificació dels reals, també ho és, encara que sigui un fet estrany.
Aquesta globalitat té una conseqüència: cada conjunt arrossega tota la
seva història interna ---els seus elements, els elements dels seus
elements, i així avall fins al buit---, i l'axioma de fonamentació
garanteix que aquesta baixada acaba.

En la teoria d'aquest llibre no hi ha cap relació \(\in\) primitiva. El
que hi ha és una pertinença \emph{recuperada}, i \emph{local}. «Ser
element de \(A\)» vol dir ser un element global \(\mathbf1\to A\), o més
generalment un element generalitzat \(X\to A\); i «pertànyer a una part
\(P\hookrightarrow A\)» vol dir factoritzar per la inclusió, cosa que la
pertinença interna del capítol~\ref{cap:parts-interna} converteix en una
aplicació \(A\to\mathbf2\) avaluada al punt. Aquesta pertinença sempre és
\emph{relativa a un conjunt ambient}: un element ho és \emph{d'un} conjunt
determinat, i la pregunta «aquest element pertany a aquell altre conjunt?»
només té sentit un cop fixat el codomini.

\begin{observacio}[Què es perd i què es guanya]\label{obs:pertinenca-local}
La pèrdua és la pertinença creuada entre conjunts arbitraris: aquí no té
sentit preguntar si \(\N\) és element de \(\R\), ni si un element de \(A\)
«és el mateix» que un element de \(B\) quan \(A\) i \(B\) no tenen cap
relació ---els elements no tenen identitat pròpia fora del conjunt que els
allotja. En la lectura estructural això no és una mancança sinó una
disciplina: reflecteix que en la pràctica matemàtica mai no comparem un
nombre natural amb un nombre real «com a conjunts», sinó a través d'una
aplicació explícita \(\N\to\R\). El guany és que desapareixen els fets
accidentals de codificació ---el \(3\in\pi\) d'abans--- que la teoria
clàssica genera i que cap matemàtic no fa servir. La informació que queda
és exactament la invariant per isomorfisme, que és la que la matemàtica
utilitza.
\end{observacio}

\begin{observacio}[Es pot recuperar una \(\in\) global?]\label{obs:recuperar-in}
Sí, i el fet que es pugui és el pont tècnic amb la teoria clàssica. A
partir d'un model dels nostres axiomes ---amb un afegit sobre la mida que
esmentarem--- es pot construir una relació \(\in\) global i un univers de
conjunts «amb elements», recuperant un model de (una part de) Zermelo. La
construcció associa a cada conjunt estructural un arbre d'aplicacions que
en codifica una versió «ben fundada», i és inversa, en el sentit adequat,
de la lectura de la secció~\ref{sec:conjunts-petits}. No la desenvolupem
aquí; el que importa és que la pertinença global no és una dada primitiva
irrenunciable, sinó una estructura \emph{derivable} quan es vol, a canvi
d'hipòtesis addicionals de reemplaçament~\cite{mclarty-catset,shulman-comparing}.
\end{observacio}

\section{Comparació amb ZF/ZFC i amb ETCS}\label{sec:comparacio}

Amb això a la mà, podem situar el llibre respecte de les dues grans
famílies d'axiomatitzacions.

Respecte de \textbf{ZF/ZFC}, la diferència de fons no és de potència sinó
de \emph{primitius}. ZF pren com a primitius els conjunts i la relació
\(\in\), i construeix les funcions com a conjunts de parells; nosaltres
prenem com a primitives les aplicacions, i recuperem la pertinença com una
noció derivada i local, tal com acabem de veure (secció~\ref{sec:pertinenca}).
La diferència no és que ZF descrigui un univers únic i la nostra teoria molts:
ZF és igualment una teoria de primer ordre amb molts models no isomorfs (pel
teorema de Löwenheim--Skolem i la incompletesa, en té de mides i tipus
d'isomorfisme diferents). El contrast correcte és de \emph{criteri d'identitat}:
en la lectura material de ZF, cada element està determinat \emph{absolutament}
per la relació \(\in\) ---dos conjunts amb els mateixos elements són el mateix,
i la pertinença fixa la identitat de cada objecte dins un model donat---,
mentre que en la lectura estructural la identitat és \emph{llevat
d'isomorfisme}: un objecte només està determinat pel seu paper en la categoria,
i cap propietat interna distingeix objectes isomorfs. La distinció, doncs, és
entre identitat rígida per \(\in\) i identitat estructural llevat
d'isomorfisme, no entre una teoria amb model únic i una amb molts models. En el
pla lògic, totes dues són teories de primer ordre, però amb signatures
diferents: la de ZF té un únic símbol de relació \(\in\); la nostra, els
predicats \(\Dom\), \(\Cod\) i \(\Com\) sobre aplicacions.

Respecte d'\textbf{ETCS}, la relació és de \emph{presentació}, no de
contingut. Com ja hem dit (observació~\ref{obs:tesi-etcs}), els nostres
axiomes són ETCS. La diferència amb l'exposició de Lawvere és que ell la
formula en el llenguatge de les categories ---parla de límits, objectes
exponencials i classificador de subobjectes--- mentre que nosaltres hem
volgut que cada axioma fos una fórmula elemental sobre aplicacions,
llegible sense saber teoria de categories, i que les construccions
categòriques apareguessin com a \emph{teoremes} i no com a vocabulari de
partida. El contingut matemàtic és el mateix; canvia qui és el públic.

\begin{observacio}[Força relativa: el lloc exacte a l'escala]\label{obs:forca-relativa}
La comparació de força mereix una mica de cura, perquè el resultat sovint
sorprèn. ETCS ---i per tant aquest llibre--- \emph{no} és tan fort com
ZFC. La seva força demostrativa és la de la teoria de Zermelo amb elecció
i amb separació restringida a fórmules \emph{acotades} (l'anomenada
\(Z\) acotada, o \(\mathrm{BZC}\)): ETCS i \(\mathrm{BZC}\) demostren
exactament els mateixos enunciats aritmètics, són \emph{equiconsistents} i
mútuament interpretables~\cite{osius1974,mclarty-catset}. El que separa ETCS de ZFC
és l'\emph{esquema de reemplaçament}, que ETCS en la seva forma bàsica no
inclou. Afegint a ETCS una versió categòrica concreta del reemplaçament ---la
formulació de reemplaçament\slash col·lecció estructural adoptada com a SR
(observació~\ref{obs:sr-rigor}, capítol~\ref{cap:alephs})--- s'obté una teoria
de la mateixa força que ZFC, i llavors la correspondència amb la teoria
clàssica és exacta en el sentit precís de la bi-interpretabilitat demostrada
a~\cite{osius1974,shulman-comparing} (no en el sentit de completesa lògica).
La correspondència depèn, doncs, de \emph{quina} versió de l'esquema
s'adopti, i és per a aquella versió que val el teorema de comparació. Així, en
l'escala de les teories de conjunts, aquest llibre se situa una mica per
sota de ZFC ---al nivell de la teoria de Zermelo amb elecció---, que és
més que suficient per a pràcticament tota la matemàtica ordinària, però
que cal saber reconèixer per no atribuir a la teoria una potència que no
té sense l'afegit del reemplaçament.
\end{observacio}

\subsection{ETCS, reemplaçament i la substitució estructural}\label{sec:etcs-sr}

La discussió anterior sobre la força relativa pren un relleu concret ara que
el llibre ha desenvolupat la part transfinita, perquè és allà on la
diferència entre ETCS i ETCS reforçat es fa tangible. Convé, doncs,
distingir amb precisió les dues teories.

La \textbf{teoria de Nivell~I} d'aquest llibre ---els capítols~1 a~18--- és
ETCS: un topos ben apuntat amb objecte dels nombres naturals i elecció, de
la força de la teoria de Zermelo acotada amb elecció. \emph{Tota} la teoria
de l'infinit fins als cardinals hi cap: Cantor, Schröder--Bernstein, bons
ordres, recursió transfinita, Zorn, Hartogs, tipus ordinals, aritmètica
ordinal i cardinal. La raó per la qual hi cap, malgrat que la presentació
clàssica de bona part d'això usa el reemplaçament, és que hem mantingut
totes les construccions \emph{acotades}: cada recursió té codomini fix, cada
suprem de tipus es pren dins un bon ordre ja construït, i la cota de Hartogs
substitueix a tot arreu la recursió \enquote{sobre tots els ordinals}. Aquest
és el contingut de la traducció que el llibre proposa: allà on la teoria
clàssica reuniria una família transfinita amb el reemplaçament, l'estructural
la manté dins un \(\Hart\) i la tracta amb la recursió de codomini fix.

La \textbf{teoria reforçada}, ETCS més la substitució estructural SR
(principi~\ref{sr:esquema}), només fa falta quan la família transfinita
\emph{no es pot acotar}: la successió \(\alpha\mapsto\aleph_\alpha\), els
beths, les jerarquies iterades \(A,P(A),P^2(A),\dots\) reunides en un sol
objecte. SR és la traducció estructural de l'\emph{esquema de reemplaçament}:
allà on ZF diu \enquote{la imatge d'un conjunt per una funció-classe és un
conjunt}, SR diu \enquote{una regla externa uniforme i invariant es
representa per una família interna}. Amb SR afegit, la teoria assoleix la
força de ZFC, i la correspondència amb la teoria clàssica esdevé exacta en el
sentit precís de la bi-interpretabilitat ---no de completesa lògica---, per a
la formulació concreta de l'esquema adoptada (observació~\ref{obs:sr-rigor})
\cite{osius1974,shulman-comparing}.

\begin{observacio}[SR és un esquema, no un axioma; i és opcional]\label{obs:sr-metateoric}
Dues precisions tanquen aquesta comparació. Primera, sobre l'\emph{estatus}:
SR no és un axioma més al costat de l'elecció, sinó un \emph{esquema} de
fórmules, com el reemplaçament clàssic; per això el llibre el presenta com un
principi a part (principi~\ref{sr:esquema}) i el marca separadament a la
caixa de dependències, no barrejat amb els axiomes. Segona, sobre
l'\emph{abast}: SR és estrictament opcional. El volum es pot llegir sencer
fins al capítol~18 sense ell; només el capítol dels alephs
(capítol~\ref{cap:alephs}) el requereix, i ho fa de manera explícita a cada
ús. Un lector que vulgui quedar-se en ETCS pur té, doncs, tota la teoria de
l'infinit numerable, ordinal i cardinal a la seva disposició; qui vulgui els
trams de la successió dels alephs i dels beths com a famílies ha d'adoptar SR, sabent que amb
això puja al nivell de força de ZFC. (El continu \(\mathfrak c=\abs{P(\N)}\) i
la hipòtesi del continu, en canvi, ja són accessibles sense SR, al Nivell~I.)
Aquesta és la línia divisòria exacta entre les dues teories, i el llibre l'ha
fet visible a cada pas en comptes d'amagar-la.
\end{observacio}

\section{Taula comparativa: ZFC, ETCS i aquest llibre}\label{sec:taula-comparativa}

Recollim en tres taules els axiomes de les tres presentacions, i en una
quarta la correspondència entre elles. Per a ZFC seguim la formulació dels
manuals~\cite{jech2003,kunen2011}; per a ETCS, la forma moderna que en donen
McLarty i Shulman~\cite{mclarty-catset,shulman-comparing}, equivalent a
l'original de Lawvere~\cite{lawvere1964,lawvere-etcs-reprint}.

\begin{center}\small
\begin{tabular}{@{}>{\raggedright\arraybackslash}p{0.30\linewidth}>{\raggedright\arraybackslash}p{0.62\linewidth}@{}}
\toprule
\multicolumn{2}{@{}l}{\textbf{ZFC} (primitius: conjunts i \(\in\))}\\
\midrule
Extensionalitat & Dos conjunts amb els mateixos elements són iguals.\\
Fonament & Tot conjunt no buit té un element disjunt d'ell.\\
Separació (esquema) & \(\{x\in A:\varphi(x)\}\) és un conjunt, per a tota fórmula \(\varphi\).\\
Parells & \(\{a,b\}\) és un conjunt.\\
Unió & \(\bigcup A\) és un conjunt.\\
Parts & \(P(A)\) és un conjunt.\\
Infinit & Hi ha un conjunt inductiu.\\
Reemplaçament (esquema) & La imatge d'un conjunt per una funció-classe és un conjunt.\\
Elecció & Tota família de conjunts no buits té una funció d'elecció.\\
\bottomrule
\end{tabular}
\end{center}

\begin{center}\small
\begin{tabular}{@{}>{\raggedright\arraybackslash}p{0.30\linewidth}>{\raggedright\arraybackslash}p{0.62\linewidth}@{}}
\toprule
\multicolumn{2}{@{}l}{\textbf{ETCS} (primitius: objectes, fletxes i composició)}\\
\midrule
Categoria & La composició és associativa i té identitats.\\
Límits finits & Hi ha objecte terminal, productes binaris i equalitzadors.\\
Exponencials & Per a tots \(A,B\) hi ha \(B^A\) amb avaluació universal.\\
Classificador & Hi ha \(\Omega\) que classifica els subobjectes.\\
Nombres naturals & Hi ha un objecte de nombres naturals.\\
Ben apuntat & \(\mathbf1\) separa les fletxes, i \(\mathbf0\not\cong\mathbf1\).\\
Elecció & Tot epimorfisme té secció.\\
\bottomrule
\end{tabular}
\end{center}

\begin{center}\small
\begin{tabular}{@{}>{\raggedright\arraybackslash}p{0.30\linewidth}>{\raggedright\arraybackslash}p{0.62\linewidth}@{}}
\toprule
\multicolumn{2}{@{}l}{\textbf{Aquest llibre} (primitius: aplicacions; \(\Dom\), \(\Cod\), \(\Com\))}\\
\midrule
Estructurals (5) & Domini i codomini, composició, associativitat, neutralitat (capítol~\ref{cap:aplicacions}).\\
Objectes distingits & Terminal, inicial, no-degeneració, separació global (capítol~\ref{cap:aplicacions}).\\
Construccions finites & Sumes, disjunció i estabilitat, productes, representació de les parts per \(\mathbf2\), coequalitzadors (capítols~\ref{cap:parts} i~\ref{cap:valors}).\\
Exactitud & Efectivitat de les relacions d'equivalència (capítol~\ref{cap:relacions}).\\
Exponencials & Existència de \(B^A\) (capítol~\ref{cap:exponencials}).\\
Nombres naturals & Recursió sobre \(\N\) (capítol~\ref{cap:naturals}).\\
Elecció & Tot epimorfisme té secció, i la numerable com a cas particular (capítols~\ref{cap:relacions} i~\ref{cap:infinitud}).\\
SR (esquema, opcional) & Reunió de famílies determinades llevat d'isomorfisme únic (capítol~\ref{cap:alephs}).\\
\bottomrule
\end{tabular}
\end{center}

La taula següent posa les tres presentacions en correspondència. La darrera
columna diu l'estatus de l'axioma del llibre respecte d'ETCS:
\emph{equivalent} (és el mateix axioma, en llenguatge d'aplicacions),
\emph{derivat} (és un teorema d'ETCS que el llibre postula perquè el
necessita abans de tenir les eines per demostrar-lo) o \emph{addicional}
(va més enllà d'ETCS).

\begin{center}\footnotesize
\begin{tabular}{@{}>{\raggedright\arraybackslash\hspace{0pt}}p{0.17\linewidth}>{\raggedright\arraybackslash\hspace{0pt}}p{0.21\linewidth}>{\raggedright\arraybackslash\hspace{0pt}}p{0.16\linewidth}>{\raggedright\arraybackslash\hspace{0pt}}p{0.18\linewidth}>{\raggedright\arraybackslash\hspace{0pt}}p{0.13\linewidth}@{}}
\toprule
Contingut & ZFC & ETCS & Aquest llibre & Estatus\\
\midrule
Composició i identitats & Implícites (funcions com a conjunts de parells) & Categoria & 5 axiomes estructurals & equivalent\\
Terminal & \(\{\varnothing\}\) (Parells) & Límits finits & Axioma & equivalent\\
Inicial & \(\varnothing\) (Separació) & Teorema en un topos & Axioma & derivat\\
Productes & Parells, Parts, Separació & Límits finits & Axioma & equivalent\\
Equalitzadors & Separació & Límits finits & Teorema (cap.~\ref{cap:parts}) & equivalent\\
Sumes; disjunció i estabilitat & Parells, Unió & Teoremes en un topos & Axiomes & derivats\\
Coequalitzadors & Parts, Separació (classes) & Teorema en un topos & Axioma & derivat\\
Efectivitat & Classes d'equivalència & Teorema en un topos & Axioma & derivat\\
Classificador & Separació i Parts (característiques) & Classificador \(\Omega\) & Representació per \(\mathbf2\) & equivalent (\(\Omega\cong\mathbf2\))\\
Exponencials, \(P(A)\) & Parts, Separació & Exponencials & Axioma; \(P(A):=\mathbf2^A\) & equivalent\\
Infinit & Infinit & Nombres naturals & Recursió sobre \(\N\) & equivalent\\
Extensionalitat & Extensionalitat & Ben apuntat & Separació global i no-degeneració & equivalent\\
Elecció & Elecció & Elecció & Elecció & equivalent\\
Separació & Esquema per a tota fórmula & Només fórmules acotades, via \(\Omega\) & Igual que ETCS & equivalent\\
Reemplaçament & Esquema & Absent & SR, opcional & addicional\\
Fonament & Fonament & Sense contrapartida & Sense contrapartida (no hi ha \(\in\) global) & ---\\
\bottomrule
\end{tabular}
\end{center}

\begin{observacio}[Lectura de la taula]\label{obs:lectura-taula}
Tres conclusions. \emph{Primera}: sense SR, els axiomes del llibre són
exactament ETCS. Els que figuren com a \enquote{derivats} ---l'inicial, les
sumes amb la seva disjunció i estabilitat, els coequalitzadors i
l'efectivitat--- són teoremes en qualsevol topos, demostrables a partir
dels límits finits, els exponencials i el classificador (vegeu el volum de
categories~\cite{perello-categories} i~\cite[cap.~IV]{maclane-moerdijk}); el
llibre els postula perquè els necessita als capítols~\ref{cap:parts}
a~\ref{cap:relacions}, abans de disposar del conjunt de parts amb què es
demostrarien. Com a teoria, doncs, el llibre no és més fort que ETCS:
només és redundant. \emph{Segona}: la diferència amb ZFC és doble. Hi ha
dos axiomes de ZFC sense contrapartida ---el \emph{fonament}, que parla
d'una pertinença global que aquí no existeix, i el \emph{reemplaçament},
que aquí és l'esquema opcional SR---, i la \emph{separació} és més feble,
restringida a les condicions que es poden expressar com a parts, és a dir,
amb quantificadors acotats per conjunts ja donats. \emph{Tercera}: que el
classificador sigui el conjunt \(\mathbf2=\mathbf1+\mathbf1\), de dos punts,
és un tret dels conjunts constants: en un topos qualsevol el classificador
\(\Omega\) pot tenir més valors de veritat, i és la combinació de ben
apuntat i elecció la que el redueix a \(\mathbf2\) (Diaconescu: l'elecció
fa booleà el topos). Per això el llibre el pot postular directament.
\end{observacio}

\begin{observacio}[Què és demostrat i què és citat]\label{obs:taula-citat}
La taula barreja dues menes d'afirmacions, que cal no confondre. Les
columnes \enquote{ZFC} i \enquote{Aquest llibre} i la correspondència entre
elles en el pla del \emph{contingut} (quin principi fa la feina de quin) són
comparacions informals que el lector pot comprovar amb el llibre a la mà.
En canvi, les afirmacions sobre \emph{teories} ---que els axiomes
\enquote{derivats} són teoremes de tot topos, que ETCS és equivalent a BZC, i
que ETCS\({}+{}\)SR és bi-interpretable amb ZFC--- són \emph{metateoremes
citats} de la literatura~\cite{osius1974,mclarty-catset,shulman-comparing,maclane-moerdijk},
no resultats demostrats en aquest volum, i la darrera depèn de la
formulació concreta de SR (observació~\ref{obs:sr-rigor}). El que el llibre
sí que demostra és la direcció directa: que tots els seus teoremes són
conseqüència dels axiomes de la taula.
\end{observacio}

\begin{observacio}[Força de les tres teories]\label{obs:forca-taula}
En força, l'escala és: \(\text{Nivell~I}=\text{ETCS}\equiv\mathrm{BZC}\)
(Zermelo amb elecció i separació acotada) \(<\) \(\mathrm{ZFC}\equiv
\text{ETCS}+\text{SR}\), on \(\equiv\) indica bi-interpretabilitat i \(<\) que
la teoria de la dreta demostra la consistència de la de l'esquerra
(observació~\ref{obs:forca-relativa};
\cite{osius1974,mclarty-catset,shulman-comparing}). La mancança de
reemplaçament no es nota en la matemàtica ordinària, i la part transfinita
del llibre s'ha mantingut al Nivell~I gràcies a les cotes de Hartogs; només
la jerarquia sencera dels alephs i dels beths l'exigeix.
\end{observacio}

\section{Límits del desenvolupament}\label{sec:limits}

Tanquem amb els límits, que un capítol de balanç honest no pot ometre.
Són de tres menes: de mida, de consistència i de força expressiva.

\subsection{Mida}\label{subsec:mida}

Els axiomes no diuen res sobre \emph{quants} conjunts hi ha ni sobre com
n'és de gran l'univers. Garanteixen que hi ha un conjunt infinit
(\(\N\)), i que a partir de qualsevol conjunt se'n poden formar el de parts
i els exponencials ---de manera que hi ha conjunts arbitràriament grans en
el sentit de Cantor---, però no postulen cap cardinal inaccessible ni cap
tancament sota famílies de mida arbitrària. En particular, sense afegir SR,
no hi ha reemplaçament: no es pot, en general, formar la unió d'una família
de conjunts indexada per un conjunt quan la família ve donada per una
\emph{fórmula} i no per una aplicació ja existent ---és exactament la
distinció entre regla externa i família interna que vam fer al
capítol~\ref{cap:families}. Aquesta és la limitació de mida més important de
la teoria de Nivell~I, i és exactament la que SR (principi~\ref{sr:esquema})
aixeca, permetent reunir les famílies transfinites no acotades del
capítol~\ref{cap:alephs}. Amb SR, la teoria té la força del reemplaçament;
sense ell, es queda ---còmodament--- en tota la matemàtica que no reuneix
col·leccions transfinites no acotades.

\subsection{Consistència relativa}\label{subsec:consistencia}

Com qualsevol teoria que interpreti l'aritmètica, la nostra no pot
demostrar la seva pròpia consistència, pel segon teorema de Gödel. El que
sí que tenim són resultats de consistència \emph{relativa}: la teoria és
consistent si ho és la teoria de Zermelo amb elecció, i recíprocament, per
la interpretació mútua esmentada. Dit d'una altra manera, no demanem al
lector cap acte de fe que no demani ja la teoria clàssica: qui accepti la
consistència de \(\mathrm{ZFC}\) ---o fins i tot només la de \(Z\) amb
elecció--- accepta la d'aquest llibre. No n'obtenim, doncs, cap garantia
absoluta de consistència; només la seguretat que no hem introduït cap risc
nou respecte del que la matemàtica ja assumeix.

\subsection{Força expressiva}\label{subsec:expressiva}

Finalment, hi ha coses que la teoria, tal com és, expressa amb dificultat
o no expressa. Les propietats que \emph{no} són invariants per isomorfisme
---els fets de codificació interna, la pertinença creuada entre conjunts
sense relació--- no s'hi poden ni tan sols formular, i això, com hem
argumentat, és volgut. Pel que fa als arguments transfinits, la situació és
més matisada del que la primera versió d'aquest capítol suggeria: la teoria
de l'ordre, els tipus ordinals i la recursió transfinita \emph{sí} que hi
són, sencers, mentre es mantinguin \emph{acotats} ---la recursió amb codomini
fix (capítol~\ref{cap:recursio}), els suprems dins un bon ordre, la cota de
Hartogs. El que queda fora de la teoria de Nivell~I, sense SR, és
exactament la reunió de col·leccions transfinites \emph{no acotades}: els
trams de la successió dels \(\aleph_\alpha\) i \(\beth_\alpha\), la jerarquia
acumulativa \(V_\alpha\) per a ordinals grans, o la recursió al llarg de
qualsevol ordinal fix per molt llarg que sigui. Tot això és accessible amb SR
(capítol~\ref{cap:alephs}), sempre tram a tram: ni amb SR no hi ha cap
objecte indexat per tots els ordinals;
i sense SR es queda fora ---no per una mancança del llenguatge, sinó per la
mateixa raó que ho queda a ETCS: manca el reemplaçament. Res d'això no és
una objecció: és la delimitació natural d'una teoria pensada per fonamentar
la matemàtica ordinària amb el mínim de compromisos, i que fa explícit,
mitjançant SR, el punt on cal pagar-ne un de més fort.

\begin{observacio}[Balanç del capítol i del llibre]\label{obs:balanc-metateoric}
Aquest capítol ha respost les tres preguntes amb què s'obria. \emph{És
realment una teoria de conjunts?} Sí: és una axiomatització d'ETCS,
autònoma, en un llenguatge d'aplicacions
(observació~\ref{obs:tesi-etcs}), i la categoria dels conjunts clàssics
n'és un model (secció~\ref{sec:conjunts-petits}). \emph{Què pot expressar?}
Tota la matemàtica que viu llevat d'isomorfisme ---que és, a la pràctica,
tota la matemàtica ordinària---, amb la força de la teoria de Zermelo amb
elecció; i, dins d'aquesta, tota la teoria clàssica de l'infinit ---Cantor,
bons ordres, recursió transfinita, ordinals i cardinals--- \emph{acotada},
sense cap principi nou (secció~\ref{sec:transfinit-balanc}). Només la reunió
de col·leccions transfinites no acotades ---els trams de la successió dels alephs i
dels beths--- demana la substitució estructural SR, que porta la teoria a la
força de ZFC (secció~\ref{sec:etcs-sr}); el continu mateix
\(\mathfrak c=\abs{P(\N)}\) i la hipòtesi del continu, en canvi, són ja de
Nivell~I. \emph{Quina relació té amb la teoria clàssica?}
És equiconsistent i mútuament interpretable amb \(Z\) acotada amb elecció,
una mica més feble que ZFC, i coincideix amb ZFC quan s'hi afegeix SR ---la
traducció estructural del reemplaçament (observació~\ref{obs:forca-relativa},
secció~\ref{sec:etcs-sr}); la diferència essencial és que substitueix la
pertinença primitiva i global per una de derivada i local
(secció~\ref{sec:pertinenca}), i que tracta ordinals i cardinals com a
\emph{tipus} llevat d'isomorfisme en comptes de com a conjunts privilegiats
(secció~\ref{sec:vonneumann}).

Amb això es tanca el recorregut començat a la introducció. Vam prometre
construir els conjunts sense posar els elements al centre, prenent les
aplicacions com a únics primitius i recuperant-ho tot ---parts,
pertinença, productes, quocients, funcions, nombres i cardinalitat--- a
partir del domini, el codomini i la composició. El recorregut ha arribat,
partint d'aquí, fins a la teoria clàssica de l'infinit sencera: mides i
formes d'ordre, cada una amb la seva aritmètica, els dos grans teoremes
diagonals ---Cantor i Burali--Forti--- que impedeixen una mida o una
longitud màximes, i, amb SR, la jerarquia dels alephs i la formulació precisa
de la hipòtesi del continu. Tot dit en un llenguatge que posa en primer pla
les relacions entre conjunts en comptes de la seva composició interna. Qui
vulgui anar més enllà ---grans cardinals, forcing, o la traducció detallada
cap a l'univers de von Neumann i tornada--- té, en la
bibliografia~\cite{mclarty-catset,shulman-comparing,lawvere-rosebrugh}, els fils
per continuar.
\end{observacio}


% ============================================================
% PART II — PRÀCTICA
% Els capítols reprodueixen els noms i l'ordre de la part teòrica
% (el capítol 20 és un balanç sense capítols paral·lels).
% ============================================================
\part{Pràctica}\label{part:practica}
\setcounter{chapter}{0}
\chapter{Aplicacions, dominis, codominis i composició}\label{cap:practica-aplicacions}

Aquest capítol acompanya el capítol homònim de la part teòrica (capítol~\ref{cap:aplicacions} de la Part~\ref{part:teoria}). Els exercicis resolts en segueixen l'ordre d'exposició i serveixen de model per als exercicis proposats: cada bloc treballa una tècnica ---traduir al llenguatge dels predicats, encadenar igualtats amb neutralitat i associativitat, raonar amb propietats universals--- que després es reutilitza. Convé intentar cada enunciat abans de llegir-ne la solució.

\section{Llenguatge formal, composició i estabilitat}\label{prac:axiomes}

\begin{exercici}
\textbf{Traduir al llenguatge dels predicats.} Escriviu, usant només els predicats \(\Dom\), \(\Cod\) i \(\Com\), la frase \enquote{\(g\circ f\) està definida i té el mateix domini que \(f\)}. No feu servir la notació \(f\colon A\to B\).
\end{exercici}
\begin{solucio}
Introduïm variables per a la composició \(h\) i per al domini comú \(d\):
\[
\exists h\,\exists d\,\bigl(\Com(f,g,h)\land\Dom(f,d)\land\Dom(h,d)\bigr).
\]
El predicat \(\Com(f,g,h)\) afirma alhora que la composició està definida i que és \(h\); la conjunció \(\Dom(f,d)\land\Dom(h,d)\) expressa que \(h\) i \(f\) tenen el mateix domini.
\end{solucio}

\begin{exercici}
\textbf{El domini d'una aplicació és un conjunt.} Useu la proposició d'estabilitat (proposició~\ref{prop:estabilitat}) per demostrar que si \(\Dom(f,d)\), aleshores \(\Dom(d,d)\).
\end{exercici}
\begin{solucio}
La proposició d'estabilitat afirma que el domini de qualsevol aplicació és un conjunt, i que un conjunt és el seu propi domini. Com que \(d\) és el domini de \(f\), és un conjunt, i per tant \(\Dom(d,d)\). En el llenguatge del capítol, \(d\) és una identitat, i l'escrivim \(d=\dom(f)\).
\end{solucio}

\begin{exercici}
\textbf{Composició amb identitats.} Sigui \(f\colon A\to B\). Demostreu que \(f\circ\id_A=f\) i \(\id_B\circ f=f\), i identifiqueu quin axioma useu a cada pas.
\end{exercici}
\begin{solucio}
Per les lleis de neutralitat (axioma~\ref{ax:neutralitat}), tota identitat és neutra per la composició quan aquesta està definida. Com que \(\dom(f)=A\), la composició \(f\circ\id_A\) està definida i val \(f\); com que \(\cod(f)=B\), la composició \(\id_B\circ f\) està definida i val \(f\). Totes dues igualtats són instàncies directes de l'axioma de neutralitat.
\end{solucio}

\begin{exercici}
\textbf{El codomini de la composició.} Siguin \(f\colon A\to B\) i \(g\colon B\to C\). Demostreu, a partir de l'axioma sobre el domini i el codomini de la composició (axioma~\ref{ax:composicio-dom-cod}), que \(\cod(g\circ f)=C\) i \(\dom(g\circ f)=A\).
\end{exercici}
\begin{solucio}
L'axioma~\ref{ax:composicio-dom-cod} diu que, quan \(g\circ f\) està definida, el seu domini és el de \(f\) i el seu codomini és el de \(g\). La composició està definida perquè \(\cod(f)=B=\dom(g)\). Aplicant l'axioma, \(\dom(g\circ f)=\dom(f)=A\) i \(\cod(g\circ f)=\cod(g)=C\). Així \(g\circ f\colon A\to C\), com anticipava la notació tipada (notació~\ref{not:aplicacions-tipades}).
\end{solucio}

\section{Identitats i conjunts}\label{prac:identitats}

\begin{exercici}
\textbf{Els conjunts són les identitats.} Demostreu que un individu \(A\) de l'univers és un conjunt si i només si \(A=\dom(A)\).
\end{exercici}
\begin{solucio}
Per definició, \(A\) és un conjunt quan és una identitat, és a dir, quan \(\Dom(A,A)\land\Cod(A,A)\). Si \(A\) és un conjunt, aleshores \(\Dom(A,A)\), i com que el domini és únic (axioma~\ref{ax:dom-cod}), \(A=\dom(A)\). Recíprocament, si \(A=\dom(A)\), aleshores \(\Dom(A,A)\); per la proposició d'estabilitat, el domini \(\dom(A)=A\) és un conjunt, de manera que \(A\) ho és.
\end{solucio}

\begin{exercici}
\textbf{Endoaplicació que no és identitat.} Doneu, en el model clàssic, un exemple d'una aplicació amb domini igual al codomini que no sigui una identitat, i expliqueu per què no contradiu la unicitat de la identitat.
\end{exercici}
\begin{solucio}
Sigui \(A\) un conjunt amb dos elements globals diferents \(a_0,a_1\), i \(c\colon A\to A\) l'aplicació constant amb valor \(a_0\), és a dir, \(c\circ a_0=c\circ a_1=a_0\). El seu domini i codomini coincideixen (tots dos són \(A\)), de manera que és una endoaplicació. No és una identitat perquè \(c\circ a_1=a_0\neq a_1\), mentre que la identitat compliria \(\id_A\circ a_1=a_1\). No hi ha contradicció: la unicitat afirma que hi ha una sola aplicació \emph{neutra} de \(A\) en \(A\), no una sola endoaplicació; les endoaplicacions poden ser moltes, i només una d'elles és neutra.
\end{solucio}

\section{Igualtat d'aplicacions i elements generalitzats}\label{prac:igualtat}

\begin{exercici}
\textbf{Igualtat per elements generalitzats.} Enuncieu el criteri d'igualtat d'aplicacions per elements generalitzats (proposició~\ref{prop:aplicacions-iguals}) i useu-lo per demostrar que si \(f\circ x=g\circ x\) per a \emph{tot} element generalitzat \(x\colon X\to A\) (amb \(X\) arbitrari), aleshores \(f=g\).
\end{exercici}
\begin{solucio}
El criteri diu que dues aplicacions paral·leles \(f,g\colon A\to B\) són iguals si i només si coincideixen en compondre-les amb tots els elements generalitzats del domini. La hipòtesi és exactament l'antecedent d'aquest criteri, per a tots els estadis \(X\); en particular, prenent \(X=A\) i \(x=\id_A\), obtenim \(f=f\circ\id_A=g\circ\id_A=g\). El criteri general permet, doncs, concloure \(f=g\) directament.
\end{solucio}

\begin{exercici}
\textbf{Separació per elements globals.} Siguin \(f,g\colon A\to B\) dues aplicacions tals que \(f\circ x=g\circ x\) per a tot element global \(x\colon\mathbf1\to A\). Enuncieu quin axioma permet concloure \(f=g\) i apliqueu-lo, tot indicant en què difereix del criteri de l'exercici anterior.
\end{exercici}
\begin{solucio}
És l'axioma de separació d'aplicacions per elements globals (axioma~\ref{ax:separacio-globals}): si dues aplicacions paral·leles coincideixen després de compondre-les amb tots els elements \emph{globals} del domini, són iguals. Com que la hipòtesi és exactament \(f\circ x=g\circ x\) per a tot \(x\colon\mathbf1\to A\), l'axioma dona \(f=g\). La diferència amb la proposició~\ref{prop:aplicacions-iguals} és l'abast: aquell criteri quantifica sobre \emph{tots} els estadis \(X\), i val sempre; aquest axioma només fa servir l'estadi \(X=\mathbf1\), i és una hipòtesi addicional que recupera l'extensionalitat clàssica.
\end{solucio}

\section{Isomorfismes}\label{prac:isomorfismes}

\begin{exercici}
\textbf{Unicitat de la inversa.} Suposeu que \(g\) i \(g'\) són totes dues inverses d'un isomorfisme \(f\colon A\to B\). Demostreu que \(g=g'\), indicant on useu l'associativitat i on la neutralitat.
\end{exercici}
\begin{solucio}
Partim de \(g\) i hi inserim \(\id_B=f\circ g'\):
\[
g=g\circ\id_B=g\circ(f\circ g')=(g\circ f)\circ g'=\id_A\circ g'=g'.
\]
La segona igualtat usa que \(g'\) és inversa de \(f\) (per la dreta); la tercera és l'associativitat (axioma~\ref{ax:associativitat}); la quarta usa que \(g\) és inversa de \(f\) (per l'esquerra); i la primera i l'última són les lleis de neutralitat. Per tant, la inversa és única, i la podem denotar \(f^{-1}\).
\end{solucio}

\begin{exercici}
\textbf{La inversa d'una composició.} Siguin \(f\colon A\to B\) i \(g\colon B\to C\) isomorfismes. Demostreu que \(g\circ f\) és un isomorfisme amb invers \(f^{-1}\circ g^{-1}\).
\end{exercici}
\begin{solucio}
Comprovem les dues igualtats de la definició d'isomorfisme. D'una banda,
\[
(f^{-1}\circ g^{-1})\circ(g\circ f)=f^{-1}\circ(g^{-1}\circ g)\circ f=f^{-1}\circ\id_B\circ f=f^{-1}\circ f=\id_A,
\]
i, de l'altra, pel mateix camí,
\[
(g\circ f)\circ(f^{-1}\circ g^{-1})=g\circ(f\circ f^{-1})\circ g^{-1}=g\circ\id_B\circ g^{-1}=\id_C.
\]
Hem usat l'associativitat per reagrupar i les definicions de \(f^{-1}\) i \(g^{-1}\). Per tant, \(g\circ f\) és un isomorfisme i \({(g\circ f)}^{-1}=f^{-1}\circ g^{-1}\).
\end{solucio}

\section{Conjunts terminals}\label{prac:terminals}

\begin{exercici}
\textbf{Unicitat del conjunt terminal.} Demostreu que si \(T_1\) i \(T_2\) són tots dos terminals, hi ha un únic isomorfisme \(T_1\to T_2\).
\end{exercici}
\begin{solucio}
Com que \(T_2\) és terminal, hi ha exactament una aplicació \(t\colon T_1\to T_2\); com que \(T_1\) és terminal, hi ha exactament una \(s\colon T_2\to T_1\). La composició \(s\circ t\colon T_1\to T_1\) és una aplicació de \(T_1\) en ell mateix; però \(T_1\) és terminal, i l'única aplicació \(T_1\to T_1\) és \(\id_{T_1}\), de manera que \(s\circ t=\id_{T_1}\). Anàlogament, \(t\circ s=\id_{T_2}\). Així, \(t\) és un isomorfisme, i és únic perquè ja ho era com a aplicació.
\end{solucio}

\begin{exercici}
\textbf{Un element global d'un terminal.} Quantes aplicacions \(\mathbf1\to\mathbf1\) hi ha? Deduïu-ne que l'únic element global de \(\mathbf1\) és \(\id_{\mathbf1}\).
\end{exercici}
\begin{solucio}
Com que \(\mathbf1\) és terminal, hi ha exactament una aplicació de qualsevol conjunt cap a \(\mathbf1\); en particular, exactament una de \(\mathbf1\) en \(\mathbf1\). Aquesta aplicació ha de ser \(\id_{\mathbf1}\), que sempre existeix. Un element global de \(\mathbf1\) és una aplicació \(\mathbf1\to\mathbf1\), i com que només n'hi ha una, l'únic element global de \(\mathbf1\) és \(\id_{\mathbf1}\).
\end{solucio}

\section{Conjunts inicials, dualitat i no-degeneració}\label{prac:inicials}

\begin{exercici}
\textbf{Invariància de la inicialitat per isomorfisme.} Sigui \(I\) un conjunt inicial i \(I\cong J\). Demostreu que \(J\) també és inicial. (Model dual de la proposició~\ref{prop:terminal-invariant-iso}.)
\end{exercici}
\begin{solucio}
Sigui \(w\colon I\to J\) un isomorfisme amb invers \(w^{-1}\), i sigui \(A\) un conjunt qualsevol. Com que \(I\) és inicial, hi ha exactament una aplicació \(u\colon I\to A\). Definim \(v:=u\circ w^{-1}\colon J\to A\); existeix, doncs, almenys una aplicació \(J\to A\). Per veure que és única, sigui \(v'\colon J\to A\) qualsevol; aleshores \(v'\circ w\colon I\to A\) és una aplicació des de l'inicial \(I\), de manera que \(v'\circ w=u\), i component per la dreta amb \(w^{-1}\) obtenim \(v'=u\circ w^{-1}=v\). Així \(J\) té exactament una aplicació cap a cada conjunt, és a dir, és inicial. L'argument és el dual exacte del de la terminalitat: s'obté invertint el sentit de totes les fletxes.
\end{solucio}

\begin{exercici}
\textbf{El conjunt inicial no té elements globals.} Demostreu que no existeix cap aplicació \(\mathbf1\to\mathbf0\).
\end{exercici}
\begin{solucio}
Suposem que existís \(x\colon\mathbf1\to\mathbf0\). Com que \(\mathbf0\) és inicial, hi ha exactament una aplicació \(\mathbf0\to\mathbf1\), diguem \(u\). La composició \(x\circ u\colon\mathbf0\to\mathbf0\) és \(\id_{\mathbf0}\) per la inicialitat, i \(u\circ x\colon\mathbf1\to\mathbf1\) és \(\id_{\mathbf1}\) per la terminalitat. Aleshores \(\mathbf0\) i \(\mathbf1\) serien isomorfs, en contradicció amb l'axioma de no-degeneració (axioma~\ref{ax:no-degeneracio}). Per tant, no hi ha cap aplicació \(\mathbf1\to\mathbf0\).
\end{solucio}

\begin{exercici}
\textbf{Distingir amb la no-degeneració.} Expliqueu per què, sense l'axioma de no-degeneració, no es podria demostrar que \(\mathbf0\) i \(\mathbf1\) són diferents.
\end{exercici}
\begin{solucio}
Els axiomes d'existència del terminal i de l'inicial només afirmen que aquests conjunts existeixen i tenen les seves propietats universals; res en aquestes propietats no impedeix que un mateix conjunt sigui alhora inicial i terminal. De fet, en un univers amb un sol conjunt (i una sola aplicació, la seva identitat) totes dues propietats es compleixen trivialment. Cal, doncs, un axioma que exclogui aquest cas degenerat: és exactament el que fa l'axioma de no-degeneració, en postular que \(\mathbf0\) i \(\mathbf1\) no són isomorfs.
\end{solucio}

\begin{exercici}
\textbf{Dualitzar formalment.} Escriviu, amb els predicats \(\Dom\), \(\Cod\) i \(\Com\), la fórmula \(\operatorname{Terminal}(T)\) de la definició~\ref{def:terminal} i apliqueu-hi la transformació dual de l'observació~\ref{obs:dualitat}. Comproveu que el resultat és la fórmula \(\operatorname{Inicial}(T)\).
\end{exercici}
\begin{solucio}
La definició diu que, per a tot conjunt \(X\), existeix una única \(f\) amb \(\Dom(f,X)\land\Cod(f,T)\). La transformació dual intercanvia \(\Dom\) i \(\Cod\) (i no toca els quantificadors ni la igualtat), de manera que la fórmula es converteix en: per a tot conjunt \(X\), existeix una única \(f\) amb \(\Cod(f,X)\land\Dom(f,T)\), és a dir, una única \(f\colon T\to X\). Això és exactament la definició~\ref{def:inicial} d'\(\operatorname{Inicial}(T)\). La dualitat no és una analogia: és una substitució mecànica de símbols.
\end{solucio}

\begin{exercici}
\textbf{La separació global no és autodual.} Escriviu l'enunciat dual de l'axioma de separació global (axioma~\ref{ax:separacio-globals}) i demostreu que, afegit als altres axiomes, implicaria que dues aplicacions paral·leles qualssevol són iguals.
\end{exercici}
\begin{solucio}
L'axioma diu: si \(f,g\colon A\to B\) i \(f\circ x=g\circ x\) per a tot \(x\colon\mathbf1\to A\), aleshores \(f=g\). El dual intercanvia el sentit de les fletxes i l'ordre de la composició, i substitueix el terminal per l'inicial: si \(f,g\colon A\to B\) i \(y\circ f=y\circ g\) per a tot \(y\colon B\to\mathbf0\), aleshores \(f=g\). Ara bé, si \(B\) no és inicial no hi ha cap aplicació \(B\to\mathbf0\) (una tal \(y\) donaria, amb \(\mathbf0_B\), un isomorfisme \(B\cong\mathbf0\) pel mateix argument de la proposició~\ref{prop:no-fletxa-u-zero}), de manera que la hipòtesi \enquote{per a tota \(y\)} es compleix buidament i l'enunciat dual conclouria \(f=g\) per a tot parell d'aplicacions paral·leles. Per això el principi de dualitat deixa de valer a partir d'aquest axioma: la teoria dels conjunts constants no és simètrica.
\end{solucio}

\section{Elements globals i aplicacions constants}\label{prac:constants}

\begin{exercici}
\textbf{Una constant és constant després de compondre.} Sigui \(c\colon A\to B\) l'aplicació constant de valor \(b\colon\mathbf1\to B\) (definició~\ref{def:constant}), i sigui \(u\colon X\to A\) qualsevol. Demostreu que \(c\circ u\) és també constant de valor \(b\).
\end{exercici}
\begin{solucio}
Per definició, \(c=b\circ{!}_A\), on \({!}_A\colon A\to\mathbf1\) és l'única aplicació cap al terminal. Aleshores
\[
c\circ u=(b\circ{!}_A)\circ u=b\circ({!}_A\circ u).
\]
Ara \({!}_A\circ u\colon X\to\mathbf1\) és una aplicació cap al terminal, i com que n'hi ha exactament una, \({!}_A\circ u={!}_X\). Per tant \(c\circ u=b\circ{!}_X\), que és, per definició, l'aplicació constant \(X\to B\) de valor \(b\). Hem usat només l'associativitat i la propietat universal del terminal.
\end{solucio}

\begin{exercici}
\textbf{El valor d'una constant sobre un punt.} Amb la notació de l'exercici anterior, demostreu que per a tot element global \(x\colon\mathbf1\to A\) es té \(c\circ x=b\). Interpreteu-ho en la lectura clàssica.
\end{exercici}
\begin{solucio}
Prenent \(X=\mathbf1\) i \(u=x\) a l'exercici anterior, \(c\circ x=b\circ{!}_{\mathbf1}\). Però \({!}_{\mathbf1}\colon\mathbf1\to\mathbf1\) és l'única aplicació del terminal en ell mateix, que és \(\id_{\mathbf1}\); per tant \(c\circ x=b\circ\id_{\mathbf1}=b\). En la lectura clàssica (secció~\ref{sec:lectura-classica}), això diu que la funció constant de valor \(b\) envia \emph{tot} punt del domini al mateix element \(b\), com és d'esperar.
\end{solucio}

\chapter{Parts i construccions universals}\label{cap:practica-parts}

Aquest capítol acompanya el capítol homònim de la part teòrica (capítol~\ref{cap:parts} de la Part~\ref{part:teoria}). Els exercicis resolts en segueixen l'ordre d'exposició --- monomorfismes i epimorfismes, parts i inclusió, pertinença generalitzada, sumes, aplicacions característiques, productes, equalitzadors, pullbacks i imatges inverses--- i serveixen de model per als exercicis proposats de la Part~\ref{part:exercicis}: cada bloc fixa una tècnica (cancel·lar per un mono, exhibir el testimoni d'una inclusió o d'una pertinença, comprovar una propietat universal per unicitat de factorització) que després es reutilitza. Convé intentar cada enunciat abans de llegir-ne la solució.

\section{Cancel·lació per un monomorfisme}\label{prac:cancel-lacio}

\begin{exercici}
\textbf{Un ús de la cancel·lació.} Sigui \(m:A\to B\) un monomorfisme i siguin \(u,v:X\to A\) amb
\[
m\circ u=m\circ v.
\]
Demostreu que \(u=v\).
\end{exercici}

\begin{solucio}
Per definició (definició~\ref{def:monomorfisme}), un monomorfisme es pot cancel·lar per l'esquerra: per a tot parell d'aplicacions paral·leles \(u,v:X\to A\),
\[
m\circ u=m\circ v\quad\Longrightarrow\quad u=v.
\]
La hipòtesi és exactament l'antecedent d'aquesta implicació; per tant, \(u=v\).
\end{solucio}

\begin{exercici}
\textbf{Dualitzeu-ho: cancel·lació per un epimorfisme.} Sigui \(e:A\to B\) un epimorfisme i siguin \(u,v:B\to Y\) amb \(u\circ e=v\circ e\). Demostreu que \(u=v\).
\end{exercici}
\begin{solucio}
Un epimorfisme (definició~\ref{def:epimorfisme}) es cancel·la per la dreta: per a tot parell d'aplicacions paral·leles \(u,v:B\to Y\), de \(u\circ e=v\circ e\) se'n segueix \(u=v\). La hipòtesi és l'antecedent, i s'obté \(u=v\). L'enunciat és el dual estricte de l'exercici anterior: s'obté intercanviant \enquote{per l'esquerra} per \enquote{per la dreta} i invertint el sentit de les fletxes.
\end{solucio}

\section{Isomorfismes, monomorfismes i epimorfismes}\label{prac:iso-mono-epi}

\begin{exercici}
\textbf{Isomorfisme implica monomorfisme.} Demostreu que tot isomorfisme és un monomorfisme, sense fer servir cap resultat que no sigui els axiomes i la definició d'inversa.
\end{exercici}
\begin{solucio}
Sigui \(f\colon A\to B\) un isomorfisme amb invers \(f^{-1}\), i suposem \(f\circ u=f\circ v\). Component per l'esquerra amb \(f^{-1}\):
\[
u=\id_A\circ u=(f^{-1}\circ f)\circ u=f^{-1}\circ(f\circ u)=f^{-1}\circ(f\circ v)=(f^{-1}\circ f)\circ v=\id_A\circ v=v.
\]
Hem usat la neutralitat, la definició d'invers i l'associativitat. Com que \(f\circ u=f\circ v\) implica \(u=v\), \(f\) és un monomorfisme. Aquest és exactament el contingut de la proposició~\ref{prop:isomorfisme-mono-epi}, que afirma que tot isomorfisme és alhora mono i epi.
\end{solucio}

\begin{exercici}
\textbf{La composició de monomorfismes és un monomorfisme.} Siguin \(m:A\rightarrowtail B\) i \(n:B\rightarrowtail C\) dos monomorfismes. Demostreu que \(n\circ m\) és un monomorfisme. (Model de la proposició~\ref{prop:composicio-mono-epi}.)
\end{exercici}
\begin{solucio}
Suposem \((n\circ m)\circ u=(n\circ m)\circ v\) amb \(u,v:X\to A\). Per associativitat (axioma~\ref{ax:associativitat}),
\[
n\circ(m\circ u)=n\circ(m\circ v).
\]
Com que \(n\) és un monomorfisme, es cancel·la per l'esquerra i \(m\circ u=m\circ v\). Com que \(m\) també ho és, \(u=v\). Per tant \(n\circ m\) és un monomorfisme. El mateix argument, dualitzat, prova que la composició d'epimorfismes és un epimorfisme.
\end{solucio}

\begin{exercici}
\textbf{Cancel·lació posterior detecta monomorfismes.} Siguin \(f:A\to B\) i \(g:B\to C\) tals que \(g\circ f\) és un monomorfisme. Demostreu que \(f\) és un monomorfisme.
\end{exercici}
\begin{solucio}
Suposem \(f\circ u=f\circ v\). Component per l'esquerra amb \(g\) i usant l'associativitat,
\[
(g\circ f)\circ u=g\circ(f\circ u)=g\circ(f\circ v)=(g\circ f)\circ v.
\]
Com que \(g\circ f\) és un monomorfisme, \(u=v\). Per tant \(f\) és un monomorfisme. Observeu que no cal cap hipòtesi sobre \(g\): només que el compost sigui mono. El resultat dual diu que si \(g\circ f\) és un epimorfisme, aleshores \(g\) ho és.
\end{solucio}

\section{Seccions, retraccions i escindits}\label{prac:seccio-retraccio}

\begin{exercici}
\textbf{Tota secció és un monomorfisme.} Sigui \(s:A\to B\) una secció, és a dir, existeix \(r:B\to A\) amb \(r\circ s=\id_A\). Demostreu que \(s\) és un monomorfisme. (Meitat de la proposició~\ref{prop:seccio-mono-retraccio-epi}.)
\end{exercici}
\begin{solucio}
Suposem \(s\circ u=s\circ v\). Component per l'esquerra amb la retracció \(r\):
\[
u=\id_A\circ u=(r\circ s)\circ u=r\circ(s\circ u)=r\circ(s\circ v)=(r\circ s)\circ v=\id_A\circ v=v.
\]
Per tant \(s\) és un monomorfisme. Dualment, tota retracció és un epimorfisme: un mono (resp.\ epi) que a més és secció (resp.\ retracció) s'anomena \emph{escindit} (observació~\ref{obs:escindits}).
\end{solucio}

\begin{exercici}
\textbf{Exhaustivitat generalitzada d'un epimorfisme escindit.} Sigui \(e:A\to B\) una retracció, amb secció \(s:B\to A\) tal que \(e\circ s=\id_B\). Demostreu que tot element generalitzat \(y:X\to B\) es factoritza com \(y=e\circ x\) per a algun \(x:X\to A\). (Model de la proposició~\ref{prop:exhaustiva-generalitzada-escindit}.)
\end{exercici}
\begin{solucio}
Preneu \(x=s\circ y:X\to A\). Aleshores, per associativitat i la relació \(e\circ s=\id_B\),
\[
e\circ x=e\circ(s\circ y)=(e\circ s)\circ y=\id_B\circ y=y.
\]
Així \(y\) factoritza per \(e\), de manera que \(e\) és exhaustiva sobre els elements generalitzats (definició~\ref{def:exhaustiva-generalitzada}). La secció proporciona explícitament una elecció de preimatge, \(x=s\circ y\).
\end{solucio}

\section{Parts i inclusió}\label{prac:parts-inclusio}

\begin{exercici}
\textbf{Unicitat del testimoni d'una inclusió.} Siguin \(m:S\rightarrowtail A\) i \(n:T\rightarrowtail A\) dues parts amb \(m\subseteq_A n\). Demostreu que l'aplicació \(k:S\to T\) amb \(m=n\circ k\) és única, i que és un monomorfisme. (Reproduïu la proposició~\ref{prop:testimoni-inclusio}.)
\end{exercici}
\begin{solucio}
Si \(m=n\circ k=n\circ k'\), com que \(n\) és un monomorfisme es cancel·la per l'esquerra i \(k=k'\); això dona la unicitat. Per veure que \(k\) és mono, suposem \(k\circ u=k\circ v\). Aleshores
\[
m\circ u=(n\circ k)\circ u=n\circ(k\circ u)=n\circ(k\circ v)=(n\circ k)\circ v=m\circ v,
\]
i, com que \(m\) és mono, \(u=v\). El testimoni d'una inclusió és, doncs, un monomorfisme únic.
\end{solucio}

\begin{exercici}
\textbf{La inclusió és reflexiva i transitiva.} Sigui \(m:S\rightarrowtail A\) una part. Demostreu que \(m\subseteq_A m\), i que si \(m\subseteq_A n\) i \(n\subseteq_A p\), aleshores \(m\subseteq_A p\).
\end{exercici}
\begin{solucio}
Per la reflexivitat, prenem \(k=\id_S\); aleshores \(m=m\circ\id_S\), i per definició d'inclusió (definició~\ref{def:inclusio}) \(m\subseteq_A m\). Per la transitivitat, siguin \(k:S\to T\) i \(l:T\to U\) els testimonis de \(m=n\circ k\) i \(n=p\circ l\). Aleshores
\[
m=n\circ k=(p\circ l)\circ k=p\circ(l\circ k),
\]
de manera que \(l\circ k\) testimonia \(m\subseteq_A p\). La relació d'inclusió és, per tant, un preordre sobre les parts de \(A\).
\end{solucio}

\section{Pertinença generalitzada}\label{prac:pertinenca}

\begin{exercici}
\textbf{Unicitat del testimoni d'una pertinença.} Siguin \(m:S\rightarrowtail A\) una part i \(x:X\to A\) amb \(x\in_X m\). Demostreu que el testimoni \(\bar x:X\to S\) amb \(x=m\circ\bar x\) és únic. (Reproduïu la proposició~\ref{prop:testimoni-pertinenca}.)
\end{exercici}
\begin{solucio}
Si \(x=m\circ\bar x=m\circ\bar x'\), com que \(m\) és un monomorfisme es cancel·la per l'esquerra i \(\bar x=\bar x'\). Per tant, quan un element generalitzat pertany a una part, ho fa \enquote{d'una sola manera}: el seu testimoni queda determinat.
\end{solucio}

\begin{exercici}
\textbf{Inclusió com a implicació de pertinences.} Siguin \(m:S\rightarrowtail A\) i \(n:T\rightarrowtail A\). Demostreu que si \(m\subseteq_A n\), aleshores per a tot \(x:X\to A\) amb \(x\in_X m\) es té \(x\in_X n\). (Una direcció de la proposició~\ref{prop:inclusio-pertinenca}.)
\end{exercici}
\begin{solucio}
Sigui \(k:S\to T\) el testimoni de \(m=n\circ k\), i \(\bar x:X\to S\) el de \(x=m\circ\bar x\). Aleshores
\[
x=m\circ\bar x=(n\circ k)\circ\bar x=n\circ(k\circ\bar x),
\]
de manera que \(k\circ\bar x:X\to T\) és un testimoni de \(x\in_X n\) (vegeu la figura~\ref{fig:fonamental} de la teoria). La inclusió de parts es tradueix, doncs, en \enquote{tot element de \(S\) és element de \(T\)}, ara en el sentit generalitzat.
\end{solucio}

\section{Sumes i definició per casos}\label{prac:sumes}

\begin{exercici}
\textbf{Unicitat de l'aplicació definida per casos.} Sigui \(A+B\) una suma amb injeccions \(i_A,i_B\), i siguin \(f:A\to C\), \(g:B\to C\). Demostreu que hi ha \emph{com a molt} una aplicació \(h:A+B\to C\) amb \(h\circ i_A=f\) i \(h\circ i_B=g\). (És la unicitat a la propietat universal; vegeu la proposició~\ref{prop:extensionalitat-suma}.)
\end{exercici}
\begin{solucio}
Suposem que \(h\) i \(h'\) compleixen totes dues les equacions: \(h\circ i_A=f=h'\circ i_A\) i \(h\circ i_B=g=h'\circ i_B\). La propietat universal de la suma (definició~\ref{def:suma-conjunts}) diu que una aplicació que surt de \(A+B\) queda determinada pels seus dos components sobre les injeccions; en altres paraules, dues aplicacions \(A+B\to C\) que coincideixen en compondre amb \(i_A\) i amb \(i_B\) són iguals. Per tant \(h=h'\). L'existència la garanteix la mateixa propietat universal; aquest exercici n'aïlla la unicitat.
\end{solucio}

\begin{exercici}
\textbf{El cograf com a retracció.} Donada \(f:A\to B\), considereu el cograf \(\kappa_f=[f,\id_B]:A+B\to B\) (definició~\ref{def:cograf}). Comproveu, a partir de les equacions de la suma, que la injecció \(i_1:B\to A+B\) admet \(\kappa_f\) com a retracció, i expliqueu per què el mateix argument no dona, sense més, que \(i_0:A\to A+B\) sigui un monomorfisme. (Model de la proposició~\ref{prop:cograf-retraccio}.)
\end{exercici}
\begin{solucio}
L'aplicació \(\kappa_f=[f,\id_B]\) existeix per la propietat universal de la suma i compleix \(\kappa_f\circ i_0=f\) i \(\kappa_f\circ i_1=\id_B\). La segona equació diu precisament que \(\kappa_f\) és una retracció de la injecció \(i_1\); per tant \(i_1\) és una secció i, en particular, un monomorfisme escindit (exercici de la secció~\ref{prac:seccio-retraccio}). Per fer el mateix amb \(i_0\) necessitaríem una aplicació \(g\colon B\to A\) amb què formar \([\id_A,g]\colon A+B\to A\), i una tal \(g\) pot no existir (per exemple, si \(B=\mathbf1\) i \(A\) no té punts). Que \emph{totes} les injeccions siguin monomorfismes és conseqüència de la disjunció de les sumes, al capítol~\ref{cap:valors}; observeu que la injecció \(i_0\) de l'exemple, \(\mathbf0\to\mathbf0+\mathbf1\), és mono per una altra raó: tota aplicació amb domini \(\mathbf0\) ho és, perquè dues aplicacions paral·leles cap a \(\mathbf0\) sempre són iguals.
\end{solucio}

\section{El conjunt \texorpdfstring{\(\mathbf2\)}{2} i aplicacions característiques}\label{prac:caracteristica}

\begin{exercici}
\textbf{Una part està determinada per la seva característica.} Sigui \(m:S\rightarrowtail A\) una part amb aplicació característica \(\chi_m:A\to\mathbf2\) (definició~\ref{def:caracteristica}). Demostreu que si dues parts \(m,n\) de \(A\) tenen la mateixa característica, \(\chi_m=\chi_n\), aleshores són equivalents. (Model de la proposició~\ref{prop:caracteristica-determina-part}.)
\end{exercici}
\begin{solucio}
La característica de \(m\) està definida per la condició \(x\in_X m\iff\chi_m\circ x=1\circ{!}_X\), per a tot estadi \(X\) i tot \(x\colon X\to A\) (notació~\ref{not:caracteristica}). Si \(\chi_m=\chi_n\), aleshores, per a tot \(x\),
\[
x\in_X m\iff\chi_m\circ x=1\circ{!}_X\iff\chi_n\circ x=1\circ{!}_X\iff x\in_X n,
\]
és a dir, \(m\) i \(n\) tenen exactament els mateixos elements generalitzats a tots els estadis. Per la proposició~\ref{prop:inclusio-pertinenca}, aplicada en els dos sentits, \(m\subseteq_A n\) i \(n\subseteq_A m\); per tant \(m\equiv_A n\) (definició~\ref{def:parts-equivalents}). Juntament amb la unicitat de la característica que dona l'axioma de representació (axioma~\ref{ax:representacio}), això fa de \(\chi_m\) un invariant complet de la classe d'equivalència d'una part.
\end{solucio}

\begin{exercici}
\textbf{Bimorfisme més representació dona isomorfisme.} Useu el context de les aplicacions característiques per explicar per què, en aquesta teoria, tot bimorfisme és un isomorfisme. (Enunciat: proposició~\ref{prop:bimorfisme-isomorfisme}.)
\end{exercici}
\begin{solucio}
Un bimorfisme (definició~\ref{def:bimorfisme}) és alhora mono i epi. Com a monomorfisme, \(f:A\rightarrowtail B\) és una part de \(B\), amb característica \(\chi_f:B\to\mathbf2\). L'argument té tres passos. Primer, \(f\) pertany a la part \(f\) amb testimoni \(\id_A\), de manera que \(\chi_f\circ f=1\circ{!}_A\); però també \((1\circ{!}_B)\circ f=1\circ{!}_A\), i com que \(f\) és un \emph{epimorfisme}, \(\chi_f=1\circ{!}_B\). Segon, això diu que \(\chi_f\circ\id_B=1\circ{!}_B\), és a dir, \(\id_B\in_B f\): existeix \(g\colon B\to A\) amb \(f\circ g=\id_B\). Tercer, \(f\circ(g\circ f)=(f\circ g)\circ f=f=f\circ\id_A\), i com que \(f\) és un \emph{monomorfisme}, \(g\circ f=\id_A\). Així \(g\) és la inversa de \(f\). Observeu on intervé cada hipòtesi: l'epimorfisme força la característica a ser la total (proposició~\ref{prop:caracteristica-total}), la representació converteix això en una secció, i el monomorfisme fa de la secció una inversa. El corol·lari~\ref{cor:isomorfisme-bimorfisme} en recull el recíproc, ja provat abans.
\end{solucio}

\section{Productes i aparellament}\label{prac:productes}

\begin{exercici}
\textbf{Les projeccions separen aplicacions.} Sigui \(A\times B\) un producte amb projeccions \(p_A,p_B\), i siguin \(u,v:X\to A\times B\) amb \(p_A\circ u=p_A\circ v\) i \(p_B\circ u=p_B\circ v\). Demostreu que \(u=v\). (Model de la proposició~\ref{prop:projeccions-separen}.)
\end{exercici}
\begin{solucio}
La propietat universal del producte (definició~\ref{def:producte}) diu que una aplicació cap a \(A\times B\) queda unívocament determinada per les seves dues components sobre les projeccions: donades \(f:X\to A\) i \(g:X\to B\), existeix una única \(\langle f,g\rangle:X\to A\times B\) amb \(p_A\circ\langle f,g\rangle=f\) i \(p_B\circ\langle f,g\rangle=g\). Ara bé, tant \(u\) com \(v\) tenen les mateixes components, \(f:=p_A\circ u=p_A\circ v\) i \(g:=p_B\circ u=p_B\circ v\); per la \emph{unicitat} de la factorització, \(u=\langle f,g\rangle=v\). Les projeccions, doncs, són conjuntament \enquote{cancel·lables}.
\end{solucio}

\begin{exercici}
\textbf{Unicitat del producte llevat d'isomorfisme.} Suposeu que \((P,p_A,p_B)\) i \((Q,q_A,q_B)\) satisfan totes dues la propietat universal del producte de \(A\) i \(B\). Esbosseu per què \(P\cong Q\). (Enunciat: proposició~\ref{prop:producte-unic}.)
\end{exercici}
\begin{solucio}
Per la propietat universal de \(Q\), les aplicacions \(p_A,p_B\) que surten de \(P\) indueixen una única \(\varphi:P\to Q\) amb \(q_A\circ\varphi=p_A\) i \(q_B\circ\varphi=p_B\). Simètricament, la propietat universal de \(P\) indueix \(\psi:Q\to P\) amb \(p_A\circ\psi=q_A\) i \(p_B\circ\psi=q_B\). El compost \(\psi\circ\varphi:P\to P\) satisfà les mateixes equacions que \(\id_P\) respecte de \(p_A,p_B\); per la unicitat de la factorització (l'exercici anterior), \(\psi\circ\varphi=\id_P\). Igualment \(\varphi\circ\psi=\id_Q\). Per tant \(\varphi\) és un isomorfisme i \(P\cong Q\). Aquest és el patró estàndard: \emph{tot objecte definit per una propietat universal és únic llevat d'isomorfisme}.
\end{solucio}

\section{Equalitzadors, pullbacks i imatges inverses}\label{prac:equalitzadors-pullbacks}

\begin{exercici}
\textbf{L'equalitzador és un monomorfisme.} Sigui \(e:E\to A\) l'equalitzador de \(f,g:A\to B\) (definició~\ref{def:equalitzador}). Demostreu que \(e\) és un monomorfisme.
\end{exercici}
\begin{solucio}
Suposem \(e\circ u=e\circ v\) amb \(u,v:X\to E\). Sigui \(w:=e\circ u=e\circ v:X\to A\). Aleshores \(f\circ w=g\circ w\) (perquè \(f\circ e=g\circ e\), i \(w\) factoritza per \(e\)); és a dir, \(w\) equalitza \(f\) i \(g\). Per la propietat universal de l'equalitzador (teorema~\ref{teo:equalitzador-universal}), existeix una \emph{única} aplicació \(t:X\to E\) amb \(e\circ t=w\). Però tant \(u\) com \(v\) compleixen \(e\circ u=w=e\circ v\); per la unicitat, \(u=t=v\). Per tant \(e\) és un monomorfisme.
\end{solucio}

\begin{exercici}
\textbf{El pullback d'un monomorfisme és un monomorfisme.} En un pullback (definició~\ref{def:pullback}) de \(f:A\to C\) al llarg de \(g:B\to C\), suposeu que \(f\) és un monomorfisme. Demostreu que el canvi de base \(g^{*}(f)\), la projecció cap a \(B\) (observació~\ref{obs:pullback-exemples}), també ho és. (Enunciat: proposició~\ref{prop:pullback-mono}.)
\end{exercici}
\begin{solucio}
Sigui el pullback \(P\) amb projeccions \(p:P\to A\) i \(q:P\to B\), tals que \(f\circ p=g\circ q\); la projecció \(q\) és la que anomenem \(g^{*}(f)\). Suposem \(q\circ u=q\circ v\) amb \(u,v:X\to P\). Volem \(u=v\). Component les equacions del quadrat, \(f\circ(p\circ u)=g\circ(q\circ u)=g\circ(q\circ v)=f\circ(p\circ v)\); com que \(f\) és mono, \(p\circ u=p\circ v\). Així, \(u\) i \(v\) tenen la mateixa composició amb totes dues projeccions \(p\) i \(q\). Per la unicitat a la propietat universal del pullback (teorema~\ref{teo:pullback-universal}), les dues aplicacions \(X\to P\) que factoritzen el mateix con són iguals: \(u=v\). Per tant \(g^{*}(f)=q\) és un monomorfisme. Aquest és el fet que fa que la imatge inversa d'una part torni a ser una part (definició~\ref{def:imatge-inversa}).
\end{solucio}

\begin{exercici}
\textbf{Pertinença a la imatge inversa.} Siguin \(g:B\to A\) i una part \(m:S\rightarrowtail A\), amb imatge inversa \(g^{-1}(m):P\rightarrowtail B\). Demostreu que, per a tot \(y:X\to B\), es té \(y\in_X g^{-1}(m)\) si i només si \(g\circ y\in_X m\). (Model de la proposició~\ref{prop:pertinenca-imatge-inversa}.)
\end{exercici}
\begin{solucio}
La imatge inversa és el pullback de \(m\) al llarg de \(g\): un quadrat cartesià amb projeccions \(q_B\colon P\to B\) (que és \(g^{-1}(m)\)) i \(q_S\colon P\to S\), amb \(g\circ q_B=m\circ q_S\). Si \(y\in_X g^{-1}(m)\), hi ha \(w\colon X\to P\) amb \(y=q_B\circ w\), i aleshores \(g\circ y=g\circ q_B\circ w=m\circ(q_S\circ w)\), de manera que \(q_S\circ w\) testimonia \(g\circ y\in_X m\). Recíprocament, si \(g\circ y=m\circ z\) per a algun \(z\colon X\to S\), el parell \((y,z)\) compleix la condició de commutativitat del pullback, i la propietat universal (teorema~\ref{teo:pullback-universal}) dona \(w\colon X\to P\) amb \(q_B\circ w=y\) i \(q_S\circ w=z\); la primera igualtat és \(y\in_X g^{-1}(m)\). En termes clàssics: \(y\) cau dins \(g^{-1}(m)\) exactament quan la seva imatge \(g\circ y\) cau dins \(m\).
\end{solucio}

\chapter[Valors de veritat i construccions finites]{Valors de veritat, parts\texorpdfstring{\\}{ }i construccions finites}\label{cap:practica-valors}

Aquest capítol acompanya el capítol homònim de la part teòrica (capítol~\ref{cap:valors} de la Part~\ref{part:teoria}). Els exercicis resolts en segueixen l'ordre d'exposició --- sumes disjuntes i estables, l'estrictesa de l'inicial, les parts extremals, la negació i la part complementària, les operacions booleanes, i les construccions finites (límits, coequalitzadors, pushouts i colímits)--- i serveixen de model per als exercicis proposats de la Part~\ref{part:exercicis}. Cada bloc fixa una tècnica: raonar amb l'axioma de sumes disjuntes, provar igualtats d'aplicacions cap a \(\mathbf2\) per extensionalitat de la definició per casos, traduir operacions de parts a operacions sobre característiques, i comprovar propietats universals per unicitat de factorització. Convé intentar cada enunciat abans de llegir-ne la solució.

\section{Sumes disjuntes i estables}\label{prac:sumes-disjuntes}

\begin{exercici}
\textbf{Monomorfismes i pullbacks.} Demostreu que una aplicació \(m:S\to A\) és un monomorfisme si i només si el quadrat format per \(S\), dues còpies de \(\id_S\) i dues còpies de \(m\) és cartesià, és a dir, si i només si \(S\), amb projeccions \(\id_S,\id_S\), és un pullback de \(m\) amb ella mateixa. Aquesta caracterització es farà servir al llarg del capítol, en particular per a les injeccions d'una suma disjunta.
\end{exercici}
\begin{solucio}
El quadrat sempre commuta: \(m\circ\id_S=m\circ\id_S\). Suposem que \(m\) és mono, i siguin \(u,v:X\to S\) amb \(m\circ u=m\circ v\), és a dir, un parell compatible per al pullback de \(m\) amb ella mateixa. Com que \(m\) és mono, \(u=v\), i l'aplicació \(w:=u:X\to S\) compleix \(\id_S\circ w=u\) i \(\id_S\circ w=v\); és única perquè \(\id_S\circ w=u\) la determina. Per tant \(S\) és el pullback. Recíprocament, si el quadrat és cartesià i \(m\circ u=m\circ v\), la propietat universal dona una única \(w\) amb \(\id_S\circ w=u\) i \(\id_S\circ w=v\), d'on \(u=w=v\): \(m\) és mono. Així, \enquote{ser mono} és una condició sobre un pullback, cosa que explica per què la disjunció d'una suma (definició~\ref{def:suma-disjunta}) es formula en termes de pullbacks i de monomorfismes alhora.
\end{solucio}

\begin{exercici}
\textbf{Estabilitat: el canvi de base preserva les sumes.} Enuncieu què afirma l'axioma de disjunció i estabilitat (axioma~\ref{ax:sumes-disjuntes}) sobre el pullback d'una suma al llarg d'una aplicació qualsevol, i expliqueu per què d'aquí se segueix que tota suma és disjunta.
\end{exercici}
\begin{solucio}
L'axioma~\ref{ax:sumes-disjuntes} diu dues coses: (i) tota suma és \emph{disjunta} (definició~\ref{def:suma-disjunta}), és a dir, el pullback de les dues injeccions és inicial i cada injecció és mono; i (ii) les sumes són \emph{estables} per canvi de base (definició~\ref{def:suma-estable}): si es fa el pullback d'una suma \(A_0\to A_0+A_1\leftarrow A_1\) al llarg de qualsevol \(f:X\to A_0+A_1\), el resultat torna a ser una suma, amb vèrtex \(X\). La disjunció és, doncs, un cas particular postulat directament; l'estabilitat és el que permet propagar la descomposició de \(A_0+A_1\) a tot element generalitzat, com fa el teorema~\ref{teo:descomposicio-canvi-base}.
\end{solucio}

\section{L'inicial estricte i la distinció \texorpdfstring{\(0\neq1\)}{0!=1}}\label{prac:zero-u}

\begin{exercici}
\textbf{Tota aplicació cap a l'inicial és un isomorfisme.} Reproduïu l'argument que mostra que \(\mathbf0\) és estricte: si \(f:X\to\mathbf0\), aleshores \(f\) és un isomorfisme. (Enunciat: proposició~\ref{prop:inicial-estricte}.)
\end{exercici}
\begin{solucio}
Les dues còpies de \(\id_{\mathbf0}\) formen una suma \(\mathbf0\to\mathbf0\leftarrow\mathbf0\) (qualsevol parell d'aplicacions que surten de l'inicial coincideix amb l'única fletxa de l'inicial). Fem el pullback d'aquesta suma al llarg de \(f:X\to\mathbf0\): el pullback de \(\id_{\mathbf0}\) al llarg de \(f\) és \(X\) amb projecció \(\id_X\). Per l'estabilitat (axioma~\ref{ax:sumes-disjuntes}), obtenim una suma \(X\to X\leftarrow X\) amb totes dues injeccions \(\id_X\); i, essent suma, és disjunta, de manera que el pullback de \(\id_X\) amb \(\id_X\) --- que és \(X\)--- és inicial. Així \(X\) és inicial, i com que l'inicial és únic llevat d'isomorfisme, hi ha \(u:X\xrightarrow{\sim}\mathbf0\). Com que de l'inicial \(X\) surt una sola aplicació cap a \(\mathbf0\), forçosament \(f=u\), i \(f\) és un isomorfisme.
\end{solucio}

\begin{exercici}
\textbf{Tot conjunt sense punts globals és inicial.} Demostreu el lema del buit (lema~\ref{lema:buit}): si \(X\) no té cap element global \(\mathbf1\to X\), aleshores \(X\cong\mathbf0\).
\end{exercici}
\begin{solucio}
Considerem les dues aplicacions \(0\circ{!}_X\) i \(1\circ{!}_X\) de \(X\) a \(\mathbf2\). Coincideixen sobre \emph{tots} els elements globals de \(X\), simplement perquè \(X\) no en té cap; per l'axioma de separació global (axioma~\ref{ax:separacio-globals}), \(0\circ{!}_X=1\circ{!}_X\). Aquesta igualtat diu que \(({!}_X,{!}_X)\) és un con sobre el pullback de \(0\) i \(1\); la propietat universal dona una aplicació \(X\to P\) cap a aquest pullback. Però la suma \(\mathbf1\xrightarrow{0}\mathbf2\xleftarrow{1}\mathbf1\) és disjunta, de manera que \(P\) és inicial; component, obtenim \(X\to\mathbf0\), i per l'estrictesa (proposició~\ref{prop:inicial-estricte}) \(X\cong\mathbf0\). Contrapositivament: tot conjunt no inicial té almenys un punt global.
\end{solucio}

\begin{exercici}
\textbf{Els dos valors de veritat són diferents.} A partir del fet que la suma \(\mathbf1\xrightarrow{0}\mathbf2\xleftarrow{1}\mathbf1\) és disjunta, demostreu que \(0\neq1\) com a elements globals de \(\mathbf2\). (Model del teorema~\ref{teo:zero-diferent-u}.)
\end{exercici}
\begin{solucio}
Com que la suma és disjunta, un pullback de \(0\) i \(1\) és inicial. Suposem, per obtenir contradicció, que \(0=1\). Aleshores aquesta aplicació comuna, essent un element global \(\mathbf1\to\mathbf2\), és un monomorfisme (tot punt \(\mathbf1\to A\) ho és, perquè \(\mathbf1\) és terminal), i el pullback d'un monomorfisme amb ell mateix és el seu domini \(\mathbf1\). Per la unicitat del pullback llevat d'isomorfisme, el pullback inicial de \(0\) i \(1\) seria isomorf a \(\mathbf1\), és a dir \(\mathbf1\cong\mathbf0\), en contradicció amb la no-degeneració (axioma~\ref{ax:no-degeneracio}). Per tant \(0\neq1\), i amb el corol·lari~\ref{cor:elements-de-2} conclourem que \(\mathbf2\) té exactament dos punts.
\end{solucio}

\section{Parts extremals: la buida i la total}\label{prac:extremals}

\begin{exercici}
\textbf{La part total té per característica la constant \(1\).} Sigui \(\id_A:A\rightarrowtail A\) la part total de \(A\) (definició~\ref{def:parts-extremals}). Demostreu que la seva aplicació característica és \(1\circ{!}_A:A\to\mathbf2\), la constant de valor \(1\).
\end{exercici}
\begin{solucio}
Comprovem que la constant \(1\circ{!}_A\) satisfà la condició característica de \(\id_A\) (notació~\ref{not:caracteristica}), i la unicitat de l'axioma de representació farà la resta. Sigui \(x:X\to A\) qualsevol. D'una banda, \(x\in_X\id_A\) sempre, amb testimoni el mateix \(x\), perquè \(\id_A\circ x=x\). De l'altra, \((1\circ{!}_A)\circ x=1\circ({!}_A\circ x)=1\circ{!}_X\), perquè \({!}_A\circ x\) és una aplicació \(X\to\mathbf1\) i només n'hi ha una. Per tant, per a tot \(x\), \(x\in_X\id_A\iff(1\circ{!}_A)\circ x=1\circ{!}_X\): totes dues condicions són sempre certes. Això és exactament dir que \(1\circ{!}_A\) és la característica de \(\id_A\). Anàlogament, la part buida \(\bot_A:\mathbf0\rightarrowtail A\) té per característica la constant \(0\circ{!}_A\) (proposició~\ref{prop:caracteristica-buida}), però aquesta direcció és més delicada: cal l'estrictesa de \(\mathbf0\) per veure que \(x\in_X\bot_A\) només passa quan \(X\) és inicial.
\end{solucio}

\begin{exercici}
\textbf{Unicitat de les parts extremals.} Demostreu que dues parts totals de \(A\) són equivalents, i que dues parts buides de \(A\) són equivalents. (Enunciat: proposició~\ref{prop:parts-extremals-uniques}.)
\end{exercici}
\begin{solucio}
Una part total és, per definició, una part \(m:S\rightarrowtail A\) que és un isomorfisme. Aleshores tot element generalitzat \(x:X\to A\) hi pertany, amb testimoni \(m^{-1}\circ x\), de manera que el mateix càlcul de l'exercici anterior dona \(\chi_m=1\circ{!}_A\). Si \(m\) i \(m'\) són totes dues totals, tenen la mateixa característica \(1\circ{!}_A\); i com que la característica determina la part llevat d'equivalència (via l'axioma de representació), \(m\equiv_A m'\). El mateix argument, amb la constant \(0\) i la proposició~\ref{prop:caracteristica-buida}, dona la unicitat de la part buida. Així, \enquote{la} part total i \enquote{la} part buida estan ben definides com a classes d'equivalència.
\end{solucio}

\section{Negació i part complementària}\label{prac:complement}

\begin{exercici}
\textbf{La negació és una involució.} Demostreu que \(\neg\circ\neg=\id_{\mathbf2}\), on \(\neg:\mathbf2\to\mathbf2\) intercanvia \(0\) i \(1\) (definició~\ref{def:negacio}). (Reproduïu la proposició~\ref{prop:negacio-involutiva}.)
\end{exercici}
\begin{solucio}
La negació es defineix per casos sobre la suma \(\mathbf1\xrightarrow{0}\mathbf2\xleftarrow{1}\mathbf1\) mitjançant \(\neg\circ0=1\) i \(\neg\circ1=0\). Comparem \(\neg\circ\neg\) amb \(\id_{\mathbf2}\) component amb les dues injeccions:
\[
(\neg\circ\neg)\circ0=\neg\circ(\neg\circ0)=\neg\circ1=0=\id_{\mathbf2}\circ0,
\qquad
(\neg\circ\neg)\circ1=\neg\circ0=1=\id_{\mathbf2}\circ1.
\]
Les dues aplicacions coincideixen sobre totes dues injeccions; per l'extensionalitat de la definició per casos (proposició~\ref{prop:extensionalitat-suma} del capítol anterior), \(\neg\circ\neg=\id_{\mathbf2}\). En particular \(\neg\) és un isomorfisme, invers de si mateix.
\end{solucio}

\begin{exercici}
\textbf{El complement és involutiu.} Sigui \(m:S\rightarrowtail A\) una part, amb complement \(m^c\) definit per \(\chi_{m^c}=\neg\circ\chi_m\) (definició~\ref{def:complementaria}). Demostreu que \((m^c)^c\equiv_A m\). (Model de la proposició~\ref{prop:complementari-involutiu}.)
\end{exercici}
\begin{solucio}
Apliquem dues vegades la definició del complement en termes de característiques i la involució de la negació:
\[
\chi_{(m^c)^c}=\neg\circ\chi_{m^c}=\neg\circ(\neg\circ\chi_m)=(\neg\circ\neg)\circ\chi_m=\id_{\mathbf2}\circ\chi_m=\chi_m.
\]
Com que dues parts amb la mateixa característica són equivalents, \((m^c)^c\equiv_A m\). Aquest càlcul redueix una propietat geomètrica de parts (complementar dues vegades torna a l'original) a una propietat algebraica d'una sola aplicació cap a \(\mathbf2\).
\end{solucio}

\begin{exercici}
\textbf{Una part i el seu complement descomponen \(A\).} Sigui \(m:S\rightarrowtail A\) amb complement \(m^c:S^c\rightarrowtail A\). Expliqueu per què \([m,m^c]:S+S^c\to A\) és un isomorfisme. (Enunciat: teorema~\ref{teo:part-complementari}.)
\end{exercici}
\begin{solucio}
El quadrat característic (teorema~\ref{teo:quadrat-caracteristic}) identifica \(m\) amb el pullback de \(1\) al llarg de \(\chi_m\), i per definició \(m^c\) és el pullback de \(0\) al llarg de la mateixa \(\chi_m\). Ara bé, \(\mathbf1\xrightarrow{1}\mathbf2\xleftarrow{0}\mathbf1\) és una suma (és la suma canònica amb els sumands intercanviats). Per l'estabilitat de les sumes (axioma~\ref{ax:sumes-disjuntes}), els seus pullbacks al llarg de \(\chi_m\) formen també una suma
\[
S\xrightarrow{\;m\;}A\xleftarrow{\;m^c\;}S^c,
\]
que a més és disjunta. La unicitat de les sumes llevat d'isomorfisme dona que l'aplicació per casos \([m,m^c]:S+S^c\to A\) és un isomorfisme. En termes clàssics: cada element de \(A\) cau exactament dins de \(S\) o dins del seu complement.
\end{solucio}

\section{Operacions booleanes amb parts}\label{prac:booleanes}

\begin{exercici}
\textbf{Característica d'una intersecció.} Siguin \(m,n\) parts de \(A\) amb característiques \(\chi_m,\chi_n:A\to\mathbf2\). Expliqueu per què la característica de la intersecció \(m\cap n\) és \(\wedge\circ\langle\chi_m,\chi_n\rangle\), on \(\wedge:\mathbf2\times\mathbf2\to\mathbf2\) és la conjunció. (Model de la proposició~\ref{prop:caracteristica-interseccio}.)
\end{exercici}
\begin{solucio}
Posem \(c:=\langle\chi_m,\chi_n\rangle:A\to\mathbf2\times\mathbf2\). La conjunció és, per definició (definició~\ref{def:conjuncio-disjuncio}), la característica de la part \(\delta_{11}=\langle1,1\rangle:\mathbf1\rightarrowtail\mathbf2\times\mathbf2\); per tant \(\wedge\circ c\) és la característica de la imatge inversa \(c^{-1}(\delta_{11})\) (proposició~\ref{prop:caracteristica-imatge-inversa}). Ara comparem elements generalitzats, no només punts: per a \(x:X\to A\), tenim \(x\in_X c^{-1}(\delta_{11})\) si i només si \(c\circ x\) es factoritza per \(\delta_{11}\), és a dir, si i només si \(c\circ x=\langle1,1\rangle\circ{!}_X\); component amb les dues projeccions, això equival a \(\chi_m\circ x=1\circ{!}_X\) i \(\chi_n\circ x=1\circ{!}_X\), és a dir, a \(x\in_X m\) i \(x\in_X n\), és a dir, a \(x\in_X(m\cap n)\) (proposició~\ref{prop:pertinenca-interseccio}). Dues parts amb els mateixos elements generalitzats són equivalents (proposició~\ref{prop:inclusio-pertinenca}) i, per tant, tenen la mateixa característica: \(\chi_{m\cap n}=\wedge\circ c\). Observeu que l'argument no ha usat la taula de veritat de \(\wedge\), sinó la seva definició com a característica; la taula serveix per als punts globals, però la igualtat de característiques demana tots els estadis.
\end{solucio}

\begin{exercici}
\textbf{Una llei booleana per càlcul de característiques.} Useu la traducció de les operacions de parts a operacions sobre característiques per demostrar la llei d'idempotència \(m\cap m\equiv_A m\). (Cas particular de les lleis booleanes, teorema~\ref{teo:lleis-booleanes}.)
\end{exercici}
\begin{solucio}
Per la proposició~\ref{prop:caracteristica-interseccio}, \(\chi_{m\cap m}=\wedge\circ\langle\chi_m,\chi_m\rangle\). La conjunció satisfà \(\wedge\circ\langle t,t\rangle=t\) per a tot valor de veritat \(t\) (la taula de \(\wedge\) val \(1\) a \((1,1)\) i \(0\) a \((0,0)\)); això es comprova sobre els dos punts \(0,1\) de \(\mathbf2\) i s'estén per extensionalitat de les taules (proposició~\ref{prop:extensionalitat-taules}). Per tant \(\chi_{m\cap m}=\chi_m\), i com que la característica determina la part, \(m\cap m\equiv_A m\). El mateix mètode --- reduir una identitat de parts a una identitat de taules de veritat sobre \(\mathbf2\)--- demostra la resta de lleis booleanes.
\end{solucio}

\begin{exercici}
\textbf{Inclusió com a conjunció.} Demostreu que \(m\subseteq_A n\) si i només si \(m\cap n\equiv_A m\). (Model de la proposició~\ref{prop:inclusio-booleana}.)
\end{exercici}
\begin{solucio}
Suposem primer \(m\subseteq_A n\). Aleshores tot element generalitzat que pertany a \(m\) pertany a \(n\) (proposició~\ref{prop:inclusio-pertinenca}); per tant les condicions \enquote{pertany a \(m\)} i \enquote{pertany a \(m\) i a \(n\)} defineixen la mateixa part, és a dir \(m\cap n\equiv_A m\). Recíprocament, si \(m\cap n\equiv_A m\), aleshores tot \(x\) que pertany a \(m\) pertany a \(m\cap n\) i, en particular, a \(n\); això és exactament \(m\subseteq_A n\) per la mateixa proposició~\ref{prop:inclusio-pertinenca}. La relació d'inclusió queda així expressada dins de l'àlgebra booleana de les parts (observació~\ref{obs:algebra-booleana-parts}).
\end{solucio}

\section{Límits, coequalitzadors, pushouts i colímits finits}\label{prac:limits-colimits}

\begin{exercici}
\textbf{El coequalitzador és un epimorfisme.} Sigui \(q:B\to Q\) el coequalitzador de \(f,g:A\to B\) (definició~\ref{def:coequalitzador}). Demostreu que \(q\) és un epimorfisme. (Reproduïu la proposició~\ref{prop:coequalitzador-epi}.)
\end{exercici}
\begin{solucio}
Suposem \(u\circ q=v\circ q\) amb \(u,v:Q\to Y\). Sigui \(w:=u\circ q=v\circ q:B\to Y\). Com que \(q\circ f=q\circ g\) (equació del coequalitzador), tenim \(w\circ f=u\circ q\circ f=u\circ q\circ g=w\circ g\); és a dir, \(w\) \enquote{coequalitza} \(f\) i \(g\). Per la propietat universal del coequalitzador (notació~\ref{not:coequalitzador}), existeix una \emph{única} aplicació \(t:Q\to Y\) amb \(t\circ q=w\). Però tant \(u\) com \(v\) compleixen \(u\circ q=w=v\circ q\); per la unicitat, \(u=t=v\). Per tant \(q\) és un epimorfisme. És el dual estricte del fet que l'equalitzador és un monomorfisme.
\end{solucio}

\begin{exercici}
\textbf{Unicitat del coequalitzador llevat d'isomorfisme.} Suposeu que \(q:B\to Q\) i \(q':B\to Q'\) són tots dos coequalitzadors de \(f,g:A\to B\). Demostreu que \(Q\cong Q'\) de manera compatible amb \(q,q'\). (Enunciat: proposició~\ref{prop:coequalitzador-unic}.)
\end{exercici}
\begin{solucio}
Com que \(q'\) coequalitza \(f,g\) i \(q\) és coequalitzador, hi ha una única \(\varphi:Q\to Q'\) amb \(\varphi\circ q=q'\). Simètricament, hi ha una única \(\psi:Q'\to Q\) amb \(\psi\circ q'=q\). Aleshores \((\psi\circ\varphi)\circ q=\psi\circ q'=q=\id_Q\circ q\); per la unicitat de la factorització a través de \(q\), \(\psi\circ\varphi=\id_Q\). Igualment \(\varphi\circ\psi=\id_{Q'}\). Per tant \(\varphi\) és un isomorfisme amb \(\varphi\circ q=q'\). És el mateix patró d'unicitat de tota construcció universal, ara en la versió \enquote{co}.
\end{solucio}

\begin{exercici}
\textbf{Un pushout com a coequalitzador d'una suma.} Descriviu com es construeix el pushout de \(f:C\to A\) i \(g:C\to B\) a partir de la suma \(A+B\) i un coequalitzador. (Model del teorema~\ref{teo:pushout-construccio}.)
\end{exercici}
\begin{solucio}
Formem la suma \(A+B\) amb injeccions \(i_A,i_B\). Les dues aplicacions \(i_A\circ f\) i \(i_B\circ g\) van de \(C\) a \(A+B\); prenem el seu coequalitzador \(q:A+B\to P\) (existeix per l'axioma~\ref{ax:coequalitzadors}). Les composicions \(q\circ i_A:A\to P\) i \(q\circ i_B:B\to P\) satisfan
\[
(q\circ i_A)\circ f=q\circ(i_A\circ f)=q\circ(i_B\circ g)=(q\circ i_B)\circ g,
\]
és a dir, formen un cocon sobre el diagrama \(A\leftarrow C\rightarrow B\). Aquest cocon és universal: donat qualsevol altre cocon \((u:A\to Y,\;v:B\to Y)\) amb \(u\circ f=v\circ g\), l'aplicació per casos \([u,v]:A+B\to Y\) coequalitza \(i_A\circ f\) i \(i_B\circ g\), i per tant es factoritza de manera única per \(q\). Així \(P\), amb \(q\circ i_A\) i \(q\circ i_B\), és el pushout (definició~\ref{def:pushout}). Aquesta construcció mostra que sumes i coequalitzadors basten per obtenir tots els colímits finits (teorema~\ref{teo:colimits-finits}).
\end{solucio}

\begin{exercici}
\textbf{Els límits finits es redueixen a productes i equalitzadors.} Expliqueu per què, tenint productes finits i equalitzadors, existeix el límit de qualsevol diagrama finit. (Model del teorema~\ref{teo:limits-finits}.)
\end{exercici}
\begin{solucio}
Sigui un diagrama finit amb conjunts \(D_1,\dots,D_n\) als vèrtexs i unes quantes aplicacions entre ells. Formem el producte \(P=D_1\times\cdots\times D_n\) (existeix per la proposició~\ref{prop:productes-finits}), amb projeccions \(p_i\). Un con sobre el diagrama és una família d'aplicacions \(x_i:X\to D_i\) compatibles amb cada fletxa \(a:D_i\to D_j\) del diagrama, és a dir \(a\circ x_i=x_j\). Cada una d'aquestes equacions de compatibilitat és una condició d'igualtat entre dues aplicacions \(P\to D_j\), a saber \(a\circ p_i\) i \(p_j\); imposar-la equival a factoritzar per l'equalitzador d'aquestes dues aplicacions. Prenent successivament l'equalitzador de totes les condicions (o l'equalitzador simultani via el producte de les \(D_j\) implicades), obtenim un monomorfisme \(e:L\rightarrowtail P\) tal que un con sobre el diagrama és exactament una aplicació \(X\to L\). Aquest \(L\), amb les projeccions \(p_i\circ e\), és el límit (definició~\ref{def:limit}). En resum: el producte recull tots els vèrtexs alhora i els equalitzadors imposen les equacions del diagrama.
\end{solucio}

\chapter{Relacions, quocients i elecció}\label{cap:practica-relacions}

Aquest capítol acompanya el capítol homònim de la part teòrica (capítol~\ref{cap:relacions} de la Part~\ref{part:teoria}). Els exercicis resolts en segueixen l'ordre d'exposició --- relacions com a parts de productes, relació inversa i diagonal, relacions d'equivalència, parell nucli, quocients i epimorfismes regulars, efectivitat, particions, i les conseqüències de l'axioma d'elecció--- i serveixen de model per als exercicis proposats de la Part~\ref{part:exercicis}. Cada bloc fixa una tècnica: traduir una relació al llenguatge de la pertinença generalitzada, construir aplicacions estructurals mitjançant la propietat universal del pullback, i raonar amb quocients i seccions. Convé intentar cada enunciat abans de llegir-ne la solució.

\section{Relacions, inversa i diagonal}\label{prac:relacions}

\begin{exercici}
\textbf{La inversa de la inversa.} Sigui \(r\colon R\rightarrowtail A\times B\) una relació. Demostreu que \(\operatorname{Inv}(\operatorname{Inv}(r))\equiv_{A\times B}r\) i que, per a tots \(x\colon X\to A\), \(y\colon X\to B\), es té \(y\mathrel{\operatorname{Inv}(r)_X}x\) si i només si \(x\mathrel{r_X}y\). (Reproduïu la proposició~\ref{prop:inversa-pertinenca}.)
\end{exercici}
\begin{solucio}
La inversa es defineix component amb l'isomorfisme de commutació \(\sigma\colon A\times B\to B\times A\), que compleix \(\sigma\circ\langle x,y\rangle=\langle y,x\rangle\). Per tant, \(\langle y,x\rangle\) es factoritza a través de \(\operatorname{Inv}(r)=\sigma\circ r\) si i només si \(\langle x,y\rangle\) es factoritza a través de \(r\); això és exactament l'equivalència \(y\mathrel{\operatorname{Inv}(r)_X}x\iff x\mathrel{r_X}y\) (notació~\ref{not:relacio}). Com que \(\sigma\circ\sigma=\id\), aplicar la inversa dues vegades retorna la relació de partida, d'on \(\operatorname{Inv}(\operatorname{Inv}(r))\equiv_{A\times B}r\).
\end{solucio}

\begin{exercici}
\textbf{La diagonal detecta la igualtat.} Sigui \(\Delta_A=\langle\id_A,\id_A\rangle\colon A\rightarrowtail A\times A\) la diagonal (definició~\ref{def:diagonal}). Demostreu que és un monomorfisme i que, per a tots \(x,x'\colon X\to A\), es té \(x\mathrel{(\Delta_A)_X}x'\) si i només si \(x=x'\). (Model de la proposició~\ref{prop:diagonal-mono}.)
\end{exercici}
\begin{solucio}
Component \(\Delta_A\) amb la primera projecció s'obté \(p_1\circ\langle\id_A,\id_A\rangle=\id_A\), que és un monomorfisme; per la proposició que diu que si \(g\circ f\) és mono aleshores \(f\) ho és (capítol~\ref{cap:parts}), \(\Delta_A\) és un monomorfisme. Ara, \(\langle x,x'\rangle\) es factoritza a través de \(\Delta_A\) si i només si existeix \(w\colon X\to A\) amb \(\langle x,x'\rangle=\langle\id_A,\id_A\rangle\circ w=\langle w,w\rangle\); igualant components, això equival a \(x=w=x'\), és a dir a \(x=x'\). La diagonal és, doncs, la relació d'igualtat en el sentit generalitzat.
\end{solucio}

\begin{exercici}
\textbf{Simetria puntual i simetria estructural.} Sigui \(r\colon R\rightarrowtail A\times A\) una relació sobre \(A\). Demostreu que \(\operatorname{Inv}(r)\equiv_{A\times A}r\) si i només si, per a tot estadi \(X\) i tots \(x,y\colon X\to A\), \(x\mathrel{r_X}y\) implica \(y\mathrel{r_X}x\). Expliqueu per què, amb la separació global, n'hi ha prou amb comprovar-ho per a punts globals.
\end{exercici}
\begin{solucio}
Per l'exercici anterior sobre la inversa, \(y\mathrel{\operatorname{Inv}(r)_X}x\iff x\mathrel{r_X}y\). Si \(\operatorname{Inv}(r)\equiv r\), les dues parts tenen els mateixos elements generalitzats, de manera que \(x\mathrel{r_X}y\) dona \(\langle y,x\rangle\in_X\operatorname{Inv}(r)\), és a dir \(\langle y,x\rangle\in_X r\), és a dir \(y\mathrel{r_X}x\). Recíprocament, si la implicació val a tots els estadis, tot element generalitzat de \(\operatorname{Inv}(r)\) ho és de \(r\) i viceversa, i dues parts amb els mateixos elements generalitzats són equivalents (proposició~\ref{prop:inclusio-pertinenca}). Quant als punts globals: dues parts \(m,n\) d'un mateix conjunt amb els mateixos punts globals són equivalents, perquè la part \(m\cap n^c\) no té cap punt i, pel lema del buit (lema~\ref{lema:buit}), és buida; com que \(m\) es descompon com la suma \((m\cap n)+(m\cap n^c)\) (pullback de la descomposició \(A\cong N+N^c\) del teorema~\ref{teo:part-complementari} al llarg de \(m\), estable per canvi de base), resulta \(m\equiv m\cap n\), és a dir, \(m\subseteq n\); i simètricament. Per tant la simetria puntual implica la simetria estructural; aquesta reducció usa la separació global, i és el motiu pel qual, en la lectura clàssica, les tres condicions sobre punts basten.
\end{solucio}

\section{El parell nucli}\label{prac:parell-nucli}

\begin{exercici}
\textbf{El parell nucli és un monomorfisme.} Sigui \(f\colon A\to B\) i \(\operatorname{Nuc}(f)=\langle k_1,k_2\rangle\colon A\times_B A\to A\times A\) el seu parell nucli (definició~\ref{def:kernel-pair}). Demostreu que \(\langle k_1,k_2\rangle\) és un monomorfisme, sense usar cap propietat de \(f\). (Reproduïu la proposició~\ref{prop:kernel-mono}.)
\end{exercici}
\begin{solucio}
Siguin \(u,v\colon X\to A\times_B A\) amb \(\langle k_1,k_2\rangle\circ u=\langle k_1,k_2\rangle\circ v\). Component amb les dues projeccions \(p_1,p_2\) del producte \(A\times A\), obtenim \(k_1\circ u=k_1\circ v\) i \(k_2\circ u=k_2\circ v\). Ara bé, per la unicitat a la propietat universal del pullback, dues aplicacions amb codomini \(A\times_B A\) que coincideixen en compondre amb \(k_1\) i amb \(k_2\) són iguals; per tant \(u=v\). Cap pas ha usat que \(f\) sigui mono: el que fa mono el parell nucli és exclusivament la unicitat del pullback (observació~\ref{obs:kernel-no-mono}).
\end{solucio}

\begin{exercici}
\textbf{Reflexivitat i simetria del parell nucli.} Amb la notació \(K=A\times_B A\), \(k_1,k_2\colon K\to A\), \(f\circ k_1=f\circ k_2\), construïu les aplicacions estructurals que testimonien que \(\operatorname{Nuc}(f)\) és reflexiu i simètric. (Part de la proposició~\ref{prop:kernel-equivalencia}.)
\end{exercici}
\begin{solucio}
\emph{Reflexivitat.} L'aplicació \(\id_A\) compleix trivialment \(f\circ\id_A=f\circ\id_A\); el parell \((\id_A,\id_A)\) és, doncs, un con sobre el pullback, i la propietat universal proporciona una única \(e\colon A\to K\) amb \(k_1\circ e=\id_A\) i \(k_2\circ e=\id_A\). Aquesta \(e\) realitza la reflexivitat (\enquote{tot element està relacionat amb ell mateix}).

\emph{Simetria.} Les projeccions intercanviades compleixen \(f\circ k_2=f\circ k_1\); per tant \((k_2,k_1)\) és un con sobre el pullback, i la propietat universal dona una única \(s\colon K\to K\) amb \(k_1\circ s=k_2\) i \(k_2\circ s=k_1\). Aquesta \(s\) realitza la simetria. (La transitivitat es construeix anàlogament sobre \(K\times_A K\).)
\end{solucio}

\section{Quocients i epimorfismes regulars}\label{prac:quocients}

\begin{exercici}
\textbf{Tot quocient és un epimorfisme regular.} Sigui \(r\) una relació d'equivalència sobre \(A\) i \(q\colon A\twoheadrightarrow A/r\) un quocient. Expliqueu per què \(q\) és un epimorfisme regular. (Model de la proposició~\ref{prop:quocient-regular}.)
\end{exercici}
\begin{solucio}
Per definició (definició~\ref{def:quocient}), el quocient \(q\) és el coequalitzador de les dues projeccions \(r_1,r_2\colon R\to A\) de la relació. Un epimorfisme regular és, precisament, una aplicació que és el coequalitzador d'algun parell d'aplicacions paral·leles (definició~\ref{def:epi-regular}); com que \(q\) ho és de \(r_1,r_2\), és un epimorfisme regular. En particular és un epimorfisme, per la proposició~\ref{prop:coequalitzador-epi} del capítol~\ref{cap:valors}. El recíproc ---que tot epimorfisme sigui regular--- no el tenim encara: el donarà l'axioma d'elecció (observació~\ref{obs:regular-escindit}).
\end{solucio}

\begin{exercici}
\textbf{Existència i unicitat del quocient.} Justifiqueu que tota relació sobre \(A\) té quocient, únic llevat d'un únic isomorfisme compatible, i que l'aplicació quocient és un epimorfisme. (Enunciat: proposició~\ref{prop:quocient-existeix}.)
\end{exercici}
\begin{solucio}
L'existència és l'axioma de coequalitzadors (axioma~\ref{ax:coequalitzadors} del capítol~\ref{cap:valors}) aplicat a les projeccions \(r_1,r_2\): el seu coequalitzador és el quocient. La unicitat llevat d'un únic isomorfisme compatible és la unicitat dels coequalitzadors (proposició~\ref{prop:coequalitzador-unic}), que ja vam provar amb el patró estàndard de les construccions universals. I que l'aplicació quocient sigui un epimorfisme és la proposició~\ref{prop:coequalitzador-epi}. Els tres ingredients provenen, doncs, de la teoria dels coequalitzadors del capítol anterior.
\end{solucio}

\section{Efectivitat i correspondència}\label{prac:efectivitat}

\begin{exercici}
\textbf{Invariància del parell nucli per isomorfisme del codomini.} Siguin \(q\colon A\to B\) i \(w\colon B\to B'\) un isomorfisme, i \(q'=w\circ q\). Demostreu que tot parell nucli de \(q\) és equivalent, com a part de \(A\times A\), a tot parell nucli de \(q'\). (Reproduïu el lema~\ref{lema:nucli-invariant}.)
\end{exercici}
\begin{solucio}
Com que \(w\) és un isomorfisme i \(q'=w\circ q\), per a qualsevol parell \(u,v\) amb codomini \(A\) es té \(q\circ u=q\circ v\iff q'\circ u=q'\circ v\) (component amb \(w\) o amb \(w^{-1}\), que es cancel·la). Per tant, un con sobre el diagrama de pullback de \(q\) amb ell mateix és exactament un con sobre el de \(q'\); les dues propietats universals determinen el mateix objecte llevat d'un únic isomorfisme, de manera que existeix \(\theta\colon K\to K'\) isomorfisme amb \(k_1'\circ\theta=k_1\) i \(k_2'\circ\theta=k_2\). Component amb l'aparellament, \(\langle k_1',k_2'\rangle\circ\theta=\langle k_1,k_2\rangle\); les dues parts es factoritzen l'una per l'altra a través d'un isomorfisme, i per tant són equivalents. Aquest lema és el que fa que el camí \enquote{relació \(\to\) quocient \(\to\) parell nucli} no depengui del quocient triat.
\end{solucio}

\begin{exercici}
\textbf{El parell nucli recupera la relació.} Expliqueu, usant l'axioma d'efectivitat, per què per a tota relació d'equivalència \(r\) i tot quocient \(q\) seu es té \(\operatorname{Nuc}(q)\equiv_{A\times A}r\). (Camí (1)$\to$(2) del teorema~\ref{teo:correspondencia}.)
\end{exercici}
\begin{solucio}
L'axioma d'efectivitat (axioma~\ref{ax:exactitud}) garanteix que existeix \emph{algun} quocient \(q_0\) de \(r\) amb un parell nucli equivalent a \(r\). Ara bé, dos quocients qualssevol de \(r\) són isomorfs de manera compatible (unicitat dels coequalitzadors, proposició~\ref{prop:coequalitzador-unic}): hi ha un isomorfisme \(w\) amb \(q=w\circ q_0\). Pel lema~\ref{lema:nucli-invariant}, el parell nucli és invariant sota postcomposició amb un isomorfisme, de manera que el parell nucli de \(q\) també és equivalent a \(r\). Aquest és el sentit precís en què \enquote{prendre quocient} i \enquote{prendre parell nucli} són operacions inverses.
\end{solucio}

\section{Particions i epimorfismes escindits}\label{prac:particions}

\begin{exercici}
\textbf{Tot epimorfisme escindit és regular.} Sigui \(q\colon A\to B\) un epimorfisme amb secció \(s\) (\(q\circ s=\id_B\)). Demostreu que \(q\) és el coequalitzador del seu parell nucli. (Reproduïu la proposició~\ref{prop:escindit-regular}.)
\end{exercici}
\begin{solucio}
Sigui \(\operatorname{Nuc}(q)\) amb projeccions \(k_1,k_2\) i \(q\circ k_1=q\circ k_2\). Cal veure que \(q\) té la propietat universal del coequalitzador de \(k_1,k_2\). Ja tenim \(q\circ k_1=q\circ k_2\). Sigui \(h\colon A\to C\) amb \(h\circ k_1=h\circ k_2\); definim \(\bar h:=h\circ s\). Considerem \(\id_A\) i \(s\circ q\): compleixen \(q\circ\id_A=q=q\circ s\circ q\) (perquè \(q\circ s=\id_B\)), de manera que el parell \((\id_A,\,s\circ q)\) es factoritza pel pullback \(\operatorname{Nuc}(q)\): hi ha \(d\colon A\to K\) amb \(k_1\circ d=\id_A\) i \(k_2\circ d=s\circ q\). Aleshores
\[
h=h\circ\id_A=h\circ k_1\circ d=h\circ k_2\circ d=h\circ s\circ q=\bar h\circ q,
\]
usant \(h\circ k_1=h\circ k_2\) al pas central. La unicitat de \(\bar h\) surt de ser \(q\) un epimorfisme. Per tant \(q\) és el coequalitzador de \(k_1,k_2\), i és regular.
\end{solucio}

\section{Conseqüències de l'axioma d'elecció}\label{prac:eleccio}

\begin{exercici}
\textbf{Amb elecció, tot epimorfisme és regular.} Demostreu que, admès l'axioma d'elecció, tot epimorfisme és un epimorfisme regular. (Reproduïu la proposició~\ref{prop:epi-regular-eleccio}.)
\end{exercici}
\begin{solucio}
L'axioma d'elecció (axioma~\ref{ax:eleccio}) afirma que tot epimorfisme és escindit, és a dir, té secció. I la proposició~\ref{prop:escindit-regular} diu que tot epimorfisme escindit és regular. Encadenant-les: tot epimorfisme té secció, per tant és escindit, per tant és regular. Aquesta és la peça que, un cop afegida l'elecció, fa que la correspondència del teorema~\ref{teo:correspondencia} s'estengui de \emph{tots} els epimorfismes amb domini \(A\), no només els regulars (corol·lari~\ref{cor:correspondencia-tots-epis}).
\end{solucio}

\begin{exercici}
\textbf{Epimorfisme i exhaustivitat generalitzada.} Demostreu que, amb l'axioma d'elecció, tot epimorfisme \(q\colon A\twoheadrightarrow B\) és exhaustiu sobre elements generalitzats: per a tot \(y\colon X\to B\) existeix \(x\colon X\to A\) amb \(q\circ x=y\). (Model de la proposició~\ref{prop:epi-exhaustiu-generalitzat}.)
\end{exercici}
\begin{solucio}
Per l'elecció, \(q\) té una secció \(s\colon B\to A\) amb \(q\circ s=\id_B\). Donat qualsevol \(y\colon X\to B\), preneu \(x:=s\circ y\). Aleshores
\[
q\circ x=q\circ(s\circ y)=(q\circ s)\circ y=\id_B\circ y=y.
\]
Per tant \(y\) factoritza per \(q\), i \(q\) és exhaustiu sobre els elements generalitzats (és el criteri de la proposició~\ref{prop:exhaustiva-generalitzada-escindit} del capítol~\ref{cap:parts}, aplicat a la secció que dona l'elecció). Prenent l'estadi \(X=\mathbf1\) se'n dedueix l'exhaustivitat sobre els punts globals; però aquesta conseqüència particular no necessita l'elecció, perquè ja es va demostrar al capítol~\ref{cap:valors} (proposició~\ref{prop:epi-exhaustiu-punts}; vegeu la proposició~\ref{prop:epi-exhaustiu-global}). El que l'elecció aporta de nou és la preimatge per a tots els estadis alhora.
\end{solucio}

\begin{exercici}
\textbf{Sistema de representants.} Demostreu que, amb l'axioma d'elecció, tot quocient \(q\colon A\twoheadrightarrow A/r\) admet una secció, i interpreteu-ho clàssicament. (Enunciat: proposició~\ref{prop:representants}.)
\end{exercici}
\begin{solucio}
El quocient \(q\) és un epimorfisme (proposició~\ref{prop:quocient-existeix}); per l'axioma d'elecció (axioma~\ref{ax:eleccio}), tot epimorfisme té secció, de manera que existeix \(s\colon A/r\to A\) amb \(q\circ s=\id_{A/r}\). En la lectura clàssica (observació~\ref{obs:seccio-representants}), aquesta secció tria per a cada classe un element que la representa: \(s([a])\) és un punt de la classe \([a]\). Així, l'elecció és exactament el que garanteix que tota partició admet un sistema de representants.
\end{solucio}

\chapter{Imatges i factoritzacions}\label{cap:practica-imatges}

Aquest capítol acompanya el capítol homònim de la part teòrica (capítol~\ref{cap:imatges} de la Part~\ref{part:teoria}). Els exercicis resolts en segueixen l'ordre d'exposició --- la factorització (epi regular)--(mono), les caracteritzacions per la factorització, la imatge directa i la seva adjunció amb la imatge inversa, i la composició de relacions amb les relacions funcionals i totals--- i serveixen de model per als exercicis proposats de la Part~\ref{part:exercicis}. Cada bloc fixa una tècnica: construir la factorització imatge via el parell nucli i el quocient, llegir monomorfismes i epimorfismes a través dels seus factors, i raonar sobre elements generalitzats amb testimonis. Convé intentar cada enunciat abans de llegir-ne la solució.

\section{La factorització imatge}\label{prac:factoritzacio}

\begin{exercici}
\textbf{Construcció de la factorització.} Sigui \(f\colon A\to B\). Descriviu com es construeixen l'epimorfisme regular \(e\) i el monomorfisme \(m\) de la factorització \(A\overset{e}{\twoheadrightarrow}\operatorname{Im}(f)\overset{m}{\hookrightarrow}B\), i justifiqueu l'existència de l'única \(m\) amb \(f=m\circ e\). (Model del teorema~\ref{teo:factoritzacio}.)
\end{exercici}
\begin{solucio}
Prenem el parell nucli \(n=\operatorname{Nuc}(f)\), que és una relació d'equivalència sobre \(A\) (proposició~\ref{prop:kernel-equivalencia}), i el seu quocient \(e\colon A\twoheadrightarrow A/n\), un epimorfisme regular (proposicions~\ref{prop:quocient-existeix} i~\ref{prop:quocient-regular}). Per definició de parell nucli, \(f\circ k_1=f\circ k_2\), on \(k_1,k_2\) són les projeccions de \(n\); és a dir, \(f\) coiguala el parell \((k_1,k_2)\). Com que \(e\) és el coequalitzador d'aquest parell (definició~\ref{def:coequalitzador}), la seva propietat universal proporciona una \emph{única} aplicació \(m\colon A/n\to B\) amb \(f=m\circ e\). Posant \(\operatorname{Im}(f):=A/n\), obtenim la factorització desitjada; falta veure que \(m\) és mono (exercici següent).
\end{solucio}

\begin{exercici}
\textbf{El factor \(m\) és un monomorfisme, sense elecció.} Demostreu que l'aplicació \(m\colon\operatorname{Im}(f)\to B\) de la factorització és un monomorfisme, usant només la separació global i l'exhaustivitat dels epimorfismes sobre punts. Doneu després l'argument alternatiu amb una secció de \(e\), i compareu-los. (Reproduïu el pas central del teorema~\ref{teo:factoritzacio}; vegeu l'observació~\ref{obs:factoritzacio-eleccio}.)
\end{exercici}
\begin{solucio}
\emph{Sense elecció.} Per la separació global, un monomorfisme és el mateix que una aplicació injectiva sobre punts (observació~\ref{obs:injectiva-mono}). Siguin \(y,y'\colon\mathbf1\to I\) amb \(m\circ y=m\circ y'\). Com que \(e\colon A\twoheadrightarrow I\) és un epimorfisme, és exhaustiu sobre punts (proposició~\ref{prop:epi-exhaustiu-punts}): \(y=e\circ a\) i \(y'=e\circ a'\) per a certs \(a,a'\colon\mathbf1\to A\). Aleshores \(f\circ a=m\circ e\circ a=m\circ y=m\circ y'=f\circ a'\), de manera que \((a,a')\) es factoritza pel parell nucli \(n\) de \(f\); i com que \(e\) coiguala les projeccions de \(n\), \(y=e\circ a=e\circ a'=y'\). Per tant \(m\) és injectiva sobre punts, i mono.

\emph{Amb elecció.} Una secció \(s\colon I\to A\) de \(e\) permet fer el mateix a qualsevol estadi: si \(m\circ x=m\circ x'\) amb \(x,x'\colon X\to I\), aleshores \(f\circ s\circ x=m\circ x=m\circ x'=f\circ s\circ x'\), de manera que \(s\circ x\) i \(s\circ x'\) estan relacionats per \(n\), i \(x=e\circ s\circ x=e\circ s\circ x'=x'\).

\emph{Comparació.} El primer argument redueix la qüestió als punts globals, que és on la separació global dona força a la teoria, i després usa que els epimorfismes tenen prou preimatges puntuals. El segon no necessita la separació global, però necessita una secció, és a dir, l'elecció. En una teoria ben apuntada el primer és el bo: no gasta un axioma que no cal.
\end{solucio}

\begin{exercici}
\textbf{Unicitat de la factorització.} Suposeu que \(f=m_1\circ e_1=m_2\circ e_2\) són dues factoritzacions (epi regular)--(mono). Demostreu que hi ha un únic isomorfisme \(u\) amb \(m_2\circ u=m_1\) i \(u\circ e_1=e_2\). (Part de unicitat del teorema~\ref{teo:factoritzacio}.)
\end{exercici}
\begin{solucio}
Per la minimalitat de la imatge, cada monomorfisme es factoritza per l'altre: hi ha \(u\colon I_1\to I_2\) i \(v\colon I_2\to I_1\) amb \(m_2\circ u=m_1\) i \(m_1\circ v=m_2\). Aleshores \(m_1\circ v\circ u=m_2\circ u=m_1\); com que \(m_1\) és mono, \(v\circ u=\id_{I_1}\), i simètricament \(u\circ v=\id_{I_2}\). Per tant \(u\) és un isomorfisme. La compatibilitat \(u\circ e_1=e_2\) surt de \(m_2\circ u\circ e_1=m_1\circ e_1=f=m_2\circ e_2\) cancel·lant el monomorfisme \(m_2\); i la unicitat de \(u\), del mateix argument de cancel·lació. És el patró habitual d'unicitat de les construccions universals.
\end{solucio}

\section{Caracteritzacions per la factorització}\label{prac:caracteritzacions}

\begin{exercici}
\textbf{Mono si i només si parell nucli trivial.} Demostreu que \(f\colon A\to B\) és un monomorfisme si i només si \(\operatorname{Nuc}(f)\equiv_{A\times A}\Delta_A\), i que en aquest cas l'aplicació exhaustiva \(e\) és un isomorfisme. (Reproduïu la proposició~\ref{prop:mono-nucli-trivial}.)
\end{exercici}
\begin{solucio}
Si \(f\) és mono, aleshores \(f\circ k_1=f\circ k_2\) implica \(k_1=k_2\), de manera que el parell nucli és la diagonal (proposició~\ref{prop:diagonal-mono}). Recíprocament, si \(\operatorname{Nuc}(f)\equiv_{A\times A}\Delta_A\) i \(x,x'\colon X\to A\) compleixen \(f\circ x=f\circ x'\), aleshores \(\langle x,x'\rangle\) es factoritza pel parell nucli, per tant per la diagonal, i això dona \(x=x'\); així \(f\) és mono. En aquest cas el quocient de \(A\) per la diagonal és el mateix \(A\), de manera que \(e\colon A\to\operatorname{Im}(f)\) és un isomorfisme: un monomorfisme \enquote{ja és} la seva pròpia imatge.
\end{solucio}

\begin{exercici}
\textbf{Epi si i només si imatge total.} Demostreu que \(f\) és un epimorfisme si i només si el monomorfisme \(m\colon\operatorname{Im}(f)\to B\) és un isomorfisme. (Model de la proposició~\ref{prop:epi-imatge-total}.)
\end{exercici}
\begin{solucio}
Si \(m\) és un isomorfisme, \(f=m\circ e\) és la composició d'un epimorfisme regular \(e\) amb un isomorfisme, i per tant és un epimorfisme (regular). Recíprocament, sigui \(f\) un epimorfisme. Com que \(f=m\circ e\), també \(m\) és un epimorfisme (si \(u\circ m=v\circ m\), component amb \(e\) s'obté \(u\circ f=v\circ f\) i \(u=v\)). Així \(m\) és un bimorfisme i, pel teorema del bimorfisme (proposició~\ref{prop:bimorfisme-isomorfisme}), un isomorfisme. Observeu que aquest argument dona de passada el corol·lari~\ref{cor:epi-regular-sense-eleccio}: tot epimorfisme és regular, perquè és un coequalitzador (\(e\)) seguit d'un isomorfisme. L'argument antic, que suposava \(f\) regular i construïa la inversa de \(m\) amb la propietat universal del coequalitzador, continua sent vàlid però ja no cal.
\end{solucio}

\begin{exercici}
\textbf{Isomorfisme via els dos factors.} Demostreu que \(f\) és un isomorfisme si i només si és alhora monomorfisme i epimorfisme. (Model del corol·lari~\ref{cor:bimorfisme-iso}.)
\end{exercici}
\begin{solucio}
Si \(f\) és un isomorfisme, és tant mono com epi. Recíprocament, si \(f\) és mono, l'aplicació exhaustiva \(e\) és un isomorfisme (proposició~\ref{prop:mono-nucli-trivial}); si a més és epi, el monomorfisme \(m\) també ho és (proposició~\ref{prop:epi-imatge-total}). Aleshores \(f=m\circ e\) és composició de dos isomorfismes, i és un isomorfisme. Aquesta és la lectura del teorema del bimorfisme a través de la factorització imatge: cada factor recull una de les dues hipòtesis.
\end{solucio}

\section{Imatge directa i adjunció}\label{prac:imatge-directa}

\begin{exercici}
\textbf{L'adjunció imatge directa\,--\,imatge inversa.} Siguin \(f\colon A\to B\), una part \(s\colon S\rightarrowtail A\) i una part \(t\colon T\rightarrowtail B\). Demostreu que \(f_{*}S\subseteq_B T\) si i només si \(S\subseteq_A f^{*}T\). (Reproduïu la proposició~\ref{prop:adjuncio-imatges}.)
\end{exercici}
\begin{solucio}
Suposem \(f_{*}S\subseteq_B T\). Un element generalitzat \(x\colon X\to A\) pertany a \(S\) quan es factoritza per \(s\); en aquest cas \(f\circ x\) es factoritza per \(f\circ s\) i, per la definició d'imatge, per \(f_{*}S\), i per hipòtesi per \(t\). Per tant \(f\circ x\) pertany a \(T\), és a dir \(x\) pertany a \(f^{*}T\) (proposició~\ref{prop:pertinenca-imatge-inversa}); això és \(S\subseteq_A f^{*}T\). Recíprocament, si \(S\subseteq_A f^{*}T\), tot element de \(S\) té imatge dins \(T\); com que \(f_{*}S\) és la part \emph{mínima} per la qual passa \(f\circ s\) i \(t\) n'és una, la minimalitat dona \(f_{*}S\subseteq_B T\). Prendre imatge directa i inversa són, doncs, operacions adjuntes entre les parts de \(A\) i les de \(B\).
\end{solucio}

\begin{exercici}
\textbf{La imatge de la part total.} Demostreu que \(f_{*}A=\operatorname{Im}(f)\), i deduïu que \(f\) és un epimorfisme si i només si \(f_{*}A\equiv_B\id_B\). (Reproduïu el corol·lari~\ref{cor:estabilitat-exhaustiva}.)
\end{exercici}
\begin{solucio}
La part total de \(A\) és \(\id_A\), i \(f\circ\id_A=f\); per definició d'imatge directa, \(f_{*}A\) és la imatge de \(f\circ\id_A=f\), és a dir \(f_{*}A=\operatorname{Im}(f)\). La segona afirmació és, aleshores, exactament la proposició~\ref{prop:epi-imatge-total}: \(f\) és epi si i només si \(m\) és iso, és a dir si i només si la imatge és tot \(B\), això és \(f_{*}A\equiv_B\id_B\).
\end{solucio}

\section{Composició de relacions i relacions funcionals}\label{prac:relacions-funcionals}

\begin{exercici}
\textbf{La composició de relacions és una relació.} Siguin \(r\) una relació de \(A\) cap a \(B\) i \(r'\) una de \(B\) cap a \(C\). Expliqueu per què \(r'\circ r\) és una relació de \(A\) cap a \(C\), i quin paper hi juga la imatge. (Model de la proposició~\ref{prop:composicio-es-relacio}.)
\end{exercici}
\begin{solucio}
Es forma el producte fibrat \(R\times_B R'\) (que encadena un testimoni de \(r\) amb un de \(r'\) sobre un mateix punt de \(B\)) i l'aplicació \(g\colon R\times_B R'\to A\times C\) que en recull els extrems, \(g=\langle r_1\circ\pi,\;r'_2\circ\pi'\rangle\). La composició \(r'\circ r\) es defineix com la \emph{imatge} d'aquesta \(g\): per la factorització (teorema~\ref{teo:factoritzacio}), la imatge és sempre un monomorfisme cap a \(A\times C\), de manera que \(r'\circ r\) és una part de \(A\times C\), és a dir una relació. La imatge és, precisament, el que introdueix el quantificador existencial \enquote{existeix \(b\) amb \(a\mathrel{r}b\) i \(b\mathrel{r'}c\)}: comprimeix tots els encadenaments possibles en el conjunt de parells efectivament connectats (observació~\ref{obs:composicio-relacions-classica}).
\end{solucio}

\begin{exercici}
\textbf{Funcionalitat \(=\) unicitat del testimoni.} Sigui \(r\colon R\rightarrowtail A\times B\) una relació amb projeccions \(r_1,r_2\). Demostreu que \(r\) és funcional (és a dir, \(r_1\) és mono) si i només si per a tot \(x\colon X\to A\) hi ha com a molt un \(y\colon X\to B\) amb \(x\mathrel{r_X}y\). (Part de la proposició~\ref{prop:funcional-lectura}.)
\end{exercici}
\begin{solucio}
Recordem que \(x\mathrel{r_X}y\) vol dir que hi ha un testimoni \(\rho\colon X\to R\) amb \(r_1\circ\rho=x\) i \(r_2\circ\rho=y\). Suposem \(r_1\) mono. Si \(x\mathrel{r_X}y\) i \(x\mathrel{r_X}y'\), amb testimonis \(\rho,\rho'\), aleshores \(r_1\circ\rho=x=r_1\circ\rho'\); com que \(r_1\) és mono, \(\rho=\rho'\), i per tant \(y=r_2\circ\rho=r_2\circ\rho'=y'\). Hi ha, doncs, com a molt un \(y\). Recíprocament, suposem la unicitat i siguin \(\rho,\rho'\) amb \(r_1\circ\rho=r_1\circ\rho'\); posant \(y=r_2\circ\rho\), \(y'=r_2\circ\rho'\), tots dos són testimonis del mateix \(x\), d'on \(y=y'\); llavors \(\rho,\rho'\) tenen les dues components iguals, i com que \(r=\langle r_1,r_2\rangle\) és mono, \(\rho=\rho'\). Això prova que \(r_1\) és mono. La funcionalitat és, literalment, la unicitat del testimoni.
\end{solucio}

\begin{exercici}
\textbf{Les aplicacions són les relacions funcionals i totals.} Demostreu que assignar a \(f\colon A\to B\) el seu graf \(\langle\id_A,f\rangle\) dona una correspondència biunívoca (llevat d'equivalència de parts) entre aplicacions \(A\to B\) i relacions funcionals i totals. (Model de la proposició~\ref{prop:aplicacio-funcional}.)
\end{exercici}
\begin{solucio}
El graf \(\langle\id_A,f\rangle\colon A\rightarrowtail A\times B\) té primera projecció \(p_1\circ\langle\id_A,f\rangle=\id_A\), un isomorfisme; per tant el graf és funcional i total (definició~\ref{def:relacio-funcional}), amb \(y=f\circ x\) l'únic element relacionat amb cada \(x\). Recíprocament, una relació funcional i total té primera projecció \(r_1\) que és alhora mono i epi, i doncs un isomorfisme (corol·lari~\ref{cor:bimorfisme-iso}); l'aplicació \(f:=r_2\circ r_1^{-1}\) té per graf una part equivalent a \(r\). Aquestes dues assignacions són mútuament inverses llevat d'equivalència. A més, sota la correspondència, la composició de relacions coincideix amb la d'aplicacions: el producte fibrat de dos grafs encadena \(x\) amb \(f\circ x\) i aquest amb \(g\circ(f\circ x)\), i la imatge d'aquest encadenament és el graf de \(g\circ f\).
\end{solucio}

\chapter{Exponencials}\label{cap:practica-exponencials}

Aquest capítol acompanya el capítol homònim de la part teòrica (capítol~\ref{cap:exponencials} de la Part~\ref{part:teoria}). Els exercicis resolts en segueixen l'ordre d'exposició --- l'avaluació i la propietat universal, les aplicacions com a punts, els exponencials de \(\mathbf0\), \(\mathbf1\) i \(\mathbf2\), l'exponencial d'una suma, la composició interna i les famílies--- i serveixen de model per als exercicis proposats de la Part~\ref{part:exercicis}. Cada bloc fixa una tècnica: comprovar propietats universals per unicitat de la transposada, calcular exponencials identificant la propietat universal d'un objecte conegut, i traduir entre aplicacions \(X\times A\to B\) i aplicacions \(X\to B^A\). Convé intentar cada enunciat abans de llegir-ne la solució.

\section{Propietat universal i transposició}\label{prac:exp-universal}

\begin{exercici}
\textbf{Unicitat de l'exponencial.} Suposeu que \((E,\operatorname{ev})\) i \((E',\operatorname{ev}')\) són tots dos exponencials de \(B\) per \(A\). Demostreu que hi ha un únic isomorfisme \(w\colon E\to E'\) amb \(\operatorname{ev}'\circ(w\times\id_A)=\operatorname{ev}\). (Reproduïu la proposició~\ref{prop:exponencial-unic}.)
\end{exercici}
\begin{solucio}
Per la propietat universal de \((E',\operatorname{ev}')\) aplicada a \(\operatorname{ev}\colon E\times A\to B\), hi ha una única \(w\colon E\to E'\) amb \(\operatorname{ev}'\circ(w\times\id_A)=\operatorname{ev}\). Simètricament, hi ha una única \(w'\colon E'\to E\) amb \(\operatorname{ev}\circ(w'\times\id_A)=\operatorname{ev}'\). Component, \(\operatorname{ev}\circ((w'\circ w)\times\id_A)=\operatorname{ev}\); però també \(\operatorname{ev}\circ(\id_E\times\id_A)=\operatorname{ev}\), i la unicitat de la transposada de \(\operatorname{ev}\) (definició~\ref{def:exponencial}) dona \(w'\circ w=\id_E\). Simètricament \(w\circ w'=\id_{E'}\). Per tant \(w\) és un isomorfisme, únic compatible amb l'avaluació. És el patró d'unicitat de tota construcció universal.
\end{solucio}

\begin{exercici}
\textbf{Descurrificar i currificar són inverses.} Per a una aplicació \(g\colon X\to B^A\), la seva descurrificació és \(\operatorname{ev}\circ(g\times\id_A)\colon X\times A\to B\). Demostreu que currificar la descurrificació torna a donar \(g\), i que descurrificar la currificació de \(h\colon X\times A\to B\) torna a donar \(h\). (Model de la definició~\ref{def:exponencial} i la notació~\ref{not:exponencial}.)
\end{exercici}
\begin{solucio}
Sigui \(h:=\operatorname{ev}\circ(g\times\id_A)\) la descurrificació de \(g\). La currificació \(\lambda(h)\) és, per definició, l'única aplicació \(X\to B^A\) amb \(\operatorname{ev}\circ(\lambda(h)\times\id_A)=h\). Però \(g\) ja compleix \(\operatorname{ev}\circ(g\times\id_A)=h\); per la \emph{unicitat} de la transposada, \(\lambda(h)=g\). Recíprocament, donada \(h\colon X\times A\to B\), la seva currificació \(\lambda(h)\) satisfà \(\operatorname{ev}\circ(\lambda(h)\times\id_A)=h\) per definició; el membre esquerre és exactament la descurrificació de \(\lambda(h)\), de manera que descurrificar \(\lambda(h)\) retorna \(h\). Les dues operacions són, doncs, mútuament inverses.
\end{solucio}

\section{Aplicacions com a punts}\label{prac:exp-punts}

\begin{exercici}
\textbf{Els punts de \(B^A\) són les aplicacions \(A\to B\).} Demostreu que hi ha una correspondència biunívoca entre els punts globals \(\mathbf1\to B^A\) i les aplicacions \(A\to B\). (Reproduïu el teorema~\ref{teo:punts-exponencial}.)
\end{exercici}
\begin{solucio}
Fixem l'isomorfisme canònic \(\iota\colon A\xrightarrow{\sim}\mathbf1\times A\) (per la propietat universal del producte amb el terminal). A un punt \(p\colon\mathbf1\to B^A\) li associem \(\Phi(p):=\operatorname{ev}\circ(p\times\id_A)\circ\iota\colon A\to B\). A una aplicació \(f\colon A\to B\) li associem \(\Psi(f):=\lambda(f\circ\iota^{-1})\colon\mathbf1\to B^A\). Que \(\Phi\) i \(\Psi\) són inverses és exactament la unicitat de la transposada: \(\Phi(\Psi(f))=f\) perquè \(\operatorname{ev}\circ(\lambda(f\circ\iota^{-1})\times\id_A)=f\circ\iota^{-1}\), i \(\Psi(\Phi(p))=p\) perquè tots dos tenen la mateixa descurrificació i la transposada és única. Així, un element de \(B^A\) \enquote{és} una aplicació \(A\to B\).
\end{solucio}

\begin{exercici}
\textbf{Els elements generalitzats són famílies.} Expliqueu per què un element generalitzat \(X\to B^A\) es correspon amb una aplicació \(X\times A\to B\), i com aquesta lectura generalitza el cas dels punts. (Model de l'observació~\ref{obs:generalitzats-exponencial} i la proposició~\ref{prop:families}.)
\end{exercici}
\begin{solucio}
La propietat universal de l'exponencial (definició~\ref{def:exponencial}) diu que les aplicacions \(X\to B^A\) es corresponen bijectivament, via la transposició, amb les aplicacions \(X\times A\to B\): a cada \(g\colon X\to B^A\) li correspon la seva descurrificació \(\operatorname{ev}\circ(g\times\id_A)\), i recíprocament. El cas dels punts és el cas particular \(X=\mathbf1\), on \(\mathbf1\times A\cong A\) i recuperem el teorema~\ref{teo:punts-exponencial}. En general, una aplicació \(X\times A\to B\) es llegeix com una \emph{família} \((f_x)_x\) d'aplicacions \(A\to B\) indexada pels elements generalitzats de \(X\), i \(g\colon X\to B^A\) n'és el \enquote{nom intern}.
\end{solucio}

\section{Exponencials de \texorpdfstring{\(\mathbf0\), \(\mathbf1\), \(\mathbf2\)}{0, 1, 2}}\label{prac:exp-basics}

\begin{exercici}
\textbf{\(A^{\mathbf1}\cong A\).} Demostreu que l'exponencial de \(A\) per \(\mathbf1\) és isomorf a \(A\). (Part de la proposició~\ref{prop:exponencials-basics}.)
\end{exercici}
\begin{solucio}
Sigui \(\rho\colon A\times\mathbf1\xrightarrow{\sim}A\) l'isomorfisme de projecció. Definim \(\alpha\colon A\to A^{\mathbf1}\) com la transposada de \(\rho\), és a dir l'única amb \(\operatorname{ev}\circ(\alpha\times\id_{\mathbf1})=\rho\). En sentit contrari, avaluant al punt únic de \(\mathbf1\), \(\beta:=\operatorname{ev}\circ\langle\id,{!}\rangle\colon A^{\mathbf1}\to A\), on \({!}\colon A^{\mathbf1}\to\mathbf1\). La composició \(\beta\circ\alpha=\operatorname{ev}\circ(\alpha\times\id_{\mathbf1})\circ\langle\id_A,{!}_A\rangle=\rho\circ\langle\id_A,{!}_A\rangle=\id_A\) es calcula directament. Per a \(\alpha\circ\beta\), que és una aplicació cap a l'exponencial, es comparen descurrificacions: \(\operatorname{ev}\circ((\alpha\circ\beta)\times\id_{\mathbf1})=\rho\circ(\beta\times\id_{\mathbf1})=\beta\circ p_1=\operatorname{ev}\circ\langle p_1,{!}\rangle=\operatorname{ev}\), i la unicitat de la transposada dona \(\alpha\circ\beta=\id_{A^{\mathbf1}}\). Per tant \(A^{\mathbf1}\cong A\): fixar l'única entrada d'una aplicació \(\mathbf1\to A\) és el mateix que triar un punt de \(A\).
\end{solucio}

\begin{exercici}
\textbf{\(A^{\mathbf0}\cong\mathbf1\) i \(\mathbf1^{A}\cong\mathbf1\).} Demostreu totes dues identificacions. (Part de la proposició~\ref{prop:exponencials-basics}.)
\end{exercici}
\begin{solucio}
Per a \(A^{\mathbf0}\): una aplicació \(X\times\mathbf0\to A\) surt del conjunt inicial, perquè \(X\times\mathbf0\cong\mathbf0\) (la projecció \(X\times\mathbf0\to\mathbf0\) és un isomorfisme per l'estrictesa de \(\mathbf0\), proposició~\ref{prop:inicial-estricte}); per tant n'hi ha exactament una. Per la propietat universal de l'exponencial, hi ha exactament una aplicació \(X\to A^{\mathbf0}\) per a cada \(X\), i això caracteritza \(A^{\mathbf0}\) com a terminal: \(A^{\mathbf0}\cong\mathbf1\). Per a \(\mathbf1^{A}\): una aplicació \(X\times A\to\mathbf1\) és única perquè \(\mathbf1\) és terminal; de nou per la propietat universal, hi ha exactament una \(X\to\mathbf1^{A}\) per a cada \(X\), i \(\mathbf1^{A}\cong\mathbf1\).
\end{solucio}

\begin{exercici}
\textbf{El conjunt de parts com a exponencial.} Expliqueu per què \(P(A):=\mathbf2^{A}\) mereix el nom de \enquote{conjunt de parts}. (Model de la definició~\ref{def:conjunt-parts}.)
\end{exercici}
\begin{solucio}
Pels punts de l'exponencial (teorema~\ref{teo:punts-exponencial}), els punts globals de \(\mathbf2^{A}\) es corresponen amb les aplicacions \(A\to\mathbf2\). Aquestes són precisament les aplicacions característiques del capítol~\ref{cap:parts}, i pel classificador de parts es corresponen, al seu torn, amb les parts de \(A\) llevat d'equivalència. Per tant, els punts de \(\mathbf2^{A}\) classifiquen les parts de \(A\); això justifica definir \(P(A):=\mathbf2^{A}\). L'existència del conjunt de parts no requereix cap axioma nou: és el cas \(B=\mathbf2\) de l'axioma d'exponencials (axioma~\ref{ax:exponencials}).
\end{solucio}

\begin{exercici}
\textbf{Els quatre punts de \(\mathbf2^{\mathbf2}\).} Demostreu que \(\mathbf2^{\mathbf2}\cong\mathbf2\times\mathbf2\) i identifiqueu els quatre punts amb les quatre aplicacions \(\mathbf2\to\mathbf2\).
\end{exercici}
\begin{solucio}
Com que \(\mathbf2=\mathbf1+\mathbf1\), l'exponencial d'una suma (proposició~\ref{prop:exponencial-suma}) dona \(\mathbf2^{\mathbf2}=\mathbf2^{\mathbf1+\mathbf1}\cong\mathbf2^{\mathbf1}\times\mathbf2^{\mathbf1}\cong\mathbf2\times\mathbf2\), usant \(A^{\mathbf1}\cong A\). Els punts de \(\mathbf2\times\mathbf2\) són els parells \(\langle a,b\rangle\) amb \(a,b\in\{0,1\}\), que són exactament les quatre injeccions \(\delta_{ab}\) del capítol~\ref{cap:valors} (proposició~\ref{prop:quatre-injeccions}); pel teorema~\ref{teo:punts-exponencial}, un punt de \(\mathbf2^{\mathbf2}\) és una aplicació \(\mathbf2\to\mathbf2\), i la correspondència envia \(\langle a,b\rangle\) a l'aplicació definida per casos \([a,b]\), que val \(a\) sobre \(0\) i \(b\) sobre \(1\). Les quatre són: les dues constants \([0,0]\) i \([1,1]\), la identitat \([0,1]\) i la negació \([1,0]\). Aquest és el primer recompte d'aplicacions del llibre, i el patró ---sumes a l'exponent es tornen productes--- és el que el capítol~\ref{cap:finits} farà servir per comptar \(B^A\) quan \(A\) i \(B\) siguin finits.
\end{solucio}

\section{L'exponencial d'una suma}\label{prac:exp-suma}

\begin{exercici}
\textbf{\(B^{A+C}\cong B^{A}\times B^{C}\).} Demostreu que l'exponencial d'una suma es descompon com un producte d'exponencials. (Reproduïu la proposició~\ref{prop:exponencial-suma}.)
\end{exercici}
\begin{solucio}
Comprovem que \(B^{A}\times B^{C}\) té la propietat universal de \(B^{A+C}\); la unicitat (proposició~\ref{prop:exponencial-unic}) donarà l'isomorfisme. Per a tot \(X\), les aplicacions \(X\to B^{A+C}\) es corresponen amb les aplicacions \(X\times(A+C)\to B\). Com que el producte distribueix sobre la suma (teorema~\ref{teo:distributivitat}), \(X\times(A+C)\cong(X\times A)+(X\times C)\), i per la propietat universal de la suma (definició~\ref{def:suma-conjunts}) una aplicació d'aquest conjunt cap a \(B\) és una parella d'aplicacions \(X\times A\to B\) i \(X\times C\to B\). Transposant cada component, aquestes corresponen a \(X\to B^{A}\) i \(X\to B^{C}\), és a dir a una única \(X\to B^{A}\times B^{C}\). Totes les correspondències són naturals en \(X\), de manera que \(B^{A}\times B^{C}\) satisfà la propietat universal de \(B^{A+C}\).
\end{solucio}

\section{Composició interna i famílies}\label{prac:exp-composicio}

\begin{exercici}
\textbf{La composició com a aplicació entre exponencials.} Demostreu que existeix una única aplicació \(c\colon C^{B}\times B^{A}\to C^{A}\) que sobre punts envia \((g,f)\) a \(g\circ f\). (Model de la proposició~\ref{prop:composicio-interna}.)
\end{exercici}
\begin{solucio}
Per la propietat universal de \(C^{A}\), definir \(c\) equival a donar una aplicació \((C^{B}\times B^{A})\times A\to C\). La prenem com
\[
(C^{B}\times B^{A})\times A\xrightarrow{\ \sim\ }C^{B}\times(B^{A}\times A)\xrightarrow{\id\times\operatorname{ev}}C^{B}\times B\xrightarrow{\operatorname{ev}}C,
\]
on la primera fletxa és l'associativitat del producte i les altres dues són avaluacions. La transposada d'aquesta aplicació és la \(c\) buscada. Sobre punts globals, seguint les avaluacions, \((g,f)\) es transforma en \(a\mapsto g(f(a))\), és a dir \(g\circ f\). La unicitat de \(c\) \emph{amb aquesta propietat sobre punts} no la dona la transposada, que només és única per a una aplicació \((C^B\times B^A)\times A\to C\) ja fixada, sinó la separació global: dues aplicacions amb domini \(C^B\times B^A\) que coincideixen sobre tots els punts, que són els parells \(\langle\bar g,\bar f\rangle\), són iguals. Anàlogament, la transposada de la projecció \(\mathbf1\times A\to A\) dona el punt \(\jmath\colon\mathbf1\to A^{A}\) que correspon a \(\id_A\). La composició, que era una operació sobre fletxes, s'ha tornat una aplicació entre conjunts.
\end{solucio}

\begin{exercici}
\textbf{Naturalitat de la transposició.} Demostreu que la correspondència entre aplicacions \(X\times A\to B\) i aplicacions \(X\to B^A\) és natural en \(X\): per a tota \(t\colon X'\to X\), currificar i precompondre amb \(t\) coincideix amb precompondre amb \(t\times\id_A\) i currificar. (Part de la proposició~\ref{prop:families}.)
\end{exercici}
\begin{solucio}
Sigui \(h\colon X\times A\to B\) amb currificació \(\widehat h=\lambda(h)\colon X\to B^A\). D'una banda, \(\widehat h\circ t\colon X'\to B^A\). De l'altra, considerem \(h\circ(t\times\id_A)\colon X'\times A\to B\) i la seva currificació \(\lambda(h\circ(t\times\id_A))\). Per veure que coincideixen n'hi ha prou de comprovar que tenen la mateixa descurrificació i usar la unicitat de la transposada:
\[
\operatorname{ev}\circ((\widehat h\circ t)\times\id_A)=\operatorname{ev}\circ(\widehat h\times\id_A)\circ(t\times\id_A)=h\circ(t\times\id_A),
\]
on hem usat \((\widehat h\circ t)\times\id_A=(\widehat h\times\id_A)\circ(t\times\id_A)\) i la definició de \(\widehat h\). Com que \(\lambda(h\circ(t\times\id_A))\) té, per definició, aquesta mateixa descurrificació, la unicitat de la transposada dona \(\widehat h\circ t=\lambda(h\circ(t\times\id_A))\). Aquesta és la naturalitat que fa que els elements generalitzats de \(B^A\) es comportin com a famílies coherents sota reindexació.
\end{solucio}

\chapter{El conjunt de parts i la pertinença interna}\label{cap:practica-parts-interna}

Aquest capítol acompanya el capítol homònim de la part teòrica (capítol~\ref{cap:parts-interna} de la Part~\ref{part:teoria}). Els exercicis resolts en segueixen l'ordre d'exposició --- la pertinença interna, les parts parametritzades, les operacions booleanes internes, la inclusió com a ordre sobre \(P(A)\), i la imatge inversa interna amb els quantificadors--- i serveixen de model per als exercicis proposats de la Part~\ref{part:exercicis}. Cada bloc fixa una tècnica: reduir una afirmació interna a la pertinença generalitzada via l'avaluació, calcular operacions internes descurrificant-les cap a \(\mathbf2\), i llegir adjuncions sobre elements generalitzats. Convé intentar cada enunciat abans de llegir-ne la solució.

\section{La pertinença interna}\label{prac:pertinenca-interna}

\begin{exercici}
\textbf{Lectura de la pertinença interna.} Siguin \(a\colon X\to A\) i \(\alpha\colon X\to P(A)\). Demostreu que \(a\mathrel{\varepsilon}\alpha\) equival a \(\operatorname{ev}\circ\langle\alpha,a\rangle=1\circ{!}_X\), i que si \(\alpha=\ulcorner m\urcorner\circ{!}_X\) és el nom constant d'una part \(m\), aleshores \(a\mathrel{\varepsilon}\alpha\iff a\in_X m\). (Reproduïu la proposició~\ref{prop:pertinenca-interna-lectura}.)
\end{exercici}
\begin{solucio}
Per la propietat característica (notació~\ref{not:caracteristica}), \(\langle a,\alpha\rangle\in_X{\in_A}\) equival a \(\chi_{\in_A}\circ\langle a,\alpha\rangle=1\circ{!}_X\). Com que \(\chi_{\in_A}=\operatorname{ev}\circ\sigma^{-1}\) i \(\sigma^{-1}\circ\langle a,\alpha\rangle=\langle\alpha,a\rangle\), això és \(\operatorname{ev}\circ\langle\alpha,a\rangle=1\circ{!}_X\). Per al cas constant, el nom \(\ulcorner m\urcorner\) descurrifica a \(\chi_m\): \(\operatorname{ev}\circ(\ulcorner m\urcorner\times\id_A)=\chi_m\circ\rho\). Component amb \(\langle{!}_X,a\rangle\) i usant \(\alpha=\ulcorner m\urcorner\circ{!}_X\), s'obté \(\operatorname{ev}\circ\langle\alpha,a\rangle=\chi_m\circ a\); per tant \(a\mathrel{\varepsilon}\alpha\iff\chi_m\circ a=1\circ{!}_X\), que és \(a\in_X m\). La pertinença interna és, doncs, la pertinença generalitzada vista dins del conjunt \(P(A)\).
\end{solucio}

\begin{exercici}
\textbf{La pertinença interna és universal.} Sigui \(\alpha\colon X\to P(A)\) i \(s_\alpha\rightarrowtail X\times A\) la part que li correspon. Demostreu que \(a\mathrel{\varepsilon}(\alpha\circ x)\iff\langle x,a\rangle\in_Y s_\alpha\), i que \(\in_A\) és la part associada a \(\id_{P(A)}\). (Reproduïu el corol·lari~\ref{cor:pertinenca-universal}.)
\end{exercici}
\begin{solucio}
Per la proposició~\ref{prop:pertinenca-interna-lectura}, \(a\mathrel{\varepsilon}(\alpha\circ x)\) equival a \(\operatorname{ev}\circ\langle\alpha\circ x,a\rangle=1\circ{!}_Y\). Com que \(\langle\alpha\circ x,a\rangle=(\alpha\times\id_A)\circ\langle x,a\rangle\), la condició és \(\operatorname{ev}\circ(\alpha\times\id_A)\circ\langle x,a\rangle=1\circ{!}_Y\); i \(\operatorname{ev}\circ(\alpha\times\id_A)\) és, per construcció, la característica de \(s_\alpha\). Això és \(\langle x,a\rangle\in_Y s_\alpha\). Prenent \(X=P(A)\) i \(\alpha=\id_{P(A)}\), la part associada té característica \(\operatorname{ev}\), que sobre \(A\times P(A)\) és la de \(\in_A\). Així, la relació de pertinença interna és el cas universal (\(\alpha=\id\)) del qual tota altra pertinença s'obté per reindexació.
\end{solucio}

\section{Parts parametritzades}\label{prac:parts-parametritzades}

\begin{exercici}
\textbf{Aplicacions cap a \(P(A)\) són parts de \(X\times A\).} Demostreu que hi ha una correspondència biunívoca, natural en \(X\), entre aplicacions \(\alpha\colon X\to P(A)\), aplicacions \(X\times A\to\mathbf2\), i parts de \(X\times A\) llevat d'equivalència. (Reproduïu el teorema~\ref{teo:parts-parametritzades}.)
\end{exercici}
\begin{solucio}
L'equivalència entre \(\alpha\colon X\to P(A)\) i \(X\times A\to\mathbf2\) és la propietat universal de l'exponencial \(P(A)=\mathbf2^{A}\) (definició~\ref{def:exponencial}): a cada \(\alpha\) li correspon la seva descurrificació \(\operatorname{ev}\circ(\alpha\times\id_A)\), biunívoca per la unicitat de la transposada; la naturalitat en \(X\) és la de la proposició~\ref{prop:families} amb \(B=\mathbf2\). L'equivalència entre \(X\times A\to\mathbf2\) i parts de \(X\times A\) és l'axioma de representació (axioma~\ref{ax:representacio}): tota aplicació cap a \(\mathbf2\) és la característica d'una única part llevat d'equivalència. Component les dues correspondències s'obté l'enunciat. En conseqüència, un element generalitzat de \(P(A)\) és una \emph{família} de parts de \(A\), no una sola part.
\end{solucio}

\section{Operacions booleanes internes}\label{prac:booleanes-internes}

\begin{exercici}
\textbf{Càlcul puntual de la intersecció interna.} Siguin \(\alpha,\beta\colon X\to P(A)\) i \(a\colon X\to A\). Demostreu que \(a\mathrel{\varepsilon}(\alpha\cap\beta)\) si i només si \(a\mathrel{\varepsilon}\alpha\) i \(a\mathrel{\varepsilon}\beta\). (Part de la proposició~\ref{prop:booleanes-puntuals}.)
\end{exercici}
\begin{solucio}
Per la proposició~\ref{prop:pertinenca-interna-lectura}, \(a\mathrel{\varepsilon}(\alpha\cap\beta)\) equival a \(\operatorname{ev}\circ\langle\alpha\cap\beta,a\rangle=1\circ{!}_X\). Descurrificant la definició de \(\cap\) (definició~\ref{def:booleanes-internes}), aquesta composició és \(\wedge\circ\langle\operatorname{ev}\circ\langle\alpha,a\rangle,\;\operatorname{ev}\circ\langle\beta,a\rangle\rangle\). Com que \(\wedge\) és la característica de \(\delta_{11}=\langle1,1\rangle\colon\mathbf1\hookrightarrow\mathbf2\times\mathbf2\), aquesta composició val \(1\circ{!}_X\) si i només si \(\langle\operatorname{ev}\circ\langle\alpha,a\rangle,\operatorname{ev}\circ\langle\beta,a\rangle\rangle\) es factoritza per \(\delta_{11}\), és a dir, si i només si tots dos components són \(1\circ{!}_X\): \(a\mathrel{\varepsilon}\alpha\) i \(a\mathrel{\varepsilon}\beta\). L'argument val a tot estadi \(X\) perquè no ha usat la taula de veritat, sinó la definició de \(\wedge\) com a característica. Compareu-ho amb la unió: \(a\mathrel{\varepsilon}(\alpha\cup\beta)\) no equival, en general, a \enquote{\(a\mathrel{\varepsilon}\alpha\) o \(a\mathrel{\varepsilon}\beta\)} (proposició~\ref{prop:booleanes-puntuals}).
\end{solucio}

\begin{exercici}
\textbf{Les operacions internes estenen les de parts.} Demostreu que sobre els noms de dues parts \(m,n\) de \(A\) es té \(\ulcorner m\urcorner\cap\ulcorner n\urcorner=\ulcorner m\cap n\urcorner\), i anàlogament per a la unió i el complement. (Part de la proposició~\ref{prop:booleanes-puntuals}.)
\end{exercici}
\begin{solucio}
Prenent \(\alpha=\ulcorner m\urcorner\circ{!}_X\) i \(\beta=\ulcorner n\urcorner\circ{!}_X\), l'exercici anterior i la proposició~\ref{prop:pertinenca-interna-lectura} converteixen l'equivalència puntual en \(a\in_X(m\cap n)\iff a\in_X m\text{ i }a\in_X n\), que és la caracterització de la intersecció de parts (proposició~\ref{prop:pertinenca-interseccio}). Com que dues parts amb la mateixa pertinença generalitzada són equivalents (proposició~\ref{prop:inclusio-pertinenca}) i parts equivalents tenen el mateix nom (proposició~\ref{prop:caracteristica-determina-part}), els punts classificadors coincideixen: \(\ulcorner m\urcorner\cap\ulcorner n\urcorner=\ulcorner m\cap n\urcorner\). Per a la unió i el complement no es pot repetir l'argument puntual, perquè la lectura amb \enquote{o} i \enquote{no} no val a tots els estadis; però no cal: les tres identitats de la proposició~\ref{prop:booleanes-puntuals} amb \(X=\mathbf1\) diuen que la descurrificació de \(\ulcorner m\urcorner\cup\ulcorner n\urcorner\) és \(\vee\circ\langle\chi_m,\chi_n\rangle=\chi_{m\cup n}\) (definició~\ref{def:operacions-parts}) i la de \(\neg\ulcorner m\urcorner\) és \(\neg\circ\chi_m=\chi_{m^c}\) (proposició~\ref{prop:caracteristica-complement}), i la unicitat del nom conclou. Les lleis booleanes de les parts es transporten així a \(P(A)\) (observació~\ref{obs:lleis-transportades}).
\end{solucio}

\section{La inclusió com a ordre}\label{prac:inclusio-ordre}

\begin{exercici}
\textbf{Lectura de la inclusió interna.} Siguin \(\alpha,\beta\colon X\to P(A)\) amb parts associades \(s_\alpha,s_\beta\hookrightarrow X\times A\). Demostreu que \(\langle\alpha,\beta\rangle\in_X{\subseteq_A}\) si i només si \(s_\alpha\subseteq s_\beta\), i que sobre els noms recupera \(m\subseteq_A n\). (Reproduïu la proposició~\ref{prop:inclusio-interna-lectura}.)
\end{exercici}
\begin{solucio}
Per definició, \(\subseteq_A=\forall_\pi(I_A)\), on \(I_A\) és la part de \(P(A)\times P(A)\times A\) classificada per la implicació puntual \(\Rightarrow\circ\langle\operatorname{ev}_1,\operatorname{ev}_2\rangle\) i \(\pi\) és la projecció que oblida \(A\). La lectura de \(\forall_\pi\) (proposició~\ref{prop:quantificadors-lectura}) diu que \(\langle\alpha,\beta\rangle\in_X\forall_\pi(I_A)\) exactament quan tot \(\langle\alpha\circ x,\beta\circ x,a\rangle\) (amb \(x\colon Y\to X\), \(a\colon Y\to A\)) pertany a \(I_A\), és a dir, quan \(\neg u\vee v=1\circ{!}_Y\) amb \(u=\operatorname{ev}\circ\langle\alpha\circ x,a\rangle\) i \(v=\operatorname{ev}\circ\langle\beta\circ x,a\rangle\). Ara, \(\neg u\vee v\) classifica \(T^c\cup U\), on \(T\) i \(U\) són les parts de \(Y\) classificades per \(u\) i \(v\); i \(T^c\cup U\) és total si i només si \(T\subseteq U\) (lleis booleanes). Com que \(T=\langle x,a\rangle^{-1}s_\alpha\) i \(U=\langle x,a\rangle^{-1}s_\beta\), la condició per a tots els \(x,a\) equival a \(s_\alpha\subseteq s_\beta\): prenent \(Y=s_\alpha\) i \(\langle x,a\rangle\) la seva inclusió s'obté la inclusió, i la monotonia de la imatge inversa dona el recíproc. Sobre els noms, \(X=\mathbf1\) i \(s_{\ulcorner m\urcorner}\) és \(m\) dins \(\mathbf1\times A\cong A\), de manera que la condició és \(m\subseteq_A n\).
\end{solucio}

\begin{exercici}
\textbf{Reflexivitat i transitivitat de la inclusió interna.} Demostreu que \(\langle\alpha,\alpha\rangle\in_X{\subseteq_A}\) i que la inclusió interna és transitiva. (Part del teorema~\ref{teo:inclusio-ordre}.)
\end{exercici}
\begin{solucio}
Per l'exercici anterior, \(\langle\alpha,\beta\rangle\in_X{\subseteq_A}\) equival a \(s_\alpha\subseteq s_\beta\) com a parts de \(X\times A\). La reflexivitat és \(s_\alpha\subseteq s_\alpha\), i la transitivitat és la de la inclusió de parts: si \(s_\alpha\subseteq s_\beta\) i \(s_\beta\subseteq s_\gamma\), tot element generalitzat de \(s_\alpha\) ho és de \(s_\gamma\) (proposició~\ref{prop:inclusio-pertinenca}). L'antisimetria, que completa l'ordre, és que dues parts incloses l'una en l'altra són equivalents i tenen la mateixa característica, d'on \(\alpha=\beta\) per la unicitat de la transposada; \(\cap\) i \(\cup\) en són l'ínfim i el suprem perquè \(s_{\alpha\cap\beta}=s_\alpha\cap s_\beta\) i \(s_{\alpha\cup\beta}=s_\alpha\cup s_\beta\).
\end{solucio}

\section{Imatge inversa interna i quantificadors}\label{prac:quantificadors}

\begin{exercici}
\textbf{Un testimoni només després d'un recobriment.} Sigui \(f=\pi_1\colon\mathbf2\times\mathbf2\to\mathbf2\) la primera projecció i \(\alpha\colon\mathbf1\to P(\mathbf2\times\mathbf2)\) el nom de la diagonal \(\Delta\). Comproveu que \(\exists_f\alpha\) és el nom de la part total de \(\mathbf2\), i que per a l'element generalitzat \(b=\id_{\mathbf2}\colon\mathbf2\to\mathbf2\) es té \(b\mathrel{\varepsilon}\exists_f\alpha\), amb testimoni \(a=\langle\id,\id\rangle\) ja a l'estadi \(\mathbf2\). Expliqueu per què, en general, cal admetre un recobriment, i per què amb l'elecció no cal.
\end{exercici}
\begin{solucio}
Sobre noms, \(\exists_f\ulcorner\Delta\urcorner=\ulcorner\pi_{1*}\Delta\urcorner\), i la imatge de la diagonal per la primera projecció és tot \(\mathbf2\), perquè \(\pi_1\circ\langle\id,\id\rangle=\id\) és un epimorfisme; per tant \(\exists_f\alpha\) és el nom de la part total. Per a \(b=\id_{\mathbf2}\), la lectura de la proposició~\ref{prop:quantificadors-lectura} demana un recobriment \(p\colon X'\twoheadrightarrow\mathbf2\) i un \(a\colon X'\to\mathbf2\times\mathbf2\) amb \(\pi_1\circ a=b\circ p\) i \(a\in\Delta\); aquí \(p=\id\) i \(a=\langle\id,\id\rangle\) serveixen. En general, el que la construcció garanteix és que \(\langle\id_X,b\rangle\) pertany a la imatge directa \((\id\times f)_*(s_\alpha)\), i pertànyer a una imatge dona la preimatge només després de fer el pullback de l'epimorfisme \(s_\alpha\twoheadrightarrow\operatorname{Im}\), que és el recobriment; recollir-la ja sobre \(X\) demana una secció d'aquest epimorfisme, i això és l'elecció. És la mateixa situació que la composició de relacions (proposició~\ref{prop:composicio-es-relacio}).
\end{solucio}

\begin{exercici}
\textbf{La imatge inversa interna estén \(f^{-1}\).} Sigui \(f\colon A\to B\) i \(P(f)\colon P(B)\to P(A)\) la imatge inversa interna. Demostreu que \(P(f)\circ\ulcorner n\urcorner=\ulcorner f^{-1}(n)\urcorner\) per a tota part \(n\) de \(B\). (Reproduïu la proposició~\ref{prop:imatge-inversa-interna}.)
\end{exercici}
\begin{solucio}
Per la definició de \(P(f)\) (definició~\ref{def:imatge-inversa-interna}) amb \(\beta=\ulcorner n\urcorner\circ{!}_X\) i la proposició~\ref{prop:pertinenca-interna-lectura}, per a tot \(a\colon X\to A\),
\[
a\mathrel{\varepsilon}\bigl(P(f)\circ\ulcorner n\urcorner\bigr)\iff(f\circ a)\in_X n\iff a\in_X f^{-1}(n),
\]
on la darrera equivalència és la proposició~\ref{prop:pertinenca-imatge-inversa}. Com que la pertinença generalitzada determina la part i el nom en determina la descurrificació, \(P(f)\circ\ulcorner n\urcorner=\ulcorner f^{-1}(n)\urcorner\). La imatge inversa interna és, doncs, la versió \enquote{com a aplicació entre conjunts de parts} de l'operació \(f^{-1}\) sobre parts.
\end{solucio}

\begin{exercici}
\textbf{La imatge inversa commuta amb el complement.} Demostreu que \(P(f)\circ\neg_B=\neg_A\circ P(f)\). (Reproduïu la proposició~\ref{prop:imatge-inversa-complement}.)
\end{exercici}
\begin{solucio}
Les dues aplicacions tenen domini \(P(B)\) i codomini \(P(A)\), de manera que basta comparar-ne les descurrificacions \(P(B)\times A\to\mathbf2\). Per a \(P(f)\circ\neg_B\): \(\operatorname{ev}\circ((P(f)\circ\neg_B)\times\id_A)=\operatorname{ev}\circ(P(f)\times\id_A)\circ(\neg_B\times\id_A)=\operatorname{ev}\circ(\id_{P(B)}\times f)\circ(\neg_B\times\id_A)\), per la definició de \(P(f)\); i com que \((\id\times f)\circ(\neg_B\times\id_A)=(\neg_B\times\id_B)\circ(\id\times f)\), això és \(\operatorname{ev}\circ(\neg_B\times\id_B)\circ(\id\times f)=\neg\circ\operatorname{ev}\circ(\id\times f)\) per la definició de \(\neg_B\). Per a \(\neg_A\circ P(f)\): \(\operatorname{ev}\circ((\neg_A\circ P(f))\times\id_A)=\neg\circ\operatorname{ev}\circ(P(f)\times\id_A)=\neg\circ\operatorname{ev}\circ(\id\times f)\). Les descurrificacions coincideixen i, per la unicitat de la transposada, les aplicacions també. Observeu que l'argument és un càlcul amb transposades i no passa per cap lectura puntual del complement, que només valdria sobre punts globals. Aquesta commutació és el que permet definir \(\forall_f\) a partir de \(\exists_f\) per dualitat.
\end{solucio}

\begin{exercici}
\textbf{Els quantificadors com a adjunts.} Siguin \(\alpha\colon X\to P(A)\) i \(\beta\colon X\to P(B)\), amb parts associades \(s_\alpha\) i \(s_\beta\). Demostreu que \(s_{\exists_f\alpha}\subseteq s_\beta\iff s_\alpha\subseteq s_{P(f)\beta}\) i que \(s_{P(f)\beta}\subseteq s_\alpha\iff s_\beta\subseteq s_{\forall_f\alpha}\). (Reproduïu la proposició~\ref{prop:quantificadors-adjuncio}.)
\end{exercici}
\begin{solucio}
Les parts associades són \(s_{\exists_f\alpha}=(\id_X\times f)_*s_\alpha\) (proposició~\ref{prop:quantificadors-lectura}, que és on intervé el lema de Beck--Chevalley) i \(s_{P(f)\beta}=(\id_X\times f)^{-1}s_\beta\) (proposició~\ref{prop:imatge-inversa-interna}). La primera equivalència és, doncs, l'adjunció imatge directa\,--\,imatge inversa del capítol~\ref{cap:imatges} (proposició~\ref{prop:adjuncio-imatges}) aplicada a \(\id_X\times f\). Per a la segona, \(s_{\forall_f\alpha}=\bigl((\id\times f)_*(s_\alpha^c)\bigr)^c\); aleshores \(s_\beta\subseteq s_{\forall_f\alpha}\) equival, per l'antimonotonia del complement, a \((\id\times f)_*(s_\alpha^c)\subseteq s_\beta^c\); per la primera adjunció, a \(s_\alpha^c\subseteq(\id\times f)^{-1}(s_\beta^c)\); i com que la imatge inversa commuta amb el complement, a \(s_\alpha^c\subseteq\bigl((\id\times f)^{-1}s_\beta\bigr)^c\), és a dir, \((\id\times f)^{-1}s_\beta\subseteq s_\alpha\). Així \(\exists_f\) és adjunt per l'esquerra i \(\forall_f\) per la dreta de la imatge inversa interna, a cada estadi \(X\).
\end{solucio}

\chapter{L'objecte dels nombres naturals}\label{cap:practica-naturals}

Aquest capítol acompanya el capítol homònim de la part teòrica (capítol~\ref{cap:naturals} de la Part~\ref{part:teoria}). Els exercicis resolts en segueixen l'ordre d'exposició --- la propietat universal de recursió, la iteració, el principi d'inducció, les equacions de Peano, l'aritmètica (suma, producte, exponenciació) i l'ordre--- i serveixen de model per als exercicis proposats de la Part~\ref{part:exercicis}. Cada bloc fixa una tècnica: definir aplicacions per recursió i provar que coincideixen per la unicitat de la recursora, demostrar igualtats aritmètiques per inducció sobre un argument, i raonar amb l'ordre via la caracterització \(m\le n\iff\exists k\,(m+k=n)\). Convé intentar cada enunciat abans de llegir-ne la solució.

\section{Propietat universal i unicitat}\label{prac:nno-universal}

\begin{exercici}
\textbf{Unicitat de l'objecte dels nombres naturals.} Suposeu que \((\N,0,s)\) i \((\N',0',s')\) compleixen totes dues la propietat universal de recursió. Demostreu que hi ha un únic isomorfisme \(\varphi\colon\N\to\N'\) amb \(\varphi\circ0=0'\) i \(\varphi\circ s=s'\circ\varphi\). (Reproduïu el teorema~\ref{teo:nno-unicitat}.)
\end{exercici}
\begin{solucio}
Per la propietat universal de \((\N,0,s)\) aplicada al sistema recursiu \((\N',0',s')\), hi ha una única \(\varphi\colon\N\to\N'\) amb \(\varphi\circ0=0'\) i \(\varphi\circ s=s'\circ\varphi\). Simètricament, hi ha una única \(\psi\colon\N'\to\N\) amb \(\psi\circ0'=0\) i \(\psi\circ s'=s\circ\psi\). La composició \(\psi\circ\varphi\) compleix \((\psi\circ\varphi)\circ0=0\) i \((\psi\circ\varphi)\circ s=s\circ(\psi\circ\varphi)\); però la identitat \(\id_\N\) compleix les mateixes dues condicions, i la propietat universal de \((\N,0,s)\) aplicada a si mateix diu que hi ha una sola aplicació així. Per tant \(\psi\circ\varphi=\id_\N\), i igualment \(\varphi\circ\psi=\id_{\N'}\). Així \(\varphi\) és un isomorfisme, únic amb les propietats requerides. És el patró d'unicitat de tota construcció universal.
\end{solucio}

\begin{exercici}
\textbf{Existència i unicitat de la iterada.} Sigui \(t\colon A\to A\). Demostreu que existeix una única \(\operatorname{it}(t)\colon\N\times A\to A\) amb \(t^{0}(x)=x\) i \(t^{\,s\circ n}(x)=t(t^{n}(x))\). (Reproduïu la proposició~\ref{prop:iterada}.)
\end{exercici}
\begin{solucio}
Considerem l'exponencial \(A^{A}\) i el sistema recursiu \((A^{A},\ulcorner\id_A\urcorner,C_t)\), on \(\ulcorner\id_A\urcorner\) és el nom de la identitat i \(C_t\colon A^{A}\to A^{A}\) és la composició interna per \(t\) per l'esquerra (el nom de \(g\mapsto t\circ g\); vegeu la proposició~\ref{prop:composicio-interna}). Per l'axioma dels naturals (axioma~\ref{ax:naturals}) hi ha una única \(\pi\colon\N\to A^{A}\) amb \(\pi\circ0=\ulcorner\id_A\urcorner\) i \(\pi\circ s=C_t\circ\pi\); intuïtivament \(\pi(n)=\ulcorner t^{n}\urcorner\). Descurrificant, \(\operatorname{it}(t):=\operatorname{ev}\circ(\pi\times\id_A)\), i les dues equacions de \(\pi\) es llegeixen, via l'avaluació, com \(t^{0}(x)=x\) i \(t^{\,s\circ n}(x)=t(t^{n}(x))\). La unicitat de \(\operatorname{it}(t)\) es redueix, per la unicitat de la transposada (definició~\ref{def:exponencial}), a la de \(\pi\), que és la de la recursora.
\end{solucio}

\section{El principi d'inducció}\label{prac:induccio}

\begin{exercici}
\textbf{Unicitat per inducció.} Siguin \(f,g\colon\N\to A\) amb \(f\circ0=g\circ0\), i suposeu que la part on \(f\) i \(g\) coincideixen és tancada pel successor. Demostreu que \(f=g\). (Reproduïu la proposició~\ref{prop:unicitat-induccio}.)
\end{exercici}
\begin{solucio}
Sigui \(e\colon E\hookrightarrow\N\) l'equalitzador de \(f\) i \(g\), és a dir la part de \(\N\) on totes dues coincideixen. La hipòtesi \(f\circ0=g\circ0\) diu que \(E\) conté el zero, i la hipòtesi de tancament diu que \(E\) és tancada pel successor. Pel principi d'inducció (teorema~\ref{teo:induccio}), tota part de \(\N\) que conté el zero i és tancada pel successor és tot \(\N\); per tant \(e\) és un isomorfisme i \(f\) i \(g\) coincideixen a tot \(\N\), és a dir \(f=g\). Aquesta és la cara \enquote{demostrativa} de la propietat universal: la part d'unicitat de la recursora és exactament el principi d'inducció (observació~\ref{obs:recursio-i-induccio}).
\end{solucio}

\section{Les equacions de Peano}\label{prac:peano}

\begin{exercici}
\textbf{El zero no és cap successor.} Demostreu que no hi ha cap \(n\colon X\to\N\) amb \(X\) no inicial tal que \(s\circ n=0\circ{!}_X\). (Part del teorema~\ref{teo:peano}.)
\end{exercici}
\begin{solucio}
Definim per recursió \(d\colon\N\to\mathbf2\) amb \(d\circ0=1\) i \(d\circ s=0\circ{!}_\N\) (la recursora del sistema \((\mathbf2,1,0\circ{!}_{\mathbf2})\)). Si fos \(s\circ n=0\circ{!}_X\) per a algun \(n\colon X\to\N\) amb \(X\) no inicial, component amb \(d\) obtindríem \(d\circ s\circ n=d\circ0\circ{!}_X\), és a dir \(0\circ{!}_X=1\circ{!}_X\). Però \(0\neq1\) (distinció dels valors de veritat, capítol~\ref{cap:valors}), i com que \(X\) no és inicial això és impossible per la separació global (axioma~\ref{ax:separacio-globals}). Per tant cap successor no és el zero.
\end{solucio}

\begin{exercici}
\textbf{El successor és injectiu.} Demostreu que \(s\colon\N\to\N\) és un monomorfisme construint un predecessor \(p\) amb \(p\circ s=\id_\N\). Expliqueu primer per què la recursora del sistema \((\N,0,\id_\N)\) \emph{no} serveix. (Part del teorema~\ref{teo:peano}.)
\end{exercici}
\begin{solucio}
La temptació és definir \(p\) per recursió amb \(p\circ0=0\) i \(p\circ s=\id_\N\). Però la forma de l'axioma és \(u\circ s=t\circ u\): amb \(t=\id_\N\) la condició és \(u\circ s=u\), no \(u\circ s=\id\), i la recursora de \((\N,0,\id_\N)\) és la constant \(0\circ{!}_\N\) (compleix les dues equacions, i és l'única). Cal recordar més: definim \(q\colon\N\to\N\times\N\) com la recursora de \((\N\times\N,\langle0,0\rangle,\langle p_2,s\circ p_2\rangle)\), de manera que \(q\circ0=\langle0,0\rangle\) i \(q\circ s=\langle p_2,s\circ p_2\rangle\circ q\); intuïtivament \(q(n)=(n-1,n)\). Primer, \(p_2\circ q=\id_\N\): tots dos compleixen \(u\circ0=0\) i \(u\circ s=s\circ u\), i la recursora de \((\N,0,s)\) és única. Després, amb \(p:=p_1\circ q\), \(p\circ s=p_1\circ\langle p_2,s\circ p_2\rangle\circ q=p_2\circ q=\id_\N\). Així \(s\) és una secció de \(p\) i, en particular, un monomorfisme (proposició~\ref{prop:seccio-mono-retraccio-epi}): si \(s\circ n=s\circ n'\), component amb \(p\) resulta \(n=n'\). La lliçó és la de l'observació~\ref{obs:predecessor-parell}: una funció que mira enrere es fa recursiva recordant el parell.
\end{solucio}

\section{Aritmètica per inducció}\label{prac:aritmetica}

\begin{exercici}
\textbf{Neutralitat de la suma per l'esquerra.} A partir de la definició recursiva \(m+0=m\), \(m+(s\circ n)=s\circ(m+n)\), demostreu que \(0+n=n\) per a tot \(n\). (Lema de la proposició~\ref{prop:suma}.)
\end{exercici}
\begin{solucio}
Fem inducció sobre \(n\) (teorema~\ref{teo:induccio}), usant la unicitat per inducció (proposició~\ref{prop:unicitat-induccio}) per garantir que la part on val la igualtat és tot \(\N\). Cas base: \(0+0=0\) directament de la definició. Pas de \(n\) a \(s\circ n\): si \(0+n=n\), aleshores
\[
0+(s\circ n)=s\circ(0+n)=s\circ n,
\]
on la primera igualtat és la definició de suma i la segona la hipòtesi d'inducció. Per inducció, \(0+n=n\) a tot \(\N\). Juntament amb \(n+0=n\), això dona la neutralitat del zero per totes dues bandes.
\end{solucio}

\begin{exercici}
\textbf{Associativitat de la suma.} Demostreu que \((k+m)+n=k+(m+n)\) per a tots \(k,m,n\). (Part de la proposició~\ref{prop:suma}.)
\end{exercici}
\begin{solucio}
Inducció sobre \(n\), amb \(k,m\) fixats. Cas base \(n=0\): \((k+m)+0=k+m=k+(m+0)\), per la definició de suma. Pas de \(n\) a \(s\circ n\):
\[
(k+m)+(s\circ n)=s\circ((k+m)+n)=s\circ(k+(m+n))=k+(m+(s\circ n)),
\]
on la primera i la tercera igualtats són la definició de la suma i la central és la hipòtesi d'inducció. Per inducció, l'associativitat val a tot \(\N\). El mateix esquema ---inducció sobre un argument amb la definició recursiva com a motor--- dona la commutativitat i la cancel·lació (proposició~\ref{prop:suma}).
\end{solucio}

\begin{exercici}
\textbf{Distributivitat del producte sobre la suma.} Demostreu que \(m\cdot(n+k)=m\cdot n+m\cdot k\). (Model de la proposició~\ref{prop:lleis-aritmetica}.)
\end{exercici}
\begin{solucio}
Inducció sobre \(k\), amb \(m,n\) fixats. Cas base \(k=0\): per les definicions de suma i producte (definicions~\ref{def:suma-naturals} i~\ref{def:producte-exp}), \(m\cdot(n+0)=m\cdot n=m\cdot n+0=m\cdot n+m\cdot0\). Pas de \(k\) a \(s\circ k\):
\[
m\cdot(n+(s\circ k))=m\cdot(s\circ(n+k))=m\cdot(n+k)+m=(m\cdot n+m\cdot k)+m,
\]
usant la definició de suma, la de producte i la hipòtesi d'inducció; i reordenant amb l'associativitat de la suma (proposició~\ref{prop:suma}),
\[
(m\cdot n+m\cdot k)+m=m\cdot n+(m\cdot k+m)=m\cdot n+m\cdot(s\circ k).
\]
Per inducció, la distributivitat val a tot \(\N\). D'ella depenen l'associativitat i la commutativitat del producte i la primera llei dels exponents.
\end{solucio}

\section{L'ordre dels naturals}\label{prac:ordre}

En aquesta secció \(k,m,n\) són punts globals de \(\N\), d'acord amb l'observació~\ref{obs:ordre-sobre-punts}: les propietats de l'ordre són afirmacions sobre parts de \(\N\times\N\), i per la separació global n'hi ha prou amb comprovar-les sobre punts, on \(m\le n\) vol dir exactament que existeix un punt \(k\) amb \(m+k=n\).

\begin{exercici}
\textbf{L'ordre és reflexiu i transitiu.} A partir de \(m\le n\iff\exists k\,(m+k=n)\), demostreu que \(n\le n\) i que \(k\le m\), \(m\le n\) impliquen \(k\le n\). (Part de la proposició~\ref{prop:ordre-total}.)
\end{exercici}
\begin{solucio}
Reflexivitat: \(n+0=n\) (definició~\ref{def:suma-naturals}), de manera que prenent \(k=0\) es té \(n\le n\). Transitivitat: si \(k+a=m\) i \(m+b=n\), aleshores per l'associativitat de la suma (proposició~\ref{prop:suma})
\[
k+(a+b)=(k+a)+b=m+b=n,
\]
de manera que \(k\le n\) amb increment \(a+b\). La caracterització \enquote{existeix un increment} redueix les propietats de l'ordre a propietats de la suma.
\end{solucio}

\begin{exercici}
\textbf{Antisimetria de l'ordre.} Demostreu que si \(m\le n\) i \(n\le m\), aleshores \(m=n\). (Part de la proposició~\ref{prop:ordre-total}.)
\end{exercici}
\begin{solucio}
Suposem \(m+k=n\) i \(n+l=m\). Substituint la primera a la segona i usant l'associativitat,
\[
m+(k+l)=(m+k)+l=n+l=m=m+0.
\]
Per la cancel·lació de la suma (proposició~\ref{prop:suma}), \(k+l=0\); i com que una suma nul·la té tots dos sumands nuls (proposició~\ref{prop:suma}), \(k=l=0\). Aleshores \(n=m+k=m+0=m\). L'antisimetria és, doncs, conseqüència directa de la cancel·lació i del fet que \(k+l=0\Rightarrow k=l=0\).
\end{solucio}

\begin{exercici}
\textbf{\(\N\) és ben ordenat.} Expliqueu l'estratègia de la demostració que tota part no inicial de \(\N\) té element mínim. (Model del teorema~\ref{teo:bon-ordre-N}.)
\end{exercici}
\begin{solucio}
La unicitat del mínim surt de l'antisimetria. Per a l'existència es prova el contrarecíproc: si una part \(P\hookrightarrow\N\) no té mínim, aleshores \(P\) és inicial. S'introdueix la part auxiliar \(R\hookrightarrow\N\) classificada per \enquote{cap natural \(\le n\) no pertany a \(P\)}, que existeix gràcies als quantificadors interns i les operacions booleanes sobre \(P(\N)\) (capítol~\ref{cap:parts-interna}), i perquè \(\le\) és una part (definició~\ref{def:ordre-naturals}). Per inducció (teorema~\ref{teo:induccio}) es veu \(R=\N\): al cas base, \(0\notin P\) (perquè si hi fos seria el mínim, ja que \(0\le k\) sempre); al pas, si cap \(k\le n\) és a \(P\), aleshores \(s\circ n\notin P\) (altrament \(s\circ n\) en seria el mínim). Com que \(R=\N\), prenent \(k=n\) resulta que cap natural pertany a \(P\), és a dir \(P\) és inicial. Aquest bon ordre és equivalent, en la matemàtica clàssica, al principi d'inducció (observació~\ref{obs:bon-ordre-classica}).
\end{solucio}

\begin{exercici}
\textbf{Per què sobre punts.} Siguin \(m,n\colon\mathbf2\to\N\) els elements generalitzats \(m=[0,s\circ0]\) i \(n=[s\circ0,0]\) (sobre el punt \(0\) de \(\mathbf2\), \(m\) val \(0\) i \(n\) val \(1\); sobre el punt \(1\), al revés). Comproveu que ni \(m\le n\) ni \(n\le m\) com a elements generalitzats, i expliqueu per què això no contradiu la totalitat de l'ordre.
\end{exercici}
\begin{solucio}
\(m\le n\) vol dir que \(\langle m,n\rangle\colon\mathbf2\to\N\times\N\) es factoritza per \(\operatorname{Leq}\); component amb el punt \(1\) de \(\mathbf2\) s'obtindria \(\langle1,0\rangle\in\operatorname{Leq}\), és a dir, \(1\le0\), fals. Simètricament, \(n\le m\) fallaria sobre el punt \(0\). L'element generalitzat \(\langle m,n\rangle\) pertany, però, a la unió \(\operatorname{Leq}\cup\operatorname{Inv}(\operatorname{Leq})\): la totalitat diu que aquesta unió és la part total de \(\N\times\N\), no que cada element generalitzat caigui sencer en un dels dos membres. És el mateix fenomen que la unió interna del capítol~\ref{cap:parts-interna}: sobre punts, \enquote{o} és \enquote{o}; sobre un estadi com \(\mathbf2\), un element es pot repartir.
\end{solucio}

\section{Recursió amb paràmetres}\label{prac:recursio-parametres}

\begin{exercici}
\textbf{Recursió amb paràmetres via exponencials.} Siguin \(a\colon X\to A\) i \(t\colon X\times A\to A\). Demostreu que existeix una única \(u\colon X\times\N\to A\) amb \(u\circ\langle x,0\rangle=a\circ x\) i \(u\circ\langle x,s\circ n\rangle=t\circ\langle x,u\circ\langle x,n\rangle\rangle\). (Reproduïu el teorema~\ref{teo:recursio-parametres}.)
\end{exercici}
\begin{solucio}
Transposem respecte de \(X\). Pel capítol~\ref{cap:exponencials}, una aplicació \(u\colon X\times\N\to A\) equival a la seva transposada \(\widehat u\colon\N\to A^{X}\), i les dues condicions sobre \(u\) es tradueixen, per la unicitat de la transposada, en condicions sobre \(\widehat u\). El valor inicial \(a\) té nom \(\ulcorner a\urcorner\colon\mathbf1\to A^{X}\), i la regla \(t\) determina una autoaplicació \(T\colon A^{X}\to A^{X}\) (el nom intern de \(g\mapsto(x\mapsto t(x,g(x)))\)). Aplicant l'axioma dels naturals (axioma~\ref{ax:naturals}) al sistema recursiu \((A^{X},\ulcorner a\urcorner,T)\) s'obté una única \(\widehat u\colon\N\to A^{X}\) amb \(\widehat u\circ0=\ulcorner a\urcorner\) i \(\widehat u\circ s=T\circ\widehat u\). La seva descurrificada \(u\) és l'aplicació buscada, única perquè ho és \(\widehat u\). Aquesta és la raó tècnica per introduir els exponencials abans que \(\N\): permeten reduir la recursió amb paràmetres a la recursió amb dades globals.
\end{solucio}

\chapter{Conjunts finits i cardinalitat finita}\label{cap:practica-finits}

Aquest capítol acompanya el capítol homònim de la part teòrica (capítol~\ref{cap:finits} de la Part~\ref{part:teoria}). Els exercicis resolts en segueixen l'ordre d'exposició --- els ordinals finits i els segments interns, l'equipotència i el conjunt de cardinals, la comparació de segments i la unicitat del cardinal, el principi dels calaixos, les parts d'un conjunt finit, i els recomptes de parts i d'aplicacions--- i serveixen de model per als exercicis proposats de la Part~\ref{part:exercicis}. Cada bloc fixa una tècnica: descompondre un segment successor com \(\seg{n}+\mathbf1\), raonar per inducció sobre el cardinal, i traduir sumes/productes/exponencials de conjunts finits en operacions aritmètiques. Convé intentar cada enunciat abans de llegir-ne la solució.

\section{Segments inicials}\label{prac:segments}

\begin{exercici}
\textbf{Descomposició successora del segment.} Demostreu que \(\seg{0}\cong\mathbf0\) i que, per a tot \(n\colon\mathbf1\to\N\), \(\seg{s\circ n}\cong\seg{n}+\mathbf1\), amb \(\seg{n}\) la part dels \(k<n\) i \(\mathbf1\) el punt \(k=n\). (Reproduïu la proposició~\ref{prop:segment-successor}.)
\end{exercici}
\begin{solucio}
\(\seg{0}=\{k:k<0\}\) és inicial, perquè \(k<0\) no es compleix per a cap \(k\): de \(k+s\circ j=0\) sortiria \(s\circ(k+j)=0\), i el zero no és cap successor (teorema~\ref{teo:peano}(2)). Per al successor, i raonant sobre punts com a tota la secció de l'ordre (observació~\ref{obs:ordre-sobre-punts}), un punt \(k\in\seg{s\circ n}\) compleix \(k<s\circ n\), que equival a \(k\le n\), és a dir \(k<n\) o bé \(k=n\), i no totes dues alhora perquè \(n\not<n\) (pas del successor de l'ordre, proposició~\ref{prop:ordre-total}). Aquesta disjunció exclusiva descompon la part \(\seg{s\circ n}\hookrightarrow\N\) en dues parts complementàries: la dels \(k<n\), que és \(\seg{n}\), i la del punt únic \(k=n\), una còpia de \(\mathbf1\). Per la caracterització de la suma disjunta com a parell de parts complementàries (definició~\ref{def:suma-disjunta}), \(\seg{s\circ n}\cong\seg{n}+\mathbf1\).
\end{solucio}

\begin{exercici}
\textbf{Els primers ordinals.} Demostreu que \([1]\cong\mathbf1\) i \([2]\cong\mathbf2\). (Model de la proposició~\ref{prop:ordinals-petits}.)
\end{exercici}
\begin{solucio}
Per la definició recursiva metateòrica dels ordinals (definició~\ref{def:ordinal-finit}), \([1]=[0]+\mathbf1=\mathbf0+\mathbf1\). Per la neutralitat del conjunt inicial respecte de la suma (capítol~\ref{cap:valors}), \(\mathbf0+\mathbf1\cong\mathbf1\); així \([1]\cong\mathbf1\). Aplicant-ho un cop més, \([2]=[1]+\mathbf1\cong\mathbf1+\mathbf1\), que és per definició el conjunt de valors de veritat \(\mathbf2\). Aquestes igualtats són un esquema sobre numerals estàndard, comprovable externament per a cada instància.
\end{solucio}

\section{Equipotència i cardinal}\label{prac:equipotencia}

\begin{exercici}
\textbf{Els punts de \(\operatorname{Isos}(A,B)\).} Demostreu que un punt global \(\mathbf1\to\operatorname{Isos}(A,B)\) és exactament un isomorfisme \(A\to B\), i que \(\operatorname{Isos}(A,B)\) és no inicial si i només si \(A\cong B\). (Reproduïu l'observació~\ref{obs:iso-punts}.)
\end{exercici}
\begin{solucio}
\(\operatorname{Isos}(A,B)\) és la imatge, per la projecció a \(B^{A}\), de la part \(P\hookrightarrow B^{A}\times A^{B}\) dels parells \((f,g)\) amb \(g\circ f=\id_A\) i \(f\circ g=\id_B\) (definició~\ref{def:iso-exponencial}). Pels punts globals de l'exponencial (teorema~\ref{teo:punts-exponencial}), un punt \(\mathbf1\to B^{A}\) és una aplicació \(f\colon A\to B\); pertany a la imatge exactament quan existeix \(g\) amb \(g\circ f=\id_A\) i \(f\circ g=\id_B\), és a dir quan \(f\) és un isomorfisme (definició~\ref{def:isomorfisme}). Per tant \(\operatorname{Isos}(A,B)\) és no inicial si i només si existeix un isomorfisme \(A\cong B\), és a dir \(A\approx B\) (definició~\ref{def:equipotencia}). La part es defineix amb un únic bloc d'operacions internes, sense triar cap isomorfisme.
\end{solucio}

\begin{exercici}
\textbf{\(\operatorname{Card}(A)\) és una part de \(\N\).} Expliqueu per què \(\operatorname{Card}(A)\) és una part interna ben definida de \(\N\), i que \(n\in\operatorname{Card}(A)\iff\seg{n}\approx A\). (Reproduïu la proposició~\ref{prop:card-interna}.)
\end{exercici}
\begin{solucio}
La part \(\mathcal E_A\hookrightarrow\N\times(A+\mathbf1)^{\N}\times\N^{A}\) és interna (definició~\ref{def:card-A}): les seves condicions (i)--(iii) són conjuncions d'equacions entre avaluacions d'exponencials \emph{fixos} i de condicions acotades per l'ordre \(<\), totes tallables amb equalitzadors, imatges inverses i quantificadors acotats ja disponibles (capítols~\ref{cap:exponencials} i~\ref{cap:parts-interna}). La imatge d'una aplicació existeix sempre (teorema~\ref{teo:factoritzacio}), de manera que \(\operatorname{Card}(A)=\operatorname{Im}(\pi)\) és una part de \(\N\). Un punt \(n\) hi pertany si i només si la fibra \((\mathcal E_A)_n\) és no inicial; i un punt d'aquesta fibra és un parell \((g,h)\) que, per (i)--(ii), exhibeix \(g|_{\seg{n}}\) i \(h\) com a inverses mútues, és a dir un isomorfisme \(\seg{n}\to A\). Recíprocament, un isomorfisme \(\varphi\colon\seg{n}\to A\) dona el triple \((n,\bar\varphi,\varphi^{-1})\). Així \(n\in\operatorname{Card}(A)\iff\seg{n}\approx A\), sense cap axioma d'elecció ni noció prèvia de finitud.
\end{solucio}

\section{Comparació de segments i unicitat del cardinal}\label{prac:comparacio}

\begin{exercici}
\textbf{Isomorfisme de segments implica igualtat.} A partir del fet que un monomorfisme \(\seg{m}\hookrightarrow\seg{n}\) implica \(m\le n\), demostreu que un isomorfisme \(\seg{m}\cong\seg{n}\) implica \(m=n\). (Part del lema~\ref{lema:comparacio-ordinals}.)
\end{exercici}
\begin{solucio}
Un isomorfisme \(\seg{m}\cong\seg{n}\) dona en particular monomorfismes en tots dos sentits: \(\seg{m}\hookrightarrow\seg{n}\) i \(\seg{n}\hookrightarrow\seg{m}\). Per la primera part del lema~\ref{lema:comparacio-ordinals}, el primer dona \(m\le n\) i el segon \(n\le m\). Per l'antisimetria de l'ordre (proposició~\ref{prop:ordre-total}), \(m=n\). Aquesta és la peça que fa que el cardinal estigui ben definit: no hi pot haver dos numerals diferents amb segments isomorfs.
\end{solucio}

\begin{exercici}
\textbf{Unicitat del cardinal finit.} Demostreu que si \(A\) és finit, la part \(\operatorname{Card}(A)\) té un únic punt, de manera que \(|A|\) està ben definit, i que dos conjunts finits són isomorfs si i només si tenen el mateix cardinal. (Reproduïu el teorema~\ref{teo:unicitat-cardinal}.)
\end{exercici}
\begin{solucio}
Si \(\seg{m}\approx A\) i \(\seg{n}\approx A\), per transitivitat de l'equipotència (definició~\ref{def:equipotencia}) \(\seg{m}\approx\seg{n}\), és a dir hi ha un isomorfisme \(\seg{m}\cong\seg{n}\); pel lema~\ref{lema:comparacio-ordinals}, \(m=n\). Així \(\operatorname{Card}(A)\) té com a molt un punt, i \(|A|\) queda ben definit. Per a la darrera afirmació: si \(A\cong B\), aleshores \(\seg{|A|}\approx A\cong B\approx\seg{|B|}\) dona \(|A|=|B|\); recíprocament, \(|A|=|B|\) dona \(A\approx\seg{|A|}=\seg{|B|}\approx B\), és a dir \(A\cong B\).
\end{solucio}

\section{El principi dels calaixos}\label{prac:calaixos}

\begin{exercici}
\textbf{Principi dels calaixos.} Demostreu que si \(A\) i \(B\) són finits amb \(|A|>|B|\), no hi ha cap monomorfisme \(A\hookrightarrow B\). (Reproduïu el teorema~\ref{teo:calaixos}.)
\end{exercici}
\begin{solucio}
Un monomorfisme \(A\hookrightarrow B\), compost amb els isomorfismes \(\seg{|A|}\approx A\) i \(B\approx\seg{|B|}\) (definició~\ref{def:conjunt-finit}), donaria un monomorfisme \(\seg{|A|}\hookrightarrow\seg{|B|}\). Pel lema~\ref{lema:comparacio-ordinals}, això forçaria \(|A|\le|B|\), contra la hipòtesi \(|A|>|B|\). Per tant no existeix cap monomorfisme \(A\hookrightarrow B\); dit altrament, tota aplicació d'un finit cap a un de cardinal estrictament menor identifica almenys dos elements.
\end{solucio}

\begin{exercici}
\textbf{Injectiu \(=\) exhaustiu \(=\) bijectiu en el cas finit.} Sigui \(f\colon A\to B\) entre finits amb \(|A|=|B|\). Demostreu que si \(f\) és un monomorfisme, aleshores \(f\) és un isomorfisme. (Part del teorema~\ref{teo:endo-finit}.)
\end{exercici}
\begin{solucio}
Per la factorització imatge (teorema~\ref{teo:factoritzacio}), \(f=m\circ e\) amb \(e\) epimorfisme regular i \(m\) monomorfisme. Com que \(f\) és mono, \(e\) és un isomorfisme (proposició~\ref{prop:mono-nucli-trivial}); així \(\operatorname{Im}(f)\cong A\) és finit amb \(|\operatorname{Im}(f)|=|A|=|B|\). Si \(m\) no fos exhaustiu, el seu complement \(C\hookrightarrow B\) seria una part no inicial, finita (lema~\ref{lema:part-finita}), i \(B\cong\operatorname{Im}(f)+C\); pel recompte de la suma (lema~\ref{lema:recompte-suma}) i \(C\) no inicial, \(|B|=|\operatorname{Im}(f)|+|C|>|\operatorname{Im}(f)|=|B|\), contradicció. Per tant \(m\) és un monomorfisme exhaustiu, és a dir un isomorfisme (proposició~\ref{prop:epi-imatge-total}, corol·lari~\ref{cor:bimorfisme-iso}); i \(f=m\circ e\) és composició de dos isomorfismes, un isomorfisme.
\end{solucio}

\begin{exercici}
\textbf{Una secció sense elecció.} Sigui \(f\colon\seg{n}\to B\) un epimorfisme. Construïu una secció de \(f\) sense l'axioma d'elecció, i deduïu-ne que la direcció (2)\(\Rightarrow\)(3) del teorema~\ref{teo:endo-finit} no necessita l'elecció.
\end{exercici}
\begin{solucio}
Per a cada punt \(b\) de \(B\), la fibra \(f^{-1}(b)\hookrightarrow\seg{n}\hookrightarrow\N\) és una part de \(\N\) no buida, perquè \(f\) és exhaustiu sobre punts (proposició~\ref{prop:epi-exhaustiu-punts}); pel bon ordre de \(\N\) (teorema~\ref{teo:bon-ordre-N}) té un mínim, únic. Això defineix \(s\) sobre punts: \(s(b):=\min f^{-1}(b)\). Per fer-ne una aplicació cal construir-la internament: la relació \(R\hookrightarrow B\times\N\) formada pels parells \((b,k)\) amb \(k<n\), \(f(k)=b\) i \(\forall j<k\,(f(j)\ne b)\) és una part (quantificador acotat, igualtats i l'ordre, com a la definició~\ref{def:card-A}); és funcional, perquè el mínim és únic, i total, perquè tota fibra en té un; per la proposició~\ref{prop:aplicacio-funcional}, \(R\) és el graf d'una aplicació \(s\colon B\to\seg{n}\), i \(f\circ s=\id_B\) perquè coincideixen sobre punts. Amb aquesta secció, la prova del teorema~\ref{teo:endo-finit} continua igual. L'elecció, doncs, no cal per als conjunts finits: el bon ordre fa la tria per nosaltres, com ja anunciava l'observació~\ref{obs:bon-ordre-classica}.
\end{solucio}

\section{Parts d'un conjunt finit}\label{prac:parts-finites}

\begin{exercici}
\textbf{Tota part d'un finit és finita.} Sigui \(S\hookrightarrow A\) una part d'un conjunt finit. Demostreu que \(S\) és finita i \(|S|\le|A|\). (Reproduïu el lema~\ref{lema:part-finita}.)
\end{exercici}
\begin{solucio}
Com que \(A\approx\seg{n}\) amb \(n=|A|\), n'hi ha prou de provar-ho per a \(A=\seg{n}\), per inducció sobre \(n\). Per a \(n=0\), \(\seg{0}\cong\mathbf0\) i l'única part és la inicial, finita de cardinal \(0\). Pas de \(n\) a \(s\circ n\): per la descomposició successora, \(\seg{s\circ n}\cong\seg{n}+\mathbf1\) (proposició~\ref{prop:segment-successor}). La imatge inversa de \(S\) per les dues injeccions la descompon en parts complementàries \(S_0\hookrightarrow\seg{n}\) i \(S_1\hookrightarrow\mathbf1\), amb \(S\cong S_0+S_1\). Per la hipòtesi d'inducció \(S_0\) és finita amb \(|S_0|\le n\); i \(S_1\) és inicial o terminal, en tot cas finita amb \(|S_1|\le1\). Pel recompte de la suma (lema~\ref{lema:recompte-suma}), \(S\) és finita i \(|S|=|S_0|+|S_1|\le n+1=|\seg{s\circ n}|\).
\end{solucio}

\begin{exercici}
\textbf{Nombre de parts d'un conjunt finit.} Demostreu que si \(|A|=n\), aleshores \(|\mathcal P(A)|=2^{n}\). (Reproduïu el teorema~\ref{teo:nombre-parts}.)
\end{exercici}
\begin{solucio}
Com que \(\mathcal P(A)\cong\mathbf2^{A}\) (definició~\ref{def:conjunt-parts}) i \(A\approx\seg{n}\), n'hi ha prou de provar \(|\mathbf2^{\seg{n}}|=2^{n}\) per inducció sobre \(n\). Per a \(n=0\), \(\mathbf2^{\seg{0}}\cong\mathbf2^{\mathbf0}\cong\mathbf1\), i \(2^{0}=1\). Pas de \(n\) a \(s\circ n\): de \(\seg{s\circ n}\cong\seg{n}+\mathbf1\) i de la conversió de l'exponencial d'una suma en producte (proposició~\ref{prop:exponencial-suma}),
\[
\mathbf2^{\seg{s\circ n}}\cong\mathbf2^{\seg{n}+\mathbf1}\cong\mathbf2^{\seg{n}}\times\mathbf2^{\mathbf1}\cong\mathbf2^{\seg{n}}\times\mathbf2.
\]
Per la hipòtesi d'inducció \(\mathbf2^{\seg{n}}\) té cardinal \(2^{n}\), i multiplicar per \(\mathbf2\) duplica el cardinal (lema~\ref{lema:recompte-producte}); així el cardinal és \(2^{n}+2^{n}=2\cdot2^{n}=2^{s\circ n}\), per la definició recursiva de l'exponenciació (definició~\ref{def:producte-exp}).
\end{solucio}

\section{Recompte d'aplicacions}\label{prac:recompte-aplicacions}

\begin{exercici}
\textbf{Nombre d'aplicacions entre finits.} Demostreu que si \(|A|=n\) i \(|B|=m\), aleshores \(|B^{A}|=m^{n}\). (Reproduïu el teorema~\ref{teo:nombre-aplicacions}.)
\end{exercici}
\begin{solucio}
Fixats \(A\approx\seg{n}\) i \(B\approx\seg{m}\), l'exponencial només depèn dels conjunts llevat d'isomorfisme, de manera que n'hi ha prou de provar \(|\seg{m}^{\seg{n}}|=m^{n}\) per inducció sobre \(n\), uniformement en \(m\). Per a \(n=0\), \(\seg{m}^{\seg{0}}\cong\seg{m}^{\mathbf0}\cong\mathbf1\), i \(m^{0}=1\). Pas de \(n\) a \(s\circ n\): de \(\seg{s\circ n}\cong\seg{n}+\mathbf1\) i de la conversió de l'exponencial d'una suma (proposició~\ref{prop:exponencial-suma}),
\[
\seg{m}^{\seg{s\circ n}}\cong\seg{m}^{\seg{n}+\mathbf1}\cong\seg{m}^{\seg{n}}\times\seg{m}^{\mathbf1}\cong\seg{m}^{\seg{n}}\times\seg{m}.
\]
Per la hipòtesi d'inducció \(\seg{m}^{\seg{n}}\) té cardinal \(m^{n}\), i pel recompte del producte (lema~\ref{lema:recompte-producte}) multiplicar per \(\seg{m}\) multiplica el cardinal per \(m\); així \(|\seg{m}^{\seg{s\circ n}}|=m^{n}\cdot m=m^{s\circ n}\), la definició recursiva de l'exponenciació. Aquest recompte justifica la notació \(|B^{A}|=m^{n}\) i l'ús de \(P(A)=\mathbf2^{A}\) com a \enquote{conjunt de \(2^{n}\) parts}.
\end{solucio}

\chapter{Infinitud, enumerabilitat i seqüències}\label{cap:practica-infinitud}

Aquest capítol acompanya el capítol homònim de la part teòrica (capítol~\ref{cap:infinitud} de la Part~\ref{part:teoria}). Els exercicis resolts en segueixen l'ordre d'exposició --- la infinitud i la infinitud de Dedekind, els conjunts enumerables i numerables, els subobjectes de \(\N\), les llistes finites, i les unions numerables--- i serveixen de model per als exercicis proposats de la Part~\ref{part:exercicis}. Cada bloc fixa una tècnica: exhibir un monomorfisme no exhaustiu, compondre enumeracions sense elecció, usar el principi del mínim en lloc d'una tria, i identificar exactament on entra l'elecció numerable. Convé intentar cada enunciat abans de llegir-ne la solució.

\section{Infinitud i infinitud de Dedekind}\label{prac:infinits}

\begin{exercici}
\textbf{\(\N\) és infinit de Dedekind.} Demostreu que el successor \(s\colon\N\to\N\) és un monomorfisme no exhaustiu, i deduïu que \(\N\) és infinit. (Reproduïu la proposició~\ref{prop:N-dedekind}.)
\end{exercici}
\begin{solucio}
Que \(s\) és un monomorfisme és el teorema~\ref{teo:peano}(1): de \(s\circ m=s\circ n\) se segueix \(m=n\). Que no és exhaustiu és el teorema~\ref{teo:peano}(2): \(0\) no és de la forma \(s\circ n\) per a cap \(n\), de manera que el punt \(0\colon\mathbf1\to\N\) no factoritza per \(s\), i la imatge de \(s\) és la part pròpia complementària del zero. Així \(s\) exhibeix \(\N\) com a infinit de Dedekind (definició~\ref{def:dedekind-infinit}). Que \(\N\) és infinit se'n segueix: si fos finit, \(s\) hauria de ser exhaustiu (teorema~\ref{teo:endo-finit}), contra el que acabem de veure.
\end{solucio}

\begin{exercici}
\textbf{Dedekind implica infinit.} Demostreu que tot conjunt infinit de Dedekind és infinit, sense cap axioma nou. (Reproduïu la proposició~\ref{prop:dedekind-implica-infinit}.)
\end{exercici}
\begin{solucio}
És el contrarecíproc del cas finit. Si un conjunt \(A\) fos alhora finit i infinit de Dedekind, tindria un monomorfisme \(A\hookrightarrow A\) no exhaustiu; però pel teorema~\ref{teo:endo-finit}, tot monomorfisme d'un conjunt finit en si mateix és exhaustiu, una contradicció. Per tant un conjunt infinit de Dedekind no pot ser finit, és a dir, és infinit. La direcció oposada ---que tot infinit és infinit de Dedekind--- sí que requereix l'elecció numerable (teorema~\ref{teo:infinit-implica-dedekind}), perquè cal triar, alhora i per a cada llargada, una llista injectiva dins \(A\); un cop triades, la còpia de \(\N\) es construeix sense més tries.
\end{solucio}

\section{Enumerables i numerables}\label{prac:enumerables}

\begin{exercici}
\textbf{Imatge i part d'un enumerable.} Sigui \(A\) enumerable. Demostreu que tota imatge exhaustiva \(g\colon A\twoheadrightarrow C\) és enumerable i que tota part de \(A\) és enumerable, sense elecció. (Part de la proposició~\ref{prop:estabilitat-enumerable}.)
\end{exercici}
\begin{solucio}
Sigui \(e\colon\N\twoheadrightarrow A\) una enumeració. \emph{Imatge:} si \(g\) és exhaustiva, la composició \(g\circ e\colon\N\twoheadrightarrow C\) també ho és (l'exhaustivitat es preserva per composició, capítol~\ref{cap:imatges}), i és una enumeració de \(C\). \emph{Part:} sigui \(m\colon P\hookrightarrow A\) una part no inicial i \(p_0\colon\mathbf1\to P\) un punt (existeix perquè \(P\) no és inicial). Definim \(r\colon A\to P\) enviant els elements de \(P\) en si mateixos i la resta a \(p_0\); és una retracció de \(m\), construïda amb la decisió del complement (capítol~\ref{cap:valors}) i sense cap tria infinita. Aleshores \(r\circ e\colon\N\twoheadrightarrow P\) és una enumeració. Els casos inicials són immediats. Cap pas no usa elecció: només es fa servir una única enumeració donada.
\end{solucio}

\begin{exercici}
\textbf{Suma i producte de dos enumerables.} Demostreu que si \(A\) i \(B\) són enumerables, aleshores \(A+B\) i \(A\times B\) també ho són. (Part de la proposició~\ref{prop:estabilitat-enumerable}.)
\end{exercici}
\begin{solucio}
Siguin \(e\colon\N\twoheadrightarrow A\) i \(f\colon\N\twoheadrightarrow B\) enumeracions. \emph{Suma:} d'un isomorfisme \(\N\cong\N+\N\) (per exemple la partició en parells i senars, obtinguda de l'aritmètica del capítol~\ref{cap:naturals}) s'obté \(\N\cong\N+\N\xrightarrow{e+f}A+B\), exhaustiva. \emph{Producte:} n'hi ha prou amb una bijecció \(\N\cong\N\times\N\) (un aparellament), perquè aleshores \(\N\cong\N\times\N\xrightarrow{e\times f}A\times B\) és exhaustiva. Un tal aparellament es construeix explícitament per recursió, per exemple amb la diagonalització \((i,j)\mapsto\tfrac12(i+j)(i+j+1)+j\), la inversa de la qual també es defineix per recursió; és una aplicació concreta, no una tria. Cap dels dos usa elecció.
\end{solucio}

\begin{exercici}
\textbf{L'aparellament per diagonals.} Sigui \(\nu\colon\N\times\N\to\N\times\N\) definida per \(\nu(0,j)=(s\circ j,0)\) i \(\nu(s\circ i',j)=(i',s\circ j)\), i \(u\colon\N\to\N\times\N\) la recursora de \((\N\times\N,\langle0,0\rangle,\nu)\). Demostreu que \(u\) és bijectiva. (Detalls de la proposició~\ref{prop:estabilitat-enumerable}(4).)
\end{exercici}
\begin{solucio}
\(\nu\) està ben definida per casos perquè tot punt \(i\) és zero o successor, de manera única (lema~\ref{lema:zero-o-successor}). \emph{\(\nu\) és injectiva:} \(\nu(i,j)=\nu(i',j')\) força el mateix cas (la segona component és \(0\) en el primer cas i un successor en el segon), i dins de cada cas la injectivitat de \(s\) dona \((i,j)=(i',j')\). \emph{\((0,0)\) no és a la imatge de \(\nu\):} en el primer cas la primera component és un successor, en el segon ho és la segona. \emph{\(u\) és injectiva:} com en la prova de les equacions de Peano, \(u\circ s=\nu\circ u\) i la part dels \(n\) tals que \(u(n)=u(n')\Rightarrow n=n'\) per a tot \(n'\) conté el zero (perquè \((0,0)\) no és a la imatge de \(\nu\), \(u(n')=u(0)\) força \(n'=0\)) i és tancada pel successor (per la injectivitat de \(\nu\)). \emph{\(u\) és exhaustiva:} per a cada \(d\), la part \(D_d\) dels parells \((i,j)\) amb \(i+j=d\) té \(d+1\) punts i \(\nu\) els recorre des de \((d,0)\) fins a \((0,d)\) i salta a \((d+1,0)\); per inducció sobre \(d\) tots els punts de \(D_d\) són a la imatge de \(u\), i tot parell és a algun \(D_d\). Per tant \(u\) és un bimorfisme, és a dir, un isomorfisme.
\end{solucio}

\section{Subobjectes de \texorpdfstring{\(\N\)}{N}}\label{prac:subobjectes}

\begin{exercici}
\textbf{Tota part de \(\N\) és enumerable.} Demostreu que tota part \(P\hookrightarrow\N\) és enumerable, i que és finita o isomorfa a \(\N\), sense elecció. (Reproduïu el teorema~\ref{teo:subobjectes-N}.)
\end{exercici}
\begin{solucio}
Si \(P\) és inicial, és enumerable per definició. Suposem \(P\) no inicial. La clau és el \emph{principi del mínim}: tota part no inicial de \(\N\) té un element mínim (teorema~\ref{teo:bon-ordre-N}), la qual cosa no usa elecció perquè el mínim és \emph{únic} i per tant canònic. Definim l'autoaplicació \(t\colon P\to P\) que envia \(x\) al mínim dels elements de \(P\) estrictament més grans que \(x\), si n'hi ha, i a \(x\) altrament (el seu graf és una part funcional i total de \(P\times P\), tallada amb quantificadors interns), i \(a\colon\N\to P\) com la recursora del sistema \((P,\min P,t)\): una recursió simple, sense paràmetres ni tries. Si per a algun \(n\) no en queda cap de més gran, la successió s'estabilitza i la seva imatge \(\{a_0<\dots<a_n\}\) és tot \(P\), que és finit. Altrament \(a\) és estrictament creixent, per tant un monomorfisme, i la seva imatge és tot \(P\) (cap element pot ser omès, perquè la successió creix almenys d'una unitat a cada pas); així \(a\colon\N\to P\) és una bijecció creixent i \(P\cong\N\).
\end{solucio}

\begin{exercici}
\textbf{Els enumerables infinits són \(\N\).} Demostreu que un conjunt és enumerable infinit si i només si és isomorf a \(\N\). (Reproduïu el corol·lari~\ref{cor:enumerable-infinit-N}.)
\end{exercici}
\begin{solucio}
Si \(A\cong\N\), és enumerable (per la identitat) i infinit (proposició~\ref{prop:N-dedekind}). Recíprocament, sigui \(A\) enumerable infinit amb enumeració \(e\colon\N\twoheadrightarrow A\). Construïm un monomorfisme \(A\hookrightarrow\N\) sense elecció: definim \(r\colon A\to\N\) enviant cada \(a\) al \emph{mínim} índex \(n\) amb \(e_n=a\); aquest mínim existeix i és únic pel principi del mínim, de manera que \(r\) està ben definida i és una secció de \(e\), en particular un monomorfisme. La seva imatge és una part de \(\N\), que pel teorema~\ref{teo:subobjectes-N} és finita o isomorfa a \(\N\); com que \(A\) és infinit, és el segon cas, i \(A\cong\N\). En conseqüència, tots els numerables infinits són isomorfs entre si, amb \(\N\) com a representant. El mínim substitueix aquí l'elecció (observació~\ref{obs:minim-vs-eleccio}).
\end{solucio}

\section{Llistes finites}\label{prac:llistes}

\begin{exercici}
\textbf{Recursió sobre llistes.} Demostreu que \((\operatorname{List}(A),\varepsilon,\operatorname{cons})\) satisfà el principi de recursió: per a tot \((B,b,t)\) amb \(b\colon\mathbf1\to B\) i \(t\colon A\times B\to B\), existeix una única \(h\colon\operatorname{List}(A)\to B\) amb \(h\circ\varepsilon=b\) i \(h\circ\operatorname{cons}=t\circ(\id_A\times h)\). (Model de la proposició~\ref{prop:recursio-llistes}.)
\end{exercici}
\begin{solucio}
\emph{Existència.} Sigui \(L=\operatorname{List}(A)\). En lloc d'una família d'aplicacions amb dominis variables, construïm per recursió simple una successió de punts de l'exponencial fix \(B^{L}\). Sigui \(\Phi\colon B^{L}\to B^{L}\) la transposada de l'aplicació \(B^{L}\times L\to B\) definida per casos sobre \(L\cong\mathbf1+A\times L\): \((h,\varepsilon)\mapsto b\) i \((h,\operatorname{cons}(a,u))\mapsto t(a,\operatorname{ev}(h,u))\). La recursora \(H\colon\N\to B^{L}\) del sistema \((B^{L},\ulcorner b\circ{!}\urcorner,\Phi)\) compleix \(H_0=b\circ{!}\) i \(H_{m+1}=\Phi(H_m)\), i per inducció sobre \(m\) les \(H_m\) coincideixen sobre les llistes de llargada \(\le m\) amb totes les posteriors. La diagonal \(h(u):=H_{\ell(u)}(u)\), és a dir \(h=\operatorname{ev}\circ\langle H\circ\ell,\id_L\rangle\), compleix \(h(\varepsilon)=H_0(\varepsilon)=b\) i \(h(\operatorname{cons}(a,u))=H_{\ell(u)+1}(\operatorname{cons}(a,u))=t(a,H_{\ell(u)}(u))=t(a,h(u))\).

\emph{Unicitat.} Si \(h'\) compleix les mateixes equacions, aleshores \(h=h'\) sobre les llistes de llargada \(0\), i si coincideixen sobre les de llargada \(\le m\), l'equació del \(\operatorname{cons}\) força la coincidència sobre les de llargada \(m+1\); per inducció sobre \(m\) (teorema~\ref{teo:induccio}), \(h=h'\). Això estableix la inicialitat de \((\operatorname{List}(A),\varepsilon,\operatorname{cons})\) a partir de la construcció, no a l'inrevés.
\end{solucio}

\begin{exercici}
\textbf{Enumerabilitat de \(\operatorname{List}(A)\).} Demostreu que si \(A\) és enumerable, aleshores \(\operatorname{List}(A)\) també ho és, sense elecció. (Model de l'observació~\ref{obs:llistes-enumerables}.)
\end{exercici}
\begin{solucio}
Per la lectura de \(\operatorname{List}(A)\) com a suma indexada \(\coprod_{n\in\N}A^{\seg{n}}\) (observació~\ref{obs:llistes-segment}), cada sumand \(A^{\seg{n}}\) és un producte finit d'enumerables i per tant enumerable (proposició~\ref{prop:estabilitat-enumerable}). La clau és que cada \(A^{\seg{n}}\) ve amb una enumeració \emph{canònica}, derivada \emph{uniformement} de la d'\(A\) i de l'aparellament \(\N\cong\N\times\N\), no triada. Amb aquesta uniformitat, l'assignació \((n,k)\mapsto\) \enquote{\(k\)-èsim element de \(A^{\seg{n}}\)} és una aplicació exhaustiva \(\N\times\N\twoheadrightarrow\operatorname{List}(A)\), que composta amb l'aparellament dona una enumeració de \(\operatorname{List}(A)\). No cal elecció precisament perquè les enumeracions dels sumands són uniformes, no arbitràries; és el contrast amb el cas de la secció següent.
\end{solucio}

\section{Unions numerables}\label{prac:unions}

\begin{exercici}
\textbf{Unió numerable d'enumerables.} Suposant l'elecció numerable, demostreu que la unió disjunta d'una família numerable \((A_n)_{n\in\N}\) d'enumerables és enumerable. (Reproduïu el teorema~\ref{teo:unio-numerable}.)
\end{exercici}
\begin{solucio}
Per a cada \(n\), \(A_n\) és enumerable: o bé inicial, o bé existeix un epimorfisme \(\N\twoheadrightarrow A_n\). El conjunt \(E'\to\N\) que sobre cada \(n\) (amb \(A_n\) no inicial) té per fibra les enumeracions de \(A_n\) és una família de conjunts no buits indexada per \(\N\). L'axioma d'elecció numerable (axioma~\ref{ax:eleccio-numerable}) en proporciona una secció, és a dir una tria simultània d'una enumeració \(e_n\colon\N\twoheadrightarrow A_n\) per a cada \(n\). Amb aquesta tria, l'aplicació \((n,k)\mapsto e_n(k)\) de \(\N\times\N\) cap a \(\coprod_n A_n\) és exhaustiva, i component-la amb un aparellament \(\N\cong\N\times\N\) s'obté una enumeració de la unió disjunta. Tota imatge seva és enumerable per la proposició~\ref{prop:estabilitat-enumerable}(1).
\end{solucio}

\begin{exercici}
\textbf{On és l'elecció, i on no.} Expliqueu per què la unió de \emph{dos} enumerables i l'enumerabilitat de \(\operatorname{List}(A)\) no necessiten elecció, mentre que la unió numerable sí. (Model de l'observació~\ref{obs:eleccio-unio}.)
\end{exercici}
\begin{solucio}
El pas irreductible del teorema~\ref{teo:unio-numerable} és la tria \emph{simultània} de les \(e_n\). Cada \(A_n\) enumerable \emph{té} alguna enumeració, però la definició d'enumerabilitat n'afirma només l'existència; res no en selecciona cap de canònica quan els \(A_n\) són arbitraris, i encadenar infinites eleccions en una sola aplicació \(n\mapsto e_n\) és exactament el contingut de l'axioma~\ref{ax:eleccio-numerable}. En canvi, la unió de dos enumerables (proposició~\ref{prop:estabilitat-enumerable}) usa només \emph{dues} enumeracions donades, i \(\operatorname{List}(A)\) (observació~\ref{obs:llistes-enumerables}) deriva totes les enumeracions dels sumands \emph{uniformement} d'una de sola. La diferència no és la mida de la família, sinó si les enumeracions vénen donades/uniformes o s'han de triar. De fet, sense elecció és consistent que una unió numerable de conjunts de dos elements no sigui enumerable.
\end{solucio}

\chapter[Schröder--Bernstein i Cantor]{Comparació de conjunts,\texorpdfstring{\\}{ }Schröder--Bernstein i Cantor}\label{cap:practica-cantor}

Aquest capítol acompanya el capítol homònim de la part teòrica (capítol~\ref{cap:cantor} de la Part~\ref{part:teoria}). Els exercicis resolts en segueixen l'ordre d'exposició ---la comparació \(\preceq\) i les seves propietats, el teorema de Schröder--Bernstein, i l'argument diagonal amb el teorema de Cantor--- i serveixen de model per als exercicis proposats de la Part~\ref{part:exercicis}. Cada bloc fixa una tècnica: manipular monomorfismes i isomorfismes com a testimonis de mida, construir un punt fix d'un operador creixent sobre \(P(A)\) sense cap tria, i explotar la diagonal per negar l'exhaustivitat. Convé intentar cada enunciat abans de llegir-ne la solució.

\section{La comparació de conjunts}\label{prac:comparacio-conjunts}

\begin{exercici}
\textbf{La comparació és un preordre.} Demostreu que \(\preceq\) és reflexiva i transitiva, i que és invariant per equipotència: si \(A\approx A'\), \(B\approx B'\) i \(A\preceq B\), aleshores \(A'\preceq B'\). (Reproduïu la proposició~\ref{prop:preceq-preordre}.)
\end{exercici}
\begin{solucio}
\emph{Reflexivitat:} \(\id_A\) és un monomorfisme (definició~\ref{def:monomorfisme}), de manera que \(A\preceq A\). \emph{Transitivitat:} si \(m\colon A\hookrightarrow B\) i \(n\colon B\hookrightarrow C\) són monomorfismes i \((n\circ m)\circ u=(n\circ m)\circ v\), aleshores \(n\circ(m\circ u)=n\circ(m\circ v)\) dona \(m\circ u=m\circ v\) (perquè \(n\) és mono), i altre cop \(u=v\) (perquè \(m\) ho és); així \(n\circ m\) és un monomorfisme i \(A\preceq C\). \emph{Invariància:} siguin \(p\colon A'\cong A\), \(q\colon B\cong B'\) isomorfismes i \(m\colon A\hookrightarrow B\) un monomorfisme; aleshores \(q\circ m\circ p\colon A'\to B'\) és una composició de monomorfismes (els isomorfismes ho són), per tant un monomorfisme, i \(A'\preceq B'\). Cap pas no usa cap axioma més enllà dels de composició.
\end{solucio}

\begin{exercici}
\textbf{Dominacions bàsiques.} Demostreu que tota part \(S\hookrightarrow A\) compleix \(S\preceq A\), i que si \(g\colon A\twoheadrightarrow C\) és exhaustiva \emph{i escindida} (té una secció), aleshores \(C\preceq A\). Comenteu on entra, o no, l'elecció.
\end{exercici}
\begin{solucio}
\emph{Parts:} la inclusió \(S\hookrightarrow A\) és un monomorfisme per definició de part (definició~\ref{def:part}), de manera que \(S\preceq A\) sense cap hipòtesi. \emph{Exhaustives escindides:} si \(s\colon C\to A\) és una secció de \(g\), és a dir \(g\circ s=\id_C\), aleshores \(s\) és un monomorfisme (té retracció \(g\)), i \(C\preceq A\). Aquí la clau és que la secció \emph{ja} ve donada: no s'hi tria res. Si només se sap que \(g\) és exhaustiva, obtenir-ne una secció requereix l'elecció (proposició~\ref{prop:epi-regular-eleccio}); és exactament el pas que aquesta versió de l'enunciat evita en donar la secció per hipòtesi.
\end{solucio}

\section{El teorema de Schröder--Bernstein}\label{prac:sb}

\begin{exercici}
\textbf{L'operador \(\Phi\) és creixent.} Donats \(f\colon A\hookrightarrow B\) i \(g\colon B\hookrightarrow A\), sigui \(\Phi(X)=A\setminus g(B\setminus f(X))\) per a \(X\in P(A)\). Demostreu que \(X\subseteq X'\) implica \(\Phi(X)\subseteq\Phi(X')\). (Pas de la prova del teorema~\ref{teo:schroder-bernstein}.)
\end{exercici}
\begin{solucio}
Suposem \(X\subseteq X'\). La imatge directa conserva la inclusió (capítol~\ref{cap:imatges}), de manera que \(f(X)\subseteq f(X')\). El complement la inverteix (capítol~\ref{cap:valors}): \(B\setminus f(X')\subseteq B\setminus f(X)\). Tornant a aplicar la imatge directa, \(g(B\setminus f(X'))\subseteq g(B\setminus f(X))\); i tornant a prendre complements, \(A\setminus g(B\setminus f(X))\subseteq A\setminus g(B\setminus f(X'))\), és a dir \(\Phi(X)\subseteq\Phi(X')\). El nombre parell d'inversions retorna el sentit de la inclusió, de manera que \(\Phi\) és creixent.
\end{solucio}

\begin{exercici}
\textbf{El menor punt fix.} Sigui \(\Phi\colon P(A)\to P(A)\) creixent qualsevol, i \(D=\bigcap\{X : \Phi(X)\subseteq X\}\). Demostreu que \(\Phi(D)=D\) i que \(D\) és el \emph{menor} punt fix. (Nucli del teorema~\ref{teo:schroder-bernstein}, i teorema de Knaster--Tarski.)
\end{exercici}
\begin{solucio}
La col·lecció de les parts \(\Phi\)-tancades (les \(X\) amb \(\Phi(X)\subseteq X\)) és no buida perquè \(A\) hi és (\(\Phi(A)\subseteq A\) sempre). \emph{\(\Phi(D)\subseteq D\):} per a tota \(X\) \(\Phi\)-tancada, \(D\subseteq X\) i \(\Phi\) creixent donen \(\Phi(D)\subseteq\Phi(X)\subseteq X\); com que això val per a tota tal \(X\), \(\Phi(D)\subseteq\bigcap X=D\). Per tant \(D\) és \(\Phi\)-tancada. \emph{\(D\subseteq\Phi(D)\):} de \(\Phi(D)\subseteq D\) i \(\Phi\) creixent surt \(\Phi(\Phi(D))\subseteq\Phi(D)\), és a dir \(\Phi(D)\) és també \(\Phi\)-tancada; per la minimalitat de \(D\), \(D\subseteq\Phi(D)\). Així \(\Phi(D)=D\). És el menor punt fix perquè tot punt fix \(X_0=\Phi(X_0)\) és en particular \(\Phi\)-tancat, d'on \(D\subseteq X_0\). Tot això viu dins \(P(A)\) i no usa cap tria: \(D\) és un ínfim, únic i canònic.
\end{solucio}

\begin{exercici}
\textbf{Schröder--Bernstein complet.} A partir de \(\Phi(D)=D\), construïu l'isomorfisme \(A\cong B\) i comproveu que és bijectiu. (Reproduïu el teorema~\ref{teo:schroder-bernstein}.)
\end{exercici}
\begin{solucio}
De \(D=\Phi(D)=A\setminus g(B\setminus f(D))\) surt, prenent complements, \(A\setminus D=g(B\setminus f(D))\). La restricció de \(g\) a \(B\setminus f(D)\) és un monomorfisme, i tot monomorfisme és un isomorfisme sobre la seva imatge (proposició~\ref{prop:mono-nucli-trivial}), que aquí és \(A\setminus D\); escrivim \(g^{-1}\colon A\setminus D\to B\setminus f(D)\) la inversa. Definim \(h:=[f|_D,g^{-1}]\colon A\cong D+(A\setminus D)\to B\) per casos (teorema~\ref{teo:part-complementari}). Comprovem que és bimorfisme \emph{sobre punts globals}, cosa que basta per la separació global (observació~\ref{obs:injectiva-mono}, observació~\ref{obs:exhaustiva-epi}); sobre punts, les alternatives \(a\in D\) o \(a\in A\setminus D\) són excloents. \emph{Injectiva:} si \(h(a)=h(a')\) amb tots dos a \(D\), \(f\) mono dona \(a=a'\); amb tots dos fora, aplicar \(g\) dona \(a=a'\); el cas mixt és impossible perquè aleshores \(h(a)\in f(D)\) i \(h(a')\notin f(D)\). \emph{Exhaustiva:} donat un punt \(b\) de \(B\), si \(b\in f(D)\) és \(b=f(a)=h(a)\) amb \(a\in D\); si \(b\notin f(D)\), aleshores \(g(b)\in A\setminus D\) i \(h(g(b))=b\). Injectiva sobre punts és mono i exhaustiva sobre punts és epi, i un bimorfisme és isomorfisme (proposició~\ref{prop:bimorfisme-isomorfisme}), així que \(A\cong B\). Observeu que la distinció de casos no es podria fer sobre elements generalitzats, que poden quedar repartits entre \(D\) i el seu complement; és la raó de passar pels punts.
\end{solucio}

\begin{exercici}
\textbf{Antisimetria i encaix.} Deduïu de Schröder--Bernstein que si \(A\preceq B\preceq A\) aleshores \(A\approx B\), i que si \(A\subseteq B\subseteq C\) amb \(A\approx C\), aleshores \(A\approx B\approx C\) (\enquote{teorema de l'entrepà}).
\end{exercici}
\begin{solucio}
La primera part és el corol·lari~\ref{cor:preceq-antisimetric}: dues dominacions oposades donen un isomorfisme (teorema~\ref{teo:schroder-bernstein}). Per a la segona, de les inclusions surt \(A\preceq B\) i \(B\preceq C\); de \(A\approx C\) surt \(C\preceq A\). Encadenant, \(B\preceq C\approx A\preceq B\), és a dir \(B\preceq A\) i \(A\preceq B\) (aquesta darrera per \(A\subseteq B\)); per Schröder--Bernstein, \(A\approx B\). I com que \(A\approx C\), també \(B\approx C\). Cap pas no usa elecció.
\end{solucio}

\section{L'argument diagonal i el teorema de Cantor}\label{prac:cantor}

\begin{exercici}
\textbf{La diagonal.} Sigui \(f\colon A\to P(A)\). Demostreu que la part \(D=\{a\in A : a\not\mathrel{\varepsilon} f(a)\}\) no és el valor de cap element, i concloeu que \(f\) no és exhaustiva. (Reproduïu el teorema~\ref{teo:diagonal}.)
\end{exercici}
\begin{solucio}
\(D\) està ben definida com a imatge inversa per \(\langle\id_A,f\rangle\) del complement de la pertinença interna \({\in_A}\subseteq A\times P(A)\) (definició~\ref{def:pertinenca-interna}), tallada amb el complement (capítol~\ref{cap:valors}). Suposem que existís \(d\colon\mathbf1\to A\) amb \(f(d)=D\). Per definició de \(D\), \(d\in D\Leftrightarrow d\not\mathrel{\varepsilon} f(d)\); i com que \(f(d)=D\), la pertinença interna a la part constant \(D\) és la pertinença generalitzada (proposició~\ref{prop:pertinenca-interna-lectura}), és a dir \(d\mathrel{\varepsilon} f(d)\Leftrightarrow d\in D\). Substituint, \(d\in D\Leftrightarrow d\notin D\), que contradiu la distinció dels valors de veritat (teorema~\ref{teo:zero-diferent-u}). Per tant cap \(d\) no compleix \(f(d)=D\), \(D\) no és a la imatge de \(f\), i \(f\) no és exhaustiva.
\end{solucio}

\begin{exercici}
\textbf{El teorema de Cantor.} Demostreu que \(A\prec P(A)\) per a tot conjunt \(A\). (Reproduïu el teorema~\ref{teo:cantor}.)
\end{exercici}
\begin{solucio}
\emph{Dominació:} el singletó \(\{-\}\colon A\to P(A)\) és un monomorfisme (proposició~\ref{prop:singletons}): \(\{a\}=\{a'\}\) força \(a=a'\), perquè \(a\mathrel{\varepsilon}\{a\}=\{a'\}\) dona \(a=a'\). Així \(A\preceq P(A)\). \emph{Estrictesa:} si hi hagués un isomorfisme \(A\cong P(A)\), en particular seria una aplicació exhaustiva \(A\to P(A)\), que el teorema~\ref{teo:diagonal} prohibeix. Per tant \(A\not\approx P(A)\), i \(A\prec P(A)\).
\end{solucio}

\begin{exercici}
\textbf{No hi ha conjunt universal.} Demostreu que no existeix cap conjunt \(U\) en què tot conjunt s'injecti. (Reproduïu el corol·lari~\ref{cor:no-universal}.)
\end{exercici}
\begin{solucio}
Si un tal \(U\) existís, \(P(U)\) seria un conjunt i s'hauria d'injectar en \(U\), és a dir \(P(U)\preceq U\). Però sempre \(U\preceq P(U)\) pels singletons (proposició~\ref{prop:singletons}). Per Schröder--Bernstein (teorema~\ref{teo:schroder-bernstein}), \(U\approx P(U)\), contra el teorema de Cantor (teorema~\ref{teo:cantor}). La contradicció mostra que cap \(U\) no domina tots els conjunts; en particular no hi ha mida màxima, perquè \(A\prec P(A)\) per a tot \(A\).
\end{solucio}

\begin{exercici}
\textbf{\(\mathbf2^{\N}\) no és enumerable.} Useu la diagonal per demostrar que no hi ha cap aplicació exhaustiva \(\N\to\mathbf2^{\N}\); concloeu que \(\mathbf2^{\N}\), i per tant \(P(\N)\), no és enumerable. (Model de l'observació~\ref{obs:cantor-fort}.)
\end{exercici}
\begin{solucio}
Com que \(P(\N)\approx\mathbf2^{\N}\) (el conjunt de parts és l'exponencial de base \(\mathbf2\), capítol~\ref{cap:exponencials}), una exhaustiva \(\N\twoheadrightarrow\mathbf2^{\N}\) donaria, composta amb l'isomorfisme, una exhaustiva \(\N\to P(\N)\), que el teorema~\ref{teo:diagonal} prohibeix. Explícitament, donada \(f\colon\N\to\mathbf2^{\N}\), la successió diagonal \(d(n)=1-f(n)(n)\) difereix de cada \(f(n)\) en la posició \(n\), de manera que no és cap \(f(n)\); \(f\) no és exhaustiva. Un conjunt enumerable admetria una tal exhaustiva des de \(\N\) (definició~\ref{def:enumerable}); com que \(\mathbf2^{\N}\) no és inicial i no n'admet cap, no és enumerable. És l'argument que separa el numerable del continu.
\end{solucio}

\chapter{Famílies, fibres i operacions indexades}\label{cap:practica-families}

Aquest capítol acompanya el capítol homònim de la part teòrica (capítol~\ref{cap:families} de la Part~\ref{part:teoria}). Els exercicis resolts en segueixen l'ordre d'exposició ---famílies i fibres, seccions i elecció, suma i producte, famílies de parts i operacions indexades, i la distinció regla externa contra família interna--- i serveixen de model per als exercicis proposats de la Part~\ref{part:exercicis}. Cada bloc fixa una tècnica: llegir una fibra com un pullback, localitzar exactament on entra l'elecció, formar unions i interseccions indexades amb els quantificadors interns, i reconèixer quan una notació \(i\mapsto X_i\) encara no representa cap objecte. Convé intentar cada enunciat abans de llegir-ne la solució.

\section{Famílies i fibres}\label{prac:families-fibres}

\begin{exercici}
\textbf{La fibra és una imatge inversa.} Sigui \(p\colon E\to I\) una família i \(i\colon\mathbf1\to I\) un índex. Demostreu que la fibra \(E_i=\mathbf1\times_I E\) és (isomorfa a) la imatge inversa \(p^{-1}(i)\hookrightarrow E\), i que la projecció \(E_i\to E\) és un monomorfisme. (Reproduïu l'observació~\ref{obs:fibra-part}.)
\end{exercici}
\begin{solucio}
Per definició, \(E_i\) és el pullback de \(i\colon\mathbf1\to I\) i \(p\colon E\to I\) (definició~\ref{def:pullback}). La imatge inversa \(p^{-1}(i)\) és, precisament, el pullback de \(p\) al llarg de la part \(i\colon\mathbf1\hookrightarrow I\) (definició~\ref{def:imatge-inversa}); com que \(i\) és un punt, tots dos pullbacks coincideixen, d'on \(E_i\cong p^{-1}(i)\). La projecció \(E_i\to E\) és el canvi de base de \(i\) al llarg de \(p\); com que \(i\) és un monomorfisme (tot punt ho és) i els monomorfismes es preserven per canvi de base (proposició~\ref{prop:pullback-mono}), la projecció \(E_i\to E\) és un monomorfisme. En la lectura clàssica (observació~\ref{obs:pullback-exemples}), els seus punts són els \(e\in E\) amb \(p(e)=i\).
\end{solucio}

\begin{exercici}
\textbf{La família constant.} Descriviu la família constant de valor \(A\) indexada per \(I\) i comproveu que totes les seves fibres són isomorfes a \(A\). (Model de la definició~\ref{def:familia-constant}.)
\end{exercici}
\begin{solucio}
La família constant és la projecció \(\pi_I\colon I\times A\to I\). La seva fibra sobre \(i\) és el pullback \(\mathbf1\times_I(I\times A)\); prenent \(C=I\), \(f=i\colon\mathbf1\to I\) i \(g=\pi_I\), el pullback és \(\{i\}\times A\), isomorf a \(A\) (observació~\ref{obs:pullback-exemples}, cas de la fibra). Per tant \(\pi_I\) representa la família \((A)_{i\in I}\) amb totes les fibres iguals a \(A\). El seu total \(I\times A\) és el producte cartesià, llegit com \enquote{\(I\) còpies de \(A\)}.
\end{solucio}

\section{Seccions i elecció}\label{prac:seccions}

\begin{exercici}
\textbf{Secció = funció d'elecció.} Demostreu que una secció \(s\colon I\to E\) de \(p\colon E\to I\) tria un element de cada fibra, i que si \(p\) és un epimorfisme regular, l'existència d'una secció és exactament l'axioma d'elecció. (Reproduïu l'observació~\ref{obs:seccio-eleccio}.)
\end{exercici}
\begin{solucio}
Si \(p\circ s=\id_I\), aleshores per a cada índex \(i\) es té \(p(s(i))=i\), de manera que \(s(i)\) és un element de la fibra \(E_i=p^{-1}(i)\); així \(s\) tria simultàniament un element a cada fibra, és a dir és una funció d'elecció (definició~\ref{def:seccio-familia}). Si \(p\) és un epimorfisme regular ---totes les fibres no inicials--- l'axioma d'elecció (axioma~\ref{ax:eleccio}) afirma exactament que \(p\) s'escindeix, és a dir que admet una secció (proposició~\ref{prop:epi-regular-eleccio}). Sense elecció, en canvi, una família de fibres totes no buides pot no tenir cap secció; la tria simultània no és automàtica. Quan cada fibra porta un element distingit canònic (un mínim, un ínfim), la secció es construeix sense elecció, com als subobjectes de \(\N\) (observació~\ref{obs:minim-vs-eleccio}).
\end{solucio}

\section{Suma i producte}\label{prac:suma-producte}

\begin{exercici}
\textbf{El producte com a conjunt de seccions.} Demostreu que \(\prod_{i\in I}E_i=\{s\in E^I : p\circ s=\id_I\}\) existeix com a part interna de \(E^I\), i que els seus elements globals són les seccions de \(p\). (Reproduïu la proposició~\ref{prop:producte-existeix}.)
\end{exercici}
\begin{solucio}
L'exponencial \(E^I\) existeix per l'axioma~\ref{ax:exponencials}. Les aplicacions \(s\mapsto p\circ s\) (postcomposició interna amb \(p\), capítol~\ref{cap:exponencials}) i \(s\mapsto\id_I\) (constant) van de \(E^I\) a \(I^I\); el seu equalitzador és una part de \(E^I\) (capítol~\ref{cap:valors}), que és per definició \(\prod_{i\in I}E_i\). Cap tria fins aquí: l'objecte existeix sempre. Un element global és una \(s\colon\mathbf1\to E^I\) amb \(p\circ s=\id_I\), que per transposició (definició~\ref{def:exponencial}) és una aplicació \(s\colon I\to E\) amb \(p\circ s=\id_I\), és a dir una secció (definició~\ref{def:seccio-familia}). Per tant el producte és no inicial si i només si \(p\) té una secció.
\end{solucio}

\begin{exercici}
\textbf{Objecte contra elements.} Expliqueu, amb el producte, la diferència entre \enquote{el producte existeix} i \enquote{el producte té elements}, i digueu on entra l'elecció. (Model de l'observació~\ref{obs:objecte-vs-elements}.)
\end{exercici}
\begin{solucio}
Que el producte \emph{existeixi com a objecte} és sempre cert: es construeix amb exponencials i un equalitzador, sense cap tria (exercici anterior). Que \emph{tingui algun element global} equival, per a fibres no buides, a l'existència d'una secció de \(p\), i això és exactament l'axioma d'elecció (proposició~\ref{prop:epi-regular-eleccio}). La notació \(\prod_{i\in I}E_i\) amaga aquesta diferència: escriure-la no compromet cap tria; afirmar que és no inicial, sí. És la versió general del producte numerable (observació~\ref{obs:producte-numerable}), on l'objecte hi és sempre però els seus punts costen l'elecció numerable.
\end{solucio}

\section{Famílies de parts i operacions indexades}\label{prac:parts-indexades}

\begin{exercici}
\textbf{Unió i intersecció indexades.} Sigui \(\alpha\colon I\to P(A)\) una família de parts de \(A\). Definiu \(\bigcup_{i}\alpha(i)\) i \(\bigcap_{i}\alpha(i)\) amb els quantificadors interns i demostreu que són, respectivament, el suprem i l'ínfim de la família a \((P(A),\subseteq)\). (Reproduïu la proposició~\ref{prop:unio-interseccio-propietats}.)
\end{exercici}
\begin{solucio}
Sigui \(R=\{(i,a) : a\mathrel{\varepsilon}\alpha(i)\}\hookrightarrow I\times A\) la relació de pertinença total (definició~\ref{def:familia-parts}) i \(\pi_A\colon I\times A\to A\) la projecció. Posem \(\bigcup_i\alpha(i)=\exists_{\pi_A}(R)\) i \(\bigcap_i\alpha(i)=\forall_{\pi_A}\) del predicat \((i,a)\in R\) (definició~\ref{def:unio-interseccio-indexada}). Per l'adjunció dels quantificadors (proposició~\ref{prop:quantificadors-adjuncio}, estadi \(\mathbf1\)), \(\exists_{\pi_A}(R)\subseteq S\) equival a \(R\subseteq\pi_A^{-1}(S)=I\times S\), és a dir, sobre punts, a \(\alpha(i)\subseteq S\) per a tot \(i\): això caracteritza \(\bigcup_i\alpha(i)\) com la menor part que conté totes les \(\alpha(i)\), el suprem. Dualment, \(\forall_{\pi_A}\) dona la major part continguda en totes, l'ínfim. Tot es fa dins \(P(A)\) amb els quantificadors interns, sense cap axioma nou.
\end{solucio}

\begin{exercici}
\textbf{De Morgan indexada.} Demostreu que \(A\setminus\bigcup_i\alpha(i)=\bigcap_i(A\setminus\alpha(i))\). (Part de la proposició~\ref{prop:unio-interseccio-propietats}.)
\end{exercici}
\begin{solucio}
Sobre punts, \(a\) és a \(A\setminus\bigcup_i\alpha(i)\) quan \emph{no} existeix cap \(i\) amb \(a\mathrel{\varepsilon}\alpha(i)\), és a dir quan \emph{per a tot} \(i\) es té \(a\mathrel{\varepsilon}(A\setminus\alpha(i))\); i això és pertànyer a \(\bigcap_i(A\setminus\alpha(i))\). Però la igualtat de parts no s'ha de comprovar sobre punts: és directament la identitat \(\neg\exists_{\pi_A}(R)=\forall_{\pi_A}(\neg R)\), que surt de \(\forall_f=\neg\exists_f\neg\) (proposició~\ref{prop:quantificadors-adjuncio}) i de les lleis de De Morgan booleanes transportades a la família (observació~\ref{obs:lleis-transportades}). Cap pas usa elecció.
\end{solucio}

\begin{exercici}
\textbf{Una partició a partir de les seves peces.} Sigui \(\chi\colon\N\to\mathbf2\) la \emph{paritat}, recursora del sistema \((\mathbf2,0,\neg)\): \(\chi(0)=0\) i \(\chi(s\circ n)=\neg\chi(n)\). Siguin \(\alpha(0):=\chi^{-1}(0)\) (els parells) i \(\alpha(1):=\chi^{-1}(1)\) (els senars), vistes com una família \(\alpha\colon\mathbf2\to P(\N)\). Comproveu les hipòtesis de la proposició~\ref{prop:familia-particio} i identifiqueu la partició \(q\colon\N\twoheadrightarrow\mathbf2\) que en resulta.
\end{exercici}
\begin{solucio}
Les dues parts són les imatges inverses dels dos punts de \(\mathbf2\) per \(\chi\), és a dir, la part classificada per \(\chi\) i el seu complement (capítol~\ref{cap:valors}): són disjuntes i la seva unió és tot \(\N\) per construcció. Són no buides: \(0\in\alpha(0)\) i \(s\circ0\in\alpha(1)\), perquè \(\chi(s\circ0)=\neg0=1\). La relació de pertinença total \(R\hookrightarrow\mathbf2\times\N\) és el graf de \(\chi\) amb els factors intercanviats, i \(r_\N\colon R\to\N\) és un isomorfisme (tot punt té exactament una paritat); la partició és \(q=r_{\mathbf2}\circ r_\N^{-1}=\chi\) mateixa, amb fibres \(\alpha(0)\) i \(\alpha(1)\). En general, una partició en dues peces és el mateix que una característica \(A\to\mathbf2\) amb les dues fibres no buides; la proposició ho generalitza a un nombre qualsevol de peces, indexat per \(I\).
\end{solucio}

\section{Regla externa contra família interna}\label{prac:regla-vs-familia}

\begin{exercici}
\textbf{\(n\mapsto\seg{n}\) com a regla externa.} Expliqueu per què l'assignació \(n\mapsto\seg{n}\) és una regla externa que \emph{sí} es pot representar per una família interna, i com. (Model de l'observació~\ref{obs:regla-externa-problema}.)
\end{exercici}
\begin{solucio}
La prescripció \(n\mapsto\seg{n}\) associa a cada natural el seu segment inicial; és una regla clara però, en si mateixa, no és cap aplicació \(E\to\N\). Es \emph{representa} per una família interna allotjant tots els \(\seg{n}\) dins un ambient comú: com que cada \(\seg{n}\) és una part de \(\N\), la regla és l'aplicació \(\N\to P(\N)\), \(n\mapsto\seg{n}\) (una família de parts d'un fix, definició~\ref{def:familia-parts}); alternativament, el total es realitza dins un exponencial fix, com al conjunt de llistes (definició~\ref{def:llistes}). La representació existeix precisament perquè hi ha un ambient (\(\N\), o \(P(\N)\)) que conté tots els valors. Sense un tal ambient comú, la reunió podria escapar-se dels axiomes.
\end{solucio}

\begin{exercici}
\textbf{Quan la regla no es representa.} Descriviu una situació en què una regla externa \(i\mapsto X_i\) \emph{no} es pot representar amb els axiomes de Nivell~I, i digueu quin principi caldria. (Model de les observacions~\ref{obs:regla-externa-problema} i~\ref{obs:sr}.)
\end{exercici}
\begin{solucio}
Si els conjunts \(X_i\) creixen sense cap cota comuna coneguda ---per exemple la regla \(n\mapsto P^n(A)\) que itera el conjunt de parts, \(A, P(A), P(P(A)),\dots\), on cada terme és estrictament més gran que l'anterior pel teorema de Cantor (teorema~\ref{teo:cantor})--- no hi ha cap conjunt ambient dins el qual tallar tots els \(X_i\) per separació, de manera que no se'n pot formar cap total \(E\to\N\) amb els axiomes bàsics. Reunir aquesta regla en una família interna demana el principi de substitució estructural SR (observació~\ref{obs:sr}), que és de Nivell~II i que no adoptem aquí. Noteu que això \emph{no} és una qüestió d'elecció: no cal triar res dins fibres, sinó garantir l'existència del total (observació~\ref{obs:sr-no-eleccio}).
\end{solucio}

\chapter{Ordres, bons ordres i segments}\label{cap:practica-bons-ordres}

Aquest capítol acompanya el capítol homònim de la part teòrica (capítol~\ref{cap:bons-ordres} de la Part~\ref{part:teoria}). Els exercicis resolts en segueixen l'ordre d'exposició ---ordres i bons ordres, segments inicials, aplicacions creixents i rigidesa, i la comparació per segments--- i serveixen de model per als exercicis proposats de la Part~\ref{part:exercicis}. Cada bloc fixa una tècnica: passar del mínim d'una part no inicial a la totalitat, identificar les parts inicials amb els segments, i explotar la rigidesa (cap punt no baixa) per obtenir la unicitat dels isomorfismes d'ordre. Convé intentar cada enunciat abans de llegir-ne la solució.

\section{Ordres i bons ordres}\label{prac:ordres}

\begin{exercici}
\textbf{Tot bon ordre és total.} Demostreu que si \({<}\) és un bon ordre sobre \(W\), aleshores dos elements distints qualssevol són comparables, i que el mínim d'una part no inicial és únic. (Reproduïu la proposició~\ref{prop:bon-ordre-total}.)
\end{exercici}
\begin{solucio}
La unicitat del mínim surt de l'antisimetria de \(\le\) (observació~\ref{obs:ordre-part-general}): dos mínims \(m,m'\) compleixen \(m\le m'\) i \(m'\le m\), d'on \(m=m'\). Per a la totalitat, siguin \(u\ne w\) dos punts i \(\{u,w\}\hookrightarrow W\) la imatge del parell \(\langle u,w\rangle\) (teorema~\ref{teo:factoritzacio}), una part no inicial. Pel bon ordre té un mínim, que és \(u\) o \(w\). Si és \(u\), aleshores \(u\le w\) i, com que \(u\ne w\), \(u<w\); si és \(w\), \(w<u\). Per tant \(u<w\) o \(w<u\).
\end{solucio}

\begin{exercici}
\textbf{\(\Z\) no és ben ordenat.} Demostreu que \(\Z\) amb l'ordre usual és total però no és un bon ordre. (Model de l'observació~\ref{obs:bon-ordre-classic}.)
\end{exercici}
\begin{solucio}
Que \(\le\) és total sobre \(\Z\) és estàndard: donats dos enters, un dels dos és menor o igual. Però \(\Z\) mateix, com a part no inicial de \(\Z\), no té element mínim: si \(m_0\) en fos el mínim, \(m_0-1<m_0\) seria un enter estrictament menor, contra la minimalitat. Equivalentment, la cadena descendent \(0>-1>-2>\cdots\) no s'atura, i la seva imatge és una part no buida sense mínim. Per tant \(\le\) sobre \(\Z\) és un ordre total que no és un bon ordre.
\end{solucio}

\section{Segments inicials}\label{prac:segments-bo}

\begin{exercici}
\textbf{Els segments són ben ordenats.} Demostreu que en un bon ordre \((W,<)\) tota part \(S\hookrightarrow W\), amb l'ordre restringit, és un bon ordre; en particular tot segment \(W_{<w}\). (Reproduïu la proposició~\ref{prop:segment-ben-ordenat}.)
\end{exercici}
\begin{solucio}
La restricció \({<}\cap(S\times S)\) és una part de \(S\times S\) (intersecció de parts, capítol~\ref{cap:valors}), irreflexiva i transitiva perquè ho és \({<}\). Donada una part no inicial \(P\hookrightarrow S\), via \(S\hookrightarrow W\) és una part no inicial de \(W\) (composició de monomorfismes, proposició~\ref{prop:preceq-preordre}), que pel bon ordre de \(W\) té un mínim \(m\); com que \(m\in P\subseteq S\) i \(m\le k\) per a tot \(k\in P\), \(m\) és el mínim de \(P\) dins \(S\). Així \({<}\!\restriction_S\) és un bon ordre, i \(S=W_{<w}\) n'és un cas.
\end{solucio}

\begin{exercici}
\textbf{Parts inicials i segments.} Demostreu que en un bon ordre tota part inicial pròpia \(S\subsetneq W\) és un segment \(W_{<w}\), amb \(w\) el mínim de \(W\setminus S\), i que el punt \(w\) és únic. (Reproduïu les proposicions~\ref{prop:parts-inicials} i~\ref{prop:segment-determina-punt}.)
\end{exercici}
\begin{solucio}
Com que \(S\) és pròpia, \(W\setminus S\) (teorema~\ref{teo:part-complementari}) és no inicial i té un mínim \(w\). \emph{\(S\subseteq W_{<w}\):} si \(u\in S\) i fos \(w\le u\), la inicialitat de \(S\) donaria \(w\in S\), contra \(w\in W\setminus S\); per tant \(u<w\). \emph{\(W_{<w}\subseteq S\):} si \(u<w\), aleshores \(u\notin W\setminus S\) perquè \(w\) n'és el mínim, d'on \(u\in S\). Així \(S=W_{<w}\). \emph{Unicitat:} si \(W_{<v}=W_{<w}\) amb \(v\ne w\), la totalitat dona (diguem) \(v<w\), d'on \(v\in W_{<w}=W_{<v}\), és a dir \(v<v\), contra la irreflexivitat.
\end{solucio}

\section{Creixents i rigidesa}\label{prac:rigidesa}

\begin{exercici}
\textbf{Cap punt no baixa.} Sigui \((W,<)\) un bon ordre i \(f\colon W\to W\) creixent i injectiva. Demostreu que \(w\le f(w)\) per a tot \(w\). (Reproduïu el lema~\ref{lema:rigidesa}.)
\end{exercici}
\begin{solucio}
Sigui \(P=\{w : f(w)<w\}\hookrightarrow W\) la part dels punts que baixen (imatge inversa per \(\langle f,\id_W\rangle\) de \({<}\), capítols~\ref{cap:parts} i~\ref{cap:relacions}). Si \(P\) no fos inicial, tindria un mínim \(w_0\) amb \(f(w_0)<w_0\). Aplicant \(f\), creixent, \(f(f(w_0))<f(w_0)\), de manera que \(w_1:=f(w_0)\) també és a \(P\) i \(w_1<w_0\), contra la minimalitat de \(w_0\). Per tant \(P\) és inicial: cap punt baixa, i \(w\le f(w)\).
\end{solucio}

\begin{exercici}
\textbf{Unicitat de l'isomorfisme d'ordre.} Demostreu que entre dos bons ordres hi ha com a molt un isomorfisme d'ordre, i que l'únic automorfisme d'ordre és la identitat. (Reproduïu el teorema~\ref{teo:unicitat-iso}.)
\end{exercici}
\begin{solucio}
Siguin \(f,g\colon W\to V\) isomorfismes d'ordre. Aleshores \(g^{-1}\circ f\) i \(f^{-1}\circ g\) són automorfismes d'ordre de \(W\), inversos l'un de l'altre. Pel lema~\ref{lema:rigidesa}, \(w\le(g^{-1}\circ f)(w)\) per a tot \(w\); i \(w\le(f^{-1}\circ g)(w)\), que component amb \(g^{-1}\circ f\) creixent dona \((g^{-1}\circ f)(w)\le w\). Les dues desigualtats forcen \((g^{-1}\circ f)(w)=w\), és a dir \(g^{-1}\circ f=\id_W\) i \(f=g\). El cas \(W=V\) dona l'automorfisme únic.
\end{solucio}

\begin{exercici}
\textbf{Cap bon ordre és isomorf a un segment propi.} Demostreu que cap bon ordre \((W,<)\) és isomorf a cap segment inicial propi seu. (Reproduïu el corol·lari~\ref{cor:no-segment-propi}.)
\end{exercici}
\begin{solucio}
Si \(f\colon W\cong W_{<w}\) fos un isomorfisme d'ordre, component amb la inclusió \(W_{<w}\hookrightarrow W\) tindríem \(f\colon W\to W\) creixent i injectiva amb \(f(w)\in W_{<w}\), és a dir \(f(w)<w\), contra el lema~\ref{lema:rigidesa} (\(w\le f(w)\)). Per tant no existeix cap tal isomorfisme; el cas d'una part inicial pròpia qualsevol es redueix a un segment (proposició~\ref{prop:parts-inicials}).
\end{solucio}

\section{Comparació de bons ordres}\label{prac:comparacio-ordres}

\begin{exercici}
\textbf{Antisimetria de la comparació.} Demostreu que si \(W\preceq_{\mathrm o}V\) i \(V\preceq_{\mathrm o}W\), aleshores \(W\cong V\). (Part de la proposició~\ref{prop:comparacio-ordres}.)
\end{exercici}
\begin{solucio}
Per transitivitat (proposició~\ref{prop:comparacio-ordres}(2)), \(W\preceq_{\mathrm o}V\preceq_{\mathrm o}W\) dona un isomorfisme d'ordre de \(W\) sobre un segment inicial de \(W\). Pel corol·lari~\ref{cor:no-segment-propi} aquest segment no pot ser propi, així que és tot \(W\), i la composició és un automorfisme, la identitat (teorema~\ref{teo:unicitat-iso}). Desfent-la, cadascun dels dos isomorfismes de partida és sobre tot l'altre bon ordre, d'on \(W\cong V\).
\end{solucio}

\begin{exercici}
\textbf{Exclusivitat de la tricotomia.} Demostreu que, entre dos bons ordres, com a molt una de les tres situacions \(W\prec_{\mathrm o}V\), \(W\cong V\), \(V\prec_{\mathrm o}W\) es dona. (Part del teorema~\ref{teo:tricotomia}.)
\end{exercici}
\begin{solucio}
Si \(W\cong V\) i \(W\prec_{\mathrm o}V\) alhora, component \(V\cong W\prec_{\mathrm o}V\) dona \(V\) isomorf a un segment propi seu, contra el corol·lari~\ref{cor:no-segment-propi}; igual amb \(W\cong V\) i \(V\prec_{\mathrm o}W\). Si \(W\prec_{\mathrm o}V\) i \(V\prec_{\mathrm o}W\) alhora, la transitivitat (proposició~\ref{prop:comparacio-ordres}) dona \(W\prec_{\mathrm o}W\), altre cop contra el corol·lari~\ref{cor:no-segment-propi}. Les tres són, doncs, mútuament excloents. (Que almenys una es doni ---la comparabilitat total--- es demostra per recursió transfinita al capítol següent, sense elecció; vegeu l'observació~\ref{obs:tricotomia-recursio}.)
\end{solucio}

\chapter{Inducció i recursió transfinites}\label{cap:practica-recursio}

Aquest capítol acompanya el capítol homònim de la part teòrica (capítol~\ref{cap:recursio} de la Part~\ref{part:teoria}). Els exercicis resolts en segueixen l'ordre d'exposició ---inducció transfinita, recursió amb codomini fix, la comparabilitat total dels bons ordres, i exemples de definició per recursió--- i serveixen de model per als exercicis proposats de la Part~\ref{part:exercicis}. Cada bloc fixa una tècnica: usar el mínim del complement per a la inducció, reunir solucions parcials compatibles dins un exponencial fix per a la recursió, i aparellar dos bons ordres pel mínim disponible sense elecció. Convé intentar cada enunciat abans de llegir-ne la solució.

\section{Inducció transfinita}\label{prac:induccio-transfinita}

\begin{exercici}
\textbf{El principi d'inducció transfinita.} Sigui \((W,<)\) un bon ordre i \(S\hookrightarrow W\) tal que \(W_{<w}\subseteq S\) implica \(w\in S\), per a tot \(w\). Demostreu que \(S=W\). (Reproduïu el teorema~\ref{teo:induccio-transfinita}.)
\end{exercici}
\begin{solucio}
Suposem \(S\ne W\); aleshores \(W\setminus S\) (teorema~\ref{teo:part-complementari}) és no inicial i, pel bon ordre (definició~\ref{def:bon-ordre}), té un mínim \(w_0\). Com que \(w_0\) és el mínim del complement, cap element \(<w_0\) no és al complement, és a dir \(W_{<w_0}\subseteq S\). Per la progressivitat aplicada a \(w_0\), \(w_0\in S\), contra \(w_0\in W\setminus S\). La contradicció força \(W\setminus S\) inicial, és a dir \(S=W\).
\end{solucio}

\begin{exercici}
\textbf{Cap cas base separat.} Expliqueu per què la inducció transfinita no necessita un cas base per al mínim global \(\bot\) de \(W\), a diferència de la inducció sobre \(\N\). (Model de l'observació~\ref{obs:induccio-propietats}.)
\end{exercici}
\begin{solucio}
El cas del mínim \(\bot\) queda inclòs a la clàusula general. En efecte, \(W_{<\bot}\) és inicial, de manera que la hipòtesi \enquote{\(W_{<\bot}\subseteq S\)} es compleix \emph{buidament} (no hi ha cap element \(<\bot\)); la progressivitat aplicada a \(\bot\) força aleshores \(\bot\in S\). Així, la mateixa clàusula \enquote{si tot el segment inicial és a \(S\), llavors \(w\) hi és} cobreix el mínim, els successors i els límits d'un sol cop. A \(\N\) es separava zero i successor perquè tot element és un o l'altre; en un bon ordre general això no val (hi ha límits sense predecessor), i la formulació amb el segment inicial els unifica.
\end{solucio}

\section{Recursió amb codomini fix}\label{prac:recursio-fix}

\begin{exercici}
\textbf{Compatibilitat de les solucions parcials.} Demostreu que dues solucions parcials \(a\colon D\to X\) i \(b\colon E\to X\) de la mateixa regla \(\Phi\) coincideixen sobre \(D\cap E\). (Pas de la prova del teorema~\ref{teo:recursio-transfinita}.)
\end{exercici}
\begin{solucio}
\(D\cap E\) és inicial (intersecció de parts inicials). Apliquem la inducció transfinita (teorema~\ref{teo:induccio-transfinita}) a \(S=\{w : w\notin D\cap E\ \text{o}\ a(w)=b(w)\}\). Sigui \(w\) amb \(W_{<w}\subseteq S\). Si \(w\notin D\cap E\), ja \(w\in S\). Si \(w\in D\cap E\), aleshores \(W_{<w}\subseteq D\cap E\) per inicialitat, i la hipòtesi dona \(a=b\) sobre \(W_{<w}\), és a dir \(a|_{W_{<w}}=b|_{W_{<w}}\); l'equació de recursió dona \(a(w)=\Phi(w,a|_{W_{<w}})=\Phi(w,b|_{W_{<w}})=b(w)\), d'on \(w\in S\). Per inducció \(S=W\), és a dir \(a=b\) sobre \(D\cap E\).
\end{solucio}

\begin{exercici}
\textbf{Existència i unicitat de \(F\).} Demostreu que per a tota regla \(\Phi\) amb valors al fix \(X\) existeix una única \(F\colon W\to X\) amb \(F(w)=\Phi(w,F|_{W_{<w}})\). (Reproduïu el teorema~\ref{teo:recursio-transfinita}.)
\end{exercici}
\begin{solucio}
\emph{Existència.} Per inducció transfinita es prova que per a cada \(w\) hi ha una solució parcial definida a \(W_{\le w}\): donat \(w\) amb solucions parcials fins a cada \(u<w\), la compatibilitat (exercici anterior) fa que la seva reunió sigui una solució parcial \(g\colon W_{<w}\to X\), i estenent-la amb \(g^+(w):=\Phi(w,g)\) s'obté una solució parcial fins a \(W_{\le w}\). La reunió de \emph{totes} les solucions parcials és, per compatibilitat, una única aplicació \(F\) sobre una part inicial \(W'\); com que hi ha solucions fins a cada \(w\), \(W'=W\), i \(F\colon W\to X\) satisfà l'equació a tot arreu. Tot viu dins l'exponencial fix \((X+\mathbf1)^W\) i dins \(P(W)\). \emph{Unicitat.} Dues solucions totals són solucions parcials amb domini \(W\); per compatibilitat coincideixen, \(F=F'\).
\end{solucio}

\begin{exercici}
\textbf{Sense elecció.} Expliqueu per què la construcció de \(F\) no usa cap elecció. (Model de l'observació~\ref{obs:recursio-sense-eleccio}.)
\end{exercici}
\begin{solucio}
Cada solució parcial està \emph{determinada} per l'equació de recursió: a cada punt del seu domini, el valor és \(\Phi(w,\cdot)\) de la història, no un element triat entre alternatives. La reunió de solucions parcials compatibles és una operació canònica sobre parts (unió indexada, capítol~\ref{cap:families}), no una selecció. I tot passa dins l'exponencial fix \((X+\mathbf1)^W\), un conjunt ja construït. No hi ha, doncs, cap punt on calgui triar simultàniament un element de cada membre d'una família; per això no intervé l'elecció. El bon ordre només s'usa, via la inducció, per garantir compatibilitat i existència fins a cada punt.
\end{solucio}

\section{La comparabilitat total}\label{prac:comparabilitat}

\begin{exercici}
\textbf{Dos bons ordres són comparables.} Demostreu que, donats dos bons ordres \(W\) i \(V\), es dona una de les tres situacions de la tricotomia, construint per recursió l'aparellament pel mínim disponible. (Reproduïu el teorema~\ref{teo:comparabilitat}.)
\end{exercici}
\begin{solucio}
Sigui \(V^+=V+\mathbf1\) amb punt exterior \(\top\), i \(F\colon W\to V^+\) definida per recursió (teorema~\ref{teo:recursio-transfinita}): \(F(w)=\min(V\setminus F(W_{<w}))\) si el complement és no buit, i \(\top\) altrament. El mínim és únic i canònic, de manera que la regla és interna i no usa elecció. Sobre \(W_0=\{w:F(w)\ne\top\}\), inicial, \(F\) és creixent i injectiva amb imatge un segment inicial de \(V\). Si \(F\) no val mai \(\top\), \(F\) és un isomorfisme de \(W\) sobre un segment de \(V\) (tot \(V\): \(W\cong V\); propi: \(W\prec_{\mathrm o}V\)). Si \(w_0\) és el mínim amb \(F(w_0)=\top\), aleshores \(F(W_{<w_0})=V\) i \(F|_{W_{<w_0}}\colon W_{<w_0}\cong V\), d'on \(V\prec_{\mathrm o}W\). En tots els casos, comparables; l'exclusivitat és el teorema~\ref{teo:tricotomia}.
\end{solucio}

\begin{exercici}
\textbf{Per què no cal elecció.} Contrasteu la comparabilitat de dos bons ordres amb la de dos conjunts qualssevol pel que fa a l'elecció. (Model de l'observació~\ref{obs:comparabilitat-vs-eleccio}.)
\end{exercici}
\begin{solucio}
A l'aparellament de dos bons ordres, a cada pas es tria el \emph{mínim} del que queda de \(V\); aquest mínim existeix i és únic pel bon ordre, de manera que no hi ha cap tria arbitrària i tota la construcció és canònica. En canvi, la comparabilitat de dos conjunts \emph{qualssevol} ---\(A\preceq B\) o \(B\preceq A\)--- no disposa de cap tria canònica sobre conjunts nus, i s'obté només via el teorema del bon ordre, que és equivalent a l'elecció. La diferència és, doncs, que un bon ordre porta incorporada a cada pas la tria que en un conjunt nu s'ha de postular.
\end{solucio}

\section{Exemples de recursió}\label{prac:exemples-recursio}

\begin{exercici}
\textbf{Successió transfinita de parts.} Donada \(g\colon P(A)\to P(A)\), definiu per recursió una successió \((S_w)_{w\in W}\) de parts d'\(A\) amb \(S_w=\bigcup_{u<w}g(S_u)\), i digueu per què és de Nivell~I. (Model de l'observació~\ref{obs:exemples-recursio}.)
\end{exercici}
\begin{solucio}
Prenem \(X=P(A)\), fix, i la regla \(\Phi(w,h)=\bigcup_{u<w}g(h(u))\), que a partir de la història \(h=S|_{W_{<w}}\) forma la unió indexada de les \(g(S_u)\) ---una part d'\(A\), obtinguda amb els quantificadors interns (capítol~\ref{cap:families}). Pel teorema~\ref{teo:recursio-transfinita} existeix una única \(S\colon W\to P(A)\) amb \(S_w=\bigcup_{u<w}g(S_u)\). És de Nivell~I perquè totes les \(S_w\) viuen dins el conjunt \emph{fix} \(P(A)\): el codomini no varia, de manera que no cal cap principi nou. La monotonia de \(g\), si es té, fa la successió creixent.
\end{solucio}

\begin{exercici}
\textbf{On cal SR.} Expliqueu per què la successió \(V_0=A\), \(V_{n+1}=P(V_n)\) \emph{no} es pot definir amb el teorema de recursió d'aquest capítol, i quin principi cal. (Model de l'observació~\ref{obs:frontera-exemple}.)
\end{exercici}
\begin{solucio}
El teorema~\ref{teo:recursio-transfinita} demana un codomini \emph{fix} \(X\) que contingui tots els valors. Aquí cada \(V_{n+1}=P(V_n)\) és estrictament més gran que \(V_n\) pel teorema de Cantor (teorema~\ref{teo:cantor}), de manera que no existeix cap conjunt \(X\) que contingui tots els \(V_n\) com a elements: el codomini hauria de créixer a cada pas. La regla \(n\mapsto V_n\) és una \emph{regla externa} (definició~\ref{def:familia-vs-regla}) i reunir-la en una família interna demana la substitució estructural SR (observació~\ref{obs:sr}), que és de Nivell~II. No és una qüestió d'elecció: no es tria res dins fibres, sinó que cal garantir l'existència del total (observació~\ref{obs:sr-no-eleccio}).
\end{solucio}

\chapter{Zorn, Hartogs i el teorema del bon ordre}\label{cap:practica-zorn}

Aquest capítol acompanya el capítol homònim de la part teòrica (capítol~\ref{cap:zorn} de la Part~\ref{part:teoria}). Els exercicis resolts en segueixen l'ordre d'exposició ---Hartogs, el teorema del bon ordre, el lema de Zorn i les seves aplicacions, i les tres equivalències--- i serveixen de model per als exercicis proposats de la Part~\ref{part:exercicis}. Cada bloc fixa una tècnica: acotar una construcció dins \(P(A\times A)\), usar \(\Hart(A)\) com a cota que atura una recursió, distingir maximal de màxim, i tancar una cadena d'equivalències amb les dependències explícites. Convé intentar cada enunciat abans de llegir-ne la solució.

\section{Hartogs}\label{prac:hartogs}

\begin{exercici}
\textbf{\(\Hart(A)\) és un bon ordre.} Demostreu que el conjunt dels tipus d'ordre dels bons ordres sobre parts d'\(A\), ordenat per segment inicial, és un bon ordre. (Reproduïu la proposició~\ref{prop:hartogs-bon-ordre}.)
\end{exercici}
\begin{solucio}
Els bons ordres sobre parts d'\(A\) formen una part \(\WO(A)\hookrightarrow P(A)\times P(A\times A)\), els parells \((S,R)\) amb \(S\hookrightarrow A\) i \(R\hookrightarrow S\times S\) bon ordre estricte sobre tot \(S\) (definició~\ref{def:bons-ordres-parts}); la relació \enquote{ser isomorf} és una equivalència interna, i pel quocient efectiu (axioma~\ref{ax:exactitud}) s'obté \(\Hart(A)=\WO(A)/\cong\). L'ordre \(<\) per segment inicial és total i exclusiu per la comparabilitat de bons ordres (teorema~\ref{teo:comparabilitat}) i la tricotomia (teorema~\ref{teo:tricotomia}); és irreflexiu perquè cap bon ordre no és isomorf a un segment propi seu (corol·lari~\ref{cor:no-segment-propi}). Per al bon ordre, donada una part no inicial \(T\) i \(\alpha_0=[(S_0,R_0)]\in T\), els tipus de \(T\) menors que \(\alpha_0\) es corresponen amb segments inicials de \((S_0,R_0)\), és a dir amb una part no inicial del suport \(S_0\) (proposició~\ref{prop:parts-inicials}), que té mínim pel bon ordre \(R_0\); aquest mínim dona el mínim de \(T\).
\end{solucio}

\begin{exercici}
\textbf{Teorema de Hartogs.} Demostreu que \(\Hart(A)\) no admet cap monomorfisme en \(A\). (Reproduïu el teorema~\ref{teo:hartogs}.)
\end{exercici}
\begin{solucio}
Si \(m\colon\Hart(A)\hookrightarrow A\) fos un monomorfisme, transportant l'ordre de \(\Hart(A)\) per \(m\) la imatge \(S_*:=m(\Hart(A))\) seria una part d'\(A\) amb un bon ordre \(R_*\hookrightarrow S_*\times S_*\) isomorf a \(\Hart(A)\); així \((S_*,R_*)\in\WO(A)\) i el seu tipus \([(S_*,R_*)]\) és un element de \(\Hart(A)\). Però el segment inicial de \(\Hart(A)\) determinat per \([(S_*,R_*)]\) és isomorf a un bon ordre d'aquell tipus (lema~\ref{lema:segment-tipus}), és a dir a \(\Hart(A)\) mateix; per tant \(\Hart(A)\) seria isomorf a un segment inicial \emph{propi} de si mateix (propi perquè \([(S_*,R_*)]\in\Hart(A)\)), contra el corol·lari~\ref{cor:no-segment-propi}. La contradicció mostra que no hi ha cap tal \(m\): \(\Hart(A)\not\preceq A\).
\end{solucio}

\section{El teorema del bon ordre}\label{prac:bon-ordre}

\begin{exercici}
\textbf{Tot conjunt admet un bon ordre.} Amb l'elecció, demostreu que tot conjunt \(A\) es pot ben ordenar, fent la recursió acotada per \(\Hart(A)\). (Reproduïu el teorema~\ref{teo:bon-ordre}.)
\end{exercici}
\begin{solucio}
Per l'elecció, sigui \(c\) una funció que tria un element de cada part no inicial d'\(A\) (proposició~\ref{prop:epi-regular-eleccio}). Amb \(A^+=A+\mathbf1\) i punt exterior \(\top\), definim per recursió (teorema~\ref{teo:recursio-transfinita}) sobre \(\Hart(A)\) una \(F\colon\Hart(A)\to A^+\): \(F(\alpha)=c(A\setminus F(\Hart(A)_{<\alpha}))\) si el complement és no buit, i \(\top\) altrament. Sobre \(\{F\ne\top\}\), \(F\) és creixent i injectiva cap a \(A\); si no valgués mai \(\top\), seria un monomorfisme \(\Hart(A)\hookrightarrow A\), contra Hartogs (teorema~\ref{teo:hartogs}). Sigui \(\alpha_0\) el mínim amb \(F(\alpha_0)=\top\): aleshores \(F(\Hart(A)_{<\alpha_0})=A\), i \(F|_{\Hart(A)_{<\alpha_0}}\) és una bijecció del segment sobre \(A\); transportant-ne l'ordre s'obté un bon ordre sobre tot \(A\). La recursió ha estat acotada per \(\Hart(A)\), no sobre cap classe.
\end{solucio}

\begin{exercici}
\textbf{Per què la cota de Hartogs.} Expliqueu quin problema evita fer la recursió sobre \(\Hart(A)\) en lloc de \enquote{sobre tots els ordinals}. (Model de l'observació~\ref{obs:cota-hartogs}.)
\end{exercici}
\begin{solucio}
La prova clàssica \enquote{va triant un element del que queda} sembla recórrer una col·lecció d'ordinals que no és a priori un conjunt, de manera que el teorema de recursió no s'hi aplicaria directament. En fer la recursió sobre el bon ordre \emph{concret} \(\Hart(A)\) ---que sí que és un conjunt (definició~\ref{def:hartogs})--- el teorema de recursió amb codomini fix (teorema~\ref{teo:recursio-transfinita}) s'aplica sense reserves, i el teorema de Hartogs garanteix que el procés s'atura abans d'esgotar \(\Hart(A)\). Així la cota converteix una recursió aparentment sense fi en una d'acotada i legítima.
\end{solucio}

\section{El lema de Zorn}\label{prac:zorn}

\begin{exercici}
\textbf{Maximal no és màxim.} Doneu un ordre parcial amb diversos maximals i cap màxim, i expliqueu la diferència. (Model de l'observació~\ref{obs:maximal-maxim}.)
\end{exercici}
\begin{solucio}
Prenem les parts pròpies d'un conjunt de dos elements \(\{0,1\}\), ordenades per inclusió: \(\varnothing\subsetneq\{0\}\) i \(\varnothing\subsetneq\{1\}\), però \(\{0\}\) i \(\{1\}\) són incomparables. Tots dos singletons són \emph{maximals} (cap part pròpia no els supera), però cap no és \emph{màxim} (no contenen l'altre). Un màxim seria \(\ge\) que tots; un maximal, només un element que cap altre no supera estrictament. En un ordre total coincideixen; en un parcial, no. El lema de Zorn dona un maximal, no un màxim.
\end{solucio}

\begin{exercici}
\textbf{El lema de Zorn.} Amb l'elecció, demostreu que un ordre parcial no inicial on tota cadena té cota superior té un element maximal. (Reproduïu el teorema~\ref{teo:zorn}.)
\end{exercici}
\begin{solucio}
Suposem que \(P\) no té maximal: aleshores tot \(x\) té un \(x'>x\), i tota cadena té una cota superior \emph{estricta}. Per l'elecció, sigui \(c\) una funció que assigna a cada cadena una cota estricta. Definim per recursió (teorema~\ref{teo:recursio-transfinita}) sobre \(\Hart(P)\) una \(F\colon\Hart(P)\to P\) amb \(F(\alpha)=c(F(\Hart(P)_{<\alpha}))\) quan la imatge és una cadena. Per inducció (teorema~\ref{teo:induccio-transfinita}), la imatge és sempre una cadena i \(F\) és estrictament creixent, per tant un monomorfisme \(\Hart(P)\hookrightarrow P\), contra Hartogs (teorema~\ref{teo:hartogs}). La contradicció ve de suposar \(P\) sense maximal; per tant en té un.
 Observeu dos punts de rigor: la tria no és d'un element d'una cadena sinó d'una \emph{cota superior estricta} seva, i es fa aplicant l'elecció a la família que sobre cada cadena té la part de les seves cotes estrictes; i la regla de recursió s'ha de completar amb una clàusula \enquote{altrament} (per exemple, un punt fix \(p_0\in P\)) perquè estigui definida per a tota història, encara que aquest cas no s'arribi a donar mai.
\end{solucio}

\begin{exercici}
\textbf{Tot espai vectorial té base.} Demostreu, amb Zorn, que tot espai vectorial té una base. (Reproduïu la proposició~\ref{prop:base}.)
\end{exercici}
\begin{solucio}
Sigui \(P\) el conjunt de les parts linealment independents de \(V\), ordenat per inclusió. Tota cadena \(C\) té cota superior \(U=\bigcup C\): \(U\) és independent perquè tota combinació lineal usa un nombre finit de vectors, tots en un mateix membre de la cadena (per ser \(C\) totalment ordenada), on són independents. Pel lema de Zorn (teorema~\ref{teo:zorn}), \(P\) té un maximal \(B\), independent per ser a \(P\). És generador: si un \(v\) no fos combinació lineal de \(B\), aleshores \(B\cup\{v\}\) seria independent i estrictament més gran, contra la maximalitat. Per tant \(B\) és una base.
\end{solucio}

\section{Les tres equivalències}\label{prac:equivalencies}

\begin{exercici}
\textbf{Bon ordre implica elecció.} Demostreu que si tot conjunt admet un bon ordre, aleshores val l'axioma d'elecció. (Part del teorema~\ref{teo:equivalencies}.)
\end{exercici}
\begin{solucio}
Sigui \(p\colon E\twoheadrightarrow I\) una família de parts no inicials. Ben ordenem \(E\). Cada fibra \(E_i=p^{-1}(i)\) és una part no inicial d'\(E\), i té un \emph{mínim} únic per al bon ordre (definició~\ref{def:bon-ordre}). L'aplicació \(s(i)=\min(E_i)\) és una secció de \(p\), definida sense cap tria addicional perquè el mínim és canònic. Per tant tota família de parts no buides té secció, que és l'elecció (axioma~\ref{ax:eleccio}). El bon ordre substitueix aquí la tria per l'operació canònica \enquote{mínim}.
\end{solucio}

\begin{exercici}
\textbf{Zorn implica bon ordre.} Demostreu que el lema de Zorn implica que tot conjunt admet un bon ordre. (Part del teorema~\ref{teo:equivalencies}.)
\end{exercici}
\begin{solucio}
Sigui \(P\) el conjunt dels bons ordres parcials \((S,R)\) d'\(A\) ---\(S\hookrightarrow A\), \(R\) bon ordre sobre \(S\)---, ordenat per extensió inicial: \((S,R)\le(S',R')\) si \(S\) és un segment inicial de \((S',R')\) i \(R=R'|_S\). Tota cadena té cota superior (la unió dels dominis i dels ordres, que és un bon ordre parcial perquè les unions de segments encaixats són ben ordenades). Pel lema de Zorn (teorema~\ref{teo:zorn}), \(P\) té un maximal \((S_*,R_*)\). Ha de ser \(S_*=A\): si quedés \(a\notin S_*\), afegint-lo al final tindríem un bon ordre parcial més gran, contra la maximalitat. Per tant \(R_*\) ben ordena tot \(A\).
\end{solucio}

\chapter{Tipus ordinals}\label{cap:practica-ordinals}

Aquest capítol acompanya el capítol homònim de la part teòrica (capítol~\ref{cap:ordinals} de la Part~\ref{part:teoria}). Els exercicis resolts en segueixen l'ordre d'exposició ---tipus ordinals i el seu ordre, zero, successor i límit, suprems acotats, i Burali--Forti--- i serveixen de model per als exercicis proposats de la Part~\ref{part:exercicis}. Cada bloc fixa una tècnica: definir sobre representants i provar la invariància, caracteritzar el successor com a menor tipus més gran, formar suprems dins un bon ordre donat, i reconèixer l'estructura diagonal de Burali--Forti. Convé intentar cada enunciat abans de llegir-ne la solució.

\section{Tipus ordinals i el seu ordre}\label{prac:tipus-ordinals}

\begin{exercici}
\textbf{L'ordre és independent del representant.} Demostreu que \(\alpha<\beta\) (definida via segments) no depèn dels representants triats de \(\alpha\) i \(\beta\). (Part de la definició~\ref{def:ordre-ordinals}.)
\end{exercici}
\begin{solucio}
Siguin \(W\cong W'\) representants de \(\alpha\) i \(V\cong V'\) de \(\beta\), amb isomorfismes d'ordre \(f\colon W\cong W'\) i \(g\colon V\cong V'\). Si \(W\) és isomorf a un segment inicial propi \(V_{<v}\) de \(V\), aleshores, transportant per \(g\), \(W\) és isomorf a \(V'_{<g(v)}\) (un isomorfisme d'ordre envia segments inicials a segments inicials), i component amb \(f^{-1}\), \(W'\) és isomorf a \(V'_{<g(v)}\), un segment propi de \(V'\). Per tant \(\alpha<\beta\) val amb un parell de representants si i només si val amb l'altre. La clau és la unicitat de l'isomorfisme d'ordre (teorema~\ref{teo:unicitat-iso}), que fa que aquests transports siguin consistents.
\end{solucio}

\begin{exercici}
\textbf{Els tipus estan totalment ordenats.} Demostreu que \(<\) sobre tipus ordinals és un ordre estricte total. (Reproduïu la proposició~\ref{prop:ordinals-total}.)
\end{exercici}
\begin{solucio}
Donats dos representants \(W\), \(V\), la tricotomia i la comparabilitat dels bons ordres (teoremes~\ref{teo:tricotomia} i~\ref{teo:comparabilitat}) diuen que es dona exactament una de: \(W\) isomorf a un segment propi de \(V\) (\(\alpha<\beta\)), \(W\cong V\) (\(\alpha=\beta\)), o \(V\) isomorf a un segment propi de \(W\) (\(\beta<\alpha\)). La transitivitat és la de la comparació per segments (proposició~\ref{prop:comparacio-ordres}) i la irreflexivitat és que cap bon ordre no és isomorf a un segment propi seu (corol·lari~\ref{cor:no-segment-propi}). Per tant \(<\) és un ordre estricte total.
\end{solucio}

\begin{exercici}
\textbf{Per què no \(\alpha\in\beta\).} Expliqueu per què en aquest llibre l'ordre ordinal es defineix per segments i no per pertinença. (Model de l'observació~\ref{obs:no-pertinença-ordinal}.)
\end{exercici}
\begin{solucio}
En la teoria de von Neumann, un ordinal és el conjunt dels ordinals menors, i \(\alpha<\beta\) es defineix com \(\alpha\in\beta\), fent servir una relació de pertinença global entre conjunts. Aquest llibre no assumeix cap pertinença global: la pertinença és sempre relativa a una part d'un conjunt fix (capítol~\ref{cap:parts-interna}). Per això \(\alpha\in\beta\) no té sentit com a definició de l'ordre ordinal; la relació correcta és \(\alpha<\beta\), definida perquè un representant de \(\alpha\) és isomorf a un segment inicial propi d'un de \(\beta\) (definició~\ref{def:ordre-ordinals}). Escriure \(\alpha\in\beta\) suggeriria una jerarquia material que el punt de vista estructural evita deliberadament.
\end{solucio}

\section{Zero, successor i límit}\label{prac:zero-successor-limit}

\begin{exercici}
\textbf{El successor és el menor tipus més gran.} Demostreu que \(\alpha<\alpha+1\) i que no hi ha cap tipus \(\gamma\) amb \(\alpha<\gamma<\alpha+1\). (Reproduïu la proposició~\ref{prop:successor-invariant}.)
\end{exercici}
\begin{solucio}
Sigui \(W\) un representant de \(\alpha\); aleshores \(W+\mathbf1\) representa \(\alpha+1\), amb el punt nou \(\ast\) després de tot \(W\). El segment inicial de \(W+\mathbf1\) determinat per \(\ast\) és \(W\), de manera que \(\alpha<\alpha+1\). Si hi hagués \(\gamma\) amb \(\alpha<\gamma<\alpha+1\), un representant \(G\) seria isomorf a un segment propi de \(W+\mathbf1\); però els segments propis de \(W+\mathbf1\) són \(W\) i els segments de \(W\) (proposició~\ref{prop:parts-inicials}), tots de tipus \(\le\alpha\), contra \(\alpha<\gamma\). Per tant \(\alpha+1\) és el menor tipus estrictament més gran que \(\alpha\).
\end{solucio}

\begin{exercici}
\textbf{Zero, successor, límit són exhaustius.} Demostreu que tot tipus ordinal és exactament un de zero, successor o límit. (Reproduïu la proposició~\ref{prop:tricotomia-successor}.)
\end{exercici}
\begin{solucio}
Sigui \(W\) un representant. Si \(W\) és inicial, el tipus és \(0\). Si no, mirem si \(W\) té element màxim ---una propietat invariant per isomorfisme d'ordre. Si en té, \(m\), aleshores \(W\cong W_{<m}+\mathbf1\), de manera que el tipus és \(\operatorname{tipus}(W_{<m})+1\), un successor. Si \(W\) no és inicial i no té màxim, el tipus no és zero ni successor (un successor té màxim, proposició~\ref{prop:successor-invariant}), i per tant és límit. Les tres classes són excloents (zero és inicial, successor té màxim, límit no) i cobreixen tots els casos.
\end{solucio}

\begin{exercici}
\textbf{\(\omega\) és el primer límit.} Demostreu que \(\omega=\operatorname{tipus}(\N)\) és un tipus límit i que és el menor. (Model de l'observació~\ref{obs:limit-classic}.)
\end{exercici}
\begin{solucio}
\(\N\) amb el seu bon ordre (teorema~\ref{teo:bon-ordre-N}) no té element màxim: per a tot \(n\) hi ha \(s(n)>n\). Per tant \(\omega\) no és successor; i no és zero perquè \(\N\) és no inicial. Així \(\omega\) és límit (definició~\ref{def:tipus-limit}). És el menor: qualsevol tipus menor que \(\omega\) és el d'un segment propi de \(\N\), és a dir d'algun \(\seg{n}\) (proposició~\ref{prop:parts-inicials}), que és finit ---zero o successor. Com que tots els tipus \(<\omega\) són zero o successors, \(\omega\) és el primer tipus límit.
\end{solucio}

\section{Suprems acotats}\label{prac:suprems}

\begin{exercici}
\textbf{Suprem d'una família acotada.} Sigui \((\alpha_i)_{i\in I}\) representada dins un bon ordre \(W\) per \(r\colon I\to W\). Demostreu que existeix \(\sup_i\alpha_i\), el menor tipus \(\ge\) tots els \(\alpha_i\). (Reproduïu la proposició~\ref{prop:suprem}.)
\end{exercici}
\begin{solucio}
Sigui \(S=\bigcup_{i}W_{<r(i)}\) la unió indexada (capítol~\ref{cap:families}) dels segments \emph{oberts} \(W_{<r(i)}\), que són els que representen els \(\alpha_i\). Com a unió de parts inicials, \(S\) és inicial, per tant un segment de \(W\) o tot \(W\) (proposició~\ref{prop:parts-inicials}), i és ben ordenat (proposició~\ref{prop:segment-ben-ordenat}). Sigui \(\sigma=\operatorname{tipus}(S)\). Cada \(\alpha_i\le\sigma\) perquè \(W_{<r(i)}\) és una part inicial de \(S\). I \(\sigma\) és el menor amb aquesta propietat: si \(\beta\ge\alpha_i\) per a tot \(i\), amb representant \(V\), cada \(W_{<r(i)}\) té un únic isomorfisme \(f_i\) sobre una part inicial de \(V\) (teorema~\ref{teo:unicitat-iso}); per la unicitat, \(f_j\) estén \(f_i\) quan \(r(i)\le r(j)\), i la reunió dels grafs és una aplicació \(f\colon S\to V\) creixent, injectiva i amb imatge inicial. Així \(\sigma\le\beta\), i \(\sigma=\sup_i\alpha_i\). Compte: amb els segments \emph{tancats} \(W_{\le r(i)}\) el resultat seria massa gran (per a una sola \(r\) igual al mínim de \(W\), \(\alpha=0\) però \(W_{\le r}\) té tipus \(1\)); vegeu l'observació~\ref{obs:suprem-oberts}.
\end{solucio}

\begin{exercici}
\textbf{Per què cal l'acotació.} Expliqueu per què no es pot prendre el suprem d'una família arbitrària de tipus sense un bon ordre comú. (Model de l'observació~\ref{obs:suprem-acotat}.)
\end{exercici}
\begin{solucio}
El suprem es construeix com la reunió de segments \emph{dins} un bon ordre \(W\) que representi tots els \(\alpha_i\). Sense un tal \(W\) ---sense un bon ordre que contingui tots els tipus de la família com a segments--- no hi ha res a reunir: caldria formar primer un bon ordre prou llarg, i si els tipus no estan acotats aquest bon ordre no existeix com a conjunt. De fet, la inexistència d'un bon ordre de \emph{tots} els tipus és el teorema de Burali--Forti (teorema~\ref{teo:burali-forti}). Per això cada suprem exigeix exhibir abans la família acotada; escriure \(\sup_i\alpha_i\) sense el \(W\) comú seria una construcció no justificada.
\end{solucio}

\section{Burali--Forti}\label{prac:burali-forti}

\begin{exercici}
\textbf{No hi ha conjunt de tots els tipus ordinals.} Demostreu que no existeix cap conjunt ben ordenat que contingui un representant de cada tipus ordinal amb la propietat de segments. (Reproduïu el teorema~\ref{teo:burali-forti}.)
\end{exercici}
\begin{solucio}
Suposem un tal \(\Omega\), ben ordenat pel segment inicial, amb \(\operatorname{tipus}(\Omega_{<\gamma})=\) el tipus representat per \(\gamma\), per a cada \(\gamma\in\Omega\). Aleshores \(\Omega\) té un tipus propi \(\omega_\Omega=\operatorname{tipus}(\Omega)\), i com que \(\Omega\) conté un representant de cada tipus, hi ha \(\gamma\in\Omega\) amb \(\operatorname{tipus}(\Omega_{<\gamma})=\omega_\Omega\). Però \(\Omega_{<\gamma}\) és un segment inicial propi de \(\Omega\), i seria isomorf a \(\Omega\) sencer, contra el corol·lari~\ref{cor:no-segment-propi}. La contradicció mostra que cap tal \(\Omega\) no existeix.
\end{solucio}

\begin{exercici}
\textbf{Burali--Forti i Cantor.} Expliqueu en quin sentit Burali--Forti és el paral·lel ordinal del teorema de Cantor. (Model de l'observació~\ref{obs:burali-cantor}.)
\end{exercici}
\begin{solucio}
Tots dos són resultats \enquote{diagonals}: la contradicció surt de fer que la construcció actuï sobre si mateixa. Cantor (teorema~\ref{teo:cantor}) dona \(A\prec P(A)\), d'on no hi ha mida màxima ni conjunt universal (corol·lari~\ref{cor:no-universal}). Burali--Forti dona que no hi ha un bon ordre de tots els tipus ordinals, d'on no hi ha longitud d'ordre màxima. La diferència de funció és que Cantor mesura cardinalitat i Burali--Forti longitud d'ordre; l'estructura de la prova ---la col·lecció de \enquote{tots els} objectes és massa gran per ser un del mateix tipus--- és la mateixa. És per això que la teoria global dels ordinals és de Nivell~II.
\end{solucio}

\chapter{Aritmètica ordinal}\label{cap:practica-aritmetica-ordinal}

Aquest capítol acompanya el capítol homònim de la part teòrica (capítol~\ref{cap:aritmetica-ordinal} de la Part~\ref{part:teoria}). Els exercicis resolts en segueixen l'ordre d'exposició ---suma, producte, exponenciació, i monotonia amb la distinció ordre/mida--- i serveixen de model per als exercicis proposats de la Part~\ref{part:exercicis}. Cada bloc fixa una tècnica: definir sobre representants i provar la invariància, calcular tipus concrets exhibint isomorfismes d'ordre, i separar amb cura l'aritmètica ordinal de la cardinal. Convé intentar cada enunciat abans de llegir-ne la solució.

\section{Suma ordinal}\label{prac:suma-ordinal}

\begin{exercici}
\textbf{La suma és un bon ordre invariant.} Demostreu que l'ordre \enquote{tot \(A\) abans de tot \(B\)} sobre \(A+B\) és un bon ordre i que el seu tipus no depèn dels representants. (Reproduïu la proposició~\ref{prop:suma-bon-ordre}.)
\end{exercici}
\begin{solucio}
És un ordre estricte total: la irreflexivitat i la transitivitat es comproven per casos (els casos mixtos són consistents perquè \enquote{\(A\) abans de \(B\)} és global), i la totalitat perquè dos elements o són al mateix sumand o a sumands diferents. És ben ordenat: donada \(P\hookrightarrow A+B\) no inicial, si \(P\cap A\) és no inicial el seu mínim (bon ordre d'\(A\)) és el mínim de \(P\); si \(P\cap A\) és inicial, \(P\subseteq B\) i el mínim és el de \(P\) en \(B\). Invariància: dos isomorfismes \(f\colon A\cong A'\), \(g\colon B\cong B'\) donen \(f+g\colon A+B\cong A'+B'\), que respecta l'ordre (teorema~\ref{teo:unicitat-iso}).
\end{solucio}

\begin{exercici}
\textbf{No commutativitat: \(1+\omega\) contra \(\omega+1\).} Demostreu que \(1+\omega=\omega\) però \(\omega+1\neq\omega\). (Reproduïu la proposició~\ref{prop:suma-no-commutativa}.)
\end{exercici}
\begin{solucio}
\(1+\omega\) és un punt seguit d'una còpia de \(\N\); l'aplicació que envia el punt afegit a \(0\) i \(n\) a \(n+1\) és un isomorfisme d'ordre amb \(\N\), d'on \(1+\omega=\omega\). En canvi \(\omega+1\) és una còpia de \(\N\) seguida d'un punt \(\ast\) posterior a tots els naturals, que és un element \emph{màxim}. Com que \(\N\) no té màxim i tenir màxim és invariant per isomorfisme d'ordre, \(\omega+1\not\cong\omega\). De fet \(\omega+1\) és el successor de \(\omega\) (definició~\ref{def:zero-successor}), estrictament més gran; per l'esquerra el punt s'absorbeix, per la dreta no.
\end{solucio}

\begin{exercici}
\textbf{Associativitat.} Demostreu que \((\alpha+\beta)+\gamma=\alpha+(\beta+\gamma)\). (Reproduïu la proposició~\ref{prop:suma-associativa}.)
\end{exercici}
\begin{solucio}
Tots dos costats són el tipus del bon ordre sobre \(A+B+C\) amb tot \(A\) abans de tot \(B\), i tot \(B\) abans de tot \(C\), i els ordres originals dins de cada sumand. L'associativitat de la suma disjunta de conjunts (capítol~\ref{cap:valors}) dona un isomorfisme \((A+B)+C\cong A+(B+C)\) que conserva aquest ordre; per la unicitat de l'isomorfisme d'ordre (teorema~\ref{teo:unicitat-iso}), els tipus coincideixen.
\end{solucio}

\section{Producte ordinal}\label{prac:producte-ordinal}

\begin{exercici}
\textbf{El producte és un bon ordre.} Demostreu que l'ordre antilexicogràfic sobre \(A\times B\) (la coordenada de \(B\) primer) és un bon ordre. (Reproduïu la proposició~\ref{prop:producte-bon-ordre}.)
\end{exercici}
\begin{solucio}
És un ordre estricte total (l'antilexicogràfic d'ordres totals ho és). Per al bon ordre, donada \(P\hookrightarrow A\times B\) no inicial: la projecció \(\pi_B(P)\) és no inicial i té mínim \(b_0\); dins la fibra sobre \(b_0\), la projecció a \(A\) té mínim \(a_0\); aleshores \((a_0,b_0)\) és el mínim de \(P\), perquè cap element amb segona coordenada més gran no el precedeix, i dins \(b_0\) mana \(a_0\). Per tant és ben ordenat.
\end{solucio}

\begin{exercici}
\textbf{\(2\cdot\omega\) contra \(\omega\cdot2\).} Demostreu que \(2\cdot\omega=\omega\) però \(\omega\cdot2=\omega+\omega\neq\omega\). (Reproduïu la proposició~\ref{prop:producte-propietats}.)
\end{exercici}
\begin{solucio}
\(2\cdot\omega\) són \(\omega\) còpies del bon ordre de dos elements: \((0,1),(2,3),\dots\), que reindexat és \(\N\); d'on \(2\cdot\omega=\omega\). \(\omega\cdot2\) són \(2\) còpies de \(\N\), l'una després de l'altra, és a dir la concatenació \(\omega+\omega\); té un element (el primer de la segona còpia) precedit per infinits altres, sense predecessor immediat, cosa que \(\N\) no té. Per tant \(\omega\cdot2\neq\omega\) i \(\omega\cdot2\neq2\cdot\omega\): el producte no és commutatiu.
\end{solucio}

\begin{exercici}
\textbf{Distributivitat per l'esquerra i fallada per la dreta.} Demostreu que \(\alpha\cdot(\beta+\gamma)=\alpha\cdot\beta+\alpha\cdot\gamma\), i doneu un contraexemple de la distributivitat per la dreta. (Reproduïu la proposició~\ref{prop:producte-propietats}.)
\end{exercici}
\begin{solucio}
\(\alpha\cdot(\beta+\gamma)\) són \((\beta+\gamma)\) còpies d'\(\alpha\), és a dir \(\beta\) còpies seguides de \(\gamma\) còpies, que és \(\alpha\cdot\beta+\alpha\cdot\gamma\); l'isomorfisme és el canònic \(A\times(B+C)\cong(A\times B)+(A\times C)\), que conserva l'antilexicogràfic. Per la dreta falla: \((1+1)\cdot\omega=2\cdot\omega=\omega\), però \(1\cdot\omega+1\cdot\omega=\omega+\omega\neq\omega\).
\end{solucio}

\section{Exponenciació ordinal}\label{prac:exponenciacio-ordinal}

\begin{exercici}
\textbf{Càlculs bàsics.} Demostreu que \(\alpha^1=\alpha\), \(\alpha^2=\alpha\cdot\alpha\), \(2^\omega=\omega\) i \(\omega^2=\omega\cdot\omega\). (Reproduïu la proposició~\ref{prop:exponenciacio-calculs}.)
\end{exercici}
\begin{solucio}
Per les equacions recursives, que valen per a \(\alpha>0\) (teorema~\ref{teo:exponenciacio-recursio}): \(\alpha^1=\alpha^{0+1}=\alpha^0\cdot\alpha=1\cdot\alpha=\alpha\); \(\alpha^2=\alpha^1\cdot\alpha=\alpha\cdot\alpha\). Per a \(2^\omega\): és el suprem dels \(2^n\) (\(n<\omega\)), tots finits, i el suprem d'una successió creixent no acotada de tipus finits és \(\omega\); així \(2^\omega=\omega\). I \(\omega^2=\omega^1\cdot\omega=\omega\cdot\omega\). En particular \(2^\omega=\omega\) és numerable, ben diferent del cardinal \(2^{\aleph_0}\), que és el continu.
\end{solucio}

\begin{exercici}
\textbf{La clàusula límit i la base zero.} Comproveu que, si la clàusula límit s'apliqués a la base \(0\), resultaria \(0^\omega=1\), i expliqueu per què la construcció per funcions de suport finit no necessita cap cota externa per als suprems. (Vegeu l'observació~\ref{obs:exponenciacio-acotacio}.)
\end{exercici}
\begin{solucio}
Amb \(0^0=1\) i \(0^{n+1}=0^n\cdot0=0\), la família \((0^n)_{n<\omega}\) és \(1,0,0,\dots\), el suprem de la qual és \(1\); però clàssicament \(0^\omega=\abs{\varnothing^{\N}}\) hauria de ser \(0\). Per això la base \(0\) es defineix a part. Per a \(\alpha>0\), \(\alpha^\beta\) és el tipus del bon ordre \(E(B,A)\) de les funcions de suport finit; si \(B'\) és una part inicial de \(B\), \(E(B',A)\) n'és una part inicial (extensió per \(0_A\)), de manera que totes les potències amb exponent \(\le\operatorname{tipus}(L)\) viuen com a parts inicials d'un sol \(E(L,A)\). Al límit, \(E(B,A)\) és la reunió de les \(E(B_{<b},A)\), i el seu tipus és el suprem: la cota és el mateix \(E(L,A)\), sense \(\Hart\) ni SR.
\end{solucio}

\section{Monotonia i ordre contra mida}\label{prac:ordre-mida}

\begin{exercici}
\textbf{Monotonia estricta per la dreta.} Demostreu que \(\beta<\gamma\) implica \(\alpha+\beta<\alpha+\gamma\), i que per l'esquerra només val \(\le\). (Reproduïu la proposició~\ref{prop:monotonia-ordinal}.)
\end{exercici}
\begin{solucio}
Si \(\beta<\gamma\), un representant de \(\beta\) és un segment inicial propi d'un de \(\gamma\); concatenant \(\alpha\) al davant, \(\alpha+\beta\) resulta un segment inicial propi de \(\alpha+\gamma\), d'on \(\alpha+\beta<\alpha+\gamma\). Per l'esquerra, \(B+A\) \emph{no} és un segment de \(C+A\) (la còpia d'\(A\) va al final), però la inclusió \(B\hookrightarrow C\) seguida de la identitat d'\(A\) és una immersió estrictament creixent \(B+A\to C+A\), i una immersió entre bons ordres compara els tipus (lema~\ref{lema:immersio-ordre}: si \(C+A\) fos més curt, seria isomorf a un segment propi de \(B+A\) i hi hauria una creixent que fa baixar un punt, contra la rigidesa). Així \(\beta+\alpha\le\gamma+\alpha\); l'estrictesa pot fallar, com mostra \(1<2\) amb \(1+\omega=2+\omega=\omega\) (proposició~\ref{prop:suma-no-commutativa}).
\end{solucio}

\begin{exercici}
\textbf{Ordre contra mida.} Demostreu que \(\omega\), \(\omega+1\), \(\omega\cdot2\) i \(\omega^2\) són tipus ordinals diferents però conjunts equipotents. (Model de l'observació~\ref{obs:ordre-vs-mida}.)
\end{exercici}
\begin{solucio}
Són tipus diferents: \(\omega<\omega+1<\omega\cdot2<\omega^2\) per la monotonia (proposició~\ref{prop:monotonia-ordinal}) i els càlculs (proposicions~\ref{prop:suma-no-commutativa} i~\ref{prop:producte-propietats}); en particular \(\omega+1\) té màxim i \(\omega\) no, etc. Però tots són \emph{numerables}: \(\omega+1=\N+\mathbf1\) és numerable (unió de dos numerables); \(\omega\cdot2=\omega+\omega\) també; i \(\omega^2=\omega\cdot\omega\) és numerable perquè \(\N\times\N\) ho és (aparellament, capítol~\ref{cap:infinitud}). Com a \emph{mides} són el mateix cardinal \(\aleph_0\); com a \emph{formes d'ordre}, tots diferents. És la distinció que motiva mantenir separades l'aritmètica ordinal i la cardinal.
\end{solucio}

\chapter{Cardinals i aritmètica cardinal}\label{cap:practica-cardinals}

Aquest capítol acompanya el capítol homònim de la part teòrica (capítol~\ref{cap:cardinals} de la Part~\ref{part:teoria}). Els exercicis resolts en segueixen l'ordre d'exposició ---aritmètica cardinal relacional, comparabilitat, cardinals inicials i König--- i serveixen de model per als exercicis proposats de la Part~\ref{part:exercicis}. Cada bloc fixa una tècnica: definir sobre classes d'equipotència i provar la invariància, obtenir la comparabilitat total del teorema del bon ordre, acotar la cerca del cardinal inicial dins un Hartogs, i reconèixer la diagonal de Cantor dins König. Convé intentar cada enunciat abans de llegir-ne la solució.

\section{Aritmètica cardinal relacional}\label{prac:aritmetica-cardinal}

\begin{exercici}
\textbf{Les operacions estan ben definides.} Demostreu que \(\abs{A}+\abs{B}\), \(\abs{A}\cdot\abs{B}\) i \(\abs{B}^{\abs{A}}\) no depenen dels representants. (Reproduïu la proposició~\ref{prop:operacions-invariants}.)
\end{exercici}
\begin{solucio}
Siguin \(f\colon A\cong A'\) i \(g\colon B\cong B'\). Aleshores \(f+g\colon A+B\cong A'+B'\) (functorialitat de la suma disjunta) i \(f\times g\colon A\times B\cong A'\times B'\) (functorialitat del producte) són isomorfismes. Per a l'exponencial, \(h\mapsto g\circ h\circ f^{-1}\) és una bijecció \(B^A\cong B'^{A'}\), amb inversa \(k\mapsto g^{-1}\circ k\circ f\) (propietat universal, definició~\ref{def:exponencial}). Com que les tres construccions són isomorfismes de conjunts, les mides coincideixen, i les operacions estan ben definides sobre classes d'equipotència.
\end{solucio}

\begin{exercici}
\textbf{Lleis de l'exponenciació cardinal.} Demostreu \(\kappa^{\mu+\nu}=\kappa^\mu\cdot\kappa^\nu\) i \((\kappa^\mu)^\nu=\kappa^{\mu\cdot\nu}\). (Part de la proposició~\ref{prop:lleis-cardinals}.)
\end{exercici}
\begin{solucio}
Amb \(\kappa=\abs{C}\), \(\mu=\abs{B}\), \(\nu=\abs{A}\): \(C^{B+A}\cong C^B\times C^A\), perquè una aplicació d'una suma disjunta \(B+A\to C\) és exactament un parell d'aplicacions \((B\to C,\,A\to C)\) (propietat universal de la suma), d'on \(\kappa^{\mu+\nu}=\kappa^\mu\cdot\kappa^\nu\). I \((C^B)^A\cong C^{B\times A}\) és la currificació (definició~\ref{def:exponencial}): una aplicació \(A\to C^B\) es correspon amb una \(B\times A\to C\), d'on \((\kappa^\mu)^\nu=\kappa^{\mu\cdot\nu}\) (l'ordre \(B\times A\) correspon a \(\mu\cdot\nu\), indiferent per commutativitat). Transportant aquests isomorfismes a mides (proposició~\ref{prop:operacions-invariants}) s'obtenen les dues lleis, amb les lletres dels representants assignades igual que a l'enunciat.
\end{solucio}

\begin{exercici}
\textbf{Cantor cardinal.} Demostreu que \(\abs{A}<2^{\abs{A}}\) per a tot conjunt \(A\). (Reproduïu el teorema~\ref{teo:cantor-cardinal}.)
\end{exercici}
\begin{solucio}
Com que \(P(A)\cong\mathbf2^A\) (capítol~\ref{cap:exponencials}), \(2^{\abs{A}}=\abs{\mathbf2^A}=\abs{P(A)}\). El teorema de Cantor (teorema~\ref{teo:cantor}) diu \(A\prec P(A)\), és a dir \(\abs{A}\le\abs{P(A)}\) i \(\abs{A}\neq\abs{P(A)}\). Per tant \(\abs{A}<2^{\abs{A}}\). Noteu el contrast amb l'ordinal \(2^\omega=\omega\): la mateixa notació, resultats oposats, perquè la cardinal oblida l'ordre i la ordinal el recorda.
\end{solucio}

\section{Comparabilitat}\label{prac:comparabilitat-cardinals}

\begin{exercici}
\textbf{Comparabilitat total.} Amb l'elecció, demostreu que dos conjunts qualssevol són comparables, i deduïu que els cardinals estan totalment ordenats mòdul equipotència. (Reproduïu el teorema~\ref{teo:comparabilitat-cardinals}.)
\end{exercici}
\begin{solucio}
Pel teorema del bon ordre (teorema~\ref{teo:bon-ordre}), \(A\) i \(B\) admeten bons ordres \(W_A\), \(W_B\). Per la comparabilitat de bons ordres (teorema~\ref{teo:comparabilitat}), un és isomorf a un segment inicial de l'altre; aquest isomorfisme d'ordre és un monomorfisme de conjunts, d'on \(A\preceq B\) o \(B\preceq A\). Combinat amb Schröder--Bernstein (teorema~\ref{teo:schroder-bernstein}), que dona l'antisimetria mòdul \(\approx\), i amb la reflexivitat i transitivitat (proposició~\ref{prop:preceq-preordre}), la comparació és un ordre total llevat d'equipotència.
\end{solucio}

\begin{exercici}
\textbf{Comparabilitat equival a elecció.} Demostreu que si dos conjunts qualssevol són comparables, aleshores val el teorema del bon ordre. (Model de l'observació~\ref{obs:comparabilitat-eleccio}.)
\end{exercici}
\begin{solucio}
Sigui \(A\) qualsevol. Per hipòtesi, \(A\) i \(\Hart(A)\) són comparables. Com que \(\Hart(A)\not\preceq A\) (teorema~\ref{teo:hartogs}), l'única possibilitat és \(A\preceq\Hart(A)\): hi ha un monomorfisme \(A\hookrightarrow\Hart(A)\). Com que \(\Hart(A)\) és un bon ordre (proposició~\ref{prop:hartogs-bon-ordre}), la imatge d'\(A\) hi és una part, que hereta un bon ordre; transportant-lo, \(A\) admet un bon ordre. Així la comparabilitat implica el teorema del bon ordre, equivalent a l'elecció (teorema~\ref{teo:equivalencies}).
\end{solucio}

\section{Cardinals inicials}\label{prac:cardinals-inicials}

\begin{exercici}
\textbf{Existència del cardinal inicial.} Amb l'elecció, demostreu que tot conjunt té un únic cardinal inicial, trobat acotadament. (Reproduïu la proposició~\ref{prop:cardinal-inicial}.)
\end{exercici}
\begin{solucio}
Treballem dins \(\Hart(A)\). La part \(E\) dels tipus \([(S,R)]\) amb \(S\approx A\) és no inicial: pel teorema del bon ordre, \(A\) admet un bon ordre \(R\), i \([(A,R)]\in E\). Com que \(\Hart(A)\) és un bon ordre (proposició~\ref{prop:hartogs-bon-ordre}), \(E\) té un mínim \(\alpha\), únic. És inicial: un bon ordre \(V\) de tipus \(<\alpha\) és isomorf a un segment propi d'un representant \((S,R)\) de \(\alpha\), és a dir, a un bon ordre sobre una part d'\(A\); si fos equipotent a \(A\), el seu tipus seria un element de \(E\) menor que \(\alpha\). La cerca no ha sortit de \(\Hart(A)\).
\end{solucio}

\begin{exercici}
\textbf{El cardinal és un tipus, no un conjunt escollit.} Expliqueu per què és més net dir que el cardinal inicial és un tipus d'ordre inicial que no pas \enquote{un conjunt}. (Model de l'observació~\ref{obs:cardinal-prudencia}.)
\end{exercici}
\begin{solucio}
El cardinal inicial de \(A\) és el menor tipus d'ordre de la seva mida, únic llevat d'un únic isomorfisme d'ordre (teorema~\ref{teo:unicitat-iso}). Dir \enquote{és un conjunt} suggeriria una tria canònica d'un bon ordre concret d'aquell tipus, que no tenim: escollir-ne un representant és una tria addicional, legítima però no canònica. Per això la teoria relacional treballa amb \(\abs{A}\) com a classe d'equipotència i només introdueix representants qui vulgui una successió distingida. A més, reunir \emph{tots} els cardinals inicials en una successió global seria de Nivell~II, i requeriria SR (observació~\ref{obs:sr}).
\end{solucio}

\section{El teorema de König}\label{prac:konig}

\begin{exercici}
\textbf{El teorema de König.} Amb l'elecció, demostreu que \(A_i\prec B_i\) per a tot \(i\) implica \(\sum_i\abs{A_i}<\prod_i\abs{B_i}\). (Reproduïu el teorema~\ref{teo:konig}.)
\end{exercici}
\begin{solucio}
\emph{\(\le\):} triant (elecció) un monomorfisme \(m_i\colon A_i\hookrightarrow B_i\) i un \(b_i\notin\operatorname{Im}(m_i)\) (existeix perquè \(A_i\not\approx B_i\)), enviem un \(a\) de la fibra \(A_i\) a la família amb \(x_i=m_i(a)\) i \(x_j=b_j\) per \(j\ne i\); és injectiva. \emph{\(\neq\):} donada \(F\colon\coprod_i A_i\to\prod_i B_i\), per a cada \(i\) l'aplicació \(F_i=\pi_i\circ F|_{A_i}\colon A_i\to B_i\) no és exhaustiva (car \(A_i\prec B_i\)), i triem \(d_i\notin\operatorname{Im}(F_i)\). La família \((d_i)_i\) no és cap valor de \(F\): si \(F(a)=(d_i)_i\) amb \(a\) a la fibra \(A_{i_0}\), la coordenada \(i_0\) donaria \(F_{i_0}(a)=d_{i_0}\), impossible. Per tant \(F\) no és exhaustiva i \(\sum_i\abs{A_i}\neq\prod_i\abs{B_i}\); amb \(\le\), l'estricta.
\end{solucio}

\begin{exercici}
\textbf{König conté Cantor.} Deduïu el teorema de Cantor \(\abs{A}<2^{\abs{A}}\) com a cas particular de König. (Model de l'observació~\ref{obs:konig-cantor}.)
\end{exercici}
\begin{solucio}
Preneu \(I=A\), \(A_i=\mathbf1\) i \(B_i=\mathbf2\) per a tot \(i\). Aleshores \(A_i=\mathbf1\prec\mathbf2=B_i\), i König dona \(\sum_{i\in A}1<\prod_{i\in A}2\), és a dir \(\abs{A}<2^{\abs{A}}\), que és Cantor cardinal (teorema~\ref{teo:cantor-cardinal}). La família diagonal \((d_i)\) de la prova de König és, en aquest cas, exactament la part diagonal de Cantor: per a cada \(i\), \(d_i\) és l'element de \(\mathbf2\) diferent del que \(F\) assigna, és a dir la negació de la pertinença. König és, doncs, la forma indexada de la diagonal.
\end{solucio}

\chapter{Alephs, cofinalitat i continu}\label{cap:practica-alephs}

Aquest capítol acompanya el capítol homònim de la part teòrica (capítol~\ref{cap:alephs} de la Part~\ref{part:teoria}), que és opcional i de Nivell~II. Els exercicis resolts en segueixen l'ordre d'exposició ---\(\aleph_0\) i el cardinal successor, la successió dels alephs, la cofinalitat, i els beths i el continu--- i serveixen de model per als exercicis proposats de la Part~\ref{part:exercicis}. Cada bloc distingeix amb cura el que és de Nivell~I (el cardinal successor, acotat per Hartogs, i també el continu \(\mathfrak c=\abs{P(\N)}\) i la mateixa hipòtesi del continu) del que demana l'esquema de reunió de famílies (els trams de la successió dels alephs i dels beths com a famílies). Convé intentar cada enunciat abans de llegir-ne la solució.

\section{\texorpdfstring{\(\aleph_0\)}{aleph 0} i el cardinal successor}\label{prac:aleph0-successor}

\begin{exercici}
\textbf{El cardinal successor existeix sense SR.} Demostreu que per a tot cardinal \(\kappa\) existeix \(\kappa^+\), acotat per Hartogs, sense usar l'esquema de reunió. (Reproduïu la proposició~\ref{prop:cardinal-successor}.)
\end{exercici}
\begin{solucio}
Sigui \(\kappa=\abs{K}\). El bon ordre \(\Hart(K)\) no s'injecta en \(K\) (teorema~\ref{teo:hartogs}), de manera que \(\abs{\Hart(K)}>\kappa\): existeix algun cardinal més gran. De fet, \(\kappa^+=\abs{\Hart(K)}\) (proposició~\ref{prop:cardinal-successor}): si \(\abs{X}>\kappa\), el teorema del bon ordre dona un bon ordre \(W\) sobre \(X\); per la comparabilitat, o bé \(W\) és isomorf a un segment propi de \(\Hart(K)\), cosa impossible perquè aquests segments representen bons ordres sobre parts de \(K\) i aleshores \(X\preceq K\), o bé \(\Hart(K)\) s'immergeix en \(W\), és a dir, \(\abs{\Hart(K)}\le\abs{X}\). I \(\Hart(K)\) és inicial: un tipus menor és el d'un segment propi, que s'injecta en \(K\). Tota la cerca ha passat dins un \(\Hart\), acotadament; no s'ha reunit cap família no acotada, de manera que no cal l'esquema de reunió (principi~\ref{sr:esquema}).
\end{solucio}

\begin{exercici}
\textbf{Un pas contra un tram de la successió.} Expliqueu per què \(\aleph_0\) i cada \(\kappa^+\) individual són de Nivell~I, però el tram \((\aleph_n)_{n<\omega}\) com a família demana l'esquema de reunió. (Model de l'observació~\ref{obs:successor-sense-sr}.)
\end{exercici}
\begin{solucio}
Cada cardinal successor \(\kappa^+\) s'obté acotadament per Hartogs (exercici anterior): és una operació que produeix \emph{un} cardinal a partir d'\emph{un} altre, dins un ambient ja construït. En canvi, reunir el tram \(\aleph_0,\aleph_1,\aleph_2,\dots\) indexat per \(\N\) en una família ---i, per a trams més llargs, prolongar-lo als límits--- és passar d'una regla externa \(n\mapsto\aleph_n\) a una família interna \(E\to\N\); i com que els \(\aleph_n\) creixen sense una cota comuna donada d'entrada, aquesta reunió no es pot fer amb els axiomes de Nivell~I i requereix l'esquema de reunió (principi~\ref{sr:esquema}). És la frontera de sempre: un pas és acotat; la successió transfinita no acotada, no.
\end{solucio}

\section{La successió dels alephs}\label{prac:alephs}

\begin{exercici}
\textbf{Els alephs són estrictament creixents.} Sota l'esquema de reunió, demostreu que \(\alpha<\beta\) implica \(\aleph_\alpha<\aleph_\beta\). (Part de la proposició~\ref{prop:alephs}.)
\end{exercici}
\begin{solucio}
Per inducció transfinita (teorema~\ref{teo:induccio-transfinita}). Al successor, \(\aleph_{\alpha+1}=(\aleph_\alpha)^+>\aleph_\alpha\) per definició del cardinal successor. Al límit \(\lambda\), \(\aleph_\lambda=\sup_{\beta<\lambda}\aleph_\beta\); per a \(\gamma<\lambda\), agafant qualsevol \(\beta\) amb \(\gamma<\beta<\lambda\) es té \(\aleph_\gamma<\aleph_\beta\le\aleph_\lambda\), i de fet \(\aleph_\lambda>\aleph_\gamma\) perquè el suprem d'una successió estrictament creixent no acotada supera cada terme. Així la successió és estrictament creixent en tot el seu recorregut.
\end{solucio}

\begin{exercici}
\textbf{Tot cardinal infinit és un aleph.} Sota elecció i l'esquema de reunió, demostreu que per a tot conjunt infinit \(A\) existeix un únic \(\alpha\) amb \(\abs{A}=\aleph_\alpha\). (Part de la proposició~\ref{prop:alephs}.)
\end{exercici}
\begin{solucio}
Per l'elecció, \(A\) té un cardinal inicial \(\kappa\) (proposició~\ref{prop:cardinal-inicial}). La successió dels alephs \(\le\kappa\) és estrictament creixent (exercici anterior) i acotada per \(\kappa\). Sigui \(\alpha\) el menor tipus amb \(\aleph_\alpha\ge\kappa\); com que cada aleph és el menor cardinal més gran que l'anterior (al successor) o el suprem (al límit), la minimalitat de \(\alpha\) força \(\aleph_\alpha=\kappa\), és a dir \(\abs{A}=\aleph_\alpha\). La unicitat és l'estricta monotonia. Sense elecció, en canvi, hi pot haver conjunts infinits no equipotents a cap aleph, de manera que la hipòtesi és essencial.
\end{solucio}

\section{Cofinalitat}\label{prac:cofinalitat}

\begin{exercici}
\textbf{\(\aleph_\omega\) és singular.} Demostreu que \(\operatorname{cf}(\aleph_\omega)=\aleph_0<\aleph_\omega\). (Part de la proposició~\ref{prop:regular-exemples}.)
\end{exercici}
\begin{solucio}
La successió \((\aleph_n)_{n<\omega}\) és estrictament creixent (proposició~\ref{prop:alephs}) i el seu suprem és \(\aleph_\omega\) per definició (definició~\ref{def:alephs}). Per tant el conjunt \(\{\aleph_n:n<\omega\}\) és cofinal en \(\aleph_\omega\): arriba fins al final, perquè qualsevol cardinal \(<\aleph_\omega\) és superat per algun \(\aleph_n\). Aquesta part cofinal té tipus d'ordre \(\omega\), és a dir cardinal \(\aleph_0\). Com que \(\aleph_0<\aleph_\omega\), la cofinalitat de \(\aleph_\omega\) és estrictament menor que ell mateix: \(\aleph_\omega\) és singular (definició~\ref{def:regular-singular}).
\end{solucio}

\begin{exercici}
\textbf{Els successors són regulars.} Demostreu que tot \(\aleph_{\alpha+1}\) és regular. (Part de la proposició~\ref{prop:regular-exemples}.)
\end{exercici}
\begin{solucio}
Suposem una part cofinal \(C\) de \(\aleph_{\alpha+1}\) de cardinal \(\le\aleph_\alpha\). Aleshores \(\aleph_{\alpha+1}\) seria la unió dels segments determinats pels elements de \(C\), és a dir una unió de \(\le\aleph_\alpha\) conjunts, cadascun de cardinal \(<\aleph_{\alpha+1}\), és a dir \(\le\aleph_\alpha\). Però una tal unió té cardinal \(\le\aleph_\alpha\cdot\aleph_\alpha=\aleph_\alpha<\aleph_{\alpha+1}\) (teorema de Hessenberg, \(\kappa\cdot\kappa=\kappa\), citat al teorema~\ref{teo:hessenberg}; la reunió usa l'elecció per triar les injeccions), contradicció. Per tant tota part cofinal té cardinal \(\aleph_{\alpha+1}\), i \(\operatorname{cf}(\aleph_{\alpha+1})=\aleph_{\alpha+1}\): el successor és regular.
\end{solucio}

\section{Beths i el continu}\label{prac:beths-continu}

\begin{exercici}
\textbf{Els beths dominen els alephs.} Sota l'esquema de reunió, demostreu que \(\aleph_\alpha\le\beth_\alpha\) per a tot \(\alpha\). (Reproduïu la proposició~\ref{prop:beths}.)
\end{exercici}
\begin{solucio}
Per inducció transfinita (teorema~\ref{teo:induccio-transfinita}). A \(\alpha=0\), \(\aleph_0=\beth_0\). Al successor, \(\aleph_{\alpha+1}=(\aleph_\alpha)^+\le 2^{\aleph_\alpha}\le 2^{\beth_\alpha}=\beth_{\alpha+1}\): la primera desigualtat perquè \((\aleph_\alpha)^+\) és el menor cardinal \(>\aleph_\alpha\) i \(2^{\aleph_\alpha}>\aleph_\alpha\) per Cantor (teorema~\ref{teo:cantor-cardinal}); la segona per monotonia de l'exponenciació cardinal i la hipòtesi d'inducció. Al límit, \(\aleph_\lambda=\sup_{\beta<\lambda}\aleph_\beta\le\sup_{\beta<\lambda}\beth_\beta=\beth_\lambda\). Així \(\aleph_\alpha\le\beth_\alpha\) sempre.
\end{solucio}

\begin{exercici}
\textbf{La hipòtesi del continu, ben formulada.} Enuncieu la hipòtesi del continu i expliqueu per què és una qüestió independent i no un teorema pendent. (Model de les observacions~\ref{obs:hipotesi-continu} i~\ref{obs:hc-independencia}.)
\end{exercici}
\begin{solucio}
El continu és \(\mathfrak c=2^{\aleph_0}=\beth_1\). Per Cantor, \(\mathfrak c>\aleph_0\) (teorema~\ref{teo:cantor-cardinal}), d'on \(\mathfrak c\ge\aleph_1\) sota elecció. La hipòtesi del continu (HC) afirma \(2^{\aleph_0}=\aleph_1\): que no hi ha cap cardinal entre \(\aleph_0\) i \(\mathfrak c\), és a dir que la jerarquia dels successors (\(\aleph\)) i la dels conjunts de parts (\(\beth\)) coincideixen al primer graó. No és un teorema pendent de demostració elemental: els resultats de Gödel i Cohen mostren que, \emph{si ZFC és consistent}, ni HC ni la seva negació no es dedueixen de ZFC; per la bi-interpretabilitat d'ETCS\({}+{}\)SR amb ZFC, tampoc no es dedueixen dels axiomes d'aquest llibre amb SR. És independència \emph{relativa a la consistència}, i és un metateorema citat, no un resultat del llibre. La teoria formula HC amb precisió ---distingint \(\aleph_1\) de \(\beth_1\)--- però no en determina la resposta; el que Cantor estableix és només la desigualtat \(\aleph_0<2^{\aleph_0}\), no on cau exactament \(2^{\aleph_0}\).
\end{solucio}


% ============================================================
% PART III — EXERCICIS PROPOSATS
% ============================================================
\part{Exercicis proposats}\label{part:exercicis}
\setcounter{chapter}{0}
\chapter{Aplicacions, dominis, codominis i composició}\label{cap:exercicis-aplicacions}

Aquests exercicis proposats corresponen al capítol homònim de la part teòrica (capítol~\ref{cap:aplicacions} de la Part~\ref{part:teoria}) i es resolen amb les tècniques dels exercicis resolts de la part pràctica (capítol~\ref{cap:practica-aplicacions} de la Part~\ref{part:practica}). Els marcats amb una indicació són els més exigents.

\section{Llenguatge formal, composició i estabilitat}

\begin{enumerate}
  \item Escriviu amb els predicats \(\Dom\), \(\Cod\) i \(\Com\) que \(h\) és la composició de \(f\) seguida de \(g\), que el domini de \(h\) coincideix amb el de \(f\) i que el codomini de \(h\) coincideix amb el de \(g\).

  \item Traduïu al llenguatge dels predicats la frase \enquote{\(f\) i \(g\) tenen el mateix codomini}. Feu-ho de dues maneres: amb una variable auxiliar per al codomini comú i sense.

  \item Siguin \(f\colon A\to B\), \(g\colon B\to C\) i \(h\colon C\to D\). Escriviu les dues maneres d'agrupar \(h\circ g\circ f\) i justifiqueu-ne la igualtat a partir de l'axioma d'associativitat.

  \item Demostreu que si \(\Com(f,g,h)\) i \(\Com(f,g,h')\), aleshores \(h=h'\). \begin{indicacio}És la unicitat de la composició (proposició~\ref{prop:composicio-unica}); cap raonament nou, només citar-la correctament.\end{indicacio}
\end{enumerate}

\section{Identitats i conjunts}

\begin{enumerate}[resume]
  \item Demostreu que per a tot conjunt \(A\) es té \(\id_A\circ\id_A=\id_A\), i que \(\id_A\) és l'única aplicació \(A\to A\) amb aquesta propietat que a més és neutra per composició.

  \item Sigui \(e\colon A\to A\) una aplicació amb \(e\circ e=e\) (un \emph{idempotent}). Doneu, en el model clàssic, un exemple d'un idempotent que no sigui la identitat. \begin{indicacio}Una aplicació constant \(A\to A\) amb valor un punt fix serveix; comproveu \(e\circ e=e\).\end{indicacio}
\end{enumerate}

\section{Igualtat d'aplicacions i elements generalitzats}

\begin{enumerate}[resume]
  \item Siguin \(f,g\colon A\to B\). Demostreu que si \(f=g\), aleshores \(f\circ x=g\circ x\) per a tot element generalitzat \(x\colon X\to A\). (És la direcció fàcil del criteri d'igualtat.)

  \item Expliqueu amb un exemple del model clàssic per què la separació per elements \emph{globals} no seria suficient si no la postuléssim: descriviu dos conjunts i dues aplicacions paral·leles que un observador que només mirés punts de \(\mathbf1\) no podria distingir sense l'axioma. \begin{indicacio}Penseu en què passa quan el domini no té elements globals, com \(\mathbf0\); relacioneu-ho amb l'observació~\ref{obs:separacio-no-degeneracio}.\end{indicacio}
\end{enumerate}

\section{Isomorfismes}

\begin{enumerate}[resume]
  \item Sigui \(f\colon A\to B\) un isomorfisme. Demostreu que \({(f^{-1})}^{-1}=f\).

  \item Demostreu que la isomorfia es comporta transitivament: de \(A\cong B\) i \(B\cong C\), deduïu \(A\cong C\). Indiqueu quin isomorfisme concret exhibeix \(A\cong C\).

  \item Demostreu que \(\cong\) és una relació d'equivalència sobre els conjunts (reflexiva, simètrica i transitiva), tot dient quins resultats del capítol useu a cada propietat.

  \item Siguin \(f,g\colon A\to B\) isomorfismes amb \(f^{-1}=g^{-1}\). Demostreu que \(f=g\). \begin{indicacio}Componeu la igualtat de les inverses convenientment amb \(f\) o \(g\) i useu la neutralitat.\end{indicacio}
\end{enumerate}

\section{Terminals, inicials i no-degeneració}

\begin{enumerate}[resume]
  \item Demostreu que si \(I_1\) i \(I_2\) són tots dos inicials, hi ha un únic isomorfisme \(I_1\to I_2\). (Dual de la unicitat del terminal.)

  \item Demostreu que tota aplicació \(A\to\mathbf1\) és \({!}_A\). Deduïu-ne que, per a cada estadi \(X\), el conjunt \(\mathbf1\) té exactament un element generalitzat en l'estadi \(X\), i compareu-ho amb el nombre d'aplicacions \(\mathbf1\to A\) per a un conjunt \(A\) qualsevol. \begin{indicacio}El que és únic és l'aplicació \emph{cap} a \(\mathbf1\); les aplicacions \emph{des de} \(\mathbf1\) són els punts d'\(A\), i poden ser molts o cap.\end{indicacio}

  \item Demostreu que no existeix cap aplicació \(\mathbf1\to\mathbf0\) reproduint l'argument del resolt corresponent, i deduïu que \(\mathbf0\) no té cap element global.

  \item Suposeu un univers alternatiu en què se satisfan tots els axiomes tret de la no-degeneració, i en què \(\mathbf0\cong\mathbf1\). Demostreu que aleshores entre cada parell de conjunts hi ha exactament una aplicació. \begin{indicacio}Vegeu l'observació~\ref{obs:zero-u-sense-separacio}; useu que tot conjunt rebria i emetria una única fletxa cap a i des de l'objecte col·lapsat.\end{indicacio}

  \item Escriviu l'enunciat dual de la proposició~\ref{prop:no-fletxa-u-zero} i digueu si és cert. \begin{indicacio}El dual és la mateixa proposició: la no-degeneració és autodual. Compareu-ho amb l'exercici resolt sobre la no-autodualitat de la separació global.\end{indicacio}

  \item \textbf{Axiomes usats.} Per a cadascun dels resultats següents, digueu quins dels nou axiomes del capítol intervenen a la demostració i quins no hi intervenen: la unicitat de la composició (proposició~\ref{prop:composicio-unica}), la unicitat de la inversa (proposició~\ref{prop:inversa-unica}), la unicitat del terminal llevat d'isomorfisme (proposició~\ref{prop:terminal-unic}), la inexistència d'aplicacions \(\mathbf1\to\mathbf0\) (proposició~\ref{prop:no-fletxa-u-zero}) i l'extensionalitat global (proposició~\ref{prop:extensionalitat-global}). \begin{indicacio}Les tres primeres només usen els axiomes estructurals i, si escau, les definicions; les dues últimes són les úniques que necessiten la no-degeneració i la separació global, respectivament.\end{indicacio}
\end{enumerate}

\section{Elements globals i aplicacions constants}

\begin{enumerate}[resume]
  \item Sigui \(b\colon\mathbf1\to B\) un punt i \(c\colon A\to B\) la constant de valor \(b\). Demostreu que si \(A\) té almenys un element global, aleshores \(c\) queda determinada pel seu valor sobre qualsevol d'ells.

  \item Siguin \(b,b'\colon\mathbf1\to B\) dos punts i \(c,c'\colon A\to B\) les constants respectives. Demostreu que \(c=c'\) si i només si \(b=b'\), suposant que \(A\) no és inicial. \begin{indicacio}Per a una direcció componeu amb un element global de \(A\); l'altra és immediata. On s'usa que \(A\) no és inicial?\end{indicacio}
\end{enumerate}

\chapter{Parts i construccions universals}\label{cap:exercicis-parts}

Aquests exercicis proposats corresponen al capítol homònim de la part teòrica (capítol~\ref{cap:parts} de la Part~\ref{part:teoria}) i es resolen amb les tècniques dels exercicis resolts de la part pràctica (capítol~\ref{cap:practica-parts} de la Part~\ref{part:practica}). Cada secció remet al bloc de resolts que en dona el model. Els marcats amb una indicació són els més exigents.

\section{Cancel·lació i composició de mono/epi}

\begin{enumerate}
  \item Demostreu que si \(m:A\rightarrowtail B\) és un monomorfisme i \(u:X\to A\) és qualsevol, aleshores \(m\circ u=m\circ u'\) implica \(u=u'\); enuncieu i demostreu la propietat dual per a un epimorfisme. (Model: secció~\ref{prac:cancel-lacio}.)

  \item Siguin \(m:A\rightarrowtail B\) i \(n:B\rightarrowtail C\) monomorfismes. Demostreu que \(n\circ m\) és un monomorfisme i, dualment, que la composició de dos epimorfismes és un epimorfisme. \begin{indicacio}Descomponeu \((n\circ m)\circ u\) amb l'associativitat i cancel·leu primer \(n\), després \(m\), com a la proposició~\ref{prop:composicio-mono-epi}.\end{indicacio}

  \item Siguin \(f:A\to B\) i \(g:B\to C\). Demostreu que si \(g\circ f\) és un epimorfisme, aleshores \(g\) és un epimorfisme. (Dual de l'exercici resolt de cancel·lació posterior.)

  \item Doneu, en el model clàssic, un exemple de \(f\) i \(g\) amb \(g\circ f\) mono però \(g\) no mono, per veure que la conclusió de l'exercici anterior no es pot enfortir a \enquote{\(g\) és mono}. \begin{indicacio}Preneu \(g\) no injectiva però injectiva sobre la imatge de \(f\).\end{indicacio}
\end{enumerate}

\section{Seccions, retraccions i exhaustivitat generalitzada}

\begin{enumerate}[resume]
  \item Demostreu que tota retracció és un epimorfisme (la meitat de la proposició~\ref{prop:seccio-mono-retraccio-epi} que no s'ha resolt al text). \begin{indicacio}Dualitzeu l'argument de \enquote{tota secció és un monomorfisme}: component per la dreta amb la secció.\end{indicacio}

  \item Sigui \(e:A\to B\) un epimorfisme escindit amb secció \(s\). Demostreu que \(s\) és un monomorfisme i que \(e\circ s=\id_B\) caracteritza \(s\) entre les possibles seccions. 

  \item Comproveu que, en el model clàssic, una aplicació exhaustiva sobre els elements globals no és necessàriament una retracció si no es disposa d'una elecció de preimatges; relacioneu-ho amb l'observació~\ref{obs:exhaustiva-generalitzada-global}. \begin{indicacio}El pas de \enquote{exhaustiva} a \enquote{escindida} requereix triar, per a cada element de \(B\), una preimatge; sense elecció, la secció pot no existir.\end{indicacio}
\end{enumerate}

\section{Parts, inclusió i pertinença}

\begin{enumerate}[resume]
  \item Demostreu que la relació d'inclusió \(\subseteq_A\) és reflexiva i transitiva sobre les parts de \(A\), i que \(m\subseteq_A n\) i \(n\subseteq_A m\) impliquen que \(m\) i \(n\) són parts equivalents. (Model: secció~\ref{prac:parts-inclusio}; per a l'antisimetria vegeu la proposició~\ref{prop:equivalencia-parts}.)

  \item Siguin \(m,n\) parts de \(A\). Demostreu la implicació que falta a la proposició~\ref{prop:inclusio-pertinenca}: si per a tot \(x:X\to A\) es té \(x\in_X m\Rightarrow x\in_X n\), aleshores \(m\subseteq_A n\). \begin{indicacio}Apliqueu la hipòtesi a l'element generalitzat \enquote{genèric} \(x=m:S\to A\), que pertany a \(m\) via \(\id_S\); el testimoni resultant és el \(k\) buscat.\end{indicacio}

  \item Demostreu que si \(x\in_X m\) i \(u:Y\to X\) és qualsevol, aleshores \(x\circ u\in_Y m\). \begin{indicacio}Si \(\bar x\) testimonia \(x\in_X m\), comproveu que \(\bar x\circ u\) testimonia \(x\circ u\in_Y m\).\end{indicacio}
\end{enumerate}

\section{Sumes, característiques i \texorpdfstring{\(\mathbf2\)}{2}}

\begin{enumerate}[resume]
  \item Donades \(f:A\to C\) i \(g:B\to C\), demostreu que l'aplicació definida per casos \([f,g]:A+B\to C\) és única amb les propietats \([f,g]\circ i_A=f\) i \([f,g]\circ i_B=g\). (Model: exercici resolt d'unicitat per casos, secció~\ref{prac:sumes}.)

  \item Demostreu que, si existeixen aplicacions \(f\colon A\to B\) i \(g\colon B\to A\), les dues injeccions \(i_0,i_1\) de la suma \(A+B\) són monomorfismes escindits. \begin{indicacio}Construïu una retracció de cada injecció com al model del cograf (proposició~\ref{prop:cograf-retraccio}); una secció és mono. Que les injeccions siguin sempre monomorfismes, sense aquesta hipòtesi, es veurà al capítol~\ref{cap:valors}.\end{indicacio}

  \item Demostreu que dues parts de \(A\) amb la mateixa aplicació característica \(\chi:A\to\mathbf2\) són equivalents, i recíprocament que parts equivalents tenen la mateixa característica. (Model: secció~\ref{prac:caracteristica}; useu l'axioma~\ref{ax:representacio}.)

  \item Comproveu que la característica de la part total \(\id_A\) és l'aplicació constant \(1\circ{!}_A\), i demostreu que una part \(m\colon S\rightarrowtail A\) amb \(\chi_m=1\circ{!}_A\) és equivalent a la total i, per tant, \(m\) és un isomorfisme. \begin{indicacio}Per a la segona part, \(\id_A\in_A m\) dona una secció de \(m\); combineu-ho amb el fet que \(m\) és mono, com a la proposició~\ref{prop:bimorfisme-isomorfisme}.\end{indicacio}
\end{enumerate}

\section{Productes, equalitzadors i pullbacks}

\begin{enumerate}[resume]
  \item Demostreu que si \(p_A\circ u=p_A\circ v\) i \(p_B\circ u=p_B\circ v\), aleshores \(u=v\), i deduïu que \(\langle p_A,p_B\rangle=\id_{A\times B}\). (Model: secció~\ref{prac:productes}.)

  \item Demostreu que el producte de \(A\) i \(B\) és únic llevat d'isomorfisme, exhibint els isomorfismes inversos entre dues realitzacions qualssevol. \begin{indicacio}Seguiu el patró de la proposició~\ref{prop:producte-unic}: induïu \(\varphi,\psi\) per les propietats universals i useu la unicitat de factorització per veure \(\psi\circ\varphi=\id\).\end{indicacio}

  \item Demostreu que l'equalitzador \(e:E\rightarrowtail A\) de \(f,g:A\to B\) és un monomorfisme, i que qualsevol altre equalitzador dels mateixos \(f,g\) hi és equivalent com a part. (Model: secció~\ref{prac:equalitzadors-pullbacks}; enunciat: teorema~\ref{teo:equalitzador-universal}.)

  \item En un pullback de \(f:A\to C\) al llarg de \(g:B\to C\), demostreu que si \(f\) és un monomorfisme, la projecció oposada és un monomorfisme, i interpreteu-ho com que la imatge inversa d'una part és una part. \begin{indicacio}Vegeu la proposició~\ref{prop:pullback-mono}; cancel·leu \(f\) i useu la unicitat de la factorització del pullback.\end{indicacio}

  \item Siguin \(g:B\to A\) i una part \(m:S\rightarrowtail A\). Demostreu que \(y\in_X g^{-1}(m)\) si i només si \(g\circ y\in_X m\), i deduïu-ne que \(g^{-1}(\id_A)\equiv_B\id_B\) i que, si \(m\subseteq_A n\), aleshores \(g^{-1}(m)\subseteq_B g^{-1}(n)\). \begin{indicacio}La primera part és la proposició~\ref{prop:pertinenca-imatge-inversa}; per a les altres dues, compareu els elements generalitzats de les parts i useu la proposició~\ref{prop:inclusio-pertinenca}.\end{indicacio}

  \item Siguin \(h\colon C\to B\), \(g\colon B\to A\) i una part \(m\) de \(A\). Demostreu que \((g\circ h)^{-1}(m)\equiv_C h^{-1}\bigl(g^{-1}(m)\bigr)\). \begin{indicacio}Les dues parts tenen els mateixos elements generalitzats: apliqueu dues vegades la proposició~\ref{prop:pertinenca-imatge-inversa}. Alternativament, useu que dos quadrats cartesians encaixats formen un quadrat cartesià.\end{indicacio}

  \item \textbf{Axiomes usats.} Digueu quins axiomes nous del capítol (sumes, representació, productes) intervenen en cadascun d'aquests resultats, i quins es demostren només amb els del capítol~\ref{cap:aplicacions}: la composició de monomorfismes és un monomorfisme; tota secció és un monomorfisme; el testimoni d'una inclusió és únic; \(0,1\colon\mathbf1\to\mathbf2\) són monomorfismes; tot bimorfisme és un isomorfisme; l'equalitzador existeix; el pullback existeix. \begin{indicacio}Els tres primers només usen els axiomes estructurals i les definicions; el quart usa les sumes; el cinquè, la representació; els dos últims, productes i representació.\end{indicacio}
\end{enumerate}

\chapter[Valors de veritat i construccions finites]{Valors de veritat, parts\texorpdfstring{\\}{ }i construccions finites}\label{cap:exercicis-valors}

Aquests exercicis proposats corresponen al capítol homònim de la part teòrica (capítol~\ref{cap:valors} de la Part~\ref{part:teoria}) i es resolen amb les tècniques dels exercicis resolts de la part pràctica (capítol~\ref{cap:practica-valors} de la Part~\ref{part:practica}). Cada secció remet al bloc de resolts que en dona el model. Els marcats amb una indicació són els més exigents.

\section{Sumes disjuntes i estables}

\begin{enumerate}
  \item Demostreu que una aplicació \(m\) és un monomorfisme si i només si és un pullback d'ella mateixa amb projeccions identitat, i useu-ho per reescriure la definició de suma disjunta (definició~\ref{def:suma-disjunta}) només en termes de pullbacks. (Model: secció~\ref{prac:sumes-disjuntes}.)

  \item Enuncieu amb precisió què vol dir que una suma és \emph{estable} per canvi de base (definició~\ref{def:suma-estable}) i expliqueu per què l'estabilitat sola no implica la disjunció. Indiqueu, per a cadascun dels resultats següents, si usa la disjunció, l'estabilitat o totes dues: l'estrictesa de \(\mathbf0\), la distinció \(0\ne1\), la descomposició per canvi de base. \begin{indicacio}El canvi de base al llarg de la identitat no dona cap informació nova. L'estrictesa usa l'estabilitat i després la disjunció; \(0\ne1\) només la disjunció; la descomposició, totes dues.\end{indicacio}

  \item Sigui \(A_0\xrightarrow{i_0}A_0+A_1\xleftarrow{i_1}A_1\) una suma i \(f:X\to A_0+A_1\) qualsevol. Descriviu la suma \(X_0\to X\leftarrow X_1\) que s'obté per canvi de base al llarg de \(f\) i interpreteu-la com una descomposició de \(X\) en dues parts complementàries. \begin{indicacio}És el teorema~\ref{teo:descomposicio-canvi-base}; els \(X_j\) són els pullbacks de les injeccions al llarg de \(f\).\end{indicacio}
\end{enumerate}

\section{Inicial estricte i distinció dels valors de veritat}

\begin{enumerate}[resume]
  \item Demostreu que tota aplicació \(f:X\to\mathbf0\) és un isomorfisme. (Model: secció~\ref{prac:zero-u}; enunciat: proposició~\ref{prop:inicial-estricte}.)

  \item Demostreu el lema del buit: un conjunt sense elements globals és inicial. \begin{indicacio}Compareu \(0\circ{!}_X\) i \(1\circ{!}_X\) i useu l'axioma de separació global (axioma~\ref{ax:separacio-globals}), com al model.\end{indicacio}

  \item Demostreu que \(0\neq1\) com a elements de \(\mathbf2\), i deduïu del corol·lari~\ref{cor:elements-de-2} que \(\mathbf2\) té exactament dos elements globals. \begin{indicacio}Suposeu \(0=1\); un punt \(\mathbf1\to\mathbf2\) és mono, i el pullback d'un mono amb ell mateix és el seu domini. Arribeu a \(\mathbf1\cong\mathbf0\).\end{indicacio}

  \item Demostreu que tot element global d'una suma \(A_0+A_1\) es factoritza a través d'exactament una de les dues injeccions. (Enunciat: teorema~\ref{teo:exhaustivitat-sumes}.)
\end{enumerate}

\section{Parts extremals}

\begin{enumerate}[resume]
  \item Demostreu que la característica de la part total \(\id_A\) és la constant \(1\), i que la de la part buida \(\mathbf0\rightarrowtail A\) és la constant \(0\). (Model: secció~\ref{prac:extremals}; proposició~\ref{prop:caracteristica-buida}.)

  \item Demostreu que dues parts totals de \(A\) són equivalents i que dues parts buides de \(A\) són equivalents, de manera que \enquote{la} part total i \enquote{la} part buida queden ben definides. (Enunciat: proposició~\ref{prop:parts-extremals-uniques}.)

  \item Verifiqueu, a partir del quadrat característic (teorema~\ref{teo:quadrat-caracteristic}), que la part buida és el pullback de \(0\) al llarg de qualsevol característica constant \(1\). \begin{indicacio}El pullback de dos punts diferents de \(\mathbf2\) és inicial, per disjunció.\end{indicacio}
\end{enumerate}

\section{Negació i complement}

\begin{enumerate}[resume]
  \item Demostreu que \(\neg\circ\neg=\id_{\mathbf2}\) i que, per tant, la negació és un isomorfisme. (Model: secció~\ref{prac:complement}; proposició~\ref{prop:negacio-involutiva}.)

  \item Demostreu que el complement és involutiu: \((m^c)^c\equiv_A m\) per a tota part \(m\). \begin{indicacio}Calculeu \(\chi_{(m^c)^c}=\neg\circ\neg\circ\chi_m\) i useu que la característica determina la part.\end{indicacio}

  \item Demostreu que una part i el seu complement descomponen \(A\): \([m,m^c]:S+S^c\to A\) és un isomorfisme. (Enunciat: teorema~\ref{teo:part-complementari}.)

  \item Demostreu que la part total i la part buida són complementàries l'una de l'altra. \begin{indicacio}Compareu les característiques: \(\neg\circ(1\circ{!}_A)=0\circ{!}_A\).\end{indicacio}
\end{enumerate}

\section{Operacions booleanes amb parts}

\begin{enumerate}[resume]
  \item Demostreu que la característica de \(m\cap n\) és \(\wedge\circ\langle\chi_m,\chi_n\rangle\) i que la de \(m\cup n\) és \(\vee\circ\langle\chi_m,\chi_n\rangle\). (Model: secció~\ref{prac:booleanes}; proposició~\ref{prop:caracteristica-interseccio} i definició~\ref{def:conjuncio-disjuncio}.)

  \item Demostreu les lleis d'idempotència \(m\cap m\equiv_A m\) i \(m\cup m\equiv_A m\) reduint-les a identitats de taules de veritat sobre \(\mathbf2\). \begin{indicacio}Useu l'extensionalitat de les taules (proposició~\ref{prop:extensionalitat-taules}).\end{indicacio}

  \item Demostreu una de les lleis de De Morgan per a parts, \((m\cap n)^c\equiv_A m^c\cup n^c\, \), traduint-la a una identitat entre \(\neg\), \(\wedge\) i \(\vee\) sobre \(\mathbf2\times\mathbf2\). \begin{indicacio}Comproveu la identitat \(\neg\circ\wedge=\vee\circ(\neg\times\neg)\) sobre els quatre punts de \(\mathbf2\times\mathbf2\); vegeu les lleis booleanes (teorema~\ref{teo:lleis-booleanes}).\end{indicacio}

  \item Demostreu que \(m\subseteq_A n\) si i només si \(m\cap n\equiv_A m\), i també si i només si \(m\cup n\equiv_A n\). (Model de la proposició~\ref{prop:inclusio-booleana}.)
\end{enumerate}

\section{Límits, coequalitzadors, pushouts i colímits finits}

\begin{enumerate}[resume]
  \item Demostreu que el coequalitzador \(q:B\to Q\) de \(f,g:A\to B\) és un epimorfisme, i que qualsevol altre coequalitzador dels mateixos \(f,g\) hi és isomorf de manera compatible. (Model: secció~\ref{prac:limits-colimits}; proposicions~\ref{prop:coequalitzador-epi} i~\ref{prop:coequalitzador-unic}.)

  \item Construïu el pushout de \(f:C\to A\) i \(g:C\to B\) com a coequalitzador de \(i_A\circ f\) i \(i_B\circ g\) sobre la suma \(A+B\), i comproveu-ne la propietat universal. \begin{indicacio}Seguiu el teorema~\ref{teo:pushout-construccio}; l'aplicació mediadora és \([u,v]\) factoritzada per \(q\).\end{indicacio}

  \item Demostreu que, tenint productes finits i equalitzadors, existeix el límit de qualsevol diagrama finit, i digueu quines aplicacions calen per imposar cada fletxa del diagrama. (Model del teorema~\ref{teo:limits-finits}.)

  \item Dualitzeu l'exercici anterior: expliqueu com, tenint sumes finites i coequalitzadors, s'obté el colímit de qualsevol diagrama finit. \begin{indicacio}És el teorema~\ref{teo:colimits-finits}; invertiu el sentit de totes les fletxes i canvieu productes/equalitzadors per sumes/coequalitzadors.\end{indicacio}

  \item Demostreu, sense usar cap axioma d'elecció, que tot epimorfisme \(f:A\to B\) és exhaustiu sobre els punts globals. \begin{indicacio}És la proposició~\ref{prop:epi-exhaustiu-punts}: per a \(y:\mathbf1\to B\), mireu la imatge inversa \(f^{-1}(y)\); si no té punts és inicial, i la seva característica \(\chi_y\circ f\) és \(0\circ{!}_A\); cancel·leu \(f\) i arribeu a \(0=1\).\end{indicacio}

  \item \textbf{Axiomes usats.} Per a cadascun dels resultats següents, digueu si es demostra amb els axiomes dels capítols~\ref{cap:aplicacions} i~\ref{cap:parts} sols, o si necessita la disjunció i estabilitat de les sumes, els coequalitzadors, o tots dos: \(0\ne1\); el lema del buit; les lleis booleanes de les parts; l'existència dels límits finits; l'existència dels pushouts; l'exhaustivitat dels epimorfismes sobre punts. \begin{indicacio}Només els límits finits es demostren amb els axiomes anteriors; els pushouts són els únics que necessiten coequalitzadors; la resta usa la disjunció i estabilitat (i el lema del buit i l'exhaustivitat, a més, la separació global).\end{indicacio}
\end{enumerate}

\chapter{Relacions, quocients i elecció}\label{cap:exercicis-relacions}

Aquests exercicis proposats corresponen al capítol homònim de la part teòrica (capítol~\ref{cap:relacions} de la Part~\ref{part:teoria}) i es resolen amb les tècniques dels exercicis resolts de la part pràctica (capítol~\ref{cap:practica-relacions} de la Part~\ref{part:practica}). Cada secció remet al bloc de resolts que en dona el model. Els marcats amb una indicació són els més exigents.

\section{Relacions, inversa i diagonal}

\begin{enumerate}
  \item Sigui \(r\colon R\rightarrowtail A\times B\) una relació. Demostreu que \(y\mathrel{\operatorname{Inv}(r)_X}x\) si i només si \(x\mathrel{r_X}y\), i deduïu que \(\operatorname{Inv}(\operatorname{Inv}(r))\equiv_{A\times B}r\). (Model: secció~\ref{prac:relacions}.)

  \item Demostreu que la diagonal \(\Delta_A\) és un monomorfisme i que \(x\mathrel{(\Delta_A)_X}x'\) equival a \(x=x'\). (Model: proposició~\ref{prop:diagonal-mono}.)

  \item Demostreu que una relació \(r\) sobre \(A\) és \emph{simètrica} en el sentit generalitzat (\(x\mathrel{r_X}y\Rightarrow y\mathrel{r_X}x\) a tots els estadis) si i només si \(\operatorname{Inv}(r)\equiv_{A\times A}r\), i que, amb la separació global, n'hi ha prou amb els punts globals. \begin{indicacio}Model: secció~\ref{prac:relacions}. Per a la segona part, dues parts amb els mateixos punts globals són equivalents pel lema del buit.\end{indicacio}

  \item Demostreu que la diagonal és simètrica: \(\operatorname{Inv}(\Delta_A)\equiv_{A\times A}\Delta_A\). \begin{indicacio}\(\sigma\circ\langle\id_A,\id_A\rangle=\langle\id_A,\id_A\rangle\).\end{indicacio}
\end{enumerate}

\section{El parell nucli}

\begin{enumerate}[resume]
  \item Demostreu que el parell nucli \(\langle k_1,k_2\rangle\colon A\times_B A\to A\times A\) de qualsevol \(f\colon A\to B\) és un monomorfisme, i observeu que la prova no usa cap propietat de \(f\). (Model: secció~\ref{prac:parell-nucli}; proposició~\ref{prop:kernel-mono}.)

  \item Construïu les aplicacions estructurals que testimonien la reflexivitat i la simetria del parell nucli \(\operatorname{Nuc}(f)\). \begin{indicacio}Per a la reflexivitat, factoritzeu \((\id_A,\id_A)\) pel pullback; per a la simetria, \((k_2,k_1)\).\end{indicacio}

  \item Completeu la prova que \(\operatorname{Nuc}(f)\) és una relació d'equivalència construint l'aplicació de transitivitat sobre \(K\times_A K\). \begin{indicacio}Seguiu la proposició~\ref{prop:kernel-equivalencia}: comproveu \(f\circ(k_1\circ p)=f\circ(k_2\circ q)\) i factoritzeu pel pullback \(K\).\end{indicacio}
\end{enumerate}

\section{Quocients i epimorfismes regulars}

\begin{enumerate}[resume]
  \item Demostreu que tota relació sobre \(A\) té quocient, únic llevat d'un únic isomorfisme compatible, i que l'aplicació quocient és un epimorfisme. (Model: secció~\ref{prac:quocients}; proposició~\ref{prop:quocient-existeix}.)

  \item Demostreu que tot quocient d'una relació d'equivalència és un epimorfisme regular. (Model: proposició~\ref{prop:quocient-regular}.)

  \item Doneu un exemple, en el model clàssic, d'un epimorfisme, i comproveu que és el quocient pel seu propi parell nucli. \begin{indicacio}Preneu una aplicació exhaustiva entre conjunts finits; el parell nucli identifica els punts amb la mateixa imatge.\end{indicacio}
\end{enumerate}

\section{Efectivitat i correspondència}

\begin{enumerate}[resume]
  \item Demostreu que el parell nucli és invariant, com a part de \(A\times A\), sota postcomposició del codomini amb un isomorfisme. (Model: secció~\ref{prac:efectivitat}; lema~\ref{lema:nucli-invariant}.)

  \item Usant l'axioma d'efectivitat, demostreu que per a tota relació d'equivalència \(r\) i tot quocient \(q\) seu es té \(\operatorname{Nuc}(q)\equiv_{A\times A}r\). (Camí (1)$\to$(2) del teorema~\ref{teo:correspondencia}.)

  \item Demostreu l'altra direcció de la correspondència: per a tot epimorfisme regular \(q\colon A\twoheadrightarrow Q\), el quocient del seu parell nucli és isomorf a \(Q\) de manera compatible amb \(q\). \begin{indicacio}Factoritzeu \(q\) per \(q'\colon A\to A/\operatorname{Nuc}(q)\) i demostreu que la mediadora \(\bar q\) és un isomorfisme, com al teorema~\ref{teo:correspondencia}.\end{indicacio}
\end{enumerate}

\section{Particions, escindits i elecció}

\begin{enumerate}[resume]
  \item Demostreu que tot epimorfisme escindit és el coequalitzador del seu parell nucli, i per tant és regular. (Model: secció~\ref{prac:particions}; proposició~\ref{prop:escindit-regular}.)

  \item Demostreu que, amb l'axioma d'elecció, tot epimorfisme és regular, i deduïu que la correspondència s'estén a tots els epimorfismes amb domini \(A\). (Model: secció~\ref{prac:eleccio}; proposició~\ref{prop:epi-regular-eleccio} i corol·lari~\ref{cor:correspondencia-tots-epis}.)

  \item Demostreu que, amb l'axioma d'elecció, tot epimorfisme \(q\colon A\twoheadrightarrow B\) és exhaustiu sobre elements generalitzats. Expliqueu per què l'exhaustivitat sobre els punts globals, en canvi, no necessita l'elecció. \begin{indicacio}Preneu una secció \(s\) i comproveu \(q\circ(s\circ y)=y\); vegeu les proposicions~\ref{prop:epi-exhaustiu-generalitzat} i~\ref{prop:epi-exhaustiu-global}.\end{indicacio}

  \item Demostreu que un epimorfisme és escindit si i només si és exhaustiu sobre elements generalitzats, i doneu un exemple clàssic (sense usar l'elecció) d'un epimorfisme escindit que no és isomorfisme. \begin{indicacio}Proposició~\ref{prop:exhaustiva-generalitzada-escindit} del capítol~\ref{cap:parts}; una projecció \(A\times B\to A\) amb \(B\) no buit té secció \(\langle\id_A,b\circ{!}_A\rangle\).\end{indicacio}

  \item \textbf{Axiomes usats.} Digueu, per a cadascun dels resultats següents, si usa l'efectivitat, l'elecció, tots dos o cap: el parell nucli és una relació d'equivalència; tot quocient és un epimorfisme regular; el parell nucli del quocient recupera la relació; tot epimorfisme escindit és regular; tot epimorfisme és regular; tot epimorfisme és exhaustiu sobre punts. \begin{indicacio}Només el tercer usa l'efectivitat i només el cinquè l'elecció; el primer, el segon i el quart no usen cap dels dos, i el sisè tampoc (capítol~\ref{cap:valors}).\end{indicacio}

  \item Demostreu que, amb l'axioma d'elecció, tot quocient \(q\colon A\twoheadrightarrow A/r\) admet una secció, i interpreteu-la com un sistema de representants de les classes. (Model: proposició~\ref{prop:representants}.)
\end{enumerate}

\chapter{Imatges i factoritzacions}\label{cap:exercicis-imatges}

Aquests exercicis proposats corresponen al capítol homònim de la part teòrica (capítol~\ref{cap:imatges} de la Part~\ref{part:teoria}) i es resolen amb les tècniques dels exercicis resolts de la part pràctica (capítol~\ref{cap:practica-imatges} de la Part~\ref{part:practica}). Cada secció remet al bloc de resolts que en dona el model. Els marcats amb una indicació són els més exigents.

\section{La factorització imatge}

\begin{enumerate}
  \item Descriviu la construcció de la factorització \(A\overset{e}{\twoheadrightarrow}\operatorname{Im}(f)\overset{m}{\hookrightarrow}B\) a partir del parell nucli i el quocient, i justifiqueu l'existència de l'única \(m\) amb \(f=m\circ e\). (Model: secció~\ref{prac:factoritzacio}; teorema~\ref{teo:factoritzacio}.)

  \item Demostreu, sense l'axioma d'elecció, que el factor \(m\) de la factorització és un monomorfisme. \begin{indicacio}Reduïu-ho als punts globals (observació~\ref{obs:injectiva-mono}); de \(m\circ y=m\circ y'\) obteniu preimatges \(a,a'\) per \(e\) (proposició~\ref{prop:epi-exhaustiu-punts}), i useu que \(e\) identifica els relacionats pel parell nucli.\end{indicacio}

  \item Deduïu de la factorització imatge i del teorema del bimorfisme que tot epimorfisme és regular, sense l'elecció. \begin{indicacio}Corol·lari~\ref{cor:epi-regular-sense-eleccio}: si \(f=m\circ e\) és epi, \(m\) és epi i mono.\end{indicacio}

  \item Demostreu que dues factoritzacions (epi regular)--(mono) d'una mateixa \(f\) són isomorfes de manera única i compatible. (Model: unicitat del teorema~\ref{teo:factoritzacio}.)

  \item Identifiqueu, sobre punts globals en el model clàssic, què és \(\operatorname{Im}(f)\) per a una aplicació concreta \(f\) entre conjunts finits, i comproveu que \(m\) n'és la inclusió i \(e\) la correstricció exhaustiva. \begin{indicacio}Vegeu l'observació~\ref{obs:imatge-classica}: la imatge és el conjunt de valors assolits.\end{indicacio}
\end{enumerate}

\section{Caracteritzacions per la factorització}

\begin{enumerate}[resume]
  \item Demostreu que \(f\) és un monomorfisme si i només si \(\operatorname{Nuc}(f)\equiv_{A\times A}\Delta_A\), i que en aquest cas \(e\) és un isomorfisme. (Model: secció~\ref{prac:caracteritzacions}; proposició~\ref{prop:mono-nucli-trivial}.)

  \item Demostreu que \(f\) és un epimorfisme si i només si el factor \(m\) és un isomorfisme. (Model: proposició~\ref{prop:epi-imatge-total}.)

  \item Demostreu que \(f\) és un isomorfisme si i només si és alhora monomorfisme i epimorfisme. \begin{indicacio}Combineu les dues caracteritzacions anteriors sobre els factors \(e\) i \(m\); vegeu el corol·lari~\ref{cor:bimorfisme-iso}.\end{indicacio}
\end{enumerate}

\section{Imatge directa i adjunció}

\begin{enumerate}[resume]
  \item Demostreu l'adjunció \(f_{*}S\subseteq_B T\iff S\subseteq_A f^{*}T\) per a parts \(S\) de \(A\) i \(T\) de \(B\). (Model: secció~\ref{prac:imatge-directa}; proposició~\ref{prop:adjuncio-imatges}.)

  \item Demostreu que \(f_{*}A=\operatorname{Im}(f)\) i que \(f\) és epimorfisme si i només si \(f_{*}A\equiv_B\id_B\). (Model: corol·lari~\ref{cor:estabilitat-exhaustiva}.)

  \item Useu l'adjunció per demostrar que \(f_{*}\) preserva les inclusions (si \(S\subseteq_A S'\) aleshores \(f_{*}S\subseteq_B f_{*}S'\)) i que \(f^{*}\) també. \begin{indicacio}De \(S\subseteq_A S'\) i \(S'\subseteq_A f^{*}f_{*}S'\) deduïu \(S\subseteq_A f^{*}f_{*}S'\), i torneu a aplicar l'adjunció.\end{indicacio}

  \item Demostreu la unitat i la counitat de l'adjunció: \(S\subseteq_A f^{*}f_{*}S\) i \(f_{*}f^{*}T\subseteq_B T\). \begin{indicacio}Preneu \(T:=f_{*}S\) i \(S:=f^{*}T\) respectivament a l'equivalència de l'adjunció.\end{indicacio}
\end{enumerate}

\section{Composició de relacions i relacions funcionals}

\begin{enumerate}[resume]
  \item Demostreu que la composició \(r'\circ r\) de dues relacions és una relació, i expliqueu quin quantificador introdueix la imatge. (Model: secció~\ref{prac:relacions-funcionals}; proposició~\ref{prop:composicio-es-relacio}.)

  \item Demostreu que \(r\) és funcional si i només si per a tot \(x\colon X\to A\) hi ha com a molt un \(y\) amb \(x\mathrel{r_X}y\), i que \(r\) és total si i només si n'hi ha almenys un (potser després d'un canvi d'estadi). (Model: proposició~\ref{prop:funcional-lectura}.)

  \item Demostreu que assignar a cada aplicació el seu graf estableix una correspondència biunívoca (llevat d'equivalència) entre aplicacions \(A\to B\) i relacions funcionals i totals, i que aquesta correspondència respecta la composició. (Model: proposició~\ref{prop:aplicacio-funcional}.)

  \item Demostreu que la relació inversa d'una relació funcional i total no té per què ser funcional, i digueu què significa, sobre punts, que \(\operatorname{Inv}(r)\) sigui total. \begin{indicacio}Traduïu: \(\operatorname{Inv}(r)\) funcional \(=\) \(r\) injectiva; \(\operatorname{Inv}(r)\) total \(=\) \(r\) exhaustiva. Doneu un \(f\) no injectiu com a contraexemple.\end{indicacio}

  \item Demostreu el lema de pertinença després d'un recobriment (lema~\ref{lema:pertinenca-recobriment}): si \(p\colon X'\twoheadrightarrow X\) és un epimorfisme i \(y\circ p\in_{X'}m\), aleshores \(y\in_Xm\). \begin{indicacio}La imatge inversa \(y^{-1}(m)\) és una part de \(X\) per la qual \(p\) es factoritza; és, doncs, mono i epi.\end{indicacio}

  \item \textbf{Axiomes usats.} Per a cadascun dels resultats següents digueu si usa la separació global, l'elecció, tots dos o cap: la construcció de \(e\) i \(m\) amb \(f=m\circ e\); que \(m\) és mono; que tot epimorfisme és regular; l'estabilitat dels epimorfismes per canvi de base; l'adjunció \(f_*\dashv f^*\); que la composició de relacions és una relació. \begin{indicacio}Cap dels sis no usa l'elecció. El segon, el tercer i el quart usen la separació global (a través de l'exhaustivitat sobre punts); el primer, el cinquè i el sisè, cap dels dos.\end{indicacio}
\end{enumerate}

\chapter{Exponencials}\label{cap:exercicis-exponencials}

Aquests exercicis proposats corresponen al capítol homònim de la part teòrica (capítol~\ref{cap:exponencials} de la Part~\ref{part:teoria}) i es resolen amb les tècniques dels exercicis resolts de la part pràctica (capítol~\ref{cap:practica-exponencials} de la Part~\ref{part:practica}). Cada secció remet al bloc de resolts que en dona el model. Els marcats amb una indicació són els més exigents.

\section{Propietat universal i transposició}

\begin{enumerate}
  \item Demostreu que l'exponencial de \(B\) per \(A\) és únic llevat d'un únic isomorfisme compatible amb l'avaluació. (Model: secció~\ref{prac:exp-universal}; proposició~\ref{prop:exponencial-unic}.)

  \item Demostreu que currificar i descurrificar són operacions mútuament inverses. \begin{indicacio}Useu la unicitat de la transposada de la definició~\ref{def:exponencial} en les dues direccions.\end{indicacio}

  \item Demostreu que la transposició és compatible amb la precomposició en \(A\): per a \(k\colon A'\to A\) i \(h\colon X\times A\to B\), definiu \(B^{k}\colon B^{A}\to B^{A'}\) com la transposada de \(\operatorname{ev}\circ(\id_{B^A}\times k)\) i demostreu que \(\lambda(h\circ(\id_X\times k))=B^{k}\circ\lambda(h)\). \begin{indicacio}Comproveu que els dos costats tenen la mateixa descurrificació i useu la unicitat de la transposada. Sobre punts, \(B^k\) envia \(f\) a \(f\circ k\).\end{indicacio}
\end{enumerate}

\section{Aplicacions com a punts i famílies}

\begin{enumerate}[resume]
  \item Demostreu que els punts globals \(\mathbf1\to B^A\) es corresponen biunívocament amb les aplicacions \(A\to B\). (Model: secció~\ref{prac:exp-punts}; teorema~\ref{teo:punts-exponencial}.)

  \item Demostreu que un element generalitzat \(X\to B^A\) es correspon amb una aplicació \(X\times A\to B\), i interpreteu-ho com una família d'aplicacions indexada per \(X\). (Model: proposició~\ref{prop:families}.)

  \item Demostreu la naturalitat de la transposició en \(X\): per a \(t\colon X'\to X\), \(\lambda(h)\circ t=\lambda(h\circ(t\times\id_A))\). \begin{indicacio}Comproveu que tots dos costats tenen la mateixa descurrificació i useu la unicitat.\end{indicacio}
\end{enumerate}

\section{Exponencials de \texorpdfstring{\(\mathbf0\), \(\mathbf1\), \(\mathbf2\)}{0, 1, 2}}

\begin{enumerate}[resume]
  \item Demostreu que \(A^{\mathbf1}\cong A\), \(A^{\mathbf0}\cong\mathbf1\) i \(\mathbf1^{A}\cong\mathbf1\). (Model: secció~\ref{prac:exp-basics}; proposició~\ref{prop:exponencials-basics}.)

  \item Demostreu que els punts de \(P(A)=\mathbf2^{A}\) es corresponen amb les parts de \(A\) llevat d'equivalència. (Model: definició~\ref{def:conjunt-parts}.)

  \item Determineu, en el model clàssic, el nombre d'aplicacions \(\mathbf2\to\mathbf2\) i identifiqueu-les amb els punts de \(\mathbf2^{\mathbf2}\). \begin{indicacio}Useu \(\mathbf2^{\mathbf2}\cong\mathbf2^{\mathbf1+\mathbf1}\cong\mathbf2\times\mathbf2\) i compteu-ne els punts.\end{indicacio}
\end{enumerate}

\section{Suma, composició interna i famílies}

\begin{enumerate}[resume]
  \item Demostreu que \(B^{A+C}\cong B^{A}\times B^{C}\). (Model: secció~\ref{prac:exp-suma}; proposició~\ref{prop:exponencial-suma}.)

  \item Demostreu que existeix una única composició interna \(c\colon C^{B}\times B^{A}\to C^{A}\) que sobre punts envia \((g,f)\) a \(g\circ f\), i una única \(\jmath\colon\mathbf1\to A^{A}\) corresponent a \(\id_A\). (Model: secció~\ref{prac:exp-composicio}; proposició~\ref{prop:composicio-interna}.)

  \item Demostreu que la composició interna és associativa i que \(\jmath\) n'és neutre, en el sentit de les igualtats internes corresponents sobre \(A^{A}\). \begin{indicacio}Traduïu l'associativitat i la neutralitat a igualtats d'aplicacions cap a l'exponencial i comproveu-les descurrificant amb l'avaluació, com a la proposició~\ref{prop:composicio-interna}.\end{indicacio}

  \item A partir de \(B^{A+C}\cong B^{A}\times B^{C}\), deduïu per iteració que l'exponencial d'una suma de \(n\) còpies de \(A\) és un producte de \(n\) còpies de \(B^{A}\). \begin{indicacio}Inducció sobre \(n\) usant la proposició~\ref{prop:exponencial-suma}; vegeu l'observació~\ref{obs:exponencial-suma-classica}. La inducció és metateòrica: \(n\) és un nombre del metallenguatge fins al capítol~\ref{cap:naturals}.\end{indicacio}

  \item Demostreu les altres dues lleis d'exponents: \((B\times C)^{A}\cong B^{A}\times C^{A}\) i \((C^{B})^{A}\cong C^{A\times B}\). \begin{indicacio}Per a la primera, una aplicació \(X\times A\to B\times C\) és un parell d'aplicacions; per a la segona, transposeu dues vegades usant l'associativitat del producte. En tots dos casos exhibiu l'avaluació, com a la prova de la proposició~\ref{prop:exponencial-suma}.\end{indicacio}

  \item \textbf{Axiomes usats.} Per a cadascun dels resultats següents, digueu si usa l'axioma d'exponencials, la separació global, tots dos o cap: la unicitat de l'exponencial; \(A^{\mathbf0}\cong\mathbf1\); la correspondència entre punts de \(B^A\) i aplicacions \(A\to B\); l'existència de la composició interna; la unicitat de la composició interna determinada pels seus valors sobre punts. \begin{indicacio}La unicitat de l'exponencial no usa cap dels dos (és el patró de les construccions universals); els tres següents usen l'axioma d'exponencials; l'últim, a més, la separació global.\end{indicacio}
\end{enumerate}

\chapter{El conjunt de parts i la pertinença interna}\label{cap:exercicis-parts-interna}

Aquests exercicis proposats corresponen al capítol homònim de la part teòrica (capítol~\ref{cap:parts-interna} de la Part~\ref{part:teoria}) i es resolen amb les tècniques dels exercicis resolts de la part pràctica (capítol~\ref{cap:practica-parts-interna} de la Part~\ref{part:practica}). Cada secció remet al bloc de resolts que en dona el model. Els marcats amb una indicació són els més exigents.

\section{La pertinença interna}

\begin{enumerate}
  \item Demostreu que \(a\mathrel{\varepsilon}\alpha\) equival a \(\operatorname{ev}\circ\langle\alpha,a\rangle=1\circ{!}_X\), i que sobre el nom constant d'una part \(m\) es recupera \(a\in_X m\). (Model: secció~\ref{prac:pertinenca-interna}; proposició~\ref{prop:pertinenca-interna-lectura}.)

  \item Demostreu que \(\in_A\) és la part de \(A\times P(A)\) associada a \(\id_{P(A)}\), i que tota pertinença \(a\mathrel{\varepsilon}(\alpha\circ x)\) s'obté reindexant \(\in_A\). (Model: corol·lari~\ref{cor:pertinenca-universal}.)

  \item Comproveu, sobre punts globals en el model clàssic, que la pertinença interna \(a\mathrel{\varepsilon}\ulcorner m\urcorner\) coincideix amb \(a\in m\) per a una part concreta \(m\) d'un conjunt finit. \begin{indicacio}Descurrifiqueu \(\ulcorner m\urcorner\) i avalueu \(\chi_m\) sobre els punts de \(A\).\end{indicacio}
\end{enumerate}

\section{Parts parametritzades}

\begin{enumerate}[resume]
  \item Demostreu la correspondència entre aplicacions \(X\to P(A)\), aplicacions \(X\times A\to\mathbf2\) i parts de \(X\times A\) llevat d'equivalència. (Model: secció~\ref{prac:parts-parametritzades}; teorema~\ref{teo:parts-parametritzades}.)

  \item Interpreteu un element generalitzat \(X\to P(A)\) com una família de parts de \(A\), i digueu què recupera el cas \(X=\mathbf1\). \begin{indicacio}Vegeu l'observació~\ref{obs:generalitzats-parts}; el cas global dona una sola part de \(\mathbf1\times A\cong A\).\end{indicacio}
\end{enumerate}

\section{Operacions booleanes internes}

\begin{enumerate}[resume]
  \item Demostreu les tres identitats de càlcul de les operacions internes (proposició~\ref{prop:booleanes-puntuals}) i deduïu-ne que \(a\mathrel{\varepsilon}(\alpha\cap\beta)\) equival a \enquote{\(a\mathrel{\varepsilon}\alpha\) i \(a\mathrel{\varepsilon}\beta\)} a tot estadi, mentre que les lectures amb \enquote{o} i \enquote{no} només valen sobre punts globals. Doneu un contraexemple amb \(X=\mathbf2\) per a la unió. (Model: secció~\ref{prac:booleanes-internes}.)

  \item Demostreu que les operacions internes estenen les de parts sobre els noms: \(\ulcorner m\urcorner\cap\ulcorner n\urcorner=\ulcorner m\cap n\urcorner\), etc. (Model: proposició~\ref{prop:booleanes-puntuals}.)

  \item Demostreu una llei de De Morgan interna, \(\neg(\alpha\cup\beta)=\neg\alpha\cap\neg\beta\), comparant descurrificacions: totes dues aplicacions \(P(A)\times P(A)\to P(A)\) es transposen a \(\neg\circ\vee\circ\langle\operatorname{ev}_1,\operatorname{ev}_2\rangle\) i a \(\wedge\circ(\neg\times\neg)\circ\langle\operatorname{ev}_1,\operatorname{ev}_2\rangle\), que coincideixen per la llei de De Morgan sobre \(\mathbf2\). \begin{indicacio}Vegeu l'observació~\ref{obs:lleis-transportades}; no cal cap lectura puntual.\end{indicacio}
\end{enumerate}

\section{La inclusió com a ordre}

\begin{enumerate}[resume]
  \item Demostreu la lectura de la inclusió interna: \(\langle\alpha,\beta\rangle\in_X{\subseteq_A}\) si i només si \(s_\alpha\subseteq s_\beta\) com a parts de \(X\times A\), i que sobre els noms recupera \(m\subseteq_A n\). (Model: secció~\ref{prac:inclusio-ordre}; proposició~\ref{prop:inclusio-interna-lectura}.)

  \item Demostreu que \(\subseteq_A\) és reflexiva, transitiva i antisimètrica, i que \(\cap\) és l'ínfim i \(\cup\) el suprem. (Model: teorema~\ref{teo:inclusio-ordre}.)

  \item Demostreu la propietat de l'ínfim: \(\langle\gamma,\alpha\cap\beta\rangle\in_X{\subseteq_A}\) equival a \(\langle\gamma,\alpha\rangle\in_X{\subseteq_A}\) i \(\langle\gamma,\beta\rangle\in_X{\subseteq_A}\). \begin{indicacio}Useu que \(s_{\alpha\cap\beta}=s_\alpha\cap s_\beta\) i que una part és inclosa en una intersecció si i només si ho és en cadascun dels factors.\end{indicacio}
\end{enumerate}

\section{Imatge inversa interna i quantificadors}

\begin{enumerate}[resume]
  \item Demostreu que \(P(f)\circ\ulcorner n\urcorner=\ulcorner f^{-1}(n)\urcorner\) i que \(P(f)\) commuta amb el complement intern. (Model: secció~\ref{prac:quantificadors}; proposicions~\ref{prop:imatge-inversa-interna} i~\ref{prop:imatge-inversa-complement}.)

  \item Demostreu el lema de Beck--Chevalley (lema~\ref{lema:beck-chevalley}) i useu-lo per veure que \(\exists_f\circ\alpha\) correspon a la part \((\id_X\times f)_*(s_\alpha)\). \begin{indicacio}Pulleu enrere la factorització imatge de \(f\circ s\) al llarg de \(h'\): el pullback d'un mono és mono i el d'un epi és epi, i la factorització és única.\end{indicacio}

  \item Demostreu que \(\exists_f\) és adjunt per l'esquerra i \(\forall_f\) adjunt per la dreta de \(P(f)\), respecte de la inclusió de les parts associades. (Model: proposició~\ref{prop:quantificadors-adjuncio}.)

  \item Demostreu que \(\exists_A:=\exists_{{!}_A}\) i \(\forall_A:=\forall_{{!}_A}\), vistes com a aplicacions \(P(A)\to P(\mathbf1)\cong\mathbf2\), envien el nom d'una part \(m\) a \(1\) exactament quan \(m\) és no buida i, respectivament, quan \(m\) és total. \begin{indicacio}Sobre noms, \(\exists_A\ulcorner m\urcorner=\ulcorner{!}_{A*}m\urcorner\), i la imatge de \(m\) dins \(\mathbf1\) és total si i només si \(S\to\mathbf1\) és epi, és a dir, si \(S\) té algun punt (capítol~\ref{cap:valors}); per a \(\forall_A\) useu el complement.\end{indicacio}

  \item \textbf{Axiomes usats.} Digueu quins dels resultats següents usen la separació global (a través de l'estabilitat dels epimorfismes o del lema del buit) i quins no: la construcció de \(\in_A\); el teorema de les parts parametritzades; les identitats de càlcul de les operacions internes; el lema de Beck--Chevalley; la lectura de \(\exists_f\) sobre punts globals; que \(\subseteq_A\) és un ordre. \begin{indicacio}Només el lema de Beck--Chevalley (estabilitat dels epis) i la lectura sobre punts (lema del buit) la usen; cap resultat del capítol no usa l'elecció.\end{indicacio}

  \item Demostreu que \(P(f)\) preserva les operacions booleanes internes: \(P(f)(\beta\cap\beta')=P(f)\beta\cap P(f)\beta'\) i anàlogament per a \(\cup\). \begin{indicacio}Comproveu-ho puntualment amb la definició~\ref{def:imatge-inversa-interna} i el càlcul de \(\cap\).\end{indicacio}
\end{enumerate}

\chapter{L'objecte dels nombres naturals}\label{cap:exercicis-naturals}

Aquests exercicis proposats corresponen al capítol homònim de la part teòrica (capítol~\ref{cap:naturals} de la Part~\ref{part:teoria}) i es resolen amb les tècniques dels exercicis resolts de la part pràctica (capítol~\ref{cap:practica-naturals} de la Part~\ref{part:practica}). Cada secció remet al bloc de resolts que en dona el model. Els marcats amb una indicació són els més exigents.

\section{Propietat universal, unicitat i iteració}

\begin{enumerate}
  \item Demostreu que l'objecte dels nombres naturals és únic llevat d'un únic isomorfisme compatible amb \(0\) i \(s\). (Model: secció~\ref{prac:nno-universal}; teorema~\ref{teo:nno-unicitat}.)

  \item Demostreu que la iterada \(\operatorname{it}(t)\colon\N\times A\to A\) d'una autoaplicació \(t\colon A\to A\) existeix i és única. \begin{indicacio}Recursió sobre el sistema \((A^{A},\ulcorner\id_A\urcorner,C_t)\) i descurrificació; vegeu la proposició~\ref{prop:iterada}.\end{indicacio}

  \item Demostreu que la recursora general \(u=\operatorname{rec}(a,t)\) s'obté iterant \(t\) a partir de \(a\): \(u\circ n=t^{n}(a)\) per a tot \(n\). \begin{indicacio}Comproveu que \(t^{(-)}(a)\) satisfà les equacions de la recursora i useu-ne la unicitat; vegeu l'observació~\ref{obs:recursio-iteracio}.\end{indicacio}
\end{enumerate}

\section{Inducció i equacions de Peano}

\begin{enumerate}[resume]
  \item Demostreu la unicitat per inducció: si \(f\circ0=g\circ0\) i la part on \(f=g\) és tancada pel successor, aleshores \(f=g\). (Model: secció~\ref{prac:induccio}; proposició~\ref{prop:unicitat-induccio}.)

  \item Demostreu que el zero no és cap successor. (Model: secció~\ref{prac:peano}; teorema~\ref{teo:peano}.)

  \item Demostreu que el successor és un monomorfisme construint el predecessor. \begin{indicacio}Recursió sobre \(\N\times\N\) amb \(q\circ0=\langle0,0\rangle\) i \(q\circ s=\langle p_2,s\circ p_2\rangle\circ q\); proveu \(p_2\circ q=\id\) per unicitat i preneu \(p=p_1\circ q\). Expliqueu per què la recursora de \((\N,0,\id)\) no serveix.\end{indicacio}

  \item Demostreu que tot punt de \(\N\) és \(0\) o un successor (lema~\ref{lema:zero-o-successor}), i deduïu-ne que \(\N\cong\mathbf1+\N\) via \([0,s]\). \begin{indicacio}La unió de la part \(0\) i de la imatge de \(s\) conté \(0\) i és tancada per \(s\); per a l'isomorfisme, \([0,s]\) és mono (el zero no és successor i \(s\) és mono) i epi (exhaustiu sobre punts pel lema).\end{indicacio}

  \item Demostreu la unicitat per inducció amb paràmetres (corol·lari~\ref{cor:induccio-parametres}) i expliqueu on intervenen els exponencials. \begin{indicacio}Transposeu \(f,g\colon X\times\N\to A\) a \(\N\to A^X\) i apliqueu la unicitat per inducció al seu equalitzador.\end{indicacio}
\end{enumerate}

\section{Aritmètica}

\begin{enumerate}[resume]
  \item Demostreu \(0+n=n\) i, a partir d'aquí, la neutralitat, l'associativitat i la commutativitat de la suma. (Model: secció~\ref{prac:aritmetica}; proposició~\ref{prop:suma}.)

  \item Demostreu que la suma és cancel·lativa (\(m+k=m+l\Rightarrow k=l\)) i que \(k+l=0\Rightarrow k=l=0\). \begin{indicacio}Inducció sobre \(m\) per a la cancel·lació, usant la injectivitat del successor; per a la suma nul·la, useu que un successor no és zero.\end{indicacio}

  \item Demostreu la distributivitat \(m\cdot(n+k)=m\cdot n+m\cdot k\) i deduïu la commutativitat del producte. \begin{indicacio}Inducció sobre \(k\) per a la distributivitat (proposició~\ref{prop:lleis-aritmetica}); per a la commutativitat, inducció amb els lemes \(0\cdot n=0\) i \((s\circ m)\cdot n=m\cdot n+n\).\end{indicacio}

  \item Demostreu les lleis dels exponents \(m^{p+q}=m^{p}\cdot m^{q}\) i \(m^{p\cdot q}=(m^{p})^{q}\). \begin{indicacio}Inducció sobre \(q\), usant la definició recursiva de l'exponenciació i les lleis del producte (definició~\ref{def:producte-exp}, proposició~\ref{prop:lleis-aritmetica}).\end{indicacio}
\end{enumerate}

\section{Ordre i bon ordre}

\begin{enumerate}[resume]
  \item Demostreu que \(\le\) és reflexiu, transitiu i antisimètric a partir de \(m\le n\iff\exists k\,(m+k=n)\). (Model: secció~\ref{prac:ordre}; proposició~\ref{prop:ordre-total}.)

  \item Demostreu que l'ordre és total: per a tots \(m,n\), \(m\le n\) o \(n\le m\). \begin{indicacio}Inducció sobre \(n\) amb \(m\) arbitrari, distingint \(n=m\) i \(n<m\); useu el pas del successor \(m\le n\iff s\circ m\le s\circ n\).\end{indicacio}

  \item Demostreu que tota part no inicial de \(\N\) té un element mínim únic. \begin{indicacio}Contrarecíproc via la part auxiliar \enquote{cap \(k\le n\) no és a \(P\)}, per inducció; vegeu el teorema~\ref{teo:bon-ordre-N}.\end{indicacio}

  \item Demostreu la recursió amb paràmetres: donades \(a\colon X\to A\) i \(t\colon X\times A\to A\), existeix una única \(u\colon X\times\N\to A\) amb les equacions corresponents. (Model: secció~\ref{prac:recursio-parametres}; teorema~\ref{teo:recursio-parametres}.)

  \item \textbf{Sobre punts o sobre elements generalitzats.} Classifiqueu els resultats següents segons si s'han enunciat i demostrat per a elements generalitzats o només per a punts globals, i digueu en cada cas per què: les lleis de l'aritmètica; \enquote{una suma nul·la té els sumands nuls}; la totalitat de l'ordre; el teorema del bon ordre. \begin{indicacio}Les igualtats d'aplicacions valen per a elements generalitzats (inducció amb paràmetres); les afirmacions amb \enquote{o} o \enquote{existeix} es demostren sobre punts, i la separació global les converteix en afirmacions sobre parts (observació~\ref{obs:ordre-sobre-punts}).\end{indicacio}
\end{enumerate}

\chapter{Conjunts finits i cardinalitat finita}\label{cap:exercicis-finits}

Aquests exercicis proposats corresponen al capítol homònim de la part teòrica (capítol~\ref{cap:finits} de la Part~\ref{part:teoria}) i es resolen amb les tècniques dels exercicis resolts de la part pràctica (capítol~\ref{cap:practica-finits} de la Part~\ref{part:practica}). Cada secció remet al bloc de resolts que en dona el model. Els marcats amb una indicació són els més exigents.

\section{Segments inicials i ordinals}

\begin{enumerate}
  \item Demostreu que \(\seg{0}\cong\mathbf0\) i que \(\seg{s\circ n}\cong\seg{n}+\mathbf1\). (Model: secció~\ref{prac:segments}; proposició~\ref{prop:segment-successor}.)

  \item Demostreu que \([1]\cong\mathbf1\) i \([2]\cong\mathbf2\), i identifiqueu la injecció \(\iota_k\) com un monomorfisme amb complement \(\ast_k\). (Model: proposició~\ref{prop:ordinals-petits}.)

  \item Demostreu que \(\seg{n}\) és una part de \(\N\) i que \(k\mathrel{\varepsilon}\seg{n}\) equival a \(k<n\). \begin{indicacio}\(\seg{n}\) és la imatge inversa de la relació \(\operatorname{Lt}\) (definició~\ref{def:ordre-naturals}) fixant el segon argument a \(n\).\end{indicacio}
\end{enumerate}

\section{Equipotència i cardinal}

\begin{enumerate}[resume]
  \item Demostreu que un punt de \(\operatorname{Isos}(A,B)\) és un isomorfisme \(A\to B\), i que \(A\approx B\) si i només si \(\operatorname{Isos}(A,B)\) és no inicial. (Model: secció~\ref{prac:equipotencia}; observació~\ref{obs:iso-punts}.)

  \item Demostreu que \(\operatorname{Card}(A)\) és una part interna de \(\N\) amb \(n\in\operatorname{Card}(A)\iff\seg{n}\approx A\), sense cap axioma d'elecció. (Model: proposició~\ref{prop:card-interna}.)

  \item Demostreu que l'equipotència \(\approx\) és reflexiva, simètrica i transitiva (al metanivell). \begin{indicacio}Exhibiu en cada cas l'isomorfisme: la identitat, la inversa i la composició; vegeu la definició~\ref{def:equipotencia}.\end{indicacio}
\end{enumerate}

\section{Comparació, unicitat i calaixos}

\begin{enumerate}[resume]
  \item Demostreu que un isomorfisme \(\seg{m}\cong\seg{n}\) implica \(m=n\), a partir del cas monomòrfic del lema. (Model: secció~\ref{prac:comparacio}; lema~\ref{lema:comparacio-ordinals}.)

  \item Demostreu la unicitat del cardinal finit i que dos finits són isomorfs si i només si tenen el mateix cardinal. (Model: teorema~\ref{teo:unicitat-cardinal}.)

  \item Demostreu el principi dels calaixos: si \(|A|>|B|\), no hi ha cap monomorfisme \(A\hookrightarrow B\). (Model: secció~\ref{prac:calaixos}; teorema~\ref{teo:calaixos}.)

  \item Demostreu que, per a \(f\colon A\to B\) entre finits amb \(|A|=|B|\), ser epimorfisme implica ser isomorfisme. \begin{indicacio}Per l'elecció, \(f\) té secció \(s\); \(s\) és mono entre finits del mateix cardinal, doncs iso per la direcció ja provada, i \(f=s^{-1}\); vegeu el teorema~\ref{teo:endo-finit}. Alternativament, construïu la secció amb el bon ordre, sense elecció (model: secció~\ref{prac:calaixos}).\end{indicacio}

  \item Expliqueu per què la inducció sobre \(n\) en l'afirmació \enquote{\(A+\seg{n}\) és finit amb cardinal \(|A|+n\)} és una inducció interna, codificant les aplicacions \(\seg{j}\to A+\seg{n}\) dins un exponencial fix. \begin{indicacio}Observació~\ref{obs:recomptes-interns}; seguiu el model de \(\mathcal E_A\) (definició~\ref{def:card-A}).\end{indicacio}

  \item \textbf{Axiomes usats.} Digueu quins dels resultats següents usen l'elecció i quins no: la construcció de \(\operatorname{Card}(A)\); la comparació de segments; la unicitat del cardinal; mono \(\Rightarrow\) iso en cardinal igual; epi \(\Rightarrow\) iso en cardinal igual; els quatre recomptes. \begin{indicacio}Només epi \(\Rightarrow\) iso l'usa tal com està demostrat, i es pot evitar amb el bon ordre.\end{indicacio}
\end{enumerate}

\section{Recomptes}

\begin{enumerate}[resume]
  \item Demostreu que tota part d'un conjunt finit és finita amb \(|S|\le|A|\), i que el seu complement també és finit. (Model: secció~\ref{prac:parts-finites}; lema~\ref{lema:part-finita}.)

  \item Demostreu els recomptes de la suma i del producte: \(|A+B|=|A|+|B|\) i \(|A\times B|=|A|\cdot|B|\). \begin{indicacio}Inducció sobre \(|B|\), amb la descomposició successora i la distributivitat; vegeu els lemes~\ref{lema:recompte-suma} i~\ref{lema:recompte-producte}.\end{indicacio}

  \item Demostreu que \(|\mathcal P(A)|=2^{n}\) per a \(|A|=n\). (Model: teorema~\ref{teo:nombre-parts}.)

  \item Demostreu que \(|B^{A}|=m^{n}\) per a \(|A|=n\), \(|B|=m\), i deduïu el nombre d'aplicacions \(\seg{3}\to\seg{2}\) i el nombre de parts de \(\seg{3}\). (Model: secció~\ref{prac:recompte-aplicacions}; teorema~\ref{teo:nombre-aplicacions}.)
\end{enumerate}

\chapter{Infinitud, enumerabilitat i seqüències}\label{cap:exercicis-infinitud}

Aquests exercicis proposats corresponen al capítol homònim de la part teòrica (capítol~\ref{cap:infinitud} de la Part~\ref{part:teoria}) i es resolen amb les tècniques dels exercicis resolts de la part pràctica (capítol~\ref{cap:practica-infinitud} de la Part~\ref{part:practica}). Cada secció remet al bloc de resolts que en dona el model. Els marcats amb una indicació són els més exigents.

\section{Infinitud i infinitud de Dedekind}

\begin{enumerate}
  \item Demostreu que \(s\colon\N\to\N\) és un monomorfisme no exhaustiu i que, per tant, \(\N\) és infinit de Dedekind i infinit. (Model: secció~\ref{prac:infinits}; proposició~\ref{prop:N-dedekind}.)

  \item Demostreu que tot conjunt infinit de Dedekind és infinit, sense cap axioma nou. (Model: proposició~\ref{prop:dedekind-implica-infinit}.)

  \item Demostreu que si \(A\) és infinit de Dedekind, aleshores existeix un monomorfisme \(\N\hookrightarrow A\). \begin{indicacio}Iteneu el monomorfisme no exhaustiu \(f\) a partir d'un punt fora de la seva imatge; useu la recursió (capítol~\ref{cap:naturals}) per definir \(n\mapsto f^{n}(a_0)\) i comproveu que és injectiva.\end{indicacio}
\end{enumerate}

\section{Enumerables i numerables}

\begin{enumerate}[resume]
  \item Demostreu que tota imatge exhaustiva i tota part d'un conjunt enumerable són enumerables, sense elecció. (Model: secció~\ref{prac:enumerables}; proposició~\ref{prop:estabilitat-enumerable}.)

  \item Demostreu que la suma i el producte de dos conjunts enumerables són enumerables. \begin{indicacio}Useu \(\N\cong\N+\N\) i un aparellament \(\N\cong\N\times\N\); vegeu la proposició~\ref{prop:estabilitat-enumerable}.\end{indicacio}

  \item Demostreu que \(\N\times\N\) és enumerable exhibint explícitament un aparellament \(\N\cong\N\times\N\). \begin{indicacio}Model: secció~\ref{prac:enumerables}; la recursora de \enquote{següent parell en l'ordre per diagonals} és bijectiva. La fórmula tancada \(\tfrac12(i+j)(i+j+1)+j\) és correcta, però pressuposa la divisió.\end{indicacio}
\end{enumerate}

\section{Subobjectes de \texorpdfstring{\(\N\)}{N}}

\begin{enumerate}[resume]
  \item Demostreu que tota part de \(\N\) és enumerable i que és finita o isomorfa a \(\N\), sense elecció. (Model: secció~\ref{prac:subobjectes}; teorema~\ref{teo:subobjectes-N}.)

  \item Demostreu que un conjunt és enumerable infinit si i només si és isomorf a \(\N\). (Model: corol·lari~\ref{cor:enumerable-infinit-N}.)

  \item Expliqueu per què el principi del mínim permet evitar l'elecció en aquests resultats, i on deixaria de funcionar per a un conjunt infinit qualsevol. \begin{indicacio}El mínim d'una part no inicial de \(\N\) és únic i canònic; en un infinit sense ordre no hi ha tria distingida (observació~\ref{obs:minim-vs-eleccio}).\end{indicacio}
\end{enumerate}

\section{Llistes i unions numerables}

\begin{enumerate}[resume]
  \item Demostreu el principi de recursió sobre \(\operatorname{List}(A)\): per a tot \((B,b,t)\) hi ha una única \(h\) amb \(h\circ\varepsilon=b\) i \(h\circ\operatorname{cons}=t\circ(\id_A\times h)\). (Model: secció~\ref{prac:llistes}; proposició~\ref{prop:recursio-llistes}.)

  \item Definiu, per recursió sobre llistes, la llargada \(\ell\), la concatenació amb una llista fixada i un plegat d'una operació binària. \begin{indicacio}Trieu \((B,b,t)\) adequats: \((\N,0,s\circ\pi_2)\) per a la llargada; itereu \(\operatorname{cons}\) per a la concatenació.\end{indicacio}

  \item Demostreu que si \(A\) és enumerable, aleshores \(\operatorname{List}(A)\) és enumerable, sense elecció. (Model: observació~\ref{obs:llistes-enumerables}.)

  \item Demostreu, amb l'elecció numerable, que la unió disjunta d'una família numerable d'enumerables és enumerable, i identifiqueu el pas que usa l'elecció. (Model: secció~\ref{prac:unions}; teorema~\ref{teo:unio-numerable} i observació~\ref{obs:eleccio-unio}.)

  \item A la prova del teorema~\ref{teo:infinit-implica-dedekind}, expliqueu per què \enquote{\(a_n\) és el primer valor de \(f_n\) que no és cap \(a_i\) amb \(i<n\)} no és una recursió de l'axioma~\ref{ax:naturals}, i comproveu amb detall que la construcció amb llistes de llargada \(2^n\) és injectiva. \begin{indicacio}La regla depèn de tota la història \(a_0,\dots,a_{n-1}\), no només de \(a_{n-1}\); amb \(V_n\) definit a partir de les \(f_{2^m-1}\), \(a_m\in V_n\) per a \(m<n\) i \(a_n\notin V_n\).\end{indicacio}

  \item Demostreu que \(\operatorname{List}(\N)\cong\N\) mitjançant la codificació \(c(\varepsilon)=0\), \(c(\operatorname{cons}(k,u))=s(\pi(k,c(u)))\), amb \(\pi\) un aparellament bijectiu. \begin{indicacio}Injectivitat per inducció sobre la llargada; exhaustivitat per inducció sobre el codi, usant que tot natural és \(0\) o \(s(\pi(k,m))\) amb \(m\) més petit que el codi.\end{indicacio}

  \item \textbf{Axiomes usats.} Per a cadascun dels resultats següents, digueu si usa l'elecció numerable o no: \(\N\) és infinit de Dedekind; Dedekind implica infinit; infinit implica Dedekind; estabilitat de l'enumerabilitat; tota part de \(\N\) és enumerable; \(\operatorname{List}(A)\) és enumerable si \(A\) ho és; la unió numerable d'enumerables és enumerable. \begin{indicacio}Només el tercer i l'últim.\end{indicacio}
\end{enumerate}

\chapter[Schröder--Bernstein i Cantor]{Comparació de conjunts,\texorpdfstring{\\}{ }Schröder--Bernstein i Cantor}\label{cap:exercicis-cantor}

Aquests exercicis proposats corresponen al capítol homònim de la part teòrica (capítol~\ref{cap:cantor} de la Part~\ref{part:teoria}) i es resolen amb les tècniques dels exercicis resolts de la part pràctica (capítol~\ref{cap:practica-cantor} de la Part~\ref{part:practica}). Cada secció remet al bloc de resolts que en dona el model. Els marcats amb una indicació són els més exigents.

\section{La comparació de conjunts}

\begin{enumerate}
  \item Demostreu que \(\preceq\) és reflexiva i transitiva i invariant per equipotència. (Model: secció~\ref{prac:comparacio-conjunts}; proposició~\ref{prop:preceq-preordre}.)

  \item Demostreu que tota part \(S\hookrightarrow A\) compleix \(S\preceq A\), i que una aplicació exhaustiva amb secció donada \(g\colon A\twoheadrightarrow C\) dona \(C\preceq A\). Indiqueu on caldria l'elecció si la secció no vingués donada. (Model: secció~\ref{prac:comparacio-conjunts}; proposició~\ref{prop:epi-regular-eleccio}.)

  \item Demostreu que \(A\preceq B\) si i només si \(A\) és equipotent a alguna part de \(B\). \begin{indicacio}Un monomorfisme \(A\hookrightarrow B\) es factoritza per la seva imatge (teorema~\ref{teo:factoritzacio}), que és una part de \(B\) isomorfa a \(A\).\end{indicacio}

  \item Demostreu que \(\N\preceq A\) si i només si \(A\) és infinit de Dedekind, i relacioneu-ho amb la infinitud sota l'elecció numerable. \begin{indicacio}Un monomorfisme \(\N\hookrightarrow A\) dona un endomonomorfisme no exhaustiu transportant el successor; vegeu el teorema~\ref{teo:infinit-implica-dedekind}.\end{indicacio}
\end{enumerate}

\section{El teorema de Schröder--Bernstein}

\begin{enumerate}[resume]
  \item Donats \(f\colon A\hookrightarrow B\) i \(g\colon B\hookrightarrow A\), demostreu que \(\Phi(X)=A\setminus g(B\setminus f(X))\) és creixent sobre \(P(A)\). (Model: secció~\ref{prac:sb}.)

  \item Demostreu que tot operador creixent \(\Phi\colon P(A)\to P(A)\) té un menor punt fix, donat per la intersecció de les seves parts \(\Phi\)-tancades. (Model: secció~\ref{prac:sb}; teorema de Knaster--Tarski.)

  \item A partir del punt fix, completeu la construcció de l'isomorfisme \(A\cong B\) i comproveu-ne la bijectivitat sobre punts; expliqueu per què és legítim comprovar-la sobre punts i no sobre elements generalitzats. (Model: teorema~\ref{teo:schroder-bernstein}.)

  \item Deduïu el \enquote{teorema de l'entrepà}: si \(A\subseteq B\subseteq C\) i \(A\approx C\), aleshores \(A\approx B\approx C\). \begin{indicacio}Les inclusions donen dominacions; encadeneu-les amb l'isomorfisme \(A\approx C\) i apliqueu Schröder--Bernstein.\end{indicacio}

  \item Demostreu que \([0,1]\approx[0,1)\) construint dues injeccions explícites i aplicant Schröder--Bernstein, sense exhibir cap bijecció directa. \begin{indicacio}\([0,1)\hookrightarrow[0,1]\) és la inclusió; \([0,1]\hookrightarrow[0,1)\) es pot donar per \(x\mapsto x/2\).\end{indicacio}
\end{enumerate}

\section{L'argument diagonal i el teorema de Cantor}

\begin{enumerate}[resume]
  \item Demostreu que la part diagonal \(D=\{a : a\not\mathrel{\varepsilon} f(a)\}\) no és el valor de cap element, i concloeu que cap \(f\colon A\to P(A)\) no és exhaustiva. (Model: secció~\ref{prac:cantor}; teorema~\ref{teo:diagonal}.)

  \item Demostreu que \(A\prec P(A)\) per a tot \(A\). (Model: teorema~\ref{teo:cantor}.)

  \item Demostreu que no existeix cap conjunt universal ni cap mida màxima. (Model: corol·lari~\ref{cor:no-universal}.)

  \item Demostreu que \(\mathbf2^{\N}\) no és enumerable mitjançant la successió diagonal, i concloeu que \(P(\N)\) tampoc. \begin{indicacio}Donada \(f\colon\N\to\mathbf2^{\N}\), la successió \(d(n)=1-f(n)(n)\) difereix de cada \(f(n)\); vegeu l'observació~\ref{obs:cantor-fort}.\end{indicacio}

  \item Demostreu que si \(A\preceq B\), aleshores \(P(A)\preceq P(B)\), i que si a més \(A\prec B\) i \(B\) no és inicial, es manté \(P(A)\prec P(B)\) sota la comparabilitat total (que assumireu). \begin{indicacio}Un monomorfisme \(A\hookrightarrow B\) indueix \(P(A)\to P(B)\) per imatge directa, injectiva; per a l'estrictesa, useu Cantor i la comparabilitat.\end{indicacio}

  \item \textbf{(Punt de reflexió.)} Expliqueu per què l'argument diagonal no produeix cap paradoxa en aquest llibre, tot i que la \enquote{col·lecció de tots els conjunts} no existeix. \begin{indicacio}La pertinença és sempre relativa a una part d'un conjunt fix; no hi ha cap relació de pertinença global sobre la qual diagonalitzar (observació~\ref{obs:cantor-classica}).\end{indicacio}

  \item \textbf{Axiomes usats.} Digueu quins axiomes intervenen en cadascun dels resultats següents, i confirmeu que cap no usa l'elecció: que \(\preceq\) és un preordre; l'existència del menor punt fix de \(\Phi\); que \(h\) és un isomorfisme; que cap \(A\to P(A)\) és exhaustiva; que no hi ha conjunt universal. \begin{indicacio}Composició; conjunt de parts i quantificador universal; separació global (per comprovar sobre punts) i teorema del bimorfisme; pertinença interna, complement i \(0\ne1\); els dos teoremes anteriors.\end{indicacio}
\end{enumerate}

\chapter{Famílies, fibres i operacions indexades}\label{cap:exercicis-families}

Aquests exercicis proposats corresponen al capítol homònim de la part teòrica (capítol~\ref{cap:families} de la Part~\ref{part:teoria}) i es resolen amb les tècniques dels exercicis resolts de la part pràctica (capítol~\ref{cap:practica-families} de la Part~\ref{part:practica}). Cada secció remet al bloc de resolts que en dona el model. Els marcats amb una indicació són els més exigents.

\section{Famílies i fibres}

\begin{enumerate}
  \item Demostreu que la fibra \(E_i=\mathbf1\times_I E\) és la imatge inversa \(p^{-1}(i)\) i que la projecció \(E_i\to E\) és un monomorfisme. (Model: secció~\ref{prac:families-fibres}; observació~\ref{obs:fibra-part}.)

  \item Descriviu la família constant \(I\times A\to I\) i comproveu que totes les fibres són isomorfes a \(A\). (Model: definició~\ref{def:familia-constant}.)

  \item Sigui \(m\colon S\hookrightarrow I\) una part, vista com a característica \(\chi_S\colon I\to\mathbf2\). Descriviu la família que representa i les seves fibres. \begin{indicacio}La fibra sobre \(i\) és \(\mathbf1\) si \(i\in S\) i \(\mathbf0\) altrament; vegeu l'observació~\ref{obs:familia-exemples}.\end{indicacio}
\end{enumerate}

\section{Seccions i elecció}

\begin{enumerate}[resume]
  \item Demostreu que una secció de \(p\colon E\to I\) tria un element de cada fibra, i que per a \(p\) epimorfisme l'existència d'una secció és l'axioma d'elecció. (Model: secció~\ref{prac:seccions}; observació~\ref{obs:seccio-eleccio}.)

  \item Demostreu que si cada fibra de \(p\colon E\to I\) porta un mínim per a un bon ordre intern donat, aleshores \(p\) té una secció \emph{sense} elecció. \begin{indicacio}Definiu \(s(i)\) com el mínim de la fibra \(E_i\); és únic i canònic, com als subobjectes de \(\N\) (observació~\ref{obs:minim-vs-eleccio}).\end{indicacio}
\end{enumerate}

\section{Suma i producte}

\begin{enumerate}[resume]
  \item Demostreu que \(\prod_{i\in I}E_i=\{s\in E^I:p\circ s=\id_I\}\) existeix com a part interna de \(E^I\) i que els seus elements globals són les seccions de \(p\). (Model: secció~\ref{prac:suma-producte}; proposició~\ref{prop:producte-existeix}.)

  \item Distingiu, amb el producte, \enquote{existeix com a objecte} de \enquote{té elements}, i localitzeu l'elecció. (Model: observació~\ref{obs:objecte-vs-elements}.)

  \item Demostreu que la suma d'una família \(p\colon E\to I\) existeix sempre i és el propi total \(E\), i que per a la família constant coincideix amb \(I\times A\). \begin{indicacio}No cal cap axioma: el total és una de les dades (observació~\ref{obs:suma-existeix}).\end{indicacio}

  \item Sigui \(I=\mathbf2\) i \(p\colon E\to\mathbf2\) amb fibres \(E_0,E_1\). Demostreu que \(\sum_{i}E_i\cong E_0+E_1\) i \(\prod_{i}E_i\cong E_0\times E_1\). \begin{indicacio}Sobre un índex de dos elements, suma i producte indexats recuperen la suma i el producte binaris.\end{indicacio}
\end{enumerate}

\section{Famílies de parts i operacions indexades}

\begin{enumerate}[resume]
  \item Definiu \(\bigcup_i\alpha(i)\) i \(\bigcap_i\alpha(i)\) amb els quantificadors interns i demostreu que són el suprem i l'ínfim a \((P(A),\subseteq)\). (Model: secció~\ref{prac:parts-indexades}; proposició~\ref{prop:unio-interseccio-propietats}.)

  \item Demostreu les dues lleis de De Morgan indexades. \begin{indicacio}\(\neg\exists_{\pi_A}(R)=\forall_{\pi_A}(\neg R)\); vegeu la proposició~\ref{prop:quantificadors-adjuncio}.\end{indicacio}

  \item Demostreu que una família \(\alpha\colon I\to P(A)\) de parts no buides, disjuntes i que cobreixen \(A\) determina una partició \(q\colon A\twoheadrightarrow I\) amb fibres \(\alpha(i)\). (Model: secció~\ref{prac:parts-indexades}; proposició~\ref{prop:familia-particio}.)

  \item Demostreu la distributivitat \(\beta\cap\bigcup_i\alpha(i)=\bigcup_i(\beta\cap\alpha(i))\) per a una part fixa \(\beta\hookrightarrow A\). \begin{indicacio}Traduïu-ho a \enquote{existeix \(i\) amb \(a\in\beta\) i \(a\in\alpha(i)\)} i useu l'adjunció de \(\exists_{\pi_A}\).\end{indicacio}
\end{enumerate}

\section{Regla externa contra família interna}

\begin{enumerate}[resume]
  \item Expliqueu per què \(n\mapsto\seg{n}\) és una regla externa representable per una família interna, i com. (Model: secció~\ref{prac:regla-vs-familia}; observació~\ref{obs:regla-externa-problema}.)

  \item Descriviu una regla externa \emph{no} representable amb els axiomes de Nivell~I i digueu quin principi caldria. (Model: observacions~\ref{obs:regla-externa-problema} i~\ref{obs:sr}.)

  \item \textbf{(Punt de reflexió.)} Argumenteu per què la representabilitat d'una regla externa \emph{no} és una qüestió d'elecció sinó de substitució. \begin{indicacio}L'elecció tria dins fibres d'una família ja existent; la substitució garanteix l'existència de la família (observació~\ref{obs:sr-no-eleccio}).\end{indicacio}

  \item \textbf{Axiomes usats.} Per a cadascun dels resultats següents, digueu si usa l'elecció, la separació global, cap dels dos, o si queda fora del Nivell~I: l'existència del producte com a objecte; la no-vacuïtat del producte d'una família de fibres no buides; la unió i la intersecció indexades; que una família de parts disjuntes que cobreixen és una partició; la representació de \(n\mapsto\seg{n}\) per una família interna; la representació d'una regla sense ambient comú. \begin{indicacio}Cap; elecció; cap; separació global; cap (exponencial fix); SR, fora del Nivell~I.\end{indicacio}
\end{enumerate}

\chapter{Ordres, bons ordres i segments}\label{cap:exercicis-bons-ordres}

Aquests exercicis proposats corresponen al capítol homònim de la part teòrica (capítol~\ref{cap:bons-ordres} de la Part~\ref{part:teoria}) i es resolen amb les tècniques dels exercicis resolts de la part pràctica (capítol~\ref{cap:practica-bons-ordres} de la Part~\ref{part:practica}). Cada secció remet al bloc de resolts que en dona el model. Els marcats amb una indicació són els més exigents.

\section{Ordres i bons ordres}

\begin{enumerate}
  \item Demostreu que tot bon ordre és total i que el mínim d'una part no inicial és únic. (Model: secció~\ref{prac:ordres}; proposició~\ref{prop:bon-ordre-total}.)

  \item Demostreu que \(\Z\) amb l'ordre usual és total però no és un bon ordre. (Model: observació~\ref{obs:bon-ordre-classic}.)

  \item Demostreu que un ordre estricte \({<}\) prohibeix les cadenes descendents finites \(w_0>w_1>\cdots>w_n>w_0\) (cap cicle), i que un bon ordre prohibeix a més tota cadena descendent infinita. \begin{indicacio}Un cicle contradiu la irreflexivitat via la transitivitat; una cadena infinita té imatge sense mínim.\end{indicacio}

  \item Sigui \({<}\) un bon ordre sobre \(W\) no inicial. Demostreu que \(W\) té un element mínim global \(\bot\), i que \(W_{<\bot}\) és inicial. \begin{indicacio}Apliqueu el bon ordre a la part total \(W\).\end{indicacio}
\end{enumerate}

\section{Segments inicials}

\begin{enumerate}[resume]
  \item Demostreu que tota part d'un bon ordre, amb l'ordre restringit, és un bon ordre. (Model: secció~\ref{prac:segments-bo}; proposició~\ref{prop:segment-ben-ordenat}.)

  \item Demostreu que en un bon ordre les parts inicials pròpies són exactament els segments \(W_{<w}\), i que el punt \(w\) és únic. (Model: proposicions~\ref{prop:parts-inicials} i~\ref{prop:segment-determina-punt}.)

  \item Demostreu que la intersecció i la unió de dues parts inicials són parts inicials, i que per a \(v\le w\) es té \(W_{<v}\subseteq W_{<w}\). \begin{indicacio}La condició de tancament cap avall es preserva per intersecció i unió; per als segments, useu la transitivitat.\end{indicacio}
\end{enumerate}

\section{Creixents i rigidesa}

\begin{enumerate}[resume]
  \item Demostreu que en un bon ordre cap punt no baixa sota una aplicació creixent injectiva: \(w\le f(w)\). (Model: secció~\ref{prac:rigidesa}; lema~\ref{lema:rigidesa}.)

  \item Demostreu que entre dos bons ordres hi ha com a molt un isomorfisme d'ordre, i que l'únic automorfisme és la identitat. (Model: teorema~\ref{teo:unicitat-iso}.)

  \item Demostreu que cap bon ordre és isomorf a un segment inicial propi seu. (Model: corol·lari~\ref{cor:no-segment-propi}.)

  \item Demostreu que sobre un ordre \emph{total} tota aplicació creixent és injectiva, i doneu un contraexemple sobre un ordre no total. \begin{indicacio}Per a la injectivitat useu la tricotomia puntual; per al contraexemple, una antisimetria fallida en un ordre parcial.\end{indicacio}
\end{enumerate}

\section{Comparació de bons ordres}

\begin{enumerate}[resume]
  \item Demostreu que la comparació \enquote{ser un segment de} és antisimètrica i transitiva sobre bons ordres. (Model: secció~\ref{prac:comparacio-ordres}; proposició~\ref{prop:comparacio-ordres}.)

  \item Demostreu que, entre dos bons ordres, com a molt una de les tres situacions de la tricotomia es dona. (Model: teorema~\ref{teo:tricotomia}.)

  \item \textbf{(Punt de reflexió.)} Expliqueu per què la comparabilitat total de dos bons ordres necessita recursió transfinita però \emph{no} elecció, a diferència de la comparabilitat de dos conjunts qualssevol. \begin{indicacio}L'aparellament es fa prenent a cada pas el mínim del que queda, únic i canònic; no hi ha cap tria (observació~\ref{obs:tricotomia-recursio}).\end{indicacio}

  \item \textbf{Punts i parts.} Per a cadascuna de les afirmacions següents, digueu si és una inclusió de parts (vàlida a tots els estadis) o una igualtat amb una unió (comprovable sobre punts): la transitivitat; la irreflexivitat; la totalitat; que una aplicació és creixent; la tricotomia. \begin{indicacio}Les tres primeres i la quarta són inclusions o disjuncions de parts; la totalitat i la tricotomia són unions totals, i es comproven sobre punts (lema~\ref{lema:transferencia-punts}).\end{indicacio}

  \item \textbf{Axiomes usats.} Confirmeu que cap resultat del capítol no usa l'elecció, i digueu on intervé la separació global: tot bon ordre és total; els segments són ben ordenats; les parts inicials pròpies són segments; la rigidesa; la unicitat de l'isomorfisme d'ordre. \begin{indicacio}La separació global intervé a la totalitat (pas de punts a parts) i, a través d'ella, a \enquote{creixent implica mono}.\end{indicacio}
\end{enumerate}

\chapter{Inducció i recursió transfinites}\label{cap:exercicis-recursio}

Aquests exercicis proposats corresponen al capítol homònim de la part teòrica (capítol~\ref{cap:recursio} de la Part~\ref{part:teoria}) i es resolen amb les tècniques dels exercicis resolts de la part pràctica (capítol~\ref{cap:practica-recursio} de la Part~\ref{part:practica}). Cada secció remet al bloc de resolts que en dona el model. Els marcats amb una indicació són els més exigents.

\section{Inducció transfinita}

\begin{enumerate}
  \item Demostreu el principi d'inducció transfinita: si \(W_{<w}\subseteq S\) implica \(w\in S\) per a tot \(w\), aleshores \(S=W\). (Model: secció~\ref{prac:induccio-transfinita}; teorema~\ref{teo:induccio-transfinita}.)

  \item Expliqueu per què la inducció transfinita no necessita cas base separat. (Model: observació~\ref{obs:induccio-propietats}.)

  \item Deduïu de la inducció transfinita el \emph{principi del mínim} per a un bon ordre: tota part no inicial té mínim. \begin{indicacio}Contrarecíproc: si una part no té mínim, useu la inducció per veure que el seu complement és tot \(W\).\end{indicacio}
\end{enumerate}

\section{Recursió amb codomini fix}

\begin{enumerate}[resume]
  \item Demostreu que dues solucions parcials d'una regla \(\Phi\) coincideixen a la intersecció dels dominis. (Model: secció~\ref{prac:recursio-fix}; teorema~\ref{teo:recursio-transfinita}.)

  \item Demostreu l'existència i unicitat de \(F\colon W\to X\) amb \(F(w)=\Phi(w,F|_{W_{<w}})\). (Model: teorema~\ref{teo:recursio-transfinita}.)

  \item Justifiqueu que la construcció de \(F\) no usa cap elecció i que tot viu dins un exponencial fix. \begin{indicacio}Cada valor és \(\Phi\) de la història; la reunió de solucions compatibles és canònica (observació~\ref{obs:recursio-sense-eleccio}).\end{indicacio}

  \item Demostreu que la recursió primitiva amb paràmetres sobre \(\N\) (teorema~\ref{teo:recursio-parametres}) és el cas particular de la recursió transfinita en què \(\Phi\) només mira el terme anterior. \begin{indicacio}Vegeu l'observació~\ref{obs:recursio-N-cas}; a \(\N\) tot no-zero és successor.\end{indicacio}
\end{enumerate}

\section{Comparabilitat total}

\begin{enumerate}[resume]
  \item Demostreu que dos bons ordres qualssevol són comparables, construint per recursió l'aparellament pel mínim disponible. (Model: secció~\ref{prac:comparabilitat}; teorema~\ref{teo:comparabilitat}.)

  \item Expliqueu per què aquesta comparabilitat no usa elecció, en contrast amb la de conjunts qualssevol. (Model: observació~\ref{obs:comparabilitat-vs-eleccio}.)

  \item Deduïu de la comparabilitat total que la comparació per segments és un ordre total (llevat d'isomorfisme d'ordre) sobre els bons ordres. \begin{indicacio}Combineu l'exclusivitat (teorema~\ref{teo:tricotomia}) amb l'existència (teorema~\ref{teo:comparabilitat}).\end{indicacio}
\end{enumerate}

\section{Exemples de recursió}

\begin{enumerate}[resume]
  \item Definiu per recursió una successió transfinita de parts d'un fix \(A\) amb \(S_w=\bigcup_{u<w}g(S_u)\) i justifiqueu que és de Nivell~I. (Model: secció~\ref{prac:exemples-recursio}; observació~\ref{obs:exemples-recursio}.)

  \item Definiu per recursió el tancament d'una part \(S_0\hookrightarrow A\) sota una operació monòtona \(g\colon P(A)\to P(A)\), i relacioneu-ho amb el menor punt fix de Schröder--Bernstein. \begin{indicacio}Itereu \(g\) transfinitament dins el fix \(P(A)\); vegeu el teorema~\ref{teo:schroder-bernstein}.\end{indicacio}

  \item Expliqueu per què \(V_0=A\), \(V_{n+1}=P(V_n)\) no es defineix amb la recursió de codomini fix, i quin principi cal. (Model: observació~\ref{obs:frontera-exemple}.)

  \item \textbf{(Punt de reflexió.)} Argumenteu per què separar en casos successor i límit és legítim només després de provar que la separació és exhaustiva per al bon ordre considerat. \begin{indicacio}La formulació general amb la història no ho requereix; la separació és una comoditat que pot amagar casos (observació~\ref{obs:successor-limit}).\end{indicacio}

  \item Comproveu que la reunió de totes les solucions parcials és una aplicació: que la relació \(G=\{(w,x):\exists a\in\mathcal A\,(a(w)=x)\}\) és funcional i té domini una part inicial. \begin{indicacio}La funcionalitat és la compatibilitat; useu la proposició~\ref{prop:aplicacio-funcional}.\end{indicacio}

  \item \textbf{Axiomes usats.} Digueu quins dels resultats següents usen elecció, i quins SR: la inducció transfinita; la recursió amb codomini fix; la comparabilitat total dels bons ordres; la successió \((P^n(A))_{n\in\N}\). \begin{indicacio}Cap dels tres primers; el quart queda fora del Nivell~I.\end{indicacio}
\end{enumerate}

\chapter{Zorn, Hartogs i el teorema del bon ordre}\label{cap:exercicis-zorn}

Aquests exercicis proposats corresponen al capítol homònim de la part teòrica (capítol~\ref{cap:zorn} de la Part~\ref{part:teoria}) i es resolen amb les tècniques dels exercicis resolts de la part pràctica (capítol~\ref{cap:practica-zorn} de la Part~\ref{part:practica}). Cada secció remet al bloc de resolts que en dona el model. Els marcats amb una indicació són els més exigents.

\section{Hartogs}

\begin{enumerate}
  \item Demostreu que \(\Hart(A)\), el conjunt de tipus d'ordre dels bons ordres sobre parts d'\(A\) ordenat per segment inicial, és un bon ordre. (Model: secció~\ref{prac:hartogs}; proposició~\ref{prop:hartogs-bon-ordre}.)

  \item Demostreu que \(\Hart(A)\) no admet cap monomorfisme en \(A\). (Model: teorema~\ref{teo:hartogs}.)

  \item Expliqueu per què tota la construcció de \(\Hart(A)\) és \emph{acotada} i no usa cap classe global d'ordinals. \begin{indicacio}Els bons ordres sobre parts d'\(A\), codificats com a parells \((S,R)\), viuen dins \(P(A)\times P(A\times A)\); el quocient per isomorfisme és efectiu (observació~\ref{obs:hartogs-acotat}).\end{indicacio}

  \item Compareu els teoremes de Hartogs i de Cantor com a resultats \enquote{diagonals}. \begin{indicacio}\(P(A)\) dona una mida més gran; \(\Hart(A)\), una longitud d'ordre més gran (observació~\ref{obs:hartogs-cantor}).\end{indicacio}
\end{enumerate}

\section{El teorema del bon ordre}

\begin{enumerate}[resume]
  \item Amb l'elecció, demostreu que tot conjunt admet un bon ordre, amb la recursió acotada per \(\Hart(A)\). (Model: secció~\ref{prac:bon-ordre}; teorema~\ref{teo:bon-ordre}.)

  \item Expliqueu quin problema evita la cota de Hartogs respecte d'una recursió \enquote{sobre tots els ordinals}. (Model: observació~\ref{obs:cota-hartogs}.)

  \item Demostreu que si \(A\) admet un bon ordre, aleshores \(A\preceq\Hart(A)\) o \(\Hart(A)\preceq A\), i deduïu-ne quin dels dos casos es dona. \begin{indicacio}Per Hartogs, \(\Hart(A)\not\preceq A\); useu la comparabilitat de bons ordres.\end{indicacio}
\end{enumerate}

\section{El lema de Zorn}

\begin{enumerate}[resume]
  \item Doneu un ordre parcial amb diversos maximals i cap màxim. (Model: secció~\ref{prac:zorn}; observació~\ref{obs:maximal-maxim}.)

  \item Amb l'elecció, demostreu el lema de Zorn. (Model: teorema~\ref{teo:zorn}.)

  \item Demostreu, amb Zorn, que tot espai vectorial té una base. (Model: proposició~\ref{prop:base}.)

  \item Demostreu, amb Zorn, que tot ideal propi d'un anell commutatiu amb unitat es conté en un ideal maximal. \begin{indicacio}Ordeneu per inclusió els ideals propis que contenen l'ideal donat; la unió d'una cadena d'ideals propis és un ideal propi (no conté \(1\)); vegeu l'observació~\ref{obs:zorn-aplicacions}.\end{indicacio}
\end{enumerate}

\section{Les tres equivalències}

\begin{enumerate}[resume]
  \item Demostreu que (bon ordre) implica (elecció). (Model: secció~\ref{prac:equivalencies}; teorema~\ref{teo:equivalencies}.)

  \item Demostreu que (Zorn) implica (bon ordre). (Model: teorema~\ref{teo:equivalencies}.)

  \item \textbf{(Punt de reflexió.)} Expliqueu per què cap de les direccions de l'equivalència no ha necessitat una recursió més forta que la del capítol~\ref{cap:recursio}. \begin{indicacio}Totes les recursions són acotades per un \(\Hart\) i de codomini fix; vegeu el final del teorema~\ref{teo:equivalencies}.\end{indicacio}

  \item Expliqueu per què la regla de recursió de la prova de Zorn s'ha de donar també en el cas en què la història no és una cadena, encara que aquest cas no es doni mai. \begin{indicacio}El teorema~\ref{teo:recursio-transfinita} demana una regla \(\Phi\) definida per a tota parella \((w,h)\); una clàusula \enquote{altrament} amb un punt fix de \(P\) hi basta.\end{indicacio}

  \item Demostreu amb detall que la unió d'una cadena de bons ordres parcials, ordenada per extensió inicial, és un bon ordre parcial. \begin{indicacio}Per a una part no buida \(T\), preneu un membre \((S,R)\) de la cadena que la toqui; el mínim de \(T\cap S\) ho és de \(T\) perquè les extensions són inicials.\end{indicacio}

  \item \textbf{Axiomes usats.} Digueu quins dels resultats següents usen l'elecció: que \(\Hart(A)\) existeix; que \(\Hart(A)\) és un bon ordre; el teorema de Hartogs; el teorema del bon ordre; el lema de Zorn; que (BO) implica (E). \begin{indicacio}Només el quart i el cinquè; el sisè construeix l'elecció a partir del bon ordre, sense suposar-la.\end{indicacio}
\end{enumerate}

\chapter{Tipus ordinals}\label{cap:exercicis-ordinals}

Aquests exercicis proposats corresponen al capítol homònim de la part teòrica (capítol~\ref{cap:ordinals} de la Part~\ref{part:teoria}) i es resolen amb les tècniques dels exercicis resolts de la part pràctica (capítol~\ref{cap:practica-ordinals} de la Part~\ref{part:practica}). Cada secció remet al bloc de resolts que en dona el model. Els marcats amb una indicació són els més exigents.

\section{Tipus ordinals i el seu ordre}

\begin{enumerate}
  \item Demostreu que \(\alpha<\beta\) no depèn dels representants triats. (Model: secció~\ref{prac:tipus-ordinals}; definició~\ref{def:ordre-ordinals}.)

  \item Demostreu que \(<\) sobre tipus ordinals és un ordre estricte total. (Model: proposició~\ref{prop:ordinals-total}.)

  \item Expliqueu per què l'ordre ordinal es defineix per segments i no per \(\alpha\in\beta\). (Model: observació~\ref{obs:no-pertinença-ordinal}.)

  \item Demostreu que la restricció de \(<\) als tipus dels bons ordres sobre parts d'un fix \(A\) ---és a dir, sobre \(\Hart(A)\)--- és un bon ordre. \begin{indicacio}És la proposició~\ref{prop:hartogs-bon-ordre}, rellegida sobre tipus.\end{indicacio}
\end{enumerate}

\section{Zero, successor i límit}

\begin{enumerate}[resume]
  \item Demostreu que \(\alpha+1\) és el menor tipus estrictament més gran que \(\alpha\), i que té element màxim. (Model: secció~\ref{prac:zero-successor-limit}; proposició~\ref{prop:successor-invariant}.)

  \item Demostreu que tot tipus ordinal és exactament un de zero, successor o límit. (Model: proposició~\ref{prop:tricotomia-successor}.)

  \item Demostreu que \(\omega=\operatorname{tipus}(\N)\) és el primer tipus límit. (Model: observació~\ref{obs:limit-classic}.)

  \item Demostreu que un tipus \(\lambda\) és límit si i només si \(\lambda=\sup_{\alpha<\lambda}\alpha\) (el suprem dels tipus estrictament menors, quan aquest està acotat). \begin{indicacio}Un successor \(\alpha+1\) no és el suprem dels menors, perquè \(\alpha\) n'és el màxim; un límit no té màxim entre els menors.\end{indicacio}
\end{enumerate}

\section{Suprems acotats}

\begin{enumerate}[resume]
  \item Demostreu que una família de tipus \(\alpha_i=\operatorname{tipus}(W_{<r(i)})\) representada dins un bon ordre \(W\) té suprem, donat per la reunió \(\bigcup_iW_{<r(i)}\) dels segments oberts, i doneu un exemple on la reunió dels segments tancats \(W_{\le r(i)}\) dona un tipus estrictament més gran. (Model: secció~\ref{prac:suprems}; proposició~\ref{prop:suprem}.)

  \item Expliqueu per què l'acotació dins un bon ordre comú és imprescindible. (Model: observació~\ref{obs:suprem-acotat}.)

  \item Demostreu que \(\sup\) és monòton: si \(\alpha_i\le\beta_i\) per a tot \(i\) i totes dues famílies estan acotades dins un mateix \(W\), aleshores \(\sup_i\alpha_i\le\sup_i\beta_i\). \begin{indicacio}Les reunions de segments respecten la inclusió.\end{indicacio}
\end{enumerate}

\section{Burali--Forti}

\begin{enumerate}[resume]
  \item Demostreu que no existeix cap conjunt ben ordenat de tots els tipus ordinals amb la propietat de segments. (Model: secció~\ref{prac:burali-forti}; teorema~\ref{teo:burali-forti}.)

  \item Expliqueu el paral·lelisme entre Burali--Forti i Cantor. (Model: observació~\ref{obs:burali-cantor}.)

  \item \textbf{(Punt de reflexió.)} Argumenteu per què la teoria \emph{global} dels tipus ordinals és de Nivell~II mentre que la teoria acotada dins un \(\Hart(A)\) és de Nivell~I. \begin{indicacio}Reunir tots els tipus demanaria SR, i ni tan sols amb SR seria un conjunt ben ordenat (observacions~\ref{obs:tipus-acotats} i~\ref{obs:burali-cantor}).\end{indicacio}

  \item Demostreu que \(\mathbf1+\N\cong\N\) com a bons ordres, mentre que \(\N+\mathbf1\not\cong\N\). Deduïu-ne que la suma de tipus no és commutativa. \begin{indicacio}El successor dona el primer isomorfisme; per al segon, \(\N+\mathbf1\) té màxim i \(\N\) no.\end{indicacio}

  \item \textbf{Axiomes usats.} Confirmeu que cap resultat del capítol no usa l'elecció, i digueu quins es poden situar dins un \(\Hart(A)\) (Nivell~I) i quins parlen de tots els tipus alhora. \begin{indicacio}Tots els resultats són de Nivell~I i sense elecció; Burali--Forti és un enunciat negatiu sobre tots els tipus, demostrat sense sortir del Nivell~I.\end{indicacio}
\end{enumerate}

\chapter{Aritmètica ordinal}\label{cap:exercicis-aritmetica-ordinal}

Aquests exercicis proposats corresponen al capítol homònim de la part teòrica (capítol~\ref{cap:aritmetica-ordinal} de la Part~\ref{part:teoria}) i es resolen amb les tècniques dels exercicis resolts de la part pràctica (capítol~\ref{cap:practica-aritmetica-ordinal} de la Part~\ref{part:practica}). Cada secció remet al bloc de resolts que en dona el model. Els marcats amb una indicació són els més exigents.

\section{Suma ordinal}

\begin{enumerate}
  \item Demostreu que la suma ordinal és un bon ordre invariant. (Model: secció~\ref{prac:suma-ordinal}; proposició~\ref{prop:suma-bon-ordre}.)

  \item Demostreu que \(1+\omega=\omega\) però \(\omega+1\neq\omega\). (Model: proposició~\ref{prop:suma-no-commutativa}.)

  \item Demostreu que la suma ordinal és associativa. (Model: proposició~\ref{prop:suma-associativa}.)

  \item Demostreu que \(\alpha+\beta=\beta+\alpha\) si i només si \(\alpha\) i \(\beta\) són \enquote{commensurables} en el sentit adequat; en particular, doneu dos tipus finits que commuten i expliqueu per què amb \(\omega\) es trenca. \begin{indicacio}Els tipus finits commuten (aritmètica de \(\N\)); amb un límit, la reabsorció per l'esquerra trenca la simetria (proposició~\ref{prop:suma-no-commutativa}).\end{indicacio}
\end{enumerate}

\section{Producte ordinal}

\begin{enumerate}[resume]
  \item Demostreu que el producte ordinal (antilexicogràfic, la coordenada de \(\beta\) primer) és un bon ordre. (Model: secció~\ref{prac:producte-ordinal}; proposició~\ref{prop:producte-bon-ordre}.)

  \item Demostreu que \(2\cdot\omega=\omega\) però \(\omega\cdot2\neq\omega\). (Model: proposició~\ref{prop:producte-propietats}.)

  \item Demostreu la distributivitat per l'esquerra i doneu un contraexemple de la distributivitat per la dreta. (Model: proposició~\ref{prop:producte-propietats}.)

  \item Demostreu que el producte ordinal és associatiu. \begin{indicacio}Tots dos costats són \(A\times B\times C\) amb la coordenada de \(C\) primer, després \(B\), després \(A\); useu l'isomorfisme canònic i la unicitat (teorema~\ref{teo:unicitat-iso}).\end{indicacio}
\end{enumerate}

\section{Exponenciació ordinal}

\begin{enumerate}[resume]
  \item Demostreu que \(\alpha^1=\alpha\), \(\alpha^2=\alpha\cdot\alpha\), \(2^\omega=\omega\) i \(\omega^2=\omega\cdot\omega\). (Model: secció~\ref{prac:exponenciacio-ordinal}; proposició~\ref{prop:exponenciacio-calculs}.)

  \item Demostreu que \(E(B,A)\cong E(C,A)\times A\) quan \(B=C+\mathbf1\), i que l'ordre antilexicogràfic correspon al del producte ordinal; deduïu \(\alpha^{\beta+1}=\alpha^\beta\cdot\alpha\). Comproveu també que la clàusula límit, aplicada a la base \(0\), donaria \(0^\omega=1\). (Model: teorema~\ref{teo:exponenciacio-recursio}, observació~\ref{obs:exponenciacio-acotacio}.)

  \item Demostreu que \(2^\omega=\omega\neq2^{\aleph_0}\), distingint l'exponenciació ordinal de la cardinal. \begin{indicacio}L'ordinal \(2^\omega\) és el suprem dels \(2^n\) finits; la cardinal \(2^{\aleph_0}\) és no numerable per Cantor (teorema~\ref{teo:cantor}).\end{indicacio}
\end{enumerate}

\section{Monotonia i ordre contra mida}

\begin{enumerate}[resume]
  \item Demostreu la monotonia estricta per la dreta i només \(\le\) per l'esquerra. (Model: secció~\ref{prac:ordre-mida}; proposició~\ref{prop:monotonia-ordinal}.)

  \item Demostreu que \(\omega\), \(\omega+1\), \(\omega\cdot2\), \(\omega^2\) són tipus diferents però tots equipotents, i que \(2^\omega=\omega\) (de manera que \(2^\omega\) no és un tipus nou, sinó el mateix \(\omega\)). (Model: observació~\ref{obs:ordre-vs-mida}.)

  \item \textbf{(Punt de reflexió.)} Expliqueu per què l'aritmètica ordinal \enquote{recorda l'ordre} i la cardinal l'\enquote{oblidarà}, i què significa exactament oblidar l'ordre. \begin{indicacio}Oblidar l'ordre és identificar tots els bons ordres equipotents; l'aritmètica cardinal opera sobre classes d'equipotència (observació~\ref{obs:ordre-vs-mida}).\end{indicacio}

  \item Demostreu que una aplicació estrictament creixent entre bons ordres \(W\to V\) implica \(\operatorname{tipus}(W)\le\operatorname{tipus}(V)\), i useu-ho per provar \(\beta\cdot\alpha\le\gamma\cdot\alpha\) quan \(\beta\le\gamma\). (Model: lema~\ref{lema:immersio-ordre}.)

  \item Demostreu que \(\omega+\omega\not\cong\omega\) trobant una propietat d'ordre invariant que distingeixi els dos bons ordres. \begin{indicacio}Un punt que no és el mínim i no té predecessor immediat.\end{indicacio}

  \item \textbf{Axiomes usats.} Confirmeu que cap resultat del capítol no usa l'elecció, i digueu si algun requereix una cota externa per mantenir-se al Nivell~I. \begin{indicacio}Cap: la construcció per funcions de suport finit proporciona les cotes; SR només caldria per a famílies de bases o exponents sense ambient comú.\end{indicacio}
\end{enumerate}

\chapter{Cardinals i aritmètica cardinal}\label{cap:exercicis-cardinals}

Aquests exercicis proposats corresponen al capítol homònim de la part teòrica (capítol~\ref{cap:cardinals} de la Part~\ref{part:teoria}) i es resolen amb les tècniques dels exercicis resolts de la part pràctica (capítol~\ref{cap:practica-cardinals} de la Part~\ref{part:practica}). Cada secció remet al bloc de resolts que en dona el model. Els marcats amb una indicació són els més exigents.

\section{Aritmètica cardinal relacional}

\begin{enumerate}
  \item Demostreu que la suma, el producte i l'exponenciació cardinals estan ben definits sobre classes d'equipotència. (Model: secció~\ref{prac:aritmetica-cardinal}; proposició~\ref{prop:operacions-invariants}.)

  \item Demostreu les lleis d'exponents \(\kappa^{\mu+\nu}=\kappa^\mu\cdot\kappa^\nu\) i \((\kappa^\mu)^\nu=\kappa^{\mu\cdot\nu}\). (Model: proposició~\ref{prop:lleis-cardinals}.)

  \item Demostreu que \(\abs{A}<2^{\abs{A}}\) per a tot \(A\), i contrasteu-ho amb \(2^\omega=\omega\) ordinal. (Model: teorema~\ref{teo:cantor-cardinal}; observació~\ref{obs:cardinal-vs-ordinal}.)

  \item Demostreu que l'aritmètica cardinal és commutativa (suma i producte), a diferència de l'ordinal. \begin{indicacio}\(A+B\cong B+A\) i \(A\times B\cong B\times A\); oblidar l'ordre elimina l'asimetria (observació~\ref{obs:cardinal-vs-ordinal}).\end{indicacio}
\end{enumerate}

\section{Comparabilitat}

\begin{enumerate}[resume]
  \item Amb l'elecció, demostreu que dos conjunts qualssevol són comparables i que els cardinals estan totalment ordenats mòdul equipotència. (Model: secció~\ref{prac:comparabilitat-cardinals}; teorema~\ref{teo:comparabilitat-cardinals}.)

  \item Demostreu que la comparabilitat total implica el teorema del bon ordre. (Model: observació~\ref{obs:comparabilitat-eleccio}.)

  \item Demostreu que, sense elecció, \(\preceq\) és igualment reflexiva, transitiva i antisimètrica mòdul \(\approx\), però potser no total. \begin{indicacio}La reflexivitat i transitivitat són la proposició~\ref{prop:preceq-preordre}; l'antisimetria és Schröder--Bernstein (corol·lari~\ref{cor:preceq-antisimetric}); la totalitat és el que costa elecció.\end{indicacio}
\end{enumerate}

\section{Cardinals inicials}

\begin{enumerate}[resume]
  \item Amb l'elecció, demostreu que tot conjunt té un únic cardinal inicial, trobat acotadament dins un Hartogs. (Model: secció~\ref{prac:cardinals-inicials}; proposició~\ref{prop:cardinal-inicial}.)

  \item Expliqueu per què el cardinal inicial és un tipus d'ordre, no un conjunt escollit. (Model: observació~\ref{obs:cardinal-prudencia}.)

  \item Demostreu que un bon ordre inicial no és equipotent a cap dels seus segments inicials propis. \begin{indicacio}Un segment propi té tipus estrictament menor; si fos equipotent, contradiria la inicialitat (definició~\ref{def:cardinal-inicial}).\end{indicacio}
\end{enumerate}

\section{El teorema de König}

\begin{enumerate}[resume]
  \item Amb l'elecció, demostreu el teorema de König: \(A_i\prec B_i\) implica \(\sum_i\abs{A_i}<\prod_i\abs{B_i}\). (Model: secció~\ref{prac:konig}; teorema~\ref{teo:konig}.)

  \item Deduïu Cantor \(\abs{A}<2^{\abs{A}}\) com a cas particular de König. (Model: observació~\ref{obs:konig-cantor}.)

  \item \textbf{(Punt de reflexió.)} Expliqueu en quin sentit el teorema de König és la forma \enquote{indexada} de la diagonal de Cantor. \begin{indicacio}La família \((d_i)\) difereix de cada valor de \(F\) en una coordenada; amb \(A_i=\mathbf1\), \(B_i=\mathbf2\), és la part diagonal (observació~\ref{obs:konig-cantor}).\end{indicacio}

  \item Doneu un contraexemple a la monotonia en l'exponent quan la base és \(0\), i expliqueu on falla la prova. \begin{indicacio}\(0^0=1\) i \(0^1=0\); per estendre una aplicació \(A\to C\) a \(B\to C\) cal un punt de \(C\).\end{indicacio}

  \item A la prova de König, identifiqueu els tres llocs on s'usa l'elecció, i escriviu cadascun com una secció d'una família interna indexada per \(I\). (Model: teorema~\ref{teo:konig}.)

  \item \textbf{Axiomes usats.} Digueu quins resultats del capítol usen l'elecció: les lleis de l'aritmètica; Cantor cardinal; la comparabilitat total; l'existència del cardinal inicial; König. \begin{indicacio}Els tres últims.\end{indicacio}
\end{enumerate}

\chapter{Alephs, cofinalitat i continu}\label{cap:exercicis-alephs}

Aquests exercicis proposats corresponen al capítol homònim de la part teòrica (capítol~\ref{cap:alephs} de la Part~\ref{part:teoria}), opcional i de Nivell~II, i es resolen amb les tècniques dels exercicis resolts de la part pràctica (capítol~\ref{cap:practica-alephs} de la Part~\ref{part:practica}). Cada secció remet al bloc de resolts que en dona el model. Els marcats amb una indicació són els més exigents. Recordeu de marcar, a cada resultat, si depèn de l'elecció, de l'esquema de reunió (principi~\ref{sr:esquema}), o de cap dels dos.

\section{\texorpdfstring{\(\aleph_0\)}{aleph 0} i el cardinal successor}

\begin{enumerate}
  \item Demostreu que \(\kappa^+\) existeix per a tot \(\kappa\), acotat per Hartogs, sense l'esquema de reunió. (Model: secció~\ref{prac:aleph0-successor}; proposició~\ref{prop:cardinal-successor}.)

  \item Expliqueu per què cada \(\kappa^+\) és de Nivell~I però la successió \((\aleph_n)_n\) demana l'esquema de reunió. (Model: observació~\ref{obs:successor-sense-sr}.)

  \item Demostreu que \(\aleph_0\) és el menor cardinal infinit. \begin{indicacio}Tot conjunt infinit conté una còpia de \(\N\) (capítol~\ref{cap:infinitud}); tot enumerable infinit és equipotent a \(\N\) (corol·lari~\ref{cor:enumerable-infinit-N}).\end{indicacio}
\end{enumerate}

\section{La successió dels alephs}

\begin{enumerate}[resume]
  \item Sota l'esquema de reunió, demostreu que els alephs són estrictament creixents. (Model: secció~\ref{prac:alephs}; proposició~\ref{prop:alephs}.)

  \item Sota elecció i l'esquema de reunió, demostreu que tot cardinal infinit és un aleph. (Model: proposició~\ref{prop:alephs}.)

  \item Identifiqueu explícitament la instància de l'esquema de reunió que s'usa per definir, fixat un ordinal \(\lambda\), la família \((\aleph_\beta)_{\beta<\lambda}\) indexada pel conjunt \(\lambda\). \begin{indicacio}La fórmula \enquote{\(\beta\mapsto\) el cardinal inicial \(\aleph_\beta\)} és funcional i invariant per isomorfisme; la notació global \(\alpha\mapsto\aleph_\alpha\) és només una regla de classe. Vegeu la definició~\ref{def:alephs}.\end{indicacio}
\end{enumerate}

\section{Cofinalitat}

\begin{enumerate}[resume]
  \item Demostreu que \(\aleph_\omega\) és singular, amb \(\operatorname{cf}(\aleph_\omega)=\aleph_0\). (Model: secció~\ref{prac:cofinalitat}; proposició~\ref{prop:regular-exemples}.)

  \item Demostreu que tot cardinal successor és regular. (Model: proposició~\ref{prop:regular-exemples}.)

  \item Deduïu del teorema de König que \(\operatorname{cf}(2^{\aleph_0})>\aleph_0\): el continu no té cofinalitat numerable. \begin{indicacio}Si \(2^{\aleph_0}\) fos suprem d'una successió numerable de cardinals menors, König donaria una contradicció (observació~\ref{obs:cofinalitat-importa}).\end{indicacio}
\end{enumerate}

\section{Beths i el continu}

\begin{enumerate}[resume]
  \item Sota l'esquema de reunió, demostreu que \(\aleph_\alpha\le\beth_\alpha\) per a tot \(\alpha\). (Model: secció~\ref{prac:beths-continu}; proposició~\ref{prop:beths}.)

  \item Enuncieu la hipòtesi del continu i la generalitzada, i expliqueu per què són qüestions independents. (Model: observacions~\ref{obs:hipotesi-continu} i~\ref{obs:hc-independencia}.)

  \item \textbf{(Punt de reflexió.)} Expliqueu per què la teoria \enquote{formula HC amb precisió però no la resol}, i què significa exactament que HC sigui independent. \begin{indicacio}La teoria distingeix \(\aleph_1\) de \(\beth_1\); Cantor només dona \(\aleph_0<2^{\aleph_0}\); Gödel i Cohen mostren que cap dels dos costats no es dedueix dels axiomes (observació~\ref{obs:hc-independencia}).\end{indicacio}

  \item Descriviu la instància de l'esquema SR que s'usa als límits de la construcció dels alephs: quin és el conjunt d'índexs, quina la fórmula, quina l'estructura, i per què l'isomorfisme entre dues realitzacions és únic. (Model: prova de la proposició~\ref{prop:alephs}.)

  \item Deduïu de König que \(\operatorname{cf}(2^{\aleph_0})>\aleph_0\). \begin{indicacio}Si \(2^{\aleph_0}=\sum_n\abs{A_n}\) amb \(\abs{A_n}<2^{\aleph_0}\), compareu amb \(\prod_n2^{\aleph_0}=2^{\aleph_0\cdot\aleph_0}\).\end{indicacio}

  \item \textbf{Axiomes usats.} Per a cadascun dels resultats següents, digueu si és de Nivell~I, si usa l'elecció, SR, o un resultat citat: \(\kappa^+=\abs{\Hart(K)}\); la família \((\aleph_\beta)_{\beta<\lambda}\); tot cardinal infinit és un aleph; la regularitat dels successors; la formulació d'HC; la independència d'HC. \begin{indicacio}Nivell~I amb elecció; SR; SR i elecció; elecció i Hessenberg (citat); Nivell~I; metateorema citat, sota la hipòtesi de consistència de ZFC.\end{indicacio}
\end{enumerate}


% ============================================================
% PART IV — TESTS
% En pantalla, un sol entorn Form per a tots els tests: un únic diccionari
% AcroForm. En mode imprimible les macros dels tests són estàtiques.
% ============================================================
\part{\NomTests}\label{part:tests}
\setcounter{chapter}{0}
\ifimprimible\else\begin{Form}\fi
\chapter{Aplicacions, dominis, codominis i composició}\label{test:cap1}

Aquest test correspon al capítol~1 de la part teòrica. Cada pregunta té una única resposta correcta.

\iniciTest{u1}{b,c,b,a,c,c,b,a,c,b,c,d}

\begin{enumerate}

\pregunta Quin és el vocabulari no lògic primitiu de la teoria?
\begin{itemize}
\opcio{a}{La pertinença, la inclusió i la igualtat.}
\opcio{b}{Els predicats \(\Dom\), \(\Cod\) i \(\Com\).}
\opcio{c}{Els conjunts, els elements i les funcions.}
\opcio{d}{La suma, el producte i la composició.}
\end{itemize}
\feedback
\justificacio{Els únics símbols no lògics del llenguatge són els dos predicats binaris \(\Dom\) i \(\Cod\) i el ternari \(\Com\). Conjunts, elements, pertinença i inclusió són nocions definides o abreviacions.}

\pregunta La composició \(g\circ f\) està definida quan\ldots
\begin{itemize}
\opcio{a}{el domini de \(f\) coincideix amb el domini de \(g\).}
\opcio{b}{el codomini de \(g\) coincideix amb el domini de \(f\).}
\opcio{c}{el codomini de \(f\) coincideix amb el domini de \(g\).}
\opcio{d}{sempre: dues aplicacions qualssevol es poden compondre.}
\end{itemize}
\feedback
\justificacio{L'axioma d'existència de la composició demana \(\cod(f)=\dom(g)\): primer s'aplica \(f\) i després \(g\). Les altres opcions confonen l'ordre o ignoren que la composició és parcial.}

\pregunta Segons la nostra teoria, un element global d'un conjunt \(A\) és\ldots
\begin{itemize}
\opcio{a}{un individu primitiu que pertany a \(A\).}
\opcio{b}{una aplicació \(\mathbf1\to A\).}
\opcio{c}{una aplicació \(A\to\mathbf1\).}
\opcio{d}{una aplicació \(A\to\mathbf2\).}
\end{itemize}
\feedback
\justificacio{No hi ha elements primitius: un punt d'\(A\) és una aplicació des del terminal, \(x\colon\mathbf1\to A\). Les aplicacions \(A\to\mathbf1\) i \(A\to\mathbf2\) són, respectivament, l'única \({!}_A\) i les futures parts d'\(A\).}

\pregunta Quantes aplicacions hi ha d'un conjunt qualsevol \(A\) cap al conjunt terminal \(\mathbf1\)?
\begin{itemize}
\opcio{a}{Exactament una.}
\opcio{b}{Tantes com elements globals té \(A\).}
\opcio{c}{Cap, si \(A\) té més d'un element global.}
\opcio{d}{Depèn del conjunt \(A\).}
\end{itemize}
\feedback
\justificacio{És la definició de terminal: des de qualsevol conjunt, inclòs \(\mathbf0\), hi ha exactament una aplicació cap a \(\mathbf1\), sigui quin sigui el nombre d'elements globals del domini.}

\pregunta Si \(f:A\to B\) és un isomorfisme, quantes inverses té?
\begin{itemize}
\opcio{a}{Almenys una, però en pot tenir més.}
\opcio{b}{Dues: una per l'esquerra i una per la dreta.}
\opcio{c}{Exactament una, per l'associativitat i les lleis de neutralitat.}
\opcio{d}{Cap: la inversa només existeix per a la identitat.}
\end{itemize}
\feedback
\justificacio{Si \(g\) i \(g'\) són inverses, \(g=g\circ(f\circ g')=(g\circ f)\circ g'=g'\): neutralitat i associativitat forcen la unicitat. Una inversa ho és alhora per l'esquerra i per la dreta.}

\pregunta Segons la proposició d'unicitat de la composició, si \(\Com(f,g,h)\) i \(\Com(f,g,h')\), aleshores\ldots
\begin{itemize}
\opcio{a}{\(h\) i \(h'\) tenen dominis diferents.}
\opcio{b}{no se'n pot dir res sense més hipòtesis.}
\opcio{c}{\(h=h'\).}
\opcio{d}{\(h\circ h'=\id\).}
\end{itemize}
\feedback
\justificacio{És exactament la proposició d'unicitat de la composició, que es demostra amb l'associativitat i les lleis de neutralitat; l'axioma només n'assegura l'existència.}

\pregunta Quan diem que \(A\) és un \emph{conjunt} en aquesta teoria?
\begin{itemize}
\opcio{a}{Quan té almenys un element global.}
\opcio{b}{Quan \(A\) és una identitat, és a dir \(\Dom(A,A)\) i \(\Cod(A,A)\).}
\opcio{c}{Quan hi ha una aplicació \(A\to\mathbf2\).}
\opcio{d}{Quan \(A\) és terminal.}
\end{itemize}
\feedback
\justificacio{Un conjunt és una identitat: una aplicació que és el seu propi domini i codomini. Tenir punts, tenir parts o ser terminal són propietats que alguns conjunts tenen i d'altres no.}

\pregunta El criteri d'igualtat d'aplicacions per elements generalitzats diu que \(f=g\) quan\ldots
\begin{itemize}
\opcio{a}{\(f\circ x=g\circ x\) per a tot element generalitzat \(x\colon X\to A\), amb \(X\) arbitrari.}
\opcio{b}{\(f\) i \(g\) tenen el mateix domini.}
\opcio{c}{\(f\circ x=g\circ x\) només per a \(x\colon\mathbf1\to A\).}
\opcio{d}{\(f\) i \(g\) són totes dues identitats.}
\end{itemize}
\feedback
\justificacio{El criteri quantifica sobre \emph{tots} els estadis \(X\); n'hi ha prou amb \(X=A\) i \(x=\id_A\). Restringir-se a \(X=\mathbf1\) és el contingut de l'axioma de separació global, no del criteri.}

\pregunta Les lleis de neutralitat afirmen que, quan les composicions estan definides,\ldots
\begin{itemize}
\opcio{a}{\(f\circ g=g\circ f\).}
\opcio{b}{tota aplicació és un isomorfisme.}
\opcio{c}{\(f\circ\id=f\) i \(\id\circ f=f\).}
\opcio{d}{\(f\circ f=f\).}
\end{itemize}
\feedback
\justificacio{Les lleis de neutralitat diuen que el domini i el codomini actuen com a neutres: \(f\circ\id_{\dom f}=f\) i \(\id_{\cod f}\circ f=f\). La composició no és commutativa ni idempotent en general.}

\pregunta Quantes aplicacions hi ha del conjunt terminal \(\mathbf1\) cap al conjunt inicial \(\mathbf0\)?
\begin{itemize}
\opcio{a}{Exactament una.}
\opcio{b}{Cap, per la no-degeneració.}
\opcio{c}{Tantes com elements té \(\mathbf0\).}
\opcio{d}{Infinites.}
\end{itemize}
\feedback
\justificacio{Una aplicació \(\mathbf1\to\mathbf0\) donaria, amb l'única \(\mathbf0\to\mathbf1\), un isomorfisme \(\mathbf0\cong\mathbf1\), prohibit per la no-degeneració. Per això \(\mathbf0\) no té cap element global.}

\pregunta Quin és l'enunciat dual de \enquote{des de tot conjunt arriba exactament una aplicació a \(T\)}?
\begin{itemize}
\opcio{a}{\enquote{Cap aplicació no arriba a \(T\)}.}
\opcio{b}{\enquote{Tota aplicació que surt de \(T\) és un isomorfisme}.}
\opcio{c}{\enquote{De \(T\) surt exactament una aplicació cap a tot conjunt}.}
\opcio{d}{\enquote{\(T\) té exactament un element global}.}
\end{itemize}
\feedback
\justificacio{La dualitat intercanvia domini i codomini: \enquote{arribar a \(T\)} es converteix en \enquote{sortir de \(T\)}, i la condició d'existència i unicitat es conserva. És la definició de conjunt inicial.}

\pregunta Quin axioma garanteix que una aplicació queda determinada pels valors que pren sobre els punts globals del seu domini?
\begin{itemize}
\opcio{a}{La llei associativa.}
\opcio{b}{L'axioma d'existència d'un conjunt terminal.}
\opcio{c}{L'axioma de no-degeneració.}
\opcio{d}{L'axioma de separació global: \(\mathbf1\) és un separador.}
\end{itemize}
\feedback
\justificacio{Que els elements \emph{generalitzats} determinin les aplicacions és un teorema; que n'hi hagi prou amb els \emph{globals} és l'axioma de separació global, que fa la teoria ben apuntada. L'existència del terminal només diu que \(\mathbf1\) existeix.}

\end{enumerate}
\clauestatica
\finalitzaTest

\chapter{Parts i construccions universals}\label{test:cap2}

Aquest test correspon al capítol~2 de la part teòrica. Cada pregunta té una única resposta correcta.

\iniciTest{u2}{a,c,b,c,b,d,b,a,c,b,c,d}

\begin{enumerate}

\pregunta Una aplicació \(m\colon A\to B\) és un \emph{monomorfisme} quan\ldots
\begin{itemize}
\opcio{a}{es pot cancel·lar per l'esquerra: \(m\circ u=m\circ v\Rightarrow u=v\).}
\opcio{b}{es pot cancel·lar per la dreta: \(u\circ m=v\circ m\Rightarrow u=v\).}
\opcio{c}{admet sempre una inversa per la dreta.}
\opcio{d}{el seu domini és el conjunt terminal.}
\end{itemize}
\feedback
\justificacio{Un monomorfisme es cancel·la per l'esquerra: és la formulació estructural de la injectivitat. La cancel·lació per la dreta defineix els epimorfismes, i tenir inversa per la dreta és ser una retracció, condició més forta.}

\pregunta La composició de dos monomorfismes\ldots
\begin{itemize}
\opcio{a}{no té per què ser un monomorfisme.}
\opcio{b}{és un epimorfisme.}
\opcio{c}{és sempre un monomorfisme.}
\opcio{d}{és un isomorfisme.}
\end{itemize}
\feedback
\justificacio{Si \((n\circ m)\circ u=(n\circ m)\circ v\), l'associativitat permet cancel·lar primer \(n\) i després \(m\). Res no obliga que la composició sigui epi ni iso.}

\pregunta Tota secció \(s\) (amb \(r\circ s=\id\)) és\ldots
\begin{itemize}
\opcio{a}{un epimorfisme.}
\opcio{b}{un monomorfisme.}
\opcio{c}{una aplicació constant.}
\opcio{d}{el conjunt terminal.}
\end{itemize}
\feedback
\justificacio{Si \(r\circ s=\id\) i \(s\circ u=s\circ v\), component amb \(r\) s'obté \(u=v\). Una secció és un monomorfisme escindit; la retracció \(r\) és l'epimorfisme.}

\pregunta Dues parts \(m,n\) de \(A\) compleixen \(m\subseteq_A n\) quan\ldots
\begin{itemize}
\opcio{a}{tenen el mateix domini.}
\opcio{b}{existeix una aplicació \(A\to A\) que les intercanvia.}
\opcio{c}{existeix \(k\) amb \(m=n\circ k\).}
\opcio{d}{tenen la mateixa característica.}
\end{itemize}
\feedback
\justificacio{La inclusió de parts és una factorització: \(m\subseteq_A n\) vol dir \(m=n\circ k\) per a alguna \(k\), que resulta única i mono. Tenir la mateixa característica és l'equivalència, no la inclusió.}

\pregunta La propietat universal del producte \(A\times B\) diu que una aplicació \(X\to A\times B\)\ldots
\begin{itemize}
\opcio{a}{existeix només si \(X\) és terminal.}
\opcio{b}{queda determinada de manera única per les seves components sobre les projeccions.}
\opcio{c}{és sempre un monomorfisme.}
\opcio{d}{coincideix amb una aplicació \(X\to A+B\).}
\end{itemize}
\feedback
\justificacio{Donades \(f\colon X\to A\) i \(g\colon X\to B\), hi ha una única \(\langle f,g\rangle\colon X\to A\times B\) amb aquestes components; per tant, dues aplicacions cap al producte amb les mateixes components són iguals.}

\pregunta L'aplicació característica \(\chi_m\colon A\to\mathbf2\) d'una part \(m\colon S\rightarrowtail A\)\ldots
\begin{itemize}
\opcio{a}{depèn de la parametrització concreta \(S\).}
\opcio{b}{només existeix quan \(S\) és terminal.}
\opcio{c}{és sempre constant.}
\opcio{d}{determina la part llevat d'equivalència.}
\end{itemize}
\feedback
\justificacio{L'axioma de representació dona una única característica per part, i dues parts amb la mateixa característica tenen els mateixos elements generalitzats, de manera que són equivalents. No depèn del domini \(S\) ni és constant en general.}

\pregunta L'equalitzador \(e\colon E\to A\) de dues aplicacions \(f,g\colon A\to B\) sempre és\ldots
\begin{itemize}
\opcio{a}{un epimorfisme.}
\opcio{b}{un monomorfisme.}
\opcio{c}{un isomorfisme.}
\opcio{d}{una aplicació constant.}
\end{itemize}
\feedback
\justificacio{Si \(e\circ u=e\circ v\), tots dos igualen \(f\) i \(g\) i, per la unicitat de la factorització, \(u=v\). L'equalitzador és una part del domini; no té per què ser epi ni iso.}

\pregunta En un pullback de \(f\colon A\to C\) al llarg de \(g\colon B\to C\), si \(f\) és un monomorfisme, la projecció oposada\ldots
\begin{itemize}
\opcio{a}{és un monomorfisme, i per això la imatge inversa d'una part és una part.}
\opcio{b}{és un epimorfisme.}
\opcio{c}{és un isomorfisme.}
\opcio{d}{no té cap propietat especial.}
\end{itemize}
\feedback
\justificacio{Si \(f\) és mono i \(q_B\circ u=q_B\circ v\), del quadrat se segueix \(f\circ q_A\circ u=f\circ q_A\circ v\), es cancel·la \(f\) i la unicitat del pullback dona \(u=v\). Per això \(f^{-1}(m)\) torna a ser una part.}

\pregunta En aquesta teoria, un bimorfisme (mono i epi alhora)\ldots
\begin{itemize}
\opcio{a}{no té per què ser un isomorfisme.}
\opcio{b}{només existeix entre conjunts terminals.}
\opcio{c}{és sempre un isomorfisme.}
\opcio{d}{és sempre una aplicació constant.}
\end{itemize}
\feedback
\justificacio{És un teorema del capítol, no una definició: un epi que és part té característica total, la representació en dona una secció, i la cancel·lació per mono la converteix en inversa.}

\pregunta Siguin \(g\colon B\to A\) i una part \(m\) de \(A\). Es té \(y\in_X g^{-1}(m)\) si i només si\ldots
\begin{itemize}
\opcio{a}{\(y\) és un element global.}
\opcio{b}{\(g\circ y\in_X m\).}
\opcio{c}{\(g\) és un isomorfisme.}
\opcio{d}{\(m\) és la part total.}
\end{itemize}
\feedback
\justificacio{És la propietat que caracteritza la imatge inversa com a pullback: \(y\) es factoritza per \(g^{-1}(m)\) exactament quan \(g\circ y\) es factoritza per \(m\).}

\pregunta Quin axioma del capítol fa falta per demostrar que tot bimorfisme és un isomorfisme?
\begin{itemize}
\opcio{a}{Cap: se segueix de les definicions de mono i epi.}
\opcio{b}{L'axioma d'existència de productes.}
\opcio{c}{L'axioma de representació de les parts.}
\opcio{d}{L'axioma de separació global.}
\end{itemize}
\feedback
\justificacio{La prova usa la característica del bimorfisme vist com a part: és l'axioma de representació. En una teoria sense classificador de parts, un bimorfisme pot no ser isomorfisme, com passa amb la inclusió dels enters en els racionals vista com a homomorfisme d'anells.}

\pregunta Quina d'aquestes construccions \emph{no} demana cap axioma nou al capítol, perquè es construeix a partir de les altres?
\begin{itemize}
\opcio{a}{La suma \(A+B\).}
\opcio{b}{El producte \(A\times B\).}
\opcio{c}{La característica \(\chi_m\) d'una part.}
\opcio{d}{El pullback \(A\times_C B\).}
\end{itemize}
\feedback
\justificacio{El pullback es construeix com un equalitzador dins del producte, i l'equalitzador com la part classificada per \(\operatorname{eq}_B\circ\langle f,g\rangle\). Sumes, productes i representació són els tres axiomes del capítol.}

\end{enumerate}
\clauestatica
\finalitzaTest

\chapter[Valors de veritat i construccions finites]{Valors de veritat, parts\texorpdfstring{\\}{ }i construccions finites}\label{test:cap3}

Aquest test correspon al capítol~3 de la part teòrica. Cada pregunta té una única resposta correcta.

\iniciTest{u3}{c,b,a,c,c,a,b,c,b,c,c,b}

\begin{enumerate}

\pregunta Que una suma \(A_0\to A_0+A_1\leftarrow A_1\) sigui \emph{disjunta} vol dir, entre altres coses, que\ldots
\begin{itemize}
\opcio{a}{els dos sumands són isomorfs.}
\opcio{b}{la suma és el conjunt terminal.}
\opcio{c}{el pullback de les dues injeccions és inicial i cada injecció és un monomorfisme.}
\opcio{d}{no existeix cap aplicació que en surti.}
\end{itemize}
\feedback
\justificacio{La disjunció diu que les dues injeccions són monomorfismes i que el seu pullback és inicial: les dues còpies no se solapen. Els sumands no han de ser isomorfs, i de la suma en surten tantes aplicacions com defineixi la propietat universal.}

\pregunta El conjunt inicial \(\mathbf0\) és \emph{estricte}. Això significa que\ldots
\begin{itemize}
\opcio{a}{no té cap element global.}
\opcio{b}{tota aplicació \(X\to\mathbf0\) és un isomorfisme.}
\opcio{c}{és isomorf al terminal.}
\opcio{d}{admet exactament dos elements.}
\end{itemize}
\feedback
\justificacio{Estricte vol dir que tota aplicació \(X\to\mathbf0\) és un isomorfisme: si alguna cosa es pot enviar a \(\mathbf0\), ja era \(\mathbf0\). Es demostra amb l'estabilitat i la disjunció de la suma \(\mathbf0+\mathbf0\). Que \(\mathbf0\) no tingui punts és una altra propietat, provada al capítol 1.}

\pregunta Els dos elements globals \(0,1\colon\mathbf1\to\mathbf2\)\ldots
\begin{itemize}
\opcio{a}{són diferents, i \(\mathbf2\) té exactament dos elements globals.}
\opcio{b}{coincideixen.}
\opcio{c}{no són monomorfismes.}
\opcio{d}{només existeixen si \(\mathbf2\) és inicial.}
\end{itemize}
\feedback
\justificacio{Com que la suma \(\mathbf1+\mathbf1\) és disjunta, si \(0=1\) el pullback de \(0\) amb ella mateixa (que és \(\mathbf1\)) seria inicial, contra la no-degeneració. I tot punt de \(\mathbf2\) ve d'un sumand: és \(0\) o \(1\).}

\pregunta Un conjunt \(X\) sense cap element global\ldots
\begin{itemize}
\opcio{a}{és necessàriament el terminal.}
\opcio{b}{té exactament un element global amagat.}
\opcio{c}{és inicial (lema del buit).}
\opcio{d}{és isomorf a \(\mathbf2\).}
\end{itemize}
\feedback
\justificacio{Si \(X\) no té punts, \(0\circ{!}_X\) i \(1\circ{!}_X\) coincideixen buidament sobre tots ells i, per separació global, són iguals; això dona una aplicació cap al pullback inicial de \(0\) i \(1\), i l'estrictesa fa \(X\cong\mathbf0\).}

\pregunta La característica de la part total \(\id_A\colon A\rightarrowtail A\) és\ldots
\begin{itemize}
\opcio{a}{la negació \(\neg\).}
\opcio{b}{una aplicació no constant.}
\opcio{c}{la constant de valor \(1\).}
\opcio{d}{la constant de valor \(0\).}
\end{itemize}
\feedback
\justificacio{Tot element generalitzat pertany a la part total, i \((1\circ{!}_A)\circ x=1\circ{!}_X\) per a tot \(x\); la unicitat de la característica fa la resta. La constant \(0\) és la característica de la part buida.}

\pregunta La negació \(\neg\colon\mathbf2\to\mathbf2\) compleix\ldots
\begin{itemize}
\opcio{a}{\(\neg\circ\neg=\id_{\mathbf2}\); és una involució.}
\opcio{b}{\(\neg\circ\neg=\neg\).}
\opcio{c}{\(\neg\) és una aplicació constant.}
\opcio{d}{\(\neg\) no és un isomorfisme.}
\end{itemize}
\feedback
\justificacio{\(\neg=[1,0]\) intercanvia els dos punts; \(\neg\circ\neg\) i \(\id_{\mathbf2}\) coincideixen sobre les dues injeccions i, per extensionalitat de la definició per casos, són iguals. Per tant \(\neg\) és un isomorfisme invers de si mateix.}

\pregunta Per a tota part \(m\colon S\rightarrowtail A\), l'aplicació \([m,m^c]\colon S+S^c\to A\)\ldots
\begin{itemize}
\opcio{a}{és sempre un monomorfisme però no exhaustiu.}
\opcio{b}{és un isomorfisme: \(m\) i el seu complement descomponen \(A\).}
\opcio{c}{és l'aplicació constant.}
\opcio{d}{només existeix quan \(S\) és buit.}
\end{itemize}
\feedback
\justificacio{\(m\) és el pullback de \(1\) al llarg de \(\chi_m\) i \(m^c\) el de \(0\); com que \(\mathbf1\xrightarrow{1}\mathbf2\xleftarrow{0}\mathbf1\) és una suma, l'estabilitat fa que els seus dos pullbacks formin una suma amb vèrtex \(A\), i la unicitat de les sumes dona l'isomorfisme.}

\pregunta La característica de la intersecció \(m\cap n\) és\ldots
\begin{itemize}
\opcio{a}{\(\neg\circ\chi_m\).}
\opcio{b}{\(\chi_m\circ\chi_n\).}
\opcio{c}{\(\wedge\circ\langle\chi_m,\chi_n\rangle\), la conjunció de les dues característiques.}
\opcio{d}{sempre la constant \(1\).}
\end{itemize}
\feedback
\justificacio{La conjunció és la característica de \(\delta_{11}=\langle1,1\rangle\), i \(\wedge\circ\langle\chi_m,\chi_n\rangle\) classifica els elements generalitzats que pertanyen alhora a \(m\) i a \(n\), és a dir, la intersecció. La composició \(\chi_m\circ\chi_n\) ni tan sols està definida.}

\pregunta El coequalitzador \(q\colon B\to Q\) de dues aplicacions \(f,g\colon A\to B\) sempre és\ldots
\begin{itemize}
\opcio{a}{un monomorfisme.}
\opcio{b}{un epimorfisme.}
\opcio{c}{un isomorfisme.}
\opcio{d}{una aplicació constant.}
\end{itemize}
\feedback
\justificacio{Dues aplicacions que coincideixen després de compondre amb \(q\) són factoritzacions de la mateixa aplicació a través de \(q\), i la unicitat de la propietat universal les iguala. És el dual del fet que l'equalitzador és mono.}

\pregunta El pushout de \(f\colon C\to A\) i \(g\colon C\to B\) es pot construir\ldots
\begin{itemize}
\opcio{a}{com un equalitzador de dues projeccions.}
\opcio{b}{com el producte \(A\times B\).}
\opcio{c}{com el coequalitzador de \(i_A\circ f\) i \(i_B\circ g\) sobre la suma \(A+B\).}
\opcio{d}{només quan \(C\) és inicial.}
\end{itemize}
\feedback
\justificacio{El pushout és el coequalitzador de \(i_A\circ f\) i \(i_B\circ g\) sobre \(A+B\): la suma ajunta \(A\) i \(B\), i el coequalitzador identifica les dues imatges de \(C\). És la construcció dual del pullback com a equalitzador dins d'un producte.}

\pregunta Quin dels resultats següents demana un axioma nou del capítol (disjunció i estabilitat, o coequalitzadors), i no es dedueix dels axiomes dels capítols anteriors?
\begin{itemize}
\opcio{a}{L'existència del límit de qualsevol diagrama finit.}
\opcio{b}{L'existència del pullback de dues aplicacions.}
\opcio{c}{La distinció \(0\ne1\) entre els dos punts de \(\mathbf2\).}
\opcio{d}{La unicitat de la característica d'una part.}
\end{itemize}
\feedback
\justificacio{Els límits finits, els pullbacks i la característica surten dels axiomes dels capítols 1 i 2. La distinció \(0\ne1\) necessita la disjunció de la suma \(\mathbf1+\mathbf1\): sense l'axioma nou, els dos punts podrien coincidir.}

\pregunta Un epimorfisme \(f\colon A\to B\) i un punt \(y\colon\mathbf1\to B\). Què es pot afirmar amb els axiomes d'aquest capítol?
\begin{itemize}
\opcio{a}{Res: cal l'axioma d'elecció per trobar una preimatge de \(y\).}
\opcio{b}{Existeix un punt \(x\colon\mathbf1\to A\) amb \(f\circ x=y\), sense elecció.}
\opcio{c}{Existeix una secció \(s\colon B\to A\) de \(f\).}
\opcio{d}{\(f\) és un isomorfisme.}
\end{itemize}
\feedback
\justificacio{La fibra \(f^{-1}(y)\) té un punt, perquè si no en tingués seria inicial i la característica \(\chi_y\circ f\) seria \(0\circ{!}_A\), cosa que, cancel·lant \(f\), donaria \(0=1\). Una secció de \(f\) ---una preimatge per a tots els elements generalitzats alhora--- sí que demana l'elecció.}

\end{enumerate}
\clauestatica
\finalitzaTest

\chapter{Relacions, quocients i elecció}\label{test:cap4}

Aquest test correspon al capítol~4 de la part teòrica. Cada pregunta té una única resposta correcta.

\iniciTest{u4}{b,c,b,a,c,b,c,a,c,b,b,c}

\begin{enumerate}

\pregunta En aquesta teoria, una \emph{relació} d'un conjunt \(A\) cap a un conjunt \(B\) és\ldots
\begin{itemize}
\opcio{a}{una aplicació \(A\to B\).}
\opcio{b}{una part d'un producte \(A\times B\).}
\opcio{c}{un element global de \(A\times B\).}
\opcio{d}{un isomorfisme \(A\to B\).}
\end{itemize}
\feedback
\justificacio{Una relació és una part \(r\colon R\rightarrowtail A\times B\); els parells relacionats són els elements generalitzats que hi pertanyen. Una aplicació \(A\to B\) és el cas particular d'una relació funcional, i es recupera com a part per la seva gràfica.}

\pregunta La diagonal \(\Delta_A=\langle\id_A,\id_A\rangle\) compleix que \(x\mathrel{(\Delta_A)_X}x'\) si i només si\ldots
\begin{itemize}
\opcio{a}{\(x\) és un element global.}
\opcio{b}{\(x\) i \(x'\) tenen codominis diferents.}
\opcio{c}{\(x=x'\).}
\opcio{d}{\(A\) és terminal.}
\end{itemize}
\feedback
\justificacio{\(\langle x,x'\rangle\) es factoritza per \(\langle\id_A,\id_A\rangle\) exactament quan \(x=w=x'\) per a alguna \(w\). La diagonal és la relació d'igualtat, i això val a tots els estadis.}

\pregunta El parell nucli \(\langle k_1,k_2\rangle\) d'una aplicació \(f\colon A\to B\) sempre és\ldots
\begin{itemize}
\opcio{a}{un epimorfisme.}
\opcio{b}{un monomorfisme, encara que \(f\) no ho sigui.}
\opcio{c}{un isomorfisme.}
\opcio{d}{un monomorfisme només si \(f\) ho és.}
\end{itemize}
\feedback
\justificacio{La prova només usa la unicitat de la factorització pel pullback: dues aplicacions cap a \(A\times_BA\) amb les mateixes projeccions són iguals. Cap propietat de \(f\) hi intervé.}

\pregunta El parell nucli de qualsevol aplicació \(f\colon A\to B\) és\ldots
\begin{itemize}
\opcio{a}{una relació d'equivalència sobre \(A\).}
\opcio{b}{una relació d'ordre.}
\opcio{c}{sempre la diagonal.}
\opcio{d}{una part de \(B\times B\).}
\end{itemize}
\feedback
\justificacio{Reflexivitat, simetria i transitivitat s'obtenen factoritzant \((\id_A,\id_A)\), \((k_2,k_1)\) i el parell adequat sobre \(K\times_AK\) pel pullback. És la diagonal només quan \(f\) és mono.}

\pregunta Un \emph{epimorfisme regular} és una aplicació que\ldots
\begin{itemize}
\opcio{a}{admet una inversa per la dreta.}
\opcio{b}{és alhora mono i epi.}
\opcio{c}{és el coequalitzador d'algun parell d'aplicacions paral·leles.}
\opcio{d}{té domini terminal.}
\end{itemize}
\feedback
\justificacio{És la definició: coequalitzador d'algun parell paral·lel. Tenir inversa per la dreta és ser escindit, condició més forta sense l'elecció; mono i epi alhora és ser bimorfisme, que ja és isomorfisme.}

\pregunta L'aplicació quocient \(q\colon A\twoheadrightarrow A/r\) d'una relació d'equivalència\ldots
\begin{itemize}
\opcio{a}{és un monomorfisme.}
\opcio{b}{és un epimorfisme regular, i és única llevat d'un únic isomorfisme compatible.}
\opcio{c}{no existeix sense l'axioma d'elecció.}
\opcio{d}{és sempre un isomorfisme.}
\end{itemize}
\feedback
\justificacio{El quocient és el coequalitzador de \(r_1,r_2\), que existeix per l'axioma de coequalitzadors i és únic llevat d'un únic isomorfisme compatible; l'elecció no hi intervé. És isomorfisme només si \(r\) és la diagonal.}

\pregunta L'axioma d'\emph{efectivitat} de les relacions d'equivalència afirma que\ldots
\begin{itemize}
\opcio{a}{tota relació és simètrica.}
\opcio{b}{tot epimorfisme té secció.}
\opcio{c}{en prendre el quocient i després el parell nucli es recupera la relació original.}
\opcio{d}{tota aplicació és un epimorfisme.}
\end{itemize}
\feedback
\justificacio{L'efectivitat diu que el parell nucli d'un quocient de \(r\) és equivalent a \(r\): quocient i parell nucli són inversos. Que tot epimorfisme tingui secció és l'elecció.}

\pregunta L'axioma d'\emph{elecció}, en la seva formulació estructural, diu que\ldots
\begin{itemize}
\opcio{a}{tot epimorfisme és escindit (té secció).}
\opcio{b}{tot monomorfisme és escindit.}
\opcio{c}{tota relació d'equivalència és efectiva.}
\opcio{d}{tot conjunt té un element global.}
\end{itemize}
\feedback
\justificacio{En forma estructural, l'elecció diu que tot epimorfisme té secció. Que tot mono sigui escindit és fals (\(\mathbf0\to\mathbf1\) no té retracció); l'efectivitat és un altre axioma.}

\pregunta Amb l'axioma d'elecció, tot epimorfisme\ldots
\begin{itemize}
\opcio{a}{és un monomorfisme.}
\opcio{b}{és un isomorfisme.}
\opcio{c}{és regular i exhaustiu sobre elements generalitzats.}
\opcio{d}{deixa d'existir.}
\end{itemize}
\feedback
\justificacio{Amb secció, un epimorfisme és regular (coequalitzador del seu parell nucli) i exhaustiu sobre elements generalitzats (\(x:=s\circ y\)). No és mono ni iso en general.}

\pregunta En la lectura clàssica, una secció \(s\) d'un quocient \(q\colon A\twoheadrightarrow A/r\)\ldots
\begin{itemize}
\opcio{a}{col·lapsa totes les classes en un punt.}
\opcio{b}{tria un representant de cada classe d'equivalència.}
\opcio{c}{és sempre la identitat.}
\opcio{d}{només existeix si \(r\) és la diagonal.}
\end{itemize}
\feedback
\justificacio{Com que \(q\circ s=\id\), el punt \(s(b)\) cau a la fibra de \(b\), és a dir, a la classe que \(b\) representa: la secció tria un element per classe.}

\pregunta Quina d'aquestes afirmacions és certa \emph{sense} l'axioma d'elecció?
\begin{itemize}
\opcio{a}{Tot epimorfisme té secció.}
\opcio{b}{Tot epimorfisme és exhaustiu sobre els punts globals.}
\opcio{c}{Tot epimorfisme és exhaustiu sobre els elements generalitzats.}
\opcio{d}{Tot quocient admet un sistema de representants.}
\end{itemize}
\feedback
\justificacio{L'exhaustivitat sobre punts es va demostrar al capítol 3 amb la separació global i la disjunció de les sumes. Les altres tres afirmacions són formulacions equivalents de l'elecció, o conseqüències directes seves.}

\pregunta Per què es postula l'efectivitat de les relacions d'equivalència com a axioma en aquest llibre?
\begin{itemize}
\opcio{a}{Perquè és independent de tots els altres axiomes del llibre.}
\opcio{b}{Perquè és falsa en el model clàssic i cal imposar-la.}
\opcio{c}{Perquè la seva demostració necessita el conjunt de parts, que encara no s'ha introduït.}
\opcio{d}{Perquè equival a l'axioma d'elecció.}
\end{itemize}
\feedback
\justificacio{En la formulació categòrica és un teorema que es demostra a partir del conjunt de parts (capítol 7); aquí es postula per poder usar-la abans. En el model clàssic és certa, i no té res a veure amb l'elecció.}

\end{enumerate}
\clauestatica
\finalitzaTest

\chapter{Imatges i factoritzacions}\label{test:cap5}

Aquest test correspon al capítol~5 de la part teòrica. Cada pregunta té una única resposta correcta.

\iniciTest{u5}{c,b,b,a,c,a,c,b,c,b,b,a}

\begin{enumerate}

\pregunta Tota aplicació \(f\colon A\to B\) admet una factorització \(A\to\operatorname{Im}(f)\to B\) on\ldots
\begin{itemize}
\opcio{a}{els dos factors són monomorfismes.}
\opcio{b}{els dos factors són epimorfismes.}
\opcio{c}{el primer és un epimorfisme regular i el segon un monomorfisme.}
\opcio{d}{el primer és un monomorfisme i el segon un isomorfisme.}
\end{itemize}
\feedback
\justificacio{La imatge és el quocient pel parell nucli, un coequalitzador, seguit de la inclusió de la imatge dins del codomini, que és un monomorfisme: epi regular seguit de mono, única llevat d'isomorfisme.}

\pregunta La imatge \(\operatorname{Im}(f)\) es pot construir com\ldots
\begin{itemize}
\opcio{a}{el producte \(A\times B\).}
\opcio{b}{el quocient de \(A\) pel parell nucli de \(f\).}
\opcio{c}{el pullback de \(f\) amb la identitat.}
\opcio{d}{la suma \(A+B\).}
\end{itemize}
\feedback
\justificacio{El parell nucli relaciona els elements amb la mateixa imatge; el seu quocient els identifica, i el que queda és exactament el conjunt de valors assolits. El producte, el pullback amb la identitat i la suma no tenen res a veure.}

\pregunta Una aplicació \(f\) és un monomorfisme si i només si el seu parell nucli és\ldots
\begin{itemize}
\opcio{a}{tot el producte \(A\times A\).}
\opcio{b}{la diagonal \(\Delta_A\).}
\opcio{c}{inicial.}
\opcio{d}{un epimorfisme.}
\end{itemize}
\feedback
\justificacio{\(f\) és mono exactament quan \(f\circ x=f\circ x'\) força \(x=x'\), és a dir, quan el parell nucli no relaciona res més que cada element amb ell mateix: la diagonal. Si fos tot \(A\times A\), \(f\) seria constant.}

\pregunta Una aplicació \(f\) és un epimorfisme regular si i només si\ldots
\begin{itemize}
\opcio{a}{el monomorfisme \(m\) de la seva factorització és un isomorfisme (la imatge és tot el codomini).}
\opcio{b}{el seu parell nucli és la diagonal.}
\opcio{c}{el seu domini és terminal.}
\opcio{d}{no té imatge.}
\end{itemize}
\feedback
\justificacio{La imatge és tot el codomini exactament quan \(m\) és isomorfisme; aleshores \(f\) és un coequalitzador seguit d'un iso, és a dir, epi regular. I si \(f\) és epi, \(m\) és epi i mono, doncs iso: per això epi i epi regular coincideixen aquí.}

\pregunta L'adjunció entre imatge directa i imatge inversa diu que \(f_{*}S\subseteq_B T\) si i només si\ldots
\begin{itemize}
\opcio{a}{\(T\subseteq_B f_{*}S\).}
\opcio{b}{\(f^{*}T\subseteq_A S\).}
\opcio{c}{\(S\subseteq_A f^{*}T\).}
\opcio{d}{\(S\) i \(T\) són parts totals.}
\end{itemize}
\feedback
\justificacio{És l'adjunció \(f_*\dashv f^*\): la imatge de \(S\) cau dins \(T\) exactament quan \(S\) cau dins la preimatge de \(T\). Les altres opcions inverteixen el sentit d'alguna inclusió.}

\pregunta La imatge directa de la part total de \(A\) és\ldots
\begin{itemize}
\opcio{a}{la imatge de \(f\).}
\opcio{b}{la part buida de \(B\).}
\opcio{c}{sempre tot \(B\).}
\opcio{d}{el parell nucli de \(f\).}
\end{itemize}
\feedback
\justificacio{\(f_*A=\operatorname{Im}(f\circ\id_A)=\operatorname{Im}(f)\). Que sigui tot \(B\) equival a que \(f\) sigui epimorfisme.}

\pregunta La composició de dues relacions \(r'\circ r\) es defineix\ldots
\begin{itemize}
\opcio{a}{com el producte \(R\times R'\).}
\opcio{b}{com la suma de \(r\) i \(r'\).}
\opcio{c}{com la imatge de l'aplicació que recull els extrems del producte fibrat \(R\times_B R'\).}
\opcio{d}{com el pullback de \(r\) amb \(r'\).}
\end{itemize}
\feedback
\justificacio{El producte fibrat \(R\times_BR'\) encadena testimonis amb el mateix punt mitjà; l'aplicació als extrems dona parells \((a,c)\), i la imatge comprimeix els encadenaments en una part de \(A\times C\): aquí apareix el \enquote{existeix \(b\)}.}

\pregunta Que una relació \(r\) sigui \emph{funcional} vol dir que la seva primera projecció \(r_1\) és\ldots
\begin{itemize}
\opcio{a}{un epimorfisme regular (existència del testimoni).}
\opcio{b}{un monomorfisme (unicitat del testimoni).}
\opcio{c}{un isomorfisme.}
\opcio{d}{una aplicació constant.}
\end{itemize}
\feedback
\justificacio{Funcional vol dir com a molt un testimoni per element de \(A\), que és la injectivitat de \(r_1\); total, almenys un, que és la seva exhaustivitat. Isomorfisme és la conjunció.}

\pregunta Sota la correspondència entre aplicacions i relacions, les aplicacions \(A\to B\) es corresponen amb les relacions que són\ldots
\begin{itemize}
\opcio{a}{simètriques.}
\opcio{b}{només funcionals (aplicacions parcials).}
\opcio{c}{funcionals i totals.}
\opcio{d}{només totals.}
\end{itemize}
\feedback
\justificacio{Una aplicació és una relació amb exactament un \(y\) per a cada \(x\): funcional (unicitat) i total (existència); la primera projecció és un isomorfisme i \(f=r_2\circ r_1^{-1}\). Només funcional dona les aplicacions parcials.}

\pregunta Una aplicació és un isomorfisme si i només si és\ldots
\begin{itemize}
\opcio{a}{un monomorfisme.}
\opcio{b}{alhora monomorfisme i epimorfisme.}
\opcio{c}{un epimorfisme regular.}
\opcio{d}{una relació total.}
\end{itemize}
\feedback
\justificacio{Mono fa de \(e\) un isomorfisme i epi fa de \(m\) un isomorfisme; el compost ho és. És el teorema del bimorfisme llegit a la factorització.}

\pregunta Quin axioma s'usa per demostrar que el factor \(m\) de la factorització imatge és un monomorfisme?
\begin{itemize}
\opcio{a}{L'axioma d'elecció: cal una secció de \(e\).}
\opcio{b}{La separació global, reduint-ho als punts, i l'exhaustivitat dels epimorfismes sobre punts.}
\opcio{c}{Cap: és immediat de la propietat universal del coequalitzador.}
\opcio{d}{L'axioma d'efectivitat sol.}
\end{itemize}
\feedback
\justificacio{Mono equival a injectiu sobre punts (separació global); si \(m\circ y=m\circ y'\), les preimatges \(a,a'\) per \(e\) estan relacionades pel parell nucli i \(e\) les identifica. L'elecció dona una prova alternativa, però no és necessària.}

\pregunta Què es pot dir, sense l'axioma d'elecció, d'un epimorfisme \(f\) en aquesta teoria?
\begin{itemize}
\opcio{a}{Que és regular i estable per canvi de base, però no necessàriament escindit.}
\opcio{b}{Que és escindit.}
\opcio{c}{Res: sense elecció els epimorfismes no tenen cap propietat especial.}
\opcio{d}{Que és un isomorfisme.}
\end{itemize}
\feedback
\justificacio{Tot epimorfisme és regular (factorització imatge i teorema del bimorfisme) i els epimorfismes són estables per canvi de base (exhaustivitat sobre punts). Que tingui secció és exactament l'elecció.}

\end{enumerate}
\clauestatica
\finalitzaTest

\chapter{Exponencials}\label{test:cap6}

Aquest test correspon al capítol~6 de la part teòrica. Cada pregunta té una única resposta correcta.

\iniciTest{u6}{b,c,b,a,c,b,c,a,b,c,b,b}

\begin{enumerate}

\pregunta Un exponencial \(B^A\) ve donat per un conjunt \(E\) i una aplicació d'avaluació \(\operatorname{ev}\colon E\times A\to B\) tal que\ldots
\begin{itemize}
\opcio{a}{\(\operatorname{ev}\) és un isomorfisme.}
\opcio{b}{tota \(h\colon X\times A\to B\) té una única transposada \(\widehat h\colon X\to E\) amb \(\operatorname{ev}\circ(\widehat h\times\id_A)=h\).}
\opcio{c}{\(E\) és el producte \(A\times B\).}
\opcio{d}{\(\operatorname{ev}\) és un epimorfisme regular.}
\end{itemize}
\feedback
\justificacio{És la propietat universal: l'avaluació és universal entre les aplicacions \(E'\times A\to B\). L'avaluació no és iso ni epi regular en general, i \(E\) no és el producte.}

\pregunta La \emph{transposada} (o currificació) de \(h\colon X\times A\to B\) és\ldots
\begin{itemize}
\opcio{a}{una aplicació \(B\to X\times A\).}
\opcio{b}{una aplicació \(X\times A\to E\).}
\opcio{c}{l'única aplicació \(\widehat h\colon X\to B^A\) que descurrifica a \(h\).}
\opcio{d}{sempre un isomorfisme.}
\end{itemize}
\feedback
\justificacio{La transposada és l'única \(\widehat h\colon X\to B^A\) tal que \(\operatorname{ev}\circ(\widehat h\times\id_A)=h\): fixa la primera variable i deixa una aplicació de la segona. Només és iso en casos especials.}

\pregunta Els punts globals \(\mathbf1\to B^A\) es corresponen biunívocament amb\ldots
\begin{itemize}
\opcio{a}{els punts de \(A\).}
\opcio{b}{les aplicacions \(A\to B\).}
\opcio{c}{les parts de \(B\).}
\opcio{d}{els punts de \(B\).}
\end{itemize}
\feedback
\justificacio{És el cas \(X=\mathbf1\) de la propietat universal, amb \(\mathbf1\times A\cong A\): les aplicacions \(\mathbf1\times A\to B\) són les aplicacions \(A\to B\).}

\pregunta L'exponencial és únic\ldots
\begin{itemize}
\opcio{a}{llevat d'un únic isomorfisme compatible amb l'avaluació.}
\opcio{b}{de manera absoluta, com a individu de la teoria.}
\opcio{c}{només quan \(A\) és terminal.}
\opcio{d}{només amb l'axioma d'elecció.}
\end{itemize}
\feedback
\justificacio{Com tota construcció universal: dos exponencials estan lligats per un únic isomorfisme compatible amb les avaluacions. No hi ha unicitat absoluta, i l'elecció no hi intervé.}

\pregunta Quin dels següents isomorfismes és correcte?
\begin{itemize}
\opcio{a}{\(A^{\mathbf1}\cong\mathbf1\).}
\opcio{b}{\(A^{\mathbf0}\cong A\).}
\opcio{c}{\(A^{\mathbf1}\cong A\).}
\opcio{d}{\(\mathbf1^{A}\cong A\).}
\end{itemize}
\feedback
\justificacio{Les aplicacions \(\mathbf1\to A\) són els punts de \(A\): \(A^{\mathbf1}\cong A\). Els altres són \(A^{\mathbf0}\cong\mathbf1\) i \(\mathbf1^A\cong\mathbf1\).}

\pregunta El conjunt de parts \(P(A)\) es defineix com\ldots
\begin{itemize}
\opcio{a}{el producte \(A\times\mathbf2\).}
\opcio{b}{l'exponencial \(\mathbf2^{A}\).}
\opcio{c}{la suma \(A+\mathbf2\).}
\opcio{d}{el parell nucli de \(A\).}
\end{itemize}
\feedback
\justificacio{Els punts de \(\mathbf2^A\) són les aplicacions \(A\to\mathbf2\), és a dir, les característiques, que classifiquen les parts d'\(A\). Per això \(P(A):=\mathbf2^A\), sense cap axioma nou.}

\pregunta L'exponencial d'una suma satisfà\ldots
\begin{itemize}
\opcio{a}{\(B^{A+C}\cong B^{A}+B^{C}\).}
\opcio{b}{\(B^{A+C}\cong B^{A\times C}\).}
\opcio{c}{\(B^{A+C}\cong B^{A}\times B^{C}\).}
\opcio{d}{\(B^{A+C}\cong B^{A}\).}
\end{itemize}
\feedback
\justificacio{Una aplicació sobre una suma és un parell d'aplicacions sobre els sumands: via la distributivitat, \(X\times(A+C)\cong X\times A+X\times C\), i transposant, \(B^{A+C}\cong B^A\times B^C\).}

\pregunta La composició interna és una aplicació\ldots
\begin{itemize}
\opcio{a}{\(c\colon C^{B}\times B^{A}\to C^{A}\) que sobre punts envia \((g,f)\) a \(g\circ f\).}
\opcio{b}{\(c\colon C^{A}\to C^{B}\times B^{A}\).}
\opcio{c}{que només existeix si \(A=B=C\).}
\opcio{d}{que coincideix amb l'avaluació.}
\end{itemize}
\feedback
\justificacio{La composició interna és la transposada de l'avaluació iterada \((C^B\times B^A)\times A\to C^B\times B\to C\); sobre punts fa \((g,f)\mapsto g\circ f\), i és única per la separació global.}

\pregunta Un element generalitzat \(X\to B^A\) es llegeix, en general, com\ldots
\begin{itemize}
\opcio{a}{un punt de \(B\).}
\opcio{b}{una família d'aplicacions \(A\to B\) parametritzada per \(X\).}
\opcio{c}{una sola aplicació \(A\to B\).}
\opcio{d}{una part de \(A\).}
\end{itemize}
\feedback
\justificacio{Un element generalitzat \(X\to B^A\) es transposa en \(X\times A\to B\): per a cada \(x\) de \(X\), una aplicació \(A\to B\). Només quan \(X=\mathbf1\) és una sola aplicació.}

\pregunta L'existència d'exponencials per a tot parell de conjunts\ldots
\begin{itemize}
\opcio{a}{es dedueix dels axiomes anteriors.}
\opcio{b}{només val per a conjunts finits.}
\opcio{c}{es postula com un axioma nou.}
\opcio{d}{és equivalent a l'axioma d'elecció.}
\end{itemize}
\feedback
\justificacio{Els exponencials es postulen; en la formulació categòrica són, amb el classificador, l'axioma que fa de la teoria un topos. No es dedueixen dels axiomes anteriors ni tenen res a veure amb l'elecció.}

\pregunta Quin axioma garanteix que la composició interna \(c\colon C^B\times B^A\to C^A\) és l'\emph{única} aplicació que sobre punts envia \((g,f)\) a \(g\circ f\)?
\begin{itemize}
\opcio{a}{L'axioma d'exponencials, per la unicitat de la transposada.}
\opcio{b}{La separació global: les aplicacions queden determinades pels punts.}
\opcio{c}{L'axioma d'elecció.}
\opcio{d}{Cap: l'existència ja implica la unicitat.}
\end{itemize}
\feedback
\justificacio{La transposada és única per a una aplicació \((C^B\times B^A)\times A\to C\) ja fixada; però la condició \enquote{sobre punts envia \((g,f)\) a \(g\circ f\)} només fixa els valors de \(c\) sobre punts globals, i és la separació global la que fa que això la determini.}

\pregunta Quants punts globals té \(\mathbf2^{\mathbf2}\), i per què?
\begin{itemize}
\opcio{a}{Dos, perquè \(\mathbf2^{\mathbf2}\cong\mathbf2\).}
\opcio{b}{Quatre, perquè \(\mathbf2^{\mathbf2}\cong\mathbf2^{\mathbf1+\mathbf1}\cong\mathbf2\times\mathbf2\), i \(\mathbf2\times\mathbf2\) té quatre punts.}
\opcio{c}{Un, perquè és terminal.}
\opcio{d}{Infinits.}
\end{itemize}
\feedback
\justificacio{Són les quatre aplicacions \(\mathbf2\to\mathbf2\): les dues constants, la identitat i la negació. L'exponencial d'una suma i \(A^{\mathbf1}\cong A\) donen \(\mathbf2\times\mathbf2\), els punts del qual són les quatre injeccions \(\delta_{ab}\).}

\end{enumerate}
\clauestatica
\finalitzaTest

\chapter{El conjunt de parts i la pertinença interna}\label{test:cap7}

Aquest test correspon al capítol~7 de la part teòrica. Cada pregunta té una única resposta correcta.

\iniciTest{u7}{b,c,b,a,c,b,a,c,b,a,c,a}

\begin{enumerate}

\pregunta La relació de pertinença interna \(\in_A\) és la part de \(A\times P(A)\) classificada per\ldots
\begin{itemize}
\opcio{a}{la negació \(\neg\).}
\opcio{b}{l'avaluació \(\operatorname{ev}\) (llegida sobre \(A\times P(A)\)).}
\opcio{c}{la diagonal de \(A\).}
\opcio{d}{el parell nucli de \(A\).}
\end{itemize}
\feedback
\justificacio{L'avaluació \(\operatorname{ev}\colon P(A)\times A\to\mathbf2\) és una aplicació cap als valors de veritat, i el classificador la converteix en la característica d'una part; aquesta part, llegida sobre \(A\times P(A)\), és \(\in_A\).}

\pregunta Que \(a\mathrel{\varepsilon}\alpha\) (amb \(\alpha=\ulcorner m\urcorner\) el nom d'una part \(m\)) equival a\ldots
\begin{itemize}
\opcio{a}{\(a\) és un element global.}
\opcio{b}{\(m\) és la part total.}
\opcio{c}{\(a\in_X m\).}
\opcio{d}{\(\alpha\) és constant.}
\end{itemize}
\feedback
\justificacio{Descurrificant el nom \(\ulcorner m\urcorner\) es recupera \(\chi_m\), de manera que \(\operatorname{ev}\circ\langle\alpha,a\rangle=\chi_m\circ a\), i això val \(1\circ{!}_X\) exactament quan \(a\in_Xm\). La pertinença interna estén la generalitzada.}

\pregunta Un element generalitzat \(X\to P(A)\) es correspon amb\ldots
\begin{itemize}
\opcio{a}{un punt de \(A\).}
\opcio{b}{una part de \(X\times A\) (una família de parts de \(A\)).}
\opcio{c}{una sola part de \(A\).}
\opcio{d}{una aplicació \(A\to X\).}
\end{itemize}
\feedback
\justificacio{Transposant \(X\to\mathbf2^A\) s'obté \(X\times A\to\mathbf2\), que classifica una part de \(X\times A\). Només quan \(X=\mathbf1\) és una sola part d'\(A\); en general és una família indexada per \(X\).}

\pregunta La intersecció interna \(\alpha\cap\beta\) es defineix com la transposada de\ldots
\begin{itemize}
\opcio{a}{\(\wedge\circ\langle\operatorname{ev}_1,\operatorname{ev}_2\rangle\).}
\opcio{b}{\(\vee\circ\langle\operatorname{ev}_1,\operatorname{ev}_2\rangle\).}
\opcio{c}{\(\neg\circ\operatorname{ev}\).}
\opcio{d}{la projecció \(P(A)\times P(A)\to P(A)\).}
\end{itemize}
\feedback
\justificacio{És la definició: la conjunció puntual de les dues avaluacions, transposada perquè actuï sobre \(P(A)\). Amb \(\vee\) s'obté la unió interna i amb \(\neg\) el complement intern.}

\pregunta Sobre els noms de parts, les operacions internes\ldots
\begin{itemize}
\opcio{a}{no tenen cap relació amb les operacions de parts.}
\opcio{b}{intercanvien unió i intersecció.}
\opcio{c}{recuperen les operacions de parts: \(\ulcorner m\urcorner\cap\ulcorner n\urcorner=\ulcorner m\cap n\urcorner\).}
\opcio{d}{sempre donen la part buida.}
\end{itemize}
\feedback
\justificacio{Les descurrificacions de \(\ulcorner m\urcorner\cap\ulcorner n\urcorner\), \(\ulcorner m\urcorner\cup\ulcorner n\urcorner\) i \(\neg\ulcorner m\urcorner\) són \(\wedge\circ\langle\chi_m,\chi_n\rangle\), \(\vee\circ\langle\chi_m,\chi_n\rangle\) i \(\neg\circ\chi_m\), que són les característiques de \(m\cap n\), \(m\cup n\) i \(m^c\).}

\pregunta La inclusió interna \(\subseteq_A\) sobre \(P(A)\) és\ldots
\begin{itemize}
\opcio{a}{una relació d'equivalència.}
\opcio{b}{un ordre parcial (reflexiu, antisimètric, transitiu).}
\opcio{c}{un isomorfisme.}
\opcio{d}{sempre buida.}
\end{itemize}
\feedback
\justificacio{Sobre les parts associades, \(\subseteq_A\) es llegeix com la inclusió de parts de \(X\times A\), que és reflexiva i transitiva; l'antisimetria ve de la unicitat de la característica i de la transposada.}

\pregunta Respecte de l'ordre d'inclusió intern, la intersecció i la unió internes són\ldots
\begin{itemize}
\opcio{a}{l'ínfim i el suprem, respectivament.}
\opcio{b}{el màxim i el mínim absoluts.}
\opcio{c}{sense relació amb l'ordre.}
\opcio{d}{tots dos l'ínfim.}
\end{itemize}
\feedback
\justificacio{\(s_{\alpha\cap\beta}=s_\alpha\cap s_\beta\) i \(s_{\alpha\cup\beta}=s_\alpha\cup s_\beta\); una part és inclosa en la intersecció si i només si ho és en totes dues, i la unió és la part més petita que les conté. Mínim i màxim són \(\varnothing_A\) i \(A_A\).}

\pregunta La imatge inversa interna \(P(f)\colon P(B)\to P(A)\), sobre els noms de parts, recupera\ldots
\begin{itemize}
\opcio{a}{la imatge directa \(f_{*}\).}
\opcio{b}{el complement.}
\opcio{c}{la imatge inversa \(f^{-1}\).}
\opcio{d}{la identitat.}
\end{itemize}
\feedback
\justificacio{La descurrificació de \(P(f)\circ\ulcorner n\urcorner\) és \(\chi_n\circ f\), que és la característica de \(f^{-1}(n)\). La imatge directa és \(\exists_f\), no \(P(f)\).}

\pregunta Els quantificadors \(\exists_f\) i \(\forall_f\) són, respectivament,\ldots
\begin{itemize}
\opcio{a}{tots dos adjunts per l'esquerra de \(P(f)\).}
\opcio{b}{adjunt per l'esquerra i adjunt per la dreta de \(P(f)\).}
\opcio{c}{iguals a \(P(f)\).}
\opcio{d}{sense relació amb \(P(f)\).}
\end{itemize}
\feedback
\justificacio{Sobre parts associades, \(\exists_f\) és la imatge directa i \(P(f)\) la imatge inversa al llarg de \(\id\times f\), i l'adjunció \(f_*\dashv f^{-1}\) del capítol 5 dona l'esquerra; \(\forall_f=\neg\exists_f\neg\) és l'adjunt dret per dualitat booleana.}

\pregunta Les operacions \(\exists_A\) i \(\forall_A\) (per \({!}_A\colon A\to\mathbf1\)) afirmen, respectivament, que una part de \(A\)\ldots
\begin{itemize}
\opcio{a}{és no buida, i que és total.}
\opcio{b}{és buida, i que és no buida.}
\opcio{c}{és un isomorfisme.}
\opcio{d}{coincideix amb la diagonal.}
\end{itemize}
\feedback
\justificacio{Per a \({!}_A\colon A\to\mathbf1\), \(\exists_A\ulcorner m\urcorner\) és el nom de la imatge de \(S\) dins \(\mathbf1\), total si i només si \(S\) té algun punt; i \(\forall_A\ulcorner m\urcorner=\neg\exists_A\ulcorner m^c\urcorner\) és \(1\) si i només si \(m^c\) és buida, és a dir, \(m\) total.}

\pregunta Per a elements generalitzats \(\alpha,\beta\colon X\to P(A)\) i \(a\colon X\to A\), quina lectura val a \emph{tots} els estadis?
\begin{itemize}
\opcio{a}{\(a\mathrel{\varepsilon}(\alpha\cup\beta)\) si i només si \(a\mathrel{\varepsilon}\alpha\) o \(a\mathrel{\varepsilon}\beta\).}
\opcio{b}{\(a\mathrel{\varepsilon}(\neg\alpha)\) si i només si no \(a\mathrel{\varepsilon}\alpha\).}
\opcio{c}{\(a\mathrel{\varepsilon}(\alpha\cap\beta)\) si i només si \(a\mathrel{\varepsilon}\alpha\) i \(a\mathrel{\varepsilon}\beta\).}
\opcio{d}{Totes tres.}
\end{itemize}
\feedback
\justificacio{La conjunció és la característica de \(\delta_{11}\), i \(\langle u,v\rangle\) es factoritza per \(\delta_{11}\) si i només si \(u=v=1\circ{!}\); això val a tot estadi. Per a \(\cup\) i \(\neg\) un element generalitzat pot quedar repartit (\(X=\mathbf2\), \(u=\id\), \(v=\neg\)): les lectures amb \enquote{o} i \enquote{no} només valen sobre punts globals.}

\pregunta Com es construeix el quantificador existencial \(\exists_f\colon P(A)\to P(B)\) com a aplicació?
\begin{itemize}
\opcio{a}{Com la transposada de la característica de la imatge directa de la pertinença \(E_A\) al llarg de \(\id\times f\).}
\opcio{b}{Com la transposada de \(\operatorname{ev}\circ(\id\times f)\).}
\opcio{c}{Postulant-lo amb un axioma nou.}
\opcio{d}{Com la composició \(P(f)\circ\neg_B\).}
\end{itemize}
\feedback
\justificacio{La imatge directa del capítol 5, aplicada a la part universal \(E_A\hookrightarrow P(A)\times A\) al llarg de \(\id\times f\), dona una part de \(P(A)\times B\); la seva característica es transposa a \(\exists_f\). L'opció (b) és \(P(f)\), la imatge inversa. No cal cap axioma: tot surt dels exponencials i de la factorització imatge.}

\end{enumerate}
\clauestatica
\finalitzaTest

\chapter{L'objecte dels nombres naturals}\label{test:cap8}

Aquest test correspon al capítol~8 de la part teòrica. Cada pregunta té una única resposta correcta.

\iniciTest{u8}{b,c,b,a,c,b,c,a,b,c,a,b}

\begin{enumerate}

\pregunta L'objecte dels nombres naturals ve donat per una terna \((\N,0,s)\) amb la propietat universal que, per a tot sistema recursiu \((A,a,t)\),\ldots
\begin{itemize}
\opcio{a}{existeixen infinites aplicacions \(\N\to A\) compatibles.}
\opcio{b}{existeix una única \(u\colon\N\to A\) amb \(u\circ0=a\) i \(u\circ s=t\circ u\).}
\opcio{c}{\(A\) és isomorf a \(\N\).}
\opcio{d}{\(t\) ha de ser un isomorfisme.}
\end{itemize}
\feedback
\justificacio{És la propietat universal de recursió: existència i unicitat de l'aplicació que respecta el zero i el successor. \(A\) i \(t\) són arbitraris; no cal cap isomorfisme.}

\pregunta L'objecte dels nombres naturals és únic\ldots
\begin{itemize}
\opcio{a}{només si s'admet l'axioma d'elecció.}
\opcio{b}{de manera absoluta.}
\opcio{c}{llevat d'un únic isomorfisme compatible amb \(0\) i \(s\).}
\opcio{d}{no és únic.}
\end{itemize}
\feedback
\justificacio{Dues ternes amb la propietat universal es corresponen per recursores mútues, les composicions de les quals són identitats per unicitat. Com sempre amb les propietats universals, la unicitat és llevat d'un únic isomorfisme compatible.}

\pregunta El principi d'inducció afirma que una part de \(\N\) que conté el zero i és tancada pel successor\ldots
\begin{itemize}
\opcio{a}{és inicial.}
\opcio{b}{és tot \(\N\).}
\opcio{c}{és un exponencial.}
\opcio{d}{és terminal.}
\end{itemize}
\feedback
\justificacio{És el principi d'inducció: el morfisme de sistemes recursius \(S\to\N\) té una recursora inversa per l'esquerra, de manera que és retracció i mono, i per tant isomorfisme.}

\pregunta El principi d'inducció es correspon amb quina meitat de la propietat universal de recursió?
\begin{itemize}
\opcio{a}{La unicitat de la recursora.}
\opcio{b}{L'existència de la recursora.}
\opcio{c}{Cap de les dues.}
\opcio{d}{Totes dues alhora.}
\end{itemize}
\feedback
\justificacio{Que dues aplicacions amb la mateixa recursió coincideixin és la unicitat de la recursora, i és exactament el que diu la inducció aplicada al seu equalitzador. L'existència és la definició per recursió.}

\pregunta Les equacions de Peano ---que el zero no és cap successor i que el successor és injectiu--- en aquest llibre són\ldots
\begin{itemize}
\opcio{a}{axiomes primitius.}
\opcio{b}{indemostrables.}
\opcio{c}{teoremes derivats de la propietat universal.}
\opcio{d}{falsos en general.}
\end{itemize}
\feedback
\justificacio{El zero no és successor es prova amb la recursora cap a \(\mathbf2\) que val \(1\) al zero i \(0\) als successors; la injectivitat, amb el predecessor construït sobre \(\N\times\N\). Totes dues surten de la propietat universal.}

\pregunta La suma \(m+n\) es defineix\ldots
\begin{itemize}
\opcio{a}{com un pullback.}
\opcio{b}{per recursió, com la iterada del successor.}
\opcio{c}{com un producte d'exponencials.}
\opcio{d}{com el parell nucli de \(s\).}
\end{itemize}
\feedback
\justificacio{\(m+n:=s^n(m)\), la iterada del successor aplicada \(n\) vegades a \(m\), que és una recursió sobre el segon argument. El producte itera la suma i l'exponenciació itera el producte.}

\pregunta Les lleis de l'aritmètica (associativitat, commutativitat, distributivitat) es demostren\ldots
\begin{itemize}
\opcio{a}{postulant-les.}
\opcio{b}{per l'axioma d'elecció.}
\opcio{c}{per inducció sobre un dels arguments, amb els altres com a paràmetres.}
\opcio{d}{només en el model clàssic.}
\end{itemize}
\feedback
\justificacio{Cada llei és una igualtat d'aplicacions \(\N^k\to\N\), i es demostra per inducció sobre un argument amb els altres com a paràmetres (inducció amb paràmetres, via l'exponencial). No cal l'elecció ni postular res.}

\pregunta L'ordre \(m\le n\) es defineix per\ldots
\begin{itemize}
\opcio{a}{la imatge de \((m,k)\mapsto(m,m+k)\): sobre punts, l'existència d'un \(k\) amb \(m+k=n\).}
\opcio{b}{\(m\) i \(n\) tenen el mateix successor.}
\opcio{c}{\(m\cdot n=0\).}
\opcio{d}{\(m=n\) exclusivament.}
\end{itemize}
\feedback
\justificacio{\(\operatorname{Leq}\) és la imatge de \((m,k)\mapsto(m,m+k)\): sobre punts, \(m\le n\) si i només si hi ha un increment \(k\) amb \(m+k=n\). És una part de \(\N\times\N\), no una fórmula externa.}

\pregunta El principi del bon ordre afirma que tota part no inicial de \(\N\)\ldots
\begin{itemize}
\opcio{a}{és tot \(\N\).}
\opcio{b}{té un element mínim.}
\opcio{c}{té un element màxim.}
\opcio{d}{és un exponencial.}
\end{itemize}
\feedback
\justificacio{És el principi del bon ordre: tota part no inicial té un mínim, únic per antisimetria. Es demostra per inducció sobre la part auxiliar \enquote{cap natural \(\le n\) no és a \(P\)}. No hi ha màxims en general.}

\pregunta La recursió amb paràmetres (dades \(a\colon X\to A\), \(t\colon X\times A\to A\)) es redueix a la recursió amb dades globals\ldots
\begin{itemize}
\opcio{a}{prenent el parell nucli.}
\opcio{b}{amb l'axioma d'elecció.}
\opcio{c}{transposant respecte de \(X\) cap a l'exponencial \(A^{X}\).}
\opcio{d}{no es pot reduir.}
\end{itemize}
\feedback
\justificacio{Una aplicació \(X\times\N\to A\) es transposa a \(\N\to A^X\), on les dades amb paràmetre esdevenen un punt \(\ulcorner a\urcorner\) i una autoaplicació \(T\) de \(A^X\); l'axioma s'hi aplica i es descurrifica. Per això els exponencials van abans que \(\N\).}

\pregunta Per què el predecessor \(p\) (amb \(p\circ s=\id_\N\)) no es pot definir com la recursora del sistema \((\N,0,\id_\N)\)?
\begin{itemize}
\opcio{a}{Perquè aquesta recursora compleix \(u\circ s=u\), i és la constant \(0\); cal una recursió sobre \(\N\times\N\) que recordi el parell \((n-1,n)\).}
\opcio{b}{Perquè \(\id_\N\) no és una autoaplicació.}
\opcio{c}{Perquè el predecessor no existeix sense l'axioma d'elecció.}
\opcio{d}{Sí que es pot: és exactament així com es fa.}
\end{itemize}
\feedback
\justificacio{La forma de l'axioma és \(u\circ s=t\circ u\); amb \(t=\id\) s'obté \(u\circ s=u\), que amb \(u\circ0=0\) força \(u=0\circ{!}\). El predecessor és la primera projecció de la recursora \(q\) de \((\N\times\N,\langle0,0\rangle,\langle p_2,s\circ p_2\rangle)\), després de veure que \(p_2\circ q=\id\).}

\pregunta La totalitat de l'ordre, \enquote{\(m\le n\) o \(n\le m\)}, s'enuncia per a punts globals. Què diu, com a afirmació sobre parts?
\begin{itemize}
\opcio{a}{Que \(\operatorname{Leq}\) és la part total de \(\N\times\N\).}
\opcio{b}{Que \(\operatorname{Leq}\cup\operatorname{Inv}(\operatorname{Leq})\) és la part total de \(\N\times\N\); un element generalitzat pot caure a la unió sense caure sencer en cap dels dos membres.}
\opcio{c}{Que \(\operatorname{Leq}\cap\operatorname{Inv}(\operatorname{Leq})\) és la diagonal.}
\opcio{d}{Que tot element generalitzat \(\langle m,n\rangle\) cau sencer a \(\operatorname{Leq}\) o a \(\operatorname{Inv}(\operatorname{Leq})\).}
\end{itemize}
\feedback
\justificacio{Per la separació global, una part de \(\N\times\N\) és total si conté tots els punts; sobre punts, pertànyer a una unió és \enquote{o}. Sobre un estadi com \(\mathbf2\) un element es pot repartir entre els dos membres, de manera que (d) és fals. L'opció (c) és l'antisimetria.}

\end{enumerate}
\clauestatica
\finalitzaTest

\chapter{Conjunts finits i cardinalitat finita}\label{test:cap9}

Aquest test correspon al capítol~9 de la part teòrica. Cada pregunta té una única resposta correcta.

\iniciTest{u9}{b,c,b,a,c,b,a,c,b,c,a,c}

\begin{enumerate}

\pregunta El segment inicial satisfà la descomposició successora\ldots
\begin{itemize}
\opcio{a}{\(\seg{s\circ n}\cong\seg{n}\times\mathbf2\).}
\opcio{b}{\(\seg{s\circ n}\cong\seg{n}+\mathbf1\).}
\opcio{c}{\(\seg{s\circ n}\cong\seg{n}\).}
\opcio{d}{\(\seg{s\circ n}\cong\mathbf1\).}
\end{itemize}
\feedback
\justificacio{Un punt \(k<s\circ n\) compleix \(k\le n\): o bé \(k<n\) o bé \(k=n\), i no totes dues coses. Les dues parts complementàries són \(\seg{n}\) i el punt \(n\), i formen la suma.}

\pregunta La finitud i la cardinalitat internes es basen en\ldots
\begin{itemize}
\opcio{a}{els ordinals metanivell \([k]\).}
\opcio{b}{un conjunt de tots els conjunts.}
\opcio{c}{la família interna dels segments \(n\mapsto\seg{n}\).}
\opcio{d}{l'axioma d'elecció.}
\end{itemize}
\feedback
\justificacio{Els segments són les fibres de la part \(\operatorname{Lt}\hookrightarrow\N\times\N\), una família interna indexada per \(\N\); els \([k]\) només existeixen per a numerals del metallenguatge, i no hi ha cap conjunt de tots els conjunts.}

\pregunta Un punt global de \(\operatorname{Isos}(A,B)\) és\ldots
\begin{itemize}
\opcio{a}{un monomorfisme \(A\to B\).}
\opcio{b}{un isomorfisme \(A\to B\).}
\opcio{c}{una part de \(A\times B\).}
\opcio{d}{un element de \(A\).}
\end{itemize}
\feedback
\justificacio{\(\operatorname{Isos}(A,B)\) és la imatge, dins \(B^A\), dels parells \((f,g)\) amb \(g\circ f=\id\) i \(f\circ g=\id\); un punt seu és una \(f\) que admet inversa, és a dir, un isomorfisme.}

\pregunta Un conjunt \(A\) és \emph{finit} quan\ldots
\begin{itemize}
\opcio{a}{\(\operatorname{Card}(A)\) és no inicial, és a dir, existeix \(n\) amb \(\seg{n}\approx A\).}
\opcio{b}{\(A\) és terminal.}
\opcio{c}{\(A\) admet un element global.}
\opcio{d}{\(A\cong\N\).}
\end{itemize}
\feedback
\justificacio{És la definició: \(\operatorname{Card}(A)\) és la imatge de la projecció de les bijeccions codificades \(\seg{n}\to A\) sobre \(\N\), i \(A\) és finit quan aquesta part té algun punt. Ser terminal o tenir un punt són condicions molt més fortes; \(\N\) no és finit.}

\pregunta El cardinal \(|A|\) d'un conjunt finit està ben definit perquè\ldots
\begin{itemize}
\opcio{a}{s'admet l'axioma d'elecció.}
\opcio{b}{\(\N\) és infinit.}
\opcio{c}{un isomorfisme \(\seg{m}\cong\seg{n}\) força \(m=n\).}
\opcio{d}{tot conjunt és finit.}
\end{itemize}
\feedback
\justificacio{El lema de comparació de segments dona \(m\le n\) de tot monomorfisme \(\seg{m}\hookrightarrow\seg{n}\); un isomorfisme dona les dues desigualtats i \(m=n\). Cap elecció no hi intervé.}

\pregunta El principi dels calaixos diu que si \(|A|>|B|\) (finits), aleshores\ldots
\begin{itemize}
\opcio{a}{\(A\cong B\).}
\opcio{b}{no hi ha cap monomorfisme \(A\hookrightarrow B\).}
\opcio{c}{no hi ha cap aplicació \(A\to B\).}
\opcio{d}{\(B\) és inicial.}
\end{itemize}
\feedback
\justificacio{Un monomorfisme \(A\hookrightarrow B\) donaria un monomorfisme \(\seg{|A|}\hookrightarrow\seg{|B|}\) i, per comparació de segments, \(|A|\le|B|\). Sí que hi ha aplicacions \(A\to B\), però cap d'injectiva.}

\pregunta Per a \(f\colon A\to B\) entre finits amb \(|A|=|B|\), ser monomorfisme, epimorfisme o isomorfisme\ldots
\begin{itemize}
\opcio{a}{són tres condicions equivalents.}
\opcio{b}{no tenen cap relació.}
\opcio{c}{només coincideixen si \(A\) és terminal.}
\opcio{d}{mai no es donen alhora.}
\end{itemize}
\feedback
\justificacio{Mono \(\Rightarrow\) iso pel recompte de la suma aplicat a \(B=\operatorname{Im}(f)+C\); epi \(\Rightarrow\) iso escindint \(f\) i aplicant l'anterior a la secció. És el tret propi del cas finit: a \(\N\) el successor és mono i no epi.}

\pregunta El nombre de parts d'un conjunt finit amb \(|A|=n\) és\ldots
\begin{itemize}
\opcio{a}{\(n\).}
\opcio{b}{\(n^{2}\).}
\opcio{c}{\(2^{n}\).}
\opcio{d}{\(2\cdot n\).}
\end{itemize}
\feedback
\justificacio{\(P(A)\cong\mathbf2^{A}\), i per inducció \(|\mathbf2^{\seg{s\circ n}}|=|\mathbf2^{\seg{n}}\times\mathbf2|=2\cdot2^n\), amb \(2^0=1\).}

\pregunta El nombre d'aplicacions \(A\to B\) entre finits amb \(|A|=n\), \(|B|=m\) és\ldots
\begin{itemize}
\opcio{a}{\(n\cdot m\).}
\opcio{b}{\(m^{n}\).}
\opcio{c}{\(n^{m}\).}
\opcio{d}{\(n+m\).}
\end{itemize}
\feedback
\justificacio{Per inducció sobre \(n\) amb l'exponencial d'una suma: \(|\seg{m}^{\seg{s\circ n}}|=|\seg{m}^{\seg{n}}\times\seg{m}|=m^n\cdot m\). L'exponent és el cardinal del domini.}

\pregunta Tota part \(S\) d'un conjunt finit \(A\)\ldots
\begin{itemize}
\opcio{a}{és sempre infinita.}
\opcio{b}{té el mateix cardinal que \(A\).}
\opcio{c}{és finita amb \(|S|\le|A|\).}
\opcio{d}{és inicial.}
\end{itemize}
\feedback
\justificacio{Per inducció sobre el cardinal d'\(A\), descomponent la part en el que cau sobre \(\seg{n}\) i el que cau sobre el punt nou; el recompte de la suma dona \(|S|\le n+1\). El complement també és finit.}

\pregunta Per què la inducció del lema de comparació de segments es fa sobre el predicat codificat \(P(m)\) i no directament sobre \enquote{existeix un monomorfisme \(\seg{m}\hookrightarrow\seg{n}\)}?
\begin{itemize}
\opcio{a}{Perquè el principi d'inducció només s'aplica a parts de \(\N\), i \enquote{existeix un monomorfisme entre dos objectes que varien amb \(m,n\)} només és una part un cop codificat dins un exponencial fix.}
\opcio{b}{Perquè la inducció necessita l'axioma d'elecció.}
\opcio{c}{Perquè els segments no són conjunts.}
\opcio{d}{Per comoditat; es podria fer directament.}
\end{itemize}
\feedback
\justificacio{Les aplicacions \(\seg{m}\to\seg{n}\) s'estenen a \(\N\to\N+\mathbf1\), i \enquote{existeix un monomorfisme} esdevé la no-inicialitat d'una fibra de la part \(\mathcal M\) de \(\N\times\N\times(\N+\mathbf1)^{\N}\): una fórmula interna en \(m\). Sense això, la inducció seria un esquema extern.}

\pregunta On intervé l'axioma d'elecció en aquest capítol?
\begin{itemize}
\opcio{a}{En la construcció de \(\operatorname{Card}(A)\), per triar un isomorfisme.}
\opcio{b}{En la unicitat del cardinal.}
\opcio{c}{Només per escindir l'epimorfisme a \enquote{epi \(\Rightarrow\) iso en cardinal igual}, i s'hi pot evitar amb el bon ordre de \(\N\).}
\opcio{d}{A tots els recomptes.}
\end{itemize}
\feedback
\justificacio{\(\operatorname{Card}(A)\) és una imatge, sense cap tria; la comparació i els recomptes són induccions. L'única secció que cal, la d'un epimorfisme entre finits, es pot definir prenent el mínim de cada fibra, que és únic.}

\end{enumerate}
\clauestatica
\finalitzaTest

\chapter{Infinitud, enumerabilitat i seqüències}\label{test:cap10}

Aquest test correspon al capítol~10 de la part teòrica. Cada pregunta té una única resposta correcta.

\iniciTest{u10}{b,c,a,b,c,b,a,c,b,c,b,b}

\begin{enumerate}

\pregunta Un conjunt \(A\) és \emph{infinit} quan\ldots
\begin{itemize}
\opcio{a}{admet un monomorfisme no exhaustiu.}
\opcio{b}{\(\operatorname{Card}(A)\) és inicial: no hi ha cap \(n\) amb \(\seg{n}\approx A\).}
\opcio{c}{és isomorf a \(\N\).}
\opcio{d}{no té elements globals.}
\end{itemize}
\feedback
\justificacio{Infinit és la negació de finit: \(\operatorname{Card}(A)\) no té cap punt. Admetre un endomonomorfisme no exhaustiu és ser infinit de Dedekind, una noció diferent fins que no s'usa l'elecció; i hi ha infinits no isomorfs a \(\N\).}

\pregunta Un conjunt \(A\) és \emph{infinit de Dedekind} quan\ldots
\begin{itemize}
\opcio{a}{no és isomorf a cap ordinal estàndard.}
\opcio{b}{tota aplicació \(A\to A\) és exhaustiva.}
\opcio{c}{existeix un monomorfisme \(A\hookrightarrow A\) no exhaustiu.}
\opcio{d}{és finit.}
\end{itemize}
\feedback
\justificacio{És la definició de Dedekind: \(A\) es posa en bijecció amb una part pròpia seva. Que tota aplicació \(A\to A\) sigui exhaustiva és fals en general; (a) és la finitud externa, i (d) és el contrari.}

\pregunta El testimoni que \(\N\) és infinit de Dedekind és\ldots
\begin{itemize}
\opcio{a}{el successor \(s\), monomorfisme no exhaustiu.}
\opcio{b}{la identitat.}
\opcio{c}{el zero.}
\opcio{d}{el parell nucli de \(s\).}
\end{itemize}
\feedback
\justificacio{El successor és injectiu (té el predecessor com a retracció) i no exhaustiu (el zero no és successor): \(\N\) es posa en bijecció amb la seva part pròpia \(\N\setminus\{0\}\).}

\pregunta La implicació \enquote{tot conjunt infinit és infinit de Dedekind}\ldots
\begin{itemize}
\opcio{a}{no usa cap axioma.}
\opcio{b}{requereix (alguna forma d')elecció, com l'elecció numerable.}
\opcio{c}{és falsa.}
\opcio{d}{és equivalent al principi del mínim.}
\end{itemize}
\feedback
\justificacio{Cal triar alhora una llista injectiva de cada llargada; sense cap elecció és consistent que hi hagi infinits que no siguin de Dedekind. La direcció contrària no usa res, i el principi del mínim no hi té res a veure.}

\pregunta Un conjunt \(A\) és \emph{enumerable} quan és inicial o bé\ldots
\begin{itemize}
\opcio{a}{admet un monomorfisme \(A\hookrightarrow\N\).}
\opcio{b}{és finit.}
\opcio{c}{existeix un epimorfisme regular \(\N\twoheadrightarrow A\).}
\opcio{d}{és isomorf a \(\mathcal P(\N)\).}
\end{itemize}
\feedback
\justificacio{Enumerable és ser imatge de \(\N\): els elements es llisten, amb repeticions. Un monomorfisme \(A\hookrightarrow\N\) dona una part de \(\N\), que també és enumerable, però la definició és per l'epimorfisme.}

\pregunta Quina d'aquestes operacions preserva l'enumerabilitat \emph{sense} elecció?
\begin{itemize}
\opcio{a}{La unió d'una família numerable arbitrària.}
\opcio{b}{La suma i el producte de dos enumerables, i les parts i imatges.}
\opcio{c}{El producte d'una família numerable.}
\opcio{d}{Cap.}
\end{itemize}
\feedback
\justificacio{Suma i producte es fan amb \(\N\cong\N+\N\) i un aparellament; parts amb una retracció i imatges per composició: tot amb les enumeracions donades. La unió i el producte d'una família numerable demanen triar.}

\pregunta Tota part \(P\hookrightarrow\N\)\ldots
\begin{itemize}
\opcio{a}{és enumerable, i finita o isomorfa a \(\N\), sense elecció.}
\opcio{b}{és sempre infinita.}
\opcio{c}{requereix l'elecció per enumerar-se.}
\opcio{d}{és isomorfa a \(\mathcal P(\N)\).}
\end{itemize}
\feedback
\justificacio{Pel bon ordre, l'enumeració creixent \(a_{n+1}=\) mínim dels elements més grans que \(a_n\) es defineix per recursió simple; o s'estabilitza (part finita) o és una bijecció amb \(\N\).}

\pregunta El fet que els subobjectes de \(\N\) s'enumerin sense elecció es deu a\ldots
\begin{itemize}
\opcio{a}{l'axioma d'exponencials.}
\opcio{b}{la finitud de \(\N\).}
\opcio{c}{el principi del mínim: cada part no inicial té un mínim únic i canònic.}
\opcio{d}{l'elecció numerable.}
\end{itemize}
\feedback
\justificacio{El mínim d'una part no inicial de \(\N\) és únic i, per tant, una tria definida: no cal cap axioma per fer-la. En un conjunt infinit sense ordre, aquesta comoditat no existeix.}

\pregunta El conjunt de llistes finites \(\operatorname{List}(A)\) satisfà un principi de recursió amb\ldots
\begin{itemize}
\opcio{a}{només un cas base.}
\opcio{b}{la llista buida \(\varepsilon\) i l'operació \(\operatorname{cons}\).}
\opcio{c}{el successor de \(\N\).}
\opcio{d}{el parell nucli.}
\end{itemize}
\feedback
\justificacio{\(\operatorname{List}(A)\) té la llista buida i \(\operatorname{cons}\), i és l'àlgebra inicial per a \(\mathbf1+A\times(-)\): tota \((B,b,t)\) rep una única \(h\). És el principi que defineix llargada, concatenació i plegats.}

\pregunta La unió numerable d'enumerables és enumerable\ldots
\begin{itemize}
\opcio{a}{sempre, sense cap axioma.}
\opcio{b}{mai.}
\opcio{c}{sota l'axioma d'elecció numerable, que tria una enumeració per a cada \(A_n\).}
\opcio{d}{només si tots els \(A_n\) són finits.}
\end{itemize}
\feedback
\justificacio{Cal una enumeració per a cada \(A_n\) alhora; la definició d'enumerable només n'afirma l'existència, i l'elecció numerable les tria. Amb les \(e_n\) triades, un aparellament enumera la unió.}

\pregunta A la prova que infinit implica Dedekind, on s'usa l'elecció numerable?
\begin{itemize}
\opcio{a}{Per veure que per a cada \(n\) existeix alguna llista injectiva de llargada \(n+1\).}
\opcio{b}{Per triar alhora una llista injectiva \(f_n\) per a cada \(n\), és a dir, una secció de la família \(E\to\N\).}
\opcio{c}{Per prendre el mínim índex nou dins una llista.}
\opcio{d}{Per construir el desplaçament del successor sobre \(\N+C\).}
\end{itemize}
\feedback
\justificacio{L'existència fibra a fibra es prova per inducció sense tries; el mínim és canònic; el desplaçament és una definició per casos. L'únic pas no canònic és la tria simultània de les \(f_n\), que és exactament una secció d'una família indexada per \(\N\).}

\pregunta Per què \(\operatorname{List}(A)\) és enumerable sense elecció quan \(A\) ho és, mentre que la unió numerable d'enumerables necessita l'elecció?
\begin{itemize}
\opcio{a}{Perquè \(\operatorname{List}(A)\) és finit.}
\opcio{b}{Perquè les enumeracions de les llistes de cada llargada deriven uniformement d'una sola enumeració d'\(A\) i d'un aparellament, mentre que en una família arbitrària les enumeracions de les fibres s'han de triar.}
\opcio{c}{Perquè \(\operatorname{List}(A)\) no és una unió.}
\opcio{d}{No és cert: també necessita l'elecció.}
\end{itemize}
\feedback
\justificacio{\(\operatorname{List}(\N)\cong\N\) per una codificació explícita, i \(\operatorname{List}(e)\) és exhaustiva prenent els mínims índexs; tot és canònic. En una família \((A_n)\) qualsevol, cada \(A_n\) té enumeracions però cap de distingida: triar-ne una per a cada \(n\) és l'elecció numerable.}

\end{enumerate}
\clauestatica
\finalitzaTest

\chapter[Schröder--Bernstein i Cantor]{Comparació de conjunts,\texorpdfstring{\\}{ }Schröder--Bernstein i Cantor}\label{test:cap11}

Aquest test correspon al capítol~11 de la part teòrica. Cada pregunta té una única resposta correcta.

\iniciTest{u11}{b,c,b,a,c,b,a,c,b,c,a,b}

\begin{enumerate}

\pregunta Escrivim \(A\preceq B\) quan\ldots
\begin{itemize}
\opcio{a}{existeix una aplicació exhaustiva \(A\to B\).}
\opcio{b}{existeix un monomorfisme \(A\hookrightarrow B\).}
\opcio{c}{\(A\) és una part interna de \(B\times B\).}
\opcio{d}{\(A\) i \(B\) tenen el mateix conjunt de parts.}
\end{itemize}
\feedback
\justificacio{Dominar és injectar-se: \(A\preceq B\) vol dir que hi ha un monomorfisme \(A\hookrightarrow B\). Una exhaustiva \(A\to B\) diu que \(B\) és \enquote{més petit}, no més gran, i les relacions de mida no són parts internes.}

\pregunta La relació \(\prec\) es defineix com\ldots
\begin{itemize}
\opcio{a}{\(A\preceq B\) i \(B\preceq A\).}
\opcio{b}{qualsevol monomorfisme no exhaustiu.}
\opcio{c}{\(A\preceq B\) però \(A\not\approx B\).}
\opcio{d}{l'existència d'un epimorfisme regular \(B\to A\).}
\end{itemize}
\feedback
\justificacio{Estrictament menor és dominar sense ser equipotent. Un monomorfisme no exhaustiu d'un conjunt en si mateix és la infinitud de Dedekind, no la comparació estricta.}

\pregunta El teorema de Schröder--Bernstein afirma que\ldots
\begin{itemize}
\opcio{a}{tot conjunt s'injecta en el seu conjunt de parts.}
\opcio{b}{si \(A\preceq B\) i \(B\preceq A\), aleshores \(A\approx B\).}
\opcio{c}{dos conjunts qualssevol són comparables.}
\opcio{d}{cap aplicació \(A\to P(A)\) no és exhaustiva.}
\end{itemize}
\feedback
\justificacio{És l'antisimetria de \(\preceq\) llevat d'equipotència: dues injeccions oposades es combinen en una bijecció. La comparabilitat de dos conjunts qualssevol és un altre teorema, que demana l'elecció.}

\pregunta A la prova de Schröder--Bernstein, la part distingida \(D\) s'obté com\ldots
\begin{itemize}
\opcio{a}{el menor punt fix d'un operador creixent \(\Phi\) sobre \(P(A)\).}
\opcio{b}{una tria d'un element a cada fibra.}
\opcio{c}{la imatge d'una enumeració de \(A\).}
\opcio{d}{el parell nucli de \(f\).}
\end{itemize}
\feedback
\justificacio{\(\Phi(X)=A\setminus g(B\setminus f(X))\) és creixent, i \(D\) és la intersecció de les parts \(\Phi\)-tancades, construïda amb el quantificador universal: el menor punt fix, canònic.}

\pregunta La prova de Schröder--Bernstein que hem donat\ldots
\begin{itemize}
\opcio{a}{requereix l'axioma d'elecció.}
\opcio{b}{requereix l'elecció numerable.}
\opcio{c}{no usa cap elecció: l'ínfim és únic i canònic.}
\opcio{d}{només val per a conjunts finits.}
\end{itemize}
\feedback
\justificacio{L'ínfim d'una família de parts és únic; no es tria res. L'elecció només faria falta per comparar dos conjunts arbitraris, no per a l'antisimetria.}

\pregunta L'argument diagonal de Cantor considera la part\ldots
\begin{itemize}
\opcio{a}{\(\{a\in A : a\mathrel{\varepsilon} f(a)\}\).}
\opcio{b}{\(\{a\in A : a\not\mathrel{\varepsilon} f(a)\}\), que no és cap valor de \(f\).}
\opcio{c}{el singletó \(\{a\}\) de cada element.}
\opcio{d}{la imatge directa \(f(A)\).}
\end{itemize}
\feedback
\justificacio{La diagonal \(D=\{a:a\not\mathrel{\varepsilon}f(a)\}\) és una part d'\(A\), tallada amb la pertinença interna i el complement; si fos \(f(d)\), tindríem \(d\in D\iff d\notin D\).}

\pregunta El teorema~\ref{teo:diagonal} (\enquote{cap \(A\to P(A)\) és exhaustiva}) és\ldots
\begin{itemize}
\opcio{a}{estrictament més fort que \enquote{no hi ha isomorfisme \(A\cong P(A)\)}.}
\opcio{b}{equivalent a l'existència d'un conjunt universal.}
\opcio{c}{una conseqüència de l'elecció.}
\opcio{d}{fals per a \(A\) infinit.}
\end{itemize}
\feedback
\justificacio{No exhaustiva prohibeix qualsevol aplicació exhaustiva, injectiva o no; no isomorf només prohibeix les bijectives. La primera és la que es transporta, per exemple, a la no-enumerabilitat de \(\mathbf2^{\N}\).}

\pregunta El monomorfisme canònic \(A\hookrightarrow P(A)\) és\ldots
\begin{itemize}
\opcio{a}{la característica del complement.}
\opcio{b}{una aplicació exhaustiva.}
\opcio{c}{el singletó \(a\mapsto\{a\}\).}
\opcio{d}{la projecció \(A\times A\to A\).}
\end{itemize}
\feedback
\justificacio{El singletó \(a\mapsto\{a\}\) és la transposada de la característica de la diagonal, i és injectiu sobre punts: \(\{a\}=\{a'\}\) força \(a=a'\).}

\pregunta El teorema de Cantor afirma, per a tot conjunt \(A\),\ldots
\begin{itemize}
\opcio{a}{\(A\approx P(A)\).}
\opcio{b}{\(A\prec P(A)\): el conjunt de parts és estrictament més gran.}
\opcio{c}{\(P(A)\preceq A\).}
\opcio{d}{\(A\) és enumerable.}
\end{itemize}
\feedback
\justificacio{Els singletons donen \(A\preceq P(A)\) i la diagonal exclou tot isomorfisme: \(A\prec P(A)\). Per a \(A=\N\), \(P(\N)\) no és enumerable.}

\pregunta Que no existeixi cap conjunt universal se segueix de\ldots
\begin{itemize}
\opcio{a}{l'axioma d'elecció.}
\opcio{b}{la definició d'enumerable.}
\opcio{c}{Cantor i Schröder--Bernstein: \(U\preceq P(U)\preceq U\) donaria \(U\approx P(U)\).}
\opcio{d}{el principi del mínim.}
\end{itemize}
\feedback
\justificacio{Un \(U\) universal tindria \(P(U)\preceq U\), i amb \(U\preceq P(U)\) Schröder--Bernstein donaria \(U\approx P(U)\), contra Cantor. No cal cap elecció.}

\pregunta A la prova de Schröder--Bernstein, per què la bijectivitat de \(h=[f|_D,g^{-1}]\) es comprova sobre punts globals?
\begin{itemize}
\opcio{a}{Perquè sobre punts les alternatives \enquote{\(a\in D\)} o \enquote{\(a\notin D\)} són excloents, mentre que un element generalitzat pot quedar repartit; i per la separació global, mono i epi es decideixen sobre punts.}
\opcio{b}{Perquè només hi ha punts en conjunts finits.}
\opcio{c}{Per comoditat; sobre elements generalitzats l'argument és idèntic.}
\opcio{d}{Perquè l'elecció ho exigeix.}
\end{itemize}
\feedback
\justificacio{La distinció de casos de la prova usa el complement com a negació, cosa que només val sobre punts. Injectiva sobre punts és mono i exhaustiva sobre punts és epi gràcies a la separació global, i un bimorfisme és isomorfisme.}

\pregunta Com es construeix la intersecció de \emph{totes} les parts \(\Phi\)-tancades, que és un ínfim no finit?
\begin{itemize}
\opcio{a}{Iterant la intersecció binària.}
\opcio{b}{Com \(\forall_\pi(R_\varphi)\), el quantificador universal al llarg de \(A\times P(A)\to A\) aplicat a la part \enquote{si \(X\) és \(\Phi\)-tancada, aleshores \(a\mathrel{\varepsilon}X\)}.}
\opcio{c}{Amb l'axioma d'elecció.}
\opcio{d}{No es pot construir; cal postular-la.}
\end{itemize}
\feedback
\justificacio{Els axiomes donen ínfims binaris; els ínfims indexats per una part de \(P(A)\) s'obtenen amb el quantificador universal del capítol 7. És el mateix mecanisme que la inclusió interna.}

\end{enumerate}
\clauestatica
\finalitzaTest

\chapter{Famílies, fibres i operacions indexades}\label{test:cap12}

Aquest test correspon al capítol~12 de la part teòrica. Cada pregunta té una única resposta correcta.

\iniciTest{u12}{b,c,b,a,c,b,c,a,b,c,a,b}

\begin{enumerate}

\pregunta Una \emph{família de conjunts} indexada per \(I\) és\ldots
\begin{itemize}
\opcio{a}{un element de \(P(I)\).}
\opcio{b}{una aplicació \(p\colon E\to I\); les fibres són els conjunts indexats.}
\opcio{c}{un isomorfisme \(E\cong I\).}
\opcio{d}{una relació d'equivalència sobre \(I\).}
\end{itemize}
\feedback
\justificacio{La família és l'aplicació \(p\), i cada conjunt indexat és una fibra \(p^{-1}(i)\), una part del total: tot viu dins \(E\), sense cap col·lecció de conjunts.}

\pregunta La fibra de \(p\colon E\to I\) sobre un índex \(i\) és\ldots
\begin{itemize}
\opcio{a}{tot el total \(E\).}
\opcio{b}{el conjunt índex \(I\).}
\opcio{c}{el producte fibrat \(\mathbf1\times_I E\), és a dir \(p^{-1}(i)\).}
\opcio{d}{l'exponencial \(E^I\).}
\end{itemize}
\feedback
\justificacio{La fibra és el pullback del punt \(i\) i de \(p\), que és la imatge inversa \(p^{-1}(i)\), una part de \(E\) els punts de la qual són els \(e\) amb \(p(e)=i\).}

\pregunta Una \emph{secció} de \(p\colon E\to I\) és\ldots
\begin{itemize}
\opcio{a}{una part de \(I\).}
\opcio{b}{una aplicació \(s\colon I\to E\) amb \(p\circ s=\id_I\): una tria a cada fibra.}
\opcio{c}{una enumeració de \(E\).}
\opcio{d}{la projecció \(I\times A\to I\).}
\end{itemize}
\feedback
\justificacio{\(p\circ s=\id\) força \(p(s(i))=i\): \(s(i)\) és a la fibra sobre \(i\). Una secció és una funció d'elecció per a la família.}

\pregunta Per a \(p\) epimorfisme (fibres no buides), l'existència d'una secció\ldots
\begin{itemize}
\opcio{a}{és exactament l'axioma d'elecció.}
\opcio{b}{no requereix cap axioma.}
\opcio{c}{és falsa sempre.}
\opcio{d}{equival al teorema de Cantor.}
\end{itemize}
\feedback
\justificacio{Fibres no buides vol dir \(p\) epi; que tot epi tingui secció és la formulació estructural de l'elecció. Sense ella, una família de fibres no buides pot no tenir secció.}

\pregunta La \emph{suma} d'una família \(p\colon E\to I\)\ldots
\begin{itemize}
\opcio{a}{no existeix sense elecció.}
\opcio{b}{és l'exponencial \(E^I\).}
\opcio{c}{és el propi total \(E\), i existeix sempre.}
\opcio{d}{és el conjunt de seccions.}
\end{itemize}
\feedback
\justificacio{Recollir totes les fibres alhora és tenir el total: la suma és \(E\), amb \(p\) com a etiqueta de cada element. No cal res.}

\pregunta El \emph{producte} \(\prod_{i\in I}E_i=\{s\in E^I:p\circ s=\id_I\}\)\ldots
\begin{itemize}
\opcio{a}{no existeix com a objecte sense elecció.}
\opcio{b}{existeix sempre com a objecte, però la seva no-vacuïtat (fibres no buides) és l'elecció.}
\opcio{c}{és sempre buit.}
\opcio{d}{coincideix amb la suma.}
\end{itemize}
\feedback
\justificacio{El producte és un equalitzador dins \(E^I\), i existeix amb exponencials i parts; els seus punts són les seccions de \(p\), que per a fibres no buides existeixen exactament amb l'elecció.}

\pregunta Una família de parts d'un fix \(A\) es codifica com\ldots
\begin{itemize}
\opcio{a}{una relació d'equivalència sobre \(A\).}
\opcio{b}{un epimorfisme \(A\to I\).}
\opcio{c}{una aplicació \(\alpha\colon I\to P(A)\), amb pertinença total \(R\hookrightarrow I\times A\).}
\opcio{d}{una secció de \(A\times A\).}
\end{itemize}
\feedback
\justificacio{Cada \(\alpha(i)\) és una part d'\(A\), i la pertinença total \(R\) és la imatge inversa de \(\in_A\): la família viu dins \(P(A)^I\) i \(I\times A\), acotada.}

\pregunta La unió i la intersecció indexades \(\bigcup_i\alpha(i)\), \(\bigcap_i\alpha(i)\) es formen amb\ldots
\begin{itemize}
\opcio{a}{els quantificadors interns \(\exists_{\pi_A}\) i \(\forall_{\pi_A}\), sense cap axioma nou.}
\opcio{b}{l'axioma d'elecció.}
\opcio{c}{una successió d'alephs.}
\opcio{d}{la substitució estructural SR.}
\end{itemize}
\feedback
\justificacio{La unió és \(\exists_{\pi_A}(R)\) i la intersecció \(\forall_{\pi_A}(R)\): suprem i ínfim per l'adjunció dels quantificadors. Cap elecció ni cap principi nou.}

\pregunta La diferència entre una \emph{família interna} i una \emph{regla externa} és que\ldots
\begin{itemize}
\opcio{a}{només difereixen en la notació.}
\opcio{b}{la família ja és una aplicació \(E\to I\); la regla \(i\mapsto X_i\) potser encara no es representa per cap aplicació.}
\opcio{c}{la regla sempre existeix i la família mai.}
\opcio{d}{la regla requereix elecció i la família no.}
\end{itemize}
\feedback
\justificacio{Una família interna és un objecte \(E\to I\); una regla externa és una prescripció del metallenguatge que pot no tenir cap aplicació que la representi quan els valors no comparteixen ambient.}

\pregunta El principi de substitució estructural SR\ldots
\begin{itemize}
\opcio{a}{és un cas de l'axioma d'elecció.}
\opcio{b}{ja s'ha adoptat com a axioma en aquest capítol.}
\opcio{c}{garanteix representar una regla externa uniforme per una família interna; és de Nivell~II i aquí només s'enuncia.}
\opcio{d}{serveix per triar elements dins fibres.}
\end{itemize}
\feedback
\justificacio{SR garanteix l'existència de la família a partir d'una regla uniforme; és de Nivell II i el llibre només l'enuncia aquí. No és l'elecció, que tria dins fibres d'una família que ja existeix.}

\pregunta Una família \(\alpha\colon I\to P(A)\) de parts no buides, disjuntes dues a dues i que cobreixen \(A\)\ldots
\begin{itemize}
\opcio{a}{és una partició de \(A\): la projecció \(R\to A\) de la pertinença total és un isomorfisme, i \(q=r_I\circ r_A^{-1}\) és l'epimorfisme amb fibres \(\alpha(i)\).}
\opcio{b}{és una secció de \(P(A)\to I\).}
\opcio{c}{no determina cap aplicació sense l'elecció.}
\opcio{d}{és el producte de la família.}
\end{itemize}
\feedback
\justificacio{Disjunció i recobriment fan \(r_A\colon R\to A\) injectiva i exhaustiva sobre punts, doncs isomorfisme per separació global; la composició amb \(r_I\) és la partició, única amb aquestes fibres. És el recíproc de les fibres d'una partició del capítol 4.}

\pregunta Per què la regla \(n\mapsto\seg{n}\) es representa per una família interna sense cap principi nou?
\begin{itemize}
\opcio{a}{Perquè l'axioma d'elecció tria un segment per a cada \(n\).}
\opcio{b}{Perquè tots els segments són parts d'un mateix conjunt, \(\N\), i la família és la relació \(\operatorname{Lt}\hookrightarrow\N\times\N\) projectada sobre \(\N\).}
\opcio{c}{Perquè SR ho garanteix.}
\opcio{d}{No es representa: és una regla externa.}
\end{itemize}
\feedback
\justificacio{Quan els valors viuen en un ambient comú, la regla és una aplicació \(I\to P(A)\) i la família és la seva pertinença total. SR només cal quan no hi ha cap ambient comú.}

\end{enumerate}
\clauestatica
\finalitzaTest

\chapter{Ordres, bons ordres i segments}\label{test:cap13}

Aquest test correspon al capítol~13 de la part teòrica. Cada pregunta té una única resposta correcta.

\iniciTest{u13}{b,c,b,a,c,b,a,c,b,c,b,b}

\begin{enumerate}

\pregunta Un \emph{ordre estricte} sobre \(W\) és una part \({<}\hookrightarrow W\times W\) que és\ldots
\begin{itemize}
\opcio{a}{reflexiva i simètrica.}
\opcio{b}{irreflexiva i transitiva.}
\opcio{c}{una relació d'equivalència.}
\opcio{d}{una aplicació \(W\to W\).}
\end{itemize}
\feedback
\justificacio{Irreflexiva i transitiva: les dues condicions es tallen amb la diagonal i la composició de relacions. La simetria és pròpia de les equivalències, i un ordre no és una aplicació sinó una part de \(W\times W\).}

\pregunta Un ordre estricte és un \emph{bon ordre} quan\ldots
\begin{itemize}
\opcio{a}{és total.}
\opcio{b}{no té elements màxims.}
\opcio{c}{tota part no inicial de \(W\) té un element mínim.}
\opcio{d}{és isomorf a \(\N\).}
\end{itemize}
\feedback
\justificacio{És la definició: tota part no inicial té mínim. La totalitat n'és una conseqüència, no la definició, i hi ha bons ordres molt més llargs que \(\N\).}

\pregunta Tot bon ordre\ldots
\begin{itemize}
\opcio{a}{és finit.}
\opcio{b}{és total (i el mínim d'una part no inicial és únic).}
\opcio{c}{té un element màxim.}
\opcio{d}{admet cadenes descendents infinites.}
\end{itemize}
\feedback
\justificacio{El mínim de \(\{u,w\}\) decideix la comparació sobre punts, i el lema de transferència la converteix en una igualtat de parts; la unicitat ve de l'antisimetria. Un bon ordre pot ser infinit i no tenir màxim.}

\pregunta El segment inicial \(W_{<w}\) és\ldots
\begin{itemize}
\opcio{a}{la part \(\{u\in W : u<w\}\).}
\opcio{b}{tot \(W\) tret de \(w\).}
\opcio{c}{el conjunt \(\{u : u>w\}\).}
\opcio{d}{el mínim de \(W\).}
\end{itemize}
\feedback
\justificacio{El segment és la part dels elements estrictament anteriors a \(w\), imatge inversa de \({<}\) al llarg de \(\langle\id,w\circ{!}\rangle\). No és el complement del punt.}

\pregunta En un bon ordre, les parts inicials pròpies\ldots
\begin{itemize}
\opcio{a}{no existeixen.}
\opcio{b}{poden no tenir mínim.}
\opcio{c}{són exactament els segments \(W_{<w}\), amb \(w=\min(W\setminus S)\).}
\opcio{d}{coincideixen amb tot \(W\).}
\end{itemize}
\feedback
\justificacio{Si \(S\) és inicial i pròpia, el mínim \(w\) del complement compleix \(S=W_{<w}\); i aquest \(w\) és únic perquè un segment determina el seu punt.}

\pregunta La \emph{rigidesa} d'un bon ordre diu que, per a \(f\) creixent i injectiva,\ldots
\begin{itemize}
\opcio{a}{\(f\) és exhaustiva.}
\opcio{b}{cap punt no baixa: \(w\le f(w)\) per a tot \(w\).}
\opcio{c}{\(f\) té un punt fix.}
\opcio{d}{\(f(w)<w\) per a algun \(w\).}
\end{itemize}
\feedback
\justificacio{Si la part dels punts que baixen fos no buida, el seu mínim \(w_0\) donaria \(f(w_0)<w_0\) i, aplicant \(f\), un punt que baixa encara més petit. Per tant cap punt no baixa.}

\pregunta Entre dos bons ordres, el nombre d'isomorfismes d'ordre és\ldots
\begin{itemize}
\opcio{a}{com a molt un.}
\opcio{b}{sempre exactament un.}
\opcio{c}{infinit.}
\opcio{d}{depèn de l'elecció.}
\end{itemize}
\feedback
\justificacio{Si \(f,g\) són isomorfismes d'ordre, \(g^{-1}\circ f\) i la seva inversa són automorfismes i cap punt no baixa en cap dels dos: la composició és la identitat i \(f=g\). Pot no haver-n'hi cap.}

\pregunta Cap bon ordre\ldots
\begin{itemize}
\opcio{a}{és total.}
\opcio{b}{té element mínim.}
\opcio{c}{és isomorf a cap segment inicial propi seu.}
\opcio{d}{admet aplicacions creixents.}
\end{itemize}
\feedback
\justificacio{Un tal isomorfisme donaria una creixent injectiva \(W\to W\) amb \(f(w)<w\), contra la rigidesa. És el que fa consistent comparar bons ordres per segments.}

\pregunta La comparació \enquote{ser un segment de} sobre bons ordres és\ldots
\begin{itemize}
\opcio{a}{només reflexiva.}
\opcio{b}{antisimètrica i transitiva (llevat d'isomorfisme d'ordre).}
\opcio{c}{simètrica.}
\opcio{d}{una relació d'equivalència.}
\end{itemize}
\feedback
\justificacio{Dues dominacions oposades componen un isomorfisme sobre un segment inicial de si mateix, que no pot ser propi; per tant és un automorfisme, la identitat. És un ordre llevat d'isomorfisme.}

\pregunta La comparabilitat total de dos bons ordres (que almenys una situació de la tricotomia es doni)\ldots
\begin{itemize}
\opcio{a}{és falsa en general.}
\opcio{b}{requereix l'axioma d'elecció.}
\opcio{c}{s'obté per recursió transfinita, sense elecció (el pas al mínim és canònic).}
\opcio{d}{és equivalent al teorema de Cantor.}
\end{itemize}
\feedback
\justificacio{L'aparellament es construeix per recursió prenent a cada pas el mínim encara lliure: una tria canònica, no una elecció. La comparabilitat de conjunts \emph{nus} sí que demana l'elecció.}

\pregunta La totalitat d'un ordre, \enquote{\(u=w\) o \(u<w\) o \(w<u\)}, s'enuncia per a punts globals. Què diu com a afirmació sobre parts?
\begin{itemize}
\opcio{a}{Que \({<}\) és tot \(W\times W\).}
\opcio{b}{Que \(\Delta\cup{<}\cup\operatorname{Inv}(<)\) és la part total de \(W\times W\); un element generalitzat pot quedar repartit entre els tres membres.}
\opcio{c}{Que \({<}\) és una relació d'equivalència.}
\opcio{d}{Que tot element generalitzat cau sencer en un dels tres membres.}
\end{itemize}
\feedback
\justificacio{Les condicions disjuntives són unions de parts, i la separació global permet comprovar-les sobre punts; però sobre un estadi general \enquote{o} no exclou, i (d) és fals. És la mateixa situació que a l'ordre de \(\N\).}

\pregunta Quina d'aquestes propietats \emph{no} val en un ordre total qualsevol, sinó només amb el bon ordre?
\begin{itemize}
\opcio{a}{Que tota aplicació creixent sigui injectiva.}
\opcio{b}{Que cap punt baixi sota una creixent injectiva (rigidesa) i, per tant, que hi hagi com a molt un isomorfisme d'ordre.}
\opcio{c}{Que l'ordre sigui transitiu.}
\opcio{d}{Que els segments inicials siguin parts.}
\end{itemize}
\feedback
\justificacio{La injectivitat de les creixents i la transitivitat valen en qualsevol ordre total; la rigidesa necessita el mínim d'una part (a \(\Z\) la translació \(n\mapsto n+1\) és creixent i injectiva i fa baixar tots els punts, en el sentit que cap punt no queda fix).}

\end{enumerate}
\clauestatica
\finalitzaTest

\chapter{Inducció i recursió transfinites}\label{test:cap14}

Aquest test correspon al capítol~14 de la part teòrica. Cada pregunta té una única resposta correcta.

\iniciTest{u14}{b,c,a,b,c,b,a,c,b,c,a,c}

\begin{enumerate}

\pregunta El principi d'inducció transfinita diu que si una part \(S\hookrightarrow W\) compleix\ldots
\begin{itemize}
\opcio{a}{\(0\in S\) i \(S\) tancada per successor, aleshores \(S=W\).}
\opcio{b}{\(W_{<w}\subseteq S\Rightarrow w\in S\) per a tot \(w\), aleshores \(S=W\).}
\opcio{c}{\(S\) no buida, aleshores \(S=W\).}
\opcio{d}{\(S\) és inicial, aleshores \(S=W\).}
\end{itemize}
\feedback
\justificacio{La progressivitat mira tot el segment inicial, no només el predecessor: no cal cas base, i cobreix els punts límit. La versió (a) és la inducció de \(\N\).}

\pregunta La prova de la inducció transfinita usa\ldots
\begin{itemize}
\opcio{a}{l'axioma d'elecció.}
\opcio{b}{una successió d'aproximacions.}
\opcio{c}{el mínim del complement \(W\setminus S\), si fos no buit.}
\opcio{d}{el teorema de Cantor.}
\end{itemize}
\feedback
\justificacio{Si \(S\neq W\), el complement té mínim \(w_0\); tot el segment \(W_{<w_0}\) és dins \(S\), i la progressivitat força \(w_0\in S\), contradicció. Cap elecció.}

\pregunta A diferència de la inducció sobre \(\N\), la transfinita\ldots
\begin{itemize}
\opcio{a}{no necessita cas base separat: el mínim hi és inclòs buidament.}
\opcio{b}{només val per a conjunts finits.}
\opcio{c}{requereix separar successor i límit sempre.}
\opcio{d}{no val per a \(\N\).}
\end{itemize}
\feedback
\justificacio{Per al mínim \(\bot\), \(W_{<\bot}\) és buit i la hipòtesi es compleix buidament, de manera que \(\bot\in S\). La separació en successor i límit és opcional.}

\pregunta En la recursió amb codomini fix, una \emph{història} en \(w\) és\ldots
\begin{itemize}
\opcio{a}{un element de \(X\).}
\opcio{b}{una aplicació \(h\colon W_{<w}\to X\) (els valors sobre el segment inicial).}
\opcio{c}{una secció d'una família.}
\opcio{d}{un bon ordre sobre \(X\).}
\end{itemize}
\feedback
\justificacio{La història és la restricció de \(F\) al segment: una aplicació \(W_{<w}\to X\), codificada com a aplicació parcial dins \((X+\mathbf1)^W\).}

\pregunta El teorema de recursió transfinita construeix \(F\colon W\to X\) com\ldots
\begin{itemize}
\opcio{a}{una tria simultània a cada fibra.}
\opcio{b}{un producte de la família.}
\opcio{c}{la reunió de les solucions parcials compatibles, dins l'exponencial fix \((X+\mathbf1)^W\).}
\opcio{d}{el quocient per una relació d'equivalència.}
\end{itemize}
\feedback
\justificacio{Les solucions parcials són compatibles (inducció), n'hi ha definides fins a cada punt (inducció) i la seva reunió és funcional i total: aquesta és \(F\). Tot dins un exponencial fix.}

\pregunta La recursió transfinita amb codomini fix\ldots
\begin{itemize}
\opcio{a}{requereix l'axioma d'elecció.}
\opcio{b}{no usa cap elecció: cada valor és \(\Phi\) de la història, i la reunió és canònica.}
\opcio{c}{només val sobre \(\N\).}
\opcio{d}{requereix la substitució estructural SR.}
\end{itemize}
\feedback
\justificacio{Cada valor està determinat per la regla, i la reunió de solucions compatibles és canònica: no hi ha cap tria. SR només caldria si el codomini no fos fix.}

\pregunta La comparabilitat total de dos bons ordres es demostra\ldots
\begin{itemize}
\opcio{a}{per recursió, aparellant-los pel mínim disponible, sense elecció.}
\opcio{b}{amb l'axioma d'elecció.}
\opcio{c}{amb el teorema de Cantor.}
\opcio{d}{només per a bons ordres finits.}
\end{itemize}
\feedback
\justificacio{\(F(w)=\min(V\setminus F(W_{<w}))\), amb un punt \(\top\) quan \(V\) s'esgota; és creixent amb imatge inicial, i els tres casos surten d'on s'atura.}

\pregunta La comparabilitat de dos bons ordres, comparada amb la de dos conjunts qualssevol,\ldots
\begin{itemize}
\opcio{a}{és més difícil i requereix més elecció.}
\opcio{b}{és equivalent.}
\opcio{c}{no usa elecció (el mínim és canònic), mentre que la de conjunts nus sí.}
\opcio{d}{és falsa.}
\end{itemize}
\feedback
\justificacio{Un bon ordre porta incorporada la tria canònica a cada pas: el mínim. En dos conjunts sense ordre no hi ha res que distingeixi cap element, i la comparabilitat demana l'elecció.}

\pregunta Una successió transfinita de parts d'un conjunt fix \(A\) (codomini \(P(A)\))\ldots
\begin{itemize}
\opcio{a}{necessita SR.}
\opcio{b}{és de Nivell~I: el codomini és fix, n'hi ha prou amb la recursió d'aquest capítol.}
\opcio{c}{requereix elecció.}
\opcio{d}{no es pot definir.}
\end{itemize}
\feedback
\justificacio{Totes les parts viuen dins el fix \(P(A)\): el codomini és \(X=P(A)\) i la recursió d'aquest capítol s'hi aplica directament.}

\pregunta La successió \(V_0=A\), \(V_{n+1}=P(V_n)\)\ldots
\begin{itemize}
\opcio{a}{es defineix amb la recursió de codomini fix.}
\opcio{b}{és de codomini fix perquè totes són iguals.}
\opcio{c}{necessita SR: els valors creixen (Cantor) i cap axioma de Nivell~I no en dona un codomini comú.}
\opcio{d}{requereix l'elecció numerable.}
\end{itemize}
\feedback
\justificacio{Per Cantor, \(V_{n+1}\) és estrictament més gran que \(V_n\), i cap axioma de Nivell~I no dona un conjunt que els contingui tots uniformement: reunir-los en una família és SR.}

\pregunta Per què la regla de recursió \(\Phi(w,h)\) s'ha de donar per a \emph{tota} parella \((w,h)\), i no només per a les històries \enquote{bones}?
\begin{itemize}
\opcio{a}{Perquè \(\Phi\) és una aplicació interna i ha d'estar definida a tot el seu domini; si la regla només cobreix un cas, cal afegir-hi una clàusula \enquote{altrament}, encara que no s'arribi a usar.}
\opcio{b}{Perquè altrament caldria l'elecció.}
\opcio{c}{No cal: n'hi ha prou amb definir-la on interessa.}
\opcio{d}{Perquè el bon ordre pot tenir punts límit.}
\end{itemize}
\feedback
\justificacio{\(\Phi\) és una aplicació d'una part de \(W\times(X+\mathbf1)^W\) cap a \(X\), i una aplicació parcial no serveix. És un detall que es repeteix a la prova del lema de Zorn.}

\pregunta On es fa servir el bon ordre a la prova de la recursió transfinita?
\begin{itemize}
\opcio{a}{Enlloc: val per a qualsevol ordre.}
\opcio{b}{Per triar un element a cada pas.}
\opcio{c}{A les dues induccions (compatibilitat i existència fins a cada punt), que depenen del mínim del complement.}
\opcio{d}{Per construir l'exponencial \((X+\mathbf1)^W\).}
\end{itemize}
\feedback
\justificacio{La inducció transfinita és el que el bon ordre aporta; en un ordre no ben fundat el \enquote{mínim del complement} podria no existir i la prova s'ensorraria. L'exponencial no depèn de l'ordre.}

\end{enumerate}
\clauestatica
\finalitzaTest

\chapter{Zorn, Hartogs i el teorema del bon ordre}\label{test:cap15}

Aquest test correspon al capítol~15 de la part teòrica. Cada pregunta té una única resposta correcta.

\iniciTest{u15}{c,b,a,c,b,a,c,b,a,c,b,b}

\begin{enumerate}

\pregunta Els bons ordres sobre parts d'un conjunt \(A\) es codifiquen com\ldots
\begin{itemize}
\opcio{a}{elements de \(A\).}
\opcio{b}{aplicacions \(A\to A\).}
\opcio{c}{parells \((S,R)\) amb \(S\hookrightarrow A\) i \(R\hookrightarrow S\times S\) un bon ordre sobre \(S\): una part de \(P(A)\times P(A\times A)\).}
\opcio{d}{seccions d'una família.}
\end{itemize}
\feedback
\justificacio{El suport ha de ser una dada: la relació buida no distingeix el bon ordre sobre un singletó del bon ordre sobre el buit. Per això \((S,R)\in P(A)\times P(A\times A)\).}

\pregunta El conjunt de Hartogs \(\Hart(A)\) és\ldots
\begin{itemize}
\opcio{a}{el conjunt de parts de \(A\).}
\opcio{b}{el conjunt de tipus d'ordre dels bons ordres sobre parts d'\(A\), ordenat per segment inicial.}
\opcio{c}{una classe global d'ordinals.}
\opcio{d}{el quocient de \(A\) per una relació.}
\end{itemize}
\feedback
\justificacio{\(\WO(A)\) es quocienta per isomorfisme d'ordre (efectivitat de les relacions d'equivalència) i s'ordena per \enquote{ser segment inicial propi de}. És un conjunt, no una classe.}

\pregunta El teorema de Hartogs afirma que\ldots
\begin{itemize}
\opcio{a}{\(\Hart(A)\) no admet cap monomorfisme en \(A\).}
\opcio{b}{\(\Hart(A)\approx A\).}
\opcio{c}{\(\Hart(A)\) és finit.}
\opcio{d}{\(\Hart(A)\hookrightarrow A\) sempre.}
\end{itemize}
\feedback
\justificacio{Si s'hi injectés, el seu tipus seria un element de \(\Hart(A)\) i \(\Hart(A)\) resultaria isomorf a un segment propi de si mateix, contra la rigidesa.}

\pregunta El teorema de Hartogs\ldots
\begin{itemize}
\opcio{a}{requereix l'axioma d'elecció.}
\opcio{b}{és fals per a \(A\) infinit.}
\opcio{c}{no usa elecció: és l'anàleg ordinal del teorema de Cantor.}
\opcio{d}{equival al lema de Zorn.}
\end{itemize}
\feedback
\justificacio{La construcció talla, quocienta i ordena dins conjunts ja fets, i la contradicció és diagonal; no hi ha cap tria. És l'anàleg, per a longituds, del que Cantor és per a mides.}

\pregunta El teorema del bon ordre diu que, amb l'elecció,\ldots
\begin{itemize}
\opcio{a}{tot conjunt és finit.}
\opcio{b}{tot conjunt admet un bon ordre.}
\opcio{c}{tot conjunt és numerable.}
\opcio{d}{tot ordre parcial té màxim.}
\end{itemize}
\feedback
\justificacio{Amb una funció d'elecció sobre les parts no buides i una recursió acotada per \(\Hart(A)\) s'exhaureix \(A\), i el bon ordre del segment es transporta. No diu res sobre finitud ni numerabilitat.}

\pregunta A la prova del teorema del bon ordre, què atura la recursió?
\begin{itemize}
\opcio{a}{La cota de Hartogs: \(\Hart(A)\not\preceq A\), i el procés esgota \(A\) abans.}
\opcio{b}{L'axioma d'exponencials.}
\opcio{c}{El teorema de Cantor.}
\opcio{d}{Res; la recursió no s'atura.}
\end{itemize}
\feedback
\justificacio{Si la recursió no s'aturés, donaria un monomorfisme \(\Hart(A)\hookrightarrow A\), prohibit per Hartogs. Per això Hartogs precedeix el teorema del bon ordre.}

\pregunta El lema de Zorn dona\ldots
\begin{itemize}
\opcio{a}{un màxim de l'ordre parcial.}
\opcio{b}{un bon ordre.}
\opcio{c}{un element \emph{maximal} (no necessàriament màxim), si tota cadena té cota superior.}
\opcio{d}{una funció d'elecció canònica.}
\end{itemize}
\feedback
\justificacio{Zorn dona un element que cap altre no supera, suposant que tota cadena té cota. Un màxim és més fort i pot no existir.}

\pregunta \enquote{Maximal} i \enquote{màxim}\ldots
\begin{itemize}
\opcio{a}{són el mateix sempre.}
\opcio{b}{difereixen: un maximal no és superat per cap altre; un màxim supera tots els altres. En un ordre parcial poden no coincidir.}
\opcio{c}{només tenen sentit en bons ordres.}
\opcio{d}{no tenen relació.}
\end{itemize}
\feedback
\justificacio{En un ordre parcial pot haver-hi diversos maximals i cap màxim; en un ordre total coincideixen. Afirmar que Zorn dona un màxim és un error habitual.}

\pregunta Una aplicació genuïna del lema de Zorn és\ldots
\begin{itemize}
\opcio{a}{que tot espai vectorial té una base.}
\opcio{b}{el teorema de Cantor.}
\opcio{c}{la construcció de \(\N\).}
\opcio{d}{el principi d'inducció.}
\end{itemize}
\feedback
\justificacio{Les parts independents ordenades per inclusió tenen cotes de cadena (la unió, perquè tota combinació lineal és finita), i un maximal és una base. És l'esquema típic d'aplicació de Zorn.}

\pregunta Sobre els altres axiomes del llibre,\ldots
\begin{itemize}
\opcio{a}{l'elecció és més forta que el teorema del bon ordre.}
\opcio{b}{el lema de Zorn és més feble que l'elecció.}
\opcio{c}{elecció, teorema del bon ordre i lema de Zorn són equivalents.}
\opcio{d}{cap dels tres no implica els altres.}
\end{itemize}
\feedback
\justificacio{(E)\(\Rightarrow\)(BO) i (E)\(\Rightarrow\)(Z) són els teoremes del capítol; (BO)\(\Rightarrow\)(E) pren el mínim de cada fibra i (Z)\(\Rightarrow\)(BO) maximalitza els bons ordres parcials.}

\pregunta A la prova del lema de Zorn, què es tria amb l'elecció?
\begin{itemize}
\opcio{a}{Un element de cada cadena.}
\opcio{b}{Una \emph{cota superior estricta} de cada cadena, aplicant l'elecció a la família que sobre cada cadena té la part de les seves cotes estrictes.}
\opcio{c}{Un bon ordre sobre \(P\).}
\opcio{d}{El maximal mateix.}
\end{itemize}
\feedback
\justificacio{Suposant que no hi ha maximal, cada cadena té cotes estrictes i la família corresponent té les fibres no buides; la secció dona \(c\). La recursió que en resulta donaria \(\Hart(P)\preceq P\), contra Hartogs.}

\pregunta Per què la unió d'una cadena de bons ordres parcials \((S,R)\), ordenats per extensió \emph{inicial}, torna a ser un bon ordre?
\begin{itemize}
\opcio{a}{Perquè tota unió d'ordres totals és un bon ordre.}
\opcio{b}{Perquè una part no buida \(T\) talla algun membre \((S,R)\), i el mínim de \(T\cap S\) ho és de tot \(T\): els punts afegits després són posteriors a tot \(S\).}
\opcio{c}{Per l'axioma d'elecció.}
\opcio{d}{Perquè la cadena és finita.}
\end{itemize}
\feedback
\justificacio{La inicialitat de les extensions és essencial: sense ella la unió pot no ser ben ordenada (una cadena d'ordres que afegeixen punts \emph{per sota} donaria una còpia de \(\Z\) negatiu).}

\end{enumerate}
\clauestatica
\finalitzaTest

\chapter{Tipus ordinals}\label{test:cap16}

Aquest test correspon al capítol~16 de la part teòrica. Cada pregunta té una única resposta correcta.

\iniciTest{u16}{b,c,b,a,c,b,c,a,b,c,c,b}

\begin{enumerate}

\pregunta Un \emph{tipus ordinal} és\ldots
\begin{itemize}
\opcio{a}{un conjunt transitiu.}
\opcio{b}{un conjunt ben ordenat considerat llevat d'isomorfisme d'ordre.}
\opcio{c}{una part de \(\N\).}
\opcio{d}{una relació d'equivalència.}
\end{itemize}
\feedback
\justificacio{Un tipus és la classe d'isomorfisme d'ordre d'un bon ordre; no es tria cap representant. Els conjunts transitius són la presentació de von Neumann, que aquest llibre no adopta.}

\pregunta L'ordre entre tipus, \(\alpha<\beta\), es defineix com\ldots
\begin{itemize}
\opcio{a}{\(\alpha\in\beta\).}
\opcio{b}{existeix una aplicació \(\alpha\to\beta\).}
\opcio{c}{un representant de \(\alpha\) és isomorf a un segment inicial propi d'un de \(\beta\).}
\opcio{d}{\(\alpha\) i \(\beta\) són equipotents.}
\end{itemize}
\feedback
\justificacio{És la comparació per segments del capítol 13, llegida sobre tipus. \(\alpha\in\beta\) no té sentit sense pertinença global, i l'equipotència compara mides, no longituds.}

\pregunta Que l'ordre \(\alpha<\beta\) no depengui dels representants es deu a\ldots
\begin{itemize}
\opcio{a}{l'axioma d'elecció.}
\opcio{b}{la unicitat de l'isomorfisme d'ordre entre bons ordres.}
\opcio{c}{l'axioma d'exponencials.}
\opcio{d}{el teorema de Cantor.}
\end{itemize}
\feedback
\justificacio{Si \(W\cong W'\) i \(V\cong V'\), els isomorfismes únics transporten segments propis en segments propis. Cap elecció no hi intervé.}

\pregunta El successor \(\alpha+1\) és el tipus de\ldots
\begin{itemize}
\opcio{a}{\(W+\mathbf1\): \(W\) amb un punt nou col·locat després de tot \(W\).}
\opcio{b}{\(P(W)\).}
\opcio{c}{\(W\times W\).}
\opcio{d}{\(\mathbf1+W\): un punt nou al principi.}
\end{itemize}
\feedback
\justificacio{El punt nou va al final, de manera que és màxim i el seu segment és \(W\). Posar-lo al principi (\(\mathbf1+W\)) dona un altre tipus en general, i per a \(W=\N\) el mateix tipus que \(W\).}

\pregunta El successor \(\alpha+1\)\ldots
\begin{itemize}
\opcio{a}{és límit.}
\opcio{b}{no té element màxim.}
\opcio{c}{és el menor tipus estrictament més gran que \(\alpha\), i té màxim.}
\opcio{d}{és igual a \(\alpha\).}
\end{itemize}
\feedback
\justificacio{\(W\) és el segment del punt nou, d'on \(\alpha<\alpha+1\); un tipus entremig seria un segment propi de \(W+\mathbf1\), és a dir \(W\) o un segment de \(W\), d'on \(\le\alpha\).}

\pregunta Un tipus és \emph{límit} quan\ldots
\begin{itemize}
\opcio{a}{és zero.}
\opcio{b}{no és zero ni successor: els seus representants són no buits i no tenen màxim.}
\opcio{c}{té element màxim.}
\opcio{d}{és finit.}
\end{itemize}
\feedback
\justificacio{Zero és el buit, successor és tenir màxim; límit és la resta. Tenir màxim és invariant per isomorfisme d'ordre, de manera que la classificació no depèn del representant.}

\pregunta El primer tipus límit és\ldots
\begin{itemize}
\opcio{a}{\(0\).}
\opcio{b}{\(1\).}
\opcio{c}{\(\omega=\operatorname{tipus}(\N)\).}
\opcio{d}{no n'hi ha.}
\end{itemize}
\feedback
\justificacio{Els tipus finits són zero o successors; \(\N\) no té màxim i no és buit, i tot tipus menor és el d'un segment \(\seg{n}\), finit.}

\pregunta El suprem d'una família de tipus \emph{representada dins un bon ordre \(W\)}\ldots
\begin{itemize}
\opcio{a}{és el tipus de la reunió \(\bigcup_iW_{<r(i)}\) dels segments oberts que representen els \(\alpha_i\).}
\opcio{b}{no existeix mai.}
\opcio{c}{requereix l'elecció.}
\opcio{d}{és el conjunt de tots els ordinals.}
\end{itemize}
\feedback
\justificacio{La reunió dels segments oberts és inicial; cada \(\alpha_i\) n'és una part inicial i qualsevol cota comuna rep la reunió, enganxant els isomorfismes únics. Amb segments tancats el resultat seria massa gran.}

\pregunta El teorema de Burali--Forti afirma que\ldots
\begin{itemize}
\opcio{a}{tot tipus té successor.}
\opcio{b}{no existeix cap bon ordre \(\Omega\) que indexi coherentment tots els tipus, és a dir, en què cada \(\gamma\) representi el tipus del segment \(\Omega_{<\gamma}\).}
\opcio{c}{\(\omega\) és el darrer tipus.}
\opcio{d}{tot bon ordre és finit.}
\end{itemize}
\feedback
\justificacio{Si \(\Omega\) existís, tindria un tipus representat per algun \(\gamma\in\Omega\), i \(\Omega\) seria isomorf al seu segment propi \(\Omega_{<\gamma}\), contra la rigidesa.}

\pregunta Burali--Forti és\ldots
\begin{itemize}
\opcio{a}{una conseqüència de l'elecció.}
\opcio{b}{independent de Cantor.}
\opcio{c}{el paral·lel ordinal del teorema de Cantor (no hi ha longitud d'ordre màxima).}
\opcio{d}{fals sense SR.}
\end{itemize}
\feedback
\justificacio{Com Cantor amb les mides, Burali--Forti diu que la col·lecció de totes les longituds és massa gran per ser un conjunt del mateix tipus; cap dels dos no usa l'elecció ni SR.}

\pregunta Amb \(\alpha_i=\operatorname{tipus}(W_{<r(i)})\), per què el suprem es pren amb els segments \emph{oberts} \(W_{<r(i)}\) i no amb els tancats \(W_{\le r(i)}\)?
\begin{itemize}
\opcio{a}{És indiferent: donen el mateix tipus.}
\opcio{b}{Perquè els tancats no són parts inicials.}
\opcio{c}{Perquè els tancats donen en general un tipus massa gran: amb un sol índex i \(r\) el mínim de \(W\), \(\alpha=0\) però \(W_{\le r}\) té tipus \(1\).}
\opcio{d}{Perquè els oberts necessiten l'elecció.}
\end{itemize}
\feedback
\justificacio{Els segments que representen els \(\alpha_i\) són els oberts, i la seva reunió és exactament el suprem; amb els tancats s'hi afegeixen els punts \(r(i)\), que poden sobrepassar-lo.}

\pregunta \(\mathbf1+\N\) (un punt nou \emph{abans} de \(\N\)) i \(\N+\mathbf1\) (un punt nou \emph{després}), com a bons ordres\ldots
\begin{itemize}
\opcio{a}{són tots dos isomorfs a \(\N\).}
\opcio{b}{el primer és isomorf a \(\N\) (via el successor) i el segon no (té màxim): la suma de tipus no és commutativa.}
\opcio{c}{cap dels dos no és isomorf a \(\N\).}
\opcio{d}{són isomorfs entre si.}
\end{itemize}
\feedback
\justificacio{\(\mathbf1+\N\cong\N\) enviant el punt nou a \(0\) i \(n\) a \(s(n)\); \(\N+\mathbf1\) té màxim, que és invariant, i per tant té tipus \(\omega+1\neq\omega\).}

\end{enumerate}
\clauestatica
\finalitzaTest

\chapter{Aritmètica ordinal}\label{test:cap17}

Aquest test correspon al capítol~17 de la part teòrica. Cada pregunta té una única resposta correcta.

\iniciTest{u17}{b,c,a,b,c,a,c,b,a,c,b,a}

\begin{enumerate}

\pregunta La suma ordinal \(\alpha+\beta\) és el tipus de\ldots
\begin{itemize}
\opcio{a}{\(A\times B\) amb l'ordre lexicogràfic.}
\opcio{b}{\(A+B\) amb tot \(A\) abans de tot \(B\).}
\opcio{c}{\(P(A+B)\).}
\opcio{d}{\(A\cup B\) sense ordre.}
\end{itemize}
\feedback
\justificacio{Concatenar és posar tot \(A\) abans de tot \(B\) dins la suma disjunta, mantenint cada ordre. L'ordre lexicogràfic sobre el producte dona el producte ordinal, no la suma.}

\pregunta La suma ordinal\ldots
\begin{itemize}
\opcio{a}{és commutativa.}
\opcio{b}{no és associativa.}
\opcio{c}{és associativa però no commutativa: \(1+\omega=\omega\neq\omega+1\).}
\opcio{d}{no és un bon ordre.}
\end{itemize}
\feedback
\justificacio{Els dos costats de l'associativitat ordenen \(A+B+C\) en el mateix ordre; però afegir un punt al davant de \(\N\) s'absorbeix i al darrere no.}

\pregunta \(1+\omega\) i \(\omega+1\)\ldots
\begin{itemize}
\opcio{a}{compleixen \(1+\omega=\omega\) i \(\omega+1>\omega\) (aquest té element màxim).}
\opcio{b}{són iguals.}
\opcio{c}{tots dos són successors de \(\omega\).}
\opcio{d}{no són numerables.}
\end{itemize}
\feedback
\justificacio{\(1+\N\cong\N\) desplaçant la numeració; \(\N+1\) té màxim, propietat invariant, de manera que \(\omega+1\neq\omega\) i és el seu successor.}

\pregunta El producte ordinal \(\alpha\cdot\beta\) és el tipus de \(A\times B\) ordenat\ldots
\begin{itemize}
\opcio{a}{només per la coordenada d'\(A\).}
\opcio{b}{primer per la coordenada de \(B\), després per la d'\(A\): \(\beta\) còpies d'\(\alpha\).}
\opcio{c}{sense ordre.}
\opcio{d}{per inclusió.}
\end{itemize}
\feedback
\justificacio{Mana la coordenada de \(B\): per a cada punt de \(B\) hi ha una còpia sencera d'\(A\). Invertir la convenció donaria \(\beta\cdot\alpha\).}

\pregunta \(2\cdot\omega\) i \(\omega\cdot2\)\ldots
\begin{itemize}
\opcio{a}{són iguals.}
\opcio{b}{no estan definits.}
\opcio{c}{compleixen \(2\cdot\omega=\omega\) però \(\omega\cdot2=\omega+\omega\neq\omega\).}
\opcio{d}{són tots dos finits.}
\end{itemize}
\feedback
\justificacio{\(\omega\) parelles una rere l'altra són \(\omega\); dues còpies de \(\omega\) tenen un punt no mínim sense predecessor immediat, que \(\N\) no té.}

\pregunta La distributivitat del producte ordinal\ldots
\begin{itemize}
\opcio{a}{val per l'esquerra: \(\alpha\cdot(\beta+\gamma)=\alpha\cdot\beta+\alpha\cdot\gamma\); per la dreta falla.}
\opcio{b}{val per tots dos costats.}
\opcio{c}{no val mai.}
\opcio{d}{val només per la dreta.}
\end{itemize}
\feedback
\justificacio{\(\beta+\gamma\) còpies d'\(\alpha\) són \(\beta\) còpies seguides de \(\gamma\) còpies; per la dreta, \((1+1)\cdot\omega=\omega\) mentre que \(\omega+\omega\neq\omega\).}

\pregunta Per a una base \(\alpha>0\), l'exponenciació ordinal satisfà les equacions\ldots
\begin{itemize}
\opcio{a}{\(\alpha^\lambda=\alpha\) als límits.}
\opcio{b}{\(\alpha^{\beta+1}=\alpha+\alpha^\beta\).}
\opcio{c}{\(\alpha^0=1\), \(\alpha^{\beta+1}=\alpha^\beta\cdot\alpha\), \(\alpha^\lambda=\sup_{\beta<\lambda}\alpha^\beta\).}
\opcio{d}{una única clàusula.}
\end{itemize}
\feedback
\justificacio{Amb la construcció per funcions de suport finit, aquestes tres equacions són un teorema per a \(\alpha>0\): unitat, multiplicar per \(\alpha\) al successor, suprem al límit.}

\pregunta La clàusula límit de l'exponenciació\ldots
\begin{itemize}
\opcio{a}{no requereix res especial.}
\opcio{b}{demana una família \((\alpha^\beta)_{\beta<\lambda}\) representada dins un bon ordre (la construcció per suport finit la proporciona), i no s'aplica a la base \(0\).}
\opcio{c}{requereix l'axioma d'elecció.}
\opcio{d}{només val per a \(\alpha\) finit.}
\end{itemize}
\feedback
\justificacio{Per a \(\alpha>0\), totes les potències fins a \(\lambda\) són parts inicials d'un sol bon ordre de funcions de suport finit, que fa de cota; per a la base \(0\) la clàusula no val. Sense un ambient comú per a bases o exponents variables caldria SR.}

\pregunta \(2^\omega\) com a tipus ordinal\ldots
\begin{itemize}
\opcio{a}{és \(\omega\) (numerable), diferent del cardinal \(2^{\aleph_0}\) (el continu).}
\opcio{b}{és el continu.}
\opcio{c}{no està definit.}
\opcio{d}{és finit.}
\end{itemize}
\feedback
\justificacio{\(2^\omega=\sup_n2^n\) i tots els \(2^n\) són finits: el suprem és \(\omega\). El cardinal \(2^{\aleph_0}\) és el del conjunt de parts de \(\N\), no numerable per Cantor.}

\pregunta La diferència clau entre aritmètica ordinal i cardinal és que\ldots
\begin{itemize}
\opcio{a}{la cardinal és més gran.}
\opcio{b}{són idèntiques.}
\opcio{c}{l'ordinal recorda l'ordre; la cardinal l'oblida (identifica bons ordres equipotents).}
\opcio{d}{l'ordinal no té exponenciació.}
\end{itemize}
\feedback
\justificacio{Tots els tipus numerables infinits tenen la mateixa mida; la cardinal els identifica, l'ordinal els distingeix per la seva forma.}

\pregunta Per què \(\beta\le\gamma\) implica \(\beta+\alpha\le\gamma+\alpha\), si \(B+A\) no és un segment de \(C+A\)?
\begin{itemize}
\opcio{a}{Perquè la suma és commutativa.}
\opcio{b}{Perquè hi ha una immersió estrictament creixent \(B+A\to C+A\), i una immersió entre bons ordres compara els tipus (rigidesa).}
\opcio{c}{Per l'axioma d'elecció.}
\opcio{d}{No és cert en general.}
\end{itemize}
\feedback
\justificacio{Si \(C+A\) fos més curt, seria isomorf a un segment propi de \(B+A\), i la composició amb la immersió faria baixar un punt, contra la rigidesa. L'estrictesa, en canvi, pot fallar.}

\pregunta Per què la base \(0\) es tracta a part en l'exponenciació ordinal?
\begin{itemize}
\opcio{a}{Perquè si s'hi apliqués la clàusula límit resultaria \(0^\omega=\sup\{1,0,0,\dots\}=1\), mentre que \(0^\beta\) ha de ser \(0\) per a \(\beta\ge1\).}
\opcio{b}{Perquè \(0^0\) no està definit.}
\opcio{c}{Perquè \(0\) no és un tipus ordinal.}
\opcio{d}{No es tracta a part.}
\end{itemize}
\feedback
\justificacio{La construcció per funcions de suport finit necessita el mínim \(0_A\), que no existeix si \(A\) és buit; i la clàusula límit, aplicada a \(0\), no reproduiria \(\abs{\varnothing^{B}}=0\). Per això \(0^0:=1\) i \(0^\beta:=0\) per a \(\beta>0\).}

\end{enumerate}
\clauestatica
\finalitzaTest

\chapter{Cardinals i aritmètica cardinal}\label{test:cap18}

Aquest test correspon al capítol~18 de la part teòrica. Cada pregunta té una única resposta correcta.

\iniciTest{u18}{b,c,a,b,c,b,a,c,b,c,b,c}

\begin{enumerate}

\pregunta El cardinal \(\abs{A}\), en el primer nivell, és\ldots
\begin{itemize}
\opcio{a}{un conjunt transitiu escollit.}
\opcio{b}{la classe d'equipotència de \(A\): el conjunt \(A\) llevat d'isomorfisme.}
\opcio{c}{el conjunt de parts de \(A\).}
\opcio{d}{un tipus d'ordre.}
\end{itemize}
\feedback
\justificacio{És la classe d'equipotència: no es tria cap representant. Els cardinals inicials, que són tipus d'ordre, apareixen més tard i amb l'elecció.}

\pregunta La suma cardinal es defineix com\ldots
\begin{itemize}
\opcio{a}{\(\abs{A}+\abs{B}=\abs{A\times B}\).}
\opcio{b}{\(\abs{A}+\abs{B}=\abs{A^B}\).}
\opcio{c}{\(\abs{A}+\abs{B}=\abs{A+B}\) (cardinal de la suma disjunta).}
\opcio{d}{el màxim de \(\abs{A}\) i \(\abs{B}\).}
\end{itemize}
\feedback
\justificacio{La suma és el cardinal de la suma disjunta; el producte, el del producte; l'exponenciació, el de l'exponencial. Prendre el màxim és un teorema per a infinits, no la definició.}

\pregunta Que les operacions cardinals estiguin ben definides vol dir que\ldots
\begin{itemize}
\opcio{a}{no depenen dels representants (invariància per equipotència).}
\opcio{b}{requereixen l'elecció.}
\opcio{c}{només valen per a conjunts finits.}
\opcio{d}{depenen de l'ordre.}
\end{itemize}
\feedback
\justificacio{Isomorfismes \(A\cong A'\), \(B\cong B'\) donen isomorfismes entre les sumes, els productes i els exponencials. Cap elecció.}

\pregunta A diferència de l'aritmètica ordinal, la cardinal\ldots
\begin{itemize}
\opcio{a}{no té exponenciació.}
\opcio{b}{és commutativa (suma i producte), perquè oblida l'ordre.}
\opcio{c}{no és associativa.}
\opcio{d}{és igual a l'ordinal.}
\end{itemize}
\feedback
\justificacio{\(A+B\cong B+A\) i \(A\times B\cong B\times A\) com a conjunts: oblidant l'ordre, l'asimetria de l'aritmètica ordinal desapareix.}

\pregunta El teorema de Cantor en forma cardinal diu\ldots
\begin{itemize}
\opcio{a}{\(\abs{A}=2^{\abs{A}}\).}
\opcio{b}{\(2^{\abs{A}}\le\abs{A}\).}
\opcio{c}{\(\abs{A}<2^{\abs{A}}\), amb \(2^{\abs{A}}=\abs{P(A)}\).}
\opcio{d}{\(\abs{A}\) és sempre finit.}
\end{itemize}
\feedback
\justificacio{\(P(A)\cong\mathbf2^A\), i Cantor dona \(A\prec P(A)\): \(\abs{A}\le2^{\abs{A}}\) i no hi ha igualtat.}

\pregunta Que dos conjunts qualssevol siguin comparables (\(A\preceq B\) o \(B\preceq A\))\ldots
\begin{itemize}
\opcio{a}{val sempre sense cap axioma.}
\opcio{b}{s'obté del teorema del bon ordre, i és equivalent a l'elecció.}
\opcio{c}{és fals.}
\opcio{d}{equival al teorema de Cantor.}
\end{itemize}
\feedback
\justificacio{Amb bons ordres sobre \(A\) i \(B\), la comparabilitat de bons ordres dona un monomorfisme en un sentit. Recíprocament, comparant \(A\) amb \(\Hart(A)\) es ben ordena \(A\): equival a l'elecció.}

\pregunta Combinant la comparabilitat amb Schröder--Bernstein, la comparació de cardinals és\ldots
\begin{itemize}
\opcio{a}{un ordre total llevat d'equipotència.}
\opcio{b}{només un preordre.}
\opcio{c}{una relació d'equivalència.}
\opcio{d}{no reflexiva.}
\end{itemize}
\feedback
\justificacio{Preordre sense elecció, antisimètric mòdul \(\approx\) per Schröder--Bernstein, i total amb l'elecció.}

\pregunta El cardinal inicial d'un conjunt és\ldots
\begin{itemize}
\opcio{a}{el conjunt de parts.}
\opcio{b}{qualsevol bon ordre equipotent.}
\opcio{c}{el menor tipus d'ordre de la seva mida (un tipus inicial, trobat acotadament).}
\opcio{d}{sempre \(\N\).}
\end{itemize}
\feedback
\justificacio{És el mínim, dins \(\Hart(A)\), dels tipus amb suport equipotent a \(A\): existeix i és únic, i cap tipus menor no té la mida d'\(A\).}

\pregunta El teorema de König afirma que, si \(A_i\prec B_i\) per a tot \(i\),\ldots
\begin{itemize}
\opcio{a}{\(\sum_i\abs{A_i}=\prod_i\abs{B_i}\).}
\opcio{b}{\(\sum_i\abs{A_i}<\prod_i\abs{B_i}\).}
\opcio{c}{\(\prod_i\abs{B_i}<\sum_i\abs{A_i}\).}
\opcio{d}{les famílies no existeixen.}
\end{itemize}
\feedback
\justificacio{La injecció de la suma en el producte dona \(\le\), i la diagonal coordenada a coordenada exclou tota exhaustiva. És una desigualtat estricta.}

\pregunta El teorema de König\ldots
\begin{itemize}
\opcio{a}{és independent de Cantor.}
\opcio{b}{contradiu Cantor.}
\opcio{c}{conté Cantor com a cas particular (\(A_i=\mathbf1\), \(B_i=\mathbf2\)); la seva prova és la diagonal.}
\opcio{d}{no usa la diagonal.}
\end{itemize}
\feedback
\justificacio{Amb \(I=A\), \(A_i=\mathbf1\), \(B_i=\mathbf2\), la suma és \(\abs{A}\) i el producte \(2^{\abs{A}}\). La família diagonal \((d_i)\) és la part diagonal de Cantor.}

\pregunta La monotonia en l'exponent, \(\kappa\le\mu\Rightarrow\nu^\kappa\le\nu^\mu\)\ldots
\begin{itemize}
\opcio{a}{val sempre.}
\opcio{b}{val per a \(\nu\neq0\); per a \(\nu=0\) falla: \(0^0=1\not\le0=0^1\).}
\opcio{c}{requereix l'elecció.}
\opcio{d}{només val per a conjunts finits.}
\end{itemize}
\feedback
\justificacio{Per estendre una aplicació \(A\to C\) a \(B\to C\) cal un valor per als punts de \(B\) fora d'\(A\), és a dir, un punt de \(C\). Amb \(C\) buit no n'hi ha, i el contraexemple és \(0^0=1\), \(0^1=0\).}

\pregunta Quantes vegades usa l'elecció la prova del teorema de König?
\begin{itemize}
\opcio{a}{Cap.}
\opcio{b}{Una, per a la diagonal.}
\opcio{c}{Tres: per triar els monomorfismes i els punts exteriors, per passar d'exhaustiva a domina, i per triar els punts de la diagonal.}
\opcio{d}{Només per definir la suma.}
\end{itemize}
\feedback
\justificacio{Les tres tries són seccions de famílies indexades per \(I\). Per a \(I\) finit es poden evitar, però en general König necessita l'elecció.}

\end{enumerate}
\clauestatica
\finalitzaTest

\chapter{Alephs, cofinalitat i continu}\label{test:cap19}

Aquest test correspon al capítol~19 de la part teòrica. Cada pregunta té una única resposta correcta.

\iniciTest{u19}{b,c,a,b,c,b,a,c,b,c,b,c}

\begin{enumerate}

\pregunta L'esquema de reunió de famílies (SR)\ldots
\begin{itemize}
\opcio{a}{és una sentència, com l'axioma d'elecció.}
\opcio{b}{és un \emph{esquema} (una instància per fórmula), no una sentència; no és un axioma més.}
\opcio{c}{és una forma d'elecció.}
\opcio{d}{ja és un dels axiomes de Nivell~I del llibre.}
\end{itemize}
\feedback
\justificacio{SR és una família d'enunciats, un per a cada fórmula funcional a isomorfisme únic, com el reemplaçament de ZFC. Es presenta com a principi a part i és opcional: el Nivell~I no l'inclou.}

\pregunta \(\aleph_0\) és\ldots
\begin{itemize}
\opcio{a}{el continu.}
\opcio{b}{un cardinal successor.}
\opcio{c}{\(\abs{\N}\), el menor cardinal infinit.}
\opcio{d}{un tipus ordinal límit qualsevol.}
\end{itemize}
\feedback
\justificacio{\(\aleph_0=\abs{\N}\) és el menor cardinal infinit: tot infinit de Dedekind conté una còpia de \(\N\), i amb l'elecció numerable tot infinit ho és. No és successor ni el continu.}

\pregunta El cardinal successor \(\kappa^+\)\ldots
\begin{itemize}
\opcio{a}{existeix sense l'esquema de reunió, acotat per Hartogs.}
\opcio{b}{requereix l'esquema de reunió.}
\opcio{c}{no existeix sempre.}
\opcio{d}{és igual a \(2^\kappa\).}
\end{itemize}
\feedback
\justificacio{\(\kappa^+=\abs{\Hart(K)}\): el menor bon ordre que no cap dins \(K\) és el menor cardinal més gran. Cap família no acotada, cap SR.}

\pregunta Què requereix l'esquema de reunió, i no només un pas?
\begin{itemize}
\opcio{a}{Calcular \(\aleph_0+\aleph_0\).}
\opcio{b}{Reunir, per a un ordinal fix \(\lambda\), el tram \((\aleph_\beta)_{\beta<\lambda}\) en una família interna indexada per \(\lambda\), especialment als passos límit.}
\opcio{c}{Obtenir un sol \(\kappa^+\).}
\opcio{d}{Comparar dos cardinals.}
\end{itemize}
\feedback
\justificacio{Cada pas és acotat; el que demana SR és reunir un tram fins a un ordinal fix en una família indexada per un conjunt, i prendre suprems als límits. La família de \emph{tots} els alephs no existeix com a objecte, ni amb SR: els tipus ordinals no formen cap conjunt d'índexs.}

\pregunta Sota elecció, tot cardinal infinit\ldots
\begin{itemize}
\opcio{a}{és \(\aleph_0\).}
\opcio{b}{és el continu.}
\opcio{c}{és un aleph \(\aleph_\alpha\) per a un únic \(\alpha\).}
\opcio{d}{no té cardinal inicial.}
\end{itemize}
\feedback
\justificacio{Amb el cardinal inicial i la família dels alephs fins a \(\kappa+1\), el mínim \(\alpha\) amb \(\aleph_\alpha\ge\kappa\) compleix \(\aleph_\alpha=\kappa\); la unicitat és la monotonia estricta.}

\pregunta Un cardinal és \emph{regular} quan\ldots
\begin{itemize}
\opcio{a}{és un cardinal límit.}
\opcio{b}{\(\operatorname{cf}(\kappa)=\kappa\): no s'arriba al seu final amb menys de \(\kappa\) passos.}
\opcio{c}{és singular.}
\opcio{d}{és igual a \(\aleph_0\).}
\end{itemize}
\feedback
\justificacio{Regular vol dir que tota part cofinal té la mateixa mida; els límits poden ser regulars o singulars.}

\pregunta \(\aleph_\omega\)\ldots
\begin{itemize}
\opcio{a}{és singular: \(\operatorname{cf}(\aleph_\omega)=\aleph_0<\aleph_\omega\).}
\opcio{b}{és regular.}
\opcio{c}{és un cardinal successor.}
\opcio{d}{és numerable.}
\end{itemize}
\feedback
\justificacio{La successió \((\aleph_n)_{n<\omega}\) és cofinal en \(\aleph_\omega\) i té longitud \(\omega\): \(\operatorname{cf}(\aleph_\omega)=\aleph_0\).}

\pregunta Els nombres beth es defineixen amb \(\beth_{\alpha+1}=\ldots\)
\begin{itemize}
\opcio{a}{\((\beth_\alpha)^+\).}
\opcio{b}{\(\beth_\alpha+1\).}
\opcio{c}{\(2^{\beth_\alpha}\), amb \(\beth_0=\aleph_0\) i suprem als límits.}
\opcio{d}{\(\beth_\alpha\cdot\aleph_0\).}
\end{itemize}
\feedback
\justificacio{Els beths pugen pel conjunt de parts; els alephs, pel successor. HCG diu que les dues jerarquies coincideixen.}

\pregunta El continu és\ldots
\begin{itemize}
\opcio{a}{\(\aleph_1\) per definició.}
\opcio{b}{\(\mathfrak c=\abs{P(\N)}=2^{\aleph_0}=\beth_1\).}
\opcio{c}{\(\aleph_0\).}
\opcio{d}{un cardinal finit.}
\end{itemize}
\feedback
\justificacio{\(\mathfrak c\) és la mida de \(P(\N)\), de Nivell~I; que sigui \(\aleph_1\) és la hipòtesi del continu, no una definició.}

\pregunta La hipòtesi del continu \(2^{\aleph_0}=\aleph_1\)\ldots
\begin{itemize}
\opcio{a}{és un teorema demostrable amb un argument diagonal.}
\opcio{b}{és falsa.}
\opcio{c}{és independent de ZFC si ZFC és consistent (i, per bi-interpretabilitat, d'aquesta teoria amb SR); la teoria la formula però no la resol.}
\opcio{d}{és el teorema de Cantor.}
\end{itemize}
\feedback
\justificacio{Gödel i Cohen mostren la independència relativa a la consistència de ZFC; com que ETCS amb SR i ZFC són bi-interpretables, el resultat es transporta. No és un teorema del llibre sinó un metateorema citat.}

\pregunta Per què l'esquema SR exigeix que cada \(X_i\) estigui determinat llevat d'un isomorfisme \emph{únic}?
\begin{itemize}
\opcio{a}{Perquè altrament l'esquema seria l'elecció.}
\opcio{b}{Perquè, sense unicitat, no hi ha cap manera canònica d'enganxar les fibres en una sola família; amb bons ordres, la rigidesa la garanteix.}
\opcio{c}{És indiferent.}
\opcio{d}{Perquè els conjunts infinits no tenen automorfismes.}
\end{itemize}
\feedback
\justificacio{Un conjunt infinit determinat només llevat d'isomorfisme té molts automorfismes, i triar-ne la identificació seria una forma d'elecció. Les instàncies del llibre porten sempre bons ordres, l'únic automorfisme dels quals és la identitat.}

\pregunta La regularitat dels cardinals successors usa\ldots
\begin{itemize}
\opcio{a}{només els axiomes de Nivell~I.}
\opcio{b}{SR.}
\opcio{c}{l'elecció i el teorema de Hessenberg, \(\kappa\cdot\kappa=\kappa\), que el llibre cita sense demostrar.}
\opcio{d}{la hipòtesi del continu.}
\end{itemize}
\feedback
\justificacio{Si una part cofinal tingués mida \(\le\aleph_\alpha\), \(\aleph_{\alpha+1}\) seria reunió de \(\le\aleph_\alpha\) conjunts de mida \(\le\aleph_\alpha\), i \(\aleph_\alpha\cdot\aleph_\alpha=\aleph_\alpha\). La reunió tria injeccions (elecció) i usa Hessenberg, que és un resultat clàssic citat.}

\end{enumerate}
\clauestatica
\finalitzaTest

\ifimprimible\else\end{Form}\fi

% ============================================================
% APÈNDIXS NUMERATS (dins \mainmatter: a \backmatter no es numeren)
% ============================================================
\appendix
\part{Apèndixs}
\chapter{Solucions dels tests}
\label{cap:solucions-tests}

\ifimprimible
Aquest apèndix reuneix les respostes i les justificacions dels tests. En aquesta edició per imprimir, són la manera de corregir-los.
\else
Aquest apèndix reuneix les respostes i les justificacions dels tests interactius. Serveixen per als lectors de PDF que no admeten la correcció automàtica i per a l'ús en paper.
\fi Es recomana consultar-les només després d'haver respost tot el test: cada justificació explica per què la resposta correcta ho és i, sovint, on falla el distractor més temptador.

\imprimeixsolucionstests

\chapter{Quasiexemples i contraexemples}\label{cap:contraexemples}

Una definició s'entén de debò quan se sap què \emph{no} la compleix. Aquest
apèndix recull, capítol per capítol, els exemples que mostren on falla una
afirmació plausible, quina hipòtesi és imprescindible o per què una
construcció no es pot simplificar. Molts provenen de la teoria clàssica de
conjunts o d'altres estructures matemàtiques: no són models dels axiomes del
llibre, sinó contrastos que fan visible el paper de cada axioma.

\section*{Aplicacions, parts i valors de veritat (capítols 1--3)}

\begin{description}
\item[Un conjunt sense punts.] El conjunt inicial \(\mathbf0\) és un conjunt
de ple dret i no té cap element global
(capítol~\ref{cap:aplicacions}). \enquote{Ser un conjunt} i \enquote{tenir
elements} són coses diferents.

\item[La separació global no és autodual.] El dual de l'axioma de separació
global (axioma~\ref{ax:separacio-globals}) diria que dues aplicacions
paral·leles coincideixen si coincideixen després de compondre-les amb
totes les aplicacions cap a \(\mathbf0\); com que d'un conjunt no inicial no
en surt cap, l'enunciat faria iguals totes les aplicacions paral·leles. A
partir d'aquest axioma la dualitat deixa de valer
(observació~\ref{obs:dualitat}).

\item[Un bimorfisme que no és isomorfisme.] En els anells, la inclusió
\(\Z\hookrightarrow\mathbb Q\) és mono i epi però no té inversa. En aquesta
teoria no pot passar: el teorema del bimorfisme
(proposició~\ref{prop:bimorfisme-isomorfisme}) usa l'axioma de
representació, que els anells no tenen.

\item[Dues parts amb el mateix domini.] Els dos punts
\(0,1\colon\mathbf1\hookrightarrow\mathbf2\) són parts de \(\mathbf2\) amb el
mateix domini i no són equivalents: una part és una fletxa, no el seu
domini.

\item[Una injecció de suma sense la retracció del cograf.] L'argument del
cograf (proposició~\ref{prop:cograf-retraccio}) fa mono la injecció
\(i_0\colon A\to A+B\) quan hi ha alguna \(B\to A\). Per a \(A=\mathbf0\) i
\(B=\mathbf1\) no n'hi ha cap; \(i_0\) és mono per una altra raó (tota
aplicació amb domini \(\mathbf0\) ho és), i en general cal la disjunció de
les sumes.

\item[Un element generalitzat repartit.] Siguin \(A=X=\mathbf2\), \(m\) el
punt \(1\colon\mathbf1\hookrightarrow\mathbf2\) i \(x=\id_{\mathbf2}\). Aleshores
\(x\) no pertany a \(m\) ni al seu complement: sobre el punt \(0\) cau al
complement i sobre el punt \(1\), a \(m\). Per això \enquote{\(x\in m^c\)
si i només si \(x\notin m\)} només val sobre punts globals
(observació~\ref{obs:complement-classic}).
\end{description}

\section*{Relacions, imatges, exponencials i parts (capítols 4--7)}

\begin{description}
\item[Un monomorfisme no escindit.] L'única aplicació \(\mathbf0\to\mathbf1\)
és mono i no té cap retracció, perquè no hi ha cap aplicació
\(\mathbf1\to\mathbf0\). L'axioma d'elecció parla dels epimorfismes, no dels
monomorfismes.

\item[Un epimorfisme sense secció.] En els conjunts constants, l'elecció ho
impedeix. Però en els conjunts que varien contínuament sobre un espai (els
feixos sobre la circumferència), la projecció del doble recobriment de la
circumferència sobre ella mateixa és un epimorfisme sense secció: la
\enquote{tria} d'un punt a cada fibra no es pot fer contínuament
\cite[pàg.~14]{bell2006}. Allà l'elecció i la separació global fallen
alhora.

\item[Un element generalitzat d'un exponencial no és una aplicació.] Un
element \(X\to B^A\) és una família d'aplicacions \(A\to B\) indexada per
\(X\); només per a \(X=\mathbf1\) n'és una de sola
(teorema~\ref{teo:punts-exponencial}).

\item[La unió interna sobre elements generalitzats.] Siguin
\(\alpha=\ulcorner1\urcorner\) i \(\beta=\ulcorner0\urcorner\), els noms dels
dos punts com a parts de \(\mathbf2\), i \(a=\id_{\mathbf2}\). Aleshores
\(a\mathrel{\varepsilon}(\alpha\cup\beta)\), perquè la unió és total, però
\(a\) no pertany ni a \(\alpha\) ni a \(\beta\). Les lectures amb \enquote{o}
i \enquote{no} només valen sobre punts
(proposició~\ref{prop:booleanes-puntuals}).
\end{description}

\section*{Nombres naturals, finitud i infinitud (capítols 8--12)}

\begin{description}
\item[El predecessor que no ho és.] La recursora del sistema
\((\N,0,\id_\N)\) compleix \(u\circ s=u\) i \(u\circ0=0\): és la constant
\(0\), no el predecessor. Cal una recursió sobre \(\N\times\N\)
(observació~\ref{obs:predecessor-parell}).

\item[Dues desigualtats que fallen alhora.] Els elements generalitzats
\(m=[0,s\circ0]\) i \(n=[s\circ0,0]\colon\mathbf2\to\N\) no compleixen ni \(m\le n\) ni
\(n\le m\), però el parell pertany a la unió de les dues relacions: la
totalitat de l'ordre és una afirmació sobre parts.

\item[Mono no iso amb la mateixa mida.] El successor \(s\colon\N\to\N\) és
mono i no epi. En els conjunts finits del mateix cardinal, mono, epi i iso
coincideixen; en els infinits, no.

\item[Una unió numerable sense elecció.] Sense l'elecció numerable és
consistent que la recta real sigui una unió numerable de conjunts
numerables (Feferman--Lévy; vegeu~\cite{jech2003}). El
teorema~\ref{teo:unio-numerable} necessita, doncs, l'elecció numerable.

\item[Dues dominacions sense isomorfisme evident.] \(\N\hookrightarrow\N+\mathbf1\)
i \(\N+\mathbf1\hookrightarrow\N\) (pel successor) són monomorfismes oposats,
cap dels dos no és isomorfisme, i tanmateix \(\N\cong\N+\mathbf1\):
Schröder--Bernstein construeix un isomorfisme \emph{nou}
(teorema~\ref{teo:schroder-bernstein}).
\end{description}

\section*{Bons ordres, ordinals i cardinals (capítols 13--19)}

\begin{description}
\item[Rigidesa en un ordre total que no és bo.] A \(\Z\), l'aplicació
\(n\mapsto n-1\) és creixent i injectiva i fa baixar tots els punts; i la
translació és un automorfisme diferent de la identitat. La rigidesa
(lema~\ref{lema:rigidesa}) depèn del mínim de les parts.

\item[Una cadena de bons ordres que no ho és.] Per a cada \(n\),
sigui \(W_n\) la part finita de \(\Z\) formada pels enters entre \(-n\) i
\(0\), amb l'ordre usual. Cada \(W_n\) és un bon ordre i estén l'anterior
\emph{per sota}; formen una cadena la unió de la qual (els enters no positius) no té
mínim. A la prova del lema de Zorn cal que les extensions siguin inicials.

\item[El suprem amb segments tancats.] Amb un sol índex i \(r\) el mínim
d'un bon ordre \(W\), el tipus de \(W_{<r}\) és \(0\) però \(W_{\le r}\) té
tipus \(1\) (observació~\ref{obs:suprem-oberts}).

\item[L'aritmètica ordinal no és commutativa.] \(1+\omega=\omega<\omega+1\)
i \(2\cdot\omega=\omega<\omega\cdot2\); i per l'esquerra la monotonia no és
estricta: \(1<2\) però \(1+\omega=2+\omega\)
(proposició~\ref{prop:monotonia-ordinal}).

\item[La base zero.] Amb la clàusula límit, \(0^\omega\) valdria
\(\sup\{1,0,0,\dots\}=1\), i no \(0\); per això la base \(0\) es tracta a
part (teorema~\ref{teo:exponenciacio-recursio}). En els cardinals, la
monotonia en l'exponent falla per la mateixa raó: \(0^0=1\not\le0=0^1\)
(proposició~\ref{prop:lleis-cardinals}).

\item[Una regla sense unicitat.] Una regla que determina cada \(X_i\) només
llevat d'isomorfisme ---per exemple, \enquote{un conjunt numerable}--- no
permet enganxar les fibres de manera canònica, perquè cada fibra té molts
automorfismes; per això l'esquema SR (principi~\ref{sr:esquema}) demana
isomorfisme únic.
\end{description}

\chapter{Tres conjunts al llarg del llibre}\label{cap:tres-conjunts}

Tres conjunts reapareixen a gairebé tots els capítols: \(\mathbf2\), \(\N\) i
el conjunt de parts \(P(A)\). Aquest apèndix recull, per a cadascun, el que
el llibre en va dient a mesura que avança, de manera que es pugui seguir el
fil d'una construcció a la següent.

\section*{\texorpdfstring{El conjunt \(\mathbf2\)}{El conjunt 2}}

\begin{description}
\item[Definició.] \(\mathbf2:=\mathbf1+\mathbf1\), la suma de dues còpies del
terminal; no es postula (capítol~\ref{cap:parts}).
\item[Dos punts.] Els seus punts són \(0\) i \(1\), són diferents i no n'hi ha
cap altre (corol·lari~\ref{cor:elements-de-2}); la distinció \(0\ne1\) usa la
disjunció de les sumes i la no-degeneració.
\item[Classificador de parts.] Cada part \(m\) d'un conjunt \(A\) té una única
característica \(\chi_m\colon A\to\mathbf2\), i cada aplicació cap a
\(\mathbf2\) classifica alguna part (axioma~\ref{ax:representacio},
teorema~\ref{teo:quadrat-caracteristic}).
\item[Lògica.] La negació, la conjunció i la disjunció són aplicacions
\(\mathbf2\to\mathbf2\) i \(\mathbf2\times\mathbf2\to\mathbf2\), i les operacions
entre parts són les seves composicions amb les característiques
(capítol~\ref{cap:valors}).
\item[Exponencials.] \(\mathbf2^{\mathbf2}\cong\mathbf2\times\mathbf2\) té quatre
punts: les dues constants, la identitat i la negació
(capítol~\ref{cap:exponencials}); i \(P(A):=\mathbf2^A\).
\item[Recompte.] \(\abs{P(A)}=2^{\abs{A}}\) per a \(A\) finit
(capítol~\ref{cap:finits}), i \(\abs{A}<2^{\abs{A}}\) en general (Cantor,
capítol~\ref{cap:cantor}).
\item[Ordinal i cardinal.] Com a tipus d'ordre, \(2^\omega=\omega\); com a
cardinal, \(2^{\aleph_0}=\mathfrak c\) és el continu
(capítols~\ref{cap:aritmetica-ordinal} i~\ref{cap:alephs}).
\end{description}

\section*{\texorpdfstring{L'objecte \(\N\)}{L'objecte N}}

\begin{description}
\item[Definició.] Una propietat universal de recursió: tot sistema
\((A,a,t)\) rep una única \(u\) amb \(u\circ0=a\) i \(u\circ s=t\circ u\)
(axioma~\ref{ax:naturals}).
\item[Inducció i Peano.] La unicitat de la recursora és la inducció; el
zero no és successor, el successor és injectiu i tot punt és zero o
successor (teorema~\ref{teo:peano}, lema~\ref{lema:zero-o-successor}).
D'aquí, \(\N\cong\mathbf1+\N\).
\item[Aritmètica i ordre.] Suma, producte i exponenciació per recursió;
\(\le\) com a part de \(\N\times\N\); bon ordre
(teorema~\ref{teo:bon-ordre-N}).
\item[Finitud.] Els segments \(\seg{n}\) són fibres d'una part de
\(\N\times\N\), i \(\operatorname{Card}(A)\) és una part de \(\N\)
(capítol~\ref{cap:finits}).
\item[Infinitud.] El successor fa de \(\N\) un infinit de Dedekind;
\(\N\cong\N\times\N\) (proposició~\ref{prop:estabilitat-enumerable}), i tota
part de \(\N\) és finita o isomorfa a \(\N\)
(teorema~\ref{teo:subobjectes-N}).
\item[Mida i forma.] \(\abs{\N}=\aleph_0\), el menor cardinal infinit;
\(\operatorname{tipus}(\N)=\omega\), el primer tipus límit; i
\(P(\N)\) no és enumerable (capítols~\ref{cap:cantor}, \ref{cap:ordinals}
i~\ref{cap:alephs}).
\end{description}

\section*{\texorpdfstring{El conjunt de parts \(P(A)\)}{El conjunt de parts P(A)}}

\begin{description}
\item[Definició.] \(P(A):=\mathbf2^A\); els seus punts són els noms de les
parts d'\(A\), i una aplicació \(X\to P(A)\) és una família de parts
(teorema~\ref{teo:parts-parametritzades}).
\item[Pertinença interna.] \(\in_A\hookrightarrow A\times P(A)\) és la part
classificada per l'avaluació (capítol~\ref{cap:parts-interna}).
\item[Lògica interna.] Operacions booleanes, imatge inversa \(P(f)\) i
quantificadors \(\exists_f\dashv P(f)\dashv\forall_f\)
(definició~\ref{def:quantificadors}, lema~\ref{lema:beck-chevalley}); la
inclusió interna en fa un reticle booleà (teorema~\ref{teo:inclusio-ordre}).
\item[Famílies i ínfims.] Les unions i interseccions indexades d'una família
de parts, i l'ínfim de les parts \(\Phi\)-tancades de Schröder--Bernstein,
s'obtenen amb els quantificadors (capítols~\ref{cap:cantor}
i~\ref{cap:families}).
\item[Mida.] \(A\prec P(A)\) (teorema~\ref{teo:cantor}); no hi ha conjunt
universal.
\item[Ordres i ordinals.] Els bons ordres sobre parts d'\(A\) viuen dins
\(P(A)\times P(A\times A)\), i el conjunt de Hartogs és un quocient d'una
part d'aquest conjunt (definició~\ref{def:hartogs}).
\item[Jerarquies.] La torre \(A,P(A),P^2(A),\dots\) com a família, i els
beths \(\beth_{\alpha+1}=2^{\beth_\alpha}\), demanen l'esquema SR
(capítols~\ref{cap:recursio} i~\ref{cap:alephs}).
\end{description}

\chapter[Correspondència amb les categories]{Correspondència amb el volum de categories}\label{cap:correspondencia}

Aquest llibre no fa servir el vocabulari de la teoria de categories, però
cada construcció seva és el cas particular, en els conjunts, d'una noció
categòrica general. La taula següent indica, per a cada construcció, el nom
categòric que rep i el capítol del volum \emph{Una introducció a la teoria de
categories}~\cite{perello-categories} on es tracta en general. Els capítols
d'aquell volum són: 1~Categories, 2~Functors, 3~Transformacions naturals,
4~Yoneda i propietats universals, 5~Límits i colímits, 6~Adjuncions,
7~Categories cartesianes tancades, 8~Subobjectes i categories regulars,
9~Topos elementals, i 10~el pont cap a la lògica categòrica.

{\small
\begin{longtable}{@{}>{\raggedright\arraybackslash\hspace{0pt}}p{0.34\linewidth}>{\raggedright\arraybackslash\hspace{0pt}}p{0.36\linewidth}>{\raggedright\arraybackslash\hspace{0pt}}p{0.20\linewidth}@{}}
\toprule
Aquest llibre & Noció categòrica & Categories\\
\midrule
\endfirsthead
\toprule
Aquest llibre & Noció categòrica & Categories\\
\midrule
\endhead
\bottomrule
\endfoot
Aplicacions, composició, identitats (cap.~\ref{cap:aplicacions}) & Categoria; la dels conjunts, \(\mathbf{Set}\) & cap.~1\\
Isomorfisme; dualitat (cap.~\ref{cap:aplicacions}) & Isomorfisme; categoria oposada & cap.~1\\
Element generalitzat \(X\to A\) & Morfisme amb codomini \(A\); punt de vista de Yoneda & cap.~4\\
Terminal i inicial (cap.~\ref{cap:aplicacions}) & Objectes terminal i inicial & cap.~4\\
Separació global & Categoria ben apuntada; \(\mathbf1\) separador & cap.~9\\
Monomorfisme, epimorfisme, secció, retracció (cap.~\ref{cap:parts}) & Les mateixes nocions en una categoria qualsevol & cap.~1\\
Part d'un conjunt; inclusió (cap.~\ref{cap:parts}) & Subobjecte; ordre de subobjectes & cap.~8\\
Suma i producte (cap.~\ref{cap:parts}) & Coproducte i producte & cap.~5\\
Equalitzador, pullback, imatge inversa (cap.~\ref{cap:parts}) & Límits finits; reindexació de subobjectes & caps.~5 i~8\\
Representació de les parts per \(\mathbf2\) (cap.~\ref{cap:parts}) & Classificador de subobjectes \(\Omega\) & cap.~9\\
Sumes disjuntes i estables (cap.~\ref{cap:valors}) & Coproductes disjunts i universals & caps.~5 i~9\\
Coequalitzador, pushout, límits i colímits finits (cap.~\ref{cap:valors}) & Colímits i límits finits & cap.~5\\
Relació d'equivalència; parell nucli; quocient (cap.~\ref{cap:relacions}) & Relació d'equivalència interna; \emph{kernel pair}; coequalitzador & caps.~5 i~8\\
Efectivitat de les relacions d'equivalència & Categoria exacta (de Barr) & cap.~8\\
Elecció: tot epi té secció & Axioma d'elecció interna; epis escindits & cap.~9\\
Factorització imatge (cap.~\ref{cap:imatges}) & Factorització epi regular\,--\,mono & cap.~8\\
Estabilitat dels epis; teoria regular & Categoria regular & cap.~8\\
Imatge directa \(\dashv\) imatge inversa (cap.~\ref{cap:imatges}) & Existencial com a adjunt de la reindexació & caps.~6 i~8\\
Exponencial \(B^A\); transposició (cap.~\ref{cap:exponencials}) & Objecte exponencial; categoria cartesiana tancada; currificació & caps.~6 i~7\\
Conjunt de parts \(P(A)=\mathbf2^A\) (cap.~\ref{cap:parts-interna}) & Objecte de parts \(\Omega^A\) & cap.~9\\
Quantificadors \(\exists_f\dashv P(f)\dashv\forall_f\); Beck--Chevalley & Hiperdoctrina de subobjectes; condició de Beck--Chevalley & caps.~8 i~9\\
Objecte dels nombres naturals (cap.~\ref{cap:naturals}) & Objecte de nombres naturals (NNO) & cap.~9\\
Conjunts finits; \(\operatorname{Card}(A)\) (cap.~\ref{cap:finits}) & Objectes finits (Kuratowski) en un topos & ---\\
Famílies \(E\to I\); fibres (cap.~\ref{cap:families}) & Categoria de fletxes \(\mathbf{Set}/I\) & caps.~4 i~9\\
Bons ordres, ordinals, cardinals (caps.~\ref{cap:bons-ordres}--\ref{cap:cardinals}) & Sense contrapartida directa al volum de categories & ---\\
Esquema SR (cap.~\ref{cap:alephs}) & Reemplaçament estructural & ---\\
Tot el sistema d'axiomes (cap.~\ref{cap:models}) & Topos ben apuntat amb NNO i elecció: ETCS & caps.~9 i~10\\
\end{longtable}
}

Dues lectures de la taula. La primera: tot el que el llibre fa fins al
capítol~\ref{cap:parts-interna} és teoria de topos aplicada a un topos
concret; el volum de categories ho fa per a un topos qualsevol, on ja no hi
ha prou punts ni elecció, i on el classificador pot tenir més de dos valors.
La segona: les files sense contrapartida ---la teoria transfinita i SR--- són
el que és propi de la teoria de \emph{conjunts}, i no de les categories: la
mida i la longitud, que només tenen sentit quan es pot comparar i ordenar
sense límit.


\backmatter
\chapter*{Balanç de la part teòrica}
\phantomsection
\addcontentsline{toc}{chapter}{Balanç de la part teòrica}
\label{cap:balanc}

Els capítols teòrics estableixen la infraestructura de la teoria estructural de conjunts, des del llenguatge d'aplicacions fins a la comparació general de mides i el llenguatge de les famílies indexades.

\begin{itemize}
  \item El capítol~\ref{cap:aplicacions} introdueix el llenguatge d'aplicacions, dominis, codominis i composició; defineix els conjunts com a identitats; i fixa els conjunts terminal i inicial, els elements generalitzats, els elements globals i la separació global.
  \item El capítol~\ref{cap:parts} desenvolupa monomorfismes, epimorfismes, parts i pertinença generalitzada. També introdueix sumes, productes, equalitzadors, pullbacks, imatges inverses i interseccions.
  \item El capítol~\ref{cap:valors} reforça les sumes amb la disjunció i l'estabilitat, distingeix els dos valors de veritat, construeix el complement i les operacions booleanes amb parts, i obté tots els límits i colímits finits.
  \item El capítol~\ref{cap:relacions} desenvolupa les relacions sobre elements generalitzats, les relacions d'equivalència en forma estructural, el parell nucli, els quocients com a coequalitzadors, els epimorfismes regulars i les particions, units en la correspondència de l'axioma d'efectivitat; i, amb l'axioma d'elecció, estableix que tot epimorfisme és escindit i regular.
  \item El capítol~\ref{cap:imatges} construeix la factorització de tota aplicació com un epimorfisme regular seguit d'un monomorfisme, defineix la imatge directa i la inversa amb la seva adjunció, demostra l'estabilitat dels epimorfismes regulars per canvi de base ---amb la qual la teoria disposa de la infraestructura d'una categoria regular--- i en treu la composició de relacions i les relacions funcionals.
  \item El capítol~\ref{cap:exponencials} postula l'existència d'exponencials, estableix la propietat universal de transposició i la unicitat de $B^A$, en calcula els casos amb $\mathbf0$, $\mathbf1$ i $\mathbf2$ ---aquest últim definint el conjunt de parts $P(A):=\mathbf2^A$---, i realitza internament la composició i les famílies d'aplicacions.
  \item El capítol~\ref{cap:parts-interna} desenvolupa el conjunt de parts com a objecte: la relació de pertinença interna, la identificació de les aplicacions cap a $P(A)$ amb les parts parametritzades, les operacions booleanes internes amb el seu càlcul puntual, la inclusió interna com a ordre parcial, i la imatge inversa, la imatge directa i el quantificador universal com a cadena d'adjuncions entre conjunts de parts.
  \item El capítol~\ref{cap:naturals} postula l'objecte dels nombres naturals mitjançant la propietat universal de recursió, en deriva el teorema de recursió, el principi d'inducció i les equacions de Peano, i construeix l'aritmètica ---suma, producte i exponenciació--- amb la recursió amb paràmetres via els exponencials.
  \item El capítol~\ref{cap:finits} converteix la finitud en una noció interna: els ordinals finits $[n]$ per recursió, el conjunt finit com el que és isomorf a algun $[n]$, la unicitat del cardinal, el principi dels calaixos i l'equivalència ---pròpia del cas finit--- entre injectivitat, exhaustivitat i bijectivitat, i el recompte de parts $2^n$ i d'aplicacions $m^n$.
  \item El capítol~\ref{cap:infinitud} estén la comparació de cardinals més enllà del cas finit: els conjunts infinits i els infinits de Dedekind i la seva relació ---amb l'elecció numerable com a bisagra---, els conjunts enumerables i numerables i la seva estabilitat sense elecció, els subobjectes de $\N$ com a sempre enumerables gràcies al mínim, les llistes finites amb el seu principi de recursió, i les unions indexades numerablement, on l'elecció numerable torna a ser essencial.
  \item El capítol~\ref{cap:cantor} obre la comparació general de mides: la relació de dominació per monomorfismes, el teorema de Schröder--Bernstein pel menor punt fix d'un operador creixent sobre $P(A)$ ---sense elecció---, l'argument diagonal de Cantor i la desigualtat estricta $A\prec P(A)$, amb la conseqüència que no hi ha conjunt universal ni mida màxima.
  \item El capítol~\ref{cap:families} fixa el llenguatge de les famílies indexades: una família és una aplicació $p\colon E\to I$ amb les fibres com a producte fibrat, les seccions com a funcions d'elecció, la suma i el producte d'una família, les famílies de parts d'un fix amb les seves unions i interseccions indexades, i la distinció central entre una família interna existent i una regla externa encara no representada, amb la substitució estructural SR enunciada com a esquema.
  \item El capítol~\ref{cap:bons-ordres} desenvolupa els ordres, els bons ordres i els segments inicials: el bon ordre com \enquote{tota part no inicial té mínim}, la rigidesa ---entre dos bons ordres hi ha com a molt un isomorfisme d'ordre--- i la comparació per segments, amb la tricotomia establerta en la seva part d'exclusivitat.
  \item El capítol~\ref{cap:recursio} estableix la inducció i la recursió transfinites amb codomini fix ---la reunió de solucions parcials dins un exponencial fix, sense elecció--- i en treu la comparabilitat total dels bons ordres.
  \item El capítol~\ref{cap:zorn} reuneix les tres formes de maximalitat: el teorema de Hartogs (un bon ordre que no s'injecta en $A$), el teorema del bon ordre i el lema de Zorn, amb totes les recursions acotades per $\Hart$, i l'equivalència entre elecció, bon ordre i Zorn.
  \item El capítol~\ref{cap:ordinals} defineix els tipus ordinals com a bons ordres llevat d'isomorfisme, amb zero, successor i límit, els suprems acotats, i el teorema de Burali--Forti com a paral·lel ordinal de Cantor.
  \item El capítol~\ref{cap:aritmetica-ordinal} construeix l'aritmètica ordinal ---suma, producte i exponenciació sobre representants---, no commutativa ($1+\omega=\omega\neq\omega+1$), amb la distinció entre ordre i mida com a fil conductor.
  \item El capítol~\ref{cap:cardinals} desenvolupa els cardinals i la seva aritmètica, relacional i sense representants, amb Cantor com $\abs{A}<2^{\abs{A}}$, la comparabilitat total sota elecció, els cardinals inicials i el teorema de König.
  \item El capítol~\ref{cap:alephs}, opcional i de Nivell~II, recorre la jerarquia dels alephs i dels beths i la cofinalitat fent servir l'esquema de reunió de famílies SR, i formula amb precisió la hipòtesi del continu com una qüestió independent ---tot i que el continu $\mathfrak c=\abs{P(\N)}$ i la mateixa equació d'HC $2^{\aleph_0}=\aleph_1$ són ja de Nivell~I, perquè $P(\N)$ i $\aleph_1=(\aleph_0)^+$ existeixen sense reunir cap família no acotada; només la successió sencera dels alephs demana SR.
  \item El capítol~\ref{cap:models} tanca el volum amb un balanç metateòric: recull tots els axiomes, caracteritza l'estructura com un model d'ETCS, distingeix la teoria de Nivell~I d'ETCS reforçat amb SR, compara l'ordinal material de von Neumann amb el tipus estructural, i en delimita els límits de mida, consistència i força expressiva.
\end{itemize}

En el metallenguatge, les classes d'equivalència de parts d'un conjunt formen una àlgebra booleana, i les relacions d'equivalència es corresponen amb les particions; internament, aquesta lògica de subconjunts es realitza com l'àlgebra d'un únic objecte $P(A)$, amb la pertinença, les operacions booleanes, l'ordre d'inclusió i els quantificadors ja construïts. L'objecte dels nombres naturals proporciona la recursió, la inducció i l'aritmètica, sobre les quals es construeixen la finitud interna, el recompte i la cardinalitat numerable, amb el paper de l'elecció numerable delimitat amb precisió. La segona part porta la construcció fins a la teoria clàssica de l'infinit: Schröder--Bernstein i Cantor obren la comparació general de mides; els bons ordres i la recursió transfinita, amb Hartogs i el teorema del bon ordre, donen la teoria de l'ordre; els tipus ordinals i els cardinals, amb les seves aritmètiques, separen la forma d'ordre de la mida; i Burali--Forti fa de paral·lel ordinal de Cantor. Tot això es fa amb els axiomes de Nivell~I ---o amb l'elecció---, mantenint les construccions acotades; només la reunió de les col·leccions transfinites no acotades ---la successió sencera dels alephs i dels beths--- demana l'esquema de substitució estructural SR, que porta la teoria a la força de ZFC. (El continu mateix, $\mathfrak c=\abs{P(\N)}$, i la hipòtesi del continu són de Nivell~I: no depenen de SR, sinó només del conjunt de parts i del cardinal successor per Hartogs.) El balanç metateòric final situa tot l'edifici, el caracteritza com un model d'ETCS, en distingeix la versió de Nivell~I de la reforçada amb SR, i el compara amb la teoria clàssica de von Neumann.

\chapter*{Glossari de símbols}
\addcontentsline{toc}{chapter}{Glossari de símbols}
\label{cap:glossari-simbols}
\markboth{Glossari de símbols}{Glossari de símbols}

Aquest glossari recull la notació pròpia del llibre, agrupada per blocs
temàtics. La darrera columna indica el capítol on el símbol s'introdueix o es
defineix.

\begingroup
\renewcommand{\arraystretch}{1.35}

\subsection*{Llenguatge bàsic: aplicacions, dominis i composició}

\noindent
\begin{longtable}{@{}p{0.20\textwidth}p{0.62\textwidth}p{0.10\textwidth}@{}}
\toprule
\textbf{Símbol} & \textbf{Significat} & \textbf{Cap.}\\
\midrule
\endfirsthead
\toprule
\textbf{Símbol} & \textbf{Significat} & \textbf{Cap.}\\
\midrule
\endhead
\bottomrule
\endfoot
\(\Dom(f,a)\) & \(a\) és el domini de l'aplicació \(f\) & \ref{cap:aplicacions}\\
\(\Cod(f,b)\) & \(b\) és el codomini de l'aplicació \(f\) & \ref{cap:aplicacions}\\
\(\Com(f,g,h)\) & \(h\) és la composició de \(f\) i després \(g\) & \ref{cap:aplicacions}\\
\(g\circ f\) & composició d'aplicacions (primer \(f\), després \(g\)) & \ref{cap:aplicacions}\\
\(\id_A\), \(\mathrm{id}\) & aplicació identitat sobre \(A\) & \ref{cap:aplicacions}\\
\(f\colon A\to B\) & aplicació de domini \(A\) i codomini \(B\) & \ref{cap:aplicacions}\\
\(x\colon X\to A\) & element generalitzat de \(A\) a l'estadi \(X\) & \ref{cap:aplicacions}\\
\(\mathbf1\) & conjunt terminal (un sol punt) & \ref{cap:aplicacions}\\
\(\mathbf0\) & conjunt inicial (buit) & \ref{cap:aplicacions}\\
\(a\colon\mathbf1\to A\) & element global o punt de \(A\) & \ref{cap:aplicacions}\\
\(A\cong B\) & \(A\) i \(B\) són isomorfs & \ref{cap:aplicacions}\\
\end{longtable}

\subsection*{Parts, lògica interna i construccions universals}

\noindent
\begin{longtable}{@{}p{0.20\textwidth}p{0.62\textwidth}p{0.10\textwidth}@{}}
\toprule
\textbf{Símbol} & \textbf{Significat} & \textbf{Cap.}\\
\midrule
\endfirsthead
\toprule
\textbf{Símbol} & \textbf{Significat} & \textbf{Cap.}\\
\midrule
\endhead
\bottomrule
\endfoot
\(S\hookrightarrow A\) & monomorfisme; \(S\) és una part de \(A\) & \ref{cap:parts}\\
\(P(A)\) & conjunt de parts de \(A\) & \ref{cap:parts-interna}\\
\(A\times B\) & producte de conjunts & \ref{cap:parts}\\
\(A+B\) & suma (coproducte) de conjunts & \ref{cap:parts}\\
\(B^A\) & exponencial: aplicacions de \(A\) cap a \(B\) & \ref{cap:exponencials}\\
\(a\mathrel{\varepsilon}\beta\) & pertinença generalitzada interna & \ref{cap:parts-interna}\\
\(\neg_A\) & complement intern sobre \(P(A)\) & \ref{cap:parts-interna}\\
\(\exists_f,\ \forall_f\) & quantificadors interns al llarg de \(f\) & \ref{cap:parts-interna}\\
\(\abs{A}\) & cardinal (mida) de \(A\) & \ref{cap:cardinals}\\
\end{longtable}

\subsection*{Naturals, ordre i transfinit}

\noindent
\begin{longtable}{@{}p{0.20\textwidth}p{0.62\textwidth}p{0.10\textwidth}@{}}
\toprule
\textbf{Símbol} & \textbf{Significat} & \textbf{Cap.}\\
\midrule
\endfirsthead
\toprule
\textbf{Símbol} & \textbf{Significat} & \textbf{Cap.}\\
\midrule
\endhead
\bottomrule
\endfoot
\(\N,\ \Z,\ \Q,\ \R\) & naturals, enters, racionals, reals & \ref{cap:infinitud}\\
\(\seg{n}\) & segment inicial \(\{0,\dots,n-1\}\) de \(\N\) & \ref{cap:infinitud}\\
\(A\preceq B\) & \(A\) s'injecta en \(B\) (mida menor o igual) & \ref{cap:cantor}\\
\(A\prec B\) & \(A\) té mida estrictament menor que \(B\) & \ref{cap:cantor}\\
\((S,R)\) & bon ordre: suport \(S\) amb relació \(R\) & \ref{cap:zorn}\\
\(\WO(A)\) & bons ordres sobre parts de \(A\) & \ref{cap:zorn}\\
\(\camp(S,R)\) & suport d'un bon ordre (és \(S\)) & \ref{cap:zorn}\\
\(\Hart(A)\) & conjunt de Hartogs de \(A\) & \ref{cap:zorn}\\
\([(S,R)],\ [R]\) & tipus d'ordre d'un bon ordre & \ref{cap:zorn}\\
\(\omega\) & tipus d'ordre dels naturals & \ref{cap:ordinals}\\
\(\aleph_\alpha\) & cardinals infinits (alephs) & \ref{cap:cardinals}\\
\(\aleph_0\) & cardinal dels naturals; \(\abs{\N}\) & \ref{cap:cardinals}\\
\(\kappa^{+}\) & cardinal successor de \(\kappa\) & \ref{cap:cardinals}\\
\(\beth_\alpha\) & cardinals beth & \ref{cap:alephs}\\
\(\mathfrak c\) & el continu; \(\abs{P(\N)}=2^{\aleph_0}\) & \ref{cap:alephs}\\
\(\operatorname{cf}(\kappa)\) & cofinalitat de \(\kappa\) & \ref{cap:alephs}\\
\end{longtable}

\endgroup

% Lectures complementàries no citades al cos (pendents d'integrar a la
% introducció i al futur capítol de conjunts variables).
\nocite{halmos,bell2006,lawvere1994,awodey,lawvere-schanuel}
\printbibliography[heading=bibintoc,title={Bibliografia}]
\printindex

\end{document}
